% Auto-generated from data/segments.json + layout.json (+ transforms) — do not edit by hand.
% volume: 12Sam01
% mode: reading
% segments: 3743
% transform_rules: 0
% toc_marks: 609

% Volume layout defaults (layout.json ``layout``).
\csromanlayoutapply{1}{1.5}{21.6pt}{5pt}{6.3pt}{65pt}{2.5em}%

% Reading physical-page layout (layout.json ``page_layout_reading_mode``).
\csromanreadingpagelayoutvolume{1}{1.5}{21.6pt}{5pt}{6.3pt}{65pt}{2.5em}%
\csromanreadingpagelayoutenable

\csromanheader{2}{\textbf{สํยุตฺตนิกาย}}
\gambhira{\textbf{สคาถาวคฺคสํยุตฺตปาฬิ}}
\namakkaram{นโม ตสฺส ภควโต อรหโต สมฺมาสมฺพุทฺธสฺส.}
\chapterhead{1. นฬวคฺค}
\csromanheader{2}{1. โอฆตรณสุตฺต}
\proseitem{1}{\csromanfolio{1}เอวํ เม สุตํ\csromandash{}เอกํ สมยํ ภควา สาวตฺถิยํ วิหรติ เชตวเน อนาถปิณฺฑิกสฺส อาราเม.\csromanspacer{} อถ โข อญฺญตรา เทวตา อภิกฺกนฺตาย รตฺติยา อภิกฺกนฺตวณฺณา เกวลกปฺปํ เชตวนํ โอภาเสตฺวา เยน ภควา เตนุปสงฺกมิ, อุปสงฺกมิตฺวา ภควนฺตํ อภิวาเทตฺวา เอกมนฺตํ อฏฺฐาสิ, เอกมนฺตํ ฐิตา โข สา เทวตา ภควนฺตํ เอตทโวจ “กถํ นุ ตฺวํ มาริส โอฆมตรี”ติ.\csromanspacer{} อปฺปติฏฺฐํ ขฺวาหํ อาวุโส อนายูหํ โอฆมตรินฺติ.\csromanspacer{} ยถากถํ ปน ตฺวํ มาริส อปฺปติฏฺฐํ อนายูหํ โอฆมตรีติ.\csromanspacer{} ยทาสฺวาหํ อาวุโส สนฺติฏฺฐามิ, ตทาสฺสุ สํสีทามิ.\csromanspacer{} ยทาสฺวาหํ อาวุโส อายูหามิ, ตทาสฺสุ นิพฺพุยฺหามิ\footnote{นิวุยฺหามิ (สฺยา, กํ, ก)}.\csromanspacer{} เอวํ ขฺวาหํ อาวุโส อปฺปติฏฺฐํ อนายูหํ โอฆมตรินฺติ.}

\csromangathameasure{“จิรสฺสํ วต ปสฺสามิ, พฺราหฺมณํ ปรินิพฺพุตํ.}
\csromangathagroup{\csromangathastanza{“จิรสฺสํ วต ปสฺสามิ, พฺราหฺมณํ ปรินิพฺพุตํ.}}

\vspace{\tipitakaparskip}
\csromancenter{อปฺปติฏฺฐํ อนายูหํ, ติณฺณํ โลเก วิสตฺติกนฺ”ติ\csromandash{}}
\prose{\csromanfolio{2}อิทมโวจ สา เทวตา.\csromanspacer{} สมนุญฺโญ สตฺถา อโหสิ.\csromanspacer{} อถ โข สา เทวตา “สมนุญฺโญ เม สตฺถา”ติ ภควนฺตํ อภิวาเทตฺวา ปทกฺขิณํ กตฺวา ตตฺเถวนฺตรธายีติ.}

\csromanheader{2}{2. นิโมกฺขสุตฺต}
\proseitem{2}{สาวตฺถินิทานํ.\csromanspacer{} อถ โข อญฺญตรา เทวตา อภิกฺกนฺตาย รตฺติยา อภิกฺกนฺตวณฺณา เกวลกปฺปํ เชตวนํ โอภาเสตฺวา เยน ภควา เตนุปสงฺกมิ, อุปสงฺกมิตฺวา ภควนฺตํ อภิวาเทตฺวา เอกมนฺตํ อฏฺฐาสิ, เอกมนฺตํ ฐิตา โข สา เทวตา ภควนฺตํ เอตทโวจ “ชานาสิ โน ตฺวํ มาริส สตฺตานํ นิโมกฺขํ ปโมกฺขํ วิเวกนฺ”ติ.\csromanspacer{} ชานามิ ขฺวาหํ อาวุโส สตฺตานํ นิโมกฺขํ ปโมกฺขํ วิเวกนฺติ.\csromanspacer{} ยถากถํ ปน ตฺวํ มาริส ชานาสิ สตฺตานํ นิโมกฺขํ ปโมกฺขํ วิเวกนฺติ.}

\prose{นนฺทีภวปริกฺขยา\footnote{นนฺทิภวปริกฺขยา (สฺยา, กํ)}, สญฺญาวิญฺญาณสงฺขยา.}

\prose{เวทนานํ นิโรธา อุปสมา, เอวํ ขฺวาหํ อาวุโส ชานามิ, สตฺตานํ นิโมกฺขํ, ปโมกฺขํ วิเวกนฺติ.}

\csromanheader{2}{3. อุปนียสุตฺต}
\proseitem{3}{สาวตฺถินิทานํ.\csromanspacer{} เอกมนฺตํ ฐิตา โข สา เทวตา ภควโต สนฺติเก}

\csromangathameasure{“อุปนียติ ชีวิตมปฺปมายุ, \\ ชรูปนีตสฺส น สนฺติ ตาณา. \\ ปุญฺญานิ กยิราถ สุขาวหานี”ติ. \\ อุปนียติ ชีวิตมปฺปมายุ, \\ ชรูปนีตสฺส น สนฺติ ตาณา. \\ โลกามิสํ ปชเห สนฺติเปกฺโขติ.}
\csromangathagroup{\csromangathastanza{“อุปนียติ ชีวิตมปฺปมายุ, \\ ชรูปนีตสฺส น สนฺติ ตาณา. \\ ปุญฺญานิ กยิราถ สุขาวหานี”ติ. \\ อุปนียติ ชีวิตมปฺปมายุ, \\ ชรูปนีตสฺส น สนฺติ ตาณา. \\ โลกามิสํ ปชเห สนฺติเปกฺโขติ.}}

\csromanheader{2}{4. อจฺเจนฺติสุตฺต}
\proseitem{4}{\csromanfolio{3}สาวตฺถินิทานํ.\csromanspacer{} เอกมนฺตํ ฐิตา โข สา เทวตา ภควโต สนฺติเก}

\csromangathameasure{“อจฺเจนฺติ กาลา ตรยนฺติ รตฺติโย, \\ วโยคุณา อนุปุพฺพํ ชหนฺติ. \\ ปุญฺญานิ กยิราถ สุขาวหานี”ติ. \\ อจฺเจนฺติ กาลา ตรยนฺติ รตฺติโย, \\ วโยคุณา อนุปุพฺพํ ชหนฺติ. \\ โลกามิสํ ปชเห สนฺติเปกฺโขติ.}
\csromangathagroup{\csromangathastanza{“อจฺเจนฺติ กาลา ตรยนฺติ รตฺติโย, \\ วโยคุณา อนุปุพฺพํ ชหนฺติ. \\ ปุญฺญานิ กยิราถ สุขาวหานี”ติ. \\ อจฺเจนฺติ กาลา ตรยนฺติ รตฺติโย, \\ วโยคุณา อนุปุพฺพํ ชหนฺติ. \\ โลกามิสํ ปชเห สนฺติเปกฺโขติ.}}

\csromanheader{2}{5. กติฉินฺทสุตฺต}
\proseitem{5}{สาวตฺถินิทานํ.\csromanspacer{} เอกมนฺตํ ฐิตา โข สา เทวตา ภควโต สนฺติเก}

\csromangathameasure{“กติ ฉินฺเท กติ ชเห, กติ จุตฺตริ ภาวเย. \\ กติ สงฺคาติโค ภิกฺขุ, ‘โอฆติณฺโณ’ติ วุจฺจตี”ติ. \\ ปญฺจ ฉินฺเท ปญฺจ ชเห, ปญฺจ จุตฺตริ ภาวเย. \\ ปญฺจ สงฺคาติโค ภิกฺขุ, “โอฆติณฺโณ”ติ วุจฺจตีติ.}
\csromangathagroup{\csromangathastanza{“กติ ฉินฺเท กติ ชเห, กติ จุตฺตริ ภาวเย. \\ กติ สงฺคาติโค ภิกฺขุ, ‘โอฆติณฺโณ’ติ วุจฺจตี”ติ.}\csromangathastanza{ปญฺจ ฉินฺเท ปญฺจ ชเห, ปญฺจ จุตฺตริ ภาวเย. \\ ปญฺจ สงฺคาติโค ภิกฺขุ, “โอฆติณฺโณ”ติ วุจฺจตีติ.}}

\csromanheader{2}{6. ชาครสุตฺต}
\proseitem{6}{สาวตฺถินิทานํ.\csromanspacer{} เอกมนฺตํ ฐิตา โข สา เทวตา ภควโต สนฺติเก}

\csromangathameasure{“กติ ชาครตํ สุตฺตา, กติ สุตฺเตสุ ชาครา. \\ กติภิ รชมาเทติ, กติภิ1 ปริสุชฺฌตี”ติ. \\ ปญฺจ ชาครตํ สุตฺตา, ปญฺจ สุตฺเตสุ ชาครา. \\ ปญฺจภิ รชมาเทติ, ปญฺจภิ2 ปริสุชฺฌตีติ.}
\csromangathagroup{\csromangathastanza{“กติ ชาครตํ สุตฺตา, กติ สุตฺเตสุ ชาครา. \\ กติภิ\footnote{กตีหิ (สี)} รชมาเทติ, กติภิ1 ปริสุชฺฌตี”ติ.}\csromangathastanza{ปญฺจ ชาครตํ สุตฺตา, ปญฺจ สุตฺเตสุ ชาครา. \\ ปญฺจภิ\footnote{ปญฺจหิ (สี)} รชมาเทติ, ปญฺจภิ2 ปริสุชฺฌตีติ.}}

\csromanheader{2}{7. อปฺปฏิวิทิตสุตฺต}
\proseitem{7}{\csromanfolio{4}สาวตฺถินิทานํ.\csromanspacer{} เอกมนฺตํ ฐิตา โข สา เทวตา ภควโต สนฺติเก}

\csromangathameasure{“เยสํ ธมฺมา อปฺปฏิวิทิตา, ปรวาเทสุ นียเร. \\ สุตฺตา เต นปฺปพุชฺฌนฺติ, กาโล เตสํ ปพุชฺฌิตุนฺ”ติ. \\ เยสํ ธมฺมา สุปฺปฏิวิทิตา, ปรวาเทสุ น นียเร. \\ เต สมฺพุทฺธา สมฺมทญฺญา, จรนฺติ วิสเม สมนฺติ.}
\csromangathagroup{\csromangathastanza{“เยสํ ธมฺมา อปฺปฏิวิทิตา, ปรวาเทสุ นียเร\footnote{นิยฺยเร (ก)}. \\ สุตฺตา เต นปฺปพุชฺฌนฺติ, กาโล เตสํ ปพุชฺฌิตุนฺ”ติ.}\csromangathastanza{เยสํ ธมฺมา สุปฺปฏิวิทิตา, ปรวาเทสุ น นียเร. \\ เต สมฺพุทฺธา สมฺมทญฺญา, จรนฺติ วิสเม สมนฺติ.}}

\csromanheader{2}{8. สุสมฺมุฏฺฐสุตฺต}
\proseitem{8}{สาวตฺถินิทานํ.\csromanspacer{} เอกมนฺตํ ฐิตา โข สา เทวตา ภควโต สนฺติเก}

\csromangathameasure{“เยสํ ธมฺมา สุสมฺมุฏฺฐา, ปรวาเทสุ นียเร. \\ สุตฺตา เต นปฺปพุชฺฌนฺติ, กาโล เตสํ ปพุชฺฌิตุนฺ”ติ. \\ เยสํ ธมฺมา อสมฺมุฏฺฐา, ปรวาเทสุ น นียเร. \\ เต สมฺพุทฺธา สมฺมทญฺญา, จรนฺติ วิสเม สมนฺติ.}
\csromangathagroup{\csromangathastanza{“เยสํ ธมฺมา สุสมฺมุฏฺฐา, ปรวาเทสุ นียเร. \\ สุตฺตา เต นปฺปพุชฺฌนฺติ, กาโล เตสํ ปพุชฺฌิตุนฺ”ติ.}\csromangathastanza{เยสํ ธมฺมา อสมฺมุฏฺฐา, ปรวาเทสุ น นียเร. \\ เต สมฺพุทฺธา สมฺมทญฺญา, จรนฺติ วิสเม สมนฺติ.}}

\csromanheader{2}{9. มานกามสุตฺต}
\proseitem{9}{สาวตฺถินิทานํ.\csromanspacer{} เอกมนฺตํ ฐิตา โข สา เทวตา ภควโต สนฺติเก}

\csromangathameasure{“น มานกามสฺส ทโม อิธตฺถิ, \\ น โมนมตฺถิ อสมาหิตสฺส. \\ เอโก อรญฺเญ วิหรํ ปมตฺโต, \\ น มจฺจุเธยฺยสฺส ตเรยฺย ปารนฺ”ติ. \\ มานํ ปหาย สุสมาหิตตฺโต, \\ สุเจตโส สพฺพธิ วิปฺปมุตฺโต. \\ เอโก อรญฺเญ วิหรํ อปฺปมตฺโต, \\ ส มจฺจุเธยฺยสฺส ตเรยฺย ปารนฺติ.}
\csromangathagroup{\csromangathastanza{“น มานกามสฺส ทโม อิธตฺถิ, \\ น โมนมตฺถิ อสมาหิตสฺส. \\ เอโก อรญฺเญ วิหรํ ปมตฺโต, \\ น มจฺจุเธยฺยสฺส ตเรยฺย ปารนฺ”ติ.}\csromangathastanza{มานํ ปหาย สุสมาหิตตฺโต, \\ สุเจตโส สพฺพธิ วิปฺปมุตฺโต. \\ เอโก อรญฺเญ วิหรํ อปฺปมตฺโต, \\ ส มจฺจุเธยฺยสฺส ตเรยฺย ปารนฺติ.}}

\csromanheader{2}{10. อรญฺญสุตฺต}
\proseitem{10}{\csromanfolio{5}สาวตฺถินิทานํ.\csromanspacer{} เอกมนฺตํ ฐิตา โข สา เทวตา ภควนฺตํ คาถาย อชฺฌภาสิ\csromandash{}}

\csromangathameasure{“อรญฺเญ วิหรนฺตานํ, สนฺตานํ พฺรหฺมจารินํ. \\ เอกภตฺตํ ภุญฺชมานานํ, เกน วณฺโณ ปสีทตี”ติ. \\ อตีตํ นานุโสจนฺติ, นปฺปชปฺปนฺติ นาคตํ. \\ ปจฺจุปฺปนฺเนน ยาเปนฺติ, เตน วณฺโณ ปสีทติ. \\ อนาคตปฺปชปฺปาย, อตีตสฺสานุโสจนา. \\ เอเตน พาลา สุสฺสนฺติ, นโฬว หริโต ลุโตติ.}
\csromangathagroup{\csromangathastanza{“อรญฺเญ วิหรนฺตานํ, สนฺตานํ พฺรหฺมจารินํ. \\ เอกภตฺตํ ภุญฺชมานานํ, เกน วณฺโณ ปสีทตี”ติ.}\csromangathastanza{อตีตํ นานุโสจนฺติ, นปฺปชปฺปนฺติ นาคตํ. \\ ปจฺจุปฺปนฺเนน ยาเปนฺติ, เตน วณฺโณ ปสีทติ.}\csromangathastanza{อนาคตปฺปชปฺปาย, อตีตสฺสานุโสจนา. \\ เอเตน พาลา สุสฺสนฺติ, นโฬว หริโต ลุโตติ.}}

\vspace{\tipitakaparskip}
\csromancenter{นฬวคฺโค ปฐโม.}
\titlehead{\textbf{ตสฺสุทฺทานํ}}
\csromangathameasure{โอฆํ นิโมกฺขํ อุปเนยฺยํ, อจฺเจนฺติ กติฉินฺทิ จ. \\ ชาครํ อปฺปฏิวิทิตา, สุสมฺมุฏฺฐา มานกามินา. \\ อรญฺเญ ทสโม วุตฺโต, วคฺโค เตน ปวุจฺจติ.}
\csromangathagroup{\csromangathastanza{โอฆํ นิโมกฺขํ อุปเนยฺยํ, อจฺเจนฺติ กติฉินฺทิ จ. \\ ชาครํ อปฺปฏิวิทิตา, สุสมฺมุฏฺฐา มานกามินา. \\ อรญฺเญ ทสโม วุตฺโต, วคฺโค เตน ปวุจฺจติ.}}

\chapterhead{2. นนฺทนวคฺค}
\csromanheader{2}{1. นนฺทนสุตฺต}
\proseitem{11}{เอวํ เม สุตํ\csromandash{}เอกํ สมยํ ภควา สาวตฺถิยํ วิหรติ เชตวเน อนาถปิณฺฑิกสฺส อาราเม.\csromanspacer{} ตตฺร โข ภควา ภิกฺขู อามนฺเตสิ “ภิกฺขโว”ติ. “ภทนฺเต”ติ เต ภิกฺขู ภควโต ปจฺจสฺโสสุํ.\csromanspacer{} ภควา เอตทโวจ\csromandash{}}

\prose{ภูตปุพฺพํ ภิกฺขเว อญฺญตรา ตาวติํสกายิกา เทวตา นนฺทเน วเน อจฺฉราสงฺฆปริวุตา ทิพฺเพหิ ปญฺจหิ กามคุเณหิ สมปฺปิตา สมงฺคีภูตา ปริจารยมานา\footnote{ปริจาริยมานา (สฺยา, กํ, ก)} ตายํ เวลายํ อิมํ คาถํ อภาสิ\csromandash{}}

\csromangathameasure{“น เต สุขํ ปชานนฺติ, เย น ปสฺสนฺติ นนฺทนํ. \\ อาวาสํ นรเทวานํ, ติทสานํ ยสสฺสินนฺ”ติ.}
\csromangathagroup{\csromangathastanza{\csromanfolio{6}“น เต สุขํ ปชานนฺติ, เย น ปสฺสนฺติ นนฺทนํ. \\ อาวาสํ นรเทวานํ, ติทสานํ ยสสฺสินนฺ”ติ.}}

\prose{เอวํ วุตฺเต ภิกฺขเว อญฺญตรา เทวตา ตํ เทวตํ คาถาย ปจฺจภาสิ\csromandash{}}

\csromangathameasure{“น ตฺวํ พาเล ปชานาสิ, ยถา อรหตํ วโจ. \\ อนิจฺจา สพฺพสงฺขารา, อุปฺปาทวยธมฺมิโน. \\ อุปฺปชฺชิตฺวา นิรุชฺฌนฺติ, เตสํ วูปสโม สุโข”ติ.}
\csromangathagroup{\csromangathastanza{“น ตฺวํ พาเล ปชานาสิ, ยถา อรหตํ วโจ. \\ อนิจฺจา สพฺพสงฺขารา\footnote{สพฺเพ สงฺขารา (สี, สฺยา, กํ)}, อุปฺปาทวยธมฺมิโน. \\ อุปฺปชฺชิตฺวา นิรุชฺฌนฺติ, เตสํ วูปสโม สุโข”ติ.}}

\csromanheader{2}{2. นนฺทติสุตฺต}
\proseitem{12}{สาวตฺถินิทานํ.\csromanspacer{} เอกมนฺตํ ฐิตา โข สา เทวตา ภควโต สนฺติเก}

\csromangathameasure{“นนฺทติ ปุตฺเตหิ ปุตฺติมา, \\ โคมา โคหิ ตเถว นนฺทติ. \\ อุปธีหิ นรสฺส นนฺทนา, \\ น หิ โส นนฺทติ โย นิรูปธี”ติ. \\ โสจติ ปุตฺเตหิ ปุตฺติมา, \\ โคมา โคหิ ตเถว โสจติ. \\ อุปธีหิ นรสฺส โสจนา, \\ น หิ โส โสจติ โย นิรูปธีติ.}
\csromangathagroup{\csromangathastanza{“นนฺทติ ปุตฺเตหิ ปุตฺติมา, \\ โคมา\footnote{โคมิโก (สี, สฺยา, กํ, อิ)} โคหิ ตเถว นนฺทติ. \\ อุปธีหิ นรสฺส นนฺทนา, \\ น หิ โส นนฺทติ โย นิรูปธี”ติ.}\csromangathastanza{โสจติ ปุตฺเตหิ ปุตฺติมา, \\ โคมา โคหิ ตเถว โสจติ. \\ อุปธีหิ นรสฺส โสจนา, \\ น หิ โส โสจติ โย นิรูปธีติ.}}

\csromanheader{2}{3. นตฺถิปุตฺตสมสุตฺต}
\proseitem{13}{สาวตฺถินิทานํ.\csromanspacer{} เอกมนฺตํ ฐิตา โข สา เทวตา ภควโต สนฺติเก}

\csromangathameasure{“นตฺถิ ปุตฺตสมํ เปมํ, นตฺถิ โคสมิตํ ธนํ. \\ นตฺถิ สูริยสมา อาภา, สมุทฺทปรมา สรา”ติ. \\ นตฺถิ อตฺตสมํ เปมํ, นตฺถิ ธญฺญสมํ ธนํ. \\ นตฺถิ ปญฺญาสมา อาภา, วุฏฺฐิ เว ปรมา สราติ.}
\csromangathagroup{\csromangathastanza{“นตฺถิ ปุตฺตสมํ เปมํ, นตฺถิ โคสมิตํ ธนํ. \\ นตฺถิ สูริยสมา\footnote{สุริยสมา (สี, สฺยา, กํ, อิ)} อาภา, สมุทฺทปรมา สรา”ติ.}\csromangathastanza{นตฺถิ อตฺตสมํ เปมํ, นตฺถิ ธญฺญสมํ ธนํ. \\ นตฺถิ ปญฺญาสมา อาภา, วุฏฺฐิ เว ปรมา สราติ.}}

\csromanheader{2}{4. ขตฺติยสุตฺต}
\csromanhangingbegin
\hangingheaditem{14}{\csromanfolio{7}ขตฺติโย ทฺวิปทํ เสฏฺโฐ, พลีพทฺโท\footnote{พลิวทฺโท (สี, อิ), พลิพทฺโท (สฺยา, กํ, ก)} จตุปฺปทํ.}
\hangingbody{โกมารี เสฏฺฐา ภริยานํ, โย จ ปุตฺตาน ปุพฺพโชติ.}
\hangingbody{สมฺพุทฺโธ ทฺวิปทํ เสฏฺโฐ, อาชานีโย จตุปฺปทํ.}
\hangingbody{สุสฺสูสา เสฏฺฐา ภริยานํ, โย จ ปุตฺตานมสฺสโวติ.}
\csromanhangingend

\csromanheader{2}{5. สณมานสุตฺต}
\csromanhangingbegin
\hangingheaditem{15}{ฐิเต มชฺฌนฺหิเก\footnote{มชฺฌนฺติเก (สพฺพตฺถ)} กาเล, สนฺนิสีเวสุ ปกฺขิสุ.}
\hangingbody{สณเตว พฺรหารญฺญํ3, ตํ ภยํ ปฏิภาติ มนฺติ.}
\hangingbody{ฐิเต มชฺฌนฺหิเก กาเล, สนฺนิสีเวสุ ปกฺขิสุ.}
\hangingbody{สณเตว พฺรหารญฺญํ, สา รติ ปฏิภาติ มนฺติ.}
\csromanhangingend

\csromanheader{2}{6. นิทฺทาตนฺทีสุตฺต}
\csromanhangingbegin
\hangingheaditem{16}{นิทฺทา ตนฺที วิชมฺภิตา\footnote{ตนฺทิ วิชมฺภิกา (สี, อิ)}, อรตี ภตฺตสมฺมโท.}
\hangingbody{เอเตน นปฺปกาสติ, อริยมคฺโค อิธ ปาณินนฺติ.}
\hangingbody{นิทฺทํ ตนฺทิํ วิชมฺภิตํ, อรติํ ภตฺตสมฺมทํ.}
\hangingbody{วีริเยน5 นํ ปณาเมตฺวา, อริยมคฺโค วิสุชฺฌตีติ.}
\csromanhangingend

\csromanheader{2}{7. ทุกฺกรสุตฺต}
\csromanhangingbegin
\hangingheaditem{17}{ทุกฺกรํ ทุตฺติติกฺขญฺจ, อพฺยตฺเตน จ สามญฺญํ.}
\hangingbody{พหูหิ ตตฺถ สมฺพาธา, ยตฺถ พาโล วิสีทตีติ.}
\hangingbody{กติหํ จเรยฺย สามญฺญํ, จิตฺตํ เจ น นิวารเย.}
\hangingbody{ปเท ปเท วิสีเทยฺย, สงฺกปฺปานํ วสานุโคติ.}
\hangingbody{กุมฺโมว องฺคานิ สเก กปาเล,}
\hangingbody{สโมทหํ ภิกฺขุ มโนวิตกฺเก.}
\hangingbody{อนิสฺสิโต อญฺญมเหฐยาโน,}
\hangingbody{ปรินิพฺพุโต นูปวเทยฺย กญฺจีติ.}
\csromanhangingend

\csromanheader{2}{8. หิรีสุตฺต}
\csromanhangingbegin
\hangingheaditem{18}{\csromanfolio{8}หิรีนิเสโธ ปุริโส, โกจิ โลกสฺมิํ วิชฺชติ.}
\hangingbody{โย นินฺทํ อปโพธติ1, อสฺโส ภโทฺร กสามิวาติ.}
\hangingbody{หิรีนิเสธา ตนุยา, เย จรนฺติ สทา สตา.}
\hangingbody{อนฺตํ ทุกฺขสฺส ปปฺปุยฺย, จรนฺติ วิสเม สมนฺติ.}
\csromanhangingend

\csromanheader{2}{9. กุฏิกาสุตฺต}
\csromanhangingbegin
\hangingheaditem{19}{กจฺจิ เต กุฏิกา นตฺถิ, กจฺจิ นตฺถิ กุลาวกา.}
\hangingbody{กจฺจิ สนฺตานกา นตฺถิ, กจฺจิ มุตฺโตสิ พนฺธนาติ.}
\hangingbody{ตคฺฆ เม กุฏิกา นตฺถิ, ตคฺฆ นตฺถิ กุลาวกา.}
\hangingbody{ตคฺฆ สนฺตานกา นตฺถิ, ตคฺฆ มุตฺโตมฺหิ พนฺธนาติ.}
\hangingbody{กินฺตาหํ กุฏิกํ พฺรูมิ, กิํ เต พฺรูมิ กุลาวกํ.}
\hangingbody{กิํ เต สนฺตานกํ พฺรูมิ, กินฺตาหํ พฺรูมิ พนฺธนนฺติ.}
\hangingbody{มาตรํ กุฏิกํ พฺรูสิ, ภริยํ พฺรูสิ กุลาวกํ.}
\hangingbody{ปุตฺเต สนฺตานเก พฺรูสิ, ตณฺหํ เม พฺรูสิ พนฺธนนฺติ.}
\hangingbody{สาหุ เต กุฏิกา นตฺถิ, สาหุ นตฺถิ กุลาวกา.}
\hangingbody{สาหุ สนฺตานกา นตฺถิ, สาหุ มุตฺโตสิ พนฺธนาติ.}
\csromanhangingend

\csromanheader{2}{10. สมิทฺธิสุตฺต}
\proseitem{20}{เอวํ เม สุตํ\csromandash{}เอกํ สมยํ ภควา ราชคเห วิหรติ ตโปทาราเม.\csromanspacer{} อถ โข อายสฺมา สมิทฺธิ รตฺติยา ปจฺจูสสมยํ ปจฺจุฏฺฐาย เยน ตโปทา เตนุปสงฺกมิ คตฺตานิ ปริสิญฺจิตุํ, ตโปเท คตฺตานิ ปริสิญฺจิตฺวา ปจฺจุตฺตริตฺวา เอกจีวโร อฏฺฐาสิ คตฺตานิ ปุพฺพาปยมาโน.\csromanspacer{} อถ โข อญฺญตรา เทวตา อภิกฺกนฺตาย รตฺติยา อภิกฺกนฺตวณฺณา เกวลกปฺปํ ตโปทํ โอภาเสตฺวา เยน อายสฺมา สมิทฺธิ เตนุปสงฺกมิ, อุปสงฺกมิตฺวา เวหาสํ ฐิตา อายสฺมนฺตํ สมิทฺธิํ คาถาย อชฺฌภาสิ\csromandash{}}

\csromangathameasure{“อภุตฺวา ภิกฺขสิ ภิกฺขุ, น หิ ภุตฺวาน ภิกฺขสิ. \\ ภุตฺวาน ภิกฺขุ ภิกฺขสฺสุ, มา ตํ กาโล อุปจฺจคา”ติ. \\ กาลํ โวหํ น ชานามิ, ฉนฺโน กาโล น ทิสฺสติ. \\ ตสฺมา อภุตฺวา ภิกฺขามิ, มา มํ กาโล อุปจฺจคาติ.}
\csromangathagroup{\csromangathastanza{\csromanfolio{9}“อภุตฺวา ภิกฺขสิ ภิกฺขุ, น หิ ภุตฺวาน ภิกฺขสิ. \\ ภุตฺวาน ภิกฺขุ ภิกฺขสฺสุ, มา ตํ กาโล อุปจฺจคา”ติ.}\csromangathastanza{กาลํ โวหํ น ชานามิ, ฉนฺโน กาโล น ทิสฺสติ. \\ ตสฺมา อภุตฺวา ภิกฺขามิ, มา มํ กาโล อุปจฺจคาติ.}}

\prose{อถ โข สา เทวตา ปถวิยํ\footnote{ปฐวิยํ (สี, สฺยา, กํ, อิ)} ปติฏฺฐหิตฺวา อายสฺมนฺตํ สมิทฺธิํ เอตทโวจ “ทหโร ตฺวํ ภิกฺขุ ปพฺพชิโต สุสุ กาฬเกโส ภเทฺรน โยพฺพเนน สมนฺนาคโต ปฐเมน วยสา อนิกฺกีฬิตาวี กาเมสุ ภุญฺช ภิกฺขุ มานุสเก กาเม, มา สนฺทิฏฺฐิกํ หิตฺวา กาลิกํ อนุธาวี”ติ.}

\prose{น ขฺวาหํ อาวุโส สนฺทิฏฺฐิกํ หิตฺวา กาลิกํ อนุธาวามิ, กาลิกญฺจ ขฺวาหํ อาวุโส หิตฺวา สนฺทิฏฺฐิกํ อนุธาวามิ, กาลิกา หิ อาวุโส กามา วุตฺตา ภควตา พหุทุกฺขา พหุปายาสา, อาทีนโว เอตฺถ ภิยฺโย.\csromanspacer{} สนฺทิฏฺฐิโก อยํ ธมฺโม อกาลิโก เอหิปสฺสิโก โอปเนยฺยิโก ปจฺจตฺตํ เวทิตพฺโพ วิญฺญูหีติ.}

\prose{กถญฺจ ภิกฺขุ กาลิกา กามา วุตฺตา ภควตา พหุทุกฺขา พหุปายาสา, อาทีนโว เอตฺถ ภิยฺโย, กถํ สนฺทิฏฺฐิโก อยํ ธมฺโม อกาลิโก เอหิปสฺสิโก โอปเนยฺยิโก ปจฺจตฺตํ เวทิตพฺโพ วิญฺญูหีติ.\csromanspacer{} อหํ โข อาวุโส นโว อจิรปพฺพชิโต อธุนาคโต อิมํ ธมฺมวินยํ, น ตาหํ\footnote{น ขฺวาหํ (สี, อิ)} สกฺโกมิ วิตฺถาเรน อาจิกฺขิตุํ, อยํ โส ภควา อรหํ สมฺมาสมฺพุทฺโธ ราชคเห วิหรติ ตโปทาราเม, ตํ ภควนฺตํ อุปสงฺกมิตฺวา เอตมตฺถํ ปุจฺฉ, ยถา เต ภควา พฺยากโรติ, ตถา นํ ธาเรยฺยาสีติ.}

\prose{น โข ภิกฺขุ สุกโร โส ภควา อเมฺหหิ อุปสงฺกมิตุํ, อญฺญาหิ มเหสกฺขาหิ เทวตาหิ ปริวุโต, สเจ โข ตฺวํ ภิกฺขุ ตํ ภควนฺตํ อุปสงฺกมิตฺวา เอตมตฺถํ ปุจฺเฉยฺยาสิ, มยมฺปิ อาคจฺเฉยฺยาม ธมฺมสฺสวนายาติ. “เอวมาวุโส”ติ โข อายสฺมา สมิทฺธิ ตสฺสา เทวตาย ปฏิสฺสุตฺวา เยน ภควา เตนุปสงฺกมิ, อุปสงฺกมิตฺวา ภควนฺตํ อภิวาเทตฺวา เอกมนฺตํ นิสีทิ, เอกมนฺตํ นิสินฺโน โข อายสฺมา สมิทฺธิ ภควนฺตํ เอตทโวจ\csromandash{}}

\prose{\csromanfolio{10}อิธาหํ ภนฺเต รตฺติยา ปจฺจูสสมยํ ปจฺจุฏฺฐาย เยน ตโปทา เตนุปสงฺกมิํ คตฺตานิ ปริสิญฺจิตุํ, ตโปเท คตฺตานิ ปริสิญฺจิตฺวา ปจฺจุตฺตริตฺวา เอกจีวโร อฏฺฐาสิํ คตฺตานิ ปุพฺพาปยมาโน.\csromanspacer{} อถ โข ภนฺเต อญฺญตรา เทวตา อภิกฺกนฺตาย รตฺติยา อภิกฺกนฺตวณฺณา เกวลกปฺปํ ตโปทํ โอภาเสตฺวา เยนาหํ เตนุปสงฺกมิ, อุปสงฺกมิตฺวา เวหาสํ ฐิตา อิมาย คาถาย อชฺฌภาสิ\csromandash{}}

\csromangathameasure{“อภุตฺวา ภิกฺขสิ ภิกฺขุ, น หิ ภุตฺวาน ภิกฺขสิ. \\ ภุตฺวาน ภิกฺขุ ภิกฺขสฺสุ, มา ตํ กาโล อุปจฺจคา”ติ.}
\csromangathagroup{\csromangathastanza{“อภุตฺวา ภิกฺขสิ ภิกฺขุ, น หิ ภุตฺวาน ภิกฺขสิ. \\ ภุตฺวาน ภิกฺขุ ภิกฺขสฺสุ, มา ตํ กาโล อุปจฺจคา”ติ.}}

\csromanheader{2}{เอวํ วุตฺเต อหํ ภนฺเต ตํ เทวตํ คาถาย ปจฺจภาสิํ\csromandash{}}
\csromangathameasure{“กาลํ โวหํ น ชานามิ, ฉนฺโน กาโล น ทิสฺสติ. \\ ตสฺมา อภุตฺวา ภิกฺขามิ, มา มํ กาโล อุปจฺจคา”ติ.}
\csromangathagroup{\csromangathastanza{“กาลํ โวหํ น ชานามิ, ฉนฺโน กาโล น ทิสฺสติ. \\ ตสฺมา อภุตฺวา ภิกฺขามิ, มา มํ กาโล อุปจฺจคา”ติ.}}

\prose{อถ โข ภนฺเต สา เทวตา ปถวิยํ ปติฏฺฐหิตฺวา มํ เอตทโวจ “ทหโร ตฺวํ ภิกฺขุ ปพฺพชิโต สุสุ กาฬเกโส ภเทฺรน โยพฺพเนน สมนฺนาคโต ปฐเมน วยสา อนิกฺกีฬิตาวี กาเมสุ ภุญฺช ภิกฺขุ มานุสเก กาเม, มา สนฺทิฏฺฐิกํ หิตฺวา กาลิกํ อนุธาวี”ติ.}

\prose{เอวํ วุตฺตาหํ ภนฺเต ตํ เทวตํ เอตทโวจํ “น ขฺวาหํ อาวุโส สนฺทิฏฺฐิกํ หิตฺวา กาลิกํ อนุธาวามิ, กาลิกญฺจ ขฺวาหํ อาวุโส หิตฺวา สนฺทิฏฺฐิกํ อนุธาวามิ, กาลิกา หิ อาวุโส กามา วุตฺตา ภควตา พหุทุกฺขา พหุปายาสา, อาทีนโว เอตฺถ ภิยฺโย.\csromanspacer{} สนฺทิฏฺฐิโก อยํ ธมฺโม อกาลิโก เอหิปสฺสิโก โอปเนยฺยิโก ปจฺจตฺตํ เวทิตพฺโพ วิญฺญูหี”ติ.}

\prose{เอวํ วุตฺเต ภนฺเต สา เทวตา มํ เอตทโวจ “กถญฺจ ภิกฺขุ กาลิกา กามา วุตฺตา ภควตา พหุทุกฺขา พหุปายาสา, อาทีนโว เอตฺถ ภิยฺโย, กถํ สนฺทิฏฺฐิโก อยํ ธมฺโม อกาลิโก เอหิปสฺสิโก โอปเนยฺยิโก ปจฺจตฺตํ เวทิตพฺโพ วิญฺญูหี”ติ.\csromanspacer{} เอวํ วุตฺตาหํ ภนฺเต ตํ เทวตํ เอตทโวจํ “อหํ โข อาวุโส นโว อจิรปพฺพชิโต อธุนาคโต อิมํ ธมฺมวินยํ, น ตาหํ สกฺโกมิ วิตฺถาเรน อาจิกฺขิตุํ, อยํ โส ภควา อรหํ สมฺมาสมฺพุทฺโธ ราชคเห วิหรติ ตโปทาราเม, ตํ ภควนฺตํ อุปสงฺกมิตฺวา เอตมตฺถํ ปุจฺฉ, ยถา เต ภควา พฺยากโรติ, ตถา นํ ธาเรยฺยาสี”ติ.}

\prose{\csromanfolio{11}เอวํ วุตฺเต ภนฺเต สา เทวตา มํ เอตทโวจ “น โข ภิกฺขุ สุกโร โส ภควา อเมฺหหิ อุปสงฺกมิตุํ, อญฺญาหิ มเหสกฺขาหิ เทวตาหิ ปริวุโต, สเจ โข ตฺวํ ภิกฺขุ ตํ ภควนฺตํ อุปสงฺกมิตฺวา เอตมตฺถํ ปุจฺเฉยฺยาสิ, มยมฺปิ อาคจฺเฉยฺยาม ธมฺมสฺสวนายา”ติ.\csromanspacer{} สเจ ภนฺเต ตสฺสา เทวตาย สจฺจํ วจนํ อิเธว สา เทวตา อวิทูเรติ.\csromanspacer{} เอวํ วุตฺเต สา เทวตา อายสฺมนฺตํ สมิทฺธิํ เอตทโวจ “ปุจฺฉ ภิกฺขุ, ปุจฺฉ ภิกฺขุ ยมหํ อนุปฺปตฺตา”ติ.\csromanspacer{} อถ โข ภควา ตํ เทวตํ คาถาหิ อชฺฌภาสิ\csromandash{}}

\csromangathameasure{“อกฺเขยฺยสญฺญิโน สตฺตา, อกฺเขยฺยสฺมิํ ปติฏฺฐิตา. \\ อกฺเขยฺยํ อปริญฺญาย, โยคมายนฺติ มจฺจุโน. \\ อกฺเขยฺยญฺจ ปริญฺญาย, อกฺขาตารํ น มญฺญติ. \\ ตํ หิ ตสฺส น โหตีติ, เยน นํ วชฺชา น ตสฺส อตฺถิ.}
\csromangathagroup{\csromangathastanza{“อกฺเขยฺยสญฺญิโน สตฺตา, อกฺเขยฺยสฺมิํ ปติฏฺฐิตา. \\ อกฺเขยฺยํ อปริญฺญาย, โยคมายนฺติ มจฺจุโน.}\csromangathastanza{อกฺเขยฺยญฺจ ปริญฺญาย, อกฺขาตารํ น มญฺญติ. \\ ตํ หิ ตสฺส น โหตีติ, เยน นํ วชฺชา น ตสฺส อตฺถิ.}}

\prose{สเจ วิชานาสิ วเทหิ ยกฺขา”ติ\footnote{ยกฺขีติ (อิ, ก)}.}

\prose{น ขฺวาหํ ภนฺเต อิมสฺส ภควตา สํขิตฺเตน ภาสิตสฺส วิตฺถาเรน อตฺถํ อาชานามิ, สาธุ เม ภนฺเต ภควา ตถา ภาสตุ, ยถาหํ อิมสฺส ภควตา สํขิตฺเตน ภาสิตสฺส วิตฺถาเรน อตฺถํ ชาเนยฺยนฺติ.}

\csromangathameasure{สโม วิเสสี อุท วา นิหีโน, \\ โย มญฺญตี โส วิวเทถ เตน. \\ ตีสุ วิธาสุ อวิกมฺปมาโน, \\ สโม วิเสสีติ น ตสฺส โหติ.}
\csromangathagroup{\csromangathastanza{สโม วิเสสี อุท วา\footnote{อถวา (สี, อิ)} นิหีโน, \\ โย มญฺญตี โส วิวเทถ\footnote{อถวา (สี, อิ)} เตน. \\ ตีสุ วิธาสุ อวิกมฺปมาโน, \\ สโม วิเสสีติ น ตสฺส โหติ.}}

\prose{สเจ วิชานาสิ วเทหิ ยกฺขาติ.}

\prose{อิมสฺสปิ ขฺวาหํ ภนฺเต ภควตา สํขิตฺเตน ภาสิตสฺส น วิตฺถาเรน อตฺถํ อาชานามิ, สาธุ เม ภนฺเต ภควา ตถา ภาสตุ, ยถาหํ อิมสฺส ภควตา สํขิตฺเตน ภาสิตสฺส วิตฺถาเรน อตฺถํ ชาเนยฺยนฺติ.}

\prose{ปหาสิ สงฺขํ น วิมานมชฺฌคา, อจฺเฉจฺฉิ\footnote{อจฺเฉชฺชิ (สฺยา, กํ, ก)} ตณฺหํ อิธ นามรูเป.}

\prose{ตํ ฉินฺนคนฺถํ อนิฆํ นิราสํ, ปริเยสมานา นาชฺฌคมุํ.}

\prose{เทวา มนุสฺสา อิธ วา หุรํ วา, สคฺเคสุ วา สพฺพนิเวสเนสุ.}

\prose{สเจ วิชานาสิ วเทหิ ยกฺขาติ.}

\prose{\csromanfolio{12}อิมสฺส ขฺวาหํ ภนฺเต ภควตา สํขิตฺเตน ภาสิตสฺส เอวํ วิตฺถาเรน อตฺถํ อาชานามิ.}

\csromangathameasure{ปาปํ น กยิรา วจสา มนสา, \\ กาเยน วา กิญฺจน สพฺพโลเก. \\ กาเม ปหาย สติมา สมฺปชาโน, \\ ทุกฺขํ น เสเวถ อนตฺถสํหิตนฺติ.}
\csromangathagroup{\csromangathastanza{ปาปํ น กยิรา วจสา มนสา, \\ กาเยน วา กิญฺจน สพฺพโลเก. \\ กาเม ปหาย สติมา สมฺปชาโน, \\ ทุกฺขํ น เสเวถ อนตฺถสํหิตนฺติ.}}

\vspace{\tipitakaparskip}
\csromancenter{นนฺทนวคฺโค ทุติโย.}
\titlehead{\textbf{ตสฺสุทฺทานํ}}
\csromangathameasure{นนฺทนา นนฺทติ เจว, นตฺถิปุตฺตสเมน จ. \\ ขตฺติโย สณมาโน จ, นิทฺทาตนฺที จ ทุกฺกรํ. \\ หิรี กุฏิกา นวโม, ทสโม วุตฺโต สมิทฺธินาติ.}
\csromangathagroup{\csromangathastanza{นนฺทนา นนฺทติ เจว, นตฺถิปุตฺตสเมน จ. \\ ขตฺติโย สณมาโน จ, นิทฺทาตนฺที จ ทุกฺกรํ. \\ หิรี กุฏิกา นวโม, ทสโม วุตฺโต สมิทฺธินาติ.}}

\chapterhead{3. สตฺติวคฺค}
\csromanheader{2}{1. สตฺติสุตฺต}
\proseitem{21}{สาวตฺถินิทานํ.\csromanspacer{} เอกมนฺตํ ฐิตา โข สา เทวตา ภควโต สนฺติเก}

\csromangathameasure{“สตฺติยา วิย โอมฏฺโฐ, ฑยฺหมาโนว มตฺถเก. \\ กามราคปฺปหานาย, สโต ภิกฺขุ ปริพฺพเช”ติ. \\ สตฺติยา วิย โอมฏฺโฐ, ฑยฺหมาโนว มตฺถเก. \\ สกฺกายทิฏฺฐิปฺปหานาย, สโต ภิกฺขุ ปริพฺพเชติ.}
\csromangathagroup{\csromangathastanza{“สตฺติยา วิย โอมฏฺโฐ, ฑยฺหมาโนว\footnote{qอยฺหมาเนว (สพฺพตฺถ)} มตฺถเก. \\ กามราคปฺปหานาย, สโต ภิกฺขุ ปริพฺพเช”ติ.}\csromangathastanza{สตฺติยา วิย โอมฏฺโฐ, ฑยฺหมาโนว มตฺถเก. \\ สกฺกายทิฏฺฐิปฺปหานาย, สโต ภิกฺขุ ปริพฺพเชติ.}}

\csromanheader{2}{2. ผุสติสุตฺต}
\csromanhangingbegin
\hangingheaditem{22}{นาผุสนฺตํ ผุสติ จ, ผุสนฺตญฺจ ตโต ผุเส.}
\hangingbody{ตสฺมา ผุสนฺตํ ผุสติ, อปฺปทุฏฺฐปโทสินนฺติ.}
\csromanhangingend

\csromangathameasure{โย อปฺปทุฏฺฐสฺส นรสฺส ทุสฺสติ, สุทฺธสฺส โปสสฺส อนงฺคณสฺส. \\ ตเมว พาลํ ปจฺเจติ ปาปํ, สุขุโม รโช ปฏิวาตํว ขิตฺโตติ.}
\csromangathagroup{\csromangathastanza{\csromanfolio{13}โย อปฺปทุฏฺฐสฺส นรสฺส ทุสฺสติ, สุทฺธสฺส โปสสฺส อนงฺคณสฺส. \\ ตเมว พาลํ ปจฺเจติ ปาปํ, สุขุโม รโช ปฏิวาตํว ขิตฺโตติ.}}

\csromanheader{2}{3. ชฏาสุตฺต}
\csromanhangingbegin
\hangingheaditem{23}{อนฺโตชฏา พหิชฏา, ชฏาย ชฏิตา ปชา.}
\hangingbody{ตํ ตํ โคตม ปุจฺฉามิ, โก อิมํ วิชฏเย ชฏนฺติ.}
\hangingbody{สีเล ปติฏฺฐาย นโร สปญฺโญ, จิตฺตํ ปญฺญญฺจ ภาวยํ.}
\hangingbody{อาตาปี นิปโก ภิกฺขุ, โส อิมํ วิชฏเย ชฏํ.}
\hangingbody{เยสํ ราโค จ โทโส จ, อวิชฺชา จ วิราชิตา.}
\hangingbody{ขีณาสวา อรหนฺโต, เตสํ วิชฏิตา ชฏา.}
\hangingbody{ยตฺถ นามญฺจ รูปญฺจ, อเสสํ อุปรุชฺฌติ.}
\hangingbody{ปฏิฆํ รูปสญฺญา จ, เอตฺเถสา ฉิชฺชเต1 ชฏาติ.}
\csromanhangingend

\csromanheader{2}{4. มโนนิวารณสุตฺต}
\csromanhangingbegin
\hangingheaditem{24}{ยโต ยโต มโน นิวารเย,}
\hangingbody{น ทุกฺขเมติ นํ ตโต ตโต.}
\hangingbody{ส สพฺพโต มโน นิวารเย,}
\hangingbody{ส สพฺพโต ทุกฺขา ปมุจฺจติ.}
\hangingbody{น สพฺพโต มโน นิวารเย,}
\hangingbody{น มโน สํยตตฺตมาคตํ.}
\hangingbody{ยโต ยโต จ ปาปกํ,}
\hangingbody{ตโต ตโต มโน นิวารเยติ.}
\csromanhangingend

\csromanheader{2}{5. อรหนฺตสุตฺต}
\csromanhangingbegin
\hangingheaditem{25}{โย โหติ ภิกฺขุ อรหํ กตาวี,}
\hangingbody{ขีณาสโว อนฺติมเทหธารี.}
\hangingbody{อหํ วทามีติปิ โส วเทยฺย,}
\hangingbody{มมํ วทนฺตีติปิ โส วเทยฺยาติ.}
\csromanhangingend

\csromangathameasure{โย โหติ ภิกฺขุ อรหํ กตาวี, \\ ขีณาสโว อนฺติมเทหธารี. \\ อหํ วทามีติปิ โส วเทยฺย, \\ มมํ วทนฺตีติปิ โส วเทยฺย. \\ โลเก สมญฺญํ กุสโล วิทิตฺวา, \\ โวหารมตฺเตน โส โวหเรยฺยาติ. \\ โย โหติ ภิกฺขุ อรหํ กตาวี, \\ ขีณาสโว อนฺติมเทหธารี. \\ มานํ นุ โข โส อุปคมฺม ภิกฺขุ, \\ อหํ วทามีติปิ โส วเทยฺย. \\ มมํ วทนฺตีติปิ โส วเทยฺยาติ. \\ ปหีนมานสฺส น สนฺติ คนฺถา, \\ วิธูปิตา มานคนฺถสฺส สพฺเพ. \\ ส วีติวตฺโต มญฺญตํ สุเมโธ, \\ อหํ วทามีติปิ โส วเทยฺย. \\ มมํ วทนฺตีติปิ โส วเทยฺย, \\ โลเก สมญฺญํ กุสโล วิทิตฺวา. \\ โวหารมตฺเตน โส โวหเรยฺยาติ.}
\csromangathagroup{\csromangathastanza{\csromanfolio{14}โย โหติ ภิกฺขุ อรหํ กตาวี, \\ ขีณาสโว อนฺติมเทหธารี. \\ อหํ วทามีติปิ โส วเทยฺย, \\ มมํ วทนฺตีติปิ โส วเทยฺย.}\csromangathastanza{โลเก สมญฺญํ กุสโล วิทิตฺวา, \\ โวหารมตฺเตน โส\footnote{ส (?)} โวหเรยฺยาติ. \\ โย โหติ ภิกฺขุ อรหํ กตาวี, \\ ขีณาสโว อนฺติมเทหธารี.}\csromangathastanza{มานํ นุ โข โส อุปคมฺม ภิกฺขุ, \\ อหํ วทามีติปิ โส วเทยฺย. \\ มมํ วทนฺตีติปิ โส วเทยฺยาติ. \\ ปหีนมานสฺส น สนฺติ คนฺถา,}\csromangathastanza{วิธูปิตา มานคนฺถสฺส สพฺเพ. \\ ส วีติวตฺโต มญฺญตํ\footnote{มานนํ (สี), มญฺญิตํ (?)} สุเมโธ, \\ อหํ วทามีติปิ โส วเทยฺย. \\ มมํ วทนฺตีติปิ โส วเทยฺย, \\ โลเก สมญฺญํ กุสโล วิทิตฺวา. \\ โวหารมตฺเตน โส โวหเรยฺยาติ.}}

\csromanheader{2}{6. ปชฺโชตสุตฺต}
\csromanhangingbegin
\hangingheaditem{26}{กติ โลกสฺมิํ ปชฺโชตา, เยหิ โลโก ปกาสติ\footnote{ปภาสติ (ก-สี)}.}
\hangingbody{ภควนฺตํ4 ปุฏฺฐุมาคมฺม, กถํ ชาเนมุ ตํ มยนฺติ.}
\hangingbody{จตฺตาโร โลเก ปชฺโชตา, ปญฺจเมตฺถ น วิชฺชติ.}
\hangingbody{ทิวา ตปติ อาทิจฺโจ, รตฺติมาภาติ จนฺทิมา.}
\hangingbody{อถ อคฺคิ ทิวารตฺติํ, ตตฺถ ตตฺถ ปกาสติ.}
\hangingbody{สมฺพุทฺโธ ตปตํ เสฏฺโฐ, เอสา อาภา อนุตฺตราติ.}
\csromanhangingend

\csromanheader{2}{7. สรสุตฺต}
\csromanhangingbegin
\hangingheaditem{27}{\csromanfolio{15}กุโต สรา นิวตฺตนฺติ, กตฺถ วฏฺฏํ น วตฺตติ.}
\hangingbody{กตฺถ นามญฺจ รูปญฺจ, อเสสํ อุปรุชฺฌตีติ.}
\hangingbody{ยตฺถ อาโป จ ปถวี, เตโช วาโย น คาธติ.}
\hangingbody{อโต สรา นิวตฺตนฺติ, เอตฺถ วฏฺฏํ น วตฺตติ.}
\hangingbody{เอตฺถ นามญฺจ รูปญฺจ, อเสสํ อุปรุชฺฌตีติ.}
\csromanhangingend

\csromanheader{2}{8. มหทฺธนสุตฺต}
\csromanhangingbegin
\hangingheaditem{28}{มหทฺธนา มหาโภคา, รฏฺฐวนฺโตปิ ขตฺติยา.}
\hangingbody{อญฺญมญฺญาภิคิชฺฌนฺติ, กาเมสุ อนลงฺกตา.}
\hangingbody{เตสุ อุสฺสุกฺกชาเตสุ, ภวโสตานุสาริสุ.}
\hangingbody{เกธ ตณฺหํ1 ปชหิํสุ2, เก โลกสฺมิํ อนุสฺสุกาติ.}
\hangingbody{หิตฺวา อคารํ ปพฺพชิตา, หิตฺวา ปุตฺตํ ปสุํ ปิยํ.}
\hangingbody{หิตฺวา ราคญฺจ โทสญฺจ, อวิชฺชญฺจ วิราชิย.}
\hangingbody{ขีณาสวา อรหนฺโต, เต โลกสฺมิํ อนุสฺสุกาติ.}
\csromanhangingend

\csromanheader{2}{9. จตุจกฺกสุตฺต}
\csromanhangingbegin
\hangingheaditem{29}{จตุจกฺกํ นวทฺวารํ, ปุณฺณํ โลเภน สํยุตํ.}
\hangingbody{ปงฺกชาตํ มหาวีร, กถํ ยาตฺรา ภวิสฺสตีติ.}
\hangingbody{เฉตฺวา นทฺธิํ วรตฺตญฺจ, อิจฺฉา โลภญฺจ ปาปกํ.}
\hangingbody{สมูลํ ตณฺหมพฺพุยฺห, เอวํ ยาตฺรา ภวิสฺสตีติ.}
\csromanhangingend

\csromanheader{2}{10. เอณิชงฺฆสุตฺต}
\csromanhangingbegin
\hangingheaditem{30}{เอณิชงฺฆํ กิสํ วีรํ, อปฺปาหารํ อโลลุปํ.}
\hangingbody{สีหํเวกจรํ นาคํ, กาเมสุ อนเปกฺขินํ.}
\hangingbody{อุปสงฺกมฺม ปุจฺฉาม, กถํ ทุกฺขา ปมุจฺจตีติ.}
\csromanhangingend

\csromangathameasure{ปญฺจกามคุณา โลเก, มโนฉฏฺฐา ปเวทิตา. \\ เอตฺถ ฉนฺทํ วิราเชตฺวา, เอวํ ทุกฺขา ปมุจฺจตีติ.}
\csromangathagroup{\csromangathastanza{\csromanfolio{16}ปญฺจกามคุณา โลเก, มโนฉฏฺฐา ปเวทิตา. \\ เอตฺถ ฉนฺทํ วิราเชตฺวา, เอวํ ทุกฺขา ปมุจฺจตีติ.}}

\vspace{\tipitakaparskip}
\csromancenter{สตฺติวคฺโค ตติโย.}
\titlehead{\textbf{ตสฺสุทฺทานํ}}
\csromangathameasure{สตฺติยา ผุสติ เจว, ชฏา มโนนิวารณา. \\ อรหนฺเตน ปชฺโชโต, สรา มหทฺธเนน จ. \\ จตุจกฺเกน นวมํ, เอณิชงฺเฆน เต ทสาติ.}
\csromangathagroup{\csromangathastanza{สตฺติยา ผุสติ เจว, ชฏา มโนนิวารณา. \\ อรหนฺเตน ปชฺโชโต, สรา มหทฺธเนน จ. \\ จตุจกฺเกน นวมํ, เอณิชงฺเฆน เต ทสาติ.}}

\chapterhead{4. สตุลฺลปกายิกวคฺค}
\csromanheader{2}{1. สพฺภิสุตฺต}
\proseitem{31}{เอวํ เม สุตํ\csromandash{}เอกํ สมยํ ภควา สาวตฺถิยํ วิหรติ เชตวเน อนาถปิณฺฑิกสฺส อาราเม.\csromanspacer{} อถ โข สมฺพหุลา สตุลฺลปกายิกา เทวตาโย อภิกฺกนฺตาย รตฺติยา อภิกฺกนฺตวณฺณา เกวลกปฺปํ เชตวนํ โอภาเสตฺวา เยน ภควา เตนุปสงฺกมิํสุ, อุปสงฺกมิตฺวา ภควนฺตํ อภิวาเทตฺวา เอกมนฺตํ อฏฺฐํสุ, เอกมนฺตํ ฐิตา โข เอกา เทวตา ภควโต สนฺติเก อิมํ คาถํ อภาสิ\csromandash{}}

\csromangathameasure{“สพฺภิเรว สมาเสถ, สพฺภิ กุพฺเพถ สนฺถวํ. \\ สตํ สทฺธมฺมมญฺญาย, เสยฺโย โหติ น ปาปิโย”ติ.}
\csromangathagroup{\csromangathastanza{“สพฺภิเรว สมาเสถ, สพฺภิ กุพฺเพถ\footnote{กฺรุพฺเพถ (ก)} สนฺถวํ. \\ สตํ สทฺธมฺมมญฺญาย, เสยฺโย โหติ น ปาปิโย”ติ.}}

\csromanheader{2}{อถ โข อปรา เทวตา ภควโต สนฺติเก อิมํ คาถํ อภาสิ\csromandash{}}
\csromangathameasure{“สพฺภิเรว สมาเสถ, สพฺภิ กุพฺเพถ สนฺถวํ. \\ สตํ สทฺธมฺมมญฺญาย, ปญฺญา ลพฺภติ นาญฺญโต”ติ.}
\csromangathagroup{\csromangathastanza{“สพฺภิเรว สมาเสถ, สพฺภิ กุพฺเพถ สนฺถวํ. \\ สตํ สทฺธมฺมมญฺญาย, ปญฺญา ลพฺภติ\footnote{ปญฺญํ ลภติ (สฺยา, กํ)} นาญฺญโต”ติ.}}

\csromanheader{2}{อถ โข อปรา เทวตา ภควโต สนฺติเก อิมํ คาถํ อภาสิ\csromandash{}}
\csromangathameasure{“สพฺภิเรว สมาเสถ, สพฺภิ กุพฺเพถ สนฺถวํ. \\ สตํ สทฺธมฺมมญฺญาย, โสกมชฺเฌ น โสจตี”ติ.}
\csromangathagroup{\csromangathastanza{“สพฺภิเรว สมาเสถ, สพฺภิ กุพฺเพถ สนฺถวํ. \\ สตํ สทฺธมฺมมญฺญาย, โสกมชฺเฌ น โสจตี”ติ.}}

\csromanheader{2}{อถ โข อปรา เทวตา ภควโต สนฺติเก อิมํ คาถํ อภาสิ\csromandash{}}
\csromangathameasure{“สพฺภิเรว สมาเสถ, สพฺภิ กุพฺเพถ สนฺถวํ. \\ สตํ สทฺธมฺมมญฺญาย, ญาติมชฺเฌ วิโรจตี”ติ.}
\csromangathagroup{\csromangathastanza{\csromanfolio{17}“สพฺภิเรว สมาเสถ, สพฺภิ กุพฺเพถ สนฺถวํ. \\ สตํ สทฺธมฺมมญฺญาย, ญาติมชฺเฌ วิโรจตี”ติ.}}

\csromanheader{2}{อถ โข อปรา เทวตา ภควโต สนฺติเก อิมํ คาถํ อภาสิ\csromandash{}}
\csromangathameasure{“สพฺภิเรว สมาเสถ, สพฺภิ กุพฺเพถ สนฺถวํ. \\ สตํ สทฺธมฺมมญฺญาย, สตฺตา คจฺฉนฺติ สุคฺคตินฺ”ติ.}
\csromangathagroup{\csromangathastanza{“สพฺภิเรว สมาเสถ, สพฺภิ กุพฺเพถ สนฺถวํ. \\ สตํ สทฺธมฺมมญฺญาย, สตฺตา คจฺฉนฺติ สุคฺคตินฺ”ติ.}}

\csromanheader{2}{อถ โข อปรา เทวตา ภควโต สนฺติเก อิมํ คาถํ อภาสิ\csromandash{}}
\csromangathameasure{“สพฺภิเรว สมาเสถ, สพฺภิ กุพฺเพถ สนฺถวํ. \\ สตํ สทฺธมฺมมญฺญาย, สตฺตา ติฏฺฐนฺติ สาตตนฺ”ติ.}
\csromangathagroup{\csromangathastanza{“สพฺภิเรว สมาเสถ, สพฺภิ กุพฺเพถ สนฺถวํ. \\ สตํ สทฺธมฺมมญฺญาย, สตฺตา ติฏฺฐนฺติ สาตตนฺ”ติ.}}

\prose{อถ โข อปรา เทวตา ภควนฺตํ เอตทโวจ “กสฺส นุ โข ภควา สุภาสิตนฺ”ติ.\csromanspacer{} สพฺพาสํ โว สุภาสิตํ ปริยาเยน, อปิ จ มมปิ สุณาถ\csromandash{}}

\csromangathameasure{สพฺภิเรว สมาเสถ, สพฺภิ กุพฺเพถ สนฺถวํ. \\ สตํ สทฺธมฺมมญฺญาย, สพฺพทุกฺขา ปมุจฺจตีติ.}
\csromangathagroup{\csromangathastanza{สพฺภิเรว สมาเสถ, สพฺภิ กุพฺเพถ สนฺถวํ. \\ สตํ สทฺธมฺมมญฺญาย, สพฺพทุกฺขา ปมุจฺจตีติ.}}

\prose{อิทมโวจ ภควา.\csromanspacer{} อตฺตมนา ตา เทวตาโย ภควนฺตํ อภิวาเทตฺวา ปทกฺขิณํ กตฺวา ตตฺเถวนฺตรธายิํสูติ.}

\csromanheader{2}{2. \textbf{มจฺฉริสุตฺต}}
\proseitem{32}{เอกํ สมยํ ภควา สาวตฺถิยํ วิหรติ เชตวเน อนาถปิณฺฑิกสฺส อาราเม.\csromanspacer{} อถ โข สมฺพหุลา สตุลฺลปกายิกา เทวตาโย อภิกฺกนฺตาย รตฺติยา อภิกฺกนฺตวณฺณา เกวลกปฺปํ เชตวนํ โอภาเสตฺวา เยน ภควา เตนุปสงฺกมิํสุ, อุปสงฺกมิตฺวา ภควนฺตํ อภิวาเทตฺวา เอกมนฺตํ อฏฺฐํสุ, เอกมนฺตํ ฐิตา โข เอกา เทวตา ภควโต สนฺติเก}

\csromangathameasure{“มจฺเฉรา จ ปมาทา จ, เอวํ ทานํ น ทียติ. \\ ปุญฺญํ อากงฺขมาเนน, เทยฺยํ โหติ วิชานตา”ติ.}
\csromangathagroup{\csromangathastanza{“มจฺเฉรา จ ปมาทา จ, เอวํ ทานํ น ทียติ\footnote{ทิยฺยติ (ก)}. \\ ปุญฺญํ อากงฺขมาเนน, เทยฺยํ โหติ วิชานตา”ติ.}}

\csromanheader{2}{อถ โข อปรา เทวตา ภควโต สนฺติเก อิมา คาถาโย อภาสิ\csromandash{}}
\csromangathameasure{“ยสฺเสว ภีโต น ททาติ มจฺฉรี, ตเทวาททโต ภยํ. \\ ชิฆจฺฉา จ ปิปาสา จ, ยสฺส ภายติ มจฺฉรี. \\ ตเมว พาลํ ผุสติ, อสฺมิํ โลเก ปรมฺหิ จ. \\ ตสฺมา วิเนยฺย มจฺเฉรํ, ทชฺชา ทานํ มลาภิภู. \\ ปุญฺญานิ ปรโลกสฺมิํ, ปติฏฺฐา โหนฺติ ปาณินนฺ”ติ.}
\csromangathagroup{\csromangathastanza{\csromanfolio{18}“ยสฺเสว ภีโต น ททาติ มจฺฉรี, ตเทวาททโต ภยํ. \\ ชิฆจฺฉา จ ปิปาสา จ, ยสฺส ภายติ มจฺฉรี.}\csromangathastanza{ตเมว พาลํ ผุสติ, อสฺมิํ โลเก ปรมฺหิ จ. \\ ตสฺมา วิเนยฺย มจฺเฉรํ, ทชฺชา ทานํ มลาภิภู. \\ ปุญฺญานิ ปรโลกสฺมิํ, ปติฏฺฐา โหนฺติ ปาณินนฺ”ติ.}}

\csromanheader{2}{อถ โข อปรา เทวตา ภควโต สนฺติเก อิมา คาถาโย อภาสิ\csromandash{}}
\csromangathameasure{“เต มเตสุ น มียนฺติ, ปนฺถานํว สหพฺพชํ. \\ อปฺปสฺมิํ เย ปเวจฺฉนฺติ, เอส ธมฺโม สนนฺตโน. \\ อปฺปสฺเมเก ปเวจฺฉนฺติ, พหุเนเก น ทิจฺฉเร. \\ อปฺปสฺมา ทกฺขิณา ทินฺนา, สหสฺเสน สมํ มิตา”ติ.}
\csromangathagroup{\csromangathastanza{“เต มเตสุ น มียนฺติ, ปนฺถานํว สหพฺพชํ. \\ อปฺปสฺมิํ เย ปเวจฺฉนฺติ, เอส ธมฺโม สนนฺตโน.}\csromangathastanza{อปฺปสฺเมเก ปเวจฺฉนฺติ, พหุเนเก น ทิจฺฉเร. \\ อปฺปสฺมา ทกฺขิณา ทินฺนา, สหสฺเสน สมํ มิตา”ติ.}}

\csromanheader{2}{อถ โข อปรา เทวตา ภควโต สนฺติเก อิมา คาถาโย อภาสิ\csromandash{}}
\csromangathameasure{“ทุทฺททํ ททมานานํ, ทุกฺกรํ กมฺม กุพฺพตํ. \\ อสนฺโต นานุกุพฺพนฺติ, สตํ ธมฺโม ทุรนฺวโย. \\ ตสฺมา สตญฺจ อสตํ, นานา โหติ อิโต คติ. \\ อสนฺโต นิรยํ ยนฺติ, สนฺโต สคฺคปรายนา”ติ. \\ อถ โข อปรา เทวตา ภควโต สนฺติเก เอตทโวจ “กสฺส นุ โข ภควา สุภาสิตนฺ”ติ. สพฺพาสํ โว สุภาสิตํ ปริยาเยน, อปิ จ มมปิ สุณาถ\csromandash{} \\ ธมฺมํ จเร โยปิ สมุญฺชกํ จเร, \\ ทารญฺจ โปสํ ททมปฺปกสฺมิํ. \\ สตํ สหสฺสานํ สหสฺสยาคินํ, \\ กลมฺปิ นาคฺฆนฺติ ตถาวิธสฺส เตติ.}
\csromangathagroup{\csromangathastanza{“ทุทฺททํ ททมานานํ, ทุกฺกรํ กมฺม กุพฺพตํ. \\ อสนฺโต นานุกุพฺพนฺติ, สตํ ธมฺโม ทุรนฺวโย\footnote{ทุรนฺนโย (สี)}.}\csromangathastanza{ตสฺมา สตญฺจ อสตํ\footnote{อสตญฺจ (สี, สฺยา, กํ)}, นานา โหติ อิโต คติ. \\ อสนฺโต นิรยํ ยนฺติ, สนฺโต สคฺคปรายนา”ติ.}\csromangathastanza{อถ โข อปรา เทวตา ภควโต สนฺติเก เอตทโวจ “กสฺส นุ โข ภควา สุภาสิตนฺ”ติ.\csromanspacer{} สพฺพาสํ โว สุภาสิตํ ปริยาเยน, อปิ จ มมปิ สุณาถ\csromandash{} \\ ธมฺมํ จเร โยปิ สมุญฺชกํ จเร, \\ ทารญฺจ โปสํ ททมปฺปกสฺมิํ. \\ สตํ สหสฺสานํ สหสฺสยาคินํ,}\csromangathastanza{กลมฺปิ นาคฺฆนฺติ ตถาวิธสฺส เตติ.}}

\csromanheader{2}{อถ โข อปรา เทวตา ภควนฺตํ คาถาย อชฺฌภาสิ\csromandash{}}
\csromangathameasure{“เกเนส ยญฺโญ วิปุโล มหคฺคโต, \\ สเมน ทินฺนสฺส น อคฺฆเมติ. \\ กถํ สตํ สหสฺสานํ สหสฺสยาคินํ, \\ กลมฺปิ นาคฺฆนฺติ ตถาวิธสฺส เต”ติ. \\ ททนฺติ เหเก วิสเม นิวิฏฺฐา, \\ เฉตฺวา วธิตฺวา อถ โสจยิตฺวา. \\ สา ทกฺขิณา อสฺสุมุขา สทณฺฑา, \\ สเมน ทินฺนสฺส น อคฺฆเมติ. \\ เอวํ สตํ สหสฺสานํ สหสฺสยาคินํ, \\ กลมฺปิ นาคฺฆนฺติ ตถาวิธสฺส เตติ.}
\csromangathagroup{\csromangathastanza{\csromanfolio{19}“เกเนส ยญฺโญ วิปุโล มหคฺคโต, \\ สเมน ทินฺนสฺส น อคฺฆเมติ. \\ กถํ\footnote{อิทํ ปทํ กตฺถจิ สีหฬ-โปตฺถเก นตฺถิ.} สตํ สหสฺสานํ สหสฺสยาคินํ, \\ กลมฺปิ นาคฺฆนฺติ ตถาวิธสฺส เต”ติ.}\csromangathastanza{ททนฺติ เหเก วิสเม นิวิฏฺฐา, \\ เฉตฺวา วธิตฺวา อถ โสจยิตฺวา. \\ สา ทกฺขิณา อสฺสุมุขา สทณฺฑา, \\ สเมน ทินฺนสฺส น อคฺฆเมติ. \\ เอวํ สตํ สหสฺสานํ สหสฺสยาคินํ, \\ กลมฺปิ นาคฺฆนฺติ ตถาวิธสฺส เตติ.}}

\csromanheader{2}{3. \textbf{สาธุสุตฺต}}
\proseitem{33}{สาวตฺถินิทานํ.\csromanspacer{} อถ โข สมฺพหุลา สตุลฺลปกายิกา เทวตาโย อภิกฺกนฺตาย รตฺติยา อภิกฺกนฺตวณฺณา เกวลกปฺปํ เชตวนํ โอภาเสตฺวา เยน ภควา เตนุปสงฺกมิํสุ, อุปสงฺกมิตฺวา ภควนฺตํ อภิวาเทตฺวา เอกมนฺตํ อฏฺฐํสุ, เอกมนฺตํ ฐิตา โข เอกา เทวตา ภควโต สนฺติเก อิมํ อุทานํ อุทาเนสิ\csromandash{}}

\prose{“สาธุ โข มาริส ทานํ.}

\csromangathameasure{มจฺเฉรา จ ปมาทา จ, เอวํ ทานํ น ทียติ. \\ ปุญฺญํ อากงฺขมาเนน, เทยฺยํ โหติ วิชานตา”ติ.}
\csromangathagroup{\csromangathastanza{มจฺเฉรา จ ปมาทา จ, เอวํ ทานํ น ทียติ. \\ ปุญฺญํ อากงฺขมาเนน, เทยฺยํ โหติ วิชานตา”ติ.}}

\csromanheader{2}{อถ โข อปรา เทวตา ภควโต สนฺติเก อิมํ อุทานํ อุทาเนสิ\csromandash{}}
\csromangathameasure{“สาธุ โข มาริส ทานํ, อปิ จ อปฺปกสฺมิมฺปิ สาหุ ทานํ. \\ อปฺปสฺเมเก ปเวจฺฉนฺติ, พหุเนเก น ทิจฺฉเร. \\ อปฺปสฺมา ทกฺขิณา ทินฺนา, สหสฺเสน สมํ มิตา”ติ.}
\csromangathagroup{\csromangathastanza{“สาธุ โข มาริส ทานํ, อปิ จ อปฺปกสฺมิมฺปิ สาหุ ทานํ. \\ อปฺปสฺเมเก ปเวจฺฉนฺติ, พหุเนเก น ทิจฺฉเร. \\ อปฺปสฺมา ทกฺขิณา ทินฺนา, สหสฺเสน สมํ มิตา”ติ.}}

\csromanheader{2}{อถ โข อปรา เทวตา ภควโต สนฺติเก อิมํ อุทานํ อุทาเนสิ\csromandash{}}
\csromangathameasure{“สาธุ โข มาริส ทานํ, อปฺปกสฺมิมฺปิ สาหุ ทานํ.}
\csromangathagroup{\csromangathastanza{\csromanfolio{20}“สาธุ โข มาริส ทานํ, อปฺปกสฺมิมฺปิ สาหุ ทานํ.}}

\prose{อปิ จ สทฺธายปิ สาหุ ทานํ.}

\csromangathameasure{ทานญฺจ ยุทฺธญฺจ สมานมาหุ, \\ อปฺปาปิ สนฺตา พหุเก ชินนฺติ. \\ อปฺปมฺปิ เจ สทฺทหาโน ททาติ, \\ เตเนว โส โหติ สุขี ปรตฺถา”ติ.}
\csromangathagroup{\csromangathastanza{ทานญฺจ ยุทฺธญฺจ สมานมาหุ, \\ อปฺปาปิ สนฺตา พหุเก ชินนฺติ. \\ อปฺปมฺปิ เจ สทฺทหาโน ททาติ, \\ เตเนว โส โหติ สุขี ปรตฺถา”ติ.}}

\csromanheader{2}{อถ โข อปรา เทวตา ภควโต สนฺติเก อิมํ อุทานํ อุทาเนสิ\csromandash{}}
\csromangathameasure{“สาธุ โข มาริส ทานํ, อปฺปกสฺมิมฺปิ สาหุ ทานํ. \\ สทฺธายปิ สาหุ ทานํ, อปิ จ ธมฺมลทฺธสฺสาปิ สาหุ ทานํ. \\ โย ธมฺมลทฺธสฺส ททาติ ทานํ, \\ อุฏฺฐานวีริยาธิคตสฺส ชนฺตุ. \\ อติกกมฺม โส เวตรณิํ ยมสฺส, \\ ทิพฺพานิ ฐานานิ อุเปติ มจฺโจ”ติ.}
\csromangathagroup{\csromangathastanza{“สาธุ โข มาริส ทานํ, อปฺปกสฺมิมฺปิ สาหุ ทานํ. \\ สทฺธายปิ สาหุ ทานํ, อปิ จ ธมฺมลทฺธสฺสาปิ สาหุ ทานํ.}\csromangathastanza{โย ธมฺมลทฺธสฺส ททาติ ทานํ, \\ อุฏฺฐานวีริยาธิคตสฺส ชนฺตุ. \\ อติกกมฺม โส เวตรณิํ ยมสฺส, \\ ทิพฺพานิ ฐานานิ อุเปติ มจฺโจ”ติ.}}

\csromanheader{2}{อถ โข อปรา เทวตา ภควโต สนฺติเก อิมํ อุทานํ อุทาเนสิ\csromandash{}}
\csromangathameasure{“สาธุ โข มาริส ทานํ, อปฺปกสฺมิมฺปิ สาหุ ทานํ. \\ สทฺธายปิ สาหุ ทานํ, ธมฺมลทฺธสฺสาปิ สาหุ ทานํ.}
\csromangathagroup{\csromangathastanza{“สาธุ โข มาริส ทานํ, อปฺปกสฺมิมฺปิ สาหุ ทานํ. \\ สทฺธายปิ สาหุ ทานํ, ธมฺมลทฺธสฺสาปิ สาหุ ทานํ.}}

\prose{อปิ จ วิเจยฺย ทานมฺปิ สาหุ ทานํ.}

\csromangathameasure{วิเจยฺย ทานํ สุคตปฺปสตฺถํ, เย ทกฺขิเณยฺยา อิธ ชีวโลเก. \\ เอเตสุ ทินฺนานิ มหปฺผลานิ, พีชานิ วุตฺตานิ ยถา สุเขตฺเต”ติ.}
\csromangathagroup{\csromangathastanza{วิเจยฺย ทานํ สุคตปฺปสตฺถํ, เย ทกฺขิเณยฺยา อิธ ชีวโลเก. \\ เอเตสุ ทินฺนานิ มหปฺผลานิ, พีชานิ วุตฺตานิ ยถา สุเขตฺเต”ติ.}}

\csromanheader{2}{อถ โข อปรา เทวตา ภควโต สนฺติเก อิมํ อุทานํ อุทาเนสิ\csromandash{}}
\csromangathameasure{“สาธุ โข มาริส ทานํ, อปฺปกสฺมิมฺปิ สาหุ ทานํ. \\ สทฺธายปิ สาหุ ทานํ, ธมฺมลทฺธสฺสาปิ สาหุ ทานํ. \\ วิเจยฺย ทานมฺปิ สาหุ ทานํ, อปิ จ ปาเณสุปิ สาธุ สํยโม. \\ โย ปาณภูตานิ อเหฐยํ จรํ, \\ ปรูปวาทา น กโรนฺติ ปาปํ. \\ ภีรุํ ปสํสนฺติ น หิ ตตฺถ สูรํ, \\ ภยา หิ สนฺโต น กโรนฺติ ปาปนฺ”ติ.}
\csromangathagroup{\csromangathastanza{“สาธุ โข มาริส ทานํ, อปฺปกสฺมิมฺปิ สาหุ ทานํ. \\ สทฺธายปิ สาหุ ทานํ, ธมฺมลทฺธสฺสาปิ สาหุ ทานํ. \\ วิเจยฺย ทานมฺปิ สาหุ ทานํ, อปิ จ ปาเณสุปิ สาธุ สํยโม.}\csromangathastanza{\csromanfolio{21}โย ปาณภูตานิ\footnote{ปาณภูเตสุ (สี, อิ)} อเหฐยํ จรํ, \\ ปรูปวาทา น กโรนฺติ ปาปํ. \\ ภีรุํ ปสํสนฺติ น หิ ตตฺถ สูรํ, \\ ภยา หิ สนฺโต น กโรนฺติ ปาปนฺ”ติ.}}

\prose{อถ โข อปรา เทวตา ภควนฺตํ เอตทโวจ “กสฺส นุ โข ภควา สุภาสิตนฺ”ติ.\csromanspacer{} สพฺพาสํ โว สุภาสิตํ ปริยาเยน, อปิ จ มมปิ สุณาถ\csromandash{} สทฺธา หิ ทานํ พหุธา ปสตฺถํ, ทานา จ โข ธมฺมปทํว เสยฺโย.\csromanspacer{} ปุพฺเพ จ หิ ปุพฺพตเร จ สนฺโต, นิพฺพานเมวชฺฌคมุํ สปญฺญาติ.}

\csromanheader{2}{4. น สนฺติสุตฺต}
\proseitem{34}{เอกํ สมยํ ภควา สาวตฺถิยํ วิหรติ เชตวเน อนาถปิณฺฑิกสฺส อาราเม.\csromanspacer{} อถ โข สมฺพหุลา สตุลฺลปกายิกา เทวตาโย อภิกฺกนฺตาย รตฺติยา อภิกฺกนฺตวณฺณา เกวลกปฺปํ เชตวนํ โอภาเสตฺวา เยน ภควา เตนุปสงฺกมิํสุ, อุปสงฺกมิตฺวา ภควนฺตํ อภิวาเทตฺวา เอกมนฺตํ อฏฺฐํสุ, เอกมนฺตํ ฐิตา โข เอกา เทวตา ภควโต สนฺติเก}

\prose{“น สนฺติ กามา มนุเชสุ นิจฺจา,}

\prose{สนฺตีธ กมนียานิ เยสุ\footnote{กาเมสุ (ก)} พทฺโธ.}

\prose{เยสุ ปมตฺโต อปุนาคมนํ,}

\prose{อนาคนฺตา ปุริโส มจฺจุเธยฺยา”ติ.}

\prose{ฉนฺทชํ อฆํ ฉนฺทชํ ทุกฺขํ, ฉนฺทวินยา อฆวินโย.}

\prose{อฆวินยา ทุกฺขวินโยติ.}

\prose{น เต กามา ยานิ จิตฺรานิ โลเก,}

\prose{สงฺกปฺปราโค ปุริสสฺส กาโม.}

\prose{ติฏฺฐนฺติ จิตฺรานิ ตเถว โลเก,}

\prose{อเถตฺถ ธีรา วินยนฺติ ฉนฺทํ.}

\prose{\csromanfolio{22}โกธํ ชเห วิปฺปชเหยฺย มานํ, สํโยชนํ สพฺพมติกฺกเมยฺย.}

\prose{ตํ นามรูปสฺมิมสชฺชมานํ, อกิญฺจนํ นานุปตนฺติ ทุกฺขา.}

\prose{ปหาสิ สงฺขํ น วิมานมชฺฌคา\footnote{น จ มานมชฺฌคา (ก-สี), น วิมานมาคา (สฺยา, กํ)}, อจฺเฉจฺฉิ ตณฺหํ อิธ นามรูเป.}

\prose{ตํ ฉินฺนคนฺถํ อนิฆํ นิราสํ, ปริเยสมานา นาชฺฌคมุํ.}

\prose{เทวา มนุสฺสา อิธ วา หุรํ วา, สคฺเคสุ วา สพฺพนิเวสเนสูติ.}

\prose{ตํ เจ หิ นาทฺทกฺขุํ ตถาวิมุตฺตํ, (อิจฺจายสฺมา โมฆราชา.)}

\prose{เทวา มนุสฺสา อิธ วา หุรํ วา.}

\prose{นรุตฺตมํ อตฺถจรํ นรานํ,}

\prose{เย ตํ นมสฺสนฺติ ปสํสิยา เตติ.}

\prose{ปสํสิยา เตปิ ภวนฺติ ภิกฺขู, (โมฆราชาติ ภควา.)}

\prose{เย ตํ นมสฺสนฺติ ตถาวิมุตฺตํ,}

\prose{อญฺญาย ธมฺมํ วิจิกิจฺฉํ ปหาย,}

\prose{สงฺคาติคา เตปิ ภวนฺติ ภิกฺขูติ.}

\csromanheader{2}{5. อุชฺฌานสญฺญิสุตฺต}
\proseitem{35}{เอกํ สมยํ ภควา สาวตฺถิยํ วิหรติ เชตวเน อนาถปิณฺฑิกสฺส อาราเม.\csromanspacer{} อถ โข สมฺพหุลา อุชฺฌานสญฺญิกา เทวตาโย อภิกฺกนฺตาย รตฺติยา อภิกฺกนฺตวณฺณา เกวลกปฺปํ เชตวนํ โอภาเสตฺวา เยน ภควา เตนุปสงฺกมิํสุ, อุปสงฺกมิตฺวา เวหาสํ อฏฺฐํสุ, เวหาสํ ฐิตา โข เอกา เทวตา ภควโต สนฺติเก อิมํ คาถํ อภาสิ\csromandash{}}

\prose{“อญฺญถา สนฺตมตฺตานํ, อญฺญถา โย ปเวทเย.}

\prose{นิกจฺจ กิตวสฺเสว, ภุตฺตํ เถยฺเยน ตสฺส ตํ.}

\prose{ยํ หิ กยิรา ตํ หิ วเท, ยํ น กยิรา น ตํ วเท.}

\prose{อกโรนฺตํ ภาสมานํ, ปริชานนฺติ ปณฺฑิตา”ติ.}

\prose{น ยิทํ ภาสิตมตฺเตน, เอกนฺตสวเนน วา.}

\prose{อนุกฺกมิตเว สกฺกา, ยายํ ปฏิปทา ทฬฺหา.}

\prose{ยาย ธีรา ปมุจฺจนฺติ, ฌายิโน มารพนฺธนา.}

\prose{\csromanfolio{23}น เว ธีรา ปกุพฺพนฺติ, วิทิตฺวา โลกปริยายํ.}

\prose{อญฺญาย นิพฺพุตา ธีรา, ติณฺณา โลเก วิสตฺติกนฺติ.}

\prose{อถ โข ตา เทวตาโย ปถวิยํ ปติฏฺฐหิตฺวา ภควโต ปาเทสุ สิรสา นิปติตฺวา ภควนฺตํ เอตทโวจุํ “สฺจฺจโย โน ภนฺเต อจฺจคมา ยถาพาลํ ยถามูฬฺหํ ยถา-อกุสลํ\footnote{ยถาพาลา ยถามูฬฺหา ยถา-อกุสลา (สพฺพตฺถ)}, ยา มยํ ภควนฺตํ อาสาเทตพฺพํ อมญฺญิมฺหา, ตาสํ โน ภนฺเต ภควา อจฺจยํ อจฺจยโต ปฏิคฺคณฺหาตุ อายติํ สํวรายา”ติ.\csromanspacer{} อถ โข ภควา สิตํ ปาตฺวากาสิ.\csromanspacer{} อถ โข ตา เทวตาโย ภิยฺโยโส มตฺตาย อุชฺฌายนฺติโย เวหาสํ อพฺภุคฺคญฺฉุํ.\csromanspacer{} เอกา เทวตา ภควโต สนฺติเก อิมํ คาถํ อภาสิ\csromandash{}}

\prose{“อจฺจยํ เทสยนฺตีนํ, โย เจ น ปฏิคณฺหติ.}

\prose{โกปนฺตโร โทสครุ, ส เวรํ ปฏิมุญฺจตี”ติ.}

\prose{อจฺจโย เจ น วิชฺเชถ, โนจิธาปคตํ\footnote{โนจีธ อปหตํ (สฺยา, กํ), โนจิธาปกตํ (?)} สิยา.}

\prose{เวรานิ น จ สมฺเมยฺยุํ, เกนีธ\footnote{เวรานิ จ สมฺเมยฺยุํ, เตนิธ (สี)} กุสโล สิยาติ.}

\prose{กสฺสจฺจยา น วิชฺชนฺติ, กสฺส นตฺถิ อปาคตํ.}

\prose{โก น สมฺโมหมาปาทิ, โก จ ธีโร\footnote{โกธ ธีโร (สฺยา, กํ)} สทา สโตติ.}

\prose{ตถาคตสฺส พุทฺธสฺส, สพฺพภูตานุกมฺปิโน.}

\prose{ตสฺสจฺจยา น วิชฺชนฺติ, ตสฺส นตฺถิ อปาคตํ.}

\prose{โส น สมฺโมหมาปาทิ, โสว\footnote{โสธ (สฺยา, กํ)} ธีโร สทา สโตติ.}

\prose{อจฺจยํ เทสยนฺตีนํ, โย เจ น ปฏิคณฺหติ.}

\prose{โกปนฺตโร โทสครุ, ส เวรํ ปฏิมุญฺจติ.}

\prose{ตํ เวรํ นาภินนฺทามิ, ปฏิคฺคณฺหามิ โวจฺจยนฺติ.}

\csromanheader{2}{6. สทฺธาสุตฺต}
\proseitem{36}{เอกํ สมยํ ภควา สาวตฺถิยํ วิหรติ เชตวเน อนาถปิณฺฑิกสฺส อาราเม.\csromanspacer{} อถ โข สมฺพหุลา สตุลฺลปกายิกา เทวตาโย อภิกฺกนฺตาย รตฺติยา อภิกฺกนฺตวณฺณา เกวลกปฺปํ \csromanfolio{24}เชตวนํ โอภาเสตฺวา เยน ภควา เตนุปสงฺกมิํสุ, อุปสงฺกมิตฺวา ภควนฺตํ อภิวาเทตฺวา เอกมนฺตํ อฏฺฐํสุ, เอกมนฺตํ ฐิตา โข เอกา เทวตา ภควโต สนฺติเก อิมํ คาถํ อภาสิ\csromandash{}}

\prose{“สทฺธา ทุติยา ปุริสสฺส โหติ,}

\prose{โน เจ อสฺสทฺธิยํ อวติฏฺฐติ.}

\prose{ยโส จ กิตฺตี จ ตตฺวสฺส โหติ,}

\prose{สคฺคํ จ โส คจฺฉติ สรีรํ วิหายา”ติ.}

\csromanheader{2}{อถ โข อปรา เทวตา ภควโต สนฺติเก อิมา คาถาโย อภาสิ\csromandash{}}
\csromangathameasure{“โกธํ ชเห วิปฺปชเหยฺย มานํ, สํโยชนํ สพฺพมติกฺกเมยฺย. \\ ตํ นามรูปสฺมิมสชฺชมานํ, อกิญฺจนํ นานุปตนฺติ สงฺคา”ติ. \\ ปมาทมนุยุญฺชนฺติ, พาลา ทุมฺเมธิโน ชนา. \\ อปฺปมาทญฺจ เมธาวี, ธนํ เสฏฺฐํว รกฺขติ. \\ มา ปมาทมนุยุญฺเชถ, มา กามรติ สนฺถวํ. \\ อปฺปมตฺโต หิ ฌายนฺโต, ปปฺโปติ ปรมํ สุขนฺติ.}
\csromangathagroup{\csromangathastanza{“โกธํ ชเห วิปฺปชเหยฺย มานํ, สํโยชนํ สพฺพมติกฺกเมยฺย. \\ ตํ นามรูปสฺมิมสชฺชมานํ, อกิญฺจนํ นานุปตนฺติ สงฺคา”ติ.}\csromangathastanza{ปมาทมนุยุญฺชนฺติ, พาลา ทุมฺเมธิโน ชนา. \\ อปฺปมาทญฺจ เมธาวี, ธนํ เสฏฺฐํว รกฺขติ.}\csromangathastanza{มา ปมาทมนุยุญฺเชถ, มา กามรติ สนฺถวํ. \\ อปฺปมตฺโต หิ ฌายนฺโต, ปปฺโปติ ปรมํ สุขนฺติ.}}

\csromanheader{2}{7. สมยสุตฺต}
\proseitem{37}{เอวํ เม สุตํ\csromandash{}เอกํ สมยํ ภควา สกฺเกสุ วิหรติ กปิลวตฺถุสฺมิํ มหาวเน มหตา ภิกฺขุสํเฆน สทฺธิํ ปญฺจมตฺเตหิ ภิกฺขุสเตหิ สพฺเพเหว อรหนฺเตหิ, ทสหิ จ โลกธาตูหิ เทวตา เยภุยฺเยน สนฺนิปติตา โหนฺติ ภควนฺตํ ทสฺสนาย ภิกฺขุสํฆญฺจ.\csromanspacer{} อถ โข จตุนฺนํ สุทฺธาวาสกายิกานํ เทวตานํ เอตทโหสิ “อยํ โข ภควา สกฺเกสุ วิหรติ กปิลวตฺถุสฺมิํ มหาวเน มหตา ภิกฺขุสํเฆน สทฺธิํ ปญฺจมตฺเตหิ ภิกฺขุสเตหิ สพฺเพเหว อรหนฺเตหิ, ทสหิ จ โลกธาตูหิ เทวตา เยภุยฺเยน สนฺนิปติตา โหนฺติ ภควนฺตํ ทสฺสนาย ภิกฺขุสํฆญฺจ.\csromanspacer{} ยํนูน มยมฺปิ เยน ภควา เตนุปสงฺกเมยฺยาม, อุปสงฺกมิตฺวา ภควโต สนฺติเก ปจฺเจกํ คาถํ\footnote{ปจฺเจกคาถํ (สี, สฺยา, กํ, อิ)} ภาเสยฺยามา”ติ.}

\prose{\csromanfolio{25}อถ โข ตา เทวตา เสยฺยถาปิ นาม พลวา ปุริโส สมิญฺชิตํ วา พาหํ ปสาเรยฺย, ปสาริตํ วา พาหํ สมิญฺเชยฺย.\csromanspacer{} เอวเมว สุทฺธาวาเสสุ เทเวสุ อนฺตรหิตา ภควโต ปุรโต ปาตุรเหสุํ.\csromanspacer{} อถ โข ตา เทวตา ภควนฺตํ อภิวาเทตฺวา เอกมนฺตํ อฏฺฐํสุ, เอกมนฺตํ ฐิตา โข เอกา เทวตา ภควโต สนฺติเก อิมํ คาถํ อภาสิ\csromandash{}}

\prose{“มหาสมโย ปวนสฺมิํ, เทวกายา สมาคตา.}

\prose{อาคตมฺห อิมํ ธมฺมสมยํ, ทกฺขิตาเย อปราชิตสํฆนฺ”ติ.}

\csromanheader{2}{อถ โข อปรา เทวตา ภควโต สนฺติเก อิมํ คาถํ อภาสิ\csromandash{}}
\prose{“ตตฺร ภิกฺขโว สมาทหํสุ, จิตฺตมตฺตโน อุชุกํ อกํสุ\footnote{อุชุกมกํสุ (สี, สฺยา, กํ, อิ)}.}

\prose{สารถีว เนตฺตานิ คเหตฺวา, อินฺทฺริยานิ รกฺขนฺติ ปณฺฑิตา”ติ.}

\csromanheader{2}{อถ โข อปรา เทวตา ภควโต สนฺติเก อิมํ คาถํ อภาสิ\csromandash{}}
\prose{“เฉตฺวา ขีลํ เฉตฺวา ปลิฆํ, อินฺทขีลํ อูหจฺจ มเนชา.}

\prose{เต จรนฺติ สุทฺธา วิมลา, จกฺขุมตา สุทนฺตา สุสุนาคา”ติ.}

\csromanheader{2}{อถ โข อปรา เทวตา ภควโต สนฺติเก อิมํ คาถํ อภาสิ\csromandash{}}
\prose{“เย เกจิ พุทฺธํ สรณํ คตาเส, น เต คมิสฺสนฺติ อปายภูมิํ.}

\prose{ปหาย มานุสํ เทหํ, เทวกายํ ปริปูเรสฺสนฺตี”ติ.}

\csromanheader{2}{8. สกลิกสุตฺต}
\proseitem{38}{เอวํ เม สุตํ\csromandash{}เอกํ สมยํ ภควา ราชคเห วิหรติ มทฺทกุจฺฉิสฺมิํ มิคทาเย.\csromanspacer{} เตน โข ปน สมเยน ภควโต ปาโท สกลิกาย\footnote{สกฺขลิกาย (ก)} ขโต โหติ.\csromanspacer{} ภุสา สุทํ ภควโต เวทนา วตฺตนฺติ สารีริกา เวทนา ทุกฺขา ติพฺพา\footnote{ติปฺปา (สี, สฺยา, กํ, อิ)} ขรา กฏุกา อสาตา อมนาปา, ตา สุทํ ภควา สโต สมฺปชาโน อธิวาเสติ อวิหญฺญมาโน.\csromanspacer{} อถ โข ภควา จตุคฺคุณํ สงฺฆาฏิํ ปญฺญาเปตฺวา ทกฺขิเณน ปสฺเสน สีหเสยฺยํ กปฺเปติ ปาเท ปาทํ อจฺจาธาย สโต สมฺปชาโน.}

\prose{\csromanfolio{26}อถ โข สตฺตสตา สตุลฺลปกายิกา เทวตาโย อภิกฺกนฺตาย รตฺติยา อภิกฺกนฺตวณฺณา เกวลกปฺปํ มทฺทกุจฺฉิํ โอภาเสตฺวา เยน ภควา เตนุปสงฺกมิํสุ, อุปสงฺกมิตฺวา ภควนฺตํ อภิวาเทตฺวา เอกมนฺตํ อฏฺฐํสุ, เอกมนฺตํ ฐิตา โข เอกา เทวตา ภควโต สนฺติเก อิมํ อุทานํ อุทาเนสิ “นาโค วต โภ สมโณ โคตโม, นาควตา จ สมุปฺปนฺนา สารีริกา เวทนา ทุกฺขา ติพฺพา ขรา กฏุกา อสาตา อมนาปา สโต สมฺปชาโน อธิวาเสติ อวิหญฺญมาโน”ติ.}

\prose{อถ โข อปรา เทวตา ภควโต สนฺติเก อิมํ อุทานํ อุทาเนสิ “สีโห วต โภ สมโณ โคตโม, สีหวตา จ สมุปฺปนฺนา สารีริกา เวทนา ทุกฺขา ติพฺพา ขรา กฏุกา อสาตา อมนาปา สโต สมฺปชาโน อธิวาเสติ อวิหญฺญมาโน”ติ.}

\prose{อถ โข อปรา เทวตา ภควโต สนฺติเก อิมํ อุทานํ อุทาเนสิ “อาชานีโย วต โภ สมโณ โคตโม, อาชานียวตา จ สมุปฺปนฺนา สารีริกา เวทนา ทุกฺขา ติพฺพา ขรา กฏุกา อสาตา อมนาปา สโต สมฺปชาโน อธิวาเสติ อวิหญฺญมาโน”ติ.}

\prose{อถ โข อปรา เทวตา ภควโต สนฺติเก อิมํ อุทานํ อุทาเนสิ “นิสโภ วต โภ สมโณ โคตโม, นิสภวตา จ สมุปฺปนฺนา สารีริกา เวทนา ทุกฺขา ติพฺพา ขรา กฏุกา อสาตา อมนาปา สโต สมฺปชาโน อธิวาเสติ อวิหญฺญมาโน”ติ.}

\prose{อถ โข อปรา เทวตา ภควโต สนฺติเก อิมํ อุทานํ อุทาเนสิ “โธรโยฺห วต โภ สมโณ โคตโม, โธรยฺหวตา จ สมุปฺปนฺนา สารีริกา เวทนา ทุกฺขา ติพฺพา ขรา กฏุกา อสาตา อมนาปา สโต สมฺปชาโน อธิวาเสติ อวิหญฺญมาโน”ติ.}

\prose{อถ โข อปรา เทวตา ภควโต สนฺติเก อิมํ อุทานํ อุทาเนสิ “ทนฺโต วต โภ สมโณ โคตโม, ทนฺตวตา จ สมุปฺปนฺนา สารีริกา เวทนา ทุกฺขา ติพฺพา ขรา กฏุกา อสาตา อมนาปา สโต สมฺปชาโน อธิวาเสติ อวิหญฺญมาโน”ติ.}

\prose{อถ โข อปรา เทวตา ภควโต สนฺติเก อิมํ อุทานํ อุทาเนสิ “ปสฺส สมาธิํ สุภาวิตํ จิตฺตญฺจ สุวิมุตฺตํ, น จาภินตํ น จาปนตํ \csromanfolio{27}น จ สสงฺขารนิคฺคยฺหวาริตคตํ\footnote{สสงฺขารนิคฺคยฺหวาริตวตํ (สี, สฺยา, กํ, อิ), สสงฺขารนิคฺคยฺหวาริวาวตํ (ก)}, โย เอวรูปํ ปุริสนาคํ ปุริสสีหํ ปุริส-อาชานียํ ปุริสนิสภํ ปุริสโธรยฺหํ ปุริสทนฺตํ อติกฺกมิตพฺพํ มญฺเญยฺย กิมญฺญตฺร อทสฺสนา”ติ.}

\prose{ปญฺจเวทา สตํ สมํ, ตปสฺสี พฺราหฺมณา จรํ.}

\prose{จิตฺตญฺจ เนสํ น สมฺมา วิมุตฺตํ, หีนตฺถรูปา น ปารํคมา เต.}

\prose{ตณฺหาธิปนฺนา วตสีลพทฺธา, ลูขํ ตปํ วสฺสสตํ จรนฺตา.}

\prose{จิตฺตญฺจ เนสํ น สมฺมา วิมุตฺตํ, หีนตฺถรูปา น ปารํคมา เต.}

\prose{น มานกามสฺส ทโม อิธตฺถิ, น โมนมตฺถิ อสมาหิตสฺส.}

\prose{เอโก อรญฺเญ วิหรํ ปมตฺโต, น มจฺจุเธยฺยสฺส ตเรยฺย ปารนฺติ.}

\prose{มานํ ปหาย สุสมาหิตตฺโต, สุเจตโส สพฺพธิ วิปฺปมุตฺโต.}

\prose{เอโก อรญฺเญ วิหรมปฺปมตฺโต, ส มจฺจุเธยฺยสฺส ตเรยฺย ปารนฺติ.}

\csromanheader{2}{9. ปฐมปชฺชุนฺนธีตุสุตฺต}
\proseitem{39}{เอวํ เม สุตํ\csromandash{}เอกํ สมยํ ภควา เวสาลิยํ วิหรติ มหาวเน กูฏาคารสาลายํ.\csromanspacer{} อถ โข โกกนทา ปชฺชุนฺนสฺส ธีตา อภิกฺกนฺตาย รตฺติยา อภิกฺกนฺตวณฺณา เกวลกปฺปํ มหาวนํ โอภาเสตฺวา เยน ภควา เตนุปสงฺกมิ, อุปสงฺกมิตฺวา ภควนฺตํ อภิวาเทตฺวา เอกมนฺตํ อฏฺฐาสิ, เอกมนฺตํ ฐิตา โข สา เทวตา โกกนทา ปชฺชุนฺนสฺส ธีตา ภควโต สนฺติเก อิมา คาถาโย อภาสิ\csromandash{}}

\prose{“เวสาลิยํ วเน วิหรนฺตํ, อคฺคํ สตฺตสฺส สมฺพุทฺธํ.}

\prose{โกกนทาหมสฺมิ อภิวนฺเท, โกกนทา ปชฺชุนฺนสฺส ธีตา.}

\prose{สุตเมว ปุเร อาสิ, ธมฺโม จกฺขุมตานุพุทฺโธ.}

\prose{สาหํ ทานิ สกฺขิ ชานามิ, มุนิโน เทสยโต สุคตสฺส.}

\prose{เย เกจิ อริยํ ธมฺมํ, วิครหนฺตา จรนฺติ ทุมฺเมธา.}

\prose{อุเปนฺติ โรรุวํ โฆรํ, จิรรตฺตํ ทุกฺขํ อนุภวนฺติ.}

\prose{เย จ โข อริเย ธมฺเม, ขนฺติยา อุปสเมน อุเปตา.}

\prose{ปหาย มานุสํ เทหํ, เทวกายํ ปริปูเรสฺสนฺตี”ติ.}

\csromanheader{2}{10. ทุติยปชฺชุนฺนธีตุสุตฺต}
\proseitem{40}{\csromanfolio{28}เอวํ เม สุตํ\csromandash{}เอกํ สมยํ ภควา เวสาลิยํ วิหรติ มหาวเน กูฏาคารสาลายํ.\csromanspacer{} อถ โข จูฬโกกนทา\footnote{จุลฺลโกกนทา (สี, สฺยา, กํ)} ปชฺชุนฺนสฺส ธีตา อภิกฺกนฺตาย รตฺติยา อภิกฺกนฺตวณฺณา เกวลกปฺปํ มหาวนํ โอภาเสตฺวา เยน ภควา เตนุปสงฺกมิ, อุปสงฺกมิตฺวา ภควนฺตํ อภิวาเทตฺวา เอกมนฺตํ อฏฺฐาสิ, เอกมนฺตํ ฐิตา โข สา เทวตา จูฬโกกนทา ปชฺชุนฺนสฺส ธีตา ภควโต สนฺติเก อิมา คาถาโย อภาสิ\csromandash{}}

\prose{“อิธาคมา วิชฺชุปภาสวณฺณา, โกกนทา ปชฺชุนฺนสฺส ธีตา.}

\prose{พุทฺธญฺจ ธมฺมญฺจ นมสฺสมานา, คาถาจิมา อตฺถวตี อภาสิ.}

\prose{พหุนาปิ โข ตํ วิภเชยฺยํ, ปริยาเยน ตาทิโส ธมฺโม.}

\prose{สํขิตฺตมตฺถํ\footnote{สํขิตฺตมตฺตํ (ก)} ลปยิสฺสามิ, ยาวตา เม มนสา ปริยตฺตํ.}

\prose{ปาปํ น กยิรา วจสา มนสา,}

\prose{กาเยน วา กิญฺจน สพฺพโลเก.}

\prose{กาเม ปหาย สติมา สมฺปชาโน,}

\prose{ทุกฺขํ น เสเวถ อนตฺถสํหิตนฺ”ติ.}

\vspace{\tipitakaparskip}
\csromancenter{สตุลฺลปกายิกวคฺโค จตุตฺโถ.}
\titlehead{\textbf{ตสฺสุทฺทานํ}}
\csromangathameasure{สพฺภิมจฺฉรินา สาธุ, น สนฺตุชฺฌานสญฺญิโน. \\ สทฺธา สมโย สกลิกํ, อุโภ ปชฺชุนฺนธีตโรติ.}
\csromangathagroup{\csromangathastanza{สพฺภิมจฺฉรินา สาธุ, น สนฺตุชฺฌานสญฺญิโน. \\ สทฺธา สมโย สกลิกํ, อุโภ ปชฺชุนฺนธีตโรติ.}}

\chapterhead{5. อาทิตฺตวคฺค}
\csromanheader{2}{1. อาทิตฺตสุตฺต}
\proseitem{41}{เอวํ เม สุตํ\csromandash{}เอกํ สมยํ ภควา สาวตฺถิยํ วิหรติ เชตวเน อนาถปิณฺฑิกสฺส อาราเม.\csromanspacer{} อถ โข อญฺญตรา เทวตา \csromanfolio{29}อภิกฺกนฺตาย รตฺติยา อภิกฺกนฺตวณฺณา เกวลกปฺปํ เชตวนํ โอภาเสตฺวา เยน ภควา เตนุปสงฺกมิ, อุปสงฺกมิตฺวา ภควนฺตํ อภิวาเทตฺวา เอกมนฺตํ อฏฺฐาสิ, เอกมนฺตํ ฐิตา โข สา เทวตา ภควโต สนฺติเก อิมา คาถาโย อภาสิ\csromandash{}}

\csromangathameasure{“อาทิตฺตสฺมิํ อคารสฺมิํ, ยํ นีหรติ ภาชนํ. \\ ตํ ตสฺส โหติ อตฺถาย, โน จ ยํ ตตฺถ ฑยฺหติ. \\ เอวํ อาทิตฺตโก โลโก, ชราย มรเณน จ. \\ นีหเรเถว ทาเนน, ทินฺนํ โหติ สุนีหตํ. \\ ทินฺนํ สุขผลํ โหติ, นาทินฺนํ โหติ ตํ ตถา. \\ โจรา หรนฺติ ราชาโน, อคฺคิ ฑหติ นสฺสติ. \\ อถ อนฺเตน ชหติ, สรีรํ สปริคฺคหํ. \\ เอตทญฺญาย เมธาวี, ภุญฺเชถ จ ทเทถ จ. \\ ทตฺวา จ ภุตฺวา จ ยถานุภาวํ, \\ อนินฺทิโต สคฺคมุเปติ ฐานนฺ”ติ.}
\csromangathagroup{\csromangathastanza{“อาทิตฺตสฺมิํ อคารสฺมิํ, ยํ นีหรติ ภาชนํ. \\ ตํ ตสฺส โหติ อตฺถาย, โน จ ยํ ตตฺถ ฑยฺหติ.}\csromangathastanza{เอวํ อาทิตฺตโก โลโก, ชราย มรเณน จ. \\ นีหเรเถว ทาเนน, ทินฺนํ โหติ สุนีหตํ.}\csromangathastanza{ทินฺนํ สุขผลํ โหติ, นาทินฺนํ โหติ ตํ ตถา. \\ โจรา หรนฺติ ราชาโน, อคฺคิ ฑหติ นสฺสติ.}\csromangathastanza{อถ อนฺเตน ชหติ, สรีรํ สปริคฺคหํ. \\ เอตทญฺญาย เมธาวี, ภุญฺเชถ จ ทเทถ จ.}\csromangathastanza{ทตฺวา จ ภุตฺวา จ ยถานุภาวํ, \\ อนินฺทิโต สคฺคมุเปติ ฐานนฺ”ติ.}}

\csromanheader{2}{2. กิํททสุตฺต}
\csromanhangingbegin
\hangingheaditem{42}{กิํทโท พลโท โหติ, กิํทโท โหติ วณฺณโท.}
\hangingbody{กิํทโท สุขโท โหติ, กิํทโท โหติ จกฺขุโท.}
\hangingbody{โก จ สพฺพทโท โหติ, ตํ เม อกฺขาหิ ปุจฺฉิโตติ.}
\hangingbody{อนฺนโท พลโท โหติ, วตฺถโท โหติ วณฺณโท.}
\hangingbody{ยานโท สุขโท โหติ, ทีปโท โหติ จกฺขุโท.}
\hangingbody{โส จ สพฺพทโท โหติ, โย ททาติ อุปสฺสยํ.}
\hangingbody{อมตํทโท จ โส โหติ, โย ธมฺมมนุสาสตีติ.}
\csromanhangingend

\csromanheader{2}{3. อนฺนสุตฺต}
\csromanhangingbegin
\hangingheaditem{43}{อนฺนเมวาภินนฺทนฺติ, อุภเย เทวมานุสา.}
\hangingbody{อถ โก นาม โส ยกฺโข, ยํ อนฺนํ นาภินนฺทตีติ.}
\csromanhangingend

\csromangathameasure{เย นํ ททนฺติ สทฺธาย, วิปฺปสนฺเนน เจตสา. \\ ตเมว อนฺนํ ภชติ, อสฺมิํ โลเก ปรมฺหิ จ. \\ ตสฺมา วิเนยฺย มจฺเฉรํ, ทชฺชา ทานํ มลาภิภู. \\ ปุญฺญานิ ปรโลกสฺมิํ, ปติฏฺฐา โหนฺติ ปาณินนฺติ.}
\csromangathagroup{\csromangathastanza{\csromanfolio{30}เย นํ ททนฺติ สทฺธาย, วิปฺปสนฺเนน เจตสา. \\ ตเมว อนฺนํ ภชติ, อสฺมิํ โลเก ปรมฺหิ จ.}\csromangathastanza{ตสฺมา วิเนยฺย มจฺเฉรํ, ทชฺชา ทานํ มลาภิภู. \\ ปุญฺญานิ ปรโลกสฺมิํ, ปติฏฺฐา โหนฺติ ปาณินนฺติ.}}

\csromanheader{2}{4. เอกมูลสุตฺต}
\csromanhangingbegin
\hangingheaditem{44}{เอกมูลํ ทฺวิราวฏฺฏํ, ติมลํ ปญฺจปตฺถรํ.}
\hangingbody{สมุทฺทํ ทฺวาทสาวฏฺฏํ, ปาตาลํ อตรี อิสีติ.}
\csromanhangingend

\csromanheader{2}{5. อโนมสุตฺต}
\csromanhangingbegin
\hangingheaditem{45}{อโนมนามํ นิปุณตฺถทสฺสิํ, ปญฺญาททํ กามาลเย อสตฺตํ.}
\hangingbody{ตํ ปสฺสถ สพฺพวิทุํ สุเมธํ,}
\hangingbody{อริเย ปเถ กมมานํ มเหสินฺติ.}
\csromanhangingend

\csromanheader{2}{6. อจฺฉราสุตฺต}
\csromanhangingbegin
\hangingheaditem{46}{อจฺฉราคณสงฺฆุฏฺฐํ, ปิสาจคณเสวิตํ.}
\hangingbody{วนนฺตํ โมหนํ นาม, กถํ ยาตฺรา ภวิสฺสตีติ.}
\hangingbody{อุชุโก นาม โส มคฺโค, อภยา นาม สา ทิสา.}
\hangingbody{รโถ อกูชโน นาม, ธมฺมจกฺเกหิ สํยุโต.}
\hangingbody{หิรี ตสฺส อปาลมฺโพ, สตฺยสฺส ปริวารณํ.}
\hangingbody{ธมฺมาหํ สารถิํ พฺรูมิ, สมฺมาทิฏฺฐิปุเรชวํ.}
\hangingbody{ยสฺส เอตาทิสํ ยานํ, อิตฺถิยา ปุริสสฺส วา.}
\hangingbody{ส เว เอเตน ยาเนน, นิพฺพานสฺเสว สนฺติเกติ.}
\csromanhangingend

\csromanheader{2}{7. วนโรปสุตฺต}
\csromanhangingbegin
\hangingheaditem{47}{เกสํ ทิวา จ รตฺโต จ, สทา ปุญฺญํ ปวฑฺฒติ.}
\hangingbody{ธมฺมฏฺฐา สีลสมฺปนฺนา, เก ชนา สคฺคคามิโนติ.}
\hangingbody{อารามโรปา วนโรปา, เย ชนา เสตุการกา.}
\hangingbody{ปปญฺจ อุทปานญฺจ, เย ททนฺติ อุปสฺสยํ.}
\csromanhangingend

\csromangathameasure{เตสํ ทิวา จ รตฺโต จ, สทา ปุญฺญํ ปวฑฺฒติ. \\ ธมฺมฏฺฐา สีลสมฺปนฺนา, เต ชนา สคฺคคามิโนติ.}
\csromangathagroup{\csromangathastanza{\csromanfolio{31}เตสํ ทิวา จ รตฺโต จ, สทา ปุญฺญํ ปวฑฺฒติ. \\ ธมฺมฏฺฐา สีลสมฺปนฺนา, เต ชนา สคฺคคามิโนติ.}}

\csromanheader{2}{8. เชตวนสุตฺต}
\csromanhangingbegin
\hangingheaditem{48}{อิทํ หิ ตํ เชตวนํ, อิสิสงฺฆนิเสวิตํ,}
\hangingbody{อาวุตฺถํ1 ธมฺมราเชน, ปีติสญฺชนนํ มม.}
\hangingbody{กมฺมํ วิชฺชา จ ธมฺโม จ, สีลํ ชีวิตมุตฺตมํ.}
\hangingbody{เอเตน มจฺจา สุชฺฌนฺติ, น โคตฺเตน ธเนน วา.}
\hangingbody{ตสฺมา หิ ปณฺฑิโต โปโส, สมฺปสฺสํ อตฺถมตฺตโน.}
\hangingbody{โยนิโส วิจิเน ธมฺมํ, เอวํ ตตฺถ วิสุชฺฌติ.}
\hangingbody{สาริปุตฺโตว ปญฺญาย, สีเลน อุปสเมน จ.}
\hangingbody{โยปิ ปารงฺคโต ภิกฺขุ, เอตาวปรโม สิยาติ.}
\csromanhangingend

\csromanheader{2}{9. มจฺฉริสุตฺต}
\csromanhangingbegin
\hangingheaditem{49}{เยธ มจฺฉริโน โลเก, กทริยา ปริภาสกา.}
\hangingbody{อญฺเญสํ ททมานานํ, อนฺตรายกรา นรา.}
\hangingbody{กีทิโส เตสํ วิปาโก, สมฺปราโย จ กีทิโส.}
\hangingbody{ภควนฺตํ ปุฏฺฐุมาคมฺม, กถํ ชาเนมุ ตํ มยนฺติ.}
\hangingbody{เยธ มจฺฉริโน โลเก, กทริยา ปริภาสกา.}
\hangingbody{อญฺเญสํ ททมานานํ, อนฺตรายกรา นรา.}
\hangingbody{นิรยํ ติรจฺฉานโยนิํ, ยมโลกํ อุปปชฺชเร.}
\hangingbody{สเจ เอนฺติ มนุสฺสตฺตํ, ทลิทฺเท ชายเร กุเล.}
\hangingbody{โจฬํ ปิณฺโฑ รตี ขิฑฺฑา, ยตฺถ กิจฺเฉน ลพฺภติ.}
\hangingbody{ปรโต อาสีสเร2 พาลา, ตมฺปิ เตสํ น ลพฺภติ.}
\hangingbody{ทิฏฺเฐ ธมฺเมส วิปาโก, สมฺปราเย3 จ ทุคฺคตีติ.}
\hangingbody{อิติเหตํ วิชานาม, อญฺญํ ปุจฺฉาม โคตม.}
\hangingbody{เยธ ลทฺธา มนุสฺสตฺตํ, วทญฺญู วีตมจฺฉรา.}
\csromanhangingend

\csromangathameasure{พุทฺเธ ปสนฺนา ธมฺเม จ, สํเฆ จ ติพฺพคารวา. \\ กีทิโส เตสํ วิปาโก, สมฺปราโย จ กีทิโส. \\ ภควนฺตํ ปุฏฺฐุมาคมฺม, กถํ ชาเนมุ ตํ มยนฺติ. \\ เยธ ลทฺธา มนุสฺสตฺตํ, วทญฺญู วีตมจฺฉรา. \\ พุทฺเธ ปสนฺนา ธมฺเม จ, สํเฆ จ ติพฺพคารวา. \\ เอเต สคฺคา ปกาสนฺติ, ยตฺถ เต อุปปชฺชเร. \\ สเจ เอนฺติ มนุสฺสตฺตํ, อฑฺเฒ อาชายเร กุเล. \\ โจฬํ ปิณฺโฑ รตี ขิฑฺฑา, ยตฺถากิจฺเฉน ลพฺภติ. \\ ปรสมฺภเตสุ โภเคสุ, วสวตฺตีว โมทเร. \\ ทิฏฺเฐ ธมฺเมส วิปาโก, สมฺปราเย จ สุคฺคตีติ.}
\csromangathagroup{\csromangathastanza{\csromanfolio{32}พุทฺเธ ปสนฺนา ธมฺเม จ, สํเฆ จ ติพฺพคารวา. \\ กีทิโส เตสํ วิปาโก, สมฺปราโย จ กีทิโส.}\csromangathastanza{ภควนฺตํ ปุฏฺฐุมาคมฺม, กถํ ชาเนมุ ตํ มยนฺติ. \\ เยธ ลทฺธา มนุสฺสตฺตํ, วทญฺญู วีตมจฺฉรา.}\csromangathastanza{พุทฺเธ ปสนฺนา ธมฺเม จ, สํเฆ จ ติพฺพคารวา. \\ เอเต สคฺคา\footnote{สคฺเค (สี, สฺยา, กํ)} ปกาสนฺติ, ยตฺถ เต อุปปชฺชเร.}\csromangathastanza{สเจ เอนฺติ มนุสฺสตฺตํ, อฑฺเฒ อาชายเร กุเล. \\ โจฬํ ปิณฺโฑ รตี ขิฑฺฑา, ยตฺถากิจฺเฉน ลพฺภติ.}\csromangathastanza{ปรสมฺภเตสุ โภเคสุ, วสวตฺตีว โมทเร. \\ ทิฏฺเฐ ธมฺเมส วิปาโก, สมฺปราเย จ สุคฺคตีติ.}}

\csromanheader{2}{10. ฆฏีการสุตฺต}
\csromanhangingbegin
\hangingheaditem{50}{อวิหํ อุปปนฺนาเส, วิมุตฺตา สตฺต ภิกฺขโว.}
\hangingbody{ราคโทสปริกฺขีณา, ติณฺณา โลเก วิสตฺติกนฺติ.}
\hangingbody{เก จ เต อตรุํ ปงฺกํ2, มจฺจุเธยฺยํ สุทุตฺตรํ.}
\hangingbody{เก หิตฺวา มานุสํ เทหํ, ทิพฺพโยคํ อุปจฺจคุนฺติ.}
\hangingbody{อุปโก ปลคณฺโฑ จ, ปุกฺกุสาติ จ เต ตโย.}
\hangingbody{ภทฺทิโย ขณฺฑเทโว จ, พาหุรคฺคิ จ สิงฺคิโย3.}
\hangingbody{เต หิตฺวา มานุสํ เทหํ, ทิพฺพโยคํ อุปจฺจคุนฺติ.}
\hangingbody{กุสลี ภาสสี เตสํ, มารปาสปฺปหายินํ.}
\hangingbody{กสฺส เต ธมฺมมญฺญาย, อจฺฉิทุํ ภวพนฺธนนฺติ.}
\hangingbody{น อญฺญตฺร ภควตา, นาญฺญตฺร ตว สาสนา.}
\hangingbody{ยสฺส เต ธมฺมมญฺญาย, อจฺฉิทุํ ภวพนฺธนํ.}
\hangingbody{ยตฺถ นามญฺจ รูปญฺจ, อเสสํ อุปรุชฺฌติ.}
\hangingbody{ตํ เต ธมฺมํ อิธญฺญาย, อจฺฉิทุํ ภวพนฺธนนฺติ.}
\hangingbody{คมฺภีรํ ภาสสี วาจํ, ทุพฺพิชานํ สุทุพฺพุธํ.}
\hangingbody{กสฺส ตฺวํ ธมฺมมญฺญาย, วาจํ ภาสสิ อีทิสนฺติ.}
\csromanhangingend

\csromangathameasure{กุมฺภกาโร ปุเร อาสิํ, เวกฬิงฺเค ฆฏีกโร. \\ มาตาเปตฺติภโร อาสิํ, กสฺสปสฺส อุปาสโก. \\ วิรโต เมถุนา ธมฺมา, พฺรหฺมจารี นิรามิโส. \\ อหุวา เต สคาเมยฺโย, อหุวา เต ปุเร สขา. \\ โสหเมเต ปชานามิ, วิมุตฺเต สตฺต ภิกฺขโว. \\ ราคโทสปริกฺขีเณ, ติณฺเณ โลเก วิสตฺติกนฺติ. \\ เอวเมตํ ตทา อาสิ, ยถา ภาสสิ ภคฺคว. \\ กุมฺภกาโร ปุเร อาสิ, เวกฬิงฺเค ฆฏีกโร. \\ มาตาเปตฺติภโร อาสิ, กสฺสปสฺส อุปาสโก. \\ วิรโต เมถุนา ธมฺมา, พฺรหฺมจารี นิรามิโส. \\ อหุวา เม สคาเมยฺโย, อหุวา เม ปุเร สขาติ. \\ เอวเมตํ ปุราณานํ, สหายานํ อหุ สงฺคโม. \\ อุภินฺนํ ภาวิตตฺตานํ, สรีรนฺติมธารินนฺติ.}
\csromangathagroup{\csromangathastanza{\csromanfolio{33}กุมฺภกาโร ปุเร อาสิํ, เวกฬิงฺเค\footnote{เวหฬิงฺเค (สี), เวภฬิงฺเค (สฺยา, กํ)} ฆฏีกโร. \\ มาตาเปตฺติภโร อาสิํ, กสฺสปสฺส อุปาสโก.}\csromangathastanza{วิรโต เมถุนา ธมฺมา, พฺรหฺมจารี นิรามิโส. \\ อหุวา เต สคาเมยฺโย, อหุวา เต ปุเร สขา.}\csromangathastanza{โสหเมเต ปชานามิ, วิมุตฺเต สตฺต ภิกฺขโว. \\ ราคโทสปริกฺขีเณ, ติณฺเณ โลเก วิสตฺติกนฺติ.}\csromangathastanza{เอวเมตํ ตทา อาสิ, ยถา ภาสสิ ภคฺคว. \\ กุมฺภกาโร ปุเร อาสิ, เวกฬิงฺเค ฆฏีกโร.}\csromangathastanza{มาตาเปตฺติภโร อาสิ, กสฺสปสฺส อุปาสโก. \\ วิรโต เมถุนา ธมฺมา, พฺรหฺมจารี นิรามิโส.}\csromangathastanza{อหุวา เม สคาเมยฺโย, อหุวา เม ปุเร สขาติ. \\ เอวเมตํ ปุราณานํ, สหายานํ อหุ สงฺคโม. \\ อุภินฺนํ ภาวิตตฺตานํ, สรีรนฺติมธารินนฺติ.}}

\vspace{\tipitakaparskip}
\csromancenter{อาทิตฺตวคฺโค ปญฺจโม.}
\titlehead{\textbf{ตสฺสุทฺทานํ}}
\csromangathameasure{อาทิตฺตํ กิํททํ อนฺนํ, เอกมูล-อโนมิยํ. \\ อจฺฉราวนโรปเชตํ, มจฺฉเรน ฆฏีกโรติ.}
\csromangathagroup{\csromangathastanza{อาทิตฺตํ กิํททํ อนฺนํ, เอกมูล-อโนมิยํ. \\ อจฺฉราวนโรปเชตํ, มจฺฉเรน ฆฏีกโรติ.}}

\chapterhead{6. ชราวคฺค}
\csromanheader{2}{1. ชราสุตฺต}
\csromanhangingbegin
\hangingheaditem{51}{กิํสุ ยาว ชรา สาธุ, กิํสุ สาธุ ปติฏฺฐิตํ.}
\hangingbody{กิํสุ นรานํ รตนํ, กิํสุ โจเรหิ ทูหรนฺติ.}
\hangingbody{สีลํ ยาว ชรา สาธุ, สทฺธา สาธุ ปติฏฺฐิตา.}
\hangingbody{ปญฺญา นรานํ รตนํ, ปุญฺญํ โจเรหิ ทูหรนฺติ.}
\csromanhangingend

\csromanheader{2}{2. อชรสาสุตฺต}
\csromanhangingbegin
\hangingheaditem{52}{\csromanfolio{34}กิํสุ อชรสา สาธุ, กิํสุ สาธุ อธิฏฺฐิตํ.}
\hangingbody{กิํสุ นรานํ รตนํ, กิํสุ โจเรตฺยหาริยนฺติ.}
\hangingbody{สีลํ อชรสา สาธุ, สทฺธา สาธุ อธิฏฺฐิตา.}
\hangingbody{ปญฺญา นรานํ รตนํ, ปุญฺญํ โจเรตฺยหาริยนฺติ.}
\csromanhangingend

\csromanheader{2}{3. มิตฺตสุตฺต}
\csromanhangingbegin
\hangingheaditem{53}{กิํสุ ปวสโต\footnote{ปถวโต (อิ, ก)} มิตฺตํ, กิํสุ มิตฺตํ สเก ฆเร.}
\hangingbody{กิํ มิตฺตํ อตฺถชาตสฺส, กิํ มิตฺตํ สมฺปรายิกนฺติ.}
\hangingbody{สตฺโถ ปวสโต มิตฺตํ, มาตา มิตฺตํ สเก ฆเร.}
\hangingbody{สหาโย อตฺถชาตสฺส, โหติ มิตฺตํ ปุนปฺปุนํ.}
\hangingbody{สยํกตานิ ปุญฺญานิ, ตํ มิตฺตํ สมฺปรายิกนฺติ.}
\csromanhangingend

\csromanheader{2}{4. วตฺถุสุตฺต}
\csromanhangingbegin
\hangingheaditem{54}{กิํสุ วตฺถุ มนุสฺสานํ, กิํสูธ ปรโม สขา.}
\hangingbody{กิํสุ ภูตา อุปชีวนฺติ, เย ปาณา ปถวิสฺสิตาติ2.}
\hangingbody{ปุตฺตา วตฺถุ มนุสฺสานํ, ภริยา จ3 ปรโม สขา.}
\hangingbody{วุฏฺฐิํ ภูตา อุปชีวนฺติ, เย ปาณา ปถวิสฺสิตาติ.}
\csromanhangingend

\csromanheader{2}{5. ปฐมชนสุตฺต}
\csromanhangingbegin
\hangingheaditem{55}{กิํสุ ชเนติ ปุริสํ, กิํสุ ตสฺส วิธาวติ.}
\hangingbody{กิํสุ สํสารมาปาทิ, กิํสุ ตสฺส มหพฺภยนฺติ.}
\hangingbody{ตณฺหา ชเนติ ปุริสํ, จิตฺตมสฺส วิธาวติ.}
\hangingbody{สตฺโต สํสารมาปาทิ, ทุกฺขมสฺส มหพฺภยนฺติ.}
\csromanhangingend

\csromanheader{2}{6. ทุติยชนสุตฺต}
\csromanhangingbegin
\hangingheaditem{56}{\csromanfolio{35}กิํสุ ชเนติ ปุริสํ, กิํสุ ตสฺส วิธาวติ.}
\hangingbody{กิํสุ สํสารมาปาทิ, กิสฺมา น ปริมุจฺจตีติ.}
\hangingbody{ตณฺหา ชเนติ ปุริสํ, จิตฺตมสฺส วิธาวติ.}
\hangingbody{สตฺโต สํสารมาปาทิ, ทุกฺขา น ปริมุจฺจติ.}
\csromanhangingend

\csromanheader{2}{7. ตติยชนสุตฺต}
\csromanhangingbegin
\hangingheaditem{57}{กิํสุ ชเนติ ปุริสํ, กิํสุ ตสฺส วิธาวติ.}
\hangingbody{กิํสุ สํสารมาปาทิ, กิํสุ ตสฺส ปรายนนฺติ.}
\hangingbody{ตณฺหา ชเนติ ปุริสํ, จิตฺตมสฺส วิธาวติ.}
\hangingbody{สตฺโต สํสารมาปาทิ, กมฺมํ ตสฺส ปรายนนฺติ.}
\csromanhangingend

\csromanheader{2}{8. อุปฺปถสุตฺต}
\csromanhangingbegin
\hangingheaditem{58}{กิํสุ อุปฺปโถ อกฺขาโต, กิํสุ รตฺตินฺทิวกฺขโย.}
\hangingbody{กิํ มลํ พฺรหฺมจริยสฺส, กิํ สินานมโนทกนฺติ.}
\hangingbody{ราโค อุปฺปโถ อกฺขาโต, วโย รตฺตินฺทิวกฺขโย.}
\hangingbody{อิตฺถี มลํ พฺรหฺมจริยสฺส, เอตฺถายํ สชฺชเต ปชา.}
\hangingbody{ตโป จ พฺรหฺมจริยญฺจ, ตํ สินานมโนทกนฺติ.}
\csromanhangingend

\csromanheader{2}{9. ทุติยสุตฺต}
\csromanhangingbegin
\hangingheaditem{59}{กิํสุ ทุติยา\footnote{ทุติยํ (สฺยา, กํ, อิ)} ปุริสสฺส โหติ, กิํสุ เจนํ ปสาสติ.}
\hangingbody{กิสฺส จาภิรโต มจฺโจ, สพฺพทุกฺขา ปมุจฺจตีติ.}
\hangingbody{สทฺธา ทุติยา ปุริสสฺส โหติ, ปญฺญา เจนํ ปสาสติ.}
\hangingbody{นิพฺพานาภิรโต มจฺโจ, สพฺพทุกฺขา ปมุจฺจตีติ.}
\csromanhangingend

\csromanheader{2}{10. กวิสุตฺต}
\proseitem{60}{กิํสุ นิทานํ คาถานํ, กิํสุ ตาสํ วิยญฺชนํ.}

\csromanheader{2}{กิํสุ สนฺนิสฺสิตา คาถา, กิํสุ คาถานมาสโยติ.}
\csromangathameasure{ฉนฺโท นิทานํ คาถานํ, อกฺขรา ตาสํ วิยญฺชนํ.}
\csromangathagroup{\csromangathastanza{\csromanfolio{36}ฉนฺโท นิทานํ คาถานํ, อกฺขรา ตาสํ วิยญฺชนํ.}}

\csromanheader{2}{นามสนฺนิสฺสิตา คาถา, กวิ คาถานมาสโยติ.}
\csromangathameasure{ชราวคฺโค ฉฏฺโฐ.}
\csromangathagroup{\csromangathastanza{ชราวคฺโค ฉฏฺโฐ.}}

\titlehead{\textbf{ตสฺสุทฺทานํ}}
\csromangathameasure{ชรา อชรสา มิตฺตํ, วตฺถุ ตีณิ ชนานิ จ. \\ อุปฺปโถ จ ทุติโย จ, กวินา ปูริโต วคฺโคติ.}
\csromangathagroup{\csromangathastanza{ชรา อชรสา มิตฺตํ, วตฺถุ ตีณิ ชนานิ จ. \\ อุปฺปโถ จ ทุติโย จ, กวินา ปูริโต วคฺโคติ.}}

\chapterhead{7. อทฺธวคฺค}
\csromanheader{2}{1. นามสุตฺต}
\csromanhangingbegin
\hangingheaditem{61}{กิํสุ สพฺพํ อทฺธภวิ\footnote{อนฺวภวิ (สี)}, กิสฺมา ภิยฺโย น วิชฺชติ.}
\hangingbody{กิสฺสสฺสุ เอกธมฺมสฺส, สพฺเพว วสมนฺวคูติ2.}
\hangingbody{นามํ สพฺพํ อทฺธภวิ, นามา ภิยฺโย น วิชฺชติ.}
\hangingbody{นามสฺส เอกธมฺมสฺส, สพฺเพว วสมนฺวคูติ.}
\csromanhangingend

\csromanheader{2}{2. จิตฺตสุตฺต}
\csromanhangingbegin
\hangingheaditem{62}{เกนสฺสุ นียติ โลโก, เกนสฺสุ ปริกสฺสติ.}
\hangingbody{กิสฺสสฺสุ เอกธมฺมสฺส, สพฺเพว วสมนฺวคูติ.}
\hangingbody{จิตฺเตน นียติ โลโก, จิตฺเตน ปริกสฺสติ.}
\hangingbody{จิตฺตสฺส เอกธมฺมสฺส, สพฺเพว วสมนฺวคูติ.}
\csromanhangingend

\csromanheader{2}{3. ตณฺหาสุตฺต}
\csromanhangingbegin
\hangingheaditem{63}{เกนสฺสุ นียติ โลโก, เกนสฺสุ ปริกสฺสติ.}
\hangingbody{กิสฺสสฺสุ เอกธมฺมสฺส, สพฺเพว วสมนฺวคูติ.}
\hangingbody{ตณฺหาย นียติ โลโก, ตณฺหาย ปริกสฺสติ.}
\hangingbody{ตณฺหาย เอกธมฺมสฺส, สพฺเพว วสมนฺวคูติ.}
\csromanhangingend

\csromanheader{2}{4. สํโยชนสุตฺต}
\csromanhangingbegin
\hangingheaditem{64}{\csromanfolio{37}กิํสุสํโยชโน โลโก, กิํสุ ตสฺส วิจารณํ.}
\hangingbody{กิสฺสสฺสุ วิปฺปหาเนน, นิพฺพานํ อิติ วุจฺจตีติ.}
\hangingbody{นนฺทีสํโยชโน1 โลโก, วิตกฺกสฺส วิจารณํ.}
\hangingbody{ตณฺหาย วิปฺปหาเนน, นิพฺพานํ อิติ วุจฺจตีติ.}
\csromanhangingend

\csromanheader{2}{5. พนฺธนสุตฺต}
\csromanhangingbegin
\hangingheaditem{65}{กิํสุสมฺพนฺธโน โลโก, กิํสุ ตสฺส วิจารณํ.}
\hangingbody{กิสฺสสฺสุ วิปฺปหาเนน, สพฺพํ ฉินฺทติ พนฺธนนฺติ.}
\hangingbody{นนฺทีสมฺพนฺธโน โลโก, วิตกฺกสฺส วิจารณํ.}
\hangingbody{ตณฺหาย วิปฺปหาเนน, สพฺพํ ฉินฺทติ พนฺธนนฺติ.}
\csromanhangingend

\csromanheader{2}{6. อตฺตหตสุตฺต}
\csromanhangingbegin
\hangingheaditem{66}{เกนสฺสุพฺภาหโต โลโก, เกนสฺสุ ปริวาริโต.}
\hangingbody{เกน สลฺเลน โอติณฺโณ, กิสฺส ธูปายิโต สทาติ.}
\hangingbody{มจฺจุนาพฺภาหโต โลโก, ชราย ปริวาริโต.}
\hangingbody{ตณฺหาสลฺเลน โอติณฺโณ, อิจฺฉาธูปายิโต สทาติ.}
\csromanhangingend

\csromanheader{2}{7. อุฑฺฑิตสุตฺต}
\csromanhangingbegin
\hangingheaditem{67}{เกนสฺสุ อุฑฺฑิโต โลโก, เกนสฺสุ ปริวาริโต.}
\hangingbody{เกนสฺสุ ปิหิโต โลโก, กิสฺมิํ โลโก ปติฏฺฐิโตติ.}
\hangingbody{ตณฺหาย อุฑฺฑิโต โลโก, ชราย ปริวาริโต.}
\hangingbody{มจฺจุนา ปิหิโต โลโก, ทุกฺเข โลโก ปติฏฺฐิโตติ.}
\csromanhangingend

\csromanheader{2}{8. ปิหิตสุตฺต}
\csromanhangingbegin
\hangingheaditem{68}{เกนสฺสุ ปิหิโต โลโก, กิสฺมิํ โลโก ปติฏฺฐิโต.}
\hangingbody{เกนสฺสุ อุฑฺฑิโต โลโก, เกนสฺสุ ปริวาริโตติ.}
\csromanhangingend

\csromangathameasure{มจฺจุนา ปิหิโต โลโก, ทุกฺเข โลโก ปติฏฺฐิโต. \\ ตณฺหาย อุฑฺฑิโต โลโก, ชราย ปริวาริโตติ.}
\csromangathagroup{\csromangathastanza{\csromanfolio{38}มจฺจุนา ปิหิโต โลโก, ทุกฺเข โลโก ปติฏฺฐิโต. \\ ตณฺหาย อุฑฺฑิโต โลโก, ชราย ปริวาริโตติ.}}

\csromanheader{2}{9. อิจฺฉาสุตฺต}
\csromanhangingbegin
\hangingheaditem{69}{เกนสฺสุ พชฺฌตี โลโก, กิสฺส วินยาย มุจฺจติ.}
\hangingbody{กิสฺสสฺสุ วิปฺปหาเนน, สพฺพํ ฉินฺทติ พนฺธนนฺติ.}
\hangingbody{อิจฺฉาย พชฺฌตี โลโก, อิจฺฉาวินยาย มุจฺจติ.}
\hangingbody{อิจฺฉาย วิปฺปหาเนน, สพฺพํ ฉินฺทติ พนฺธนนฺติ.}
\csromanhangingend

\csromanheader{2}{10. โลกสุตฺต}
\csromanhangingbegin
\hangingheaditem{70}{กิสฺมิํ โลโก สมุปฺปนฺโน, กิสฺมิํ กุพฺพติ สนฺถวํ.}
\hangingbody{กิสฺส โลโก อุปาทาย, กิสฺมิํ โลโก วิหญฺญตีติ.}
\hangingbody{ฉสุ โลโก สมุปฺปนฺโน, ฉสุ กุพฺพติ สนฺถวํ.}
\hangingbody{ฉนฺนเมว อุปาทาย, ฉสุ โลโก วิหญฺญตีติ.}
\csromanhangingend

\vspace{\tipitakaparskip}
\csromancenter{อทฺธวคฺโค\footnote{อนฺววคฺโค (สี)} สตฺตโม.}
\titlehead{\textbf{ตสฺสุทฺทานํ}}
\csromangathameasure{นามํ จิตฺตญฺจ ตณฺหา จ, สํโยชนญฺจ พนฺธนา. \\ อพฺภาหตุฑฺฑิโต ปิหิโต, อิจฺฉา โลเกน เต ทสาติ.}
\csromangathagroup{\csromangathastanza{นามํ จิตฺตญฺจ ตณฺหา จ, สํโยชนญฺจ พนฺธนา. \\ อพฺภาหตุฑฺฑิโต ปิหิโต, อิจฺฉา โลเกน เต ทสาติ.}}

\chapterhead{8. เฉตฺวาวคฺค}
\csromanheader{2}{1. เฉตฺวาสุตฺต}
\proseitem{71}{สาวตฺถินิทานํ.\csromanspacer{} เอกมนฺตํ ฐิตา โข สา เทวตา ภควนฺตํ คาถาย อชฺฌภาสิ\csromandash{}}

\csromangathameasure{กิํสุ เฉตฺวา สุขํ เสติ, กิํสุ เฉตฺวา น โสจติ. \\ กิสฺสสฺสุ เอกธมฺมสฺส, วธํ โรเจสิ โคตมาติ. \\ โกธํ เฉตฺวา สุขํ เสติ, โกธํ เฉตฺวา น โสจติ. \\ โกธสฺส วิสมูลสฺส, มธุรคฺคสฺส เทวเต.}
\csromangathagroup{\csromangathastanza{กิํสุ เฉตฺวา\footnote{ฌตฺวา (สี), ฆตฺวา (สฺยา, กํ) เอวมุปริปิ.} สุขํ เสติ, กิํสุ เฉตฺวา น โสจติ. \\ กิสฺสสฺสุ เอกธมฺมสฺส, วธํ โรเจสิ โคตมาติ.}\csromangathastanza{\csromanfolio{39}โกธํ เฉตฺวา สุขํ เสติ, โกธํ เฉตฺวา น โสจติ. \\ โกธสฺส วิสมูลสฺส, มธุรคฺคสฺส เทวเต.}}

\prose{วธํ อริยา ปสํสนฺติ, ตํ หิ เฉตฺวา น โสจตีติ. \textbf{2.}\csromanspacer{}\textbf{ รถสุตฺต}}

\proseitem{72}{กิํสุ รถสฺส ปญฺญาณํ, กิํสุ ปญฺญาณมคฺคิโน.}

\csromangathameasure{กิํสุ รฏฺฐสฺส ปญฺญาณํ, กิํสุ ปญฺญาณมิตฺถิยาติ. \\ ธโช รถสฺส ปญฺญาณํ, ธูโม ปญฺญาณมคฺคิโน.}
\csromangathagroup{\csromangathastanza{กิํสุ รฏฺฐสฺส ปญฺญาณํ, กิํสุ ปญฺญาณมิตฺถิยาติ. \\ ธโช รถสฺส ปญฺญาณํ, ธูโม ปญฺญาณมคฺคิโน.}}

\prose{ราชา รฏฺฐสฺส ปญฺญาณํ, ภตฺตา ปญฺญาณมิตฺถิยาติ. \textbf{3.}\csromanspacer{}\textbf{ วิตฺตสุตฺต}}

\proseitem{73}{กิํสูธ วิตฺตํ ปุริสสฺส เสฏฺฐํ,}

\prose{กิํสุ สุจิณฺโณ สุขมาวหติ.}

\csromangathameasure{กิํสุ หเว สาทุตรํ รสานํ, \\ กถํชีวิํ ชีวิตมาหุ เสฏฺฐนฺติ. \\ สทฺธีธ วิตฺตํ ปุริสสฺส เสฏฺฐํ, \\ ธมฺโม สุจิณฺโณ สุขมาวหติ. \\ สจฺจํ หเว สาทุตรํ รสานํ,}
\csromangathagroup{\csromangathastanza{กิํสุ หเว สาทุตรํ\footnote{สาธุตรํ (ก)} รสานํ, \\ กถํชีวิํ\footnote{สาธุตรํ (ก)} ชีวิตมาหุ เสฏฺฐนฺติ. \\ สทฺธีธ วิตฺตํ ปุริสสฺส เสฏฺฐํ, \\ ธมฺโม สุจิณฺโณ สุขมาวหติ.}\csromangathastanza{สจฺจํ หเว สาทุตรํ รสานํ,}}

\csromanheader{2}{ปญฺญาชีวิํ ชีวิตมาหุ เสฏฺฐนฺติ. \textbf{4. วุฏฺฐิสุตฺต}}
\proseitem{74}{กิํสุ อุปฺปตตํ เสฏฺฐํ, กิํสุ นิปตตํ วรํ,}

\csromangathameasure{กิํสุ ปวชมานานํ, กิํสุ ปวทตํ วรนฺติ. \\ พีชํ อุปฺปตตํ เสฏฺฐํ, วุฏฺฐิ นิปตตํ วรา. \\ คาโว ปวชมานานํ, ปุตฺโต ปวทตํ วโรติ. \\ วิชฺชา อุปฺปตตํ เสฏฺฐา, อวิชฺชา นิปตตํ วรา.}
\csromangathagroup{\csromangathastanza{กิํสุ ปวชมานานํ, กิํสุ ปวทตํ วรนฺติ. \\ พีชํ อุปฺปตตํ เสฏฺฐํ, วุฏฺฐิ นิปตตํ วรา.}\csromangathastanza{คาโว ปวชมานานํ, ปุตฺโต ปวทตํ วโรติ. \\ วิชฺชา อุปฺปตตํ เสฏฺฐา, อวิชฺชา นิปตตํ วรา.}}

\prose{สํโฆ ปวชมานานํ, พุทฺโธ ปวทตํ วโรติ. \textbf{5.}\csromanspacer{}\textbf{ ภีตาสุตฺต}}

\csromanhangingbegin
\hangingheaditem{75}{กิํสูธ ภีตา ชนตา อเนกา,}
\hangingbody{มคฺโค จเนกายตนปฺปวุตฺโต.}
\hangingbody{ปุจฺฉามิ ตํ โคตม ภูริปญฺญ,}
\hangingbody{กิสฺมิํ ฐิโต ปรโลกํ น ภาเยติ.}
\csromanhangingend

\csromangathameasure{วาจํ มนญฺจ ปณิธาย สมฺมา, \\ กาเยน ปาปานิ อกุพฺพมาโน. \\ พหฺวนฺนปานํ ฆรมาวสนฺโต, \\ สทฺโธ มุทู สํวิภาคี วทญฺญู. \\ เอเตสุ ธมฺเมสุ ฐิโต จตูสุ. \\ ธมฺเม ฐิโต ปรโลกํ น ภาเยติ.}
\csromangathagroup{\csromangathastanza{\csromanfolio{40}วาจํ มนญฺจ ปณิธาย สมฺมา, \\ กาเยน ปาปานิ อกุพฺพมาโน. \\ พหฺวนฺนปานํ ฆรมาวสนฺโต, \\ สทฺโธ มุทู สํวิภาคี วทญฺญู. \\ เอเตสุ ธมฺเมสุ ฐิโต จตูสุ. \\ ธมฺเม ฐิโต ปรโลกํ น ภาเยติ.}}

\csromanheader{2}{6. นชีรติสุตฺต}
\csromanhangingbegin
\hangingheaditem{76}{กิํ ชีรติ กิํ น ชีรติ, กิํสุ อุปฺปโถติ วุจฺจติ.}
\hangingbody{กิํสุ ธมฺมานํ ปริปนฺโถ, กิํสุ รตฺตินฺทิวกฺขโย.}
\hangingbody{กิํ มลํ พฺรหฺมจริยสฺส, กิํ สินานมโนทกํ.}
\hangingbody{กติ โลกสฺมิํ ฉิทฺทานิ, ยตฺถ วิตฺตํ1 น ติฏฺฐติ.}
\hangingbody{ภควนฺตํ ปุฏฺฐุมาคมฺม, กถํ ชาเนมุ ตํ มยนฺติ.}
\hangingbody{รูปํ ชีรติ มจฺจานํ, นามโคตฺตํ น ชีรติ.}
\hangingbody{ราโค อุปฺปโถติ วุจฺจติ.}
\hangingbody{โลโภ ธมฺมานํ ปริปนฺโถ, วโย รตฺตินฺทิวกฺขโย.}
\hangingbody{อิตฺถี มลํ พฺรหฺมจริยสฺส, เอตฺถายํ สชฺชเต ปชา.}
\hangingbody{ตโป จ พฺรหฺมจริยญฺจ, ตํ สินานมโนทกํ.}
\hangingbody{ฉ โลกสฺมิํ ฉิทฺทานิ, ยตฺถ วิตฺตํ น ติฏฺฐติ.}
\hangingbody{อาลสฺยญฺจ2 ปมาโท จ, อนุฏฺฐานํ อสํยโม.}
\hangingbody{นิทฺทา ตนฺที3 จ เต ฉิทฺเท, สพฺพโส ตํ วิวชฺชเยติ.}
\csromanhangingend

\csromanheader{2}{7. อิสฺสริยสุตฺต}
\csromanhangingbegin
\hangingheaditem{77}{กิํสุ อิสฺสริยํ โลเก, กิํสุ ภณฺฑานมุตฺตมํ.}
\hangingbody{กิํสุ สตฺถมลํ โลเก, กิํสุ โลกสฺมิมพฺพุทํ.}
\hangingbody{กิํสุ หรนฺตํ วาเรนฺติ, หรนฺโต ปน โก ปิโย.}
\hangingbody{กิํสุ ปุนปฺปุนายนฺตํ, อภินนฺทนฺติ ปณฺฑิตาติ.}
\csromanhangingend

\csromangathameasure{วโส อิสฺสริยํ โลเก, อิตฺถี ภณฺฑานมุตฺตมํ. \\ โกโธ สตฺถมลํ โลเก, โจรา โลกสฺมิมพฺพุทา. \\ โจรํ หรนฺตํ วาเรนฺติ, หรนฺโต สมโณ ปิโย. \\ สมณํ ปุนปฺปุนายนฺตํ, อภินนฺทนฺติ ปณฺฑิตาติ.}
\csromangathagroup{\csromangathastanza{\csromanfolio{41}วโส อิสฺสริยํ โลเก, อิตฺถี ภณฺฑานมุตฺตมํ. \\ โกโธ สตฺถมลํ โลเก, โจรา โลกสฺมิมพฺพุทา.}\csromangathastanza{โจรํ หรนฺตํ วาเรนฺติ, หรนฺโต สมโณ ปิโย. \\ สมณํ ปุนปฺปุนายนฺตํ, อภินนฺทนฺติ ปณฺฑิตาติ.}}

\csromanheader{2}{8. กามสุตฺต}
\csromanhangingbegin
\hangingheaditem{78}{กิมตฺถกาโม น ทเท, กิํ มจฺโจ น ปริจฺจเช.}
\hangingbody{กิํสุ มุญฺเจยฺย กลฺยาณํ, ปาปิกํ น จ โมจเยติ.}
\hangingbody{อตฺตานํ น ทเท โปโส, อตฺตานํ น ปริจฺจเช.}
\hangingbody{วาจํ มุญฺเจยฺย กลฺยาณํ, ปาปิกญฺจ น โมจเยติ.}
\csromanhangingend

\csromanheader{2}{9. ปาเถยฺยสุตฺต}
\csromanhangingbegin
\hangingheaditem{79}{กิํสุ พนฺธติ ปาเถยฺยํ, กิํสุ โภคานมาสโย.}
\hangingbody{กิํสุ นรํ ปริกสฺสติ, กิํสุ โลกสฺมิ ทุชฺชหํ.}
\hangingbody{กิสฺมิํ พทฺธา ปุถู สตฺตา, ปาเสน สกุณี ยถาติ.}
\hangingbody{สทฺธา พนฺธติ ปาเถยฺยํ, สิรี โภคานมาสโย.}
\hangingbody{อิจฺฉา นรํ ปริกสฺสติ, อิจฺฉา โลกสฺมิ ทุชฺชหา.}
\hangingbody{อิจฺฉาพทฺธา ปุถู สตฺตา, ปาเสน สกุณี ยถาติ.}
\csromanhangingend

\csromanheader{2}{10. ปชฺโชตสุตฺต}
\csromanhangingbegin
\hangingheaditem{80}{กิํสุ โลกสฺมิ ปชฺโชโต, กิํสุ โลกสฺมิ ชาคโร.}
\hangingbody{กิํสุ กมฺเม สชีวานํ, กิมสฺส อิริยาปโถ.}
\hangingbody{กิํสุ อลสํ อนลสญฺจ1, มาตา ปุตฺตํว โปสติ.}
\hangingbody{กิํ ภูตา อุปชีวนฺติ, เย ปาณา ปถวิสฺสิตาติ.}
\hangingbody{ปญฺญา โลกสฺมิ ปชฺโชโต, สติ โลกสฺมิ ชาครา.}
\hangingbody{คาโว กมฺเม สชีวานํ, สีตสฺส อิริยาปโถ.}
\hangingbody{วุฏฺฐิ อลสํ อนลสญฺจ, มาตา ปุตฺตํว โปสติ.}
\hangingbody{วุฏฺฐิํ ภูตา อุปชีวนฺติ, เย ปาณา ปถวิสฺสิตาติ.}
\csromanhangingend

\csromanheader{2}{11. อรณสุตฺต}
\csromanhangingbegin
\hangingheaditem{81}{\csromanfolio{42}เกสูธ อรณา โลเก, เกสํ วุสิตํ น นสฺสติ.}
\hangingbody{เกธ อิจฺฉํ ปริชานนฺติ, เกสํ โภชิสฺสิยํ สทา.}
\hangingbody{กิํสุ มาตา ปิตา ภาตา, วนฺทนฺติ นํ ปติฏฺฐิตํ.}
\hangingbody{กิํสุ อิธ ชาติหีนํ, อภิวาเทนฺติ ขตฺติยาติ.}
\hangingbody{สมณีธ อรณา โลเก, สมณานํ วุสิตํ น นสฺสติ.}
\hangingbody{สมณา อิจฺฉํ ปริชานนฺติ, สมณานํ โภชิสฺสิยํ สทา.}
\hangingbody{สมณํ มาตา ปิตา ภาตา, วนฺทนฺติ นํ ปติฏฺฐิตํ.}
\hangingbody{สมณีธ ชาติหีนํ, อภิวาเทนฺติ ขตฺติยาติ.}
\csromanhangingend

\vspace{\tipitakaparskip}
\csromancenter{เฉตฺวาวคฺโค อฏฺฐโม.}
\titlehead{\textbf{ตสฺสุทฺทานํ}}
\csromangathameasure{เฉตฺวา รถญฺจ จิตฺตญฺจ, วุฏฺฐิ ภีตา นชีรติ. \\ อิสฺสรํ กามํ ปาเถยฺยํ, ปชฺโชโต อรเณน จาติ. \\ \textbf{เทวตาสํยุตฺตํ สมตฺตํ.}}
\csromangathagroup{\csromangathastanza{เฉตฺวา รถญฺจ จิตฺตญฺจ, วุฏฺฐิ ภีตา นชีรติ. \\ อิสฺสรํ กามํ ปาเถยฺยํ, ปชฺโชโต อรเณน จาติ. \\ \textbf{เทวตาสํยุตฺตํ สมตฺตํ.}}}

\csromanheader{2}{2. เทวปุตฺตสํยุตฺต}
\csromanheader{1}{1. ปฐมวคฺค}
\csromanheader{2}{1. ปฐมกสฺสปสุตฺต}
\proseitem{82}{\csromanfolio{43}เอวํ เม สุตํ\csromandash{}เอกํ สมยํ ภควา สาวตฺถิยํ วิหรติ เชตวเน อนาถปิณฺฑิกสฺส อาราเม.\csromanspacer{} อถ โข กสฺสโป เทวปุตฺโต อภิกฺกนฺตาย รตฺติยา อภิกฺกนฺตวณฺโณ เกวลกปฺปํ เชตวนํ โอภาเสตฺวา เยน ภควา เตนุปสงฺกมิ, อุปสงฺกมิตฺวา ภควนฺตํ อภิวาเทตฺวา เอกมนฺตํ อฏฺฐาสิ, เอกมนฺตํ ฐิโต โข กสฺสโป เทวปุตฺโต ภควนฺตํ เอตทโวจ “ภิกฺขุํ ภควา ปกาเสสิ, โน จ ภิกฺขุโน อนุสาสนฺ”ติ.\csromanspacer{} เตน หิ กสฺสป ตญฺเญเวตฺถ ปฏิภาตูติ.}

\csromangathameasure{สุภาสิตสฺส สิกฺเขถ, สมณูปาสนสฺส จ. \\ เอกาสนสฺส จ รโห, จิตฺตวูปสมสฺส จาติ.}
\csromangathagroup{\csromangathastanza{สุภาสิตสฺส สิกฺเขถ, สมณูปาสนสฺส จ. \\ เอกาสนสฺส จ รโห, จิตฺตวูปสมสฺส จาติ.}}

\prose{อิทมโวจ กสฺสโป เทวปุตฺโต, สมนุญฺโญ สตฺถา อโหสิ.\csromanspacer{} อถ โข กสฺสโป เทวปุตฺโต “สมนุญฺโญ เม สตฺถา”ติ ภควนฺตํ อภิวาเทตฺวา ปทกฺขิณํ กตฺวา ตตฺเถวนฺตรธายีติ.}

\csromanheader{2}{2. ทุติยกสฺสปสุตฺต}
\proseitem{83}{สาวตฺถินิทานํ.\csromanspacer{} เอกมนฺตํ ฐิโต โข กสฺสโป เทวปุตฺโต ภควโต สนฺติเก อิมํ คาถํ อภาสิ\csromandash{}}

\csromangathameasure{ภิกฺขุ สิยา ฌายี วิมุตฺตจิตฺโต, \\ อากงฺเข เจ หทยสฺสานุปตฺติํ. \\ โลกสฺส ญตฺวา อุทยพฺพยญฺจ, \\ สุเจตโส อนิสฺสิโต ตทานิสํโสติ.}
\csromangathagroup{\csromangathastanza{ภิกฺขุ สิยา ฌายี วิมุตฺตจิตฺโต, \\ อากงฺเข เจ หทยสฺสานุปตฺติํ. \\ โลกสฺส ญตฺวา อุทยพฺพยญฺจ, \\ สุเจตโส อนิสฺสิโต ตทานิสํโสติ.}}

\csromanheader{2}{3. มาฆสุตฺต}
\proseitem{84}{สาวตฺถินิทานํ.\csromanspacer{} อถ โข มาโฆ เทวปุตฺโต อภิกฺกนฺตาย รตฺติยา อภิกฺกนฺตวณฺโณ เกวลกปฺปํ เชตวนํ โอภาเสตฺวา เยน ภควา เตนุปสงฺกมิ, อุปสงฺกมิตฺวา ภควนฺตํ อภิวาเทตฺวา เอกมนฺตํ \csromanfolio{44}อฏฺฐาสิ, เอกมนฺตํ ฐิโต โข มาโฆ เทวปุตฺโต ภควนฺตํ คาถาย อชฺฌภาสิ\csromandash{}}

\csromangathameasure{กิํสุ เฉตฺวา สุขํ เสติ, กิํสุ เฉตฺวา น โสจติ. \\ กิสฺสสฺสุ เอกธมฺมสฺส, วธํ โรเจสิ โคตมาติ. \\ โกธํ เฉตฺวา สุขํ เสติ, โกธํ เฉตฺวา น โสจติ. \\ โกธสฺส วิสมูลสฺส, มธุรคฺคสฺส วตฺรภู. \\ วธํ อริยา ปสํสนฺติ, ตํ หิ เฉตฺวา น โสจตีติ.}
\csromangathagroup{\csromangathastanza{กิํสุ เฉตฺวา สุขํ เสติ, กิํสุ เฉตฺวา น โสจติ. \\ กิสฺสสฺสุ เอกธมฺมสฺส, วธํ โรเจสิ โคตมาติ.}\csromangathastanza{โกธํ เฉตฺวา สุขํ เสติ, โกธํ เฉตฺวา น โสจติ. \\ โกธสฺส วิสมูลสฺส, มธุรคฺคสฺส วตฺรภู. \\ วธํ อริยา ปสํสนฺติ, ตํ หิ เฉตฺวา น โสจตีติ.}}

\csromanheader{2}{4. มาคธสุตฺต}
\proseitem{85}{สาวตฺถินิทานํ.\csromanspacer{} เอกมนฺตํ ฐิโต โข มาคโธ เทวปุตฺโต ภควนฺตํ คาถาย อชฺฌภาสิ\csromandash{}}

\csromangathameasure{กติ โลกสฺมิํ ปชฺโชตา, เยหิ โลโก ปกาสติ. \\ ภวนฺตํ ปุฏฺฐุมาคมฺม, กถํ ชาเนมุ ตํ มยนฺติ. \\ จตฺตาโร โลเก ปชฺโชตา, ปญฺจเมตฺถ น วิชฺชติ. \\ ทิวา ตปติ อาทิจฺโจ, รตฺติมาภาติ จนฺทิมา. \\ อถ อคฺคิ ทิวารตฺติํ, ตตฺถ ตตฺถ ปกาสติ. \\ สมฺพุทฺโธ ตปตํ เสฏฺโฐ, เอสา อาภา อนุตฺตราติ.}
\csromangathagroup{\csromangathastanza{กติ โลกสฺมิํ ปชฺโชตา, เยหิ โลโก ปกาสติ. \\ ภวนฺตํ ปุฏฺฐุมาคมฺม, กถํ ชาเนมุ ตํ มยนฺติ.}\csromangathastanza{จตฺตาโร โลเก ปชฺโชตา, ปญฺจเมตฺถ น วิชฺชติ. \\ ทิวา ตปติ อาทิจฺโจ, รตฺติมาภาติ จนฺทิมา.}\csromangathastanza{อถ อคฺคิ ทิวารตฺติํ, ตตฺถ ตตฺถ ปกาสติ. \\ สมฺพุทฺโธ ตปตํ เสฏฺโฐ, เอสา อาภา อนุตฺตราติ.}}

\csromanheader{2}{5. ทามลิสุตฺต}
\proseitem{86}{สาวตฺถินิทานํ.\csromanspacer{} อถ โข ทามลิ เทวปุตฺโต อภิกฺกนฺตาย รตฺติยา อภิกฺกนฺตวณฺโณ เกวลกปฺปํ เชตวนํ โอภาเสตฺวา เยน ภควา เตนุปสงฺกมิ, อุปสงฺกมิตฺวา ภควนฺตํ อภิวาเทตฺวา เอกมนฺตํ อฏฺฐาสิ, เอกมนฺตํ ฐิโต โข ทามลิ เทวปุตฺโต ภควโต สนฺติเก อิมํ คาถํ อภาสิ\csromandash{}}

\csromangathameasure{กรณียเมตํ พฺราหฺมเณน, ปธานํ อกิลาสุนา. \\ กามานํ วิปฺปหาเนน, น เตนาสีสเต ภวํติ.}
\csromangathagroup{\csromangathastanza{กรณียเมตํ พฺราหฺมเณน, ปธานํ อกิลาสุนา. \\ กามานํ วิปฺปหาเนน, น เตนาสีสเต ภวํติ.}}

\vspace{\tipitakaparskip}
\csromancenter{นตฺถิ กิจฺจํ พฺราหฺมณสฺส, (ทามลีติ ภควา)}
\csromancenter{กตกิจฺโจ หิ พฺราหฺมโณ.}
\csromanheader{2}{2. เทวปุตฺตสํยุตฺต}
\csromangathameasure{ยาว น คาธํ ลภติ นทีสุ, \\ อายูหติ สพฺพคตฺเตภิ ชนฺตุ. \\ คาธญฺจ ลทฺธาน ถเล ฐิโต โย, \\ นายูหตี ปารคโต หิ โสว. \\ เอสูปมา ทามลิ พฺราหฺมณสฺส, \\ ขีณาสวสฺส นิปกสฺส ฌายิโน. \\ ปปฺปุยฺย ชาติมรณสฺส อนฺตํ, \\ นายูหตี ปารคโต หิ โสติ.}
\csromangathagroup{\csromangathastanza{\csromanfolio{45}ยาว น คาธํ ลภติ นทีสุ, \\ อายูหติ สพฺพคตฺเตภิ ชนฺตุ. \\ คาธญฺจ ลทฺธาน ถเล ฐิโต โย, \\ นายูหตี ปารคโต หิ โสว\footnote{โสติ (สี, อิ, ก), โหติ (สฺยา, กํ), โส (?)}.}\csromangathastanza{เอสูปมา ทามลิ พฺราหฺมณสฺส, \\ ขีณาสวสฺส นิปกสฺส ฌายิโน. \\ ปปฺปุยฺย ชาติมรณสฺส อนฺตํ, \\ นายูหตี ปารคโต หิ โสติ\footnote{โหตีติ (สฺยา, กํ)}.}}

\csromanheader{2}{6. กามทสุตฺต}
\proseitem{87}{สาวตฺถินิทานํ.\csromanspacer{} เอกมนฺตํ ฐิโต โข กามโท เทวปุตฺโต ภควนฺตํ เอตทโวจ “ทุกฺกรํ ภควา, สุทุกฺกรํ ภควา”ติ.}

\csromanheader{2}{ทุกฺกรํ วาปิ กโรนฺติ (กามทาติ ภควา)}
\prose{เสขา สีลสมาหิตา.}

\csromangathameasure{ฐิตตฺตา อนคาริยุเปตสฺส, \\ ตุฏฺฐิ โหติ สุขาวหาติ.}
\csromangathagroup{\csromangathastanza{ฐิตตฺตา อนคาริยุเปตสฺส, \\ ตุฏฺฐิ โหติ สุขาวหาติ.}}

\prose{“ทุลฺลภา ภควา ยทิทํ ตุฏฺฐิ” ติ.}

\csromanmatikaanchor{mtk.0}
\csromanmatikaanchor{mtk.157}
\csromantocmarknopage{chapter}{นิทานวคฺคสํยุตฺตปาฬิ}{mtk.0}
\bookmark[dest={mtk.0},level=0]{นิทานวคฺคสํยุตฺตปาฬิ}
\prose{ทุลฺลภํ วาปิ ลภนฺติ, (กามทาติ ภควา)}

\prose{จิตฺตวูปสเม รตา.}

\csromangathameasure{เยสํ ทิวา จ รตฺโต จ, \\ ภาวนาย รโต มโนติ.}
\csromangathagroup{\csromangathastanza{เยสํ ทิวา จ รตฺโต จ, \\ ภาวนาย รโต มโนติ.}}

\prose{“ทุสฺสมาทหํ ภควา ยทิทํ จิตฺตนฺ” ติ.}

\vspace{\tipitakaparskip}
\csromancenter{ทุสฺสมาทหํ วาปิ สมาทหนฺติ, (กามทาติ ภควา)}
\prose{อินฺทฺริยูปสเม รตา.}

\csromangathameasure{เต เฉตฺวา มจฺจุโน ชาลํ, \\ อริยา คจฺฉนฺติ กามทาติ.}
\csromangathagroup{\csromangathastanza{เต เฉตฺวา มจฺจุโน ชาลํ, \\ อริยา คจฺฉนฺติ กามทาติ.}}

\prose{\csromanfolio{46}“ทุคฺคโม ภควา วิสโม มคฺโค”ติ.}

\csromangathameasure{ทุคฺคเม วิสเม วาปิ, อริยา คจฺฉนฺติ กามท. \\ อนริยา วิสเม มคฺเค, ปปตนฺติ อวํสิรา. \\ อริยานํ สโม มคฺโค, อริยา หิ วิสเม สมาติ.}
\csromangathagroup{\csromangathastanza{ทุคฺคเม วิสเม วาปิ, อริยา คจฺฉนฺติ กามท. \\ อนริยา วิสเม มคฺเค, ปปตนฺติ อวํสิรา. \\ อริยานํ สโม มคฺโค, อริยา หิ วิสเม สมาติ.}}

\csromanheader{2}{7. ปญฺจาลจณฺฑสุตฺต}
\proseitem{88}{สาวตฺถินิทานํ.\csromanspacer{} เอกมนฺตํ ฐิโต โข ปญฺจาลจณฺโฑ เทวปุตฺโต ภควโต สนฺติเก อิมํ คาถํ อภาสิ\csromandash{}}

\csromangathameasure{สมฺพาเธ วต โอกาสํ, อวินฺทิ ภูริเมธโส. \\ โย ฌานมพุชฺฌิ พุทฺโธ, ปฏิลีนนิสโภ มุนีติ.}
\csromangathagroup{\csromangathastanza{สมฺพาเธ วต โอกาสํ, อวินฺทิ ภูริเมธโส. \\ โย ฌานมพุชฺฌิ\footnote{ฌานมพุธา (ก-สี), ฌานมพุทฺธิ (สฺยา, กํ, อิ, ก)} พุทฺโธ, ปฏิลีนนิสโภ มุนีติ.}}

\vspace{\tipitakaparskip}
\csromancenter{สมฺพาเธ วาปิ วินฺทนฺติ, (ปญฺจาลจณฺฑาติ ภควา)}
\prose{ธมฺมํ นิพฺพานปตฺติยา.}

\csromangathameasure{เย สติํ ปจฺจลตฺถํสุ, \\ สมฺมา เต สุสมาหิตาติ.}
\csromangathagroup{\csromangathastanza{เย สติํ ปจฺจลตฺถํสุ, \\ สมฺมา เต สุสมาหิตาติ.}}

\csromanheader{2}{8. ตายนสุตฺต}
\proseitem{89}{สาวตฺถินิทานํ.\csromanspacer{} อถ โข ตายโน เทวปุตฺโต ปุราณติตฺถกโร อภิกฺกนฺตาย รตฺติยา อภิกฺกนฺตาวณฺโณ เกวลกปฺปํ เชตวนํ โอภาเสตฺวา เยน ภควา เตนุปสงฺกมิ, อุปสงฺกมิตฺวา ภควนฺตํ อภิวาเทตฺวา เอกมนฺตํ อฏฺฐาสิ, เอกมนฺตํ ฐิโต โข ตายโน เทวปุตฺโต ภควโต สนฺติเก อิมา คาถาโย อภาสิ\csromandash{}}

\csromangathameasure{ฉินฺท โสตํ ปรกมฺม, กาเม ปนุท พฺราหฺมณ. \\ นปฺปหาย มุนี กาเม, เนกตฺตมุปปชฺชติ. \\ กยิรา เจ กยิราเถนํ, ทฬฺหเมนํ ปรกฺกเม. \\ สิถิโล หิ ปริพฺพาโช, ภิยฺโย อากิรเต รชํ. \\ อกตํ ทุกฺกฏํ เสยฺโย, ปจฺฉา ตปติ ทุกฺกฏํ. \\ กตญฺจ สุกตํ เสยฺโย, ยํ กตฺวา นานุตปฺปติ.}
\csromangathagroup{\csromangathastanza{ฉินฺท โสตํ ปรกมฺม, กาเม ปนุท พฺราหฺมณ. \\ นปฺปหาย มุนี กาเม, เนกตฺตมุปปชฺชติ.}\csromangathastanza{กยิรา เจ กยิราเถนํ, ทฬฺหเมนํ ปรกฺกเม. \\ สิถิโล หิ ปริพฺพาโช, ภิยฺโย อากิรเต รชํ.}\csromangathastanza{อกตํ ทุกฺกฏํ\footnote{ทุกฺกตํ (สี, อิ)} เสยฺโย, ปจฺฉา ตปติ ทุกฺกฏํ. \\ กตญฺจ สุกตํ เสยฺโย, ยํ กตฺวา นานุตปฺปติ.}}

\csromanheader{2}{2. เทวปุตฺตสํยุตฺต}
\csromangathameasure{กุโส ยถา ทุคฺคหิโต, หตฺถเมวานุกนฺตติ. \\ สามญฺญํ ทุปฺปรามฏฺฐํ, นิรยายูปกฑฺฒติ. \\ ยํ กิญฺจิ สิถิลํ กมฺมํ, สํกิลิฏฺฐญฺจ ยํ วตํ. \\ สงฺกสฺสรํ พฺรหฺมจริยํ, น ตํ โหติ มหปฺผลนฺติ.}
\csromangathagroup{\csromangathastanza{\csromanfolio{47}กุโส ยถา ทุคฺคหิโต, หตฺถเมวานุกนฺตติ. \\ สามญฺญํ ทุปฺปรามฏฺฐํ, นิรยายูปกฑฺฒติ.}\csromangathastanza{ยํ กิญฺจิ สิถิลํ กมฺมํ, สํกิลิฏฺฐญฺจ ยํ วตํ. \\ สงฺกสฺสรํ พฺรหฺมจริยํ, น ตํ โหติ มหปฺผลนฺติ.}}

\prose{อิทมโวจ ตายโน เทวปุตฺโต, อิทํ วตฺวา ภควนฺตํ อภิวาเทตฺวา ปทกฺขิณํ กตฺวา ตตฺเถวนฺตรธายีติ.}

\prose{อถ โข ภควา ตสฺสา รตฺติยา อจฺจเยน ภิกฺขู อามนฺเตสิ “อิมํ ภิกฺขเว รตฺติํ ตายโน นาม เทวปุตฺโต ปุราณติตฺถกโร อภิกฺกนฺตาย รตฺติยา อภิกฺกนฺตวณฺโณ เกวลกปฺปํ เชตวนํ โอภาเสตฺวา เยนาหํ เตนุปสงฺกมิ, อุปสงฺกมิตฺวา มํ อภิวาเทตฺวา เอกมนฺตํ อฏฺฐาสิ, เอกมนฺตํ ฐิโต โข ภิกฺขเว ตายโน เทวปุตฺโต มม สนฺติเก อิมา คาถาโย อภาสิ\csromandash{}}

\csromangathameasure{ฉินฺท โสตํ ปรกฺกมฺม, กาเม ปนุท พฺราหฺมณ. \\ นปฺปหาย มุนี กาเม, เนกตฺตมุปปชฺชติ. \\ กยิรา เจ กยิราเถนํ, ทฬฺหเมนํ ปรกฺกเม. \\ สิถิโล หิ ปริพฺพาโช, ภิยฺโย อากิรเต รชํ. \\ อกตํ ทุกฺกฏํ เสยฺโย, ปจฺฉา ตปติ ทุกฺกฏํ. \\ กตญฺจ สุกตํ เสยฺโย, ยํ กตฺวา นานุตปฺปติ. \\ กุโส ยถา ทุคฺคหิโต, หตฺถเมวานุกนฺตติ. \\ สามญฺญํ ทุปฺปรามฏฺฐํ, นิรยายูปกฑฺฒติ. \\ ยํ กิญฺจิ สิถิลํ กมฺมํ, สํกิลิฏฺฐญฺจ ยํ วตํ. \\ สงฺกสฺสรํ พฺรหฺมจริยํ, น ตํ โหติ มหปฺผลนฺติ.}
\csromangathagroup{\csromangathastanza{ฉินฺท โสตํ ปรกฺกมฺม, กาเม ปนุท พฺราหฺมณ. \\ นปฺปหาย มุนี กาเม, เนกตฺตมุปปชฺชติ.}\csromangathastanza{กยิรา เจ กยิราเถนํ, ทฬฺหเมนํ ปรกฺกเม. \\ สิถิโล หิ ปริพฺพาโช, ภิยฺโย อากิรเต รชํ.}\csromangathastanza{อกตํ ทุกฺกฏํ เสยฺโย, ปจฺฉา ตปติ ทุกฺกฏํ. \\ กตญฺจ สุกตํ เสยฺโย, ยํ กตฺวา นานุตปฺปติ.}\csromangathastanza{กุโส ยถา ทุคฺคหิโต, หตฺถเมวานุกนฺตติ. \\ สามญฺญํ ทุปฺปรามฏฺฐํ, นิรยายูปกฑฺฒติ.}\csromangathastanza{ยํ กิญฺจิ สิถิลํ กมฺมํ, สํกิลิฏฺฐญฺจ ยํ วตํ. \\ สงฺกสฺสรํ พฺรหฺมจริยํ, น ตํ โหติ มหปฺผลนฺติ.}}

\prose{อิทมโวจ ภิกฺขเว ตายโน เทวปุตฺโต, อิทํ วตฺวา มํ อภิวาเทตฺวา ปทกฺขิณํ กตฺวา ตตฺเถวนฺตรธายิ.\csromanspacer{} อุคฺคณฺหาถ ภิกฺขเว ตายนคาถา, ปริยาปุณาถ ภิกฺขเว ตายนคาถา, ธาเรถ ภิกฺขเว ตายนคาถา, อตฺถสํหิตา ภิกฺขเว ตายนคาถา อาทิพฺรหฺมจริยิกา”ติ.}

\csromanheader{2}{9. จนฺทิมสุตฺต}
\proseitem{90}{\csromanfolio{48}สาวตฺถินิทานํ.\csromanspacer{} เตน โข ปน สมเยน จนฺทิมา เทวปุตฺโต ราหุนา อสุรินฺเทน คหิโต โหติ.\csromanspacer{} อถ โข จนฺทิมา เทวปุตฺโต ภควนฺตํ อนุสฺสรมาโน ตายํ เวลายํ อิมํ คาถํ อภาสิ\csromandash{}}

\csromangathameasure{นโม เต พุทฺธ วีรตฺถุ, วิปฺปมุตฺโตสิ สพฺพธิ. \\ สมฺพาธปฏิปนฺโนสฺมิ, ตสฺส เม สรณํ ภวาติ.}
\csromangathagroup{\csromangathastanza{นโม เต พุทฺธ วีรตฺถุ, วิปฺปมุตฺโตสิ สพฺพธิ. \\ สมฺพาธปฏิปนฺโนสฺมิ, ตสฺส เม สรณํ ภวาติ.}}

\prose{อถ โข ภควา จนฺทิมํ เทวปุตฺตํ อารพฺภ ราหุํ อสุรินฺทํ คาถาย อชฺฌภาสิ\csromandash{}}

\csromangathameasure{ตถาคตํ อรหนฺตํ, จนฺทิมา สรณํ คโต. \\ ราหุ จนฺทํ ปมุญฺจสฺสุ, พุทฺธา โลกานุกมฺปกาติ.}
\csromangathagroup{\csromangathastanza{ตถาคตํ อรหนฺตํ, จนฺทิมา สรณํ คโต. \\ ราหุ จนฺทํ ปมุญฺจสฺสุ, พุทฺธา โลกานุกมฺปกาติ.}}

\prose{อถ โข ราหุ อสุรินฺโท จนฺทิมํ เทวปุตฺตํ มุญฺจิตฺวา ตรมานรูโป เยน เวปจิตฺติ อสุรินฺโท เตนุปสงฺกมิ, อุปสงฺกมิตฺวา สํวิคฺโค โลมหฏฺฐชาโต เอกมนฺตํ อฏฺฐาสิ, เอกมนฺตํ ฐิตํ โข ราหุํ อสุรินฺทํ เวปจิตฺติ อสุรินฺโท คาถาย อชฺฌภาสิ\csromandash{}}

\csromangathameasure{กิํ นุ สนฺตรมาโนว, ราหุ จนฺทํ ปมุญฺจสิ. \\ สํวิคฺครูโป อาคมฺม, กิํ นุ ภีโตว ติฏฺฐสีติ. \\ สตฺตธา เม ผเล มุทฺธา, ชีวนฺโต น สุขํ ลเภ. \\ พุทฺธคาถาภิคีโตมฺหิ, โน เจ มุญฺเจยฺย จนฺทิมนฺติ.}
\csromangathagroup{\csromangathastanza{กิํ นุ สนฺตรมาโนว, ราหุ จนฺทํ ปมุญฺจสิ. \\ สํวิคฺครูโป อาคมฺม, กิํ นุ ภีโตว ติฏฺฐสีติ.}\csromangathastanza{สตฺตธา เม ผเล มุทฺธา, ชีวนฺโต น สุขํ ลเภ. \\ พุทฺธคาถาภิคีโตมฺหิ, โน เจ มุญฺเจยฺย จนฺทิมนฺติ.}}

\csromanheader{2}{10. สูริยสุตฺต}
\proseitem{91}{สาวตฺถินิทานํ.\csromanspacer{} เตน โข ปน สมเยน สูริโย เทวปุตฺโต ราหุนา อสุรินฺเทน คหิโต โหติ.\csromanspacer{} อถ โข สูริโย เทวปุตฺโต ภควนฺตํ อนุสฺสรมาโน ตายํ เวลายํ อิมํ คาถํ อภาสิ\csromandash{}}

\csromangathameasure{นโม เต พุทฺธ วีรตฺถุ, วิปฺปมุตฺโตสิ สพฺพธิ. \\ สมฺพาธปฏิปนฺโนสฺมิ, ตสฺส เม สรณํ ภวาติ.}
\csromangathagroup{\csromangathastanza{นโม เต พุทฺธ วีรตฺถุ, วิปฺปมุตฺโตสิ สพฺพธิ. \\ สมฺพาธปฏิปนฺโนสฺมิ, ตสฺส เม สรณํ ภวาติ.}}

\prose{อถ โข ภควา สูริยํ เทวปุตฺตํ อารพฺภ ราหุํ อสุรินฺทํ คาถาหิ อชฺฌภาสิ\csromandash{}}

\csromanheader{2}{2. เทวปุตฺตสํยุตฺต}
\csromangathameasure{ตถาคตํ อรหนฺตํ, สูริโย สรณํ คโต. \\ ราหุ สูริยํ ปมุญฺจสฺสุ, พุทฺธา โลกานุกมฺปกา. \\ โย อนฺธกาเร ตมสิ ปภงฺกโร, \\ เวโรจโน มณฺฑลี อุคฺคเตโช. \\ มา ราหุ คิลี จรมนฺตลิกฺเข, \\ ปชํ มมํ ราหุ ปมุญฺจ สูริยนฺติ.}
\csromangathagroup{\csromangathastanza{\csromanfolio{49}ตถาคตํ อรหนฺตํ, สูริโย สรณํ คโต. \\ ราหุ สูริยํ\footnote{สุริยํ (สี, สฺยา, กํ, อิ)} ปมุญฺจสฺสุ, พุทฺธา โลกานุกมฺปกา.}\csromangathastanza{โย อนฺธกาเร ตมสิ ปภงฺกโร, \\ เวโรจโน มณฺฑลี อุคฺคเตโช. \\ มา ราหุ คิลี จรมนฺตลิกฺเข, \\ ปชํ มมํ ราหุ ปมุญฺจ สูริยนฺติ.}}

\prose{อถ โข ราหุ อสุรินฺโท สูริยํ เทวปุตฺตํ มุญฺจิตฺวา ตรมานรูโป เยน เวปจิตฺติ อสุรินฺโท เตนุปสงฺกมิ, อุปสงฺกมิตฺวา สํวิคฺโค โลมหฏฺฐชาโต เอกมนฺตํ อฏฺฐาสิ, เอกมนฺตํ ฐิตํ โข ราหุํ อสุรินฺทํ เวปจิตฺติ อสุรินฺโท คาถาย อชฺฌภาสิ\csromandash{}}

\csromangathameasure{กิํ นุ สนฺตรมาโนว, ราหุ สูริยํ ปมุญฺจสิ. \\ สํวิคฺครูโป อาคมฺม, กิํ นุ ภีโตว ติฏฺฐสีติ. \\ สตฺตธา เม ผเล มุทฺธา, ชีวนฺโต น สุขํ ลเภ. \\ พุทฺธคาถาภิคีโตมฺหิ, โน เจ มุญฺเจยฺย สูริยนฺติ.}
\csromangathagroup{\csromangathastanza{กิํ นุ สนฺตรมาโนว, ราหุ สูริยํ ปมุญฺจสิ. \\ สํวิคฺครูโป อาคมฺม, กิํ นุ ภีโตว ติฏฺฐสีติ.}\csromangathastanza{สตฺตธา เม ผเล มุทฺธา, ชีวนฺโต น สุขํ ลเภ. \\ พุทฺธคาถาภิคีโตมฺหิ, โน เจ มุญฺเจยฺย สูริยนฺติ.}}

\vspace{\tipitakaparskip}
\csromancenter{ปฐโม วคฺโค.}
\titlehead{\textbf{ตสฺสุทฺทานํ}}
\csromangathameasure{เทฺว กสฺสปา จ มาโฆ จ, มาคโธ ทามลิ กามโท. \\ ปญฺจาลจณฺโฑ ตายโน, จนฺทิมสูริเยน เต ทสาติ.}
\csromangathagroup{\csromangathastanza{เทฺว กสฺสปา จ มาโฆ จ, มาคโธ ทามลิ กามโท. \\ ปญฺจาลจณฺโฑ ตายโน, จนฺทิมสูริเยน เต ทสาติ.}}

\chapterhead{2. อนาถปิณฺฑิกวคฺค}
\csromanheader{2}{1. จนฺทิมสสุตฺต}
\proseitem{92}{สาวตฺถินิทานํ.\csromanspacer{} อถ โข จนฺทิมโส\footnote{จนฺทิมาโส (ก)} เทวปุตฺโต อภิกฺกนฺตาย รตฺติยา อภิกฺกนฺตวณฺโณ เกวลกปฺปํ เชตวนํ โอภาเสตฺวา เยน ภควา เตนุปสงฺกมิ, อุปสงฺกมิตฺวา ภควนฺตํ อภิวาเทตฺวา เอกมนฺตํ \csromanfolio{50}อฏฺฐาสิ, เอกมนฺตํ ฐิโต โข จนฺทิมโส เทวปุตฺโต ภควโต สนฺติเก อิมํ คาถํ อภาสิ\csromandash{}}

\vspace{\tipitakaparskip}
\csromancenter{เต หิ โสตฺถิํ คมิสฺสนฺติ กจฺเฉวามกเส มคา.}
\vspace{0.5\baselineskip}
\csromangathameasure{ฌานานิ อุปสมฺปชฺช, เอโกทิ นิปกา สตาติ. \\ เต หิ ปารํ คมิสฺสนฺติ, เฉตฺวา ชาลํว อมฺพุโช. \\ ฌานานิ อุปสมฺปชฺช, อปฺปมตฺตา รณญฺชหาติ.}
\csromangathagroup{\csromangathastanza{ฌานานิ อุปสมฺปชฺช, เอโกทิ นิปกา สตาติ. \\ เต หิ ปารํ คมิสฺสนฺติ, เฉตฺวา ชาลํว อมฺพุโช. \\ ฌานานิ อุปสมฺปชฺช, อปฺปมตฺตา รณญฺชหาติ.}}

\csromanheader{2}{2. เวณฺฑุสุตฺต}
\proseitem{93}{เอกมนฺตํ ฐิโต โข เวณฺฑุ\footnote{เวณฺหุ (สี)} เทวปุตฺโต ภควโต สนฺติเก อิมํ คาถํ อภาสิ\csromandash{}}

\csromangathameasure{สุขิตาว เต มนุชา, สุคตํ ปยิรุปาสิย. \\ ยุญฺชํ โคตมสาสเน, อปฺปมตฺตา นุ สิกฺขเรติ.}
\csromangathagroup{\csromangathastanza{สุขิตาว เต\footnote{สุขิตา วต เต (สี, สฺยา, กํ)} มนุชา, สุคตํ ปยิรุปาสิย. \\ ยุญฺชํ\footnote{สุขิตา วต เต (สี, สฺยา, กํ)} โคตมสาสเน, อปฺปมตฺตา นุ สิกฺขเรติ.}}

\prose{เย เม ปวุตฺเต สิฏฺฐิปเท\footnote{สตฺถิปเท (สี, สฺยา, กํ, อิ)}, (เวณฺฑูติ ภควา)}

\prose{อนุสิกฺขนฺติ ฌายิโน.}

\csromangathameasure{กาเล เต อปฺปมชฺชนฺตา, \\ น มจฺจุวสคา สิยุนฺติ.}
\csromangathagroup{\csromangathastanza{กาเล เต อปฺปมชฺชนฺตา, \\ น มจฺจุวสคา สิยุนฺติ.}}

\csromanheader{2}{3. ทีฆลฏฺฐิสุตฺต}
\proseitem{94}{เอวํ เม สุตํ\csromandash{}เอกํ สมยํ ภควา ราชคเห วิหรติ เวฬุวเน กลนฺทกนิวาเป.\csromanspacer{} อถ โข ทีฆลฏฺฐิ เทวปุตฺโต อภิกฺกนฺตาย รตฺติยา อภิกฺกนฺตวณฺโณ เกวลกปฺปํ เวฬุวนํ โอภาเสตฺวา เยน ภควา เตนุปสงฺกมิ, อุปสงฺกมิตฺวา ภควนฺตํ อภิวาเทตฺวา เอกมนฺตํ อฏฺฐาสิ, เอกมนฺตํ ฐิโต โข ทีฆลฏฺฐิ เทวปุตฺโต ภควโต สนฺติเก อิมํ คาถํ อภาสิ\csromandash{}}

\csromangathameasure{ภิกฺขุ สิยา ฌายี วิมุตฺตจิตฺโต, \\ อากงฺเข เจ หทยสฺสานุปตฺติํ. \\ โลกสฺส ญตฺวา อุทยพฺพยญฺจ, \\ สุเจตโส อนิสฺสิโต ตทานิสํโสติ.}
\csromangathagroup{\csromangathastanza{ภิกฺขุ สิยา ฌายี วิมุตฺตจิตฺโต, \\ อากงฺเข เจ หทยสฺสานุปตฺติํ. \\ โลกสฺส ญตฺวา อุทยพฺพยญฺจ, \\ สุเจตโส อนิสฺสิโต ตทานิสํโสติ.}}

\csromanheader{2}{2. เทวปุตฺตสํยุตฺต}
\csromanheader{2}{4. นนฺทนสุตฺต}
\proseitem{95}{\csromanfolio{51}เอกมนฺตํ ฐิโต โข นนฺทโน เทวปุตฺโต ภควนฺตํ คาถาย อชฺฌภาสิ\csromandash{}}

\csromangathameasure{ปุจฺฉามิ ตํ โคตม ภูริปญฺญ, \\ อนาวฏํ ภควโต ญาณทสฺสนํ. \\ กถํวิธํ สีลวนฺตํ วทนฺติ, \\ กถํวิธํ ปญฺญวนฺตํ วทนฺติ. \\ กถํวิโธ ทุกฺขมติจฺจ อิริยติ, \\ กถํวิธํ เทวตา ปูชยนฺตีติ. \\ โย สีลวา ปญฺญวา ภาวิตตฺโต, \\ สมาหิโต ฌานรโต สตีมา. \\ สพฺพสฺส โสกา วิคตา ปหีนา, \\ ขีณาสโว อนฺติมเทหธารี. \\ ตถาวิธํ สีลวนฺตํ วทนฺติ, \\ ตถาวิธํ ปญฺญวนฺตํ วทนฺติ. \\ ตถาวิโธ ทุกฺขมติจฺจ อิริยติ, \\ ตถาวิธํ เทวตา ปูชยนฺตีติ.}
\csromangathagroup{\csromangathastanza{ปุจฺฉามิ ตํ โคตม ภูริปญฺญ, \\ อนาวฏํ ภควโต ญาณทสฺสนํ. \\ กถํวิธํ สีลวนฺตํ วทนฺติ, \\ กถํวิธํ ปญฺญวนฺตํ วทนฺติ.}\csromangathastanza{กถํวิโธ ทุกฺขมติจฺจ อิริยติ, \\ กถํวิธํ เทวตา ปูชยนฺตีติ. \\ โย สีลวา ปญฺญวา ภาวิตตฺโต, \\ สมาหิโต ฌานรโต สตีมา.}\csromangathastanza{สพฺพสฺส โสกา วิคตา ปหีนา, \\ ขีณาสโว อนฺติมเทหธารี. \\ ตถาวิธํ สีลวนฺตํ วทนฺติ, \\ ตถาวิธํ ปญฺญวนฺตํ วทนฺติ. \\ ตถาวิโธ ทุกฺขมติจฺจ อิริยติ, \\ ตถาวิธํ เทวตา ปูชยนฺตีติ.}}

\csromanheader{2}{5. จนฺทนสุตฺต}
\proseitem{96}{เอกมนฺตํ ฐิโต โข จนฺทโน เทวปุตฺโต ภควนฺตํ คาถาย อชฺฌภาสิ\csromandash{}}

\csromangathameasure{“กถํสุ ตรติ โอฆํ, รตฺตินฺทิวมตนฺทิโต. \\ อปฺปติฏฺเฐ อนาลมฺเพ, โก คมฺภีเร น สีทตี”ติ. \\ สพฺพทา สีลสมฺปนฺโน, ปญฺญวา สุสมาหิโต. \\ อารทฺธวีริโย ปหิตตฺโต, โอฆํ ตรติ ทุตฺตรํ. \\ วิรโต กามสญฺญาย, รูปสํโยชนาติโค. \\ นนฺทีราคปริกฺขีโณ, โส คมฺภีเร น สีทตีติ.}
\csromangathagroup{\csromangathastanza{“กถํสุ\footnote{โกสุธ (สี)} ตรติ โอฆํ, รตฺตินฺทิวมตนฺทิโต. \\ อปฺปติฏฺเฐ อนาลมฺเพ, โก คมฺภีเร น สีทตี”ติ.}\csromangathastanza{สพฺพทา สีลสมฺปนฺโน, ปญฺญวา สุสมาหิโต. \\ อารทฺธวีริโย ปหิตตฺโต, โอฆํ ตรติ ทุตฺตรํ.}\csromangathastanza{วิรโต กามสญฺญาย, รูปสํโยชนาติโค. \\ นนฺทีราคปริกฺขีโณ, โส คมฺภีเร น สีทตีติ.}}

\csromanheader{2}{6. วาสุทตฺตสุตฺต}
\proseitem{97}{\csromanfolio{52}เอกมนฺตํ ฐิโต โข วาสุทตฺโต เทวปุตฺโต ภควโต สนฺติเก อิมํ คาถํ อภาสิ\csromandash{}}

\csromangathameasure{“สตฺติยา วิย โอมฏฺโฐ, ฑยฺหมาโนว มตฺถเก. \\ กามราคปฺปหานาย, สโต ภิกฺขุ ปริพฺพเช”ติ. \\ สตฺติยา วิย โอมฏฺโฐ, ฑยฺหมาโนว มตฺถเก. \\ สกฺกายทิฏฺฐิปฺปหานาย, สโต ภิกฺขุ ปริพฺพเชติ.}
\csromangathagroup{\csromangathastanza{“สตฺติยา วิย โอมฏฺโฐ, ฑยฺหมาโนว\footnote{qอยฺหมาเนว (สพฺพตฺถ)} มตฺถเก. \\ กามราคปฺปหานาย, สโต ภิกฺขุ ปริพฺพเช”ติ.}\csromangathastanza{สตฺติยา วิย โอมฏฺโฐ, ฑยฺหมาโนว มตฺถเก. \\ สกฺกายทิฏฺฐิปฺปหานาย, สโต ภิกฺขุ ปริพฺพเชติ.}}

\csromanheader{2}{7. สุพฺรหฺมสุตฺต}
\proseitem{98}{เอกมนฺตํ ฐิโต โข สุพฺรหฺมา เทวปุตฺโต ภควนฺตํ คาถาย อชฺฌภาสิ\csromandash{}}

\csromangathameasure{“นิจฺจํ อุตฺรสฺตมิทํ จิตฺตํ, นิจฺจํ อุพฺพิคฺคมิทํ มโน. \\ อนุปฺปนฺเนสุ กิจฺเฉสุ3, อโถ อุปฺปติเตสุ จ. \\ สเจ อตฺถิ อนุตฺรสฺตํ, ตํ เม อกฺขาหิ ปุจฺฉิโต”ติ. \\ นาญฺญตฺร โพชฺฌา ตปสา, นาญฺญตฺรินฺทฺริยสํวรา. \\ นาญฺญตฺร สพฺพนิสฺสคฺคา, โสตฺถิํ ปสฺสามิ ปาณินนฺติ.}
\csromangathagroup{\csromangathastanza{“นิจฺจํ อุตฺรสฺตมิทํ จิตฺตํ, นิจฺจํ อุพฺพิคฺคมิทํ\footnote{อุพฺพิคฺคิทํ (มหาสติปฏฺฐานสุตฺตวณฺณนายํ) 3. กิจฺเจสุ (พหูสุ)} มโน. \\ อนุปฺปนฺเนสุ กิจฺเฉสุ3, อโถ อุปฺปติเตสุ จ.}\csromangathastanza{สเจ อตฺถิ อนุตฺรสฺตํ, ตํ เม อกฺขาหิ ปุจฺฉิโต”ติ. \\ นาญฺญตฺร โพชฺฌา ตปสา\footnote{โพชฺฌงฺคตปสา (สี, สฺยา, กํ, อิ) 5. เนว (สี, สฺยา, กํ)}, นาญฺญตฺรินฺทฺริยสํวรา. \\ นาญฺญตฺร สพฺพนิสฺสคฺคา, โสตฺถิํ ปสฺสามิ ปาณินนฺติ.}}

\prose{อิทมโวจ ฯเปฯ ตตฺเถวนฺตรธายีติ.}

\csromanheader{2}{8. กกุธสุตฺต}
\proseitem{99}{เอวํ เม สุตํ\csromandash{}เอกํ สมยํ ภควา สาเกเต วิหรติ อญฺชนวเน มิคทาเย.\csromanspacer{} อถ โข กกุโธ เทวปุตฺโต อภิกฺกนฺตาย รตฺติยา อภิกฺกนฺตวณฺโณ เกวลกปฺปํ อญฺชนวนํ โอภาเสตฺวา เยน ภควา เตนุปสงฺกมิ, อุปสงฺกมิตฺวา ภควนฺตํ อภิวาเทตฺวา เอกมนฺตํ อฏฺฐาสิ, เอกมนฺตํ ฐิโต โข กกุโธ เทวปุตฺโต ภควนฺตํ เอตทโวจ “นนฺทสิ สมณา”ติ.\csromanspacer{} กิํ ลทฺธา อาวุโสติ.\csromanspacer{} เตน หิ สมณ โสจสีติ.\csromanspacer{} กิํ ชียิตฺถ อาวุโสติ.\csromanspacer{} เตน หิ สมณ เนว นนฺทสิ น จ5 โสจสีติ.\csromanspacer{} เอวมาวุโสติ.}

\csromanheader{2}{2. เทวปุตฺตสํยุตฺต}
\csromangathameasure{กจฺจิ ตฺวํ อนโฆ ภิกฺขุ, กจฺจิ นนฺที น วิชฺชติ. \\ กจฺจิ ตํ เอกมาสีนํ, อรตี นาภิกีรตีติ. \\ อนโฆ เว อหํ ยกฺข, อโถ นนฺที น วิชฺชติ. \\ อโถ มํ เอกมาสีนํ, อรตี นาภิกีรตีติ. \\ กถํ ตฺวํ อนโฆ ภิกฺขุ, กถํ นนฺที น วิชฺชติ. \\ กถํ ตํ เอกมาสีนํ, อรตี นาภิกีรตีติ. \\ อฆชาตสฺส เว นนฺที, นนฺทีชาตสฺส เว อฆํ. \\ อนนฺที อนโฆ ภิกฺขุ, เอวํ ชานาหิ อาวุโสติ. \\ จิรสฺสํ วต ปสฺสามิ, พฺราหฺมณํ ปรินิพฺพุตํ. \\ อนนฺทิํ อนฆํ ภิกฺขุํ, ติณฺณํ โลเก วิสตฺติกนฺติ.}
\csromangathagroup{\csromangathastanza{\csromanfolio{53}กจฺจิ ตฺวํ อนโฆ\footnote{อนิโฆ (สพฺพตฺถ)} ภิกฺขุ, กจฺจิ นนฺที\footnote{นนฺทิ (สี, สฺยา, กํ)} น วิชฺชติ. \\ กจฺจิ ตํ เอกมาสีนํ, อรตี นาภิกีรตีติ.}\csromangathastanza{อนโฆ เว อหํ ยกฺข, อโถ นนฺที น วิชฺชติ. \\ อโถ มํ เอกมาสีนํ, อรตี นาภิกีรตีติ.}\csromangathastanza{กถํ ตฺวํ อนโฆ ภิกฺขุ, กถํ นนฺที น วิชฺชติ. \\ กถํ ตํ เอกมาสีนํ, อรตี นาภิกีรตีติ.}\csromangathastanza{อฆชาตสฺส เว นนฺที, นนฺทีชาตสฺส เว อฆํ. \\ อนนฺที อนโฆ ภิกฺขุ, เอวํ ชานาหิ อาวุโสติ.}\csromangathastanza{จิรสฺสํ วต ปสฺสามิ, พฺราหฺมณํ ปรินิพฺพุตํ. \\ อนนฺทิํ อนฆํ ภิกฺขุํ, ติณฺณํ โลเก วิสตฺติกนฺติ.}}

\csromanheader{2}{9. อุตฺตรสุตฺต}
\proseitem{100}{ราชคหนิทานํ.\csromanspacer{} เอกมนฺตํ ฐิโต โข อุตฺตโร เทวปุตฺโต ภควโต สนฺติเก อิมํ คาถํ อภาสิ\csromandash{}}

\csromangathameasure{“อุปนียติ ชีวิตมปฺปมายุ, \\ ชรูปนีตสฺส น สนฺติ ตาณา. \\ ปุญฺญานิ กยิราถ สุขาวหานี”ติ. \\ อุปนียติ ชีวิตมปฺปมายุ, \\ ชรูปนีตสฺส น สนฺติ ตาณา. \\ โลกามิสํ ปชเห สนฺติเปกฺโขติ.}
\csromangathagroup{\csromangathastanza{“อุปนียติ ชีวิตมปฺปมายุ, \\ ชรูปนีตสฺส น สนฺติ ตาณา. \\ ปุญฺญานิ กยิราถ สุขาวหานี”ติ. \\ อุปนียติ ชีวิตมปฺปมายุ, \\ ชรูปนีตสฺส น สนฺติ ตาณา. \\ โลกามิสํ ปชเห สนฺติเปกฺโขติ.}}

\csromanheader{2}{10. อนาถปิณฺฑิกสุตฺต}
\proseitem{101}{เอกมนฺตํ ฐิโต โข อนาถปิณฺฑิโก เทวปุตฺโต ภควโต สนฺติเก อิมา คาถาโย อภาสิ\csromandash{}}

\csromangathameasure{“อิทํ หิ ตํ เชตวนํ, อิสิสํฆนิเสวิตํ. \\ อาวุตฺถํ ธมฺมราเชน, ปีติสญฺชนนํ มม. \\ กมฺมํ วิชฺชา จ ธมฺโม จ, สีลํ ชีวิตมุตฺตมํ. \\ เอเตน มจฺจา สุชฺฌนฺติ, น โคตฺเตน ธเนน วา. \\ ตสฺมา หิ ปณฺฑิโต โปโส, สมฺปสฺสํ อตฺถมตฺตโน. \\ โยนิโส วิจิเน ธมฺมํ, เอวํ ตตฺถ วิสุชฺฌติ. \\ สาริปุตฺโตว ปญฺญาย, สีเลน อุปสเมน จ. \\ โยปิ ปารงฺคโต ภิกฺขุ, เอตาวปรโม สิยา”ติ.}
\csromangathagroup{\csromangathastanza{“อิทํ หิ ตํ เชตวนํ, อิสิสํฆนิเสวิตํ. \\ อาวุตฺถํ ธมฺมราเชน, ปีติสญฺชนนํ มม.}\csromangathastanza{\csromanfolio{54}กมฺมํ วิชฺชา จ ธมฺโม จ, สีลํ ชีวิตมุตฺตมํ. \\ เอเตน มจฺจา สุชฺฌนฺติ, น โคตฺเตน ธเนน วา.}\csromangathastanza{ตสฺมา หิ ปณฺฑิโต โปโส, สมฺปสฺสํ อตฺถมตฺตโน. \\ โยนิโส วิจิเน ธมฺมํ, เอวํ ตตฺถ วิสุชฺฌติ.}\csromangathastanza{สาริปุตฺโตว ปญฺญาย, สีเลน อุปสเมน จ. \\ โยปิ ปารงฺคโต ภิกฺขุ, เอตาวปรโม สิยา”ติ.}}

\prose{อิทมโวจ อนาถปิณฺฑิโก เทวปุตฺโต อิทํ วตฺวา ภควนฺตํ อภิวาเทตฺวา ปทกฺขิณํ กตฺวา ตตฺเถวนฺตรธายิ.}

\prose{อถ โข ภควา ตสฺสา รตฺติยา อจฺจเยน ภิกฺขู อามนฺเตสิ “อิมํ ภิกฺขเว รตฺติํ อญฺญตโร เทวปุตฺโต อภิกฺกนฺตาย รตฺติยา อภิกฺกนฺตวณฺโณ เกวลกปฺปํ เชตวนํ โอภาเสตฺวา เยนาหํ เตนุปสงฺกมิ, อุปสงฺกมิตฺวา มํ อภิวาเทตฺวา เอกมนฺตํ อฏฺฐาสิ, เอกมนฺตํ ฐิโต โข ภิกฺขเว โส เทวปุตฺโต มม สนฺติเก อิมา คาถาโย อภาสิ\csromandash{}}

\csromangathameasure{‘อิทํ หิ ตํ เชตวนํ, อิสิสํฆนิเสวิตํ. \\ อาวุตฺถํ ธมฺมราเชน, ปีติสญฺชนนํ มม. \\ กมฺมํ วิชฺชา จ ธมฺโม จ, สีลํ ชีวิตมุตฺตมํ. \\ เอเตน มจฺจา สุชฺฌนฺติ, น โคตฺเตน ธเนน วา. \\ ตสฺมา หิ ปณฺฑิโต โปโส, สมฺปสฺสํ อตฺถมตฺตโน. \\ โยนิโส วิจิเน ธมฺมํ, เอวํ ตตฺถ วิสุชฺฌติ. \\ สาริปุตฺโตว ปญฺญาย, สีเลน อุปสเมน จ. \\ โยปิ ปารงฺคโต ภิกฺขุ, เอตาวปรโม สิยา’ติ.}
\csromangathagroup{\csromangathastanza{‘อิทํ หิ ตํ เชตวนํ, อิสิสํฆนิเสวิตํ. \\ อาวุตฺถํ ธมฺมราเชน, ปีติสญฺชนนํ มม.}\csromangathastanza{กมฺมํ วิชฺชา จ ธมฺโม จ, สีลํ ชีวิตมุตฺตมํ. \\ เอเตน มจฺจา สุชฺฌนฺติ, น โคตฺเตน ธเนน วา.}\csromangathastanza{ตสฺมา หิ ปณฺฑิโต โปโส, สมฺปสฺสํ อตฺถมตฺตโน. \\ โยนิโส วิจิเน ธมฺมํ, เอวํ ตตฺถ วิสุชฺฌติ.}\csromangathastanza{สาริปุตฺโตว ปญฺญาย, สีเลน อุปสเมน จ. \\ โยปิ ปารงฺคโต ภิกฺขุ, เอตาวปรโม สิยา’ติ.}}

\prose{อิทมโวจ ภิกฺขเว โส เทวปุตฺโต อิทํ วตฺวา มํ อภิวาเทตฺวา ปทกฺขิณํ กตฺวา ตตฺเถวนฺตรธายี”ติ.}

\prose{เอวํ วุตฺเต อายสฺมา อานนฺโท ภควนฺตํ เอตทโวจ “โส หิ นูน ภนฺเต อนาถปิณฺฑิโก เทวปุตฺโต ภวิสฺสติ, อนาถปิณฺฑิโก คหปติ อายสฺมนฺเต สาริปุตฺเต อภิปฺปสนฺโน อโหสิ” ติ.\csromanspacer{} สาธุ}

\csromanheader{2}{2. เทวปุตฺตสํยุตฺต}
\prose{\csromanfolio{55}สาธุ อานนฺท, ยาวตกํ โข อานนฺท ตกฺกาย ปตฺตพฺพํ, อนุปฺปตฺตํ ตํ ตยา, อนาถปิณฺฑิโก หิ โส อานนฺท เทวปุตฺโตติ.}

\vspace{\tipitakaparskip}
\csromancenter{อนาถปิณฺฑิกวคฺโค ทุติโย.}
\titlehead{\textbf{ตสฺสุทฺทานํ}}
\csromangathameasure{จนฺทิมโส จ เวณฺฑุ จ, ทีฆลฏฺฐิ จ นนฺทโน. \\ จนฺทโน วาสุทตฺโต จ, สุพฺรหฺมา กกุเธน จ. \\ อุตฺตโร นวโม วุตฺโต, ทสโม อนาถปิณฺฑิโกติ.}
\csromangathagroup{\csromangathastanza{จนฺทิมโส\footnote{จนฺทิมาโส (อิ, ก)} จ เวณฺฑุ\footnote{เวณฺหุ (สี, ก)} จ, ทีฆลฏฺฐิ จ นนฺทโน. \\ จนฺทโน วาสุทตฺโต จ, สุพฺรหฺมา กกุเธน จ. \\ อุตฺตโร นวโม วุตฺโต, ทสโม อนาถปิณฺฑิโกติ.}}

\chapterhead{3. นานาติตฺถิยวคฺค}
\csromanheader{2}{1. สิวสุตฺต}
\proseitem{102}{เอวํ เม สุตํ\csromandash{}เอกํ สมยํ ภควา สาวตฺถิยํ วิหรติ เชตวเน อนาถปิณฺฑิกสฺส อาราเม.\csromanspacer{} อถ โข สิโว เทวปุตฺโต อภิกฺกนฺตาย รตฺติยา อภิกฺกนฺตวณฺโณ เกวลกปฺปํ เชตวนํ โอภาเสตฺวา เยน ภควา เตนุปสงฺกมิ, อุปสงฺกมิตฺวา ภควนฺตํ อภิวาเทตฺวา เอกมนฺตํ อฏฺฐาสิ, เอกมนฺตํ ฐิโต โข สิโว เทวปุตฺโต ภควโต สนฺติเก อิมา คาถาโย อภาสิ\csromandash{}}

\csromangathameasure{“สพฺภิเรว สมาเสถ, สพฺภิ กุพฺเพถ สนฺถวํ. \\ สตํ สทฺธมฺมมญฺญาย, เสยฺโย โหติ น ปาปิโย. \\ สพฺภิเรว สมาเสถ, สพฺภิ กุพฺเพถ สนฺถวํ. \\ สตํ สทฺธมฺมมญฺญาย, ปญฺญา ลพฺภติ นาญฺญโต. \\ สพฺภิเรว สมาเสถ, สพฺภิ กุพฺเพถ สนฺถวํ. \\ สตํ สทฺธมฺมมญฺญาย, โสกมชฺเฌ น โสจติ. \\ สพฺภิเรว สมาเสถ, สพฺภิ กุพฺเพถ สนฺถวํ. \\ สตํ สทฺธมฺมมญฺญาย, ญาติมชฺเฌ วิโรจติ. \\ สพฺภิเรว สมาเสถ, สพฺภิ กุพฺเพถ สนฺถวํ. \\ สตํ สทฺธมฺมมญฺญาย, สตฺตา คจฺฉนฺติ สุคฺคติํ. \\ สพฺภิเรว สมาเสถ, สพฺภิ กุพฺเพถ สนฺถวํ. \\ สตํ สทฺธมฺมมญฺญาย, สตฺตา ติฏฺฐนฺติ สาตตนฺ”ติ.}
\csromangathagroup{\csromangathastanza{“สพฺภิเรว สมาเสถ, สพฺภิ กุพฺเพถ สนฺถวํ. \\ สตํ สทฺธมฺมมญฺญาย, เสยฺโย โหติ น ปาปิโย.}\csromangathastanza{สพฺภิเรว สมาเสถ, สพฺภิ กุพฺเพถ สนฺถวํ. \\ สตํ สทฺธมฺมมญฺญาย, ปญฺญา ลพฺภติ นาญฺญโต.}\csromangathastanza{สพฺภิเรว สมาเสถ, สพฺภิ กุพฺเพถ สนฺถวํ. \\ สตํ สทฺธมฺมมญฺญาย, โสกมชฺเฌ น โสจติ.}\csromangathastanza{สพฺภิเรว สมาเสถ, สพฺภิ กุพฺเพถ สนฺถวํ. \\ สตํ สทฺธมฺมมญฺญาย, ญาติมชฺเฌ วิโรจติ.}\csromangathastanza{\csromanfolio{56}สพฺภิเรว สมาเสถ, สพฺภิ กุพฺเพถ สนฺถวํ. \\ สตํ สทฺธมฺมมญฺญาย, สตฺตา คจฺฉนฺติ สุคฺคติํ.}\csromangathastanza{สพฺภิเรว สมาเสถ, สพฺภิ กุพฺเพถ สนฺถวํ. \\ สตํ สทฺธมฺมมญฺญาย, สตฺตา ติฏฺฐนฺติ สาตตนฺ”ติ.}}

\csromanheader{2}{อถ โข ภควา สิวํ เทวปุตฺตํ คาถาย ปจฺจภาสิ\csromandash{}}
\csromangathameasure{“สพฺภิเรว สมาเสถ, สพฺภิ กุพฺเพถ สนฺถวํ. \\ สตํ สทฺธมฺมมญฺญาย, สพฺพทุกฺขา ปมุจฺจตี”ติ.}
\csromangathagroup{\csromangathastanza{“สพฺภิเรว สมาเสถ, สพฺภิ กุพฺเพถ สนฺถวํ. \\ สตํ สทฺธมฺมมญฺญาย, สพฺพทุกฺขา ปมุจฺจตี”ติ.}}

\csromanheader{2}{2. เขมสุตฺต}
\proseitem{103}{เอกมนฺตํ ฐิโต โข เขโม เทวปุตฺโต ภควโต สนฺติเก อิมา คาถาโย อภาสิ\csromandash{}}

\csromangathameasure{“จรนฺติ พาลา ทุมฺเมธา, อมิตฺเตเนว อตฺตนา. \\ กโรนฺตา ปาปกํ กมฺมํ, ยํ โหติ กฏุกปฺผลํ. \\ น ตํ กมฺมํ กตํ สาธุ, ยํ กตฺวา อนุตปฺปติ. \\ ยสฺส อสฺสุมุโข โรทํ, วิปากํ ปฏิเสวติ. \\ ตญฺจ กมฺมํ กตํ สาธุ, ยํ กตฺวา นานุตปฺปติ. \\ ยสฺส ปตีโต สุมโน, วิปากํ ปฏิเสวติ. \\ ปฏิกจฺเจว ตํ กยิรา, ยํ ชญฺญา หิตมตฺตโน. \\ น สากฏิกจินฺตาย, มนฺตา ธีโร ปรกฺกเม. \\ ยถา สากฏิโก มฏฺฐํ, สมํ หิตฺวา มหาปถํ. \\ วิสมํ มคฺคมารุยฺห, อกฺขจฺฉินฺโนว ฌายติ. \\ เอวํ ธมฺมา อปฺปกฺกมฺม, อธมฺมมนุวตฺติย. \\ มนฺโท มจฺจุมุขํ ปตฺโต, อกฺขจฺฉินฺโนว ฌายตี”ติ.}
\csromangathagroup{\csromangathastanza{“จรนฺติ พาลา ทุมฺเมธา, อมิตฺเตเนว อตฺตนา. \\ กโรนฺตา ปาปกํ กมฺมํ, ยํ โหติ กฏุกปฺผลํ.}\csromangathastanza{น ตํ กมฺมํ กตํ สาธุ, ยํ กตฺวา อนุตปฺปติ. \\ ยสฺส อสฺสุมุโข โรทํ, วิปากํ ปฏิเสวติ.}\csromangathastanza{ตญฺจ กมฺมํ กตํ สาธุ, ยํ กตฺวา นานุตปฺปติ. \\ ยสฺส ปตีโต สุมโน, วิปากํ ปฏิเสวติ.}\csromangathastanza{ปฏิกจฺเจว\footnote{ปฏิคจฺเจว (สี)} ตํ กยิรา, ยํ ชญฺญา หิตมตฺตโน. \\ น สากฏิกจินฺตาย, มนฺตา ธีโร ปรกฺกเม.}\csromangathastanza{ยถา สากฏิโก มฏฺฐํ\footnote{ปนฺถํ (สี), ปสตฺถํ (สฺยา, กํ)}, สมํ หิตฺวา มหาปถํ. \\ วิสมํ มคฺคมารุยฺห, อกฺขจฺฉินฺโนว ฌายติ.}\csromangathastanza{เอวํ ธมฺมา อปฺปกฺกมฺม, อธมฺมมนุวตฺติย. \\ มนฺโท มจฺจุมุขํ ปตฺโต, อกฺขจฺฉินฺโนว ฌายตี”ติ.}}

\csromanheader{2}{3. เสรีสุตฺต}
\vspace{\tipitakaparskip}
\csromancenter{เอกมนฺตํ ฐิโต โข เสรี เทวปุตฺโต ภควนฺตํ คาถาย อชฺฌภาสิ\csromandash{}}
\csromanheader{2}{2. เทวปุตฺตสํยุตฺต}
\csromangathameasure{“อนฺนเมวาภินนฺทนฺติ, อุภเย เทวมานุสา. \\ อถ โก นาม โส ยกฺโข, ยํ อนฺนํ นาภินนฺทตี”ติ. \\ เย นํ ททนฺติ สทฺธาย, วิปฺปสนฺเนน เจตสา. \\ ตเมว อนฺนํ ภชติ, อสฺมิํ โลเก ปรมฺหิ จ. \\ ตสฺมา วิเนยฺย มจฺเฉรํ, ทชฺชา ทานํ มลาภิภู. \\ ปุญฺญานิ ปรโลกสฺมิํ, ปติฏฺฐา โหนฺติ ปาณินนฺติ.}
\csromangathagroup{\csromangathastanza{\csromanfolio{57}“อนฺนเมวาภินนฺทนฺติ, อุภเย เทวมานุสา. \\ อถ โก นาม โส ยกฺโข, ยํ อนฺนํ นาภินนฺทตี”ติ.}\csromangathastanza{เย นํ ททนฺติ สทฺธาย, วิปฺปสนฺเนน เจตสา. \\ ตเมว อนฺนํ ภชติ, อสฺมิํ โลเก ปรมฺหิ จ.}\csromangathastanza{ตสฺมา วิเนยฺย มจฺเฉรํ, ทชฺชา ทานํ มลาภิภู. \\ ปุญฺญานิ ปรโลกสฺมิํ, ปติฏฺฐา โหนฺติ ปาณินนฺติ.}}

\prose{อจฺฉริยํ ภนฺเต, อพฺภุตํ ภนฺเต, ยาวสุภาสิตมิทํ ภนฺเต ภควตา\csromandash{}}

\csromangathameasure{“เย นํ ททนฺติ สทฺธาย, วิปฺปสนฺเนน เจตสา. \\ ตเมว อนฺนํ ภชติ, อสฺมิํ โลเก ปรมฺหิ จ. \\ ตสฺมา วิเนยฺย มจฺเฉรํ, ทชฺชา ทานํ มลาภิภู. \\ ปุญฺญานิ ปรโลกสฺมิํ, ปติฏฺฐา โหนฺติ ปาณินนฺ”ติ.}
\csromangathagroup{\csromangathastanza{“เย นํ ททนฺติ สทฺธาย, วิปฺปสนฺเนน เจตสา. \\ ตเมว อนฺนํ ภชติ, อสฺมิํ โลเก ปรมฺหิ จ.}\csromangathastanza{ตสฺมา วิเนยฺย มจฺเฉรํ, ทชฺชา ทานํ มลาภิภู. \\ ปุญฺญานิ ปรโลกสฺมิํ, ปติฏฺฐา โหนฺติ ปาณินนฺ”ติ.}}

\prose{ภูตปุพฺพาหํ ภนฺเต สิรี\footnote{เสรี (สี, สฺยา, กํ, อิ)} นาม ราชา อโหสิํ ทายโก ทานปติ ทานสฺส วณฺณวาที, ตสฺส มยฺหํ ภนฺเต จตูสุ ทฺวาเรสุ ทานํ ทียิตฺถ สมณ พฺราหฺมณ กปณทฺธิก วนิพฺพก ยาจกานํ.\csromanspacer{} อถ โข มํ ภนฺเต อิตฺถาคารํ อุปสงฺกมิตฺวา เอตทโวจ\footnote{อิตฺถาคารา อุปสงฺกมิตฺวา เอตทโวจุํ (ก)} “เทวสฺส โข\footnote{เทวสฺเสว โข (ก-สี)} ทานํ ทียติ, อมฺหากํ ทานํ น ทียติ.\csromanspacer{} สาธุ มยมฺปิ เทวํ นิสฺสาย ทานานิ ทเทยฺยาม, ปุญฺญานิ กเรยฺยามา”ติ.\csromanspacer{} ตสฺส มยฺหํ ภนฺเต เอตทโหสิ “อหํ โขสฺมิ ทายโก ทานปติ ทานสฺส วณฺณวาที, ‘ทานํ ทสฺสามา’ติ วทนฺเต กินฺติ วเทยฺยนฺ”ติ.\csromanspacer{} โส ขฺวาหํ ภนฺเต ปฐมํ ทฺวารํ อิตฺถาคารสฺส อทาสิํ.\csromanspacer{} ตตฺถ อิตฺถาคารสฺส ทานํ ทียิตฺถ, มม ทานํ ปฏิกฺกมิ.}

\prose{อถ โข มํ ภนฺเต ขตฺติยา อนุยนฺตา อุปสงฺกมิตฺวา เอตทโวจุํ “เทวสฺส โข ทานํ ทียติ, อิตฺถาคารสฺส ทานํ ทียติ, อมฺหากํ ทานํ น ทียติ.\csromanspacer{} สาธุ มยมฺปิ เทวํ นิสฺสาย ทานานิ ทเทยฺยาม, ปุญฺญานิ กเรยฺยามา”ติ.\csromanspacer{} ตสฺส มยฺหํ ภนฺเต เอตทโหสิ “อหํ โขสฺมิ ทายโก ทานปติ ทานสฺส วณฺณวาที, ‘ทานํ ทสฺสามา’ติ วทนฺเต กินฺติ}

\prose{\csromanfolio{58}วเทยฺยนฺ”ติ.\csromanspacer{} โส ขฺวาหํ ภนฺเต ทุติยํ ทฺวารํ ขตฺติยานํ อนุยนฺตานํ อทาสิํ.\csromanspacer{} ตตฺถ ขตฺติยานํ อนุยนฺตานํ ทานํ ทียิตฺถ, มม ทานํ ปฏิกฺกมิ.}

\prose{อถ โข มํ ภนฺเต พลกาโย อุปสงฺกมิตฺวา เอตทโวจ “เทวสฺส โข ทานํ ทียติ, อิตฺถาคารสฺส ทานํ ทียติ, ขตฺติยานํ อนุยนฺตานํ ทานํ ทียติ, อมฺหากํ ทานํ น ทียติ.\csromanspacer{} สาธุ มยมฺปิ เทวํ นิสฺสาย ทานานิ ทเทยฺยาม, ปุญฺญานิ กเรยฺยามา”ติ.\csromanspacer{} ตสฺส มยฺหํ ภนฺเต เอตทโหสิ “อหํ โขสฺมิ ทายโก ทานปติ ทานสฺส วณฺณวาที, ‘ทานํ ทสฺสามา’ติ วทนฺเต กินฺติ วเทยฺยนฺ”ติ.\csromanspacer{} โส ขฺวาหํ ภนฺเต ตติยํ ทฺวารํ พลกายสฺส อทาสิํ.\csromanspacer{} ตตฺถ พลกายสฺส ทานํ ทียิตฺถ, มม ทานํ ปฏิกฺกมิ.}

\prose{อถ โข มํ ภนฺเต พฺราหฺมณคหปติกา อุปสงฺกมิตฺวา เอตทโวจุํ “เทวสฺส โข ทานํ ทียติ, อิตฺถาคารสฺส ทานํ ทียติ, ขตฺติยานํ อนุยนฺตานํ ทานํ ทียติ, พลกายสฺส ทานํ ทียติ, อมฺหากํ ทานํ น ทียติ.\csromanspacer{} สาธุ มยมฺปิ เทวํ นิสฺสาย ทานานิ ทเทยฺยาม, ปุญฺญานิ กเรยฺยามา”ติ.\csromanspacer{} ตสฺส มยฺหํ ภนฺเต เอตทโหสิ “อหํ โขสฺมิ ทายโก ทานปติ ทานสฺส วณฺณวาที, ‘ทานํ ทสฺสามา’ติ วทนฺเต กินฺติ วเทยฺยนฺ”ติ.\csromanspacer{} โส ขฺวาหํ ภนฺเต จตุตฺถํ ทฺวารํ พฺราหฺมณคหปติกานํ อทาสิํ.\csromanspacer{} ตตฺถ พฺราหฺมณคหปติกานํ ทานํ ทียิตฺถ, มม ทานํ ปฏิกฺกมิ.}

\prose{อถ โข มํ ภนฺเต ปุริสา อุปสงฺกมิตฺวา เอตทโวจุํ “น โข ทานิ เทวสฺส โกจิ ทานํ ทียตี”ติ.\csromanspacer{} เอวํ วุตฺตาหํ ภนฺเต เต ปุริเส เอตทโวจํ “เตน หิ ภเณ โย พาหิเรสุ ชนปเทสุ อาโย สญฺชายติ, ตโต อุปฑฺฒํ อนฺเตปุเร ปเวเสถ, อุปฑฺฒํ ตตฺเถว ทานํ เทถ สมณพฺราหฺมณกปณทฺธิกวนิพฺพกยาจกานนฺ”ติ.\csromanspacer{} โส ขฺวาหํ ภนฺเต เอวํ ทีฆรตฺตํ กตานํ ปุญฺญานํ เอวํ ทีฆรตฺตํ กตานํ กุสลานํ ธมฺมานํ ปริยนฺตํ นาธิคจฺฉามิ “เอตฺตกํ ปุญฺญนฺ”ติ วา “เอตฺตโก ปุญฺญวิปาโก”ติ วา “เอตฺตกํ สคฺเค ฐาตพฺพนฺ”ติ วาติ.\csromanspacer{} อจฺฉริยํ ภนฺเต, อพฺภุตํ ภนฺเต, ยาวสุภาสิตมิทํ ภนฺเต ภควตา\csromandash{}}

\csromangathameasure{“เย นํ ททนฺติ สทฺธาย, วิปฺปสนฺเนน เจตสา. \\ ตเมว อนฺนํ ภชติ, อสฺมิํ โลเก ปรมฺหิ จ. \\ ตสฺมา วิเนยฺย มจฺเฉรํ, ทชฺชา ทานํ มลาภิภู. \\ ปุญฺญานิ ปรโลกสฺมิํ, ปติฏฺฐา โหนฺติ ปาณินนฺ”ติ.}
\csromangathagroup{\csromangathastanza{“เย นํ ททนฺติ สทฺธาย, วิปฺปสนฺเนน เจตสา. \\ ตเมว อนฺนํ ภชติ, อสฺมิํ โลเก ปรมฺหิ จ.}\csromangathastanza{ตสฺมา วิเนยฺย มจฺเฉรํ, ทชฺชา ทานํ มลาภิภู. \\ ปุญฺญานิ ปรโลกสฺมิํ, ปติฏฺฐา โหนฺติ ปาณินนฺ”ติ.}}

\csromanheader{2}{2. เทวปุตฺตสํยุตฺต}
\csromanheader{2}{4. ฆฏีการสุตฺต}
\csromanheader{2}{105. เอกมนฺตํ ฐิโต โข ฆฏีกาโร เทวปุตฺโต ภควโต สนฺติเก}
\csromangathameasure{“อวิหํ อุปปนฺนาเส, วิมุตฺตา สตฺต ภิกฺขโว. \\ ราคโทสปริกฺขีณา, ติณฺณา โลเก วิสตฺติกนฺ”ติ. \\ เก จ เต อตรุํ ปงฺกํ, มจฺจุเธยฺยํ สุทุตฺตรํ. \\ เก หิตฺวา มานุสํ เทหํ, ทิพฺพโยคํ อุปจฺจคุนฺติ. \\ อุปโก ปลคณฺโฑ จ, ปุกฺกุสาติ จ เต ตโย. \\ ภทฺทิโย ขณฺฑเทโว จ, พาหุรคฺคิ จ สงฺคิโย. \\ เต หิตฺวา มานุสํ เทหํ, ทิพฺพโยคํ อุปจฺจคุนฺติ. \\ กุสลี ภาสสี เตสํ, มารปาสปฺปหายินํ. \\ กสฺส เต ธมฺมมญฺญาย, อจฺฉิทุํ ภวพนฺธนนฺติ. \\ น อญฺญตฺร ภควตา, นาญฺญตฺร ตว สาสนา. \\ ยสฺส เต ธมฺมมญฺญาย, อจฺฉิทุํ ภวพนฺธนํ. \\ ยตฺถ นามญฺจ รูปญฺจ, อเสสํ อุปรุชฺฌติ. \\ ตํ เต ธมฺมํ อิธญฺญาย, อจฺฉิทุํ ภวพนฺธนนฺติ. \\ คมฺภีรํ ภาสสี วาจํ, ทุพฺพิชานํ สุทุพฺพุธํ. \\ กสฺส ตฺวํ ธมฺมมญฺญาย, วาจํ ภาสสิ อีทิสนฺติ. \\ กุมฺภกาโร ปุเร อาสิํ, เวกฬิงฺเค ฆฏีกโร. \\ มาตาเปตฺติภโร อาสิํ, กสฺสปสฺส อุปาสโก. \\ วิรโต เมถุนา ธมฺมา, พฺรหฺมจารี นิรามิโส. \\ อหุวา เต สคาเมยฺโย, อหุวา เต ปุเร สขา. \\ โสหเมเต ปชานามิ, วิมุตฺเต สตฺต ภิกฺขโว. \\ ราคโทสปริกฺขีเณ, ติณฺเณ โลเก วิสตฺติกนฺติ. \\ เอวเมตํ ตทา อาสิ, ยถา ภาสสิ ภคฺคว. \\ กุมฺภกาโร ปุเร อาสิ, เวกฬิงฺเค ฆฏีกโร. \\ มาตาเปตฺติภโร อาสิ, กสฺสปสฺส อุปาสโก. \\ วิรโต เมถุนา ธมฺมา, พฺรหฺมจารี นิรามิโส. \\ อหุวา เม สคาเมยฺโย, อหุวา เม ปุเร สขาติ. \\ เอวเมตํ ปุราณานํ, สหายานํ อหุ สงฺคโม. \\ อุภินฺนํ ภาวิตตฺตานํ, สรีรนฺติมธารินนฺติ.}
\csromangathagroup{\csromangathastanza{\csromanfolio{59}“อวิหํ อุปปนฺนาเส, วิมุตฺตา สตฺต ภิกฺขโว. \\ ราคโทสปริกฺขีณา, ติณฺณา โลเก วิสตฺติกนฺ”ติ.}\csromangathastanza{เก จ เต อตรุํ ปงฺกํ, มจฺจุเธยฺยํ สุทุตฺตรํ. \\ เก หิตฺวา มานุสํ เทหํ, ทิพฺพโยคํ อุปจฺจคุนฺติ.}\csromangathastanza{อุปโก ปลคณฺโฑ\footnote{ผลคณฺโฑ (ก)} จ, ปุกฺกุสาติ จ เต ตโย. \\ ภทฺทิโย ขณฺฑเทโว จ, พาหุรคฺคิ จ สงฺคิโย\footnote{ผลคณฺโฑ (ก)}.}\csromangathastanza{เต หิตฺวา มานุสํ เทหํ, ทิพฺพโยคํ อุปจฺจคุนฺติ. \\ กุสลี ภาสสี เตสํ, มารปาสปฺปหายินํ.}\csromangathastanza{กสฺส เต ธมฺมมญฺญาย, อจฺฉิทุํ ภวพนฺธนนฺติ. \\ น อญฺญตฺร ภควตา, นาญฺญตฺร ตว สาสนา.}\csromangathastanza{ยสฺส เต ธมฺมมญฺญาย, อจฺฉิทุํ ภวพนฺธนํ. \\ ยตฺถ นามญฺจ รูปญฺจ, อเสสํ อุปรุชฺฌติ.}\csromangathastanza{ตํ เต ธมฺมํ อิธญฺญาย, อจฺฉิทุํ ภวพนฺธนนฺติ. \\ คมฺภีรํ ภาสสี วาจํ, ทุพฺพิชานํ สุทุพฺพุธํ.}\csromangathastanza{กสฺส ตฺวํ ธมฺมมญฺญาย, วาจํ ภาสสิ อีทิสนฺติ. \\ กุมฺภกาโร ปุเร อาสิํ, เวกฬิงฺเค ฆฏีกโร.}\csromangathastanza{มาตาเปตฺติภโร อาสิํ, กสฺสปสฺส อุปาสโก. \\ วิรโต เมถุนา ธมฺมา, พฺรหฺมจารี นิรามิโส.}\csromangathastanza{อหุวา เต สคาเมยฺโย, อหุวา เต ปุเร สขา. \\ โสหเมเต ปชานามิ, วิมุตฺเต สตฺต ภิกฺขโว.}\csromangathastanza{ราคโทสปริกฺขีเณ, ติณฺเณ โลเก วิสตฺติกนฺติ. \\ เอวเมตํ ตทา อาสิ, ยถา ภาสสิ ภคฺคว.}\csromangathastanza{กุมฺภกาโร ปุเร อาสิ, เวกฬิงฺเค ฆฏีกโร. \\ มาตาเปตฺติภโร อาสิ, กสฺสปสฺส อุปาสโก.}\csromangathastanza{\csromanfolio{60}วิรโต เมถุนา ธมฺมา, พฺรหฺมจารี นิรามิโส. \\ อหุวา เม สคาเมยฺโย, อหุวา เม ปุเร สขาติ.}\csromangathastanza{เอวเมตํ ปุราณานํ, สหายานํ อหุ สงฺคโม. \\ อุภินฺนํ ภาวิตตฺตานํ, สรีรนฺติมธารินนฺติ.}}

\csromanheader{2}{5. ชนฺตุสุตฺต}
\proseitem{106}{เอวํ เม สุตํ\csromandash{}เอกํ สมยํ สมฺพหุลา ภิกฺขู โกสเลสุ วิหรนฺติ หิมวนฺตปสฺเส อรญฺญกุฏิกาย อุทฺธตา อุนฺนฬา จปลา มุขรา วิกิณฺณวาจา มุฏฺฐสฺสติโน อสมฺปชานา อสมาหิตา วิพฺภนฺตจิตฺตา ปากตินฺทฺริยา.}

\prose{อถ โข ชนฺตุ เทวปุตฺโต ตทหุโปสเถ ปนฺนรเส เยน เต ภิกฺขู เตนุปสงฺกมิ, อุปสงฺกมิตฺวา เต ภิกฺขู คาถาหิ อชฺฌภาสิ\csromandash{}}

\csromangathameasure{“สุขชีวิโน ปุเร อาสุํ, ภิกฺขู โคตมสาวกา. \\ อนิจฺฉา ปิณฺฑเมสนา, อนิจฺฉา สยนาสนํ. \\ โลเก อนิจฺจตํ ญตฺวา, ทุกฺขสฺสนฺตํ อกํสุ เต. \\ ทุปฺโปสํ กตฺวา อตฺตานํ, คาเม คามณิกา วิย. \\ ภุตฺวา ภุตฺวา นิปชฺชนฺติ, ปราคาเรสุ มุจฺฉิตา. \\ สงฺฆสฺส อญฺชลิํ กตฺวา, อิเธกจฺเจ วทามหํ. \\ อปวิทฺธา อนาถา เต, ยถา เปตา ตเถว เต. \\ เย โข ปมตฺตา วิหรนฺติ, เต เม สนฺธาย ภาสิตํ. \\ เย อปฺปมตฺตา วิหรนฺติ, นโม เตสํ กโรมหนฺ”ติ.}
\csromangathagroup{\csromangathastanza{“สุขชีวิโน ปุเร อาสุํ, ภิกฺขู โคตมสาวกา. \\ อนิจฺฉา ปิณฺฑเมสนา\footnote{ปิณฺฑเมสานา (?)}, อนิจฺฉา สยนาสนํ.}\csromangathastanza{โลเก อนิจฺจตํ ญตฺวา, ทุกฺขสฺสนฺตํ อกํสุ เต. \\ ทุปฺโปสํ กตฺวา อตฺตานํ, คาเม คามณิกา วิย.}\csromangathastanza{ภุตฺวา ภุตฺวา นิปชฺชนฺติ, ปราคาเรสุ มุจฺฉิตา. \\ สงฺฆสฺส อญฺชลิํ กตฺวา, อิเธกจฺเจ วทามหํ\footnote{วนฺทามหํ (ก)}.}\csromangathastanza{อปวิทฺธา อนาถา เต, ยถา เปตา ตเถว เต\footnote{ตเถว จ (สี)}. \\ เย โข ปมตฺตา วิหรนฺติ, เต เม สนฺธาย ภาสิตํ. \\ เย อปฺปมตฺตา วิหรนฺติ, นโม เตสํ กโรมหนฺ”ติ.}}

\csromanheader{2}{6. โรหิตสฺสสุตฺต}
\proseitem{107}{สาวตฺถินิทานํ.\csromanspacer{} เอกมนฺตํ ฐิโต โข โรหิตสฺโส เทวปุตฺโต ภควนฺตํ เอตทโวจ “ยตฺถ นุ โข ภนฺเต น ชายติ น ชียติ น มียติ\footnote{น ชิยฺยติ น มิยฺยติ (สฺยา, กํ, ก)} น จวติ น อุปปชฺชติ, สกฺกา นุ โข โส ภนฺเต คมเนน}

\csromanheader{2}{2. เทวปุตฺตสํยุตฺต}
\prose{\csromanfolio{61}โลกสฺส อนฺโต ญาตุํ วา ทฏฺฐุํ วา ปาปุณิตุํ วา”ติ.\csromanspacer{} ยตฺถ โข อาวุโส น ชายติ น ชียติ น มียติ น จวติ น อุปปชฺชติ, นาหํ ตํ คมเนน โลกสฺส อนฺตํ ญาเตยฺยํ ทฏฺเฐยฺยํ ปตฺเตยฺยนฺติ วทามีติ.}

\prose{อจฺฉริยํ ภนฺเต, อพฺภุตํ ภนฺเต, ยาวสุภาสิตมิทํ ภนฺเต ภควตา “ยตฺถ โข อาวุโส น ชายติ น ชียติ น มียติ น จวติ น อุปปชฺชติ, นาหํ ตํ คมเนน โลกสฺส อนฺตํ ญาเตยฺยํ ทฏฺเฐยฺยํ ปตฺเตยฺยนฺติ วทามี”ติ.}

\prose{ภูตปุพฺพาหํ ภนฺเต โรหิตสฺโส นาม อิสิ อโหสิํ โภชปุตฺโต อิทฺธิมา เวหาสงฺคโม.\csromanspacer{} ตสฺส มยฺหํ ภนฺเต เอวรูโป ชโว อโหสิ\csromandash{}เสยฺยถาปิ นาม ทฬฺหธมฺมา\footnote{ทฬฺหธมฺโม (สพฺพตฺถ) ฏีกา จ โมคฺคลฺลานพฺยากรณํ จ โอโลเกตพฺพํ.} ธนุคฺคโห สุสิกฺขิโต กตหตฺโถ กตโยคฺโค กตูปาสโน ลหุเกน อสเนน อปฺปกสิเรเนว ติริยํ ตาลจฺฉายํ อติปาเตยฺย.\csromanspacer{} ตสฺส มยฺหํ ภนฺเต เอวรูโป ปทวีติหาโร อโหสิ\csromandash{}เสยฺยถาปิ นาม ปุรตฺถิมา สมุทฺทา ปจฺฉิโม สมุทฺโท.\csromanspacer{} ตสฺส มยฺหํ ภนฺเต เอวรูปํ อิจฺฉาคตํ อุปฺปชฺชิ “อหํ คมเนน โลกสฺส อนฺตํ ปาปุณิสฺสามี”ติ.\csromanspacer{} โส ขฺวาหํ ภนฺเต เอวรูเปน ชเวน สมนฺนาคโต เอวรูเปน จ ปทวีติหาเรน, อญฺญเตฺรว อสิต ปีตขายิต สายิตา อญฺญตฺร อุจฺจาร ปสฺสาว กมฺมา อญฺญตฺร นิทฺทากิลมถปฏิวิโนทนา วสฺสสตายุโก วสฺสสตชีวี วสฺสสตํ คนฺตฺวา อปฺปตฺวาว โลกสฺส อนฺตํ อนฺตราว กาลงฺกโต.}

\prose{อจฺฉริยํ ภนฺเต, อพฺภุตํ ภนฺเต, ยาวสุภาสิตมิทํ ภนฺเต ภควตา “ยตฺถ โข อาวุโส น ชายติ น ชียติ น มียติ น จวติ น อุปปชฺชติ, นาหํ ตํ คมเนน โลกสฺส อนฺตํ ญาเตยฺยํ ทฏฺเฐยฺยํ ปตฺเตยฺยนฺติ วทามี”ติ.}

\prose{น โข ปนาหํ อาวุโส อปฺปตฺวา โลกสฺส อนฺตํ ทุกฺขสฺส อนฺตกิริยํ วทามิ, อปิ จ ขฺวาหํ อาวุโส อิมสฺมิํเยว พฺยามมตฺเต กเฬวเร สสญฺญิมฺหิ สมนเก โลกญฺจ ปญฺญเปมิ โลกสมุทยญฺจ โลกนิโรธญฺจ โลกนิโรธคามินิญฺจ ปฏิปทนฺติ.}

\csromangathameasure{คมเนน น ปตฺตพฺโพ, โลกสฺสนฺโต กุทาจนํ. \\ น จ อปฺปตฺวา โลกนฺตํ, ทุกฺขา อตฺถิ ปโมจนํ.}
\csromangathagroup{\csromangathastanza{คมเนน น ปตฺตพฺโพ, โลกสฺสนฺโต กุทาจนํ. \\ น จ อปฺปตฺวา โลกนฺตํ, ทุกฺขา อตฺถิ ปโมจนํ.}}

\prose{\csromanfolio{62}ตสฺมา หเว โลกวิทู สุเมโธ.}

\csromangathameasure{โลกนฺตคู วุสิตพฺรหฺมจริโย, \\ โลกสฺส อนฺตํ สมิตาวิ ญตฺวา. \\ นาสีสติ โลกมิมํ ปรญฺจาติ.}
\csromangathagroup{\csromangathastanza{โลกนฺตคู วุสิตพฺรหฺมจริโย, \\ โลกสฺส อนฺตํ สมิตาวิ ญตฺวา. \\ นาสีสติ โลกมิมํ ปรญฺจาติ.}}

\csromanheader{2}{7. นนฺทสุตฺต}
\proseitem{108}{เอกมนฺตํ ฐิโต โข นนฺโท เทวปุตฺโต ภควโต สนฺติเก อิมํ คาถํ อภาสิ\csromandash{}}

\csromangathameasure{“อจฺเจนฺติ กาลา ตรยนฺติ รตฺติโย, \\ วโยคุณา อนุปุพฺพํ ชหนฺติ. \\ ปุญฺญานิ กยิราถ สุขาวหานี”ติ. \\ อจฺเจนฺติ กาลา ตรยนฺติ รตฺติโย, \\ วโยคุณา อนุปุพฺพํ ชหนฺติ. \\ โลกามิสํ ปชเห สนฺติเปกฺโขติ.}
\csromangathagroup{\csromangathastanza{“อจฺเจนฺติ กาลา ตรยนฺติ รตฺติโย, \\ วโยคุณา อนุปุพฺพํ ชหนฺติ. \\ ปุญฺญานิ กยิราถ สุขาวหานี”ติ. \\ อจฺเจนฺติ กาลา ตรยนฺติ รตฺติโย, \\ วโยคุณา อนุปุพฺพํ ชหนฺติ. \\ โลกามิสํ ปชเห สนฺติเปกฺโขติ.}}

\csromanheader{2}{8. นนฺทิวิสาลสุตฺต}
\proseitem{109}{เอกมนฺตํ ฐิโต โข นนฺทิวิสาโล เทวปุตฺโต ภควนฺตํ คาถาย อชฺฌภาสิ\csromandash{}}

\csromangathameasure{“จตุจกฺกํ นวทฺวารํ, ปุณฺณํ โลเภน สํยุตํ. \\ ปงฺกชาตํ มหาวีร, กถํ ยาตฺรา ภวิสฺสตี”ติ. \\ เฉตฺวา นทฺธิํ วรตฺตญฺจ, อิจฺฉาโลภญฺจ ปาปกํ. \\ สมูลํ ตณฺหมพฺพุยฺห, เอวํ ยาตฺรา ภวิสฺสตีติ.}
\csromangathagroup{\csromangathastanza{“จตุจกฺกํ นวทฺวารํ, ปุณฺณํ โลเภน สํยุตํ. \\ ปงฺกชาตํ มหาวีร, กถํ ยาตฺรา ภวิสฺสตี”ติ.}\csromangathastanza{เฉตฺวา นทฺธิํ วรตฺตญฺจ, อิจฺฉาโลภญฺจ ปาปกํ. \\ สมูลํ ตณฺหมพฺพุยฺห, เอวํ ยาตฺรา ภวิสฺสตีติ.}}

\csromanheader{2}{9. สุสิมสุตฺต}
\proseitem{110}{สาวตฺถินิทานํ.\csromanspacer{} อถ โข อายสฺมา อานนฺโท เยน ภควา เตนุปสงฺกมิ, อุปสงฺกมิตฺวา ภควนฺตํ อภิวาเทตฺวา เอกมนฺตํ นิสีทิ, เอกมนฺตํ นิสินฺนํ โข อายสฺมนฺตํ อานนฺทํ ภควา เอตทโวจ “ตุยฺหมฺปิ โน อานนฺท สาริปุตฺโต รุจฺจตี”ติ.}

\csromanheader{2}{2. เทวปุตฺตสํยุตฺต}
\prose{\csromanfolio{63}กสฺส หิ นาม ภนฺเต อพาลสฺส อทุฏฺฐสฺส อมูฬฺหสฺส อวิปลฺลตฺถจิตฺตสฺส อายสฺมา สาริปุตฺโต น รุจฺเจยฺย, ปณฺฑิโต ภนฺเต อายสฺมา สาริปุตฺโต, มหาปญฺโญ ภนฺเต อายสฺมา สาริปุตฺโต, ปุถุปญฺโญ ภนฺเต อายสฺมา สาริปุตฺโต, หาสปญฺโญ\footnote{หาสุปญฺโญ (สี)} ภนฺเต อายสฺมา สาริปุตฺโต, ชวนปญฺโญ ภนฺเต อายสฺมา สาริปุตฺโต, ติกฺขปญฺโญ ภนฺเต อายสฺมา สาริปุตฺโต, นิพฺเพธิกปญฺโญ ภนฺเต อายสฺมา สาริปุตฺโต, อปฺปิจฺโฉ ภนฺเต อายสฺมา สาริปุตฺโต, สนฺตุฏฺโฐ ภนฺเต อายสฺมา สาริปุตฺโต, ปวิวิตฺโต ภนฺเต อายสฺมา สาริปุตฺโต, อสํสฏฺโฐ ภนฺเต อายสฺมา สาริปุตฺโต, อารทฺธวีริโย ภนฺเต อายสฺมา สาริปุตฺโต, วตฺตา ภนฺเต อายสฺมา สาริปุตฺโต, วจนกฺขโม ภนฺเต อายสฺมา สาริปุตฺโต, โจทโก ภนฺเต อายสฺมา สาริปุตฺโต, ปาปครหี ภนฺเต อายสฺมา สาริปุตฺโต.\csromanspacer{} กสฺส หิ นาม ภนฺเต อพาลสฺส อทุฏฺฐสฺส อมูฬฺหสฺส อวิปลฺลตฺถจิตฺตสฺส อายสฺมา สาริปุตฺโต น รุจฺเจยฺยาติ.}

\prose{เอวเมตํ อานนฺท, เอวเมตํ อานนฺท, กสฺส หิ นาม อานนฺท อพาลสฺส อทุฏฺฐสฺส อมูฬฺหสฺส อวิปลฺลตฺถจิตฺตสฺส สาริปุตฺโต น รุจฺเจยฺย, ปณฺฑิโต อานนฺท สาริปุตฺโต, มหาปญฺโญ อานนฺท สาริปุตฺโต, ปุถุปญฺโญ อานนฺท สาริปุตฺโต, หาสปญฺโญ อานนฺท สาริปุตฺโต, ชวนปญฺโญ อานนฺท สาริปุตฺโต, ติกฺขปญฺโญ อานนฺท สาริปุตฺโต, นิพฺเพธิกปญฺโญ อานนฺท สาริปุตฺโต, อปฺปิจฺโฉ อานนฺท สาริปุตฺโต, สนฺตุฏฺโฐ อานนฺท สาริปุตฺโต, ปวิวิตฺโต อานนฺท สาริปุตฺโต, อสํสฏฺโฐ อานนฺท สาริปุตฺโต, อารทฺธวีริโย อานนฺท สาริปุตฺโต, วตฺตา อานนฺท สาริปุตฺโต, วจนกฺขโม อานนฺท สาริปุตฺโต, โจทโก อานนฺท สาริปุตฺโต, ปาปครหี อานนฺท สาริปุตฺโต.\csromanspacer{} กสฺส หิ นาม อานนฺท อพาลสฺส อทุฏฺฐสฺส อมูฬฺหสฺส อวิปลฺลตฺถจิตฺตสฺส สาริปุตฺโต น รุจฺเจยฺยาติ.}

\prose{อถ โข สุสิโม\footnote{สุสีโม (สี)} เทวปุตฺโต อายสฺมโต สาริปุตฺตสฺส วณฺเณ ภญฺญมาเน มหติยา เทวปุตฺตปริสาย ปริวุโต เยน ภควา เตนุปสงฺกมิ, อุปสงฺกมิตฺวา ภควนฺตํ อภิวาเทตฺวา เอกมนฺตํ อฏฺฐาสิ, เอกมนฺตํ ฐิโต โข สุสิโม เทวปุตฺโต ภควนฺตํ เอตทโวจ}

\prose{\csromanfolio{64}“เอวเมตํ ภควา, เอวเมตํ สุคต, กสฺส หิ นาม ภนฺเต อพาลสฺส อทุฏฺฐสฺส อมูฬฺหสฺส อวิปลฺลตฺถจิตฺตสฺส อายสฺมา สาริปุตฺโต น รุจฺเจยฺย, ปณฺฑิโต ภนฺเต อายสฺมา สาริปุตฺโต, มหาปญฺโญ ภนฺเต, ปุถุปญฺโญ ภนฺเต, หาสปญฺโญ ภนฺเต, ชวนปญฺโญ ภนฺเต, ติกฺขปญฺโญ ภนฺเต, นิพฺเพธิกปญฺโญ ภนฺเต, อปฺปิจฺโฉ ภนฺเต, สนฺตุฏฺโฐ ภนฺเต, ปวิวิตฺโต ภนฺเต, อสํสฏฺโฐ ภนฺเต, อารทฺธวีริโย ภนฺเต, วตฺตา ภนฺเต, วจนกฺขโม ภนฺเต, โจทโก ภนฺเต, ปาปครหี ภนฺเต อายสฺมา สาริปุตฺโต.\csromanspacer{} กสฺส หิ นาม ภนฺเต อพาลสฺส อทุฏฺฐสฺส อมูฬฺหสฺส อวิปลฺลตฺถจิตฺตสฺส อายสฺมา สาริปุตฺโต น รุจฺเจยฺย.}

\prose{อหมฺปิ หิ ภนฺเต ยญฺญเทว เทวปุตฺตปริสํ อุปสงฺกมิํ, เอตเทว พหุลํ สทฺทํ สุณามิ “ปณฺฑิโต อายสฺมา สาริปุตฺโต, มหาปญฺโญ อายสฺมา, ปุถุปญฺโญ อายสฺมา, หาสปญฺโญ อายสฺมา, ชวนปญฺโญ อายสฺมา, ติกฺขปญฺโญ อายสฺมา, นิพฺเพธิกปญฺโญ อายสฺมา, อปฺปิจฺโฉ อายสฺมา, สนฺตุฏฺโฐ อายสฺมา, ปวิวิตฺโต อายสฺมา, อสํสฏฺโฐ อายสฺมา, อารทฺธวีริโย อายสฺมา, วตฺตา อายสฺมา, วจนกฺขโม อายสฺมา, โจทโก อายสฺมา, ปาปครหี อายสฺมา สาริปุตฺโต”ติ.\csromanspacer{} กสฺส หิ นาม ภนฺเต อพาลสฺส อทุฏฺฐสฺส อมูฬฺหสฺส อวิปลฺลตฺถจิตฺตสฺส อายสฺมา สาริปุตฺโต น รุจฺเจยฺยาติ.}

\prose{อถ โข สุสิมสฺส เทวปุตฺตสฺส เทวปุตฺตปริสา อายสฺมโต สาริปุตฺตสฺส วณฺเณ ภญฺญมาเน อตฺตมนา ปมุทิตา ปีติโสมนสฺสชาตา อุจฺจาวจา วณฺณนิภา อุปทํเสติ, เสยฺยถาปิ นาม มณิ เวฬุริโย สุโภ ชาติมา อฏฺฐํโส สุปริกมฺมกโต ปณฺฑุกมฺพเล นิกฺขิตฺโต ภาสเต จ ตปเต จ วิโรจติ จ, เอวเมวํ สุสิมสฺส เทวปุตฺตสฺส เทวปุตฺตปริสา อายสฺมโต สาริปุตฺตสฺส วณฺเณ ภญฺญมาเน อตฺตมนา ปมุทิตา ปีติโสมนสฺสชาตา อุจฺจาวจา วณฺณนิภา อุปทํเสติ.}

\prose{เสยฺยาถาปิ นาม นิกฺขํ ชมฺโพนทํ ทกฺข กมฺมารปุตฺต อุกฺกามุข สุกุสลสมฺปหฏฺฐํ ปณฺฑุกมฺพเล นิกฺขิตฺตํ ภาสเต จ ตปเต จ วิโรจติ จ, เอวเมวํ สุสิมสฺส เทวปุตฺตสฺส เทวปุตฺตปริสา อายสฺมโต สาริปุตฺตสฺส วณฺเณ ภญฺญมาเน อตฺตมนา ปมุทิตา ปีติโสมนสฺสชาตา อุจฺจาวจา วณฺณนิภา อุปทํเสติ.}

\csromanheader{2}{2. เทวปุตฺตสํยุตฺต}
\prose{\csromanfolio{65}เสยฺยถาปิ นาม สรทสมเย วิทฺเธ วิคตวลาหเก เทเว รตฺติยา ปจฺจูสสมยํ โอสธิตารกา ภาสเต จ ตปเต จ วิโรจติ จ, เอวเมวํ สุสิมสฺส เทวปุตฺตสฺส เทวปุตฺตปริสา อายสฺมโต สาริปุตฺตสฺส วณฺเณ ภญฺญมาเน อตฺตมนา ปมุทิตา ปีติโสมนสฺสชาตา อุจฺจาวจา วณฺณนิภา อุปทํเสติ.}

\prose{เสยฺยถาปิ นาม สรทสมเย วิทฺเธ วิคตวลาหเก เทเว อาทิจฺโจ นภํ อพฺภุสฺสกฺกมาโน\footnote{อพฺภุสฺสุกฺกมาโน (สี, สฺยา, กํ, อิ), อพฺภุคฺคมมาโน (ที 2. 150 ปิฏฺเฐ)} สพฺพํ อากาสคตํ ตมคตํ อภิวิหจฺจ ภาสเต จ ตปเต จ วิโรจติ จ, เอวเมวํ สุสิมสฺส เทวปุตฺตสฺส เทวปุตฺตปริสา อายสฺมโต สาริปุตฺตสฺส วณฺเณ ภญฺญมาเน อตฺตมนา ปมุทิตา ปีติโสมนสฺสชาตา อุจฺจาวจา วณฺณนิภา อุปทํเสติ.}

\prose{อถ โข สุสิโม เทวปุตฺโต อายสฺมนฺตํ สาริปุตฺตํ อารพฺภ ภควโต สนฺติเก อิมํ คาถํ อภาสิ\csromandash{}}

\csromangathameasure{“ปณฺฑิโตติ สมญฺญาโต, สาริปุตฺโต อโกธโน. \\ อปฺปิจฺโฉ โสรโต ทนฺโต, สตฺถุวณฺณาภโต อิสี”ติ.}
\csromangathagroup{\csromangathastanza{“ปณฺฑิโตติ สมญฺญาโต, สาริปุตฺโต อโกธโน. \\ อปฺปิจฺโฉ โสรโต ทนฺโต, สตฺถุวณฺณาภโต อิสี”ติ.}}

\prose{อถ โข ภควา อายสฺมนฺตํ สาริปุตฺตํ อารพฺภ สุสิมํ เทวปุตฺตํ คาถาย ปจฺจภาสิ\csromandash{}}

\csromangathameasure{“ปณฺฑิโตติ สมญฺญาโต, สาริปุตฺโต อโกธโน. \\ อปฺปิจฺโฉ โสรโต ทนฺโต, กาลํ กงฺขติ สุทนฺโต” ติ.}
\csromangathagroup{\csromangathastanza{“ปณฺฑิโตติ สมญฺญาโต, สาริปุตฺโต อโกธโน. \\ อปฺปิจฺโฉ โสรโต ทนฺโต, กาลํ กงฺขติ สุทนฺโต”\footnote{กาลํ กงฺขติ ภตโก สุทนฺโต (สี), กาลํ กงฺขติ ภาวิโต สุทนฺโต (สฺยา, กํ), กาลํ กงฺขติ ภติโก สุทนฺโต (อิ)} ติ.}}

\csromanheader{2}{10. นานาติตฺถิยสาวกสุตฺต}
\proseitem{111}{เอวํ เม สุตํ\csromandash{}เอกํ สมยํ ภควา ราชคเห วิหรติ เวฬุวเน กลนฺทกนิวาเป.\csromanspacer{} อถ โข สมฺพหุลา นานาติตฺถิยสาวกา เทวปุตฺตา อสโม จ สหลิ\footnote{สหลี (สี, สฺยา, กํ, อิ)} จ นีโก\footnote{นิงฺโก (สี,อิ), นิโก (สฺยา, กํ)} จ อาโกฏโก จ เวคพฺภริ จ\footnote{เวฏมฺพรี จ (สี, สฺยา, กํ, อิ)} มาณวคามิโย จ อภิกฺกนฺตาย รตฺติยา อภิกฺกนฺตวณฺณา เกวลกปฺปํ เวฬุวนํ โอภาเสตฺวา เยน ภควา เตนุปสงฺกมิํสุ, อุปสงฺกมิตฺวา ภควนฺตํ \csromanfolio{66}อภิวาเทตฺวา เอกมนฺตํ อฏฺฐํสุ, เอกมนฺตํ ฐิโต โข อสโม เทวปุตฺโต ปูรณํ กสฺสปํ อารพฺภ ภควโต สนฺติเก อิมํ คาถํ อภาสิ\csromandash{}}

\csromangathameasure{“อิธ ฉินฺทิตมาริเต, หตชานีสุ กสฺสโป. \\ น ปาปํ สมนุปสฺสติ, ปุญฺญํ วา ปน อตฺตโน. \\ ส เว วิสฺสาสมาจิกฺขิ, สตฺถา อรหติ มานนนฺ”ติ.}
\csromangathagroup{\csromangathastanza{“อิธ ฉินฺทิตมาริเต, หตชานีสุ กสฺสโป. \\ น ปาปํ สมนุปสฺสติ, ปุญฺญํ วา ปน อตฺตโน. \\ ส เว วิสฺสาสมาจิกฺขิ, สตฺถา อรหติ มานนนฺ”ติ.}}

\prose{อถ โข สหลิ เทวปุตฺโต มกฺขลิํ โคสาลํ อารพฺภ ภควโต สนฺติเก อิมํ คาถํ อภาสิ\csromandash{}}

\csromangathameasure{“ตโปชิคุจฺฉาย สุสํวุตตฺโต, \\ วาจํ ปหาย กลหํ ชเนน. \\ สโมสวชฺชา วิรโต สจฺจวาที, \\ น หิ นูน ตาทิสํ กโรติ ปาปนฺ”ติ.}
\csromangathagroup{\csromangathastanza{“ตโปชิคุจฺฉาย สุสํวุตตฺโต, \\ วาจํ ปหาย กลหํ ชเนน. \\ สโมสวชฺชา วิรโต สจฺจวาที, \\ น หิ นูน ตาทิสํ กโรติ\footnote{น ห นุน ตาที ปกโรติ (สี, สฺยา, กํ)} ปาปนฺ”ติ.}}

\prose{อถ โข นีโก เทวปุตฺโต นิคณฺฐํ นาฏปุตฺตํ\footnote{นาถปุตฺตํ (สี)} อารพฺภ ภควโต สนฺติเก อิมํ คาถํ อภาสิ\csromandash{}}

\csromangathameasure{“เชคุจฺฉี นิปโก ภิกฺขุ, จาตุยามสุสํวุโต. \\ ทิฏฺฐํ สุตญฺจ อาจิกฺขํ, น หิ นูน กิพฺพิสี สิยา”ติ.}
\csromangathagroup{\csromangathastanza{“เชคุจฺฉี นิปโก ภิกฺขุ, จาตุยามสุสํวุโต. \\ ทิฏฺฐํ สุตญฺจ อาจิกฺขํ, น หิ นูน กิพฺพิสี สิยา”ติ.}}

\prose{อถ โข อาโกฏโก เทวปุตฺโต นานาติตฺถิเย อารพฺภ ภควโต สนฺติเก}

\csromangathameasure{“ปกุธโก กาติยาโน นิคณฺโฐ, \\ เย จาปิเม มกฺขลิปูรณาเส. \\ คณสฺส สตฺถาโร สามญฺญปฺปตฺตา, \\ น หิ นูน เต สปฺปุริเสหิ ทูเร”ติ.}
\csromangathagroup{\csromangathastanza{“ปกุธโก กาติยาโน นิคณฺโฐ, \\ เย จาปิเม มกฺขลิปูรณาเส. \\ คณสฺส สตฺถาโร สามญฺญปฺปตฺตา, \\ น หิ นูน เต สปฺปุริเสหิ ทูเร”ติ.}}

\prose{อถ โข เวคพฺภริ เทวปุตฺโต อาโกฏกํ เทวปุตฺตํ คาถาย ปจฺจภาสิ\csromandash{}}

\csromangathameasure{“สหาจริเตน ฉโว สิคาโล, \\ น โกตฺถุโก สีหสโม กทาจิ. \\ นคฺโค มุสาวาที คณสฺส สตฺถา, \\ สงฺกสฺสราจาโร น สตํ สริกฺโข”ติ.}
\csromangathagroup{\csromangathastanza{“สหาจริเตน\footnote{สหารเวนาปิ (ก-สี), สคารเวนาปิ (อิ)} ฉโว สิคาโล\footnote{สิงฺคาโล (ก)}, \\ น โกตฺถุโก สีหสโม กทาจิ. \\ นคฺโค มุสาวาที คณสฺส สตฺถา, \\ สงฺกสฺสราจาโร น สตํ สริกฺโข”ติ.}}

\csromanheader{2}{2. เทวปุตฺตสํยุตฺต}
\prose{\csromanfolio{67}อถ โข มาโร ปาปิมา เวคพฺภริํ เทวปุตฺตํ อนฺวาวิสิตฺวา ภควโต สนฺติเก อิมํ คาถํ อภาสิ\csromandash{}}

\csromangathameasure{“ตโปชิคุจฺฉาย อายุตฺตา, ปาลยํ ปวิเวกิยํ. \\ รูเป จ เย นิวิฏฺฐาเส, เทวโลกาภินนฺทิโน. \\ เต เว สมฺมานุสาสนฺติ, ปรโลกาย มาติยา”ติ.}
\csromangathagroup{\csromangathastanza{“ตโปชิคุจฺฉาย อายุตฺตา, ปาลยํ ปวิเวกิยํ. \\ รูเป จ เย นิวิฏฺฐาเส, เทวโลกาภินนฺทิโน. \\ เต เว สมฺมานุสาสนฺติ, ปรโลกาย มาติยา”ติ.}}

\prose{อถ โข ภควา “มาโร อยํ ปาปิมา” อิติ วิทิตฺวา มารํ ปาปิมนฺตํ คาถาย ปจฺจภาสิ\csromandash{}}

\csromangathameasure{“เย เกจิ รูปา อิธ วา หุรํ วา, \\ เย จนฺตลิกฺขสฺมิํ ปภาสวณฺณา. \\ สพฺเพว เต เต นมุจิปฺปสตฺถา, \\ อามิสํว มจฺฉานํ วธาย ขิตฺตา”ติ.}
\csromangathagroup{\csromangathastanza{“เย เกจิ รูปา อิธ วา หุรํ วา, \\ เย จนฺตลิกฺขสฺมิํ ปภาสวณฺณา. \\ สพฺเพว เต เต นมุจิปฺปสตฺถา, \\ อามิสํว มจฺฉานํ วธาย ขิตฺตา”ติ.}}

\prose{อถ โข มาณวคามิโย เทวปุตฺโต ภควนฺตํ อารพฺภ ภควโต สนฺติเก อิมา คาถาโย อภาสิ\csromandash{}}

\csromangathameasure{“วิปุโล ราชคหียานํ, คิริเสฏฺโฐ ปวุจฺจติ. \\ เสโต หิมวตํ เสฏฺโฐ, อาทิจฺโจ อฆคามินํ. \\ สมุทฺโท อุทธินํ เสฏฺโฐ, นกฺขตฺตานญฺจ จนฺทิมา. \\ สเทวกสฺส โลกสฺส, พุทฺโธ อคฺโค ปวุจฺจตี”ติ.}
\csromangathagroup{\csromangathastanza{“วิปุโล ราชคหียานํ, คิริเสฏฺโฐ ปวุจฺจติ. \\ เสโต หิมวตํ เสฏฺโฐ, อาทิจฺโจ อฆคามินํ.}\csromangathastanza{สมุทฺโท อุทธินํ เสฏฺโฐ, นกฺขตฺตานญฺจ จนฺทิมา\footnote{นกฺขตฺตานํว จนฺทิมา (ก)}. \\ สเทวกสฺส โลกสฺส, พุทฺโธ อคฺโค ปวุจฺจตี”ติ.}}

\vspace{\tipitakaparskip}
\csromancenter{นานาติตฺถิยวคฺโค ตติโย.}
\titlehead{\textbf{ตสฺสุทฺทานํ}}
\csromangathameasure{สิโว เขโม จ เสรี จ, ฆฏี ชนฺตุ จ โรหิโต. \\ นนฺโท นนฺทิวิสาโล จ, สุสิโม นานาติตฺถิเยน เต ทสาติ. \\ \textbf{เทวปุตฺตสํยุตฺตํ สมตฺตํ.}}
\csromangathagroup{\csromangathastanza{สิโว เขโม จ เสรี จ, ฆฏี ชนฺตุ จ โรหิโต. \\ นนฺโท นนฺทิวิสาโล จ, สุสิโม นานาติตฺถิเยน เต ทสาติ. \\ \textbf{เทวปุตฺตสํยุตฺตํ สมตฺตํ.}}}

\csromanheader{2}{3. โกสลสํยุตฺต}
\chapterhead{1. ปฐมวคฺค}
\csromanheader{2}{1. ทหรสุตฺต}
\proseitem{112}{\csromanfolio{68}เอวํ เม สุตํ\csromandash{}เอกํ สมยํ ภควา สาวตฺถิยํ วิหรติ เชตวเน อนาถปิณฺฑิกสฺส อาราเม.\csromanspacer{} อถ โข ราชา ปเสนทิ โกสโล เยน ภควา เตนุปสงฺกมิ, อุปสงฺกมิตฺวา ภควตา สทฺธิํ สมฺโมทิ, สมฺโมทนียํ กถํ สารณียํ วีติสาเรตฺวา เอกมนฺตํ นิสีทิ, เอกมนฺตํ นิสินฺโน โข ราชา ปเสนทิ โกสโล ภควนฺตํ เอตทโวจ “ภวมฺปิ โน โคตโม ‘อนุตฺตรํ สมฺมาสมฺโพธิํ อภิสมฺพุทฺโธ’ติ ปฏิชานาตี” ติ.\csromanspacer{} ยํ หิ ตํ มหาราช สมฺมา วทมาโน วเทยฺย “อนุตฺตรํ สมฺมาสมฺโพธิํ อภิสมฺพุทฺโธ”ติ, มเมว\footnote{มมํ (สพฺพตฺถ)} ตํ สมฺมา วทมาโน วเทยฺย, อหํ หิ มหาราช อนุตฺตรํ สมฺมาสมฺโพธิํ อภิสมฺพุทฺโธติ.}

\prose{เยปิ เต โภ โคตม สมณพฺราหฺมณา สงฺฆิโน คณิโน คณาจริยา ญาตา ยสสฺสิโน ติตฺถกรา สาธุสมฺมตา พหุชนสฺส.\csromanspacer{} เสยฺยถิทํ, ปูรโณ กสฺสโป มกฺขลิ โคสาโล นิคณฺโฐ นาฏปุตฺโต สญฺจโย เพลฏฺฐปุตฺโต ปกุโธ กจฺจายโน อชิโต เกสกมฺพโล, เตปิ มยา “อนุตฺตรํ สมฺมาสมฺโพธิํ อภิสมฺพุทฺโธติ ปฏิชานาถา”ติ ปุฏฺฐา สมานา “อนุตฺตรํ สมฺมาสมฺโพธิํ อภิสมฺพุทฺโธ”ติ น ปฏิชานนฺติ, กิํ ปน ภวํ โคตโม ทหโร เจว ชาติยา, นโว จ ปพฺพชฺชายาติ.}

\prose{จตฺตาโร โข เม มหาราช “ทหรา”ติ น อุญฺญาตพฺพา “ทหรา”ติ น ปริโภตพฺพา.\csromanspacer{} กตเม จตฺตาโร, ขตฺติโย โข มหาราช “ทหโร”ติ น อุญฺญาตพฺโพ “ทหโร”ติ น ปริโภตพฺโพ, อุรโค โข มหาราช “ทหโร”ติ น อุญฺญาตพฺโพ “ทหโร”ติ น ปริโภตพฺโพ, อคฺคิ โข มหาราช “ทหโร”ติ น อุญฺญาตพฺโพ “ทหโร”ติ น ปริโภตพฺโพ, ภิกฺขุ โข มหาราช “ทหโร”ติ น อุญฺญาตพฺโพ “ทหโร”ติ น ปริโภตพฺโพ.\csromanspacer{} อิเม โข มหาราช จตฺตาโร “ทหรา”ติ น อุญฺญาตพฺพา “ทหรา”ติ น ปริโภตพฺพาติ.}

\csromanheader{2}{3. โกสลสํยุตฺต}
\prose{\csromanfolio{69}อิทมโวจ ภควา, อิทํ วตฺวาน สุคโต อถาปรํ เอตทโวจ สตฺถา\csromandash{}}

\csromangathameasure{ขตฺติยํ ชาติสมฺปนฺนํ, อภิชาตํ ยสสฺสินํ. \\ “ทหโร”ติ นาวชาเนยฺย, น นํ ปริภเว นโร. \\ ฐานํ หิ โส มนุชินฺโท, รชฺชํ ลทฺธาน ขตฺติโย. \\ โส กุทฺโธ ราชทณฺเฑน, ตสฺมิํ ปกฺกมเต ภุสํ. \\ ตสฺมา ตํ ปริวชฺเชยฺย, รกฺขํ ชีวิตมตฺตโน. \\ คาเม วา ยทิ วา รญฺเญ, ยตฺถ ปสฺเส ภุชงฺคมํ. \\ “ทหโร”ติ นาวชาเนยฺย, น นํ ปริภเว นโร. \\ อุจฺจาวเจหิ วณฺเณหิ, อุรโค จรติ เตชสี. \\ โส อาสชฺช ฑํเส พาลํ, นรํ นาริํ จ เอกทา. \\ ตสฺมา ตํ ปริวชฺเชยฺย, รกฺขํ ชีวิตมตฺตโน. \\ ปหูตภกฺขํ ชาลินํ, ปาวกํ กณฺหวตฺตนิํ. \\ ทหโรติ นาวชาเนยฺย, น นํ ปริภเว นโร. \\ ลทฺธา หิ โส อุปาทานํ, มหา หุตฺวาน ปาวโก. \\ โส อาสชฺช ฑเห พาลํ, นรํ นาริํ จ เอกทา. \\ ตสฺมา ตํ ปริวชฺเชยฺย, รกฺขํ ชีวิตมตฺตโน. \\ วนํ ยทคฺคิ ฑหติ, ปาวโก กณฺหวตฺตนี. \\ ชายนฺติ ตตฺถ ปาโรหา, อโหรตฺตานมจฺจเย. \\ ยญฺจ โข สีลสมฺปนฺโน, ภิกฺขุ ฑหติ เตชสา. \\ น ตสฺส ปุตฺตา ปสโว, ทายาทา วินฺทเร ธนํ. \\ อนปจฺจา อทายาทา, ตาลาวตฺถู ภวนฺติ เต. \\ ตสฺมา หิ ปณฺฑิโต โปโส, สมฺปสฺสํ อตฺถมตฺตโน. \\ ภุชงฺคมํ ปาวกญฺจ, ขตฺติยญฺจ ยสสฺสินํ. \\ ภิกฺขุญฺจ สีลสมฺปนฺนํ, สมฺมเทว สมาจเรติ.}
\csromangathagroup{\csromangathastanza{ขตฺติยํ ชาติสมฺปนฺนํ, อภิชาตํ ยสสฺสินํ. \\ “ทหโร”ติ นาวชาเนยฺย, น นํ ปริภเว นโร.}\csromangathastanza{ฐานํ หิ โส มนุชินฺโท, รชฺชํ ลทฺธาน ขตฺติโย. \\ โส กุทฺโธ ราชทณฺเฑน, ตสฺมิํ ปกฺกมเต ภุสํ.}\csromangathastanza{ตสฺมา ตํ ปริวชฺเชยฺย, รกฺขํ ชีวิตมตฺตโน. \\ คาเม วา ยทิ วา รญฺเญ, ยตฺถ ปสฺเส ภุชงฺคมํ.}\csromangathastanza{“ทหโร”ติ นาวชาเนยฺย, น นํ ปริภเว นโร. \\ อุจฺจาวเจหิ วณฺเณหิ, อุรโค จรติ เตชสี\footnote{เตชสา (สี, ก), เตชสิ (อิ, ก)}.}\csromangathastanza{โส อาสชฺช ฑํเส พาลํ, นรํ นาริํ จ เอกทา. \\ ตสฺมา ตํ ปริวชฺเชยฺย, รกฺขํ ชีวิตมตฺตโน.}\csromangathastanza{ปหูตภกฺขํ ชาลินํ, ปาวกํ กณฺหวตฺตนิํ. \\ ทหโรติ นาวชาเนยฺย, น นํ ปริภเว นโร.}\csromangathastanza{ลทฺธา หิ โส อุปาทานํ, มหา หุตฺวาน ปาวโก. \\ โส อาสชฺช ฑเห\footnote{ทเห} พาลํ, นรํ นาริํ จ เอกทา.}\csromangathastanza{ตสฺมา ตํ ปริวชฺเชยฺย, รกฺขํ ชีวิตมตฺตโน. \\ วนํ ยทคฺคิ ฑหติ\footnote{ทหติ (ก)}, ปาวโก กณฺหวตฺตนี.}\csromangathastanza{ชายนฺติ ตตฺถ ปาโรหา, อโหรตฺตานมจฺจเย. \\ ยญฺจ โข สีลสมฺปนฺโน, ภิกฺขุ ฑหติ เตชสา.}\csromangathastanza{น ตสฺส ปุตฺตา ปสโว, ทายาทา วินฺทเร ธนํ. \\ อนปจฺจา อทายาทา, ตาลาวตฺถู ภวนฺติ เต.}\csromangathastanza{ตสฺมา หิ ปณฺฑิโต โปโส, สมฺปสฺสํ อตฺถมตฺตโน. \\ ภุชงฺคมํ ปาวกญฺจ, ขตฺติยญฺจ ยสสฺสินํ. \\ ภิกฺขุญฺจ สีลสมฺปนฺนํ, สมฺมเทว สมาจเรติ.}}

\prose{เอวํ วุตฺเต ราชา ปเสนทิ โกสโล ภควนฺตํ เอตทโวจ “อภิกฺกนฺตํ ภนฺเต อภิกฺกนฺตํ ภนฺเต, เสยฺยถาปิ ภนฺเต นิกฺกุชฺชิตํ\footnote{นิกุชฺชิตํ (?)} วา}

\prose{\csromanfolio{70}อุกฺกุชฺเชยฺย, ปฏิจฺฉนฺนํ วา วิวเรยฺย, มูฬฺหสฺส วา มคฺคํ อาจิกฺเขยฺย, อนฺธกาเร วา เตลปชฺโชตํ ธาเรยฺย ‘จกฺขุมนฺโต รูปานิ ทกฺขนฺตี’ติ.\csromanspacer{} เอวเมวํ ภควตา อเนกปริยาเยน ธมฺโม ปกาสิโต, เอสาหํ ภนฺเต ภควนฺตํ สรณํ คจฺฉามิ ธมฺมญฺจ ภิกฺขุสํฆญฺจ, อุปาสกํ มํ ภนฺเต ภควา ธาเรตุ อชฺชตคฺเค ปาณุเปตํ สรณํ คตนฺ”ติ.}

\csromanheader{2}{2. ปุริสสุตฺต}
\proseitem{113}{สาวตฺถินิทานํ.\csromanspacer{} อถ โข ราชา ปเสนทิ โกสโล เยน ภควา เตนุปสงฺกมิ, อุปสงฺกมิตฺวา ภควนฺตํ อภิวาเทตฺวา เอกมนฺตํ นิสีทิ, เอกมนฺตํ นิสินฺโน โข ราชา ปเสนทิ โกสโล ภควนฺตํ เอตทโวจ “กติ นุ โข ภนฺเต ปุริสสฺส ธมฺมา อชฺฌตฺตํ อุปฺปชฺชมานา อุปฺปชฺชนฺติ อหิตาย ทุกฺขาย อผาสุวิหารายา”ติ.}

\prose{ตโย โข มหาราช ปุริสสฺส ธมฺมา อชฺฌตฺตํ อุปฺปชฺชมานา อุปฺปชฺชนฺติ อหิตาย ทุกฺขาย อผาสุวิหาราย.\csromanspacer{} กตเม ตโย, โลโภ โข มหาราช ปุริสสฺส ธมฺโม อชฺฌตฺตํ อุปฺปชฺชมาโน อุปฺปชฺชติ อหิตาย ทุกฺขาย อผาสุวิหาราย.\csromanspacer{} โทโส โข มหาราช ปุริสสฺส ธมฺโม อชฺฌตฺตํ อุปฺปชฺชมาโน อุปฺปชฺชติ อหิตาย ทุกฺขาย อผาสุวิหาราย.\csromanspacer{} โมโห โข มหาราช ปุริสสฺส ธมฺโม อชฺฌตฺตํ อุปฺปชฺชมาโน อุปฺปชฺชติ อหิตาย ทุกฺขาย อผาสุวิหาราย.\csromanspacer{} อิเม โข มหาราช ตโย ปุริสสฺส ธมฺมา อชฺฌตฺตํ อุปฺปชฺชมานา อุปฺปชฺชนฺติ อหิตาย ทุกฺขาย อผาสุวิหารายาติ.\csromanspacer{} อิทมโวจ ฯเปฯ.}

\csromangathameasure{โลโภ โทโส จ โมโห จ, ปุริสํ ปาปเจตสํ. \\ หิํสนฺติ อตฺตสมฺภูตา, ตจสารํว สมฺผลนฺติ.}
\csromangathagroup{\csromangathastanza{โลโภ โทโส จ โมโห จ, ปุริสํ ปาปเจตสํ. \\ หิํสนฺติ อตฺตสมฺภูตา, ตจสารํว สมฺผลนฺติ\footnote{สปฺผลนฺติ (สฺยา, กํ)}.}}

\csromanheader{2}{3. ชรามรณสุตฺต}
\proseitem{114}{สาวตฺถินิทานํ.\csromanspacer{} เอกมนฺตํ นิสินฺโน โข ราชา ปเสนทิ โกสโล ภควนฺตํ เอตทโวจ “อตฺถิ นุ โข ภนฺเต ชาตสฺส อญฺญตฺร ชรามรณา”ติ.\csromanspacer{} นตฺถิ โข มหาราช ชาตสฺส อญฺญตฺร ชรามรณา.\csromanspacer{} เยปิ เต มหาราช ขตฺติยมหาสาลา อฑฺฒา มหทฺธนา มหาโภคา}

\vspace{\tipitakaparskip}
\csromancenter{โกสลสํยุตฺต ปหูตชาตรูปรชตา ปหูตวิตฺตูปกรณา ปหูตธนธญฺญา, เตสมฺปิ ชาตานํ นตฺถิ อญฺญตฺร ชรามรณา.\csromanspacer{} เยปิ เต มหาราช พฺราหฺมณมหาสาลา ฯเปฯ คหปติมหาสาลา อฑฺฒา มหทฺธนา มหาโภคา ปหูตชาตรูปรชตา ปหูตวิตฺตูปกรณา ปหูตธนธญฺญา, เตสมฺปิ ชาตานํ นตฺถิ อญฺญตฺร ชรามรณา.\csromanspacer{} เยปิ เต มหาราช ภิกฺขู อรหนฺโต ขีณาสวา วุสิตวนฺโต กตกรณียา โอหิตภารา อนุปฺปตฺตสทตฺถา ปริกฺขีณภวสํโยชนา สมฺมทญฺญาวิมุตฺตา, เตสํปายํ กาโย เภทนธมฺโม นิกฺเขปนธมฺโมติ.\csromanspacer{} อิทมโวจ ฯเปฯ.}
\vspace{0.5\baselineskip}
\csromangathameasure{ชีรนฺติ เว ราชรถา สุจิตฺตา, อโถ สรีรมฺปิ ชรํ อุเปติ. \\ สตญฺจ ธมฺโม น ชรํ อุเปติ, สนฺโต หเว สพฺภิ ปเวทยนฺตีติ.}
\csromangathagroup{\csromangathastanza{\csromanfolio{71}ชีรนฺติ เว ราชรถา สุจิตฺตา, อโถ สรีรมฺปิ ชรํ อุเปติ. \\ สตญฺจ ธมฺโม น ชรํ อุเปติ, สนฺโต หเว สพฺภิ ปเวทยนฺตีติ.}}

\csromanheader{2}{4. ปิยสุตฺต}
\proseitem{115}{สาวตฺถินิทานํ.\csromanspacer{} เอกมนฺตํ นิสินฺโน โข ราชา ปเสนทิ โกสโล ภควนฺตํ เอตทโวจ\csromandash{}อิธ มยฺหํ ภนฺเต รโหคตสฺส ปฏิสลฺลีนสฺส เอวํ เจตโส ปริวิตกฺโก อุทปาทิ “เกสํ นุ โข ปิโย อตฺตา เกสํ อปฺปิโย อตฺตา”ติ.\csromanspacer{} ตสฺส มยฺหํ ภนฺเต เอตทโหสิ “เย โข เกจิ กาเยน ทุจฺจริตํ จรนฺติ, วาจาย ทุจฺจริตํ จรนฺติ, มนสา ทุจฺจริตํ จรนฺติ, เตสํ อปฺปิโย อตฺตา.\csromanspacer{} กิญฺจาปิ เต เอวํ วเทยฺยุํ ‘ปิโย โน อตฺตา’ติ.\csromanspacer{} อถ โข เตสํ อปฺปิโย อตฺตา.\csromanspacer{} ตํ กิสฺส เหตุ, ยํ หิ อปฺปิโย อปฺปิยสฺส กเรยฺย, ตํ เต อตฺตนาว อตฺตโน กโรนฺติ.\csromanspacer{} ตสฺมา เตสํ อปฺปิโย อตฺตา.\csromanspacer{} เย จ โข เกจิ กาเยน สุจริตํ จรนฺติ, วาจาย สุจริตํ จรนฺติ, มนสา สุจริตํ จรนฺติ, เตสํ ปิโย อตฺตา.\csromanspacer{} กิญฺจาปิ เต เอวํ วเทยฺยุํ ‘อปฺปิโย โน อตฺตา’ติ.\csromanspacer{} อถ โข เตสํ ปิโย อตฺตา.\csromanspacer{} ตํ กิสฺส เหตุ, ยํ หิ ปิโย ปิยสฺส กเรยฺย, ตํ เต อตฺตนาว อตฺตโน กโรนฺติ.\csromanspacer{} ตสฺมา เตสํ ปิโย อตฺตา”ติ.}

\prose{เอวเมตํ มหาราช เอวเมตํ มหาราช, เย หิ เกจิ มหาราช กาเยน ทุจฺจริตํ จรนฺติ, วาจาย ทุจฺจริตํ จรนฺติ, มนสา ทุจฺจริตํ จรนฺติ, เตสํ อปฺปิโย อตฺตา.\csromanspacer{} กิญฺจาปิ เต เอวํ วเทยฺยุํ “ปิโย โน อตฺตา”ติ.\csromanspacer{} อถ โข เตสํ อปฺปิโย อตฺตา.\csromanspacer{} ตํ กิสฺส เหตุ, ยํ หิ มหาราช อปฺปิโย อปฺปิยสฺส กเรยฺย, ตํ เต อตฺตนาว อตฺตโน กโรนฺติ.\csromanspacer{} ตสฺมา เตสํ อปฺปิโย อตฺตา.\csromanspacer{} เย จ โข เกจิ มหาราช \csromanfolio{72}กาเยน สุจริตํ จรนฺติ, วาจาย สุจริตํ จรนฺติ, มนสา สุจริตํ จรนฺติ, เตสํ ปิโย อตฺตา.\csromanspacer{} กิญฺจาปิ เต เอวํ วเทยฺยุํ “อปฺปิโย โน อตฺตา”ติ.\csromanspacer{} อถ โข เตสํ ปิโย อตฺตา.\csromanspacer{} ตํ กิสฺส เหตุ, ยํ หิ มหาราช ปิโย ปิยสฺส กเรยฺย, ตํ เต อตฺตนาว อตฺตโน กโรนฺติ.\csromanspacer{} ตสฺมา เตสํ ปิโย อตฺตาติ.\csromanspacer{} อิทมโวจ ฯเปฯ.}

\csromangathameasure{อตฺตานญฺเจ ปิยํ ชญฺญา, น นํ ปาเปน สํยุเช. \\ น หิ ตํ สุลภํ โหติ, สุขํ ทุกฺกฏการินา. \\ อนฺตเกนาธิปนฺนสฺส, ชหโต มานุสํ ภวํ. \\ กิํ หิ ตสฺส สกํ โหติ, กิญฺจ อาทาย คจฺฉติ. \\ กิญฺจสฺส อนุคํ โหติ, ฉายาว อนปายินี. \\ อุโภ ปุญฺญญฺจ ปาปญฺจ, ยํ มจฺโจ กุรุเต อิธ. \\ ตํ หิ ตสฺส สกํ โหติ, ตญฺจ อาทาย คจฺฉติ. \\ ตญฺจสฺส อนุคํ โหติ, ฉายาว อนปายินี. \\ ตสฺมา กเรยฺย กลฺยาณํ, นิจยํ สมฺปรายิกํ. \\ ปุญฺญานิ ปรโลกสฺมิํ, ปติฏฺฐา โหนฺติ ปาณินนฺติ.}
\csromangathagroup{\csromangathastanza{อตฺตานญฺเจ ปิยํ ชญฺญา, น นํ ปาเปน สํยุเช. \\ น หิ ตํ สุลภํ โหติ, สุขํ ทุกฺกฏการินา.}\csromangathastanza{อนฺตเกนาธิปนฺนสฺส, ชหโต มานุสํ ภวํ. \\ กิํ หิ ตสฺส สกํ โหติ, กิญฺจ อาทาย คจฺฉติ.}\csromangathastanza{กิญฺจสฺส อนุคํ โหติ, ฉายาว อนปายินี\footnote{อนุปายินี (สฺยา, กํ, ก)}. \\ อุโภ ปุญฺญญฺจ ปาปญฺจ, ยํ มจฺโจ กุรุเต อิธ.}\csromangathastanza{ตํ หิ ตสฺส สกํ โหติ, ตญฺจ\footnote{ตํว (?)} อาทาย คจฺฉติ. \\ ตญฺจสฺส\footnote{ตํว (?)} อนุคํ โหติ, ฉายาว อนปายินี.}\csromangathastanza{ตสฺมา กเรยฺย กลฺยาณํ, นิจยํ สมฺปรายิกํ. \\ ปุญฺญานิ ปรโลกสฺมิํ, ปติฏฺฐา โหนฺติ ปาณินนฺติ.}}

\csromanheader{2}{5. อตฺตรกฺขิตสุตฺต}
\proseitem{116}{สาวตฺถินิทานํ.\csromanspacer{} เอกมนฺตํ นิสินฺโน โข ราชา ปเสนทิ โกสโล ภควนฺตํ เอตทโวจ\csromandash{}อิธ มยฺหํ ภนฺเต รโหคตสฺส ปฏิสลฺลีนสฺส เอวํ เจตโส ปริวิตกฺโก อุทปาทิ “เกสํ นุ โข รกฺขิโต อตฺตา เกสํ อรกฺขิโต อตฺตา”ติ.\csromanspacer{} ตสฺส มยฺหํ ภนฺเต เอตทโหสิ “เย โข เกจิ กาเยน ทุจฺจริตํ จรนฺติ, วาจาย ทุจฺจริตํ จรนฺติ, มนสา ทุจฺจริตํ จรนฺติ, เตสํ อรกฺขิโต อตฺตา.\csromanspacer{} กิญฺจาปิ เต หตฺถิกาโย วา รกฺเขยฺย, อสฺสกาโย วา รกฺเขยฺย, รถกาโย วา รกฺเขยฺย, ปตฺติกาโย วา รกฺเขยฺย, อถ โข เตสํ อรกฺขิโต อตฺตา.\csromanspacer{} ตํ กิสฺส เหตุ, พาหิรา เหสา รกฺขา, เนสา รกฺขา อชฺฌตฺติกา, ตสฺมา เตสํ อรกฺขิโต อตฺตา.\csromanspacer{} เย จ โข เกจิ กาเยน สุจริตํ จรนฺติ, วาจาย สุจริตํ จรนฺติ, มนสา สุจริตํ จรนฺติ, เตสํ รกฺขิโต อตฺตา.\csromanspacer{} กิญฺจาปิ เต เนว หตฺถิกาโย รกฺเขยฺย, น อสฺสกาโย รกฺเขยฺย, น รถกาโย}

\vspace{\tipitakaparskip}
\csromancenter{โกสลสํยุตฺต รกฺเขยฺย, น ปตฺติกาโย รกฺเขยฺย, อถ โข เตสํ รกฺขิโต อตฺตา.\csromanspacer{} ตํ กิสฺส เหตุ, อชฺฌตฺติกา เหสา รกฺขา, เนสา รกฺขา พาหิรา, ตสฺมา เตสํ รกฺขิโต อตฺตา”ติ.}
\prose{\csromanfolio{73}เอวเมตํ มหาราช เอวเมตํ มหาราช, เย หิ เกจิ มหาราช กาเยน ทุจฺจริตํ จรนฺติ ฯเปฯ เตสํ อรกฺขิโต อตฺตา.\csromanspacer{} ตํ กิสฺส เหตุ, พาหิรา เหสา มหาราช รกฺขา, เนสา รกฺขา อชฺฌตฺติกา, ตสฺมา เตสํ อรกฺขิโต อตฺตา.\csromanspacer{} เย จ โข เกจิ มหาราช กาเยน สุจริตํ จรนฺติ, วาจาย สุจริตํ จรนฺติ, มนสา สุจริตํ จรนฺติ, เตสํ รกฺขิโต อตฺตา.\csromanspacer{} กิญฺจาปิ เต เนว หตฺถิกาโย รกฺเขยฺย, น อสฺสกาโย รกฺเขยฺย, น รถกาโย รกฺเขยฺย, น ปตฺติกาโย รกฺเขยฺย, อถ โข เตสํ รกฺขิโต อตฺตา.\csromanspacer{} ตํ กิสฺส เหตุ, อชฺฌตฺติกา เหสา มหาราช รกฺขา, เนสา รกฺขา พาหิรา, ตสฺมา เตสํ รกฺขิโต อตฺตาติ.\csromanspacer{} อิทมโวจ ฯเปฯ.}

\csromangathameasure{กาเยน สํวโร สาธุ, สาธุ วาจาย สํวโร. \\ มนสา สํวโร สาธุ, สาธุ สพฺพตฺถ สํวโร. \\ สพฺพตฺถ สํวุโต ลชฺชี, รกฺขิโตติ ปวุจฺจตีติ.}
\csromangathagroup{\csromangathastanza{กาเยน สํวโร สาธุ, สาธุ วาจาย สํวโร. \\ มนสา สํวโร สาธุ, สาธุ สพฺพตฺถ สํวโร. \\ สพฺพตฺถ สํวุโต ลชฺชี, รกฺขิโตติ ปวุจฺจตีติ.}}

\csromanheader{2}{6. อปฺปกสุตฺต}
\proseitem{117}{สาวตฺถินิทานํ.\csromanspacer{} เอกมนฺตํ นิสินฺโน โข ราชา ปเสนทิ โกสโล ภควนฺตํ เอตทโวจ “อิธ มยฺหํ ภนฺเต รโหคตสฺส ปฏิสลฺลีนสฺส เอวํ เจตโส ปริวิตกฺโก อุทปาทิ ‘อปฺปกา เต สตฺตา โลกสฺมิํ, เย อุฬาเร อุฬาเร โภเค ลภิตฺวา น เจว มชฺชนฺติ น จ ปมชฺชนฺติ, น จ กาเมสุ เคธํ อาปชฺชนฺติ, น จ สตฺเตสุ วิปฺปฏิปชฺชนฺติ.\csromanspacer{} อถ โข เอเตว พหุตรา สตฺตา โลกสฺมิํ, เย อุฬาเร อุฬาเร โภเค ลภิตฺวา มชฺชนฺติ เจว ปมชฺชนฺติ จ, กาเมสุ จ เคธํ อาปชฺชนฺติ, สตฺเตสุ จ วิปฺปฏิปชฺชนฺตี’ติ”.}

\prose{เอวเมตํ มหาราช เอวเมตํ มหาราช, อปฺปกา เต มหาราช สตฺตา โลกสฺมิํ, เย อุฬาเร อุฬาเร โภเค ลภิตฺวา น เจว มชฺชนฺติ น จ ปมชฺชนฺติ, น จ กาเมสุ เคธํ อาปชฺชนฺติ, น จ สตฺเตสุ วิปฺปฏิปชฺชนฺติ.\csromanspacer{} อถ โข เอเตว พหุตรา สตฺตา โลกสฺมิํ, เย อุฬาเร อุฬาเร โภเค \csromanfolio{74}ลภิตฺวา มชฺชนฺติ เจว ปมชฺชนฺติ จ, กาเมสุ จ เคธํ อาปชฺชนฺติ, สตฺเตสุ จ วิปฺปฏิปชฺชนฺตีติ.\csromanspacer{} อิทมโวจ ฯเปฯ.}

\csromangathameasure{สารตฺตา กามโภเคสุ, คิทฺธา กาเมสุ มุจฺฉิตา. \\ อติสารํ น พุชฺฌนฺติ, มิคา กูฏํว โอฑฺฑิตํ. \\ ปจฺฉา’สํ กฏุกํ โหติ, วิปาโก หิสฺส ปาปโกติ.}
\csromangathagroup{\csromangathastanza{สารตฺตา กามโภเคสุ, คิทฺธา กาเมสุ มุจฺฉิตา. \\ อติสารํ น พุชฺฌนฺติ, มิคา กูฏํว โอฑฺฑิตํ. \\ ปจฺฉา’สํ กฏุกํ โหติ, วิปาโก หิสฺส ปาปโกติ.}}

\csromanheader{2}{7. อฑฺฑกรณสุตฺต}
\proseitem{118}{สาวตฺถินิทานํ.\csromanspacer{} เอกมนฺตํ นิสินฺโน โข ราชา ปเสนทิ โกสโล ภควนฺตํ เอตทโวจ “อิธาหํ ภนฺเต อฑฺฑกรเณ\footnote{อตฺถกรเณ (สี, สฺยา, กํ, อิ)} นิสินฺโน ปสฺสามิ ขตฺติยมหาสาเลปิ พฺราหฺมณมหาสาเลปิ คหปติมหาสาเลปิ อฑฺเฒ มหทฺธเน มหาโภเค ปหูตชาตรูปรชเต ปหูตวิตฺตูปกรเณ ปหูตธนธญฺเญ กามเหตุ กามนิทานํ กามาธิกรณํ สมฺปชานมุสา ภาสนฺเต.\csromanspacer{} ตสฺส มยฺหํ ภนฺเต เอตทโหสิ ‘อลํ ทานิ เม อฑฺฑกรเณน, ภทฺรมุโข ทานิ อฑฺฑกรเณน ปญฺญายิสฺสตี’ติ”.}

\prose{(เอวเมตํ มหาราช เอวเมตํ มหาราช,)\footnote{( ) สี-อิ-โปตฺถเกสุ นตฺถิ.} เยปิ เต มหาราช ขตฺติยมหาสาลา พฺราหฺมณมหาสาลา คหปติมหาสาลา อฑฺฒา มหทฺธนา มหาโภคา ปหูตชาตรูปรชตา ปหูตวิตฺตูปกรณา ปหูตธนธญฺญา กามเหตุ กามนิทานํ กามาธิกรณํ สมฺปชานมุสา ภาสนฺติ, เตสํ ตํ ภวิสฺสติ ทีฆรตฺตํ อหิตาย ทุกฺขายาติ.\csromanspacer{} อิทมโวจ ฯเปฯ.}

\csromangathameasure{สารตฺตา กามโภเคสุ, คิทฺธา กาเมสุ มุจฺฉิตา. \\ อติสารํ น พุชฺฌนฺติ, มจฺฉา ขิปฺปํว โอฑฺฑิตํ. \\ ปจฺฉา’สํ กฏุกํ โหติ, วิปาโก หิสฺส ปาปโกติ.}
\csromangathagroup{\csromangathastanza{สารตฺตา กามโภเคสุ, คิทฺธา กาเมสุ มุจฺฉิตา. \\ อติสารํ น พุชฺฌนฺติ, มจฺฉา ขิปฺปํว โอฑฺฑิตํ. \\ ปจฺฉา’สํ กฏุกํ โหติ, วิปาโก หิสฺส ปาปโกติ.}}

\csromanheader{2}{8. มลฺลิกาสุตฺต}
\proseitem{119}{สาวตฺถินิทานํ.\csromanspacer{} เตน โข ปน สมเยน ราชา ปเสนทิ โกสโล มลฺลิกาย เทวิยา สทฺธิํ อุปริปาสาทวรคโต โหติ.\csromanspacer{} อถ โข ราชา ปเสนทิ โกสโล มลฺลิกํ เทวิํ เอตทโวจ “อตฺถิ นุ โข เต มลฺลิเก โกจญฺโญ อตฺตนา ปิยตโร”ติ.\csromanspacer{} นตฺถิ โข เม มหาราช โกจญฺโญ อตฺตนา ปิยตโร.\csromanspacer{} ตุยฺหํ ปน มหาราช}

\vspace{\tipitakaparskip}
\csromancenter{โกสลสํยุตฺต อตฺถญฺโญ โกจิ อตฺตนา ปิยตโรติ.\csromanspacer{} มยฺหมฺปิ โข มลฺลิเก นตฺถญฺโญ โกจิ อตฺตนา ปิยตโรติ.}
\prose{\csromanfolio{75}อถ โข ราชา ปเสนทิ โกสโล ปาสาทา โอโรหิตฺวา เยน ภควา เตนุปสงฺกมิ, อุปสงฺกมิตฺวา ภควนฺตํ อภิวาเทตฺวา เอกมนฺตํ นิสีทิ, เอกมนฺตํ นิสินฺโน โข ราชา ปเสนทิ โกสโล ภควนฺตํ เอตทโวจ “อิธาหํ ภนฺเต มลฺลิกาย เทวิยา สทฺธิํ อุปริปาสาทวรคโต มลฺลิกํ เทวิํ เอตทโวจํ ‘อตฺถิ นุ โข เต มลฺลิเก โกจญฺโญ อตฺตนา ปิยตโร’ติ.\csromanspacer{} เอวํ วุตฺเต ภนฺเต มลฺลิกา เทวี มํ เอตทโวจ ‘นตฺถิ โข เม มหาราช โกจญฺโญ อตฺตนา ปิยตโร.\csromanspacer{} ตุยฺหํ ปน มหาราช อตฺถญฺโญ โกจิ อตฺตนา ปิยตโร’ติ.\csromanspacer{} เอวํ วุตฺตาหํ ภนฺเต มลฺลิกํ เทวิํ เอตทโวจํ ‘มยฺหมฺปิ โข มลฺลิเก นตฺถญฺโญ โกจิ อตฺตนา ปิยตโร’ติ”.}

\prose{อถ โข ภควา เอตมตฺถํ วิทิตฺวา ตายํ เวลายํ อิมํ คาถํ อภาสิ\csromandash{}}

\csromangathameasure{“สพฺพา ทิสา อนุปริคมฺม เจตสา, \\ เนวชฺฌคา ปิยตรมตฺตนา กฺวจิ. \\ เอวํ ปิโย ปุถุ อตฺตา ปเรสํ, \\ ตสฺมา น หิํเส ปรมตฺตกาโม”ติ.}
\csromangathagroup{\csromangathastanza{“สพฺพา ทิสา อนุปริคมฺม เจตสา, \\ เนวชฺฌคา ปิยตรมตฺตนา กฺวจิ. \\ เอวํ ปิโย ปุถุ อตฺตา ปเรสํ, \\ ตสฺมา น หิํเส ปรมตฺตกาโม”ติ.}}

\csromanheader{2}{9. ยญฺญสุตฺต}
\proseitem{120}{สาวตฺถินิทานํ.\csromanspacer{} เตน โข ปน สมเยน รญฺโญ ปเสนทิสฺส โกสลสฺส มหายญฺโญ ปจฺจุปฏฺฐิโต โหติ, ปญฺจ จ อุสภสตานิ ปญฺจ จ วจฺฉตรสตานิ ปญฺจ จ วจฺฉตริสตานิ ปญฺจ จ อชสตานิ ปญฺจ จ อุรพฺภสตานิ ถูณูปนีตานิ โหนฺติ ยญฺญตฺถาย.\csromanspacer{} เยปิสฺส เต โหนฺติ ทาสาติ วา เปสฺสาติ วา กมฺมกราติ วา, เตปิ ทณฺฑตชฺชิตา ภยตชฺชิตา อสฺสุมุขา รุทมานา ปริกมฺมานิ กโรนฺติ.}

\prose{อถ โข สมฺพหุลา ภิกฺขู ปุพฺพณฺหสมยํ นิวาเสตฺวา ปตฺตจีวรมาทาย สาวตฺถิํ ปิณฺฑาย ปวิสิํสุ, สาวตฺถิยํ ปิณฺฑาย จริตฺวา ปจฺฉาภตฺตํ ปิณฺฑปาตปฏิกฺกนฺตา เยน ภควา เตนุปสงฺกมิํสุ, อุปสงฺกมิตฺวา ภควนฺตํ อภิวาเทตฺวา เอกมนฺตํ นิสีทิํสุ, เอกมนฺตํ นิสินฺนา โข เต ภิกฺขู ภควนฺตํ เอตทโวจุํ “อิธ ภนฺเต รญฺโญ ปเสนทิสฺส โกสลสฺส มหายญฺโญ \csromanfolio{76}ปจฺจุปฏฺฐิโต โหติ, ปญฺจ จ อุสภสตานิ ปญฺจ จ วจฺฉตรสตานิ ปญฺจ จ วจฺฉตริสตานิ ปญฺจ จ อชสตานิ ปญฺจ จ อุรพฺภสตานิ ถูณูปนีตานิ โหนฺติ ยญฺญตฺถาย.\csromanspacer{} เยปิสฺส เต โหนฺติ ทาสาติ วา เปสฺสาติ วา กมฺมกราติ วา, เตปิ ทณฺฑตชฺชิตา ภยตชฺชิตา อสฺสุมุขา รุทมานา ปริกมฺมานิ กโรนฺตี”ติ.}

\prose{อถ โข ภควา เอตมตฺถํ วิทิตฺวา ตายํ เวลายํ อิมา คาถาโย อภาสิ\csromandash{}}

\prose{“อสฺสเมธํ ปุริสเมธํ, สมฺมาปาสํ วาชเปยฺยํ นิรคฺคฬํ.}

\prose{มหายญฺญา มหารมฺภา\footnote{วาชเปยฺยํ. นิรคฺคฬํ มหารมฺภา (ก)}, น เต โหนฺติ มหปฺผลา.}

\prose{อเชฬกา จ คาโว จ, วิวิธา ยตฺถ หญฺญเร.}

\prose{น ตํ สมฺมคฺคตา ยญฺญํ, อุปยนฺติ มเหสิโน.}

\prose{เย จ ยญฺญา นิรารมฺภา, ยชนฺติ อนุกุลํ สทา.}

\prose{อเชฬกา จ คาโว จ, วิวิธา เนตฺถ หญฺญเร.}

\prose{เอตํ สมฺมคฺคตา ยญฺญํ, อุปยนฺติ มเหสิโน.}

\prose{เอตํ ยเชถ เมธาวี, เอโส ยญฺโญ มหปฺผโล.}

\prose{เอตํ หิ ยชมานสฺส, เสยฺโย โหติ น ปาปิโย.}

\prose{ยญฺโญ จ วิปุโล โหติ, ปสีทนฺติ จ เทวตา”ติ.}

\csromanheader{2}{10. พนฺธนสุตฺต}
\proseitem{121}{เตน โข ปน สมเยน รญฺญา ปเสนทินา โกสเลน มหาชนกาโย พนฺธาปิโต โหติ, อปฺเปกจฺเจ รชฺชูหิ อปฺเปกจฺเจ อนฺทูหิ อปฺเปกจฺเจ สงฺขลิกาหิ.}

\prose{อถ โข สมฺพหุลา ภิกฺขู ปุพฺพณฺหสมยํ นิวาเสตฺวา ปตฺตจีวรมาทาย สาวตฺถิํ ปิณฺฑาย ปวิสิํสุ, สาวตฺถิยํ ปิณฺฑาย จริตฺวา ปจฺฉาภตฺตํ ปิณฺฑปาตปฏิกฺกนฺตา เยน ภควา เตนุปสงฺกมิํสุ, อุปสงฺกมิตฺวา ภควนฺตํ อภิวาเทตฺวา เอกมนฺตํ นิสีทิํสุ, เอกมนฺตํ นิสินฺนา โข เต ภิกฺขู ภควนฺตํ เอตทโวจุํ “อิธ ภนฺเต รญฺญา ปเสนทินา โกสเลน มหาชนกาโย พนฺธาปิโต อปฺเปกจฺเจ รชฺชูหิ อปฺเปกจฺเจ อนฺทูหิ อปฺเปกจฺเจ สงฺขลิกาหี”ติ.}

\csromanheader{2}{3. โกสลสํยุตฺต}
\prose{\csromanfolio{77}อถ โข ภควา เอตมตฺถํ วิทิตฺวา ตายํ เวลายํ อิมา คาถาโย อภาสิ\csromandash{}}

\csromangathameasure{“น ตํ ทฬฺหํ พนฺธนมาหุ ธีรา, \\ ยทายสํ ทารุชํ ปพฺพชญฺจ. \\ สารตฺตรตฺตา มณิกุณฺฑเลสุ, \\ ปุตฺเตสุ ทาเรสุ จ ยา อเปกฺขา. \\ เอตํ ทฬฺหํ พนฺธนมาหุ ธีรา, \\ โอหารินํ สิถิลํ ทุปฺปมุญฺจํ. \\ เอตมฺปิ เฉตฺวาน ปริพฺพชนฺติ, \\ อนเปกฺขิโน กามสุขํ ปหายา”ติ.}
\csromangathagroup{\csromangathastanza{“น ตํ ทฬฺหํ พนฺธนมาหุ ธีรา, \\ ยทายสํ ทารุชํ ปพฺพชญฺจ. \\ สารตฺตรตฺตา มณิกุณฺฑเลสุ, \\ ปุตฺเตสุ ทาเรสุ จ ยา อเปกฺขา.}\csromangathastanza{เอตํ ทฬฺหํ พนฺธนมาหุ ธีรา, \\ โอหารินํ สิถิลํ ทุปฺปมุญฺจํ. \\ เอตมฺปิ เฉตฺวาน ปริพฺพชนฺติ, \\ อนเปกฺขิโน กามสุขํ ปหายา”ติ.}}

\vspace{\tipitakaparskip}
\csromancenter{ปฐโม วคฺโค.}
\titlehead{\textbf{ตสฺสุทฺทานํ}}
\csromangathameasure{ทหโร ปุริโส ชรา, ปิยํ อตฺตานรกฺขิโต. \\ อปฺปกา อฑฺฑกรณํ, มลฺลิกา ยญฺญพนฺธนนฺติ.}
\csromangathagroup{\csromangathastanza{ทหโร ปุริโส ชรา, ปิยํ อตฺตานรกฺขิโต. \\ อปฺปกา อฑฺฑกรณํ, มลฺลิกา ยญฺญพนฺธนนฺติ.}}

\chapterhead{2. ทุติยวคฺค}
\csromanheader{2}{1. สตฺตชฏิลสุตฺต}
\proseitem{122}{เอกํ สมยํ ภควา สาวตฺถิยํ วิหรติ ปุพฺพาราเม มิคารมาตุปาสาเท.\csromanspacer{} เตน โข ปน สมเยน ภควา สายนฺหสมยํ ปฏิสลฺลานา วุฏฺฐิโต พหิทฺวารโกฏฺฐเก นิสินฺโน โหติ.\csromanspacer{} อถ โข ราชา ปเสนทิ โกสโล เยน ภควา เตนุปสงฺกมิ, อุปสงฺกมิตฺวา ภควนฺตํ อภิวาเทตฺวา เอกมนฺตํ นิสีทิ.}

\prose{เตน โข ปน สมเยน สตฺต จ ชฏิลา สตฺต จ นิคณฺฐา สตฺต จ อเจลกา สตฺต จ เอกสาฏกา สตฺต จ ปริพฺพาชกา ปรูฬฺหกจฺฉนขโลมา ขาริวิวิธมาทาย\footnote{ขาริวิธํ อาทาย (อิ) ที 1. 95 ปิฏฺเฐปิ. ตทฏฺฐกถาปิ โอโลเกตพฺพา.} ภควโต อวิทูเร อติกฺกมนฺติ.\csromanspacer{} อถ โข ราชา ปเสนทิ โกสโล อุฏฺฐายาสนา เอกํสํ อุตฺตราสงฺคํ กริตฺวา \csromanfolio{78}ทกฺขิณชาณุมณฺฑลํ ปถวิยํ นิหนฺตฺวา เยน เต สตฺต จ ชฏิลา สตฺต จ นิคณฺฐา สตฺต จ อเจลกา สตฺต จ เอกสาฏกา สตฺต จ ปริพฺพาชกา เตนญฺชลิํ ปณาเมตฺวา ติกฺขตฺตุํ นามํ สาเวสิ “ราชาหํ ภนฺเต ปเสนทิ โกสโล ฯเปฯ ราชาหํ ภนฺเต ปเสนทิ โกสโล”ติ.}

\prose{อถ โข ราชา ปเสนทิ โกสโล อจิรปกฺกนฺเตสุ เตสุ สตฺตสุ จ ชฏิเลสุ สตฺตสุ จ นิคณฺเฐสุ สตฺตสุ จ อเจลเกสุ สตฺตสุ จ เอกสาฏเกสุ สตฺตสุ จ ปริพฺพาชเกสุ เยน ภควา เตนุปสงฺกมิ, อุปสงฺกมิตฺวา ภควนฺตํ อภิวาเทตฺวา เอกมนฺตํ นิสีทิ, เอกมนฺตํ นิสินฺโน โข ราชา ปเสนทิ โกสโล ภควนฺตํ เอตทโวจ “เย เต ภนฺเต โลเก อรหนฺโต วา อรหตฺตมคฺคํ วา สมาปนฺนา, เอเต เตสํ อญฺญตรา”ติ.\csromanspacer{} ทุชฺชานํ โข เอตํ มหาราช ตยา คิหินา กามโภคินา ปุตฺตสมฺพาธสยนํ อชฺฌาวสนฺเตน กาสิกจนฺทนํ ปจฺจนุโภนฺเตน มาลาคนฺธวิเลปนํ ธารยนฺเตน ชาตรูปรชตํ สาทิยนฺเตน “อิเม วา อรหนฺโต, อิเม วา อรหตฺตมคฺคํ สมาปนฺนา”ติ.}

\prose{สํวาเสน โข มหาราช สีลํ เวทิตพฺพํ, ตญฺจ โข ทีเฆน อทฺธุนา น อิตฺตรํ, มนสิกโรตา โน อมนสิกโรตา, ปญฺญวตา โน ทุปฺปญฺเญน.\csromanspacer{} สํโวหาเรน โข มหาราช โสเจยฺยํ เวทิตพฺพํ, ตญฺจ โข ทีเฆน อทฺธุนา น อิตฺตรํ, มนสิกโรตา โน อมนสิกโรตา, ปญฺญวตา โน ทุปฺปญฺเญน.\csromanspacer{} อาปทาสุ โข มหาราช ถาโม เวทิตพฺโพ, โส จ โข ทีเฆน อทฺธุนา น อิตฺตรํ, มนสิกโรตา โน อมนสิกโรตา, ปญฺญวตา โน ทุปฺปญฺเญน.\csromanspacer{} สากจฺฉาย โข มหาราช ปญฺญา เวทิตพฺพา, สา จ โข ทีเฆน อทฺธุนา น อิตฺตรํ, มนสิกโรตา โน อมนสิกโรตา, ปญฺญวตา โน ทุปฺปญฺเญนาติ.}

\prose{อจฺฉริยํ ภนฺเต อพฺภุตํ ภนฺเต, ยาว สุภาสิตมิทํ ภนฺเต ภควตา\csromandash{}“ทุชฺชานํ โข เอตํ มหาราช ตยา คิหินา กามโภคินา ปุตฺตสมฺพาธสยนํ อชฺฌาวสนฺเตน กาสิกจนฺทนํ ปจฺจนุโภนฺเตน มาลาคนฺธวิเลปนํ ธารยนฺเตน ชาตรูปรชตํ สาทิยนฺเตน ‘อิเม วา อรหนฺโต, อิเม วา อรหตฺตมคฺคํ สมาปนฺนา’ติ.\csromanspacer{} สํวาเสน โข มหาราช สีลํ เวทิตพฺพํ, ตญฺจ โข ทีเฆน อทฺธุนา น อิตฺตรํ, มนสิกโรตา โน อมนสิกโรตา, ปญฺญวตา โน ทุปฺปญฺเญน.\csromanspacer{} สํโวหาเรน โข}

\vspace{\tipitakaparskip}
\csromancenter{โกสลสํยุตฺต มหาราช โสเจยฺยํ เวทิตพฺพํ, ตญฺจ โข ทีเฆน อทฺธุนา น อิตฺตรํ, มนสิกโรตา โน อมนสิกโรตา, ปญฺญวตา โน ทุปฺปญฺเญน.\csromanspacer{} อาปทาสุ โข มหาราช ถาโม เวทิตพฺโพ, โส จ โข ทีเฆน อทฺธุนา น อิตฺตรํ, มนสิกโรตา โน อมนสิกโรตา, ปญฺญวตา โน ทุปฺปญฺเญน.\csromanspacer{} สากจฺฉาย โข มหาราช ปญฺญา เวทิตพฺพา, สา จ โข ทีเฆน อทฺธุนา น อิตฺตรํ, มนสิกโรตา โน อมนสิกโรตา, ปญฺญวตา โน ทุปฺปญฺเญนา”ติ.}
\prose{\csromanfolio{79}เอเต ภนฺเต มม ปุริสา จรา โอจรกา ชนปทํ โอจริตฺวา อาคจฺฉนฺติ, เตหิ ปฐมํ โอจิณฺณํ อหํ ปจฺฉา โอสาปยิสฺสามิ\footnote{โอยายิสฺสามิ (สี), โอหยิสฺสามิ (สฺยา, กํ)}.\csromanspacer{} อิทานิ เต ภนฺเต ตํ รโชชลฺลํ ปวาเหตฺวา สุนฺหาตา สุวิลิตฺตา กปฺปิตเกสมสฺสู โอทาตวตฺถา\footnote{โอทาตวตฺถวสนา (สี)} ปญฺจหิ กามคุเณหิ สมปฺปิตา สมงฺคีภูตา ปริจาเรสฺสนฺตีติ.}

\prose{อถ โข ภควา เอตมตฺถํ วิทิตฺวา ตายํ เวลายํ อิมา คาถาโย อภาสิ\csromandash{}}

\csromangathameasure{“น วณฺณรูเปน นโร สุชาโน, \\ น วิสฺสเส อิตฺตรทสฺสเนน. \\ สุสญฺญตานํ หิ วิยญฺชเนน, \\ อสญฺญตา โลกมิมํ จรนฺติ. \\ ปติรูปโก มตฺติกากุณฺฑโลว, \\ โลหฑฺฒมาโสว สุวณฺณฉนฺโน. \\ จรนฺติ โลเก ปริวารฉนฺนา, \\ อนฺโต อสุทฺธา พหิ โสภมานา”ติ.}
\csromangathagroup{\csromangathastanza{“น วณฺณรูเปน นโร สุชาโน, \\ น วิสฺสเส อิตฺตรทสฺสเนน. \\ สุสญฺญตานํ หิ วิยญฺชเนน, \\ อสญฺญตา โลกมิมํ จรนฺติ.}\csromangathastanza{ปติรูปโก มตฺติกากุณฺฑโลว, \\ โลหฑฺฒมาโสว สุวณฺณฉนฺโน. \\ จรนฺติ โลเก\footnote{เอเก (สี, อิ)} ปริวารฉนฺนา, \\ อนฺโต อสุทฺธา พหิ โสภมานา”ติ.}}

\csromanheader{2}{2. ปญฺจราชสุตฺต}
\proseitem{123}{สาวตฺถินิทานํ.\csromanspacer{} เตน โข ปน สมเยน ปญฺจนฺนํ ราชูนํ ปเสนทิปมุขานํ ปญฺจหิ กามคุเณหิ สมปฺปิตานํ สมงฺคีภูตา- นํ ปริจารยมานานํ อยมนฺตรากถา อุทปาทิ “กิํ นุ โข กามานํ อคฺคนฺ”ติ.\csromanspacer{} ตเตฺรกจฺเจ\footnote{ตเตฺรเก (สี, อิ)} เอวมาหํสุ “รูปา กามานํ อคฺคนฺ”ติ.\csromanspacer{} เอกจฺเจ \csromanfolio{80}เอวมาหํสุ “สทฺทา กามานํ อคฺคนฺ”ติ.\csromanspacer{} เอกจฺเจ เอวมาหํสุ “คนฺธา กามานํ อคฺคนฺ”ติ.\csromanspacer{} เอกจฺเจ เอวมาหํสุ “รสา กามานํ อคฺคนฺ”ติ.\csromanspacer{} เอกจฺเจ เอวมาหํสุ “โผฏฺฐพฺพา กามานํ อคฺคนฺ”ติ.\csromanspacer{} ยโต โข เต ราชาโน นาสกฺขิํสุ อญฺญมญฺญํ สญฺญาเปตุํ.}

\prose{อถ โข ราชา ปเสนทิ โกสโล เต ราชาโน เอตทโวจ “อายาม มาริสา เยน ภควา เตนุปสงฺกมิสฺสาม, อุปสงฺกมิตฺวา ภควนฺตํ เอตมตฺถํ ปฏิปุจฺฉิสฺสาม, ยถา โน ภควา พฺยากริสฺสติ, ตถา นํ ธาเรสฺสามา”ติ\footnote{ธาเรยฺยามาติ (สี, สฺยา, กํ, อิ)}. “เอวํ มาริสา”ติ โข เต ราชาโน รญฺโญ ปเสนทิสฺส โกสลสฺส ปจฺจสฺโสสุํ.}

\prose{อถ โข เต ปญฺจ ราชาโน ปเสนทิปมุขา เยน ภควา เตนุปสงฺกมิํสุ, อุปสงฺกมิตฺวา ภควนฺตํ อภิวาเทตฺวา เอกมนฺตํ นิสีทิํสุ, เอกมนฺตํ นิสินฺโน โข ราชา ปเสนทิ โกสโล ภควนฺตํ เอตทโวจ “อิธ ภนฺเต อมฺหากํ ปญฺจนฺนํ ราชูนํ ปญฺจหิ กามคุเณหิ สมปฺปิตานํ สมงฺคีภูตานํ ปริจารยมานานํ อยมนฺตรากถา อุทปาทิ ‘กิํ นุ โข กามานํ อคฺคนฺ’ติ.\csromanspacer{} เอกจฺเจ เอวมาหํสุ ‘รูปา กามานํ อคฺคนฺ’ติ.\csromanspacer{} เอกจฺเจ เอวมาหํสุ ‘สทฺทา กามานํ อคฺคนฺ’ติ.\csromanspacer{} เอกจฺเจ เอวมาหํสุ ‘คนฺธา กามานํ อคฺคนฺ’ติ.\csromanspacer{} เอกจฺเจ เอวมาหํสุ ‘รสา กามานํ อคฺคนฺ’ติ.\csromanspacer{} เอกจฺเจ เอวมาหํสุ ‘โผฏฺฐพฺพา กามานํ อคฺคนฺ’ติ.\csromanspacer{} กิํ นุ โข ภนฺเต กามานํ อคฺคนฺ’ติ.}

\prose{“มนาปปริยนฺตํ ขฺวาหํ มหาราช ปญฺจสุ กามคุเณสุ อคฺคนฺ”ติ วทามิ, เตว\footnote{เต จ (สี, อิ, ก), เย จ (สฺยา, กํ)} มหาราช รูปา เอกจฺจสฺส มนาปา โหนฺติ, เตว\footnote{เต จ (สี, อิ, ก)} รูปา เอกจฺจสฺส อมนาปา โหนฺติ.\csromanspacer{} เยหิ จ โย รูเปหิ อตฺตมโน โหติ ปริปุณฺณสงฺกปฺโป, โส เตหิ รูเปหิ อญฺญํ รูปํ อุตฺตริตรํ วา ปณีตตรํ วา น ปตฺเถติ.\csromanspacer{} เต ตสฺส รูปา ปรมา โหนฺติ, เต ตสฺส รูปา อนุตฺตรา โหนฺติ.}

\prose{เตว มหาราช สทฺทา เอกจฺจสฺส มนาปา โหนฺติ, เตว สทฺทา เอกจฺจสฺส อมนาปา โหนฺติ.\csromanspacer{} เยหิ จ โย สทฺเทหิ อตฺตมโน โหติ ปริปุณฺณสงฺกปฺโป, โส เตหิ สทฺเทหิ อญฺญํ สทฺทํ อุตฺตริตรํ}

\vspace{\tipitakaparskip}
\csromancenter{โกสลสํยุตฺต วา ปณีตตรํ วา น ปตฺเถติ.\csromanspacer{} เต ตสฺส สทฺทา ปรมา โหนฺติ, เต ตสฺส สทฺทา อนุตฺตรา โหนฺติ.}
\prose{\csromanfolio{81}เตว มหาราช คนฺธา เอกจฺจสฺส มนาปา โหนฺติ, เตว คนฺธา เอกจฺจสฺส อมนาปา โหนฺติ.\csromanspacer{} เยหิ จ โย คนฺเธหิ อตฺตมโน โหติ ปริปุณฺณสงฺกปฺโป, โส เตหิ คนฺเธหิ อญฺญํ คนฺธํ อุตฺตริตรํ วา ปณีตตรํ วา น ปตฺเถติ.\csromanspacer{} เต ตสฺส คนฺธา ปรมา โหนฺติ, เต ตสฺส คนฺธา อนุตฺตรา โหนฺติ.}

\prose{เตว มหาราช รสา เอกจฺจสฺส มนาปา โหนฺติ, เตว รสา เอกจฺจสฺส อมนาปา โหนฺติ.\csromanspacer{} เยหิ จ โย รเสหิ อตฺตมโน โหติ ปริปุณฺณสงฺกปฺโป, โส เตหิ รเสหิ อญฺญํ รสํ อุตฺตริตรํ วา ปณีตตรํ วา น ปตฺเถติ.\csromanspacer{} เต ตสฺส รสา ปรมา โหนฺติ, เต ตสฺส รสา อนุตฺตรา โหนฺติ.}

\prose{เตว มหาราช โผฏฺฐพฺพา เอกจฺจสฺส มนาปา โหนฺติ, เตว โผฏฺฐพฺพา เอกจฺจสฺส อมนาปา โหนฺติ.\csromanspacer{} เยหิ จ โย โผฏฺฐพฺเพหิ อตฺตมโน โหติ ปริปุณฺณสงฺกปฺโป, โส เตหิ โผฏฺฐพฺเพหิ อญฺญํ โผฏฺฐพฺพํ อุตฺตริตรํ วา ปณีตตรํ วา น ปตฺเถติ.\csromanspacer{} เต ตสฺส โผฏฺฐพฺพา ปรมา โหนฺติ, เต ตสฺส โผฏฺฐพฺพา อนุตฺตรา โหนฺตีติ.}

\prose{เตน โข ปน สมเยน จนฺทนงฺคลิโก อุปาสโก ตสฺสํ ปริสายํ นิสินฺโน โหติ.\csromanspacer{} อถ โข จนฺทนงฺคลิโก อุปาสโก อุฏฺฐายาสนา เอกํสํ อุตฺตราสงฺคํ กริตฺวา เยน ภควา เตนญฺชลิํ ปณาเมตฺวา ภควนฺตํ เอตทโวจ “ปฏิภาติ มํ ภควา, ปฏิภาติ มํ สุคตา”ติ. “ปฏิภาตุ ตํ จนฺทนงฺคลิกา”ติ ภควา อโวจ. อถ โข จนฺทนงฺคลิโก อุปาสโก ภควโต สมฺมุขา ตทนุรูปาย คาถาย อภิตฺถวิ\csromandash{}}

\csromangathameasure{“ปทุมํ ยถา โกกนทํ สุคนฺธํ, \\ ปาโต สิยา ผุลฺลมวีตคนฺธํ. \\ องฺคีรสํ ปสฺส วิโรจมานํ, \\ ตปนฺตมาทิจฺจมิวนฺตลิกฺเข”ติ.}
\csromangathagroup{\csromangathastanza{“ปทุมํ ยถา โกกนทํ สุคนฺธํ, \\ ปาโต สิยา ผุลฺลมวีตคนฺธํ. \\ องฺคีรสํ ปสฺส วิโรจมานํ, \\ ตปนฺตมาทิจฺจมิวนฺตลิกฺเข”ติ.}}

\prose{\csromanfolio{82}อถ โข เต ปญฺจ ราชาโน จนฺทนงฺคลิกํ อุปาสกํ ปญฺจหิ อุตฺตราสงฺเคหิ อจฺฉาเทสุํ.\csromanspacer{} อถ โข จนฺทนงฺคลิโก อุปาสโก เตหิ ปญฺจหิ อุตฺตราสงฺเคหิ ภควนฺตํ อจฺฉาเทสีติ.}

\csromanheader{2}{3. โทณปากสุตฺต}
\proseitem{124}{สาวตฺถินิทานํ.\csromanspacer{} เตน โข ปน สมเยน ราชา ปเสนทิ โกสโล โทณปากกุรํ\footnote{โทณปากสุทํ (สี), โทณปากํ สุทํ (อิ)} ภุญฺชติ.\csromanspacer{} อถ โข ราชา ปเสนทิ โกสโล ภุตฺตาวี มหสฺสาสี เยน ภควา เตนุปสงฺกมิ, อุปสงฺกมิตฺวา ภควนฺตํ อภิวาเทตฺวา เอกมนฺตํ นิสีทิ.}

\prose{อถ โข ภควา ราชานํ ปเสนทิํ โกสลํ ภุตฺตาวิํ มหสฺสาสิํ วิทิตฺวา ตายํ เวลายํ อิมํ คาถํ อภาสิ\csromandash{}}

\csromangathameasure{“มนุชสฺส สทา สตีมโต, \\ มตฺตํ ชานโต ลทฺธโภชเน. \\ ตนุกสฺส ภวนฺติ เวทนา, \\ สณิกํ ชีรติ อายุปาลยนฺ”ติ.}
\csromangathagroup{\csromangathastanza{“มนุชสฺส สทา สตีมโต, \\ มตฺตํ ชานโต ลทฺธโภชเน. \\ ตนุกสฺส\footnote{ตนุ ตสฺส (สี, อิ) 3. ( ) สี-สฺยา-กํ-อิ-โปตฺถเกสุ นตฺถิ.} ภวนฺติ เวทนา, \\ สณิกํ ชีรติ อายุปาลยนฺ”ติ.}}

\prose{เตน โข ปน สมเยน สุทสฺสโน มาณโว รญฺโญ ปเสนทิสฺส โกสลสฺส ปิฏฺฐิโต ฐิโต โหติ.\csromanspacer{} อถ โข ราชา ปเสนทิ โกสโล สุทสฺสนํ มาณวํ อามนฺเตสิ “เอหิ ตฺวํ ตาต สุทสฺสน ภควโต สนฺติเก อิมํ คาถํ ปริยาปุณิตฺวา มม ภตฺตาภิหาเร (ภตฺตาภิหาเร)3 ภาส, อหญฺจ เต เทวสิกํ กหาปณสตํ (กหาปณสตํ)\footnote{( ) สี-สฺยา-กํ-โปตฺถเกสุ นตฺถิ.} นิจฺจํ ภิกฺขํ ปวตฺตยิสฺสามี”ติ. “เอวํ เทวา”ติ โข สุทสฺสโน มาณโว รญฺโญ ปเสนทิสฺส โกสลสฺส ปฏิสฺสุตฺวา ภควโต สนฺติเก อิมํ คาถํ ปริยาปุณิตฺวา รญฺโญ ปเสนทิสฺส โกสลสฺส ภตฺตาภิหาเร สุทํ ภาสติ\csromandash{}}

\csromangathameasure{“มนุชสฺส สทา สตีมโต, \\ มตฺตํ ชานโต ลทฺธโภชเน. \\ ตนุกสฺส ภวนฺติ เวทนา, \\ สณิกํ ชีรติ อายุปาลยนฺ”ติ.}
\csromangathagroup{\csromangathastanza{“มนุชสฺส สทา สตีมโต, \\ มตฺตํ ชานโต ลทฺธโภชเน. \\ ตนุกสฺส ภวนฺติ เวทนา, \\ สณิกํ ชีรติ อายุปาลยนฺ”ติ.}}

\csromanheader{2}{3. โกสลสํยุตฺต}
\prose{\csromanfolio{83}อถ โข ราชา ปเสนทิ โกสโล อนุปุพฺเพน นาฬิโกทนปรมตาย\footnote{นาฬิโกทนมตฺตาย (ก)} สณฺฐาสิ.\csromanspacer{} อถ โข ราชา ปเสนทิ โกสโล อปเรน สมเยน สุสลฺลิขิตคตฺโต ปาณินา คตฺตานิ อนุมชฺชนฺโต ตายํ เวลายํ อิมํ อุทานํ อุทาเนสิ “อุภเยน วต มํ โส ภควา อตฺเถน อนุกมฺปิ ทิฏฺฐธมฺมิเกน เจว อตฺเถน สมฺปรายิเกน จา”ติ.}

\csromanheader{2}{4. ปฐมสงฺคามสุตฺต}
\proseitem{125}{สาวตฺถินิทานํ.\csromanspacer{} อถ โข ราชา มาคโธ อชาตสตฺตุ เวเทหิปุตฺโต จตุรงฺคินิํ เสนํ สนฺนยฺหิตฺวา ราชานํ ปเสนทิํ โกสลํ อพฺภุยฺยาสิ เยน กาสิ.\csromanspacer{} อสฺโสสิ โข ราชา ปเสนทิ โกสโล “ราชา กิร มาคโธ อชาตสตฺตุ เวเทหิปุตฺโต จตุรงฺคินิํ เสนํ สนฺนยฺหิตฺวา มมํ อพฺภุยฺยาโต เยน กาสี”ติ.\csromanspacer{} อถ โข ราชา ปเสนทิ โกสโล จตุรงฺคินิํ เสนํ สนฺนยฺหิตฺวา ราชานํ มาคธํ อชาตสตฺตุํ เวเทหิปุตฺตํ ปจฺจุยฺยาสิ เยน กาสิ.\csromanspacer{} อถ โข ราชา จ มาคโธ อชาตสตฺตุ เวเทหิปุตฺโต ราชา จ ปเสนทิ โกสโล สงฺคาเมสุํ.\csromanspacer{} เตสฺมิํ โข ปน สงฺคาเม ราชา มาคโธ อชาตสตฺตุ เวเทหิปุตฺโต ราชานํ ปเสนทิํ โกสลํ ปราเชสิ, ปราชิโต จ ราชา ปเสนทิ โกสโล สกเมว\footnote{สงฺคามา (ก)} ราชธานิํ สาวตฺถิํ ปจฺจุยฺยาสิ\footnote{ปายาสิ (สี, อิ)}.}

\prose{อถ โข สมฺพหุลา ภิกฺขู ปุพฺพณฺหสมยํ นิวาเสตฺวา ปตฺตจีวรมาทาย สาวตฺถิํ ปิณฺฑาย ปวิสิํสุ, สาวตฺถิยํ ปิณฺฑาย จริตฺวา ปจฺฉาภตฺตํ ปิณฺฑปาตปฏิกฺกนฺตา เยน ภควา เตนุปสงฺกมิํสุ, อุปสงฺกมิตฺวา ภควนฺตํ อภิวาเทตฺวา เอกมนฺตํ นิสีทิํสุ, เอกมนฺตํ นิสินฺนา โข เต ภิกฺขู ภควนฺตํ เอตทโวจุํ\csromandash{}}

\prose{อิธ ภนฺเต ราชา มาคโธ อชาตสตฺตุ เวเทหิปุตฺโต จตุรงฺคินิํ เสนํ สนฺนยฺหิตฺวา ราชานํ ปเสนทิํ โกสลํ อพฺภุยฺยาสิ เยน กาสิ.\csromanspacer{} อสฺโสสิ โข ภนฺเต ราชา ปเสนทิ โกสโล “ราชา กิร มาคโธ อชาตสตฺตุ เวเทหิปุตฺโต จตุรงฺคินิํ เสนํ สนฺนยฺหิตฺวา มมํ อพฺภุยฺยาโต เยน กาสี”ติ.\csromanspacer{} อถ โข ภนฺเต ราชา ปเสนทิ \csromanfolio{84}โกสโล จตุรงฺคินิํ เสนํ สนฺนยฺหิตฺวา ราชานํ มาคธํ อชาตสตฺตุํ เวเทหิปุตฺตํ ปจฺจุยฺยาสิ เยน กาสิ.\csromanspacer{} อถ โข ภนฺเต ราชา จ มาคโธ อชาตสตฺตุ เวเทหิปุตฺโต ราชา จ ปเสนทิ โกสโล สงฺคาเมสุํ.\csromanspacer{} ตสฺมิํ โข ปน ภนฺเต สงฺคาเม ราชา มาคโธ อชาตสตฺตุ เวเทหิปุตฺโต ราชานํ ปเสนทิํ โกสลํ ปราเชสิ, ปราชิโต จ ภนฺเต ราชา ปเสนทิ โกสโล สกเมว ราชธานิํ สาวตฺถิํ ปจฺจุยฺยาสีติ.}

\prose{ราชา ภิกฺขเว มาคโธ อชาตสตฺตุ เวเทหิปุตฺโต ปาปมิตฺโต ปาปสหาโย ปาปสมฺปวงฺโก, ราชา จ โข ภิกฺขเว ปเสนทิ โกสโล กลฺยาณมิตฺโต กลฺยาณสหาโย กลฺยาณสมฺปวงฺโก, อชฺเชว\footnote{อชฺชตญฺจ (สี, อิ), อชฺเชวํ (สฺยา, กํ)} ภิกฺขเว ราชา ปเสนทิ โกสโล อิมํ รตฺติํ ทุกฺขํ เสติ ปราชิโตติ.\csromanspacer{} อิทมโวจ ฯเปฯ.}

\csromangathameasure{ชยํ เวรํ ปสวติ, ทุกฺขํ เสติ ปราชิโต. \\ อุปสนฺโต สุขํ เสติ, หิตฺวา ชยปราชยนฺติ.}
\csromangathagroup{\csromangathastanza{ชยํ เวรํ ปสวติ, ทุกฺขํ เสติ ปราชิโต. \\ อุปสนฺโต สุขํ เสติ, หิตฺวา ชยปราชยนฺติ.}}

\csromanheader{2}{5. ทุติยสงฺคามสุตฺต}
\proseitem{126}{2อถ โข ราชา มาคโธ อชาตสตฺตุ เวเทหิปุตฺโต จตุรงฺคินิํ เสนํ สนฺนยฺหิตฺวา ราชานํ ปเสนทิํ โกสลํ อพฺภุยฺยาสิ เยน กาสิ.\csromanspacer{} อสฺโสสิ โข ราชา ปเสนทิ โกสโล “ราชา กิร มาคโธ อชาตสตฺตุ เวเทหิปุตฺโต จตุรงฺคินิํ เสนํ สนฺนยฺหิตฺวา มมํ อพฺภุยฺยาโต เยน กาสี”ติ.\csromanspacer{} อถ โข ราชา ปเสนทิ โกสโล จตุรงฺคินิํ เสนํ สนฺนยฺหิตฺวา ราชานํ มาคธํ อชาตสตฺตุํ เวเทหิปุตฺตํ ปจฺจุยฺยาสิ เยน กาสิ.\csromanspacer{} อถ โข ราชา จ มาคโธ อชาตสตฺตุ เวเทหิปุตฺโต ราชา จ ปเสนทิ โกสโล สงฺคาเมสุํ.\csromanspacer{} ตสฺมิํ โข ปน สงฺคเม ราชา ปเสนทิ โกสโล ราชานํ มาคธํ อชาตสตฺตุํ เวเทหิปุตฺตํ ปราเชสิ, ชีวคฺคาหญฺจ นํ อคฺคเหสิ.\csromanspacer{} อถ โข รญฺโญ ปเสนทิสฺส โกสลสฺส เอตทโหสิ “กิญฺจาปิ โข}

\vspace{\tipitakaparskip}
\csromancenter{โกสลสํยุตฺต มฺยายํ ราชา มาคโธ อชาตสตฺตุ เวเทหิปุตฺโต อทุพฺภนฺตสฺส ทุพฺภติ, อถ จ ปน เม ภาคิเนยฺโย โหติ.\csromanspacer{} ยํนูนาหํ รญฺโญ มาคธสฺส อชาตสตฺตุโน เวเทหิปุตฺตสฺส สพฺพํ หตฺถิกายํ ปริยาทิยิตฺวา สพฺพํ อสฺสกายํ ปริยาทิยิตฺวา สพฺพํ รถกายํ ปริยาทิยิตฺวา สพฺพํ ปตฺติกายํ ปริยาทิยิตฺวา ชีวนฺตเมว นํ โอสชฺเชยฺยนฺ”ติ\footnote{โอสฺสชฺเชยฺยนฺติ (สี, สฺยา, กํ, อิ)}.}
\prose{\csromanfolio{85}อถ โข ราชา ปเสนทิ โกสโล รญฺโญ มาคธสฺส อชาตสตฺตุโน เวเทหิปุตฺตสฺส สพฺพํ หตฺถิกายํ ปริยาทิยิตฺวา สพฺพํ อสฺสกายํ ปริยาทิยิตฺวา สพฺพํ รถกายํ ปริยาทิยิตฺวา สพฺพํ ปตฺติกายํ ปริยาทิยิตฺวา ชีวนฺตเมว นํ โอสชฺชิ\footnote{โอสฺสชิ (สี), โอสฺสชฺชิ (สฺยา, กํ, อิ)}.}

\prose{อถ โข สมฺพหุลา ภิกฺขู ปุพฺพณฺหสมยํ นิวาเสตฺวา ปตฺตจีวรมาทาย สาวตฺถิํ ปิณฺฑาย ปวิสิํสุ, สาวตฺถิยํ ปิณฺฑาย จริตฺวา ปจฺฉาภตฺตํ ปิณฺฑปาตปฏิกฺกนฺตา เยน ภควา เตนุปสงฺกมิํสุ, อุปสงฺกมิตฺวา ภควนฺตํ อภิวาเทตฺวา เอกมนฺตํ นิสีทิํสุ, เอกมนฺตํ นิสินฺนา โข เต ภิกฺขู ภควนฺตํ เอตทโวจุํ\csromandash{}}

\prose{อิธ ภนฺเต ราชา มาคโธ อชาตสตฺตุ เวเทหิปุตฺโต จตุรงฺคินิํ เสนํ สนฺนยฺหิตฺวา ราชานํ ปเสนทิํ โกสลํ อพฺภุยฺยาสิ เยน กาสิ.\csromanspacer{} อสฺโสสิ โข ภนฺเต ราชา ปเสนทิ โกสโล “ราชา กิร มาคโธ อชาตสตฺตุ เวเทหิปุตฺโต จตุรงฺคินิํ เสนํ สนฺนยฺหิตฺวา มมํ อพฺภุยฺยาโต เยน กาสี”ติ.\csromanspacer{} อถ โข ภนฺเต ราชา ปเสนทิ โกสโล จตุรงฺคินิํ เสนํ สนฺนยฺหิตฺวา ราชานํ มาคธํ อชาตสตฺตุํ เวเทหิปุตฺตํ ปจฺจุยฺยาสิ เยน กาสิ.\csromanspacer{} อถ โข ภนฺเต ราชา จ มาคโธ อชาตสตฺตุ เวเทหิปุตฺโต ราชาจ ปเสนทิ โกสโล สงฺคาเมสุํ.\csromanspacer{} ตสฺมิํ โข ปน ภนฺเต สงฺคาเม ราชา ปเสนทิ โกสโล ราชานํ มาคธํ อชาตสตฺตุํ เวเทหิปุตฺตํ ปราเชสิ, ชีวคฺคาหญฺจ นํ อคฺคเหสิ.\csromanspacer{} อถ โข ภนฺเต รญฺโญ ปเสนทิสฺส โกสลสฺส เอตทโหสิ “กิญฺจาปิ โข มฺยายํ ราชา มาคโธ อชาตสตฺตุ เวเทหิปุตฺโต อทุพฺภนฺตสฺส ทุพฺภติ, อถ จ ปน เม ภาคิเนยฺโย โหติ.\csromanspacer{} ยํนูนาหํ รญฺโญ มาคธสฺส อชาตสตฺตุโน เวเทหิปุตฺตสฺส สพฺพํ หตฺถิกายํ ปริยาทิยิตฺวา สพฺพํ อสฺสกายํ สพฺพํ รถกายํ สพฺพํ ปตฺติกายํ ปริยาทิยิตฺวา ชีวนฺตเมว นํ โอสชฺเชยฺยนฺ”ติ.}

\prose{\csromanfolio{86}อถ โข ภนฺเต ราชา ปเสนทิ โกสโล รญฺโญ มาคธสฺส อชาตสตฺตุโน เวเทหิปุตฺตสฺส สพฺพํ หตฺถิกายํ ปริยาทิยิตฺวา สพฺพํ อสฺสกายํ ปริยาทิยิตฺวา สพฺพํ รถกายํ ปริยาทิยิตฺวา สพฺพํ ปตฺติกายํ ปริยาทิยิตฺวา ชีวนฺตเมว นํ โอสชฺชีติ.}

\prose{อถ โข ภควา เอตมตฺถํ วิทิตฺวา ตายํ เวลายํ อิมา คาถาโย อภาสิ\csromandash{}}

\csromangathameasure{วิลุมฺปเตว ปุริโส, ยาวสฺส อุปกปฺปติ. \\ ยทา จญฺเญ วิลุมฺปนฺติ, โส วิลุตฺโต วิลุปฺปติ. \\ ฐานํ หิ มญฺญติ พาโล, ยาว ปาปํ น ปจฺจติ, \\ ยทา จ ปจฺจติ ปาปํ, \\ อถ ทุกฺขํ นิคจฺฉติ. \\ หนฺตา ลภติ หนฺตารํ, เชตารํ ลภเต ชยํ, \\ อกฺโกสโก จ อกฺโกสํ, โรเสตารญฺจ โรสโก. \\ อถ กมฺมวิวฏฺเฏน, โส วิลุตฺโต วิลุปฺปตีติ.}
\csromangathagroup{\csromangathastanza{วิลุมฺปเตว ปุริโส, ยาวสฺส อุปกปฺปติ. \\ ยทา จญฺเญ วิลุมฺปนฺติ, โส วิลุตฺโต วิลุปฺปติ\footnote{วิลุมฺปติ (สี, อิ, ก)}.}\csromangathastanza{ฐานํ หิ มญฺญติ พาโล, ยาว ปาปํ น ปจฺจติ, \\ ยทา จ ปจฺจติ ปาปํ, \\ อถ ทุกฺขํ นิคจฺฉติ. \\ หนฺตา ลภติ\footnote{ลภติ หนฺตา (สี, สฺยา, กํ)} หนฺตารํ, เชตารํ ลภเต ชยํ, \\ อกฺโกสโก จ อกฺโกสํ, โรเสตารญฺจ โรสโก. \\ อถ กมฺมวิวฏฺเฏน, โส วิลุตฺโต วิลุปฺปตีติ.}}

\csromanheader{2}{6. มลฺลิกาสุตฺต}
\proseitem{127}{สาวตฺถินิทานํ.\csromanspacer{} อถ โข ราชา ปเสนทิ โกสโล เยน ภควา เตนุปสงฺกมิ, อุปสงฺกมิตฺวา ภควนฺตํ อภิวาเทตฺวา เอกมนฺตํ นิสีทิ.\csromanspacer{} อถ โข อญฺญตโร ปุริโส เยน ราชา ปเสนทิ โกสโล เตนุปสงฺกมิ, อุปสงฺกมิตฺวา รญฺโญ ปเสนทิสฺส โกสลสฺส อุปกณฺณเก อาโรเจสิ “มลฺลิกา เทว เทวี ธีตรํ วิชาตา”ติ.\csromanspacer{} เอวํ วุตฺเต ราชา ปเสนทิ โกสโล อนตฺตมโน อโหสิ.}

\prose{อถ โข ภควา ราชานํ ปเสนทิํ โกสลํ อนตฺตมนตํ วิทิตฺวา ตายํ เวลายํ อิมา คาถาโย อภาสิ\csromandash{}}

\csromangathameasure{อิตฺถีปิ หิ เอกจฺจิยา, เสยฺยา โปส ชนาธิป. \\ เมธาวินี สีลวตี, สสฺสุเทวา ปติพฺพตา. \\ ตสฺสา โย ชายติ โปโส, สูโร โหติ ทิสมฺปติ. \\ ตาทิสา สุภคิยา ปุตฺโต, รชฺชมฺปิ อนุสาสตีติ.}
\csromangathagroup{\csromangathastanza{อิตฺถีปิ หิ เอกจฺจิยา, เสยฺยา โปส ชนาธิป. \\ เมธาวินี สีลวตี, สสฺสุเทวา ปติพฺพตา.}\csromangathastanza{ตสฺสา โย ชายติ โปโส, สูโร โหติ ทิสมฺปติ. \\ ตาทิสา สุภคิยา\footnote{สุภริยาปุตฺโต (ก)} ปุตฺโต, รชฺชมฺปิ อนุสาสตีติ.}}

\csromanheader{2}{3. โกสลสํยุตฺต \textbf{7. อปฺปมาทสุตฺต}}
\proseitem{128}{\csromanfolio{87}สาวตฺถินิทานํ.\csromanspacer{} เอกมนฺตํ นิสีทิ, เอกมนฺตํ นิสินฺโน โข ราชา ปเสนทิ โกสโล ภควนฺตํ เอตทโวจ “อตฺถิ นุ โข ภนฺเต เอโก ธมฺโม, โย อุโภ อตฺเถ สมธิคฺคยฺห ติฏฺฐติ ทิฏฺฐธมฺมิกญฺเจว อตฺถํ สมฺปรายิกญฺจา”ติ.\csromanspacer{} อตฺถิ โข มหาราช เอโก ธมฺโม, โย อุโภ อตฺเถ สมธิคฺคยฺห ติฏฺฐติ ทิฏฺฐธมฺมิกญฺเจว อตฺถํ สมฺปรายิกญฺจาติ.}

\prose{กตโม ปน ภนฺเต เอโก ธมฺโม, โย อุโภ อตฺเถ สมธิคฺคยฺห ติฏฺฐติ ทิฏฺฐธมฺมิกญฺเจว อตฺถํ สมฺปรายิกญฺจาติ.\csromanspacer{} อปฺปมาโท โข มหาราช เอโก ธมฺโม, โย อุโภ อตฺเถ สมธิคฺคยฺห ติฏฺฐติ ทิฏฺฐธมฺมิกญฺเจว อตฺถํ สมฺปรายิกญฺจาติ.\csromanspacer{} เสยฺยถาปิ มหาราช ยานิ กานิจิ ชงฺคลานํ\footnote{ชงฺคมานํ (สี, อิ)} ปาณานํ ปทชาตานิ, สพฺพานิ ตานิ หตฺถิปเท สโมธานํ คจฺฉนฺติ, หตฺถิปทํ เตสํ อคฺคมกฺขายติ ยทิทํ มหนฺตตฺเตน.\csromanspacer{} เอวเมว โข มหาราช อปฺปมาโท เอโก ธมฺโม, โย อุโภ อตฺเถ สมธิคฺคยฺห ติฏฺฐติ ทิฏฺฐธมฺมิกญฺเจว อตฺถํ สมฺปรายิกญฺจาติ.\csromanspacer{} อิทมโวจ ฯเปฯ.}

\csromangathameasure{อายุํ อโรคิยํ วณฺณํ, สคฺคํ อุจฺจากุลีนตํ. \\ รติโย ปตฺถยนฺเตน, อุฬารา อปราปรา. \\ อปฺปมาทํ ปสํสนฺติ, ปุญฺญกิริยาสุ ปณฺฑิตา. \\ อปฺปมตฺโต อุโภ อตฺเถ, อธิคฺคณฺหาติ ปณฺฑิโต. \\ ทิฏฺเฐ ธมฺเม จ โย อตฺโถ, โย จตฺโถ สมฺปรายิโก. \\ อตฺถาภิสมยา ธีโร, ปณฺฑิโตติ ปวุจฺจตีติ.}
\csromangathagroup{\csromangathastanza{อายุํ อโรคิยํ วณฺณํ, สคฺคํ อุจฺจากุลีนตํ. \\ รติโย ปตฺถยนฺเตน, อุฬารา อปราปรา.}\csromangathastanza{อปฺปมาทํ ปสํสนฺติ, ปุญฺญกิริยาสุ ปณฺฑิตา. \\ อปฺปมตฺโต อุโภ อตฺเถ, อธิคฺคณฺหาติ ปณฺฑิโต.}\csromangathastanza{ทิฏฺเฐ ธมฺเม จ โย อตฺโถ, โย จตฺโถ สมฺปรายิโก. \\ อตฺถาภิสมยา ธีโร, ปณฺฑิโตติ ปวุจฺจตีติ.}}

\csromanheader{2}{8. กลฺยาณมิตฺตสุตฺต}
\proseitem{129}{สาวตฺถินิทานํ.\csromanspacer{} เอกมนฺตํ นิสินฺโน โข ราชา ปเสนทิ โกสโล ภควนฺตํ เอตทโวจ “อิธ มยฺหํ ภนฺเต รโหคตสฺส ปฏิสลฺลีนสฺส เอวํ เจตโส ปริวิตกฺโก อุทปาทิ ‘สฺวากฺขาโต ภควตา ธมฺโม, โส จ โข กลฺยาณมิตฺตสฺส กลฺยาณสหายสฺส กลฺยาณสมฺปวงฺกสฺส, โน ปาปมิตฺตสฺส โน ปาปสหายสฺส โน ปาปสมฺปวงฺกสฺสา’ติ”.\csromanspacer{} เอวเมตํ มหาราช, เอวเมตํ มหาราช, \csromanfolio{88}สฺวากฺขาโต มหาราช มยา ธมฺโม, โส จ โข กลฺยาณมิตฺตสฺส กลฺยาณสหายสฺส กลฺยาณสมฺปวงฺกสฺส, โน ปาปมิตฺตสฺส โน ปาปสหายสฺส โน ปาปสมฺปวงฺกสฺสาติ.}

\prose{เอกมิทาหํ มหาราช สมยํ สกฺเกสุ วิหรามิ นครกํ นาม สกฺยานํ นิคโม.\csromanspacer{} อถ โข มหาราช อานนฺโท ภิกฺขุ เยนาหํ เตนุปสงฺกมิ, อุปสงฺกมิตฺวา มํ อภิวาเทตฺวา เอกมนฺตํ นิสีทิ, เอกมนฺตํ นิสินฺโน โข มหาราช อานนฺโท ภิกฺขุ มํ เอตทโวจ “อุปฑฺฒมิทํ ภนฺเต พฺรหฺมจริยสฺส, ยทิทํ กลฺยาณมิตฺตตา กลฺยาณสหายตา กลฺยาณสมฺปวงฺกตา”ติ.}

\prose{เอวํ วุตฺตาหํ มหาราช อานนฺทํ ภิกฺขุํ เอตทโวจํ “มา เหวํ อานนฺท, มา เหวํ อานนฺท, สกลเมว หิทํ อานนฺท พฺรหฺมจริยํ, ยทิทํ กลฺยาณมิตฺตตา กลฺยาณสหายตา กลฺยาณสมฺปวงฺกตา, กลฺยาณมิตฺตสฺเสตํ อานนฺท ภิกฺขุโน ปาฏิกงฺขํ กลฺยาณสหายสฺส กลฺยาณสมฺปวงฺกสฺส, อริยํ อฏฺฐงฺคิกํ มคฺคํ ภาเวสฺสติ, อริยํ อฏฺฐงฺคิกํ มคฺคํ พหุลีกริสฺสติ.}

\prose{กถญฺจ อานนฺท ภิกฺขุ กลฺยาณมิตฺโต กลฺยาณสหาโย กลฺยาณสมฺปวงฺโก อริยํ อฏฺฐงฺคิกํ มคฺคํ ภาเวติ, อริยํ อฏฺฐงฺคิกํ มคฺคํ พหุลีกโรติ.\csromanspacer{} อิธานนฺท ภิกฺขุ สมฺมาทิฏฺฐิํ ภาเวติ วิเวกนิสฺสิตํ วิราคนิสฺสิตํ นิโรธนิสฺสิตํ โวสฺสคฺคปริณามิํ.\csromanspacer{} สมฺมาสงฺกปฺปํ ภาเวติ.\csromanspacer{} สมฺมาวาจํ ภาเวติ.\csromanspacer{} สมฺมากมฺมนฺตํ ภาเวติ.\csromanspacer{} สมฺมา-อาชีวํ ภาเวติ.\csromanspacer{} สมฺมาวายามํ ภาเวติ.\csromanspacer{} สมฺมาสติํ ภาเวติ.\csromanspacer{} สมฺมาสมาธิํ ภาเวติ วิเวกนิสฺสิตํ วิราคนิสฺสิตํ นิโรธนิสฺสิตํ โวสฺสคฺคปริณามิํ.\csromanspacer{} เอวํ โข อานนฺท ภิกฺขุ กลฺยาณมิตฺโต กลฺยาณสหาโย กลฺยาณสมฺปวงฺโก อริยํ อฏฺฐงฺคิกํ มคฺคํ ภาเวติ, อริยํ อฏฺฐงฺคิกํ มคฺคํ พหุลีกโรติ.\csromanspacer{} ตทมินาเปตํ อานนฺท ปริยาเยน เวทิตพฺพํ, ยถา สกลเมวิทํ พฺรหฺมจริยํ, ยทิทํ กลฺยาณมิตฺตตา กลฺยาณสหายตา กลฺยาณสมฺปวงฺกตาติ.}

\prose{มมํ หิ อานนฺท กลฺยาณมิตฺตํ อาคมฺม ชาติธมฺมา สตฺตา ชาติยา ปริมุจฺจนฺติ, ชราธมฺมา สตฺตา ชราย ปริมุจฺจนฺติ, พฺยาธิธมฺมา สตฺตา พฺยาธิโต ปริมุจฺจนฺติ, มรณธมฺมา สตฺตา มรเณน ปริมุจฺจนฺติ, โสกปริเทวทุกฺขโทมนสฺสุปายาสธมฺมา สตฺตา โสกปริเทวทุกฺขโทมนสฺสุปายาเสหิ ปริมุจฺจนฺติ.\csromanspacer{} อิมินา โข เอตํ อานนฺท ปริยาเยน เวทิตพฺพํ, ยถา}

\vspace{\tipitakaparskip}
\csromancenter{โกสลสํยุตฺต สกลเมวิทํ พฺรหฺมจริยํ, ยทิทํ กลฺยาณมิตฺตตา กลฺยาณสหายตา กลฺยาณสมฺปวงฺกตาติ.}
\prose{\csromanfolio{89}ตสฺมาติห เต มหาราช เอวํ สิกฺขิตพฺพํ “กลฺยาณมิตฺโต ภวิสฺสามิ กลฺยาณสหาโย กลฺยาณสมฺปวงฺโก”ติ, เอวํ หิ เต มหาราช สิกฺขิตพฺพํ.}

\prose{กลฺยาณมิตฺตสฺส เต มหาราช กลฺยาณสหายสฺส กลฺยาณสมฺปวงฺกสฺส อยํ เอโก ธมฺโม อุปนิสฺสาย วิหาตพฺโพ อปฺปมาโท กุสเลสุ ธมฺเมสุ.}

\prose{อปฺปมตฺตสฺส เต มหาราช วิหรโต อปฺปมาทํ อุปนิสฺสาย อิตฺถาคารสฺส อนุยนฺตสฺส เอวํ ภวิสฺสติ “ราชา โข อปฺปมตฺโต วิหรติ อปฺปมาทํ อุปนิสฺสาย, หนฺท มยมฺปิ อปฺปมตฺตา วิหราม อปฺปมาทํ อุปนิสฺสายา”ติ.}

\prose{อปฺปมตฺตสฺส เต มหาราช วิหรโต อปฺปมาทํ อุปนิสฺสาย ขตฺติยานมฺปิ อนุยนฺตานํ เอวํ ภวิสฺสติ “ราชา โข อปฺปมตฺโต วิหรติ อปฺปมาทํ อุปนิสฺสาย, หนฺท มยมฺปิ อปฺปมตฺตา วิหราม อปฺปมาทํ อุปนิสฺสายา”ติ.}

\prose{อปฺปมตฺตสฺส เต มหาราช วิหรโต อปฺปมาทํ อุปนิสฺสาย พลกายสฺสปิ เอวํ ภวิสฺสติ “ราชา โข อปฺปมตฺโต วิหรติ อปฺปมาทํ อุปนิสฺสาย, หนฺท มยมฺปิ อปฺปมตฺตา วิหราม อปฺปมาทํ อุปนิสฺสายา”ติ.}

\prose{อปฺปมตฺตสฺส เต มหาราช วิหรโต อปฺปมาทํ อุปนิสฺสาย เนคมชานปทสฺสปิ เอวํ ภวิสฺสติ “ราชา โข อปฺปมตฺโต วิหรติ อปฺปมาทํ อุปนิสฺสาย, หนฺท มยมฺปิ อปฺปมตฺตา วิหราม อปฺปมาทํ อุปนิสฺสายา”ติ.}

\prose{อปฺปมตฺตสฺส เต มหาราช วิหรโต อปฺปมาทํ อุปนิสฺสาย อตฺตาปิ คุตฺโต รกฺขิโต ภวิสฺสติ, อิตฺถาคารมฺปิ คุตฺตํ รกฺขิตํ ภวิสฺสติ, โกสโกฏฺฐาคารมฺปิ คุตฺตํ รกฺขิตํ ภวิสฺสตีติ.\csromanspacer{} อิทมโวจ ฯเปฯ.}

\csromangathameasure{โภเค ปตฺถยมาเนน, อุฬาเร อปราปเร. \\ อปฺปมาทํ ปสํสนฺติ, ปุญฺญกิริยาสุ ปณฺฑิตา. \\ อปฺปมตฺโต อุโภ อตฺเถ, อธิคฺคณฺหาติ ปณฺฑิโต. \\ ทิฏฺเฐ ธมฺเม จ โย อตฺโถ, โย จตฺโถ สมฺปรายิโก. \\ อตฺถาภิสมยา ธีโร, ปณฺฑิโตติ ปวุจฺจตีติ.}
\csromangathagroup{\csromangathastanza{โภเค ปตฺถยมาเนน, อุฬาเร อปราปเร. \\ อปฺปมาทํ ปสํสนฺติ, ปุญฺญกิริยาสุ ปณฺฑิตา.}\csromangathastanza{อปฺปมตฺโต อุโภ อตฺเถ, อธิคฺคณฺหาติ ปณฺฑิโต. \\ ทิฏฺเฐ ธมฺเม จ โย อตฺโถ, โย จตฺโถ สมฺปรายิโก. \\ อตฺถาภิสมยา ธีโร, ปณฺฑิโตติ ปวุจฺจตีติ.}}

\csromanheader{2}{9. ปฐม-อปุตฺตกสุตฺต}
\proseitem{130}{\csromanfolio{90}สาวตฺถินิทานํ.\csromanspacer{} อถ โข ราชา ปเสนทิ โกสโล ทิวา ทิวสฺส เยน ภควา เตนุปสงฺกมิ, อุปสงฺกมิตฺวา ภควนฺตํ อภิวาเทตฺวา เอกมนฺตํ นิสีทิ, เอกมนฺตํ นิสินฺนํ โข ราชานํ ปเสนทิํ โกสลํ ภควา เอตทโวจ “หนฺท กุโต นุ ตฺวํ มหาราช อาคจฺฉสิ ทิวา ทิวสฺสา”ติ.}

\prose{อิธ ภนฺเต สาวตฺถิยํ เสฏฺฐิ คหปติ กาลงฺกโต, ตมหํ อปุตฺตกํ สาปเตยฺยํ ราชนฺเตปุรํ อติหริตฺวา อาคจฺฉามิ.\csromanspacer{} อสีติ ภนฺเต สตสหสฺสานิ หิรญฺญสฺเสว, โก ปน วาโท รูปิยสฺส.\csromanspacer{} ตสฺส โข ปน ภนฺเต เสฏฺฐิสฺส คหปติสฺส เอวรูโป ภตฺตโภโค อโหสิ “กณาชกํ ภุญฺชติ พิลงฺคทุติยํ”.\csromanspacer{} เอวรูโป วตฺถโภโค อโหสิ “สาณํ ธาเรติ ติปกฺขวสนํ”.\csromanspacer{} เอวรูโป ยานโภโค อโหสิ “ชชฺชรรถเกน ยาติ ปณฺณฉตฺตเกน ธาริยมาเนนา”ติ.}

\prose{เอวเมตํ มหาราช, เอวเมตํ มหาราช, อสปฺปุริโส โข มหาราช อุฬาเร โภเค ลภิตฺวา เนวตฺตานํ สุเขติ ปีเณติ, น มาตาปิตโร สุเขติ ปีเณติ, น ปุตฺตทารํ สุเขติ ปีเณติ, น ทาสกมฺมกรโปริเส สุเขติ ปีเณติ, น มิตฺตามจฺเจ สุเขติ ปีเณติ, น สมณพฺราหฺมเณสุ อุทฺธคฺคิกํ ทกฺขิณํ ปติฏฺฐาเปติ โสวคฺคิกํ สุขวิปากํ สคฺคสํวตฺตนิกํ.\csromanspacer{} ตสฺส เต โภเค เอวํ สมฺมา อปริภุญฺชิยมาเน\footnote{อปริภุญฺชมาโน (สพฺพตฺถ)} ราชาโน วา หรนฺติ, โจรา วา หรนฺติ, อคฺคิ วา ฑหติ, อุทกํ วา วหติ, อปฺปิยา วา ทายาทา หรนฺติ.\csromanspacer{} เอวํส เต\footnote{เอวํ สนฺเต (สี, อิ)} มหาราช โภคา สมฺมา อปริภุญฺชิยมานา ปริกฺขยํ คจฺฉนฺติ, โน ปริโภคํ.}

\prose{เสยฺยถาปิ มหาราช อมนุสฺสฏฺฐาเน โปกฺขรณี อจฺโฉทกา สีโตทกา สาโตทกา เสโตทกา สุปติตฺถา รมณียา, ตํ ชโน เนว หเรยฺย น ปิเวยฺย น นหาเยยฺย น ยถาปจฺจยํ วา กเรยฺย, เอวํ หิ ตํ มหาราช อุทกํ สมฺมา อปริภุญฺชิยมานํ\footnote{อปริภุญฺชมานํ (สฺยา, กํ)} ปริกฺขยํ}

\vspace{\tipitakaparskip}
\csromancenter{โกสลสํยุตฺต คจฺเฉยฺย, โน ปริโภคํ.\csromanspacer{} เอวเมว โข มหาราช อสปฺปุริโส อุฬาเร โภเค ลภิตฺวา เนวตฺตานํ สุเขติ ปีเณติ, น มาตาปิตโร สุเขติ ปีเณติ, น ปุตฺตทารํ สุเขติ ปีเณติ, น ทาสกมฺมกรโปริเส สุเขติ ปีเณติ, น มิตฺตามจฺเจ สุเขติ ปีเณติ, น สมณพฺราหฺมเณสุ อุทฺธคฺคิกํ ทกฺขิณํ ปติฏฺฐาเปติ โสวคฺคิกํ สุขวิปากํ สคฺคสํวตฺตนิกํ.\csromanspacer{} ตสฺส เต โภเค เอวํ สมฺมา อปริภุญฺชิยมาเน ราชาโน วา หรนฺติ, โจรา วา หรนฺติ, อคฺคิ วา ฑหติ, อุทกํ วา วหติ, อปฺปิยา วา ทายาทา หรนฺติ.\csromanspacer{} เอวํส เต\footnote{เอวํ สนฺเต (สี, อิ)} มหาราช โภคา สมฺมา อปริภุญฺชิยมานา ปริกฺขยํ คจฺฉนฺติ, โน ปริโภคํ.}
\prose{\csromanfolio{91}สปฺปุริโส จ โข มหาราช อุฬาเร โภเค ลภิตฺวา อตฺตานํ สุเขติ ปีเณติ, มาตาปิตโร สุเขติ ปีเณติ, ปุตฺตทารํ สุเขติ ปีเณติ, ทาสกมฺมกรโปริเส สุเขติ ปีเณติ, มิตฺตามจฺเจ สุเขติ ปีเณติ, สมณพฺราหฺมเณสุ อุทฺธคฺคิกํ ทกฺขิณํ ปติฏฺฐาเปติ โสวคฺคิกํ สุขวิปากํ สคฺคสํวตฺตนิกํ.\csromanspacer{} ตสฺส เต โภเค เอวํ สมฺมา ปริภุญฺชิยมาเน เนว ราชาโน หรนฺติ, น โจรา หรนฺติ, น อคฺคิ ฑหติ, น อุทกํ วหติ, น อปฺปิยา ทายาทา หรนฺติ.\csromanspacer{} เอวํส เต มหาราช โภคา สมฺมา ปริภุญฺชิยมานา ปริโภคํ คจฺฉนฺติ, โน ปริกฺขยํ.}

\prose{เสยฺยถาปิ มหาราช คามสฺส วา นิคมสฺส วา อวิทูเร โปกฺขรณี อจฺโฉทกา สีโตทกา สาโตทกา เสโตทกา สุปติตฺถา รมณียา, ตญฺจ อุทกํ ชโน หเรยฺยปิ ปิเวยฺยปิ นหาเยยฺยปิ ยถาปจฺจยมฺปิ กเรยฺย.\csromanspacer{} เอวํ หิ ตํ มหาราช อุทกํ สมฺมา ปริภุญฺชิยมานํ ปริโภคํ คจฺเฉยฺย, โน ปริกฺขยํ.\csromanspacer{} เอวเมว โข มหาราช สปฺปุริโส อุฬาเร โภเค ลภิตฺวา อตฺตานํ สุเขติ ปีเณติ, มาตาปิตโร สุเขติ ปีเณติ, ปุตฺตทารํ สุเขติ ปีเณติ, ทาสกมฺมกรโปริเส สุเขติ ปีเณติ, มิตฺตามจฺเจ สุเขติ ปีเณติ.\csromanspacer{} สมณพฺราหฺมเณสุ อุทฺธคฺคิกํ ทกฺขิณํ ปติฏฺฐาเปติ โสวคฺคิกํ สุขวิปากํ สคฺคสํวตฺตนิกํ.\csromanspacer{} ตสฺส เต โภเค เอวํ สมฺมา ปริภุญฺชิยมาเน เนว ราชาโน หรนฺติ, น โจรา หรนฺติ, น อคฺคิ ฑหติ, น อุทกํ วหติ, น อปฺปิยา ทายาทา หรนฺติ.\csromanspacer{} เอวํส เต มหาราช โภคา สมฺมา ปริภุญฺชิยมานา ปริโภคํ คจฺฉนฺติ, โน ปริกฺขยนฺติ.}

\csromangathameasure{อมนุสฺสฏฺฐาเน อุทกํว สีตํ, \\ ตทเปยฺยมานํ ปริโสสเมติ. \\ เอวํ ธนํ กาปุริโส ลภิตฺวา, \\ เนวตฺตนา ภุญฺชติ โน ททาติ. \\ ธีโร จ วิญฺญู อธิคมฺม โภเค, \\ โส ภุญฺชติ กิจฺจกโร จ โหติ. \\ โส ญาติสํฆํ นิสโภ ภริตฺวา, \\ อนินฺทิโต สคฺคมุเปติ ฐานนฺติ.}
\csromangathagroup{\csromangathastanza{\csromanfolio{92}อมนุสฺสฏฺฐาเน อุทกํว สีตํ, \\ ตทเปยฺยมานํ ปริโสสเมติ. \\ เอวํ ธนํ กาปุริโส ลภิตฺวา, \\ เนวตฺตนา ภุญฺชติ โน ททาติ.}\csromangathastanza{ธีโร จ วิญฺญู อธิคมฺม โภเค, \\ โส ภุญฺชติ กิจฺจกโร จ โหติ. \\ โส ญาติสํฆํ นิสโภ ภริตฺวา, \\ อนินฺทิโต สคฺคมุเปติ ฐานนฺติ.}}

\csromanheader{2}{10. ทุติย-อปุตฺตกสุตฺต}
\proseitem{131}{อถ โข ราชา ปเสนทิ โกสโล ทิวา ทิวสฺส เยน ภควา เตนุปสงฺกมิ, อุปสงฺกมิตฺวา เอกมนฺตํ นิสินฺนํ โข ราชานํ ปเสนทิํ โกสลํ ภควา เอตทโวจ “หนฺท กุโต นุ ตฺวํ มหาราช อาคจฺฉสิ ทิวา ทิวสฺสา”ติ.}

\prose{อิธ ภนฺเต สาวตฺถิยํ เสฏฺฐิ คหปติ กาลงฺกโต, ตมหํ อปุตฺตกํ สาปเตยฺยํ ราชนฺเตปุรํ อติหริตฺวา อาคจฺฉามิ.\csromanspacer{} สตํ ภนฺเต สตสหสฺสานิ หิรญฺญสฺเสว, โก ปน วาโท รูปิยสฺส.\csromanspacer{} ตสฺส โข ปน ภนฺเต เสฏฺฐิสฺส คหปติสฺส เอวรูโป ภตฺตโภโค อโหสิ “กณาชกํ ภุญฺชติ พิลงฺคทุติยํ”.\csromanspacer{} เอวรูโป วตฺถโภโค อโหสิ “สาณํ ธาเรติ ติปกฺขวสนํ”.\csromanspacer{} เอวรูโป ยานโภโค อโหสิ “ชชฺชรรถเกน ยาติ ปณฺณฉตฺเตเกน ธาริยมาเนนา”ติ.}

\prose{เอวเมตํ มหาราช, เอวเมตํ มหาราช, ภูตปุพฺพํ โส มหาราช เสฏฺฐิ คหปติ ตคฺครสิขิํ นาม ปจฺเจกสมฺพุทฺธํ ปิณฺฑปาเตน ปฏิปาเทสิ, “เทถ สมณสฺส ปิณฺฑนฺ”ติ วตฺวา อุฏฺฐายาสนา ปกฺกามิ.\csromanspacer{} ทตฺวา จ ปน ปจฺฉา วิปฺปฏิสารี อโหสิ “วรเมตํ ปิณฺฑปาตํ ทาสา วา กมฺมกรา วา ภุญฺเชยฺยุนฺ”ติ.\csromanspacer{} ภาตุ จ ปน เอกปุตฺตกํ สาปเตยฺยสฺส การณา ชีวิตา โวโรเปสิ.}

\prose{ยํ โข โส มหาราช เสฏฺฐิ คหปติ ตคฺครสิขิํ ปจฺเจกสมฺพุทฺธํ ปิณฺฑปาเตน ปฏิปาเทสิ, ตสฺส กมฺมสฺส วิปาเกน สตฺตกฺขตฺตุํ สุคติํ สคฺคํ โลกํ อุปปชฺชิ, ตสฺเสว กมฺมสฺส วิปากาวเสเสน อิมิสฺสาเยว}

\vspace{\tipitakaparskip}
\csromancenter{โกสลสํยุตฺต สาวตฺถิยา สตฺตกฺขตฺตุํ เสฏฺฐิตฺตํ กาเรสิ.\csromanspacer{} ยํ โข โส มหาราช เสฏฺฐิ คหปติ ทตฺวา ปจฺฉา วิปฺปฏิสารี อโหสิ “วรเมตํ ปิณฺฑปาตํ ทาสา วา กมฺมกรา วา ภุญฺเชยฺยุนฺ”ติ, ตสฺส กมฺมสฺส วิปาเกน นาสฺสุฬาราย ภตฺตโภคาย จิตฺตํ นมติ, นาสฺสุฬาราย วตฺถโภคาย จิตฺตํ นมติ, นาสฺสุฬาราย ยานโภคาย จิตฺตํ นมติ, นาสฺสุฬารานํ ปญฺจนฺนํ กามคุณานํ โภคาย จิตฺตํ นมติ.\csromanspacer{} ยํ โข โส มหาราช เสฏฺฐิ คหปติ ภาตุ จ ปน เอกปุตฺตกํ สาปเตยฺยสฺส การณา ชีวิตา โวโรเปสิ, ตสฺส กมฺมสฺส วิปาเกน พหูนิ วสฺสานิ พหูนิ วสฺสสตานิ พหูนิ วสฺสสหสฺสานิ พหูนิ วสฺสสตสหสฺสานิ นิรเย ปจฺจิตฺถ, ตสฺเสว กมฺมสฺส วิปากาวเสเสน อิทํ สตฺตมํ อปุตฺตกํ สาปเตยฺยํ ราชโกสํ ปเวเสติ.\csromanspacer{} ตสฺส โข มหาราช เสฏฺฐิสฺส คหปติสฺส ปุราณญฺจ ปุญฺญํ ปริกฺขีณํ นวญฺจ ปุญฺญํ อนุปจิตํ, อชฺช ปน มหาราช เสฏฺฐิ คหปติ มหาโรรุเว นิรเย ปจฺจตีติ.\csromanspacer{} เอวํ ภนฺเต เสฏฺฐิ คหปติ มหาโรรุวํ นิรยํ อุปปนฺโนติ.\csromanspacer{} เอวํ มหาราช เสฏฺฐิ คหปติ มหาโรรุวํ นิรยํ อุปปนฺโนติ.\csromanspacer{} อิทมโวจ ฯเปฯ.}
\prose{\csromanfolio{93}ธญฺญํ ธนํ รชตํ ชาตรูปํ, ปริคฺคหํ วาปิ ยทตฺถิ กิญฺจิ.}

\prose{ทาสา กมฺมกรา เปสฺสา, เย จสฺส อนุชีวิโน.}

\prose{สพฺพํ นาทาย คนฺตพฺพํ, สพฺพํ นิกฺขิปฺปคามินํ\footnote{นิกฺขีปคามินํ (สฺยา, กํ, ก)}.}

\prose{ยญฺจ กโรติ กาเยน, วาจาย อุท เจตสา.}

\prose{ตํ หิ ตสฺส สกํ โหติ, ตญฺจ อาทาย คจฺฉติ.}

\prose{ตญฺจสฺส อนุคํ โหติ, ฉายาว อนปายินี.}

\prose{ตสฺมา กเรยฺย กลฺยาณํ, นิจยํ สมฺปรายิกํ.}

\prose{ปุญฺญานิ ปรโลกสฺมิํ, ปติฏฺฐา โหนฺติ ปาณินนฺติ.}

\vspace{\tipitakaparskip}
\csromancenter{ทุติโย วคฺโค.}
\titlehead{\textbf{ตสฺสุทฺทานํ}}
\csromangathameasure{ชฏิลา ปญฺจ ราชาโน, โทณปากกุเรน จ. \\ สงฺคาเมน เทฺว วุตฺตานิ, มลฺลิกา เทฺว อปฺปมาเทน จ. \\ อปุตฺตเกน เทฺว วุตฺตา, วคฺโค เตน ปวุจฺจตีติ.}
\csromangathagroup{\csromangathastanza{ชฏิลา ปญฺจ ราชาโน, โทณปากกุเรน จ. \\ สงฺคาเมน เทฺว วุตฺตานิ, มลฺลิกา\footnote{ธีตรา (พหูสุ)} เทฺว อปฺปมาเทน จ. \\ อปุตฺตเกน เทฺว วุตฺตา, วคฺโค เตน ปวุจฺจตีติ.}}

\chapterheadpageread{3. ตติยวคฺค}
\csromanheader{2}{1. ปุคฺคลสุตฺต}
\proseitem{132}{\csromanfolio{94}สาวตฺถินิทานํ.\csromanspacer{} อถ โข ราชา ปเสนทิ โกสโล เยน ภควา เตนุปสงฺกมิ, อุปสงฺกมิตฺวา ภควนฺตํ อภิวาเทตฺวา เอกมนฺตํ นิสีทิ, เอกมนฺตํ นิสินฺนํ โข ราชานํ ปเสนทิํ โกสลํ ภควา เอตทโวจ “จตฺตาโรเม มหาราช ปุคฺคลา สนฺโต สํวิชฺชมานา โลกสฺมิํ.\csromanspacer{} กตเม จตฺตาโร, ตโมตมปรายโน ตโมโชติปรายโน โชติตมปรายโน โชติโชติปรายโน.}

\prose{กถญฺจ มหาราช ปุคฺคโล ตโมตมปรายโน โหติ, อิธ มหาราช เอกจฺโจ ปุคฺคโล นีเจ กุเล ปจฺจาชาโต โหติ จณฺฑาลกุเล วา เวนกุเล\footnote{เวณกุเล (สี, สฺยา, กํ, อิ)} วา เนสาทกุเล วา รถการกุเล วา ปุกฺกุสกุเล วา ทลิทฺเท อปฺปนฺนปานโภชเน กสิรวุตฺติเก, ยตฺถ กสิเรน ฆาสจฺฉาโท ลพฺภติ.\csromanspacer{} โส จ โหติ ทุพฺพณฺโณ ทุทฺทสิโก โอโกฏิมโก พวฺหาพาโธ\footnote{พหฺวาพาโธ (ก)} กาโณ วา กุณี วา ขญฺโช วา ปกฺขหโต วา, น ลาภี อนฺนสฺส ปานสฺส วตฺถสฺส ยานสฺส มาลาคนฺธวิเลปนสฺส เสยฺยาวสถปทีเปยฺยสฺส.\csromanspacer{} โส กาเยน ทุจฺจริตํ จรติ, วาจาย ทุจฺจริตํ จรติ, มนสา ทุจฺจริตํ จรติ.\csromanspacer{} โส กาเยน ทุจฺจริตํ จริตฺวา วาจาย ทุจฺจริตํ จริตฺวา มนสา ทุจฺจริตํ จริตฺวา กายสฺส เภทา ปรํ มรณา อปายํ ทุคฺคติํ วินิปาตํ นิรยํ อุปปชฺชติ.}

\prose{เสยฺยถาปิ มหาราช ปุริโส อนฺธการา วา อนฺธการํ คจฺเฉยฺย, ตมา วา ตมํ คจฺเฉยฺย, โลหิตมลา วา โลหิตมลํ คจฺเฉยฺย.\csromanspacer{} ตถูปมาหํ มหาราช อิมํ ปุคฺคลํ วทามิ.\csromanspacer{} เอวํ โข มหาราช ปุคฺคโล ตโมตมปรายโน โหติ.}

\prose{กถญฺจ มหาราช ปุคฺคโล ตโมโชติปรายโน โหติ, อิธ มหาราช เอกจฺโจ ปุคฺคโล นีเจ กุเล ปจฺจาชาโต โหติ จณฺฑาลกุเล วา เวนกุเล วา เนสาทกุเล วา รถการกุเล วา ปุกฺกุสกุเล วา ทลิทฺเท อปฺปนฺนปานโภชเน กสิรวุตฺติเก, ยตฺถ กสิเรน ฆาสจฺฉาโท ลพฺภติ.\csromanspacer{} โส จ โข โหติ ทุพฺพณฺโณ}

\vspace{\tipitakaparskip}
\csromancenter{โกสลสํยุตฺต ทุทฺทสิโก โอโกฏิมโก พวฺหาพาโธ กาโณ วา กุณี วา ขญฺโช วา ปกฺขหโต วา, น ลาภี อนฺนสฺส ปานสฺส วตฺถสฺส ยานสฺส มาลาคนฺธวิเลปนสฺส เสยฺยาวสถปทีเปยฺยสฺส.\csromanspacer{} โส กาเยน สุจริตํ จรติ, วาจาย สุจริตํ จรติ, มนสา สุจริตํ จรติ.\csromanspacer{} โส กาเยน สุจริตํ จริตฺวา วาจาย สุจริตํ จริตฺวา มนสา สุจริตํ จริตฺวา กายสฺส เภทา ปรํ มรณา สุคติํ สคฺคํ โลกํ อุปปชฺชติ.}
\prose{\csromanfolio{95}เสยฺยถาปิ มหาราช ปุริโส ปถวิยา วา ปลฺลงฺกํ อาโรเหยฺย, ปลฺลงฺกา วา อสฺสปิฏฺฐิํ อาโรเหยฺย, อสฺสปิฏฺฐิยา วา หตฺถิกฺขนฺธํ อาโรเหยฺย, หตฺถิกฺขนฺธา วา ปาสาทํ อาโรเหยฺย.\csromanspacer{} ตถูปมาหํ มหาราช อิมํ ปุคฺคลํ วทามิ.\csromanspacer{} เอวํ โข มหาราช ปุคฺคโล ตโมโชติปรายโน โหติ.}

\prose{กถญฺจ มหาราช ปุคฺคโล โชติตมปรายโน โหติ, อิธ มหาราช เอกจฺโจ ปุคฺคโล อุจฺเจ กุเล ปจฺจาชาโต โหติ ขตฺติยมหาสาลกุเล วา พฺราหฺมณมหาสาลกุเล วา คหปติมหาสาลกุเล วา อฑฺเฒ มหทฺธเน มหาโภเค ปหูตชาตรูปรชเต ปหูตวิตฺตูปกรเณ ปหูตธนธญฺเญ, โส จ โหติ อภิรูโป ทสฺสนีโย ปาสาทิโก ปรมาย วณฺณโปกฺขรตาย สมนฺนาคโต, ลาภี อนฺนสฺส ปานสฺส วตฺถสฺส ยานสฺส มาลาคนฺธวิเลปนสฺส เสยฺยาวสถปทีเปยฺยสฺส.\csromanspacer{} โส กาเยน ทุจฺจริตํ จรติ, วาจาย ทุจฺจริตํ จรติ, มนสา ทุจฺจริตํ จรติ.\csromanspacer{} โส กาเยน ทุจฺจริตํ จริตฺวา วาจาย ทุจฺจริตํ จริตฺวา มนสา ทุจฺจริตํ จริตฺวา กายสฺส เภทา ปรํ มรณา อปายํ ทุคฺคติํ วินิปาตํ นิรยํ อุปปชฺชติ.}

\prose{เสยฺยถาปิ มหาราช ปุริโส ปาสาทา วา หตฺถิกฺขนฺธํ โอโรเหยฺย, หตฺถิกฺขนฺธา วา อสฺสปิฏฺฐิํ โอโรเหยฺย, อสฺสปิฏฺฐิยา วา ปลฺลงฺกํ โอโรเหยฺย.\csromanspacer{} ปลฺลงฺกา วา ปถวิํ โอโรเหยฺย, ปถวิยา วา อนฺธการํ ปวิเสยฺย, ตถูปมาหํ มหาราช อิมํ ปุคฺคลํ วทามิ.\csromanspacer{} เอวํ โข มหาราช ปุคฺคโล โชติตมปรายโน โหติ.}

\prose{กถญฺจ มหาราช ปุคฺคโล โชติโชติปรายโน โหติ, อิธ มหาราช เอกจฺโจ ปุคฺคโล อุจฺเจ กุเล ปจฺจาชาโต โหติ ขตฺติยมหาสาลกุเล วา พฺราหฺมณมหาสาลกุเล วา \csromanfolio{96}คหปติมหาสาลกุเล วา อฑฺเฒ มหทฺธเน มหาโภเค ปหูตชาตรูปรชเต ปหูตวิตฺตูปกรเณ ปหูตธนธญฺเญ, โส จ โหติ อภิรูโป ทสฺสนีโย ปาสาทิโก ปรมาย วณฺณโปกฺขรตาย สมนฺนาคโต, ลาภี อนฺนสฺส ปานสฺส วตฺถสฺส ยานสฺส มาลาคนฺธวิเลปนสฺส เสยฺยาวสถปทีเปยฺยสฺส.\csromanspacer{} โส กาเยน สุจริตํ จรติ, วาจาย สุจริตํ จรติ, มนสา สุจริตํ จรติ.\csromanspacer{} โส กาเยน สุจริตํ จริตฺวา วาจาย สุจริตํ จริตฺวา มนสา สุจริตํ จริตฺวา กายสฺส เภทา ปรํ มรณา สุคติํ สคฺคํ โลกํ อุปปชฺชติ.}

\prose{เสยฺยถาปิ มหาราช ปุริโส ปลฺลงฺกา วา ปลฺลงฺกํ สงฺกเมยฺย, อสฺสปิฏฺฐิยา วา อสฺสปิฏฺฐิํ สงฺกเมยฺย, หตฺถิกฺขนฺธา วา หตฺถิกฺขนฺธํ สงฺกเมยฺย, ปาสาทา วา ปาสาทํ สงฺกเมยฺย.\csromanspacer{} ตถูปมาหํ มหาราช อิมํ ปุคฺคลํ วทามิ.\csromanspacer{} เอวํ โข มหาราช ปุคฺคโล โชติโชติปรายโน โหติ.\csromanspacer{} อิเม โข มหาราช จตฺตาโร ปุคฺคลา สนฺโต สํวิชฺชมานา โลกสฺมินฺติ.\csromanspacer{} อิทมโวจ ฯเปฯ.}

\csromangathameasure{ทลิทฺโท ปุริโส ราช, อสฺสทฺโธ โหติ มจฺฉรี. \\ กทริโย ปาปสงฺกปฺโป, มิจฺฉาทิฏฺฐิ อนาทโร. \\ สมเณ พฺราหฺมเณ วาปิ, อญฺเญ วาปิ วนิพฺพเก. \\ อกฺโกสติ ปริภาสติ, นตฺถิโก โหติ โรสโก. \\ ททมานํ นิวาเรติ, ยาจมานาน โภชนํ. \\ ตาทิโส ปุริโส ราช, มียมาโน ชนาธิป. \\ อุเปติ นิรยํ โฆรํ, ตโมตมปรายโน. \\ ทลิทฺโท ปุริโส ราช, สทฺโธ โหติ อมจฺฉรี. \\ ททาติ เสฏฺฐสงฺกปฺโป, อพฺยคฺคมนโส นโร. \\ สมเณ พฺราหฺมเณ วาปิ, อญฺเญ วาปิ วนิพฺพเก. \\ อุฏฺฐาย อภิวาเทติ, สมจริยาย สิกฺขติ. \\ ททมานํ น วาเรติ, ยาจมานาน โภชนํ. \\ ตาทิโส ปุริโส ราช, มียมาโน ชนาธิป. \\ อุเปติ ติทิวํ ฐานํ, ตโมโชติปรายโน.}
\csromangathagroup{\csromangathastanza{ทลิทฺโท ปุริโส ราช, อสฺสทฺโธ โหติ มจฺฉรี. \\ กทริโย ปาปสงฺกปฺโป, มิจฺฉาทิฏฺฐิ อนาทโร.}\csromangathastanza{สมเณ พฺราหฺมเณ วาปิ, อญฺเญ วาปิ วนิพฺพเก. \\ อกฺโกสติ ปริภาสติ, นตฺถิโก โหติ โรสโก.}\csromangathastanza{ททมานํ นิวาเรติ, ยาจมานาน โภชนํ. \\ ตาทิโส ปุริโส ราช, มียมาโน ชนาธิป.}\csromangathastanza{อุเปติ นิรยํ โฆรํ, ตโมตมปรายโน. \\ ทลิทฺโท ปุริโส ราช, สทฺโธ โหติ อมจฺฉรี.}\csromangathastanza{ททาติ เสฏฺฐสงฺกปฺโป, อพฺยคฺคมนโส นโร. \\ สมเณ พฺราหฺมเณ วาปิ, อญฺเญ วาปิ วนิพฺพเก.}\csromangathastanza{อุฏฺฐาย อภิวาเทติ, สมจริยาย สิกฺขติ. \\ ททมานํ น วาเรติ\footnote{น นิวาเรติ (สี)}, ยาจมานาน โภชนํ.}\csromangathastanza{ตาทิโส ปุริโส ราช, มียมาโน ชนาธิป. \\ อุเปติ ติทิวํ ฐานํ, ตโมโชติปรายโน.}}

\csromanheader{2}{3. โกสลสํยุตฺต}
\csromangathameasure{อฑฺโฒ เจ ปุริโส ราช, อสฺสทฺโธ โหติ มจฺฉรี. \\ กทริโย ปาปสงฺกปฺโป, มิจฺฉาทิฏฺฐิ อนาทโร. \\ สมเณ พฺราหฺมเณ วาปิ, อญฺเญ วาปิ วนิพฺพเก. \\ อกฺโกสติ ปริภาสติ, นตฺถิโก โหติ โรสโก. \\ ททมานํ นิวาเรติ, ยาจมานาน โภชนํ. \\ ตาทิโส ปุริโส ราช, มียมาโน ชนาธิป. \\ อุเปติ นิรยํ โฆรํ, โชติตมปรายโน. \\ อฑฺโฒ เจ ปุริโส ราช, สทฺโธ โหติ อมจฺฉรี. \\ ททาติ เสฏฺฐสงฺกปฺโป, อพฺยคฺคมนโส นโร. \\ สมเณ พฺราหฺมเณ วาปิ, อญฺเญ วาปิ วนิพฺพเก. \\ อุฏฺฐาย อภิวาเทติ, สมจริยาย สิกฺขติ. \\ ททมานํ น วาเรติ, ยาจมานาน โภชนํ. \\ ตาทิโส ปุริโส ราช, มียมาโน ชนาธิป. \\ อุเปติ ติทิวํ ฐานํ, โชติโชติปรายโนติ.}
\csromangathagroup{\csromangathastanza{\csromanfolio{97}อฑฺโฒ เจ\footnote{อฑฺโฒ เว (อิ, ก)} ปุริโส ราช, อสฺสทฺโธ โหติ มจฺฉรี. \\ กทริโย ปาปสงฺกปฺโป, มิจฺฉาทิฏฺฐิ อนาทโร.}\csromangathastanza{สมเณ พฺราหฺมเณ วาปิ, อญฺเญ วาปิ วนิพฺพเก. \\ อกฺโกสติ ปริภาสติ, นตฺถิโก โหติ โรสโก.}\csromangathastanza{ททมานํ นิวาเรติ, ยาจมานาน โภชนํ. \\ ตาทิโส ปุริโส ราช, มียมาโน ชนาธิป.}\csromangathastanza{อุเปติ นิรยํ โฆรํ, โชติตมปรายโน. \\ อฑฺโฒ เจ ปุริโส ราช, สทฺโธ โหติ อมจฺฉรี.}\csromangathastanza{ททาติ เสฏฺฐสงฺกปฺโป, อพฺยคฺคมนโส นโร. \\ สมเณ พฺราหฺมเณ วาปิ, อญฺเญ วาปิ วนิพฺพเก.}\csromangathastanza{อุฏฺฐาย อภิวาเทติ, สมจริยาย สิกฺขติ. \\ ททมานํ น วาเรติ, ยาจมานาน โภชนํ.}\csromangathastanza{ตาทิโส ปุริโส ราช, มียมาโน ชนาธิป. \\ อุเปติ ติทิวํ ฐานํ, โชติโชติปรายโนติ.}}

\csromanheader{2}{2. อยฺยิกาสุตฺต}
\proseitem{133}{สาวตฺถินิทานํ.\csromanspacer{} เอกมนฺตํ นิสินฺนํ โข ราชานํ ปเสนทิํ โกสลํ ภควา เอตทโวจ “หนฺท กุโต นุ ตฺวํ มหาราช อาคจฺฉสิ ทิวา ทิวสฺสา”ติ.}

\prose{อยฺยิกา เม ภนฺเต กาลงฺกตา ชิณฺณา วุฑฺฒา มหลฺลิกา อทฺธคตา วโย- อนุปฺปตฺตา วีสวสฺสสติกา ชาติยา.\csromanspacer{} อยฺยิกา โข ปน เม ภนฺเต ปิยา โหติ มนาปา.\csromanspacer{} หตฺถิรตเนน เจปาหํ ภนฺเต ลเภยฺยํ “มา เม อยฺยิกา กาลมกาสี”ติ, หตฺถิรตนมฺปาหํ ทเทยฺยํ “มา เม อยฺยิกา กาลมกาสี”ติ.\csromanspacer{} อสฺสรตเนน เจปาหํ ภนฺเต ลเภยฺยํ “มา เม อยฺยิกา กาลมกาสี”ติ, อสฺสรตนมฺปาหํ ทเทยฺยํ “มา เม อยฺยิกา กาลมกาสี”ติ.\csromanspacer{} คามวเรน เจปาหํ ภนฺเต ลเภยฺยํ “มา เม อยฺยิกา กาลมกาสี”ติ, คามวรมฺปาหํ ทเทยฺยํ “มา เม อยฺยิกา กาลมกาสี”ติ.\csromanspacer{} ชนปทปเทเสน\footnote{ชนปเทน (สี, สฺยา, อิ)} เจปาหํ ภนฺเต \csromanfolio{98}ลเภยฺยํ “มา เม อยฺยิกา กาลมกาสี”ติ, ชนปทปเทสมฺปาหํ ทเทยฺยํ “มา เม อยฺยิกา กาลมกาสี”ติ.\csromanspacer{} สพฺเพ สตฺตา มหาราช มรณธมฺมา มรณปริโยสานา มรณํ อนตีตาติ.\csromanspacer{} อจฺฉริยํ ภนฺเต, อพฺภุตํ ภนฺเต, ยาวสุภาสิตมิทํ ภนฺเต ภควตา “สพฺเพ สตฺตา มรณธมฺมา มรณปริโยสานา มรณํ อนตีตา”ติ.}

\prose{เอวเมตํ มหาราช, เอวเมตํ มหาราช, สพฺเพ สตฺตา มรณธมฺมา มรณปริโยสานา มรณํ อนตีตา.\csromanspacer{} เสยฺยถาปิ มหาราช ยานิ กานิจิ กุมฺภการภาชนานิ อามกานิ เจว ปกฺกานิ จ, สพฺพานิ ตานิ เภทนธมฺมานิ เภทนปริโยสานานิ เภทนํ อนตีตานิ.\csromanspacer{} เอวเมว โข มหาราช สพฺเพ สตฺตา มรณธมฺมา มรณปริโยสานา มรณํ อนตีตาติ.\csromanspacer{} อิทมโวจ ฯเปฯ.}

\csromangathameasure{สพฺเพ สตฺตา มริสฺสนฺติ, มรณนฺตํ หิ ชีวิตํ. \\ ยถากมฺมํ คมิสฺสนฺติ, ปุญฺญปาปผลูปคา. \\ นิรยํ ปาปกมฺมนฺตา, ปุญฺญกมฺมา จ สุคฺคติํ. \\ ตสฺมา กเรยฺย กลฺยาณํ, นิจยํ สมฺปรายิกํ. \\ ปุญฺญานิ ปรโลกสฺมิํ, ปติฏฺฐา โหนฺติ ปาณินนฺติ.}
\csromangathagroup{\csromangathastanza{สพฺเพ สตฺตา มริสฺสนฺติ, มรณนฺตํ หิ ชีวิตํ. \\ ยถากมฺมํ คมิสฺสนฺติ, ปุญฺญปาปผลูปคา.}\csromangathastanza{นิรยํ ปาปกมฺมนฺตา, ปุญฺญกมฺมา จ สุคฺคติํ. \\ ตสฺมา กเรยฺย กลฺยาณํ, นิจยํ สมฺปรายิกํ. \\ ปุญฺญานิ ปรโลกสฺมิํ, ปติฏฺฐา โหนฺติ ปาณินนฺติ.}}

\csromanheader{2}{3. โลกสุตฺต}
\proseitem{134}{สาวตฺถินิทานํ.\csromanspacer{} เอกมนฺตํ นิสินฺโน โข ราชา ปเสนทิ โกสโล ภควนฺตํ เอตทโวจ “กติ นุ โข ภนฺเต โลกสฺส ธมฺมา อุปฺปชฺชมานา อุปฺปชฺชนฺติ อหิตาย ทุกฺขาย อผาสุวิหารายา”ติ.\csromanspacer{} ตโย โข มหาราช โลกสฺส ธมฺมา อุปฺปชฺชมานา อุปฺปชฺชนฺติ อหิตาย ทุกฺขาย อผาสุวิหาราย.\csromanspacer{} กตเม ตโย, โลโภ โข มหาราช โลกสฺส ธมฺโม อุปฺปชฺชมาโน อุปฺปชฺชติ อหิตาย ทุกฺขาย อผาสุวิหาราย.\csromanspacer{} โทโส โข มหาราช โลกสฺส ธมฺโม อุปฺปชฺชมาโน อุปฺปชฺชติ อหิตาย ทุกฺขาย อผาสุวิหาราย.\csromanspacer{} โมโห โข มหาราช โลกสฺส ธมฺโม อุปฺปชฺชมาโน อุปฺปชฺชติ อหิตาย ทุกฺขาย อผาสุวิหาราย.\csromanspacer{} อิเม โข มหาราช ตโย โลกสฺส ธมฺมา อุปฺปชฺชมานา อุปฺปชฺชนฺติ อหิตาย ทุกฺขาย อผาสุวิหารายาติ.\csromanspacer{} อิทมโวจ ฯเปฯ.}

\csromanheader{2}{3. โกสลสํยุตฺต}
\csromangathameasure{โลโภ โทโส จ โมโห จ, ปุริสํ ปาปเจตสํ. \\ หิํสนฺติ อตฺตสมฺภูตา, ตจสารํว สมฺผลนฺติ.}
\csromangathagroup{\csromangathastanza{\csromanfolio{99}โลโภ โทโส จ โมโห จ, ปุริสํ ปาปเจตสํ. \\ หิํสนฺติ อตฺตสมฺภูตา, ตจสารํว สมฺผลนฺติ.}}

\csromanheader{2}{4. อิสฺสตฺตสุตฺต}
\proseitem{135}{สาวตฺถินิทานํ.\csromanspacer{} เอกมนฺตํ นิสินฺโน โข ราชา ปเสนทิ โกสโล ภควนฺตํ เอตทโวจ “กตฺถ นุ โข ภนฺเต ทานํ ทาตพฺพนฺ”ติ.\csromanspacer{} ยตฺถ โข มหาราช จิตฺตํ ปสีทตีติ.\csromanspacer{} กตฺถ ปน ภนฺเต ทินฺนํ มหปฺผลนฺติ.\csromanspacer{} อญฺญํ โข เอตํ มหาราช กตฺถ ทานํ ทาตพฺพํ, อญฺญํ ปเนตํ กตฺถ ทินฺนํ มหปฺผลนฺติ.\csromanspacer{} สีลวโต โข มหาราช ทินฺนํ มหปฺผลํ, โน ตถา ทุสฺสีเล.\csromanspacer{} เตน หิ มหาราช ตญฺเญเวตฺถ ปฏิปุจฺฉิสฺสามิ.\csromanspacer{} ยถา เต ขเมยฺย, ตถา นํ พฺยากเรยฺยาสิ.\csromanspacer{} ตํ กิํ มญฺญสิ มหาราช, อิธ ตฺยสฺส ยุทฺธํ ปจฺจุปฏฺฐิตํ สงฺคาโม สมุปพฺยูโฬฺห\footnote{สมูปพฺพูโฬฺห (สี), สมุปพฺพุโฬฺห (อิ)}.\csromanspacer{} อถ อาคจฺเฉยฺย ขตฺติยกุมาโร อสิกฺขิโต อกตหตฺโถ อกตโยคฺโค อกตูปาสโน ภีรุ ฉมฺภี อุตฺราสี ปลายี, ภเรยฺยาสิ ตํ ปุริสํ, อตฺโถ จ เต ตาทิเสน ปุริเสนาติ.\csromanspacer{} นาหํ ภนฺเต ภเรยฺยํ ตํ ปุริสํ, น จ เม อตฺโถ ตาทิเสน ปุริเสนาติ.\csromanspacer{} อถ อาคจฺเฉยฺย พฺราหฺมณกุมาโร อสิกฺขิโต ฯเปฯ.\csromanspacer{} อถ อาคจฺเฉยฺย เวสฺสกุมาโร อสิกฺขิโต ฯเปฯ.\csromanspacer{} อถ อาคจฺเฉยฺย สุทฺทกุมาโร อสิกฺขิโต ฯเปฯ น จ เม อตฺโถ ตาทิเสน ปุริเสนาติ.}

\prose{ตํ กิํ มญฺญสิ มหาราช, อิธ ตฺยสฺส ยุทฺธํ ปจฺจุปฏฺฐิตํ สงฺคาโม สมุปพฺยูโฬฺห.\csromanspacer{} อถ อาคจฺเฉยฺย ขตฺติยกุมาโร สุสิกฺขิโต กตหตฺโถ กตโยคฺโค กตูปาสโน อภีรุ อจฺฉมฺภี อนุตฺราสี อปลายี, ภเรยฺยาสิ ตํ ปุริสํ, อตฺโถ จ เต ตาทิเสน ปุริเสนาติ.\csromanspacer{} ภเรยฺยาหํ ภนฺเต ตํ ปุริสํ, อตฺโถ จ เม ตาทิเสน ปุริเสนาติ.\csromanspacer{} อถ อาคจฺเฉยฺย พฺราหฺมณกุมาโร ฯเปฯ.\csromanspacer{} อถ อาคจฺเฉยฺย เวสฺสกุมาโร ฯเปฯ.\csromanspacer{} อถ อาคจฺเฉยฺย สุทฺทกุมาโร สุสิกฺขิโต กตหตฺโถ กตโยคฺโค กตูปาสโน อภีรุ อจฺฉมฺภี อนุตฺราสี อปลายี, ภเรยฺยาสิ ตํ ปุริสํ, อตฺโถ จ เต ตาทิเสน ปุริเสนาติ.\csromanspacer{} ภเรยฺยาหํ ภนฺเต ตํ ปุริสํ, อตฺโถ จ เม ตาทิเสน ปุริเสนาติ.}

\prose{เอวเมว โข มหาราช ยสฺมา กสฺมา เจปิ\footnote{ยสฺมา เจปิ (สี, สฺยา, กํ, ก)} กุลา อคารสฺมา อนคาริยํ ปพฺพชิโต โหติ, โส จ โหติ ปญฺจงฺควิปฺปหีโน ปญฺจงฺคสมนฺนาคโต, \csromanfolio{100}ตสฺมิํ ทินฺนํ มหปฺผลํ โหติ.\csromanspacer{} กตมานิ ปญฺจงฺคานิ ปหีนานิ โหนฺติ.\csromanspacer{} กามจฺฉนฺโท ปหีโน โหติ, พฺยาปาโท ปหีโน โหติ, ถินมิทฺธํ ปหีนํ โหติ, อุทฺธจฺจกุกฺกุจฺจํ ปหีนํ โหติ, วิจิกิจฺฉา ปหีนา โหติ.\csromanspacer{} อิมานิ ปญฺจงฺคานิ ปหีนานิ โหนฺติ.\csromanspacer{} กตเมหิ ปญฺจหงฺเคหิ สมนฺนาคโต โหติ.\csromanspacer{} อเสกฺเขน สีลกฺขนฺเธน สมนฺนาคโต โหติ, อเสกฺเขน สมาธิกฺขนฺเธน สมนฺนาคโต โหติ, อเสกฺเขน ปญฺญากฺขนฺเธน สมนฺนาคโต โหติ, อเสกฺเขน วิมุตฺติกฺขนฺเธน สมนฺนาคโต โหติ, อเสกฺเขน วิมุตฺติญาณทสฺสนกฺขนฺเธน สมนฺนาคโต โหติ.\csromanspacer{} อิเมหิ ปญฺจหงฺเคหิ สมนฺนาคโต โหติ.\csromanspacer{} อิติ ปญฺจงฺควิปฺปหีเน ปญฺจงฺคสมนฺนาคเต ทินฺนํ มหปฺผลนฺติ.\csromanspacer{} อิทมโวจ ภควา ฯเปฯ สตฺถา\csromandash{}}

\csromangathameasure{“อิสฺสตฺตํ พลวีริยญฺจ, ยสฺมิํ วิชฺเชถ มาณเว. \\ ตํ ยุทฺธตฺโถ ภเร ราชา, นาสูรํ ชาติปจฺจยา. \\ ตเถว ขนฺติโสรจฺจํ, ธมฺมา ยสฺมิํ ปติฏฺฐิตา. \\ อริยวุตฺติํ เมธาวิํ, หีนชจฺจมฺปิ ปูชเย. \\ การเย อสฺสเม รมฺเม, วาสเยตฺถ พหุสฺสุเต. \\ ปปญฺจ วิวเน กยิรา, ทุคฺเค สงฺกมนานิ จ. \\ อนฺนํ ปานํ ขาทนียํ, วตฺถเสนาสนานิ จ. \\ ทเทยฺย อุชุภูเตสุ, วิปฺปสนฺเนน เจตสา. \\ ยถา หิ เมโฆ ถนยํ, วิชฺชุมาลี สตกฺกกุ. \\ ถลํ นินฺนญฺจ ปูเรติ, อภิวสฺสํ วสุนฺธรํ. \\ ตเถว สทฺโธ สุตวา, อภิสงฺขจฺจ โภชนํ. \\ วนิพฺพเก ตปฺปยติ, อนฺนปาเนน ปณฺฑิโต. \\ อาโมทมาโน ปกิเรติ, เทถ เทถาติ ภาสติ. \\ ตํ หิสฺส คชฺชิตํ โหติ, เทวสฺเสว ปวสฺสโต. \\ สา ปุญฺญธารา วิปุลา, ทาตารํ อภิวสฺสตี”ติ.}
\csromangathagroup{\csromangathastanza{“อิสฺสตฺตํ\footnote{อิสฺสตฺถํ (สี, สฺยา, กํ)} พลวีริยญฺจ\footnote{พลวิริยญฺจ (สี, สฺยา, กํ, อิ)}, ยสฺมิํ วิชฺเชถ มาณเว. \\ ตํ ยุทฺธตฺโถ ภเร ราชา, นาสูรํ ชาติปจฺจยา.}\csromangathastanza{ตเถว ขนฺติโสรจฺจํ, ธมฺมา ยสฺมิํ ปติฏฺฐิตา. \\ อริยวุตฺติํ เมธาวิํ, หีนชจฺจมฺปิ ปูชเย.}\csromangathastanza{การเย อสฺสเม รมฺเม, วาสเยตฺถ พหุสฺสุเต. \\ ปปญฺจ วิวเน กยิรา, ทุคฺเค สงฺกมนานิ จ.}\csromangathastanza{อนฺนํ ปานํ ขาทนียํ, วตฺถเสนาสนานิ จ. \\ ทเทยฺย อุชุภูเตสุ, วิปฺปสนฺเนน เจตสา.}\csromangathastanza{ยถา หิ เมโฆ ถนยํ, วิชฺชุมาลี สตกฺกกุ. \\ ถลํ นินฺนญฺจ ปูเรติ, อภิวสฺสํ วสุนฺธรํ.}\csromangathastanza{ตเถว สทฺโธ สุตวา, อภิสงฺขจฺจ โภชนํ. \\ วนิพฺพเก ตปฺปยติ, อนฺนปาเนน ปณฺฑิโต.}\csromangathastanza{อาโมทมาโน ปกิเรติ, เทถ เทถาติ ภาสติ. \\ ตํ หิสฺส คชฺชิตํ โหติ, เทวสฺเสว ปวสฺสโต. \\ สา ปุญฺญธารา วิปุลา, ทาตารํ อภิวสฺสตี”ติ.}}

\csromanheader{2}{3. โกสลสํยุตฺต \textbf{5. ปพฺพตูปมสุตฺต}}
\proseitem{136}{\csromanfolio{101}สาวตฺถินิทานํ.\csromanspacer{} เอกมนฺตํ นิสินฺนํ โข ราชานํ ปเสนทิํ โกสลํ ภควา เอตทโวจ “กนฺท กุโต นุ ตฺวํ มหาราช อาคจฺฉสิ ทิวา ทิวสฺสา”ติ.\csromanspacer{} ยานิ ตานิ ภนฺเต รญฺญํ ขตฺติยานํ มุทฺธาวสิตฺตานํ อิสฺสริยมทมตฺตานํ กามเคธปริยุฏฺฐิตานํ ชนปทตฺถาวริยปฺปตฺตานํ มหนฺตํ ปถวิมณฺฑลํ อภิวิชิย อชฺฌาวสนฺตานํ ราชกรณียานิ ภวนฺติ, เตสุ ขฺวาหํ เอตรหิ อุสฺสุกฺกมาปนฺโนติ.}

\prose{ตํ กิํ มญฺญสิ มหาราช, อิธ เต ปุริโส อาคจฺเฉยฺย ปุรตฺถิมาย ทิสาย สทฺธายิโก ปจฺจยิโก, โส ตํ อุปสงฺกมิตฺวา เอวํ วเทยฺย “ยคฺเฆ มหาราช ชาเนยฺยาสิ, อหํ อาคจฺฉามิ ปุรตฺถิมาย ทิสาย, ตตฺถทฺทสํ มหนฺตํ ปพฺพตํ อพฺภสมํ, สพฺเพ ปาเณ นิปฺโปเถนฺโต อาคจฺฉติ.\csromanspacer{} ยํ เต มหาราช กรณียํ, ตํ กโรหี”ติ.\csromanspacer{} อถ ทุติโย ปุริโส อาคจฺเฉยฺย ปจฺฉิมาย ทิสาย ฯเปฯ.\csromanspacer{} อถ ตติโย ปุริโส อาคจฺเฉยฺย อุตฺตราย ทิสาย.\csromanspacer{} อถ จตุตฺโถ ปุริโส อาคจฺเฉยฺย ทกฺขิณาย ทิสาย สทฺธายิโก ปจฺจยิโก, โส ตํ อุปสงฺกมิตฺวา เอวํ วเทยฺย “ยคฺเฆ มหาราช ชาเนยฺยาสิ, อหํ อาคจฺฉามิ ทกฺขิณาย ทิสาย ตตฺถทฺทสํ มหนฺตํ ปพฺพตํ อพฺภสมํ, สพฺเพ ปาเณ นิปฺโปเถนฺโต อาคจฺฉติ.\csromanspacer{} ยํ เต มหาราช กรณียํ, ตํ กโรหี”ติ.\csromanspacer{} เอวรูเป เต มหาราช มหติ มหพฺภเย สมุปฺปนฺเน ทารุเณ มนุสฺสกฺขเย\footnote{มนุสฺสกาเย (ก)} ทุลฺลเภ มนุสฺสตฺเต กิมสฺส กรณียนฺติ.}

\prose{เอวรูเป เม ภนฺเต มหติ มหพฺภเย สมุปฺปนฺเน ทารุเณ มนุสฺสกฺขเย ทุลฺลเภ มนุสฺสตฺเต กิมสฺส กรณียํ อญฺญตฺร ธมฺมจริยาย อญฺญตฺร สมจริยาย อญฺญตฺร กุสลกิริยาย อญฺญตฺร ปุญฺญกิริยายาติ.}

\prose{อาโรเจมิ โข เต มหาราช, ปฏิเวเทมิ โข เต มหาราช, อธิวตฺตติ โข ตํ มหาราช ชรามรณํ, อธิวตฺตมาเน เจ เต มหาราช ชรามรเณ กิมสฺส กรณียนฺติ.\csromanspacer{} อธิวตฺตมาเน จ เม ภนฺเต ชรามรเณ กิมสฺส กรณียํ อญฺญตฺร ธมฺมจริยาย สมจริยาย กุสลกิริยาย ปุญฺญกิริยาย.\csromanspacer{} ยานิ ตานิ ภนฺเต รญฺญํ ขตฺติยานํ มุทฺธาวสิตฺตานํ \csromanfolio{102}อิสฺสริยมทมตฺตานํ กามเคธปริยุฏฺฐิตานํ ชนปทตฺถาวริยปฺปตฺตานํ มหนฺตํ ปถวิมณฺฑลํ อภิวิชิย อชฺฌาวสนฺตานํ หตฺถิยุทฺธานิ ภวนฺติ, เตสมฺปิ ภนฺเต หตฺถิยุทฺธานํ นตฺถิ คติ นตฺถิ วิสโย อธิวตฺตมาเน ชรามรเณ.\csromanspacer{} ยานิปิ ตานิ ภนฺเต รญฺญํ ขตฺติยานํ มุทฺธาวสิตฺตานํ ฯเปฯ อชฺฌาวสนฺตานํ อสฺสยุทฺธานิ ภวนฺติ ฯเปฯ รถยุทฺธานิ ภวนฺติ ฯเปฯ ปตฺติยุทฺธานิ ภวนฺติ, เตสมฺปิ ภนฺเต ปตฺติยุทฺธานํ นตฺถิ คติ นตฺถิ วิสโย อธิวตฺตมาเน ชรามรเณ.\csromanspacer{} สนฺติ โข ปน ภนฺเต อิมสฺมิํ ราชกุเล มนฺติโน มหามตฺตา, เย ปโหนฺติ\footnote{เยสํ โหนฺติ (ก)} อาคเต ปจฺจตฺถิเก มนฺเตหิ เภทยิตุํ, เตสมฺปิ ภนฺเต มนฺตยุทฺธานํ นตฺถิ คติ นตฺถิ วิสโย อธิวตฺตมาเน ชรามรเณ.\csromanspacer{} สํวิชฺชติ โข ปน ภนฺเต อิมสฺมิํ ราชกุเล ปหูตํ หิรญฺญสุวณฺณํ ภูมิคตญฺเจว เวหาสฏฺฐญฺจ, เยน มยํ ปโหม อาคเต ปจฺจตฺถิเก ธเนน อุปลาเปตุํ, เตสมฺปิ ภนฺเต ธนยุทฺธานํ นตฺถิ คติ นตฺถิ วิสโย อธิวตฺตมาเน ชรามรเณ.\csromanspacer{} อธิวตฺตมาเน จ เม ภนฺเต ชรามรเณ กิมสฺส กรณียํ อญฺญตฺร ธมฺมจริยาย สมจริยาย กุสลกิริยาย ปุญฺญกิริยายาติ.}

\prose{เอวเมตํ มหาราช, เอวเมตํ มหาราช, อธิวตฺตมาเน ชรามรเณ กิมสฺส กรณียํ อญฺญตฺร ธมฺมจริยาย สมจริยาย กุสลกิริยาย ปุญฺญกิริยายาติ.\csromanspacer{} อิทมโวจ ภควา ฯเปฯ สตฺถา\csromandash{}}

\csromangathameasure{“ยถาปิ เสลา วิปุลา, นภํ อาหจฺจ ปพฺพตา. \\ สมนฺตานุปริยาเยยฺยุํ, นิปฺโปเถนฺโต จตุทฺทิสา. \\ เอวํ ชรา จ มจฺจุ จ, อธิวตฺตนฺติ ปาณิเน. \\ ขตฺติเย พฺราหฺมเณ เวสฺเส, สุทฺเท จณฺฑาลปุกฺกุเส. \\ น กิญฺจิ ปริวชฺเชติ, สพฺพเมวาภิมทฺทติ. \\ น ตตฺถ หตฺถีนํ ภูมิ, น รถานํ น ปตฺติยา. \\ น จาปิ มนฺตยุทฺเธน, สกฺกา เชตุํ ธเนน วา. \\ ตสฺมา หิ ปณฺฑิโต โปโส, สมฺปสฺสํ อตฺถมตฺตโน. \\ พุทฺเธ ธมฺเม จ สํเฆ จ, ธีโร สทฺธํ นิเวสเย.}
\csromangathagroup{\csromangathastanza{“ยถาปิ เสลา วิปุลา, นภํ อาหจฺจ ปพฺพตา. \\ สมนฺตานุปริยาเยยฺยุํ, นิปฺโปเถนฺโต จตุทฺทิสา.}\csromangathastanza{เอวํ ชรา จ มจฺจุ จ, อธิวตฺตนฺติ ปาณิเน\footnote{ปาณิโน (สี, สฺยา, กํ, อิ)}. \\ ขตฺติเย พฺราหฺมเณ เวสฺเส, สุทฺเท จณฺฑาลปุกฺกุเส.}\csromangathastanza{น กิญฺจิ\footnote{น กญฺจิ (?)} ปริวชฺเชติ, สพฺพเมวาภิมทฺทติ. \\ น ตตฺถ หตฺถีนํ ภูมิ, น รถานํ น ปตฺติยา.}\csromangathastanza{น จาปิ มนฺตยุทฺเธน, สกฺกา เชตุํ ธเนน วา. \\ ตสฺมา หิ ปณฺฑิโต โปโส, สมฺปสฺสํ อตฺถมตฺตโน. \\ พุทฺเธ ธมฺเม จ สํเฆ จ, ธีโร สทฺธํ นิเวสเย.}}

\csromanheader{2}{3. โกสลสํยุตฺต}
\csromangathameasure{โย ธมฺมํ จริ กาเยน, วาจาย อุท เจตสา. \\ อิเธว นํ ปสํสนฺติ, เปจฺจ สคฺเค ปโมทตี”ติ.}
\csromangathagroup{\csromangathastanza{\csromanfolio{103}โย ธมฺมํ จริ\footnote{ธมฺมจารี (สี, สฺยา, กํ, อิ)} กาเยน, วาจาย อุท เจตสา. \\ อิเธว นํ ปสํสนฺติ, เปจฺจ สคฺเค ปโมทตี”ติ.}}

\vspace{\tipitakaparskip}
\csromancenter{ตติโย วคฺโค.}
\titlehead{\textbf{ตสฺสุทฺทานํ}}
\csromangathameasure{ปุคฺคโล อยฺยิกา โลโก, อิสฺสตฺตํ ปพฺพตูปมา. \\ เทสิตํ พุทฺธเสฏฺเฐน, อิมํ โกสลปญฺจกนฺติ. \\ \textbf{โกสลสํยุตฺตํ สมตฺตํ.}}
\csromangathagroup{\csromangathastanza{ปุคฺคโล อยฺยิกา โลโก, อิสฺสตฺตํ\footnote{อิสฺสตฺถํ (สี, สฺยา, กํ)} ปพฺพตูปมา. \\ เทสิตํ พุทฺธเสฏฺเฐน, อิมํ โกสลปญฺจกนฺติ. \\ \textbf{โกสลสํยุตฺตํ สมตฺตํ.}}}

\csromanheader{2}{4. มารสํยุตฺต}
\chapterhead{1. ปฐมวคฺค}
\csromanheader{2}{1. ตโปกมฺมสุตฺต}
\proseitem{137}{\csromanfolio{104}เอวํ เม สุตํ\csromandash{}เอกํ สมยํ ภควา อุรุเวลายํ วิหรติ นชฺชา เนรญฺชราย ตีเร อชปาลนิโคฺรธมูเล ปฐมาภิสมฺพุทฺโธ.\csromanspacer{} อถ โข ภควโต รโหคตสฺส ปฏิสลฺลีนสฺส เอวํ เจตโส ปริวิตกฺโก อุทปาทิ “มุตฺโต วตมฺหิ ตาย ทุกฺกรการิกาย, สาธุ มุตฺโต วตมฺหิ ตาย อนตฺถสํหิตาย ทุกฺกรการิกาย, สาธุ วตมฺหิ มุตฺโต โพธิํ สมชฺฌคนฺติ\footnote{สาธุ ฐิโต สโต โพธิํ สมชฺฌคนฺติ (สี, อิ), สาธุ วตมฺหิ สตฺโต โพธิสมชฺฌคูติ (สฺยา, กํ)}.}

\prose{อถ โข มาโร ปาปิมา ภควโต เจตสา เจโตปริวิตกฺกมญฺญาย เยน ภควา เตนุปสงฺกมิ, อุปสงฺกมิตฺวา ภควนฺตํ คาถาย อชฺฌภาสิ\csromandash{}}

\csromangathameasure{“ตโปกมฺมา อปกฺกมฺม, เย น สุชฺฌนฺติ มาณวา. \\ อสุทฺโธ มญฺญสิ สุทฺโธ, สุทฺธิมคฺคา อปรทฺโธ” ติ.}
\csromangathagroup{\csromangathastanza{“ตโปกมฺมา อปกฺกมฺม, เย น สุชฺฌนฺติ มาณวา. \\ อสุทฺโธ มญฺญสิ สุทฺโธ, สุทฺธิมคฺคา อปรทฺโธ”\footnote{สุทฺธิมคฺคมปรทฺโธ (สี, สฺยา, กํ, อิ)} ติ.}}

\prose{อถ โข ภควา มาโร อยํ ปาปิมา อิติ วิทิตฺวา มารํ ปาปิมนฺตํ คาถาหิ อชฺฌภาสิ\csromandash{}}

\csromangathameasure{“อนตฺถสํหิตํ ญตฺวา, ยํ กิญฺจิ อมรํ ตปํ. \\ สพฺพํ นตฺถาวหํ โหติ, ผิยาริตฺตํว ธมฺมนิ. \\ สีลํ สมาธิ ปญฺญญฺจ, มคฺคํ โพธาย ภาวยํ. \\ ปตฺโตสฺมิ ปรมํ สุทฺธิํ, นิหโต ตฺวมสิ อนฺตกา”ติ.}
\csromangathagroup{\csromangathastanza{“อนตฺถสํหิตํ ญตฺวา, ยํ กิญฺจิ อมรํ ตปํ\footnote{อปรํ ตปํ (ก)}. \\ สพฺพํ นตฺถาวหํ โหติ, ผิยาริตฺตํว ธมฺมนิ\footnote{อปรํ ตปํ (ก)}.}\csromangathastanza{สีลํ สมาธิ ปญฺญญฺจ, มคฺคํ โพธาย ภาวยํ. \\ ปตฺโตสฺมิ ปรมํ สุทฺธิํ, นิหโต ตฺวมสิ อนฺตกา”ติ.}}

\prose{อถ โข มาโร ปาปิมา “ชานาติ มํ ภควา, ชานาติ มํ สุคโต”ติ ทุกฺขี ทุมฺมโน ตตฺเถวนฺตรธายีติ.}

\csromanheader{2}{4. มารสํยุตฺต \textbf{2. หตฺถิราชวณฺณสุตฺต}}
\proseitem{138}{\csromanfolio{105}เอวํ เม สุตํ\csromandash{}เอกํ สมยํ ภควา อุรุเวลายํ วิหรติ นชฺชา เนรญฺชราย ตีเร อชปาลนิโคฺรธมูเล ปฐมาภิสมฺพุทฺโธ.\csromanspacer{} เตน โข ปน สมเยน ภควา รตฺตนฺธการติมิสายํ อพฺโภกาเส นิสินฺโน โหติ, เทโว จ เอกเมกํ ผุสายติ.\csromanspacer{} อถ โข มาโร ปาปิมา ภควโต ภยํ ฉมฺภิตตฺตํ โลมหํสํ อุปฺปาเทตุกาโม มหนฺตํ หตฺถิราชวณฺณํ อภินิมฺมินิตฺวา เยน ภควา เตนุปสงฺกมิ.\csromanspacer{} เสยฺยถาปิ นาม มหา-อริฏฺฐโก มณิ, เอวมสฺส สีสํ โหติ.\csromanspacer{} เสยฺยถาปิ นาม สุทฺธํ รูปิยํ, เอวมสฺส ทนฺตา โหนฺติ.\csromanspacer{} เสยฺยถาปิ นาม มหตี นงฺคลีสา\footnote{นงฺคลสีสา (อิ, ก)}, เอวมสฺส โสณฺโฑ โหติ.\csromanspacer{} อถ โข ภควา มาโร อยํ ปาปิมา อิติ วิทิตฺวา มารํ ปาปิมนฺตํ คาถาย อชฺฌภาสิ\csromandash{}}

\csromangathameasure{“สํสรํ ทีฆมทฺธานํ, วณฺณํ กตฺวา สุภาสุภํ. \\ อลํ เต เตน ปาปิม, นิหโต ตฺวมสิ อนฺตกา”ติ.}
\csromangathagroup{\csromangathastanza{“สํสรํ ทีฆมทฺธานํ, วณฺณํ กตฺวา สุภาสุภํ. \\ อลํ เต เตน ปาปิม, นิหโต ตฺวมสิ อนฺตกา”ติ.}}

\prose{อถ โข มาโร ปาปิมา “ชานาติ มํ ภควา, ชานาติ มํ สุคโต”ติ ทุกฺขี ทุมฺมโน ตตฺเถวนฺตรธายีติ.}

\csromanheader{2}{3. สุภสุตฺต}
\proseitem{139}{เอวํ เม สุตํ\csromandash{}เอกํ สมยํ ภควา อุรุเวลายํ วิหรติ นชฺชา เนรญฺชราย ตีเร อชปาลนิโคฺรธมูเล ปฐมาภิสมฺพุทฺโธ.\csromanspacer{} เตน โข ปน สมเยน ภควา รตฺตนฺธการติมิสายํ อพฺโภกาเส นิสินฺโน โหติ, เทโว จ เอกเมกํ ผุสายติ.\csromanspacer{} อถ โข มาโร ปาปิมา ภควโต ภยํ ฉมฺภิตตฺตํ โลมหํสํ อุปฺปาเทตุกาโม เยน ภควา เตนุปสงฺกมิ, อุปสงฺกมิตฺวา ภควโต อวิทูเร อุจฺจาวจา วณฺณนิภา อุปทํเสติ สุภา เจว อสุภา จ.\csromanspacer{} อถ โข ภควา มาโร อยํ ปาปิมา อิติ วิทิตฺวา มารํ ปาปิมนฺตํ คาถาหิ อชฺฌภาสิ\csromandash{}}

\csromangathameasure{“สํสรํ ทีฆมทฺธานํ, วณฺณํ กตฺวา สุภาสุภํ. \\ อลํ เต เตน ปาปิม, นิหโต ตฺวมสิ อนฺตก. \\ เย จ กาเยน วาจาย, มนสา จ สุสํวุตา. \\ น เต มารวสานุคา, น เต มารสฺส พทฺธคู” ติ.}
\csromangathagroup{\csromangathastanza{“สํสรํ ทีฆมทฺธานํ, วณฺณํ กตฺวา สุภาสุภํ. \\ อลํ เต เตน ปาปิม, นิหโต ตฺวมสิ อนฺตก.}\csromangathastanza{\csromanfolio{106}เย จ กาเยน วาจาย, มนสา จ สุสํวุตา. \\ น เต มารวสานุคา, น เต มารสฺส พทฺธคู”\footnote{พทฺธภู (ก), ปจฺจคู (สี, สฺยา, กํ, อิ)} ติ.}}

\prose{อถ โข มาโร ฯเปฯ ตตฺเถวนฺตรธายีติ.}

\csromanheader{2}{4. ปฐม มารปาสสุตฺต}
\proseitem{140}{เอวํ เม สุตํ\csromandash{}เอกํ สมยํ ภควา พาราณสิยํ วิหรติ อิสิปตเน มิคทาเย.\csromanspacer{} ตตฺร โข ภควา ภิกฺขู อามนฺเตสิ “ภิกฺขโว”ติ. “ภทนฺเต”ติ เต ภิกฺขู ภควโต ปจฺจสฺโสสุํ.\csromanspacer{} ภควา เอตทโวจ\csromandash{}}

\prose{มยฺหํ โข ภิกฺขเว โยนิโส มนสิการา โยนิโส สมฺมปฺปธานา อนุตฺตรา วิมุตฺติ อนุปฺปตฺตา, อนุตฺตรา วิมุตฺติ สจฺฉิกตา.\csromanspacer{} ตุเมฺหปิ ภิกฺขเว โยนิโส มนสิการา โยนิโส สมฺมปฺปธานา อนุตฺตรํ วิมุตฺติํ อนุปาปุณาถ, อนุตฺตรํ วิมุตฺติํ สจฺฉิกโรถาติ.\csromanspacer{} อถ โข มาโร ปาปิมา เยน ภควา เตนุปสงฺกมิ, อุปสงฺกมิตฺวา ภควนฺตํ คาถาย อชฺฌภาสิ\csromandash{}}

\csromangathameasure{“พทฺโธสิ มารปาเสน, เย ทิพฺพา เย จ มานุสา. \\ มารพนฺธนพทฺโธสิ, น เม สมณ โมกฺขสี”ติ. \\ มุตฺตาหํ มารปาเสน, เย ทิพฺพา เย จ มานุสา. \\ มารพนฺธนมุตฺโตมฺหิ, นิหโต ตฺวมสิ อนฺตกาติ.}
\csromangathagroup{\csromangathastanza{“พทฺโธสิ มารปาเสน, เย ทิพฺพา เย จ มานุสา. \\ มารพนฺธนพทฺโธสิ, น เม สมณ โมกฺขสี”ติ.}\csromangathastanza{มุตฺตาหํ\footnote{มุตฺโตหํ (สี, สฺยา, กํ, อิ)} มารปาเสน, เย ทิพฺพา เย จ มานุสา. \\ มารพนฺธนมุตฺโตมฺหิ, นิหโต ตฺวมสิ อนฺตกาติ.}}

\prose{อถ โข มาโร ปาปิมา ฯเปฯ ตตฺเถวนฺตรธายีติ.}

\csromanheader{2}{5. ทุติยมารปาสสุตฺต}
\proseitem{141}{เอกํ สมยํ ภควา พาราณสิยํ วิหรติ อิสิปตเน มิคทาเย.\csromanspacer{} ตตฺร โข ภควา ภิกฺขู อามนฺเตสิ “ภิกฺขโว”ติ. “ภทนฺเต”ติ เต ภิกฺขู ภควโต ปจฺจสฺโสสุํ.\csromanspacer{} ภควา เอตทโวจ\csromandash{}}

\prose{มุตฺตาหํ ภิกฺขเว สพฺพปาเสหิ เย ทิพฺพา เย จ มานุสา.\csromanspacer{} ตุเมฺหปิ ภิกฺขเว มุตฺตา สพฺพปาเสหิ เย ทิพฺพา เย จ มานุสา.\csromanspacer{} จรถ ภิกฺขเว จาริกํ พหุชนหิตาย พหุชนสุขาย โลกานุกมฺปาย อตฺถาย หิตาย สุขาย เทวมนุสฺสานํ, มา เอเกน เทฺว อคมิตฺถ, เทเสถ ภิกฺขเว ธมฺมํ อาทิกลฺยาณํ มชฺเฌกลฺยาณํ ปริโยสานกลฺยาณํ สาตฺถํ}

\vspace{\tipitakaparskip}
\csromancenter{มารสํยุตฺต สพฺยญฺชนํ เกวลปริปุณฺณํ ปริสุทฺธํ พฺรหฺมจริยํ ปกาเสถ, สนฺติ สตฺตา อปฺปรชกฺขชาติกา อสฺสวนตา ธมฺมสฺส ปริหายนฺติ, ภวิสฺสนฺติ ธมฺมสฺส อญฺญาตาโร.\csromanspacer{} อหมฺปิ ภิกฺขเว เยน อุรุเวลา เสนานิคโม เตนุปสงฺกมิสฺสามิ ธมฺมเทสนายาติ.\csromanspacer{} อถ โข มาโร ปาปิมา เยน ภควา เตนุปสงฺกมิ, อุปสงฺกมิตฺวา ภควนฺตํ คาถาย อชฺฌภาสิ\csromandash{}}
\vspace{0.5\baselineskip}
\csromangathameasure{“พทฺโธสิ สพฺพปาเสหิ, เย ทิพฺพา เย จ มานุสา. \\ มหาพนฺธนพทฺโธสิ, น เม สมณ โมกฺขสี”ติ. \\ มุตฺตาหํ สพฺพปาเสหิ, เย ทิพฺพา เย จ มานุสา, \\ มหาพนฺธนมุตฺโตมฺหิ, \\ นิหโต ตฺวมสิ อนฺตกาติ.}
\csromangathagroup{\csromangathastanza{\csromanfolio{107}“พทฺโธสิ สพฺพปาเสหิ, เย ทิพฺพา เย จ มานุสา. \\ มหาพนฺธนพทฺโธสิ, น เม สมณ โมกฺขสี”ติ.}\csromangathastanza{มุตฺตาหํ สพฺพปาเสหิ, เย ทิพฺพา เย จ มานุสา, \\ มหาพนฺธนมุตฺโตมฺหิ, \\ นิหโต ตฺวมสิ อนฺตกาติ.}}

\prose{อถ โข มาโร ปาปิมา ฯเปฯ ตตฺเถวนฺตรธายีติ.}

\csromanheader{2}{6. สปฺปสุตฺต}
\proseitem{142}{เอวํ เม สุตํ\csromandash{}เอกํ สมยํ ภควา ราชคเห วิหรติ เวฬุวเน กลนฺทกนิวาเป.\csromanspacer{} เตน โข ปน สมเยน ภควา รตฺตนฺธการติมิสายํ อพฺโภกาเส นิสินฺโน โหติ, เทโว จ เอกเมกํ ผุสายติ.}

\prose{อถ โข มาโร ปาปิมา ภควโต ภยํ ฉมฺภิตตฺตํ โลมหํสํ อุปฺปาเทตุกาโม มหนฺตํ สปฺปราชวณฺณํ อภินิมฺมินิตฺวา เยน ภควา เตนุปสงฺกมิ.\csromanspacer{} เสยฺยถาปิ นาม มหตี เอกรุกฺขิกา นาวา, เอวมสฺส กาโย โหติ.\csromanspacer{} เสยฺยถาปิ นาม มหนฺตํ โสณฺฑิกากิฬญฺชํ, เอวมสฺส ผโณ โหติ.\csromanspacer{} เสยฺยถาปิ นาม มหตี โกสลิกา กํสปาติ, เอวมสฺส อกฺขีนิ ภวนฺติ.\csromanspacer{} เสยฺยถาปิ นาม เทเว คฬคฬายนฺเต วิชฺชุลฺลตา นิจฺฉรนฺติ, เอวมสฺส มุขโต ชิวฺหา นิจฺฉรติ.\csromanspacer{} เสยฺยถาปิ นาม กมฺมารคคฺคริยา ธมมานาย สทฺโท โหติ, เอวมสฺส อสฺสาสปสฺสาสานํ สทฺโท โหติ.}

\prose{อถ โข ภควา มาโร อยํ ปาปิมา อิติ วิทิตฺวา มารํ ปาปิมนฺตํ คาถาหิ อชฺฌภาสิ\csromandash{}}

\csromangathameasure{“โย สุญฺญเคหานิ เสวติ, \\ เสยฺยา โส มุนิ อตฺตสญฺญโต. \\ โวสฺสชฺช จเรยฺย ตตฺถ โส, \\ ปติรูปํ หิ ตถาวิธสฺส ตํ. \\ จรกา พหู เภรวา พหู, \\ อโถ ฑํสสรีสปา พหู. \\ โลมมฺปิ น ตตฺถ อิญฺชเย, \\ สุญฺญาคารคโต มหามุนิ. \\ นภํ ผเลยฺย ปถวี จเลยฺย, \\ สพฺเพปิ ปาณา อุท สนฺตเสยฺยุํ. \\ สลฺลมฺปิ เจ อุรสิ ปกปฺปเยยฺยุํ, \\ อุปธีสุ ตาณํ น กโรนฺติ พุทฺธา”ติ.}
\csromangathagroup{\csromangathastanza{“โย สุญฺญเคหานิ เสวติ, \\ เสยฺยา โส มุนิ อตฺตสญฺญโต. \\ โวสฺสชฺช จเรยฺย ตตฺถ โส, \\ ปติรูปํ หิ ตถาวิธสฺส ตํ.}\csromangathastanza{\csromanfolio{108}จรกา พหู เภรวา พหู, \\ อโถ ฑํสสรีสปา\footnote{qอํส สิริํสปา (สี, สฺยา, กํ, อิ)} พหู. \\ โลมมฺปิ น ตตฺถ อิญฺชเย, \\ สุญฺญาคารคโต มหามุนิ.}\csromangathastanza{นภํ ผเลยฺย ปถวี จเลยฺย, \\ สพฺเพปิ ปาณา อุท สนฺตเสยฺยุํ. \\ สลฺลมฺปิ เจ อุรสิ ปกปฺปเยยฺยุํ, \\ อุปธีสุ ตาณํ น กโรนฺติ พุทฺธา”ติ.}}

\prose{อถ โข มาโร ปาปิมา “ชานาติ มํ ภควา, ชานาติ มํ สุคโต”ติ ทุกฺขี ทุมฺมโน ตตฺเถวนฺตรธายีติ.}

\csromanheader{2}{7. สุปติสุตฺต}
\proseitem{143}{เอกํ สมยํ ภควา ราชคเห วิหรติ เวฬุวเน กลนฺทกนิวาเป.\csromanspacer{} อถ โข ภควา พหุเทวรตฺติํ อพฺโภกาเส จงฺกมิตฺวา รตฺติยา ปจฺจูสสมยํ ปาเท ปกฺขาเลตฺวา วิหารํ ปวิสิตฺวา ทกฺขิเณน ปสฺเสน สีหเสยฺยํ กปฺเปสิ ปาเท ปาทํ อจฺจาธาย สโต สมฺปชาโน อุฏฺฐานสญฺญํ มนสิ กริตฺวา.\csromanspacer{} อถ โข มาโร ปาปิมา เยน ภควา เตนุปสงฺกมิ, อุปสงฺกมิตฺวา ภควนฺตํ คาถาย อชฺฌภาสิ\csromandash{}}

\csromangathameasure{“กิํ โสปฺปสิ กิํ นุ โสปฺปสิ, \\ กิมิทํ โสปฺปสิ ทุพฺภโค วิย. \\ สุญฺญมคารนฺติ โสปฺปสิ, \\ กิมิทํ โสปฺปสิ สูริเย อุคฺคเต”ติ. \\ ยสฺส ชาลินี วิสตฺติกา, \\ ตณฺหา นตฺถิ กุหิญฺจิ เนตเว. \\ สพฺพูปธิปริกฺขยา พุทฺโธ, \\ โสปฺปติ กิํ ตเวตฺถ มาราติ.}
\csromangathagroup{\csromangathastanza{“กิํ โสปฺปสิ กิํ นุ โสปฺปสิ, \\ กิมิทํ โสปฺปสิ ทุพฺภโค\footnote{ทุพฺภโต (สฺยา, กํ), ทุพฺภโย (อิ)} วิย. \\ สุญฺญมคารนฺติ โสปฺปสิ, \\ กิมิทํ โสปฺปสิ สูริเย อุคฺคเต”ติ.}\csromangathastanza{ยสฺส ชาลินี วิสตฺติกา, \\ ตณฺหา นตฺถิ กุหิญฺจิ เนตเว. \\ สพฺพูปธิปริกฺขยา พุทฺโธ, \\ โสปฺปติ กิํ ตเวตฺถ มาราติ.}}

\prose{อถ โข มาโร ปาปิมา ฯเปฯ ตตฺเถวนฺตรธายีติ.}

\csromanheader{2}{4. มารสํยุตฺต \textbf{8. นนฺทติสุตฺต}}
\proseitem{144}{\csromanfolio{109}เอวํ เม สุตํ\csromandash{}เอกํ สมยํ ภควา สาวตฺถิยํ วิหรติ เชตวเน อนาถปิณฺฑิกสฺส อาราเม.\csromanspacer{} อถ โข มาโร ปาปิมา เยน ภควา เตนุปสงฺกมิ, อุปสงฺกมิตฺวา ภควโต สนฺติเก อิมํ คาถํ อภาสิ\csromandash{}}

\csromangathameasure{“นนฺทติ ปุตฺเตหิ ปุตฺติมา, โคมา โคภิ ตเถว นนฺทติ. \\ อุปธีหิ นรสฺส นนฺทนา, น หิ โส นนฺทติ โย นิรูปธี”ติ. \\ โสจติ ปุตฺเตหิ ปุตฺติมา, โคมา โคภิ ตเถว โสจติ. \\ อุปธีหิ นรสฺส โสจนา, น หิ โส โสจติ โย นิรูปธีติ.}
\csromangathagroup{\csromangathastanza{“นนฺทติ ปุตฺเตหิ ปุตฺติมา, โคมา โคภิ ตเถว นนฺทติ. \\ อุปธีหิ นรสฺส นนฺทนา, น หิ โส นนฺทติ โย นิรูปธี”ติ.}\csromangathastanza{โสจติ ปุตฺเตหิ ปุตฺติมา, โคมา โคภิ ตเถว โสจติ. \\ อุปธีหิ นรสฺส โสจนา, น หิ โส โสจติ โย นิรูปธีติ.}}

\prose{อถ โข มาโร ปาปิมา “ชานาติ มํ ภควา, ชานาติ มํ สุคโต”ติ ทุกฺขี ทุมฺมโน ตตฺเถวนฺตรธายีติ.}

\csromanheader{2}{9. ปฐม-อายุสุตฺต}
\proseitem{145}{เอวํ เม สุตํ\csromandash{}เอกํ สมยํ ภควา ราชคเห วิหรติ เวฬุวเน กลนฺทกนิวาเป.\csromanspacer{} ตตฺร โข ภควา ภิกฺขู อามนฺเตสิ “ภิกฺขโว”ติ. “ภทนฺเต”ติ เต ภิกฺขู ภควโต ปจฺจสฺโสสุํ.\csromanspacer{} ภควา เอตทโวจ\csromandash{}}

\prose{อปฺปมิทํ ภิกฺขเว มนุสฺสานํ อายุ, คมนีโย สมฺปราโย กตฺตพฺพํ กุสลํ, จริตพฺพํ พฺรหฺมจริยํ, นตฺถิ ชาตสฺส อมรณํ.\csromanspacer{} โย ภิกฺขเว จิรํ ชีวติ, โส วสฺสสตํ อปฺปํ วา ภิยฺโยติ.}

\prose{อถ โข มาโร ปาปิมา เยน ภควา เตนุปสงฺกมิ, อุปสงฺกมิตฺวา ภควนฺตํ คาถาย อชฺฌภาสิ\csromandash{}}

\csromangathameasure{“ทีฆมายุ มนุสฺสานํ, น นํ หีเฬ สุโปริโส. \\ จเรยฺย ขีรมตฺโตว, นตฺถิ มจฺจุสฺส อาคโม”ติ. \\ อปฺปมายุ มนุสฺสานํ, หีเฬยฺย นํ สุโปริโส. \\ จเรยฺยาทิตฺตสีโสว, นตฺถิ มจฺจุสฺส นาคโมติ.}
\csromangathagroup{\csromangathastanza{“ทีฆมายุ มนุสฺสานํ, น นํ หีเฬ สุโปริโส. \\ จเรยฺย ขีรมตฺโตว, นตฺถิ มจฺจุสฺส อาคโม”ติ.}\csromangathastanza{อปฺปมายุ มนุสฺสานํ, หีเฬยฺย นํ สุโปริโส. \\ จเรยฺยาทิตฺตสีโสว, นตฺถิ มจฺจุสฺส นาคโมติ.}}

\prose{อถ โข มาโร ฯเปฯ ตตฺเถวนฺตรธายีติ.}

\csromanheader{2}{10. ทุติย-อายุสุตฺต}
\proseitem{146}{\csromanfolio{110}เอวํ เม สุตํ\csromandash{}เอกํ สมยํ ภควา ราชคเห วิหรติ เวฬุวเน กลนฺทกนิวาเป.\csromanspacer{} ตตฺร โข ภควา ฯเปฯ เอตทโวจ\csromandash{}}

\prose{อปฺปมิทํ ภิกฺขเว มนุสฺสานํ อายุ, คมนีโย สมฺปราโย, กตฺตพฺพํ กุสลํ, จริตพฺพํ พฺรหฺมจริยํ, นตฺถิ ชาตสฺส อมรณํ.\csromanspacer{} โย ภิกฺขเว จิรํ ชีวติ, โส วสฺสสตํ อปฺปํ วา ภิยฺโยติ.}

\prose{อถ โข มาโร ปาปิมา เยน ภควา เตนุปสงฺกมิ, อุปสงฺกมิตฺวา ภควนฺตํ คาถาย อชฺฌภาสิ\csromandash{}}

\csromangathameasure{“นาจฺจยนฺติ อโหรตฺตา, ชีวิตํ นูปรุชฺฌติ. \\ อายุ อนุปริยายติ มจฺจานํ, เนมีว รถกุพฺพรนฺ”ติ. \\ อจฺจยนฺติ อโหรตฺตา, ชีวิตํ อุปรุชฺฌติ. \\ อายุ ขียติ มจฺจานํ, กุนฺนทีนํว โอทกนฺติ.}
\csromangathagroup{\csromangathastanza{“นาจฺจยนฺติ อโหรตฺตา, ชีวิตํ นูปรุชฺฌติ. \\ อายุ อนุปริยายติ มจฺจานํ, เนมีว รถกุพฺพรนฺ”ติ.}\csromangathastanza{อจฺจยนฺติ อโหรตฺตา, ชีวิตํ อุปรุชฺฌติ. \\ อายุ ขียติ มจฺจานํ, กุนฺนทีนํว โอทกนฺติ.}}

\prose{อถ โข มาโร ปาปิมา “ชานาติ มํ ภควา, ชานาติ มํ สุคโต”ติ ทุกฺขี ทุมฺมโน ตตฺเถวนฺตรธายีติ.}

\vspace{\tipitakaparskip}
\csromancenter{ปฐโม วคฺโค.}
\titlehead{\textbf{ตสฺสุทฺทานํ}}
\csromangathameasure{ตโปกมฺมญฺจ นาโค จ, สุภํ ปาเสน เต ทุเว. \\ สปฺโป สุปติ นนฺทนํ, อายุนา อปเร ทุเวติ.}
\csromangathagroup{\csromangathastanza{ตโปกมฺมญฺจ นาโค จ, สุภํ ปาเสน เต ทุเว. \\ สปฺโป สุปติ นนฺทนํ, อายุนา อปเร ทุเวติ.}}

\chapterhead{2. \textbf{ทุติยวคฺค}}
\csromanheader{2}{1. ปาสาณสุตฺต}
\proseitem{147}{เอกํ สมยํ ภควา ราชคเห วิหรติ คิชฺฌกูเฏ ปพฺพเต.\csromanspacer{} เตน โข ปน สมเยน ภควา รตฺตนฺธการติมิสายํ อพฺโภกาเส นิสินฺโน โหติ, เทโว จ เอกเมกํ ผุสายติ.\csromanspacer{} อถ โข มาโร ปาปิมา ภควโต ภยํ ฉมฺภิตตฺตํ โลมหํสํ อุปฺปาเทตุกาโม เยน}

\vspace{\tipitakaparskip}
\csromancenter{มารสํยุตฺต ภควา เตนุปสงฺกมิ, อุปสงฺกมิตฺวา ภควโต อวิทูเร มหนฺเต ปาสาเณ ปทาเลสิ.}
\prose{\csromanfolio{111}อถ โข ภควา มาโร อยํ ปาปิมา อิติ วิทิตฺวา มารํ ปาปิมนฺตํ คาถาย อชฺฌภาสิ\csromandash{}}

\csromangathameasure{“สเจปิ เกวลํ สพฺพํ, คิชฺฌกูฏํ จเลสฺสสิ. \\ เนว สมฺมาวิมุตฺตานํ, พุทฺธานํ อตฺถิ อิญฺชิตนฺ”ติ.}
\csromangathagroup{\csromangathastanza{“สเจปิ เกวลํ สพฺพํ, คิชฺฌกูฏํ จเลสฺสสิ\footnote{คเฬยฺยสิ (สฺยา, กํ), จเลยฺยาสิ (ก)}. \\ เนว สมฺมาวิมุตฺตานํ, พุทฺธานํ อตฺถิ อิญฺชิตนฺ”ติ.}}

\prose{อถ โข มาโร ปาปิมา “ชานาติ มํ ภควา, ชานาติ มํ สุคโต”ติ ทุกฺขี ทุมฺมโน ตตฺเถวนฺตรธายีติ.}

\csromanheader{2}{2. กินฺนุสีหสุตฺต}
\proseitem{148}{เอกํ สมยํ ภควา สาวตฺถิยํ วิหรติ เชตวเน อนาถปิณฺฑิกสฺส อาราเม.\csromanspacer{} เตน โข ปน สมเยน ภควา มหติยา ปริสาย ปริวุโต ธมฺมํ เทเสติ.}

\prose{อถ โข มารสฺส ปาปิมโต เอตทโหสิ “อยํ โข สมโณ โคตโม มหติยา ปริสาย ปริวุโต ธมฺมํ เทเสติ, ยํนูนาหํ เยน สมโณ โคตโม เตนุปสงฺกเมยฺยํ วิจกฺขุกมฺมายา”ติ.\csromanspacer{} อถ โข มาโร ปาปิมา เยน ภควา เตนุปสงฺกมิ, อุปสงฺกมิตฺวา ภควนฺตํ คาถาย อชฺฌภาสิ\csromandash{}}

\csromangathameasure{“กินฺนุ สีโหว นทสิ, ปริสายํ วิสารโท. \\ ปฏิมลฺโล หิ เต อตฺถิ, วิชิตาวี นุ มญฺญสี”ติ. \\ นทนฺติ เว มหาวีรา, ปริสาสุ วิสารทา. \\ ตถาคตา พลปฺปตฺตา, ติณฺณา โลเก วิสตฺติกนฺติ.}
\csromangathagroup{\csromangathastanza{“กินฺนุ สีโหว นทสิ, ปริสายํ วิสารโท. \\ ปฏิมลฺโล หิ เต อตฺถิ, วิชิตาวี นุ มญฺญสี”ติ.}\csromangathastanza{นทนฺติ เว มหาวีรา, ปริสาสุ วิสารทา. \\ ตถาคตา พลปฺปตฺตา, ติณฺณา โลเก วิสตฺติกนฺติ.}}

\prose{อถ โข มาโร ปาปิมา “ชานาติ มํ ภควา, ชานาติ มํ สุคโต”ติ ทุกฺขี ทุมฺมโน ตตฺเถวนฺตรธายีติ.}

\csromanheader{2}{3. สกลิกสุตฺต}
\proseitem{149}{เอวํ เม สุตํ\csromandash{}เอกํ สมยํ ภควา ราชคเห วิหรติ มทฺทกุจฺฉิสฺมิํ มิคทาเย.\csromanspacer{} เตน โข ปน สมเยน ภควโต ปาโท \csromanfolio{112}สกลิกาย ขโต โหติ, ภุสา สุทํ ภควโต เวทนา วตฺตนฺติ สารีริกา ทุกฺขา ติพฺพา ขรา กฏุกา อสาตา อมนาปา.\csromanspacer{} ตา สุทํ ภควา สโต สมฺปชาโน อธิวาเสติ อวิหญฺญมาโน.\csromanspacer{} อถ โข ภควา จตุคฺคุณํ สงฺฆาฏิํ ปญฺญเปตฺวา ทกฺขิเณน ปสฺเสน สีหเสยฺยํ กปฺเปสิ ปาเท ปาทํ อจฺจาธาย สโต สมฺปชาโน.\csromanspacer{} อถ โข มาโร ปาปิมา เยน ภควา เตนุปสงฺกมิ, อุปสงฺกมิตฺวา ภควนฺตํ คาถาย อชฺฌภาสิ\csromandash{}}

\csromangathameasure{“มนฺทิยา นุ โข เสสิ อุทาหุ กาเวยฺยมตฺโต, \\ อตฺถา นุ เต สมฺปจุรา น สนฺติ. \\ เอโก วิวิตฺเต สยนาสนมฺหิ, \\ นิทฺทามุโข กิมิทํ โสปฺปเส วา”ติ. \\ น มนฺทิยา สยามิ นาปิ กาเวยฺยมตฺโต, \\ อตฺถํ สเมจฺจาหมเปตโสโก. \\ เอโก วิวิตฺเต สยนาสนมฺหิ, \\ สยามหํ สพฺพภูตานุกมฺปี. \\ เยสมฺปิ สลฺลํ อุรสิ ปวิฏฺฐํ, \\ มุหุํ มุหุํ หทยํ เวธมานํ. \\ เตปีธ โสปฺปํ ลภเร สสลฺลา, \\ ตสฺมา อหํ น สุเป วีตสลฺโล. \\ ชคฺคํ น สงฺเก นปิ เภมิ โสตฺตุํ, \\ รตฺตินฺทิวา นานุตปนฺติ มามํ. \\ หานิํ น ปสฺสามิ กุหิญฺจิ โลเก, \\ ตสฺมา สุเป สพฺพภูตานุกมฺปีติ.}
\csromangathagroup{\csromangathastanza{“มนฺทิยา นุ โข เสสิ อุทาหุ กาเวยฺยมตฺโต, \\ อตฺถา นุ เต สมฺปจุรา น สนฺติ. \\ เอโก วิวิตฺเต สยนาสนมฺหิ, \\ นิทฺทามุโข กิมิทํ โสปฺปเส วา”ติ.}\csromangathastanza{น มนฺทิยา สยามิ นาปิ กาเวยฺยมตฺโต, \\ อตฺถํ สเมจฺจาหมเปตโสโก. \\ เอโก วิวิตฺเต สยนาสนมฺหิ, \\ สยามหํ สพฺพภูตานุกมฺปี.}\csromangathastanza{เยสมฺปิ สลฺลํ อุรสิ ปวิฏฺฐํ, \\ มุหุํ มุหุํ หทยํ เวธมานํ. \\ เตปีธ โสปฺปํ ลภเร สสลฺลา, \\ ตสฺมา อหํ น สุเป วีตสลฺโล.}\csromangathastanza{ชคฺคํ น สงฺเก นปิ เภมิ โสตฺตุํ, \\ รตฺตินฺทิวา นานุตปนฺติ มามํ. \\ หานิํ น ปสฺสามิ กุหิญฺจิ โลเก, \\ ตสฺมา สุเป สพฺพภูตานุกมฺปีติ.}}

\prose{อถ โข มาโร ปาปิมา “ชานาติ มํ ภควา, ชานาติ มํ สุคโต”ติ ทุกฺขี ทุมฺมโน ตตฺเถวนฺตรธายีติ.}

\csromanheader{2}{4. ปติรูปสุตฺต}
\proseitem{150}{เอกํ สมยํ ภควา โกสเลสุ วิหรติ เอกสาลายํ พฺราหฺมณคาเม.\csromanspacer{} เตน โข ปน สมเยน ภควา มหติยา คิหิปริสาย ปริวุโต ธมฺมํ เทเสติ.}

\csromanheader{2}{4. มารสํยุตฺต}
\prose{\csromanfolio{113}อถ โข มารสฺส ปาปิมโต เอตทโหสิ “อยํ โข สมโณ โคตโม มหติยา คิหิปริสาย ปริวุโต ธมฺมํ เทเสติ, ยํนูนาหํ เยน สมโณ โคตโม เตนุปสงฺกเมยฺยํ วิจกฺขุกมฺมายา”ติ.\csromanspacer{} อถ โข มาโร ปาปิมา เยน ภควา เตนุปสงฺกมิ, อุปสงฺกมิตฺวา ภควนฺตํ คาถาย อชฺฌภาสิ\csromandash{}}

\csromangathameasure{“เนตํ ตว ปติรูปํ, ยทญฺญมนุสาสสิ. \\ อนุโรธวิโรเธสุ, มา สชฺชิตฺโถ ตทาจรนฺ”ติ. \\ หิตานุกมฺปี สมฺพุทฺโธ, ยทญฺญมนุสาสติ. \\ อนุโรธวิโรเธหิ, วิปฺปมุตฺโต ตถาคโตติ.}
\csromangathagroup{\csromangathastanza{“เนตํ ตว ปติรูปํ, ยทญฺญมนุสาสสิ. \\ อนุโรธวิโรเธสุ, มา สชฺชิตฺโถ ตทาจรนฺ”ติ.}\csromangathastanza{หิตานุกมฺปี สมฺพุทฺโธ, ยทญฺญมนุสาสติ. \\ อนุโรธวิโรเธหิ, วิปฺปมุตฺโต ตถาคโตติ.}}

\prose{อถ โข มาโร ปาปิมา “ชานาติ มํ ภควา, ชานาติ มํ สุคโต”ติ ทุกฺขี ทุมฺมโน ตตฺเถวนฺตรธายีติ.}

\csromanheader{2}{5. มานสสุตฺต}
\proseitem{151}{เอวํ เม สุตํ\csromandash{}เอกํ สมยํ ภควา สาวตฺถิยํ วิหรติ เชตวเน อนาถปิณฺฑิกสฺส อาราเม.\csromanspacer{} อถ โข มาโร ปาปิมา เยน ภควา เตนุปสงฺกมิ, อุปสงฺกมิตฺวา ภควนฺตํ คาถาย อชฺฌภาสิ\csromandash{}}

\csromangathameasure{“อนฺตลิกฺขจโร ปาโส, ยฺวายํ จรติ มานโส. \\ เตน ตํ พาธยิสฺสามิ, น เม สมณ โมกฺขสี”ติ. \\ รูปา สทฺทา รสา คนฺธา, โผฏฺฐพฺพา จ มโนรมา. \\ เอตฺถ เม วิคโต ฉนฺโท, นิหโต ตฺวมสิ อนฺตกาติ.}
\csromangathagroup{\csromangathastanza{“อนฺตลิกฺขจโร ปาโส, ยฺวายํ จรติ มานโส. \\ เตน ตํ พาธยิสฺสามิ, น เม สมณ โมกฺขสี”ติ.}\csromangathastanza{รูปา สทฺทา รสา คนฺธา, โผฏฺฐพฺพา จ มโนรมา. \\ เอตฺถ เม วิคโต ฉนฺโท, นิหโต ตฺวมสิ อนฺตกาติ.}}

\prose{อถ โข มาโร ปาปิมา “ชานาติ มํ ภควา, ชานาติ มํ สุคโต”ติ ทุกฺขี ทุมฺมโน ตตฺเถวนฺตรธายีติ.}

\csromanheader{2}{6. ปตฺตสุตฺต}
\proseitem{152}{สาวตฺถินิทานํ.\csromanspacer{} เตน โข ปน สมเยน ภควา ปญฺจนฺนํ อุปาทานกฺขนฺธานํ อุปาทาย ภิกฺขูนํ ธมฺมิยา กถาย สนฺทสฺเสติ สมาทเปติ\footnote{สมาทาเปติ (?)} สมุตฺเตเชติ สมฺปหํเสติ.\csromanspacer{} เต จ ภิกฺขู อฏฺฐิํ กตฺวา\footnote{อฏฺฐิกตฺวา (สี, สฺยา, กํ, อิ)} \csromanfolio{114}มนสิ กตฺวา สพฺพเจตสา\footnote{สพฺพเจตโส (สี, สฺยา, กํ, อิ), สพฺพํ เจตสา (ก)} สมนฺนาหริตฺวา โอหิตโสตา ธมฺมํ สุณนฺติ.}

\prose{อถ โข มารสฺส ปาปิมโต เอตทโหสิ “อยํ โข สมโณ โคตโม ปญฺจนฺนํ อุปาทานกฺขนฺธานํ อุปาทาย ภิกฺขูนํ ธมฺมิยา กถาย สนฺทสฺเสติ สมาทเปติ สมุตฺเตเชติ สมฺปหํเสติ.\csromanspacer{} เต จ ภิกฺขู อฏฺฐิํ กตฺวา มนสิ กตฺวา สพฺพเจตสา สมนฺนาหริตฺวา โอหิตโสตา ธมฺมํ สุณนฺติ.\csromanspacer{} ยํนูนาหํ เยน สมโณ โคตโม เตนุปสงฺกเมยฺยํ วิจกฺขุกมฺมายา”ติ.}

\prose{เตน โข ปน สมเยน สมฺพหุลา ปตฺตา อพฺโภกาเส นิกฺขิตฺตา โหนฺติ.\csromanspacer{} อถ โข มาโร ปาปิมา พลีพทฺทวณฺณํ อภินิมฺมินิตฺวา เยน เต ปตฺตา เตนุปสงฺกมิ.\csromanspacer{} อถ โข อญฺญตโร ภิกฺขุ อญฺญตรํ ภิกฺขุํ เอตทโวจ “ภิกฺขุ ภิกฺขุ เอโส พลีพทฺโท ปตฺเต ภินฺเทยฺยา”ติ.\csromanspacer{} เอวํ วุตฺเต ภควา ตํ ภิกฺขุํ เอตทโวจ “น โส ภิกฺขุ พลีพทฺโท, มาโร เอโส ปาปิมา ตุมฺหากํ วิจกฺขุกมฺมาย อาคโต”ติ.\csromanspacer{} อถ โข ภควา มาโร อยํ ปาปิมา อิติ วิทิตฺวา มารํ ปาปิมนฺตํ คาถาย อชฺฌภาสิ\csromandash{}}

\csromangathameasure{“รูปํ เวทยิตํ สญฺญา, วิญฺญาณํ ยญฺจ สงฺขตํ. \\ เนโสหมสฺมิ เนตํ เม, เอวํ ตตฺถ วิรชฺชติ. \\ เอวํ วิรตฺตํ เขมตฺตํ, สพฺพสํโยชนาติคํ. \\ อเนฺวสํ สพฺพฏฺฐาเนสุ, มารเสนาปิ นาชฺฌคา”ติ.}
\csromangathagroup{\csromangathastanza{“รูปํ เวทยิตํ สญฺญา, วิญฺญาณํ ยญฺจ สงฺขตํ. \\ เนโสหมสฺมิ เนตํ เม, เอวํ ตตฺถ วิรชฺชติ.}\csromangathastanza{เอวํ วิรตฺตํ เขมตฺตํ, สพฺพสํโยชนาติคํ. \\ อเนฺวสํ สพฺพฏฺฐาเนสุ, มารเสนาปิ นาชฺฌคา”ติ.}}

\prose{อถ โข มาโร ปาปิมา ฯเปฯ ตตฺเถวนฺตรธายีติ.}

\csromanheader{2}{7. ฉผสฺสายตนสุตฺต}
\proseitem{153}{เอกํ สมยํ ภควา เวสาลิยํ วิหรติ มหาวเน กูฏาคารสาลายํ.\csromanspacer{} เตน โข ปน สมเยน ภควา ฉนฺนํ ผสฺสายตนานํ อุปาทาย ภิกฺขูนํ ธมฺมิยา กถาย สนฺทสฺเสติ สมาทเปติ สมุตฺเตเชติ สมฺปหํเสติ.\csromanspacer{} เต จ ภิกฺขู อฏฺฐิํ กตฺวา มนสิ กตฺวา สพฺพเจตสา สมนฺนาหริตฺวา โอหิตโสตา ธมฺมํ สุณนฺติ.}

\csromanheader{2}{4. มารสํยุตฺต}
\prose{\csromanfolio{115}อถ โข มารสฺส ปาปิมโต เอตทโหสิ “อยํ โข สมโณ โคตโม ฉนฺนํ ผสฺสายตนานํ อุปาทาย ภิกฺขูนํ ธมฺมิยา กถาย สนฺทสฺเสติ สมาทเปติ สมุตฺเตเชติ สมฺปหํเสติ.\csromanspacer{} เต จ ภิกฺขู อฏฺฐิํ กตฺวา มนสิ กตฺวา สพฺพเจตสา สมนฺนาหริตฺวา โอหิตโสตา ธมฺมํ สุณนฺติ.\csromanspacer{} ยํนูนาหํ เยน สมโณ โคตโม เตนุปสงฺกเมยฺยํ วิจกฺขุกมฺมายา”ติ.\csromanspacer{} อถ โข มาโร ปาปิมา เยน ภควา เตนุปสงฺกมิ, อุปสงฺกมิตฺวา ภควโต อวิทูเร มหนฺตํ ภยเภรวํ สทฺทมกาสิ, อปิสฺสุทํ ปถวี มญฺเญ อุนฺทฺรียติ\footnote{อุทฺรียติ (สี, สฺยา, กํ, อิ), อุ + ทร + ย + ติ = อุทฺรียติ.}.\csromanspacer{} อถ โข อญฺญตโร ภิกฺขุ อญฺญตรํ ภิกฺขุํ เอตทโวจ “ภิกฺขุ ภิกฺขุ เอสา ปถวี มญฺเญ อุนฺทฺรียตี”ติ.\csromanspacer{} เอวํ วุตฺเต ภควา ตํ ภิกฺขุํ เอตทโวจ “เนสา ภิกฺขุ ปถวี อุนฺทฺรียติ, มาโร เอโส ปาปิมา ตุมฺหากํ วิจกฺขุกมฺมาย อาคโต”ติ.\csromanspacer{} อถ โข ภควา มาโร อยํ ปาปิมา อิติ วิทิตฺวา มารํ ปาปิมนฺตํ คาถาย อชฺฌภาสิ\csromandash{}}

\csromangathameasure{“รูปา สทฺทา รสา คนฺธา, ผสฺสา ธมฺมา จ เกวลา. \\ เอตํ โลกามิสํ โฆรํ, เอตฺถ โลโก วิมุจฺฉิโต. \\ เอตญฺจ สมติกฺกมฺม, สโต พุทฺธสฺส สาวโก. \\ มารเธยฺยํ อติกฺกมฺม, อาทิจฺโจว วิโรจตี”ติ.}
\csromangathagroup{\csromangathastanza{“รูปา สทฺทา รสา คนฺธา, ผสฺสา ธมฺมา จ เกวลา. \\ เอตํ โลกามิสํ โฆรํ, เอตฺถ โลโก วิมุจฺฉิโต.}\csromangathastanza{เอตญฺจ สมติกฺกมฺม, สโต พุทฺธสฺส สาวโก. \\ มารเธยฺยํ อติกฺกมฺม, อาทิจฺโจว วิโรจตี”ติ.}}

\prose{อถ โข มาโร ปาปิมา ฯเปฯ ตตฺเถวนฺตรธายีติ.}

\csromanheader{2}{8. ปิณฺฑสุตฺต}
\proseitem{154}{เอกํ สมยํ ภควา มคเธสุ วิหรติ ปญฺจสาลายํ พฺราหฺมณคาเม.\csromanspacer{} เตน โข ปน สมเยน ปญฺจสาลายํ พฺราหฺมณคาเม กุมาริกานํ ปาหุนกานิ ภวนฺติ.\csromanspacer{} อถ โข ภควา ปุพฺพณฺหสมยํ นิวาเสตฺวา ปตฺตจีวรมาทาย ปญฺจสาลํ พฺราหฺมณคามํ ปิณฺฑาย ปาวิสิ.\csromanspacer{} เตน โข ปน สมเยน ปญฺจสาเลยฺยกา พฺราหฺมณคหปติกา มาเรน ปาปิมตา อนฺวาวิฏฺฐา ภวนฺติ “มา สมโณ โคตโม ปิณฺฑมลตฺถา”ติ.}

\prose{อถ โข ภควา ยถาโธเตน ปตฺเตน ปญฺจสาลํ พฺราหฺมณคามํ ปิณฺฑาย ปาวิสิ, ตถาโธเตน\footnote{ยถาโธเตน (?)} ปตฺเตน ปฏิกฺกมิ.\csromanspacer{} อถ โข มาโร ปาปิมา เยน ภควา เตนุปสงฺกมิ, อุปสงฺกมิตฺวา ภควนฺตํ เอตทโวจ \csromanfolio{116}“อปิ ตฺวํ สมณ ปิณฺฑมลตฺถา”ติ.\csromanspacer{} ตถา นุ ตฺวํ ปาปิม อกาสิ, ยถาหํ ปิณฺฑํ น ลเภยฺยนฺติ.\csromanspacer{} เตน หิ ภนฺเต ภควา ทุติยมฺปิ ปญฺจสาลํ พฺราหฺมณคามํ ปิณฺฑาย ปวิสตุ, ตถาหํ กริสฺสามิ, ยถา ภควา ปิณฺฑํ ลจฺฉตีติ.}

\csromangathameasure{อปุญฺญํ ปสวิ มาโร, อาสชฺช นํ ตถาคตํ. \\ กิํ นุ มญฺญสิ ปาปิม, น เม ปาปํ วิปจฺจติ. \\ สุสุขํ วต ชีวาม, เยสํ โน นตฺถิ กิญฺจนํ. \\ ปีติภกฺขา ภวิสฺสาม, เทวา อาภสฺสรา ยถาติ.}
\csromangathagroup{\csromangathastanza{อปุญฺญํ ปสวิ มาโร, อาสชฺช นํ ตถาคตํ. \\ กิํ นุ มญฺญสิ ปาปิม, น เม ปาปํ วิปจฺจติ.}\csromangathastanza{สุสุขํ วต ชีวาม, เยสํ โน นตฺถิ กิญฺจนํ. \\ ปีติภกฺขา ภวิสฺสาม, เทวา อาภสฺสรา ยถาติ.}}

\prose{อถ โข มาโร ปาปิมา “ชานาติ มํ ภควา, ชานาติ มํ สุคโต”ติ ทุกฺขี ทุมฺมโน ตตฺเถวนฺตรธายีติ.}

\csromanheader{2}{9. กสฺสกสุตฺต}
\proseitem{155}{สาวตฺถินิทานํ.\csromanspacer{} เตน โข ปน สมเยน ภควา ภิกฺขูนํ นิพฺพานปฏิสํยุตฺตาย ธมฺมิยา กถาย สนฺทสฺเสติ สมาทเปติ สมุตฺเตเชติ สมฺปหํเสติ.\csromanspacer{} เต จ ภิกฺขู อฏฺฐิํ กตฺวา มนสิ กตฺวา สพฺพเจตสา สมนฺนาหริตฺวา โอหิตโสตา ธมฺมํ สุณนฺติ.}

\prose{อถ โข มารสฺส ปาปิมโต เอตทโหสิ “อยํ โข สมโณ โคตโม ภิกฺขูนํ นิพฺพานปฏิสํยุตฺตาย ธมฺมิยา กถาย ฯเปฯ.\csromanspacer{} ยํนูนาหํ เยน สมโณ โคตโม เตนุปสงฺกเมยฺยํ วิจกฺขุกมฺมายา”ติ.\csromanspacer{} อถ โข มาโร ปาปิมา กสฺสกวณฺณํ อภินิมฺมินิตฺวา มหนฺตํ นงฺคลํ ขนฺเธ กริตฺวา ทีฆปาจนยฏฺฐิํ คเหตฺวา หฏหฏเกโส สาณสาฏินิวตฺโถ กทฺทมมกฺขิเตหิ ปาเทหิ เยน ภควา เตนุปสงฺกมิ, อุปสงฺกมิตฺวา ภควนฺตํ เอตทโวจ “อปิ สมณ พลีพทฺเท อทฺทสา”ติ.\csromanspacer{} กิํ ปน ปาปิม เต พลีพทฺเทหีติ.\csromanspacer{} มเมว สมณ จกฺขุ, มม รูปา, มม จกฺขุสมฺผสฺสวิญฺญาณายตนํ, กุหิํ เม สมณ คนฺตฺวา โมกฺขสิ.\csromanspacer{} มเมว สมณ โสตํ, มม สทฺทา ฯเปฯ.\csromanspacer{} มเมว สมณ ฆานํ, มม คนฺธา.\csromanspacer{} มเมว สมณ ชิวฺหา, มม รสา.\csromanspacer{} มเมว สมณ กาโย, มม โผฏฺฐพฺพา.\csromanspacer{} มเมว สมณ มโน, มม ธมฺมา, มม มโนสมฺผสฺสวิญฺญาณายตนํ, กุหิํ เม สมณ คนฺตฺวา โมกฺขสีติ.}

\csromanheader{2}{4. มารสํยุตฺต}
\prose{\csromanfolio{117}ตเวว ปาปิม จกฺขุ, ตว รูปา, ตว จกฺขุสมฺผสฺสวิญฺญาณายตนํ.\csromanspacer{} ยตฺถ จ โข ปาปิม นตฺถิ จกฺขุ, นตฺถิรูปา,นตฺถิจกฺขุสมฺผสฺสวิญฺญาณายตนํ, อคติ ตว ตตฺถ ปาปิม.\csromanspacer{} ตเวว ปาปิม โสตํ, ตว สทฺทา, ตว โสตสมฺผสฺสวิญฺญาณายตนํ.\csromanspacer{} ยตฺถ จ โข ปาปิม นตฺถิ โสตํ, นตฺถิ สทฺทา, นตฺถิ โสตสมฺผสฺสวิญฺญาณายตนํ, อคติ ตว ตตฺถ ปาปิม.\csromanspacer{} ตเวว ปาปิม ฆานํ, ตว คนฺธา, ตว ฆานสมฺผสฺสวิญฺญาณายตนํ.\csromanspacer{} ยตฺถ จ โข ปาปิม นตฺถิ ฆานํ, นตฺถิ คนฺธา, นตฺถิ ฆานสมฺผสฺสวิญฺญาณายตนํ, อคติ ตว ตตฺถ ปาปิม.\csromanspacer{} ตเวว ปาปิม ชิวฺหา, ตว รสา, ตว ชิวฺหาสมฺผสฺสวิญฺญาณายตนํ ฯเปฯ.\csromanspacer{} ตเวว ปาปิม กาโย, ตว โผฏฺฐพฺพา, ตว กายสมฺผสฺสวิญฺญาณายตนํ ฯเปฯ.\csromanspacer{} ตเวว ปาปิม มโน, ตว ธมฺมา, ตว มโนสมฺผสฺสวิญฺญาณายตนํ.\csromanspacer{} ยตฺถ จ โข ปาปิม นตฺถิ มโน, นตฺถิ ธมฺมา, นตฺถิ มโนสมฺผสฺสวิญฺญาณายตนํ, อคติ ตว ตตฺถ ปาปิมาติ.}

\csromangathameasure{ยํ วทนฺติ มม ยิทนฺติ, เย วทนฺติ มมนฺติ จ. \\ เอตฺถ เจ เต มโน อตฺถิ, น เม สมณ โมกฺขสีติ. \\ ยํ วทนฺติ น ตํ มยฺหํ, เย วทนฺติ น เต อหํ. \\ เอวํ ปาปิม ชานาหิ, น เม มคฺคมฺปิ ทกฺขสีติ.}
\csromangathagroup{\csromangathastanza{ยํ วทนฺติ มม ยิทนฺติ, เย วทนฺติ มมนฺติ จ. \\ เอตฺถ เจ เต มโน อตฺถิ, น เม สมณ โมกฺขสีติ.}\csromangathastanza{ยํ วทนฺติ น ตํ มยฺหํ, เย วทนฺติ น เต อหํ. \\ เอวํ ปาปิม ชานาหิ, น เม มคฺคมฺปิ ทกฺขสีติ.}}

\prose{อถ โข มาโร ปาปิมา ฯเปฯ ตตฺเถวนฺตรธายีติ.}

\csromanheader{2}{10. รชฺชสุตฺต}
\proseitem{156}{เอกํ สมยํ ภควา โกสเลสุ วิหรติ หิมวนฺตปเทเส\footnote{หิมวนฺตปสฺเส (สี)} อรญฺญกุฏิกายํ.\csromanspacer{} อถ โข ภควโต รโหคตสฺส ปฏิสลฺลีนสฺส เอวํ เจตโส ปริวิตกฺโก อุทปาทิ “สกฺกา นุ โข รชฺชํ กาเรตุํ อหนํ อฆาตยํ อชินํ อชาปยํ อโสจํ อโสจาปยํ ธมฺเมนา”ติ.}

\prose{อถ โข มาโร ปาปิมา ภควโต เจตสา เจโตปริวิตกฺกมญฺญาย เยน ภควา เตนุปสงฺกมิ, อุปสงฺกมิตฺวา ภควนฺตํ เอตทโวจ “กาเรตุ ภนฺเต ภควา รชฺชํ, กาเรตุ สุคโต รชฺชํ อหนํ อฆาตยํ อชินํ อชาปยํ อโสจํ อโสจาปยํ ธมฺเมนา”ติ.\csromanspacer{} กิํ ปน เม ตฺวํ ปาปิม ปสฺสสิ, ยํ มํ ตฺวํ เอวํ วเทสิ “กาเรตุ ภนฺเต ภควา รชฺชํ, กาเรตุ สุคโต รชฺชํ อหนํ อฆาตยํ อชินํ อชาปยํ อโสจํ \csromanfolio{118}อโสจาปยํ ธมฺเมนา”ติ.\csromanspacer{} ภควตา โข ภนฺเต จตฺตาโร อิทฺธิปาทา ภาวิตา พหุลีกตา ยานีกตา วตฺถุกตา อนุฏฺฐิตา ปริจิตา สุสมารทฺธา, อากงฺขมาโน จ ภนฺเต ภควา หิมวนฺตํ ปพฺพตราชํ สุวณฺณํเตฺวว อธิมุจฺเจยฺย, สุวณฺณญฺจ ปนสฺสาติ\footnote{สุวณฺณปพฺพตสฺสาติ (สี, สฺยา, กํ), สุวณฺณญฺจ ปพฺพตสฺสาติ (อิ)}.}

\csromangathameasure{ปพฺพตสฺส สุวณฺณสฺส, ชาตรูปสฺส เกวโล. \\ ทฺวิตฺตาว นาลเมกสฺส, อิติ วิทฺวา สมญฺจเร. \\ โย ทุกฺขมทฺทกฺขิ ยโตนิทานํ, \\ กาเมสุ โส ชนฺตุ กถํ นเมยฺย. \\ อุปธิํ วิทิตฺวา สงฺโคติ โลเก, \\ ตสฺเสว ชนฺตุ วินยาย สิกฺเขติ.}
\csromangathagroup{\csromangathastanza{ปพฺพตสฺส สุวณฺณสฺส, ชาตรูปสฺส เกวโล. \\ ทฺวิตฺตาว นาลเมกสฺส, อิติ วิทฺวา สมญฺจเร.}\csromangathastanza{โย ทุกฺขมทฺทกฺขิ ยโตนิทานํ, \\ กาเมสุ โส ชนฺตุ กถํ นเมยฺย. \\ อุปธิํ วิทิตฺวา สงฺโคติ โลเก, \\ ตสฺเสว ชนฺตุ วินยาย สิกฺเขติ.}}

\prose{อถ โข มาโร ปาปิมา “ชานาติ มํ ภควา, ชานาติ มํ สุคโต”ติ ทุกฺขี ทุมฺมโน ตตฺเถวนฺตรธายีติ.}

\vspace{\tipitakaparskip}
\csromancenter{ทุติโย วคฺโค.}
\titlehead{\textbf{ตสฺสุทฺทานํ}}
\csromangathameasure{ปาสาโณ สีโห สกลิกํ, ปติรูปญฺจ มานสํ. \\ ปตฺตํ อายตนํ ปิณฺฑํ, กสฺสกํ รชฺเชน เต ทสาติ.}
\csromangathagroup{\csromangathastanza{ปาสาโณ สีโห สกลิกํ\footnote{สกฺขลิกํ (ก)}, ปติรูปญฺจ มานสํ. \\ ปตฺตํ อายตนํ ปิณฺฑํ, กสฺสกํ รชฺเชน เต ทสาติ.}}

\chapterhead{3. ตติยวคฺค}
\csromanheader{2}{1. สมฺพหุลสุตฺต}
\proseitem{157}{เอวํ เม สุตํ\csromandash{}เอกํ สมยํ ภควา สกฺเกสุ วิหรติ สิลาวติยํ.\csromanspacer{} เตน โข ปน สมเยน สมฺพหุลา ภิกฺขู ภควโต อวิทูเร อปฺปมตฺตา อาตาปิโน ปหิตตฺตา วิหรนฺติ.\csromanspacer{} อถ โข มาโร ปาปิมา พฺราหฺมณวณฺณํ อภินิมฺมินิตฺวา มหนฺเตน ชฏณฺฑุเวน อชินกฺขิปนิวตฺโถ ชิณฺโณ โคปานสิวงฺโก ฆุรุฆุรุปสฺสาสี อุทุมฺพรทณฺฑํ คเหตฺวา เยน เต ภิกฺขู เตนุปสงฺกมิ, อุปสงฺกมิตฺวา เต ภิกฺขู เอตทโวจ “ทหรา}

\vspace{\tipitakaparskip}
\csromancenter{มารสํยุตฺต ภวนฺโต ปพฺพชิตา สุสู กาฬเกสา ภเทฺรน โยพฺพเนน สมนฺนาคตา ปฐเมน วยสา อนิกฺกีฬิตาวิโน กาเมสุ, ภุญฺชนฺตุ ภวนฺโต มานุสเก กาเม, มา สนฺทิฏฺฐิกํ หิตฺวา กาลิกํ อนุธาวิตฺถา”ติ.\csromanspacer{} น โข มยํ พฺราหฺมณ สนฺทิฏฺฐิกํ หิตฺวา กาลิกํ อนุธาวาม, กาลิกญฺจ โข มยํ พฺราหฺมณ หิตฺวา สนฺทิฏฺฐิกํ อนุธาวาม, กาลิกา หิ พฺราหฺมณ กามา วุตฺตา ภควตา พหุทุกฺขา พหุปายาสา, อาทีนโว เอตฺถ ภิยฺโย, สนฺทิฏฺฐิโก อยํ ธมฺโม อกาลิโก เอหิปสฺสิโก โอปเนยฺยิโก ปจฺจตฺตํ เวทิตพฺโพ วิญฺญูหีติ.\csromanspacer{} เอวํ วุตฺเต มาโร ปาปิมา สีสํ โอกมฺเปตฺวา ชิวฺหํ นิลฺลาเลตฺวา ติวิสาขํ นลาเฏ นลาฏิกํ วุฏฺฐาเปตฺวา ทณฺฑโมลุพฺภ ปกฺกามิ.}
\prose{\csromanfolio{119}อถ โข เต ภิกฺขู เยน ภควา เตนุปสงฺกมิํสุ, อุปสงฺกมิตฺวา ภควนฺตํ อภิวาเทตฺวา เอกมนฺตํ นิสีทิํสุ, เอกมนฺตํ นิสินฺนา โข เต ภิกฺขู ภควนฺตํ เอตทโวจุํ “อิธ มยํ ภนฺเต ภควโต อวิทูเร อปฺปมตฺตา อาตาปิโน ปหิตตฺตา วิหราม, อถ โข ภนฺเต อญฺญตโร พฺราหฺมโณ มหนฺเตน ชฏณฺฑุเวน อชินกฺขิปนิวตฺโถ ชิณฺโณ โคปานสิวงฺโก ฆุรุฆุรุปสฺสาสี อุทุมฺพรทณฺฑํ คเหตฺวา เยน มยํ เตนุปสงฺกมิ, อุปสงฺกมิตฺวา อเมฺห เอตทโวจ ‘ทหรา ภวนฺโต ปพฺพชิตา สุสู กาฬเกสา ภเทฺรน โยพฺพเนน สมนฺนาคตา ปฐเมน วยสา อนิกฺกีฬิตาวิโน กาเมสุ, ภุญฺชนฺตุ ภวนฺโต มานุสเก กาเม, มา สนฺทิฏฺฐิกํ หิตฺวา กาลิกํ อนุธาวิตฺถา’ติ”.\csromanspacer{} เอวํ วุตฺเต มยํ ภนฺเต ตํ พฺราหฺมณํ เอตทโวจุมฺห “น โข มยํ พฺราหฺมณ สนฺทิฏฺฐิกํ หิตฺวา กาลิกํ อนุธาวาม, กาลิกญฺจ โข มยํ พฺราหฺมณ หิตฺวา สนฺทิฏฺฐิกํ อนุธาวาม, กาลิกา หิ พฺราหฺมณ กามา วุตฺตา ภควตา พหุทุกฺขา พหุปายาสา, อาทีนโว เอตฺถ ภิยฺโย, สนฺทิฏฺฐิโก อยํ ธมฺโม อกาลิโก เอหิปสฺสิโก โอปเนยฺยิโก ปจฺจตฺตํ เวทิตพฺโพ วิญฺญูหี”ติ.\csromanspacer{} เอวํ วุตฺเต ภนฺเต โส พฺราหฺมโณ สีสํ โอกมฺเปตฺวา ชิวฺหํ นิลฺลาเลตฺวา ติวิสาขํ นลาเฏ นลาฏิกํ วุฏฺฐาเปตฺวา ทณฺฑโมลุพฺภ ปกฺกนฺโตติ.}

\prose{เนโส ภิกฺขเว พฺราหฺมโณ, มาโร เอโส ปาปิมา ตุมฺหากํ วิจกฺขุกมฺมาย อาคโตติ.\csromanspacer{} อถ โข ภควา เอตมตฺถํ วิทิตฺวา ตายํ เวลายํ อิมํ คาถํ อภาสิ\csromandash{}}

\csromangathameasure{“โย ทุกฺขมทฺทกฺขิ ยโตนิทานํ, \\ กาเมสุ โส ชนฺตุ กถํ นเมยฺย. \\ อุปธิํ วิทิตฺวา สงฺโคติ โลเก, \\ ตสฺเสว ชนฺตุ วินยาย สิกฺเข”ติ.}
\csromangathagroup{\csromangathastanza{\csromanfolio{120}“โย ทุกฺขมทฺทกฺขิ ยโตนิทานํ, \\ กาเมสุ โส ชนฺตุ กถํ นเมยฺย. \\ อุปธิํ วิทิตฺวา สงฺโคติ โลเก, \\ ตสฺเสว ชนฺตุ วินยาย สิกฺเข”ติ.}}

\csromanheader{2}{2. สมิทฺธิสุตฺต}
\proseitem{158}{เอกํ สมยํ ภควา สกฺเกสุ วิหรติ สิลาวติยํ.\csromanspacer{} เตน โข ปน สมเยน อายสฺมา สมิทฺธิ ภควโต อวิทูเร อปฺปมตฺโต อาตาปี ปหิตตฺโต วิหรติ.\csromanspacer{} อถ โข อายสฺมโต สมิทฺธิสฺส รโหคตสฺส ปฏิสลฺลีนสฺส เอวํ เจตโส ปริวิตกฺโก อุทปาทิ “ลาภา วต เม, สุลทฺธํ วต เม, ยสฺส เม สตฺถา อรหํ สมฺมาสมฺพุทฺโธ, ลาภา วต เม, สุลทฺธํ วต เม, ยฺวาหํ เอวํ สฺวากฺขาเต ธมฺมวินเย ปพฺพชิโต, ลาภา วต เม, สุลทฺธํ วต เม, ยสฺส เม สพฺรหฺมจาริโน สีลวนฺโต กลฺยาณธมฺมา”ติ.\csromanspacer{} อถ โข มาโร ปาปิมา อายสฺมโต สมิทฺธิสฺส เจตสา เจโตปริวิตกฺกมญฺญาย เยนายสฺมา สมิทฺธิ เตนุปสงฺกมิ, อุปสงฺกมิตฺวา อายสฺมโต สมิทฺธิสฺส อวิทูเร มหนฺตํ ภยเภรวํ สทฺทมกาสิ.\csromanspacer{} อปิสฺสุทํ ปถวี มญฺเญ อุนฺทฺรียติ.}

\prose{อถ โข อายสฺมา สมิทฺธิ เยน ภควา เตนุปสงฺกมิ, อุปสงฺกมิตฺวา ภควนฺตํ อภิวาเทตฺวา เอกมนฺตํ นิสีทิ, เอกมนฺตํ นิสินฺโน อายสฺมา สมิทฺธิ ภควนฺตํ เอตทโวจ “อิธาหํ ภนฺเต ภควโต อวิทูเร อปฺปมตฺโต อาตาปี ปหิตตฺโต วิหรามิ, ตสฺส มยฺหํ ภนฺเต รโหคตสฺส ปฏิสลฺลีนสฺส เอวํ เจตโส ปริวิตกฺโก อุทปาทิ ‘ลาภา วต เม, สุลทฺธํ วต เม, ยสฺส เม สตฺถา อรหํ สมฺมาสมฺพุทฺโธ, ลาภา วต เม, สุลทฺธํ วต เม, ยฺวาหํ เอวํ สฺวากฺขาเต ธมฺมวินเย ปพฺพชิโต, ลาภา วต เม, สุลทฺธํ วต เม, ยสฺส เม สพฺรหฺมจาริโน สีลวนฺโต กลฺยาณธมฺมา’ติ.\csromanspacer{} ตสฺส มยฺหํ ภนฺเต อวิทูเร มหาภยเภรวสทฺโท อโหสิ.\csromanspacer{} อปิสฺสุทํ ปถวี มญฺเญ อุนฺทฺรียตี”ติ.}

\prose{เนสา สมิทฺธิ ปถวี อุนฺทฺรียติ, มาโร เอโส ปาปิมา ตุยฺหํ วิจกฺขุกมฺมาย อาคโต.\csromanspacer{} คจฺฉ ตฺวํ สมิทฺธิ ตตฺเถว อปฺปมตฺโต อาตาปี ปหิตตฺโต วิหราหีติ. “เอวํ ภนฺเต”ติ โข อายสฺมา สมิทฺธิ ภควโต}

\vspace{\tipitakaparskip}
\csromancenter{มารสํยุตฺต ปฏิสฺสุตฺวา อุฏฺฐายาสนา ภควนฺตํ อภิวาเทตฺวา ปทกฺขิณํ กตฺวา ปกฺกามิ.\csromanspacer{} ทุติยมฺปิ โข อายสฺมา สมิทฺธิ ตตฺเถว อปฺปมตฺโต อาตาปี ปหิตตฺโต วิหาสิ.\csromanspacer{} ทุติยมฺปิ โข อายสฺมโต สมิทฺธิสฺส รโหคตสฺส ปฏิสลฺลีนสฺส ฯเปฯ.\csromanspacer{} ทุติยมฺปิ โข มาโร ปาปิมา อายสฺมโต สมิทฺธิสฺส เจตสา เจโตปริวิตกฺกมญฺญาย ฯเปฯ.\csromanspacer{} อปิสฺสุทํ ปถวี มญฺเญ อุนฺทฺรียติ.\csromanspacer{} อถ โข อายสฺมา สมิทฺธิ มารํ ปาปิมนฺตํ คาถาย อชฺฌภาสิ\csromandash{}}
\vspace{0.5\baselineskip}
\csromangathameasure{“สทฺธายาหํ ปพฺพชิโต, อคารสฺมา อนคาริยํ. \\ สติ ปญฺญา จ เม พุทฺธา, จิตฺตญฺจ สุสมาหิตํ. \\ กามํ กรสฺสุ รูปานิ, เนว มํ พฺยาธยิสฺสสี”ติ.}
\csromangathagroup{\csromangathastanza{\csromanfolio{121}“สทฺธายาหํ ปพฺพชิโต, อคารสฺมา อนคาริยํ. \\ สติ ปญฺญา จ เม พุทฺธา, จิตฺตญฺจ สุสมาหิตํ. \\ กามํ กรสฺสุ รูปานิ, เนว มํ พฺยาธยิสฺสสี”ติ.}}

\prose{อถ โข มาโร ปาปิมา “ชานาติ มํ สมิทฺธิ ภิกฺขู”ติ ทุกฺขี ทุมฺมโน ตตฺเถวนฺตรธายีติ.}

\csromanheader{2}{3. โคธิกสุตฺต}
\proseitem{159}{เอวํ เม สุตํ\csromandash{}เอกํ สมยํ ภควา ราชคเห วิหรติ เวฬุวเน กลนฺทกนิวาเป.\csromanspacer{} เตน โข ปน สมเยน อายสฺมา โคธิโก อิสิคิลิปสฺเส วิหรติ กาฬสิลายํ.\csromanspacer{} อถ โข อายสฺมา โคธิโก อปฺปมตฺโต อาตาปี ปหิตตฺโต วิหรนฺโต สามยิกํ เจโตวิมุตฺติํ ผุสิ.\csromanspacer{} อถ โข อายสฺมา โคธิโก ตมฺหา สามยิกาย เจโตวิมุตฺติยา ปริหายิ.\csromanspacer{} ทุติยมฺปิ โข อายสฺมา โคธิโก อปฺปมตฺโต อาตาปี ปหิตตฺโต วิหรนฺโต สามยิกํ เจโตวิมุตฺติํ ผุสิ.\csromanspacer{} ทุติยมฺปิ โข อายสฺมา โคธิโก ตมฺหา สามยิกาย เจโตวิมุตฺติยา ปริหายิ.\csromanspacer{} ตติยมฺปิ โข อายสฺมา โคธิโก อปฺปมตฺโต อาตาปี ปหิตตฺโต วิหรนฺโต สามยิกํ เจโตวิมุตฺติํ ผุสิ.\csromanspacer{} ตติยมฺปิ โข อายสฺมา โคธิโก ตมฺหา ฯเปฯ ปริหายิ.\csromanspacer{} จตุตฺถมฺปิ โข อายสฺมา โคธิโก อปฺปมตฺโต ฯเปฯ วิมุตฺติํ ผุสิ.\csromanspacer{} จตุตฺถมฺปิ โข อายสฺมา โคธิโก ตมฺหา ฯเปฯ ปริหายิ.\csromanspacer{} ปญฺจมมฺปิ โข อายสฺมา โคธิโก ฯเปฯ เจโตวิมุตฺติํ ผุสิ.\csromanspacer{} ปญฺจมมฺปิ โข อายสฺมา ฯเปฯ วิมุตฺติยา ปริหายิ.\csromanspacer{} ฉฏฺฐมฺปิ โข อายสฺมา โคธิโก อปฺปมตฺโต อาตาปี ปหิตตฺโต วิหรนฺโต สามยิกํ เจโตวิมุตฺติํ ผุสิ.\csromanspacer{} ฉฏฺฐมฺปิ โข อายสฺมา โคธิโก ตมฺหา สามยิกาย เจโตวิมุตฺติยา ปริหายิ. \csromanfolio{122}สตฺตมมฺปิ โข อายสฺมา โคธิโก อปฺปมตฺโต อาตาปี ปหิตตฺโต วิหรนฺโต สามยิกํ เจโตวิมุตฺติํ ผุสิ.}

\prose{อถ โข อายสฺมโต โคธิกสฺส เอตทโหสิ “ยาว ฉฏฺฐํ ขฺวาหํ สามยิกาย เจโตวิมุตฺติยา ปริหีโน, ยํนูนาหํ สตฺถํ อาหเรยฺยนฺ”ติ.\csromanspacer{} อถ โข มาโร ปาปิมา อายสฺมโต โคธิกสฺส เจตสา เจโตปริวิตกฺกมญฺญาย เยน ภควา เตนุปสงฺกมิ, อุปสงฺกมิตฺวา ภควนฺตํ คาถาหิ อชฺฌภาสิ\csromandash{}}

\csromangathameasure{“มหาวีร มหาปญฺญ, อิทฺธิยา ยสสา ชล. \\ สพฺพเวรภยาตีต, ปาเท วนฺทามิ จกฺขุม. \\ สาวโก เต มหาวีร, มรณํ มรณาภิภู. \\ อากงฺขติ เจตยติ, ตํ นิเสธ ชุตินฺธร. \\ กถํ หิ ภควา ตุยฺหํ, สาวโก สาสเน รโต. \\ อปฺปตฺตมานโส เสกฺโข, กาลํ กยิรา ชเนสุตา”ติ.}
\csromangathagroup{\csromangathastanza{“มหาวีร มหาปญฺญ, อิทฺธิยา ยสสา ชล. \\ สพฺพเวรภยาตีต, ปาเท วนฺทามิ จกฺขุม.}\csromangathastanza{สาวโก เต มหาวีร, มรณํ มรณาภิภู. \\ อากงฺขติ เจตยติ, ตํ นิเสธ ชุตินฺธร.}\csromangathastanza{กถํ หิ ภควา ตุยฺหํ, สาวโก สาสเน รโต. \\ อปฺปตฺตมานโส เสกฺโข, กาลํ กยิรา ชเนสุตา”ติ.}}

\prose{เตน โข ปน สมเยน อายสฺมตา โคธิเกน สตฺถํ อาหริตํ โหติ.\csromanspacer{} อถ โข ภควา มาโร อยํ ปาปิมา อิติ วิทิตฺวา มารํ ปาปิมนฺตํ คาถาย อชฺฌภาสิ\csromandash{}}

\csromangathameasure{“เอวํ หิ ธีรา กุพฺพนฺติ, นาวกงฺขนฺติ ชีวิตํ. \\ สมูลํ ตณฺหมพฺพุยฺห, โคธิโก ปรินิพฺพุโต”ติ.}
\csromangathagroup{\csromangathastanza{“เอวํ หิ ธีรา กุพฺพนฺติ, นาวกงฺขนฺติ ชีวิตํ. \\ สมูลํ ตณฺหมพฺพุยฺห, โคธิโก ปรินิพฺพุโต”ติ.}}

\prose{อถ โข ภควา ภิกฺขู อามนฺเตสิ “อายาม ภิกฺขเว เยน อิสิคิลิปสฺสํ กาฬสิลา เตนุปสงฺกมิสฺสาม, ยตฺถ โคธิเกน กุลปุตฺเตน สตฺถํ อาหริตนฺ”ติ. “เอวํ ภนฺเต”ติ โข เต ภิกฺขู ภควโต ปจฺจสฺโสสุํ.}

\prose{อถ โข ภควา สมฺพหุเลหิ ภิกฺขูหิ สทฺธิํ เยน อิสิคิลิปสฺสํ กาฬสิลา เตนุปสงฺกมิ.\csromanspacer{} อทฺทสา โข ภควา อายสฺมนฺตํ โคธิกํ ทูรโตว มญฺจเก วิวตฺตกฺขนฺธํ เสมานํ\footnote{เสยฺยมานํ (สฺยา, กํ), โสปฺปมานํ (ก)}.\csromanspacer{} เตน โข ปน สมเยน ธูมายิตตฺตํ ติมิรายิตตฺตํ คจฺฉเตว ปุริมํ ทิสํ, คจฺฉติ ปจฺฉิมํ ทิสํ, คจฺฉติ อุตฺตรํ ทิสํ, คจฺฉติ ทกฺขิณํ ทิสํ, คจฺฉติ อุทฺธํ, คจฺฉติ อโธ, คจฺฉติ อนุทิสํ.}

\csromanheader{2}{4. มารสํยุตฺต}
\prose{\csromanfolio{123}อถ โข ภควา ภิกฺขู อามนฺเตสิ “ปสฺสถ โน ตุเมฺห ภิกฺขเว เอตํ ธูมายิตตฺตํ ติมิรายิตตฺตํ คจฺฉเตว ปุริมํ ทิสํ, คจฺฉติ ปจฺฉิมํ ทิสํ, คจฺฉติ อุตฺตรํ ทิสํ, คจฺฉติ ทกฺขิณํ ทิสํ, คจฺฉติ อุทฺธํ, คจฺฉติ อโธ, คจฺฉติ อนุทิสนฺ”ติ.\csromanspacer{} เอวํ ภนฺเต.\csromanspacer{} เอโส โข ภิกฺขเว มาโร ปาปิมา โคธิกสฺส กุลปุตฺตสฺส วิญฺญาณํ สมเนฺวสติ “กตฺถ โคธิกสฺส กุลปุตฺตสฺส วิญฺญาณํ ปติฏฺฐิตนฺ”ติ.\csromanspacer{} อปฺปติฏฺฐิเตน จ ภิกฺขเว วิญฺญาเณน โคธิโก กุลปุตฺโต ปรินิพฺพุโตติ.\csromanspacer{} อถ โข มาโร ปาปิมา เพลุวปณฺฑุวีณํ อาทาย เยน ภควา เตนุปสงฺกมิ, อุปสงฺกมิตฺวา ภควนฺตํ คาถาย อชฺฌภาสิ\csromandash{}}

\csromangathameasure{“อุทฺธํ อโธ จ ติริยํ, ทิสา อนุทิสา สฺวหํ. \\ อเนฺวสํ นาธิคจฺฉามิ, โคธิโก โส กุหิํ คโต”ติ. \\ โย ธีโร ธิติสมฺปนฺโน, ฌายี ฌานรโต สทา. \\ อโหรตฺตํ อนุยุญฺชํ, ชีวิตํ อนิกามยํ. \\ เชตฺวาน มจฺจุโน เสนํ, อนาคนฺตฺวา ปุนพฺภวํ. \\ สมูลํ ตณฺหมพฺพุยฺห, โคธิโก ปรินิพฺพุโตติ. \\ ตสฺส โสกปเรตสฺส, วีณา กจฺฉา อภสฺสถ. \\ ตโต โส ทุมฺมโน ยกฺโข, ตตฺเถวนฺตรธายถาติ.}
\csromangathagroup{\csromangathastanza{“อุทฺธํ อโธ จ ติริยํ, ทิสา อนุทิสา สฺวหํ. \\ อเนฺวสํ นาธิคจฺฉามิ, โคธิโก โส กุหิํ คโต”ติ.}\csromangathastanza{โย\footnote{โส (สี, อิ)} ธีโร ธิติสมฺปนฺโน, ฌายี ฌานรโต สทา. \\ อโหรตฺตํ อนุยุญฺชํ, ชีวิตํ อนิกามยํ.}\csromangathastanza{เชตฺวาน มจฺจุโน\footnote{เภตฺวา นมุจิโน (สี)} เสนํ, อนาคนฺตฺวา ปุนพฺภวํ. \\ สมูลํ ตณฺหมพฺพุยฺห, โคธิโก ปรินิพฺพุโตติ.}\csromangathastanza{ตสฺส โสกปเรตสฺส, วีณา กจฺฉา อภสฺสถ. \\ ตโต โส ทุมฺมโน ยกฺโข, ตตฺเถวนฺตรธายถาติ\footnote{ตตฺเถวนฺตรธายิถาติ (สฺยา, กํ), ตตฺเถว อนฺตรธายีติ (ก)}.}}

\csromanheader{2}{4. สตฺตวสฺสานุพนฺธสุตฺต}
\proseitem{160}{เอวํ เม สุตํ\csromandash{}เอกํ สมยํ ภควา อุรุเวลายํ วิหรติ นชฺชา เนรญฺชราย ตีเร อชปาลนิโคฺรเธ.\csromanspacer{} เตน โข ปน สมเยน มาโร ปาปิมา สตฺตวสฺสานิ ภควนฺตํ อนุพนฺโธ โหติ โอตาราเปกฺโข โอตารํ อลภมาโน.\csromanspacer{} อถ โข มาโร ปาปิมา เยน ภควา เตนุปสงฺกมิ, อุปสงฺกมิตฺวา ภควนฺตํ คาถาย อชฺฌภาสิ\csromandash{}}

\csromangathameasure{“โสกาวติณฺโณ นุ วนมฺหิ ฌายสิ, \\ วิตฺตํ นุ ชีโน อุท ปตฺถยาโน. \\ อาคุํ นุ คามสฺมิมกาสิ กิญฺจิ, \\ กสฺมา ชเนน น กโรสิ สกฺขิํ. \\ สกฺขี น สมฺปชฺชติ เกนจิ เต”ติ. \\ โสกสฺส มูลํ ปลิขาย สพฺพํ, \\ อนาคุ ฌายามิ อโสจมาโน. \\ เฉตฺวาน สพฺพํ ภวโลภชปฺปํ, \\ อนาสโว ฌายามิ ปมตฺตพนฺธูติ. \\ ยํ วทนฺติ มม ยิทนฺติ, \\ เย วทนฺติ มมนฺติ จ. \\ เอตฺถ เจ เต มโน อตฺถิ, \\ น เม สมณ โมกฺขสีติ. ยํ วทนฺติ น ตํ มยฺหํ, \\ เย วทนฺติ น เต อหํ. เอวํ ปาปิม ชานาหิ, \\ น เม มคฺคมฺปิ ทกฺขสีติ. สเจ มคฺคํ อนุพุทฺธํ, \\ เขมํ อมตคามินํ. อเปหิ คจฺฉ ตฺวเมเวโก, \\ กิมญฺญมนุสาสสีติ. อมจฺจุเธยฺยํ ปุจฺฉนฺติ, \\ เย ชนา ปารคามิโน. เตสาหํ ปุฏฺโฐ อกฺขามิ, \\ ยํ สจฺจํ ตํ นิรูปธินฺติ.}
\csromangathagroup{\csromangathastanza{“โสกาวติณฺโณ นุ วนมฺหิ ฌายสิ, \\ วิตฺตํ นุ ชีโน อุท ปตฺถยาโน. \\ อาคุํ นุ คามสฺมิมกาสิ กิญฺจิ, \\ กสฺมา ชเนน น กโรสิ สกฺขิํ.}\csromangathastanza{สกฺขี น สมฺปชฺชติ เกนจิ เต”ติ. \\ โสกสฺส มูลํ ปลิขาย สพฺพํ, \\ อนาคุ ฌายามิ อโสจมาโน. \\ เฉตฺวาน สพฺพํ ภวโลภชปฺปํ,}\csromangathastanza{\csromanfolio{124}อนาสโว ฌายามิ ปมตฺตพนฺธูติ. \\ ยํ วทนฺติ มม ยิทนฺติ, \\ เย วทนฺติ มมนฺติ จ. \\ เอตฺถ เจ เต มโน อตฺถิ, \\ น เม สมณ โมกฺขสีติ. ยํ วทนฺติ น ตํ มยฺหํ, \\ เย วทนฺติ น เต อหํ. เอวํ ปาปิม ชานาหิ,}\csromangathastanza{น เม มคฺคมฺปิ ทกฺขสีติ. สเจ มคฺคํ อนุพุทฺธํ, \\ เขมํ อมตคามินํ. อเปหิ คจฺฉ ตฺวเมเวโก,}\csromangathastanza{กิมญฺญมนุสาสสีติ. อมจฺจุเธยฺยํ ปุจฺฉนฺติ, \\ เย ชนา ปารคามิโน. เตสาหํ ปุฏฺโฐ อกฺขามิ, \\ ยํ สจฺจํ ตํ นิรูปธินฺติ.}}

\prose{เสยฺยถาปิ ภนฺเต คามสฺส วา นิคมสฺส วา อวิทูเร โปกฺขรณี, ตตฺรสฺส กกฺกฏโก.\csromanspacer{} อถ โข ภนฺเต สมฺพหุลา กุมารกา วา กุมาริกาโย วา ตมฺหา คามา วา นิคมา วา นิกฺขมิตฺวา เยน สา โปกฺขรณี เตนุปสงฺกเมยฺยุํ, อุปสงฺกมิตฺวา ตํ กกฺกฏกํ อุทกา อุทฺธริตฺวา ถเล ปติฏฺฐเปยฺยุํ.\csromanspacer{} ยํ ยเทว หิ โส ภนฺเต กกฺกฏโก อฬํ อภินินฺนาเมยฺย.\csromanspacer{} ตํ ตเทว เต กุมารกา วา กุมาริกาโย วา กฏฺเฐน วา กถลาย วา สญฺฉินฺเทยฺยุํ สมฺภญฺเชยฺยุํ สมฺปลิภญฺเชยฺยุํ, เอวํ หิ โส ภนฺเต กกฺกฏโก สพฺเพหิ อเฬหิ สญฺฉินฺเนหิ สมฺภคฺเคหิ สมฺปลิภคฺเคหิ อภพฺโพ ตํ โปกฺขรณิํ โอตริตุํ.\csromanspacer{} เอวเมว โข ภนฺเต ยานิ กานิจิ วิสูกายิกานิ\footnote{ยานิ วิสุกายิกานิ (สี, อิ, ก)} วิเสวิตานิ วิปฺผนฺทิตานิ สพฺพานิ ตานิ\footnote{กานิจิ กานิจิ สพฺพานิ (สี, อิ, ก)} ภควตา สญฺฉินฺนานิ สมฺภคฺคานิ สมฺปลิภคฺคานิ, อภพฺโพ ทานาหํ ภนฺเต ปุน ภควนฺตํ อุปสงฺกมิตุํ, ยทิทํ โอตาราเปกฺโขติ.\csromanspacer{} อถ โข มาโร ปาปิมา ภควโต สนฺติเก อิมา นิพฺเพชนียา คาถาโย อภาสิ\csromandash{}}

\csromangathameasure{“เมทวณฺณญฺจ ปาสาณํ, วายโส อนุปริยคา. \\ อเปตฺถ มุทุํ วินฺเทม, อปิ อสฺสาทนา สิยา.}
\csromangathagroup{\csromangathastanza{“เมทวณฺณญฺจ ปาสาณํ, วายโส อนุปริยคา. \\ อเปตฺถ มุทุํ วินฺเทม, อปิ อสฺสาทนา สิยา.}}

\csromanheader{2}{4. มารสํยุตฺต}
\csromangathameasure{อลทฺธา ตตฺถ อสฺสาทํ, วายเสตฺโต อปกฺกเม. \\ กาโกว เสลมาสชฺช, นิพฺพิชฺชาเปม โคตมา”ติ.}
\csromangathagroup{\csromangathastanza{\csromanfolio{125}อลทฺธา ตตฺถ อสฺสาทํ, วายเสตฺโต อปกฺกเม. \\ กาโกว เสลมาสชฺช, นิพฺพิชฺชาเปม โคตมา”ติ.}}

\csromanheader{2}{5. มารธีตุสุตฺต}
\proseitem{161}{อถ โข มาโร ปาปิมา ภควโต สนฺติเก อิมา นิพฺเพชนียา คาถาโย อภาสิตฺวา ตมฺหา ฐานา อปกฺกมฺม ภควโต อวิทูเร ปถวิยํ ปลฺลงฺเกน นิสีทิ ตุณฺหีภูโต มงฺกุภูโต ปตฺตกฺขนฺโธ อโธมุโข ปชฺฌายนฺโต อปฺปฏิภาโน กฏฺเฐน ภูมิํ วิลิขนฺโต.\csromanspacer{} อถ โข ตณฺหา จ อรติ จ รคา จ มารธีตโร เยน มาโร ปาปิมา เตนุปสงฺกมิํสุ, อุปสงฺกมิตฺวา มารํ ปาปิมนฺตํ คาถาย อชฺฌภาสิํสุ\csromandash{}}

\csromangathameasure{“เกนาสิ ทุมฺมโน ตาต, ปุริสํ กํ นุ โสจสิ. \\ มยํ ตํ ราคปาเสน, อารญฺญมิว กุญฺชรํ. \\ พนฺธิตฺวา อานยิสฺสาม, วสโค เต ภวิสฺสตี”ติ. \\ อรหํ สุคโต โลเก, น ราเคน สุวานโย. \\ มารเธยฺยํ อติกฺกนฺโต, ตสฺมา โสจามหํ ภุสนฺติ.}
\csromangathagroup{\csromangathastanza{“เกนาสิ ทุมฺมโน ตาต, ปุริสํ กํ นุ โสจสิ. \\ มยํ ตํ ราคปาเสน, อารญฺญมิว กุญฺชรํ.}\csromangathastanza{พนฺธิตฺวา อานยิสฺสาม, วสโค เต ภวิสฺสตี”ติ. \\ อรหํ สุคโต โลเก, น ราเคน สุวานโย. \\ มารเธยฺยํ อติกฺกนฺโต, ตสฺมา โสจามหํ ภุสนฺติ.}}

\prose{อถ โข ตณฺหา จ อรติ จ รคา จ มารธีตโร เยน ภควา เตนุปสงฺกมิํสุ, อุปสงฺกมิตฺวา ภควนฺตํ เอตทโวจุํ “ปาเท เต สมณ ปริจาเรมา”ติ.\csromanspacer{} อถ โข ภควา น มนสากาสิ, ยถา ตํ อนุตฺตเร อุปธิสงฺขเย วิมุตฺโต.}

\prose{อถ โข ตณฺหา จ อรติ จ รคา จ มารธีตโร เอกมนฺตํ อปกฺกมฺม เอวํ สมจินฺเตสุํ “อุจฺจาวจา โข ปุริสานํ อธิปฺปายา, ยํนูน มยํ เอกสตํ เอกสตํ กุมาริวณฺณสตํ อภินิมฺมิเนยฺยามา”ติ.\csromanspacer{} อถ โข ตณฺหา จ อรติ จ รคา จ มารธีตโร เอกสตํ เอกสตํ กุมาริวณฺณสตํ อภินิมฺมินิตฺวา เยน ภควา เตนุปสงฺกมิํสุ, อุปสงฺกมิตฺวา ภควนฺตํ เอตทโวจุํ “ปาเท เต สมณ ปริจาเรมา”ติ.\csromanspacer{} ตมฺปิ ภควา น มนสากาสิ, ยถา ตํ อนุตฺตเร อุปธิสงฺขเย วิมุตฺโต.}

\prose{อถ โข ตณฺหา จ อรติ จ รคา จ มารธีตโร เอกมนฺตํ อปกฺกมฺม เอวํ สมจินฺเตสุํ “อุจฺจาวจา โข ปุริสานํ อธิปฺปายา, ยํนูน มยํ เอกสตํ \csromanfolio{126}เอกสตํ อวิชาตวณฺณสตํ อภินิมฺมิเนยฺยามา”ติ.\csromanspacer{} อถ โข ตณฺหา จ อรติ จ รคา จ มารธีตโร เอกสตํ เอกสตํ อวิชาตวณฺณสตํ อภินิมฺมินิตฺวา เยน ภควา เตนุปสงฺกมิํสุ, อุปสงฺกมิตฺวา ภควนฺตํ เอตทโวจุํ “ปาเท เต สมณ ปริจาเรมา”ติ.\csromanspacer{} ตมฺปิ ภควา น มนสากาสิ, ยถา ตํ อนุตฺตเร อุปธิสงฺขเย วิมุตฺโต.}

\prose{อถ โข ตณฺหา จ ฯเปฯ ยํนูน มยํ เอกสตํ เอกสตํ สกิํ วิชาตวณฺณสตํ อภินิมฺมิเนยฺยามาติ.\csromanspacer{} อถ โข ตณฺหา จ ฯเปฯ สกิํ วิชาตวณฺณสตํ อภินิมฺมินิตฺวา เยน ภควา เตนุปสงฺกมิํสุ, อุปสงฺกมิตฺวา ภควนฺตํ เอตทโวจุํ “ปาเท เต สมณ ปริจาเรมา”ติ.\csromanspacer{} ตมฺปิ ภควา น มนสากาสิ, ยถา ตํ อนุตฺตเร อุปธิสงฺขเย วิมุตฺโต.}

\prose{อถ โข ตณฺหา จ ฯเปฯ ยํนูน มยํ เอกสตํ เอกสตํ ทุวิชาตวณฺณสตํ อภินิมฺมิเนยฺยามาติ.\csromanspacer{} อถ โข ตณฺหา จ ฯเปฯ ทุวิชาตวณฺณสตํ อภินิมฺมินิตฺวา เยน ภควา ฯเปฯ ยถา ตํ อนุตฺตเร อุปธิสงฺขเย วิมุตฺโต.\csromanspacer{} อถ โข ตณฺหา จ ฯเปฯ มชฺฌิมิตฺถิวณฺณสตํ อภินิมฺมิเนยฺยามาติ.\csromanspacer{} อถ โข ตณฺหา จ ฯเปฯ มชฺฌิมิตฺถิวณฺณสตํ อภินิมฺมินิตฺวา ฯเปฯ อนุตฺตเร อุปธิสงฺขเย วิมุตฺโต. อถ โข ตณฺหา จ ฯเปฯ มหิตฺถิวณฺณสตํ อภินิมฺมิเนยฺยามาติ.\csromanspacer{} อถ โข ตณฺหา จ ฯเปฯ มหิตฺถิวณฺณสตํ อภินิมฺมินิตฺวา เยน ภควา ฯเปฯ อนุตฺตเร อุปธิสงฺขเย วิมุตฺโต.\csromanspacer{} อถ โข ตณฺหา จ อรติ จ รคา จ มารธีตโร เอกมนฺตํ อปกฺกมฺม เอตทโวจุํ “สจฺจํ กิร โน ปิตา อโวจ\csromandash{}}

\csromangathameasure{‘อรหํ สุคโต โลเก, น ราเคน สุวานโย. \\ มารเธยฺยํ อติกฺกนฺโต, ตสฺมา โสจามหํ ภุสนฺ’ติ”.}
\csromangathagroup{\csromangathastanza{‘อรหํ สุคโต โลเก, น ราเคน สุวานโย. \\ มารเธยฺยํ อติกฺกนฺโต, ตสฺมา โสจามหํ ภุสนฺ’ติ”.}}

\prose{ยํ หิ มยํ สมณํ วา พฺราหฺมณํ วา อวีตราคํ อิมินา อุปกฺกเมน อุปกฺกเมยฺยาม, หทยํ วาสฺส ผเลยฺย, อุณฺหํ โลหิตํ วา มุขโต อุคฺคจฺเฉยฺย, อุมฺมาทํ วา ปาปุเณยฺย จิตฺตกฺเขปํ วา.\csromanspacer{} เสยฺยถา วา ปน นโฬ หริโต ลุโต อุสฺสุสฺสติ วิสุสฺสติ มิลายติ.\csromanspacer{} เอวเมว อุสฺสุสฺเสยฺย วิสุสฺเสยฺย มิลาเยยฺยาติ.}

\csromanheader{2}{4. มารสํยุตฺต}
\prose{\csromanfolio{127}อถ โข ตณฺหา จ อรติ จ รคา จ มารธีตโร เยน ภควา เตนุปสงฺกมิํสุ, อุปสงฺกมิตฺวา เอกมนฺตํ อฏฺฐํสุ, เอกมนฺตํ ฐิตา โข ตณฺหา มารธีตา ภควนฺตํ คาถาย อชฺฌภาสิ\csromandash{}}

\csromangathameasure{“โสกาวติณฺโณ นุ วนมฺหิ ฌายสิ, \\ วิตฺตํ นุ ชีโน อุท ปตฺถยาโน. \\ อาคุํ นุ คามสฺมิมกาสิ กิญฺจิ, \\ กสฺมา ชเนน น กโรสิ สกฺขิํ. \\ สกฺขี น สมฺปชฺชติ เกนจิ เต”ติ. \\ อตฺถสฺส ปตฺติํ หทยสฺส สนฺติํ, \\ เชตฺวาน เสนํ ปิยสาตรูปํ. \\ เอโกหํ ฌายํ สุขมนุโพธิํ, \\ ตสฺมา ชเนน น กโรมิ สกฺขิํ. \\ สกฺขี น สมฺปชฺชติ เกนจิ เมติ.}
\csromangathagroup{\csromangathastanza{“โสกาวติณฺโณ นุ วนมฺหิ ฌายสิ, \\ วิตฺตํ นุ ชีโน อุท ปตฺถยาโน. \\ อาคุํ นุ คามสฺมิมกาสิ กิญฺจิ, \\ กสฺมา ชเนน น กโรสิ สกฺขิํ.}\csromangathastanza{สกฺขี น สมฺปชฺชติ เกนจิ เต”ติ. \\ อตฺถสฺส ปตฺติํ หทยสฺส สนฺติํ, \\ เชตฺวาน เสนํ ปิยสาตรูปํ. \\ เอโกหํ\footnote{เอกาหํ (สฺยา, กํ, อิ, ก)} ฌายํ สุขมนุโพธิํ, \\ ตสฺมา ชเนน น กโรมิ สกฺขิํ. \\ สกฺขี น สมฺปชฺชติ เกนจิ เมติ.}}

\csromanheader{2}{อถ โข อรติ\footnote{อรติ จ (ก) 3. กถํ ฌายํ (สฺยา, กํ, อิ), กถชฺฌายํ (ก) 4. น กุปฺปตี นสฺสรตี น ถีโน (สี)} มารธีตา ภควนฺตํ คาถาย อชฺฌภาสิ\csromandash{}}
\csromangathameasure{“กถํวิหารีพหุโลธ ภิกฺขุ, \\ ปญฺโจฆติณฺโณ อตรีธ ฉฏฺฐํ. \\ กถํ ฌายิํ3 พหุลํ กามสญฺญา, \\ ปริพาหิรา โหนฺติ อลทฺธ โย ตนฺ”ติ. \\ ปสฺสทฺธกาโย สุวิมุตฺตจิตฺโต, \\ อสงฺขราโน สติมา อโนโก. \\ อญฺญาย ธมฺมํ อวิตกฺกฌายี, \\ น กุปฺปติ น สรติ น ถิโน4. \\ เอวํวิหารีพหุโลธ ภิกฺขุ, \\ ปญฺโจฆติณฺโณ อตรีธ ฉฏฺฐํ. \\ เอวํ ฌายิํ พหุลํ กามสญฺญา, \\ ปริพาหิรา โหนฺติ อลทฺธ โย ตนฺติ. \\ อถ โข รคา มารธีตา ภควโต สนฺติเก คาถาย อชฺฌภาสิ\csromandash{} \\ “อจฺเฉชฺช ตณฺหํ คณสงฺฆจารี, \\ อทฺธา จริสฺสนฺติ พหู จ สทฺธา. \\ พหุํ วตายํ ชนตํ อโนโก, \\ อจฺเฉชฺช เนสฺสติ มจฺจุราชสฺส ปารนฺติ. \\ นยนฺติ เว มหาวีรา, \\ สทฺธมฺเมน ตถาคตา. \\ ธมฺเมน นยมานานํ, \\ กา อุสูยา วิชานตนฺ”ติ.}
\csromangathagroup{\csromangathastanza{“กถํวิหารีพหุโลธ ภิกฺขุ, \\ ปญฺโจฆติณฺโณ อตรีธ ฉฏฺฐํ. \\ กถํ ฌายิํ3 พหุลํ กามสญฺญา, \\ ปริพาหิรา โหนฺติ อลทฺธ โย ตนฺ”ติ.}\csromangathastanza{ปสฺสทฺธกาโย สุวิมุตฺตจิตฺโต, \\ อสงฺขราโน สติมา อโนโก. \\ อญฺญาย ธมฺมํ อวิตกฺกฌายี, \\ น กุปฺปติ น สรติ น ถิโน4.}\csromangathastanza{เอวํวิหารีพหุโลธ ภิกฺขุ, \\ ปญฺโจฆติณฺโณ อตรีธ ฉฏฺฐํ. \\ เอวํ ฌายิํ พหุลํ กามสญฺญา, \\ ปริพาหิรา โหนฺติ อลทฺธ โย ตนฺติ.}\csromangathastanza{\csromanfolio{128}อถ โข รคา\footnote{รคา จ (ก)} มารธีตา ภควโต สนฺติเก คาถาย อชฺฌภาสิ\csromandash{} \\ “อจฺเฉชฺช ตณฺหํ คณสงฺฆจารี, \\ อทฺธา จริสฺสนฺติ\footnote{รคา จ (ก)} พหู จ สทฺธา. \\ พหุํ วตายํ ชนตํ อโนโก,}\csromangathastanza{อจฺเฉชฺช เนสฺสติ มจฺจุราชสฺส ปารนฺติ. \\ นยนฺติ เว มหาวีรา, \\ สทฺธมฺเมน ตถาคตา. \\ ธมฺเมน นยมานานํ,}\csromangathastanza{กา อุสูยา วิชานตนฺ”ติ.}}

\prose{อถ โข ตณฺหา จ อรติ จ รคา จ มารธีตโร เยน มาโร ปาปิมา เตนุปสงฺกมิํสุ.\csromanspacer{} อทฺทสา โข มาโร ปาปิมา ตณฺหญฺจ อรติญฺจ รคญฺจ มารธีตโร ทูรโตว อาคจฺฉนฺติโย, ทิสฺวาน คาถาหิ อชฺฌภาสิ\csromandash{}}

\csromangathameasure{“พาลา กุมุทนาเฬหิ, ปพฺพตํ อภิมตฺถถ. \\ คิริํ นเขน ขนถ, อโย ทนฺเตหิ ขาทถ. \\ เสลํว สิรสูหจฺจ, ปาตาเล คาธเมสถ. \\ ขาณุํว อุรสาสชฺช, นิพฺพิชฺชาเปถ โคตมา”ติ. \\ ททฺทลฺลมานา อาคญฺฉุํ, ตณฺหา จ อรตี รคา. \\ ตา ตตฺถ ปนุที สตฺถา, ตูลํ ภฏฺฐํว มาลุโตติ.}
\csromangathagroup{\csromangathastanza{“พาลา กุมุทนาเฬหิ, ปพฺพตํ อภิมตฺถถ\footnote{อภิมนฺถถ (สี)}. \\ คิริํ นเขน ขนถ, อโย ทนฺเตหิ ขาทถ.}\csromangathastanza{เสลํว สิรสูหจฺจ\footnote{สิรสิ อูหจฺจ (สี), สิรสิ โอหจฺจ (สฺยา, กํ)}, ปาตาเล คาธเมสถ. \\ ขาณุํว อุรสาสชฺช, นิพฺพิชฺชาเปถ โคตมา”ติ.}\csromangathastanza{ททฺทลฺลมานา อาคญฺฉุํ, ตณฺหา จ อรตี รคา. \\ ตา ตตฺถ ปนุที สตฺถา, ตูลํ ภฏฺฐํว มาลุโตติ.}}

\vspace{\tipitakaparskip}
\csromancenter{ตติโย วคฺโค.}
\titlehead{\textbf{ตสฺสุทฺทานํ}}
\csromangathameasure{สมฺพหุลา สมิทฺธิ จ, โคธิกํ สตฺตวสฺสานิ. \\ ธีตรํ เทสิตํ พุทฺธ, เสฏฺเฐน อิมํ มารปญฺจกนฺติ. \\ \textbf{มารสํยุตฺตํ สมตฺตํ.}}
\csromangathagroup{\csromangathastanza{สมฺพหุลา สมิทฺธิ จ, โคธิกํ สตฺตวสฺสานิ. \\ ธีตรํ เทสิตํ พุทฺธ, เสฏฺเฐน อิมํ มารปญฺจกนฺติ. \\ \textbf{มารสํยุตฺตํ สมตฺตํ.}}}

\csromanheader{2}{5. ภิกฺขุนีสํยุตฺต}
\csromanheader{2}{1. อาฬวิกาสุตฺต}
\proseitem{162}{\csromanfolio{129}เอวํ เม สุตํ\csromandash{}เอกํ สมยํ ภควา สาวตฺถิยํ วิหรติ เชตวเน อนาถปิณฺฑิกสฺส อาราเม.\csromanspacer{} อถ โข อาฬวิกา ภิกฺขุนี ปุพฺพณฺหสมยํ นิวาเสตฺวา ปตฺตจีวรมาทาย สาวตฺถิํ ปิณฺฑาย ปาวิสิ, สาวตฺถิยํ ปิณฺฑาย จริตฺวา ปจฺฉาภตฺตํ ปิณฺฑปาตปฏิกฺกนฺตา เยน อนฺธวนํ เตนุปสงฺกมิ วิเวกตฺถินี.\csromanspacer{} อถ โข มาโร ปาปิมา อาฬวิกาย ภิกฺขุนิยา ภยํ ฉมฺภิตตฺตํ โลมหํสํ อุปฺปาเทตุกาโม วิเวกมฺหา จาเวตุกาโม เยน อาฬวิกา ภิกฺขุนี เตนุปสงฺกมิ, อุปสงฺกมิตฺวา อาฬวิกํ ภิกฺขุนิํ คาถาย อชฺฌภาสิ\csromandash{}}

\csromangathameasure{“นตฺถิ นิสฺสรณํ โลเก, กิํ วิเวเกน กาหสิ. \\ ภุญฺชสฺสุ กามรติโย, มาหุ ปจฺฉานุตาปินี”ติ.}
\csromangathagroup{\csromangathastanza{“นตฺถิ นิสฺสรณํ โลเก, กิํ วิเวเกน กาหสิ. \\ ภุญฺชสฺสุ กามรติโย, มาหุ ปจฺฉานุตาปินี”ติ.}}

\prose{อถ โข อาฬวิกาย ภิกฺขุนิยา เอตทโหสิ “โก นุ ขฺวายํ, มนุสฺโส วา อมนุสฺโส วา คาถํ ภาสตี”ติ.\csromanspacer{} อถ โข อาฬวิกาย ภิกฺขุนิยา เอตทโหสิ “มาโร โข อยํ ปาปิมา มม ภยํ ฉมฺภิตตฺตํ โลมหํสํ อุปฺปาเทตุกาโม วิเวกมฺหา จาเวตุกาโม คาถํ ภาสตี”ติ.\csromanspacer{} อถ โข อาฬวิกา ภิกฺขุนี มาโร อยํ ปาปิมา อิติ วิทิตฺวา มารํ ปาปิมนฺตํ คาถาหิ ปจฺจภาสิ\csromandash{}}

\csromangathameasure{“อตฺถิ นิสฺสรณํ โลเก, ปญฺญาย เม สุผุสฺสิตํ. \\ ปมตฺตพนฺธุ ปาปิม, น ตฺวํ ชานาสิ ตํ ปทํ. \\ สตฺติสูลูปมา กามา, ขนฺธาสํ อธิกุฏฺฏนา. \\ ยํ ตฺวํ กามรติํ พฺรูสิ, อรติ มยฺห สา อหู”ติ.}
\csromangathagroup{\csromangathastanza{“อตฺถิ นิสฺสรณํ โลเก, ปญฺญาย เม สุผุสฺสิตํ\footnote{สุผสฺสิตํ (สี, อิ)}. \\ ปมตฺตพนฺธุ ปาปิม, น ตฺวํ ชานาสิ ตํ ปทํ.}\csromangathastanza{สตฺติสูลูปมา กามา, ขนฺธาสํ อธิกุฏฺฏนา. \\ ยํ ตฺวํ กามรติํ พฺรูสิ, อรติ มยฺห สา อหู”ติ.}}

\prose{อถ โข มาโร ปาปิมา “ชานาติ มํ อาฬวิกา ภิกฺขุนี”ติ ทุกฺขี ทุมฺมโน ตตฺเถวนฺตรธายีติ.}

\csromanheader{2}{2. โสมาสุตฺต}
\proseitem{163}{\csromanfolio{130}สาวตฺถินิทานํ.\csromanspacer{} อถ โข โสมา ภิกฺขุนี ปุพฺพณฺหสมยํ นิวาเสตฺวา ปตฺตจีวรมาทาย สาวตฺถิํ ปิณฺฑาย ปาวิสิ, สาวตฺถิยํ ปิณฺฑาย จริตฺวา ปจฺฉาภตฺตํ ปิณฺฑปาตปฏิกฺกนฺตา เยน อนฺธวนํ เตนุปสงฺกมิ ทิวาวิหาราย.\csromanspacer{} อนฺธวนํ อชฺโฌคาเหตฺวา อญฺญตรสฺมิํ รุกฺขมูเล ทิวาวิหารํ นิสีทิ.\csromanspacer{} อถ โข มาโร ปาปิมา โสมาย ภิกฺขุนิยา ภยํ ฉมฺภิตตฺตํ โลมหํสํ อุปฺปาเทตุกาโม สมาธิมฺหา จาเวตุกาโม เยน โสมา ภิกฺขุนี เตนุปสงฺกมิ, อุปสงฺกมิตฺวา โสมํ ภิกฺขุนิํ คาถาย อชฺฌภาสิ\csromandash{}}

\csromangathameasure{“ยํ ตํ อิสีหิ ปตฺตพฺพํ, ฐานํ ทุรภิสมฺภวํ. \\ น ตํ ทฺวงฺคุลปญฺญาย, สกฺกา ปปฺโปตุมิตฺถิยา”ติ.}
\csromangathagroup{\csromangathastanza{“ยํ ตํ อิสีหิ ปตฺตพฺพํ, ฐานํ ทุรภิสมฺภวํ. \\ น ตํ ทฺวงฺคุลปญฺญาย, สกฺกา ปปฺโปตุมิตฺถิยา”ติ.}}

\prose{อถ โข โสมาย ภิกฺขุนิยา เอตทโหสิ “โก นุ ขฺวายํ, มนุสฺโส วา อมนุสฺโส วา คาถํ ภาสตี”ติ.\csromanspacer{} อถ โข โสมาย ภิกฺขุนิยา เอตทโหสิ “มาโร โข อยํ ปาปิมา มม ภยํ ฉมฺภิตตฺตํ โลมหํสํ อุปฺปาเทตุกาโม สมาธิมฺหา จาเวตุกาโม คาถํ ภาสตี”ติ.\csromanspacer{} อถ โข โสมา ภิกฺขุนี มาโร อยํ ปาปิมา อิติ วิทิตฺวา มารํ ปาปิมนฺตํ คาถาหิ ปจฺจภาสิ\csromandash{}}

\csromangathameasure{“อิตฺถิภาโว กิํ กยิรา, จิตฺตมฺหิ สุสมาหิเต. \\ ญาณมฺหิ วตฺตมานมฺหิ, สมฺมา ธมฺมํ วิปสฺสโต. \\ ยสฺส นูน สิยา เอวํ, อิตฺถาหํ ปุริโสติ วา. \\ กิญฺจิ วา ปน อญฺญสฺมิ, ตํ มาโร วตฺตุมรหตี”ติ.}
\csromangathagroup{\csromangathastanza{“อิตฺถิภาโว กิํ กยิรา, จิตฺตมฺหิ สุสมาหิเต. \\ ญาณมฺหิ วตฺตมานมฺหิ, สมฺมา ธมฺมํ วิปสฺสโต.}\csromangathastanza{ยสฺส นูน สิยา เอวํ, อิตฺถาหํ ปุริโสติ วา. \\ กิญฺจิ วา ปน อญฺญสฺมิ\footnote{อสฺมีติ (สฺยา, กํ, อิ)}, ตํ มาโร วตฺตุมรหตี”ติ.}}

\prose{อถ โข มาโร ปาปิมา “ชานาติ มํ โสมา ภิกฺขุนี”ติ ทุกฺขี ทุมฺมโน ตตฺเถวนฺตรธายีติ.}

\csromanheader{2}{3. กิสาโคตมีสุตฺต}
\proseitem{164}{สาวตฺถินิทานํ.\csromanspacer{} อถ โข กิสาโคตมี ภิกฺขุนี ปุพฺพณฺหสมยํ นิวาเสตฺวา ปตฺตจีวรมาทาย สาวตฺถิํ ปิณฺฑาย ปาวิสิ, สาวตฺถิยํ ปิณฺฑาย จริตฺวา ปจฺฉาภตฺตํ ปิณฺฑปาตปฏิกฺกนฺตา เยน อนฺธวนํ}

\csromanheader{2}{5. ภิกฺขุนีสํยุตฺต}
\prose{\csromanfolio{131}เตนุปสงฺกมิ ทิวาวิหาราย.\csromanspacer{} อนฺธวนํ อชฺโฌคาเหตฺวา อญฺญตรสฺมิํ รุกฺขมูเล ทิวาวิหารํ นิสีทิ.\csromanspacer{} อถ โข มาโร ปาปิมา กิสาโคตมิยา ภิกฺขุนิยา ภยํ ฉมฺภิตตฺตํ โลมหํสํ อุปฺปาเทตุกาโม สมาธิมฺหา จาเวตุกาโม เยน กิสาโคตมี ภิกฺขุนี เตนุปสงฺกมิ, อุปสงฺกมิตฺวา กิสาโคตมิํ ภิกฺขุนิํ คาถาย อชฺฌภาสิ\csromandash{}}

\csromangathameasure{“กิํ นุ ตฺวํ มตปุตฺตาว, เอกมาสิ รุทมฺมุขี. \\ วนมชฺฌคตา เอกา, ปุริสํ นุ คเวสสี”ติ.}
\csromangathagroup{\csromangathastanza{“กิํ นุ ตฺวํ มตปุตฺตาว, เอกมาสิ รุทมฺมุขี. \\ วนมชฺฌคตา เอกา, ปุริสํ นุ คเวสสี”ติ.}}

\prose{อถ โข กิสาโคตมิยา ภิกฺขุนิยา เอตทโหสิ “โก นุ ขฺวายํ, มนุสฺโส วา อมนุสฺโส วา คาถํ ภาสตี”ติ.\csromanspacer{} อถ โข กิสาโคตมิยา ภิกฺขุนิยา เอตทโหสิ “มาโร โข อยํ ปาปิมา มม ภยํ ฉมฺภิตตฺตํ โลมหํสํ อุปฺปาเทตุกาโม สมาธิมฺหา จาเวตุกาโม คาถํ ภาสตี”ติ.}

\prose{อถ โข กิสาโคตมี ภิกฺขุนี มาโร อยํ ปาปิมา อิติ วิทิตฺวา มารํ ปาปิมนฺตํ คาถาหิ ปจฺจภาสิ\csromandash{}}

\csromangathameasure{“อจฺจนฺตํ มตปุตฺตามฺหิ, ปุริสา เอตทนฺติกา. \\ น โสจามิ น โรทามิ, น ตํ ภายามิ อาวุโส. \\ สพฺพตฺถ วิหตา นนฺที, ตโมกฺขนฺโธ ปทาลิโต. \\ เชตฺวาน มจฺจุโน เสนํ, วิหรามิ อนาสวา”ติ.}
\csromangathagroup{\csromangathastanza{“อจฺจนฺตํ มตปุตฺตามฺหิ, ปุริสา เอตทนฺติกา. \\ น โสจามิ น โรทามิ, น ตํ ภายามิ อาวุโส.}\csromangathastanza{สพฺพตฺถ วิหตา นนฺที, ตโมกฺขนฺโธ ปทาลิโต. \\ เชตฺวาน มจฺจุโน\footnote{เชตฺวา นมุจิโน (สี)} เสนํ, วิหรามิ อนาสวา”ติ.}}

\prose{อถ โข มาโร ปาปิมา “ชานาติ มํ กิสาโคตมี ภิกฺขุนี”ติ ทุกฺขี ทุมฺมโน ตตฺเถวนฺตรธายีติ.}

\csromanheader{2}{4. วิชยาสุตฺต}
\proseitem{165}{สาวตฺถินิทานํ.\csromanspacer{} อถ โข วิชยา ภิกฺขุนี ปุพฺพณฺหสมยํ นิวาเสตฺวา ฯเปฯ อญฺญตรสฺมิํ รุกฺขมูเล ทิวาวิหารํ นิสีทิ.\csromanspacer{} อถ โข มาโร ปาปิมา วิชยาย ภิกฺขุนิยา ภยํ ฉมฺภิตตฺตํ โลมหํสํ อุปฺปาเทตุกาโม สมาธิมฺหา จาเวตุกาโม เยน วิชยา ภิกฺขุนี เตนุปสงฺกมิ, อุปสงฺกมิตฺวา วิชยํ ภิกฺขุนิํ คาถาย อชฺฌภาสิ\csromandash{}}

\csromangathameasure{“ทหรา ตฺวํ รูปวตี, อหญฺจ ทหโร สุสุ. \\ ปญฺจงฺคิเกน ตุริเยน, เอหยฺเยภิรมามเส”ติ.}
\csromangathagroup{\csromangathastanza{\csromanfolio{132}“ทหรา ตฺวํ รูปวตี, อหญฺจ ทหโร สุสุ. \\ ปญฺจงฺคิเกน ตุริเยน, เอหยฺเยภิรมามเส”ติ\footnote{เอหิ อยฺเย รมามเสติ (สี)}.}}

\prose{อถ โข วิชยาย ภิกฺขุนิยา เอตทโหสิ “โก นุ ขฺวายํ, มนุสฺโส วา อมนุสฺโส วา คาถํ ภาสตี”ติ.\csromanspacer{} อถ โข วิชยาย ภิกฺขุนิยา เอตทโหสิ “มาโร โข อยํ ปาปิมา มม ภยํ ฉมฺภิตตฺตํ โลมหํสํ อุปฺปาเทตุกาโม สมาธิมฺหา จาเวตุกาโม คาถํ ภาสตี”ติ.\csromanspacer{} อถ โข วิชยา ภิกฺขุนี “มาโร อยํ ปาปิมา” อิติ วิทิตฺวา มารํ ปาปิมนฺตํ คาถาหิ ปจฺจภาสิ\csromandash{}}

\csromangathameasure{รูปา สทฺทา รสา คนฺธา, โผฏฺฐพฺพา จ มโนรมา. \\ นิยฺยาตยามิ ตุเยฺหว, มาร นาหํ เตนตฺถิกา. \\ อิมินา ปูติกาเยน, ภินฺทเนน ปภงฺคุนา. \\ อฏฺฏียามิ หรายามิ, กามตณฺหา สมูหตา. \\ เย จ รูปูปคา สตฺตา, เย จ อรูปฏฺฐายิโน. \\ ยา จ สนฺตา สมาปตฺติ, สพฺพตฺถ วิหโต ตโมติ.}
\csromangathagroup{\csromangathastanza{รูปา สทฺทา รสา คนฺธา, โผฏฺฐพฺพา จ มโนรมา. \\ นิยฺยาตยามิ ตุเยฺหว, มาร นาหํ เตนตฺถิกา.}\csromangathastanza{อิมินา ปูติกาเยน, ภินฺทเนน ปภงฺคุนา. \\ อฏฺฏียามิ หรายามิ, กามตณฺหา สมูหตา.}\csromangathastanza{เย จ รูปูปคา สตฺตา, เย จ อรูปฏฺฐายิโน\footnote{อารุปฺปฏฺฐายิโน (สี, อิ)}. \\ ยา จ สนฺตา สมาปตฺติ, สพฺพตฺถ วิหโต ตโมติ.}}

\prose{อถ โข มาโร ปาปิมา “ชานาติ มํ วิชยา ภิกฺขุนี”ติ ทุกฺขี ทุมฺมโน ตตฺเถวนฺตรธายีติ.}

\csromanheader{2}{5. อุปฺปลวณฺณาสุตฺต}
\proseitem{166}{สาวตฺถินิทานํ.\csromanspacer{} อถ โข อุปฺปลวณฺณา ภิกฺขุนี ปุพฺพณฺหสมยํ นิวาเสตฺวา ฯเปฯ อญฺญตรสฺมิํ สุปุปฺผิตสาลรุกฺขมูเล อฏฺฐาสิ.\csromanspacer{} อถ โข มาโร ปาปิมา อุปฺปลวณฺณาย ภิกฺขุนิยา ภยํ ฉมฺภิตตฺตํ โลมหํสํ อุปฺปาเทตุกาโม สมาธิมฺหา จาเวตุกาโม เยน อุปฺปลวณฺณา ภิกฺขุนี เตนุปสงฺกมิ, อุปสงฺกมิตฺวา อุปฺปลวณฺณํ ภิกฺขุนิํ คาถาย อชฺฌภาสิ\csromandash{}}

\csromangathameasure{สุปุปฺผิตคฺคํ อุปคมฺม ภิกฺขุนิ, \\ เอกา ตุวํ ติฏฺฐสิ สาลมูเล. \\ น จตฺถิ เต ทุติยา วณฺณธาตุ, \\ พาเล น ตฺวํ ภายสิ ธุตฺตกานนฺติ.}
\csromangathagroup{\csromangathastanza{สุปุปฺผิตคฺคํ อุปคมฺม ภิกฺขุนิ, \\ เอกา ตุวํ ติฏฺฐสิ สาลมูเล. \\ น จตฺถิ เต ทุติยา วณฺณธาตุ, \\ พาเล น ตฺวํ ภายสิ ธุตฺตกานนฺติ.}}

\csromanheader{2}{5. ภิกฺขุนีสํยุตฺต}
\prose{\csromanfolio{133}อถ โข อุปฺปลวณฺณาย ภิกฺขุนิยา เอตทโหสิ “โก นุ ขฺวายํ, มนุสฺโส วา อมนุสฺโส วา คาถํ ภาสตี”ติ.\csromanspacer{} อถ โข อุปฺปลวณฺณาย ภิกฺขุนิยา เอตทโหสิ “มาโร โข อยํ ปาปิมา มม ภยํ ฉมฺภิตตฺตํ โลมหํสํ อุปฺปาเทตุกาโม สมาธิมฺหา จาเวตุกาโม คาถํ ภาสตี”ติ.\csromanspacer{} อถ โข อุปฺปลวณฺณา ภิกฺขุนี “มาโร อยํ ปาปิมา” อิติ วิทิตฺวา มารํ ปาปิมนฺตํ คาถาหิ ปจฺจภาสิ\csromandash{}}

\csromangathameasure{สตํ สหสฺสานิปิ ธุตฺตกานํ, \\ อิธาคตา ตาทิสกา ภเวยฺยุํ. \\ โลมํ น อิญฺชามิ น สนฺตสามิ, \\ น มาร ภายามิ ตเมกิกาปิ. \\ เอสา อนฺตรธายามิ, กุจฺฉิํ วา ปวิสามิ เต. \\ ปขุมนฺตริกายมฺปิ, ติฏฺฐนฺติํ มํ น ทกฺขสิ. \\ จิตฺตสฺมิํ วสีภูตามฺหิ, อิทฺธิปาทา สุภาวิตา. \\ สพฺพพนฺธนมุตฺตามฺหิ, น ตํ ภายามิ อาวุโสติ.}
\csromangathagroup{\csromangathastanza{สตํ สหสฺสานิปิ ธุตฺตกานํ, \\ อิธาคตา ตาทิสกา ภเวยฺยุํ. \\ โลมํ น อิญฺชามิ น สนฺตสามิ, \\ น มาร ภายามิ ตเมกิกาปิ. \\ เอสา อนฺตรธายามิ, กุจฺฉิํ วา ปวิสามิ เต. \\ ปขุมนฺตริกายมฺปิ, ติฏฺฐนฺติํ มํ น ทกฺขสิ.}\csromangathastanza{จิตฺตสฺมิํ วสีภูตามฺหิ, อิทฺธิปาทา สุภาวิตา. \\ สพฺพพนฺธนมุตฺตามฺหิ, น ตํ ภายามิ อาวุโสติ.}}

\prose{อถ โข มาโร ปาปิมา “ชานาติ มํ อุปฺปลวณฺณา ภิกฺขุนี”ติ ทุกฺขี ทุมฺมโน ตตฺเถวนฺตรธายีติ.}

\csromanheader{2}{6. จาลาสุตฺต}
\proseitem{167}{สาวตฺถินิทานํ.\csromanspacer{} อถ โข จาลา ภิกฺขุนี ปุพฺพณฺหสมยํ นิวาเสตฺวา ฯเปฯ อญฺญตรสฺมิํ รุกฺขมูเล ทิวาวิหารํ นิสีทิ.\csromanspacer{} อถ โข มาโร ปาปิมา เยน จาลา ภิกฺขุนี เตนุปสงฺกมิ, อุปสงฺกมิตฺวา จาลํ ภิกฺขุนิํ เอตทโวจ “กิํ นุ ตฺวํ ภิกฺขุนิ น โรเจสี”ติ.\csromanspacer{} ชาติํ ขฺวาหํ อาวุโส น โรเจมีติ.}

\csromangathameasure{กิํ นุ ชาติํ น โรเจสิ, ชาโต กามานิ ภุญฺชติ. \\ โก นุ ตํ อิทมาทปยิ, ชาติํ มา โรจ ภิกฺขุนีติ. \\ ชาตสฺส มรณํ โหติ, ชาโต ทุกฺขานิ ผุสฺสติ. \\ พนฺธํ วธํ ปริเกฺลสํ, ตสฺมา ชาติํ น โรจเย. \\ พุทฺโธ ธมฺมมเทเสสิ, ชาติยา สมติกฺกมํ. \\ สพฺพทุกฺขปฺปหานาย, โส มํ สจฺเจ นิเวสยิ. \\ เย จ รูปูปคา สตฺตา, เย จ อรูปฏฺฐายิโน. \\ นิโรธํ อปฺปชานนฺตา, อาคนฺตาโร ปุนพฺภวนฺติ.}
\csromangathagroup{\csromangathastanza{กิํ นุ ชาติํ น โรเจสิ, ชาโต กามานิ ภุญฺชติ. \\ โก นุ ตํ อิทมาทปยิ, ชาติํ มา โรจ\footnote{มา โรเจสิ (สี, อิ)} ภิกฺขุนีติ.}\csromangathastanza{ชาตสฺส มรณํ โหติ, ชาโต ทุกฺขานิ ผุสฺสติ\footnote{ปสฺสติ (สี, อิ)}. \\ พนฺธํ วธํ ปริเกฺลสํ, ตสฺมา ชาติํ น โรจเย.}\csromangathastanza{\csromanfolio{134}พุทฺโธ ธมฺมมเทเสสิ, ชาติยา สมติกฺกมํ. \\ สพฺพทุกฺขปฺปหานาย, โส มํ สจฺเจ นิเวสยิ.}\csromangathastanza{เย จ รูปูปคา สตฺตา, เย จ อรูปฏฺฐายิโน. \\ นิโรธํ อปฺปชานนฺตา, อาคนฺตาโร ปุนพฺภวนฺติ.}}

\prose{อถ โข มาโร ปาปิมา “ชานาติ มํ จาลา ภิกฺขุนี”ติ ทุกฺขี ทุมฺมโน ตตฺเถวนฺตรธายีติ.}

\csromanheader{2}{7. อุปจาลาสุตฺต}
\proseitem{168}{สาวตฺถินิทานํ.\csromanspacer{} อถ โข อุปจาลา ภิกฺขุนี ปุพฺพณฺหสมยํ นิวาเสตฺวา ฯเปฯ อญฺญตรสฺมิํ รุกฺขมูเล ทิวาวิหารํ นิสีทิ.\csromanspacer{} อถ โข มาโร ปาปิมา เยน อุปจาลา ภิกฺขุนี เตนุปสงฺกมิ, อุปสงฺกมิตฺวา อุปจาลํ ภิกฺขุนิํ เอตทโวจ “กตฺถ นุ ตฺวํ ภิกฺขุนิ อุปฺปชฺชิตุกามา”ติ.\csromanspacer{} น ขฺวาหํ อาวุโส กตฺถจิ อุปฺปชฺชิตุกามาติ.}

\csromangathameasure{ตาวติํสา จ ยามา จ, ตุสิตา จาปิ เทวตา. \\ นิมฺมานรติโน เทวา, เย เทวา วสวตฺติโน. \\ ตตฺถ จิตฺตํ ปณิเธหิ, รติํ ปจฺจนุโภสฺสสีติ. \\ ตาวติํสา จ ยามา จ, ตุสิตา จาปิ เทวตา. \\ นิมฺมานรติโน เทวา, เย เทวา วสวตฺติโน. \\ กามพนฺธนพทฺธา เต, เอนฺติ มารวสํ ปุน. \\ สพฺโพ อาทีปิโต โลโก, สพฺโพ โลโก ปธูปิโต. \\ สพฺโพ ปชฺชาลิโต โลโก, สพฺโพ โลโก ปกมฺปิโต. \\ อกมฺปิตํ อปชฺชลิตํ, อปุถุชฺชนเสวิตํ. \\ อคติ ยตฺถ มารสฺส, ตตฺถ เม นิรโต มโนติ.}
\csromangathagroup{\csromangathastanza{ตาวติํสา จ ยามา จ, ตุสิตา จาปิ เทวตา. \\ นิมฺมานรติโน เทวา, เย เทวา วสวตฺติโน.}\csromangathastanza{ตตฺถ จิตฺตํ ปณิเธหิ, รติํ ปจฺจนุโภสฺสสีติ. \\ ตาวติํสา จ ยามา จ, ตุสิตา จาปิ เทวตา.}\csromangathastanza{นิมฺมานรติโน เทวา, เย เทวา วสวตฺติโน. \\ กามพนฺธนพทฺธา เต, เอนฺติ มารวสํ ปุน.}\csromangathastanza{สพฺโพ อาทีปิโต\footnote{สพฺโพว อาทิตฺโต (สฺยา, กํ)} โลโก, สพฺโพ โลโก ปธูปิโต. \\ สพฺโพ ปชฺชาลิโต\footnote{สพฺโพว อาทิตฺโต (สฺยา, กํ)} โลโก, สพฺโพ โลโก ปกมฺปิโต.}\csromangathastanza{อกมฺปิตํ อปชฺชลิตํ\footnote{อจลิตํ (สี, สฺยา, กํ, อิ)}, อปุถุชฺชนเสวิตํ. \\ อคติ ยตฺถ มารสฺส, ตตฺถ เม นิรโต มโนติ.}}

\prose{อถ โข มาโร ปาปิมา “ชานาติ มํ อุปจาลา ภิกฺขุนี”ติ ทุกฺขี ทุมฺมโน ตตฺเถวนฺตรธายีติ.}

\csromanheader{2}{5. ภิกฺขุนีสํยุตฺต}
\csromanheader{2}{8. สีสุปจาลาสุตฺต}
\proseitem{169}{\csromanfolio{135}สาวตฺถินิทานํ.\csromanspacer{} อถ โข สีสุปจาลา\footnote{สีสูปจาลา (สี)} ภิกฺขุนี ปุพฺพณฺหสมยํ นิวาเสตฺวา ฯเปฯ อญฺญตรสฺมิํ รุกฺขมูเล ทิวาวิหารํ นิสีทิ.\csromanspacer{} อถ โข มาโร ปาปิมา เยน สีสุปจาลา ภิกฺขุนี เตนุปสงฺกมิ, อุปสงฺกมิตฺวา สีสุปจาลํ ภิกฺขุนิํ เอตทโวจ “กสฺส นุ ตฺวํ ภิกฺขุนิ ปาสณฺฑํ โรเจสี”ติ.\csromanspacer{} น ขฺวาหํ อาวุโส กสฺสจิ ปาสณฺฑํ โรเจมีติ.}

\csromangathameasure{กํ นุ อุทฺทิสฺส มุณฺฑาสิ, สมณี วิย ทิสฺสสิ, \\ น จ โรเจสิ ปาสณฺฑํ, \\ กิมิว จรสิ โมมูหาติ. \\ อิโต พหิทฺธา ปาสณฺฑา, \\ ทิฏฺฐีสุ ปสีทนฺติ เต. น เตสํ ธมฺมํ โรเจมิ, \\ เต ธมฺมสฺส อโกวิทา. อตฺถิ สกฺยกุเล ชาโต, \\ พุทฺโธ อปฺปฏิปุคฺคโล. สพฺพาภิภู มารนุโท, \\ สพฺพตฺถมปราชิโต. สพฺพตฺถ มุตฺโต อสิโต, \\ สพฺพํ ปสฺสติ จกฺขุมา. สพฺพกมฺมกฺขยํ ปตฺโต, \\ วิมุตฺโต อุปธิสงฺขเย. โส มยฺหํ ภควา สตฺถา, \\ ตสฺส โรเจมิ สาสนนฺติ.}
\csromangathagroup{\csromangathastanza{กํ นุ อุทฺทิสฺส มุณฺฑาสิ, สมณี วิย ทิสฺสสิ, \\ น จ โรเจสิ ปาสณฺฑํ, \\ กิมิว จรสิ โมมูหาติ. \\ อิโต พหิทฺธา ปาสณฺฑา, \\ ทิฏฺฐีสุ ปสีทนฺติ เต. น เตสํ ธมฺมํ โรเจมิ, \\ เต ธมฺมสฺส อโกวิทา. อตฺถิ สกฺยกุเล ชาโต,}\csromangathastanza{พุทฺโธ อปฺปฏิปุคฺคโล. สพฺพาภิภู มารนุโท, \\ สพฺพตฺถมปราชิโต. สพฺพตฺถ มุตฺโต อสิโต,}\csromangathastanza{สพฺพํ ปสฺสติ จกฺขุมา. สพฺพกมฺมกฺขยํ ปตฺโต, \\ วิมุตฺโต อุปธิสงฺขเย. โส มยฺหํ ภควา สตฺถา, \\ ตสฺส โรเจมิ สาสนนฺติ.}}

\prose{อถ โข มาโร ปาปิมา “ชานาติ มํ สีสุปจาลา ภิกฺขุนี”ติ ทุกฺขี ทุมฺมโน ตตฺเถวนฺตรธายีติ.}

\csromanheader{2}{9. เสลาสุตฺต}
\proseitem{170}{สาวตฺถินิทานํ.\csromanspacer{} อถ โข เสลา ภิกฺขุนี ปุพฺพณฺหสมยํ นิวาเสตฺวา ฯเปฯ อญฺญตรสฺมิํ รุกฺขมูเล ทิวาวิหารํ นิสีทิ.\csromanspacer{} อถ โข มาโร ปาปิมา เสลาย ภิกฺขุนิยา ภยํ ฉมฺภิตตฺตํ โลมหํสํ อุปฺปาเทตุกาโม ฯเปฯ เสลํ ภิกฺขุนิํ คาถาย อชฺฌภาสิ\csromandash{}}

\csromangathameasure{เกนิทํ ปกตํ พิมฺพํ, กฺวนุ พิมฺพสฺส การโก. \\ กฺวนุ พิมฺพํ สมุปฺปนฺนํ, กฺวนุ พิมฺพํ นิรุชฺฌตีติ.}
\csromangathagroup{\csromangathastanza{เกนิทํ ปกตํ พิมฺพํ, กฺวนุ\footnote{กฺวนฺนุ (สี, อิ), กฺวจิ (สฺยา, กํ, ก)} พิมฺพสฺส การโก. \\ กฺวนุ พิมฺพํ สมุปฺปนฺนํ, กฺวนุ พิมฺพํ นิรุชฺฌตีติ.}}

\prose{อถ โข เสลาย ภิกฺขุนิยา เอตทโหสิ “โก นุ ขฺวายํ, มนุสฺโส วา อมนุสฺโส วา คาถํ ภาสตี”ติ.\csromanspacer{} อถ โข เสลาย \csromanfolio{136}ภิกฺขุนิยา เอตทโหสิ “มาโร โข อยํ ปาปิมา มม ภยํ ฉมฺภิตตฺตํ โลมหํสํ อุปฺปาเทตุกาโม สมาธิมฺหา จาเวตุกาโม คาถํ ภาสตี”ติ.\csromanspacer{} อถ โข เสลา ภิกฺขุนี “มาโร อยํ ปาปิมา” อิติ วิทิตฺวา มารํ ปาปิมนฺตํ คาถาหิ ปจฺจภาสิ\csromandash{}}

\csromangathameasure{นยิทํ อตฺตกตํ พิมฺพํ, นยิทํ ปรกตํ1 อฆํ. \\ เหตุํ ปฏิจฺจ สมฺภูตํ, เหตุภงฺคา นิรุชฺฌติ. \\ ยถา อญฺญตรํ พีชํ, เขตฺเต วุตฺตํ วิรูหติ. \\ ปถวีรสญฺจาคมฺม, สิเนหญฺจ ตทูภยํ. \\ เอวํ ขนฺธา จ ธาตุโย, ฉ จ อายตนา อิเม. \\ เหตุํ ปฏิจฺจ สมฺภูตา, เหตุภงฺคา นิรุชฺฌเรติ.}
\csromangathagroup{\csromangathastanza{นยิทํ อตฺตกตํ\footnote{นยิทํ ปกตํ (สฺยา, กํ)} พิมฺพํ, นยิทํ ปรกตํ1 อฆํ. \\ เหตุํ ปฏิจฺจ สมฺภูตํ, เหตุภงฺคา นิรุชฺฌติ.}\csromangathastanza{ยถา อญฺญตรํ พีชํ, เขตฺเต วุตฺตํ วิรูหติ. \\ ปถวีรสญฺจาคมฺม, สิเนหญฺจ ตทูภยํ.}\csromangathastanza{เอวํ ขนฺธา จ ธาตุโย, ฉ จ อายตนา อิเม. \\ เหตุํ ปฏิจฺจ สมฺภูตา, เหตุภงฺคา นิรุชฺฌเรติ.}}

\prose{อถ โข มาโร ปาปิมา “ชานาติ มํ เสลา ภิกฺขุนี”ติ ทุกฺขี ทุมฺมโน ตตฺเถวนฺตรธายีติ.}

\csromanheader{2}{10. วชิราสุตฺต}
\proseitem{171}{สาวตฺถินิทานํ.\csromanspacer{} อถ โข วชิรา ภิกฺขุนี ปุพฺพณฺหสมยํ นิวาเสตฺวา ปตฺตจีวรมาทาย สาวตฺถิํ ปิณฺฑาย ปาวิสิ, สาวตฺถิยํ ปิณฺฑาย จริตฺวา ปจฺฉาภตฺตํ ปิณฺฑปาตปฏิกฺกนฺตา เยน อนฺธวนํ เตนุปสงฺกมิ ทิวาวิหาราย, อนฺธวนํ อชฺโฌคาเหตฺวา อญฺญตรสฺมิํ รุกฺขมูเล ทิวาวิหารํ นิสีทิ.\csromanspacer{} อถ โข มาโร ปาปิมา วชิราย ภิกฺขุนิยา ภยํ ฉมฺภิตตฺตํ โลมหํสํ อุปฺปาเทตุกาโม สมาธิมฺหา จาเวตุกาโม เยน วชิรา ภิกฺขุนี เตนุปสงฺกมิ, อุปสงฺกมิตฺวา วชิรํ ภิกฺขุนิํ คาถาย อชฺฌภาสิ\csromandash{}}

\csromangathameasure{เกนายํ ปกโต สตฺโต, กุวํ สตฺตสฺส การโก. \\ กุวํ สตฺโต สมุปฺปนฺโน, กุวํ สตฺโต นิรุชฺฌตีติ.}
\csromangathagroup{\csromangathastanza{เกนายํ ปกโต สตฺโต, กุวํ สตฺตสฺส การโก. \\ กุวํ สตฺโต สมุปฺปนฺโน, กุวํ สตฺโต นิรุชฺฌตีติ.}}

\prose{อถ โข วชิราย ภิกฺขุนิยา เอตทโหสิ “โก นุ ขฺวายํ, มนุสฺโส วา อมนุสฺโส วา คาถํ ภาสตี”ติ.\csromanspacer{} อถ โข วชิราย ภิกฺขุนิยา เอตทโหสิ “มาโร โข อยํ ปาปิมา มม ภยํ ฉมฺภิตตฺตํ โลมหํสํ อุปฺปาเทตุกาโม สมาธิมฺหา จาเวตุกาโม คาถํ ภาสตี”ติ.\csromanspacer{} อถ}

\csromanheader{2}{5. ภิกฺขุนีสํยุตฺต}
\prose{\csromanfolio{137}โข วชิรา ภิกฺขุนี “มาโร อยํ ปาปิมา” อิติ วิทิตฺวา มารํ ปาปิมนฺตํ คาถาหิ ปจฺจภาสิ\csromandash{}}

\csromangathameasure{กิํ นุ สตฺโตติ ปจฺเจสิ, มาร ทิฏฺฐิคตํ นุ เต. \\ สุทฺธสงฺขารปุญฺโชยํ, นยิธ สตฺตุปลพฺภติ. \\ ยถา หิ องฺคสมฺภารา, โหติ สทฺโท รโถ อิติ. \\ เอวํ ขนฺเธสุ สนฺเตสุ, โหติ สตฺโตติ สมฺมุติ. \\ ทุกฺขเมว หิ สมฺโภติ, ทุกฺขํ ติฏฺฐติ เวติ จ. \\ นาญฺญตฺร ทุกฺขา สมฺโภติ, นาญฺญํ ทุกฺขา นิรุชฺฌตีติ.}
\csromangathagroup{\csromangathastanza{กิํ นุ สตฺโตติ ปจฺเจสิ, มาร ทิฏฺฐิคตํ นุ เต. \\ สุทฺธสงฺขารปุญฺโชยํ, นยิธ สตฺตุปลพฺภติ.}\csromangathastanza{ยถา หิ องฺคสมฺภารา, โหติ สทฺโท รโถ อิติ. \\ เอวํ ขนฺเธสุ สนฺเตสุ, โหติ สตฺโตติ สมฺมุติ\footnote{สมฺมติ (สฺยา, กํ)}.}\csromangathastanza{ทุกฺขเมว หิ สมฺโภติ, ทุกฺขํ ติฏฺฐติ เวติ จ. \\ นาญฺญตฺร ทุกฺขา สมฺโภติ, นาญฺญํ ทุกฺขา นิรุชฺฌตีติ.}}

\prose{อถ โข มาโร ปาปิมา “ชานาติ มํ วชิรา ภิกฺขุนี”ติ ทุกฺขี ทุมฺมโน ตตฺเถวนฺตรธายีติ.}

\vspace{\tipitakaparskip}
\csromancenter{\textbf{ภิกฺขุนีสํยุตฺตํ สมตฺตํ.}}
\titlehead{\textbf{ตสฺสุทฺทานํ}}
\csromangathameasure{อาฬวิกา จ โสมา จ, โคตมี วิชยา สห. \\ อุปฺปลวณฺณา จ จาลา, อุปจาลา สีสุปจาลา จ. \\ เสลา วชิราย เต ทสาติ.}
\csromangathagroup{\csromangathastanza{อาฬวิกา จ โสมา จ, โคตมี วิชยา สห. \\ อุปฺปลวณฺณา จ จาลา, อุปจาลา สีสุปจาลา จ. \\ เสลา วชิราย เต ทสาติ.}}

\csromanheader{2}{6. พฺรหฺมสํยุตฺต}
\chapterhead{1. ปฐมวคฺค}
\csromanheader{2}{1. พฺรหฺมายาจนสุตฺต}
\proseitem{172}{\csromanfolio{138}เอวํ เม สุตํ\csromandash{}เอกํ สมยํ ภควา อุรุเวลายํ วิหรติ นชฺชา เนรญฺชราย ตีเร อชปาลนิโคฺรธมูเล ปฐมาภิสมฺพุทฺโธ.\csromanspacer{} อถ โข ภควโต รโหคตสฺส ปฏิสลฺลีนสฺส เอวํ เจตโส ปริวิตกฺโก อุทปาทิ “อธิคโต โข มฺยายํ ธมฺโม คมฺภีโร ทุทฺทโส ทุรนุโพโธ สนฺโต ปณีโต อตกฺกาวจโร นิปุโณ ปณฺฑิตเวทนีโย, อาลยรามา โข ปนายํ ปชา อาลยรตา อาลยสมฺมุทิตา, อาลยรามาย โข ปน ปชาย อาลยรตาย อาลยสมฺมุทิตาย ทุทฺทสํ อิทํ ฐานํ ยทิทํ อิทปฺปจฺจยตาปฏิจฺจสมุปฺปาโท, อิทมฺปิ โข ฐานํ ทุทฺทสํ ยทิทํ สพฺพสงฺขารสมโถ สพฺพูปธิปฏินิสฺสคฺโค ตณฺหากฺขโย วิราโค นิโรโธ นิพฺพานํ.\csromanspacer{} อหญฺเจว โข ปน ธมฺมํ เทเสยฺยํ, ปเร จ เม น อาชาเนยฺยุํ, โส มมสฺส กิลมโถ, สา มมสฺส วิเหสา”ติ.\csromanspacer{} อปิสฺสุ ภควนฺตํ อิมา อนจฺฉริยา คาถาโย ปฏิภํสุ ปุพฺเพ อสฺสุตปุพฺพา\csromandash{}}

\csromangathameasure{กิจฺเฉน เม อธิคตํ, หลํ ทานิ ปกาสิตุํ. \\ ราคโทสปเรเตหิ, นายํ ธมฺโม สุสมฺพุโธ. \\ ปฏิโสตคามิํ นิปุณํ, คมฺภีรํ ทุทฺทสํ อณุํ. \\ ราครตฺตา น ทกฺขนฺติ, ตโมขนฺเทน อาวุฏาติ.}
\csromangathagroup{\csromangathastanza{กิจฺเฉน เม อธิคตํ, หลํ ทานิ ปกาสิตุํ. \\ ราคโทสปเรเตหิ, นายํ ธมฺโม สุสมฺพุโธ.}\csromangathastanza{ปฏิโสตคามิํ นิปุณํ, คมฺภีรํ ทุทฺทสํ อณุํ. \\ ราครตฺตา น ทกฺขนฺติ, ตโมขนฺเทน อาวุฏาติ\footnote{ตโมกฺขนฺเธน อาวุตาติ (สี, สฺยา, กํ, อิ)}.}}

\prose{อิติห ภควโต ปฏิสญฺจิกฺขโต อปฺโปสฺสุกฺกตาย จิตฺตํ นมติ, โน ธมฺมเทสนาย.}

\prose{อถ โข พฺรหฺมุโน สหมฺปติสฺส ภควโต เจตสา เจโตปริวิตกฺกมญฺญาย เอตทโหสิ “นสฺสติ วต โภ โลโก, วินสฺสติ วต โภ โลโก, ยตฺร หิ นาม ตถาคตสฺส อรหโต สมฺมาสมฺพุทฺธสฺส อปฺโปสฺสุกฺกตาย จิตฺตํ นมติ\footnote{นมิสฺสติ (?)}, โน ธมฺมเทสนายา”ติ.}

\vspace{\tipitakaparskip}
\csromancenter{พฺรหฺมสํยุตฺต อถ โข พฺรหฺมา สหมฺปติ เสยฺยถาปิ นาม พลวา ปุริโส สมิญฺชิตํ\footnote{สมฺมิญฺชิตํ (สี, สฺยา, กํ, อิ)} วา พาหํ ปสาเรยฺย, ปสาริตํ วา พาหํ สมิญฺเชยฺย, เอวเมว พฺรหฺมโลเก อนฺตรหิโต ภควโต ปุรโต ปาตุรโหสิ.\csromanspacer{} อถ โข พฺรหฺมา สหมฺปติ เอกํสํ อุตฺตราสงฺคํ กริตฺวา ทกฺขิณชาณุมณฺฑลํ ปถวิยํ นิหนฺตฺวา เยน ภควา เตนญฺชลิํ ปณาเมตฺวา ภควนฺตํ เอตทโวจ “เทเสตุ ภนฺเต ภควา ธมฺมํ, เทเสตุ สุคโต ธมฺมํ, สนฺติ สตฺตา อปฺปรชกฺขชาติกา, อสฺสวนตา ธมฺมสฺส ปริหายนฺติ, ภวิสฺสนฺติ ธมฺมสฺส อญฺญาตาโร”ติ.\csromanspacer{} อิทมโวจ พฺรหฺมา สหมฺปติ, อิทํ วตฺวา อถาปรํ เอตทโวจ\csromandash{}}
\vspace{0.5\baselineskip}
\csromangathameasure{ปาตุรโหสิ มคเธสุ ปุพฺเพ, \\ ธมฺโม อสุทฺโธ สมเลหิ จินฺติโต. \\ อปาปุเรตํ อมตสฺส ทฺวารํ, \\ สุณนฺตุ ธมฺมํ วิมเลนานุพุทฺธํ.}
\csromangathagroup{\csromangathastanza{\csromanfolio{139}ปาตุรโหสิ มคเธสุ ปุพฺเพ, \\ ธมฺโม อสุทฺโธ สมเลหิ จินฺติโต. \\ อปาปุเรตํ\footnote{อวาปุเรตํ (สี)} อมตสฺส ทฺวารํ, \\ สุณนฺตุ ธมฺมํ วิมเลนานุพุทฺธํ.}}

\nitthitam{เสเล ยถา ปพฺพตมุทฺธนิฏฺฐิโต,}
\prose{ยถาปิ ปสฺเส ชนตํ สมนฺตโต.}

\csromangathameasure{ตถูปมํ ธมฺมมยํ สุเมธ, \\ ปาสาทมารุยฺห สมนฺตจกฺขุ. \\ โสกาวติณฺณํ ชนตมเปตโสโก, \\ อเวกฺขสฺสุ ชาติชราภิภูตํ. \\ อุฏฺเฐหิ วีร วิชิตสงฺคาม, \\ สตฺถวาห อนณ4 วิจร โลเก. \\ เทสสฺสุ ภควา ธมฺมํ, \\ อญฺญาตาโร ภวิสฺสนฺตีติ.}
\csromangathagroup{\csromangathastanza{ตถูปมํ ธมฺมมยํ สุเมธ, \\ ปาสาทมารุยฺห สมนฺตจกฺขุ. \\ โสกาวติณฺณํ\footnote{โสกาวกิณฺณํ (สี) 4. อณณ (รูปสิทฺธิฏีกา)} ชนตมเปตโสโก, \\ อเวกฺขสฺสุ ชาติชราภิภูตํ.}\csromangathastanza{อุฏฺเฐหิ วีร วิชิตสงฺคาม, \\ สตฺถวาห อนณ4 วิจร โลเก. \\ เทสสฺสุ\footnote{เทเสตุ (สฺยา, กํ, อิ, ก)} ภควา ธมฺมํ, \\ อญฺญาตาโร ภวิสฺสนฺตีติ.}}

\prose{อถ โข ภควา พฺรหฺมุโน จ อชฺเฌสนํ วิทิตฺวา สตฺเตสุ จ การุญฺญตํ ปฏิจฺจ พุทฺธจกฺขุนา โลกํ โวโลเกสิ, อทฺทสา โข ภควา พุทฺธจกฺขุนา โลกํ โวโลเกนฺโต สตฺเต อปฺปรชกฺเข มหารชกฺเข ติกฺขินฺทฺริเย มุทินฺทฺริเย สฺวากาเร ทฺวากาเร สุวิญฺญาปเย ทุวิญฺญาปเย อปฺเปกจฺเจ}

\prose{\csromanfolio{140}ปรโลกวชฺชภยทสฺสาวิเน\footnote{ทสฺสาวิโน (สี, สฺยา, กํ, อิ)} วิหรนฺเต อปฺเปกจฺเจ น ปรโลกวชฺชภยทสฺสาวิเน วิหรนฺเต, เสยฺยถาปิ นาม อุปฺปลินิยํ วา ปทุมินิยํ วา ปุณฺฑรีกินิยํ วา อปฺเปกจฺจานิ อุปฺปลานิ วา ปทุมานิ วา ปุณฺฑรีกานิ วา อุทเก ชาตานิ อุทเก สํวฑฺฒานิ อุทกานุคฺคตานิ อนฺโต นิมุคฺคโปสีนิ, อปฺเปกจฺจานิ อุปฺปลานิ วา ปทุมานิ วา ปุณฺฑรีกานิ วา อุทเก ชาตานิ อุทเก สํวฑฺฒานิ สโมทกํ ฐิตานิ, อปฺเปกจฺจานิ อุปฺปลานิ วา ปทุมานิ วา ปุณฺฑรีกานิ วา อุทเก ชาตานิ อุทเก สํวฑฺฒานิ อุทกา อจฺจุคฺคมฺม ฐิตานิ\footnote{ติฏฺฐนฺติ (สี, สฺยา, กํ, อิ)} อนุปลิตฺตานิ อุทเกน, เอวเมว ภควา พุทฺธจกฺขุนา โลกํ โวโลเกนฺโต อทฺทส สตฺเต อปฺปรชกฺเข มหารชกฺเข ติกฺขินฺทฺริเย มุทินฺทฺริเย สฺวากาเร ทฺวากาเร สุวิญฺญาปเย ทุวิญฺญาปเย อปฺเปกจฺเจ ปรโลกวชฺชภยทสฺสาวิเน วิหรนฺเต อปฺเปกจฺเจ น ปรโลกวชฺชภยทสฺสาวิเน วิหรนฺเต, ทิสฺวาน พฺรหฺมานํ สหมฺปติํ คาถาย ปจฺจภาสิ\csromandash{}}

\csromangathameasure{อปารุตา เตสํ อมตสฺส ทฺวารา, \\ เย โสตวนฺโต ปมุญฺจนฺตุ สทฺธํ. \\ วิหิํสสญฺญี ปคุณํ น ภาสิํ, \\ ธมฺมํ ปณีตํ มนุเชสุ พฺรหฺเมติ.}
\csromangathagroup{\csromangathastanza{อปารุตา เตสํ อมตสฺส ทฺวารา, \\ เย โสตวนฺโต ปมุญฺจนฺตุ สทฺธํ. \\ วิหิํสสญฺญี ปคุณํ น ภาสิํ, \\ ธมฺมํ ปณีตํ มนุเชสุ พฺรหฺเมติ.}}

\prose{อถ โข พฺรหฺมา สหมฺปติ “กตาวกาโส โขมฺหิ ภควตา ธมฺมเทสนายา”ติ ภควนฺตํ อภิวาเทตฺวา ปทกฺขิณํ กตฺวา ตตฺเถวนฺตรธายีติ.}

\csromanheader{2}{2. คารวสุตฺต}
\proseitem{173}{เอวํ เม สุตํ\csromandash{}เอกํ สมยํ ภควา อุรุเวลายํ วิหรติ นชฺชา เนรญฺชราย ตีเร อชปาลนิโคฺรธมูเล ปฐมาภิสมฺพุทฺโธ.\csromanspacer{} อถ โข ภควโต รโหคตสฺส ปฏิสลฺลีนสฺส เอวํ เจตโส ปริวิตกฺโก อุทปาทิ “ทุกฺขํ โข อคารโว วิหรติ อปฺปติสฺโส, กํ นุ ขฺวาหํ สมณํ วา พฺราหฺมณํ วา สกฺกตฺวา ครุํ กตฺวา\footnote{ครุกตฺวา (สี, สฺยา, กํ, อิ)} อุปนิสฺสาย วิหเรยฺยนฺ”ติ.}

\prose{อถ โข ภควโต เอตทโหสิ “อปริปุณฺณสฺส โข สีลกฺขนฺธสฺส ปาริปูริยา อญฺญํ สมณํ วา พฺราหฺมณํ วา สกฺกตฺวา ครุํ กตฺวา อุปนิสฺสาย}

\vspace{\tipitakaparskip}
\csromancenter{พฺรหฺมสํยุตฺต วิหเรยฺยํ, น โข ปนาหํ ปสฺสามิ สเทวเก โลเก สมารเก สพฺรหฺมเก สสฺสมณพฺราหฺมณิยา ปชาย สเทวมนุสฺสาย อตฺตนา สีลสมฺปนฺนตรํ อญฺญํ สมณํ วา พฺราหฺมณํ วา, ยมหํ สกฺกตฺวา ครุํ กตฺวา อุปนิสฺสาย วิหเรยฺยํ.}
\prose{\csromanfolio{141}อปริปุณฺณสฺส โข สมาธิกฺขนฺธสฺส ปาริปูริยา อญฺญํ สมณํ วา พฺราหฺมณํ วา สกฺกตฺวา ครุํ กตฺวา อุปนิสฺสาย วิหเรยฺยํ, น โข ปนาหํ ปสฺสามิ สเทวเก โลเก ฯเปฯ อตฺตนา สมาธิสมฺปนฺนตรํ อญฺญํ สมณํ วา พฺราหฺมณํ วา, ยมหํ สกฺกตฺวา ครุํ กตฺวา อุปนิสฺสาย วิหเรยฺยํ.}

\prose{อปริปุณฺณสฺส ปญฺญากฺขนฺธสฺส ปาริปูริยา อญฺญํ สมณํ วา พฺราหฺมณํ วา สกฺกตฺวา ครุํ กตฺวา อุปนิสฺสาย วิหเรยฺยํ, น โข ปนาหํ ปสฺสามิ สเทวเก ฯเปฯ อตฺตนา ปญฺญาสมฺปนฺนตรํ อญฺญํ สมณํ วา พฺราหฺมณํ วา, ยมหํ สกฺกตฺวา ครุํ กตฺวา อุปนิสฺสาย วิหเรยฺยํ.}

\prose{อปริปุณฺณสฺส โข วิมุตฺติกฺขนฺธสฺส ปาริปูริยา อญฺญํ สมณํ วา พฺราหฺมณํ วา สกฺกตฺวา ครุํ กตฺวา อุปนิสฺสาย วิหเรยฺยํ, น โข ปนาหํ ปสฺสามิ สเทวเก ฯเปฯ อตฺตนา วิมุตฺติสมฺปนฺนตรํ อญฺญํ สมณํ วา พฺราหฺมณํ วา, ยมหํ สกฺกตฺวา ครุํ กตฺวา อุปนิสฺสาย วิหเรยฺยํ.}

\prose{อปริปุณฺณสฺส โข วิมุตฺติญาณทสฺสนกฺขนฺธสฺส ปาริปูริยา อญฺญํ สมณํ วา พฺราหฺมณํ วา สกฺกตฺวา ครุํ กตฺวา อุปนิสฺสาย วิหเรยฺยํ.\csromanspacer{} น โข ปนาหํ ปสฺสามิ สเทวเก โลเก สมารเก สพฺรหฺมเก สสฺสมณพฺราหฺมณิยา ปชาย สเทวมนุสฺสาย อตฺตนา วิมุตฺติญาณทสฺสนสมฺปนฺนตรํ อญฺญํ สมณํ วา พฺราหฺมณํ วา, ยมหํ สกฺกตฺวา ครุํ กตฺวา อุปนิสฺสาย วิหเรยฺยํ.\csromanspacer{} ยํนูนาหํ ยฺวายํ ธมฺโม มยา อภิสมฺพุทฺโธ, ตเมว ธมฺมํ สกฺกตฺวา ครุํ กตฺวา อุปนิสฺสาย วิหเรยฺยนฺ”ติ.}

\prose{อถ โข พฺรหฺมา สหมฺปติ ภควโต เจตสา เจโตปริวิตกฺกมญฺญาย เสยฺยถาปิ นาม พลวา ปุริโส สมิญฺชิตํ วา พาหํ ปสาเรยฺย, ปสาริตํ วา พาหํ สมิญฺเชยฺย, เอวเมว พฺรหฺมโลเก อนฺตรหิโต ภควโต ปุรโต ปาตุรโหสิ.\csromanspacer{} อถ โข พฺรหฺมา สหมฺปติ เอกํสํ อุตฺตราสงฺคํ กริตฺวา เยน ภควา เตนญฺชลิํ ปณาเมตฺวา ภควนฺตํ เอตทโวจ “เอวเมตํ ภควา, เอวเมตํ สุคต, เยปิ เต ภนฺเต อเหสุํ อตีตมทฺธานํ อรหนฺโต สมฺมาสมฺพุทฺธา, \csromanfolio{142}เตปิ ภควนฺโต ธมฺมญฺเญว สกฺกตฺวา ครุํ กตฺวา อุปนิสฺสาย วิหริํสุ.\csromanspacer{} เยปิ เต ภนฺเต ภวิสฺสนฺติ อนาคตมทฺธานํ อรหนฺโต สมฺมาสมฺพุทฺธา, เตปิ ภควนฺโต ธมฺมญฺเญว สกฺกตฺวา ครุํ กตฺวา อุปนิสฺสาย วิหริสฺสนฺติ.\csromanspacer{} ภควาปิ ภนฺเต เอตรหิ อรหํ สมฺมาสมฺพุทฺโธ ธมฺมญฺเญว สกฺกตฺวา ครุํ กตฺวา อุปนิสฺสาย วิหรตูติ.\csromanspacer{} อิทมโวจ พฺรหฺมา สหมฺปติ, อิทํ วตฺวา อถาปรํ เอตทโวจ\csromandash{}}

\csromangathameasure{เย จ อตีตา สมฺพุทฺธา, เย จ พุทฺธา อนาคตา. \\ โย เจตรหิ สมฺพุทฺโธ, พหูนํ โสกนาสโน. \\ สพฺเพ สทฺธมฺมครุโน, วิหํสุ วิหรนฺติ จ. \\ ตถาปิ วิหริสฺสนฺติ, เอสา พุทฺธาน ธมฺมตา. \\ ตสฺมา หิ อตฺตกาเมน, มหตฺตมภิกงฺขตา. \\ สทฺธมฺโม ครุกาตพฺโพ, สรํ พุทฺธาน สาสนนฺติ.}
\csromangathagroup{\csromangathastanza{เย จ อตีตา สมฺพุทฺธา, เย จ พุทฺธา อนาคตา. \\ โย เจตรหิ สมฺพุทฺโธ, พหูนํ\footnote{พหุนฺนํ (สี, สฺยา, กํ, อิ)} โสกนาสโน.}\csromangathastanza{สพฺเพ สทฺธมฺมครุโน, วิหํสุ\footnote{วิหริํสุ (สี, สฺยา,กํ, อิ)} วิหรนฺติ จ. \\ ตถาปิ วิหริสฺสนฺติ, เอสา พุทฺธาน ธมฺมตา.}\csromangathastanza{ตสฺมา หิ อตฺตกาเมน\footnote{อตฺถกาเมน (สี, อิ, ก)}, มหตฺตมภิกงฺขตา. \\ สทฺธมฺโม ครุกาตพฺโพ, สรํ พุทฺธาน สาสนนฺติ.}}

\csromanheader{2}{3. พฺรหฺมเทวสุตฺต}
\proseitem{174}{เอวํ เม สุตํ\csromandash{}เอกํ สมยํ ภควา สาวตฺถิยํ วิหรติ เชตวเน อนาถปิณฺฑิกสฺส อาราเม.\csromanspacer{} เตน โข ปน สมเยน อญฺญตริสฺสา พฺราหฺมณิยา พฺรหฺมเทโว นาม ปุตฺโต ภควโต สนฺติเก อคารสฺมา อนคาริยํ ปพฺพชิโต โหติ.}

\prose{อถ โข อายสฺมา พฺรหฺมเทโว เอโก วูปกฏฺโฐ อปฺปมตฺโต อาตาปี ปหิตตฺโต วิหรนฺโต นจิรสฺเสว ยสฺสตฺถาย กุลปุตฺตา สมฺมเทว อคารสฺมา อนคาริยํ ปพฺพชนฺติ, ตทนุตฺตรํ พฺรหฺมจริยปริโยสานํ ทิฏฺเฐว ธมฺเม สยํ อภิญฺญา สจฺฉิกตฺวา อุปสมฺปชฺช วิหาสิ, “ขีณา ชาติ, วุสิตํ พฺรหฺมจริยํ, กตํ กรณียํ, นาปรํ อิตฺถตฺตายา”ติ อพฺภญฺญาสิ, อญฺญตโร จ ปนายสฺมา พฺรหฺมเทโว อรหตํ อโหสิ.}

\prose{อถ โข อายสฺมา พฺรหฺมเทโว ปุพฺพณฺหสมยํ นิวาเสตฺวา ปตฺตจีวรมาทาย สาวตฺถิํ ปิณฺฑาย ปาวิสิ, สาวตฺถิยํ สปทานํ ปิณฺฑาย จรมาโน เยน สกมาตุ นิเวสนํ เตนุปสงฺกมิ.\csromanspacer{} เตน โข ปน สมเยน อายสฺมโต พฺรหฺมเทวสฺส มาตา พฺราหฺมณี พฺรหฺมุโน อาหุติํ นิจฺจํ}

\vspace{\tipitakaparskip}
\csromancenter{พฺรหฺมสํยุตฺต ปคฺคณฺหาติ.\csromanspacer{} อถ โข พฺรหฺมุโน สหมฺปติสฺส เอตทโหสิ “อยํ โข อายสฺมโต พฺรหฺมเทวสฺส มาตา พฺราหฺมณี พฺรหฺมุโน อาหุติํ นิจฺจํ ปคฺคณฺหาติ, ยํนูนาหํ ตํ อุปสงฺกมิตฺวา สํเวเชยฺยนฺ”ติ.\csromanspacer{} อถ โข พฺรหฺมา สหมฺปติ เสยฺยถาปิ นาม พลวา ปุริโส สมิญฺชิตํ วา พาหํ ปสาเรยฺย, ปสาริตํ วา พาหํ สมิญฺเชยฺย, เอวเมว พฺรหฺมโลเก อนฺตรหิโต อายสฺมโต พฺรหฺมเทวสฺส มาตุ นิเวสเน ปาตุรโหสิ.\csromanspacer{} อถ โข พฺรหฺมา สหมฺปติ เวหาสํ ฐิโต อายสฺมโต พฺรหฺมเทวสฺส มาตรํ พฺราหฺมณิํ คาถาย อชฺฌภาสิ\csromandash{}}
\vspace{0.5\baselineskip}
\csromangathameasure{ทูเร อิโต พฺราหฺมณิ พฺรหฺมโลโก, \\ ยสฺสาหุติํ ปคฺคณฺหาสิ นิจฺจํ. \\ เนตาทิโส พฺราหฺมณิ พฺรหฺมภกฺโข, \\ กิํ ชปฺปสิ พฺรหฺมปถํ อชานํ. \\ เอโส หิ เต พฺราหฺมณิ พฺรหฺมเทโว, \\ นิรูปธิโก อติเทวปตฺโต. \\ อกิญฺจโน ภิกฺขุ อนญฺญโปสี, \\ โย เต โส ปิณฺฑาย ฆรํ ปวิฏฺโฐ. \\ อาหุเนยฺโย เวทคุ ภาวิตตฺโต, \\ นรานํ เทวานญฺจ ทกฺขิเณยฺโย. \\ พาหิตฺวา ปาปานิ อนูปลิตฺโต, \\ ฆาเสสนํ อิริยติ สีติภูโต. \\ น ตสฺส ปจฺฉา น ปุรตฺถมตฺถิ, \\ สนฺโต วิธูโม อนิโฆ นิราโส. \\ นิกฺขิตฺตทณฺโฑ ตสถาวเรสุ, \\ โส ตฺยาหุติํ ภุญฺชตุ อคฺคปิณฺฑํ. \\ วิเสนิภูโต อุปสนฺตจิตฺโต, \\ นาโคว ทนฺโต จรติ อเนโช. \\ ภิกฺขุ สุสีโล สุวิมุตฺตจิตฺโต, \\ โส ตฺยาหุติํ ภุญฺชตุ อคฺคปิณฺฑํ. \\ ตสฺมิํ ปสนฺนา อวิกมฺปมานา, \\ ปติฏฺฐเปหิ ทกฺขิณํ ทกฺขิเณยฺเย. \\ กโรหิ ปุญฺญํ สุขมายติกํ, \\ ทิสฺวา มุนิํ พฺราหฺมณิ โอฆติณฺณนฺติ. \\ ตสฺมิํ ปสนฺนา อวิกมฺปมานา, \\ ปติฏฺฐเปสิ ทกฺขิณํ ทกฺขิเณยฺเย. \\ อกาสิ ปุญฺญํ สุขมายติกํ, \\ ทิสฺวา มุนิํ พฺราหฺมณี โอฆติณฺณนฺติ.}
\csromangathagroup{\csromangathastanza{\csromanfolio{143}ทูเร อิโต พฺราหฺมณิ พฺรหฺมโลโก, \\ ยสฺสาหุติํ ปคฺคณฺหาสิ นิจฺจํ. \\ เนตาทิโส พฺราหฺมณิ พฺรหฺมภกฺโข, \\ กิํ ชปฺปสิ พฺรหฺมปถํ อชานํ\footnote{อชานนฺตี (สี, อิ, ก)}.}\csromangathastanza{เอโส หิ เต พฺราหฺมณิ พฺรหฺมเทโว, \\ นิรูปธิโก อติเทวปตฺโต. \\ อกิญฺจโน ภิกฺขุ อนญฺญโปสี, \\ โย เต โส\footnote{เต โส (สี, อิ), โย เต ส (?)} ปิณฺฑาย ฆรํ ปวิฏฺโฐ.}\csromangathastanza{อาหุเนยฺโย เวทคุ ภาวิตตฺโต, \\ นรานํ เทวานญฺจ ทกฺขิเณยฺโย. \\ พาหิตฺวา ปาปานิ อนูปลิตฺโต, \\ ฆาเสสนํ อิริยติ สีติภูโต.}\csromangathastanza{น ตสฺส ปจฺฉา น ปุรตฺถมตฺถิ, \\ สนฺโต วิธูโม อนิโฆ นิราโส. \\ นิกฺขิตฺตทณฺโฑ ตสถาวเรสุ, \\ โส ตฺยาหุติํ ภุญฺชตุ อคฺคปิณฺฑํ.}\csromangathastanza{วิเสนิภูโต อุปสนฺตจิตฺโต, \\ นาโคว ทนฺโต จรติ อเนโช. \\ ภิกฺขุ สุสีโล สุวิมุตฺตจิตฺโต, \\ โส ตฺยาหุติํ ภุญฺชตุ อคฺคปิณฺฑํ.}\csromangathastanza{\csromanfolio{144}ตสฺมิํ ปสนฺนา อวิกมฺปมานา, \\ ปติฏฺฐเปหิ ทกฺขิณํ ทกฺขิเณยฺเย. \\ กโรหิ ปุญฺญํ สุขมายติกํ, \\ ทิสฺวา มุนิํ พฺราหฺมณิ โอฆติณฺณนฺติ.}\csromangathastanza{ตสฺมิํ ปสนฺนา อวิกมฺปมานา, \\ ปติฏฺฐเปสิ ทกฺขิณํ ทกฺขิเณยฺเย. \\ อกาสิ ปุญฺญํ สุขมายติกํ, \\ ทิสฺวา มุนิํ พฺราหฺมณี โอฆติณฺณนฺติ.}}

\csromanheader{2}{4. พกพฺรหฺมสุตฺต}
\proseitem{175}{เอวํ เม สุตํ\csromandash{}เอกํ สมยํ ภควา สาวตฺถิยํ วิหรติ เชตวเน อนาถปิณฺฑิกสฺส อาราเม.\csromanspacer{} เตน โข ปน สมเยน พกสฺส พฺรหฺมุโน เอวรูปํ ปาปกํ ทิฏฺฐิคตํ อุปฺปนฺนํ โหติ “อิทํ นิจฺจํ อิทํ ธุวํ อิทํ สสฺสตํ อิทํ เกวลํ อิทํ อจวนธมฺมํ, อิทํ หิ น ชายติ น ชียติ น มียติ น จวติ น อุปปชฺชติ, อิโต จ ปนญฺญํ อุตฺตริ\footnote{อุตฺตริํ (สี, สฺยา, กํ, อิ)} นิสฺสรณํ นตฺถี”ติ.}

\prose{อถ โข ภควา พกสฺส พฺรหฺมุโน เจตสา เจโตปริวิตกฺกมญฺญาย เสยฺยถาปิ นาม พลวา ปุริโส สมิญฺชิตํ วา พาหํ ปสาเรยฺย, ปสาริตํ วา พาหํ สมิญฺเชยฺย, เอวเมว เชตวเน อนฺตรหิโต ตสฺมิํ พฺรหฺมโลเก ปาตุรโหสิ.\csromanspacer{} อทฺทสา โข พโก พฺรหฺมา ภควนฺตํ ทูรโตว อาคจฺฉนฺตํ, ทิสฺวาน ภควนฺตํ เอตทโวจ “เอหิ โข มาริส, สฺวาคตํ เต มาริส, จิรสฺสํ โข มาริส อิมํ ปริยายมกาสิ, ยทิทํ อิธาคมนาย.\csromanspacer{} อิทํ หิ มาริส นิจฺจํ อิทํ ธุวํ อิทํ สสฺสตํ อิทํ เกวลํ อิทํ อจวนธมฺมํ, อิทํ หิ น ชายติ น ชียติ น มียติ น จวติ น อุปปชฺชติ, อิโต จ ปนญฺญํ อุตฺตริ นิสฺสรณํ นตฺถี”ติ.}

\prose{เอวํ วุตฺเต ภควา พกํ พฺรหฺมานํ เอตทโวจ “อวิชฺชาคโต วต โภ พโก พฺรหฺมา, อวิชฺชาคโต วต โภ พโก พฺรหฺมา,}

\vspace{\tipitakaparskip}
\csromancenter{พฺรหฺมสํยุตฺต ยตฺร หิ นาม อนิจฺจํเยว สมานํ ‘นิจฺจนฺ’ติ วกฺขติ, อธุวํเยว สมานํ ‘ธุวนฺ’ติ วกฺขติ, อสสฺสตํเยว สมานํ ‘สสฺสตนฺ’ติ วกฺขติ, อเกวลํเยว สมานํ ‘เกวลนฺ’ติ วกฺขติ, จวนธมฺมํเยว สมานํ ‘อจวนธมฺมนฺ’ติ วกฺขติ.\csromanspacer{} ยตฺถ จ ปน ชายติ จ ชียติ จ มียติ จ จวติ จ อุปปชฺชติ จ, ตญฺจ ตถา วกฺขติ.\csromanspacer{} อิทํ หิ น ชายติ น ชียติ น มียติ น จวติ น อุปปชฺชติ, สนฺตญฺจ ปนญฺญํ อุตฺตริ นิสฺสรณํ, นตฺถญฺญํ อุตฺตริ นิสฺสรณนฺติ วกฺขตี”ติ.}
\vspace{0.5\baselineskip}
\csromangathameasure{ทฺวาสตฺตติ โคตม ปุญฺญกมฺมา, \\ วสวตฺติโน ชาติชรํ อตีตา. \\ อยมนฺติมา เวทคู พฺรหฺมุปปตฺติ, \\ อสฺมาภิชปฺปนฺติ ชนา อเนกาติ. \\ อปฺปํ หิ เอตํ น หิ ทีฆมายุ, \\ ยํ ตฺวํ พก มญฺญสิ ทีฆมายุํ. \\ สตํ สหสฺสานํ นิรพฺพุทานํ, \\ อายุํ ปชานามิ ตวาหํ พฺรหฺเมติ. \\ อนนฺตทสฺสี ภควาหมสฺมิ, \\ ชาติชรํ โสกมุปาติวตฺโต. \\ กิํ เม ปุราณํ วตสีลวตฺตํ, \\ อาจิกฺข เม ตํ ยมหํ วิชญฺญาติ. \\ ยํ ตฺวํ อปาเยสิ พหู มนุสฺเส, \\ ปิปาสิเต ฆมฺมนิ สมฺปเรเต. \\ ตํ เต ปุราณํ วตสีลวตฺตํ, \\ สุตฺตปฺปพุทฺโธว อนุสฺสรามิ. \\ ยํ เอณิกูลสฺมิํ ชนํ คหีตํ, \\ อโมจยี คยฺหกํ นียมานํ. \\ ตํ เต ปุราณํ วตสีลวตฺตํ, \\ สุตฺตปฺปพุทฺโธว อนุสฺสรามิ. \\ คงฺคาย โสตสฺมิํ คหีตนาวํ, \\ ลุทฺเทน นาเคน มนุสฺสกมฺยา. \\ ปโมจยิตฺถ พลสา ปสยฺห, \\ ตํ เต ปุราณํ วตสีลวตฺตํ. \\ สุตฺตปฺปพุทฺโธว อนุสฺสรามิ. \\ กปฺโป จ เต พทฺธจโร อโหสิํ, \\ สมฺพุทฺธิมนฺตํ วตินํ อมญฺญิ. \\ ตํ เต ปุราณํ วตสีลวตฺตํ, \\ สุตฺตปฺปพุทฺโธว อนุสฺสรามีติ. \\ อทฺธา ปชานาสิ มเมตมายุํ, \\ อญฺเญปิ ชานาสิ ตถา หิ พุทฺโธ. \\ ตถา หิ ตฺยายํ ชลิตานุภาโว, \\ โอภาสยํ ติฏฺฐติ พฺรหฺมโลกนฺติ.}
\csromangathagroup{\csromangathastanza{\csromanfolio{145}ทฺวาสตฺตติ โคตม ปุญฺญกมฺมา, \\ วสวตฺติโน ชาติชรํ อตีตา. \\ อยมนฺติมา เวทคู พฺรหฺมุปปตฺติ, \\ อสฺมาภิชปฺปนฺติ ชนา อเนกาติ.}\csromangathastanza{อปฺปํ หิ เอตํ น หิ ทีฆมายุ, \\ ยํ ตฺวํ พก มญฺญสิ ทีฆมายุํ. \\ สตํ สหสฺสานํ\footnote{สหสฺสาน (สฺยา, กํ)} นิรพฺพุทานํ, \\ อายุํ ปชานามิ ตวาหํ พฺรหฺเมติ.}\csromangathastanza{อนนฺตทสฺสี ภควาหมสฺมิ, \\ ชาติชรํ โสกมุปาติวตฺโต. \\ กิํ เม ปุราณํ วตสีลวตฺตํ, \\ อาจิกฺข เม ตํ ยมหํ วิชญฺญาติ.}\csromangathastanza{ยํ ตฺวํ อปาเยสิ พหู มนุสฺเส, \\ ปิปาสิเต ฆมฺมนิ สมฺปเรเต. \\ ตํ เต ปุราณํ วตสีลวตฺตํ, \\ สุตฺตปฺปพุทฺโธว อนุสฺสรามิ.}\csromangathastanza{ยํ เอณิกูลสฺมิํ ชนํ คหีตํ, \\ อโมจยี คยฺหกํ นียมานํ. \\ ตํ เต ปุราณํ วตสีลวตฺตํ, \\ สุตฺตปฺปพุทฺโธว อนุสฺสรามิ.}\csromangathastanza{\csromanfolio{146}คงฺคาย โสตสฺมิํ คหีตนาวํ, \\ ลุทฺเทน นาเคน มนุสฺสกมฺยา. \\ ปโมจยิตฺถ พลสา ปสยฺห, \\ ตํ เต ปุราณํ วตสีลวตฺตํ.}\csromangathastanza{สุตฺตปฺปพุทฺโธว อนุสฺสรามิ. \\ กปฺโป จ เต พทฺธจโร อโหสิํ, \\ สมฺพุทฺธิมนฺตํ\footnote{สมฺพุทฺธิวนฺตํ (พหูสุ)} วตินํ อมญฺญิ. \\ ตํ เต ปุราณํ วตสีลวตฺตํ,}\csromangathastanza{สุตฺตปฺปพุทฺโธว อนุสฺสรามีติ. \\ อทฺธา ปชานาสิ มเมตมายุํ, \\ อญฺเญปิ\footnote{อญฺญมฺปิ (สี, อิ)} ชานาสิ ตถา หิ พุทฺโธ. \\ ตถา หิ ตฺยายํ ชลิตานุภาโว,}\csromangathastanza{โอภาสยํ ติฏฺฐติ พฺรหฺมโลกนฺติ.}}

\csromanheader{2}{5. อญฺญตรพฺรหฺมสุตฺต}
\proseitem{176}{สาวตฺถินิทานํ.\csromanspacer{} เตน โข ปน สมเยน อญฺญตรสฺส พฺรหฺมุโน เอวรูปํ ปาปกํ ทิฏฺฐิคตํ อุปฺปนฺนํ โหติ “นตฺถิ โส สมโณ วา พฺราหฺมโณ วา, โย อิธ อาคจฺเฉยฺยา”ติ.\csromanspacer{} อถ โข ภควา ตสฺส พฺรหฺมุโน เจตสา เจโตปริวิตกฺกมญฺญาย เสยฺยถาปิ นาม พลวา ปุริโส ฯเปฯ ตสฺมิํ พฺรหฺมโลเก ปาตุรโหสิ.\csromanspacer{} อถ โข ภควา ตสฺส พฺรหฺมุโน อุปริ เวหาสํ ปลฺลงฺเกน นิสีทิ เตโชธาตุํ สมาปชฺชิตฺวา.}

\prose{อถ โข อายสฺมโต มหาโมคฺคลฺลานสฺส เอตทโหสิ “กหํ นุ โข ภควา เอตรหิ วิหรตี”ติ.\csromanspacer{} อทฺทสา โข อายสฺมา มหาโมคฺคลฺลาโน\footnote{มหาโมคฺคลาโน (ก)} ภควนฺตํ ทิพฺเพน จกฺขุนา วิสุทฺเธน อติกฺกนฺตมานุสเกน ตสฺส พฺรหฺมุโน อุปริเวหาสํ ปลฺลงฺเกน นิสินฺนํ เตโชธาตุํ สมาปนฺนํ, ทิสฺวาน เสยฺยถาปิ นาม พลวา ปุริโส สมิญฺชิตํ วา พาหํ ปสาเรยฺย, ปสาริตํ วา พาหํ สมิญฺเชยฺย, เอวเมว เชตวเน อนฺตรหิโต ตสฺมิํ พฺรหฺมโลเก ปาตุรโหสิ.\csromanspacer{} อถ โข อายสฺมา}

\vspace{\tipitakaparskip}
\csromancenter{พฺรหฺมสํยุตฺต มหาโมคฺคลฺลาโน ปุรตฺถิมํ ทิสํ นิสฺสาย\footnote{อุปนิสฺสาย (สี)} ตสฺส พฺรหฺมุโน อุปริ เวหาสํ ปลฺลงฺเกน นิสีทิ เตโชธาตุํ สมาปชฺชิตฺวา นีจตรํ ภควโต.}
\prose{\csromanfolio{147}อถ โข อายสฺมโต มหากสฺสปสฺส เอตทโหสิ “กหํ นุ โข ภควา เอตรหิ วิหรตี”ติ.\csromanspacer{} อทฺทสา โข อายสฺมา มหากสฺสโป ภควนฺตํ ทิพฺเพน จกฺขุนา ฯเปฯ ทิสฺวาน เสยฺยถาปิ นาม พลวา ปุริโส ฯเปฯ เอวเมว เชตวเน อนฺตรหิโต ตสฺมิํ พฺรหฺมโลเก ปาตุรโหสิ.\csromanspacer{} อถ โข อายสฺมา มหากสฺสโป ทกฺขิณํ ทิสํ นิสฺสาย ตสฺส พฺรหฺมุโน อุปริ เวหาสํ ปลฺลงฺเกน นิสีทิ เตโชธาตุํ สมาปชฺชิตฺวา นีจตรํ ภควโต.}

\prose{อถ โข อายสฺมโต มหากปฺปินสฺส เอตทโหสิ “กหํ นุ โข ภควา เอตรหิ วิหรตี”ติ.\csromanspacer{} อทฺทสา โข อายสฺมา มหากปฺปิโน ภควนฺตํ ทิพฺเพน จกฺขุนา ฯเปฯ เตโชธาตุํ สมาปนฺนํ, ทิสฺวาน เสยฺยถาปิ นาม พลวา ปุริโส ฯเปฯ เอวเมว เชตวเน อนฺตรหิโต ตสฺมิํ พฺรหฺมโลเก ปาตุรโหสิ.\csromanspacer{} อถ โข อายสฺมา มหากปฺปิโน ปจฺฉิมํ ทิสํ นิสฺสาย ตสฺส พฺรหฺมุโน อุปริ เวหาสํ ปลฺลงฺเกน นิสีทิ เตโชธาตุํ สมาปชฺชิตฺวา นีจตรํ ภควโต.}

\prose{อถ โข อายสฺมโต อนุรุทฺธสฺส เอตทโหสิ “กหํ นุ โข ภควา เอตรหิ วิหรตี”ติ.\csromanspacer{} อทฺทสา โข อายสฺมา อนุรุทฺโธ ฯเปฯ เตโชธาตุํ สมาปนฺนํ, ทิสฺวาน เสยฺยถาปิ นาม พลวา ปุริโส ฯเปฯ ตสฺมิํ พฺรหฺมโลเก ปาตุรโหสิ.\csromanspacer{} อถ โข อายสฺมา อนุรุทฺโธ อุตฺตรํ ทิสํ นิสฺสาย ตสฺส พฺรหฺมุโน อุปริ เวหาสํ ปลฺลงฺเกน นิสีทิ เตโชธาตุํ สมาปชฺชิตฺวา นีจตรํ ภควโต.}

\prose{อถ โข อายสฺมา มหาโมคฺคลฺลาโน ต พฺรหฺมานํ คาถาย อชฺฌภาสิ\csromandash{}}

\csromangathameasure{อชฺชาปิ เต อาวุโส สา ทิฏฺฐิ, ยา เต ทิฏฺฐิ ปุเร อหุ. \\ ปสฺสสิ วีติวตฺตนฺตํ, พฺรหฺมโลเก ปภสฺสรนฺติ. \\ น เม มาริส สา ทิฏฺฐิ, ยา เม ทิฏฺฐิ ปุเร อหุ. \\ ปสฺสามิ วีติวตฺตนฺตํ, พฺรหฺมโลเก ปภสฺสรํ. \\ สฺวาหํ อชฺช กถํ วชฺชํ, อหํ นิจฺโจมฺหิ สสฺสโตติ.}
\csromangathagroup{\csromangathastanza{อชฺชาปิ เต อาวุโส สา ทิฏฺฐิ, ยา เต ทิฏฺฐิ ปุเร อหุ. \\ ปสฺสสิ วีติวตฺตนฺตํ, พฺรหฺมโลเก ปภสฺสรนฺติ.}\csromangathastanza{น เม มาริส สา ทิฏฺฐิ, ยา เม ทิฏฺฐิ ปุเร อหุ. \\ ปสฺสามิ วีติวตฺตนฺตํ, พฺรหฺมโลเก ปภสฺสรํ. \\ สฺวาหํ อชฺช กถํ วชฺชํ, อหํ นิจฺโจมฺหิ สสฺสโตติ.}}

\prose{\csromanfolio{148}อถ โข ภควา ตํ พฺรหฺมานํ สํเวเชตฺวา เสยฺยถาปิ นาม พลวา ปุริโส สมิญฺชิตํ วา พาหํ ปสาเรยฺย, ปสาริตํ วา พาหํ สมิญฺเชยฺย, เอวเมว ตสฺมิํ พฺรหฺมโลเก อนฺตรหิโต เชตวเน ปาตุรโหสิ.\csromanspacer{} อถ โข โส พฺรหฺมา อญฺญตรํ พฺรหฺมปาริสชฺชํ อามนฺเตสิ “เอหิ ตฺวํ มาริส เยนายสฺมา มหาโมคฺคลฺลาโน เตนุปสงฺกม, อุปสงฺกมิตฺวา อายสฺมนฺตํ มหาโมคฺคลฺลานํ เอวํ วเทหิ ‘อตฺถิ นุ โข มาริส โมคฺคลฺลาน อญฺเญปิ ตสฺส ภควโต สาวกา เอวํมหิทฺธิกา เอวํมหานุภาวา เสยฺยถาปิ ภวํ โมคฺคลฺลาโน กสฺสโป กปฺปิโน อนุรุทฺโธ’ติ”. “เอวํ มาริสา”ติ โข โส พฺรหฺมปาริสชฺโช ตสฺส พฺรหฺมุโน ปฏิสฺสุตฺวา เยนายสฺมา มหาโมคฺคลฺลาโน เตนุปสงฺกมิ, อุปสงฺกมิตฺวา อายสฺมนฺตํ มหาโมคฺคลฺลานํ เอตทโวจ “อตฺถิ นุ โข มาริส โมคฺคลฺลาน อญฺเญปิ ตสฺส ภควโต สาวกา เอวํมหิทฺธิกา เอวํมหานุภาวา เสยฺยถาปิ ภวํ โมคฺคลฺลาโน กสฺสโป กปฺปิโน อนุรุทฺโธ”ติ.\csromanspacer{} อถ โข อายสฺมา มหาโมคฺคลฺลาโน ตํ พฺรหฺมปาริสชฺชํ คาถาย อชฺฌภาสิ\csromandash{}}

\csromangathameasure{เตวิชฺชา อิทฺธิปตฺตา จ, เจโตปริยายโกวิทา. \\ ขีณาสวา อรหนฺโต, พหู พุทฺธสฺส สาวกาติ.}
\csromangathagroup{\csromangathastanza{เตวิชฺชา อิทฺธิปตฺตา จ, เจโตปริยายโกวิทา. \\ ขีณาสวา อรหนฺโต, พหู พุทฺธสฺส สาวกาติ.}}

\prose{อถ โข โส พฺรหฺมปาริสชฺโช อายสฺมโต มหาโมคฺคลฺลานสฺส ภาสิตํ อภินนฺทิตฺวา อนุโมทิตฺวา เยน โส พฺรหฺมา เตนุปสงฺกมิ, อุปสงฺกมิตฺวา ตํ พฺรหฺมานํ เอตทโวจ “อายสฺมา มาริส มหาโมคฺคลฺลาโน เอวมาห\csromandash{}}

\csromangathameasure{“เตวิชฺชา อิทฺธิปตฺตา จ, เจโตปริยายโกวิทา. \\ ขีณาสวา อรหนฺโต, พหู พุทฺธสฺส สาวกา”ติ.}
\csromangathagroup{\csromangathastanza{“เตวิชฺชา อิทฺธิปตฺตา จ, เจโตปริยายโกวิทา. \\ ขีณาสวา อรหนฺโต, พหู พุทฺธสฺส สาวกา”ติ.}}

\prose{อิทมโวจ โส พฺรหฺมปาริสชฺโช.\csromanspacer{} อตฺตมโน จ โส พฺรหฺมา ตสฺส พฺรหฺมปาริสชฺชสฺส ภาสิตํ อภินนฺทีติ.}

\csromanheader{2}{6. พฺรหฺมโลกสุตฺต}
\proseitem{177}{สาวตฺถินิทานํ.\csromanspacer{} เตน โข ปน สมเยน ภควา ทิวาวิหารคโต โหติ ปฏิสลฺลีโน.\csromanspacer{} อถ โข สุพฺรหฺมา จ ปจฺเจกพฺรหฺมา}

\vspace{\tipitakaparskip}
\csromancenter{พฺรหฺมสํยุตฺต สุทฺธาวาโส จ ปจฺเจกพฺรหฺมา เยน ภควา เตนุปสงฺกมิํสุ, อุปสงฺกมิตฺวา ปจฺเจกํ ทฺวารพาหํ\footnote{ปจฺเจกทฺวารพาหํ (อิ, ก)} อุปนิสฺสาย อฏฺฐํสุ.\csromanspacer{} อถ โข สุพฺรหฺมา ปจฺเจกพฺรหฺมา สุทฺธาวาสํ ปจฺเจกพฺรหฺมานํ เอตทโวจ “อกาโล โข ตาว มาริส ภควนฺตํ ปยิรุปาสิตุํ, ทิวาวิหารคโต ภควา ปฏิสลฺลีโน จ, อสุโก จ พฺรหฺมโลโก อิทฺโธ เจว ผีโต จ, พฺรหฺมา จ ตตฺร ปมาทวิหารํ วิหรติ, อายาม มาริส เยน โส พฺรหฺมโลโก เตนุปสงฺกมิสฺสาม, อุปสงฺกมิตฺวา ตํ พฺรหฺมานํ สํเวเชยฺยามา”ติ. “เอวํ มาริสา”ติ โข สุทฺธาวาโส ปจฺเจกพฺรหฺมา สุพฺรหฺมุโน ปจฺเจกพฺรหฺมุโน ปจฺจสฺโสสิ.}
\prose{\csromanfolio{149}อถ โข สุพฺรหฺมา จ ปจฺเจกพฺรหฺมา สุทฺธาวาโส จ ปจฺเจกพฺรหฺมา เสยฺยถาปิ นาม พลวา ปุริโส ฯเปฯ เอวเมว ภควโต ปุรโต อนฺตรหิตา ตสฺมิํ พฺรหฺมโลเก ปาตุรเหสุํ.\csromanspacer{} อทฺทสา โข โส พฺรหฺมา เต พฺรหฺมาโน ทูรโตว อาคจฺฉนฺเต, ทิสฺวาน เต พฺรหฺมาโน เอตทโวจ “หนฺท กุโต นุ ตุเมฺห มาริสา อาคจฺฉถา”ติ.\csromanspacer{} อาคตา โข มยํ มาริส อมฺห ตสฺส ภควโต สนฺติกา อรหโต สมฺมาสมฺพุทฺธสฺส, คจฺเฉยฺยาสิ ปน ตฺวํ มาริส ตสฺส ภควโต อุปฏฺฐานํ อรหโต สมฺมาสมฺพุทฺธสฺสาติ.}

\prose{เอวํ วุตฺโต\footnote{เอวํ วุตฺเต (สี, สฺยา, กํ)} โข โส พฺรหฺมา ตํ วจนํ อนธิวาเสนฺโต สหสฺสกฺขตฺตุํ อตฺตานํ อภินิมฺมินิตฺวา สุพฺรหฺมานํ ปจฺเจกพฺรหฺมานํ เอตทโวจ “ปสฺสสิ เม โน ตฺวํ มาริส เอวรูปํ อิทฺธานุภาวนฺ”ติ.\csromanspacer{} ปสฺสามิ โข ตฺยาหํ มาริส เอวรูปํ อิทฺธานุภาวนฺติ.\csromanspacer{} โส ขฺวาหํ มาริส เอวํมหิทฺธิโก เอวํมหานุภาโว กสฺส อญฺญสฺส สมณสฺส วา พฺราหฺมณสฺส วา อุปฏฺฐานํ คมิสฺสามีติ.}

\prose{อถ โข สุพฺรหฺมา ปจฺเจกพฺรหฺมา ทฺวิสหสฺสกฺขตฺตุํ อตฺตานํ อภินิมฺมินิตฺวา ตํ พฺรหฺมานํ เอตทโวจ “ปสฺสสิ เม โน ตฺวํ มาริส เอวรูปํ อิทฺธานุภาวนฺ”ติ.\csromanspacer{} ปสฺสามิ โข ตฺยาหํ มาริส เอวรูปํ อิทฺธานุภาวนฺติ.\csromanspacer{} ตยา จ โข มาริส มยา จ เสฺวว ภควา มหิทฺธิกตโร เจว มหานุภาวตโร จ, คจฺเฉยฺยาสิ ตฺวํ มาริส ตสฺส ภควโต อุปฏฺฐานํ อรหโต สมฺมาสมฺพุทฺธสฺสาติ.\csromanspacer{} อถ โข โส พฺรหฺมา สุพฺรหฺมานํ ปจฺเจกพฺรหฺมานํ คาถาย อชฺฌภาสิ\csromandash{}}

\csromangathameasure{“ตโย สุปณฺณา จตุโร จ หํสา, \\ พฺยคฺฆีนิสา ปญฺจสตา จ ฌายิโน. \\ ตยิทํ วิมานํ ชลเต จ พฺรหฺเม, \\ โอภาสยํ อุตฺตรสฺสํ ทิสายนฺ”ติ. \\ กิญฺจาปิ เต ตํ ชลเต วิมานํ, \\ โอภาสยํ อุตฺตรสฺสํ ทิสายํ. \\ รูเป รณํ ทิสฺวา สทา ปเวธิตํ, \\ ตสฺมา น รูเป รมตี สุเมโธติ.}
\csromangathagroup{\csromangathastanza{\csromanfolio{150}“ตโย สุปณฺณา จตุโร จ หํสา, \\ พฺยคฺฆีนิสา ปญฺจสตา จ ฌายิโน. \\ ตยิทํ วิมานํ ชลเต จ\footnote{ชลเตว (อิ, ก)} พฺรหฺเม, \\ โอภาสยํ อุตฺตรสฺสํ ทิสายนฺ”ติ.}\csromangathastanza{กิญฺจาปิ เต ตํ ชลเต วิมานํ, \\ โอภาสยํ อุตฺตรสฺสํ ทิสายํ. \\ รูเป รณํ ทิสฺวา สทา ปเวธิตํ, \\ ตสฺมา น รูเป รมตี สุเมโธติ.}}

\prose{อถ โข สุพฺรหฺมา จ ปจฺเจกพฺรหฺมา สุทฺธาวาโส จ ปจฺเจกพฺรหฺมา ตํ พฺรหฺมานํ สํเวเชตฺวา ตตฺเถวนฺตรธายิํสุ.\csromanspacer{} อคมาสิ จ โข โส พฺรหฺมา อปเรน สมเยน ภควโต อุปฏฺฐานํ อรหโต สมฺมาสมฺพุทฺธสฺสาติ.}

\csromanheader{2}{7. โกกาลิกสุตฺต}
\proseitem{178}{สาวตฺถินิทานํ.\csromanspacer{} เตน โข ปน สมเยน ภควา ทิวาวิหารคโต โหติ ปฏิสลฺลีโน.\csromanspacer{} อถ โข สุพฺรหฺมา จ ปจฺเจกพฺรหฺมา สุทฺธาวาโส จ ปจฺเจกพฺรหฺมา เยน ภควา เตนุปสงฺกมิํสุ, อุปสงฺกมิตฺวา ปจฺเจกํ ทฺวารพาหํ นิสฺสาย อฏฺฐํสุ.\csromanspacer{} อถ โข สุพฺรหฺมา ปจฺเจกพฺรหฺมา โกกาลิกํ ภิกฺขุํ อารพฺภ ภควโต สนฺติเก}

\csromangathameasure{“อปฺปเมยฺยํ ปมินนฺโต, โกธ วิทฺวา วิกปฺปเย. \\ อปฺปเมยฺยํ ปมายินํ, นิวุตํ ตํ มญฺเญ ปุถุชฺชนนฺ”ติ.}
\csromangathagroup{\csromangathastanza{“อปฺปเมยฺยํ ปมินนฺโต, โกธ วิทฺวา วิกปฺปเย. \\ อปฺปเมยฺยํ ปมายินํ, นิวุตํ ตํ มญฺเญ ปุถุชฺชนนฺ”ติ.}}

\csromanheader{2}{8. กตโมทกติสฺสสุตฺต}
\proseitem{179}{สาวตฺถินิทานํ.\csromanspacer{} เตน โข ปน สมเยน ภควา ทิวาวิหารคโต โหติ ปฏิสลฺลีโน.\csromanspacer{} อถ โข สุพฺรหฺมา จ ปจฺเจกพฺรหฺมา สุทฺธาวาโส จ ปจฺเจกพฺรหฺมา เยน ภควา เตนุปสงฺกมิํสุ, อุปสงฺกมิตฺวา ปจฺเจกํ ทฺวารพาหํ นิสฺสาย อฏฺฐํสุ.\csromanspacer{} อถ โข สุทฺธาวาโส ปจฺเจกพฺรหฺมา กตโมทกติสฺสกํ\footnote{กตโมรกติสฺสกํ (สี, สฺยา, กํ)} ภิกฺขุํ อารพฺภ ภควโต สนฺติเก อิมํ คาถํ อภาสิ\csromandash{}}

\csromanheader{2}{6. พฺรหฺมสํยุตฺต}
\csromangathameasure{“อปฺปเมยฺยํ ปมินนฺโต, โกธ วิทฺวา วิกปฺปเย. \\ อปฺปเมยฺยํ ปมายินํ, นิวุตํ ตํ มญฺเญ อกิสฺสวนฺ”ติ.}
\csromangathagroup{\csromangathastanza{\csromanfolio{151}“อปฺปเมยฺยํ ปมินนฺโต, โกธ วิทฺวา วิกปฺปเย. \\ อปฺปเมยฺยํ ปมายินํ, นิวุตํ ตํ มญฺเญ อกิสฺสวนฺ”ติ.}}

\csromanheader{2}{9. ตุรูพฺรหฺมสุตฺต}
\proseitem{180}{สาวตฺถินิทานํ.\csromanspacer{} เตน โข ปน สมเยน โกกาลิโก ภิกฺขุ อาพาธิโก โหติ ทุกฺขิโต พาฬฺหคิลาโน.\csromanspacer{} อถ โข ตุรู\footnote{ตุทุ (สี, สฺยา, กํ, อิ)} ปจฺเจกพฺรหฺมา อภิกฺกนฺตาย รตฺติยา อภิกฺกนฺตวณฺโณ เกวลกปฺปํ เชตวนํ โอภาเสตฺวา เยน โกกาลิโก ภิกฺขุ เตนุปสงฺกมิ, อุปสงฺกมิตฺวา เวหาสํ ฐิโต โกกาลิกํ ภิกฺขุํ เอตทโวจ “ปสาเทหิ โกกาลิก สาริปุตฺตโมคฺคลฺลาเนสุ จิตฺตํ, เปสลา สาริปุตฺตโมคฺคลฺลานา”ติ.\csromanspacer{} โกสิ ตฺวํ อาวุโสติ.\csromanspacer{} อหํ ตุรู ปจฺเจกพฺรหฺมาติ.\csromanspacer{} นนุ ตฺวํ อาวุโส ภควตา อนาคามี พฺยากโต, อถ กิญฺจรหิ อิธาคโต, ปสฺส ยาวญฺจ เต อิทํ อปรทฺธนฺติ.}

\csromangathameasure{ปุริสสฺส หิ ชาตสฺส, กุฐารี ชายเต มุเข. \\ ยาย ฉินฺทติ อตฺตานํ, พาโล ทุพฺภาสิตํ ภณํ. \\ โย นินฺทิยํ ปสํสติ, \\ ตํ วา นินฺทติ โย ปสํสิโย. \\ วิจินาติ มุเขน โส กลิํ, \\ กลินา เตน สุขํ น วินฺทติ. \\ อปฺปมตฺตโก อยํ กลิ, \\ โย อกฺเขสุ ธนปราชโย. \\ สพฺพสฺสาปิ สหาปิ อตฺตนา, \\ อยเมว มหนฺตตโร กลิ. \\ โย สุคเตสุ มนํ ปโทสเย. \\ สตํ สหสฺสานํ นิรพฺพุทานํ, \\ ฉตฺติํสติ ปญฺจ จ อพฺพุทานิ. \\ ยมริยครหี นิรยํ อุเปติ, \\ วาจํ มนญฺจ ปณิธาย ปาปกนฺติ.}
\csromangathagroup{\csromangathastanza{ปุริสสฺส หิ ชาตสฺส, กุฐารี\footnote{กุธารี (สฺยา, กํ, ก)} ชายเต มุเข. \\ ยาย ฉินฺทติ อตฺตานํ, พาโล ทุพฺภาสิตํ ภณํ.}\csromangathastanza{โย นินฺทิยํ ปสํสติ, \\ ตํ วา นินฺทติ โย ปสํสิโย. \\ วิจินาติ มุเขน โส กลิํ, \\ กลินา เตน สุขํ น วินฺทติ.}\csromangathastanza{อปฺปมตฺตโก อยํ กลิ, \\ โย อกฺเขสุ ธนปราชโย. \\ สพฺพสฺสาปิ สหาปิ อตฺตนา, \\ อยเมว มหนฺตตโร กลิ.}\csromangathastanza{โย สุคเตสุ มนํ ปโทสเย. \\ สตํ สหสฺสานํ นิรพฺพุทานํ, \\ ฉตฺติํสติ ปญฺจ จ อพฺพุทานิ. \\ ยมริยครหี\footnote{ยมริเย ครหี (สฺยา, กํ), ยมริยํ ครหํ (ก)} นิรยํ อุเปติ,}\csromangathastanza{วาจํ มนญฺจ ปณิธาย ปาปกนฺติ.}}

\csromanheader{2}{10. โกกาลิกสุตฺต}
\proseitem{181}{\csromanfolio{152}สาวตฺถินิทานํ.\csromanspacer{} อถ โข โกกาลิโก ภิกฺขุ เยน ภควา เตนุปสงฺกมิ, อุปสงฺกมิตฺวา ภควนฺตํ อภิวาเทตฺวา เอกมนฺตํ นิสีทิ, เอกมนฺตํ นิสินฺโน โข โกกาลิโก ภิกฺขุ ภควนฺตํ เอตทโวจ “ปาปิจฺฉา ภนฺเต สาริปุตฺตโมคฺคลฺลานา, ปาปิกานํ อิจฺฉานํ วสํ คตา”ติ.\csromanspacer{} เอวํ วุตฺเต ภควา โกกาลิกํ ภิกฺขุํ เอตทโวจ “มา เหวํ โกกาลิก อวจ, มา เหวํ โกกาลิก อวจ, ปสาเทหิ โกกาลิก สาริปุตฺตโมคฺคลฺลาเนสุ จิตฺตํ, เปสลา สาริปุตฺตโมคฺคลฺลานา”ติ.\csromanspacer{} ทุติยมฺปิ โข โกกาลิโก ภิกฺขุ ภควนฺตํ เอตทโวจ “กิญฺจาปิ เม ภนฺเต ภควา สทฺธายิโก ปจฺจยิโก, อถ โข ปาปิจฺฉาว ภนฺเต สาริปุตฺตโมคฺคลฺลานา, ปาปิกานํ อิจฺฉานํ วสํ คตา”ติ.\csromanspacer{} ทุติยมฺปิ โข ภควา โกกาลิกํ ภิกฺขุํ เอตทโวจ “มา เหวํ โกกาลิก อวจ, มา เหวํ โกกาลิก อวจ, ปสาเทหิ โกกาลิก สาริปุตฺตโมคฺคลฺลาเนสุ จิตฺตํ, เปสลา สาริปุตฺตโมคฺคลฺลานา”ติ.\csromanspacer{} ตติยมฺปิ โข โกกาลิโก ภิกฺขุ ภควนฺตํ เอตทโวจ “กิญฺจาปิ ฯเปฯ อิจฺฉานํ วสํ คตา”ติ.\csromanspacer{} ตติยมฺปิ โข ภควา โกกาลิกํ ภิกฺขุํ เอตทโวจ “มา เหวํ ฯเปฯ เปสลา สาริปุตฺตโมคฺคลฺลานา”ติ.}

\prose{อถ โข โกกาลิโก ภิกฺขุ อุฏฺฐายาสนา ภควนฺตํ อภิวาเทตฺวา ปทกฺขิณํ กตฺวา ปกฺกามิ.\csromanspacer{} อจิรปกฺกนฺตสฺส จ โกกาลิกสฺส ภิกฺขุโน สาสปมตฺตีหิ ปีฬกาหิ\footnote{ปิฬกาหิ (สี, อิ)} สพฺโพ กาโย ผุโฏ อโหสิ, สาสปมตฺติโย หุตฺวา มุคฺคมตฺติโย อเหสุํ, มุคฺคมตฺติโย หุตฺวา กลายมตฺติโย อเหสุํ, กลายมตฺติโย หุตฺวา โกลฏฺฐิมตฺติโย อเหสุํ, โกลฏฺฐิมตฺติโย หุตฺวา โกลมตฺติโย อเหสุํ, โกลมตฺติโย หุตฺวา อามลกมตฺติโย อเหสุํ, อามลกมตฺติโย หุตฺวา เพลุวสลาฏุกมตฺติโย อเหสุํ, เพลุวสลาฏุกมตฺติโย หุตฺวา พิลฺลมตฺติโย อเหสุํ, พิลฺลมตฺติโย หุตฺวา ปภิชฺชิํสุ, ปุพฺพญฺจ โลหิตญฺจ ปคฺฆริํสุ.\csromanspacer{} อถ โข โกกาลิโก ภิกฺขุ เตเนว อาพาเธน กาลมกาสิ, กาลงฺกโต จ โกกาลิโก ภิกฺขุ ปทุมํ นิรยํ อุปปชฺชิ สาริปุตฺตโมคฺคลฺลาเนสุ จิตฺตํ อาฆาเตตฺวา.}

\csromanheader{2}{6. พฺรหฺมสํยุตฺต}
\prose{\csromanfolio{153}อถ โข พฺรหฺมา สหมฺปติ อภิกฺกนฺตาย รตฺติยา อภิกฺกนฺตวณฺโณ เกวลกปฺปํ เชตวนํ โอภาเสตฺวา เยน ภควา เตนุปสงฺกมิ, อุปสงฺกมิตฺวา ภควนฺตํ อภิวาเทตฺวา เอกมนฺตํ อฏฺฐาสิ, เอกมนฺตํ ฐิโต โข พฺรหฺมา สหมฺปติ ภควนฺตํ เอตทโวจ “โกกาลิโก ภนฺเต ภิกฺขุ กาลงฺกโต, กาลงฺกโต จ ภนฺเต โกกาลิโก ภิกฺขุ ปทุมํ นิรยํ อุปปนฺโน สาริปุตฺตโมคฺคลฺลาเนสุ จิตฺตํ อาฆาเตตฺวา”ติ.\csromanspacer{} อิทมโวจ พฺรหฺมา สหมฺปติ, อิทํ วตฺวา ภควนฺตํ อภิวาเทตฺวา ปทกฺขิณํ กตฺวา ตตฺเถวนฺตรธายีติ.}

\prose{อถ โข ภควา ตสฺสา รตฺติยา อจฺจเยน ภิกฺขู อามนฺเตสิ “อิมํ ภิกฺขเว รตฺติํ พฺรหฺมา สหมฺปติ อภิกฺกนฺตาย รตฺติยา อภิกฺกนฺตวณฺโณ เกวลกปฺปํ เชตวนํ โอภาเสตฺวา เยนาหํ เตนุปสงฺกมิ, อุปสงฺกมิตฺวา มํ อภิวาเทตฺวา เอกมนฺตํ อฏฺฐาสิ, เอกมนฺตํ ฐิโต โข ภิกฺขเว พฺรหฺมา สหมฺปติ มํ เอตทโวจ ‘โกกาลิโก ภนฺเต ภิกฺขุ กาลงฺกโต, กาลงฺกโต จ ภนฺเต โกกาลิโก ภิกฺขุ ปทุมํ นิรยํ อุปปนฺโน สาริปุตฺตโมคฺคลฺลาเนสุ จิตฺตํ อาฆาเตตฺวา’ติ.\csromanspacer{} อิทมโวจ ภิกฺขเว พฺรหฺมา สหมฺปติ, อิทํ วตฺวา มํ อภิวาเทตฺวา ปทกฺขิณํ กตฺวา ตตฺเถวนฺตรธายี”ติ.}

\prose{เอวํ วุตฺเต อญฺญตโร ภิกฺขุ ภควนฺตํ เอตทโวจ “กีวทีฆํ นุ โข ภนฺเต ปทุเม นิรเย อายุปฺปมาณนฺ”ติ.\csromanspacer{} ทีฆํ โข ภิกฺขุ ปทุเม นิรเย อายุปฺปมาณํ, ตํ น สุกรํ สงฺขาตุํ\csromandash{}เอตฺตกานิ วสฺสานิ อิติ วา, เอตฺตกานิ วสฺสสตานิ อิติ วา, เอตฺตกานิ วสฺสสหสฺสานิ อิติ วา, เอตฺตกานิ วสฺสสตสหสฺสานิ อิติ วาติ.\csromanspacer{} สกฺกา ปน ภนฺเต อุปมํ กาตุนฺติ. “สกฺกา ภิกฺขู”ติ ภควา อโวจ\csromandash{}}

\prose{เสยฺยถาปิ ภิกฺขุ วีสติขาริโก โกสลโก ติลวาโห, ตโต ปุริโส วสฺสสตสฺส วสฺสสตสฺส อจฺจเยน เอกเมกํ ติลํ อุทฺธเรยฺย, ขิปฺปตรํ โข โส ภิกฺขุ วีสติขาริโก โกสลโก ติลวาโห อิมินา อุปกฺกเมน ปริกฺขยํ ปริยาทานํ คจฺเฉยฺย, น เตฺวว เอโก อพฺพุโท นิรโย.\csromanspacer{} เสยฺยถาปิ ภิกฺขุ วีสติ อพฺพุทา นิรยา, เอวเมโก นิรพฺพุทนิรโย.\csromanspacer{} เสยฺยถาปิ ภิกฺขุ วีสติ นิรพฺพุทา นิรยา, เอวเมโก อพโพ นิรโย.\csromanspacer{} เสยฺยถาปิ ภิกฺขุ วีสติ อพพา นิรยา,}

\prose{\csromanfolio{154}เอวเมโก อฏโฏ นิรโย.\csromanspacer{} เสยฺยถาปิ ภิกฺขุ วีสติ อฏฏา นิรยา, เอวเมโก อหโห นิรโย.\csromanspacer{} เสยฺยถาปิ ภิกฺขุ วีสติ อหหา นิรยา, เอวเมโก กุมุโท นิรโย.\csromanspacer{} เสยฺยถาปิ ภิกฺขุ วีสติ กุมุทา นิรยา, เอวเมโก โสคนฺธิโก นิรโย.\csromanspacer{} เสยฺยถาปิ ภิกฺขุ วีสติ โสคนฺธิกา นิรยา, เอวเมโก อุปฺปลนิรโย.\csromanspacer{} เสยฺยถาปิ ภิกฺขุ วีสติ อุปฺปลา นิรยา, เอวเมโก ปุณฺฑริโก นิรโย.\csromanspacer{} เสยฺยถาปิ ภิกฺขุ วีสติ ปุณฺฑริกา นิรยา, เอวเมโก ปทุโม นิรโย.\csromanspacer{} ปทุเม ปน ภิกฺขุ นิรเย โกกาลิโก ภิกฺขุ อุปปนฺโน สาริปุตฺตโมคฺคลฺลาเนสุ จิตฺตํ อาฆาเตตฺวาติ.\csromanspacer{} อิทมโวจ ภควา, อิทํ วตฺวาน สุคโต อถาปรํ เอตทโวจ สตฺถา\csromandash{}}

\csromangathameasure{“ปุริสสฺส หิ ชาตสฺส, \\ กุฐารี ชายเต มุเข. \\ ยาย ฉินฺทติ อตฺตานํ, \\ พาโล ทุพฺภาสิตํ ภณํ. \\ โย นินฺทิยํ ปสํสติ, \\ ตํ วา นินฺทติ โย ปสํสิโย. \\ วิจินาติ มุเขน โส กลิํ, \\ กลินา เตน สุขํ น วินฺทติ. \\ อปฺปมตฺตโก อยํ กลิ, \\ โย อกฺเขสุ ธนปราชโย. \\ สพฺพสฺสาปิ สหาปิ อตฺตนา, \\ อยเมว มหนฺตตโร กลิ. \\ โย สุคเตสุ มนํ ปโทสเย. \\ สตํ สหสฺสานํ นิรพฺพุทานํ, \\ ฉตฺติํสติ ปญฺจ จ อพฺพุทานิ. \\ ยมริยครหี นิรยํ อุเปติ, \\ วาจํ มนญฺจ ปณิธาย ปาปกนฺ”ติ.}
\csromangathagroup{\csromangathastanza{“ปุริสสฺส หิ ชาตสฺส, \\ กุฐารี ชายเต มุเข. \\ ยาย ฉินฺทติ อตฺตานํ, \\ พาโล ทุพฺภาสิตํ ภณํ.}\csromangathastanza{โย นินฺทิยํ ปสํสติ, \\ ตํ วา นินฺทติ โย ปสํสิโย. \\ วิจินาติ มุเขน โส กลิํ, \\ กลินา เตน สุขํ น วินฺทติ.}\csromangathastanza{อปฺปมตฺตโก อยํ กลิ, \\ โย อกฺเขสุ ธนปราชโย. \\ สพฺพสฺสาปิ สหาปิ อตฺตนา, \\ อยเมว มหนฺตตโร กลิ.}\csromangathastanza{โย สุคเตสุ มนํ ปโทสเย. \\ สตํ สหสฺสานํ นิรพฺพุทานํ, \\ ฉตฺติํสติ ปญฺจ จ อพฺพุทานิ. \\ ยมริยครหี นิรยํ อุเปติ,}\csromangathastanza{วาจํ มนญฺจ ปณิธาย ปาปกนฺ”ติ.}}

\vspace{\tipitakaparskip}
\csromancenter{ปฐโม วคฺโค.}
\titlehead{6. พฺรหฺมสํยุตฺต \textbf{ตสฺสุทฺทานํ}}
\csromangathameasure{อายาจนํ คารโว พฺรหฺมเทโว, \\ พโก จ พฺรหฺมา อปรา จ ทิฏฺฐิ. \\ ปมาทโกกาลิกติสฺสโก จ, \\ ตุรู จ พฺรหฺมา อปโร จ โกกาลิโกติ.}
\csromangathagroup{\csromangathastanza{\csromanfolio{155}อายาจนํ คารโว พฺรหฺมเทโว, \\ พโก จ พฺรหฺมา อปรา จ ทิฏฺฐิ. \\ ปมาทโกกาลิกติสฺสโก จ, \\ ตุรู จ พฺรหฺมา อปโร จ โกกาลิโกติ.}}

\chapterhead{2. ทุติยวคฺค}
\csromanheader{2}{1. สนงฺกุมารสุตฺต}
\proseitem{182}{เอวํ เม สุตํ\csromandash{}เอกํ สมยํ ภควา ราชคเห วิหรติ สปฺปินีตีเร.\csromanspacer{} อถ โข พฺรหฺมา สนงฺกุมาโร อภิกฺกนฺตาย รตฺติยา อภิกฺกนฺตวณฺโณ เกวลกปฺปํ สปฺปินีตีรํ โอภาเสตฺวา เยน ภควา เตนุปสงฺกมิ, อุปสงฺกมิตฺวา ภควนฺตํ อภิวาเทตฺวา เอกมนฺตํ อฏฺฐาสิ, เอกมนฺตํ ฐิโต โข พฺรหฺมา สนงฺกุมาโร ภควโต สนฺติเก อิมํ คาถํ อภาสิ\csromandash{}}

\csromangathameasure{“ขตฺติโย เสฏฺโฐ ชเนตสฺมิํ, เย โคตฺตปฏิสาริโน. \\ วิชฺชาจรณสมฺปนฺโน, โส เสฏฺโฐ เทวมานุเส”ติ.}
\csromangathagroup{\csromangathastanza{“ขตฺติโย เสฏฺโฐ ชเนตสฺมิํ, เย โคตฺตปฏิสาริโน. \\ วิชฺชาจรณสมฺปนฺโน, โส เสฏฺโฐ เทวมานุเส”ติ.}}

\prose{อิทมโวจ พฺรหฺมา สนงฺกุมาโร, สมนุญฺโญ สตฺถา อโหสิ.\csromanspacer{} อถ โข พฺรหฺมา สนงฺกุมาโร “สมนุญฺโญ เม สตฺถา”ติ ภควนฺตํ อภิวาเทตฺวา ปทกฺขิณํ กตฺวา ตตฺเถวนฺตรธายีติ.}

\csromanheader{2}{2. เทวทตฺตสุตฺต}
\proseitem{183}{เอวํ เม สุตํ\csromandash{}เอกํ สมยํ ภควา ราชคเห วิหรติ คิชฺฌกูเฏ ปพฺพเต อจิรปกฺกนฺเต เทวทตฺเต.\csromanspacer{} อถ โข พฺรหฺมา สหมฺปติ อภิกฺกนฺตาย รตฺติยา อภิกฺกนฺตวณฺโณ เกวลกปฺปํ คิชฺฌกูฏํ ปพฺพตํ โอภาเสตฺวา เยน ภควา เตนุปสงฺกมิ, อุปสงฺกมิตฺวา ภควนฺตํ อภิวาเทตฺวา เอกมนฺตํ อฏฺฐาสิ, เอกมนฺตํ ฐิโต โข พฺรหฺมา สหมฺปติ เทวทตฺตํ อารพฺภ ภควโต สนฺติเก อิมํ คาถํ อภาสิ\csromandash{}}

\csromangathameasure{“ผลํ เว กทลิํ หนฺติ, ผลํ เวฬุํ ผลํ นฬํ. \\ สกฺกาโร กาปุริสํ หนฺติ, คพฺโภ อสฺสตริํ ยถา”ติ.}
\csromangathagroup{\csromangathastanza{\csromanfolio{156}“ผลํ เว กทลิํ หนฺติ, ผลํ เวฬุํ ผลํ นฬํ. \\ สกฺกาโร กาปุริสํ หนฺติ, คพฺโภ อสฺสตริํ ยถา”ติ.}}

\csromanheader{2}{3. อนฺธกวินฺทสุตฺต}
\proseitem{184}{เอกํ สมยํ ภควา มาคเธสุ วิหรติ อนฺธกวินฺเท.\csromanspacer{} เตน โข ปน สมเยน ภควา รตฺตนฺธการติมิสายํ อพฺโภกาเส นิสินฺโน โหติ, เทโว จ เอกเมกํ ผุสายติ.\csromanspacer{} อถ โข พฺรหฺมา สหมฺปติ อภิกฺกนฺตาย รตฺติยา อภิกฺกนฺตวณฺโณ เกวลกปฺปํ อนฺธกวินฺทํ โอภาเสตฺวา เยน ภควา เตนุปสงฺกมิ, อุปสงฺกมิตฺวา ภควนฺตํ อภิวาเทตฺวา เอกมนฺตํ อฏฺฐาสิ, เอกมนฺตํ ฐิโต โข พฺรหฺมา สหมฺปติ ภควโต สนฺติเก อิมา คาถาโย อภาสิ\csromandash{}}

\csromangathameasure{“เสเวถ ปนฺตานิ เสนาสนานิ, \\ จเรยฺย สํโยชนวิปฺปโมกฺขา. \\ สเจ รติํ นาธิคจฺเฉยฺย ตตฺถ, \\ สํเฆ วเส รกฺขิตตฺโต สตีมา. \\ กุลากุลํ ปิณฺฑิกาย จรนฺโต, \\ อินฺทฺริยคุตฺโต นิปโก สติมา. \\ เสเวถ ปนฺตานิ เสนาสนานิ, \\ ภยา ปมุตฺโต อภเย วิมุตฺโต. \\ ยตฺถ เภรวา สรีสปา, \\ วิชฺชุ สญฺจรติ ถนยติ เทโว. \\ อนฺธการติมิสาย รตฺติยา, \\ นิสีทิ ตตฺถ ภิกฺขุ วิคตโลมหํโส. \\ อิทํ หิ ชาตุ เม ทิฏฺฐํ, น ยิทํ อิติหีติหํ. \\ เอกสฺมิํ พฺรหฺมจริยสฺมิํ, สหสฺสํ มจฺจุหายินํ. \\ ภิยฺโย ปญฺจสตา เสกฺขา, ทสา จ ทสธา ทส. \\ สพฺเพ โสตสมาปนฺนา, อติรจฺฉานคามิโน.}
\csromangathagroup{\csromangathastanza{“เสเวถ ปนฺตานิ เสนาสนานิ, \\ จเรยฺย สํโยชนวิปฺปโมกฺขา. \\ สเจ รติํ นาธิคจฺเฉยฺย ตตฺถ, \\ สํเฆ วเส รกฺขิตตฺโต สตีมา.}\csromangathastanza{กุลากุลํ ปิณฺฑิกาย จรนฺโต, \\ อินฺทฺริยคุตฺโต นิปโก สติมา. \\ เสเวถ ปนฺตานิ เสนาสนานิ, \\ ภยา ปมุตฺโต อภเย วิมุตฺโต.}\csromangathastanza{ยตฺถ เภรวา สรีสปา\footnote{สิริํสปา (สี, สฺยา, กํ, อิ)}, \\ วิชฺชุ สญฺจรติ ถนยติ เทโว. \\ อนฺธการติมิสาย รตฺติยา, \\ นิสีทิ ตตฺถ ภิกฺขุ วิคตโลมหํโส. \\ อิทํ หิ ชาตุ เม ทิฏฺฐํ, น ยิทํ อิติหีติหํ. \\ เอกสฺมิํ พฺรหฺมจริยสฺมิํ, สหสฺสํ มจฺจุหายินํ.}\csromangathastanza{ภิยฺโย\footnote{ภีโย (สี, สฺยา, กํ, อิ)} ปญฺจสตา เสกฺขา, ทสา จ ทสธา ทส. \\ สพฺเพ โสตสมาปนฺนา, อติรจฺฉานคามิโน.}}

\csromanheader{2}{6. พฺรหฺมสํยุตฺต}
\csromangathameasure{อถายํ อิตรา ปชา, ปุญฺญภาคาติ เม มโน. \\ สงฺขาตุํ โนปิ สกฺโกมิ, มุสาวาทสฺส โอตฺตปนฺ”ติ.}
\csromangathagroup{\csromangathastanza{\csromanfolio{157}อถายํ\footnote{อตฺถายํ-อิติปิ ที 2 (176) ปิฏฺเฐ.} อิตรา ปชา, ปุญฺญภาคาติ เม มโน. \\ สงฺขาตุํ โนปิ สกฺโกมิ, มุสาวาทสฺส โอตฺตปนฺ”ติ\footnote{อตฺถายํ-อิติปิ ที 2 (176) ปิฏฺเฐ.}.}}

\csromanheader{2}{4. อรุณวตีสุตฺต}
\proseitem{185}{เอวํ เม สุตํ\csromandash{}เอกํ สมยํ ภควา สาวตฺถิยํ วิหรติ ฯเปฯ.\csromanspacer{} ตตฺร โข ภควา ภิกฺขู อามนฺเตสิ “ภิกฺขโว”ติ. “ภทนฺเต”ติ เต ภิกฺขู ภควโต ปจฺจสฺโสสุํ.\csromanspacer{} ภควา เอตทโวจ\csromandash{}}

\prose{ภูตปุพฺพํ ภิกฺขเว ราชา อโหสิ อรุณวา นาม.\csromanspacer{} รญฺโญ โข ปน ภิกฺขเว อรุณวโต อรุณวตี นาม ราชธานี อโหสิ.\csromanspacer{} อรุณวติํ โข ปน ภิกฺขเว ราชธานิํ\footnote{อรุณวติยํ โข ปน ภิกฺขเว ราชธานิยํ (อิ, ก)} สิขี ภควา อรหํ สมฺมาสมฺพุทฺโธ อุปนิสฺสาย วิหาสิ.\csromanspacer{} สิขิสฺส โข ปน ภิกฺขเว ภควโต อรหโต สมฺมาสมฺพุทฺธสฺส อภิภูสมฺภวํ นาม สาวกยุคํ อโหสิ อคฺคํ ภทฺทยุคํ.\csromanspacer{} อถ โข ภิกฺขเว สิขี ภควา อรหํ สมฺมาสมฺพุทฺโธ อภิภุํ ภิกฺขุํ อามนฺเตสิ “อายาม พฺราหฺมณ เยน อญฺญตโร พฺรหฺมโลโก เตนุปสงฺกมิสฺสาม ยาว ภตฺตสฺส กาโล ภวิสฺสตี”ติ. “เอวํ ภนฺเต”ติ โข ภิกฺขเว อภิภู ภิกฺขุ สิขิสฺส ภควโต อรหโต สมฺมาสมฺพุทฺธสฺส ปจฺจสฺโสสิ.\csromanspacer{} อถ โข ภิกฺขเว สิขี ภควา อรหํ สมฺมาสมฺพุทฺโธ อภิภู จ ภิกฺขุ เสยฺยถาปิ นาม พลวา ปุริโส สมิญฺชิตํ วา พาหํ ปสาเรยฺย, ปสาริตํ วา พาหํ สมิญฺเชยฺย, เอวเมว อรุณวติยา ราชธานิยา อนฺตรหิตา ตสฺมิํ พฺรหฺมโลเก ปาตุรเหสุํ.}

\prose{อถ โข ภิกฺขเว สิขี ภควา อรหํ สมฺมาสมฺพุทฺโธ อภิภุํ ภิกฺขุํ อามนฺเตสิ “ปฏิภาตุ พฺราหฺมณ ตํ พฺรหฺมุโน จ พฺรหฺมปริสาย จ พฺรหฺมปาริสชฺชานญฺจ ธมฺมี กถา”ติ. “เอวํ ภนฺเต”ติ โข ภิกฺขเว อภิภู ภิกฺขุ สิขิสฺส ภควโต อรหโต สมฺมาสมฺพุทฺธสฺส ปฏิสฺสุตฺวา พฺรหฺมานญฺจ พฺรหฺมปริสญฺจ พฺรหฺมปาริสชฺเช จ ธมฺมิยา กถาย สนฺทสฺเสสิ สมาทเปสิ สมุตฺเตเชสิ สมฺปหํเสสิ.\csromanspacer{} ตตฺร สุทํ ภิกฺขเว พฺรหฺมา จ พฺรหฺมปริสา จ พฺรหฺมปาริสชฺชา จ อุชฺฌายนฺติ ขิยฺยนฺติ\footnote{ขียนฺติ (สี, สฺยา, กํ, อิ)} วิปาเจนฺติ “อจฺฉริยํ \csromanfolio{158}วต โภ, อพฺภุตํ วต โภ, กถํ หิ นาม สตฺถริ สมฺมุขีภูเต สาวโก ธมฺมํ เทเสสฺสตี”ติ.}

\prose{อถ โข ภิกฺขเว สิขี ภควา อรหํ สมฺมาสมฺพุทฺโธ อภิภุํ ภิกฺขุํ อามนฺเตสิ “อุชฺฌายนฺติ โข เต พฺราหฺมณ พฺรหฺมา จ พฺรหฺมปริสา จ พฺรหฺมปาริสชฺชา จ ‘อจฺฉริยํ วต โภ, อพฺภุตํ วต โภ, กถํ หิ นาม สตฺถริ สมฺมุขีภูเต สาวโก ธมฺมํ เทเสสฺสตี’ติ.\csromanspacer{} เตน หิ ตฺวํ พฺราหฺมณ ภิยฺโยโส มตฺตาย พฺรหฺมานญฺจ พฺรหฺมปริสญฺจ พฺรหฺมปาริสชฺเช จ สํเวเชหี”ติ. “เอวํ ภนฺเต”ติ โข ภิกฺขเว อภิภู ภิกฺขุ สิขิสฺส ภควโต อรหโต สมฺมา- สมฺพุทฺธสฺส ปฏิสฺสุตฺวา ทิสฺสมาเนนปิ กาเยน ธมฺมํ เทเสสิ, อทิสฺสมาเนนปิ กาเยน ธมฺมํ เทเสสิ, ทิสฺสมาเนนปิ เหฏฺฐิเมน อุปฑฺฒกาเยน อทิสฺสมาเนน อุปริเมน อุปฑฺฒกาเยน ธมฺมํ เทเสสิ, ทิสฺสมาเนนปิ อุปริเมน อุปฑฺฒกาเยน อทิสฺสมาเนน เหฏฺฐิเมน อุปฑฺฒกาเยน ธมฺมํ เทเสสิ.\csromanspacer{} ตตฺร สุทํ ภิกฺขเว พฺรหฺมา จ พฺรหฺมปริสา จ พฺรหฺมปาริสชฺชา จ อจฺฉริยพฺภุตจิตฺตชาตา อเหสุํ “อจฺฉริยํ วต โภ, อพฺภุตํ วต โภ สมณสฺส มหิทฺธิกตา มหานุภาวตา”ติ.}

\prose{อถ โข อภิภู ภิกฺขุ สิขิํ ภควนฺตํ อรหนฺตํ สมฺมาสมฺพุทฺธํ เอตทโวจ “อภิชานามิ ขฺวาหํ ภนฺเต ภิกฺขุสํฆสฺส มชฺเฌ เอวรูปิํ วาจํ ภาสิตา ‘ปโหมิ ขฺวาหํ อาวุโส พฺรหฺมโลเก ฐิโต สหสฺสิโลกธาตุํ\footnote{สหสฺสีโลกธาตุํ (สี, สฺยา, กํ, อิ)} สเรน วิญฺญาเปตุนฺ’ติ”.\csromanspacer{} เอตสฺส พฺราหฺมณ กาโล, เอตสฺส พฺราหฺมณ กาโล, ยํ ตฺวํ พฺราหฺมณ พฺรหฺมโลเก ฐิโต สหสฺสิโลกธาตุํ สเรน วิญฺญาเปยฺยาสีติ. “เอวํ ภนฺเต”ติ โข ภิกฺขเว อภิภู ภิกฺขุ สิขิสฺส ภควโต อรหโต สมฺมาสมฺพุทฺธสฺส ปฏิสฺสุตฺวา พฺรหฺมโลเก ฐิโต อิมา คาถาโย อภาสิ\csromandash{}}

\csromangathameasure{“อารมฺภถ นิกฺกมถ, ยุญฺชถ พุทฺธสาสเน. \\ ธุนาถ มจฺจุโน เสนํ, นฬาคารํว กุญฺชโร. \\ โย อิมสฺมิํ ธมฺมวินเย, อปฺปมตฺโต วิหสฺสติ. \\ ปหาย ชาติสํสารํ, ทุกฺขสฺสนฺตํ กริสฺสตี”ติ.}
\csromangathagroup{\csromangathastanza{“อารมฺภถ\footnote{อารพฺภถ (สพฺพตฺถ)} นิกฺกมถ\footnote{นิกฺขมถ (สี, อิ)}, ยุญฺชถ พุทฺธสาสเน. \\ ธุนาถ มจฺจุโน เสนํ, นฬาคารํว กุญฺชโร.}\csromangathastanza{โย อิมสฺมิํ ธมฺมวินเย, อปฺปมตฺโต วิหสฺสติ. \\ ปหาย ชาติสํสารํ, ทุกฺขสฺสนฺตํ กริสฺสตี”ติ.}}

\csromanheader{2}{6. พฺรหฺมสํยุตฺต}
\prose{\csromanfolio{159}อถ โข ภิกฺขเว สิขี จ ภควา อรหํ สมฺมาสมฺพุทฺโธ อภิภู จ ภิกฺขุ พฺรหฺมานญฺจ พฺรหฺมปริสญฺจ พฺรหฺมปาริสชฺเช จ สํเวเชตฺวา เสยฺยถาปิ นาม ฯเปฯ ตสฺมิํ พฺรหฺมโลเก อนฺตรหิตา อรุณวติยา ราชธานิยา ปาตุรเหสุํ.\csromanspacer{} อถ โข ภิกฺขเว สิขี ภควา อรหํ สมฺมาสมฺพุทฺโธ ภิกฺขู อามนฺเตสิ “อสฺสุตฺถ โน ตุเมฺห ภิกฺขเว อภิภุสฺส ภิกฺขุโน พฺรหฺมโลเก ฐิตสฺส คาถาโย ภาสมานสฺสา”ติ.\csromanspacer{} อสฺสุมฺห โข มยํ ภนฺเต อภิภุสฺส ภิกฺขุโน พฺรหฺมโลเก ฐิตสฺส คาถาโย ภาสมานสฺสาติ.\csromanspacer{} ยถา กถํ ปน ตุเมฺห ภิกฺขเว อสฺสุตฺถ อภิภุสฺส ภิกฺขุโน พฺรหฺมโลเก ฐิตสฺส คาถาโย ภาสมานสฺสาติ.\csromanspacer{} เอวํ โข มยํ ภนฺเต อสฺสุมฺห อภิภุสฺส ภิกฺขุโน พฺรหฺมโลเก ฐิตสฺส คาถาโย ภาสมานสฺส\csromandash{}}

\csromangathameasure{“อารมฺภถ นิกฺกมถ, ยุญฺชถ พุทฺธสาสเน. \\ ธุนาถ มจฺจุโน เสนํ, นฬาคารํว กุญฺชโร. \\ โย อิมสฺมิํ ธมฺมวินเย, อปฺปมตฺโต วิหสฺสติ. \\ ปหาย ชาติสํสารํ, ทุกฺขสฺสนฺตํ กริสฺสตี”ติ.}
\csromangathagroup{\csromangathastanza{“อารมฺภถ นิกฺกมถ, ยุญฺชถ พุทฺธสาสเน. \\ ธุนาถ มจฺจุโน เสนํ, นฬาคารํว กุญฺชโร.}\csromangathastanza{โย อิมสฺมิํ ธมฺมวินเย, อปฺปมตฺโต วิหสฺสติ. \\ ปหาย ชาติสํสารํ, ทุกฺขสฺสนฺตํ กริสฺสตี”ติ.}}

\prose{เอวํ โข มยํ ภนฺเต อสฺสุมฺห อภิภุสฺส ภิกฺขุโน พฺรหฺมโลเก ฐิตสฺส คาถาโย ภาสมานสฺสาติ.\csromanspacer{} สาธุ สาธุ ภิกฺขเว, สาธุ โข ตุเมฺห ภิกฺขเว อสฺสุตฺถ อภิภุสฺส ภิกฺขุโน พฺรหฺมโลเก ฐิตสฺส คาถาโย ภาสมานสฺสาติ.}

\prose{อิทมโวจ ภควา.\csromanspacer{} อตฺตมนา เต ภิกฺขู ภควโต ภาสิตํ อภินนฺทุนฺติ.}

\csromanheader{2}{5. ปรินิพฺพานสุตฺต}
\proseitem{186}{เอกํ สมยํ ภควา กุสินารายํ วิหรติ อุปวตฺตเน มลฺลานํ สาลวเน อนฺตเรน ยมกสาลานํ ปรินิพฺพานสมเย.\csromanspacer{} อถ โข ภควา ภิกฺขู อามนฺเตสิ “หนฺท ทานิ ภิกฺขเว อามนฺตยามิ โว, วยธมฺมา สงฺขารา, อปฺปมาเทน สมฺปาเทถา”ติ.\csromanspacer{} อยํ ตถาคตสฺส ปจฺฉิมา วาจา.}

\prose{\csromanfolio{160}อถ โข ภควา ปฐมํ ฌานํ\footnote{ปฐมชฺฌานํ (สฺยา, กํ) เอวํ ทุติยํ ฌานํ อิจฺจาทีสุปิ.} สมาปชฺชิ, ปฐมา ฌานา\footnote{ปฐมชฺฌานา (สฺยา, กํ) เอวํ ทุติยา ฌานา อิจฺจาทีสุปิ.} วุฏฺฐหิตฺวา ทุติยํ ฌานํ สมาปชฺชิ, ทุติยา ฌานา วุฏฺฐหิตฺวา ตติยํ ฌานํ สมาปชฺชิ, ตติยา ฌานา วุฏฺฐหิตฺวา จตุตฺถํ ฌานํ สมาปชฺชิ, จตุตฺถา ฌานา วุฏฺฐหิตฺวา อากาสานญฺจายตนํ สมาปชฺชิ, อากาสานญฺจายตนา วุฏฺฐหิตฺวา วิญฺญาณญฺจายตนํ สมาปชฺชิ, วิญฺญาณญฺจายตนา วุฏฺฐหิตฺวา อากิญฺจญฺญายตนํ สมาปชฺชิ, อากิญฺจญฺญายตนา วุฏฺฐหิตฺวา เนวสญฺญานาสญฺญายตนํ สมาปชฺชิ, เนวสญฺญานาสญฺญายตนา วุฏฺฐหิตฺวา สญฺญาเวทยิตนิโรธํ สมาปชฺชิ.}

\prose{สญฺญาเวทยิตนิโรธา วุฏฺฐหิตฺวา เนวสญฺญานาสญฺญายตนํ สมาปชฺชิ, เนวสญฺญานาสญฺญายตนา วุฏฺฐหิตฺวา อากิญฺจญฺญายตนํ สมาปชฺชิ, อากิญฺจญฺญายตนา วุฏฺฐหิตฺวา วิญฺญาณญฺจายตนํ สมาปชฺชิ, วิญฺญาณญฺจายตนา วุฏฺฐหิตฺวา อากาสานญฺจายตนํ สมาปชฺชิ, อากาสานญฺจายตนา วุฏฺฐหิตฺวา จตุตฺถํ ฌานํ สมาปชฺชิ, จตุตฺถา ฌานา วุฏฺฐหิตฺวา ตติยํ ฌานํ สมาปชฺชิ, ตติยา ฌานา วุฏฺฐหิตฺวา ทุติยํ ฌานํ สมาปชฺชิ, ทุติยา ฌานา วุฏฺฐหิตฺวา ปฐมํ ฌานํ สมาปชฺชิ, ปฐมา ฌานา วุฏฺฐหิตฺวา ทุติยํ ฌานํ สมาปชฺชิ, ทุติยา ฌานา วุฏฺฐหิตฺวา ตติยํ ฌานํ สมาปชฺชิ, ตติยา ฌานา วุฏฺฐหิตฺวา จตุตฺถํ ฌานํ สมาปชฺชิ, จตุตฺถา ฌานา วุฏฺฐหิตฺวา สมนนฺตรํ ภควา ปรินิพฺพายิ, ปรินิพฺพุเต ภควติ สห ปรินิพฺพานา พฺรหฺมา สหมฺปติ}

\csromangathameasure{“สพฺเพว นิกฺขิปิสฺสนฺติ, ภูตา โลเก สมุสฺสยํ. \\ ยตฺถ เอตาทิโส สตฺถา, โลเก อปฺปฏิปุคฺคโล. \\ ตถาคโต พลปฺปตฺโต, สมฺพุทฺโธ ปรินิพฺพุโต”ติ.}
\csromangathagroup{\csromangathastanza{“สพฺเพว นิกฺขิปิสฺสนฺติ, ภูตา โลเก สมุสฺสยํ. \\ ยตฺถ เอตาทิโส สตฺถา, โลเก อปฺปฏิปุคฺคโล. \\ ตถาคโต พลปฺปตฺโต, สมฺพุทฺโธ ปรินิพฺพุโต”ติ.}}

\prose{ปรินิพฺพุเต ภควติ สห ปรินิพฺพานา สกฺโก เทวานมินฺโท อิมํ คาถํ อภาสิ\csromandash{}}

\csromangathameasure{“อนิจฺจา วต สงฺขารา, อุปฺปาทวยธมฺมิโน. \\ อุปฺปชฺชิตฺวา นิรุชฺฌนฺติ, เตสํ วูปสโม สุโข”ติ.}
\csromangathagroup{\csromangathastanza{“อนิจฺจา วต สงฺขารา, อุปฺปาทวยธมฺมิโน. \\ อุปฺปชฺชิตฺวา นิรุชฺฌนฺติ, เตสํ วูปสโม สุโข”ติ.}}

\csromanheader{2}{6. พฺรหฺมสํยุตฺต}
\prose{\csromanfolio{161}ปรินิพฺพุเต ภควติ สห ปรินิพฺพานา อายสฺมา อานนฺโท อิมํ คาถํ อภาสิ\csromandash{}}

\csromangathameasure{“ตทาสิ ยํ ภิํสนกํ, ตทาสิ โลมหํสนํ. \\ สพฺพาการวรูเปเต, สมฺพุทฺเธ ปรินิพฺพุเต”ติ.}
\csromangathagroup{\csromangathastanza{“ตทาสิ ยํ ภิํสนกํ, ตทาสิ โลมหํสนํ. \\ สพฺพาการวรูเปเต, สมฺพุทฺเธ ปรินิพฺพุเต”ติ.}}

\prose{ปรินิพฺพุเต ภควติ สห ปรินิพฺพานา อายสฺมา อนุรุทฺโธ อิมา คาถาโย อภาสิ\csromandash{}}

\csromangathameasure{“นาหุ อสฺสาสปสฺสาโส, ฐิตจิตฺตสฺส ตาทิโน. \\ อเนโช สนฺติมารพฺภ, จกฺขุมา ปรินิพฺพุโต. \\ อสลฺลีเนน จิตฺเตน, เวทนํ อชฺฌวาสยิ. \\ ปชฺโชตสฺเสว นิพฺพานํ, วิโมกฺโข เจตโส อหู”ติ.}
\csromangathagroup{\csromangathastanza{“นาหุ อสฺสาสปสฺสาโส, ฐิตจิตฺตสฺส ตาทิโน. \\ อเนโช สนฺติมารพฺภ, จกฺขุมา ปรินิพฺพุโต\footnote{ยํ กาลมกรี มุนิ (มหาปรินิพฺพานสุตฺเต)}.}\csromangathastanza{อสลฺลีเนน จิตฺเตน, เวทนํ อชฺฌวาสยิ. \\ ปชฺโชตสฺเสว นิพฺพานํ, วิโมกฺโข เจตโส อหู”ติ.}}

\vspace{\tipitakaparskip}
\csromancenter{ทุติโย วคฺโค\textbf{.}}
\titlehead{\textbf{ตสฺสุทฺทานํ}}
\csromangathameasure{พฺรหฺมาสนํ เทวทตฺโต, อนฺธกวินฺโท อรุณวตี. \\ ปรินิพฺพาเนน จ เทสิตํ, อิทํ พฺรหฺมปญฺจกนฺติ. \\ \textbf{พฺรหฺมสํยุตฺตํ สมตฺตํ}\textbf{.}}
\csromangathagroup{\csromangathastanza{พฺรหฺมาสนํ เทวทตฺโต, อนฺธกวินฺโท อรุณวตี. \\ ปรินิพฺพาเนน จ เทสิตํ, อิทํ พฺรหฺมปญฺจกนฺติ. \\ \textbf{พฺรหฺมสํยุตฺตํ สมตฺตํ}\footnote{อิโต ปรํ มรมฺม-โปตฺถเกสุ เอวมฺปิ ทิสฺสติ\csromandash{} พฺรหฺมายาจนํ อคารวญฺจ, พฺรหฺมเทโว พโก จ พฺรหฺมา. อญฺญตโร จ พฺรหฺมา โกกาลิกญฺจ, ติสฺสกญฺจ ตุรู จ. พฺรหฺมา โกกาลิกภิกฺขุ, สนงฺกุมาเรน เทวทตฺตํ. อนฺธกวินฺทํ อรุณวติ, ปรินิพฺพาเนน ปนฺนรสาติ.}\textbf{.}}}

\csromanheader{2}{7. พฺราหฺมณสํยุตฺต}
\chapterhead{1. อรหนฺตวคฺค}
\csromanheader{2}{1. ธนญฺชานีสุตฺต}
\proseitem{187}{\csromanfolio{162}เอวํ เม สุตํ\csromandash{}เอกํ สมยํ ภควา ราชคเห วิหรติ เวฬุวเน กลนฺทกนิวาเป.\csromanspacer{} เตน โข ปน สมเยน อญฺญตรสฺส ภารทฺวาชโคตฺตสฺส พฺราหฺมณสฺส ธนญฺชานี\footnote{ธานญฺชานี (อิ, สี-ฏฺฐ)} นาม พฺราหฺมณี อภิปฺปสนฺนา โหติ พุทฺเธ จ ธมฺเม จ สํเฆ จ.\csromanspacer{} อถ โข ธนญฺชานี พฺราหฺมณี ภารทฺวาชโคตฺตสฺส พฺราหฺมณสฺส ภตฺตํ อุปสํหรนฺตี อุปกฺขลิตฺวา ติกฺขตฺตุํ อุทานํ อุทาเนสิ\csromandash{}}

\csromanheader{2}{“นโม ตสฺส ภควโต อรหโต สมฺมาสมฺพุทฺธสฺส.}
\namakkaram{นโม ตสฺส ภควโต อรหโต สมฺมาสมฺพุทฺธสฺส.}
\namakkaram{นโม ตสฺส ภควโต อรหโต สมฺมาสมฺพุทฺธสฺสา”ติ.}
\prose{เอวํ วุตฺเต ภารทฺวาชโคตฺโต พฺราหฺมโณ ธนญฺชานิํ พฺราหฺมณิํ เอตทโวจ “เอวเมวํ ปนายํ วสลี ยสฺมิํ วา ตสฺมิํ วา ตสฺส มุณฺฑกสฺส สมณสฺส วณฺณํ ภาสติ, อิทานิ ตฺยาหํ วสลิ ตสฺส สตฺถุโน วาทํ อาโรเปสฺสามี”ติ.\csromanspacer{} น ขฺวาหํ ตํ พฺราหฺมณ ปสฺสามิ สเทวเก โลเก สมารเก สพฺรหฺมเก สสฺสมณพฺราหฺมณิยา ปชาย สเทวมนุสฺสาย, โย ตสฺส ภควโต วาทํ อาโรเปยฺย อรหโต สมฺมาสมฺพุทฺธสฺส.\csromanspacer{} อปิ จ ตฺวํ พฺราหฺมณ คจฺฉ, คนฺตฺวา วิชานิสฺสสีติ\footnote{คนฺตฺวาปิ ชานิสฺสสีติ (สฺยา, กํ)}.}

\prose{อถ โข ภารทฺวาชโคตฺโต พฺราหฺมโณ กุปิโต อนตฺตมโน เยน ภควา เตนุปสงฺกมิ, อุปสงฺกมิตฺวา ภควตา สทฺธิํ สมฺโมทิ, สมฺโมทนียํ กถํ สารณียํ วีติสาเรตฺวา เอกมนฺตํ นิสีทิ, เอกมนฺตํ นิสินฺโน โข ภารทฺวาชโคตฺโต พฺราหฺมโณ ภควนฺตํ คาถาย อชฺฌภาสิ\csromandash{}}

\csromangathameasure{“กิํสุ เฉตฺวา สุขํ เสติ, กิํสุ เฉตฺวา น โสจติ. \\ กิสฺสสฺสุ เอกธมฺมสฺส, วธํ โรเจสิ โคตมา”ติ.}
\csromangathagroup{\csromangathastanza{“กิํสุ เฉตฺวา สุขํ เสติ, กิํสุ เฉตฺวา น โสจติ. \\ กิสฺสสฺสุ เอกธมฺมสฺส, วธํ โรเจสิ โคตมา”ติ.}}

\csromanheader{2}{7. พฺราหฺมณสํยุตฺต}
\csromangathameasure{โกธํ เฉตฺวา สุขํ เสติ, โกธํ เฉตฺวา น โสจติ. \\ โกธสฺส วิสมูลสฺส, มธุรคฺคสฺส พฺราหฺมณ. \\ วธํ อริยา ปสํสนฺติ, ตํ หิ เฉตฺวา น โสจตีติ.}
\csromangathagroup{\csromangathastanza{\csromanfolio{163}โกธํ เฉตฺวา สุขํ เสติ, โกธํ เฉตฺวา น โสจติ. \\ โกธสฺส วิสมูลสฺส, มธุรคฺคสฺส พฺราหฺมณ. \\ วธํ อริยา ปสํสนฺติ, ตํ หิ เฉตฺวา น โสจตีติ.}}

\prose{เอวํ วุตฺเต ภารทฺวาชโคตฺโต พฺราหฺมโณ ภควนฺตํ เอตทโวจ “อภิกฺกนฺตํ โภ โคตม, อภิกฺกนฺตํ โภ โคตม, เสยฺยถาปิ โภ โคตม นิกฺกุชฺชิตํ วา อุกฺกุชฺเชยฺย, ปฏิจฺฉนฺนํ วา วิวเรยฺย, มูฬฺหสฺส วา มคฺคํ อาจิกฺเขยฺย, อนฺธกาเร วา เตลปชฺโชตํ ธาเรยฺย ‘จกฺขุมนฺโต รูปานิ ทกฺขนฺตี’ติ.\csromanspacer{} เอวเมวํ โภตา โคตเมน อเนกปริยาเยน ธมฺโม ปกาสิโต, เอสาหํ ภนฺเต ภควนฺตํ โคตมํ สรณํ คจฺฉามิ ธมฺมญฺจ ภิกฺขุสํฆญฺจ, ลเภยฺยาหํ โภโต โคตมสฺส สนฺติเก ปพฺพชฺชํ, ลเภยฺยํ อุปสมฺปทนฺ”ติ.}

\prose{อลตฺถ โข ภารทฺวาชโคตฺโต พฺราหฺมโณ ภควโต สนฺติเก ปพฺพชฺชํ, อลตฺถ อุปสมฺปทํ.\csromanspacer{} อจิรูปสมฺปนฺโน โข ปนายสฺมา ภารทฺวาโช เอโก วูปกฏฺโฐ อปฺปมตฺโต อาตาปี ปหิตตฺโต วิหรนฺโต นจิรสฺเสว ยสฺสตฺถาย กุลปุตฺตา สมฺมเทว อคารสฺมา อนคาริยํ ปพฺพชนฺติ, ตทนุตฺตรํ พฺรหฺมจริยปริโยสานํ ทิฏฺเฐว ธมฺเม สยํ อภิญฺญา สจฺฉิกตฺวา อุปสมฺปชฺช วิหาสิ, “ขีณา ชาติ, วุสิตํ พฺรหฺมจริยํ, กตํ กรณียํ, นาปรํ อิตฺถตฺตายา”ติ อพฺภญฺญาสิ.\csromanspacer{} อญฺญตโร จ ปนายสฺมา ภารทฺวาโช อรหตํ อโหสีติ.}

\csromanheader{2}{2. อกฺโกสสุตฺต}
\proseitem{188}{เอกํ สมยํ ภควา ราชคเห วิหรติ เวฬุวเน กลนฺทกนิวาเป.\csromanspacer{} อสฺโสสิ โข อกฺโกสกภารทฺวาโช พฺราหฺมโณ “ภารทฺวาชโคตฺโต กิร พฺราหฺมโณ สมณสฺส โคตมสฺส สนฺติเก อคารสฺมา อนคาริยํ ปพฺพชิโต”ติ, กุปิโต อนตฺตมโน เยน ภควา เตนุปสงฺกมิ, อุปสงฺกมิตฺวา ภควนฺตํ อสพฺภาหิ ผรุสาหิ วาจาหิ อกฺโกสติ ปริภาสติ.}

\prose{เอวํ วุตฺเต ภควา อกฺโกสกภารทฺวาชํ พฺราหฺมณํ เอตทโวจ “ตํ กิํ มญฺญสิ พฺราหฺมณ, อปิ นุ โข เต อาคจฺฉนฺติ มิตฺตามจฺจา ญาติสาโลหิตา อติถิโย\footnote{อติถโย (?)}”ติ.\csromanspacer{} อปฺเปกทา เม โภ โคตม อาคจฺฉนฺติ \csromanfolio{164}มิตฺตามจฺจา ญาติสาโลหิตา อติ ถิโยติ.\csromanspacer{} ตํ กิํ มญฺญสิ พฺราหฺมณ, อปิ นุ เตสํ อนุปฺปเทสิ ขาทนียํ วา โภชนียํ วา สายนียํ วาติ.\csromanspacer{} อปฺเปกทา เนสาหํ โภ โคตม อนุปฺปเทมิ ขาทนียํ วา โภชนียํ วา สายนียํ วาติ.\csromanspacer{} สเจ โข ปน เต พฺราหฺมณ นปฺปฏิคฺคณฺหนฺติ, กสฺส ตํ โหตีติ.\csromanspacer{} สเจ เต โภ โคตม นปฺปฏิคฺคณฺหนฺติ, อมฺหากเมว ตํ โหตีติ.\csromanspacer{} เอวเมว โข พฺราหฺมณ ยํ ตฺวํ อเมฺห อนกฺโกสนฺเต อกฺโกสสิ, อโรเสนฺเต โรเสสิ, อภณฺฑนฺเต ภณฺฑสิ, ตํ เต มยํ นปฺปฏิคฺคณฺหาม, ตเวเวตํ พฺราหฺมณ โหติ, ตเวเวตํ พฺราหฺมณ โหติ.}

\prose{โย โข พฺราหฺมณ อกฺโกสนฺตํ ปจฺจกฺโกสติ, โรเสนฺตํ ปฏิโรเสติ, ภณฺฑนฺตํ ปฏิภณฺฑติ, อยํ วุจฺจติ พฺราหฺมณ สมฺภุญฺชติ วีติหรตีติ.\csromanspacer{} เต มยํ ตยา เนว สมฺภุญฺชาม น วีติหราม.\csromanspacer{} ตเวเวตํ พฺราหฺมณ โหติ, ตเวเวตํ พฺราหฺมณ โหตีติ.\csromanspacer{} ภวนฺตํ โข โคตมํ สราชิกา ปริสา เอวํ ชานาติ “อรหํ สมโณ โคตโม”ติ, อถ จ ปน ภวํ โคตโม กุชฺฌตีติ.}

\csromangathameasure{อกฺโกธสฺส กุโต โกโธ, ทนฺตสฺส สมชีวิโน. \\ สมฺมทญฺญา วิมุตฺตสฺส, อุปสนฺตสฺส ตาทิโน. \\ ตสฺเสว เตน ปาปิโย, โย กุทฺธํ ปฏิกุชฺฌติ. \\ กุทฺธํ อปฺปฏิกุชฺฌนฺโต, สงฺคามํ เชติ ทุชฺชยํ. \\ อุภินฺนมตฺถํ จรติ, อตฺตโน จ ปรสฺส จ. \\ ปรํ สงฺกุปิตํ ญตฺวา, โย สโต อุปสมฺมติ. \\ อุภินฺนํ ติกิจฺฉนฺตานํ, อตฺตโน จ ปรสฺส จ. \\ ชนา มญฺญนฺติ พาโลติ, เย ธมฺมสฺส อโกวิทาติ.}
\csromangathagroup{\csromangathastanza{อกฺโกธสฺส กุโต โกโธ, ทนฺตสฺส สมชีวิโน. \\ สมฺมทญฺญา วิมุตฺตสฺส, อุปสนฺตสฺส ตาทิโน.}\csromangathastanza{ตสฺเสว เตน ปาปิโย, โย กุทฺธํ ปฏิกุชฺฌติ. \\ กุทฺธํ อปฺปฏิกุชฺฌนฺโต, สงฺคามํ เชติ ทุชฺชยํ.}\csromangathastanza{อุภินฺนมตฺถํ จรติ, อตฺตโน จ ปรสฺส จ. \\ ปรํ สงฺกุปิตํ ญตฺวา, โย สโต อุปสมฺมติ.}\csromangathastanza{อุภินฺนํ ติกิจฺฉนฺตานํ, อตฺตโน จ ปรสฺส จ. \\ ชนา มญฺญนฺติ พาโลติ, เย ธมฺมสฺส อโกวิทาติ.}}

\prose{เอวํ วุตฺเต อกฺโกสกภารทฺวาโช พฺราหฺมโณ ภควนฺตํ เอตทโวจ “อภิกฺกนฺตํ โภ โคตม ฯเปฯ เอสาหํ ภวนฺตํ โคตมํ สรณํ คจฺฉามิ ธมฺมญฺจ ภิกฺขุสํฆญฺจ, ลเภยฺยาหํ ภนฺเต โภโต โคตมสฺส สนฺติเก ปพฺพชฺชํ, ลเภยฺยํ อุปสมฺปทนฺ”ติ.}

\prose{อลตฺถ โข อกฺโกสกภารทฺวาโช พฺราหฺมโณ ภควโต สนฺติเก ปพฺพชฺชํ, อลตฺถ อุปสมฺปทํ.\csromanspacer{} อจิรูปสมฺปนฺโน โข ปนายสฺมา อกฺโกสกภารทฺวาโช เอโก วูปกฏฺโฐ อปฺปมตฺโต อาตาปี ปหิตตฺโต}

\csromanheader{2}{7. พฺราหฺมณสํยุตฺต}
\prose{\csromanfolio{165}วิหรนฺโต นจิรสฺเสว ยสฺสตฺถาย กุลปุตฺตา สมฺมเทว อคารสฺมา อนคาริยํ ปพฺพชนฺติ, ตทนุตฺตรํ พฺรหฺมจริยปริโยสานํ ทิฏฺเฐว ธมฺเม สยํ อภิญฺญา สจฺฉิกตฺวา อุปสมฺปชฺช วิหาสิ, “ขีณา ชาติ, วุสิตํ พฺรหฺมจริยํ, กตํ กรณียํ, นาปรํ อิตฺถตฺตายา”ติ อพฺภญฺญาสิ.\csromanspacer{} อญฺญตโร จ ปนายสฺมา ภารทฺวาโช อรหตํ อโหสีติ.}

\csromanheader{2}{3. อสุรินฺทกสุตฺต}
\proseitem{189}{เอกํ สมยํ ภควา ราชคเห วิหรติ เวฬุวเน กลนฺทกนิวาเป.\csromanspacer{} อสฺโสสิ โข อสุรินฺทกภารทฺวาโช พฺราหฺมโณ “ภารทฺวาชโคตฺโต พฺราหฺมโณ กิร สมณสฺส โคตมสฺส สนฺติเก อคารสฺมา อนคาริยํ ปพฺพชิโต”ติ กุปิโต อนตฺตมโน เยน ภควา เตนุปสงฺกมิ, อุปสงฺกมิตฺวา ภควนฺตํ อสพฺภาหิ ผรุสาหิ วาจาหิ อกฺโกสติ ปริภาสติ.\csromanspacer{} เอวํ วุตฺเต ภควา ตุณฺหี อโหสิ.\csromanspacer{} อถ โข อสุรินฺทกภารทฺวาโช พฺราหฺมโณ ภควนฺตํ เอตทโวจ “ชิโตสิ สมณ, ชิโหสิ สมณา”ติ.}

\csromangathameasure{ชยํ เว มญฺญติ พาโล, วาจาย ผรุสํ ภณํ. \\ ชยญฺเจวสฺส ตํ โหติ, ยา ติติกฺขา วิชานโต. \\ ตสฺเสว เตน ปาปิโย, โย กุทฺธํ ปฏิกุชฺฌติ. \\ กุทฺธํ อปฺปฏิกุชฺฌนฺโต, สงฺคามํ เชติ ทุชฺชยํ. \\ อุภินฺนมตฺถํ จรติ, อตฺตโน จ ปรสฺส จ. \\ ปรํ สงฺกุปิตํ ญตฺวา, โย สโต อุปสมฺมติ. \\ อุภินฺนํ ติกิจฺฉนฺตานํ, อตฺตโน จ ปรสฺส จ. \\ ชนา มญฺญนฺติ พาโลติ, เย ธมฺมสฺส อโกวิทาติ.}
\csromangathagroup{\csromangathastanza{ชยํ เว มญฺญติ พาโล, วาจาย ผรุสํ ภณํ. \\ ชยญฺเจวสฺส ตํ โหติ, ยา ติติกฺขา วิชานโต.}\csromangathastanza{ตสฺเสว เตน ปาปิโย, โย กุทฺธํ ปฏิกุชฺฌติ. \\ กุทฺธํ อปฺปฏิกุชฺฌนฺโต, สงฺคามํ เชติ ทุชฺชยํ.}\csromangathastanza{อุภินฺนมตฺถํ จรติ, อตฺตโน จ ปรสฺส จ. \\ ปรํ สงฺกุปิตํ ญตฺวา, โย สโต อุปสมฺมติ.}\csromangathastanza{อุภินฺนํ ติกิจฺฉนฺตานํ, อตฺตโน จ ปรสฺส จ. \\ ชนา มญฺญนฺติ พาโลติ, เย ธมฺมสฺส อโกวิทาติ.}}

\prose{เอวํ วุตฺเต อสุรินฺทกภารทฺวาโช พฺราหฺมโณ ภควนฺตํ เอตทโวจ อภิกฺกนฺตํ โภ โคตม ฯเปฯ อพฺภญฺญาสิ.\csromanspacer{} อญฺญตโร จ ปนายสฺมา ภารทฺวาโช อรหตํ อโหสีติ.}

\csromanheader{2}{4. พิลงฺคิกสุตฺต}
\proseitem{190}{\csromanfolio{166}เอกํ สมยํ ภควา ราชคเห วิหรติ เวฬุวเน กลนฺทกนิวาเป.\csromanspacer{} อสฺโสสิ โข พิลงฺคิกภารทฺวาโช พฺราหฺมโณ “ภารทฺวาชโคตฺโต กิร พฺราหฺมโณ สมณสฺส โคตมสฺส สนฺติเก อคารสฺมา อนคาริยํ ปพฺพชิโต”ติ กุปิโต อนตฺตมโน เยน ภควา เตนุปสงฺกมิ, อุปสงฺกมิตฺวา ตุณฺหีภูโต เอกมนฺตํ อฏฺฐาสิ.\csromanspacer{} อถ โข ภควา พิลงฺคิกสฺส ภารทฺวาชสฺส พฺราหฺมณสฺส เจตสา เจโตปริวิตกฺกมญฺญาย พิลงฺคิกํ ภารทฺวาชํ พฺราหฺมณํ คาถาย อชฺฌภาสิ\csromandash{}}

\csromangathameasure{“โย อปฺปทุฏฺฐสฺส นรสฺส ทุสฺสติ, \\ สุทฺธสฺส โปสสฺส อนงฺคณสฺส. \\ ตเมว พาลํ ปจฺเจติ ปาปํ, \\ สุขุโม รโช ปฏิวาตํว ขิตฺโต”ติ.}
\csromangathagroup{\csromangathastanza{“โย อปฺปทุฏฺฐสฺส นรสฺส ทุสฺสติ, \\ สุทฺธสฺส โปสสฺส อนงฺคณสฺส. \\ ตเมว พาลํ ปจฺเจติ ปาปํ, \\ สุขุโม รโช ปฏิวาตํว ขิตฺโต”ติ.}}

\prose{เอวํ วุตฺเต พิลงฺคิกภารทฺวาโช พฺราหฺมโณ ภควนฺตํ เอตทโวจ อภิกฺกนฺตํ โภ โคตม ฯเปฯ อพฺภญฺญาสิ.\csromanspacer{} อญฺญตโร จ ปนายสฺมา ภารทฺวาโช อรหตํ อโหสีติ.}

\csromanheader{2}{5. อหิํสกสุตฺต}
\proseitem{191}{สาวตฺถินิทานํ.\csromanspacer{} อถ โข อหิํสกภารทฺวาโช พฺราหฺมโณ เยน ภควา เตนุปสงฺกมิ, อุปสงฺกมิตฺวา ภควตา สทฺธิํ สมฺโมทิ, สมฺโมทนียํ กถํ สารณียํ วีติสาเรตฺวา เอกมนฺตํ นิสีทิ, เอกมนฺตํ นิสินฺโน โข อหิํสกภารทฺวาโช พฺราหฺมโณ ภควนฺตํ เอตทโวจ “อหิํสกาหํ โภ โคตม, อหิํสกาหํ โภ โคตมา”ติ.}

\csromangathameasure{ยถา นามํ ตถา จสฺส, สิยา โข ตฺวํ อหิํสโก. \\ โย จ กาเยน วาจาย, มนสา จ น หิํสติ. \\ ส เว อหิํสโก โหติ, โย ปรํ น วิหิํสตีติ.}
\csromangathagroup{\csromangathastanza{ยถา นามํ ตถา จสฺส, สิยา โข ตฺวํ อหิํสโก. \\ โย จ กาเยน วาจาย, มนสา จ น หิํสติ. \\ ส เว อหิํสโก โหติ, โย ปรํ น วิหิํสตีติ.}}

\prose{เอวํ วุตฺเต อหิํสกภารทฺวาโช พฺราหฺมโณ ภควนฺตํ เอตทโวจ อภิกฺกนฺตํ โภ โคตม ฯเปฯ อพฺภญฺญาสิ.\csromanspacer{} อญฺญตโร จ ปนายสฺมา อหิํสกภารทฺวาโช อรหตํ อโหสีติ.}

\csromanheader{2}{7. พฺราหฺมณสํยุตฺต}
\csromanheader{2}{6. ชฏาสุตฺต}
\proseitem{192}{\csromanfolio{167}สาวตฺถินิทานํ.\csromanspacer{} อถ โข ชฏาภารทฺวาโช พฺราหฺมโณ เยน ภควา เตนุปสงฺกมิ, อุปสงฺกมิตฺวา ภควตา สทฺธิํ สมฺโมทิ, สมฺโมทนียํ กถํ สารณียํ วีติสาเรตฺวา เอกมนฺตํ นิสีทิ, เอกมนฺตํ นิสินฺโน โข ชฏาภารทฺวาโช พฺราหฺมโณ ภควนฺตํ คาถาย อชฺฌภาสิ\csromandash{}}

\csromangathameasure{“อนฺโตชฏา พหิชฏา, ชฏาย ชฏิตา ปชา. \\ ตํ ตํ โคตม ปุจฺฉามิ, โก อิมํ วิชฏเย ชฏนฺ”ติ. \\ สีเล ปติฏฺฐาย นโร สปญฺโญ, จิตฺตํ ปญฺญญฺจ ภาวยํ. \\ อาตาปี นิปโก ภิกฺขุ, โส อิมํ วิชฏเย ชฏํ. \\ เยสํ ราโค จ โทโส จ, อวิชฺชา จ วิราชิตา. \\ ขีณาสวา อรหนฺโต, เตสํ วิชฏิตา ชฏา. \\ ยตฺถ นามญฺจ รูปญฺจ, อเสสํ อุปรุชฺฌติ. \\ ปฏิฆํ รูปสญฺญา จ, เอตฺเถสา ฉิชฺชเต ชฏาติ.}
\csromangathagroup{\csromangathastanza{“อนฺโตชฏา พหิชฏา, ชฏาย ชฏิตา ปชา. \\ ตํ ตํ โคตม ปุจฺฉามิ, โก อิมํ วิชฏเย ชฏนฺ”ติ.}\csromangathastanza{สีเล ปติฏฺฐาย นโร สปญฺโญ, จิตฺตํ ปญฺญญฺจ ภาวยํ. \\ อาตาปี นิปโก ภิกฺขุ, โส อิมํ วิชฏเย ชฏํ.}\csromangathastanza{เยสํ ราโค จ โทโส จ, อวิชฺชา จ วิราชิตา. \\ ขีณาสวา อรหนฺโต, เตสํ วิชฏิตา ชฏา.}\csromangathastanza{ยตฺถ นามญฺจ รูปญฺจ, อเสสํ อุปรุชฺฌติ. \\ ปฏิฆํ รูปสญฺญา จ, เอตฺเถสา ฉิชฺชเต ชฏาติ.}}

\prose{เอวํ วุตฺเต ชฏาภารทฺวาโช ภควนฺตํ เอตทโวจ อภิกฺกนฺตํ โภ โคตม ฯเปฯ.\csromanspacer{} อญฺญตโร จ ปนายสฺมา ภารทฺวาโช อรหตํ อโหสีติ.}

\csromanheader{2}{7. สุทฺธิกสุตฺต}
\proseitem{193}{สาวตฺถินิทานํ.\csromanspacer{} อถ โข สุทฺธิกภารทฺวาโช พฺราหฺมโณ เยน ภควา เตนุปสงฺกมิ, อุปสงฺกมิตฺวา ภควตา สทฺธิํ สมฺโมทิ, สมฺโมทนียํ กถํ สารณียํ วีติสาเรตฺวา เอกมนฺตํ นิสีทิ, เอกมนฺตํ นิสินฺโน โข สุทฺธิกภารทฺวาโช พฺราหฺมโณ ภควโต สนฺติเก อิมํ คาถํ อชฺฌภาสิ\csromandash{}}

\csromangathameasure{“น พฺราหฺมโณ สุชฺฌติ โกจิ, โลเก สีลวาปิ ตโปกรํ. \\ วิชฺชาจรณสมฺปนฺโน, โส สุชฺฌติ น อญฺญา อิตรา ปชา”ติ. \\ พหุมฺปิ ปลปํ ชปฺปํ, น ชจฺจา โหติ พฺราหฺมโณ. \\ อนฺโตกสมฺพุ สํกิลิฏฺโฐ, กุหนํ อุปนิสฺสิโต. \\ ขตฺติโย พฺราหฺมโณ เวสฺโส, สุทฺโท จณฺฑาลปุกฺกุโส. \\ อารทฺธวีริโย ปหิตตฺโต, นิจฺจํ ทฬฺหปรกฺกโม. \\ ปปฺโปติ ปรมํ สุทฺธิํ, เอวํ ชานาหิ พฺราหฺมณาติ.}
\csromangathagroup{\csromangathastanza{“น พฺราหฺมโณ\footnote{นาพฺราหฺมโณ (?)} สุชฺฌติ โกจิ, โลเก สีลวาปิ ตโปกรํ. \\ วิชฺชาจรณสมฺปนฺโน, โส สุชฺฌติ น อญฺญา อิตรา ปชา”ติ.}\csromangathastanza{พหุมฺปิ ปลปํ ชปฺปํ, น ชจฺจา โหติ พฺราหฺมโณ. \\ อนฺโตกสมฺพุ สํกิลิฏฺโฐ, กุหนํ อุปนิสฺสิโต.}\csromangathastanza{\csromanfolio{168}ขตฺติโย พฺราหฺมโณ เวสฺโส, สุทฺโท จณฺฑาลปุกฺกุโส. \\ อารทฺธวีริโย ปหิตตฺโต, นิจฺจํ ทฬฺหปรกฺกโม. \\ ปปฺโปติ ปรมํ สุทฺธิํ, เอวํ ชานาหิ พฺราหฺมณาติ.}}

\prose{เอวํ วุตฺเต สุทฺธิกภารทฺวาโช พฺราหฺมโณ ภควนฺตํ เอตทโวจ อภิกฺกนฺตํ โภ โคตม ฯเปฯ.\csromanspacer{} อญฺญตโร จ ปนายสฺมา ภารทฺวาโช อรหตํ อโหสีติ.}

\csromanheader{2}{8. อคฺคิกสุตฺต}
\proseitem{194}{เอกํ สมยํ ภควา ราชคเห วิหรติ เวฬุวเน กลนฺทกนิวาเป.\csromanspacer{} เตน โข ปน สมเยน อคฺคิกภารทฺวาชสฺส พฺราหฺมณสฺส สปฺปินา ปายโส สนฺนิหิโต โหติ “อคฺคิํ ชุหิสฺสามิ, อคฺคิหุตฺตํ ปริจริสฺสามี”ติ.}

\prose{อถ โข ภควา ปุพฺพณฺหสมยํ นิวาเสตฺวา ปตฺตจีวรมาทาย ราชคหํ ปิณฺฑาย ปาวิสิ.\csromanspacer{} ราชคเห สปทานํ ปิณฺฑาย จรมาโน เยน อคฺคิกภารทฺวาชสฺส พฺราหฺมณสฺส นิเวสนํ เตนุปสงฺกมิ, อุปสงฺกมิตฺวา เอกมนฺตํ อฏฺฐาสิ.\csromanspacer{} อทฺทสา โข อคฺคิกภารทฺวาโช พฺราหฺมโณ ภควนฺตํ ปิณฺฑาย ฐิตํ, ทิสฺวาน ภควนฺตํ คาถาย อชฺฌภาสิ\csromandash{}}

\csromangathameasure{“ตีหิ วิชฺชาหิ สมฺปนฺโน, ชาติมา สุตวา พหู. \\ วิชฺชาจรณสมฺปนฺโน, โสมํ ภุญฺเชยฺย ปายสนฺ”ติ. \\ พหุมฺปิ ปลปํ ชปฺปํ, น ชจฺจา โหติ พฺราหฺมโณ. \\ อนฺโตกสมฺพุ สํกิลิฏฺโฐ, กุหนาปริวาริโต. \\ ปุพฺเพนิวาสํ โย เวที, สคฺคาปายญฺจ ปสฺสติ. \\ อโถ ชาติกฺขยํ ปตฺโต, อภิญฺญาโวสิโต มุนิ. \\ เอตาหิ ตีหิ วิชฺชาหิ, เตวิชฺโช โหติ พฺราหฺมโณ. \\ วิชฺชาจรณสมฺปนฺโน, โสมํ ภุญฺเชยฺย ปายสนฺติ.}
\csromangathagroup{\csromangathastanza{“ตีหิ วิชฺชาหิ สมฺปนฺโน, ชาติมา สุตวา พหู. \\ วิชฺชาจรณสมฺปนฺโน, โสมํ ภุญฺเชยฺย ปายสนฺ”ติ.}\csromangathastanza{พหุมฺปิ ปลปํ ชปฺปํ, น ชจฺจา โหติ พฺราหฺมโณ. \\ อนฺโตกสมฺพุ สํกิลิฏฺโฐ, กุหนาปริวาริโต.}\csromangathastanza{ปุพฺเพนิวาสํ โย เวที, สคฺคาปายญฺจ ปสฺสติ. \\ อโถ ชาติกฺขยํ ปตฺโต, อภิญฺญาโวสิโต มุนิ.}\csromangathastanza{เอตาหิ ตีหิ วิชฺชาหิ, เตวิชฺโช โหติ พฺราหฺมโณ. \\ วิชฺชาจรณสมฺปนฺโน, โสมํ ภุญฺเชยฺย ปายสนฺติ.}}

\prose{ภุญฺชตุ ภวํ โคตโม, พฺราหฺมโณ ภวนฺติ.}

\csromanheader{2}{7. พฺราหฺมณสํยุตฺต}
\csromangathameasure{คาถาภิคีตํ เม อโภชเนยฺยํ, \\ สมฺปสฺสตํ พฺราหฺมณ เนส ธมฺโม. \\ คาถาภิคีตํ ปนุทนฺติ พุทฺธา, \\ ธมฺเม สติ พฺราหฺมณ วุตฺติเรสา. \\ อญฺเญน จ เกวลินํ มเหสิํ, \\ ขีณาสวํ กุกฺกุจฺจวูปสนฺตํ. \\ อนฺเนน ปาเนน อุปฏฺฐหสฺสุ, \\ เขตฺตํ หิ ตํ ปุญฺญเปกฺขสฺส โหตีติ.}
\csromangathagroup{\csromangathastanza{\csromanfolio{169}คาถาภิคีตํ เม อโภชเนยฺยํ, \\ สมฺปสฺสตํ พฺราหฺมณ เนส ธมฺโม. \\ คาถาภิคีตํ ปนุทนฺติ พุทฺธา, \\ ธมฺเม สติ พฺราหฺมณ วุตฺติเรสา.}\csromangathastanza{อญฺเญน จ เกวลินํ มเหสิํ, \\ ขีณาสวํ กุกฺกุจฺจวูปสนฺตํ. \\ อนฺเนน ปาเนน อุปฏฺฐหสฺสุ, \\ เขตฺตํ หิ ตํ ปุญฺญเปกฺขสฺส โหตีติ.}}

\prose{เอวํ วุตฺเต อคฺคิกภารทฺวาโช พฺราหฺมโณ ภควนฺตํ เอตทโวจ อภิกฺกนฺตํ โภ โคตม ฯเปฯ.\csromanspacer{} อญฺญตโร จ ปนายสฺมา อคฺคิกภารทฺวาโช อรหตํ อโหสีติ.}

\csromanheader{2}{9. สุนฺทริกสุตฺต}
\proseitem{195}{เอกํ สมยํ ภควา โกสเลสุ วิหรติ สุนฺทริกาย นทิยา ตีเร.\csromanspacer{} เตน โข ปน สมเยน สุนฺทริกภารทฺวาโช พฺราหฺมโณ สุนฺทริกาย นทิยา ตีเร อคฺคิํ ชุหติ อคฺคิหุตฺตํ ปริจรติ.\csromanspacer{} อถ โข สุนฺทริกภารทฺวาโช พฺราหฺมโณ อคฺคิํ ชุหิตฺวา อคฺคิหุตฺตํ ปริจริตฺวา อุฏฺฐายาสนา สมนฺตา จตุทฺทิสา อนุวิโลเกสิ “โก นุ โข อิมํ หพฺยเสสํ ภุญฺเชยฺยา”ติ.\csromanspacer{} อทฺทสา โข สุนฺทริกภารทฺวาโช พฺราหฺมโณ ภควนฺตํ อญฺญตรสฺมิํ รุกฺขมูเล สสีสํ ปารุตํ นิสินฺนํ, ทิสฺวาน วาเมน หตฺเถน หพฺยเสสํ คเหตฺวา ทกฺขิเณน หตฺเถน กมณฺฑลุํ คเหตฺวา เยน ภควา เตนุปสงฺกมิ.\csromanspacer{} อถ โข ภควา สุนฺทริกภารทฺวาชสฺส พฺราหฺมณสฺส ปทสทฺเทน สีสํ วิวริ.\csromanspacer{} อถ โข สุนฺทริกภารทฺวาโช พฺราหฺมโณ “มุณฺโฑ อยํ ภวํ, มุณฺฑโก อยํ ภวนฺ”ติ ตโตว ปุน นิวตฺติตุกาโม อโหสิ.\csromanspacer{} อถ โข สุนฺทริกภารทฺวาชสฺส พฺราหฺมณสฺส เอตทโหสิ “มุณฺฑาปิ หิ อิเธกจฺเจ พฺราหฺมณา ภวนฺติ, ยํนูนาหํ ตํ อุปสงฺกมิตฺวา ชาติํ ปุจฺเฉยฺยนฺ”ติ.}

\prose{อถ โข สุนฺทริกภารทฺวาโช พฺราหฺมโณ เยน ภควา เตนุปสงฺกมิ, อุปสงฺกมิตฺวา ภควนฺตํ เอตทโวจ “กิํ ชจฺโจ ภวนฺ”ติ.}

\csromangathameasure{มา ชาติํ ปุจฺฉ จรณญฺจ ปุจฺฉ, \\ กฏฺฐา หเว ชายติ ชาตเวโท. \\ นีจากุลีโนปิ มุนิ ธิติมา, \\ อาชานีโย โหติ หิรีนิเสโธ. \\ สจฺเจน ทนฺโต ทมสา อุเปโต, \\ เวทนฺตคู วุสิตพฺรหฺมจริโย. \\ ยญฺโญปนีโต ตมุปวฺหเยถ, \\ กาเลน โส ชุหติ ทกฺขิเณยฺเยติ. \\ อทฺธา สุยิฏฺฐํ สุหุตํ มม ยิทํ, \\ ยํ ตาทิสํ เวทคุมทฺทสามิ. \\ ตุมฺหาทิสานํ หิ อทสฺสเนน, \\ อญฺโญ ชโน ภุญฺชติ หพฺยเสสนฺติ.}
\csromangathagroup{\csromangathastanza{\csromanfolio{170}มา ชาติํ ปุจฺฉ จรณญฺจ ปุจฺฉ, \\ กฏฺฐา หเว ชายติ ชาตเวโท. \\ นีจากุลีโนปิ มุนิ ธิติมา, \\ อาชานีโย โหติ หิรีนิเสโธ.}\csromangathastanza{สจฺเจน ทนฺโต ทมสา อุเปโต, \\ เวทนฺตคู วุสิตพฺรหฺมจริโย. \\ ยญฺโญปนีโต ตมุปวฺหเยถ, \\ กาเลน โส ชุหติ ทกฺขิเณยฺเยติ.}\csromangathastanza{อทฺธา สุยิฏฺฐํ สุหุตํ มม ยิทํ, \\ ยํ ตาทิสํ เวทคุมทฺทสามิ. \\ ตุมฺหาทิสานํ หิ อทสฺสเนน, \\ อญฺโญ ชโน ภุญฺชติ หพฺยเสสนฺติ.}}

\prose{ภุญฺชตุ ภวํ โคตโม, พฺราหฺมโณ ภวนฺติ.}

\csromangathameasure{คาถาภิคีตํ เม อโภชเนยฺยํ, \\ สมฺปสฺสตํ พฺราหฺมณ เนส ธมฺโม. \\ คาถาภิคีตํ ปนุทนฺติ พุทฺธา, \\ ธมฺเม สติ พฺราหฺมณ วุตฺติเรสา. \\ อญฺเญน จ เกวลินํ มเหสิํ, \\ ขีณาสวํ กุกฺกุจฺจวูปสนฺตํ. \\ อนฺเนน ปาเนน อุปฏฺฐหสฺสุ, \\ เขตฺตํ หิ ตํ ปุญฺญเปกฺขสฺส โหตีติ.}
\csromangathagroup{\csromangathastanza{คาถาภิคีตํ เม อโภชเนยฺยํ, \\ สมฺปสฺสตํ พฺราหฺมณ เนส ธมฺโม. \\ คาถาภิคีตํ ปนุทนฺติ พุทฺธา, \\ ธมฺเม สติ พฺราหฺมณ วุตฺติเรสา.}\csromangathastanza{อญฺเญน จ เกวลินํ มเหสิํ, \\ ขีณาสวํ กุกฺกุจฺจวูปสนฺตํ. \\ อนฺเนน ปาเนน อุปฏฺฐหสฺสุ, \\ เขตฺตํ หิ ตํ ปุญฺญเปกฺขสฺส โหตีติ.}}

\prose{อถ กสฺส จาหํ โภ โคตม อิมํ หพฺยเสสํ ทมฺมีติ.\csromanspacer{} น ขฺวาหํ พฺราหฺมณ ปสฺสามิ สเทวเก โลเก สมารเก สพฺรหฺมเก สสฺสมณพฺราหฺมณิยา ปชาย สเทวมนุสฺสาย, ยสฺเสโส หพฺยเสโส ภุตฺโต สมฺมา ปริณามํ คจฺเฉยฺย, อญฺญตฺร พฺราหฺมณ ตถาคตสฺส วา ตถาคตสาวกสฺส วา.\csromanspacer{} เตน หิ ตฺวํ พฺราหฺมณ ตํ หพฺยเสสํ อปฺปหริเต วา ฉฏฺเฏหิ อปฺปาณเก วา อุทเก โอปิลาเปหีติ.}

\prose{อถ โข สุนฺทริกภารทฺวาโช พฺราหฺมโณ ตํ หพฺยเสสํ อปฺปาณเก อุทเก โอปิลาเปสิ.\csromanspacer{} อถ โข โส หพฺยเสโส อุทเก ปกฺขิตฺโต}

\csromanheader{2}{7. พฺราหฺมณสํยุตฺต}
\prose{\csromanfolio{171}จิจฺจิฏายติ จิฏิจิฏายติ สนฺธูปายติ สมฺปธูปายติ.\csromanspacer{} เสยฺยถาปิ นาม ผาโล\footnote{โลโห (ก)} ทิวสํสนฺตตฺโต\footnote{ทิวสสนฺตตฺโต (สี, สฺยา, กํ, อิ)} อุทเก ปกฺขิตฺโต จิจฺจิฏายติ จิฏิจิฏายติ สนฺธูปายติ สมฺปธูปายติ, เอวเมว โส หพฺยเสโส อุทเก ปกฺขิตฺโต จิจฺจิฏายติ จิฏิจิฏายติ สนฺธูปายติ สมฺปธูปายติ.}

\prose{อถ โข สุนฺทริกภารทฺวาโช พฺราหฺมโณ สํวิคฺโค โลมหฏฺฐชาโต เยน ภควา เตนุปสงฺกมิ, อุปสงฺกมิตฺวา เอกมนฺตํ อฏฺฐาสิ, เอกมนฺตํ ฐิตํ โข สุนฺทริกภารทฺวาชํ พฺราหฺมณํ ภควา คาถาหิ อชฺฌภาสิ\csromandash{}}

\csromangathameasure{“มา พฺราหฺมณ ทารุ สมาทหาโน, \\ สุทฺธิํ อมญฺญิ พหิทฺธา หิ เอตํ. \\ น หิ เตน สุทฺธิํ กุสลา วทนฺติ, \\ โย พาหิเรน ปริสุทฺธิมิจฺเฉ. \\ หิตฺวา อหํ พฺราหฺมณ ทารุทาหํ, \\ อชฺฌตฺตเมวุชฺชลยามิ โชติํ. \\ นิจฺจคฺคินี นิจฺจสมาหิตตฺโต, \\ อรหํ อหํ พฺรหฺมจริยํ จรามิ. \\ มาโน หิ เต พฺราหฺมณ ขาริภาโร, \\ โกโธ ธุโม ภสฺมนิ โมสวชฺชํ. \\ ชิวฺหา สุชา หทยํ โชติฐานํ, \\ อตฺตา สุทนฺโต ปุริสสฺส โชติ. \\ ธมฺโม รหโท พฺราหฺมณ สีลติตฺโถ, \\ อนาวิโล สพฺภิ สตํ ปสตฺโถ. \\ ยตฺถ หเว เวทคุโน สินาตา, \\ อนลฺลคตฺตาว ตรนฺติ ปารํ. \\ สจฺจํ ธมฺโม สํยโม พฺรหฺมจริยํ, \\ มชฺเฌ สิตา พฺราหฺมณ พฺรหฺมปตฺติ. \\ ส ตุชฺชุภูเตสุ นโม กโรหิ, \\ ตมหํ นรํ ธมฺมสารีติ พฺรูมี”ติ.}
\csromangathagroup{\csromangathastanza{“มา พฺราหฺมณ ทารุ สมาทหาโน, \\ สุทฺธิํ อมญฺญิ พหิทฺธา หิ เอตํ. \\ น หิ เตน สุทฺธิํ กุสลา วทนฺติ, \\ โย พาหิเรน ปริสุทฺธิมิจฺเฉ.}\csromangathastanza{หิตฺวา อหํ พฺราหฺมณ ทารุทาหํ, \\ อชฺฌตฺตเมวุชฺชลยามิ\footnote{อชฺฌตฺตเมว ชลยามิ (สี, สฺยา, กํ, อิ)} โชติํ. \\ นิจฺจคฺคินี นิจฺจสมาหิตตฺโต, \\ อรหํ อหํ พฺรหฺมจริยํ จรามิ.}\csromangathastanza{มาโน หิ เต พฺราหฺมณ ขาริภาโร, \\ โกโธ ธุโม ภสฺมนิ โมสวชฺชํ. \\ ชิวฺหา สุชา หทยํ โชติฐานํ, \\ อตฺตา สุทนฺโต ปุริสสฺส โชติ.}\csromangathastanza{ธมฺโม รหโท พฺราหฺมณ สีลติตฺโถ, \\ อนาวิโล สพฺภิ สตํ ปสตฺโถ. \\ ยตฺถ หเว เวทคุโน สินาตา, \\ อนลฺลคตฺตาว\footnote{อนลฺลีนคตฺตาว (สี, อิ, ก)} ตรนฺติ ปารํ.}\csromangathastanza{สจฺจํ ธมฺโม สํยโม พฺรหฺมจริยํ, \\ มชฺเฌ สิตา พฺราหฺมณ พฺรหฺมปตฺติ. \\ ส ตุชฺชุภูเตสุ นโม กโรหิ, \\ ตมหํ นรํ ธมฺมสารีติ พฺรูมี”ติ.}}

\prose{\csromanfolio{172}เอวํ วุตฺเต สุนฺทริกภารทฺวาโช พฺราหฺมโณ ภควนฺตํ เอตทโวจ อภิกฺกนฺตํ โภ โคตม ฯเปฯ.\csromanspacer{} อญฺญตโร จ ปนายสฺมา ภารทฺวาโช อรหตํ อโหสีติ.}

\csromanheader{2}{10. พหุธีตรสุตฺต}
\proseitem{196}{เอกํ สมยํ ภควา โกสเลสุ วิหรติ อญฺญตรสฺมิํ วนสณฺเฑ.\csromanspacer{} เตน โข ปน สมเยน อญฺญตรสฺส ภารทฺวาชโคตฺตสฺส พฺราหฺมณสฺส จตุทฺทส พลีพทฺทา นฏฺฐา โหนฺติ.\csromanspacer{} อถ โข ภารทฺวาชโคตฺโต พฺราหฺมโณ เต พลีพทฺเท คเวสนฺโต เยน โส วนสณฺโฑ เตนุปสงฺกมิ, อุปสงฺกมิตฺวา อทฺทส ภควนฺตํ ตสฺมิํ วนสณฺเฑ นิสินฺนํ ปลฺลงฺกํ อาภุชิตฺวา อุชุํ กายํ ปณิธาย ปริมุขํ สติํ อุปฏฺฐเปตฺวา, ทิสฺวาน เยน ภควา เตนุปสงฺกมิ, อุปสงฺกมิตฺวา ภควโต สนฺติเก อิมา คาถาโย อภาสิ\csromandash{}}

\csromangathameasure{“น หิ นูนิมสฺส สมณสฺส, พลีพทฺทา จตุทฺทส. \\ อชฺชสฏฺฐิํ น ทิสฺสนฺติ, เตนายํ สมโณ สุขี. \\ น หิ นูนิมสฺส สมณสฺส, ติลาเขตฺตสฺมิ ปาปกา. \\ เอกปณฺณา ทุปณฺณา จ, เตนายํ สมโณ สุขี. \\ น หิ นูนิมสฺส สมณสฺส, ตุจฺฉโกฏฺฐสฺมิ มูสิกา. \\ อุสฺโสฬฺหิกาย นจฺจนฺติ, เตนายํ สมโณ สุขี. \\ น หิ นูนิมสฺส สมณสฺส, สนฺถาโร สตฺตมาสิโก. \\ อุปฺปาฏเกหิ สญฺฉนฺโน, เตนายํ สมโณ สุขี. \\ น หิ นูนิมสฺส สมณสฺส, วิธวา สตฺต ธีตโร. \\ เอกปุตฺตา ทุปุตฺตา จ, เตนายํ สมโณ สุขี. \\ น หิ นูนิมสฺส สมณสฺส, ปิงฺคลา ติลกาหตา. \\ โสตฺตํ ปาเทน โพเธติ, เตนายํ สมโณ สุขี. \\ น หิ นูนิมสฺส สมณสฺส, ปจฺจูสมฺหิ อิณายิกา. \\ เทถ เทถาติ โจเทนฺติ, เตนายํ สมโณ สุขี”ติ.}
\csromangathagroup{\csromangathastanza{“น หิ นูนิมสฺส\footnote{นหนูนิมสฺส (สี, สฺยา, กํ)} สมณสฺส, พลีพทฺทา จตุทฺทส. \\ อชฺชสฏฺฐิํ น ทิสฺสนฺติ, เตนายํ สมโณ สุขี.}\csromangathastanza{น หิ นูนิมสฺส สมณสฺส, ติลาเขตฺตสฺมิ ปาปกา. \\ เอกปณฺณา ทุปณฺณา\footnote{ทฺวิปณฺณา (สี, อิ)} จ, เตนายํ สมโณ สุขี.}\csromangathastanza{น หิ นูนิมสฺส สมณสฺส, ตุจฺฉโกฏฺฐสฺมิ มูสิกา. \\ อุสฺโสฬฺหิกาย นจฺจนฺติ, เตนายํ สมโณ สุขี.}\csromangathastanza{น หิ นูนิมสฺส สมณสฺส, สนฺถาโร สตฺตมาสิโก. \\ อุปฺปาฏเกหิ สญฺฉนฺโน, เตนายํ สมโณ สุขี.}\csromangathastanza{น หิ นูนิมสฺส สมณสฺส, วิธวา สตฺต ธีตโร. \\ เอกปุตฺตา ทุปุตฺตา\footnote{ทฺวิปุตฺตา (สี, อิ)} จ, เตนายํ สมโณ สุขี.}\csromangathastanza{น หิ นูนิมสฺส สมณสฺส, ปิงฺคลา ติลกาหตา. \\ โสตฺตํ ปาเทน โพเธติ, เตนายํ สมโณ สุขี.}\csromangathastanza{น หิ นูนิมสฺส สมณสฺส, ปจฺจูสมฺหิ อิณายิกา. \\ เทถ เทถาติ โจเทนฺติ, เตนายํ สมโณ สุขี”ติ.}}

\csromanheader{2}{7. พฺราหฺมณสํยุตฺต}
\csromangathameasure{น หิ มยฺหํ พฺราหฺมณ, พลีพทฺทา จตุทฺทส. \\ อชฺชสฏฺฐิํ น ทิสฺสนฺติ, เตนาหํ พฺราหฺมณา สุขี. \\ น หิ มยฺหํ พฺราหฺมณ, ติลาเขตฺตสฺมิ ปาปกา. \\ เอกปณฺณา ทุปณฺณา จ, เตนาหํ พฺราหฺมณา สุขี. \\ น หิ มยฺหํ พฺราหฺมณ, ตุจฺฉโกฏฺฐสฺมิ มูสิกา. \\ อุสฺโสฬฺหิกาย นจฺจนฺติ, เตนาหํ พฺราหฺมณา สุขี. \\ น หิ มยฺหํ พฺราหฺมณ, สนฺถาโร สตฺตมาสิโก. \\ อุปฺปาฏเกหิ สญฺฉนฺโน, เตนาหํ พฺราหฺมณา สุขี. \\ น หิ มยฺหํ พฺราหฺมณ, วิธวา สตฺต ธีตโร. \\ เอกปุตฺตา ทุปุตฺตา จ, เตนาหํ พฺราหฺมณา สุขี. \\ น หิ มยฺหํ พฺราหฺมณ, ปิงฺคลา ติลกาหตา. \\ โสตฺตํ ปาเทน โพเธติ, เตนาหํ พฺราหฺมณา สุขี. \\ น หิ มยฺหํ พฺราหฺมณ, ปจฺจูสมฺหิ อิณายิกา. \\ เทถ เทถาติ โจเทนฺติ, เตนาหํ พฺราหฺมณา สุขีติ.}
\csromangathagroup{\csromangathastanza{\csromanfolio{173}น หิ มยฺหํ พฺราหฺมณ, พลีพทฺทา จตุทฺทส. \\ อชฺชสฏฺฐิํ น ทิสฺสนฺติ, เตนาหํ พฺราหฺมณา สุขี.}\csromangathastanza{น หิ มยฺหํ พฺราหฺมณ, ติลาเขตฺตสฺมิ ปาปกา. \\ เอกปณฺณา ทุปณฺณา จ, เตนาหํ พฺราหฺมณา สุขี.}\csromangathastanza{น หิ มยฺหํ พฺราหฺมณ, ตุจฺฉโกฏฺฐสฺมิ มูสิกา. \\ อุสฺโสฬฺหิกาย นจฺจนฺติ, เตนาหํ พฺราหฺมณา สุขี.}\csromangathastanza{น หิ มยฺหํ พฺราหฺมณ, สนฺถาโร สตฺตมาสิโก. \\ อุปฺปาฏเกหิ สญฺฉนฺโน, เตนาหํ พฺราหฺมณา สุขี.}\csromangathastanza{น หิ มยฺหํ พฺราหฺมณ, วิธวา สตฺต ธีตโร. \\ เอกปุตฺตา ทุปุตฺตา จ, เตนาหํ พฺราหฺมณา สุขี.}\csromangathastanza{น หิ มยฺหํ พฺราหฺมณ, ปิงฺคลา ติลกาหตา. \\ โสตฺตํ ปาเทน โพเธติ, เตนาหํ พฺราหฺมณา สุขี.}\csromangathastanza{น หิ มยฺหํ พฺราหฺมณ, ปจฺจูสมฺหิ อิณายิกา. \\ เทถ เทถาติ โจเทนฺติ, เตนาหํ พฺราหฺมณา สุขีติ.}}

\prose{เอวํ วุตฺเต ภารทฺวาชโคตฺโต พฺราหฺมโณ ภควนฺตํ เอตทโวจ “อภิกฺกนฺตํ โภ โคตม, อภิกฺกนฺตํ โภ โคตม, เสยฺยถาปิ โภ โคตม นิกฺกุชฺชิตํ วา อุกฺกุชฺเชยฺย, ปฏิจฺฉนฺนํ วา วิวเรยฺย, มูฬฺหสฺส วา มคฺคํ อาจิกฺเขยฺย, อนฺธกาเร วา เตลปชฺโชตํ ธาเรยฺย ‘จกฺขุมนฺโต รูปานิ ทกฺกนฺตี’ติ.\csromanspacer{} เอวเมว โภตา โคตเมน อเนกปริยาเยน ธมฺโม ปกาสิโต, เอสาหํ ภวนฺตํ โคตมํ สรณํ คจฺฉามิ ธมฺมญฺจ ภิกฺขุสํฆญฺจ, ลเภยฺยาหํ โภโต โคตมสฺส สนฺติเก ปพฺพชฺชํ, ลเภยฺยํ อุปสมฺปทนฺ”ติ.}

\prose{อลตฺถ โข ภารทฺวาชโคตฺโต พฺราหฺมโณ ภควโต สนฺติเก ปพฺพชฺชํ, อลตฺถ อุปสมฺปทํ.\csromanspacer{} อจิรูปสมฺปนฺโน ปนายสฺมา ภารทฺวาโช เอโก วูปกฏฺโฐ อปฺปมตฺโต อาตาปี ปหิตตฺโต วิหรนฺโต นจิรสฺเสว, ยสฺสตฺถาย กุลปุตฺตา สมฺมเทว อคารสฺมา อนคาริยํ ปพฺพชนฺติ, ตทนุตฺตรํ พฺรหฺมจริยปริโยสานํ ทิฏฺเฐว ธมฺเม สยํ อภิญฺญา สจฺฉิกตฺวา อุปสมฺปชฺช วิหาสิ, “ขีณา ชาติ, วุสิตํ พฺรหฺมจริยํ, กตํ กรณียํ,}

\prose{\csromanfolio{174}นาปรํ อิตฺถตฺตายา”ติ อพฺภญฺญาสิ.\csromanspacer{} อญฺญตโร จ ปนายสฺมา ภารทฺวาโช อรหตํ อโหสีติ.}

\vspace{\tipitakaparskip}
\csromancenter{อรหนฺตวคฺโค ปฐโม.}
\titlehead{\textbf{ตสฺสุทฺทานํ}}
\csromangathameasure{ธนญฺชานี จ อกฺโกสํ, อสุรินฺทํ พิลงฺคิกํ. \\ อหิํสกํ ชฏา เจว, สุทฺธิกํ เจว อคฺคิกา. \\ สุนฺทริกํ พหุธีตเรน จ เต ทสาติ.}
\csromangathagroup{\csromangathastanza{ธนญฺชานี จ อกฺโกสํ, อสุรินฺทํ พิลงฺคิกํ. \\ อหิํสกํ ชฏา เจว, สุทฺธิกํ เจว อคฺคิกา. \\ สุนฺทริกํ พหุธีตเรน จ เต ทสาติ.}}

\chapterhead{2. อุปาสกวคฺค}
\csromanheader{2}{1. กสิภารทฺวาชสุตฺต}
\proseitem{197}{เอวํ เม สุตํ\csromandash{}เอกํ สมยํ ภควา มคเธสุ วิหรติ ทกฺขิณาคิริสฺมิํ เอกนาฬายํ พฺราหฺมณคาเม.\csromanspacer{} เตน โข ปน สมเยน กสิภารทฺวาชสฺส\footnote{กสิกภารทฺวาชสฺส (ก)} พฺราหฺมณสฺส ปญฺจมตฺตานิ นงฺคลสตานิ ปยุตฺตานิ โหนฺติ วปฺปกาเล.\csromanspacer{} อถ โข ภควา ปุพฺพณฺหสมยํ นิวาเสตฺวา ปตฺตจีวรมาทาย เยน กสิภารทฺวาชสฺส พฺราหฺมณสฺส กมฺมนฺโต เตนุปสงฺกมิ.}

\prose{เตน โข ปน สมเยน กสิภารทฺวาชสฺส พฺราหฺมณสฺส ปริเวสนา วตฺตติ.\csromanspacer{} อถ โข ภควา เยน ปริเวสนา เตนุปสงฺกมิ, อุปสงฺกมิตฺวา เอกมนฺตํ อฏฺฐาสิ.\csromanspacer{} อทฺทสา โข กสิภารทฺวาโช พฺราหฺมโณ ภควนฺตํ ปิณฺฑาย ฐิตํ ทิสฺวา ภควนฺตํ เอตทโวจ “อหํ โข สมณ กสามิ จ วปามิ จ, กสิตฺวา จ วปิตฺวา จ ภุญฺชามิ, ตฺวมฺปิ สมณ กสสฺสุ จ วปสฺสุ จ, กสิตฺวา จ วปิตฺวา จ ภุญฺชสฺสู”ติ.\csromanspacer{} อหมฺปิ โข พฺราหฺมณ กสามิ จ วปามิ จ, กสิตฺวา จ วปิตฺวา จ ภุญฺชามีติ.\csromanspacer{} น โข มยํ ปสฺสาม โภโต โคตมสฺส ยุคํ วา นงฺคลํ วา ผาลํ วา ปาจนํ วา พลีพทฺเท วา, อถ จ ปน ภวํ โคตโม เอวมาห “อหมฺปิ โข พฺราหฺมณ กสามิ จ วปามิ จ, กสิตฺวา จ วปิตฺวา จ}

\csromanheader{2}{7. พฺราหฺมณสํยุตฺต}
\prose{\csromanfolio{175}ภุญฺชามี”ติ.\csromanspacer{} อถ โข กสิภารทฺวาโช พฺราหฺมโณ ภควนฺตํ คาถาย อชฺฌภาสิ\csromandash{}}

\csromangathameasure{“กสฺสโก ปฏิชานาสิ, น จ ปสฺสามิ เต กสิํ. \\ กสฺสโก ปุจฺฉิโต พฺรูหิ, กถํ ชาเนมุ ตํ กสินฺ”ติ. \\ สทฺทา พีชํ ตโป วุฏฺฐิ, ปญฺญา เม ยุคนงฺคลํ. \\ หิรี อีสา มโน โยตฺตํ, สติ เม ผาลปาจนํ. \\ กายคุตฺโต วจีคุตฺโต, อาหาเร อุทเร ยโต. \\ สจฺจํ กโรมิ นิทฺทานํ, โสรจฺจํ เม ปโมจนํ. \\ วีริยํ เม ธุรโธรยฺหํ, โยคกฺเขมาธิวาหนํ. \\ คจฺฉติ อนิวตฺตนฺตํ, ยตฺถ คนฺตฺวา น โสจติ. \\ เอวเมสา กสี กฏฺฐา, สา โหติ อมตปฺผลา. \\ เอตํ กสิํ กสิตฺวาน, สพฺพทุกฺขา ปมุจฺจตีติ.}
\csromangathagroup{\csromangathastanza{“กสฺสโก ปฏิชานาสิ, น จ ปสฺสามิ เต กสิํ. \\ กสฺสโก ปุจฺฉิโต พฺรูหิ, กถํ ชาเนมุ ตํ กสินฺ”ติ.}\csromangathastanza{สทฺทา พีชํ ตโป วุฏฺฐิ, ปญฺญา เม ยุคนงฺคลํ. \\ หิรี อีสา มโน โยตฺตํ, สติ เม ผาลปาจนํ.}\csromangathastanza{กายคุตฺโต วจีคุตฺโต, อาหาเร อุทเร ยโต. \\ สจฺจํ กโรมิ นิทฺทานํ, โสรจฺจํ เม ปโมจนํ.}\csromangathastanza{วีริยํ เม ธุรโธรยฺหํ, โยคกฺเขมาธิวาหนํ. \\ คจฺฉติ อนิวตฺตนฺตํ, ยตฺถ คนฺตฺวา น โสจติ.}\csromangathastanza{เอวเมสา กสี กฏฺฐา, สา โหติ อมตปฺผลา. \\ เอตํ กสิํ กสิตฺวาน, สพฺพทุกฺขา ปมุจฺจตีติ.}}

\prose{ภุญฺชตุ ภวํ โคตโม กสฺสโก ภวํ, ยํ หิ ภวํ โคตโม อมตปฺผลมฺปิ กสิํ กสตีติ\footnote{ภาสตีติ (ก)}.}

\csromangathameasure{คาถาภิคีตํ เม อโภชเนยฺยํ, \\ สมฺปสฺสตํ พฺราหฺมณ เนส ธมฺโม. \\ คาถาภิคีตํ ปนุทนฺติ พุทฺธา, \\ ธมฺเม สติ พฺราหฺมณ วุตฺติเรสา. \\ อญฺเญน จ เกวลินํ มเหสิํ, \\ ขีณาสวํ กุกฺกุจฺจวูปสนฺตํ. \\ อนฺเนน ปาเนน อุปฏฺฐหสฺสุ, \\ เขตฺตํ หิ ตํ ปุญฺญเปกฺขสฺส โหตีติ.}
\csromangathagroup{\csromangathastanza{คาถาภิคีตํ เม อโภชเนยฺยํ, \\ สมฺปสฺสตํ พฺราหฺมณ เนส ธมฺโม. \\ คาถาภิคีตํ ปนุทนฺติ พุทฺธา, \\ ธมฺเม สติ พฺราหฺมณ วุตฺติเรสา.}\csromangathastanza{อญฺเญน จ เกวลินํ มเหสิํ, \\ ขีณาสวํ กุกฺกุจฺจวูปสนฺตํ. \\ อนฺเนน ปาเนน อุปฏฺฐหสฺสุ, \\ เขตฺตํ หิ ตํ ปุญฺญเปกฺขสฺส โหตีติ.}}

\prose{เอวํ วุตฺเต กสิภารทฺวาโช พฺราหฺมโณ ภควนฺตํ เอตทโวจ “อภิกฺกนฺตํ โภ โคตม ฯเปฯ อชฺชตคฺเค ปาณุเปตํ สรณํ คตนฺ”ติ.}

\csromanheader{2}{2. อุทยสุตฺต}
\proseitem{198}{สาวตฺถินิทานํ.\csromanspacer{} อถ โข ภควา ปุพฺพณฺหสมยํ นิวาเสตฺวา ปตฺตจีวรมาทาย เยน อุทยสฺส พฺราหฺมณสฺส นิเวสนํ เตนุปสงฺกมิ. \csromanfolio{176}อถ โข อุทโย พฺราหฺมโณ ภควโต ปตฺตํ โอทเนน ปูเรสิ.\csromanspacer{} ทุติยมฺปิ โข ภควา ปุพฺพณฺหสมยํ นิวาเสตฺวา ปตฺตจีวรมาทาย เยน อุทยสฺส พฺราหฺมณสฺส นิเวสนํ เตนุปสงฺกมิ ฯเปฯ.\csromanspacer{} ตติยมฺปิ โข อุทโย พฺราหฺมโณ ภควโต ปตฺตํ โอทเนน ปูเรตฺวา ภควนฺตํ เอตทโวจ “ปกฏฺฐโกยํ สมโณ โคตโม ปุนปฺปุนํ อาคจฺฉตี”ติ.}

\prose{ปุนปฺปุนํ เจว วปนฺติ พีชํ,}

\prose{ปุนปฺปุนํ วสฺสติ เทวราชา.}

\prose{ปุนปฺปุนํ เขตฺตํ กสนฺติ กสฺสกา,}

\prose{ปุนปฺปุนํ ธญฺญมุเปติ รฏฺฐํ.}

\prose{ปุนปฺปุนํ ยาจกา ยาจยนฺติ, ปุนปฺปุนํ ทานปตี ททนฺติ.}

\prose{ปุนปฺปุนํ ทานปตี ททิตฺวา, ปุนปฺปุนํ สคฺคมุเปนฺติ ฐานํ.}

\prose{ปุนปฺปุนํ ขีรนิกา ทุหนฺติ, ปุนปฺปุนํ วจฺโฉ อุเปติ มาตรํ.}

\prose{ปุนปฺปุนํ กิลมติ ผนฺทติ จ, ปุนปฺปุนํ คพฺภมุเปติ มนฺโท.}

\prose{ปุนปฺปุนํ ชายติ มียติ จ, ปุนปฺปุนํ สิวถิกํ\footnote{สีวถิกํ (สี, สฺยา, กํ, อิ)} หรนฺติ.}

\prose{มคฺคญฺจ ลทฺธา อปุนพฺภวาย, น ปุนปฺปุนํ ชายติ ภูริปญฺโญติ\footnote{ปุนปฺปุนํ ชายติ ภูริปญฺโญติ (สฺยา, กํ, ก)}.}

\prose{เอวํ วุตฺเต อุทโย พฺราหฺมโณ ภควนฺตํ เอตทโวจ “อภิกฺกนฺตํ โภ โคตม ฯเปฯ อุปาสกํ มํ ภวํ โคตโม ธาเรตุ อชฺชตคฺเค ปาณุเปตํ สรณํ คตนฺ”ติ.}

\csromanheader{2}{3. เทวหิตสุตฺต}
\proseitem{199}{สาวตฺถินิทานํ.\csromanspacer{} เตน โข ปน สมเยน ภควา วาเตหาพาธิโก โหติ.\csromanspacer{} อายสฺมา จ อุปวาโณ ภควโต อุปฏฺฐาโก โหติ.\csromanspacer{} อถ โข ภควา อายสฺมนฺตํ อุปวาณํ อามนฺเตสิ “อิงฺฆ เม ตฺวํ อุปวาณ อุโณฺหทกํ ชานาหี”ติ. “เอวํ ภนฺเต”ติ โข อายสฺมา อุปวาโณ ภควโต ปฏิสฺสุตฺวา นิวาเสตฺวา ปตฺตจีวรมาทาย เยน เทวหิตสฺส พฺราหฺมณสฺส นิเวสนํ เตนุปสงฺกมิ, อุปสงฺกมิตฺวา ตุณฺหีภูโต เอกมนฺตํ อฏฺฐาสิ.\csromanspacer{} อทฺทสา โข เทวหิโต พฺราหฺมโณ อายสฺมนฺตํ อุปวาณํ ตุณฺหีภูตํ เอกมนฺตํ ฐิตํ, ทิสฺวาน อายสฺมนฺตํ อุปวาณํ คาถาย อชฺฌภาสิ\csromandash{}}

\csromanheader{2}{7. พฺราหฺมณสํยุตฺต}
\csromangathameasure{“ตุณฺหีภูโต ภวํ ติฏฺฐํ, มุณฺโฑ สงฺฆาฏิปารุโต. \\ กิํ ปตฺถยาโน กิํ เอสํ, กิํ นุ ยาจิตุมาคโต”ติ. \\ อรหํ สุคโต โลเก, วาเตหาพาธิโก มุนิ. \\ สเจ อุโณฺหทกํ อตฺถิ, มุนิโน เทหิ พฺราหฺมณ. \\ ปูชิโต ปูชเนยฺยานํ, สกฺกเรยฺยาน สกฺกโต. \\ อปจิโต อปเจยฺยานํ, ตสฺส อิจฺฉามิ หาตเวติ.}
\csromangathagroup{\csromangathastanza{\csromanfolio{177}“ตุณฺหีภูโต ภวํ ติฏฺฐํ, มุณฺโฑ สงฺฆาฏิปารุโต. \\ กิํ ปตฺถยาโน กิํ เอสํ, กิํ นุ ยาจิตุมาคโต”ติ.}\csromangathastanza{อรหํ สุคโต โลเก, วาเตหาพาธิโก มุนิ. \\ สเจ อุโณฺหทกํ อตฺถิ, มุนิโน เทหิ พฺราหฺมณ.}\csromangathastanza{ปูชิโต ปูชเนยฺยานํ, สกฺกเรยฺยาน สกฺกโต. \\ อปจิโต อปเจยฺยานํ\footnote{อปจิเนยฺยานํ (สี, สฺยา, กํ) ฏีกา โอโลเกตพฺพา.}, ตสฺส อิจฺฉามิ หาตเวติ.}}

\prose{อถ โข เทวหิโต พฺราหฺมโณ อุโณฺหทกสฺส กาชํ ปุริเสน คาหาเปตฺวา ผาณิตสฺส จ ปุฏํ อายสฺมโต อุปวาณสฺส ปาทาสิ.\csromanspacer{} อถ โข อายสฺมา อุปวาโณ เยน ภควา เตนุปสงฺกมิ, อุปสงฺกมิตฺวา ภควนฺตํ อุโณฺหทเกน นฺหาเปตฺวา\footnote{นหาเปตฺวา (สี, อิ)} อุโณฺหทเกน ผาณิตํ อาโลเลตฺวา ภควโต ปาทาสิ.\csromanspacer{} อถ โข ภควโต อาพาโธ ปฏิปฺปสฺสมฺภิ.}

\prose{อถ โข เทวหิโต พฺราหฺมโณ เยน ภควา เตนุปสงฺกมิ, อุปสงฺกมิตฺวา ภควตา สทฺธิํ สมฺโมทิ, สมฺโมทนียํ กถํ สารณียํ วีติสาเรตฺวา เอกมนฺตํ นิสีทิ, เอกมนฺตํ นิสินฺโน โข เทวหิโต พฺราหฺมโณ ภควนฺตํ คาถาย อชฺฌภาสิ\csromandash{}}

\csromangathameasure{“กตฺถ ทชฺชา เทยฺยธมฺมํ, กตฺถ ทินฺนํ มหปฺผลํ. \\ กถํ หิ ยชมานสฺส, กถํ อิชฺฌติ ทกฺขิณา”ติ. \\ ปุพฺเพนิวาสํ โย เวที, สคฺคาปายญฺจ ปสฺสติ. \\ อโถ ชาติกฺขยํ ปตฺโต, อภิญฺญาโวสิโต มุนิ. \\ เอตฺถ ทชฺชา เทยฺยธมฺมํ, เอตฺถ ทินฺนํ มหปฺผลํ. \\ เอวํ หิ ยชมานสฺส, เอวํ อิชฺฌติ ทกฺขิณาติ.}
\csromangathagroup{\csromangathastanza{“กตฺถ ทชฺชา เทยฺยธมฺมํ, กตฺถ ทินฺนํ มหปฺผลํ. \\ กถํ หิ ยชมานสฺส, กถํ อิชฺฌติ ทกฺขิณา”ติ.}\csromangathastanza{ปุพฺเพนิวาสํ โย เวที, สคฺคาปายญฺจ ปสฺสติ. \\ อโถ ชาติกฺขยํ ปตฺโต, อภิญฺญาโวสิโต มุนิ.}\csromangathastanza{เอตฺถ ทชฺชา เทยฺยธมฺมํ, เอตฺถ ทินฺนํ มหปฺผลํ. \\ เอวํ หิ ยชมานสฺส, เอวํ อิชฺฌติ ทกฺขิณาติ.}}

\prose{เอวํ วุตฺเต เทวหิโต พฺราหฺมโณ ภควนฺตํ เอตทโวจ “อภิกฺกนฺตํ โภ โคตม ฯเปฯ อุปาสกํ มํ ภวํ โคตโม ธาเรตุ อชฺชตคฺเค ปาณุเปตํ สรณํ คตนฺ”ติ.}

\csromanheader{2}{4. มหาสาลสุตฺต}
\proseitem{200}{\csromanfolio{178}สาวตฺถินิทานํ.\csromanspacer{} อถ โข อญฺญตโร พฺราหฺมณมหาสาโล ลูโข ลูขปาวุรโณ\footnote{ลูขปาปุรโณ (สี, สฺยา, กํ, อิ)} เยน ภควา เตนุปสงฺกมิ, อุปสงฺกมิตฺวา ภควตา สทฺธิํ สมฺโมทิ, สมฺโมทนียํ กถํ สารณียํ วีติสาเรตฺวา เอกมนฺตํ นิสีทิ, เอกมนฺตํ นิสินฺนํ โข ตํ พฺราหฺมณมหาสาลํ ภควา เอตทโวจ “กินฺนุ ตฺวํ พฺราหฺมณ ลูโข ลูขปาวุรโณ”ติ.\csromanspacer{} อิธ เม โภ โคตม จตฺตาโร ปุตฺตา เต มํ ทาเรหิ สํปุจฺฉ ฆรา นิกฺขาเมนฺตีติ.\csromanspacer{} เตน หิ ตฺวํ พฺราหฺมณ อิมา คาถาโย ปริยาปุณิตฺวา สภายํ มหาชนกาเย สนฺนิปติเต ปุตฺเตสุ จ สนฺนิสินฺเนสุ ภาสสฺสุ\csromandash{}}

\csromangathameasure{“เยหิ ชาเตหิ นนฺทิสฺสํ, เยสญฺจ ภวมิจฺฉิสํ. \\ เต มํ ทาเรหิ สํปุจฺฉ, สาว วาเรนฺติ สูกรํ. \\ อสนฺตา กิร มํ ชมฺมา, ตาต ตาตาติ ภาสเร. \\ รกฺขสา ปุตฺตรูเปน, เต ชหนฺติ วโยคตํ. \\ อสฺโสว ชิณฺโณ นิพฺโภโค, ขาทนา อปนียติ. \\ พาลกานํ ปิตา เถโร, ปราคาเรสุ ภิกฺขติ. \\ ทณฺโฑว กิร เม เสยฺโย, ยญฺเจ ปุตฺตา อนสฺสวา. \\ จณฺฑมฺปิ โคณํ วาเรติ, อโถ จณฺฑมฺปิ กุกฺกุรํ. \\ อนฺธกาเร ปุเร โหติ, คมฺภีเร คาธเมธติ. \\ ทณฺฑสฺส อานุภาเวน, ขลิตฺวา ปติติฏฺฐตี”ติ.}
\csromangathagroup{\csromangathastanza{“เยหิ ชาเตหิ นนฺทิสฺสํ, เยสญฺจ ภวมิจฺฉิสํ. \\ เต มํ ทาเรหิ สํปุจฺฉ, สาว วาเรนฺติ สูกรํ.}\csromangathastanza{อสนฺตา กิร มํ ชมฺมา, ตาต ตาตาติ ภาสเร. \\ รกฺขสา ปุตฺตรูเปน, เต ชหนฺติ วโยคตํ.}\csromangathastanza{อสฺโสว ชิณฺโณ นิพฺโภโค, ขาทนา อปนียติ. \\ พาลกานํ ปิตา เถโร, ปราคาเรสุ ภิกฺขติ.}\csromangathastanza{ทณฺโฑว กิร เม เสยฺโย, ยญฺเจ ปุตฺตา อนสฺสวา. \\ จณฺฑมฺปิ โคณํ วาเรติ, อโถ จณฺฑมฺปิ กุกฺกุรํ.}\csromangathastanza{อนฺธกาเร ปุเร โหติ, คมฺภีเร คาธเมธติ. \\ ทณฺฑสฺส อานุภาเวน, ขลิตฺวา ปติติฏฺฐตี”ติ.}}

\prose{อถ โข โส พฺราหฺมณมหาสาโล ภควโต สนฺติเก อิมา คาถาโย ปริยาปุณิตฺวา สภายํ มหาชนกาเย สนฺนิปติเต ปุตฺเตสุ จ สนฺนิสินฺเนสุ อภาสิ\csromandash{}}

\csromangathameasure{“เยหิ ชาเตหิ นนฺทิสฺสํ, เยสญฺจ ภวมิจฺฉิสํ. \\ เต มํ ทาเรหิ สํปุจฺฉ, สาว วาเรนฺติ สูกรํ. \\ อสนฺตา กิร มํ ชมฺมา, ตาต ตาตาติ ภาสเร. \\ รกฺขสา ปุตฺตรูเปน, เต ชหนฺติ วโยคตํ.}
\csromangathagroup{\csromangathastanza{“เยหิ ชาเตหิ นนฺทิสฺสํ, เยสญฺจ ภวมิจฺฉิสํ. \\ เต มํ ทาเรหิ สํปุจฺฉ, สาว วาเรนฺติ สูกรํ.}\csromangathastanza{อสนฺตา กิร มํ ชมฺมา, ตาต ตาตาติ ภาสเร. \\ รกฺขสา ปุตฺตรูเปน, เต ชหนฺติ วโยคตํ.}}

\csromanheader{2}{7. พฺราหฺมณสํยุตฺต}
\csromangathameasure{อสฺโสว ชิณฺโณ นิพฺโภโค, ขาทนา อปนียติ. \\ พาลกานํ ปิตา เถโร, ปราคาเรสุ ภิกฺขติ. \\ ทณฺโฑว กิร เม เสยฺโย, ยญฺเจ ปุตฺตา อนสฺสวา. \\ จณฺฑมฺปิ โคณํ วาเรติ, อโถ จณฺฑมฺปิ กุกฺกุรํ. \\ อนฺธกาเร ปุเร โหติ, คมฺภีเร คาธเมธติ. \\ ทณฺฑสฺส อานุภาเวน, ขลิตฺวา ปติติฏฺฐตี”ติ.}
\csromangathagroup{\csromangathastanza{\csromanfolio{179}อสฺโสว ชิณฺโณ นิพฺโภโค, ขาทนา อปนียติ. \\ พาลกานํ ปิตา เถโร, ปราคาเรสุ ภิกฺขติ.}\csromangathastanza{ทณฺโฑว กิร เม เสยฺโย, ยญฺเจ ปุตฺตา อนสฺสวา. \\ จณฺฑมฺปิ โคณํ วาเรติ, อโถ จณฺฑมฺปิ กุกฺกุรํ.}\csromangathastanza{อนฺธกาเร ปุเร โหติ, คมฺภีเร คาธเมธติ. \\ ทณฺฑสฺส อานุภาเวน, ขลิตฺวา ปติติฏฺฐตี”ติ.}}

\prose{อถ โข นํ พฺราหฺมณมหาสาลํ ปุตฺตา ฆรํ เนตฺวา นฺหาเปตฺวา ปจฺเจกํ ทุสฺสยุเคน อจฺฉาเทสุํ.\csromanspacer{} อถ โข โส พฺราหฺมณมหาสาโล เอกํ ทุสฺสยุคํ อาทาย เยน ภควา เตนุปสงฺกมิ, อุปสงฺกมิตฺวา ภควตา สทฺธิํ สมฺโมทิ, สมฺโมทนียํ กถํ สารณียํ วีติสาเรตฺวา เอกมนฺตํ นิสีทิ, เอกมนฺตํ นิสินฺโน โข พฺราหฺมณมหาสาโล ภควนฺตํ เอตทโวจ “มยํ โภ โคตม พฺราหฺมณา นาม อาจริยสฺส อาจริยธนํ ปริเยสาม, ปฏิคฺคณฺหตุ เม ภวํ โคตโม อาจริยธนนฺ”ติ.\csromanspacer{} ปฏิคฺคเหสิ ภควา อนุกมฺปํ อุปาทาย.\csromanspacer{} อถ โข โส พฺราหฺมณมหาสาโล ภควนฺตํ เอตทโวจ “อภิกฺกนฺตํ โภ โคตม ฯเปฯ อุปาสกํ มํ ภวํ โคตโม ธาเรตุ อชฺชตคฺเค ปาณุเปตํ สรณํ คตนฺ”ติ.}

\csromanheader{2}{5. มานตฺถทฺธสุตฺต}
\proseitem{201}{สาวตฺถินิทานํ.\csromanspacer{} เตน โข ปน สมเยน มานตฺถทฺโธ นาม พฺราหฺมโณ สาวตฺถิยํ ปฏิวสติ, โส เนว มาตรํ อภิวาเทติ, น ปิตรํ อภิวาเทติ, น อาจริยํ อภิวาเทติ, น เชฏฺฐภาตรํ อภิวาเทติ.\csromanspacer{} เตน โข ปน สมเยน ภควา มหติยา ปริสาย ปริวุโต ธมฺมํ เทเสติ.\csromanspacer{} อถ โข มานตฺถทฺธสฺส พฺราหฺมณสฺส เอตทโหสิ “อยํ โข สมโณ โคตโม มหติยา ปริสาย ปริวุโต ธมฺมํ เทเสติ, ยํนูนาหํ เยน สมโณ โคตโม เตนุปสงฺกเมยฺยํ, สเจ มํ สมโณ โคตโม อาลปิสฺสติ, อหมฺปิ ตํ อาลปิสฺสามิ.\csromanspacer{} โน เจ มํ สมโณ โคตโม อาลปิสฺสติ, อหมฺปิ นาลปิสฺสามี”ติ.\csromanspacer{} อถ โข มานตฺถทฺโธ พฺราหฺมโณ เยน ภควา เตนุปสงฺกมิ, อุปสงฺกมิตฺวา ตุณฺหีภูโต เอกมนฺตํ อฏฺฐาสิ.\csromanspacer{} อถ โข ภควา ตํ นาลปิ.\csromanspacer{} อถ โข มานตฺถทฺโธ พฺราหฺมโณ “นายํ \csromanfolio{180}สมโณ โคตโม กิญฺจิ ชานาตี”ติ ตโตว ปุน นิวตฺติตุกาโม อโหสิ.\csromanspacer{} อถ โข ภควา มานตฺถทฺธสฺส พฺราหฺมณสฺส เจตสา เจโตปริวิตกฺกมญฺญาย มานตฺถทฺธํ พฺราหฺมณํ คาถาย อชฺฌภาสิ\csromandash{}}

\csromangathameasure{“น มานํ พฺราหฺมณ สาธุ, อตฺถิกสฺสีธ พฺราหฺมณ. \\ เยน อตฺเถน อาคจฺฉิ, ตเมวมนุพฺรูหเย”ติ.}
\csromangathagroup{\csromangathastanza{“น มานํ พฺราหฺมณ สาธุ, อตฺถิกสฺสีธ พฺราหฺมณ. \\ เยน อตฺเถน อาคจฺฉิ, ตเมวมนุพฺรูหเย”ติ.}}

\prose{อถ โข มานตฺถทฺโธ พฺราหฺมโณ “จิตฺตํ เม สมโณ โคตโม ชานาตี”ติ ตตฺเถว ภควโต ปาเทสุ สิรสา นิปติตฺวา ภควโต ปาทานิ มุเขน จ ปริจุมฺพติ ปาณีหิ จ ปริสมฺพาหติ, นามญฺจ สาเวติ “มานตฺถทฺธาหํ โภ โคตม, มานตฺถทฺธาหํ โภ โคตมา”ติ.\csromanspacer{} อถ โข สา ปริสา อพฺภุตวิตฺตชาตา\footnote{อพฺภุตจิตฺตชาตา (สี, สฺยา, กํ, อิ), อจฺฉริยพฺภุตจิตฺตชาตา (ก)} อโหสิ “อจฺฉริยํ วต โภ อพฺภุตํ วต โภ, อยํ หิ มานตฺถทฺโธ พฺราหฺมโณ เนว มาตรํ อภิวาเทติ, น ปิตรํ อภิวาเทติ, น อาจริยํ อภิวาเทติ, น เชฏฺฐภาตรํ อภิวาเทติ.\csromanspacer{} อถ จ ปน สมเณ โคตเม เอวรูปํ ปรมนิปจฺจการํ กโรตี”ติ.\csromanspacer{} อถ โข ภควา มานตฺถทฺธํ พฺราหฺมณํ เอตทโวจ “อลํ พฺราหฺมณ อุฏฺเฐหิ สเก อาสเน นิสีท, ยโต เต มยิ จิตฺตํ ปสนฺนนฺ”ติ.\csromanspacer{} อถ โข มานตฺถทฺโธ พฺราหฺมโณ สเก อาสเน นิสีทิตฺวา ภควนฺตํ คาถาย อชฺฌภาสิ\csromandash{}}

\csromangathameasure{“เกสุ น มานํ กยิราถ, เกสุ จสฺส สคารโว. \\ กฺยสฺส อปจิตา อสฺสุ, กฺยสฺสุ สาธุ สุปูชิตา”ติ. \\ มาตริ ปิตริ จาปิ, อโถ เชฏฺฐมฺหิ ภาตริ.}
\csromangathagroup{\csromangathastanza{“เกสุ น มานํ กยิราถ, เกสุ จสฺส สคารโว. \\ กฺยสฺส อปจิตา อสฺสุ, กฺยสฺสุ สาธุ สุปูชิตา”ติ. \\ มาตริ ปิตริ จาปิ, อโถ เชฏฺฐมฺหิ ภาตริ.}}

\prose{อาจริเย จตุตฺถมฺหิ.}

\csromangathameasure{เตสุ น มานํ กยิราถ, \\ เตสุ อสฺส สคารโว. \\ ตฺยสฺส อปจิตา อสฺสุ, \\ ตฺยสฺสุ สาธุ สุปูชิตา. \\ อรหนฺเต สีตีภูเต, กตกิจฺเจ อนาสเว. \\ นิหจฺจ มานํ อถทฺโธ, เต นมสฺเส อนุตฺตเรติ.}
\csromangathagroup{\csromangathastanza{เตสุ น มานํ กยิราถ, \\ เตสุ อสฺส สคารโว. \\ ตฺยสฺส อปจิตา อสฺสุ, \\ ตฺยสฺสุ สาธุ สุปูชิตา. \\ อรหนฺเต สีตีภูเต, กตกิจฺเจ อนาสเว. \\ นิหจฺจ มานํ อถทฺโธ, เต นมสฺเส อนุตฺตเรติ.}}

\csromanheader{2}{7. พฺราหฺมณสํยุตฺต}
\prose{\csromanfolio{181}เอวํ วุตฺเต มานตฺถทฺโธ พฺราหฺมโณ ภควนฺตํ เอตทโวจ “อภิกฺกนฺตํ โภ โคตม ฯเปฯ อุปาสกํ มํ ภวํ โคตโม ธาเรตุ อชฺชตคฺเค ปาณุเปตํ สรณํ คตนฺ”ติ.}

\csromanheader{2}{6. ปจฺจนีกสุตฺต}
\proseitem{202}{สาวตฺถินิทานํ.\csromanspacer{} เตน โข ปน สมเยน ปจฺจนีกสาโต นาม พฺราหฺมโณ สาวตฺถิยํ ปฏิวสติ.\csromanspacer{} อถ โข ปจฺจนีกสาตสฺส พฺราหฺมณสฺส เอตทโหสิ “ยํนูนาหํ เยน สมโณ โคตโม เตนุปสงฺกเมยฺยํ, ยํ ยเทว สมโณ โคตโม ภาสิสฺสติ, ตํ ตเทวสฺสาหํ\footnote{ตเทว สาหํ (ก)} ปจฺจนีกาสฺสนฺ”ติ\footnote{ปจฺจนีกสฺสนฺติ (อิ), ปจฺจนีกสาตนฺติ (ก)}.\csromanspacer{} เตน โข ปน สมเยน ภควา อพฺโภกาเส จงฺกมติ.\csromanspacer{} อถ โข ปจฺจนีกสาโต พฺราหฺมโณ เยน ภควา เตนุปสงฺกมิ, อุปสงฺกมิตฺวา ภควนฺตํ จงฺกมนฺตํ เอตทโวจ “ภณ สมณธมฺมนฺ”ติ.}

\csromangathameasure{น ปจฺจนีกสาเตน, สุวิชานํ สุภาสิตํ. \\ อุปกฺกิลิฏฺฐจิตฺเตน, สารมฺภพหุเลน จ. \\ โย จ วิเนยฺย สารมฺภํ, อปฺปสาทญฺจ เจตโส. \\ อาฆาตํ ปฏินิสฺสชฺช, ส เว ชญฺญา สุภาสิตนฺติ.}
\csromangathagroup{\csromangathastanza{น ปจฺจนีกสาเตน, สุวิชานํ สุภาสิตํ. \\ อุปกฺกิลิฏฺฐจิตฺเตน, สารมฺภพหุเลน จ.}\csromangathastanza{โย จ วิเนยฺย สารมฺภํ, อปฺปสาทญฺจ เจตโส. \\ อาฆาตํ ปฏินิสฺสชฺช, ส เว\footnote{สเจ (สฺยา, กํ, ก)} ชญฺญา สุภาสิตนฺติ.}}

\prose{เอวํ วุตฺเต ปจฺจนีกสาโต พฺราหฺมโณ ภควนฺตํ เอตทโวจ “อภิกฺกนฺตํ โภ โคตม, อภิกฺกนฺตํ โภ โคตม ฯเปฯ อุปาสกํ มํ ภวํ โคตโม ธาเรตุ อชฺชตคฺเค ปาณุเปตํ สรณํ คตนฺ”ติ.}

\csromanheader{2}{7. นวกมฺมิกสุตฺต}
\proseitem{203}{เอกํ สมยํ ภควา โกสเลสุ วิหรติ อญฺญตรสฺมิํ วนสณฺเฑ.\csromanspacer{} เตน โข ปน สมเยน นวกมฺมิกภารทฺวาโช พฺราหฺมโณ ตสฺมิํ วนสณฺเฑ กมฺมนฺตํ การาเปติ.\csromanspacer{} อทฺทสา โข นวกมฺมิกภารทฺวาโช พฺราหฺมโณ ภควนฺตํ อญฺญรตสฺมิํ สาลรุกฺขมูเล นิสินฺนํ ปลฺลงฺกํ อาภุชิตฺวา อุชุํ กายํ ปณิธาย ปริมุขํ สติํ อุปฏฺฐเปตฺวา, ทิสฺวานสฺส เอตทโหสิ “อหํ โข อิมสฺมิํ วนสณฺเฑ กมฺมนฺตํ การาเปนฺโต รมามิ, อยํ สมโณ โคตโม กิํ การาเปนฺโต รมตี”ติ. \csromanfolio{182}อถ โข นวกมฺมิกภารทฺวาโช พฺราหฺมโณ เยน ภควา เตนุปสงฺกมิ, อุปสงฺกมิตฺวา ภควนฺตํ คาถาย อชฺฌภาสิ\csromandash{}}

\csromangathameasure{“เก นุ กมฺมนฺตา กรียนฺติ, ภิกฺขุ สาลวเน ตว. \\ ยเทกโก อรญฺญสฺมิํ, รติํ วินฺทติ โคตโม”ติ. \\ น เม วนสฺมิํ กรณียมตฺถิ, \\ อุจฺฉินฺนมูลํ เม วนํ วิสูกํ. \\ สฺวาหํ วเน นิพฺพนโถ วิสลฺโล, \\ เอโก รเม อรติํ วิปฺปหายาติ.}
\csromangathagroup{\csromangathastanza{“เก นุ กมฺมนฺตา กรียนฺติ, ภิกฺขุ สาลวเน ตว. \\ ยเทกโก อรญฺญสฺมิํ, รติํ วินฺทติ โคตโม”ติ.}\csromangathastanza{น เม วนสฺมิํ กรณียมตฺถิ, \\ อุจฺฉินฺนมูลํ เม วนํ วิสูกํ. \\ สฺวาหํ วเน นิพฺพนโถ วิสลฺโล, \\ เอโก รเม อรติํ วิปฺปหายาติ.}}

\prose{เอวํ วุตฺเต นวกมฺมิกภารทฺวาโช พฺราหฺมโณ ภควนฺตํ เอตทโวจ “อภิกฺกนฺตํ โภ โคตม ฯเปฯ อุปาสกํ มํ ภวํ โคตโม ธาเรตุ อชฺชตคฺเค ปาณุเปตํ สรณํ คตนฺ”ติ.}

\csromanheader{2}{8. กฏฺฐหารสุตฺต}
\proseitem{204}{เอกํ สมยํ ภควา โกสเลสุ วิหรติ อญฺญตรสฺมิํ วนสณฺเฑ.\csromanspacer{} เตน โข ปน สมเยน อญฺญตรสฺส ภารทฺวาชโคตฺตสฺส พฺราหฺมณสฺส สมฺพหุลา อนฺเตวาสิกา กฏฺฐหารกา มาณวกา เยน วนสณฺโฑ เตนุปสงฺกมิํสุ, อุปสงฺกมิตฺวา อทฺทสํสุ ภควนฺตํ ตสฺมิํ วนสณฺเฑ นิสินฺนํ ปลฺลงฺกํ อาภุชิตฺวา อุชุํ กายํ ปณิธาย ปริมุขํ สติํ อุปฏฺฐเปตฺวา, ทิสฺวาน เยน ภารทฺวาชโคตฺโต พฺราหฺมโณ เตนุปสงฺกมิํสุ, อุปสงฺกมิตฺวา ภารทฺวาชโคตฺตํ พฺราหฺมณํ เอตทโวจุํ “ยคฺเฆ ภวํ ชาเนยฺยาสิ, อสุกสฺมิํ วนสณฺเฑ สมโณ นิสินฺโน ปลฺลงฺกํ อาภุชิตฺวา อุชุํ กายํ ปณิธาย ปริมุขํ สติํ อุปฏฺฐเปตฺวา.\csromanspacer{} อถ โข ภารทฺวาชโคตฺโต พฺราหฺมโณ เตหิ มาณวเกหิ สทฺธิํ เยน โส วนสณฺโฑ เตนุปสงฺกมิ.\csromanspacer{} อทฺทสา โข ภควนฺตํ ตสฺมิํ วนสณฺเฑ นิสินฺนํ ปลฺลงฺกํ อาภุชิตฺวา อุชุํ กายํ ปณิธาย ปริมุขํ สติํ อุปฏฺฐเปตฺวา, ทิสฺวาน เยน ภควา เตนุปสงฺกมิ, อุปสงฺกมิตฺวา ภควนฺตํ คาถาย อชฺฌภาสิ\csromandash{}}

\csromangathameasure{“คมฺภีรรูเป พหุเภรเว วเน, \\ สุญฺญํ อรญฺญํ วิชนํ วิคาหิย. \\ อนิญฺชมาเนน ฐิเตน วคฺคุนา, \\ สุจารุรูปํ วต ภิกฺขุ ฌายสิ.}
\csromangathagroup{\csromangathastanza{“คมฺภีรรูเป พหุเภรเว วเน, \\ สุญฺญํ อรญฺญํ วิชนํ วิคาหิย. \\ อนิญฺชมาเนน ฐิเตน วคฺคุนา, \\ สุจารุรูปํ วต ภิกฺขุ ฌายสิ.}}

\csromanheader{2}{7. พฺราหฺมณสํยุตฺต}
\csromangathameasure{น ยตฺถ คีตํ นปิ ยตฺถ วาทิตํ, \\ เอโก อรญฺเญ วนวสฺสิโต มุนิ. \\ อจฺเฉรรูปํ ปฏิภาติ มํ อิทํ, \\ ยเทกโก ปีติมโน วเน วเส. \\ มญฺญามหํ โลกาธิปติสหพฺยตํ, \\ อากงฺขมาโน ติทิวํ อนุตฺตรํ. \\ กสฺมา ภวํ วิชนมรญฺญมสฺสิโต, \\ ตโป อิธ กฺรุพฺพสิ พฺรหฺมปตฺติยา”ติ. \\ ยา กาจิ กงฺขา อภินนฺทนา วา, \\ อเนกธาตูสุ ปุถู สทาสิตา. \\ อญฺญาณมูลปฺปภวา ปชปฺปิตา, \\ สพฺพา มยา พฺยนฺติกตา สมูลิกา. \\ สฺวาหํ อกงฺโข อสิโต อนูปโย, \\ สพฺเพสุ ธมฺเมสุ วิสุทฺธทสฺสโน. \\ ปปฺปุยฺย สมฺโพธิมนุตฺตรํ สิวํ, \\ ฌายามหํ พฺรหฺม รโห วิสารโทติ.}
\csromangathagroup{\csromangathastanza{\csromanfolio{183}น ยตฺถ คีตํ นปิ ยตฺถ วาทิตํ, \\ เอโก อรญฺเญ วนวสฺสิโต มุนิ. \\ อจฺเฉรรูปํ ปฏิภาติ มํ อิทํ, \\ ยเทกโก ปีติมโน วเน วเส.}\csromangathastanza{มญฺญามหํ โลกาธิปติสหพฺยตํ, \\ อากงฺขมาโน ติทิวํ อนุตฺตรํ. \\ กสฺมา ภวํ วิชนมรญฺญมสฺสิโต, \\ ตโป อิธ กฺรุพฺพสิ พฺรหฺมปตฺติยา”ติ.}\csromangathastanza{ยา กาจิ กงฺขา อภินนฺทนา วา, \\ อเนกธาตูสุ ปุถู สทาสิตา. \\ อญฺญาณมูลปฺปภวา ปชปฺปิตา, \\ สพฺพา มยา พฺยนฺติกตา สมูลิกา.}\csromangathastanza{สฺวาหํ อกงฺโข อสิโต อนูปโย, \\ สพฺเพสุ ธมฺเมสุ วิสุทฺธทสฺสโน. \\ ปปฺปุยฺย สมฺโพธิมนุตฺตรํ สิวํ, \\ ฌายามหํ พฺรหฺม รโห วิสารโทติ.}}

\prose{เอวํ วุตฺเต ภารทฺวาชโคตฺโต พฺราหฺมโณ ภควนฺตํ เอตทโวจ “อภิกฺกนฺตํ โภ โคตม, อภิกฺกนฺตํ โภ โคตม ฯเปฯ อชฺชตคฺเค ปาณุเปตํ สรณํ คตนฺ”ติ.}

\csromanheader{2}{9. มาตุโปสกสุตฺต}
\proseitem{205}{สาวตฺถินิทานํ.\csromanspacer{} อถ โข มาตุโปสโก พฺราหฺมโณ เยน ภควา เตนุปสงฺกมิ, อุปสงฺกมิตฺวา ภควตา สทฺธิํ สมฺโมทิ, สมฺโมทนียํ กถํ สารณียํ วีติสาเรตฺวา เอกมนฺตํ นิสีทิ, เอกมนฺตํ นิสินฺโน โข มาตุโปสโก พฺราหฺมโณ ภควนฺตํ เอตทโวจ “อหํ หิ โภ โคตม ธมฺเมน ภิกฺขํ ปริเยสามิ, ธมฺเมน ภิกฺขํ ปริเยสิตฺวา มาตาปิตโร โปเสมิ, กจฺจาหํ โภ โคตม เอวํการี กิจฺจการี โหมี”ติ.\csromanspacer{} ตคฺฆ ตฺวํ พฺราหฺมณ เอวํการี กิจฺจการี โหสิ, โย โข พฺราหฺมณ ธมฺเมน ภิกฺขํ ปริเยสติ, ธมฺเมน ภิกฺขํ ปริเยสิตฺวา มาตาปิตโร โปเสติ, พหุํ โส ปุญฺญํ ปสวตีติ.}

\csromangathameasure{โย มาตรํ ปิตรํ วา, มจฺโจ ธมฺเมน โปสติ. \\ ตาย นํ ปาริจริยาย, มาตาปิตูสุ ปณฺฑิตา. \\ อิเธว นํ ปสํสนฺติ, เปจฺจ สคฺเค ปโมทตีติ.}
\csromangathagroup{\csromangathastanza{\csromanfolio{184}โย มาตรํ ปิตรํ วา, มจฺโจ ธมฺเมน โปสติ. \\ ตาย นํ ปาริจริยาย, มาตาปิตูสุ ปณฺฑิตา. \\ อิเธว นํ ปสํสนฺติ, เปจฺจ สคฺเค ปโมทตีติ.}}

\prose{เอวํ วุตฺเต มาตุโปสโก พฺราหฺมโณ ภควนฺตํ เอตทโวจ “อภิกฺกนฺตํ โภ โคตม, อภิกฺกนฺตํ โภ โคตม ฯเปฯ อุปาสกํ มํ ภวํ โคตโม ธาเรตุ อชฺชตคฺเค ปาณุเปตํ สรณํ คตนฺ”ติ.}

\csromanheader{2}{10. ภิกฺขกสุตฺต}
\proseitem{206}{สาวตฺถินิทานํ.\csromanspacer{} อถ โข ภิกฺขโก พฺราหฺมโณ เยน ภควา เตนุปสงฺกมิ, อุปสงฺกมิตฺวา ภควตา สทฺธิํ สมฺโมทิ, สมฺโมทนียํ กถํ สารณียํ วีติสาเรตฺวา เอกมนฺตํ นิสีทิ, เอกมนฺตํ นิสินฺโน โข ภิกฺขโก พฺราหฺมโณ ภควนฺตํ เอตทโวจ “อหมฺปิ โข โภ โคตม ภิกฺขโก, ภวมฺปิ ภิกฺขโก, อิธ โน กิํ นานากรณนฺ”ติ.}

\csromangathameasure{น เตน ภิกฺขโก โหติ, ยาวตา ภิกฺขเต ปเร. \\ วิสฺสํ ธมฺมํ สมาทาย, ภิกฺขุ โหติ น ตาวตา. \\ โยธ ปุญฺญญฺจ ปาปญฺจ, พาหิตฺวา พฺรหฺมจริยํ. \\ สงฺขาย โลเก จรติ, ส เว ภิกฺขูติ วุจฺจตีติ.}
\csromangathagroup{\csromangathastanza{น เตน ภิกฺขโก โหติ, ยาวตา ภิกฺขเต ปเร. \\ วิสฺสํ ธมฺมํ สมาทาย, ภิกฺขุ โหติ น ตาวตา.}\csromangathastanza{โยธ ปุญฺญญฺจ ปาปญฺจ, พาหิตฺวา พฺรหฺมจริยํ. \\ สงฺขาย โลเก จรติ, ส เว ภิกฺขูติ วุจฺจตีติ.}}

\prose{เอวํ วุตฺเต ภิกฺขโก พฺราหฺมโณ ภควนฺตํ เอตทโวจ “อภิกฺกนฺตํ โภ โคตม, อภิกฺกนฺตํ โภ โคตม ฯเปฯ อุปาสกํ มํ ภวํ โคตโม ธาเรตุ อชฺชตคฺเค ปาณุเปตํ สรณํ คตนฺ”ติ.}

\csromanheader{2}{11. สงฺคารวสุตฺต}
\proseitem{207}{สาวตฺถินิทานํ.\csromanspacer{} เตน โข ปน สมเยน สงฺคารโว นาม พฺราหฺมโณ สาวตฺถิยํ ปฏิวสติ, อุทกสุทฺธิโก อุทเกน ปริสุทฺธิํ ปจฺเจติ, สายํ ปาตํ อุทโกโรหนานุโยคมนุยุตฺโต วิหรติ.\csromanspacer{} อถ โข อายสฺมา อานนฺโท ปุพฺพณฺหสมยํ นิวาเสตฺวา ปตฺตจีวรมาทาย สาวตฺถิํ ปิณฺฑาย ปาวิสิ, สาวตฺถิยํ ปิณฺฑาย จริตฺวา ปจฺฉาภตฺตํ ปิณฺฑปาตปฏิกฺกนฺโต เยน ภควา เตนุปสงฺกมิ, อุปสงฺกมิตฺวา ภควนฺตํ อภิวาเทตฺวา เอกมนฺตํ นิสีทิ, เอกมนฺตํ นิสินฺโน โข อายสฺมา อานนฺโท ภควนฺตํ เอตทโวจ “อิธ ภนฺเต สงฺคารโว นาม พฺราหฺมโณ สาวตฺถิยํ ปฏิวสติ,}

\csromanheader{2}{7. พฺราหฺมณสํยุตฺต}
\prose{\csromanfolio{185}อุทกสุทฺธิโก อุทเกน สุทฺธิํ ปจฺเจติ, สายํ ปาตํ อุทโกโรหนานุโยคมนุยุตฺโต วิหรติ.\csromanspacer{} สาธุ ภนฺเต ภควา เยน สงฺคารวสฺส พฺราหฺมณสฺส นิเวสนํ เตนุปสงฺกมตุ อนุกมฺปํ อุปาทายา”ติ.\csromanspacer{} อธิวาเสสิ ภควา ตุณฺหีภาเวน.}

\prose{อถ โข ภควา ปุพฺพณฺหสมยํ นิวาเสตฺวา ปตฺตจีวรมาทาย เยน สงฺคารวสฺส พฺราหฺมณสฺส นิเวสนํ เตนุปสงฺกมิ, อุปสงฺกมิตฺวา ปญฺญตฺเต อาสเน นิสีทิ.\csromanspacer{} อถ โข สงฺคารโว พฺราหฺมโณ เยน ภควา เตนุปสงฺกมิ, อุปสงฺกมิตฺวา ภควตา สทฺธิํ สมฺโมทิ, สมฺโมทนียํ กถํ สารณียํ วีติสาเรตฺวา เอกมนฺตํ นิสีทิ, เอกมนฺตํ นิสินฺนํ โข สงฺคารวํ พฺราหฺมณํ ภควา เอตทโวจ “สจฺจํ กิร ตฺวํ พฺราหฺมณ อุทกสุทฺธิโก อุทเกน สุทฺธิํ ปจฺเจสิ, สายํ ปาตํ อุทโกโรหนานุโยคมนุยุตฺโต วิหรสี”ติ.\csromanspacer{} เอวํ โภ โคตม.\csromanspacer{} กิํ ปน ตฺวํ พฺราหฺมณ อตฺถวสํ สมฺปสฺสมาโน อุทกสุทฺธิโก อุทกสุทฺธิํ ปจฺเจสิ, สายํ ปาตํ อุทโกโรหนานุโยคมนุยุตฺโต วิหรสีติ.\csromanspacer{} อิธ เม โภ โคตม\footnote{อิธ เม โภ โคตม อหํ (อิ, ก)} ยํ ทิวา ปาปกมฺมํ กตํ โหติ, ตํ สายํ นฺหาเนน\footnote{นหาเนน (สี, สฺยา, กํ, อิ)} ปวาเหมิ.\csromanspacer{} ยํ รตฺติํ ปาปกมฺมํ กตํ โหติ, ตํ ปาตํ นฺหาเนน ปวาเหมิ.\csromanspacer{} อิมํ ขฺวาหํ โภ โคตม อตฺถวสํ สมฺปสฺสมาโน อุทกสุทฺธิโก อุทเกน สุทฺธิํ ปจฺเจมิ, สายํ ปาตํ อุทโกโรหนานุโยคมนุยุตฺโต วิหรามีติ.}

\csromangathameasure{ธมฺโม รหโท พฺราหฺมณ สีลติตฺโถ, \\ อนาวิโล สพฺภิ สตํ ปสตฺโถ. \\ ยตฺถ หเว เวทคุโน สินาตา, \\ อนลฺลคตฺตาว ตรนฺติ ปารนฺติ.}
\csromangathagroup{\csromangathastanza{ธมฺโม รหโท พฺราหฺมณ สีลติตฺโถ, \\ อนาวิโล สพฺภิ สตํ ปสตฺโถ. \\ ยตฺถ หเว เวทคุโน สินาตา, \\ อนลฺลคตฺตาว\footnote{อนลฺลีนคตฺตาว (ก)} ตรนฺติ ปารนฺติ.}}

\prose{เอวํ วุตฺเต สงฺคารโว พฺราหฺมโณ ภควนฺตํ เอตทโวจ “อภิกฺกนฺตํ โภ โคตม, อภิกฺกนฺตํ โภ โคตม ฯเปฯ อุปาสกํ มํ ภวํ โคตโม ธาเรตุ อชฺชตคฺเค ปาณุเปตํ สรณํ คตนฺ”ติ.}

\csromanheader{2}{12. โขมทุสฺสสุตฺต}
\proseitem{208}{เอวํ เม สุตํ\csromandash{}เอกํ สมยํ ภควา สกฺเกสุ วิหรติ โขมทุสฺสํ นาม สกฺยานํ นิคโม.\csromanspacer{} อถ โข ภควา ปุพฺพณฺหสมยํ \csromanfolio{186}นิวาเสตฺวา ปตฺตจีวรมาทาย โขมทุสฺสํ นิคมํ ปิณฺฑาย ปาวิสิ.\csromanspacer{} เตน โข ปน สมเยน โขมทุสฺสกา พฺราหฺมณคหปติกา สภายํ สนฺนิปติตา โหนฺติ เกนจิเทว กรณีเยน, เทโว จ เอกเมกํ ผุสายติ.\csromanspacer{} อถ โข ภควา เยน สา สภา เตนุปสงฺกมิ.\csromanspacer{} อทฺทสํสุ โขมทุสฺสกา พฺราหฺมณคหปติกา ภควนฺตํ ทูรโตว อาคจฺฉนฺตํ, ทิสฺวาน เอตทโวจุํ “เก จ มุณฺฑกา สมณกา, เก จ สภาธมฺมํ ชานิสฺสนฺตี”ติ.\csromanspacer{} อถ โข ภควา โขมทุสฺสเก พฺราหฺมณคหปติเก คาถาย อชฺฌภาสิ\csromandash{}}

\csromangathameasure{“เนสา สภา ยตฺถ น สนฺติ สนฺโต, \\ สนฺโต น เต เย น วทนฺติ ธมฺมํ. \\ ราคญฺจ โทสญฺจ ปหาย โมหํ, \\ ธมฺมํ วทนฺตา จ ภวนฺติ สนฺโต”ติ.}
\csromangathagroup{\csromangathastanza{“เนสา สภา ยตฺถ น สนฺติ สนฺโต, \\ สนฺโต น เต เย น วทนฺติ ธมฺมํ. \\ ราคญฺจ โทสญฺจ ปหาย โมหํ, \\ ธมฺมํ วทนฺตา จ ภวนฺติ สนฺโต”ติ.}}

\prose{เอวํ วุตฺเต โขมทุสฺสกา พฺราหฺมณคหปติกา ภควนฺตํ เอตทโวจุํ “อภิกฺกนฺตํ โภ โคตม, อภิกฺกนฺตํ โภ โคตม, เสยฺยถาปิ โภ โคตม นิกฺกุชฺชิตํ วา อุกฺกุชฺเชยฺย, ปฏิจฺฉนฺนํ วา วิวเรยฺย, มูฬฺหสฺส วา มคฺคํ อาจิกฺเขยฺย, อนฺธกาเร วา เตลปชฺโชตํ ธาเรยฺย ‘จกฺขุมนฺโต รูปานิ ทกฺขนฺตี’ติ.\csromanspacer{} เอวเมวํ โภตา โคตเมน อเนกปริยาเยน ธมฺโม ปกาสิโต, เอเต มยํ ภวนฺตํ โคตมํ สรณํ คจฺฉาม ธมฺมญฺจ ภิกฺขุสํฆญฺจ, อุปาสเก โน ภวํ โคตโม ธาเรตุ อชฺชตคฺเค ปาณุเปเต สรณํ คเต”ติ.}

\vspace{\tipitakaparskip}
\csromancenter{อุปาสกวคฺโค ทุติโย.}
\titlehead{\textbf{ตสฺสุทฺทานํ}}
\csromangathameasure{กสิ อุทโย เทวหิโต, อญฺญตรมหาสาลํ. \\ มานถทฺธํ ปจฺจนีกํ, นวกมฺมิ กฏฺฐหารํ. \\ มาตุโปสกํ ภิกฺขโก, สงฺคารโว จ โขมทุสฺเสน ทฺวาทสาติ. \\ \textbf{พฺราหฺมณสํยุตฺตํ สมตฺตํ.}}
\csromangathagroup{\csromangathastanza{กสิ อุทโย เทวหิโต, อญฺญตรมหาสาลํ. \\ มานถทฺธํ ปจฺจนีกํ, นวกมฺมิ กฏฺฐหารํ.}\csromangathastanza{มาตุโปสกํ ภิกฺขโก, สงฺคารโว จ โขมทุสฺเสน ทฺวาทสาติ. \\ \textbf{พฺราหฺมณสํยุตฺตํ สมตฺตํ.}}}

\csromanheader{2}{8. วงฺคีสสํยุตฺต}
\csromanheader{2}{1. นิกฺขนฺตสุตฺต}
\proseitem{209}{\csromanfolio{187}เอวํ เม สุตํ\csromandash{}เอกํ สมยํ อายสฺมา วงฺคีโส อาฬวิยํ วิหรติ อคฺคาฬเว เจติเย อายสฺมตา นิโคฺรธกปฺเปน อุปชฺฌาเยน สทฺธิํ.\csromanspacer{} เตน โข ปน สมเยน อายสฺมา วงฺคีโส นวโก โหติ อจิรปพฺพชิโต โอหิยฺยโก วิหารปาโล.\csromanspacer{} อถ โข สมฺพหุลา อิตฺถิโย สมลงฺกริตฺวา เยน อคฺคาฬวโก อาราโม เตนุปสงฺกมิํสุ วิหารเปกฺขิกาโย.\csromanspacer{} อถ โข อายสฺมโต วงฺคีสสฺส ตา อิตฺถิโย ทิสฺวา อนภิรติ อุปฺปชฺชติ, ราโค จิตฺตํ อนุทฺธํเสติ.\csromanspacer{} อถ โข อายสฺมโต วงฺคีสสฺส เอตทโหสิ “อลาภา วต เม, น วต เม ลาภา, ทุลฺลทฺธํ วต เม, น วต เม สุลทฺธํ.\csromanspacer{} ยสฺส เม อนภิรติ อุปฺปนฺนา, ราโค จิตฺตํ อนุทฺธํเสติ.\csromanspacer{} ตํ กุเตตฺถ ลพฺภา, ยํ เม ปโร อนภิรติํ วิโนเทตฺวา อภิรติํ อุปฺปาเทยฺย.\csromanspacer{} ยํนูนาหํ อตฺตนาว อตฺตโน อนภิรติํ วิโนเทตฺวา อภิรติํ อุปฺปาเทยฺยนฺ”ติ.\csromanspacer{} อถ โข อายสฺมา วงฺคีโส อตฺตนาว อตฺตโน อนภิรติํ วิโนเทตฺวา อภิรติํ อุปฺปาเทตฺวา ตายํ เวลายํ อิมา คาถาโย อภาสิ\csromandash{}}

\csromangathameasure{“นิกฺขนฺตํ วต มํ สนฺตํ, อคารสฺมานคาริยํ. \\ วิตกฺกา อุปธาวนฺติ, ปคพฺภา กณฺหาโต อิเม. \\ อุคฺคปุตฺตา มหิสฺสาสา, สิกฺขิตา ทฬฺหธมฺมิโน. \\ สมนฺตา ปริกิเรยฺยุํ, สหสฺสํ อปลายินํ. \\ สเจปิ เอตโต ภิยฺโย, อาคมิสฺสนฺติ อิตฺถิโย. \\ เนว มํ พฺยาธยิสฺสนฺติ, ธมฺเม สมฺหิ ปติฏฺฐิตํ. \\ สกฺขี หิ เม สุตํ เอตํ, พุทฺธสฺสาทิจฺจพนฺธุโน. \\ นิพฺพานคมนํ มคฺคํ, ตตฺถ เม นิรโต มโน. \\ เอวญฺเจ มํ วิหรนฺตํ, ปาปิม อุปคจฺฉสิ. \\ ตถา มจฺจุ กริสฺสามิ, น เม มคฺคมฺปิ ทกฺขสี”ติ.}
\csromangathagroup{\csromangathastanza{“นิกฺขนฺตํ วต มํ สนฺตํ, อคารสฺมานคาริยํ. \\ วิตกฺกา อุปธาวนฺติ, ปคพฺภา กณฺหาโต อิเม.}\csromangathastanza{อุคฺคปุตฺตา มหิสฺสาสา, สิกฺขิตา ทฬฺหธมฺมิโน. \\ สมนฺตา ปริกิเรยฺยุํ, สหสฺสํ อปลายินํ.}\csromangathastanza{สเจปิ เอตโต\footnote{เอตฺตโต (สี, อิ, ก), เอตฺตกา (สฺยา, กํ)} ภิยฺโย, อาคมิสฺสนฺติ อิตฺถิโย. \\ เนว มํ พฺยาธยิสฺสนฺติ\footnote{เอตฺตโต (สี, อิ, ก), เอตฺตกา (สฺยา, กํ)}, ธมฺเม สมฺหิ ปติฏฺฐิตํ.}\csromangathastanza{สกฺขี หิ เม สุตํ เอตํ, พุทฺธสฺสาทิจฺจพนฺธุโน. \\ นิพฺพานคมนํ มคฺคํ, ตตฺถ เม นิรโต มโน.}\csromangathastanza{เอวญฺเจ มํ วิหรนฺตํ, ปาปิม อุปคจฺฉสิ. \\ ตถา มจฺจุ กริสฺสามิ, น เม มคฺคมฺปิ ทกฺขสี”ติ.}}

\csromanheader{2}{2. อรตีสุตฺต}
\proseitem{210}{\csromanfolio{188}เอกํ สมยํ ฯเปฯ อายสฺมา วงฺคีโส อาฬวิยํ วิหรติ อคฺคาฬเว เจติเย อายสฺมตา นิโคฺรธกปฺเปน อุปชฺฌาเยน สทฺธิํ.\csromanspacer{} เตน โข ปน สมเยน อายสฺมา นิโคฺรธกปฺโป ปจฺฉาภตฺตํ ปิณฺฑปาตปฏิกฺกนฺโต วิหารํ ปวิสติ, สายํ วา นิกฺขมติ อปฺปรชฺชุ วา กาเล.\csromanspacer{} เตน โข ปน สมเยน อายสฺมโต วงฺคีสสฺส อนภิรติ อุปฺปนฺนา โหติ, ราโค จิตฺตํ อนุทฺธํเสติ.\csromanspacer{} อถ โข อายสฺมโต วงฺคีสสฺส เอตทโหสิ “อลาภา วต เม, น วต เม ลาภา.\csromanspacer{} ทุลฺลทฺธํ วต เม, น วต เม สุลทฺธํ.\csromanspacer{} ยสฺส เม อนภิรติ อุปฺปนฺนา, ราโค จิตฺตํ อนุทฺธํเสติ.\csromanspacer{} ตํ กุเตตฺถ ลพฺภา, ยํ เม ปโร อนภิรติํ วิโนเทตฺวา อภิรติํ อุปฺปาเทยฺย.\csromanspacer{} ยํนูนาหํ อตฺตนาว อตฺตโน อนภิรติํ วิโนเทตฺวา อภิรติํ อุปฺปาเทยฺยนฺ”ติ.\csromanspacer{} อถ โข อายสฺมา วงฺคีโส อตฺตนาว อตฺตโน อนภิรติํ วิโนเทตฺวา อภิรติํ อุปฺปาเทตฺวา ตายํ เวลายํ อิมา คาถาโย อภาสิ\csromandash{}}

\prose{อรติญฺจ รติญฺจ ปหาย, สพฺพโส เคหสิตญฺจ วิตกฺกํ.}

\prose{วนถํ น กเรยฺย กุหิญฺจิ, นิพฺพนโถ อรโต ส หิ ภิกฺขุ\footnote{ส ภิกฺขุ (ก)}.}

\prose{ยมิธ ปถวิญฺจ เวหาสํ, รูปคตญฺจ ชคโตคธํ.}

\prose{กิญฺจิ ปริชียติ สพฺพมนิจฺจํ, เอวํ สเมจฺจ จรนฺติ มุตตฺตา.}

\prose{อุปธีสุ ชนา คธิตาเส\footnote{คถิตาเส (สี)}, ทิฏฺฐสุเต ปฏิเฆ จ มุเต จ.}

\prose{เอตฺถ วิโนทย ฉนฺทมเนโช, โย เอตฺถ น ลิมฺปติ ตํ มุนิมาหุ.}

\prose{อถ สฏฺฐินิสฺสิตา สวิตกฺกา, ปุถู ชนตาย อธมฺมา นิวิฏฺฐา.}

\prose{น จ วคฺคคตสฺส กุหิญฺจิ, โน ปน ทุฏฺฐุลฺลภาณี ส ภิกฺขุ.}

\prose{ทพฺโพ จิรรตฺตสมาหิโต, อกุหโก นิปโก อปิหาลุ.}

\prose{สนฺตํ ปทํ อชฺฌคมา มุนิ, ปฏิจฺจ ปรินิพฺพุโต กงฺขติ กาลนฺติ.}

\csromanheader{2}{3. เปสลสุตฺต}
\proseitem{211}{เอกํ สมยํ อายสฺมา วงฺคีโส อาฬวิยํ วิหรติ อคฺคาฬเว เจติเย อายสฺมตา นิโคฺรธกปฺเปน อุปชฺฌาเยน สทฺธิํ.\csromanspacer{} เตน โข ปน}

\vspace{\tipitakaparskip}
\csromancenter{วงฺคีสสํยุตฺต สมเยน อายสฺมา วงฺคีโส อตฺตโน ปฏิภาเนน อญฺเญ เปสเล ภิกฺขู อติมญฺญติ.\csromanspacer{} อถ โข อายสฺมโต วงฺคีสสฺส เอตทโหสิ “อลาภา วต เม, น วต เม ลาภา.\csromanspacer{} ทุลฺลทฺธํ วต เม, น วต เม สุลทฺธํ.\csromanspacer{} ยฺวาหํ อตฺตโน ปฏิภาเนน อญฺเญ เปสเล ภิกฺขู อติมญฺญามี”ติ.\csromanspacer{} อถ โข อายสฺมา วงฺคีโส อตฺตนาว อตฺตโน วิปฺปฏิสารํ อุปฺปาเทตฺวา ตายํ เวลายํ อิมา คาถาโย อภาสิ\csromandash{}}
\vspace{0.5\baselineskip}
\csromangathameasure{มานํ ปชหสฺสุ โคตม, มานปถญฺจ ปชหสฺสุ อเสสํ. \\ มานปถสฺมิํ ส มุจฺฉิโต, วิปฺปฏิสารีหุวา จิรรตฺตํ. \\ มกฺเขน มกฺขิตา ปชา, มานหตา นิรยํ ปปตนฺติ. \\ โสจนฺติ ชนา จิรรตฺตํ, มานหตา นิรยํ อุปปนฺนา. \\ น หิ โสจติ ภิกฺขุ กทาจิ, มคฺคชิโน สมฺมาปฏิปนฺโน. \\ กิตฺติญฺจ สุขญฺจ อนุโภติ, ธมฺมทโสติ ตมาหุ ปหิตตฺตํ. \\ ตสฺมา อขิโลธ ปธานวา, นีวรณานิ ปหาย วิสุทฺโธ. \\ มานญฺจ ปหาย อเสสํ, วิชฺชายนฺตกโร สมิตาวีติ.}
\csromangathagroup{\csromangathastanza{\csromanfolio{189}มานํ ปชหสฺสุ โคตม, มานปถญฺจ ปชหสฺสุ อเสสํ. \\ มานปถสฺมิํ ส มุจฺฉิโต, วิปฺปฏิสารีหุวา จิรรตฺตํ.}\csromangathastanza{มกฺเขน มกฺขิตา ปชา, มานหตา นิรยํ ปปตนฺติ. \\ โสจนฺติ ชนา จิรรตฺตํ, มานหตา นิรยํ อุปปนฺนา.}\csromangathastanza{น หิ โสจติ ภิกฺขุ กทาจิ, มคฺคชิโน สมฺมาปฏิปนฺโน. \\ กิตฺติญฺจ สุขญฺจ อนุโภติ, ธมฺมทโสติ ตมาหุ ปหิตตฺตํ.}\csromangathastanza{ตสฺมา อขิโลธ ปธานวา, นีวรณานิ ปหาย วิสุทฺโธ. \\ มานญฺจ ปหาย อเสสํ, วิชฺชายนฺตกโร สมิตาวีติ.}}

\csromanheader{2}{4. อานนฺทสุตฺต}
\proseitem{212}{เอกํ สมยํ อายสฺมา อานนฺโท สาวตฺถิยํ วิหรติ เชตวเน อนาถปิณฺฑิกสฺส อาราเม.\csromanspacer{} อถ โข อายสฺมา อานนฺโท ปุพฺพณฺหสมยํ นิวาเสตฺวา ปตฺตจีวรมาทาย สาวตฺถิํ ปิณฺฑาย ปาวิสิ อายสฺมตา วงฺคีเสน ปจฺฉาสมเณน.\csromanspacer{} เตน โข ปน สมเยน อายสฺมโต วงฺคีสสฺส อนภิรติ อุปฺปนฺนา โหติ, ราโค จิตฺตํ อนุทฺธํเสติ.\csromanspacer{} อถ โข อายสฺมา วงฺคีโส อายสฺมนฺตํ อานนฺทํ คาถาย อชฺฌภาสิ\csromandash{}}

\csromangathameasure{กามราเคน ฑยฺหามิ, จิตฺตํ เม ปริฑยฺหติ. \\ สาธุ นิพฺพาปนํ พฺรูหิ, อนุกมฺปาย โคตมาติ. \\ สญฺญาย วิปริเยสา, จิตฺตํ เต ปริฑยฺหติ. \\ นิมิตฺตํ ปริวชฺเชหิ, สุภํ ราคูปสํหิตํ. \\ สงฺขาเร ปรโต ปสฺส, ทุกฺขโต มา จ อตฺตโต. \\ นิพฺพาเปหิ มหาราคํ, มา ฑยฺหิตฺโถ ปุนปฺปุนํ. \\ อสุภาย จิตฺตํ ภาเวหิ, เอกคฺคํ สุสมาหิตํ. \\ สติ กายคตา ตฺยตฺถุ, นิพฺพิทาพหุโล ภว. \\ อนิมิตฺตญฺจ ภาเวหิ, มานานุสยมุชฺชห. \\ ตโต มานาภิสมยา, อุปสนฺโต จริสฺสสีติ.}
\csromangathagroup{\csromangathastanza{กามราเคน ฑยฺหามิ, จิตฺตํ เม ปริฑยฺหติ. \\ สาธุ นิพฺพาปนํ พฺรูหิ, อนุกมฺปาย โคตมาติ.}\csromangathastanza{สญฺญาย วิปริเยสา, จิตฺตํ เต ปริฑยฺหติ. \\ นิมิตฺตํ ปริวชฺเชหิ, สุภํ ราคูปสํหิตํ.}\csromangathastanza{สงฺขาเร ปรโต ปสฺส, ทุกฺขโต มา จ อตฺตโต. \\ นิพฺพาเปหิ มหาราคํ, มา ฑยฺหิตฺโถ ปุนปฺปุนํ.}\csromangathastanza{\csromanfolio{190}อสุภาย จิตฺตํ ภาเวหิ, เอกคฺคํ สุสมาหิตํ. \\ สติ กายคตา ตฺยตฺถุ, นิพฺพิทาพหุโล ภว.}\csromangathastanza{อนิมิตฺตญฺจ ภาเวหิ, มานานุสยมุชฺชห. \\ ตโต มานาภิสมยา, อุปสนฺโต จริสฺสสีติ.}}

\csromanheader{2}{5. สุภาสิตสุตฺต}
\proseitem{213}{สาวตฺถินิทานํ.\csromanspacer{} ตตฺร โข ภควา ภิกฺขู อามนฺเตสิ “ภิกฺขโว”ติ. “ภทนฺเต”ติ เต ภิกฺขู ภควโต ปจฺจสฺโสสุํ.\csromanspacer{} ภควา เอตทโวจ\csromandash{}}

\prose{จตูหิ ภิกฺขเว องฺเคหิ สมนฺนาคตา วาจา สุภาสิตา โหติ, โน ทุพฺภาสิตา, อนวชฺชา จ อนนุวชฺชา จ วิญฺญูนํ.\csromanspacer{} กตเมหิ จตูหิ, อิธ ภิกฺขเว ภิกฺขุ สุภาสิตํเยว ภาสติ, โน ทุพฺภาสิตํ.\csromanspacer{} ธมฺมํเยว ภาสติ, โน อธมฺมํ.\csromanspacer{} ปิยํเยว ภาสติ, โน อปฺปิยํ.\csromanspacer{} สจฺจํเยว ภาสติ, โน อลิกํ.\csromanspacer{} อิเมหิ โข ภิกฺขเว จตูหิ องฺเคหิ สมนฺนาคตา วาจา สุภาสิตา โหติ, โน ทุพฺภาสิตา, อนวชฺชา จ อนนุวชฺชา จ วิญฺญูนนฺติ.\csromanspacer{} อิทมโวจ ภควา, อิทํ วตฺวาน สุคโต อถาปรํ เอตทโวจ สตฺถา\csromandash{}}

\csromangathameasure{“สุภาสิตํ อุตฺตมมาหุ สนฺโต, \\ ธมฺมํ ภเณ นาธมฺมํ ตํ ทุติยํ. \\ ปิยํ ภเณ นาปฺปิยํ ตํ ตติยํ, \\ สจฺจํ ภเณ นาลิกํ ตํ จตุตฺถนฺ”ติ.}
\csromangathagroup{\csromangathastanza{“สุภาสิตํ อุตฺตมมาหุ สนฺโต, \\ ธมฺมํ ภเณ นาธมฺมํ ตํ ทุติยํ. \\ ปิยํ ภเณ นาปฺปิยํ ตํ ตติยํ, \\ สจฺจํ ภเณ นาลิกํ ตํ จตุตฺถนฺ”ติ.}}

\prose{อถ โข อายสฺมา วงฺคีโส อุฏฺฐายาสนา เอกํสํ อุตฺตราสงฺคํ กริตฺวา เยน ภควา เตนญฺชลิํ ปณาเมตฺวา ภควนฺตํ เอตทโวจ “ปฏิภาติ มํ ภควา, ปฏิภาติ มํ สุคตา”ติ. “ปฏิภาตุ ตํ วงฺคีสา”ติ ภควา อโวจ.\csromanspacer{} อถ โข อายสฺมา วงฺคีโส ภควนฺตํ สมฺมุขา สารุปฺปาหิ คาถาหิ อภิตฺถวิ\csromandash{}}

\csromangathameasure{“ตเมว วาจํ ภาเสยฺย, ยายตฺตานํ น ตาปเย. \\ ปเร จ น วิหิํเสยฺย, สา เว วาจา สุภาสิตา. \\ ปิยวาจํว ภาเสยฺย, ยา วาจา ปฏินนฺทิตา. \\ ยํ อนาทาย ปาปานิ, ปเรสํ ภาสเต ปิยํ.}
\csromangathagroup{\csromangathastanza{“ตเมว วาจํ ภาเสยฺย, ยายตฺตานํ น ตาปเย. \\ ปเร จ น วิหิํเสยฺย, สา เว วาจา สุภาสิตา.}\csromangathastanza{ปิยวาจํว ภาเสยฺย, ยา วาจา ปฏินนฺทิตา. \\ ยํ อนาทาย ปาปานิ, ปเรสํ ภาสเต ปิยํ.}}

\csromanheader{2}{8. วงฺคีสสํยุตฺต}
\csromangathameasure{สจฺจํ เว อมตา วาจา, เอส ธมฺโม สนนฺตโน, \\ สจฺเจ อตฺเถ จ ธมฺเม จ, \\ อาหุ สนฺโต ปติฏฺฐิตา. \\ ยํ พุทฺโธ ภาสเต วาจํ, \\ เขมํ นิพฺพานปตฺติยา. ทุกฺขสฺสนฺตกิริยาย, \\ สา เว วาจานมุตฺตมา”ติ.}
\csromangathagroup{\csromangathastanza{\csromanfolio{191}สจฺจํ เว อมตา วาจา, เอส ธมฺโม สนนฺตโน, \\ สจฺเจ อตฺเถ จ ธมฺเม จ, \\ อาหุ สนฺโต ปติฏฺฐิตา. \\ ยํ พุทฺโธ ภาสเต วาจํ, \\ เขมํ นิพฺพานปตฺติยา. ทุกฺขสฺสนฺตกิริยาย, \\ สา เว วาจานมุตฺตมา”ติ.}}

\csromanheader{2}{6. สาริปุตฺตสุตฺต}
\proseitem{214}{เอกํ สมยํ อายสฺมา สาริปุตฺโต สาวตฺถิยํ วิหรติ เชตวเน อนาถปิณฺฑิกสฺส อาราเม.\csromanspacer{} เตน โข ปน สมเยน อายสฺมา สาริปุตฺโต ภิกฺขู ธมฺมิยา กถาย สนฺทสฺเสติ สมาทเปติ สมุตฺเตเชติ สมฺปหํเสติ โปริยา วาจาย วิสฺสฏฺฐาย อเนลคลาย\footnote{อเนฬคลาย (สี, ก), อเนลคฬาย (สฺยา, กํ, อิ)} อตฺถสฺส วิญฺญาปนิยา, เต จ ภิกฺขู อฏฺฐิํกตฺวา มนสิ กตฺวา สพฺพเจตสา สมนฺนาหริตฺวา โอหิตโสตา ธมฺมํ สุณนฺติ.\csromanspacer{} อถ โข อายสฺมโต วงฺคีสสฺส เอตทโหสิ “อยํ โข อายสฺมา สาริปุตฺโต ภิกฺขู ธมฺมิยา กถาย สนฺทสฺเสติ สมาทเปติ สมุตฺเตเชติ สมฺปหํเสติ โปริยา วาจาย วิสฺสฏฺฐาย อเนลคลาย อตฺถสฺส วิญฺญาปนิยา, เต จ ภิกฺขู อฏฺฐิํ กตฺวา มนสิ กตฺวา สพฺพเจตสา สมนฺนาหริตฺวา โอหิตโสตา ธมฺมํ สุณนฺติ.\csromanspacer{} ยํนูนาหํ อายสฺมนฺตํ สาริปุตฺตํ สมฺมุขา สารุปฺปาหิ คาถาหิ อภิตฺถเวยฺยนฺ”ติ.}

\prose{อถ โข อายสฺมา วงฺคีโส อุฏฺฐายาสนา เอกํสํ อุตฺตราสงฺคํ กริตฺวา เยนายสฺมา สาริปุตฺโต เตนญฺชลิํ ปณาเมตฺวา อายสฺมนฺตํ สาริปุตฺตํ เอตทโวจ “ปฏิภาติ มํ อาวุโส สาริปุตฺต, ปฏิภาติ มํ อาวุโส สาริปุตฺตา”ติ. “ปฏิภาตุ ตํ อาวุโส วงฺคีสา”ติ.\csromanspacer{} อถ โข อายสฺมา วงฺคีโส อายสฺมนฺตํ สาริปุตฺตํ สมฺมุขา สารุปฺปาหิ คาถาหิ อภิตฺถวิ\csromandash{}}

\csromangathameasure{คมฺภีรปญฺโญ เมธาวี, มคฺคามคฺคสฺส โกวิโท. \\ สาริปุตฺโต มหาปญฺโญ, ธมฺมํ เทเสติ ภิกฺขุนํ. \\ สํขิตฺเตนปิ เทเสติ, วิตฺถาเรนปิ ภาสติ. \\ สาฬิกายิว นิคฺโฆโส, ปฏิภานํ อุทีรยิ. \\ ตสฺส ตํ เทสยนฺตสฺส, สุณนฺติ มธุรํ คิรํ. \\ สเรน รชนีเยน, สวนีเยน วคฺคุนา. \\ อุทคฺคจิตฺตา มุทิตา, โสตํ โอเธนฺติ ภิกฺขโวติ.}
\csromangathagroup{\csromangathastanza{คมฺภีรปญฺโญ เมธาวี, มคฺคามคฺคสฺส โกวิโท. \\ สาริปุตฺโต มหาปญฺโญ, ธมฺมํ เทเสติ ภิกฺขุนํ.}\csromangathastanza{สํขิตฺเตนปิ เทเสติ, วิตฺถาเรนปิ ภาสติ. \\ สาฬิกายิว นิคฺโฆโส, ปฏิภานํ อุทีรยิ\footnote{อุทีริยิ (สฺยา, กํ) อุทีริยติ (สามญฺญผลสุตฺตฏีกานุรูปํ)}.}\csromangathastanza{\csromanfolio{192}ตสฺส ตํ เทสยนฺตสฺส, สุณนฺติ มธุรํ คิรํ. \\ สเรน รชนีเยน, สวนีเยน วคฺคุนา. \\ อุทคฺคจิตฺตา มุทิตา, โสตํ โอเธนฺติ ภิกฺขโวติ.}}

\csromanheader{2}{7. ปวารณาสุตฺต}
\proseitem{215}{เอกํ สมยํ ภควา สาวตฺถิยํ วิหรติ ปุพฺพาราเม มิคารมาตุปาสาเท มหตา ภิกฺขุสํเฆน สทฺธิํ ปญฺจมตฺเตหิ ภิกฺขุสเตหิ สพฺเพเหว อรหนฺเตหิ.\csromanspacer{} เตน โข ปน สมเยน ภควา ตทหุโปสเถ ปนฺนรเส ปวารณาย ภิกฺขุสํฆปริวุโต อพฺโภกาเส นิสินฺโน โหติ.\csromanspacer{} อถ โข ภควา ตุณฺหีภูตํ ภิกฺขุสํฆํ อนุวิโลเกตฺวา ภิกฺขู อามนฺเตสิ “หนฺท ทานิ ภิกฺขเว ปวาเรมิ โว, น จ เม กิญฺจิ ครหถ กายิกํ วา วาจสิกํ วา”ติ.}

\prose{เอวํ วุตฺเต อายสฺมา สาริปุตฺโต อุฏฺฐายาสนา เอกํสํ อุตฺตราสงฺคํ กริตฺวา เยน ภควา เตนญฺชลิํ ปณาเมตฺวา ภควนฺตํ เอตทโวจ “น โข มยํ ภนฺเต ภควโต กิญฺจิ ครหาม กายิกํ วา วาจสิกํ วา.\csromanspacer{} ภควา หิ ภนฺเต อนุปฺปนฺนสฺส มคฺคสฺส อุปฺปาเทตา, อสญฺชาตสฺส มคฺคสฺส สญฺชเนตา, อนกฺขาตสฺส มคฺคสฺส อกฺขาตา, มคฺคญฺญู มคฺควิทู มคฺคโกวิโท.\csromanspacer{} มคฺคานุคา จ ภนฺเต เอตรหิ สาวกา วิหรนฺติ ปจฺฉา สมนฺนาคตา, อหญฺจ โข ภนฺเต ภควนฺตํ ปวาเรมิ, น จ เม ภควา กิญฺจิ ครหติ กายิกํ วา วาจสิกํ วา”ติ.}

\prose{น ขฺวาหํ เต สาริปุตฺต กิญฺจิ ครหามิ กายิกํ วา วาจสิกํ วา, ปณฺฑิโต ตฺวํ สาริปุตฺต, มหาปญฺโญ ตฺวํ สาริปุตฺต, ปุถุปญฺโญ ตฺวํ สาริปุตฺต, หาสปญฺโญ ตฺวํ สาริปุตฺต, ชวนปญฺโญ ตฺวํ สาริปุตฺต, ติกฺขปญฺโญ ตฺวํ สาริปุตฺต, นิพฺเพธิกปญฺโญ ตฺวํ สาริปุตฺต.\csromanspacer{} เสยฺยถาปิ สาริปุตฺต รญฺโญ จกฺกวตฺติสฺส เชฏฺฐปุตฺโต ปิตรา ปวตฺติตํ จกฺกํ สมฺมเทว อนุปฺปวตฺเตติ.\csromanspacer{} เอวเมว โข ตฺวํ สาริปุตฺต มยา อนุตฺตรํ ธมฺมจกฺกํ ปวตฺติตํ สมฺมเทว อนุปฺปวตฺเตสีติ.}

\prose{โน เจ กิร เม ภนฺเต ภควา กิญฺจิ ครหติ กายิกํ วา วาจสิกํ วา, อิเมสํ ปน ภนฺเต ภควา ปญฺจนฺนํ ภิกฺขุสตานํ น กิญฺจิ ครหติ กายิกํ วา วาจสิกํ วาติ.\csromanspacer{} อิเมสมฺปิ ขฺวาหํ สาริปุตฺต ปญฺจนฺนํ ภิกฺขุสตานํ}

\vspace{\tipitakaparskip}
\csromancenter{วงฺคีสสํยุตฺต น กิญฺจิ ครหามิ กายิกํ วา วาจสิกํ วา, อิเมสํ หิ สาริปุตฺต ปญฺจนฺนํ ภิกฺขุสตานํ สฏฺฐิ ภิกฺขู เตวิชฺชา สฏฺฐิ ภิกฺขู ฉฬภิญฺญา สฏฺฐิ ภิกฺขู อุภโตภาควิมุตฺตา, อถ อิตเร ปญฺญาวิมุตฺตาติ.}
\prose{\csromanfolio{193}อถ โข อายสฺมา วงฺคีโส อุฏฺฐายาสนา เอกํสํ อุตฺตราสงฺคํ กริตฺวา เยน ภควา เตนญฺชลิํ ปณาเมตฺวา ภควนฺตํ เอตทโวจ “ปฏิภาติ มํ ภควา, ปฏิภาติ มํ สุคตา”ติ. “ปฏิภาตุ ตํ วงฺคีสา”ติ ภควา อโวจ.\csromanspacer{} อถ โข อายสฺมา วงฺคีโส ภควนฺตํ สมฺมุขา สารุปฺปาหิ คาถาหิ อภิตฺถวิ\csromandash{}}

\csromangathameasure{“อชฺช ปนฺนรเส วิสุทฺธิยา, ภิกฺขู ปญฺจสตา สมาคตา. \\ สํโยชนพนฺธนจฺฉิทา, อนีฆา ขีณปุนพฺภวา อิสี. \\ จกฺกวตฺตี ยถา ราชา, อมจฺจปริวาริโต. \\ สมนฺตา อนุปริเยติ, สาครนฺตํ มหิํ อิมํ. \\ เอวํ วิชิตสงฺคามํ, สตฺถวาหํ อนุตฺตรํ. \\ สาวกา ปยิรุปาสนฺติ, เตวิชฺชา มจฺจุหายิโน. \\ สพฺเพ ภควโต ปุตฺตา, ปลาเปตฺถ น วิชฺชติ. \\ ตณฺหาสลฺลสฺส หนฺตารํ, วนฺเท อาทิจฺจพนฺธุนนฺ”ติ.}
\csromangathagroup{\csromangathastanza{“อชฺช ปนฺนรเส วิสุทฺธิยา, ภิกฺขู ปญฺจสตา สมาคตา. \\ สํโยชนพนฺธนจฺฉิทา, อนีฆา ขีณปุนพฺภวา อิสี.}\csromangathastanza{จกฺกวตฺตี ยถา ราชา, อมจฺจปริวาริโต. \\ สมนฺตา อนุปริเยติ, สาครนฺตํ มหิํ อิมํ.}\csromangathastanza{เอวํ วิชิตสงฺคามํ, สตฺถวาหํ อนุตฺตรํ. \\ สาวกา ปยิรุปาสนฺติ, เตวิชฺชา มจฺจุหายิโน.}\csromangathastanza{สพฺเพ ภควโต ปุตฺตา, ปลาเปตฺถ น วิชฺชติ. \\ ตณฺหาสลฺลสฺส หนฺตารํ, วนฺเท อาทิจฺจพนฺธุนนฺ”ติ.}}

\csromanheader{2}{8. ปโรสหสฺสสุตฺต}
\proseitem{216}{เอกํ สมยํ ภควา สาวตฺถิยํ วิหรติ เชตวเน อนาถปิณฺฑิกสฺส อาราเม มหตา ภิกฺขุสํเฆน สทฺธิํ อฑฺฒเตลเสหิ ภิกฺขุสเตหิ.\csromanspacer{} เตน โข ปน สมเยน ภควา ภิกฺขู นิพฺพานปฏิสํยุตฺตาย ธมฺมิยา กถาย สนฺทสฺเสติ สมาทเปติ สมุตฺเตเชติ สมฺปหํเสติ, เต จ ภิกฺขู อฏฺฐิํ กตฺวา มนสิ กตฺวา สพฺพเจตสา สมนฺนาหริตฺวา โอหิตโสตา ธมฺมํ สุณนฺติ.\csromanspacer{} อถ โข อายสฺมโต วงฺคีสสฺส เอตทโหสิ “อยํ โข ภควา ภิกฺขู นิพฺพานปฏิสํยุตฺตาย ธมฺมิยา กถาย สนฺทสฺเสติ สมาทเปติ สมุตฺเตเชติ สมฺปหํเสติ, เต จ ภิกฺขู อฏฺฐิํ กตฺวา มนสิ กตฺวา สพฺพเจตสา สมนฺนาหริตฺวา โอหิตโสตา ธมฺมํ สุณนฺติ.\csromanspacer{} ยํนูนาหํ ภควนฺตํ สมฺมุขา สารุปฺปาหิ คาถาหิ อภิตฺถเวยฺยนฺ”ติ.}

\prose{\csromanfolio{194}อถ โข อายสฺมา วงฺคีโส อุฏฺฐายาสนา เอกํสํ อุตฺตราสงฺคํ กริตฺวา เยน ภควา เตนญฺชลิํ ปณาเมตฺวา ภควนฺตํ เอตทโวจ “ปฏิภาติ มํ ภควา, ปฏิภาติ มํ สุคตา”ติ. “ปฏิภาตุ ตํ วงฺคีสา”ติ ภควา อโวจ.\csromanspacer{} อถ โข อายสฺมา วงฺคีโส ภควนฺตํ สมฺมุขา สารุปฺปาหิ คาถาหิ อภิตฺถวิ\csromandash{}}

\csromangathameasure{“ปโรสหสฺสํ ภิกฺขูนํ, สุคตํ ปยิรุปาสติ. \\ เทเสนฺตํ วิรชํ ธมฺมํ, นิพฺพานํ อกุโตภยํ. \\ สุณนฺติ ธมฺมํ วิมลํ, สมฺมาสมฺพุทฺธเทสิตํ. \\ โสภติ วต สมฺพุทฺโธ, ภิกฺขุสํฆปุรกฺขโต. \\ นาคนาโมสิ ภควา, อิสีนํ อิสิสตฺตโม. \\ มหาเมโฆว หุตฺวาน, สาวเก อภิวสฺสติ. \\ ทิวาวิหารา นิกฺขมฺม, สตฺถุทสฺสนกมฺยตา. \\ สาวโก เต มหาวีร, ปาเท วนฺทติ วงฺคิโส”ติ.}
\csromangathagroup{\csromangathastanza{“ปโรสหสฺสํ ภิกฺขูนํ, สุคตํ ปยิรุปาสติ. \\ เทเสนฺตํ วิรชํ ธมฺมํ, นิพฺพานํ อกุโตภยํ.}\csromangathastanza{สุณนฺติ ธมฺมํ วิมลํ, สมฺมาสมฺพุทฺธเทสิตํ. \\ โสภติ วต สมฺพุทฺโธ, ภิกฺขุสํฆปุรกฺขโต.}\csromangathastanza{นาคนาโมสิ ภควา, อิสีนํ อิสิสตฺตโม. \\ มหาเมโฆว หุตฺวาน, สาวเก อภิวสฺสติ.}\csromangathastanza{ทิวาวิหารา นิกฺขมฺม, สตฺถุทสฺสนกมฺยตา\footnote{สตฺถุทสฺสนกามตา (สี, สฺยา, กํ)}. \\ สาวโก เต มหาวีร, ปาเท วนฺทติ วงฺคิโส”ติ.}}

\prose{กิํ นุ เต วงฺคีส อิมา คาถาโย ปุพฺเพ ปริวิตกฺกิตา, อุทาหุ ฐานโสว ตํ ปฏิภนฺตีติ.\csromanspacer{} น โข เม ภนฺเต อิมา คาถาโย ปุพฺเพ ปริวิตกฺกิตา, อถ โข ฐานโสว มํ ปฏิภนฺตีติ.\csromanspacer{} เตน หิ ตํ วงฺคีส ภิยฺโยโส มตฺตาย ปุพฺเพ อปริวิตกฺกิตา คาถาโย ปฏิภนฺตูติ. “เอวํ ภนฺเต”ติ โข อายสฺมา วงฺคีโส ภควโต ปฏิสฺสุตฺวา ภิยฺโยโส มตฺตาย ภควนฺตํ ปุพฺเพ อปริวิตกฺกิตาหิ คาถาหิ อภิตฺถวิ\csromandash{}}

\csromangathameasure{“อุมฺมคฺคปถํ มารสฺส อภิภุยฺย, จรติ ปภิชฺช ขิลานิ. \\ ตํ ปสฺสถ พนฺธปมุญฺจกรํ, อสิตํ ภาคโส ปวิภชํ. \\ โอฆสฺส นิตฺถรณตฺถํ, อเนกวิหิตํ มคฺคํ อกฺขาสิ. \\ ตสฺมิญฺจ อมเต อกฺขาเต, ธมฺมทฺทสา ฐิตา อสํหีรา. \\ ปชฺโชตกโร อติวิชฺฌ, สพฺพฏฺฐิตีนํ อติกฺกมมทฺทส. \\ ญตฺวา จ สจฺฉิกตฺวา จ, อคฺคํ โส เทสยิ ทสทฺธานํ.}
\csromangathagroup{\csromangathastanza{“อุมฺมคฺคปถํ\footnote{อุมฺมคฺคสตํ (สฺยา, กํ, ก)} มารสฺส อภิภุยฺย, จรติ ปภิชฺช ขิลานิ. \\ ตํ ปสฺสถ พนฺธปมุญฺจกรํ, อสิตํ ภาคโส ปวิภชํ.}\csromangathastanza{โอฆสฺส นิตฺถรณตฺถํ, อเนกวิหิตํ มคฺคํ อกฺขาสิ. \\ ตสฺมิญฺจ อมเต อกฺขาเต, ธมฺมทฺทสา ฐิตา อสํหีรา.}\csromangathastanza{ปชฺโชตกโร อติวิชฺฌ\footnote{อติวิชฺฌ ธมฺมํ (สี, สฺยา, กํ)}, สพฺพฏฺฐิตีนํ อติกฺกมมทฺทส. \\ ญตฺวา จ สจฺฉิกตฺวา จ, อคฺคํ โส เทสยิ ทสทฺธานํ.}}

\csromanheader{2}{8. วงฺคีสสํยุตฺต}
\csromangathameasure{เอวํ สุเทสิเต ธมฺเม, \\ โก ปมาโท วิชานตํ ธมฺมํ. \\ ตสฺมา หิ ตสฺส ภควโต สาสเน, \\ อปฺปมตฺโต สทา นมสฺสมนุสิกฺเข”ติ.}
\csromangathagroup{\csromangathastanza{\csromanfolio{195}เอวํ สุเทสิเต ธมฺเม, \\ โก ปมาโท วิชานตํ ธมฺมํ\footnote{โก ปมาโท วิชานตํ (สี, สฺยา, กํ)}. \\ ตสฺมา หิ ตสฺส ภควโต สาสเน, \\ อปฺปมตฺโต สทา นมสฺสมนุสิกฺเข”ติ.}}

\csromanheader{2}{9. โกณฺฑญฺญสุตฺต}
\proseitem{217}{เอกํ สมยํ ภควา ราชคเห วิหรติ เวฬุวเน กลนฺทกนิวาเป.\csromanspacer{} อถ โข อายสฺมา อญฺญาสิโกณฺฑญฺโญ\footnote{อญฺญาโกณฺฑญฺโญ (สี, สฺยา, กํ)} สุจิรสฺเสว เยน ภควา เตนุปสงฺกมิ, อุปสงฺกมิตฺวา ภควโต ปาเทสุ สิรสา นิปติตฺวา ภควโต ปาทานิ มุเขน จ ปริจุมฺพติ, ปาณีหิ จ ปริสมฺพาหติ, นามญฺจ สาเวติ “โกณฺฑญฺโญหํ ภควา, โกณฺฑญฺโญหํ สุคตา”ติ.\csromanspacer{} อถ โข อายสฺมโต วงฺคีสสฺส เอตทโหสิ “อยํ โข อายสฺมา อญฺญาสิโกณฺฑญฺโญ สุจิรสฺเสว เยน ภควา เตนุปสงฺกมิ, อุปสงฺกมิตฺวา ภควโต ปาเทสุ สิรสา นิปติตฺวา ภควโต ปาทานิ มุเขน จ ปริจุมฺพติ, ปาณีหิ จ ปริสมฺพาหติ, นามญฺจ สาเวติ ‘โกณฺฑญฺโญหํ ภควา, โกณฺฑญฺโญหํ สุคตา’ติ.\csromanspacer{} ยํนูนาหํ อายสฺมนฺตํ อญฺญาสิโกณฺฑญฺญํ ภควโต สมฺมุขา สารุปฺปาหิ คาถาหิ อภิตฺถเวยฺยนฺ”ติ.}

\prose{อถ โข อายสฺมา วงฺคีโส อุฏฺฐายาสนา เอกํสํ อุตฺตราสงฺคํ กริตฺวา เยน ภควา เตนญฺชลิํ ปณาเมตฺวา ภควนฺตํ เอตทโวจ “ปฏิภาติ มํ ภควา, ปฏิภาติ มํ สุคตา”ติ. “ปฏิภาตุ ตํ วงฺคีสา”ติ ภควา อโวจ.\csromanspacer{} อถ โข อายสฺมา วงฺคีโส อายสฺมนฺตํ อญฺญาสิโกณฺฑญฺญํ ภควโต สมฺมุขา สารุปฺปาหิ คาถาหิ อภิตฺถวิ\csromandash{}}

\csromangathameasure{“พุทฺธานุพุทฺโธ โส เถโร, โกณฺฑญฺโญ ติพฺพนิกฺกโม. \\ ลาภี สุขวิหารานํ, วิเวกานํ อภิณฺหโส. \\ ยํ สาวเกน ปตฺตพฺพํ, สตฺถุสาสนการินา. \\ สพฺพสฺส ตํ อนุปฺปตฺตํ, อปฺปมตฺตสฺส สิกฺขโต. \\ มหานุภาโว เตวิชฺโช, เจโตปริยายโกวิโท. \\ โกณฺฑญฺโญ พุทฺธทายาโท, ปาเท วนฺทติ สตฺถุโน”ติ.}
\csromangathagroup{\csromangathastanza{“พุทฺธานุพุทฺโธ โส เถโร, โกณฺฑญฺโญ ติพฺพนิกฺกโม. \\ ลาภี สุขวิหารานํ, วิเวกานํ อภิณฺหโส.}\csromangathastanza{ยํ สาวเกน ปตฺตพฺพํ, สตฺถุสาสนการินา. \\ สพฺพสฺส ตํ อนุปฺปตฺตํ, อปฺปมตฺตสฺส สิกฺขโต.}\csromangathastanza{มหานุภาโว เตวิชฺโช, เจโตปริยายโกวิโท. \\ โกณฺฑญฺโญ พุทฺธทายาโท\footnote{พุทฺธสาวโก (อิ)}, ปาเท วนฺทติ สตฺถุโน”ติ.}}

\csromanheader{2}{10. โมคฺคลฺลานสุตฺต}
\proseitem{218}{\csromanfolio{196}เอกํ สมยํ ภควา ราชคเห วิหรติ อิสิคิลิปสฺเส กาฬสิลายํ มหตา ภิกฺขุสํเฆน สทฺธิํ ปญฺจมตฺเตหิ ภิกฺขุสเตหิ สพฺเพเหว อรหนฺเตหิ, เตสํ สุทํ อายสฺมา มหาโมคฺคลฺลาโน เจตสา จิตฺตํ สมนฺเนสติ\footnote{สมเนฺวสติ (สฺยา-ฏฺฐ)} วิปฺปมุตฺตํ นิรุปธิํ.\csromanspacer{} อถ โข อายสฺมโต วงฺคีสสฺส เอตทโหสิ “อยํ โข ภควา ราชคเห วิหรติ อิสิคิลิปสฺเส กาฬสิลายํ มหตา ภิกฺขุสํเฆน สทฺธิํ ปญฺจมตฺเตหิ ภิกฺขุสเตหิ สพฺเพเหว อรหนฺเตหิ, เตสํ สุทํ อายสฺมา มหาโมคฺคลฺลาโน เจตสา จิตฺตํ สมนฺเนสติ วิปฺปมุตฺตํ นิรุปธิํ.\csromanspacer{} ยํนูนาหํ อายสฺมนฺตํ มหาโมคฺคลฺลานํ ภควโต สมฺมุขา สารุปฺปาหิ คาถาหิ อภิตฺถเวยฺยนฺ”ติ.}

\prose{อถ โข อายสฺมา วงฺคีโส อุฏฺฐายาสนา เอกํสํ อุตฺตราสงฺคํ กริตฺวา เยน ภควา เตนญฺชลิํ ปณาเมตฺวา ภควนฺตํ เอตทโวจ “ปฏิภาติ มํ ภควา, ปฏิภาติ มํ สุคตา”ติ. “ปฏิภาตุ ตํ วงฺคีสา”ติ ภควา อโวจ.\csromanspacer{} อถ โข อายสฺมา วงฺคีโส อายสฺมนฺตํ มหาโมคฺคลฺลานํ ภควโต สมฺมุขา สารุปฺปาหิ คาถาหิ อภิตฺถวิ\csromandash{}}

\csromangathameasure{“นคสฺส ปสฺเส อาสีนํ, มุนิํ ทุกฺขสฺส ปารคุํ. \\ สาวกา ปยิรุปาสนฺติ, เตวิชฺชา มจฺจุหายิโน. \\ เต เจตสา อนุปริเยติ, โมคฺคลฺลาโน มหิทฺธิโก. \\ จิตฺตํ เนสํ สมนฺเนสํ, วิปฺปมุตฺตํ นิรูปธิํ. \\ เอวํ สพฺพงฺคสมฺปนฺนํ, มุนิํ ทุกฺขสฺส ปารคุํ. \\ อเนกาการสมฺปนฺนํ, ปยิรุปาสนฺติ โคตมนฺ”ติ.}
\csromangathagroup{\csromangathastanza{“นคสฺส ปสฺเส อาสีนํ, มุนิํ ทุกฺขสฺส ปารคุํ. \\ สาวกา ปยิรุปาสนฺติ, เตวิชฺชา มจฺจุหายิโน.}\csromangathastanza{เต เจตสา อนุปริเยติ\footnote{อนุปริเยสติ (สี, สฺยา, กํ)}, โมคฺคลฺลาโน มหิทฺธิโก. \\ จิตฺตํ เนสํ สมนฺเนสํ\footnote{อนุปริเยสติ (สี, สฺยา, กํ)}, วิปฺปมุตฺตํ นิรูปธิํ.}\csromangathastanza{เอวํ สพฺพงฺคสมฺปนฺนํ, มุนิํ ทุกฺขสฺส ปารคุํ. \\ อเนกาการสมฺปนฺนํ, ปยิรุปาสนฺติ โคตมนฺ”ติ.}}

\csromanheader{2}{11. คคฺคราสุตฺต}
\proseitem{219}{เอกํ สมยํ ภควา จมฺปายํ วิหรติ คคฺคราย โปกฺขรณิยา ตีเร มหตา ภิกฺขุสํเฆน สทฺธิํ ปญฺจมตฺเตหิ ภิกฺขุสเตหิ สตฺตหิ จ อุปาสกสเตหิ สตฺตหิ จ อุปาสิกาสเตหิ อเนเกหิ จ เทวตาสหสฺเสหิ, ตฺยาสฺสุทํ ภควา อติโรจติ\footnote{อติวิโรจติ (ก)} วณฺเณน เจว ยสสา จ.\csromanspacer{} อถ โข อายสฺมโต วงฺคีสสฺส เอตทโหสิ “อยํ}

\vspace{\tipitakaparskip}
\csromancenter{วงฺคีสสํยุตฺต โข ภควา จมฺปายํ วิหรติ คคฺคราย โปกฺขรณิยา ตีเร มหตา ภิกฺขุสํเฆน สทฺธิํ ปญฺจมตฺเตหิ ภิกฺขุสเตหิ สตฺตหิ จ อุปาสกสเตหิ สตฺตหิ จ อุปาสิกาสเตหิ อเนเกหิ จ เทวตาสหสฺเสหิ, ตฺยาสฺสุทํ ภควา อติโรจติ วณฺเณน เจว ยสสา จ, ยํนูนาหํ ภควนฺตํ สมฺมุขา สารุปฺปาย คาถาย อภิตฺถเวยฺยนฺ”ติ.}
\prose{\csromanfolio{197}อถ โข อายสฺมา วงฺคีโส อุฏฺฐายาสนา เอกํสํ อุตฺตราสงฺคํ กริตฺวา เยน ภควา เตนญฺชลิํ ปณาเมตฺวา ภควนฺตํ เอตทโวจ “ปฏิภาติ มํ ภควา, ปฏิภาติ มํ สุคตา”ติ. “ปฏิภาตุ ตํ วงฺคีสา”ติ ภควา อโวจ.\csromanspacer{} อถ โข อายสฺมา วงฺคีโส ภควนฺตํ สมฺมุขา สารุปฺปาย คาถาย อภิตฺถวิ\csromandash{}}

\csromangathameasure{“จนฺโท ยถา วิคตวลาหเก นเภ, \\ วิโรจติ วิคตมโลว ภาณุมา. \\ เอวมฺปิ องฺคีรส ตฺวํ มหามุนิ, \\ อติโรจสิ ยสสา สพฺพโลกนฺ”ติ.}
\csromangathagroup{\csromangathastanza{“จนฺโท ยถา วิคตวลาหเก นเภ, \\ วิโรจติ วิคตมโลว ภาณุมา. \\ เอวมฺปิ องฺคีรส ตฺวํ มหามุนิ, \\ อติโรจสิ ยสสา สพฺพโลกนฺ”ติ.}}

\csromanheader{2}{12. วงฺคีสสุตฺต}
\proseitem{220}{เอกํ สมยํ อายสฺมา วงฺคีโส สาวตฺถิยํ วิหรติ เชตวเน อนาถปิณฺฑิกสฺส อาราเม.\csromanspacer{} เตน โข ปน สมเยน อายสฺมา วงฺคีโส อจิร- อรหตฺตปฺปตฺโต หุตฺวา\footnote{โหติ (สี, สฺยา, กํ)} วิมุตฺติสุขํ ปฏิสํเวที\footnote{วิมุตฺติสุขปฏิสํเวที (สี, อิ)} ตายํ เวลายํ อิมา คาถาโย อภาสิ\csromandash{}}

\csromangathameasure{“กาเวยฺยมตฺตา วิจริมฺห ปุพฺเพ, คามา คามํ ปุรา ปุรํ. \\ อถทฺทสาม สมฺพุทฺธํ, สทฺธา โน อุปปชฺชถ. \\ โส เม ธมฺมมเทเสสิ, ขนฺธายตนธาตุโย. \\ ตสฺสาหํ ธมฺมํ สุตฺวาน, ปพฺพชิํ อนคาริยํ. \\ พหุนฺนํ วต อตฺถาย, โพธิํ อชฺฌคมา มุนิ. \\ ภิกฺขูนํ ภิกฺขุนีนญฺจ, เย นิยามคตทฺทสา. \\ สฺวาคตํ วต เม อาสิ, มม พุทฺธสฺส สนฺติเก. \\ ติสฺโส วิชฺชา อนุปฺปตฺตา, กตํ พุทฺธสฺส สาสนํ. \\ ปุพฺเพนิวาสํ ชานามิ, ทิพฺพจกฺขุํ วิโสธิตํ. \\ เตวิชฺโช อิทฺธิปตฺโตมฺหิ, เจโตปริยายโกวิโท”ติ.}
\csromangathagroup{\csromangathastanza{“กาเวยฺยมตฺตา วิจริมฺห ปุพฺเพ, คามา คามํ ปุรา ปุรํ. \\ อถทฺทสาม สมฺพุทฺธํ, สทฺธา โน อุปปชฺชถ.}\csromangathastanza{โส เม ธมฺมมเทเสสิ, ขนฺธายตนธาตุโย\footnote{ขนฺเธ อายตนานิ ธาตุโย (สฺยา, กํ, อิ, ก)}. \\ ตสฺสาหํ ธมฺมํ สุตฺวาน, ปพฺพชิํ อนคาริยํ.}\csromangathastanza{พหุนฺนํ วต อตฺถาย, โพธิํ อชฺฌคมา มุนิ. \\ ภิกฺขูนํ ภิกฺขุนีนญฺจ, เย นิยามคตทฺทสา.}\csromangathastanza{\csromanfolio{198}สฺวาคตํ วต เม อาสิ, มม พุทฺธสฺส สนฺติเก. \\ ติสฺโส วิชฺชา อนุปฺปตฺตา, กตํ พุทฺธสฺส สาสนํ.}\csromangathastanza{ปุพฺเพนิวาสํ ชานามิ, ทิพฺพจกฺขุํ วิโสธิตํ. \\ เตวิชฺโช อิทฺธิปตฺโตมฺหิ, เจโตปริยายโกวิโท”ติ.}}

\vspace{\tipitakaparskip}
\csromancenter{\textbf{วงฺคีสสํยุตฺตํ สมตฺตํ.}}
\titlehead{\textbf{ตสฺสุทฺทานํ}}
\csromangathameasure{นิกฺขนฺตํ อรติ เจว, เปสลา อติมญฺญนา. \\ อานนฺเทน สุภาสิตา, สาริปุตฺตปวารณา. \\ ปโรสหสฺสํ โกณฺฑญฺโญ, โมคฺคลฺลาเนน คคฺครา. \\ วงฺคีเสน ทฺวาทสาติ.}
\csromangathagroup{\csromangathastanza{นิกฺขนฺตํ อรติ เจว, เปสลา อติมญฺญนา. \\ อานนฺเทน สุภาสิตา, สาริปุตฺตปวารณา.}\csromangathastanza{ปโรสหสฺสํ โกณฺฑญฺโญ, โมคฺคลฺลาเนน คคฺครา. \\ วงฺคีเสน ทฺวาทสาติ.}}

\csromanheader{2}{9. วนสํยุตฺต}
\csromanheader{2}{1. วิเวกสุตฺต}
\proseitem{221}{\csromanfolio{199}เอวํ เม สุตํ\csromandash{}เอกํ สมยํ อญฺญตโร ภิกฺขุ โกสเลสุ วิหรติ อญฺญตรสฺมิํ วนสณฺเฑ.\csromanspacer{} เตน โข ปน สมเยน โส ภิกฺขุ ทิวาวิหารคโต ปาปเก อกุสเล วิตกฺเก วิตกฺเกติ เคหนิสฺสิเต.\csromanspacer{} อถ โข ยา ตสฺมิํ วนสณฺเฑ อธิวตฺถา เทวตา ตสฺส ภิกฺขุโน อนุกมฺปิกา อตฺถกามา ตํ ภิกฺขุํ สํเวเชตุกามา เยน โส ภิกฺขุ เตนุปสงฺกมิ, อุปสงฺกมิตฺวา ตํ ภิกฺขุํ คาถาหิ อชฺฌภาสิ\csromandash{}}

\csromangathameasure{วิเวกกาโมสิ วนํ ปวิฏฺโฐ, \\ อถ เต มโน นิจฺฉรตี พหิทฺธา. \\ ชโน ชนสฺมิํ วินยสฺสุ ฉนฺทํ, \\ ตโต สุขี โหหิสิ วีตราโค. \\ อรติํ ปชหาสิ สโต, ภวาสิ สตํ ตํ สารยามเส. \\ ปาตาลรโช หิ ทุตฺตโร, มา ตํ กามรโช อวาหริ. \\ สกุโณ ยถา ปํสุกุนฺถิโต, วิธุนํ ปาตยติ สิตํ รชํ. \\ เอวํ ภิกฺขุ ปธานวา สติมา, วิธุนํ ปาตยติ สิตํ รชนฺติ.}
\csromangathagroup{\csromangathastanza{วิเวกกาโมสิ วนํ ปวิฏฺโฐ, \\ อถ เต มโน นิจฺฉรตี พหิทฺธา. \\ ชโน ชนสฺมิํ วินยสฺสุ ฉนฺทํ, \\ ตโต สุขี โหหิสิ วีตราโค. \\ อรติํ ปชหาสิ สโต, ภวาสิ สตํ ตํ สารยามเส. \\ ปาตาลรโช หิ ทุตฺตโร, มา ตํ กามรโช อวาหริ.}\csromangathastanza{สกุโณ ยถา ปํสุกุนฺถิโต\footnote{ปํสุกุณฺฐิโต (ก), ปํสุกุณฺฑิโต (สี, สฺยา, กํ, อิ)}, วิธุนํ ปาตยติ สิตํ รชํ. \\ เอวํ ภิกฺขุ ปธานวา สติมา, วิธุนํ ปาตยติ สิตํ รชนฺติ.}}

\prose{อถ โข โส ภิกฺขุ ตาย เทวตาย สํเวชิโต สํเวคมาปาทีติ.}

\csromanheader{2}{2. อุปฏฺฐานสุตฺต}
\proseitem{222}{เอกํ สมยํ อญฺญตโร ภิกฺขุ โกสเลสุ วิหรติ อญฺญตรสฺมิํ วนสณฺเฑ.\csromanspacer{} เตน โข ปน สมเยน โส ภิกฺขุ ทิวาวิหารคโต สุปติ.\csromanspacer{} อถ โข ยา ตสฺมิํ วนสณฺเฑ อธิวตฺถา เทวตา ตสฺส ภิกฺขุโน อนุกมฺปิกา อตฺถกามา ตํ ภิกฺขุํ สํเวเชตุกามา เยน โส ภิกฺขุ เตนุปสงฺกมิ, อุปสงฺกมิตฺวา ตํ ภิกฺขุํ คาถาหิ อชฺฌภาสิ\csromandash{}}

\csromangathameasure{อุฏฺเฐหิ ภิกฺขุ กิํ เสสิ, โก อตฺโถ สุปิเตน เต. \\ อาตุรสฺส หิ กา นิทฺทา, สลฺลวิทฺธสฺส รุปฺปโต. \\ ยาย สทฺธาย ปพฺพชิโต, อคารสฺมานคาริยํ. \\ ตเมว สทฺธํ พฺรูเหหิ, มา นิทฺทาย วสํ คมีติ. \\ อนิจฺจา อทฺธุวา กามา, เยสุ มนฺโทว มุจฺฉิโต. \\ พทฺเธสุ มุตฺตํ อสิตํ, กสฺมา ปพฺพชิตํ ตเป. \\ ฉนฺทราคสฺส วินยา, อวิชฺชาสมติกฺกมา. \\ ตํ ญาณํ ปรโมทานํ, กสฺมา ปพฺพชิตํ ตเป. \\ เฉตฺวา อวิชฺชํ วิชฺชาย, อาสวานํ ปริกฺขยา. \\ อโสกํ อนุปายาสํ, กสฺมา ปพฺพชิตํ ตเป. \\ อารทฺธวีริยํ ปหิตตฺตํ, นิจฺจํ ทฬฺหปรกฺกมํ. \\ นิพฺพานํ อภิกงฺขนฺตํ, กสฺมา ปพฺพชิตํ ตเป”ติ.}
\csromangathagroup{\csromangathastanza{อุฏฺเฐหิ ภิกฺขุ กิํ เสสิ, โก อตฺโถ สุปิเตน\footnote{สุปิเนน (สี)} เต. \\ อาตุรสฺส หิ กา นิทฺทา, สลฺลวิทฺธสฺส รุปฺปโต.}\csromangathastanza{\csromanfolio{200}ยาย สทฺธาย ปพฺพชิโต\footnote{ยาย สทฺธาปพฺพชิโต (สี, สฺยา, กํ)}, อคารสฺมานคาริยํ. \\ ตเมว สทฺธํ พฺรูเหหิ, มา นิทฺทาย วสํ คมีติ.}\csromangathastanza{อนิจฺจา อทฺธุวา กามา, เยสุ มนฺโทว มุจฺฉิโต. \\ พทฺเธสุ\footnote{ขนฺเธสุ (สี)} มุตฺตํ อสิตํ, กสฺมา ปพฺพชิตํ ตเป.}\csromangathastanza{ฉนฺทราคสฺส วินยา, อวิชฺชาสมติกฺกมา. \\ ตํ ญาณํ ปรโมทานํ\footnote{ปริโยทาตํ (สี, อิ), ปรโมทาตํ (สฺยา, กํ), ปรมโวทานํ (สี-ฏฺฐ)}, กสฺมา ปพฺพชิตํ ตเป.}\csromangathastanza{เฉตฺวา\footnote{เภตฺวา (สี, สฺยา, กํ, อิ)} อวิชฺชํ วิชฺชาย, อาสวานํ ปริกฺขยา. \\ อโสกํ อนุปายาสํ, กสฺมา ปพฺพชิตํ ตเป.}\csromangathastanza{อารทฺธวีริยํ ปหิตตฺตํ, นิจฺจํ ทฬฺหปรกฺกมํ. \\ นิพฺพานํ อภิกงฺขนฺตํ, กสฺมา ปพฺพชิตํ ตเป”ติ.}}

\csromanheader{2}{3. กสฺสปโคตฺตสุตฺต}
\proseitem{223}{เอกํ สมยํ อายสฺมา กสฺสปโคตฺโต โกสเลสุ วิหรติ อญฺญตรสฺมิํ วนสณฺเฑ.\csromanspacer{} เตน โข ปน สมเยน อายสฺมา กสฺสปโคตฺโต ทิวาวิหารคโต อญฺญตรํ เฉตํ โอวทติ.\csromanspacer{} อถ โข ยา ตสฺมิํ วนสณฺเฑ อธิวตฺถา เทวตา อายสฺมนฺตํ กสฺสปโคตฺตํ สํเวเชตุกามา เยนายสฺมา กสฺสปโคตฺโต เตนุปสงฺกมิ, อุปสงฺกมิตฺวา อายสฺมนฺตํ กสฺสปโคตฺตํ คาถาหิ อชฺฌภาสิ\csromandash{}}

\csromangathameasure{“คิริทุคฺคจรํ เฉตํ, อปฺปปญฺญํ อเจตสํ. \\ อกาเล โอวทํ ภิกฺขุ, มนฺโทว ปฏิภาติ มํ. \\ สุณาติ น วิชานาติ, อาโลเกติ น ปสฺสติ. \\ ธมฺมสฺมิํ ภญฺญมานสฺมิํ, อตฺถํ พาโล น พุชฺฌติ. \\ สเจปิ ทส ปชฺโชเต, ธารยิสฺสสิ กสฺสป. \\ เนว ทกฺขติ รูปานิ, จกฺขุ หิสฺส น วิชฺชตี”ติ.}
\csromangathagroup{\csromangathastanza{“คิริทุคฺคจรํ เฉตํ, อปฺปปญฺญํ อเจตสํ. \\ อกาเล โอวทํ ภิกฺขุ, มนฺโทว ปฏิภาติ มํ.}\csromangathastanza{สุณาติ น วิชานาติ, อาโลเกติ น ปสฺสติ. \\ ธมฺมสฺมิํ ภญฺญมานสฺมิํ, อตฺถํ พาโล น พุชฺฌติ.}\csromangathastanza{สเจปิ ทส ปชฺโชเต, ธารยิสฺสสิ กสฺสป. \\ เนว ทกฺขติ รูปานิ, จกฺขุ หิสฺส น วิชฺชตี”ติ.}}

\prose{อถ โข อายสฺมา กสฺสปโคตฺโต ตาย เทวตาย สํเวชิโต สํเวคมาปาทีติ.}

\csromanheader{2}{9. วนสํยุตฺต \textbf{4. สมฺพหุลสุตฺต}}
\proseitem{224}{\csromanfolio{201}เอกํ สมยํ สมฺพหุลา ภิกฺขู โกสเลสุ วิหรนฺติ อญฺญตรสฺมิํ วนสณฺเฑ.\csromanspacer{} อถ โข เต ภิกฺขู วสฺสํวุฏฺฐา\footnote{วสฺสํวุตฺถา (สี, สฺยา)} เตมาสจฺจเยน จาริกํ ปกฺกมิํสุ.\csromanspacer{} อถ โข ยา ตสฺมิํ วนสณฺเฑ อธิวตฺถา เทวตา เต ภิกฺขู อปสฺสนฺตี ปริเทวมานา ตายํ เวลายํ อิมํ คาถํ อภาสิ\csromandash{}}

\csromangathameasure{“อรติ วิย เมชฺช ขายติ, \\ พหุเก ทิสฺวาน วิวิตฺเต อาสเน. \\ เต จิตฺตกถา พหุสฺสุตา, \\ โกเม โคตมสาวกา คตา”ติ.}
\csromangathagroup{\csromangathastanza{“อรติ วิย เมชฺช ขายติ, \\ พหุเก ทิสฺวาน วิวิตฺเต อาสเน. \\ เต จิตฺตกถา พหุสฺสุตา, \\ โกเม โคตมสาวกา คตา”ติ.}}

\csromanheader{2}{เอวํ วุตฺเต อญฺญตรา เทวตา ตํ เทวตํ คาถาย ปจฺจภาสิ\csromandash{}}
\csromangathameasure{“มาคธํ คตา โกสลํ คตา, เอกจฺจิยา ปน วชฺชิภูมิยา. \\ มคา วิย อสงฺคจาริโน, อนิเกตา วิหรนฺติ ภิกฺขโว”ติ.}
\csromangathagroup{\csromangathastanza{“มาคธํ คตา โกสลํ คตา, เอกจฺจิยา ปน วชฺชิภูมิยา. \\ มคา วิย อสงฺคจาริโน, อนิเกตา วิหรนฺติ ภิกฺขโว”ติ.}}

\csromanheader{2}{5. อานนฺทสุตฺต}
\proseitem{225}{เอกํ สมยํ อายสฺมา อานนฺโท โกสเลสุ วิหรติ อญฺญตรสฺมิํ วินสณฺเฑ.\csromanspacer{} เตน โข ปน สมเยน อายสฺมา อานนฺโท อติเวลํ คิหิสญฺญตฺติพหุโล วิหรติ.\csromanspacer{} อถ โข ยา ตสฺมิํ วนสณฺเฑ อธิวตฺถา เทวตา อายสฺมโต อานนฺทสฺส อนุกมฺปิกา อตฺถกามา อายสฺมนฺตํ อานนฺทํ สํเวเชตุกามา เยนายสฺมา อานนฺโท เตนุปสงฺกมิ, อุปสงฺกมิตฺวา อายสฺมนฺตํ อานนฺทํ คาถาย อชฺฌภาสิ\csromandash{}}

\csromangathameasure{“รุกฺขมูลคหนํ ปสกฺกิย, นิพฺพานํ หทยสฺมิํ โอปิย. \\ ฌาย โคตม มา ปมาโท, กิํ เต พิฬิพิฬิกา กริสฺสตี”ติ.}
\csromangathagroup{\csromangathastanza{“รุกฺขมูลคหนํ ปสกฺกิย, นิพฺพานํ หทยสฺมิํ โอปิย. \\ ฌาย โคตม มา ปมาโท\footnote{มา จ ปมาโท (สี, อิ)}, กิํ เต พิฬิพิฬิกา กริสฺสตี”ติ.}}

\prose{อถ โข อายสฺมา อานนฺโท ตาย เทวตาย สํเวชิโต สํเวคมาปาทีติ.}

\csromanheader{2}{6. อนุรุทฺธสุตฺต}
\proseitem{226}{เอกํ สมยํ อายสฺมา อนุรุทฺโธ โกสเลสุ วิหรติ อญฺญตรสฺมิํ วนสณฺเฑ.\csromanspacer{} อถ โข อญฺญตรา ตาวติํสกายิกา เทวตา \csromanfolio{202}ชาลินี นาม อายสฺมโต อนุรุทฺธสฺส ปุราณทุติยิกา เยนายสฺมา อนุรุทฺโธ เตนุปสงฺกมิ, อุปสงฺกมิตฺวา อายสฺมนฺตํ อนุรุทฺธํ คาถาย อชฺฌภาสิ\csromandash{}}

\csromangathameasure{“ตตฺถ จิตฺตํ ปณิเธหิ, ยตฺถ เต วุสิตํ ปุเร. \\ ตาวติํเสสุ เทเวสุ, สพฺพกามสมิทฺธิสุ. \\ ปุรกฺขโต ปริวุโต, เทวกญฺญาหิ โสภสี”ติ. \\ ทุคฺคตา เทวกญฺญาโย, สกฺกายสฺมิํ ปติฏฺฐิตา. \\ เต จาปิ ทุคฺคตา สตฺตา, เทวกญฺญาหิ ปตฺถิตาติ. \\ น เต สุขํ ปชานนฺติ, เย น ปสฺสนฺติ นนฺทนํ. \\ อาวาสํ นรเทวานํ, ติทสานํ ยสสฺสินนฺติ. \\ น ตฺวํ พาเล วิชานาสิ, ยถา อรหตํ วโจ. \\ อนิจฺจา สพฺพสงฺขารา, อุปฺปาทวยธมฺมิโน. \\ อุปฺปชฺชิตฺวา นิรุชฺฌนฺติ, เตสํ วูปสโม สุโข. \\ นตฺถิ ทานิ ปุนาวาโส, เทวกายสฺมิ ชาลินิ. \\ วิกฺขีโณ ชาติสํสาโร, นตฺถิ ทานิ ปุนพฺภโวติ.}
\csromangathagroup{\csromangathastanza{“ตตฺถ จิตฺตํ ปณิเธหิ, ยตฺถ เต วุสิตํ ปุเร. \\ ตาวติํเสสุ เทเวสุ, สพฺพกามสมิทฺธิสุ.}\csromangathastanza{ปุรกฺขโต ปริวุโต, เทวกญฺญาหิ โสภสี”ติ. \\ ทุคฺคตา เทวกญฺญาโย, สกฺกายสฺมิํ ปติฏฺฐิตา.}\csromangathastanza{เต จาปิ ทุคฺคตา สตฺตา, เทวกญฺญาหิ ปตฺถิตาติ. \\ น เต สุขํ ปชานนฺติ, เย น ปสฺสนฺติ นนฺทนํ.}\csromangathastanza{อาวาสํ นรเทวานํ, ติทสานํ ยสสฺสินนฺติ. \\ น ตฺวํ พาเล วิชานาสิ, ยถา อรหตํ วโจ.}\csromangathastanza{อนิจฺจา สพฺพสงฺขารา, อุปฺปาทวยธมฺมิโน. \\ อุปฺปชฺชิตฺวา นิรุชฺฌนฺติ, เตสํ วูปสโม สุโข.}\csromangathastanza{นตฺถิ ทานิ ปุนาวาโส, เทวกายสฺมิ ชาลินิ. \\ วิกฺขีโณ ชาติสํสาโร, นตฺถิ ทานิ ปุนพฺภโวติ.}}

\csromanheader{2}{7. นาคทตฺตสุตฺต}
\proseitem{227}{เอกํ สมยํ อายสฺมา นาคทตฺโต โกสเลสุ วิหรติ อญฺญตรสฺมิํ วนสณฺเฑ.\csromanspacer{} เตน โข ปน สมเยน อายสฺมา นาคทตฺโต อติกาเลน คามํ ปวิสติ, อติทิวา ปฏิกฺกมติ.\csromanspacer{} อถ โข ยา ตสฺมิํ วนสณฺเฑ อธิวตฺถา เทวตา อายสฺมโต นาคทตฺตสฺส อนุกมฺปิกา อตฺถกามา อายสฺมนฺตํ นาคทตฺตํ สํเวเชตุกามา เยนายสฺมา นาคทตฺโต เตนุปสงฺกมิ, อุปสงฺกมิตฺวา อายสฺมนฺตํ นาคทตฺตํ คาถาหิ อชฺฌภาสิ\csromandash{}}

\csromangathameasure{“กาเล ปวิส นาคทตฺต, ทิวา จ อาคนฺตฺวา อติเวลจารี. \\ สํสฏฺโฐ คหฏฺเฐหิ, สมานสุขทุกฺโข. \\ ภายามิ นาคทตฺตํ สุปฺปคพฺภํ, กุเลสุ วินิพทฺธํ. \\ มา เหว มจฺจุรญฺโญ พลวโต, อนฺตกสฺส วสํ อุเปสี”ติ.}
\csromangathagroup{\csromangathastanza{“กาเล ปวิส นาคทตฺต, ทิวา จ อาคนฺตฺวา อติเวลจารี. \\ สํสฏฺโฐ คหฏฺเฐหิ, สมานสุขทุกฺโข.}\csromangathastanza{ภายามิ นาคทตฺตํ สุปฺปคพฺภํ, กุเลสุ วินิพทฺธํ. \\ มา เหว มจฺจุรญฺโญ พลวโต, อนฺตกสฺส วสํ อุเปสี”ติ\footnote{วสเมยฺยาติ (สี, อิ), วสเมสีติ (สฺยา, กํ)}.}}

\csromanheader{2}{9. วนสํยุตฺต}
\vspace{\tipitakaparskip}
\csromancenter{อถ โข อายสฺมา นาคทตฺโต ตาย เทวตาย สํเวชิโต สํเวคมาปาทีติ.}
\csromanheader{2}{8. กุลฆรณีสุตฺต}
\proseitem{228}{\csromanfolio{203}เอกํ สมยํ อญฺญตโร ภิกฺขุ โกสเลสุ วิหรติ อญฺญตรสฺมิํ วนสณฺเฑ.\csromanspacer{} เตน โข ปน สมเยน โส ภิกฺขุ อญฺญตรสฺมิํ กุเล อติเวลํ อชฺโฌคาฬฺหปฺปตฺโต วิหรติ.\csromanspacer{} อถ โข ยา ตสฺมิํ วนสณฺเฑ อธิวตฺถา เทวตา ตสฺส ภิกฺขุโน อนุกมฺปิกา อตฺถกามา ตํ ภิกฺขุํ สํเวเชตุกามา ยา ตสฺมิํ กุเล กุลฆรณี, ตสฺสา วณฺณํ อภินิมฺมินิตฺวา เยน โส ภิกฺขุ เตนุปสงฺกมิ, อุปสงฺกมิตฺวา ตํ ภิกฺขุํ คาถาย อชฺฌภาสิ\csromandash{}}

\prose{“นทีตีเรสุ สณฺฐาเน, สภาสุ รถิยาสุ จ.}

\prose{ชนา สงฺคมฺม มนฺเตนฺติ, มญฺจ ตญฺจ\footnote{ตฺวญฺจ (ก)} กิมนฺตรนฺ”ติ.}

\prose{พหูหิ สทฺทา ปจฺจูหา, ขมิตพฺพา ตปสฺสินา.}

\prose{น เตน มงฺกุ โหตพฺพํ, น หิ เตน กิลิสฺสติ.}

\prose{โย จ สทฺทปริตฺตาสี, วเน วาตมิโค ยถา.}

\prose{ลหุจิตฺโตติ ตํ อาหุ, นาสฺส สมฺปชฺชเต วตนฺติ.}

\csromanheader{2}{9. วชฺชิปุตฺตสุตฺต}
\proseitem{229}{เอกํ สมยํ อญฺญตโร วชฺชิปุตฺตโก ภิกฺขุ เวสาลิยํ วิหรติ อญฺญตรสฺมิํ วนสณฺเฑ.\csromanspacer{} เตน โข ปน สมเยน เวสาลิยํ วชฺชิปุตฺตโก สพฺพรตฺติจาโร โหติ.\csromanspacer{} อถ โข โส ภิกฺขุ เวสาลิยา ตูริย ตาฬิต วาทิต นิคฺโฆสสทฺทํ สุตฺวา ปริเทวมาโน ตายํ เวลายํ อิมํ คาถํ อภาสิ\csromandash{}}

\prose{“เอกกา มยํ อรญฺเญ วิหราม, อปวิทฺธํว\footnote{อปวิฏฺฐํว (สฺยา, กํ)} วนสฺมิํ ทารุกํ.}

\prose{เอตาทิสิกาย รตฺติยา, โก สุ นามเมฺหหิ\footnote{นาม อเมฺหหิ (สี, อิ)} ปาปิโย”ติ.}

\prose{อถ โข ยา ตสฺมิํ วนสณฺเฑ อธิวตฺถา เทวตา ตสฺส ภิกฺขุโน อนุกมฺปิกา อตฺถกามา ตํ ภิกฺขุํ สํเวเชตุกามา เยน โส ภิกฺขุ เตนุปสงฺกมิ, อุปสงฺกมิตฺวา ตํ ภิกฺขุํ คาถาย อชฺฌภาสิ\csromandash{}}

\csromangathameasure{“เอกโกว ตฺวํ อรญฺเญ วิหรสิ, \\ อปวิทฺธํว วนสฺมิํ ทารุกํ. \\ ตสฺส เต พหุกา ปิหยนฺติ, \\ เนรยิกา วิย สคฺคคามินนฺ”ติ.}
\csromangathagroup{\csromangathastanza{\csromanfolio{204}“เอกโกว ตฺวํ อรญฺเญ วิหรสิ, \\ อปวิทฺธํว วนสฺมิํ ทารุกํ. \\ ตสฺส เต พหุกา ปิหยนฺติ, \\ เนรยิกา วิย สคฺคคามินนฺ”ติ.}}

\prose{อถ โข โส ภิกฺขุ ตาย เทวตาย สํเวชิโต สํเวคมาปาทีติ. \textbf{10.}\csromanspacer{}\textbf{ สชฺฌายสุตฺต}}

\proseitem{230}{เอกํ สมยํ อญฺญตโร ภิกฺขุ โกสเลสุ วิหรติ อญฺญตรสฺมิํ วนสณฺเฑ.\csromanspacer{} เตน โข ปน สมเยน โส ภิกฺขุ ยํ สุทํ ปุพฺเพ อติเวลํ สชฺฌายพหุโล วิหรติ, โส อปเรน สมเยน อปฺโปสฺสุกฺโก ตุณฺหีภูโต สงฺกสายติ.\csromanspacer{} อถ โข ยา ตสฺมิํ วนสณฺเฑ อธิวตฺถา เทวตา ตสฺส ภิกฺขุโน ธมฺมํ อสุณนฺตี เยน โส ภิกฺขุ เตนุปสงฺกมิ, อุปสงฺกมิตฺวา ตํ ภิกฺขุํ คาถาย อชฺฌภาสิ\csromandash{}}

\csromangathameasure{“กสฺมา ตุวํ ธมฺมปทานิ ภิกฺขุ, \\ นาธียสิ ภิกฺขุหิ สํวสนฺโต. \\ สุตฺวาน ธมฺมํ ลภติปฺปสาทํ, \\ ทิฏฺเฐว ธมฺเม ลภติปฺปสํสนฺ”ติ. \\ อหุ ปุเร ธมฺมปเทสุ ฉนฺโท, \\ ยาว วิราเคน สมาคมิมฺห. \\ ยโต วิราเคน สมาคมิมฺห, \\ ยํ กิญฺจิ ทิฏฺฐํว สุตํ มุตํ วา.}
\csromangathagroup{\csromangathastanza{“กสฺมา ตุวํ ธมฺมปทานิ ภิกฺขุ, \\ นาธียสิ ภิกฺขุหิ สํวสนฺโต. \\ สุตฺวาน ธมฺมํ ลภติปฺปสาทํ, \\ ทิฏฺเฐว ธมฺเม ลภติปฺปสํสนฺ”ติ.}\csromangathastanza{อหุ ปุเร ธมฺมปเทสุ ฉนฺโท, \\ ยาว วิราเคน สมาคมิมฺห. \\ ยโต วิราเคน สมาคมิมฺห, \\ ยํ กิญฺจิ ทิฏฺฐํว สุตํ มุตํ วา.}}

\csromanheader{2}{อญฺญาย นิกฺเขปนมาหุ สนฺโตติ. \textbf{11. อกุสลวิตกฺกสุตฺต}}
\proseitem{231}{เอกํ สมยํ อญฺญตโร ภิกฺขุ โกสเลสุ วิหรติ อญฺญตรสฺมิํ วนสณฺเฑ.\csromanspacer{} เตน โข ปน สมเยน โส ภิกฺขุ ทิวาวิหารคโต ปาปเก อกุสเล วิตกฺเก วิตกฺเกติ.\csromanspacer{} เสยฺยถิทํ, กามวิตกฺกํ พฺยาปาทวิตกฺกํ วิหิํสาวิตกฺกํ.\csromanspacer{} อถ โข ยา ตสฺมิํ วนสณฺเฑ อธิวตฺถา เทวตา ตสฺส ภิกฺขุโน อนุกมฺปิกา อตฺถกามา ตํ ภิกฺขุํ สํเวเชตุกามา เยน โส ภิกฺขุ เตนุปสงฺกมิ, อุปสงฺกมิตฺวา ตํ ภิกฺขุํ คาถาหิ อชฺฌภาสิ\csromandash{}}

\csromangathameasure{“อโยนิโส มนสิการา, โส วิตกฺเกหิ ขชฺชสิ. \\ อโยนิโส ปฏินิสฺสชฺช, โยนิโส อนุจินฺตย.}
\csromangathagroup{\csromangathastanza{“อโยนิโส มนสิการา, โส วิตกฺเกหิ ขชฺชสิ. \\ อโยนิโส\footnote{อโยนิํ (อิ, ก)} ปฏินิสฺสชฺช, โยนิโส อนุจินฺตย.}}

\csromanheader{2}{9. วนสํยุตฺต}
\csromangathameasure{สตฺถารํ ธมฺมมารพฺภ, สํฆํ สีลานิ อตฺตโน. \\ อธิคจฺฉสิ ปาโมชฺชํ, ปีติสุขมสํสยํ. \\ ตโต ปาโมชฺชพหุโล, ทุกฺขสฺสนฺตํ กริสฺสสี”ติ.}
\csromangathagroup{\csromangathastanza{\csromanfolio{205}สตฺถารํ ธมฺมมารพฺภ, สํฆํ สีลานิ อตฺตโน. \\ อธิคจฺฉสิ ปาโมชฺชํ, ปีติสุขมสํสยํ. \\ ตโต ปาโมชฺชพหุโล, ทุกฺขสฺสนฺตํ กริสฺสสี”ติ.}}

\prose{อถ โข โส ภิกฺขุ ตาย เทวตาย สํเวชิโต สํเวคมาปาทีติ.}

\csromanheader{2}{12. มชฺฌนฺหิกสุตฺต}
\proseitem{232}{เอกํ สมยํ อญฺญตโร ภิกฺขุ โกสเลสุ วิหรติ อญฺญตรสฺมิํ วนสณฺเฑ.\csromanspacer{} อถ โข ตสฺมิํ วนสณฺเฑ อธิวตฺถา เทวตา เยน โส ภิกฺขุ เตนุปสงฺกมิ, อุปสงฺกมิตฺวา ตสฺส ภิกฺขุโน สนฺติเก อิมํ คาถํ อภาสิ\csromandash{}}

\csromangathameasure{“ฐิเต มชฺฌนฺหิเก กาเล, สนฺนิสีเวสุ ปกฺขิสุ. \\ สณเตว พฺรหารญฺญํ, ตํ ภยํ ปฏิภาติ มํ. \\ ฐิเต มชฺฌนฺหิเก กาเล, สนฺนิสีเวสุ ปกฺขิสุ. \\ สณเตว พฺรหารญฺญํ, สา รติ ปฏิภาติ มนฺ”ติ.}
\csromangathagroup{\csromangathastanza{“ฐิเต มชฺฌนฺหิเก กาเล, สนฺนิสีเวสุ\footnote{สนฺนิสินฺเนสุ (สฺยา, กํ, อิ)} ปกฺขิสุ. \\ สณเตว พฺรหารญฺญํ, ตํ ภยํ ปฏิภาติ มํ.}\csromangathastanza{ฐิเต มชฺฌนฺหิเก กาเล, สนฺนิสีเวสุ ปกฺขิสุ. \\ สณเตว พฺรหารญฺญํ, สา รติ ปฏิภาติ มนฺ”ติ.}}

\csromanheader{2}{13. ปากตินฺทฺริยสุตฺต}
\proseitem{233}{เอกํ สมยํ สมฺพหุลา ภิกฺขู โกสเลสุ วิหรนฺติ อญฺญตรสฺมิํ วนสณฺเฑ อุทฺธตา อุนฺนฬา จปลา มุขรา วิกิณฺณวาจา มุฏฺฐสฺสติโน อสมฺปชานา อสมาหิตา วิพฺภนฺตจิตฺตา ปากตินฺทฺริยา.\csromanspacer{} อถ โข ยา ตสฺมิํ วนสณฺเฑ อธิวตฺถา เทวตา เตสํ ภิกฺขูนํ อนุกมฺปิกา อตฺถกามา เต ภิกฺขู สํเวเชตุกามา เยน เต ภิกฺขู เตนุปสงฺกมิ, อุปสงฺกมิตฺวา เต ภิกฺขู คาถาหิ อชฺฌภาสิ\csromandash{}}

\csromangathameasure{“สุขชีวิโน ปุเร อาสุํ, ภิกฺขู โคตมสาวกา. \\ อนิจฺฉา ปิณฺฑเมสนา, อนิจฺฉา สยนาสนํ. \\ โลเก อนิจฺจตํ ญตฺวา, ทุกฺขสฺสนฺตํ อกํสุ เต. \\ ทุปฺโปสํ กตฺวา อตฺตานํ, คาเม คามณิกา วิย. \\ ภุตฺวา ภุตฺวา นิปชฺชนฺติ, ปราคาเรสุ มุจฺฉิตา. \\ สํฆสฺส อญฺชลิํ กตฺวา, อิเธกจฺเจ วทามหํ. \\ อปวิทฺธา อนาถา เต, ยถา เปตา ตเถว เต. \\ เย โข ปมตฺตา วิหรนฺติ, เต เม สนฺธาย ภาสิตํ. \\ เย อปฺปมตฺตา วิหรนฺติ, นโม เตสํ กโรมหนฺ”ติ.}
\csromangathagroup{\csromangathastanza{“สุขชีวิโน ปุเร อาสุํ, ภิกฺขู โคตมสาวกา. \\ อนิจฺฉา ปิณฺฑเมสนา, อนิจฺฉา สยนาสนํ.}\csromangathastanza{โลเก อนิจฺจตํ ญตฺวา, ทุกฺขสฺสนฺตํ อกํสุ เต. \\ ทุปฺโปสํ กตฺวา อตฺตานํ, คาเม คามณิกา วิย.}\csromangathastanza{ภุตฺวา ภุตฺวา นิปชฺชนฺติ, ปราคาเรสุ มุจฺฉิตา. \\ สํฆสฺส อญฺชลิํ กตฺวา, อิเธกจฺเจ วทามหํ.}\csromangathastanza{\csromanfolio{206}อปวิทฺธา\footnote{อปวิฏฺฐา (สฺยา, กํ)} อนาถา เต, ยถา เปตา ตเถว เต. \\ เย โข ปมตฺตา วิหรนฺติ, เต เม สนฺธาย ภาสิตํ. \\ เย อปฺปมตฺตา วิหรนฺติ, นโม เตสํ กโรมหนฺ”ติ.}}

\prose{อถ โข เต ภิกฺขู ตาย เทวตาย สํเวชิตา สํเวคมาปาทุนฺติ.}

\csromanheader{2}{14. คนฺธตฺเถนสุตฺต}
\proseitem{234}{เอกํ สมยํ อญฺญตโร ภิกฺขุ โกสเลสุ วิหรติ อญฺญตรสฺมิํ วนสณฺเฑ.\csromanspacer{} เตน โข ปน สมเยน โส ภิกฺขุ ปจฺฉาภตฺตํ ปิณฺฑปาตปฏิกฺกนฺโต โปกฺขรณิํ โอคาเหตฺวา ปทุมํ อุปสิงฺฆติ.\csromanspacer{} อถ โข ยา ตสฺมิํ วนสณฺเฑ อธิวตฺถา เทวตา ตสฺส ภิกฺขุโน อนุกมฺปิกา อตฺถกามา ตํ ภิกฺขุํ สํเวเชตุกามา เยน โส ภิกฺขุ เตนุปสงฺกมิ, อุปสงฺกมิตฺวา ตํ ภิกฺขุํ คาถาย อชฺฌภาสิ\csromandash{}}

\csromangathameasure{“ยเมตํ วาริชํ ปุปฺผํ, อทินฺนํ อุปสิงฺฆสิ. \\ เอกงฺคเมตํ เถยฺยานํ, คนฺธตฺเถโนสิ มาริสา”ติ. \\ น หรามิ น ภญฺชามิ, อารา สิงฺฆามิ วาริชํ. \\ อถ เกน นุ วณฺเณน, คนฺธตฺเถโนติ วุจฺจติ. \\ ยฺวายํ ภิสานิ ขนติ, ปุณฺฑรีกานิ ภญฺชติ. \\ เอวํ อากิณฺณกมฺมนฺโต, กสฺมา เอโส น วุจฺจตีติ. \\ อากิณฺณลุทฺโท ปุริโส, ธาติเจลํว มกฺขิโต. \\ ตสฺมิํ เม วจนํ นตฺถิ, ตฺวญฺจารหามิ วตฺตเว. \\ อนงฺคณสฺส โปสสฺส, นิจฺจํ สุจิคเวสิโน. \\ วาลคฺคมตฺตํ ปาปสฺส, อพฺภามตฺตํว ขายตีติ. \\ อทฺธา มํ ยกฺข ชานาสิ, อโถ เม อนุกมฺปสิ. \\ ปุนปิ ยกฺข วชฺชาสิ, ยทา ปสฺสสิ เอทิสนฺติ.}
\csromangathagroup{\csromangathastanza{“ยเมตํ วาริชํ ปุปฺผํ, อทินฺนํ อุปสิงฺฆสิ. \\ เอกงฺคเมตํ เถยฺยานํ, คนฺธตฺเถโนสิ มาริสา”ติ.}\csromangathastanza{น หรามิ น ภญฺชามิ, อารา สิงฺฆามิ วาริชํ. \\ อถ เกน นุ วณฺเณน, คนฺธตฺเถโนติ วุจฺจติ.}\csromangathastanza{ยฺวายํ ภิสานิ ขนติ, ปุณฺฑรีกานิ ภญฺชติ. \\ เอวํ อากิณฺณกมฺมนฺโต, กสฺมา เอโส น วุจฺจตีติ.}\csromangathastanza{อากิณฺณลุทฺโท ปุริโส, ธาติเจลํว มกฺขิโต. \\ ตสฺมิํ เม วจนํ นตฺถิ, ตฺวญฺจารหามิ วตฺตเว.}\csromangathastanza{อนงฺคณสฺส โปสสฺส, นิจฺจํ สุจิคเวสิโน. \\ วาลคฺคมตฺตํ ปาปสฺส, อพฺภามตฺตํว ขายตีติ.}\csromangathastanza{อทฺธา มํ ยกฺข ชานาสิ, อโถ เม อนุกมฺปสิ. \\ ปุนปิ ยกฺข วชฺชาสิ, ยทา ปสฺสสิ เอทิสนฺติ.}}

\csromanheader{2}{9. วนสํยุตฺต}
\csromangathameasure{เนว ตํ อุปชีวาม, นปิ เต ภตกามฺหเส. \\ ตฺวเมว ภิกฺขุ ชาเนยฺย, เยน คจฺเฉยฺย สุคฺคตินฺติ.}
\csromangathagroup{\csromangathastanza{\csromanfolio{207}เนว ตํ อุปชีวาม, นปิ เต ภตกามฺหเส. \\ ตฺวเมว ภิกฺขุ ชาเนยฺย, เยน คจฺเฉยฺย สุคฺคตินฺติ.}}

\prose{อถ โข โส ภิกฺขุ ตาย เทวตาย สํเวชิโต สํเวคมาปาทีติ.}

\vspace{\tipitakaparskip}
\csromancenter{\textbf{วนสํยุตฺตํ สมตฺตํ.}}
\titlehead{\textbf{ตสฺสุทฺทานํ}}
\csromangathameasure{วิเวกํ อุปฏฺฐานญฺจ, กสฺสปโคตฺเตน สมฺพหุลา. \\ อานนฺโท อนุรุทฺโธ จ, นาคทตฺตญฺจ กุลฆรณี. \\ วชฺชิปุตฺโต จ เวสาลี, สชฺฌาเยน อโยนิโส. \\ มชฺฌนฺหิกาลมฺหิ ปากตินฺทฺริย-ปทุมปุปฺเผน จุทฺทส ภเวติ.}
\csromangathagroup{\csromangathastanza{วิเวกํ อุปฏฺฐานญฺจ, กสฺสปโคตฺเตน สมฺพหุลา. \\ อานนฺโท อนุรุทฺโธ จ, นาคทตฺตญฺจ กุลฆรณี.}\csromangathastanza{วชฺชิปุตฺโต จ เวสาลี, สชฺฌาเยน อโยนิโส. \\ มชฺฌนฺหิกาลมฺหิ ปากตินฺทฺริย-ปทุมปุปฺเผน จุทฺทส ภเวติ.}}

\csromanheader{2}{10. ยกฺขสํยุตฺต}
\csromanheader{2}{1. อินฺทกสุตฺต}
\proseitem{235}{\csromanfolio{208}เอวํ เม สุตํ\csromandash{}เอกํ สมยํ ภควา ราชคเห วิหรติ อินฺทกูเฏ ปพฺพเต อินฺทกสฺส ยกฺขสฺส ภวเน.\csromanspacer{} อถ โข อินฺทโก ยกฺโข เยน ภควา เตนุปสงฺกมิ, อุปสงฺกมิตฺวา ภควนฺตํ คาถาย อชฺฌภาสิ\csromandash{}}

\csromangathameasure{“รูปํ น ชีวนฺติ วทนฺติ พุทฺธา, \\ กถํ นฺวยํ วินฺทติมํ สรีรํ. \\ กุตสฺส อฏฺฐียกปิณฺฑเมติ, \\ กถํ นฺวยํ สชฺชติ คพฺภรสฺมินฺ”ติ. \\ ปฐมํ กลลํ โหติ, กลลา โหติ อพฺพุทํ. \\ อพฺพุทา ชายเต เปสิ, เปสิ นิพฺพตฺตตี ฆโน. \\ ฆนา ปสาขา ชายนฺติ, เกสา โลมา นขาปิ จ. \\ ยญฺจสฺส ภุญฺชตี มาตา, อนฺนํ ปานญฺจ โภชนํ.}
\csromangathagroup{\csromangathastanza{“รูปํ น ชีวนฺติ วทนฺติ พุทฺธา, \\ กถํ นฺวยํ วินฺทติมํ สรีรํ. \\ กุตสฺส อฏฺฐียกปิณฺฑเมติ, \\ กถํ นฺวยํ สชฺชติ คพฺภรสฺมินฺ”ติ. \\ ปฐมํ กลลํ โหติ, กลลา โหติ อพฺพุทํ. \\ อพฺพุทา ชายเต เปสิ, เปสิ นิพฺพตฺตตี ฆโน.}\csromangathastanza{ฆนา ปสาขา ชายนฺติ, เกสา โลมา นขาปิ จ. \\ ยญฺจสฺส ภุญฺชตี มาตา, อนฺนํ ปานญฺจ โภชนํ.}}

\vspace{\tipitakaparskip}
\csromancenter{เตน โส ตตฺถ ยาเปติ, มาตุกุจฺฉิคโต นโรติ. \textbf{2.}\csromanspacer{}\textbf{ สกฺกนามสุตฺต}}
\proseitem{236}{เอกํ สมยํ ภควา ราชคเห วิหรติ คิชฺฌกูเฏ ปพฺพเต.\csromanspacer{} อถ โข สกฺกนามโก ยกฺโข เยน ภควา เตนุปสงฺกมิ, อุปสงฺกมิตฺวา ภควนฺตํ คาถาย อชฺฌภาสิ\csromandash{}}

\csromangathameasure{“สพฺพคนฺถปฺปหีนสฺส, วิปฺปมุตฺตสฺส เต สโต. \\ สมณสฺส น ตํ สาธุ, ยทญฺญมนุสาสสี”ติ. \\ เยน เกนจิ วณฺเณน, สํวาโส สกฺก ชายติ. \\ น ตํ อรหติ สปฺปญฺโญ, มนสา อนุกมฺปิตุํ. \\ มนสา เจ ปสนฺเนน, ยทญฺญมนุสาสติ.}
\csromangathagroup{\csromangathastanza{“สพฺพคนฺถปฺปหีนสฺส, วิปฺปมุตฺตสฺส เต สโต. \\ สมณสฺส น ตํ สาธุ, ยทญฺญมนุสาสสี”ติ\footnote{ยทญฺญมนุสาสตีติ (สี, สฺยา, กํ, อิ)}.}\csromangathastanza{เยน เกนจิ วณฺเณน, สํวาโส สกฺก ชายติ. \\ น ตํ อรหติ สปฺปญฺโญ, มนสา อนุกมฺปิตุํ. \\ มนสา เจ ปสนฺเนน, ยทญฺญมนุสาสติ.}}

\vspace{\tipitakaparskip}
\csromancenter{น เตน โหติ สํยุตฺโต, ยานุกมฺปา\footnote{สานุกมฺปา (สี, อิ)} อนุทฺทยาติ. \textbf{3.}\csromanspacer{}\textbf{ สูจิโลมสุตฺต}}
\proseitem{237}{เอกํ สมยํ ภควา คยายํ วิหรติ ฏงฺกิตมญฺเจ สูจิโลมสฺส ยกฺขสฺส ภวเน.\csromanspacer{} เตน โข ปน สมเยน ขโร จ ยกฺโข}

\csromanheader{2}{10. ยกฺขสํยุตฺต}
\prose{\csromanfolio{209}สูจิโลโม จ ยกฺโข ภควโต อวิทูเร อติกฺกมนฺติ.\csromanspacer{} อถ โข ขโร ยกฺโข สูจิโลมํ ยกฺขํ เอตทโวจ “เอโส สมโณ”ติ.\csromanspacer{} เนโส สมโณ, สมณโก เอโส, ยาว ชานามิ ยทิ วา โส สมโณ ยทิ วา ปน โส สมณโกติ.}

\prose{อถ โข สูจิโลโม ยกฺโข เยน ภควา เตนุปสงฺกมิ, อุปสงฺกมิตฺวา ภควโต กายํ อุปนาเมสิ.\csromanspacer{} อถ โข ภควา กายํ อปนาเมสิ.\csromanspacer{} อถ โข สูจิโลโม ยกฺโข ภควนฺตํ เอตทโวจ “ภายสิ มํ สมณา”ติ.\csromanspacer{} น ขฺวาหํ ตํ อาวุโส ภายามิ, อปิ จ เต สมฺผสฺโส ปาปโกติ.\csromanspacer{} ปญฺหํ ตํ สมณ ปุจฺฉิสฺสามิ, สเจ เม น พฺยากริสฺสสิ, จิตฺตํ วา เต ขิปิสฺสามิ, หทยํ วา เต ผาเลสฺสามิ, ปาเทสุ วา คเหตฺวา ปารคงฺคาย\footnote{ปารํ คงฺคาย (ก)} ขิปิสฺสามีติ.\csromanspacer{} น ขฺวาหํ ตํ อาวุโส ปสฺสามิ สเทวเก โลเก สมารเก สพฺรหฺมเก สสฺสมณพฺราหฺมณิยา ปชาย สเทวมนุสฺสาย, โย เม จิตฺตํ วา ขิเปยฺย หทยํ วา ผาเลยฺย ปาเทสุ วา คเหตฺวา ปารคงฺคาย ขิเปยฺย.\csromanspacer{} อปิ จ ตฺวํ อาวุโส ปุจฺฉ ยทา กงฺขสีติ. ( )\footnote{(อถ โข สูจิโลโม ยกฺโข ภควนฺตํ คาถาย อชฺฌภาสิ.) (สี)}}

\csromangathameasure{ราโค จ โทโส จ กุโตนิทานา, \\ อรตี รตี โลมหํโส กุโตชา. \\ กุโต สมุฏฺฐาย มโนวิตกฺกา, \\ กุมารกา ธงฺกมิโวสฺสชนฺตีติ. \\ ราโค จ โทโส จ อิโตนิทานา, \\ อรตี รตี โลมหํโส อิโตชา. \\ อิโต สมุฏฺฐาย มโนวิตกฺกา, \\ กุมารกา ธงฺกมิโวสฺสชนฺติ. \\ สฺเนหชา อตฺตสมฺภูตา, \\ นิโคฺรธสฺเสว ขนฺธชา. \\ ปุถู วิสตฺตา กาเมสุ, \\ มาลุวาว วิตตา วเน. \\ เย นํ ปชานนฺติ ยโตนิทานํ, \\ เต นํ วิโนเทนฺติ สุโณหิ ยกฺข. \\ เต ทุตฺตรํ โอฆมิมํ ตรนฺติ, \\ อติณฺณปุพฺพํ อปุนพฺภวายาติ.}
\csromangathagroup{\csromangathastanza{ราโค จ โทโส จ กุโตนิทานา, \\ อรตี รตี โลมหํโส กุโตชา. \\ กุโต สมุฏฺฐาย มโนวิตกฺกา, \\ กุมารกา ธงฺกมิโวสฺสชนฺตีติ.}\csromangathastanza{ราโค จ โทโส จ อิโตนิทานา, \\ อรตี รตี โลมหํโส อิโตชา. \\ อิโต สมุฏฺฐาย มโนวิตกฺกา, \\ กุมารกา ธงฺกมิโวสฺสชนฺติ.}\csromangathastanza{สฺเนหชา อตฺตสมฺภูตา, \\ นิโคฺรธสฺเสว ขนฺธชา. \\ ปุถู วิสตฺตา กาเมสุ, \\ มาลุวาว วิตตา\footnote{วิตฺถตา (สฺยา, กํ)} วเน.}\csromangathastanza{\csromanfolio{210}เย นํ ปชานนฺติ ยโตนิทานํ, \\ เต นํ วิโนเทนฺติ สุโณหิ ยกฺข. \\ เต ทุตฺตรํ โอฆมิมํ ตรนฺติ, \\ อติณฺณปุพฺพํ อปุนพฺภวายาติ.}}

\csromanheader{2}{4. มณิภทฺทสุตฺต}
\proseitem{238}{เอกํ สมยํ ภควา มคเธสุ วิหรติ มณิมาลิเก เจติเย มณิภทฺทสฺส ยกฺขสฺส ภวเน.\csromanspacer{} อถ โข มณิภทฺโท ยกฺโข เยน ภควา เตนุปสงฺกมิ, อุปสงฺกมิตฺวา ภควโต สนฺติเก อิมํ คาถํ อภาสิ\csromandash{}}

\csromangathameasure{“สตีมโต สทา ภทฺทํ, สติมา สุขเมธติ. \\ สตีมโต สุเว เสยฺโย, เวรา จ ปริมุจฺจตี”ติ. \\ สตีมโต สทา ภทฺทํ, สติมา สุขเมธติ. \\ สตีมโต สุเว เสยฺโย, เวรา น ปริมุจฺจติ. \\ ยสฺส สพฺพมโหรตฺตํ, อหิํสาย รโต มโน. \\ เมตฺตํโส สพฺพภูเตสุ, เวรํ ตสฺส น เกนจีติ.}
\csromangathagroup{\csromangathastanza{“สตีมโต สทา ภทฺทํ, สติมา สุขเมธติ. \\ สตีมโต สุเว เสยฺโย, เวรา จ ปริมุจฺจตี”ติ.}\csromangathastanza{สตีมโต สทา ภทฺทํ, สติมา สุขเมธติ. \\ สตีมโต สุเว เสยฺโย, เวรา น ปริมุจฺจติ.}\csromangathastanza{ยสฺส สพฺพมโหรตฺตํ\footnote{รตฺติํ (สฺยา, กํ, ก)}, อหิํสาย รโต มโน. \\ เมตฺตํโส สพฺพภูเตสุ, เวรํ ตสฺส น เกนจีติ.}}

\csromanheader{2}{5. สานุสุตฺต}
\proseitem{239}{เอกํ สมยํ ภควา สาวตฺถิยํ วิหรติ เชตวเน อนาถปิณฺฑิกสฺส อาราเม.\csromanspacer{} เตน โข ปน สมเยน อญฺญตริสฺสา อุปาสิกาย สานุ นาม ปุตฺโต ยกฺเขน คหิโต โหติ.\csromanspacer{} อถ โข สา อุปาสิกา ปริเทวมานา ตายํ เวลายํ อิมา คาถาโย อภาสิ\csromandash{}}

\csromangathameasure{“จาตุทฺทสิํ ปญฺจทสิํ, ยา จ ปกฺขสฺส อฏฺฐมี. \\ ปาฏิหาริยปกฺขญฺจ, อฏฺฐงฺคสุสมาคตํ. \\ อุโปสถํ อุปวสนฺติ, พฺรหฺมจริยํ จรนฺติ เย. \\ น เตหิ ยกฺขา กีฬนฺติ, อิติ เม อรหตํ สุตํ. \\ สา ทานิ อชฺช ปสฺสามิ, ยกฺขา กีฬนฺติ สานุนา”ติ.}
\csromangathagroup{\csromangathastanza{“จาตุทฺทสิํ ปญฺจทสิํ, ยา จ ปกฺขสฺส อฏฺฐมี. \\ ปาฏิหาริยปกฺขญฺจ, อฏฺฐงฺคสุสมาคตํ.}\csromangathastanza{อุโปสถํ อุปวสนฺติ, พฺรหฺมจริยํ จรนฺติ เย. \\ น เตหิ ยกฺขา กีฬนฺติ, อิติ เม อรหตํ สุตํ. \\ สา ทานิ อชฺช ปสฺสามิ, ยกฺขา กีฬนฺติ สานุนา”ติ.}}

\csromanheader{2}{10. ยกฺขสํยุตฺต}
\csromangathameasure{จาตุทฺทสิํ ปญฺจทสิํ, ยา จ ปกฺขสฺส อฏฺฐมี. \\ ปาฏิหาริยปกฺขญฺจ, อฏฺฐงฺคสุสมาคตํ. \\ อุโปสถํ อุปวสนฺติ, พฺรหฺมจริยํ จรนฺติ เย. \\ น เตหิ ยกฺขา กีฬนฺติ, สาหุ เต อรหตํ สุตํ. \\ สานุํ ปพุทฺธํ วชฺชาสิ, ยกฺขานํ วจนํ อิทํ. \\ มากาสิ ปาปกํ กมฺมํ, อาวิ วา ยทิ วา รโห. \\ สเจ จ ปาปกํ กมฺมํ, กริสฺสสิ กโรสิ วา. \\ น เต ทุกฺขา ปมุตฺยตฺถิ, อุปฺปจฺจาปิ ปลายโตติ. \\ มตํ วา อมฺม โรทนฺติ, โย วา ชีวํ น ทิสฺสติ. \\ ชีวนฺตํ อมฺม ปสฺสนฺตี, กสฺมา มํ อมฺม โรทสีติ. \\ มตํ วา ปุตฺต โรทนฺติ, โย วา ชีวํ น ทิสฺสติ. \\ โย จ กาเม จชิตฺวาน, ปุนราคจฺฉเต อิธ. \\ ตํ วาปิ ปุตฺต โรทนฺติ, ปุน ชีวํ มโต หิ โส. \\ กุกฺกุฬา อุพฺภโต ตาต, กุกฺกุฬํ ปติตุมิจฺฉสิ. \\ นรกา อุพฺภโต ตาต, นรกํ ปติตุมิจฺฉสิ. \\ อภิธาวถ ภทฺทนฺเต, กสฺส อุชฺฌาปยามเส. \\ อาทิตฺตา นีหตํ ภณฺฑํ, ปุน ฑยฺหิตุมิจฺฉสีติ.}
\csromangathagroup{\csromangathastanza{\csromanfolio{211}จาตุทฺทสิํ ปญฺจทสิํ, ยา จ ปกฺขสฺส อฏฺฐมี. \\ ปาฏิหาริยปกฺขญฺจ, อฏฺฐงฺคสุสมาคตํ.}\csromangathastanza{อุโปสถํ อุปวสนฺติ, พฺรหฺมจริยํ จรนฺติ เย. \\ น เตหิ ยกฺขา กีฬนฺติ, สาหุ เต อรหตํ สุตํ.}\csromangathastanza{สานุํ ปพุทฺธํ วชฺชาสิ, ยกฺขานํ วจนํ อิทํ. \\ มากาสิ ปาปกํ กมฺมํ, อาวิ วา ยทิ วา รโห.}\csromangathastanza{สเจ จ\footnote{สเจว (สฺยา, กํ, อิ, ก), ยญฺเจว (สี)} ปาปกํ กมฺมํ, กริสฺสสิ กโรสิ วา. \\ น เต ทุกฺขา ปมุตฺยตฺถิ, อุปฺปจฺจาปิ ปลายโตติ.}\csromangathastanza{มตํ วา อมฺม โรทนฺติ, โย วา ชีวํ น ทิสฺสติ. \\ ชีวนฺตํ อมฺม ปสฺสนฺตี, กสฺมา มํ อมฺม โรทสีติ.}\csromangathastanza{มตํ วา ปุตฺต โรทนฺติ, โย วา ชีวํ น ทิสฺสติ. \\ โย จ กาเม จชิตฺวาน, ปุนราคจฺฉเต อิธ.}\csromangathastanza{ตํ วาปิ ปุตฺต โรทนฺติ, ปุน ชีวํ มโต หิ โส. \\ กุกฺกุฬา อุพฺภโต ตาต, กุกฺกุฬํ\footnote{กุกฺกุเฬ (สี)} ปติตุมิจฺฉสิ.}\csromangathastanza{นรกา อุพฺภโต ตาต, นรกํ ปติตุมิจฺฉสิ. \\ อภิธาวถ ภทฺทนฺเต, กสฺส อุชฺฌาปยามเส. \\ อาทิตฺตา นีหตํ\footnote{นิพฺภตํ (สฺยา, กํ, ก), นิภตํ (อิ, ก)} ภณฺฑํ, ปุน ฑยฺหิตุมิจฺฉสีติ.}}

\csromanheader{2}{6. ปิยงฺกรสุตฺต}
\proseitem{240}{เอกํ สมยํ อายสฺมา อนุรุทฺโธ สาวตฺถิยํ วิหรติ เชตวเน อนาถปิณฺฑิกสฺส อาราเม.\csromanspacer{} เตน โข ปน สมเยน อายสฺมา อนุรุทฺโธ รตฺติยา ปจฺจูสสมยํ ปจฺจุฏฺฐาย ธมฺมปทานิ ภาสติ.\csromanspacer{} อถ โข ปิยงฺกรมาตา ยกฺขินี ปุตฺตกํ เอวํ โตเสสิ\csromandash{}}

\csromangathameasure{“มา สทฺทํ กริ ปิยงฺกร, ภิกฺขุ ธมฺมปทานิ ภาสติ. \\ อปิจ ธมฺมปทํ วิชานิย, ปฏิปชฺเชม หิตาย โน สิยา. \\ ปาเณสุ จ สํยมามเส, สมฺปชานมุสา น ภณามเส. \\ สิกฺเขม สุสีลฺยมตฺตโน, อปิ มุจฺเจม ปิสาจโยนิยา”ติ.}
\csromangathagroup{\csromangathastanza{“มา สทฺทํ กริ ปิยงฺกร, ภิกฺขุ ธมฺมปทานิ ภาสติ. \\ อปิจ\footnote{อปิ (สี)} ธมฺมปทํ วิชานิย, ปฏิปชฺเชม หิตาย โน สิยา.}\csromangathastanza{ปาเณสุ จ สํยมามเส, สมฺปชานมุสา น ภณามเส. \\ สิกฺเขม สุสีลฺยมตฺตโน\footnote{สุสีลมตฺตโน (สี, ก)}, อปิ มุจฺเจม ปิสาจโยนิยา”ติ.}}

\csromanheader{2}{7. ปุนพฺพสุสุตฺต}
\proseitem{241}{\csromanfolio{212}เอกํ สมยํ ภควา สาวตฺถิยํ วิหรติ เชตวเน อนาถปิณฺฑิกสฺส อาราเม.\csromanspacer{} เตน โข ปน สมเยน ภควา ภิกฺขู นิพฺพานปฏิสํยุตฺตาย ธมฺมิยา กถาย สนฺทสฺเสติ สมาทเปติ สมุตฺเตเชติ สมฺปหํเสติ, เต จ ภิกฺขู อฏฺฐิํ กตฺวา มนสิ กตฺวา สพฺพเจตสา สมนฺนาหริตฺวา โอหิตโสตา ธมฺมํ สุณนฺติ.\csromanspacer{} อถ โข ปุนพฺพสุมาตา ยกฺขินี ปุตฺตเก เอวํ โตเสสิ\csromandash{}}

\csromangathameasure{“ตุณฺหี อุตฺตริเก โหหิ, ตุณฺหี โหหิ ปุนพฺพสุ. \\ ยาวาหํ พุทฺธเสฏฺฐสฺส, ธมฺมํ โสสฺสามิ สตฺถุโน. \\ นิพฺพานํ ภควา อาห, สพฺพคนฺถปฺปโมจนํ. \\ อติเวลา จ เม โหติ, อสฺมิํ ธมฺเม ปิยายนา. \\ ปิโย โลเก สโก ปุตฺโต, \\ ปิโย โลเก สโก ปติ. \\ ตโต ปิยตรา มยฺหํ, \\ อสฺส ธมฺมสฺส มคฺคนา. \\ น หิ ปุตฺโต ปติ วาปิ, ปิโย ทุกฺขา ปโมจเย. \\ ยถา สทฺธมฺมสฺสวนํ, ทุกฺขา โมเจติ ปาณินํ. \\ โลเก ทุกฺขปเรตสฺมิํ, ชรามรณสํยุเต. \\ ชรามรณโมกฺขาย, ยํ ธมฺมํ อภิสมฺพุธํ. \\ ตํ ธมฺมํ โสตุมิจฺฉามิ, ตุณฺหี โหหิ ปุนพฺพสู”ติ. \\ อมฺมา น พฺยาหริสฺสามิ, ตุณฺหีภูตายมุตฺตรา. \\ ธมฺมเมว นิสาเมหิ, สทฺธมฺมสฺสวนํ สุขํ. \\ สทฺธมฺมสฺส อนญฺญาย, อมฺมา ทุกฺขํ จรามเส. \\ เอส เทวมนุสฺสานํ, สมฺมูฬฺหานํ ปภงฺกโร. \\ พุทฺโธ อนฺติมสารีโร, ธมฺมํ เทเสติ จกฺขุมาติ. \\ สาธุ โข ปณฺฑิโต นาม, ปุตฺโต ชาโต อุเรสโย. \\ ปุตฺโต เม พุทฺธเสฏฺฐสฺส, ธมฺมํ สุทฺธํ ปิยายติ. \\ ปุนพฺพสุ สุขี โหหิ, อชฺชาหมฺหิ สมุคฺคตา. \\ ทิฏฺฐานิ อริยสจฺจานิ, อุตฺตราปิ สุณาตุ เมติ.}
\csromangathagroup{\csromangathastanza{“ตุณฺหี อุตฺตริเก โหหิ, ตุณฺหี โหหิ ปุนพฺพสุ. \\ ยาวาหํ พุทฺธเสฏฺฐสฺส, ธมฺมํ โสสฺสามิ สตฺถุโน.}\csromangathastanza{นิพฺพานํ ภควา อาห, สพฺพคนฺถปฺปโมจนํ. \\ อติเวลา จ เม โหติ, อสฺมิํ ธมฺเม ปิยายนา.}\csromangathastanza{ปิโย โลเก สโก ปุตฺโต, \\ ปิโย โลเก สโก ปติ. \\ ตโต ปิยตรา มยฺหํ, \\ อสฺส ธมฺมสฺส มคฺคนา. \\ น หิ ปุตฺโต ปติ วาปิ, ปิโย ทุกฺขา ปโมจเย. \\ ยถา สทฺธมฺมสฺสวนํ, ทุกฺขา โมเจติ ปาณินํ.}\csromangathastanza{โลเก ทุกฺขปเรตสฺมิํ, ชรามรณสํยุเต. \\ ชรามรณโมกฺขาย, ยํ ธมฺมํ อภิสมฺพุธํ.}\csromangathastanza{ตํ ธมฺมํ โสตุมิจฺฉามิ, ตุณฺหี โหหิ ปุนพฺพสู”ติ. \\ อมฺมา น พฺยาหริสฺสามิ, ตุณฺหีภูตายมุตฺตรา.}\csromangathastanza{ธมฺมเมว นิสาเมหิ, สทฺธมฺมสฺสวนํ สุขํ. \\ สทฺธมฺมสฺส อนญฺญาย, อมฺมา ทุกฺขํ จรามเส.}\csromangathastanza{เอส เทวมนุสฺสานํ, สมฺมูฬฺหานํ ปภงฺกโร. \\ พุทฺโธ อนฺติมสารีโร, ธมฺมํ เทเสติ จกฺขุมาติ.}\csromangathastanza{สาธุ โข ปณฺฑิโต นาม, ปุตฺโต ชาโต อุเรสโย. \\ ปุตฺโต เม พุทฺธเสฏฺฐสฺส, ธมฺมํ สุทฺธํ ปิยายติ.}\csromangathastanza{ปุนพฺพสุ สุขี โหหิ, อชฺชาหมฺหิ สมุคฺคตา. \\ ทิฏฺฐานิ อริยสจฺจานิ, อุตฺตราปิ สุณาตุ เมติ.}}

\csromanheader{2}{10. ยกฺขสํยุตฺต}
\csromanheader{2}{8. สุทตฺตสุตฺต}
\proseitem{242}{\csromanfolio{213}เอกํ สมยํ ภควา ราชคเห วิหรติ สีตวเน.\csromanspacer{} เตน โข ปน สมเยน อนาถปิณฺฑิโก คหปติ ราชคหํ อนุปฺปตฺโต โหติ เกนจิเทว กรณีเยน, อสฺโสสิ โข อนาถปิณฺฑิโก คหปติ “พุทฺโธ กิร โลเก อุปฺปนฺโน”ติ.\csromanspacer{} ตาวเทว จ ปน ภควนฺตํ ทสฺสนาย อุปสงฺกมิตุกาโม โหติ, อถสฺส อนาถปิณฺฑิกสฺส คหปติสฺส เอตทโหสิ “อกาโล โข อชฺช ภควนฺตํ ทสฺสนาย อุปสงฺกมิตุํ, เสฺว ทานาหํ กาเลน ภควนฺตํ ทสฺสนาย คมิสฺสามี”ติ พุทฺธคตาย สติยา นิปชฺชิ, รตฺติยา สุทํ ติกฺขตฺตุํ วุฏฺฐาสิ ปภาตนฺติ มญฺญมาโน.\csromanspacer{} อถ โข อนาถปิณฺฑิโก คหปติ เยน สิวถิกทฺวารํ\footnote{สีวถิกทฺวารํ (สี, สฺยา, กํ, อิ)} เตนุปสงฺกมิ, อมนุสฺสา ทฺวารํ วิวริํสุ.\csromanspacer{} อถ โข อนาถปิณฺฑิกสฺส คหปติสฺส นครมฺหา นิกฺขมนฺตสฺส อาโลโก อนฺตรธายิ, อนฺธกาโร ปาตุรโหสิ, ภยํ ฉมฺภิตตฺตํ โลมหํโส อุทปาทิ, ตโตว ปุน นิวตฺติตุกาโม อโหสิ.\csromanspacer{} อถ โข สิวโก\footnote{สีวโก (สี, อิ)} ยกฺโข อนฺตรหิโต สทฺทมนุสฺสาเวสิ\csromandash{}}

\csromangathameasure{“สตํ หตฺถี สตํ อสฺสา, สตํ อสฺสตรีรถา. \\ สตํ กญฺญาสหสฺสานิ, อามุกฺกมณิกุณฺฑลา. \\ เอกสฺส ปทวีติหารสฺส, \\ กลํ นาคฺฆนฺติ โสฬสิํ. \\ อภิกฺกม คหปติ, อภิกฺกม คหปติ, \\ อภิกฺกมนํ เต เสยฺโย โน ปฏิกฺกมนนฺ”ติ.}
\csromangathagroup{\csromangathastanza{“สตํ หตฺถี สตํ อสฺสา, สตํ อสฺสตรีรถา. \\ สตํ กญฺญาสหสฺสานิ, อามุกฺกมณิกุณฺฑลา.}\csromangathastanza{เอกสฺส ปทวีติหารสฺส, \\ กลํ นาคฺฆนฺติ โสฬสิํ. \\ อภิกฺกม คหปติ, อภิกฺกม คหปติ, \\ อภิกฺกมนํ เต เสยฺโย โน ปฏิกฺกมนนฺ”ติ.}}

\prose{อถ โข อนาถปิณฺฑิกสฺส คหปติสฺส อนฺธกาโร อนฺตรธายิ, อาโลโก ปาตุรโหสิ, ยํ อโหสิ ภยํ ฉมฺภิตตฺตํ โลมหํโส, โส ปฏิปฺปสฺสมฺภิ.\csromanspacer{} ทุติยมฺปิ โข อนาถปิณฺฑิกสฺส คหปติสฺส อาโลโก อนฺตรธายิ, อนฺธกาโร ปาตุรโหสิ, ภยํ ฉมฺภิตตฺตํ โลมหํโส อุทปาทิ, ตโตว ปุน นิวตฺติตุกาโม อโหสิ.\csromanspacer{} ทุติยมฺปิ โข สิวโก ยกฺโข อนฺตรหิโต สทฺทมนุสฺสาเวสิ\csromandash{}}

\prose{\csromanfolio{214}“สตํ หตฺถี สตํ อสฺสา ฯเปฯ กลํ นาคฺฆนฺติ โสฬสิํ.}

\csromangathameasure{อภิกฺกม คหปติ, อภิกฺกม คหปติ, \\ อภิกฺกมนํ เต เสยฺโย โน ปฏิกฺกมนนฺ”ติ.}
\csromangathagroup{\csromangathastanza{อภิกฺกม คหปติ, อภิกฺกม คหปติ, \\ อภิกฺกมนํ เต เสยฺโย โน ปฏิกฺกมนนฺ”ติ.}}

\prose{อถ โข อนาถปิณฺฑิกสฺส คหปติสฺส อนฺธกาโร อนฺตรธายิ, อาโลโก ปาตุรโหสิ, ยํ อโหสิ ภยํ ฉมฺภิตตฺตํ โลมหํโส, โส ปฏิปฺปสฺสมฺภิ.\csromanspacer{} ตติยมฺปิ โข อนาถปิณฺฑิกสฺส คหปติสฺส อาโลโก อนฺตรธายิ, อนฺธกาโร ปาตุรโหสิ, ภยํ ฉมฺภิตตฺตํ โลมหํโส อุทปาทิ, ตโตว ปุน นิวตฺติตุกาโม อโหสิ.\csromanspacer{} ตติยมฺปิ โข สิวโก ยกฺโข อนฺตรหิโต สทฺทมนุสฺสาเวสิ\csromandash{}}

\prose{“สตํ หตฺถี สตํ อสฺสา ฯเปฯ กลํ นาคฺฆนฺติ โสฬสิํ.}

\csromangathameasure{อภิกฺกม คหปติ, อภิกฺกม คหปติ, \\ อภิกฺกมนํ เต เสยฺโย โน ปฏิกฺกมนนฺ”ติ.}
\csromangathagroup{\csromangathastanza{อภิกฺกม คหปติ, อภิกฺกม คหปติ, \\ อภิกฺกมนํ เต เสยฺโย โน ปฏิกฺกมนนฺ”ติ.}}

\prose{อถ โข อนาถปิณฺฑิกสฺส คหปติสฺส อนฺธกาโร อนฺตรธายิ, อาโลโก ปาตุรโหสิ, ยํ อโหสิ ภยํ ฉมฺภิตตฺตํ โลมหํโส, โส ปฏิปฺปสฺสมฺภิ.\csromanspacer{} อถ โข อนาถปิณฺฑิโก คหปติ เยน สีตวนํ เยน ภควา เตนุปสงฺกมิ.}

\prose{เตน โข ปน สมเยน ภควา รตฺติยา ปจฺจูสสมยํ ปจฺจุฏฺฐาย อพฺโภกาเส จงฺกมติ.\csromanspacer{} อทฺทสา โข ภควา อนาถปิณฺฑิกํ คหปติํ ทูรโตว อาคจฺฉนฺตํ, ทิสฺวาน จงฺกมา โอโรหิตฺวา ปญฺญตฺเต อาสเน นิสีทิ, นิสชฺช โข ภควา อนาถปิณฺฑิกํ คหปติํ เอตทโวจ “เอหิ สุทตฺตา”ติ.\csromanspacer{} อถ โข อนาถปิณฺฑิโก คหปติ นาเมน มํ ภควา อาลปตีติ หฏฺโฐ อุทคฺโค ตตฺเถว ภควโต ปาเทสุ สิรสา นิปติตฺวา ภควนฺตํ เอตทโวจ “กจฺจิ ภนฺเต ภควา สุขมสยิตฺถา”ติ.}

\csromangathameasure{สพฺพทา เว สุขํ เสติ, พฺราหฺมโณ ปรินิพฺพุโต. \\ โย น ลิมฺปติ กาเมสุ, สีติภูโต นิรูปธิ. \\ สพฺพา อาสตฺติโย เฉตฺวา, วิเนยฺย หทเย ทรํ. \\ อุปสนฺโต สุขํ เสติ, สนฺติํ ปปฺปุยฺย เจตสาติ.}
\csromangathagroup{\csromangathastanza{สพฺพทา เว สุขํ เสติ, พฺราหฺมโณ ปรินิพฺพุโต. \\ โย น ลิมฺปติ กาเมสุ, สีติภูโต นิรูปธิ.}\csromangathastanza{สพฺพา อาสตฺติโย เฉตฺวา, วิเนยฺย หทเย ทรํ. \\ อุปสนฺโต สุขํ เสติ, สนฺติํ ปปฺปุยฺย เจตสาติ\footnote{เจตโสติ (สี)}.}}

\csromanheader{2}{10. ยกฺขสํยุตฺต}
\csromanheader{2}{9. ปฐมสุกฺกาสุตฺต}
\proseitem{243}{\csromanfolio{215}เอกํ สมยํ ภควา ราชคเห วิหรติ เวฬุวเน กลนฺทกนิวาเป.\csromanspacer{} เตน โข ปน สมเยน สุกฺกา ภิกฺขุนี มหติยา ปริสาย ปริวุตา ธมฺมํ เทเสติ.\csromanspacer{} อถ โข สุกฺกาย ภิกฺขุนิยา อภิปฺปสนฺโน ยกฺโข ราชคเห รถิกาย รถิกํ สิงฺฆาฏเกน สิงฺฆาฏกํ อุปสงฺกมิตฺวา ตายํ เวลายํ อิมา คาถาโย อภาสิ\csromandash{}}

\csromangathameasure{“กิํ เม กตา ราชคเห มนุสฺสา, มธุปีตาว เสยเร. \\ เย สุกฺกํ น ปยิรุปาสนฺติ, เทเสนฺติํ อมตํ ปทํ. \\ ตญฺจ ปน อปฺปฏิวานียํ, อเสจนกโมชวํ. \\ ปิวนฺติ มญฺเญ สปฺปญฺญา, วลาหกมิว ปนฺถคู”ติ.}
\csromangathagroup{\csromangathastanza{“กิํ เม กตา ราชคเห มนุสฺสา, มธุปีตาว เสยเร. \\ เย สุกฺกํ น ปยิรุปาสนฺติ, เทเสนฺติํ อมตํ ปทํ.}\csromangathastanza{ตญฺจ ปน อปฺปฏิวานียํ, อเสจนกโมชวํ. \\ ปิวนฺติ มญฺเญ สปฺปญฺญา, วลาหกมิว ปนฺถคู”ติ\footnote{วลาหกมิวทฺธคูติ (สี)}.}}

\csromanheader{2}{10. ทุติยสุกฺกาสุตฺต}
\proseitem{244}{เอกํ สมยํ ภควา ราชคเห วิหรติ เวฬุวเน กลนฺทกนิวาเป.\csromanspacer{} เตน โข ปน สมเยน อญฺญตโร อุปาสโก สุกฺกาย ภิกฺขุนิยา โภชนํ อทาสิ.\csromanspacer{} อถ โข สุกฺกาย ภิกฺขุนิยา อภิปฺปสนฺโน ยกฺโข ราชคเห รถิกาย รถิกํ สิงฺฆาฏเกน สิงฺฆาฏกํ อุปสงฺกมิตฺวา ตายํ เวลายํ อิมํ คาถํ อภาสิ\csromandash{}}

\csromangathameasure{“ปุญฺญํ วต ปสวิ พหุํ, สปฺปญฺโญ วตายํ อุปาสโก. \\ โย สุกฺกาย อทาสิ โภชนํ, สพฺพคนฺเถหิ วิปฺปมุตฺติยา”ติ.}
\csromangathagroup{\csromangathastanza{“ปุญฺญํ วต ปสวิ พหุํ, สปฺปญฺโญ วตายํ อุปาสโก. \\ โย สุกฺกาย อทาสิ โภชนํ, สพฺพคนฺเถหิ วิปฺปมุตฺติยา”ติ\footnote{วิปฺปมุตฺตายาติ (สฺยา, กํ)}.}}

\csromanheader{2}{11. จีราสุตฺต}
\proseitem{245}{เอวํ เม สุตํ\csromandash{}เอกํ สมยํ ภควา ราชคเห วิหรติ เวฬุวเน กลนฺทกนิวาเป.\csromanspacer{} เตน โข ปน สมเยน อญฺญตโร อุปาสโก จีราย\footnote{จิราย (ก)} ภิกฺขุนิยา จีวรํ อทาสิ.\csromanspacer{} อถ โข จีราย ภิกฺขุนิยา อภิปฺปสนฺโน ยกฺโข ราชคเห รถิกาย รถิกํ สิงฺฆาฏเกน สิงฺฆาฏกํ อุปสงฺกมิตฺวา ตายํ เวลายํ อิมํ คาถํ อภาสิ\csromandash{}}

\csromangathameasure{“ปุญฺญํ วต ปสวิ พหุํ, สปฺปญฺโญ วตายํ อุปาสโก. \\ โย จีราย อทาสิ จีวรํ, สพฺพโยเคหิ วิปฺปมุตฺติยา”ติ.}
\csromangathagroup{\csromangathastanza{\csromanfolio{216}“ปุญฺญํ วต ปสวิ พหุํ, สปฺปญฺโญ วตายํ อุปาสโก. \\ โย จีราย อทาสิ จีวรํ, สพฺพโยเคหิ วิปฺปมุตฺติยา”ติ\footnote{วิปฺปมุตฺตายาติ (สฺยา, กํ)}.}}

\csromanheader{2}{12. อาฬวกสุตฺต}
\proseitem{246}{เอวํ เม สุตํ\csromandash{}เอกํ สมยํ ภควา อาฬวิยํ วิหรติ อาฬวกสฺส ยกฺขสฺส ภวเน.\csromanspacer{} อถ โข อาฬวโก ยกฺโข เยน ภควา เตนุปสงฺกมิ, อุปสงฺกมิตฺวา ภควนฺตํ เอตทโวจ “นิกฺขม สมณา”ติ. “สาธาวุโส”ติ ภควา นิกฺขมิ.\csromanspacer{} ปวิส สมณาติ. “สาธาวุโส”ติ ภควา ปาวิสิ.\csromanspacer{} ทุติยมฺปิ โข อาฬวโก ยกฺโข ภควนฺตํ เอตทโวจ “นิกฺขม สมณา”ติ. “สาธาวุโส”ติ ภควา นิกฺขมิ.\csromanspacer{} ปวิส สมณาติ. “สาธาวุโส”ติ ภควา ปาวิสิ.\csromanspacer{} ตติยมฺปิ โข อาฬวโก ยกฺโข ภควนฺตํ เอตทโวจ “นิกฺขม สมณา”ติ. “สาธาวุโส”ติ ภควา นิกฺขมิ.\csromanspacer{} ปวิส สมณาติ. “สาธาวุโส”ติ ภควา ปาวิสิ.\csromanspacer{} จตุตฺถมฺปิ โข อาฬวโก ยกฺโข ภควนฺตํ เอตทโวจ “นิกฺขม สมณา”ติ.\csromanspacer{} น ขฺวาหํ ตํ อาวุโส นิกฺขมิสฺสามิ, ยํ เต กรณียํ, ตํ กโรหีติ.\csromanspacer{} ปญฺหํ ตํ สมณ ปุจฺฉิสฺสามิ, สเจ เม น พฺยากริสฺสสิ, จิตฺตํ วา เต ขิปิสฺสามิ, หทยํ วา เต ผาเลสฺสามิ, ปาเทสุ วา คเหตฺวา ปารคงฺคาย ขิปิสฺสามีติ.\csromanspacer{} น ขฺวาหํ ตํ อาวุโส ปสฺสามิ สเทวเก โลเก สมารเก สพฺรหฺมเก สสฺสมณพฺราหฺมณิยา ปชาย สเทวมนุสฺสาย, โย เม จิตฺตํ วา ขิเปยฺย หทยํ วา ผาเลยฺย ปาเทสุ วา คเหตฺวา ปารคงฺคาย ขิเปยฺย.\csromanspacer{} อปิ จ ตฺวํ อาวุโส ปุจฺฉ ยทา กงฺขสีติ. ( )\footnote{(อถ โข อาฬวโก ยกฺโข ภควนฺตํ คาถาย อชฺฌภาสิ.) (สี)}}

\prose{กิํสูธ วิตฺตํ ปุริสสฺส เสฏฺฐํ, กิํสุ สุจิณฺณํ สุขมาวหาติ.}

\prose{กิํสุ หเว สาทุตรํ รสานํ, กถํชีวิํ ชีวิตมาหุ เสฏฺฐนฺติ.}

\prose{สทฺธีธ วิตฺตํ ปุริสสฺส เสฏฺฐํ, ธมฺโม สุจิณฺโณ สุขมาวหาติ.}

\prose{สจฺจํ หเว สาทุตรํ รสานํ, ปญฺญาชีวิํ ชีวิตมาหุ เสฏฺฐนฺติ.}

\prose{กถํสุ ตรติ โอฆํ, กถํสุ ตรติ อณฺณวํ.}

\prose{กถํสุ ทุกฺขมจฺเจติ, กถํสุ ปริสุชฺฌตีติ.}

\csromanheader{2}{10. ยกฺขสํยุตฺต}
\csromangathameasure{สทฺธาย ตรติ โอฆํ, อปฺปมาเทน อณฺณวํ. \\ วีริเยน ทุกฺขมจฺเจติ, ปญฺญาย ปริสุชฺฌตีติ. \\ กถํสุ ลภเต ปญฺญํ, กถํสุ วินฺทเต ธนํ. \\ กถํสุ กิตฺติํ ปปฺโปติ, กถํ มิตฺตานิ คนฺถติ. \\ อสฺมา โลกา ปรํ โลกํ, กถํ เปจฺจ น โสจตีติ. \\ สทฺทหาโน อรหตํ, ธมฺมํ นิพฺพานปตฺติยา. \\ สุสฺสูสํ ลภเต ปญฺญํ, อปฺปมตฺโต วิจกฺขโณ. \\ ปติรูปการี ธุรวา, อุฏฺฐาตา วินฺทเต ธนํ. \\ สจฺเจน กิตฺติํ ปปฺโปติ, ททํ มิตฺตานิ คนฺถติ. \\ อสฺมา โลกา ปรํ โลกํ, เอวํ เปจฺจ น โสจติ. \\ ยสฺเสเต จตุโร ธมฺมา, สทฺธสฺส ฆรเมสิโน. \\ สจฺจํ ทมฺโม ธิติ จาโค, ส เว เปจฺจ น โสจติ. \\ อิงฺฆ อญฺเญปิ ปุจฺฉสฺสุ, ปุถู สมณพฺราหฺมเณ. \\ ยทิ สจฺจา ทมฺมา จาคา, ขนฺตฺยา ภิยฺโยธ วิชฺชตีติ. \\ กถํ นุ ทานิ ปุจฺเฉยฺยํ, ปุถู สมณพฺราหฺมเณ. \\ โยหํ อชฺช ปชานามิ, โย อตฺโถ สมฺปรายิโก. \\ อตฺถาย วต เม พุทฺโธ, วาสายาฬวิมาคมา. \\ โยหํ อชฺช ปชานามิ, ยตฺถ ทินฺนํ มหปฺผลํ. \\ โส อหํ วิจริสฺสามิ, คามา คามํ ปุรา ปุรํ. \\ นมสฺสมาโน สมฺพุทฺธํ, ธมฺมสฺส จ สุธมฺมตนฺติ.}
\csromangathagroup{\csromangathastanza{\csromanfolio{217}สทฺธาย ตรติ โอฆํ, อปฺปมาเทน อณฺณวํ. \\ วีริเยน ทุกฺขมจฺเจติ, ปญฺญาย ปริสุชฺฌตีติ.}\csromangathastanza{กถํสุ ลภเต ปญฺญํ, กถํสุ วินฺทเต ธนํ. \\ กถํสุ กิตฺติํ ปปฺโปติ, กถํ มิตฺตานิ คนฺถติ.}\csromangathastanza{อสฺมา โลกา ปรํ โลกํ, กถํ เปจฺจ น โสจตีติ. \\ สทฺทหาโน อรหตํ, ธมฺมํ นิพฺพานปตฺติยา.}\csromangathastanza{สุสฺสูสํ\footnote{สุสฺสูสา (สี, อิ)} ลภเต ปญฺญํ, อปฺปมตฺโต วิจกฺขโณ. \\ ปติรูปการี ธุรวา, อุฏฺฐาตา วินฺทเต ธนํ.}\csromangathastanza{สจฺเจน กิตฺติํ ปปฺโปติ, ททํ มิตฺตานิ คนฺถติ. \\ อสฺมา โลกา ปรํ โลกํ, เอวํ เปจฺจ น โสจติ.}\csromangathastanza{ยสฺเสเต จตุโร ธมฺมา, สทฺธสฺส ฆรเมสิโน. \\ สจฺจํ ทมฺโม ธิติ จาโค, ส เว เปจฺจ น โสจติ.}\csromangathastanza{อิงฺฆ อญฺเญปิ ปุจฺฉสฺสุ, ปุถู สมณพฺราหฺมเณ. \\ ยทิ สจฺจา ทมฺมา จาคา, ขนฺตฺยา ภิยฺโยธ วิชฺชตีติ.}\csromangathastanza{กถํ นุ ทานิ ปุจฺเฉยฺยํ, ปุถู สมณพฺราหฺมเณ. \\ โยหํ\footnote{โสหํ (สี), สฺวาหํ (ก)} อชฺช ปชานามิ, โย อตฺโถ สมฺปรายิโก.}\csromangathastanza{อตฺถาย วต เม พุทฺโธ, วาสายาฬวิมาคมา\footnote{มาคโต (อิ, ก)}. \\ โยหํ\footnote{มาคโต (อิ, ก)} อชฺช ปชานามิ, ยตฺถ ทินฺนํ มหปฺผลํ.}\csromangathastanza{โส อหํ วิจริสฺสามิ, คามา คามํ ปุรา ปุรํ. \\ นมสฺสมาโน สมฺพุทฺธํ, ธมฺมสฺส จ สุธมฺมตนฺติ.}}

\vspace{\tipitakaparskip}
\csromancenter{\textbf{ยกฺขสํยุตฺตํ สมตฺตํ.}}
\titlehead{\textbf{ตสฺสุทฺทานํ}}
\csromangathameasure{อินฺทโก สกฺก สูจิ จ, มณิภทฺโท จ สานุ จ. \\ ปิยงฺกร ปุนพฺพสุ สุทตฺโต จ, เทฺว สุกฺกา จีร-อาฬวีติ ทฺวาทส.}
\csromangathagroup{\csromangathastanza{อินฺทโก สกฺก สูจิ จ, มณิภทฺโท จ สานุ จ. \\ ปิยงฺกร ปุนพฺพสุ สุทตฺโต จ, เทฺว สุกฺกา จีร-อาฬวีติ ทฺวาทส.}}

\csromanheader{2}{11. สกฺกสํยุตฺต}
\chapterhead{1. ปฐมวคฺค}
\csromanheader{2}{1. สุวีรสุตฺต}
\proseitem{247}{\csromanfolio{218}เอวํ เม สุตํ\csromandash{}เอกํ สมยํ ภควา สาวตฺถิยํ วิหรติ เชตวเน อนาถปิณฺฑิกสฺส อาราเม.\csromanspacer{} ตตฺร โข ภควา ภิกฺขู อามนฺเตสิ “ภิกฺขโว”ติ. “ภทนฺเต”ติ เต ภิกฺขู ภควโต ปจฺจสฺโสสุํ.\csromanspacer{} ภควา เอตทโวจ\csromandash{}}

\prose{ภูตปุพฺพํ ภิกฺขเว อสุรา เทเว อภิยํสุ.\csromanspacer{} อถ โข ภิกฺขเว สกฺโก เทวานมินฺโท สุวีรํ เทวปุตฺตํ อามนฺเตสิ “เอเต ตาต สุวีร อสุรา เทเว อภิยนฺติ, คจฺฉ ตาต สุวีร อสุเร ปจฺจุยฺยาหี”ติ. “เอวํ ภทฺทนฺตวา”ติ โข ภิกฺขเว สุวีโร เทวปุตฺโต สกฺกสฺส เทวานมินฺทสฺส ปฏิสฺสุตฺวา ปมาทํ อาปาเทสิ\footnote{อาหเรสิ (กตฺถจิ) นวงฺคุตฺตเร สีหนาทสุตฺเตปิ.}.\csromanspacer{} ทุติยมฺปิ โข ภิกฺขเว สกฺโก เทวานมินฺโท สุวีรํ เทวปุตฺตํ อามนฺเตสิ “เอเต ตาต สุวีร อสุรา เทเว อภิยนฺติ, คจฺฉ ตาต สุวีร อสุเร ปจฺจุยฺยาหี”ติ. “เอวํ ภทฺทนฺตวา”ติ โข ภิกฺขเว สุวีโร เทวปุตฺโต สกฺกสฺส เทวานมินฺทสฺส ปฏิสฺสุตฺวา ทุติยมฺปิ ปมาทํ อาปาเทสิ.\csromanspacer{} ตติยมฺปิ โข ภิกฺขเว สกฺโก เทวานมินฺโท สุวีรํ เทวปุตฺตํ อามนฺเตสิ “เอเต ตาต สุวีร อสุรา เทเว อภิยนฺติ, คจฺฉ ตาต สุวีร อสุเร ปจฺจุยฺยาหี”ติ. “เอวํ ภทฺทนฺตวา”ติ โข ภิกฺขเว สุวีโร เทวปุตฺโต สกฺกสฺส เทวานมินฺทสฺส ปฏิสฺสุตฺวา ตติยมฺปิ ปมาทํ อาปาเทสิ.\csromanspacer{} อถ โข ภิกฺขเว สกฺโก เทวานมินฺโท สุวีรํ เทวปุตฺตํ คาถาย อชฺฌภาสิ\csromandash{}}

\csromangathameasure{“อนุฏฺฐหํ อวายามํ, สุขํ ยตฺราธิคจฺฉติ. \\ สุวีร ตตฺถ คจฺฉาหิ, มญฺจ ตตฺเถว ปาปยา”ติ. \\ อลสฺวสฺส อนุฏฺฐาตา, น จ กิจฺจานิ การเย. \\ สพฺพกามสมิทฺธสฺส, ตํ เม สกฺก วรํ ทิสาติ.}
\csromangathagroup{\csromangathastanza{“อนุฏฺฐหํ อวายามํ, สุขํ ยตฺราธิคจฺฉติ. \\ สุวีร ตตฺถ คจฺฉาหิ, มญฺจ ตตฺเถว ปาปยา”ติ.}\csromangathastanza{อลสฺวสฺส\footnote{อลส’สฺส (สี, อิ), อลสฺวายํ (สฺยา, กํ)} อนุฏฺฐาตา, น จ กิจฺจานิ การเย. \\ สพฺพกามสมิทฺธสฺส, ตํ เม สกฺก วรํ ทิสาติ.}}

\csromanheader{2}{11. สกฺกสํยุตฺต}
\csromangathameasure{ยตฺถาลโส อนุฏฺฐาตา, อจฺจนฺตํ สุขเมธติ. \\ สุวีร ตตฺถ คจฺฉาหิ, มญฺจ ตตฺเถว ปาปยาติ. \\ อกมฺมุนา เทวเสฏฺฐ, สกฺก วินฺเทมุ ยํ สุขํ. \\ อโสกํ อนุปายาสํ, ตํ เม สกฺก วรํ ทิสาติ. \\ สเจ อตฺถิ อกมฺเมน, โกจิ กฺวจิ น ชีวติ. \\ นิพฺพานสฺส หิ โส มคฺโค, สุวีร ตตฺถ คจฺฉาหิ.}
\csromangathagroup{\csromangathastanza{\csromanfolio{219}ยตฺถาลโส อนุฏฺฐาตา, อจฺจนฺตํ สุขเมธติ. \\ สุวีร ตตฺถ คจฺฉาหิ, มญฺจ ตตฺเถว ปาปยาติ.}\csromangathastanza{อกมฺมุนา\footnote{อกมฺมนา (สี, อิ)} เทวเสฏฺฐ, สกฺก วินฺเทมุ ยํ สุขํ. \\ อโสกํ อนุปายาสํ, ตํ เม สกฺก วรํ ทิสาติ.}\csromangathastanza{สเจ อตฺถิ อกมฺเมน, โกจิ กฺวจิ น ชีวติ. \\ นิพฺพานสฺส หิ โส มคฺโค, สุวีร ตตฺถ คจฺฉาหิ.}}

\prose{มญฺจ ตตฺเถว ปาปยาติ.}

\prose{โส หิ นาม ภิกฺขเว สกฺโก เทวานมินฺโท สกํ ปุญฺญผลํ อุปชีวมาโน เทวานํ ตาวติํสานํ อิสฺสริยาธิปจฺจํ รชฺชํ กาเรนฺโต อุฏฺฐานวีริยสฺส วณฺณวาที ภวิสฺสติ.\csromanspacer{} อิธ โข ตํ ภิกฺขเว โสเภถ, ยํ ตุเมฺห เอวํ สฺวากฺขาเต ธมฺมวินเย ปพฺพชิตา สมานา อุฏฺฐเหยฺยาถ ฆเฏยฺยาถ วายเมยฺยาถ อปฺปตฺตสฺส ปตฺติยา อนธิคตสฺส อธิคมาย อสจฺฉิกตสฺส สจฺฉิกิริยายาติ.}

\csromanheader{2}{2. สุสีมสุตฺต}
\proseitem{248}{สาวตฺถิยํ.\csromanspacer{} ตตฺร โข ภควา ภิกฺขู อามนฺเตสิ “ภิกฺขโว”ติ. “ภทนฺเต”ติ เต ภิกฺขู ภควโต ปจฺจสฺโสสุํ.\csromanspacer{} ภควา เอตทโวจ\csromandash{}}

\prose{ภูตปุพฺพํ ภิกฺขเว อสุรา เทเว อภิยํสุ.\csromanspacer{} อถ โข ภิกฺขเว สกฺโก เทวานมินฺโท สุสีมํ\footnote{สุสิมํ (สฺยา, กํ, ก)} เทวปุตฺตํ อามนฺเตสิ “เอเต ตาต สุสีม อสุรา เทเว อภิยนฺติ, คจฺฉ ตาต สุสีม อสุเร ปจฺจุยฺยาหี”ติ. “เอวํ ภทฺทนฺตวา”ติ โข ภิกฺขเว สุสีโม เทวปุตฺโต สกฺกสฺส เทวานมินฺทสฺส ปฏิสฺสุตฺวา ปมาทํ อาปาเทสิ.\csromanspacer{} ทุติยมฺปิ โข ภิกฺขเว สกฺโก เทวานมินฺโท สุสีมํ เทวปุตฺตํ อามนฺเตสิ ฯเปฯ.\csromanspacer{} ทุติยมฺปิ ปมาทํ อาปาเทสิ.\csromanspacer{} ตติยมฺปิ โข ภิกฺขเว สกฺโก เทวานมินฺโท สุสีมํ เทวปุตฺตํ อามนฺเตสิ ฯเปฯ.\csromanspacer{} ตติยมฺปิ ปมาทํ อาปาเทสิ.\csromanspacer{} อถ โข ภิกฺขเว สกฺโก เทวานมินฺโท สุสีมํ เทวปุตฺตํ คาถาย อชฺฌภาสิ\csromandash{}}

\csromangathameasure{“อนุฏฺฐหํ อวายามํ, สุขํ ยตฺราธิคจฺฉติ. \\ สุสีม ตตฺถ คจฺฉาหิ, มญฺจ ตตฺเถว ปาปยา”ติ. \\ อลสฺวสฺส อนุฏฺฐาตา, น จ กิจฺจานิ การเย. \\ สพฺพกามสมิทฺธสฺส, ตํ เม สกฺก วรํ ทิสาติ. \\ ยตฺถาลโส อนุฏฺฐาตา, อจฺจนฺตํ สุขเมธติ. \\ สุสีม ตตฺถ คจฺฉาหิ, มญฺจ ตตฺเถว ปาปยาติ. \\ อกมฺมุนา เทวเสฏฺฐ, สกฺก วินฺเทมุ ยํ สุขํ. \\ อโสกํ อนุปายาสํ, ตํ เม สกฺก วรํ ทิสาติ. \\ สเจ อตฺถิ อกมฺเมน, โกจิ กฺวจิ น ชีวติ. \\ นิพฺพานสฺส หิ โส มคฺโค, สุสีม ตตฺถ คจฺฉาหิ.}
\csromangathagroup{\csromangathastanza{“อนุฏฺฐหํ อวายามํ, สุขํ ยตฺราธิคจฺฉติ. \\ สุสีม ตตฺถ คจฺฉาหิ, มญฺจ ตตฺเถว ปาปยา”ติ.}\csromangathastanza{\csromanfolio{220}อลสฺวสฺส อนุฏฺฐาตา, น จ กิจฺจานิ การเย. \\ สพฺพกามสมิทฺธสฺส, ตํ เม สกฺก วรํ ทิสาติ.}\csromangathastanza{ยตฺถาลโส อนุฏฺฐาตา, อจฺจนฺตํ สุขเมธติ. \\ สุสีม ตตฺถ คจฺฉาหิ, มญฺจ ตตฺเถว ปาปยาติ.}\csromangathastanza{อกมฺมุนา เทวเสฏฺฐ, สกฺก วินฺเทมุ ยํ สุขํ. \\ อโสกํ อนุปายาสํ, ตํ เม สกฺก วรํ ทิสาติ.}\csromangathastanza{สเจ อตฺถิ อกมฺเมน, โกจิ กฺวจิ น ชีวติ. \\ นิพฺพานสฺส หิ โส มคฺโค, สุสีม ตตฺถ คจฺฉาหิ.}}

\prose{มญฺจ ตตฺเถว ปาปยาติ.}

\prose{โส หิ นาม ภิกฺขเว สกฺโก เทวานมินฺโท สกํ ปุญฺญผลํ อุปชีวมาโน เทวานํ ตาวติํสานํ อิสฺสริยาธิปจฺจํ รชฺชํ กาเรนฺโต อุฏฺฐานวีริยสฺส วณฺณวาที ภวิสฺสติ.\csromanspacer{} อิธ โข ตํ ภิกฺขเว โสเภถ, ยํ ตุเมฺห เอวํ สฺวากฺขาเต ธมฺมวินเย ปพฺพชิตา สมานา อุฏฺฐเหยฺยาถ ฆเฏยฺยาถ วายเมยฺยาถ อปฺปตฺตสฺส ปตฺติยา อนธิคตสฺส อธิคมาย อสจฺฉิกตสฺส สจฺฉิกิริยายาติ.}

\csromanheader{2}{3. ธชคฺคสุตฺต}
\proseitem{249}{สาวตฺถิยํ.\csromanspacer{} ตตฺร โข ภควา ภิกฺขู อามนฺเตสิ “ภิกฺขโว”ติ. “ภทนฺเต”ติ เต ภิกฺขู ภควโต ปจฺจสฺโสสุํ.\csromanspacer{} ภควา เอตทโวจ\csromandash{}}

\prose{ภูตปุพฺพํ ภิกฺขเว เทวาสุรสงฺคาโม สมุปพฺยูโฬฺห อโหสิ.\csromanspacer{} อถ โข ภิกฺขเว สกฺโก เทวานมินฺโท เทเว ตาวติํเส อามนฺเตสิ\csromandash{}}

\prose{สเจ มาริสา เทวานํ สงฺคามคตานํ อุปฺปชฺเชยฺย ภยํ วา ฉมฺภิตตฺตํ วา โลมหํโส วา, มเมว ตสฺมิํ สมเย ธชคฺคํ อุลฺโลเกยฺยาถ, มมํ หิ โว ธชคฺคํ อุลฺโลกยตํ ยํ ภวิสฺสติ ภยํ วา ฉมฺภิตตฺตํ วา โลมหํโส วา, โส ปหียิสฺสติ.}

\prose{โน เจ เม ธชคฺคํ อุลฺโลเกยฺยาถ, อถ ปชาปติสฺส เทวราชสฺส ธชคฺคํ อุลฺโลเกยฺยาถ, ปชาปติสฺส หิ โว เทวราชสฺส ธชคฺคํ อุลฺโลกยตํ ยํ ภวิสฺสติ ภยํ วา ฉมฺภิตตฺตํ วา โลมหํโส วา, โส ปหียิสฺสติ.}

\csromanheader{2}{11. สกฺกสํยุตฺต}
\prose{\csromanfolio{221}โน เจ ปชาปติสฺส เทวราชสฺส ธชคฺคํ อุลฺโลเกยฺยาถ, อถ วรุณสฺส เทวราชสฺส ธชคฺคํ อุลฺโลเกยฺยาถ, วรุณสฺส หิ โว เทวราชสฺส ธชคฺคํ อุลฺโลกยตํ ยํ ภวิสฺสติ ภยํ วา ฉมฺภิตตฺตํ วา โลมหํโส วา, โส ปหียิสฺสติ.}

\prose{โน เจ วรุณสฺส เทวราชสฺส ธชคฺคํ อุลฺโลเกยฺยาถ, อถ ¢สานสฺส เทวราชสฺส ธชคฺคํ อุลฺโลเกยฺยาถ, ¢สานสฺส หิ โว เทวราชสฺส ธชคฺคํ อุลฺโลกยตํ ยํ ภวิสฺสติ ภยํ วา ฉมฺภิตตฺตํ วา โลมหํโส วา, โส ปหียิสฺสตีติ.}

\prose{ตํ โข ปน ภิกฺขเว สกฺกสฺส วา เทวานมินฺทสฺส ธชคฺคํ อุลฺโลกยตํ, ปชาปติสฺส วา เทวราชสฺส ธชคฺคํ อุลฺโลกยตํ, วรุณสฺส วา เทวราชสฺส ธชคฺคํ อุลฺโลกยตํ, ¢สานสฺส วา เทวราชสฺส ธชคฺคํ อุลฺโลกยตํ ยํ ภวิสฺสติ ภยํ วา ฉมฺภิตตฺตํ วา โลมหํโส วา, โส ปหีเยถาปิ โนปิ ปหีเยถ\footnote{โน ปหีเยถ (ก)}.}

\prose{ตํ กิสฺส เหตุ, สกฺโก หิ ภิกฺขเว เทวานมินฺโท อวีตราโค อวีตโทโส อวีตโมโห ภีรุ ฉมฺภี อุตฺราสี ปลายีติ.}

\prose{อหญฺจ โข ภิกฺขเว เอวํ วทามิ\csromandash{}สเจ ตุมฺหากํ ภิกฺขเว อรญฺญคตานํ วา รุกฺขมูลคตานํ วา สุญฺญคารคตานํ วา อุปฺปชฺเชยฺย ภยํ วา ฉมฺภิตตฺตํ วา โลมหํโส วา, มเมว ตสฺมิํ สมเย อนุสฺสเรยฺยาถ\csromandash{}}

\prose{“อิติปิ โส ภควา อรหํ สมฺมาสมฺพุทฺโธ วิชฺชาจรณสมฺปนฺโน สุคโต โลกวิทู อนุตฺตโร ปุริสทมฺมสารถิ สตฺถา เทวมนุสฺสานํ พุทฺโธ ภควา”ติ.\csromanspacer{} มมํ หิ โว ภิกฺขเว อนุสฺสรตํ ยํ ภวิสฺสติ ภยํ วา ฉมฺภิตตฺตํ วา โลมหํโส วา, โส ปหียิสฺสติ.}

\prose{โน เจ มํ อนุสฺสเรยฺยาถ, อถ ธมฺมํ อนุสฺสเรยฺยาถ “สฺวากฺขาโต ภควตา ธมฺโม สนฺทิฏฺฐิโก อกาลิโก เอหิปสฺสิโก โอปเนยฺยิโก ปจฺจตฺตํ เวทิตพฺโพ วิญฺญูหี”ติ.\csromanspacer{} ธมฺมํ หิ โว ภิกฺขเว อนุสฺสรตํ ยํ ภวิสฺสติ ภยํ วา ฉมฺภิตตฺตํ วา โลมหํโส วา, โส ปหียิสฺสติ.}

\prose{โน เจ ธมฺมํ อนุสฺสเรยฺยาถ, อถ สํฆํ อนุสฺสเรยฺยาถ “สุปฺปฏิปนฺโน ภควโต สาวกสํโฆ, อุชุปฺปฏิปนฺโน ภควโต}

\prose{\csromanfolio{222}สาวกสํโฆ, ญายปฺปฏิปนฺโน ภควโต สาวกสํโฆ, สามีจิปฺปฏิปนฺโน ภควโต สาวกสํโฆ, ยทิทํ จตฺตาริ ปุริสยุคานิ อฏฺฐ ปุริสปุคฺคลา เอส ภควโต สาวกสํโฆ, อาหุเนยฺโย ปาหุเนยฺโย ทกฺขิเณยฺโย อญฺชลิกรณีโย อนุตฺตรํ ปุญฺญกฺเขตฺตํ โลกสฺสา”ติ.\csromanspacer{} สํฆํ หิ โว ภิกฺขเว อนุสฺสรตํ ยํ ภวิสฺสติ ภยํ วา ฉมฺภิตตฺตํ วา โลมหํโส วา, โส ปหียิสฺสติ.}

\prose{ตํ กิสฺส เหตุ, ตถาคโต หิ ภิกฺขเว อรหํ สมฺมาสมฺพุทฺโธ วีตราโค วีตโทโส วีตโมโห อภีรุ อจฺฉมฺภี อนุตฺราสี อปลายีติ.\csromanspacer{} อิทมโวจ ภควา, อิทํ วตฺวาน สุคโต อถาปรํ เอตทโวจ สตฺถา\csromandash{}}

\csromangathameasure{“อรญฺเญ รุกฺขมูเล วา, สุญฺญาคาเรว ภิกฺขโว. \\ อนุสฺสเรถ สมฺพุทฺธํ, ภยํ ตุมฺหาก โน สิยา. \\ โน เจ พุทฺธํ สเรยฺยาถ, โลกเชฏฺฐํ นราสภํ. \\ อถ ธมฺมํ สเรยฺยาถ, นิยฺยานิกํ สุเทสิตํ. \\ โน เจ ธมฺมํ สเรยฺยาถ, นิยฺยานิกํ สุเทสิตํ. \\ อถ สํฆํ สเรยฺยาถ, ปุญฺญกฺเขตฺตํ อนุตฺตรํ. \\ เอวํ พุทฺธํ สรนฺตานํ, ธมฺมํ สํฆญฺจ ภิกฺขโว. \\ ภยํ วา ฉมฺภิตตฺตํ วา, โลมหํโส น เหสฺสตี”ติ.}
\csromangathagroup{\csromangathastanza{“อรญฺเญ รุกฺขมูเล วา, สุญฺญาคาเรว ภิกฺขโว. \\ อนุสฺสเรถ\footnote{อนุสฺสเรยฺยาถ (ก) ปทสิทฺธิ ปน จินฺเตตพฺพา.} สมฺพุทฺธํ, ภยํ ตุมฺหาก โน สิยา.}\csromangathastanza{โน เจ พุทฺธํ สเรยฺยาถ, โลกเชฏฺฐํ นราสภํ. \\ อถ ธมฺมํ สเรยฺยาถ, นิยฺยานิกํ สุเทสิตํ.}\csromangathastanza{โน เจ ธมฺมํ สเรยฺยาถ, นิยฺยานิกํ สุเทสิตํ. \\ อถ สํฆํ สเรยฺยาถ, ปุญฺญกฺเขตฺตํ อนุตฺตรํ.}\csromangathastanza{เอวํ พุทฺธํ สรนฺตานํ, ธมฺมํ สํฆญฺจ ภิกฺขโว. \\ ภยํ วา ฉมฺภิตตฺตํ วา, โลมหํโส น เหสฺสตี”ติ.}}

\csromanheader{2}{4. เวปจิตฺติสุตฺต}
\proseitem{250}{สาวตฺถินิทานํ.\csromanspacer{} ภูตปุพฺพํ ภิกฺขเว เทวาสุรสงฺคาโม สมุปพฺยูโฬฺห อโหสิ.\csromanspacer{} อถ โข ภิกฺขเว เวปจิตฺติ อสุรินฺโท อสุเร อามนฺเตสิ “สเจ มาริสา เทวานํ อสุรสงฺคาเม สมุปพฺยูเฬฺห อสุรา ชิเนยฺยุํ, เทวา ปราชิเนยฺยุํ\footnote{ปราเชยฺยุํ (สี, อิ)}.\csromanspacer{} เยน นํ สกฺกํ เทวานมินฺทํ กณฺฐปญฺจเมหิ พนฺธเนหิ พนฺธิตฺวา มม สนฺติเก อาเนยฺยาถ อสุรปุรนฺ”ติ.\csromanspacer{} สกฺโกปิ โข ภิกฺขเว เทวานมินฺโท เทเว ตาวติํเส อามนฺเตสิ “สเจ มาริสา เทวานํ อสุรสงฺคาเม สมุปพฺยูเฬฺห เทวา ชิเนยฺยุํ, อสุรา ปราชิเนยฺยุํ.\csromanspacer{} เยน นํ เวปจิตฺติํ อสุรินฺทํ กณฺฐปญฺจเมหิ พนฺธเนหิ พนฺธิตฺวา มม สนฺติเก อาเนยฺยาถ สุธมฺมสภนฺ”ติ.\csromanspacer{} ตสฺมิํ โข ปน ภิกฺขเว สงฺคาเม เทวา}

\vspace{\tipitakaparskip}
\csromancenter{สกฺกสํยุตฺต ชินิํสุ, อสุรา ปราชินิํสุ\footnote{ปราชิํสุ (สี, อิ)}.\csromanspacer{} อถ โข ภิกฺขเว เทวา ตาวติํสา เวปจิตฺติํ อสุรินฺทํ กณฺฐปญฺจเมหิ พนฺธเนหิ พนฺธิตฺวา สกฺกสฺส เทวานมินฺทสฺส สนฺติเก อาเนสุํ สุธมฺมสภํ.\csromanspacer{} ตตฺร สุทํ ภิกฺขเว เวปจิตฺติ อสุรินฺโท กณฺฐปญฺจเมหิ พนฺธเนหิ พทฺโธ สกฺกํ เทวานมินฺทํ สุธมฺมสภํ ปวิสนฺตญฺจ นิกฺขมนฺตญฺจ อสพฺภาหิ ผรุสาหิ วาจาหิ อกฺโกสติ ปริภาสติ.\csromanspacer{} อถ โข ภิกฺขเว มาตลิ สงฺคาหโก สกฺกํ เทวานมินฺทํ คาถาหิ อชฺฌภาสิ\csromandash{}}
\vspace{0.5\baselineskip}
\csromangathameasure{“ภยา นุ มฆวา สกฺก, ทุพฺพลฺยา โน ติติกฺขสิ. \\ สุณนฺโต ผรุสํ วาจํ, สมฺมุขา เวปจิตฺติโน”ติ. \\ นาหํ ภยา น ทุพฺพลฺยา, ขมามิ เวปจิตฺติโน. \\ กถํ หิ มาทิโส วิญฺญู, พาเลน ปฏิสํยุเชติ. \\ ภิยฺโย พาลา ปภิชฺเชยฺยุํ, โน จสฺส ปฏิเสธโก. \\ ตสฺมา ภุเสน ทณฺเฑน, ธีโร พาลํ นิเสธเยติ. \\ เอตเทว อหํ มญฺเญ, พาลสฺส ปฏิเสธนํ. \\ ปรํ สงฺกุปิตํ ญตฺวา, โย สโต อุปสมฺมตีติ. \\ เอตเทว ติติกฺขาย, วชฺชํ ปสฺสามิ วาสว. \\ ยทา นํ มญฺญติ พาโล, ภยา มฺยายํ ติติกฺขติ. \\ อชฺฌารุหติ ทุมฺเมโธ, โคว ภิยฺโย ปลายินนฺติ. \\ กามํ มญฺญตุ วา มา วา, ภยา มฺยายํ ติติกฺขติ. \\ สทตฺถปรมา อตฺถา, ขนฺตฺยา ภิยฺโย น วิชฺชติ. \\ โย หเว พลวา สนฺโต, ทุพฺพลสฺส ติติกฺขติ. \\ ตมาหุ ปรมํ ขนฺติํ, นิจฺจํ ขมติ ทุพฺพโล. \\ อพลํ ตํ พลํ อาหุ, ยสฺส พาลพลํ พลํ. \\ พลสฺส ธมฺมคุตฺตสฺส, ปฏิวตฺตา น วิชฺชติ. \\ ตสฺเสว เตน ปาปิโย, โย กุทฺธํ ปฏิกุชฺฌติ. \\ กุทฺธํ อปฺปฏิกุชฺฌนฺโต, สงฺคามํ เชติ ทุชฺชยํ. \\ อุภินฺนมตฺถํ จรติ, อตฺตโน จ ปรสฺส จ. \\ ปรํ สงฺกุปิตํ ญตฺวา, โย สโต อุปสมฺมติ. \\ อุภินฺนํ ติกิจฺฉนฺตานํ, อตฺตโน จ ปรสฺส จ. \\ ชนา มญฺญนฺติ พาโลติ, เย ธมฺมสฺส อโกวิทาติ.}
\csromangathagroup{\csromangathastanza{\csromanfolio{223}“ภยา นุ มฆวา สกฺก, ทุพฺพลฺยา โน ติติกฺขสิ. \\ สุณนฺโต ผรุสํ วาจํ, สมฺมุขา เวปจิตฺติโน”ติ.}\csromangathastanza{นาหํ ภยา น ทุพฺพลฺยา, ขมามิ เวปจิตฺติโน. \\ กถํ หิ มาทิโส วิญฺญู, พาเลน ปฏิสํยุเชติ.}\csromangathastanza{ภิยฺโย พาลา ปภิชฺเชยฺยุํ, โน จสฺส ปฏิเสธโก. \\ ตสฺมา ภุเสน ทณฺเฑน, ธีโร พาลํ นิเสธเยติ.}\csromangathastanza{เอตเทว อหํ มญฺเญ, พาลสฺส ปฏิเสธนํ. \\ ปรํ สงฺกุปิตํ ญตฺวา, โย สโต อุปสมฺมตีติ.}\csromangathastanza{เอตเทว ติติกฺขาย, วชฺชํ ปสฺสามิ วาสว. \\ ยทา นํ มญฺญติ พาโล, ภยา มฺยายํ ติติกฺขติ.}\csromangathastanza{อชฺฌารุหติ ทุมฺเมโธ, โคว ภิยฺโย ปลายินนฺติ. \\ กามํ มญฺญตุ วา มา วา, ภยา มฺยายํ ติติกฺขติ.}\csromangathastanza{สทตฺถปรมา อตฺถา, ขนฺตฺยา ภิยฺโย น วิชฺชติ. \\ โย หเว พลวา สนฺโต, ทุพฺพลสฺส ติติกฺขติ.}\csromangathastanza{ตมาหุ ปรมํ ขนฺติํ, นิจฺจํ ขมติ ทุพฺพโล. \\ อพลํ ตํ พลํ อาหุ, ยสฺส พาลพลํ พลํ.}\csromangathastanza{พลสฺส ธมฺมคุตฺตสฺส, ปฏิวตฺตา น วิชฺชติ. \\ ตสฺเสว เตน ปาปิโย, โย กุทฺธํ ปฏิกุชฺฌติ.}\csromangathastanza{กุทฺธํ อปฺปฏิกุชฺฌนฺโต, สงฺคามํ เชติ ทุชฺชยํ. \\ อุภินฺนมตฺถํ จรติ, อตฺตโน จ ปรสฺส จ.}\csromangathastanza{\csromanfolio{224}ปรํ สงฺกุปิตํ ญตฺวา, โย สโต อุปสมฺมติ. \\ อุภินฺนํ ติกิจฺฉนฺตานํ, อตฺตโน จ ปรสฺส จ. \\ ชนา มญฺญนฺติ พาโลติ, เย ธมฺมสฺส อโกวิทาติ.}}

\prose{โส หิ นาม ภิกฺขเว สกฺโก เทวานมินฺโท สกํ ปุญฺญผลํ อุปชีวมาโน เทวานํ ตาวติํสานํ อิสฺสริยาธิปจฺจํ รชฺชํ กาเรนฺโต ขนฺติโสรจฺจสฺส วณฺณวาที ภวิสฺสติ.\csromanspacer{} อิธ โข ตํ ภิกฺขเว โสเภถ, ยํ ตุเมฺห เอวํ สฺวากฺขาเต ธมฺมวินเย ปพฺพชิตา สมานา ขมา จ ภเวยฺยาถ โสรตา จาติ.}

\csromanheader{2}{5. สุภาสิตชยสุตฺต}
\proseitem{251}{สาวตฺถินิทานํ.\csromanspacer{} ภูตปุพฺพํ ภิกฺขเว เทวาสุรสงฺคาโม สมุปพฺยูโฬฺห อโหสิ.\csromanspacer{} อถ โข ภิกฺขเว เวปจิตฺติ อสุรินฺโท สกฺกํ เทวานมินฺทํ เอตทโวจ “โหตุ เทวานมินฺท สุภาสิเตน ชโย”ติ.\csromanspacer{} โหตุ เวปจิตฺติ สุภาสิเตน ชโยติ.\csromanspacer{} อถ โข ภิกฺขเว เทวา จ อสุรา จ ปาริสชฺเช ฐเปสุํ “อิเม โน สุภาสิตทุพฺภาสิตํ อาชานิสฺสนฺตี”ติ.\csromanspacer{} อถ โข ภิกฺขเว เวปจิตฺติ อสุรินฺโท สกฺกํ เทวานมินฺทํ เอตทโวจ “ภณ เทวานมินฺท คาถนฺ”ติ.\csromanspacer{} เอวํ วุตฺเต ภิกฺขเว สกฺโก เทวานมินฺโท เวปจิตฺติํ อสุรินฺทํ เอตทโวจ “ตุเมฺห เขฺวตฺถ เวปจิตฺติ ปุพฺพเทวา ภณ เวปจิตฺติ คาถนฺ”ติ.\csromanspacer{} เอวํ วุตฺเต ภิกฺขเว เวปจิตฺติ อสุรินฺโท อิมํ คาถํ อภาสิ\csromandash{}}

\csromangathameasure{“ภิยฺโย พาลา ปภิชฺเชยฺยุํ, โน จสฺส ปฏิเสธโก. \\ ตสฺมา ภุเสน ทณฺเฑน, ธีโร พาลํ นิเสธเย”ติ.}
\csromangathagroup{\csromangathastanza{“ภิยฺโย พาลา ปภิชฺเชยฺยุํ, โน จสฺส ปฏิเสธโก. \\ ตสฺมา ภุเสน ทณฺเฑน, ธีโร พาลํ นิเสธเย”ติ.}}

\prose{ภาสิตาย โข ปน ภิกฺขเว เวปจิตฺตินา อสุรินฺเทน คาถาย อสุรา อนุโมทิํสุ, เทวา ตุณฺหี อเหสุํ.\csromanspacer{} อถ โข ภิกฺขเว เวปจิตฺติ อสุรินฺโท สกฺกํ เทวานมินฺทํ เอตทโวจ “ภณ เทวานมินฺท คาถนฺ”ติ.\csromanspacer{} เอวํ วุตฺเต ภิกฺขเว สกฺโก เทวานมินฺโท อิมํ คาถํ อภาสิ\csromandash{}}

\csromangathameasure{“เอตเทว อหํ มญฺเญ, พาลสฺส ปฏิเสธนํ. \\ ปรํ สงฺกุปิตํ ญตฺวา, โย สโต อุปสมฺมตี”ติ.}
\csromangathagroup{\csromangathastanza{“เอตเทว อหํ มญฺเญ, พาลสฺส ปฏิเสธนํ. \\ ปรํ สงฺกุปิตํ ญตฺวา, โย สโต อุปสมฺมตี”ติ.}}

\csromanheader{2}{11. สกฺกสํยุตฺต}
\prose{\csromanfolio{225}ภาสิตาย โข ปน ภิกฺขเว สกฺเกน เทวานมินฺเทน คาถาย เทวา อนุโมทิํสุ, อสุรา ตุณฺหี อเหสุํ.\csromanspacer{} อถ โข ภิกฺขเว สกฺโก เทวานมินฺโท เวปจิตฺติํ อสุรินฺทํ เอตทโวจ “ภณ เวปจิตฺติ คาถนฺ”ติ.\csromanspacer{} เอวํ วุตฺเต ภิกฺขเว เวปจิตฺติ อสุรินฺโท อิมํ คาถํ อภาสิ\csromandash{}}

\csromangathameasure{“เอตเทว ติติกฺขาย, วชฺชํ ปสฺสามิ วาสว. \\ ยทา นํ มญฺญติ พาโล, ภยา มฺยายํ ติติกฺขติ. \\ อชฺฌารุหติ ทุมฺเมโธ, โคว ภิยฺโย ปลายินนฺ”ติ.}
\csromangathagroup{\csromangathastanza{“เอตเทว ติติกฺขาย, วชฺชํ ปสฺสามิ วาสว. \\ ยทา นํ มญฺญติ พาโล, ภยา มฺยายํ ติติกฺขติ. \\ อชฺฌารุหติ ทุมฺเมโธ, โคว ภิยฺโย ปลายินนฺ”ติ.}}

\prose{ภาสิตาย โข ปน ภิกฺขเว เวปจิตฺตินา อสุรินฺเทน คาถาย อสุรา อนุโมทิํสุ, เทวา ตุณฺหี อเหสุํ.\csromanspacer{} อถ โข ภิกฺขเว เวปจิตฺติ อสุรินฺโท สกฺกํ เทวานมินฺทํ เอตทโวจ “ภณ เทวานมินฺท คาถนฺ”ติ.\csromanspacer{} เอวํ วุตฺเต ภิกฺขเว สกฺโก เทวานมินฺโท อิมา คาถาโย อภาสิ\csromandash{}}

\csromangathameasure{“กามํ มญฺญตุ วา มา วา, ภยา มฺยายํ ติติกฺขติ. \\ สทตฺถปรมา อตฺถา, ขนฺตฺยา ภิยฺโย น วิชฺชติ. \\ โย หเว พลวา สนฺโต, ทุพฺพลสฺส ติติกฺขติ. \\ ตมาหุ ปรมํ ขนฺติํ, นิจฺจํ ขมติ ทุพฺพโล. \\ อพลํ ตํ พลํ อาหุ, ยสฺส พาลพลํ พลํ. \\ พลสฺส ธมฺมคุตฺตสฺส, ปฏิวตฺตา น วิชฺชติ. \\ ตสฺเสว เตน ปาปิโย, โย กุทฺธํ ปฏิกุชฺฌติ. \\ กุทฺธํ อปฺปฏิกุชฺฌนฺโต, สงฺคามํ เชติ ทุชฺชยํ. \\ อุภินฺนมตฺถํ จรติ, อตฺตโน จ ปรสฺส จ. \\ ปรํ สงฺกุปิตํ ญตฺวา, โย สโต อุปสมฺมติ. \\ อุภินฺนํ ติกิจฺฉนฺตานํ, อตฺตโน จ ปรสฺส จ. \\ ชนา มญฺญนฺติ พาโลติ, เย ธมฺมสฺส อโกวิทา”ติ.}
\csromangathagroup{\csromangathastanza{“กามํ มญฺญตุ วา มา วา, ภยา มฺยายํ ติติกฺขติ. \\ สทตฺถปรมา อตฺถา, ขนฺตฺยา ภิยฺโย น วิชฺชติ.}\csromangathastanza{โย หเว พลวา สนฺโต, ทุพฺพลสฺส ติติกฺขติ. \\ ตมาหุ ปรมํ ขนฺติํ, นิจฺจํ ขมติ ทุพฺพโล.}\csromangathastanza{อพลํ ตํ พลํ อาหุ, ยสฺส พาลพลํ พลํ. \\ พลสฺส ธมฺมคุตฺตสฺส, ปฏิวตฺตา น วิชฺชติ.}\csromangathastanza{ตสฺเสว เตน ปาปิโย, โย กุทฺธํ ปฏิกุชฺฌติ. \\ กุทฺธํ อปฺปฏิกุชฺฌนฺโต, สงฺคามํ เชติ ทุชฺชยํ.}\csromangathastanza{อุภินฺนมตฺถํ จรติ, อตฺตโน จ ปรสฺส จ. \\ ปรํ สงฺกุปิตํ ญตฺวา, โย สโต อุปสมฺมติ.}\csromangathastanza{อุภินฺนํ ติกิจฺฉนฺตานํ, อตฺตโน จ ปรสฺส จ. \\ ชนา มญฺญนฺติ พาโลติ, เย ธมฺมสฺส อโกวิทา”ติ.}}

\prose{ภาสิตาสุ โข ปน ภิกฺขเว สกฺเกน เทวานมินฺเทน คาถาสุ เทวา อนุโมทิํสุ, อสุรา ตุณฺหี อเหสุํ.\csromanspacer{} อถ โข ภิกฺขเว เทวานญฺจ อสุรานญฺจ ปาริสชฺชา เอตทโวจุํ “ภาสิตา โข เวปจิตฺตินา อสุรินฺเทน คาถาโย, ตา จ โข สทณฺฑาวจรา สสตฺถาวจรา, อิติ ภณฺฑนํ}

\prose{\csromanfolio{226}อิติ วิคฺคโห อิติ กลโห.\csromanspacer{} ภาสิตา โข\footnote{ภาสิตา โข ปน (สี)} สกฺเกน เทวานมินฺเทน คาถาโย, ตา จ โข อทณฺฑาวจรา อสตฺถาวจรา, อิติ อภณฺฑนํ อิติ อวิคฺคโห อิติ อกลโห, สกฺกสฺส เทวานมินฺทสฺส สุภาสิเตน ชโย”ติ.\csromanspacer{} อิติ โข ภิกฺขเว สกฺกสฺส เทวานมินฺทสฺส สุภาสิเตน ชโย อโหสีติ.}

\csromanheader{2}{6. กุลาวกสุตฺต}
\proseitem{252}{สาวตฺถิยํ.\csromanspacer{} ภูตปุพฺพํ ภิกฺขเว เทวาสุรสงฺคาโม สมุปพฺยูโฬฺห อโหสิ.\csromanspacer{} ตสฺมิํ โข ปน ภิกฺขเว สงฺคาเม อสุรา ชินิํสุ, เทวา ปราชินิํสุ.\csromanspacer{} ปราชิตา จ โข ภิกฺขเว เทวา อปายํเสฺวว อุตฺตเรนมุขา, อภิยํเสฺวว เน อสุรา.\csromanspacer{} อถ โข ภิกฺขเว สกฺโก เทวานมินฺโท มาตลิ สงฺคาหกํ คาถาย อชฺฌภาสิ\csromandash{}}

\csromangathameasure{“กุลาวกา มาตลิ สิมฺพลิสฺมิํ, \\ ¢สามุเขน ปริวชฺชยสฺสุ. \\ กามํ จชาม อสุเรสุ ปาณํ, \\ มายิเม ทิชา วิกุลาวกา อเหสุนฺ”ติ.}
\csromangathagroup{\csromangathastanza{“กุลาวกา มาตลิ สิมฺพลิสฺมิํ, \\ ¢สามุเขน ปริวชฺชยสฺสุ. \\ กามํ จชาม อสุเรสุ ปาณํ, \\ มายิเม ทิชา วิกุลาวกา\footnote{วิกุลาวา (สฺยา, กํ, ก)} อเหสุนฺ”ติ.}}

\prose{“เอวํ ภทฺทนฺตวา”ติ โข ภิกฺขเว มาตลิ สงฺคาหโก สกฺกสฺส เทวานมินฺทสฺส ปฏิสฺสุตฺวา สหสฺสยุตฺตํ อาชญฺญรถํ ปจฺจุทาวตฺเตสิ.\csromanspacer{} อถ โข ภิกฺขเว อสุรานํ เอตทโหสิ “ปจฺจุทาวตฺโต โข ทานิ สกฺกสฺส เทวานมินฺทสฺส สหสฺสยุตฺโต อาชญฺญรโถ, ทุติยมฺปิ โข เทวา อสุเรหิ สงฺคาเมสฺสนฺตี”ติ ภีตา อสุรปุรเมว ปาวิสิํสุ.\csromanspacer{} อิติ โข ภิกฺขเว สกฺกสฺส เทวานมินฺทสฺส ธมฺเมน ชโย อโหสีติ.}

\csromanheader{2}{7. นทุพฺภิยสุตฺต}
\proseitem{253}{สาวตฺถิยํ.\csromanspacer{} ภูตปุพฺพํ ภิกฺขเว สกฺกสฺส เทวานมินฺทสฺส รโหคตสฺส ปฏิสลฺลีนสฺส เอวํ เจตโส ปริวิตกฺโก อุทปาทิ “โยปิ เม อสฺส สุปจฺจตฺถิโก ตสฺสปาหํ น ทุพฺเภยฺยนฺ”ติ.\csromanspacer{} อถ โข ภิกฺขเว เวปจิตฺติ อสุรินฺโท สกฺกสฺส เทวานมินฺทสฺส เจตสา เจโตปริวิตกฺกมญฺญาย เยน สกฺโก เทวานมินฺโท เตนุปสงฺกมิ.\csromanspacer{} อทฺทสา โข}

\vspace{\tipitakaparskip}
\csromancenter{สกฺกสํยุตฺต ภิกฺขเว สกฺโก เทวานมินฺโท เวปจิตฺติํ อสุรินฺทํ ทูรโตว อาคจฺฉนฺตํ, ทิสฺวาน เวปจิตฺติํ อสุรินฺทํ เอตทโวจ “ติฏฺฐ เวปจิตฺติ คหิโตสี”ติ.}
\vspace{0.5\baselineskip}
\csromangathameasure{ยเทว เต มาริส ปุพฺเพ จิตฺตํ, ตเทว ตฺวํ มา ปชหาสีติ.}
\csromangathagroup{\csromangathastanza{\csromanfolio{227}ยเทว เต มาริส ปุพฺเพ จิตฺตํ, ตเทว ตฺวํ มา ปชหาสีติ\footnote{ตเทว ตฺวํ มาริส ปหาสีติ (สี, สฺยา, กํ) 2. อทฺรุพฺภาย (ก)}.}}

\prose{สปสฺสุ จ เม เวปจิตฺติ อทุพฺภายาติ2.}

\csromangathameasure{ยํ มุสา ภณโต ปาปํ, ยํ ปาปํ อริยูปวาทิโน. \\ มิตฺตทฺทุโน จ ยํ ปาปํ, ยํ ปาปํ อกตญฺญุโน. \\ ตเมว ปาปํ ผุสตุ, โย เต ทุพฺเภ สุชมฺปตีติ.}
\csromangathagroup{\csromangathastanza{ยํ มุสา ภณโต ปาปํ, ยํ ปาปํ อริยูปวาทิโน. \\ มิตฺตทฺทุโน จ ยํ ปาปํ, ยํ ปาปํ อกตญฺญุโน. \\ ตเมว ปาปํ ผุสตุ\footnote{ผุสติ (สี, อิ)}, โย เต ทุพฺเภ สุชมฺปตีติ.}}

\csromanheader{2}{8. เวโรจน-อสุรินฺทสุตฺต}
\proseitem{254}{สาวตฺถิยํ เชตวเน.\csromanspacer{} เตน โข ปน สมเยน ภควา ทิวาวิหารคโต โหติ ปฏิสลฺลีโน.\csromanspacer{} อถ โข สกฺโก จ เทวานมินฺโท เวโรจโน จ อสุรินฺโท เยน ภควา เตนุปสงฺกมิํสุ, อุปสงฺกมิตฺวา ปจฺเจกํ ทฺวารพาหํ นิสฺสาย อฏฺฐํสุ.\csromanspacer{} อถ โข เวโรจโน อสุรินฺโท ภควโต สนฺติเก อิมํ คาถํ อภาสิ\csromandash{}}

\csromangathameasure{“วายเมเถว ปุริโส, ยาว อตฺถสฺส นิปฺผทา. \\ นิปฺผนฺนโสภโน อตฺโถ5, เวโรจนวโจ อิทนฺ”ติ. \\ วายเมเถว ปุริโส, ยาว อตฺถสฺส นิปฺผทา. \\ นิปฺผนฺนโสภโน อตฺโถ, ขนฺตฺยา ภิยฺโย น วิชฺชตีติ. \\ สพฺเพ สตฺตา อตฺถชาตา, ตตฺถ ตตฺถ ยถารหํ. \\ สํโยคปรมา เตฺวว, สมฺโภคา สพฺพปาณินํ. \\ นิปฺผนฺนโสภโน อตฺโถ, เวโรจนวโจ อิทนฺติ. \\ สพฺเพ สตฺตา อตฺถชาตา, ตตฺถ ตตฺถ ยถารหํ. \\ สํโยคปรมา เตฺวว, สมฺโภคา สพฺพปาณินํ. \\ นิปฺผนฺนโสภโน อตฺโถ, ขนฺตฺยา ภิยฺโย น วิชฺชตีติ.}
\csromangathagroup{\csromangathastanza{“วายเมเถว ปุริโส, ยาว อตฺถสฺส นิปฺผทา. \\ นิปฺผนฺนโสภโน\footnote{โสภิโน (สี), โสภโณ (อิ, ก) 5. อตฺถา (สี)} อตฺโถ5, เวโรจนวโจ อิทนฺ”ติ.}\csromangathastanza{วายเมเถว ปุริโส, ยาว อตฺถสฺส นิปฺผทา. \\ นิปฺผนฺนโสภโน อตฺโถ\footnote{นิปฺผนฺนโสภิโน อตฺถา (สี, สฺยา, กํ)}, ขนฺตฺยา ภิยฺโย น วิชฺชตีติ.}\csromangathastanza{สพฺเพ สตฺตา อตฺถชาตา, ตตฺถ ตตฺถ ยถารหํ. \\ สํโยคปรมา เตฺวว, สมฺโภคา สพฺพปาณินํ.}\csromangathastanza{นิปฺผนฺนโสภโน อตฺโถ, เวโรจนวโจ อิทนฺติ. \\ สพฺเพ สตฺตา อตฺถชาตา, ตตฺถ ตตฺถ ยถารหํ.}\csromangathastanza{สํโยคปรมา เตฺวว, สมฺโภคา สพฺพปาณินํ. \\ นิปฺผนฺนโสภโน อตฺโถ, ขนฺตฺยา ภิยฺโย น วิชฺชตีติ.}}

\csromanheader{2}{9. อรญฺญายตน-อิสิสุตฺต}
\proseitem{255}{\csromanfolio{228}สาวตฺถิยํ.\csromanspacer{} ภูตปุพฺพํ ภิกฺขเว สมฺพหุลา อิสโย สีลวนฺโต กลฺยาณธมฺมา อรญฺญายตเน ปณฺณกุฏีสุ สมฺมนฺติ.\csromanspacer{} อถ โข ภิกฺขเว สกฺโก จ เทวานมินฺโท เวปจิตฺติ จ อสุรินฺโท เยน เต อิสโย สีลวนฺโต กลฺยาณธมฺมา เตนุปสงฺกมิํสุ.\csromanspacer{} อถ โข ภิกฺขเว เวปจิตฺติ อสุรินฺโท ปฏลิโย\footnote{อฏลิโย (สี, สฺยา, กํ, อิ), อาฏลิโย (ก) ม 2 (362) ปิฏฺเฐปิ.} อุปาหนา อาโรหิตฺวา ขคฺคํ โอลคฺเคตฺวา ฉตฺเตน ธาริยมาเนน อคฺคทฺวาเรน อสฺสมํ ปวิสิตฺวา เต อิสโย สีลวนฺเต กลฺยาณธมฺเม อปพฺยามโต กริตฺวา อติกฺกมิ.\csromanspacer{} อถ โข ภิกฺขเว สกฺโก เทวานมินฺโท ปฏลิโย อุปาหนา โอโรหิตฺวา ขคฺคํ อญฺเญสํ ทตฺวา ฉตฺตํ อปนาเมตฺวา ทฺวาเรเนว อสฺสมํ ปวิสิตฺวา เต อิสโย สีลวนฺเต กลฺยาณธมฺเม อนุวาตํ ปญฺชลิโก นมสฺสมาโน อฏฺฐาสิ.\csromanspacer{} อถ โข ภิกฺขเว เต อิสโย สีลวนฺโต กลฺยาณธมฺมา สกฺกํ เทวานมินฺทํ คาถาย อชฺฌภาสิํสุ\csromandash{}}

\csromangathameasure{“คนฺโธ อิสีนํ จิรทิกฺขิตานํ, \\ กายา จุโต คจฺฉติ มาลุเตน. \\ อิโต ปฏิกฺกมฺม สหสฺสเนตฺต, \\ คนฺโธ อิสีนํ อสุจิ เทวราชา”ติ. \\ คนฺโธ อิสีนํ จิรทิกฺขิตานํ, \\ กายา จุโต คจฺฉตุ มาลุเตน. \\ สุจิตฺรปุปฺผํ สิรสฺมิํว มาลํ, \\ คนฺธํ เอตํ ปฏิกงฺขาม ภนฺเต. \\ น เหตฺถ เทวา ปฏิกูลสญฺญิโนติ.}
\csromangathagroup{\csromangathastanza{“คนฺโธ อิสีนํ จิรทิกฺขิตานํ, \\ กายา จุโต คจฺฉติ มาลุเตน. \\ อิโต ปฏิกฺกมฺม สหสฺสเนตฺต, \\ คนฺโธ อิสีนํ อสุจิ เทวราชา”ติ.}\csromangathastanza{คนฺโธ อิสีนํ จิรทิกฺขิตานํ, \\ กายา จุโต คจฺฉตุ\footnote{คจฺฉติ (สี, สฺยา, กํ)} มาลุเตน. \\ สุจิตฺรปุปฺผํ สิรสฺมิํว มาลํ, \\ คนฺธํ เอตํ ปฏิกงฺขาม ภนฺเต.}\csromangathastanza{น เหตฺถ เทวา ปฏิกูลสญฺญิโนติ.}}

\csromanheader{2}{10. สมุทฺทกสุตฺต}
\proseitem{256}{สาวตฺถิยํ.\csromanspacer{} ภูตปุพฺพํ ภิกฺขเว สมฺพหุลา อิสโย สีลวนฺโต กลฺยาณธมฺมา สมุทฺทตีเร ปณฺณกุฏีสุ สมฺมนฺติ.\csromanspacer{} เตน โข ปน สมเยน เทวาสุรสงฺคาโม สมุปพฺยูโฬฺห อโหสิ.\csromanspacer{} อถ โข ภิกฺขเว เตสํ อิสีนํ สีลวนฺตานํ กลฺยาณธมฺมานํ เอตทโหสิ “ธมฺมิกา โข เทวา อธมฺมิกา อสุรา, สิยาปิ โน อสุรโต ภยํ.\csromanspacer{} ยํนูน มยํ สมฺพรํ อสุรินฺทํ}

\vspace{\tipitakaparskip}
\csromancenter{สกฺกสํยุตฺต อุปสงฺกมิตฺวา อภยทกฺขิณํ ยาเจยฺยามา”ติ.\csromanspacer{} อถ โข ภิกฺขเว เต อิสโย สีลวนฺโต กลฺยาณธมฺมา.\csromanspacer{} เสยฺยถาปิ นาม พลวา ปุริโส สมิญฺชิตํ วา พาหํ ปสาเรยฺย, ปสาริตํ วา พาหํ สมิญฺเชยฺย.\csromanspacer{} เอวเมว สมุทฺทตีเร ปณฺณกุฏีสุ อนฺตรหิตา สมฺพรสฺส อสุรินฺทสฺส สมฺมุเข ปาตุรเหสุํ.\csromanspacer{} อถ โข ภิกฺขเว เต อิสโย สีลวนฺโต กลฺยาณธมฺมา สมฺพรํ อสุรินฺทํ คาถาย อชฺฌภาสิํสุ\csromandash{}}
\vspace{0.5\baselineskip}
\csromangathameasure{“อิสโย สมฺพรํ ปตฺตา, ยาจนฺติ อภยทกฺขิณํ. \\ กามํกโร หิ เต ทาตุํ, ภยสฺส อภยสฺส วา”ติ. \\ อิสีนํ อภยํ นตฺถิ, ทุฏฺฐานํ สกฺกเสวินํ. \\ อภยํ ยาจมานานํ, ภยเมว ททามิ โวติ. \\ อภยํ ยาจมานานํ, ภยเมว ททาสิ โน. \\ ปฏิคฺคณฺหาม เต เอตํ, อกฺขยํ โหตุ เต ภยํ. \\ ยาทิสํ วปเต พีชํ, ตาทิสํ หรเต ผลํ. \\ กลฺยาณการี กลฺยาณํ, ปาปการี จ ปาปกํ. \\ ปวุตฺตํ ตาต เต พีชํ, ผลํ ปจฺจนุโภสฺสสีติ.}
\csromangathagroup{\csromangathastanza{\csromanfolio{229}“อิสโย สมฺพรํ ปตฺตา, ยาจนฺติ อภยทกฺขิณํ. \\ กามํกโร หิ เต ทาตุํ, ภยสฺส อภยสฺส วา”ติ.}\csromangathastanza{อิสีนํ อภยํ นตฺถิ, ทุฏฺฐานํ สกฺกเสวินํ. \\ อภยํ ยาจมานานํ, ภยเมว ททามิ โวติ.}\csromangathastanza{อภยํ ยาจมานานํ, ภยเมว ททาสิ โน. \\ ปฏิคฺคณฺหาม เต เอตํ, อกฺขยํ โหตุ เต ภยํ.}\csromangathastanza{ยาทิสํ วปเต พีชํ, ตาทิสํ หรเต ผลํ. \\ กลฺยาณการี กลฺยาณํ, ปาปการี จ ปาปกํ. \\ ปวุตฺตํ ตาต เต พีชํ, ผลํ ปจฺจนุโภสฺสสีติ.}}

\prosecont{อถ โข ภิกฺขเว เต อิสโย สีลวนฺโต กลฺยาณธมฺมา สมฺพรํ อสุรินฺทํ อภิสปิตฺวา เสยฺยถาปิ นาม พลวา ปุริโส สมิญฺชิตํ วา พาหํ ปสาเรยฺย, ปสาริตํ วา พาหํ สมิญฺเชยฺย.\csromanspacer{} เอวเมว สมฺพรสฺส อสุรินฺทสฺส สมฺมุเข อนฺตรหิตา สมุทฺทตีเร ปณฺณกุฏีสุ ปาตุรเหสุํ.\csromanspacer{} อถ โข ภิกฺขเว สมฺพโร อสุรินฺโท เตหิ อิสีหิ สีลวนฺเตหิ กลฺยาณธมฺเมหิ อภิสปิโต รตฺติยา สุทํ ติกฺขตฺตุํ อุพฺพิชฺชีติ.}

\vspace{\tipitakaparskip}
\csromancenter{ปฐโม วคฺโค.}
\titlehead{\textbf{ตสฺสุทฺทานํ}}
\csromangathameasure{สุวีรํ สุสีมํ เจว, ธชคฺคํ เวปจิตฺติโน. \\ สุภาสิตํ ชยํ เจว, กุลาวกํ นทุพฺภิยํ. \\ เวโรจน อสุรินฺโท, อิสโย อรญฺญกํ เจว. \\ อิสโย จ สมุทฺทกาติ.}
\csromangathagroup{\csromangathastanza{สุวีรํ สุสีมํ เจว, ธชคฺคํ เวปจิตฺติโน. \\ สุภาสิตํ ชยํ เจว, กุลาวกํ นทุพฺภิยํ.}\csromangathastanza{เวโรจน อสุรินฺโท, อิสโย อรญฺญกํ เจว. \\ อิสโย จ สมุทฺทกาติ.}}

\chapterheadpageread{2. ทุติยวคฺค}
\csromanheader{2}{1. วตปทสุตฺต}
\proseitem{257}{\csromanfolio{230}สาวตฺถิยํ.\csromanspacer{} สกฺกสฺส ภิกฺขเว เทวานมินฺทสฺส ปุพฺเพ มนุสฺสภูตสฺส สตฺต วตปทานิ\footnote{วตฺตปทานิ (ก)} สมตฺตานิ สมาทินฺนานิ อเหสุํ, เยสํ สมาทินฺนตฺตา สกฺโก สกฺกตฺตํ อชฺฌคา.\csromanspacer{} กตมานิ สตฺต วตปทานิ, ยาวชีวํ มาตาเปตฺติภโร อสฺสํ, ยาวชีวํ กุเล เชฏฺฐาปจายี อสฺสํ, ยาวชีวํ สณฺหวาโจ อสฺสํ, ยาวชีวํ อปิสุณวาโจ อสฺสํ, ยาวชีวํ วิคตมลมจฺเฉเรน เจตสา อคารํ อชฺฌาวเสยฺยํ, มุตฺตจาโค ปยตปาณิ โวสฺสคฺครโต ยาจโยโค ทานสํวิภาครโต, ยาวชีวํ สจฺจวาโจ อสฺสํ, ยาวชีวํ อกฺโกธโน อสฺสํ “สเจปิ เม โกโธ อุปฺปชฺเชยฺย, ขิปฺปเมว นํ ปฏิวิเนยฺยนฺ”ติ.\csromanspacer{} สกฺกสฺส ภิกฺขเว เทวานมินฺทสฺส ปุพฺเพ มนุสฺสภูตสฺส อิมานิ สตฺต วตปทานิ สมตฺตานิ สมาทินฺนานิ อเหสุํ, เยสํ สมาทินฺนตฺตา สกฺโก สกฺกตฺตํ อชฺฌคาติ.}

\csromangathameasure{มาตาเปตฺติภรํ ชนฺตุํ, กุเล เชฏฺฐาปจายินํ. \\ สณฺหํ สขิลสมฺภาสํ, เปสุเณยฺยปฺปหายินํ. \\ มจฺเฉรวินเย ยุตฺตํ, สจฺจํ โกธาภิภุํ นรํ. \\ ตํ เว เทวา ตาวติํสา, อาหุ “สปฺปุริโส” อิตีติ.}
\csromangathagroup{\csromangathastanza{มาตาเปตฺติภรํ ชนฺตุํ, กุเล เชฏฺฐาปจายินํ. \\ สณฺหํ สขิลสมฺภาสํ, เปสุเณยฺยปฺปหายินํ.}\csromangathastanza{มจฺเฉรวินเย ยุตฺตํ, สจฺจํ โกธาภิภุํ นรํ. \\ ตํ เว เทวา ตาวติํสา, อาหุ “สปฺปุริโส” อิตีติ.}}

\csromanheader{2}{2. สกฺกนามสุตฺต}
\proseitem{258}{สาวตฺถิยํ เชตวเน.\csromanspacer{} ตตฺร โข ภควา ภิกฺขู เอตทโวจ “สกฺโก ภิกฺขเว เทวานมินฺโท ปุพฺเพ มนุสฺสภูโต สมาโน มโฆ นาม มาณโว อโหสิ, ตสฺมา มฆวา”ติ วุจฺจติ.}

\prose{สกฺโก ภิกฺขเว เทวานมินฺโท ปุพฺเพ มนุสฺสภูโต สมาโน ปุเร\footnote{ปุเร ปุเร (สี, อิ)} ทานํ อทาสิ, ตสฺมา “ปุรินฺทโท”ติ วุจฺจติ.}

\prose{สกฺโก ภิกฺขเว เทวานมินฺโท ปุพฺเพ มนุสฺสภูโต สมาโน สกฺกจฺจํ ทานํ อทาสิ, ตสฺมา “สกฺโก”ติ วุจฺจติ.}

\csromanheader{2}{11. สกฺกสํยุตฺต}
\prose{\csromanfolio{231}สกฺโก ภิกฺขเว เทวานมินฺโท ปุพฺเพ มนุสฺสภูโต สมาโน อาวสถํ อทาสิ, ตสฺมา “วาสโว”ติ วุจฺจติ.}

\prose{สกฺโก ภิกฺขเว เทวานมินฺโท สหสฺสมฺปิ อตฺถานํ มุหุตฺเตน จินฺเตติ, ตสฺมา “สหสฺสกฺโข”ติ วุจฺจติ.}

\prose{สกฺกสฺส ภิกฺขเว เทวานมินฺทสฺส สุชา นาม อสุรกญฺญา ปชาปติ, ตสฺมา “สุชมฺปตี”ติ วุจฺจติ.}

\prose{สกฺโก ภิกฺขเว เทวานมินฺโท เทวานํ ตาวติํสานํ อิสฺสริยาธิปจฺจํ รชฺชํ กาเรติ, ตสฺมา “เทวานมินฺโท”ติ วุจฺจติ.}

\prose{สกฺกสฺส ภิกฺขเว เทวานมินฺทสฺส ปุพฺเพ มนุสฺสภูตสฺส สตฺต วตปทานิ สมตฺตานิ สมาทินฺนานิ อเหสุํ, เยสํ สมาทินฺนตฺตา สกฺโก สกฺกตฺตํ อชฺฌคา.\csromanspacer{} กตมานิ สตฺต วตปทานิ, ยาวชีวํ มาตาเปตฺติภโร อสฺสํ, ยาวชีวํ กุเล เชฏฺฐาปจายี อสฺสํ, ยาวชีวํ สณฺหวาโจ อสฺสํ, ยาวชีวํ อปิสุณวาโจ อสฺสํ, ยาวชีวํ วิคตมลมจฺเฉเรน เจตสา อคารํ อชฺฌาวเสยฺยํ, มุตฺตจาโค ปยตปาณิ โวสฺสคฺครโต ยาจโยโค ทานสํวิภาครโต, ยาวชีวํ สจฺจวาโจ อสฺสํ, ยาวชีวํ อกฺโกธโน อสฺสํ “สเจปิ เม โกโธ อุปฺปชฺเชยฺย, ขิปฺปเมว นํ ปฏิวิเนยฺยนฺ”ติ.\csromanspacer{} สกฺกสฺส ภิกฺขเว เทวานมินฺทสฺส ปุพฺเพ มนุสฺสภูตสฺส อิมานิ สตฺต วตปทานิ สมตฺตานิ สมาทินฺนานิ อเหสุํ, เยสํ สมาทินฺนตฺตา สกฺโก สกฺกตฺตํ อชฺฌคาติ.}

\csromangathameasure{มาตาเปตฺติภรํ ชนฺตุํ, กุเล เชฏฺฐาปจายินํ. \\ สณฺหํ สขิลสมฺภาสํ, เปสุเณยฺยปฺปหายินํ. \\ มจฺเฉรวินเย ยุตฺตํ, สจฺจํ โกธาภิภุํ นรํ. \\ ตํ เว เทวา ตาวติํสา, อาหุ “สปฺปุริโส” อิตีติ.}
\csromangathagroup{\csromangathastanza{มาตาเปตฺติภรํ ชนฺตุํ, กุเล เชฏฺฐาปจายินํ. \\ สณฺหํ สขิลสมฺภาสํ, เปสุเณยฺยปฺปหายินํ.}\csromangathastanza{มจฺเฉรวินเย ยุตฺตํ, สจฺจํ โกธาภิภุํ นรํ. \\ ตํ เว เทวา ตาวติํสา, อาหุ “สปฺปุริโส” อิตีติ.}}

\csromanheader{2}{3. มหาลิสุตฺต}
\proseitem{259}{เอวํ เม สุตํ\csromandash{}เอกํ สมยํ ภควา เวสาลิยํ วิหรติ มหาวเน กูฏาคารสาลายํ.\csromanspacer{} อถ โข มหาลิ ลิจฺฉวี เยน ภควา เตนุปสงฺกมิ, อุปสงฺกมิตฺวา ภควนฺตํ อภิวาเทตฺวา เอกมนฺตํ นิสีทิ, เอกมนฺตํ นิสินฺโน โข มหาลิ ลิจฺฉวี ภควนฺตํ เอตทโวจ “ทิฏฺโฐ \csromanfolio{232}โข ภนฺเต ภควตา สกฺโก เทวานมินฺโท”ติ.\csromanspacer{} ทิฏฺโฐ โข เม มหาลิ สกฺโก เทวานมินฺโทติ.\csromanspacer{} โส หิ นุน ภนฺเต สกฺกปติรูปโก ภวิสฺสติ, ทุทฺทโส หิ ภนฺเต สกฺโก เทวานมินฺโทติ.\csromanspacer{} สกฺกญฺจ ขฺวาหํ มหาลิ ปชานามิ สกฺกกรเณ จ ธมฺเม, เยสํ ธมฺมานํ สมาทินฺนตฺตา สกฺโก สกฺกตฺตํ อชฺฌคา, ตญฺจ ปชานามิ.}

\prose{สกฺโก มหาลิ เทวานมินฺโท ปุพฺเพ มนุสฺสภูโต สมาโน มโฆ นาม มาณโว อโหสิ, ตสฺมา “มฆวา”ติ วุจฺจติ.}

\prose{สกฺโก มหาลิ เทวานมินฺโท ปุพฺเพ มนุสฺสภูโต สมาโน สกฺกจฺจํ ทานํ อทาสิ, ตสฺมา “สกฺโก”ติ วุจฺจติ.}

\prose{สกฺโก มหาลิ เทวานมินฺโท ปุพฺเพ มนุสฺสภูโต สมาโน ปุเร ทานํ อทาสิ, ตสฺมา “ปุรินฺทโท”ติ วุจฺจติ.}

\prose{สกฺโก มหาลิ เทวานมินฺโท ปุพฺเพ มนุสฺสภูโต สมาโน อาวสถํ อทาสิ, ตสฺมา “วาสโว”ติ วุจฺจติ.}

\prose{สกฺโก มหาลิ เทวานมินฺโท สหสฺสมฺปิ อตฺถานํ มุหุตฺเตน จินฺเตติ, ตสฺมา “สหสฺสกฺโข”ติ วุจฺจติ.}

\prose{สกฺกสฺส มหาลิ เทวานมินฺทสฺส สุชา นาม อสุรกญฺญา ปชาปติ, ตสฺมา “สุชมฺปตี”ติ วุจฺจติ.}

\prose{สกฺโก มหาลิ เทวานมินฺโท เทวานํ ตาวติํสานํ อิสฺสริยาธิปจฺจํ รชฺชํ กาเรติ, ตสฺมา “เทวานมินฺโท”ติ วุจฺจติ.}

\prose{สกฺกสฺส มหาลิ เทวานมินฺทสฺส ปุพฺเพ มนุสฺสภูตสฺส สตฺต วตปทานิ สมตฺตานิ สมาทินฺนานิ อเหสุํ, เยสํ สมาทินฺนตฺตา สกฺโก สกฺกตฺตํ อชฺฌคา.\csromanspacer{} กตมานิ สตฺต วตปทานิ, ยาวชีวํ มาตาเปตฺติภโร อสฺสํ, ยาวชีวํ กุเล เชฏฺฐาปจายี อสฺสํ, ยาวชีวํ สณฺหวาโจ อสฺสํ, ยาวชีวํ อปิสุณวาโจ อสฺสํ, ยาวชีวํ วิคตมลมจฺเฉเรน เจตสา อคารํ อชฺฌาวเสยฺยํ, มุตฺตจาโค ปยตปาณิ โวสฺสคฺครโต ยาจโยโค ทานสํวิภาครโต, ยาวชีวํ สจฺจวาโจ อสฺสํ, ยาวชีวํ อกฺโกธโน อสฺสํ “สเจปิ เม โกโธ อุปฺปชฺเชยฺย, ขิปฺปเมว นํ ปฏิวิเนยฺยนฺ”ติ.\csromanspacer{} สกฺกสฺส มหาลิ เทวานมินฺทสฺส ปุพฺเพ มนุสฺสภูตสฺส}

\vspace{\tipitakaparskip}
\csromancenter{สกฺกสํยุตฺต อิมานิ สตฺต วตปทานิ สมตฺตานิ สมาทินฺนานิ อเหสุํ, เยสํ สมาทินตฺตา สกฺโก สกฺกตฺตํ อชฺฌคาติ.}
\vspace{0.5\baselineskip}
\csromangathameasure{มาตาเปตฺติภรํ ชนฺตุํ, กุเล เชฏฺฐาปจายินํ. \\ สณฺหํ สขิลสมฺพาสํ, เปสุเณยฺยปฺปหายินํ. \\ มจฺเฉรวินเย ยุตฺตํ, สจฺจํ โกธาภิภุํ นรํ. \\ ตํ เว เทวา ตาวติํสา, อาหุ “สปฺปุริโส” อิตีติ.}
\csromangathagroup{\csromangathastanza{\csromanfolio{233}มาตาเปตฺติภรํ ชนฺตุํ, กุเล เชฏฺฐาปจายินํ. \\ สณฺหํ สขิลสมฺพาสํ, เปสุเณยฺยปฺปหายินํ.}\csromangathastanza{มจฺเฉรวินเย ยุตฺตํ, สจฺจํ โกธาภิภุํ นรํ. \\ ตํ เว เทวา ตาวติํสา, อาหุ “สปฺปุริโส” อิตีติ.}}

\csromanheader{2}{4. ทลิทฺทสุตฺต}
\proseitem{260}{เอกํ สมยํ ภควา ราชคเห วิหรติ เวฬุวเน กลนฺทกนิวาเป.\csromanspacer{} ตตฺร โข ภควา ภิกฺขู อามนฺเตสิ “ภิกฺขโว”ติ. “ภทนฺเต”ติ เต ภิกฺขู ภควโต ปจฺจสฺโสสุํ.\csromanspacer{} ภควา เอตทโวจ\csromandash{}}

\prose{ภูตปุพฺพํ ภิกฺขเว อญฺญตโร ปุริโส อิมสฺมิํเยว ราชคเห มนุสฺสทลิทฺโท\footnote{มนุสฺสทฬิทฺโท (สี, สฺยา, กํ)} อโหสิ มนุสฺสกปโณ มนุสฺสวราโก, โส ตถาคตปฺปเวทิเต ธมฺมวินเย สทฺธํ สมาทิยิ, สีลํ สมาทิยิ, สุตํ สมาทิยิ, จาคํ สมาทิยิ, ปญฺญํ สมาทิยิ.\csromanspacer{} โส ตถาคตปฺปเวทิเต ธมฺมวินเย สทฺธํ สมาทิยิตฺวา สีลํ สมาทิยิตฺวา สุตํ สมาทิยิตฺวา จาคํ สมาทิยิตฺวา ปญฺญํ สมาทิยิตฺวา กายสฺส เภทา ปรํ มรณา สุคติํ สคฺคํ โลกํ อุปปชฺชิ เทวานํ ตาวติํสานํ สหพฺยตํ.\csromanspacer{} โส อญฺเญ เทเว อติโรจติ วณฺเณน เจว ยสสา จ.\csromanspacer{} ตตฺร สุทํ ภิกฺขเว เทวา ตาวติํสา อุชฺฌายนฺติ ขิยฺยนฺติ วิปาเจนฺติ “อจฺฉริยํ วต โภ, อพฺภุตํ วต โภ, อยํ หิ เทวปุตฺโต ปุพฺเพ มนุสฺสภูโต สมาโน มนุสฺสทลิทฺโท อโหสิ มนุสฺสกปโณ มนุสฺสวราโก, โส กายสฺส เภทา ปรํ มรณา สุคติํ สคฺคํ โลกํ อุปปนฺโน เทวานํ ตาวติํสานํ สหพฺยตํ.\csromanspacer{} โส อญฺเญ เทเว อติโรจติ วณฺเณน เจว ยสสา จา”ติ.}

\prose{อถ โข ภิกฺขเว สกฺโก เทวานมินฺโท เทเว ตาวติํเส อามนฺเตสิ “มา โข ตุเมฺห มาริสา เอตสฺส เทวปุตฺตสฺส อุชฺฌายิตฺถ เอโส โข มาริสา เทวปุตฺโต ปุพฺเพ มนุสฺสภูโต สมาโน \csromanfolio{234}ตถาคตปฺปเวทิเต ธมฺมวินเย สทฺธํ สมาทิยิ, สีลํ สมาทิยิ, สุตํ สมาทิยิ, จาคํ สมาทิยิ, ปญฺญํ สมาทิยิ.\csromanspacer{} โส ตถาคตปฺปเวทิเต ธมฺมวินเย สทฺธํ สมาทิยิตฺวา สีลํ สมาทิยิตฺวา สุตํ สมาทิยิตฺวา จาคํ สมาทิยิตฺวา ปญฺญํ สมาทิยิตฺวา กายสฺส เภทา ปรํ มรณา สุคติํ สคฺคํ โลกํ อุปปนฺโน เทวานํ ตาวติํสานํ สหพฺยตํ.\csromanspacer{} โส อญฺเญ เทเว อติโรจติ วณฺเณน เจว ยสสา จา”ติ.\csromanspacer{} อถ โข ภิกฺขเว สกฺโก เทวานมินฺโท เทเว ตาวติํเส อนุนยมาโน ตายํ เวลายํ อิมา คาถาโย อภาสิ\csromandash{}}

\csromangathameasure{“ยสฺส สทฺธา ตถาคเต, อจลา สุปฺปติฏฺฐิตา. \\ สีลญฺจ ยสฺส กลฺยาณํ, อริยกนฺตํ ปสํสิตํ. \\ สํเฆ ปสาโท ยสฺสตฺถิ, อุชุภูตญฺจ ทสฺสนํ. \\ อทลิทฺโทติ ตํ อาหุ, อโมฆํ ตสฺส ชีวิตํ. \\ ตสฺมา สทฺธญฺจ สีลญฺจ, ปสาทํ ธมฺมทสฺสนํ. \\ อนุยุญฺเชถ เมธาวี, สรํ พุทฺธาน สาสนนฺ”ติ.}
\csromangathagroup{\csromangathastanza{“ยสฺส สทฺธา ตถาคเต, อจลา สุปฺปติฏฺฐิตา. \\ สีลญฺจ ยสฺส กลฺยาณํ, อริยกนฺตํ ปสํสิตํ.}\csromangathastanza{สํเฆ ปสาโท ยสฺสตฺถิ, อุชุภูตญฺจ ทสฺสนํ. \\ อทลิทฺโทติ ตํ อาหุ, อโมฆํ ตสฺส ชีวิตํ.}\csromangathastanza{ตสฺมา สทฺธญฺจ สีลญฺจ, ปสาทํ ธมฺมทสฺสนํ. \\ อนุยุญฺเชถ เมธาวี, สรํ พุทฺธาน สาสนนฺ”ติ.}}

\csromanheader{2}{5. รามเณยฺยกสุตฺต}
\proseitem{261}{สาวตฺถิยํ เชตวเน.\csromanspacer{} อถ โข สกฺโก เทวานมินฺโท เยน ภควา เตนุปสงฺกมิ, อุปสงฺกมิตฺวา ภควนฺตํ อภิวาเทตฺวา เอกมนฺตํ อฏฺฐาสิ, เอกมนฺตํ ฐิโต โข สกฺโก เทวานมินฺโท ภควนฺตํ เอตทโวจ “กิํ นุ โข ภนฺเต ภูมิรามเณยฺยกนฺ”ติ.}

\csromangathameasure{อารามเจตฺยา วนเจตฺยา, โปกฺขรญฺโญ สุนิมฺมิตา. \\ มนุสฺสรามเณยฺยสฺส, กลํ นาคฺฆนฺติ โสฬสิํ. \\ คาเม วา ยทิ วารญฺเญ, นินฺเน วา ยทิ วา ถเล. \\ ยตฺถ อรหนฺโต วิหรนฺติ, ตํ ภูมิรามเณยฺยกนฺติ.}
\csromangathagroup{\csromangathastanza{อารามเจตฺยา วนเจตฺยา, โปกฺขรญฺโญ สุนิมฺมิตา. \\ มนุสฺสรามเณยฺยสฺส, กลํ นาคฺฆนฺติ โสฬสิํ.}\csromangathastanza{คาเม วา ยทิ วารญฺเญ, นินฺเน วา ยทิ วา ถเล. \\ ยตฺถ อรหนฺโต วิหรนฺติ, ตํ ภูมิรามเณยฺยกนฺติ.}}

\csromanheader{2}{6. ยชมานสุตฺต}
\proseitem{262}{เอกํ สมยํ ภควา ราชคเห วิหรติ คิชฺฌกูเฏ ปพฺพเต.\csromanspacer{} อถ โข สกฺโก เทวานมินฺโท เยน ภควา เตนุปสงฺกมิ, อุปสงฺกมิตฺวา}

\vspace{\tipitakaparskip}
\csromancenter{สกฺกสํยุตฺต ภควนฺตํ อภิวาเทตฺวา เอกมนฺตํ อฏฺฐาสิ, เอกมนฺตํ ฐิโต โข สกฺโก เทวานมินฺโท ภควนฺตํ คาถาย อชฺฌภาสิ\csromandash{}}
\vspace{0.5\baselineskip}
\csromangathameasure{“ยชมานานํ มนุสฺสานํ, ปุญฺญเปกฺขาน ปาณินํ. \\ กโรตํ โอปธิกํ ปุญฺญํ, กตฺถ ทินฺนํ มหปฺผลนฺ”ติ. \\ จตฺตาโร จ ปฏิปนฺนา, จตฺตาโร จ ผเล ฐิตา. \\ เอส สํโฆ อุชุภูโต, ปญฺญาสีลสมาหิโต. \\ ยชมานานํ มนุสฺสานํ, ปุญฺญเปกฺขาน ปาณินํ. \\ กโรตํ โอปธิกํ ปุญฺญํ, สํเฆ ทินฺนํ มหปฺผลนฺติ.}
\csromangathagroup{\csromangathastanza{\csromanfolio{235}“ยชมานานํ มนุสฺสานํ, ปุญฺญเปกฺขาน ปาณินํ. \\ กโรตํ โอปธิกํ ปุญฺญํ, กตฺถ ทินฺนํ มหปฺผลนฺ”ติ.}\csromangathastanza{จตฺตาโร จ ปฏิปนฺนา, จตฺตาโร จ ผเล ฐิตา. \\ เอส สํโฆ อุชุภูโต, ปญฺญาสีลสมาหิโต.}\csromangathastanza{ยชมานานํ มนุสฺสานํ, ปุญฺญเปกฺขาน ปาณินํ. \\ กโรตํ โอปธิกํ ปุญฺญํ, สํเฆ ทินฺนํ มหปฺผลนฺติ.}}

\csromanheader{2}{7. พุทฺธวนฺทนาสุตฺต}
\proseitem{263}{สาวตฺถิยํ เชตวเน.\csromanspacer{} เตน โข ปน สมเยน ภควา ทิวาวิหารคโต โหติ ปฏิสลฺลีโน.\csromanspacer{} อถ โข สกฺโก จ เทวานมินฺโท พฺรหฺมา จ สหมฺปติ เยน ภควา เตนุปสงฺกมิํสุ, อุปสงฺกมิตฺวา ปจฺเจกํ ทฺวารพาหํ นิสฺสาย อฏฺฐํสุ.\csromanspacer{} อถ โข สกฺโก เทวานมินฺโท ภควโต สนฺติเก อิมํ คาถํ อภาสิ\csromandash{}}

\csromangathameasure{“อุฏฺเฐหิ วีร วิชิตสงฺคาม, \\ ปนฺนภาร อนณ วิจร โลเก. \\ จิตฺตญฺจ เต สุวิมุตฺตํ, \\ จนฺโท ยถา ปนฺนรสาย รตฺตินฺ”ติ.}
\csromangathagroup{\csromangathastanza{“อุฏฺเฐหิ วีร วิชิตสงฺคาม, \\ ปนฺนภาร อนณ วิจร โลเก. \\ จิตฺตญฺจ เต สุวิมุตฺตํ, \\ จนฺโท ยถา ปนฺนรสาย รตฺตินฺ”ติ.}}

\prose{น โข เทวานมินฺท ตถาคตา เอวํ วนฺทิตพฺพา, เอวญฺจ โข เทวานมินฺท ตถาคตา วนฺทิตพฺพา.}

\csromangathameasure{“อุฏฺเฐหิ วีร วิชิตสงฺคาม, \\ สตฺถวาห อนณ วิจร โลเก. \\ เทสสฺสุ ภควา ธมฺมํ, \\ อญฺญาตาโร ภวิสฺสนฺตี”ติ.}
\csromangathagroup{\csromangathastanza{“อุฏฺเฐหิ วีร วิชิตสงฺคาม, \\ สตฺถวาห อนณ วิจร โลเก. \\ เทสสฺสุ ภควา ธมฺมํ, \\ อญฺญาตาโร ภวิสฺสนฺตี”ติ.}}

\csromanheader{2}{8. คหฏฺฐวนฺทนาสุตฺต}
\proseitem{264}{สาวตฺถิยํ.\csromanspacer{} ตตฺร ฯเปฯ เอตทโวจ\csromandash{}ภูตปุพฺพํ ภิกฺขเว สกฺโก เทวานมินฺโท มาตลิํ สงฺคาหกํ อามนฺเตสิ “โยเชหิ สมฺม มาตลิ \csromanfolio{236}สหสฺสยุตฺตํ อาชญฺญรถํ อุยฺยานภูมิํ คจฺฉาม สุภูมิํ ทสฺสนายา”ติ. “เอวํ ภทฺทนฺตวา”ติ โข ภิกฺขเว มาตลิ สงฺคาหโก สกฺกสฺส เทวานมินฺทสฺส ปฏิสฺสุตฺวา สหสฺสยุตฺตํ อาชญฺญรถํ โยเชตฺวา สกฺกสฺส เทวานมินฺทสฺส ปฏิเวเทสิ “ยุตฺโต โข เต มาริส สหสฺสยุตฺโต อาชญฺญรโถ, ยสฺส ทานิ กาลํ มญฺญสี”ติ.\csromanspacer{} อถ โข ภิกฺขเว สกฺโก เทวานมินฺโท เวชยนฺตปาสาทา โอโรหนฺโต อญฺชลิํ กตฺวา\footnote{ปญฺชลิโก (อิ), ปญฺชลิํ กตฺวา (ก)} สุทํ ปุถุทฺทิสา นมสฺสติ.\csromanspacer{} อถ โข ภิกฺขเว มาตลิ สงฺคาหโก สกฺกํ เทวานมินฺทํ คาถาย อชฺฌภาสิ\csromandash{}}

\csromangathameasure{“ตํ นมสฺสนฺติ เตวิชฺชา, สพฺเพ ภุมฺมา จ ขตฺติยา. \\ จตฺตาโร จ มหาราชา, ติทสา จ ยสสฺสิโน. \\ อถ โก นาม โส ยกฺโข, ยํ ตฺวํ สกฺก นมสฺสสี”ติ. \\ มํ นมสฺสนฺติ เตวิชฺชา, สพฺเพ ภุมฺมา จ ขตฺติยา. \\ จตฺตาโร จ มหาราชา, ติทสา จ ยสสฺสิโน. \\ อหญฺจ สีลสมฺปนฺเน, จิรรตฺตสมาหิเต. \\ สมฺมาปพฺพชิเต วนฺเท, พฺรหฺมจริยปรายเน. \\ เย คหฏฺฐา ปุญฺญกรา, สีลวนฺโต อุปาสกา. \\ ธมฺเมน ทารํ โปเสนฺติ, เต นมสฺสามิ มาตลีติ. \\ เสฏฺฐา หิ กิร โลกสฺมิํ, เย ตฺวํ สกฺก นมสฺสสิ. \\ อหมฺปิ เต นมสฺสามิ, เย นมสฺสสิ วาสวาติ. \\ อิทํ วตฺวาน มฆวา, เทวราชา สุชมฺปติ. \\ ปุถุทฺทิสา นมสฺสิตฺวา, ปมุโข รถมารุหีติ.}
\csromangathagroup{\csromangathastanza{“ตํ นมสฺสนฺติ เตวิชฺชา, สพฺเพ ภุมฺมา จ ขตฺติยา. \\ จตฺตาโร จ มหาราชา, ติทสา จ ยสสฺสิโน.}\csromangathastanza{อถ โก นาม โส ยกฺโข, ยํ ตฺวํ สกฺก นมสฺสสี”ติ. \\ มํ นมสฺสนฺติ เตวิชฺชา, สพฺเพ ภุมฺมา จ ขตฺติยา.}\csromangathastanza{จตฺตาโร จ มหาราชา, ติทสา จ ยสสฺสิโน. \\ อหญฺจ สีลสมฺปนฺเน, จิรรตฺตสมาหิเต.}\csromangathastanza{สมฺมาปพฺพชิเต วนฺเท, พฺรหฺมจริยปรายเน. \\ เย คหฏฺฐา ปุญฺญกรา, สีลวนฺโต อุปาสกา.}\csromangathastanza{ธมฺเมน ทารํ โปเสนฺติ, เต นมสฺสามิ มาตลีติ. \\ เสฏฺฐา หิ กิร โลกสฺมิํ, เย ตฺวํ สกฺก นมสฺสสิ.}\csromangathastanza{อหมฺปิ เต นมสฺสามิ, เย นมสฺสสิ วาสวาติ. \\ อิทํ วตฺวาน มฆวา, เทวราชา สุชมฺปติ. \\ ปุถุทฺทิสา นมสฺสิตฺวา, ปมุโข รถมารุหีติ.}}

\csromanheader{2}{9. สตฺถารวนฺทนาสุตฺต}
\proseitem{265}{สาวตฺถิยํ เชตวเน.\csromanspacer{} ภูตปุพฺพํ ภิกฺขเว สกฺโก เทวานมินฺโท มาตลิํ สงฺคาหกํ อามนฺเตสิ “โยเชหิ สมฺม มาตลิ สหสฺสยุตฺตํ อาชญฺญรถํ อุยฺยานภูมิํ คจฺฉาม สุภูมิํ ทสฺสนายา”ติ. “เอวํ ภทฺทนฺตวา”ติ โข ภิกฺขเว มาตลิ สงฺคาหโก สกฺกสฺส เทวานมินฺทสฺส}

\vspace{\tipitakaparskip}
\csromancenter{สกฺกสํยุตฺต ปฏิสฺสุตฺวา สหสฺสยุตฺตํ อาชญฺญรถํ โยเชตฺวา สกฺกสฺส เทวานมินฺทสฺส ปฏิเวเทสิ “ยุตฺโต โข เต มาริส สหสฺสยุตฺโต อาชญฺญรโถ, ยสฺส ทานิ กาลํ มญฺญสี”ติ.\csromanspacer{} อถ โข ภิกฺขเว สกฺโก เทวานมินฺโท เวชยนฺตปาสาทา โอโรหนฺโต อญฺชลิํ กตฺวา สุทํ ภควนฺตํ นมสฺสติ.\csromanspacer{} อถ โข ภิกฺขเว มาตลิ สงฺคาหโก สกฺกํ เทวานมินฺทํ คาถาย อชฺฌภาสิ\csromandash{}}
\vspace{0.5\baselineskip}
\csromangathameasure{“ยํ หิ เทวา มนุสฺสา จ, ตํ นมสฺสนฺติ วาสว. \\ อถ โก นาม โส ยกฺโข, ยํ ตฺวํ สกฺก นมสฺสสี”ติ. \\ โย อิธ สมฺมาสมฺพุทฺโธ, อสฺมิํ โลเก สเทวเก. \\ อโนมนามํ สตฺถารํ, ตํ นมสฺสามิ มาตลิ. \\ เยสํ ราโค จ โทโส จ, อวิชฺชา จ วิราชิตา. \\ ขีณาสวา อรหนฺโต, เต นมสฺสามิ มาตลิ. \\ เย ราคโทสวินยา, อวิชฺชาสมติกฺกมา. \\ เสกฺขา อปจยารามา, อปฺปมตฺตานุสิกฺขเร.}
\csromangathagroup{\csromangathastanza{\csromanfolio{237}“ยํ หิ เทวา มนุสฺสา จ, ตํ นมสฺสนฺติ วาสว. \\ อถ โก นาม โส ยกฺโข, ยํ ตฺวํ สกฺก นมสฺสสี”ติ.}\csromangathastanza{โย อิธ สมฺมาสมฺพุทฺโธ, อสฺมิํ โลเก สเทวเก. \\ อโนมนามํ สตฺถารํ, ตํ นมสฺสามิ มาตลิ.}\csromangathastanza{เยสํ ราโค จ โทโส จ, อวิชฺชา จ วิราชิตา. \\ ขีณาสวา อรหนฺโต, เต นมสฺสามิ มาตลิ.}\csromangathastanza{เย ราคโทสวินยา, อวิชฺชาสมติกฺกมา. \\ เสกฺขา อปจยารามา, อปฺปมตฺตานุสิกฺขเร.}}

\prose{เต นมสฺสามิ มาตลีติ.}

\csromangathameasure{เสฏฺฐา หิ กิร โลกสฺมิํ, เย ตฺวํ สกฺก นมสฺสสิ. \\ อหมฺปิ เต นมสฺสามิ, เย นมสฺสสิ วาสวาติ. \\ อิทํ วตฺวาน มฆวา, เทวราชา สุชมฺปติ. \\ ภควนฺตํ นมสฺสิตฺวา, ปมุโข รถมารุหีติ.}
\csromangathagroup{\csromangathastanza{เสฏฺฐา หิ กิร โลกสฺมิํ, เย ตฺวํ สกฺก นมสฺสสิ. \\ อหมฺปิ เต นมสฺสามิ, เย นมสฺสสิ วาสวาติ.}\csromangathastanza{อิทํ วตฺวาน มฆวา, เทวราชา สุชมฺปติ. \\ ภควนฺตํ นมสฺสิตฺวา, ปมุโข รถมารุหีติ.}}

\csromanheader{2}{10. สํฆวนฺทนาสุตฺต}
\proseitem{266}{สาวตฺถิยํ เชตวเน.\csromanspacer{} ตตฺร โข ฯเปฯ เอตทโวจ\csromandash{}ภูตปุพฺพํ ภิกฺขเว สกฺโก เทวานมินฺโท มาตลิํ สงฺคาหกํ อามนฺเตสิ “โยเชหิ สมฺม มาตลิ สหสฺสยุตฺตํ อาชญฺญรถํ อุยฺยานภูมิํ คจฺฉาม สุภูมิํ ทสฺสนายา”ติ. “เอวํ ภทฺทนฺตวา”ติ โข ภิกฺขเว มาตลิ สงฺคาหโก สกฺกสฺส เทวานมินฺทสฺส ปฏิสฺสุตฺวา สหสฺสยุตฺตํ อาชญฺญรถํ โยเชตฺวา สกฺกสฺส เทวานมินฺทสฺส ปฏิเวเทสิ “ยุตฺโต โข เต มาริส สหสฺสยุตฺโต อาชญฺญรโถ, ยสฺส ทานิ กาลํ มญฺญสี”ติ.\csromanspacer{} อถ โข \csromanfolio{238}ภิกฺขเว สกฺโก เทวานมินฺโท เวชยนฺตปาสาทา โอโรหนฺโต อญฺชลิํ กตฺวา สุทํ ภิกฺขุสํฆํ นมสฺสติ.\csromanspacer{} อถ โข ภิกฺขเว มาตลิ สงฺคาหโก สกฺกํ เทวานมินฺทํ คาถาย อชฺฌภาสิ\csromandash{}}

\csromangathameasure{“ตํ หิ เอเต นมสฺเสยฺยุํ, ปูติเทหสยา นรา. \\ นิมุคฺคา กุณปเมฺหเต, ขุปฺปิปาสสมปฺปิตา. \\ กิํ นุ เตสํ ปิหยสิ, อนาคาราน วาสว. \\ อาจารํ อิสินํ พฺรูหิ, ตํ สุโณม วโจ ตวา”ติ. \\ เอตํ เตสํ ปิหยามิ, อนาคาราน มาตลิ. \\ ยมฺหา คามา ปกฺกมนฺติ, อนเปกฺขา วชนฺติ เต. \\ น เตสํ โกฏฺเฐ โอเปนฺติ, น กุมฺภิ น กโฬปิยํ.}
\csromangathagroup{\csromangathastanza{“ตํ หิ เอเต นมสฺเสยฺยุํ, ปูติเทหสยา นรา. \\ นิมุคฺคา กุณปเมฺหเต, ขุปฺปิปาสสมปฺปิตา.}\csromangathastanza{กิํ นุ เตสํ ปิหยสิ, อนาคาราน วาสว. \\ อาจารํ อิสินํ พฺรูหิ, ตํ สุโณม วโจ ตวา”ติ.}\csromangathastanza{เอตํ เตสํ ปิหยามิ, อนาคาราน มาตลิ. \\ ยมฺหา คามา ปกฺกมนฺติ, อนเปกฺขา วชนฺติ เต. \\ น เตสํ โกฏฺเฐ โอเปนฺติ, น กุมฺภิ\footnote{น กุมฺภา (สฺยา, กํ, อิ, ก)} น กโฬปิยํ\footnote{ขโฬปิยํ (สี)}.}}

\nitthitam{ปรนิฏฺฐิตเมสานา\footnote{ปรนิฏฺฐิตเมสนา (สฺยา, กํ, ก)}, เตน ยาเปนฺติ สุพฺพตา.}
\csromangathameasure{สุมนฺตมนฺติโน ธีรา, ตุณฺหีภูตา สมญฺจรา. \\ เทวา วิรุทฺธา อสุเรหิ, ปุถุ มจฺจา จ มาตลิ. \\ อวิรุทฺธา วิรุทฺเธสุ, อตฺตทณฺเฑสุ นิพฺพุตา. \\ สาทาเนสุ อนาทานา, เต นมสฺสามิ มาตลีติ. \\ เสฏฺฐา หิ กิร โลกสฺมิํ, เย ตฺวํ สกฺก นมสฺสสิ. \\ อหมฺปิ เต นมสฺสามิ, เย นมสฺสสิ วาสวาติ. \\ อิทํ วตฺวาน มฆวา, เทวราชา สุชมฺปติ. \\ ภิกฺขุสํฆํ นมสฺสิตฺวา, ปมุโข รถมารุหีติ.}
\csromangathagroup{\csromangathastanza{สุมนฺตมนฺติโน ธีรา, ตุณฺหีภูตา สมญฺจรา. \\ เทวา วิรุทฺธา อสุเรหิ, ปุถุ มจฺจา จ มาตลิ.}\csromangathastanza{อวิรุทฺธา วิรุทฺเธสุ, อตฺตทณฺเฑสุ นิพฺพุตา. \\ สาทาเนสุ อนาทานา, เต นมสฺสามิ มาตลีติ.}\csromangathastanza{เสฏฺฐา หิ กิร โลกสฺมิํ, เย ตฺวํ สกฺก นมสฺสสิ. \\ อหมฺปิ เต นมสฺสามิ, เย นมสฺสสิ วาสวาติ.}\csromangathastanza{อิทํ วตฺวาน มฆวา, เทวราชา สุชมฺปติ. \\ ภิกฺขุสํฆํ นมสฺสิตฺวา, ปมุโข รถมารุหีติ.}}

\vspace{\tipitakaparskip}
\csromancenter{ทุติโย วคฺโค.}
\titlehead{\textbf{ตสฺสุทฺทานํ}}
\csromangathameasure{เทวา ปน ตโย วุตฺตา, ทลิทฺทญฺจ รามเณยฺยกํ. \\ ยชมานญฺจ วนฺทนา, ตโย สกฺกนมสฺสนาติ.}
\csromangathagroup{\csromangathastanza{เทวา ปน\footnote{วตปเทน (สี, สฺยา, กํ)} ตโย วุตฺตา, ทลิทฺทญฺจ รามเณยฺยกํ. \\ ยชมานญฺจ วนฺทนา, ตโย สกฺกนมสฺสนาติ.}}

\csromanheader{4}{11. สกฺกสํยุตฺต \textbf{3. ตติยวคฺค}}
\csromanheader{2}{1. เฉตฺวาสุตฺต}
\proseitem{267}{\csromanfolio{239}สาวตฺถิยํ เชตวเน.\csromanspacer{} อถ โข สกฺโก เทวานมินฺโท เยน ภควา เตนุปสงฺกมิ, อุปสงฺกมิตฺวา ภควนฺตํ อภิวาเทตฺวา เอกมนฺตํ อฏฺฐาสิ, เอกมนฺตํ ฐิโต โข สกฺโก เทวานมินฺโท ภควนฺตํ คาถาย อชฺฌภาสิ\csromandash{}}

\csromangathameasure{“กิํสุ เฉตฺวา สุขํ เสติ, กิํสุ เฉตฺวา น โสจติ. \\ กิสฺสสฺสุ เอกธมฺมสฺส, วธํ โรเจสิ โคตมาติ. \\ โกธํ เฉตฺวา สุขํ เสติ, โกธํ เฉตฺวา น โสจติ. \\ โกธสฺส วิสมูลสฺส, มธุรคฺคสฺส วาสว. \\ วธํ อริยา ปสํสนฺติ, ตํ หิ เฉตฺวา น โสจตี”ติ.}
\csromangathagroup{\csromangathastanza{“กิํสุ เฉตฺวา สุขํ เสติ, กิํสุ เฉตฺวา น โสจติ. \\ กิสฺสสฺสุ เอกธมฺมสฺส, วธํ โรเจสิ โคตมาติ.}\csromangathastanza{โกธํ เฉตฺวา สุขํ เสติ, โกธํ เฉตฺวา น โสจติ. \\ โกธสฺส วิสมูลสฺส, มธุรคฺคสฺส วาสว. \\ วธํ อริยา ปสํสนฺติ, ตํ หิ เฉตฺวา น โสจตี”ติ.}}

\csromanheader{2}{2. ทุพฺพณฺณิยสุตฺต}
\proseitem{268}{สาวตฺถิยํ เชตวเน.\csromanspacer{} ตตฺร โข ฯเปฯ เอตทโวจ\csromandash{}ภูตปุพฺพํ ภิกฺขเว อญฺญตโร ยกฺโข ทุพฺพณฺโณ โอโกฏิมโก สกฺกสฺส เทวานมินฺทสฺส อาสเน นิสินฺโน อโหสิ.\csromanspacer{} ตตฺร สุทํ ภิกฺขเว เทวา ตาวติํสา อุชฺฌายนฺติ ขิยฺยนฺติ วิปาเจนฺติ “อจฺฉริยํ วต โภ, อพฺภุตํ วต โภ, อยํ ยกฺโข ทุพฺพณฺโณ โอโกฏิมโก สกฺกสฺส เทวานมินฺทสฺส อาสเน นิสินฺโน”ติ.\csromanspacer{} ยถา ยถา โข ภิกฺขเว เทวา ตาวติํสา อุชฺฌายนฺติ ขิยฺยนฺติ วิปาเจนฺติ, ตถา ตถา โส ยกฺโข อภิรูปตโร เจว โหติ ทสฺสนียตโร จ ปาสาทิกตโร จ.}

\prose{อถ โข ภิกฺขเว เทวา ตาวติํสา เยน สกฺโก เทวานมินฺโท เตนุปสงฺกมิํสุ, อุปสงฺกมิตฺวา สกฺกํ เทวานมินฺทํ เอตทโวจุํ “อิธ เต มาริส อญฺญตโร ยกฺโข ทุพฺพณฺโณ โอโกฏิมโก สกฺกสฺส เทวานมินฺทสฺส อาสเน นิสินฺโน.\csromanspacer{} ตตฺร สุทํ มาริส เทวา ตาวติํสา อุชฺฌายนฺติ ขิยฺยนฺติ วิปาเจนฺติ ‘อจฺฉริยํ วต โภ, อพฺภุตํ วต โภ, อยํ ยกฺโข ทุพฺพณฺโณ โอโกฏิมโก สกฺกสฺส เทวานมินฺทสฺส อาสเน นิสินฺโน’ติ.\csromanspacer{} ยถา ยถา โข มาริส เทวา อุชฺฌายนฺติ ขิยฺยนฺติ วิปาเจนฺติ, ตถา ตถา โส ยกฺโข อภิรูปตโร เจว โหติ \csromanfolio{240}ทสฺสนียตโร จ ปาสาทิกตโร จาติ.\csromanspacer{} โส หิ นุน มาริส โกธภกฺโข ยกฺโข ภวิสฺสตี”ติ.}

\prose{อถ โข ภิกฺขเว สกฺโก เทวานมินฺโท เยน โส โกธภกฺโข ยกฺโข เตนุปสงฺกมิ, อุปสงฺกมิตฺวา เอกํสํ อุตฺตราสงฺคํ กริตฺวา ทกฺขิณชาณุมณฺฑลํ ปถวิยํ นิหนฺตฺวา เยน โส โกธภกฺโข ยกฺโข เตนญฺชลิํ ปณาเมตฺวา ติกฺขตฺตุํ นามํ สาเวติ “สกฺโกหํ มาริส เทวานมินฺโท, สกฺโกหํ มาริส เทวานมินฺโท”ติ.\csromanspacer{} ยถา ยถา โข ภิกฺขเว สกฺโก เทวานมินฺโท นามํ สาเวสิ, ตถา ตถา โส ยกฺโข ทุพฺพณฺณตโร เจว อโหสิ โอโกฏิมกตโร จ.\csromanspacer{} ทุพฺพณฺณตโร เจว หุตฺวา โอโกฏิมกตโร จ ตตฺเถวนฺตรธายีติ.\csromanspacer{} อถ โข ภิกฺขเว สกฺโก เทวานมินฺโท สเก อาสเน นิสีทิตฺวา เทเว ตาวติํเส อนุนยมาโน ตายํ เวลายํ อิมา คาถาโย อภาสิ\csromandash{}}

\csromangathameasure{“น สูปหตจิตฺโตมฺหิ, นาวตฺเตน สุวานโย. \\ น โว จิราหํ กุชฺฌามิ, โกโธ มยิ นาวติฏฺฐติ. \\ กุทฺธาหํ น ผรุสํ พฺรูมิ, น จ ธมฺมานิ กิตฺตเย. \\ สนฺนิคฺคณฺหามิ อตฺตานํ, สมฺปสฺสํ อตฺถมตฺตโน”ติ.}
\csromangathagroup{\csromangathastanza{“น สูปหตจิตฺโตมฺหิ, นาวตฺเตน สุวานโย. \\ น โว จิราหํ กุชฺฌามิ, โกโธ มยิ นาวติฏฺฐติ.}\csromangathastanza{กุทฺธาหํ น ผรุสํ พฺรูมิ, น จ ธมฺมานิ กิตฺตเย. \\ สนฺนิคฺคณฺหามิ อตฺตานํ, สมฺปสฺสํ อตฺถมตฺตโน”ติ.}}

\csromanheader{2}{3. สมฺพริมายาสุตฺต}
\proseitem{269}{สาวตฺถิยํ ฯเปฯ ภควา เอตทโวจ\csromandash{}ภูตปุพฺพํ ภิกฺขเว เวปจิตฺติ อสุรินฺโท อาพาธิโก อโหสิ ทุกฺขิโต พาฬฺหคิลาโน, อถ โข ภิกฺขเว สกฺโก เทวานมินฺโท เยน เวปจิตฺติ อสุรินฺโท เตนุปสงฺกมิ คิลานปุจฺฉโก.\csromanspacer{} อทฺทสา โข ภิกฺขเว เวปจิตฺติ อสุรินฺโท สกฺกํ เทวานมินฺทํ ทูรโตว อาคจฺฉนฺตํ, ทิสฺวาน สกฺกํ เทวานมินฺทํ เอตทโวจ “ติกิจฺฉ มํ เทวานมินฺทา”ติ.\csromanspacer{} วาเจหิ มํ เวปจิตฺติ สมฺพริมายนฺติ.\csromanspacer{} น ตาวาหํ วาเจมิ, ยาวาหํ มาริส อสุเร ปฏิปุจฺฉามีติ.\csromanspacer{} อถ โข ภิกฺขเว เวปจิตฺติ อสุรินฺโท อสุเร ปฏิปุจฺฉิ “วาเจมหํ มาริสา สกฺกํ เทวานมินฺทํ สมฺพริมายนฺ”ติ.\csromanspacer{} มา โข ตฺวํ มาริส วาเจสิ สกฺกํ เทวานมินฺทํ สมฺพริมายนฺติ.\csromanspacer{} อถ โข ภิกฺขเว เวปจิตฺติ อสุรินฺโท สกฺกํ เทวานมินฺทํ คาถาย อชฺฌภาสิ\csromandash{}}

\csromanheader{2}{11. สกฺกสํยุตฺต}
\csromangathameasure{“มายาวี มฆวา สกฺก, เทวราช สุชมฺปติ. \\ อุเปติ นิรยํ โฆรํ, สมฺพโรว สตํ สมนฺ”ติ.}
\csromangathagroup{\csromangathastanza{\csromanfolio{241}“มายาวี มฆวา สกฺก, เทวราช สุชมฺปติ. \\ อุเปติ นิรยํ โฆรํ, สมฺพโรว สตํ สมนฺ”ติ.}}

\csromanheader{2}{4. อจฺจยสุตฺต}
\proseitem{270}{สาวตฺถิยํ ฯเปฯ อาราเม.\csromanspacer{} เตน โข ปน สมเยน เทฺว ภิกฺขู สมฺปโยเชสุํ, ตเตฺรโก ภิกฺขุ อจฺจสรา.\csromanspacer{} อถ โข โส ภิกฺขุ ตสฺส ภิกฺขุโน สนฺติเก อจฺจยํ อจฺจยโต เทเสติ, โส ภิกฺขุ นปฺปฏิคฺคณฺหาติ.\csromanspacer{} อถ โข สมฺพหุลา ภิกฺขู เยน ภควา เตนุปสงฺกมิํสุ, อุปสงฺกมิตฺวา ภควนฺตํ อภิวาเทตฺวา เอกมนฺตํ นิสีทิํสุ, เอกมนฺตํ นิสินฺนา โข เต ภิกฺขู ภควนฺตํ เอตทโวจุํ “อิธ ภนฺเต เทฺว ภิกฺขู สมฺปโยเชสุํ, ตเตฺรโก ภิกฺขุ อจฺจสรา.\csromanspacer{} อถ โข โส ภนฺเต ภิกฺขุ ตสฺส ภิกฺขุโน สนฺติเก อจฺจยํ อจฺจยโต เทเสติ, โส ภิกฺขุ นปฺปฏิคฺคณฺหาตี”ติ.}

\prose{เทฺวเม ภิกฺขเว พาลา, โย จ อจฺจยํ อจฺจยโต น ปสฺสติ, โย จ อจฺจยํ เทเสนฺตสฺส ยถาธมฺมํ นปฺปฏิคฺคณฺหาติ.\csromanspacer{} อิเม โข ภิกฺขเว เทฺว พาลา.\csromanspacer{} เทฺวเม ภิกฺขเว ปณฺฑิตา, โย จ อจฺจยํ อจฺจยโต ปสฺสติ, โย จ อจฺจยํ เทเสนฺตสฺส ยถาธมฺมํ ปฏิคฺคณฺหาติ.\csromanspacer{} อิเม โข ภิกฺขเว เทฺว ปณฺฑิตา.}

\prose{ภูตปุพฺพํ ภิกฺขเว สกฺโก เทวานมินฺโท สุธมฺมายํ สภายํ เทเว ตาวติํเส อนุนยมาโน ตายํ เวลายํ อิมํ คาถํ อภาสิ\csromandash{}}

\csromangathameasure{“โกโธ โว วสมายาตุ, มา จ มิตฺเตหิ โว ชรา. \\ อครหิยํ มา ครหิตฺถ, มา จ ภาสิตฺถ เปสุณํ. \\ อถ ปาปชนํ โกโธ, ปพฺพโตวาภิมทฺทตี”ติ.}
\csromangathagroup{\csromangathastanza{“โกโธ โว วสมายาตุ, มา จ มิตฺเตหิ โว ชรา. \\ อครหิยํ มา ครหิตฺถ, มา จ ภาสิตฺถ เปสุณํ. \\ อถ ปาปชนํ โกโธ, ปพฺพโตวาภิมทฺทตี”ติ.}}

\csromanheader{2}{5. อกฺโกธสุตฺต}
\proseitem{271}{เอวํ เม สุตํ\csromandash{}เอกํ สมยํ ภควา สาวตฺถิยํ วิหรติ เชตวเน อนาถปิณฺฑิกสฺส อาราเม.\csromanspacer{} ตตฺร โข ภควา ภิกฺขู ฯเปฯ ภควา เอตทโวจ\csromandash{}ภูตปุพฺพํ ภิกฺขเว สกฺโก เทวานมินฺโท สุธมฺมายํ สภายํ เทเว ตาวติํเส อนุนยมาโน ตายํ เวลายํ อิมํ คาถํ อภาสิ\csromandash{}}

\csromangathameasure{“มา โว โกโธ อชฺฌภวิ, มา จ กุชฺฌิตฺถ กุชฺฌตํ. \\ อกฺโกโธ อวิหิํสา จ, อริเยสุ จ ปฏิปทา. \\ อถ ปาปชนํ โกโธ, ปพฺพโตวาภิมทฺทตี”ติ.}
\csromangathagroup{\csromangathastanza{\csromanfolio{242}“มา โว โกโธ อชฺฌภวิ, มา จ กุชฺฌิตฺถ กุชฺฌตํ. \\ อกฺโกโธ อวิหิํสา จ, อริเยสุ จ ปฏิปทา\footnote{วสตี สทา (สี, สฺยา, กํ, อิ)}. \\ อถ ปาปชนํ โกโธ, ปพฺพโตวาภิมทฺทตี”ติ.}}

\vspace{\tipitakaparskip}
\csromancenter{ตติโย วคฺโค.}
\titlehead{\textbf{ตสฺสุทฺทานํ}}
\csromangathameasure{เฉตฺวา ทุพฺพณฺณิย มายา, อจฺจเยน อโกธโน. \\ เทสิตํ พุทฺธเสฏฺเฐน, อิทํ หิ สกฺกปญฺจกนฺติ. \\ สกฺกสํยุตฺตํ สมตฺตํ. สคาถาวคฺโค ปฐโม.}
\csromangathagroup{\csromangathastanza{เฉตฺวา ทุพฺพณฺณิย มายา, อจฺจเยน อโกธโน. \\ เทสิตํ พุทฺธเสฏฺเฐน, อิทํ หิ สกฺกปญฺจกนฺติ. \\ สกฺกสํยุตฺตํ สมตฺตํ. สคาถาวคฺโค ปฐโม.}}

\titlehead{\textbf{ตสฺสุทฺทานํ}}
\csromangathameasure{เทวตา เทวปุตฺโต จ, ราชา มาโร จ ภิกฺขุนี. \\ พฺรหฺมา พฺราหฺมณ วงฺคีโส, วนยกฺเขน วาสโวติ.}
\csromangathagroup{\csromangathastanza{เทวตา เทวปุตฺโต จ, ราชา มาโร จ ภิกฺขุนี. \\ พฺรหฺมา พฺราหฺมณ วงฺคีโส, วนยกฺเขน วาสโวติ.}}

\nitthitam{\textbf{สคาถาวคฺคสํยุตฺตปาฬิ นิฏฺฐิตา.}}
\csromanheader{2}{\textbf{สํยุตฺตนิกาย}}
\gambhira{\textbf{นิทานวคฺคสํยุตฺตปาฬิ}}
\namakkaram{นโม ตสฺส ภควโต อรหโต สมฺมาสมฺพุทฺธสฺส.}
\chapterhead{1. นิทานสํยุตฺต}
\csromanheader{1}{1. พุทฺธวคฺค}
\csromanheader{2}{1. ปฏิจฺจสมุปฺปาทสุตฺต}
\proseitem{1}{\csromanfolio{243}เอวํ เม สุตํ\csromandash{}เอกํ สมยํ ภควา สาวตฺถิยํ วิหรติ เชตวเน อนาถปิณฺฑิกสฺส อาราเม.\csromanspacer{} ตตฺร โข ภควา ภิกฺขู อามนฺเตสิ “ภิกฺขโว”ติ. “ภทนฺเต”ติ เต ภิกฺขู ภควโต ปจฺจสฺโสสุํ.\csromanspacer{} ภควา เอตทโวจ “ปฏิจฺจสมุปฺปาทํ โว ภิกฺขเว เทเสสฺสามิ, ตํ สุณาถ สาธุกํ มนสิ กโรถ, ภาสิสฺสามี”ติ. “เอวํ ภนฺเต”ติ โข เต ภิกฺขู ภควโต ปจฺจสฺโสสุํ.\csromanspacer{} ภควา เอตทโวจ\csromandash{}}

\prose{กตโม จ ภิกฺขเว ปฏิจฺจสมุปฺปาโท.\csromanspacer{} อวิชฺชาปจฺจยา ภิกฺขเว สงฺขารา, สงฺขารปจฺจยา วิญฺญาณํ, วิญฺญาณปจฺจยา นามรูปํ, นามรูปปจฺจยา สฬายตนํ, สฬายตนปจฺจยา ผสฺโส, ผสฺสปจฺจยา เวทนา, เวทนาปจฺจยา ตณฺหา, ตณฺหาปจฺจยา อุปาทานํ, อุปาทานปจฺจยา ภโว, ภวปจฺจยา ชาติ, ชาติปจฺจยา ชรามรณํ โสกปริเทวทุกฺขโทมนสฺสุปายาสา สมฺภวนฺติ, เอวเมตสฺส เกวลสฺส ทุกฺขกฺขนฺธสฺส สมุทโย โหติ.\csromanspacer{} อยํ วุจฺจติ ภิกฺขเว ปฏิจฺจสมุปฺปาโท.}

\csromanmatikaanchor{mtk.1}
\csromanmatikaanchor{mtk.214}
\csromantocmarknopage{chapter}{1. นิทานสํยุตฺต}{mtk.1}
\bookmark[dest={mtk.1},level=0]{1. นิทานสํยุตฺต}
\gambhira{นิทานวคฺคสํยุตฺตปาฬิ}
\prose{\csromanfolio{244}อวิชฺชาย เตฺวว อเสสวิราคนิโรธา สงฺขารนิโรโธ, สงฺขารนิโรธา วิญฺญาณนิโรโธ, วิญฺญาณนิโรธา นามรูปนิโรโธ, นามรูปนิโรธา สฬายตนนิโรโธ, สฬายตนนิโรธา ผสฺสนิโรโธ, ผสฺสนิโรธา เวทนานิโรโธ, เวทนานิโรธา ตณฺหานิโรโธ, ตณฺหานิโรธา อุปาทานนิโรโธ, อุปาทานนิโรธา ภวนิโรโธ, ภวนิโรธา ชาตินิโรโธ, ชาตินิโรธา ชรามรณํ โสกปริเทวทุกฺขโทมนสฺสุปายาสา นิรุชฺฌนฺติ, เอวเมตสฺส เกวลสฺส ทุกฺขกฺขนฺธสฺส นิโรโธ โหตีติ.\csromanspacer{} อิทมโวจ ภควา.\csromanspacer{} อตฺตมนา เต ภิกฺขู ภควโต ภาสิตํ อภินนฺทุนฺติ.\csromanspacer{}\csromanspacer{}ปฐมํ.}

\csromanheader{2}{2. วิภงฺคสุตฺต}
\proseitem{2}{สาวตฺถิยํ วิหรติ ฯเปฯ.\csromanspacer{} ปฏิจฺจสมุปฺปาทํ โว ภิกฺขเว เทเสสฺสามิ วิภชิสฺสามิ, ตํ สุณาถ สาธุกํ มนสิ กโรถ, ภาสิสฺสามีติ. “เอวํ ภนฺเต”ติ โข เต ภิกฺขู ภควโต ปจฺจสฺโสสุํ.\csromanspacer{} ภควา เอตทโวจ\csromandash{}}

\prose{กตโม จ ภิกฺขเว ปฏิจฺจสมุปฺปาโท.\csromanspacer{} อวิชฺชาปจฺจยา ภิกฺขเว สงฺขารา, สงฺขารปจฺจยา วิญฺญาณํ, วิญฺญาณปจฺจยา นามรูปํ, นามรูปปจฺจยา สฬายตนํ, สฬายตนปจฺจยา ผสฺโส, ผสฺสปจฺจยา เวทนา, เวทนาปจฺจยา ตณฺหา, ตณฺหาปจฺจยา อุปาทานํ, อุปาทานปจฺจยา ภโว, ภวปจฺจยา ชาติ, ชาติปจฺจยา ชรามรณํ โสกปริเทวทุกฺขโทมนสฺสุปายาสา สมฺภวนฺติ, เอวเมตสฺส เกวลสฺส ทุกฺขกฺขนฺธสฺส สมุทโย โหติ.}

\prose{กตมญฺจ ภิกฺขเว ชรามรณํ.\csromanspacer{} ยา เตสํ เตสํ สตฺตานํ ตมฺหิ ตมฺหิ สตฺตนิกาเย ชรา ชีรณตา ขณฺฑิจฺจํ ปาลิจฺจํ วลิตฺตจตา อายุโน สํหานิ อินฺทฺริยานํ ปริปาโก, อยํ วุจฺจติ ชรา.\csromanspacer{} ยา เตสํ เตสํ สตฺตานํ ตมฺหา ตมฺหา สตฺตนิกายา จุติ จวนตา เภโท อนฺตรธานํ มจฺจุ มรณํ กาลกิริยา ขนฺธานํ เภโท กเฬวรสฺส นิกฺเขโป ( )\footnote{(ชีวิตินฺทฺริยสฺส อุปจฺเฉโท) (สฺยา, กํ) เอวมุปริปิ, อฏฺฐกถายํ ปน น ทิสฺสติ.}, อิทํ วุจฺจติ มรณํ.\csromanspacer{} อิติ อยญฺจ ชรา อิทญฺจ มรณํ, อิทํ วุจฺจติ ภิกฺขเว ชรามรณํ.}

\prose{กตมา จ ภิกฺขเว ชาติ.\csromanspacer{} ยา เตสํ เตสํ สตฺตานํ ตมฺหิ ตมฺหิ สตฺตนิกาเย ชาติ สญฺชาติ โอกฺกนฺติ อภินิพฺพตฺติ ขนฺธานํ ปาตุภาโว อายตนานํ ปฏิลาโภ, อยํ วุจฺจติ ภิกฺขเว ชาติ.}

\csromanheader{2}{1. นิทานสํยุตฺต}
\prose{\csromanfolio{245}กตโม จ ภิกฺขเว ภโว.\csromanspacer{} ตโยเม ภิกฺขเว ภวา, กามภโว รูปภโว อรูปภโว, อยํ วุจฺจติ ภิกฺขเว ภโว.}

\prose{กตมญฺจ ภิกฺขเว อุปาทานํ.\csromanspacer{} จตฺตาริมานิ ภิกฺขเว อุปาทานานิ, กามุปาทานํ ทิฏฺฐุปาทานํ สีลพฺพตุปาทานํ อตฺตวาทุปาทานํ, อิทํ วุจฺจติ ภิกฺขเว อุปาทานํ.}

\prose{กตมา จ ภิกฺขเว ตณฺหา.\csromanspacer{} ฉยิเม ภิกฺขเว ตณฺหากายา, รูปตณฺหา สทฺทตณฺหา คนฺธตณฺหา รสตณฺหา โผฏฺฐพฺพตณฺหา ธมฺมตณฺหา, อยํ วุจฺจติ ภิกฺขเว ตณฺหา.}

\prose{กตมา จ ภิกฺขเว เวทนา.\csromanspacer{} ฉยิเม ภิกฺขเว เวทนากายา, จกฺขุสมฺผสฺสชา เวทนา โสตสมฺผสฺสชา เวทนา ฆานสมฺผสฺสชา เวทนา ชิวฺหาสมฺผสฺสชา เวทนา กายสมฺผสฺสชา เวทนา มโนสมฺผสฺสชา เวทนา, อยํ วุจฺจติ ภิกฺขเว เวทนา.}

\prose{กตโม จ ภิกฺขเว ผสฺโส.\csromanspacer{} ฉยิเม ภิกฺขเว ผสฺสกายา, จกฺขุสมฺผสฺโส โสตสมฺผสฺโส ฆานสมฺผสฺโส ชิวฺหาสมฺผสฺโส กายสมฺผสฺโส มโนสมฺผสฺโส, อยํ วุจฺจติ ภิกฺขเว ผสฺโส.}

\prose{กตมญฺจ ภิกฺขเว สฬายตนํ.\csromanspacer{} จกฺขายตนํ โสตายตนํ ฆานายตนํ ชิวฺหายตนํ กายายตนํ มนายตนํ, อิทํ วุจฺจติ ภิกฺขเว สฬายตนํ.}

\prose{กตมญฺจ ภิกฺขเว นามรูปํ.\csromanspacer{} เวทนา สญฺญา เจตนา ผสฺโส มนสิกาโร, อิทํ วุจฺจติ นามํ.\csromanspacer{} จตฺตาโร จ มหาภูตา จตุนฺนญฺจ มหาภูตานํ อุปาทายรูปํ, อิทํ วุจฺจติ รูปํ.\csromanspacer{} อิติ อิทญฺจ นามํ อิทญฺจ รูปํ, อิทํ วุจฺจติ ภิกฺขเว นามรูปํ.}

\prose{กตมญฺจ ภิกฺขเว วิญฺญาณํ.\csromanspacer{} ฉยิเม ภิกฺขเว วิญฺญาณกายา, จกฺขุวิญฺญาณํ โสตวิญฺญาณํ ฆานวิญฺญาณํ ชิวฺหาวิญฺญาณํ กายวิญฺญาณํ มโนวิญฺญาณํ, อิทํ วุจฺจติ ภิกฺขเว วิญฺญาณํ.}

\prose{กตเม จ ภิกฺขเว สงฺขารา.\csromanspacer{} ตโยเม ภิกฺขเว สงฺขารา, กายสงฺขาโร วจีสงฺขาโร จิตฺตสงฺขาโร, อิเม วุจฺจนฺติ ภิกฺขเว สงฺขารา.}

\prose{กตมา จ ภิกฺขเว อวิชฺชา.\csromanspacer{} ยํ โข ภิกฺขเว ทุกฺเข อญฺญาณํ ทุกฺขสมุทเย อญฺญาณํ ทุกฺขนิโรเธ อญฺญาณํ ทุกฺขนิโรธคามินิยา ปฏิปทาย อญฺญาณํ, อยํ วุจฺจติ ภิกฺขเว อวิชฺชา.}

\gambhira{นิทานวคฺคสํยุตฺตปาฬิ}
\prose{\csromanfolio{246}อิติ โข ภิกฺขเว อวิชฺชาปจฺจยา สงฺขารา, สงฺขารปจฺจยา วิญฺญาณํ ฯเปฯ เอวเมตสฺส เกวลสฺส ทุกฺขกฺขนฺธสฺส สมุทโย โหติ.\csromanspacer{} อวิชฺชายเตฺวว อเสสวิราคนิโรธา สงฺขารนิโรโธ, สงฺขารนิโรโธ วิญฺญาณนิโรโธ ฯเปฯ เอวเมตสฺส เกวลสฺส ทุกฺขกฺขนฺธสฺส นิโรโธ โหตีติ.\csromanspacer{}\csromanspacer{}ทุติยํ.}

\csromanheader{2}{3. ปฏิปทาสุตฺต}
\proseitem{3}{สาวตฺถิยํ วิหรติ ฯเปฯ.\csromanspacer{} มิจฺฉาปฏิปทญฺจ โว ภิกฺขเว เทเสสฺสามิ สมฺมาปฏิปทญฺจ, ตํ สุณาถ สาธุกํ มนสิ กโรถ, ภาสิสฺสามีติ. “เอวํ ภนฺเต”ติ โข เต ภิกฺขู ภควโต ปจฺจสฺโสสุํ.\csromanspacer{} ภควา เอตทโวจ\csromandash{}}

\prose{กตมา จ ภิกฺขเว มิจฺฉาปฏิปทา.\csromanspacer{} อวิชฺชาปจฺจยา ภิกฺขเว สงฺขารา, สงฺขารปจฺจยา วิญฺญาณํ ฯเปฯ เอวเมตสฺส เกวลสฺส ทุกฺขกฺขนฺธสฺส สมุทโย โหติ, อยํ วุจฺจติ ภิกฺขเว มิจฺฉาปฏิปทา.}

\prose{กตมา จ ภิกฺขเว สมฺมาปฏิปทา.\csromanspacer{} อวิชฺชาย เตฺวว อเสสวิราคนิโรธา สงฺขารนิโรโธ, สงฺขารนิโรธา วิญฺญาณนิโรโธ ฯเปฯ เอวเมตสฺส เกวลสฺส ทุกฺขกฺขนฺธสฺส นิโรโธ โหติ, อยํ วุจฺจติ ภิกฺขเว สมฺมาปฏิปทาติ.\csromanspacer{}\csromanspacer{}ตติยํ.}

\csromanheader{2}{4. วิปสฺสีสุตฺต}
\proseitem{4}{สาวตฺถิยํ วิหรติ ฯเปฯ.\csromanspacer{} วิปสฺสิสฺส ภิกฺขเว ภควโต อรหโต สมฺมาสมฺพุทฺธสฺส ปุพฺเพว สมฺโพธา อนภิสมฺพุทฺธสฺส โพธิสตฺตสฺเสว สโต เอตทโหสิ “กิจฺฉํ วตายํ โลโก อาปนฺโน ชายติ จ ชียติ จ มียติ จ จวติ จ อุปปชฺชติ จ, อถ จ ปนิมสฺส ทุกฺขสฺส นิสฺสรณํ นปฺปชานาติ ชรามรณสฺส.\csromanspacer{} กุทาสฺสุ นาม อิมสฺส ทุกฺขสฺส นิสฺสรณํ ปญฺญายิสฺสติ ชรามรณสฺสา”ติ.}

\prose{อถ โข ภิกฺขเว วิปสฺสิสฺส โพธิสตฺตสฺส เอตทโหสิ “กิมฺหิ นุ โข สติ ชรามรณํ โหติ, กิํปจฺจยา ชรามรณนฺ”ติ.\csromanspacer{} อถ โข ภิกฺขเว วิปสฺสิสฺส โพธิสตฺตสฺส โยนิโส มนสิการา อหุ ปญฺญาย อภิสมโย “ชาติยา โข สติ ชรามรณํ โหติ, ชาติปจฺจยา ชรามรณนฺ”ติ.}

\csromanheader{2}{1. นิทานสํยุตฺต}
\prose{\csromanfolio{247}อถ โข ภิกฺขเว วิปสฺสิสฺส โพธิสตฺตสฺส เอตทโหสิ “กิมฺหิ นุ โข สติ ชาติ โหติ, กิํปจฺจยา ชาตี”ติ.\csromanspacer{} อถ โข ภิกฺขเว วิปสฺสิสฺส โพธิสตฺตสฺส โยนิโส มนสิการา อหุ ปญฺญาย อภิสมโย “ภเว โข สติ ชาติ โหติ, ภวปจฺจยา ชาตี”ติ.}

\prose{อถ โข ภิกฺขเว วิปสฺสิสฺส โพธิสตฺตสฺส เอตทโหสิ “กิมฺหิ นุ โข สติ ภโว โหติ, กิํปจฺจยา ภโว”ติ.\csromanspacer{} อถ โข ภิกฺขเว วิปสฺสิสฺส โพธิสตฺตสฺส โยนิโส มนสิการา อหุ ปญฺญาย อภิสมโย “อุปาทาเน โข สติ ภโว โหติ, อุปาทานปจฺจยา ภโว”ติ.}

\prose{อถ โข ภิกฺขเว วิปสฺสิสฺส โพธิสตฺตสฺส เอตทโหสิ “กิมฺหิ นุ โข สติ อุปาทานํ โหติ, กิํปจฺจยา อุปาทานนฺ”ติ.\csromanspacer{} อถ โข ภิกฺขเว วิปสฺสิสฺส โพธิสตฺตสฺส โยนิโส มนสิการา อหุ ปญฺญาย อภิสมโย “ตณฺหาย โข สติ อุปาทานํ โหติ, ตณฺหาปจฺจยา อุปาทานนฺ”ติ.}

\prose{อถ โข ภิกฺขเว วิปสฺสิสฺส โพธิสตฺตสฺส เอตทโหสิ “กิมฺหิ นุ โข สติ ตณฺหา โหติ, กิํปจฺจยา ตณฺหา”ติ.\csromanspacer{} อถ โข ภิกฺขเว วิปสฺสิสฺส โพธิสตฺตสฺส โยนิโส มนสิการา อหุ ปญฺญาย อภิสมโย “เวทนาย โข สติ ตณฺหา โหติ, เวทนาปจฺจยา ตณฺหา”ติ.}

\prose{อถ โข ภิกฺขเว วิปสฺสิสฺส โพธิสตฺตสฺส เอตทโหสิ “กิมฺหิ นุ โข สติ เวทนา โหติ, กิํปจฺจยา เวทนา”ติ.\csromanspacer{} อถ โข ภิกฺขเว วิปสฺสิสฺส โพธิสตฺตสฺส โยนิโส มนสิการา อหุ ปญฺญาย อภิสมโย “ผสฺเส โข สติ เวทนา โหติ, ผสฺสปจฺจยา เวทนา”ติ.}

\prose{อถ โข ภิกฺขเว วิปสฺสิสฺส โพธิสตฺตสฺส เอตทโหสิ “กิมฺหิ นุ โข สติ ผสฺโส โหติ, กิํปจฺจยา ผสฺโส”ติ.\csromanspacer{} อถ โข ภิกฺขเว วิปสฺสิสฺส โพธิสตฺตสฺส โยนิโส มนสิการา อหุ ปญฺญาย อภิสมโย “สฬายตเน โข สติ ผสฺโส โหติ, สฬายตนปจฺจยา ผสฺโส”ติ.}

\csromanmatikaanchor{mtk.2}
\csromantocmarknopage{section}{1. พุทฺธวคฺค}{mtk.2}
\bookmark[dest={mtk.2},level=1]{1. พุทฺธวคฺค}
\gambhira{นิทานวคฺคสํยุตฺตปาฬิ}
\csromanmatikaanchor{mtk.3}
\csromantocmarkref{subsection}{1. ปฏิจฺจสมุปฺปาทสุตฺต}{mtk.3}
\bookmark[dest={mtk.3},level=2]{1. ปฏิจฺจสมุปฺปาทสุตฺต}
\prose{\csromanfolio{248}อถ โข ภิกฺขเว วิปสฺสิสฺส โพธิสตฺตสฺส เอตทโหสิ “กิมฺหิ นุ โข สติ สฬายตนํ โหติ, กิํปจฺจยา สฬายตนนฺ”ติ.\csromanspacer{} อถ โข ภิกฺขเว วิปสฺสิสฺส โพธิสตฺตสฺส โยนิโส มนสิการา อหุ ปญฺญาย อภิสมโย “นามรูเป โข สติ สฬายตนํ โหติ, นามรูปปจฺจยา สฬายตนนฺ”ติ.}

\prose{อถ โข ภิกฺขเว วิปสฺสิสฺส โพธิสตฺตสฺส เอตทโหสิ “กิมฺหิ นุ โข สติ นามรูปํ โหติ, กิํปจฺจยา นามรูปนฺ”ติ.\csromanspacer{} อถ โข ภิกฺขเว วิปสฺสิสฺส โพธิสตฺตสฺส โยนิโส มนสิการา อหุ ปญฺญาย อภิสมโย “วิญฺญาเณ โข สติ นามรูปํ โหติ, วิญฺญาณปจฺจยา นามรูปนฺ”ติ.}

\prose{อถ โข ภิกฺขเว วิปสฺสิสฺส โพธิสตฺตสฺส เอตทโหสิ “กิมฺหิ นุ โข สติ วิญฺญาณํ โหติ, กิํปจฺจยา วิญฺญาณนฺ”ติ.\csromanspacer{} อถ โข ภิกฺขเว วิปสฺสิสฺส โพธิสตฺตสฺส โยนิโส มนสิการา อหุ ปญฺญาย อภิสมโย “สงฺขาเรสุ โข สติ วิญฺญาณํ โหติ, สงฺขารปจฺจยา วิญฺญาณนฺ”ติ.}

\prose{อถ โข ภิกฺขเว วิปสฺสิสฺส โพธิสตฺตสฺส เอตทโหสิ “กิมฺหิ นุ โข สติ สงฺขารา โหนฺติ, กิํปจฺจยา สงฺขารา”ติ.\csromanspacer{} อถ โข ภิกฺขเว วิปสฺสิสฺส โพธิสตฺตสฺส โยนิโส มนสิการา อหุ ปญฺญาย อภิสมโย “อวิชฺชาย โข สติ สงฺขารา โหนฺติ, อวิชฺชาปจฺจยา สงฺขารา”ติ.}

\prose{อิติ หิทํ อวิชฺชาปจฺจยา สงฺขารา, สงฺขารปจฺจยา วิญฺญาณํ ฯเปฯ เอวเมตสฺส เกวลสฺส ทุกฺขกฺขนฺธสฺส สมุทโย โหติ. “สมุทโย สมุทโย”ติ โข ภิกฺขเว วิปสฺสิสฺส โพธิสตฺตสฺส ปุพฺเพ อนนุสฺสุเตสุ ธมฺเมสุ จกฺขุํ อุทปาทิ, ญาณํ อุทปาทิ, ปญฺญา อุทปาทิ, วิชฺชา อุทปาทิ, อาโลโก อุทปาทิ.}

\csromanmatikaanchor{mtk.4}
\csromantocmarkref{subsection}{2. วิภงฺคสุตฺต}{mtk.4}
\bookmark[dest={mtk.4},level=2]{2. วิภงฺคสุตฺต}
\prose{อถ โข ภิกฺขเว วิปสฺสิสฺส โพธิสตฺตสฺส เอตทโหสิ “กิมฺหิ นุ โข อสติ ชรามรณํ น โหติ, กิสฺส นิโรธา ชรามรณนิโรโธ”ติ.\csromanspacer{} อถ โข ภิกฺขเว วิปสฺสิสฺส โพธิสตฺตสฺส โยนิโส มนสิการา อหุ ปญฺญาย อภิสมโย “ชาติยา โข อสติ ชรามรณํ น โหติ, ชาตินิโรธา ชรามรณนิโรโธ”ติ.}

\csromanheader{2}{1. นิทานสํยุตฺต}
\prose{\csromanfolio{249}อถ โข ภิกฺขเว วิปสฺสิสฺส โพธิสตฺตสฺส เอตทโหสิ “กิมฺหิ นุ โข อสติ ชาติ น โหติ, กิสฺส นิโรธา ชาตินิโรโธ”ติ.\csromanspacer{} อถ โข ภิกฺขเว วิปสฺสิสฺส โพธิสตฺตสฺส โยนิโส มนสิการา อหุ ปญฺญาย อภิสมโย “ภเว โข อสติ ชาติ น โหติ, ภวนิโรธา ชาตินิโรโธ”ติ.}

\prose{อถ โข ภิกฺขเว วิปสฺสิสฺส โพธิสตฺตสฺส เอตทโหสิ “กิมฺหิ นุ โข อสติ ภโว น โหติ, กิสฺส นิโรธา ภวนิโรโธ”ติ.\csromanspacer{} อถ โข ภิกฺขเว วิปสฺสิสฺส โพธิสตฺตสฺส โยนิโส มนสิการา อหุ ปญฺญาย อภิสมโย “อุปาทาเน โข อสติ ภโว น โหติ, อุปาทานนิโรธา ภวนิโรโธ”ติ.}

\prose{อถ โข ภิกฺขเว วิปสฺสิสฺส โพธิสตฺตสฺส เอตทโหสิ “กิมฺหิ นุ โข อสติ อุปาทานํ น โหติ, กิสฺส นิโรธา อุปาทานนิโรโธ”ติ.\csromanspacer{} อถ โข ภิกฺขเว วิปสฺสิสฺส โพธิสตฺตสฺส โยนิโส มนสิการา อหุ ปญฺญาย อภิสมโย “ตณฺหาย โข อสติ อุปาทานํ น โหติ, ตณฺหานิโรธา อุปาทานนิโรโธ”ติ.}

\prose{อถ โข ภิกฺขเว วิปสฺสิสฺส โพธิสตฺตสฺส เอตทโหสิ “กิมฺหิ นุ โข อสติ ตณฺหา น โหติ, กิสฺส นิโรธา ตณฺหานิโรโธ”ติ.\csromanspacer{} อถ โข ภิกฺขเว วิปสฺสิสฺส โพธิสตฺตสฺส โยนิโส มนสิการา อหุ ปญฺญาย อภิสมโย “เวทนาย โข อสติ ตณฺหา น โหติ, เวทนานิโรธา ตณฺหานิโรโธ”ติ.}

\prose{อถ โข ภิกฺขเว วิปสฺสิสฺส โพธิสตฺตสฺส เอตทโหสิ “กิมฺหิ นุ โข อสติ เวทนา น โหติ, กิสฺส นิโรธา เวทนานิโรโธ”ติ.\csromanspacer{} อถ โข ภิกฺขเว วิปสฺสิสฺส โพธิสตฺตสฺส โยนิโส มนสิการา อหุ ปญฺญาย อภิสมโย “ผสฺเส โข อสติ เวทนา น โหติ, ผสฺสนิโรธา เวทนานิโรโธ”ติ.}

\prose{อถ โข ภิกฺขเว วิปสฺสิสฺส โพธิสตฺตสฺส เอตทโหสิ “กิมฺหิ นุ โข อสติ ผสฺโส น โหติ, กิสฺส นิโรธา ผสฺสนิโรโธ”ติ.\csromanspacer{} อถ โข ภิกฺขเว วิปสฺสิสฺส โพธิสตฺตสฺส โยนิโส มนสิการา}

\vspace{\tipitakaparskip}
\csromancenter{นิทานวคฺคสํยุตฺตปาฬิ อหุ ปญฺญาย อภิสมโย “สฬายตเน โข อสติ ผสฺโส น โหติ, สฬายตนนิโรธา ผสฺสนิโรโธ”ติ.}
\prose{\csromanfolio{250}อถ โข ภิกฺขเว วิปสฺสิสฺส โพธิสตฺตสฺส เอตทโหสิ “กิมฺหิ นุ โข อสติ สฬายตนํ น โหติ, กิสฺส นิโรธา สฬายตนนิโรโธ”ติ.\csromanspacer{} อถ โข ภิกฺขเว วิปสฺสิสฺส โพธิสตฺตสฺส โยนิโส มนสิการา อหุ ปญฺญาย อภิสมโย “นามรูเป โข อสติ สฬายตนํ น โหติ, นามรูปนิโรธา สฬายตนนิโรโธ”ติ.}

\prose{อถ โข ภิกฺขเว วิปสฺสิสฺส โพธิสตฺตสฺส เอตทโหสิ “กิมฺหิ นุ โข อสติ นามรูปํ น โหติ, กิสฺส นิโรธา นามรูปนิโรโธ”ติ.\csromanspacer{} อถ โข ภิกฺขเว วิปสฺสิสฺส โพธิสตฺตสฺส โยนิโส มนสิการา อหุ ปญฺญาย อภิสมโย “วิญฺญาเณ โข อสติ นามรูปํ น โหติ, วิญฺญาณนิโรธา นามรูปนิโรโธ”ติ.}

\prose{อถ โข ภิกฺขเว วิปสฺสิสฺส โพธิสตฺตสฺส เอตทโหสิ “กิมฺหิ นุ โข อสติ วิญฺญาณํ น โหติ, กิสฺส นิโรธา วิญฺญาณนิโรโธ”ติ.\csromanspacer{} อถ โข ภิกฺขเว วิปสฺสิสฺส โพธิสตฺตสฺส โยนิโส มนสิการา อหุ ปญฺญาย อภิสมโย “สงฺขาเรสุ โข อสติ วิญฺญาณํ น โหติ, สงฺขารนิโรธา วิญฺญาณนิโรโธ”ติ.}

\prose{อถ โข ภิกฺขเว วิปสฺสิสฺส โพธิสตฺตสฺส เอตทโหสิ “กิมฺหิ นุ โข อสติ สงฺขารา น โหนฺติ, กิสฺส นิโรธา สงฺขารนิโรโธ”ติ.\csromanspacer{} อถ โข ภิกฺขเว วิปสฺสิสฺส โพธิสตฺตสฺส โยนิโส มนสิการา อหุ ปญฺญาย อภิสมโย “อวิชฺชาย โข อสติ สงฺขารา น โหนฺติ, อวิชฺชานิโรธา สงฺขารนิโรโธ”ติ.}

\prose{อิติ หิทํ อวิชฺชานิโรธา สงฺขารนิโรโธ, สงฺขารนิโรธา วิญฺญาณนิโรโธ ฯเปฯ เอวเมตสฺส เกวลสฺส ทุกฺขกฺขนฺธสฺส นิโรโธ โหตีติ. “นิโรโธ นิโรโธ”ติ โข ภิกฺขเว วิปสฺสิสฺส โพธิสตฺตสฺส ปุพฺเพ อนนุสฺสุเตสุ ธมฺเมสุ จกฺขุํ อุทปาทิ, ญาณํ อุทปาทิ, ปญฺญา อุทปาทิ, วิชฺชา อุทปาทิ, อาโลโก อุทปาทีติ.\csromanspacer{}\csromanspacer{}จตุตฺถํ. (สตฺตนฺนมฺปิ พุทฺธานํ เอวํ วิตฺถาเรตพฺโพ.)}

\csromanheader{2}{1. นิทานสํยุตฺต \textbf{5. สิขีสุตฺต}}
\vspace{\tipitakaparskip}
\csromancenter{สิขิสฺส ภิกฺขเว ภควโต อรหโต สมฺมาสมฺพุทฺธสฺส ฯเปฯ.}
\csromanheader{2}{6. เวสฺสภูสุตฺต}
\vspace{\tipitakaparskip}
\csromancenter{เวสฺสภุสฺส ภิกฺขเว ภควโต อรหโต สมฺมาสมฺพุทฺธสฺส ฯเปฯ.}
\csromanmatikaanchor{mtk.5}
\csromantocmarkref{subsection}{3. ปฏิปทาสุตฺต}{mtk.5}
\bookmark[dest={mtk.5},level=2]{3. ปฏิปทาสุตฺต}
\csromanheader{2}{7. กกุสนฺธสุตฺต}
\vspace{\tipitakaparskip}
\csromancenter{กกุสนฺธสฺส ภิกฺขเว ภควโต อรหโต สมฺมาสมฺพุทฺธสฺส ฯเปฯ.}
\csromanheader{2}{8. โกณาคมนสุตฺต}
\vspace{\tipitakaparskip}
\csromancenter{โกณาคมนสฺส ภิกฺขเว ภควโต อรหโต สมฺมาสมฺพุทฺธสฺส ฯเปฯ.}
\csromanmatikaanchor{mtk.6}
\csromantocmarkref{subsection}{4. วิปสฺสีสุตฺต}{mtk.6}
\bookmark[dest={mtk.6},level=2]{4. วิปสฺสีสุตฺต}
\csromanheader{2}{9. กสฺสปสุตฺต}
\vspace{\tipitakaparskip}
\csromancenter{กสฺสปสฺส ภิกฺขเว ภควโต อรหโต สมฺมาสมฺพุทฺธสฺส ฯเปฯ.}
\csromanheader{2}{10. โคตมสุตฺต}
\proseitem{10}{\csromanfolio{251}ปุพฺเพว เม ภิกฺขเว สมฺโพธา อนภิสมฺพุทฺธสฺส โพธิสตฺตสฺเสว สโต เอตทโหสิ “กิจฺฉํ วตายํ โลโก อาปนฺโน ชายติ จ ชียติ จ มียติ จ จวติ จ อุปปชฺชติ จ, อถ จ ปนิมสฺส ทุกฺขสฺส นิสฺสรณํ นปฺปชานาติ ชรามรณสฺส.\csromanspacer{} กุทาสฺสุ นาม อิมสฺส ทุกฺขสฺส นิสฺสรณํ ปญฺญายิสฺสติ ชรามรณสฺสา”ติ.}

\prose{ตสฺส มยฺหํ ภิกฺขเว เอตทโหสิ “กิมฺหิ นุ โข สติ ชรามรณํ โหติ, กิํปจฺจยา ชรามรณนฺ”ติ.\csromanspacer{} ตสฺส มยฺหํ ภิกฺขเว โยนิโส มนสิการา อหุ ปญฺญาย อภิสมโย “ชาติยา โข สติ ชรามรณํ โหติ, ชาติปจฺจยา ชรามรณนฺ”ติ.}

\prose{ตสฺส มยฺหํ ภิกฺขเว เอตทโหสิ “กิมฺหิ นุ โข สติ ชาติ โหติ ฯเปฯ ภโว.\csromanspacer{} อุปาทานํ.\csromanspacer{} ตณฺหา.\csromanspacer{} เวทนา.\csromanspacer{} ผสฺโส.\csromanspacer{} สฬายตนํ.\csromanspacer{} นามรูปํ.\csromanspacer{} วิญฺญาณํ.\csromanspacer{} สงฺขารา โหนฺติ, กิํปจฺจยา สงฺขารา”ติ.\csromanspacer{} ตสฺส มยฺหํ}

\vspace{\tipitakaparskip}
\csromancenter{นิทานวคฺคสํยุตฺตปาฬิ ภิกฺขเว โยนิโส มนสิการา อหุ ปญฺญาย อภิสมโย “อวิชฺชาย โข สติ สงฺขารา โหนฺติ, อวิชฺชาปจฺจยา สงฺขารา”ติ.}
\csromanmatikaanchor{mtk.7}
\csromanmatikaanchor{mtk.13}
\csromantocmarkref{subsection}{5. สิขีสุตฺต}{mtk.7}
\bookmark[dest={mtk.7},level=2]{5. สิขีสุตฺต}
\prose{\csromanfolio{252}อิติ หิทํ อวิชฺชาปจฺจยา สงฺขารา, สงฺขารปจฺจยา วิญฺญาณํ ฯเปฯ เอวเมตสฺส เกวลสฺส ทุกฺขกฺขนฺธสฺส สมุทโย โหติ. “สมุทโย สมุทโย”ติ โข เม ภิกฺขเว ปุพฺเพ อนนุสฺสุเตสุ ธมฺเมสุ จกฺขุํ อุทปาทิ, ญาณํ อุทปาทิ, ปญฺญา อุทปาทิ, วิชฺชา อุทปาทิ, อาโลโก อุทปาทิ.}

\prose{ตสฺส มยฺหํ ภิกฺขเว เอตทโหสิ “กิมฺหิ นุ โข อสติ ชรามรณํ น โหติ, กิสฺส นิโรธา ชรามรณนิโรโธ”ติ.\csromanspacer{} ตสฺส มยฺหํ ภิกฺขเว โยนิโส มนสิการา อหุ ปญฺญาย อภิสมโย “ชาติยา โข อสติ ชรามรณํ น โหติ, ชาตินิโรธา ชรามรณนิโรโธ”ติ.}

\prose{ตสฺส มยฺหํ ภิกฺขเว เอตทโหสิ “กิมฺหิ นุ โข อสติ ชาติ น โหติ ฯเปฯ ภโว.\csromanspacer{} อุปาทานํ.\csromanspacer{} ตณฺหา.\csromanspacer{} เวทนา.\csromanspacer{} ผสฺโส.\csromanspacer{} สฬายตนํ.\csromanspacer{} นามรูปํ.\csromanspacer{} วิญฺญาณํ.\csromanspacer{} สงฺขารา น โหนฺติ, กิสฺส นิโรธา สงฺขารนิโรโธ”ติ.\csromanspacer{} ตสฺส มยฺหํ ภิกฺขเว โยนิโส มนสิการา อหุ ปญฺญาย อภิสมโย “อวิชฺชาย โข อสติ สงฺขารา น โหนฺติ, อวิชฺชานิโรธา สงฺขารนิโรโธ”ติ.}

\prose{อิติ หิทํ อวิชฺชานิโรธา สงฺขารนิโรโธ, สงฺขารนิโรธา วิญฺญาณนิโรโธ ฯเปฯ เอวเมตสฺส เกวลสฺส ทุกฺขกฺขนฺธสฺส นิโรโธ โหติ. “นิโรโธ นิโรโธ”ติ โข เม ภิกฺขเว ปุพฺเพ อนนุสฺสุเตสุ ธมฺเมสุ จกฺขุํ อุทปาทิ, ญาณํ อุทปาทิ, ปญฺญา อุทปาทิ, วิชฺชา อุทปาทิ, อาโลโก อุทปาทีติ.\csromanspacer{}\csromanspacer{}ทสมํ.}

\vspace{\tipitakaparskip}
\csromancenter{พุทฺธวคฺโค ปฐโม.}
\titlehead{\textbf{ตสฺสุทฺทานํ}}
\csromangathameasure{เทสนา วิภงฺคปฏิปทา จ, \\ วิปสฺสี สิขี จ เวสฺสภู. \\ กกุสนฺโธ โกณาคมโน กสฺสโป, \\ มหาสกฺยมุนิ จ โคตโมติ.}
\csromangathagroup{\csromangathastanza{เทสนา วิภงฺคปฏิปทา จ, \\ วิปสฺสี สิขี จ เวสฺสภู. \\ กกุสนฺโธ โกณาคมโน กสฺสโป, \\ มหาสกฺยมุนิ จ โคตโมติ.}}

\csromanheader{1}{1. นิทานสํยุตฺต \textbf{2. อาหารวคฺค}}
\csromanheader{2}{1. อาหารสุตฺต}
\proseitem{11}{\csromanfolio{253}เอวํ เม สุตํ\csromandash{}เอกํ สมยํ ภควา สาวตฺถิยํ วิหรติ เชตวเน อนาถปิณฺฑิกสฺส อาราเม.\csromanspacer{} จตฺตาโรเม ภิกฺขเว อาหารา ภูตานํ วา สตฺตานํ ฐิติยา สมฺภเวสีนํ วา อนุคฺคหาย.\csromanspacer{} กตเม จตฺตาโร.\csromanspacer{} กพฬีกาโร\footnote{กพฬิํกาโร (สี, อิ), กวฬีกาโร (สฺยา, กํ)} อาหาโร โอฬาริโก วา สุขุโม วา, ผสฺโส ทุติโย, มโนสญฺเจตนา ตติยา, วิญฺญาณํ จตุตฺถํ.\csromanspacer{} อิเม โข ภิกฺขเว จตฺตาโร อาหารา ภูตานํ วา สตฺตานํ ฐิติยา สมฺภเวสีนํ วา อนุคฺคหาย.}

\prose{อิเม ภิกฺขเว จตฺตาโร อาหารา กิํนิทานา กิํสมุทยา กิํชาติกา กิํปภวา.\csromanspacer{} อิเม จตฺตาโร อาหารา ตณฺหานิทานา ตณฺหาสมุทยา ตณฺหาชาติกา ตณฺหาปภวา.\csromanspacer{} ตณฺหา จายํ ภิกฺขเว กิํนิทานา กิํสมุทยา กิํชาติกา กิํปภวา.\csromanspacer{} ตณฺหา เวทนานิทานา เวทนาสมุทยา เวทนาชาติกา เวทนาปภวา.\csromanspacer{} เวทนา จายํ ภิกฺขเว กิํนิทานา กิํสมุทยา กิํชาติกา กิํปภวา.\csromanspacer{} เวทนา ผสฺสนิทานา ผสฺสสมุทยา ผสฺสชาติกา ผสฺสปภวา.\csromanspacer{} ผสฺโส จายํ ภิกฺขเว กิํนิทาโน กิํสมุทโย กิํชาติโก กิํปภโว.\csromanspacer{} ผสฺโส สฬายตนนิทาโน สฬายตนสมุทโย สฬายตนชาติโก สฬายตนปภโว.\csromanspacer{} สฬายตนญฺจิทํ ภิกฺขเว กิํนิทานํ กิํสมุทยํ กิํชาติกํ กิํปภวํ.\csromanspacer{} สฬายตนํ นามรูปนิทานํ นามรูปสมุทยํ นามรูปชาติกํ นามรูปปภวํ.\csromanspacer{} นามรูปญฺจิทํ ภิกฺขเว กิํนิทานํ กิํสมุทยํ กิํชาติกํ กิํปภวํ.\csromanspacer{} นามรูปํ วิญฺญาณนิทานํ วิญฺญาณสมุทยํ วิญฺญาณชาติกํ วิญฺญาณปภวํ.\csromanspacer{} วิญฺญาณญฺจิทํ ภิกฺขเว กิํนิทานํ กิํสมุทยํ กิํชาติกํ กิํปภวํ.\csromanspacer{} วิญฺญาณํ สงฺขารนิทานํ สงฺขารสมุทยํ สงฺขารชาติกํ สงฺขารปภวํ.\csromanspacer{} สงฺขารา จิเม ภิกฺขเว กิํนิทานา กิํสมุทยา กิํชาติกา กิํปภวา.\csromanspacer{} สงฺขารา อวิชฺชานิทานา อวิชฺชาสมุทยา อวิชฺชาชาติกา อวิชฺชาปภวา.}

\prose{อิติ โข ภิกฺขเว อวิชฺชาปจฺจยา สงฺขารา, สงฺขารปจฺจยา วิญฺญาณํ ฯเปฯ เอวเมตสฺส เกวลสฺส ทุกฺขกฺขนฺธสฺส สมุทโย โหติ.\csromanspacer{} อวิชฺชาย เตฺวว อเสสวิราคนิโรธา สงฺขารนิโรโธ, สงฺขารนิโรธา}

\vspace{\tipitakaparskip}
\csromancenter{นิทานวคฺคสํยุตฺตปาฬิ วิญฺญาณนิโรโธ ฯเปฯ เอวเมตสฺส เกวลสฺส ทุกฺขกฺขนฺธสฺส นิโรโธ โหตีติ.\csromanspacer{}\csromanspacer{}ปฐมํ.}
\csromanheader{2}{2. โมฬิยผคฺคุนสุตฺต}
\proseitem{12}{\csromanfolio{254}สาวตฺถิยํ วิหรติ.\csromanspacer{} จตฺตาโรเม ภิกฺขเว อาหารา ภูตานํ วา สตฺตานํ ฐิติยา สมฺภเวสีนํ วา อนุคฺคหาย.\csromanspacer{} กตเม จตฺตาโร.\csromanspacer{} กพฬีกาโร อาหาโร โอฬาริโก วา สุขุโม วา, ผสฺโส ทุติโย, มโนสญฺเจตนา ตติยา, วิญฺญาณํ จตุตฺถํ.\csromanspacer{} อิเม โข ภิกฺขเว จตฺตาโร อาหารา ภูตานํ วา สตฺตานํ ฐิติยา สมฺภเวสีนํ วา อนุคฺคหายาติ.}

\prose{เอวํ วุตฺเต อายสฺมา โมฬิยผคฺคุโน ภควนฺตํ เอตทโวจ “โก นุ โข ภนฺเต วิญฺญาณาหารํ อาหาเรตี”ติ. “โน กลฺโล ปโญฺห”ติ ภควา อโวจ, “อาหาเรตี”ติ อหํ น วทามิ. “อาหาเรตี”ติ จาหํ วเทยฺยํ, ตตฺรสฺส กลฺโล ปโญฺห “โก นุ โข ภนฺเต อาหาเรตี”ติ.\csromanspacer{} เอวญฺจาหํ น วทามิ, เอวํ มํ อวทนฺตํ โย เอวํ ปุจฺเฉยฺย “กิสฺส นุ โข ภนฺเต วิญฺญาณาหาโร”ติ.\csromanspacer{} เอส กลฺโล ปโญฺห, ตตฺร กลฺลํ เวยฺยากรณํ “วิญฺญาณาหาโร อายติํ ปุนพฺภวาภินิพฺพตฺติยา ปจฺจโย, ตสฺมิํ ภูเต สติ สฬายตนํ, สฬายตนปจฺจยา ผสฺโส”ติ.}

\prose{โก นุ โข ภนฺเต ผุสตีติ. “โน กลฺโล ปโญฺห”ติ ภควา อโวจ, “ผุสตี”ติ อหํ น วทามิ. “ผุสตี”ติ จาหํ วเทยฺยํ, ตตฺรสฺส กลฺโล ปโญฺห “โก นุ โข ภนฺเต ผุสตี”ติ.\csromanspacer{} เอวญฺจาหํ น วทามิ, เอวํ มํ อวทนฺตํ โย เอวํ ปุจฺเฉยฺย “กิํ ปจฺจยา นุ โข ภนฺเต ผสฺโส”ติ.\csromanspacer{} เอส กลฺโล ปโญฺห, ตตฺร กลฺลํ เวยฺยากรณํ สฬายตนปจฺจยา ผสฺโส, ผสฺสปจฺจยา เวทนา”ติ.}

\prose{โก นุ โข ภนฺเต เวทยตีติ\footnote{เวทิยตีติ (สี, อิ, ก)}. “โน กลฺโล ปโญฺห”ติ ภควา อโวจ, “เวทยตี”ติ อหํ น วทามิ. “เวทยตี”ติ จาหํ วเทยฺยํ, ตตฺรสฺส กลฺโล ปโญฺห “โก นุ โข ภนฺเต เวทยตี”ติ.\csromanspacer{} เอวญฺจาหํ น วทามิ, เอวํ มํ อวทนฺตํ โย เอวํ ปุจฺเฉยฺย “กิํ ปจฺจยา นุ โข ภนฺเต เวทนา”ติ.\csromanspacer{} เอส กลฺโล ปโญฺห, ตตฺร กลฺลํ เวยฺยากรณํ “ผสฺสปจฺจยา เวทนา, เวทนาปจฺจยา ตณฺหา”ติ.}

\csromanheader{2}{1. นิทานสํยุตฺต}
\prose{\csromanfolio{255}โก นุ โข ภนฺเต ตสตีติ\footnote{ตณฺหียตีติ (สี, สฺยา, กํ)}. “โน กลฺโล ปโญฺห”ติ ภควา อโวจ, “ตสตี”ติ อหํ น วทามิ. “ตสตี”ติ จาหํ วเทยฺยํ, ตตฺรสฺส กลฺโล ปโญฺห “โก นุ โข ภนฺเต ตสตี”ติ.\csromanspacer{} เอวญฺจาหํ น วทามิ, เอวํ มํ อวทนฺตํ โย เอวํ ปุจฺเฉยฺย “กิํปจฺจยา นุ โข ภนฺเต ตณฺหา”ติ.\csromanspacer{} เอส กลฺโล ปโญฺห, ตตฺร กลฺลํ เวยฺยากรณํ “เวทนาปจฺจยา ตณฺหา, ตณฺหาปจฺจยา อุปาทานนฺ”ติ.}

\prose{โก นุ โข ภนฺเต อุปาทิยตีติ. “โน กลฺโล ปโญฺห”ติ ภควา อโวจ, “อุปาทิยตี”ติ อหํ น วทามิ. “อุปาทิยตี”ติ จาหํ วเทยฺยํ, ตตฺรสฺส กลฺโล ปโญฺห “โก นุ โข ภนฺเต อุปาทิยตี”ติ.\csromanspacer{} เอวญฺจาหํ น วทามิ, เอวํ มํ อวทนฺตํ โย เอวํ ปุจฺเฉยฺย “กิํ ปจฺจยา นุ โข ภนฺเต อุปาทานนฺ”ติ.\csromanspacer{} เอส กลฺโล ปโญฺห, ตตฺร กลฺลํ เวยฺยากรณํ “ตณฺหาปจฺจยา อุปาทานํ, อุปาทานปจฺจยา ภโว”ติ ฯเปฯ เอวเมตสฺส เกวลสฺส ทุกฺขกฺขนฺธสฺส สมุทโย โหติ.}

\prose{ฉนฺทํ เตฺวว ผคฺคุน ผสฺสายตนานํ อเสสวิราคนิโรธา ผสฺสนิโรโธ, ผสฺสนิโรธา เวทนานิโรโธ, เวทนานิโรธา ตณฺหานิโรโธ, ตณฺหานิโรธา อุปาทานนิโรโธ, อุปาทานนิโรธา ภวนิโรโธ, ภวนิโรธา ชาตินิโรโธ, ชาตินิโรธา ชรามรณํ โสกปริเทวทุกฺขโทมนสฺสุปายาสา นิรุชฺฌนฺติ, เอวเมตสฺส เกวลสฺส ทุกฺขกฺขนฺธสฺส นิโรโธ โหตีติ.\csromanspacer{}\csromanspacer{}ทุติยํ.}

\csromanheader{2}{3. สมณพฺราหฺมณสุตฺต}
\proseitem{13}{สาวตฺถิยํ วิหรติ.\csromanspacer{} เย หิ เกจิ ภิกฺขเว สมณา วา พฺราหฺมณา วา ชรามรณํ นปฺปชานนฺติ, ชรามรณสมุทยํ นปฺปชานนฺติ, ชรามรณนิโรธํ นปฺปชานนฺติ, ชรามรณนิโรธคามินิํ ปฏิปทํ นปฺปชานนฺติ.\csromanspacer{} ชาติํ ฯเปฯ.\csromanspacer{} ภวํ.\csromanspacer{} อุปาทานํ.\csromanspacer{} ตณฺหํ.\csromanspacer{} เวทนํ.\csromanspacer{} ผสฺสํ.\csromanspacer{} อฬายตนํ.\csromanspacer{} นามรูปํ.\csromanspacer{} วิญฺญาณํ.\csromanspacer{} สงฺขาเร นปฺปชานนฺติ, สงฺขารสมุทยํ นปฺปชานนฺติ, สงฺขารนิโรธํ นปฺปชานนฺติ, สงฺขารนิโรธคามินิํ ปฏิปทํ นปฺปชานนฺติ.\csromanspacer{} น เม เต ภิกฺขเว สมณา วา พฺราหฺมณา วา สมเณสุ วา สมณสมฺมตา, พฺราหฺมเณสุ วา พฺราหฺมณสมฺมตา, น จ ปน เต อายสฺมนฺโต สามญฺญตฺถํ วา}

\vspace{\tipitakaparskip}
\csromancenter{นิทานวคฺคสํยุตฺตปาฬิ พฺรหฺมญฺญตฺถํ\footnote{พฺราหฺมญฺญตฺถํ (สฺยา, กํ) โมคฺคลฺลานพฺยากรณํ โอโลเกตพฺพํ.} วา ทิฏฺเฐว ธมฺเม สยํ อภิญฺญา สจฺฉิกตฺวา อุปสมฺปชฺช วิหรนฺติ.}
\csromanmatikaanchor{mtk.8}
\csromantocmarkref{subsection}{6. เวสฺสภูสุตฺต}{mtk.8}
\bookmark[dest={mtk.8},level=2]{6. เวสฺสภูสุตฺต}
\prose{\csromanfolio{256}เย จ โข เกจิ ภิกฺขเว สมณา วา พฺราหฺมณา วา ชรามรณํ ปชานนฺติ, ชรามรณสมุทยํ ปชานนฺติ, ชรามรณนิโรธํ ปชานนฺติ, ชรามรณนิโรธคามินิํ ปฏิปทํ ปชานนฺติ.\csromanspacer{} ชาติํ ฯเปฯ.\csromanspacer{} ภวํ.\csromanspacer{} อุปาทานํ.\csromanspacer{} ตณฺหํ.\csromanspacer{} เวทนํ.\csromanspacer{} ผสฺสํ.\csromanspacer{} สฬายตนํ.\csromanspacer{} นามรูปํ.\csromanspacer{} วิญฺญาณํ.\csromanspacer{} สงฺขาเร ปชานนฺติ, สงฺขารสมุทยํ ปชานนฺติ, สงฺขารนิโรธํ ปชานนฺติ, สงฺขารนิโรธคามินิํ ปฏิปทํ ปชานนฺติ.\csromanspacer{} เต โข เม ภิกฺขเว สมณา วา พฺราหฺมณา วา สมเณสุ เจว สมณสมฺมตา, พฺราหฺมเณสุ จ พฺราหฺมณสมฺมตา, เต จ ปนายสฺมนฺโต สามญฺญตฺถญฺจ พฺรหฺมญฺญตฺถญฺจ ทิฏฺเฐว ธมฺเม สยํ อภิญฺญา สจฺฉิกตฺวา อุปสมฺปชฺช วิหรนฺตีติ.\csromanspacer{}\csromanspacer{}ตติยํ.}

\csromanheader{2}{4. ทุติยสมณพฺราหฺมณสุตฺต}
\csromanmatikaanchor{mtk.9}
\csromantocmarkref{subsection}{7. กกุสนฺธสุตฺต}{mtk.9}
\bookmark[dest={mtk.9},level=2]{7. กกุสนฺธสุตฺต}
\proseitem{14}{สาวตฺถิยํ วิหรติ.\csromanspacer{} เย หิ เกจิ ภิกฺขเว สมณา วา พฺราหฺมณา วา อิเม ธมฺเม นปฺปชานนฺติ, อิเมสํ ธมฺมานํ สมุทยํ นปฺปชานนฺติ, อิเมสํ ธมฺมานํ นิโรธํ นปฺปชานนฺติ, อิเมสํ ธมฺมานํ นิโรธคามินิํ ปฏิปทํ นปฺปชานนฺติ.\csromanspacer{} กตเม ธมฺเม นปฺปชานนฺติ, กตเมสํ ธมฺมานํ สมุทยํ นปฺปชานนฺติ, กตเมสํ ธมฺมานํ นิโรธํ นปฺปชานนฺติ, กตเมสํ ธมฺมานํ นิโรธคามินิํ ปฏิปทํ นปฺปชานนฺติ.}

\prose{ชรามรณํ นปฺปชานนฺติ, ชรามรณสมุทยํ นปฺปชานนฺติ, ชรามรณนิโรธํ นปฺปชานนฺติ, ชรามรณนิโรธคามินิํ ปฏิปทํ นปฺปชานนฺติ.\csromanspacer{} ชาติํ ฯเปฯ.\csromanspacer{} ภวํ.\csromanspacer{} อุปาทานํ.\csromanspacer{} ตณฺหํ.\csromanspacer{} เวทนํ.\csromanspacer{} ผสฺสํ.\csromanspacer{} สฬายตนํ.\csromanspacer{} นามรูปํ.\csromanspacer{} วิญฺญาณํ.\csromanspacer{} สงฺขาเร นปฺปชานนฺติ, สงฺขารสมุทยํ นปฺปชานนฺติ, สงฺขารนิโรธํ นปฺปชานนฺติ, สงฺขารนิโรธคามินิํ ปฏิปทํ นปฺปชานนฺติ.\csromanspacer{} อิเม ธมฺเม นปฺปชานนฺติ, อิเมสํ ธมฺมานํ สมุทยํ นปฺปชานนฺติ, อิเมสํ ธมฺมานํ นิโรธํ นปฺปชานนฺติ, อิเมสํ ธมฺมานํ นิโรธคามินิํ ปฏิปทํ นปฺปชานนฺติ.\csromanspacer{} น เม เต ภิกฺขเว สมณา วา พฺราหฺมณา วา สมเณสุ วา สมณสมฺมตา, พฺราหฺมเณสุ วา พฺราหฺมณสมฺมตา, น จ ปน เต อายสฺมนฺโต สมญฺญตฺถํ วา พฺรหฺมญฺญตฺถํ วา ทิฏฺเฐว ธมฺเม สยํ อภิญฺญา สจฺฉิกตฺวา อุปสมฺปชฺช วิหรนฺติ.}

\csromanmatikaanchor{mtk.10}
\csromantocmarkref{subsection}{8. โกณาคมนสุตฺต}{mtk.10}
\bookmark[dest={mtk.10},level=2]{8. โกณาคมนสุตฺต}
\csromanheader{2}{1. นิทานสํยุตฺต}
\prose{\csromanfolio{257}เย จ โข เกจิ ภิกฺขเว สมณา วา พฺราหฺมณา วา อิเม ธมฺเม ปชานนฺติ, อิเมสํ ธมฺมานํ สมุทยํ ปชานนฺติ, อิเมสํ ธมฺมานํ นิโรธํ ปชานนฺติ, อิเมสํ ธมฺมานํ นิโรธคามินิํ ปฏิปทํ ปชานนฺติ.\csromanspacer{} กตเม ธมฺเม ปชานนฺติ, กตเมสํ ธมฺมานํ สมุทยํ ปชานนฺติ, กตเมสํ ธมฺมานํ นิโรธํ ปชานนฺติ, กตเมสํ ธมฺมานํ นิโรธคามินิํ ปฏิปทํ ปชานนฺติ.}

\csromanmatikaanchor{mtk.11}
\csromantocmarkref{subsection}{9. กสฺสปสุตฺต}{mtk.11}
\bookmark[dest={mtk.11},level=2]{9. กสฺสปสุตฺต}
\prose{ชรามรณํ ปชานนฺติ, ชรามรณสมุทยํ ปชานนฺติ, ชรามรณนิโรธํ ปชานนฺติ, ชรามรณนิโรธคามินิํ ปฏิปทํ ปชานนฺติ.\csromanspacer{} ชาติํ ฯเปฯ.\csromanspacer{} ภวํ.\csromanspacer{} อุปาทานํ.\csromanspacer{} ตณฺหํ.\csromanspacer{} เวทนํ.\csromanspacer{} ผสฺสํ.\csromanspacer{} สฬายตนํ.\csromanspacer{} นามรูปํ.\csromanspacer{} วิญฺญาณํ.\csromanspacer{} สงฺขาเร ปชานนฺติ, สงฺขารสมุทยํ ปชานนฺติ, สงฺขารนิโรธํ ปชานนฺติ, สงฺขารนิโรธคามินิํ ปฏิปทํ ปชานนฺติ.\csromanspacer{} อิเม ธมฺเม ปชานนฺติ, อิเมสํ ธมฺมานํ สมุทยํ ปชานนฺติ, อิเมสํ ธมฺมานํ นิโรธํ ปชานนฺติ, อิเมสํ ธมฺมานํ นิโรธคามินิํ ปฏิปทํ ปชานนฺติ.\csromanspacer{} เต โข เม ภิกฺขเว สมณา วา พฺราหฺมณา วา สมเณสุ เจว สมณสมฺมตา, พฺราหฺมเณสุ จ พฺราหฺมณสมฺมตา, เต จ ปนายสฺมนฺโต สามญฺญตฺถญฺจ พฺรหฺมญฺญตฺถญฺจ ทิฏฺเฐว ธมฺเม สยํ อภิญฺญา สจฺฉิกตฺวา อุปสมฺปชฺช วิหรนฺตีติ.\csromanspacer{}\csromanspacer{}จตุตฺถํ.}

\csromanheader{2}{5. กจฺจานโคตฺตสุตฺต}
\csromanmatikaanchor{mtk.12}
\csromantocmarkref{subsection}{10. โคตมสุตฺต}{mtk.12}
\bookmark[dest={mtk.12},level=2]{10. โคตมสุตฺต}
\csromantocmarkref{section}{อุทฺทานคาถา}{mtk.13}
\bookmark[dest={mtk.13},level=1]{อุทฺทานคาถา}
\proseitem{15}{สาวตฺถิยํ คิหรติ.\csromanspacer{} อถ โข อายสฺมา กจฺจานโคตฺโต เยน ภควา เตนุปสงฺกมิ, อุปสงฺกมิตฺวา ภควนฺตํ อภิวาเทตฺวา เอกมนฺตํ นิสีทิ, เอกมนฺตํ นิสินฺโน โข อายสฺมา กจฺจานโคตฺโต ภควนฺตํ เอตทโวจ “สมฺมาทิฏฺฐิ สมฺมาทิฏฺฐีติ ภนฺเต วุจฺจติ, กิตฺตาวตา นุ โข ภนฺเต สมฺมาทิฏฺฐิ โหตี”ติ.}

\prose{ทฺวยนิสฺสิโต ขฺวายํ กจฺจาน โลโก เยภุยฺเยน อตฺถิตญฺเจว นตฺถิตญฺจ.\csromanspacer{} โลกสมุทยํ โข กจฺจาน ยถาภูตํ สมฺมปฺปญฺญาย ปสฺสโต ยา โลเก นตฺถิตา, สา น โหติ.\csromanspacer{} โลกนิโรธํ โข กจฺจาน ยถาภูตํ สมฺมปฺปญฺญาย ปสฺสโต ยา โลเก อตฺถิตา, สา น โหติ.\csromanspacer{} อุปยุปาทานาภินิเวสวินิพนฺโธ\footnote{อุปายุปาทานาภินิเวสวินิพนฺโธ (สี, สฺยา, กํ, อิ)} ขฺวายํ กจฺจาน โลโก เยภุยฺเยน, ตญฺจายํ อุปยุปาทานํ เจตโส อธิฏฺฐานํ อภินิเวสานุสยํ น อุเปติ น อุปาทิยติ นาธิฏฺฐาติ “อตฺตา เม”ติ. “ทุกฺขเมว}

\vspace{\tipitakaparskip}
\csromancenter{นิทานวคฺคสํยุตฺตปาฬิ อุปฺปชฺชมานํ อุปฺปชฺชติ, ทุกฺขํ นิรุชฺฌมานํ นิรุชฺฌตี”ติ น กงฺขติ น วิจิกิจฺฉติ, อปรปจฺจยา ญาณเมวสฺส เอตฺถ โหติ.\csromanspacer{} เอตฺตาวตา โข กจฺจาน สมฺมาทิฏฺฐิ โหติ.}
\prose{\csromanfolio{258}“สพฺพมตฺถี”ติ โข กจฺจาน อยเมโก อนฺโต, “สพฺพํ นตฺถี”ติ อยํ ทุติโย อนฺโต, เอเต เต กจฺจาน อุโภ อนฺเต อนุปคมฺม มชฺเฌน ตถาคโต ธมฺมํ เทเสติ\csromandash{}อวิชฺชาปจฺจยา สงฺขารา, สงฺขารปจฺจยา วิญฺญาณํ ฯเปฯ เอวเมตสฺส เกวลสฺส ทุกฺขกฺขนฺธสฺส สมุทโย โหติ.\csromanspacer{} อวิชฺชาย เตฺวว อเสสวิราคนิโรธา สงฺขารนิโรโธ, สงฺขารนิโรธา วิญฺญาณนิโรโธ ฯเปฯ เอวเมตสฺส เกวลสฺส ทุกฺขกฺขนฺธสฺส นิโรโธ โหตีติ.\csromanspacer{}\csromanspacer{}ปญฺจมํ.}

\csromanheader{2}{6. ธมฺมกถิกสุตฺต}
\proseitem{16}{สาวตฺถิยํ.\csromanspacer{} อถ โข อญฺญตโร ภิกฺขุ เยน ภควา เตนุปสงฺกมิ, อุปสงฺกมิตฺวา ภควนฺตํ อภิวาเทตฺวา เอกมนฺตํ นิสีทิ, เอกมนฺตํ นิสินฺโน โข โส ภิกฺขุ ภควนฺตํ เอตทโวจ “ธมฺมกถิโก ธมฺมกถิโกติ ภนฺเต วุจฺจติ, กิตฺตาวตา นุ โข ภนฺเต ธมฺมกถิโก โหตี”ติ.}

\prose{ชรามรณสฺส เจ ภิกฺขุ นิพฺพิทาย วิราคาย นิโรธาย ธมฺมํ เทเสติ, “ธมฺมกถิโก ภิกฺขู”ติ อลํ วจนาย.\csromanspacer{} ชรามรณสฺส เจ ภิกฺขุ นิพฺพิทาย วิราคาย นิโรธาย ปฏิปนฺโน โหติ, “ธมฺมานุธมฺมปฺปฏิปนฺโน ภิกฺขู”ติ อลํ วจนาย.\csromanspacer{} ชรามรณสฺส เจ ภิกฺขุ นิพฺพิทา วิราคา นิโรธา อนุปาทาวิมุตฺโต โหติ, “ทิฏฺฐธมฺมนิพฺพานปฺปตฺโต ภิกฺขู”ติ อลํ วจนาย.}

\prose{ชาติยา เจ ภิกฺขุ ฯเปฯ.\csromanspacer{} ภวสฺส เจ ภิกฺขุ.\csromanspacer{} อุปาทานสฺส เจ ภิกฺขุ.\csromanspacer{} ตณฺหาย เจ ภิกฺขุ.\csromanspacer{} เวทนาย เจ ภิกฺขุ.\csromanspacer{} ผสฺสสฺส เจ ภิกฺขุ.\csromanspacer{} สฬายตนสฺส เจ ภิกฺขุ.\csromanspacer{} นามรูปสฺส เจ ภิกฺขุ.\csromanspacer{} วิญฺญาณสฺส เจ ภิกฺขุ.\csromanspacer{} สงฺขารานํ เจ ภิกฺขุ.\csromanspacer{} อวิชฺชาย เจ ภิกฺขุ นิพฺพิทาย วิราคาย นิโรธาย ธมฺมํ เทเสติ, “ธมฺมกถิโก ภิกฺขู”ติ อลํ วจนาย.\csromanspacer{} อวิชฺชาย เจ ภิกฺขุ นิพฺพิทาย วิราคาย นิโรธาย ปฏิปนฺโน โหติ, “ธมฺมานุธมฺมปฺปฏิปนฺโน ภิกฺขู”ติ อลํ วจนาย.\csromanspacer{} อวิชฺชาย เจ ภิกฺขุ นิพฺพิทา วิราคา นิโรธา อนุปาทาวิมุตฺโต โหติ, “ทิฏฺฐธมฺมนิพฺพานปฺปตฺโต ภิกฺขู”ติ อลํ วจนายาติ.\csromanspacer{}\csromanspacer{}ฉฏฺฐํ.}

\csromanheader{2}{1. นิทานสํยุตฺต \textbf{7. อเจลกสฺสปสุตฺต}}
\proseitem{17}{\csromanfolio{259}เอวํ เม สุตํ\csromandash{}เอกํ สมยํ ภควา ราชคเห วิหรติ เวฬุวเน กลนฺทกนิวาเป.\csromanspacer{} อถ โข ภควา ปุพฺพณฺหสมยํ นิวาเสตฺวา ปตฺตจีวรมาทาย ราชคหํ ปิณฺฑาย ปาวิสิ.\csromanspacer{} อทฺทสา โข อเจโล กสฺสโป ภควนฺตํ ทูรโตว อาคจฺฉนฺตํ, ทิสฺวาน เยน ภควา เตนุปสงฺกมิ, อุปสงฺกมิตฺวา ภควตา สทฺธิํ สมฺโมทิ, สมฺโมทนียํ กถํ สารณียํ วีติสาเรตฺวา เอกมนฺตํ อฏฺฐาสิ, เอกมนฺตํ ฐิโต โข อเจโล กสฺสโป ภควนฺตํ เอตทโวจ “ปุจฺเฉยฺยาม มยํ ภวนฺตํ โคตมํ กญฺจิเทว\footnote{กิญฺจิเทว (ก)} เทสํ, สเจ โน ภวํ โคตโม โอกาสํ กโรติ ปญฺหสฺส เวยฺยากรณายา”ติ.}

\prose{อกาโล โข ตาว กสฺสป ปญฺหสฺส, อนฺตรฆรํ ปวิฏฺฐมฺหาติ.\csromanspacer{} ทุติยมฺปิ โข อเจโล กสฺสโป ภควนฺตํ เอตทโวจ “ปุจฺเฉยฺยาม มยํ ภวนฺตํ โคตมํ กญฺจิเทว เทสํ, สเจ โน ภวํ โคตโม โอกาสํ กโรติ ปญฺหสฺส เวยฺยากรณายา”ติ.\csromanspacer{} อกาโล โข ตาว กสฺสป ปญฺหสฺส, อนฺตรฆรํ ปวิฏฺฐมฺหาติ.\csromanspacer{} ตติยมฺปิ โข อเจโล กสฺสโป ฯเปฯ อานฺตรฆรํ ปวิฏฺฐมฺหาติ.\csromanspacer{} เอวํ วุตฺเต อเจโล กสฺสโป ภควนฺตํ เอตทโวจ “น โข ปน มยํ ภวนฺตํ โคตมํ พหุเทว ปุจฺฉิตุกามา”ติ.\csromanspacer{} ปุจฺฉ กสฺสป ยทากงฺขสีติ.}

\prose{กิํ นุ โข โภ โคตม สยํกตํ ทุกฺขนฺติ. “มา เหวํ กสฺสปา”ติ ภควา อโวจ.\csromanspacer{} กิํ ปน โภ โคตม ปรํกตํ ทุกฺขนฺติ. “มา เหวํ กสฺสปา”ติ ภควา อโวจ.\csromanspacer{} กิํ นุ โข โภ โคตม สยํกตญฺจ ปรํกตญฺจ ทุกฺขนฺติ. “มา เหวํ กสฺสปา”ติ ภควา อโวจ.\csromanspacer{} กิํ ปน โภ โคตม อสยํการํ อปรํการํ อธิจฺจสมุปฺปนฺนํ ทุกฺขนฺติ. “มา เหวํ กสฺสปา”ติ ภควา อโวจ.\csromanspacer{} กิํ นุ โข โภ โคตม นตฺถิ ทุกฺขนฺติ.\csromanspacer{} น โข กสฺสป นตฺถิ ทุกฺขํ, อตฺถิ โข กสฺสป ทุกฺขนฺติ.\csromanspacer{} เตน หิ ภวํ โคตโม ทุกฺขํ น ชานาติ น ปสฺสตีติ.\csromanspacer{} น ขฺวาหํ กสฺสป ทุกฺขํ น ชานามิ น ปสฺสามิ, ชานามิ ขฺวาหํ กสฺสป ทุกฺขํ, ปสฺสามิ ขฺวาหํ กสฺสป ทุกฺขนฺติ.}

\csromanmatikaanchor{mtk.14}
\csromantocmarknopage{section}{2. อาหารวคฺค}{mtk.14}
\bookmark[dest={mtk.14},level=1]{2. อาหารวคฺค}
\prose{“กิํ นุ โข โภ โคตม สยํกตํ ทุกฺขนฺ”ติ อิติ ปุฏฺโฐ สมาโน “มา เหวํ กสฺสปา”ติ วเทสิ. “กิํ ปน โภ โคตม ปรํกตํ ทุกฺขนฺ”ติ}

\csromanmatikaanchor{mtk.15}
\csromantocmarkref{subsection}{1. อาหารสุตฺต}{mtk.15}
\bookmark[dest={mtk.15},level=2]{1. อาหารสุตฺต}
\vspace{\tipitakaparskip}
\csromancenter{นิทานวคฺคสํยุตฺตปาฬิ อิติ ปุฏฺโฐ สมาโน “มา เหวํ กสฺสปา”ติ วเทสิ. “กิํ นุ โข โภ โคตม สยํกตญฺจ ปรํกตญฺจ ทุกฺขนฺ”ติ อิติ ปุฏฺโฐ สมาโน “มา เหวํ กสฺสปา”ติ วเทสิ. “กิํ ปน โภ โคตม อสยํการํ อปรํการํ อธิจฺจสมุปฺปนฺนํ ทุกฺขนฺ”ติ อิติ ปุฏฺโฐ สมาโน “มา เหวํ กสฺสปา”ติ วเทสิ. “กิํ นุ โข โภ โคตม นตฺถิ ทุกฺขนฺ”ติ อิติ ปุฏฺโฐ สมาโน “น โข กสฺสป นตฺถิ ทุกฺขํ, อตฺถิ โข กสฺสป ทุกฺขนฺ”ติ วเทสิ. “เตน หิ ภวํ โคตโม ทุกฺขํ น ชานาติ น ปสฺสตี”ติ อิติ ปุฏฺโฐ สมาโน “น ขฺวาหํ กสฺสป ทุกฺขํ น ชานามิ น ปสฺสามิ, ชานามิ ขฺวาหํ กสฺสป ทุกฺขํ, ปสฺสามิ ขฺวาหํ กสฺสป ทุกฺขนฺ”ติ วเทสิ.\csromanspacer{} อาจิกฺขตุ จ\footnote{อยํ จกาโร สี-โปตฺถเก นตฺถิ.} เม ภนฺเต ภควา ทุกฺขํ, เทเสตุ จ1 เม ภนฺเต ภควา ทุกฺขนฺติ.}
\prose{\csromanfolio{260}“โส กโรติ โส ปฏิสํเวทยตี”ติ\footnote{ปฏิสํเวทิยตีติ (สี, อิ, ก)} โข กสฺสป อาทิโต สโต “สยํกตํ ทุกฺขนฺ”ติ อิติ วทํ สสฺสตํ เอตํ ปเรติ. “อญฺโญ กโรติ อญฺโญ ปฏิสํเวทยตี”ติ โข กสฺสป เวทนาภิตุนฺนสฺส สโต “ปรํกตํ ทุกฺขนฺ”ติ อิติ วทํ อุจฺเฉทํ เอตํ ปเรติ.\csromanspacer{} เอเต เต กสฺสป อุโภ อนฺเต อนุปคมฺม มชฺเฌน ตถาคโต ธมฺมํ เทเสติ\csromandash{}อวิชฺชาปจฺจยา สงฺขารา, สงฺขารปจฺจยา วิญฺญาณํ ฯเปฯ เอวเมตสฺส เกวลสฺส ทุกฺขกฺขนฺธสฺส สมุทโย โหติ.\csromanspacer{} อวิชฺชาย เตฺวว อเสสวิราคนิโรธา สงฺขารนิโรโธ, สงฺขารนิโรธา วิญฺญาณนิโรโธ ฯเปฯ เอวเมตสฺส เกวลสฺส ทุกฺขกฺขนฺธสฺส นิโรโธ โหตีติ.}

\prose{เอวํ วุตฺเต อเจโล กสฺสโป ภควนฺตํ เอตทโวจ “อภิกฺกนฺตํ ภนฺเต, อภิกฺกนฺตํ ภนฺเต, เสยฺยถาปิ ภนฺเต นิกฺกุชฺชิตํ วา อุกฺกุชฺเชยฺย ฯเปฯ ‘จกฺขุมนฺโต รูปานิ ทกฺขนฺตี’ติ, เอวเมวํ ภควตา อเนกปริยาเยน ธมฺโม ปกาสิโต, เอสาหํ ภนฺเต ภควนฺตํ สรณํ คจฺฉามิ ธมฺมญฺจ ภิกฺขุสํฆญฺจ, ลเภยฺยาหํ ภนฺเต ภควโต สนฺติเก ปพฺพชฺชํ, ลเภยฺยํ อุปสมฺปทนฺ”ติ.}

\prose{โย โข กสฺสป อญฺญติตฺถิยปุพฺโพ อิมสฺมิํ ธมฺมวินเย อากงฺขติ ปพฺพชฺชํ อากงฺขติ อุปสมฺปทํ, โส จตฺตาโร มาเส ปริวสติ, จตุนฺนํ มาสานํ อจฺจเยน\footnote{อจฺจเยน ปริวุฏฺฐปริวาสํ (สฺยา, กํ, อิ, ก)} อารทฺธจิตฺตา ภิกฺขู\footnote{ภิกฺขู อากงฺขมานา (สฺยา, กํ, อิ, ก)} ปพฺพาเชนฺติ อุปสมฺปาเทนฺติ ภิกฺขุภาวาย.\csromanspacer{} อปิ จ มยา ปุคฺคลเวมตฺตตา วิทิตาติ.}

\csromanheader{2}{1. นิทานสํยุตฺต}
\csromanmatikaanchor{mtk.16}
\csromantocmarkref{subsection}{2. โมฬิยผคฺคุณสุตฺต}{mtk.16}
\bookmark[dest={mtk.16},level=2]{2. โมฬิยผคฺคุณสุตฺต}
\prose{\csromanfolio{261}สเจ ภนฺเต อญฺญติตฺถิยปุพฺโพ อิมสฺมิํ ธมฺมวินเย อากงฺขติ ปพฺพชฺชํ อากงฺขติ อุปสมฺปทํ, จตฺตาโร มาเส ปริวสติ, จตุนฺนํ มาสานํ อจฺจเยน\footnote{อจฺจเยน ปริวุฏฺฐปริวาสํ (สฺยา, กํ, อิ, ก)} อารทฺธจิตฺตา ภิกฺขู\footnote{ภิกฺขู อากงฺขมานา (สฺยา, กํ, อิ, ก)} ปพฺพาเชนฺติ อุปสมฺปาเทนฺติ ภิกฺขุภาวาย, อหํ จตฺตาริ วสฺสานิ ปริวสิสฺสามิ, จตุนฺนํ วสฺสานํ อจฺจเยน1 อารทฺธจิตฺตา ภิกฺขู ปพฺพาเชนฺตุ อุปสมฺปาเทนฺตุ ภิกฺขุภาวายาติ.}

\prose{อลตฺถ โข อเจโล กสฺสโป ภควโต สนฺติเก ปพฺพชฺชํ, อลตฺถ อุปสมฺปทํ, อจิรูปสมฺปนฺโน จ ปนายสฺมา กสฺสโป เอโก วูปกฏฺโฐ อปฺปมตฺโต อาตาปี ปหิตตฺโต วิหรนฺโต นจิรสฺเสว ยสฺสตฺถาย กุลปุตฺตา สมฺมเทว อคารสฺมา อนคาริยํ ปพฺพชนฺติ, ตทนุตฺตรํ พฺรหฺมจริยปริโยสานํ ทิฏฺเฐว ธมฺเม สยํ อภิญฺญา สจฺฉิกตฺวา อุปสมฺปชฺช วิหาสิ, “ขีณา ชาติ, วุสิตํ พฺรหฺมจริยํ, กตํ กรณียํ, นาปรํ อิตฺถตฺตายา”ติ อพฺภญฺญาสิ.\csromanspacer{} อญฺญตโร จ ปนายสฺมา กสฺสโป อรหตํ อโหสีติ.\csromanspacer{}\csromanspacer{}สตฺตมํ.}

\csromanheader{2}{8. ติมฺพรุกสุตฺต}
\proseitem{18}{สาวตฺถิยํ วิหรติ.\csromanspacer{} อถ โข ติมฺพรุโก ปริพฺพาชโก เยน ภควา เตนุปสงฺกมิ, อุปสงฺกมิตฺวา ภควตา สทฺธิํ สมฺโมทิ, สมฺโมทนียํ กถํ สารณียํ วีติสาเรตฺวา เอกมนฺตํ นิสีทิ, เอกมนฺตํ นิสินฺโน โข ติมฺพรุโก ปริพฺพาชโก ภควนฺตํ เอตทโวจ\csromandash{}}

\prose{กิํ นุ โข โภ โคตม สยํกตํ สุขทุกฺขนฺติ. “มา เหวํ ติมฺพรุกา”ติ ภควา อโวจ.\csromanspacer{} กิํ ปน โภ โคตม ปรํกตํ สุขทุกฺขนฺติ. “มา เหวํ ติมฺพรุกา”ติ ภควา อโวจ.\csromanspacer{} กิํ นุ โข โภ โคตม สยํกตญฺจ ปรํกตญฺจ สุขทุกฺขนฺติ. “มา เหวํ ติมฺพรุกา”ติ ภควา อโวจ.\csromanspacer{} กิํ ปน โภ โคตม อสยํการํ อปรํการํ อธิจฺจสมุปฺปนฺนํ สุขทุกฺขนฺติ. “มา เหวํ ติมฺพรุกา”ติ ภควา อโวจ.\csromanspacer{} กิํ นุ โข โภ โคตม นตฺถิ สุขทุกฺขนฺติ.\csromanspacer{} น โข ติมฺพรุก นตฺถิ สุขทุกฺขํ, อตฺถิ โข ติมฺพรุก สุขทุกฺขนฺติ.\csromanspacer{} เตน หิ ภวํ โคตโม สุขทุกฺขํ น ชานาติ น ปสฺสตีติ.\csromanspacer{} น ขฺวาหํ ติมฺพรุก สุขทุกฺขํ น ชานามิ น ปสฺสามิ, ชานามิ ขฺวาหํ ติมฺพรุก สุขทุกฺขํ ปสฺสามิ ขฺวาหํ ติมฺพรุก สุขทุกฺขนฺติ.}

\gambhira{นิทานวคฺคสํยุตฺตปาฬิ}
\prose{\csromanfolio{262}“กิํ นุ โข โภ โคตม สยํกตํ สุขทุกฺขนฺ”ติ อิติ ปุฏฺโฐ สมาโน “มา เหวํ ติมฺพรุกา”ติ วเทสิ. “กิํ ปน โภ โคตม ปรํกตํ สุขทุกฺขนฺ”ติ อิติ ปุฏฺโฐ สมาโน “มา เหวํ ติมฺพรุกา”ติ วเทสิ. “กิํ นุ โข โภ โคตม สยํกตญฺจ ปรํกตญฺจ สุขทุกฺขนฺ”ติ อิติ ปุฏฺโฐ สมาโน “มา เหวํ ติมฺพรุกา”ติ วเทสิ. “กิํ ปน โภ โคตม อสยํการํ อปรํการํ อธิจฺจสมุปฺปนฺนํ สุขทุกฺขนฺ”ติ อิติ ปุฏฺโฐ สมาโน “มา เหวํ ติมฺพรุกา”ติ วเทสิ. “กิํ นุ โข โภ โคตม นตฺถิ สุขทุกฺขนฺ”ติ อิติ ปุฏฺโฐ สมาโน “น โข ติมฺพรุก นตฺถิ สุขทุกฺขํ, อตฺถิ โข ติมฺพรุก สุขทุกฺขนฺ”ติ วเทสิ. “เตน หิ ภวํ โคตโม สุขทุกฺขํ น ชานาติ น ปสฺสตี”ติ อิติ ปุฏฺโฐ สมาโน “น ขฺวาหํ ติมฺพรุก สุขทุกฺขํ น ชานามิ น ปสฺสามิ, ชานามิ ขฺวาหํ ติมฺพรุก สุขทุกฺขํ, ปสฺสามิ ขฺวาหํ ติมฺพรุก สุขทุกฺขนฺ”ติ วเทสิ.\csromanspacer{} อาจิกฺขตุ จ เม ภวํ โคตโม สุขทุกฺขํ, เทเสตุ จ เม ภวํ โคตโม สุขทุกฺขนฺติ.}

\prose{“สา เวทนา โส เวทยตี”ติ โข ติมฺพรุก อาทิโต สโต “สยํกตํ สุขทุกฺขนฺ”ติ เอวมฺปาหํ น วทามิ. “อญฺญา เวทนา อญฺโญ เวทยตี”ติ โข ติมฺพรุก เวทนาภิตุนฺนสฺส สโต “ปรํกตํ สุขทุกฺขนฺ”ติ เอวมฺปาหํ น วทามิ.\csromanspacer{} เอเต เต ติมฺพรุก อุโภ อนฺเต อนุปคมฺม มชฺเฌน ตถาคโต ธมฺมํ เทเสติ\csromandash{}อวิชฺชาปจฺจยา สงฺขารา, สงฺขารปจฺจยา วิญฺญาณํ ฯเปฯ เอวเมตสฺส เกวลสฺส ทุกฺขกฺขนฺธสฺส สมุทโย โหติ.\csromanspacer{} อวิชฺชาย เตฺวว อเสสวิราคนิโรธา สงฺขารนิโรโธ, สงฺขารนิโรธา วิญฺญาณนิโรโธ ฯเปฯ เอวเมตสฺส เกวลสฺส ทุกฺขกฺขนฺธสฺส นิโรโธ โหตีติ.}

\prose{เอวํ วุตฺเต ติมฺพรุโก ปริพฺพาชโก ภควนฺตํ เอตทโวจ “อภิกฺกนฺตํ โภ โคตม ฯเปฯ เอสาหํ ภวนฺตํ โคตมํ สรณํ คจฺฉามิ ธมฺมญฺจ ภิกฺขุสํฆญฺจ, อุปาสกํ มํ ภวํ โคตโม ธาเรตุ อชฺชตคฺเค ปาณุเปตํ สรณํ คตนฺ”ติ.\csromanspacer{}\csromanspacer{}อฏฺฐมํ.}

\csromanmatikaanchor{mtk.17}
\csromantocmarkref{subsection}{3. สมณพฺราหฺมณสุตฺต}{mtk.17}
\bookmark[dest={mtk.17},level=2]{3. สมณพฺราหฺมณสุตฺต}
\csromanheader{2}{9. พาลปณฺฑิตสุตฺต}
\proseitem{19}{สาวตฺถิยํ วิหรติ.\csromanspacer{} อวิชฺชานีวรณสฺส ภิกฺขเว พาลสฺส ตณฺหาย สมฺปยุตฺตสฺส เอวมยํ กาโย สมุทาคโต.\csromanspacer{} อิติ อยญฺเจว กาโย พหิทฺธา จ นามรูปํ, อิตฺเถตํ ทฺวยํ, ทฺวยํ ปฏิจฺจ ผสฺโส}

\vspace{\tipitakaparskip}
\csromancenter{นิทานสํยุตฺต สเฬวายตนานิ\footnote{สฬายตนานิ (ก)}, เยหิ ผุฏฺโฐ พาโล สุขทุกฺขํ ปฏิสํเวทยติ เอเตสํ วา อญฺญตเรน.}
\prose{\csromanfolio{263}อวิชฺชานีวรณสฺส ภิกฺขเว ปณฺฑิตสฺส ตณฺหาย สมฺปยุตฺตสฺส เอวมยํ กาโย สมุทาคโต.\csromanspacer{} อิติ อยญฺเจว กาโย พหิทฺธา จ นามรูปํ, อิตฺเถตํ ทฺวยํ, ทฺวยํ ปฏิจฺจ ผสฺโส สเฬวายตนานิ, เยหิ ผุฏฺโฐ ปณฺฑิโต สุขทุกฺขํ ปฏิสํเวทยติ เอเตสํ วา อญฺญตเรน.}

\csromanmatikaanchor{mtk.18}
\csromantocmarkref{subsection}{4. ทุติยสมณพฺราหฺมณสุตฺต}{mtk.18}
\bookmark[dest={mtk.18},level=2]{4. ทุติยสมณพฺราหฺมณสุตฺต}
\prose{ตตฺร ภิกฺขเว โก วิเสโส โก อธิปฺปยาโส\footnote{อธิปฺปาโย (สี, อิ, ก), อธิปฺปายโส (สฺยา, กํ), อธิ + ป + ยสุ + ณ + สิ = อธิปฺปยาโส.} กิํ นานากรณํ ปณฺฑิตสฺส พาเลนาติ.\csromanspacer{} ภควํมูลกา โน ภนฺเต ธมฺมา ภควํเนตฺติกา ภควํปฏิสรณา, สาธุ วต ภนฺเต ภควนฺตํเยว ปฏิภาตุ เอตสฺส ภาสิตสฺส อตฺโถ, ภควโต สุตฺวา ภิกฺขู ธาเรสฺสนฺตีติ.}

\prose{เตน หิ ภิกฺขเว สุณาถ สาธุกํ มนสิ กโรถ, ภาสิสฺสามีติ. “เอวํ ภนฺเต”ติ โข เต ภิกฺขู ภควโต ปจฺจสฺโสสุํ.\csromanspacer{} ภควา เอตทโวจ\csromandash{}}

\prose{ยาย จ ภิกฺขเว อวิชฺชาย นิวุตสฺส พาลสฺส ยาย จ ตณฺหาย สมฺปยุตฺตสฺส อยํ กาโย สมุทาคโต, สา เจว อวิชฺชา พาลสฺส อปฺปหีนา, สา จ ตณฺหา อปริกฺขีณา.\csromanspacer{} ตํ กิสฺส เหตุ, น ภิกฺขเว พาโล อจริ พฺรหฺมจริยํ สมฺมา ทุกฺขกฺขยาย.\csromanspacer{} ตสฺมา พาโล กายสฺส เภทา กายูปโค โหติ, โส กายูปโค สมาโน น ปริมุจฺจติ ชาติยา ชรามรเณน โสเกหิ ปริเทเวหิ ทุกฺเขหิ โทมนสฺเสหิ อุปายาเสหิ, “น ปริมุจฺจติ ทุกฺขสฺมา”ติ วทามิ.}

\prose{ยาย จ ภิกฺขเว อวิชฺชาย นิวุตสฺส ปณฺฑิตสฺส ยาย จ ตณฺหาย สมฺปยุตฺตสฺส อยํ กาโย สมุทาคโต, สา เจว อวิชฺชา ปณฺฑิตสฺส ปหีนา, สา จ ตณฺหา ปริกฺขีณา.\csromanspacer{} ตํ กิสฺส เหตุ, อจริ ภิกฺขเว ปณฺฑิโต พฺรหฺมจริยํ สมฺมา ทุกฺขกฺขยาย.\csromanspacer{} ตสฺมา ปณฺฑิโต กายสฺส เภทา น กายูปโค โหติ, โส อกายูปโค สมาโน ปริมุจฺจติ ชาติยา ชรามรเณน โสเกหิ ปริเทเวหิ ทุกฺเขหิ โทมนสฺเสหิ อุปายาเสหิ, “ปริมุจฺจติ ทุกฺขสฺมา”ติ วทามิ.\csromanspacer{} อยํ โข ภิกฺขเว}

\vspace{\tipitakaparskip}
\csromancenter{นิทานวคฺคสํยุตฺตปาฬิ วิเสโส อยํ อธิปฺปยาโส อิทํ นานากรณํ ปณฺฑิตสฺส พาเลน ยทิทํ พฺรหฺมจริยวาโสติ.\csromanspacer{}\csromanspacer{}นวมํ.}
\csromanheader{2}{10. ปจฺจยสุตฺต}
\csromanmatikaanchor{mtk.19}
\csromantocmarkref{subsection}{5. กจฺจานโคตฺตสุตฺต}{mtk.19}
\bookmark[dest={mtk.19},level=2]{5. กจฺจานโคตฺตสุตฺต}
\proseitem{20}{\csromanfolio{264}สาวตฺถิยํ วิหรติ.\csromanspacer{} ปฏิจฺจสมุปฺปาทญฺจ โว ภิกฺขเว เทเสสฺสามิ ปฏิจฺจสมุปฺปนฺเน จ ธมฺเม, ตํ สุณาถ สาธุกํ มนสิ กโรถ, ภาสิสฺสามีติ. “เอวํ ภนฺเต”ติ โข เต ภิกฺขู ภควโต ปจฺจสฺโสสุํ.\csromanspacer{} ภควา เอตทโวจ\csromandash{}}

\prose{กตโม จ ภิกฺขเว ปฏิจฺจสมุปฺปาโท.\csromanspacer{} ชาติปจฺจยา ภิกฺขเว ชรามรณํ, อุปฺปาทา วา ตถาคตานํ อนุปฺปาทา วา ตถาคตานํ ฐิตาว สา ธาตุ ธมฺมฏฺฐิตตา ธมฺมนิยามตา อิทปฺปจฺจยตา, ตํ ตถาคโต อภิสมฺพุชฺฌติ อภิสเมติ, อภิสมฺพุชฺฌิตฺวา อภิสเมตฺวา อาจิกฺขติ เทเสติ ปญฺญาเปติ ปฏฺฐเปติ วิวรติ วิภชติ อุตฺตานีกโรติ “ปสฺสถา”ติ จาห.\csromanspacer{} ชาติปจฺจยา ภิกฺขเว ชรามรณํ.}

\prose{ภวปจฺจยา ภิกฺขเว ชาติ ฯเปฯ.\csromanspacer{} อุปาทานปจฺจยา ภิกฺขเว ภโว.\csromanspacer{} ตณฺหาปจฺจยา ภิกฺขเว อุปาทานํ.\csromanspacer{} เวทนาปจฺจยา ภิกฺขเว ตณฺหา.\csromanspacer{} ผสฺสปจฺจยา ภิกฺขเว เวทนา.\csromanspacer{} สฬายตนปจฺจยา ภิกฺขเว ผสฺโส.\csromanspacer{} นามรูปปจฺจยา ภิกฺขเว สฬายตนํ.\csromanspacer{} วิญฺญาณปจฺจยา ภิกฺขเว นามรูปํ.\csromanspacer{} สงฺขารปจฺจยา ภิกฺขเว วิญฺญาณํ.\csromanspacer{} อวิชฺชาปจฺจยา ภิกฺขเว สงฺขารา อุปฺปาทา วา ตถาคตานํ อนุปฺปาทา วา ตถาคตานํ ฐิตาว สา ธาตุ ธมฺมฏฺฐิตตา ธมฺมนิยามตา อิทปฺปจฺจยตา, ตํ ตถาคโต อภิสมฺพุชฺฌติ อภิสเมติ, อภิสมฺพุชฺฌิตฺวา อภิสเมตฺวา อาจิกฺขติ เทเสติ ปญฺญาเปติ ปฏฺฐเปติ วิวรติ วิภชติ อุตฺตานีกโรติ “ปสฺสถา”ติ จาห.\csromanspacer{} อวิชฺชาปจฺจยา ภิกฺขเว สงฺขารา.\csromanspacer{} อิติ โข ภิกฺขเว ยา ตตฺร ตถตา อวิตถตา อนญฺญถตา อิทปฺปจฺจยตา.\csromanspacer{} อยํ วุจฺจติ ภิกฺขเว ปฏิจฺจสมุปฺปาโท.}

\prose{กตเม จ ภิกฺขเว ปฏิจฺจสมุปฺปนฺนา ธมฺมา.\csromanspacer{} ชรามรณํ ภิกฺขเว อนิจฺจํ สงฺขตํ ปฏิจฺจสมุปฺปนฺนํ ขยธมฺมํ วยธมฺมํ วิราคธมฺมํ นิโรธธมฺมํ.\csromanspacer{} ชาติ ภิกฺขเว อนิจฺจา สงฺขตา ปฏิจฺจสมุปฺปนฺนา ขยธมฺมา วยธมฺมา วิราคธมฺมา นิโรธธมฺมา.\csromanspacer{} ภโว ภิกฺขเว อนิจฺโจ สงฺขโต ปฏิจฺจสมุปฺปนฺโน ขยธมฺโม วยธมฺโม วิราคธมฺโม นิโรธธมฺโม.\csromanspacer{} อุปาทานํ ภิกฺขเว ฯเปฯ.\csromanspacer{} ตณฺหา ภิกฺขเว.\csromanspacer{} เวทนา ภิกฺขเว.\csromanspacer{} ผสฺโส ภิกฺขเว.\csromanspacer{} สฬายตนํ ภิกฺขเว.\csromanspacer{} นามรูปํ ภิกฺขเว.}

\vspace{\tipitakaparskip}
\csromancenter{นิทานสํยุตฺต วิญฺญาณํ ภิกฺขเว.\csromanspacer{} สงฺขารา ภิกฺขเว.\csromanspacer{} อวิชฺชา ภิกฺขเว อนิจฺจา สงฺขตา ปฏิจฺจสมุปฺปนฺนา ขยธมฺมา วยธมฺมา วิราคธมฺมา นิโรธธมฺมา.\csromanspacer{} อิเม วุจฺจนฺติ ภิกฺขเว ปฏิจฺจสมุปฺปนฺนา ธมฺมา.}
\csromanmatikaanchor{mtk.20}
\csromanmatikaanchor{mtk.25}
\csromantocmarkref{subsection}{6. ธมฺมกถิกสุตฺต}{mtk.20}
\bookmark[dest={mtk.20},level=2]{6. ธมฺมกถิกสุตฺต}
\prose{\csromanfolio{265}ยโต โข ภิกฺขเว อริยสาวกสฺส อยญฺจ ปฏิจฺจสมุปฺปาโท อิเม จ ปฏิจฺจสมุปฺปนฺนา ธมฺมา ยถาภูตํ สมฺมปฺปญฺญาย สุทิฏฺฐา โหนฺติ, โส วต ปุพฺพนฺตํ วา ปฏิธาวิสฺสติ “อโหสิํ นุ โข อหํ\footnote{นุ ขฺวาหํ (สฺยา, กํ, อิ, ก)} อตีตมทฺธานํ, นนุ โข อโหสิํ อตีตมทฺธานํ, กิํ นุ โข อโหสิํ อตีตมทฺธานํ, กถํ นุ โข อโหสิํ อตีตมทฺธานํ, กิํ หุตฺวา กิํ อโหสิํ นุ โข อหํ อตีตมทฺธานนฺ”ติ.\csromanspacer{} อปรนฺตํ วา อุปธาวิสฺสติ\footnote{อปธาวิสฺสติ (ก)} “ภวิสฺสามิ นุ โข อหํ อนาคตมทฺธานํ, นนุ โข ภวิสฺสามิ อนาคตมทฺธานํ, กิํ นุ โข ภวิสฺสามิ อนาคตมทฺธานํ, กถํ นุ โข ภวิสฺสามิ อนาคตมทฺธานํ, กิํ หุตฺวา กิํ ภวิสฺสามิ นุ โข อหํ อนาคตมทฺธานนฺ”ติ.\csromanspacer{} เอตรหิ วา ปจฺจุปฺปนฺนํ อทฺธานํ อชฺฌตฺตํ กถํกถี ภวิสฺสติ “อหํ นุ โขสฺมิ, โน นุ โขสฺมิ, กิํ นุ โขสฺมิ, กถํ นุ โขสฺมิ, อยํ นุ โข สตฺโต กุโต อาคโต, โส กุหิํ คมิสฺสตี”ติ เนตํ ฐานํ วิชฺชติ.\csromanspacer{} ตํ กิสฺส เหตุ, ตถา หิ ภิกฺขเว อริยสาวกสฺส อยญฺจ ปฏิจฺจสมุปฺปาโท อิเม จ ปฏิจฺจสมุปฺปนฺนา ธมฺมา ยถาภูตํ สมฺมปฺปญฺญาย สุทิฏฺฐาติ.\csromanspacer{}\csromanspacer{}ทสมํ.}

\vspace{\tipitakaparskip}
\csromancenter{อาหารวคฺโค ทุติโย.}
\titlehead{\textbf{ตสฺสุทฺทานํ}}
\csromangathameasure{อาหารํ ผคฺคุโน เจว, เทฺว จ สมณพฺราหฺมณา. \\ กจฺจานโคตฺต ธมฺมกถิกํ, อเจลํ ติมฺพรุเกน จ. \\ พาลปณฺฑิตโต เจว, ทสโม ปจฺจเยน จาติ.}
\csromangathagroup{\csromangathastanza{อาหารํ ผคฺคุโน เจว, เทฺว จ สมณพฺราหฺมณา. \\ กจฺจานโคตฺต ธมฺมกถิกํ, อเจลํ ติมฺพรุเกน จ. \\ พาลปณฺฑิตโต เจว, ทสโม ปจฺจเยน จาติ.}}

\csromanmatikaanchor{mtk.21}
\csromantocmarkref{subsection}{7. อเจลกสฺสปสุตฺต}{mtk.21}
\bookmark[dest={mtk.21},level=2]{7. อเจลกสฺสปสุตฺต}
\csromanheader{1}{3. ทสพลวคฺค}
\csromanheader{2}{1. ทสพลสุตฺต}
\proseitem{21}{สาวตฺถิยํ วิหรติ.\csromanspacer{} ทสพลสมนฺนาคโต ภิกฺขเว ตถาคโต จตูหิ จ เวสารชฺเชหิ สมนฺนาคโต อาสภํ ฐานํ ปฏิชานาติ,}

\vspace{\tipitakaparskip}
\csromancenter{นิทานวคฺคสํยุตฺตปาฬิ ปริสาสุ สีหนาทํ นทติ, พฺรหฺมจกฺกํ ปวตฺเตติ\csromandash{}อิติ รูปํ อิติ รูปสฺส สมุทโย อิติ รูปสฺส อตฺถงฺคโม, อิติ เวทนา อิติ เวทนาย สมุทโย อิติ เวทนาย อตฺถงฺคโม, อิติ สญฺญา อิติ สญฺญาย สมุทโย อิติ สญฺญาย อตฺถงฺคโม, อิติ สงฺขารา อิติ สงฺขารานํ สมุทโย อิติ สงฺขารานํ อตฺถงฺคโม, อิติ วิญฺญาณํ อิติ วิญฺญาณสฺส สมุทโย อิติ วิญฺญาณสฺส อตฺถงฺคโม, อิติ อิมสฺมิํ สติ อิทํ โหติ, อิมสฺสุปฺปาทา อิทํ อุปฺปชฺชติ, อิมสฺมิํ อสติ อิทํ น โหติ, อิมสฺส นิโรธา อิทํ นิรุชฺฌติ, ยทิทํ อวิชฺชาปจฺจยา สงฺขารา, สงฺขารปจฺจยา วิญฺญาณํ ฯเปฯ เอวเมตสฺส เกวลสฺส ทุกฺขกฺขนฺธสฺส สมุทโย โหติ.\csromanspacer{} อวิชฺชาย เตฺวว อเสสวิราคนิโรธา สงฺขารนิโรโธ, สงฺขารนิโรธา วิญฺญาณนิโรโธ ฯเปฯ เอวเมตสฺส เกวลสฺส ทุกฺขกฺขนฺธสฺส นิโรโธ โหตีติ.\csromanspacer{}\csromanspacer{}ปฐมํ.}
\csromanheader{2}{2. ทุติยทสพลสุตฺต}
\proseitem{22}{\csromanfolio{266}สาวตฺถิยํ วิหรติ.\csromanspacer{} ทสพลสมนฺนาคโต ภิกฺขเว ตถาคโต จตูหิ จ เวสารชฺเชหิ สมนฺนาคโต อาสภํ ฐานํ ปฏิชานาติ, ปริสาสุ สีหนาทํ นทติ, พฺรหฺมจกฺกํ ปวตฺเตติ\csromandash{}อิติ รูปํ อิติ รูปสฺส สมุทโย อิติ รูปสฺส อตฺถงฺคโม, อิติ เวทนา อิติ เวทนาย สมุทโย อิติ เวทนาย อตฺถงฺคโม, อิติ สญฺญา อิติ สญฺญาย สมุทโย อิติ สญฺญาย อตฺถงฺคโม, อิติ สงฺขารา อิติ สงฺขารานํ สมุทโย อิติ สงฺขารานํ อตฺถงฺคโม, อิติ วิญฺญาณํ อิติ วิญฺญาณสฺส สมุทโย อิติ วิญฺญาณสฺส อตฺถงฺคโม, อิติ อิมสฺมิํ สติ อิทํ โหติ, อิมสฺสุปฺปาทา อิทํ อุปฺปชฺชติ, อิมสฺมิํ อสติ อิทํ น โหติ, อิมสฺส นิโรธา อิทํ นิรุชฺฌติ, ยทิทํ อวิชฺชาปจฺจยา สงฺขารา, สงฺขารปจฺจยา วิญฺญาณํ ฯเปฯ เอวเมตสฺส เกวลสฺส ทุกฺขกฺขนฺธสฺส สมุทโย โหติ.\csromanspacer{} อวิชฺชาย เตฺวว อเสสวิราคนิโรธา สงฺขารนิโรโธ, สงฺขารนิโรธา วิญฺญาณนิโรโธ ฯเปฯ เอวเมตสฺส เกวลสฺส ทุกฺขกฺขนฺธสฺส นิโรโธ โหติ.}

\prose{เอวํ สฺวากฺขาโต ภิกฺขเว มยา ธมฺโม อุตฺตาโน วิวโฏ ปกาสิโต ฉินฺนปิโลติโก.\csromanspacer{} เอวํ สฺวากฺขาเต โข ภิกฺขเว มยา ธมฺเม อุตฺตาเน วิวเฏ ปกาสิเต ฉินฺนปิโลติเก อลเมว สทฺธาปพฺพชิเตน กุลปุตฺเตน วีริยํ อารภิตุํ “กามํ ตโจ จ นฺหารุ\footnote{นหารุ (สี, สฺยา, กํ, อิ)} จ}

\vspace{\tipitakaparskip}
\csromancenter{นิทานสํยุตฺต อฏฺฐิ จ อวสิสฺสตุ สรีเร อุปสุสฺสตุ\footnote{อวสุสฺสตุ ม 2 (146) ปิฏฺเฐ.} มํสโลหิตํ.\csromanspacer{} ยํ ตํ ปุริสถาเมน ปุริสวีริเยน ปุริสปรกฺกเมน ปตฺตพฺพํ, น ตํ อปาปุณิตฺวา วีริยสฺส สณฺฐานํ ภวิสฺสตี”ติ.}
\prose{\csromanfolio{267}ทุกฺขํ ภิกฺขเว กุสีโต วิหรติ โวกิณฺโณ ปาปเกหิ อกุสเลหิ ธมฺเมหิ, มหนฺตญฺจ สทตฺถํ ปริหาเปติ, อารทฺธวีริโย จ โข ภิกฺขเว สุขํ วิหรติ ปวิวิตฺโต ปาปเกหิ อกุสเลหิ ธมฺเมหิ, มหนฺตญฺจ สทตฺถํ ปริปูเรติ.\csromanspacer{} น ภิกฺขเว หีเนน อคฺคสฺส ปตฺติ โหติ, อคฺเคน จ โข ภิกฺขเว อคฺคสฺส ปตฺติ โหติ.\csromanspacer{} มณฺฑเปยฺยมิทํ ภิกฺขเว พฺรหฺมจริยํ, สตฺถา สมฺมุขีภูโต.\csromanspacer{} ตสฺมาติห ภิกฺขเว วีริยํ อารภถ อปฺปตฺตสฺส ปตฺติยา อนธิคตสฺส อธิคมาย อสจฺฉิกตสฺส สจฺฉิกิริยาย, “เอวํ โน อยํ อมฺหากํ ปพฺพชฺชา อวญฺฌา ภวิสฺสติ สผลา ส-อุทฺรยา.\csromanspacer{} เยสญฺจ\footnote{เยสํ (สี, สฺยา, กํ), เยสํ หิ (อิ, ก)} มยํ ปริภุญฺชาม จีวรปิณฺฑปาตเสนาสนคิลานปฺปจฺจยเภสชฺชปริกฺขารํ.\csromanspacer{} เตสํ เต การา อเมฺหสุ มหปฺผลา ภวิสฺสนฺติ มหานิสํสา”ติ.\csromanspacer{} เอวํ หิ โว ภิกฺขเว สิกฺขิตพฺพํ.\csromanspacer{} อตฺตตฺถํ วา หิ ภิกฺขเว สมฺปสฺสมาเนน อลเมว อปฺปมาเทน สมฺปาเทตุํ, ปรตฺถํ วา หิ ภิกฺขเว สมฺปสฺสมาเนน อลเมว อปฺปมาเทน สมฺปาเทตุํ, อุภยตฺถํ วา หิ ภิกฺขเว สมฺปสฺสมาเนน อลเมว อปฺปมาเทน สมฺปาเทตุนฺติ.\csromanspacer{}\csromanspacer{}ทุติยํ.}

\csromanheader{2}{3. อุปนิสสุตฺต}
\proseitem{23}{สาวตฺถิยํ วิหรติ.\csromanspacer{} ชานโต อหํ ภิกฺขเว ปสฺสโต อาสวานํ ขยํ วทามิ, โน อชานโต โน อปสฺสโต.\csromanspacer{} กิญฺจ ภิกฺขเว ชานโต กิํ ปสฺสโต อาสวานํ ขโย โหติ.\csromanspacer{} อิติ รูปํ อิติ รูปสฺส สมุทโย อิติ รูปสฺส อตฺถงฺคโม, อิติ เวทนา ฯเปฯ อิติ สญฺญา.\csromanspacer{} อิติ สงฺขารา.\csromanspacer{} อิติ วิญฺญาณํ อิติ วิญฺญาณสฺส สมุทโย อิติ วิญฺญาณสฺส อตฺถงฺคโมติ.\csromanspacer{} เอวํ โข ภิกฺขเว ชานโต เอวํ ปสฺสโต อาสวานํ ขโย โหติ.}

\prose{ยมฺปิสฺส ตํ ภิกฺขเว ขยสฺมิํ ขเยญาณํ, ตมฺปิ ส-อุปนิสํ วทามิ, โน อนุปนิสํ.\csromanspacer{} กา จ ภิกฺขเว ขเยญาณสฺส อุปนิสา, “วิมุตฺตี”ติสฺส วจนียํ.\csromanspacer{} วิมุตฺติมฺปาหํ\footnote{วิมุตฺติมฺปหํ (สี, สฺยา, กํ)} ภิกฺขเว ส-อุปนิสํ วทามิ, โน อนุปนิสํ.\csromanspacer{} กา จ ภิกฺขเว}

\csromanmatikaanchor{mtk.22}
\csromantocmarkref{subsection}{8. ติมฺพรุกสุตฺต}{mtk.22}
\bookmark[dest={mtk.22},level=2]{8. ติมฺพรุกสุตฺต}
\vspace{\tipitakaparskip}
\csromancenter{นิทานวคฺคสํยุตฺตปาฬิ วิมุตฺติยา อุปนิสา, “วิราโค”ติสฺส วจนียํ.\csromanspacer{} วิราคมฺปาหํ ภิกฺขเว ส- อุปนิสํ วทามิ, โน อนุปนิสํ.\csromanspacer{} กา จ ภิกฺขเว วิราคสฺส อุปนิสา, “นิพฺพิทา”ติสฺส วจนียํ.\csromanspacer{} นิพฺพิทมฺปาหํ ภิกฺขเว ส-อุปนิสํ วทามิ, โน อนุปนิสํ.\csromanspacer{} กา จ ภิกฺขเว นิพฺพิทาย อุปนิสา, “ยถาภูตญาณทสฺสนนฺ”ติสฺส วจนียํ.\csromanspacer{} ยถาภูตญาณทสฺสนมฺปาหํ ภิกฺขเว ส-อุปนิสํ วทามิ, โน อนุปนิสํ.\csromanspacer{} กา จ ภิกฺขเว ยถาภูตญาณทสฺสนสฺส อุปนิสา, “สมาธี”ติสฺส วจนียํ.\csromanspacer{} สมาธิมฺปาหํ ภิกฺขเว ส-อุปนิสํ วทามิ, โน อนุปนิสํ.}
\prose{\csromanfolio{268}กา จ ภิกฺขเว สมาธิสฺส อุปนิสา, “สุขนฺ”ติสฺส วจนียํ.\csromanspacer{} สุขมฺปาหํ ภิกฺขเว ส-อุปนิสํ วทามิ, โน อนุปนิสํ.\csromanspacer{} กา จ ภิกฺขเว สุขสฺส อุปนิสา, “ปสฺสทฺธี”ติสฺส วจนียํ.\csromanspacer{} ปสฺสทฺธิมฺปาหํ ภิกฺขเว ส-อุปนิสํ วทามิ, โน อนุปนิสํ.\csromanspacer{} กา จ ภิกฺขเว ปสฺสทฺธิยา อุปนิสา, “ปีตี”ติสฺส วจนียํ.\csromanspacer{} ปีติมฺปาหํ ภิกฺขเว ส-อุปนิสํ วทามิ, โน อนุปนิสํ.\csromanspacer{} กา จ ภิกฺขเว ปีติยา อุปนิสา, “ปาโมชฺชนฺ”ติสฺส วจนียํ.\csromanspacer{} ปาโมชฺชมฺปาหํ ภิกฺขเว ส- อุปนิสํ วทามิ, โน อนุปนิสํ.\csromanspacer{} กา จ ภิกฺขเว ปาโมชฺชสฺส อุปนิสา, “สทฺธา”ติสฺส วจนียํ.\csromanspacer{} สทฺธมฺปาหํ ภิกฺขเว ส-อุปนิสํ วทามิ, โน อนุปนิสํ.}

\prose{กา จ ภิกฺขเว สทฺธาย อุปนิสา, “ทุกฺข”นฺตสฺส วจนียํ.\csromanspacer{} ทุกฺขมฺปาหํ ภิกฺขเว ส-อุปนิสํ วทามิ, โน อนุปนิสํ.\csromanspacer{} กา จ ภิกฺขเว ทุกฺขสฺส อุปนิสา, “ชาตี”ติสฺส วจนียํ.\csromanspacer{} ชาติมฺปาหํ ภิกฺขเว ส-อุปนิสํ วทามิ, โน อนุปนิสํ.\csromanspacer{} กา จ ภิกฺขเว ชาติยา อุปนิสา, “ภโว”ติสฺส วจนียํ.\csromanspacer{} ภวมฺปาหํ ภิกฺขเว ส-อุปนิสํ วทามิ, โน อนุปนิสํ.\csromanspacer{} กา จ ภิกฺขเว ภวสฺส อุปนิสา, “อุปาทานนฺ”ติสฺส วจนียํ.\csromanspacer{} อุปาทานมฺปาหํ ภิกฺขเว ส-อุปนิสํ วทามิ, โน อนุปนิสํ.\csromanspacer{} กา จ ภิกฺขเว อุปาทานสฺส อุปนิสา, “ตณฺหา”ติสฺส วจนียํ.\csromanspacer{} ตณฺหมฺปาหํ ภิกฺขเว ส-อุปนิสํ วทามิ, โน อนุปนิสํ.}

\prose{กา จ ภิกฺขเว ตณฺหาย อุปนิสา, “เวทนา”ติสฺส วจนียํ ฯเปฯ “ผสฺโส”ติสฺส วจนียํ. “สฬายตนนฺ”ติสฺส วจนียํ. “นามรูปนฺ”ติสฺส วจนียํ. “วิญฺญาณนฺ”ติสฺส วจนียํ. “สงฺขารา”ติสฺส วจนียํ.\csromanspacer{} สงฺขาเรปาหํ ภิกฺขเว ส-อุปนิเส วทามิ, โน อนุปนิเส.\csromanspacer{} กา จ ภิกฺขเว สงฺขารานํ อุปนิสา, “อวิชฺชา”ติสฺส วจนียํ.}

\csromanheader{2}{1. นิทานสํยุตฺต}
\prose{\csromanfolio{269}อิติ โข ภิกฺขเว อวิชฺชูปนิสา สงฺขารา, สงฺขารูปนิสํ วิญฺญาณํ, วิญฺญาณูปนิสํ นามรูปํ, นามรูปูปนิสํ สฬายตนํ, สฬายตนูปนิโส ผสฺโส, ผสฺสูปนิสา เวทนา, เวทนูปนิสา ตณฺหา, ตณฺหูปนิสํ อุปาทานํ, อุปาทานูปนิโส ภโว, ภวูปนิสา ชาติ, ชาตูปนิสํ ทุกฺขํ, ทุกฺขูปนิสา สทฺธา, สทฺธูปนิสํ ปาโมชฺชํ, ปาโมชฺชูปนิสา ปีติ, ปีตูปนิสา ปสฺสทฺธิ, ปสฺสทฺธูปนิสํ สุขํ, สุขูปนิโส สมาธิ, สมาธูปนิสํ ยถาภูตญาณทสฺสนํ, ยถาภูตญาณทสฺสนูปนิสา นิพฺพิทา, นิพฺพิทูปนิโส วิราโค, วิราคูปนิสา วิมุตฺติ, วิมุตฺตูปนิสํ ขเยญาณํ.}

\prose{เสยฺยถาปิ ภิกฺขเว อุปริปพฺพเต ถุลฺลผุสิตเก เทเว วสฺสนฺเต ตํ อุทกํ ยถานินฺนํ ปวตฺตมานํ ปพฺพตกนฺทร ปทร สาขา ปริปูเรติ, ปพฺพตกนฺทรปทร สาขา ปริปูรา กุโสพฺเภ\footnote{กุสฺสุพฺเภ (สี, สฺยา, กํ), กุสุพฺเภ (อิ) ณฺวาทิ (129) สุตฺตํ โอโลเกตพฺพํ.} ปริปูเรนฺติ, กุโสพฺภา ปริปูรา มหาโสพฺเภ ปริปูเรนฺติ, มหาโสพฺภา ปริปูรา กุนฺนทิโย ปริปูเรนฺติ, กุนฺนทิโย ปริปูรา มหานทิโย ปริปูเรนฺติ, มหานทิโย ปริปูรา มหาสมุทฺทํ ปริปูเรนฺติ.}

\csromanmatikaanchor{mtk.23}
\csromantocmarkref{subsection}{9. พาลปณฺฑิตสุตฺต}{mtk.23}
\bookmark[dest={mtk.23},level=2]{9. พาลปณฺฑิตสุตฺต}
\prose{เอวเมว โข ภิกฺขเว อวิชฺชูปนิสา สงฺขารา, สงฺขารูปนิสํ วิญฺญาณํ, วิญฺญาณูปนิสํ นามรูปํ, นามรูปูปนิสํ สฬายตนํ, สฬายตนูปนิโส ผสฺโส, ผสฺสูปนิสา เวทนา, เวทนูปนิสา ตณฺหา, ตณฺหูปนิสํ อุปาทานํ, อุปาทานูปนิโส ภโว, ภวูปนิสา ชาติ, ชาตูปนิสํ ทุกฺขํ, ทุกฺขูปนิสา สทฺธา, สทฺธูปนิสํ ปาโมชฺชํ, ปาโมชฺชูปนิสา ปีติ, ปีตูปนิสา ปสฺสทฺธิ, ปสฺสทฺธูปนิสํ สุขํ, สุขูปนิโส สมาธิ, สมาธูปนิสํ ยถาภูตญาณทสฺสนํ, ยถาภูตญาณทสฺสนูปนิสา นิพฺพิทา, นิพฺพิทูปนิโส วิราโค, วิราคูปนิสา วิมุตฺติ, วิมุตฺตูปนิสํ ขเยญาณนฺติ.\csromanspacer{}\csromanspacer{}ตติยํ.}

\csromanheader{2}{4. อญฺญติตฺถิยสุตฺต}
\proseitem{24}{ราชคเห วิหรติ เวฬุวเน.\csromanspacer{} อถ โข อายสฺมา สาริปุตฺโต ปุพฺพณฺหสมยํ นิวาเสตฺวา ปตฺตจีวรมาทาย ราชคหํ ปิณฺฑาย ปาวิสิ.\csromanspacer{} อถ โข อายสฺมโต สาริปุตฺตสฺส เอตทโหสิ “อติปฺปโค โข ตาว ราชคเห ปิณฺฑาย จริตุํ, ยํนูนาหํ เยน อญฺญติตฺถิยานํ ปริพฺพาชกานํ อาราโม เตนุปสงฺกเมยฺยนฺ”ติ}

\gambhira{นิทานวคฺคสํยุตฺตปาฬิ}
\prose{\csromanfolio{270}อถ โข อายสฺมา สาริปุตฺโต เยน อญฺญติตฺถิยานํ ปริพฺพาชกานํ อาราโม เตนุปสงฺกมิ, อุปสงฺกมิตฺวา เตหิ อญฺญติตฺถิเยหิ ปริพฺพาชเกหิ สทฺธิํ สมฺโมทิ, สมฺโมทนียํ กถํ สารณียํ วีติสาเรตฺวา เอกมนฺตํ นิสีทิ, เอกมนฺตํ นิสินฺนํ โข อายสฺมนฺตํ สาริปุตฺตํ เต อญฺญติตฺถิยา ปริพฺพาชกา เอตทโวจุํ\csromandash{}}

\prose{“สนฺตาวุโส สาริปุตฺต เอเก สมณพฺราหฺมณา กมฺมวาทา สยํกตํ ทุกฺขํ ปญฺญเปนฺติ, สนฺติ ปนาวุโส สาริปุตฺต เอเก สมณพฺราหฺมณา กมฺมวาทา ปรํกตํ ทุกฺขํ ปญฺญเปนฺติ, สนฺตาวุโส สาริปุตฺต เอเก สมณพฺราหฺมณา กมฺมวาทา สยํกตญฺจ ปรํกตญฺจ ทุกฺขํ ปญฺญเปนฺติ, สนฺติ ปนาวุโส สาริปุตฺต เอเก สมณพฺราหฺมณา กมฺมวาทา อสยํการํ อปรํการํ อธิจฺจสมุปฺปนฺนํ ทุกฺขํ ปญฺญเปนฺติ.\csromanspacer{} อิธ ปนาวุโส สาริปุตฺต สมโณ โคตโม กิํวาที กิมกฺขายี, กถํ พฺยากรมานา จ มยํ วุตฺตวาทิโน เจว สมณสฺส โคตมสฺส อสฺสาม, น จ สมณํ โคตมํ อภูเตน อพฺภาจิกฺเขยฺยาม, ธมฺมสฺส จานุธมฺมํ พฺยากเรยฺยาม, น จ โกจิ สหธมฺมิโก วาทานุปาโต\footnote{วาทานุวาโท (ก) ที 1 (153) ปิฏฺเฐปิ.} คารยฺหํ ฐานํ อาคจฺเฉยฺยา”ติ.}

\prose{ปฏิจฺจสมุปฺปนฺนํ โข อาวุโส ทุกฺขํ วุตฺตํ ภควตา.\csromanspacer{} กิํ ปฏิจฺจ, ผสฺสํ ปฏิจฺจ.\csromanspacer{} อิติ วทํ วุตฺตวาที เจว ภควโต อสฺส, น จ ภควนฺตํ อภูเตน อพฺภาจิกฺเขยฺย, ธมฺมสฺส จานุธมฺมํ พฺยากเรยฺย, น จ โกจิ สหธมฺมิโก วาทานุปาโต คารยฺหํ ฐานํ อาคจฺเฉยฺย.}

\prose{ตตฺราวุโส เย เต สมณพฺราหฺมณา กมฺมวาทา สยํกตํ ทุกฺขํ ปญฺญเปนฺติ, ตทปิ ผสฺสปจฺจยา.\csromanspacer{} เยปิ เต สมณพฺราหฺมณา กมฺมวาทา ปรํกตํ ทุกฺขํ ปญฺญเปนฺติ, ตทปิ ผสฺสปจฺจยา.\csromanspacer{} เยปิ เต สมณพฺราหฺมณา กมฺมวาทา สยํกตญฺจ ปรํกตญฺจ ทุกฺขํ ปญฺญเปนฺติ, ตทปิ ผสฺสปจฺจยา.\csromanspacer{} เยปิ เต สมณพฺราหฺมณา กมฺมวาทา อสยํการํ อปรํการํ อธิจฺจสมุปฺปนฺนํ ทุกฺขํ ปญฺญเปนฺติ, ตทปิ ผสฺสปจฺจยา.}

\prose{ตตฺราวุโส เย เต สมณพฺราหฺมณา กมฺมวาทา สยํกตํ ทุกฺขํ ปญฺญเปนฺติ, เต วต อญฺญตฺร ผสฺสา ปฏิสํเวทิสฺสนฺตีติ เนตํ ฐานํ วิชฺชติ.\csromanspacer{} เยปิ เต สมณพฺราหฺมณา กมฺมวาทา ปรํกตํ ทุกฺขํ ปญฺญเปนฺติ, เต วต อญฺญตฺร ผสฺสา ปฏิสํเวทิสฺสนฺตีติ เนตํ ฐานํ วิชฺชติ.\csromanspacer{} เยปิ}

\csromanmatikaanchor{mtk.24}
\csromantocmarkref{subsection}{10. ปจฺจยสุตฺต}{mtk.24}
\bookmark[dest={mtk.24},level=2]{10. ปจฺจยสุตฺต}
\csromantocmarkref{section}{อุทฺทานคาถา}{mtk.25}
\bookmark[dest={mtk.25},level=1]{อุทฺทานคาถา}
\vspace{\tipitakaparskip}
\csromancenter{นิทานสํยุตฺต เต สมณพฺราหฺมณา กมฺมวาทา สยํกตญฺจ ปรํกตญฺจ ทุกฺขํ ปญฺญเปนฺติ, เต วต อญฺญตฺร ผสฺสา ปฏิสํเวทิสฺสนฺตีติ เนตํ ฐานํ วิชฺชติ.\csromanspacer{} เยปิ เต สมณพฺราหฺมณา กมฺมวาทา อสยํการํ อปรํการํ อธิจฺจสมุปฺปนฺนํ ทุกฺขํ ปญฺญเปนฺติ, เต วต อญฺญตฺร ผสฺสา ปฏิสํเวทิสฺสนฺตีติ เนตํ ฐานํ วิชฺชตีติ.}
\prose{\csromanfolio{271}อสฺโสสิ โข อายสฺมา อานนฺโท อายสฺมโต สาริปุตฺตสฺส เตหิ อญฺญติตฺถิเยหิ ปริพฺพาชเกหิ สทฺธิํ อิมํ กถาสลฺลาปํ.\csromanspacer{} อถ โข อายสฺมา อานนฺโท ราชคเห ปิณฺฑาย จริตฺวา ปจฺฉาภตฺตํ ปิณฺฑปาตปฏิกฺกนฺโต เยน ภควา เตนุปสงฺกมิ, อุปสงฺกมิตฺวา ภควนฺตํ อภิวาเทตฺวา เอกมนฺตํ นิสีทิ, เอกมนฺตํ นิสินฺโน โข อายสฺมา อานนฺโท ยาวตโก อายสฺมโต สาริปุตฺตสฺส เตหิ อญฺญติตฺถิเยหิ ปริพฺพาชเกหิ สทฺธิํ อโหสิ กถาสลฺลาโป, ตํ สพฺพํ ภควโต อาโรเจสิ.}

\prose{สาธุ สาธุ อานนฺท, ยถา ตํ สาริปุตฺโต สมฺมา พฺยากรมาโน พฺยากเรยฺย.\csromanspacer{} ปฏิจฺจสมุปฺปนฺนํ โข อานนฺท ทุกฺขํ วุตฺตํ มยา.\csromanspacer{} กิํ ปฏิจฺจ, ผสฺสํ ปฏิจฺจ.\csromanspacer{} อิติ วทํ วุตฺตวาที เจว เม อสฺส, น จ มํ อภูเตน อพฺภาจิกฺเขยฺย, ธมฺมสฺส จานุธมฺมํ พฺยากเรยฺย, น จ โกจิ สหธมฺมิโก วาทานุปาโต คารยฺหํ ฐานํ อาคจฺเฉยฺย.}

\prose{ตตฺรานนฺท เย เต สมณพฺราหฺมณา กมฺมวาทา สยํกตํ ทุกฺขํ ปญฺญเปนฺติ, ตทปิ ผสฺสปจฺจยา.\csromanspacer{} เยปิ เต ฯเปฯ.\csromanspacer{} เยปิ เต ฯเปฯ.\csromanspacer{} เยปิ เต สมณพฺราหฺมณา กมฺมวาทา อสยํการํ อปรํการํ อธิจฺจสมุปฺปนฺนํ ทุกฺขํ ปญฺญเปนฺติ, ตทปิ ผสฺสปจฺจยา.}

\prose{ตตฺรานนฺท เยปิ เต สมณพฺราหฺมณา กมฺมวาทา สยํกตํ ทุกฺขํ ปญฺญเปนฺติ, เต วต อญฺญตฺร ผสฺสา ปฏิสํเวทิสฺสนฺตีติ เนตํ ฐานํ วิชฺชติ.\csromanspacer{} เยปิ เต ฯเปฯ.\csromanspacer{} เยปิ เต ฯเปฯ.\csromanspacer{} เยปิ เต สมณพฺราหฺมณา กมฺมวาทา อสยํการํ อปรํการํ อธิจฺจสมุปฺปนฺนํ ทุกฺขํ ปญฺญเปนฺติ, เต วต อญฺญตฺร ผสฺสา ปฏิสํเวทิสฺสนฺตีติ เนตํ ฐานํ วิชฺชติ.}

\prose{เอกมิทาหํ อานนฺท สมยํ อิเธว ราชคเห วิหรามิ เวฬุวเน กลนฺทกนิวาเป.\csromanspacer{} อถ ขฺวาหํ อานนฺท ปุพฺพณฺหสมยํ นิวาเสตฺวา ปตฺตจีวรมาทาย ราชคหํ ปิณฺฑาย ปาวิสิํ, ตสฺส มยฺหํ อานนฺท เอตทโหสิ}

\vspace{\tipitakaparskip}
\csromancenter{นิทานวคฺคสํยุตฺตปาฬิ “อติปฺปโค โข ตาว ราชคเห ปิณฺฑาย จริตุํ, ยํนูนาหํ เยน อญฺญติตฺถิยานํ ปริพฺพาชกานํ อาราโม เตนุปสงฺกเมยฺยนฺ”ติ.}
\prose{\csromanfolio{272}อถ ขฺวาหํ อานนฺท เยน อญฺญติตฺถิยานํ ปริพฺพาชกานํ อาราโม เตนุปสงฺกมิํ, อุปสงฺกมิตฺวา เตหิ อญฺญติตฺถิเยหิ ปริพฺพาชเกหิ สทฺธิํ สมฺโมทิํ, สมฺโมทนียํ กถํ สารณียํ วีติสาเรตฺวา เอกมนฺตํ นิสีทิํ, เอกมนฺตํ นิสินฺนํ โข มํ อานนฺท เต อญฺญติตฺถิยา ปริพฺพาชกา เอตทโวจุํ\csromandash{}}

\prose{“สนฺตาวุโส โคตม เอเก สมณพฺราหฺมณา กมฺมวาทา สยํกตํ ทุกฺขํ ปญฺญเปนฺติ, สนฺติ ปนาวุโส โคตม เอเก สมณพฺราหฺมณา กมฺมวาทา ปรํกตํ ทุกฺขํ ปญฺญเปนฺติ, สนฺตาวุโส โคตม เอเก สมณพฺราหฺมณา กมฺมวาทา สยํกตญฺจ ปรํกตญฺจ ทุกฺขํ ปญฺญเปนฺติ, สนฺติ ปนาวุโส โคตม เอเก สมณพฺราหฺมณา กมฺมวาทา อสยํการํ อปรํการํ อธิจฺจสมุปฺปนฺนํ ทุกฺขํ ปญฺญเปนฺติ.\csromanspacer{} อิธ โน อายสฺมา โคตโม กิํวาที กิมกฺขายี, กถํ พฺยากรมานา จ มยํ วุตฺตวาทิโน เจว อายสฺมโต โคตมสฺส อสฺสาม, น จ อายสฺมนฺตํ โคตมํ อภูเตน อพฺภาจิกฺเขยฺยาม, ธมฺมสฺส จานุธมฺมํ พฺยากเรยฺยาม, น จ โกจิ สหธมฺมิโก วาทานุปาโต คารยฺหํ ฐานํ อาคจฺเฉยฺยา”ติ.}

\prose{เอวํ วุตฺตาหํ อานนฺท เต อญฺญติตฺถิเย ปริพฺพาชเก เอตทโวจํ “ปฏิจฺจสมุปฺปนฺนํ โข อาวุโส ทุกฺขํ วุตฺตํ มยา.\csromanspacer{} กิํ ปฏิจฺจ, ผสฺสํ ปฏิจฺจ.\csromanspacer{} อิติ วทํ วุตฺตวาที เจว เม อสฺส, น จ มํ อภูเตน อพฺภาจิกฺเขยฺย, ธมฺมสฺส จานุธมฺมํ พฺยากเรยฺย, น จ โกจิ สหธมฺมิโก วาทานุปาโต คารยฺหํ ฐานํ อาคจฺเฉยฺยา”ติ.}

\csromanmatikaanchor{mtk.26}
\csromantocmarknopage{section}{3. ทสพลวคฺค}{mtk.26}
\bookmark[dest={mtk.26},level=1]{3. ทสพลวคฺค}
\prose{ตตฺราวุโส เย เต สมณพฺราหฺมณา กมฺมวาทา สยํกตํ ทุกฺขํ ปญฺญเปนฺติ, ตทปิ ผสฺสปจฺจยา.\csromanspacer{} เยปิ เต ฯเปฯ.\csromanspacer{} เยปิ เต ฯเปฯ.\csromanspacer{} เยปิ เต สมณพฺราหฺมณา กมฺมวาทา อสยํการํ อปรํการํ อธิจฺจสมุปฺปนฺนํ ทุกฺขํ ปญฺญเปนฺติ, ตทปิ ผสฺสปจฺจยา.}

\csromanmatikaanchor{mtk.27}
\csromantocmarkref{subsection}{1. ทสพลสุตฺต}{mtk.27}
\bookmark[dest={mtk.27},level=2]{1. ทสพลสุตฺต}
\prose{ตตฺราวุโส เย เต สมณพฺราหฺมณา กมฺมวาทา สยํกตํ ทุกฺขํ ปญฺญเปนฺติ, เต วต อญฺญตฺร ผสฺสา ปฏิสํเวทิสฺสนฺตีติ เนตํ ฐานํ วิชฺชติ.\csromanspacer{} เยปิ เต ฯเปฯ.\csromanspacer{} เยปิ เต ฯเปฯ.\csromanspacer{} เยปิ เต สมณพฺราหฺมณา กมฺมวาทา อสยํการํ อปรํการํ อธิจฺจสมุปฺปนฺนํ ทุกฺขํ ปญฺญเปนฺติ, เต วต อญฺญตฺร ผสฺสา ปฏิสํเวทิสฺสนฺตีติ เนตํ ฐานํ วิชฺชตีติ.\csromanspacer{} อจฺฉริยํ ภนฺเต อพฺภุตํ}

\vspace{\tipitakaparskip}
\csromancenter{นิทานสํยุตฺต ภนฺเต, ยตฺร หิ นาม เอเกน ปเทน สพฺโพ อตฺโถ วุตฺโต ภวิสฺสติ.\csromanspacer{} สิยา นุ โข ภนฺเต เอเสวตฺโถ วิตฺถาเรน วุจฺจมาโน คมฺภีโร เจว อสฺส คมฺภีราวภาโส จาติ.}
\prose{\csromanfolio{273}เตน หานนฺท ตญฺเญเวตฺถ ปฏิภาตูติ.\csromanspacer{} สเจ มํ ภนฺเต เอวํ ปุจฺเฉยฺยุํ “ชรามรณํ อาวุโส อานนฺท กิํนิทานํ กิํสมุทยํ กิํชาติกํ กิํปภวนฺ”ติ.\csromanspacer{} เอวํ ปุฏฺโฐหํ ภนฺเต เอวํ พฺยากเรยฺยํ “ชรามรณํ โข อาวุโส ชาตินิทานํ ชาติสมุทยํ ชาติชาติกํ ชาติปภวนฺ”ติ.\csromanspacer{} เอวํ ปุฏฺโฐหํ ภนฺเต เอวํ พฺยากเรยฺยํ.}

\csromanmatikaanchor{mtk.28}
\csromantocmarkref{subsection}{2. ทุติยทสพลสุตฺต}{mtk.28}
\bookmark[dest={mtk.28},level=2]{2. ทุติยทสพลสุตฺต}
\prose{สเจ มํ ภนฺเต เอวํ ปุจฺเฉยฺยุํ “ชาติ ปนาวุโส อานนฺท กิํนิทานา กิํสมุทยา กิํชาติกา กิํปภวา”ติ.\csromanspacer{} เอวํ ปุฏฺโฐหํ ภนฺเต เอวํ พฺยากเรยฺยํ “ชาติ โข อาวุโส ภวนิทานา ภวสมุทยา ภวชาติกา ภวปฺปภวา”ติ.\csromanspacer{} เอวํ ปุฏฺโฐหํ ภนฺเต เอวํ พฺยากเรยฺยํ.}

\prose{สเจ มํ ภนฺเต เอวํ ปุจฺเฉยฺยุํ “ภโว ปนาวุโส อานนฺท กิํนิทาโน กิํสมุทโย กิํชาติโก กิํปภโว”ติ.\csromanspacer{} เอวํ ปุฏฺโฐหํ ภนฺเต เอวํ พฺยากเรยฺยํ “ภโว โข อาวุโส อุปาทานนิทาโน อุปาทานสมุทโย อุปาทานชาติโก อุปาทานปฺปภโว”ติ.\csromanspacer{} เอวํ ปุฏฺโฐหํ ภนฺเต เอวํ พฺยากเรยฺยํ.}

\prose{สเจ มํ ภนฺเต เอวํ ปุจฺเฉยฺยํ “อุปาทานํ ปนาวุโส ฯเปฯ ตณฺหา ปนาวุโส ฯเปฯ เวทนา ปนาวุโส ฯเปฯ.\csromanspacer{} สเจ มํ ภนฺเต เอวํ ปุจฺเฉยฺยุํ “ผสฺโส ปนาวุโส อานนฺท กิํนิทาโน กิํสมุทโย กิํชาติโก กิํปภโว”ติ.\csromanspacer{} เอวํ ปุฏฺโฐหํ ภนฺเต เอวํ พฺยากเรยฺยํ “ผสฺโส โข อาวุโส สฬายตนนิทาโน สฬายตนสมุทโย สฬายตนชาติโก สฬายตนปฺปภโว”ติ. “ฉนฺนํเตฺวว อาวุโส ผสฺสายตนานํ อเสสวิราคนิโรธา ผสฺสนิโรโธ, ผสฺสนิโรธา เวทนานิโรโธ, เวทนานิโรธา ตณฺหานิโรโธ, ตณฺหานิโรธา อุปาทานนิโรโธ, อุปาทานนิโรธา ภวนิโรโธ, ภวนิโรธา ชาตินิโรโธ, ชาตินิโรธา ชรามรณํ โสกปริเทวทุกฺขโทมนสฺสุปายาสา นิรุชฺฌนฺติ, เอวเมตสฺส เกวลสฺส ทุกฺขกฺขนฺธสฺส นิโรโธ โหตี”ติ.\csromanspacer{} เอวํ ปุฏฺโฐหํ ภนฺเต เอวํ พฺยากเรยฺยนฺติ.\csromanspacer{}\csromanspacer{}จตุตฺถํ.}

\csromanheader{2}{นิทานวคฺคสํยุตฺตปาฬิ \textbf{5. ภูมิชสุตฺต}}
\proseitem{25}{\csromanfolio{274}สาวตฺถิยํ วิหรติ.\csromanspacer{} อถ โข อายสฺมา ภูมิโช สายนฺหสมยํ ปฏิสลฺลานา วุฏฺฐิโต เยนายสฺมา สาริปุตฺโต เตนุปสงฺกมิ, อุปสงฺกมิตฺวา อายสฺมตา สาริปุตฺเตน สทฺธิํ สมฺโมทิ, สมฺโมทนียํ กถํ สารณียํ วีติสาเรตฺวา เอกมนฺตํ นิสีทิ, เอกมนฺตํ นิสินฺโน โข อายสฺมา ภูมิโช อายสฺมนฺตํ สาริปุตฺตํ เอตทโวจ\csromandash{}}

\csromanmatikaanchor{mtk.29}
\csromantocmarkref{subsection}{3. อุปนิสสุตฺต}{mtk.29}
\bookmark[dest={mtk.29},level=2]{3. อุปนิสสุตฺต}
\prose{“สนฺตาวุโส สาริปุตฺต เอเก สมณพฺราหฺมณา กมฺมวาทา สยํกตํ สุขทุกฺขํ ปญฺญเปนฺติ, สนฺติ ปนาวุโส สาริปุตฺต เอเก สมณพฺราหฺมณา กมฺมวาทา ปรํกตํ สุขทุกฺขํ ปญฺญเปนฺติ, สนฺตาวุโส สาริปุตฺต เอเก สมณพฺราหฺมณา กมฺมวาทา สยํกตญฺจ ปรํกตญฺจ สุขทุกฺขํ ปญฺญเปนฺติ, สนฺติ ปนาวุโส สาริปุตฺต เอเก สมณพฺราหฺมณา กมฺมวาทา อสยํการํ อปรํการํ อธิจฺจสมุปฺปนฺนํ สุขทุกฺขํ ปญฺญเปนฺติ.\csromanspacer{} อิธ โน อาวุโส สาริปุตฺต ภควา กิํวาที กิมกฺขายี, กถํ พฺยากรมานา จ มยํ วุตฺตวาทิโน เจว ภควโต อสฺสาม, น จ ภควนฺตํ อภูเตน อพฺภาจิกฺเขยฺยาม, ธมฺมสฺส จานุธมฺมํ พฺยากเรยฺยาม, น จ โกจิ สหธมฺมิโก วาทานุปาโต คารยฺหํ ฐานํ อาคจฺเฉยฺยา”ติ.}

\prose{ปฏิจฺจสมุปฺปนฺนํ โข อาวุโส สุขทุกฺขํ วุตฺตํ ภควตา.\csromanspacer{} กิํ ปฏิจฺจ, ผสฺสํ ปฏิจฺจ.\csromanspacer{} อิติ วทํ วุตฺตวาที เจว ภควโต อสฺส, น จ ภควนฺตํ อภูเตน อพฺภาจิกฺเขยฺย, ธมฺมสฺส จานุธมฺมํ พฺยากเรยฺย, น จ โกจิ สหธมฺมิโก วาทานุปาโต คารยฺหํ ฐานํ อาคจฺเฉยฺย.}

\prose{ตตฺราวุโส เย เต สมณพฺราหฺมณา กมฺมวาทา สยํกตํ สุขทุกฺขํ ปญฺญเปนฺติ, ตทปิ ผสฺสปจฺจยา.\csromanspacer{} เยปิ เต ฯเปฯ.\csromanspacer{} เยปิ เต ฯเปฯ.\csromanspacer{} เยปิ เต สมณพฺราหฺมณา กมฺมวาทา อสยํการํ อปรํการํ อธิจฺจสมุปฺปนฺนํ สุขทุกฺขํ ปญฺญเปนฺติ, ตทปิ ผสฺสปจฺจยา.}

\prose{ตตฺราวุโส เย เต สมณพฺราหฺมณา กมฺมวาทา สยํกตํ สุขทุกฺขํ ปญฺญเปนฺติ, เต วต อญฺญตฺร ผสฺสา ปฏิสํเวทิสฺสนฺตีติ เนตํ ฐานํ วิชฺชติ.\csromanspacer{} เยปิ เต ฯเปฯ.\csromanspacer{} เยปิ เต ฯเปฯ.\csromanspacer{} เยปิ เต สมณพฺราหฺมณา กมฺมวาทา อสยํการํ อปรํการํ อธิจฺจสมุปฺปนฺนํ สุขทุกฺขํ ปญฺญเปนฺติ, เต วต อญฺญตฺร ผสฺสา ปฏิสํเวทิสฺสนฺตีติ เนตํ ฐานํ วิชฺชตีติ.}

\csromanheader{2}{1. นิทานสํยุตฺต}
\prose{\csromanfolio{275}อสฺโสสิ โข อายสฺมา อานนฺโท อายสฺมโต สาริปุตฺตสฺส อายสฺมตา ภูมิเชน สทฺธิํ อิมํ กถาสลฺลาปํ.\csromanspacer{} อถ โข อายสฺมา อานนฺโท เยน ภควา เตนุปสงฺกมิ, อุปสงฺกมิตฺวา ภควนฺตํ อภิวาเทตฺวา เอกมนฺตํ นิสีทิ, เอกมนฺตํ นิสินฺโน โข อายสฺมา อานนฺโท ยาวตโก อายสฺมโต สาริปุตฺตสฺส อายสฺมตา ภูมิเชน สทฺธิํ อโหสิ กถาสลฺลาโป, ตํ สพฺพํ ภควโต อาโรเจสิ.}

\prose{สาธุ สาธุ อานนฺท, ยถา ตํ สาริปุตฺโต สมฺมา พฺยากรมาโน พฺยากเรยฺย.\csromanspacer{} ปฏิจฺจสมุปฺปนฺนํ โข อานนฺท สุขทุกฺขํ วุตฺตํ มยา.\csromanspacer{} กิํ ปฏิจฺจ, ผสฺสํ ปฏิจฺจ.\csromanspacer{} อิติ วทํ วุตฺตวาที เจว เม อสฺส, น จ มํ อภูเตน อพฺภาจิกฺเขยฺย, ธมฺมสฺส จานุธมฺมํ พฺยากเรยฺย, น จ โกจิ สหธมฺมิโก วาทานุปาโต คารยฺหํ ฐานํ อาคจฺเฉยฺย.}

\prose{ตตฺรานนฺท เย เต สมณพฺราหฺมณา กมฺมวาทา สยํกตํ สุขทุกฺขํ ปญฺญเปนฺติ, ตทปิ ผสฺสปจฺจยา.\csromanspacer{} เยปิ เต ฯเปฯ.\csromanspacer{} เยปิ เต ฯเปฯ.\csromanspacer{} เยปิ เต สมณพฺราหฺมณา กมฺมวาทา อสยํการํ อปรํการํ อธิจฺจสมุปฺปนฺนํ สุขทุกฺขํ ปญฺญเปนฺติ, ตทปิ ผสฺสปจฺจยา.}

\prose{ตตฺรานนฺท เย เต สมณพฺราหฺมณา กมฺมวาทา สยํกตํ สุขทุกฺขํ ปญฺญเปนฺติ, เต วต อญฺญตฺร ผสฺสา ปฏิสํเวทิสฺสนฺตีติ เนตํ ฐานํ วิชฺชติ.\csromanspacer{} เยปิ เต ฯเปฯ.\csromanspacer{} เยปิ เต ฯเปฯ.\csromanspacer{} เยปิ เต สมณพฺราหฺมณา กมฺมวาทา อสยํการํ อปรํการํ อธิจฺจสมุปฺปนฺนํ สุขทุกฺขํ ปญฺญเปนฺติ, เต วต อญฺญตฺร ผสฺสา ปฏิสํเวทิสฺสนฺตีติ เนตํ ฐานํ วิชฺชติ.}

\prose{กาเย วา หานนฺท สติ กายสญฺเจตนาเหตุ อุปฺปชฺชติ อชฺฌตฺตํ สุขทุกฺขํ.\csromanspacer{} วาจาย วา หานนฺท สติ วจีสญฺเจตนาเหตุ อุปฺปชฺชติ อชฺฌตฺตํ สุขทุกฺขํ.\csromanspacer{} มเน วา หานนฺท สติ มโนสญฺเจตนาเหตุ อุปฺปชฺชติ อชฺฌตฺตํ สุขทุกฺขํ, อวิชฺชาปจฺจยา จ.}

\prose{สามํ วา ตํ อานนฺท กายสงฺขารํ อภิสงฺขโรติ, ยํปจฺจยา’สฺส\footnote{ยํปจฺจยาย (สฺยา, กํ), ยํปจฺจยา ยํ (ก)} ตํ อุปฺปชฺชติ อชฺฌตฺตํ สุขทุกฺขํ.\csromanspacer{} ปเร วา ตํ\footnote{ปเร วาสฺส ตํ (สี, อิ), ปเร วายตํ (สฺยา, กํ)} อานนฺท กายสงฺขารํ อภิสงฺขโรนฺติ, ยํปจฺจยา’สฺส ตํ อุปฺปชฺชติ อชฺฌตฺตํ สุขทุกฺขํ.\csromanspacer{} สมฺปชาโน วา ตํ อานนฺท กายสงฺขารํ อภิสงฺขโรติ, ยํปจฺจยา’สฺส ตํ อุปฺปชฺชติ อชฺฌตฺตํ สุขทุกฺขํ.}

\csromanmatikaanchor{mtk.30}
\csromantocmarkref{subsection}{4. อญฺญติตฺถิยสุตฺต}{mtk.30}
\bookmark[dest={mtk.30},level=2]{4. อญฺญติตฺถิยสุตฺต}
\vspace{\tipitakaparskip}
\csromancenter{นิทานวคฺคสํยุตฺตปาฬิ อสมฺปชาโน วา ตํ อานนฺท กายสงฺขารํ อภิสงฺขโรติ, ยํปจฺจยา’สฺส ตํ อุปฺปชฺชติ อชฺฌตฺตํ สุขทุกฺขํ.}
\prose{\csromanfolio{276}สามํ วา ตํ อานนฺท วจีสงฺขารํ อภิสงฺขโรติ, ยํปจฺจยา’สฺส ตํ อุปฺปชฺชติ อชฺฌตฺตํ สุขทุกฺขํ.\csromanspacer{} ปเร วา ตํ อานนฺท วจีสงฺขารํ อภิสงฺขโรนฺติ, ยํปจฺจยา’สฺส ตํ อุปฺปชฺชติ อชฺฌตฺตํ สุขทุกฺขํ.\csromanspacer{} สมฺปชาโน วา ตํ อานนฺท ฯเปฯ.\csromanspacer{} อสมฺปชาโน วา ตํ อานนฺท วจีสงฺขารํ อภิสงฺขโรติ, ยํปจฺจยา’สฺส ตํ อุปฺปชฺชติ อชฺฌตฺตํ สุขทุกฺขํ.}

\prose{สามํ วา ตํ อานนฺท มโนสงฺขารํ อภิสงฺขโรติ, ยํปจฺจยา’สฺส ตํ อุปฺปชฺชติ อชฺฌตฺตํ สุขทุกฺขํ.\csromanspacer{} ปเร วา ตํ อานนฺท มโนสงฺขารํ อภิสงฺขเรนฺติ, ยํปจฺจยา’สฺส ตํ อุปฺปชฺชติ อชฺฌตฺตํ สุขทุกฺขํ.\csromanspacer{} สมฺปชาโน วา ตํ อานนฺท ฯเปฯ.\csromanspacer{} อสมฺปชาโน วา ตํ อานนฺท มโนสงฺขารํ อภิสงฺขโรติ, ยํปจฺจยา’สฺส ตํ อุปฺปชฺชติ อชฺฌตฺตํ สุขทุกฺขํ.}

\prose{อิเมสุ อานนฺท ธมฺเมสุ อวิชฺชา อนุปติตา, อวิชฺชาย เตฺวว อานนฺท อเสสวิราคนิโรธา โส กาโย น โหติ, ยํปจฺจยา’สฺส ตํ อุปฺปชฺชติ อชฺฌตฺตํ สุขทุกฺขํ.\csromanspacer{} สา วาจา น โหติ, ยํปจฺจยา’สฺส ตํ อุปฺปชฺชติ อชฺฌตฺตํ สุขทุกฺขํ.\csromanspacer{} โส มโน น โหติ, ยํปจฺจยา’สฺส ตํ อุปฺปชฺชติ อชฺฌตฺตํ สุขทุกฺขํ.\csromanspacer{} เขตฺตํ ตํ น โหติ ฯเปฯ.\csromanspacer{} วตฺถุ ตํ น โหติ ฯเปฯ.\csromanspacer{} อายตนํ ตํ น โหติ ฯเปฯ.\csromanspacer{} อธิกรณํ ตํ น โหติ, ยํปจฺจยา’สฺส ตํ อุปฺปชฺชติ อชฺฌตฺตํ สุขทุกฺขนฺติ.\csromanspacer{}\csromanspacer{}ปญฺจมํ.}

\csromanheader{2}{6. อุปวาณสุตฺต}
\proseitem{26}{สาวตฺถิยํ วิหรติ.\csromanspacer{} อถ โข อายสฺมา อุปวาโณ เยน ภควา เตนุปสงฺกมิ, อุปสงฺกมิตฺวา ภควนฺตํ อภิวาเทตฺวา เอกมนฺตํ นิสีทิ, เอกมนฺตํ นิสินฺโน โข อายสฺมา อุปวาโณ ภควนฺตํ เอตทโวจ\csromandash{}}

\prose{“สนฺติ ภนฺเต เอเก สมณพฺราหฺมณา สยํกตํ ทุกฺขํ ปญฺญเปนฺติ, สนฺติ ปน ภนฺเต เอเก สมณพฺราหฺมณา ปรํกตํ ทุกฺขํ ปญฺญเปนฺติ, สนฺติ ปน ภนฺเต เอเก สมณพฺราหฺมณา สยํกตญฺจ ปรํกตญฺจ ทุกฺขํ ปญฺญเปนฺติ, สนฺติ ปน ภนฺเต เอเก สมณพฺราหฺมณา อสยํการํ อปรํการํ อธิจฺจสมุปฺปนฺนํ ทุกฺขํ ปญฺญเปนฺติ.\csromanspacer{} อิธ โน ภนฺเต ภควา กิํวาที กิมกฺขายี, กถํ พฺยากรมานา จ มยํ วุตฺตวาทิโน เจว ภควโต อสฺสาม, น จ ภควนฺตํ}

\vspace{\tipitakaparskip}
\csromancenter{นิทานสํยุตฺต อภูเตน อพฺภาจิกฺเขยฺยาม, ธมฺมสฺส จานุธมฺมํ พฺยากเรยฺยาม, น จ โกจิ สหธมฺมิโก วาทานุปาโต คารยฺหํ ฐานํ อาคจฺเฉยฺยา”ติ.}
\prose{\csromanfolio{277}ปฏิจฺจสมุปฺปนฺนํ โข อุปวาณ ทุกฺขํ วุตฺตํ มยา.\csromanspacer{} กิํ ปฏิจฺจ, ผสฺสํ ปฏิจฺจ.\csromanspacer{} อิติ วทํ วุตฺตวาที เจว เม อสฺส, น จ มํ อภูเตน อพฺภาจิกฺเขยฺย, ธมฺมสฺส จานุธมฺมํ พฺยากเรยฺย, น จ โกจิ สหธมฺมิโก วาทานุปาโต คารยฺหํ ฐานํ อาคจฺเฉยฺย.}

\prose{ตตฺร อุปวาณ เย เต สมณพฺราหฺมณา สยํกตํ ทุกฺขํ ปญฺญเปนฺติ, ตทปิ ผสฺสปจฺจยา.\csromanspacer{} เยปิ เต ฯเปฯ.\csromanspacer{} เยปิ เต ฯเปฯ.\csromanspacer{} เยปิ เต สมณพฺราหฺมณา อสยํการํ อปรํการํ อธิจฺจสมุปฺปนฺนํ ทุกฺขํ ปญฺญเปนฺติ, ตทปิ ผสฺสปจฺจยา.}

\prose{ตตฺร อุปวาณ เย เต สมณพฺราหฺมณา สยํกตํ ทุกฺขํ ปญฺญเปนฺติ, เต วต อญฺญตฺร ผสฺสา ปฏิสํเวทิสฺสนฺตีติ เนตํ ฐานํ วิชฺชติ.\csromanspacer{} เยปิ เต ฯเปฯ.\csromanspacer{} เยปิ เต ฯเปฯ.\csromanspacer{} เยปิ เต สมณพฺราหฺมณา อสยํการํ อปรํการํ อธิจฺจสมุปฺปนฺนํ ทุกฺขํ ปญฺญเปนฺติ, เต วต อญฺญตฺร ผสฺสา ปฏิสํเวทิสฺสนฺตีติ เนตํ ฐานํ วิชฺชตีติ.\csromanspacer{}\csromanspacer{}ฉฏฺฐํ.}

\csromanheader{2}{7. ปจฺจยสุตฺต}
\proseitem{27}{สาวตฺถิยํ วิหรติ.\csromanspacer{} อวิชฺชาปจฺจยา ภิกฺขเว สงฺขารา, สงฺขารปจฺจยา วิญฺญาณํ ฯเปฯ เอวเมตสฺส เกวลสฺส ทุกฺขกฺขนฺธสฺส สมุทโย โหติ.}

\prose{กตมญฺจ ภิกฺขเว ชรามรณํ, ยา เตสํ เตสํ สตฺตานํ ตมฺหิ ตมฺหิ สตฺตนิกาเย ชรา ชีรณตา ขณฺฑิจฺจํ ปาลิจฺจํ วลิตฺตจตา อายุโน สํหานิ อินฺทฺริยานํ ปริปาโก.\csromanspacer{} อยํ วุจฺจติ ชรา.\csromanspacer{} ยา เตสํ เตสํ สตฺตานํ ตมฺหา ตมฺหา สตฺตนิกายา จุติ จวนตา เภโท อนฺตรธานํ มจฺจุ มรณํ กาลกิริยา ขนฺธานํ เภโท กเฬวรสฺส นิกฺเขโป, อิทํ วุจฺจติ มรณํ, อิติ อยญฺจ ชรา อิทญฺจ มรณํ, อิทํ วุจฺจติ ภิกฺขเว ชรามรณํ.\csromanspacer{} ชาติสมุทยา ชรามรณสมุทโย, ชาตินิโรธา ชรามรณนิโรโธ.\csromanspacer{} อยเมว อริโย อฏฺฐงฺคิโก มคฺโค ชรามรณนิโรธคามินี ปฏิปทา.\csromanspacer{} เสยฺยถิทํ, สมฺมาทิฏฺฐิ สมฺมาสงฺกปฺโป สมฺมาวาจา สมฺมากมฺมนฺโต สมฺมา-อาชีโว สมฺมาวายาโม สมฺมาสติ สมฺมาสมาธิ.}

\gambhira{นิทานวคฺคสํยุตฺตปาฬิ}
\prose{\csromanfolio{278}กตมา จ ภิกฺขเว ชาติ ฯเปฯ.\csromanspacer{} กตโม จ ภิกฺขเว ภโว.\csromanspacer{} กตมญฺจ ภิกฺขเว อุปาทานํ.\csromanspacer{} กตมา จ ภิกฺขเว ตณฺหา.\csromanspacer{} กตมา จ ภิกฺขเว เวทนา.\csromanspacer{} กตโม จ ภิกฺขเว ผสฺโส.\csromanspacer{} กตมญฺจ ภิกฺขเว สฬายตนํ.\csromanspacer{} กตมญฺจ ภิกฺขเว นามรูปํ.\csromanspacer{} กตมญฺจ ภิกฺขเว วิญฺญาณํ.}

\prose{กตเม จ ภิกฺขเว สงฺขารา.\csromanspacer{} ตโยเม ภิกฺขเว สงฺขารา, กายสงฺขาโร วจีสงฺขาโร จิตฺตสงฺขาโร.\csromanspacer{} อิเม วุจฺจนฺติ ภิกฺขเว สงฺขารา.\csromanspacer{} อวิชฺชาสมุทยา สงฺขารสมุทโย, อวิชฺชานิโรธา สงฺขารนิโรโธ, อยเมว อริโย อฏฺฐงฺคิโก มคฺโค สงฺขารนิโรธคามินี ปฏิปทา.\csromanspacer{} เสยฺยถิทํ, สมฺมาทิฏฺฐิ ฯเปฯ สมฺมาสมาธิ.}

\prose{ยโต โข ภิกฺขเว อริยสาวโก เอวํ ปจฺจยํ ปชานาติ, เอวํ ปจฺจยสมุทยํ ปชานาติ, เอวํ ปจฺจยนิโรธํ ปชานาติ, เอวํ ปจฺจยนิโรธคามินิํ ปฏิปทํ ปชานาติ.\csromanspacer{} อยํ วุจฺจติ ภิกฺขเว อริยสาวโก ทิฏฺฐิสมฺปนฺโน อิติปิ, ทสฺสนสมฺปนฺโน อิติปิ, อาคโต อิมํ สทฺธมฺมํ อิติปิ, ปสฺสติ อิมํ สทฺธมฺมํ อิติปิ, เสกฺเขน ญาเณน สมนฺนาคโต อิติปิ, เสกฺขาย วิชฺชาย สมนฺนาคโต อิติปิ, ธมฺมโสตํ สมาปนฺโน อิติปิ, อริโย นิพฺเพธิกปญฺโญ อิติปิ, อมตทฺวารํ อาหจฺจ ติฏฺฐติ อิติปีติ.\csromanspacer{}\csromanspacer{}สตฺตมํ.}

\csromanheader{2}{8. ภิกฺขุสุตฺต}
\proseitem{28}{สาวตฺถิยํ วิหรติ.\csromanspacer{} ตตฺร โข ฯเปฯ.\csromanspacer{} อิธ ภิกฺขเว ภิกฺขุ ชรามรณํ ปชานาติ, ชรามรณสมุทยํ ปชานาติ, ชรามรณนิโรธํ ปชานาติ, ชรามรณนิโรธคามินิํ ปฏิปทํ ปชานาติ.\csromanspacer{} ชาติํ ปชานาติ ฯเปฯ.\csromanspacer{} ภวํ ปชานาติ.\csromanspacer{} อุปาทานํ ปชานาติ.\csromanspacer{} ตณฺหํ ปชานาติ.\csromanspacer{} เวทนํ ปชานาติ.\csromanspacer{} ผสฺสํ ปชานาติ.\csromanspacer{} สฬายตนํ ปชานาติ.\csromanspacer{} นามรูปํ ปชานาติ.\csromanspacer{} วิญฺญาณํ ปชานาติ.\csromanspacer{} สงฺขาเร ปชานาติ, สงฺขารสมุทยํ ปชานาติ, สงฺขารนิโรธํ ปชานาติ, สงฺขารนิโรธคามินิํ ปฏิปทํ ปชานาติ.}

\prose{กตมญฺจ ภิกฺขเว ชรามรณํ.\csromanspacer{} ยา เตสํ เตสํ สตฺตานํ ตมฺหิ ตมฺหิ สตฺตนิกาเย ชรา ชีรณตา ขณฺฑิจฺจํ ปาลิจฺจํ วลิตฺตจตา อายุโน สํหานิ อินฺทฺริยานํ ปริปาโก.\csromanspacer{} อยํ วุจฺจติ ชรา.\csromanspacer{} ยา เตสํ เตสํ สตฺตานํ ตมฺหา ตมฺหา สตฺตนิกายา จุติ จวนตา เภโท อนฺตรธานํ มจฺจุ มรณํ กาลกิริยา ขนฺธานํ เภโท กเฬวรสฺส นิกฺเขโป, อิทํ}

\vspace{\tipitakaparskip}
\csromancenter{นิทานสํยุตฺต วุจฺจติ มรณํ.\csromanspacer{} อิติ อยญฺจ ชรา อิทญฺจ มรณํ, อิทํ วุจฺจติ ภิกฺขเว ชรามรณํ.\csromanspacer{} ชาติสมุทยา ชรามรณสมุทโย, ชาตินิโรธา ชรามรณนิโรโธ, อยเมว อริโย อฏฺฐงฺคิโก มคฺโค ชรามรณนิโรธคามินี ปฏิปทา.\csromanspacer{} เสยฺยถิทํ, สมฺมาทิฏฺฐิ ฯเปฯ สมฺมาสมาธิ.}
\prose{\csromanfolio{279}กตมา จ ภิกฺขเว ชาติ ฯเปฯ.\csromanspacer{} กตโม จ ภิกฺขเว ภโว.\csromanspacer{} กตมญฺจ ภิกฺขเว อุปาทานํ.\csromanspacer{} กตมา จ ภิกฺขเว ตณฺหา.\csromanspacer{} เวทนา.\csromanspacer{} ผสฺโส.\csromanspacer{} สฬายตนํ.\csromanspacer{} นามรูปํ.\csromanspacer{} วิญฺญาณํ.}

\prose{กตเม จ ภิกฺขเว สงฺขารา.\csromanspacer{} ตโยเม ภิกฺขเว สงฺขารา, กายสงฺขาโร วจีสงฺขาโร จิตฺตสงฺขาโร.\csromanspacer{} อิเม วุจฺจนฺติ ภิกฺขเว สงฺขารา.\csromanspacer{} อวิชฺชาสมุทยา สงฺขารสมุทโย, อวิชฺชานิโรธา สงฺขารนิโรโธ, อยเมว อริโย อฏฺฐงฺคิโก มคฺโค สงฺขารนิโรธคามินี ปฏิปทา.\csromanspacer{} เสยฺยถิทํ, สมฺมาทิฏฺฐิ ฯเปฯ สมฺมาสมาธิ.}

\prose{ยโต โข ภิกฺขเว ภิกฺขุ เอวํ ชรามรณํ ปชานาติ, เอวํ ชรามรณสมุทยํ ปชานาติ, เอวํ ชรามรณนิโรธํ ปชานาติ, เอวํ ชรามรณนิโรธคามินิํ ปฏิปทํ ปชานาติ.\csromanspacer{} เอวํ ชาติํ ปชานาติ ฯเปฯ ภวํ.\csromanspacer{} อุปาทานํ.\csromanspacer{} ตณฺหํ.\csromanspacer{} เวทนํ.\csromanspacer{} ผสฺสํ.\csromanspacer{} สฬายตนํ.\csromanspacer{} นามรูปํ.\csromanspacer{} วิญฺญาณํ.\csromanspacer{} สงฺขาเร.\csromanspacer{} สงฺขารสมุทยํ.\csromanspacer{} สงฺขารนิโรธํ.\csromanspacer{} เอวํ สงฺขารนิโรธคามินิํ ปฏิปทํ ปชานาติ.\csromanspacer{} อยํ วุจฺจติ ภิกฺขเว ภิกฺขุ ทิฏฺฐิสมฺปนฺโน อิติปิ, ทสฺสนสมฺปนฺโน อิติปิ, อาคโต อิมํ สทฺธมฺมํ อิติปิ, ปสฺสติ อิมํ สทฺธมฺมํ อิติปิ, เสกฺเขน ญาเณน สมนฺนาคโต อิติปิ, เสกฺขาย วิชฺชาย สมนฺนาคโต อิติปิ, ธมฺมโสตํ สมาปนฺโน อิติปิ, อริโย นิพฺเพธิกปญฺโญ อิติปิ, อมตทฺวารํ อาหจฺจ ติฏฺฐติ อิติปีติ.\csromanspacer{}\csromanspacer{}อฏฺฐมํ.}

\csromanmatikaanchor{mtk.31}
\csromantocmarkref{subsection}{5. ภูมิชสุตฺต}{mtk.31}
\bookmark[dest={mtk.31},level=2]{5. ภูมิชสุตฺต}
\csromanheader{2}{9. สมณพฺราหฺมณสุตฺต}
\proseitem{29}{สาวตฺถิยํ วิหรติ.\csromanspacer{} ตตฺร โข ฯเปฯ.\csromanspacer{} เย หิ เกจิ ภิกฺขเว สมณา วา พฺราหฺมณา วา ชรามรณํ น ปริชานนฺติ, ชรามรณสมุทยํ น ปริชานนฺติ, ชรามรณนิโรธํ น ปริชานนฺติ, ชรามรณนิโรโคมินิํ ปฏิปทํ น ปริชานนฺติ.\csromanspacer{} ชาติํ น ปริชานนฺติ ฯเปฯ.\csromanspacer{} ภวํ.\csromanspacer{} อุปาทานํ.\csromanspacer{} ตณฺหํ.\csromanspacer{} เวทนํ.\csromanspacer{} ผสฺสํ.\csromanspacer{} สฬายตนํ.\csromanspacer{} นามรูปํ.\csromanspacer{} วิญฺญาณํ.\csromanspacer{} สงฺขาเร.\csromanspacer{} สงฺขารสมุทยํ.\csromanspacer{} สงฺขารนิโรธํ.\csromanspacer{} สงฺขารนิโรธคามินิํ ปฏิปทํ น ปริชานนฺติ.\csromanspacer{} น เมเต ภิกฺขเว สมณา วา พฺราหฺมณา วา สมเณสุ วา สมณสมฺมตา, พฺราหฺมเณสุ วา พฺราหฺมณสมฺมตา.\csromanspacer{} น จ}

\vspace{\tipitakaparskip}
\csromancenter{นิทานวคฺคสํยุตฺตปาฬิ ปเนเต อายสฺมนฺโต สามญฺญตฺถํ วา พฺรหฺมญฺญตฺถํ วา ทิฏฺเฐว ธมฺเม สยํ อภิญฺญา สจฺฉิกตฺวา อุปสมฺปชฺช วิหรนฺติ.}
\prose{\csromanfolio{280}เย จ โข เกจิ ภิกฺขเว สมณา วา พฺราหฺมณา วา ชรามรณํ ปริชานนฺติ, ชรามรณสมุทยํ ปริชานนฺติ, ชรามรณนิโรธํ ปริชานนฺติ, ชรามรณนิโรธคามินิํ ปฏิปทํ ปริชานนฺติ.\csromanspacer{} ชาติํ ปริชานนฺติ ฯเปฯ.\csromanspacer{} ภวํ.\csromanspacer{} อุปาทานํ.\csromanspacer{} ตณฺหํ.\csromanspacer{} เวทนํ.\csromanspacer{} ผสฺสํ.\csromanspacer{} สฬายตนํ.\csromanspacer{} นามรูปํ.\csromanspacer{} วิญฺญาณํ.\csromanspacer{} สงฺขาเร ปริชานนฺติ, สงฺขาร- สมุทยํ ปริชานนฺติ, สงฺขารนิโรธํ ปริชานนฺติ, สงฺขารนิโรธ- คามินิํ ปฏิปทํ ปริชานนฺติ.\csromanspacer{} เต โขเม ภิกฺขเว สมณา วา พฺราหฺมณา วา สมเณสุ เจว สมณสมฺมตา, พฺราหฺมเณสุ จ พฺราหฺมณสมฺมตา.\csromanspacer{} เต จ ปนายสฺมนฺโต สามญฺญตฺถญฺจ พฺรหฺมญฺญตฺถญฺจ ทิฏฺเฐว ธมฺเม สยํ อภิญฺญา สจฺฉิกตฺวา อุปสมฺปชฺช วิหรนฺตีติ.\csromanspacer{}\csromanspacer{}นวมํ.}

\csromanheader{2}{10. ทุติยสมณพฺราหฺมณสุตฺต}
\proseitem{30}{สาวตฺถิยํ วิหรติ.\csromanspacer{} ตตฺร โข ฯเปฯ.\csromanspacer{} เย หิ เกจิ ภิกฺขเว สมณา วา พฺราหฺมณา วา ชรามรณํ นปฺปชานนฺติ, ชรามรณสมุทยํ นปฺปชานนฺติ, ชรามรณนิโรธํ นปฺปชานนฺติ, ชรามรณนิโรธคามินิํ ปฏิปทํ นปฺปชานนฺติ.\csromanspacer{} เต วต ชรามรณํ สมติกฺกมฺม ฐสฺสนฺตีติ เนตํ ฐานํ วิชฺชติ.\csromanspacer{} ชาติํ นปฺปชานนฺติ ฯเปฯ.\csromanspacer{} ภวํ.\csromanspacer{} อุปาทานํ.\csromanspacer{} ตณฺหํ.\csromanspacer{} เวทนํ.\csromanspacer{} ผสฺสํ.\csromanspacer{} สฬายตนํ.\csromanspacer{} นามรูปํ.\csromanspacer{} วิญฺญาณํ.\csromanspacer{} สงฺขาเร นปฺปชานนฺติ, สงฺขารสมุทยํ นปฺปชานนฺติ, สงฺขารนิโรธํ นปฺปชานนฺติ, สงฺขารนิโรธคามินิํ ปฏิปทํ นปฺปชานนฺติ.\csromanspacer{} เต วต สงฺขาเร สมติกฺกมฺม ฐสฺสนฺตีติ เนตํ ฐานํ วิชฺชติ.}

\prose{เย จ โข เกจิ ภิกฺขเว สมณา วา พฺราหฺมณา วา ชรามรณํ ปชานนฺติ, ชรามรณสมุทยํ ปชานนฺติ, ชรามรณนิโรธํ ปชานนฺติ, ชรามรณนิโรธคามินิํ ปฏิปทํ ปชานนฺติ.\csromanspacer{} เต วต ชรามรณํ สมติกฺกมฺม ฐสฺสนฺตีติ ฐานเมตํ วิชฺชติ.\csromanspacer{} ชาติํ ปชานนฺติ ฯเปฯ.\csromanspacer{} ภวํ.\csromanspacer{} อุปาทานํ.\csromanspacer{} ตณฺหํ.\csromanspacer{} เวทนํ.\csromanspacer{} ผสฺสํ.\csromanspacer{} สฬายตนํ.\csromanspacer{} นามรูปํ.\csromanspacer{} วิญฺญาณํ.\csromanspacer{} สงฺขาเร ปชานนฺติ, สงฺขารสมุทยํ ปชานนฺติ, สงฺขาร- นิโรธํ ปชานนฺติ, สงฺขารนิโรธคามินิํ ปฏิปทํ ปชานนฺติ.\csromanspacer{} เต วต สงฺขาเร สมติกฺกมฺม ฐสฺสนฺตีติ ฐานเมตํ วิชฺชตีติ.\csromanspacer{}\csromanspacer{}ทสมํ.}

\vspace{\tipitakaparskip}
\csromancenter{ทสพลวคฺโค ตติโย.}
\titlehead{1. นิทานสํยุตฺต \textbf{ตสฺสุทฺทานํ}}
\csromanmatikaanchor{mtk.32}
\csromanmatikaanchor{mtk.37}
\csromantocmarkref{subsection}{6. อุปวาณสุตฺต}{mtk.32}
\bookmark[dest={mtk.32},level=2]{6. อุปวาณสุตฺต}
\csromangathameasure{เทฺว ทสพลา อุปนิสา จ, อญฺญติตฺถิยภูมิโช. \\ อุปวาโณ ปจฺจโย ภิกฺขุ, เทฺว จ สมณพฺราหฺมณาติ.}
\csromangathagroup{\csromangathastanza{\csromanfolio{281}เทฺว ทสพลา อุปนิสา จ, อญฺญติตฺถิยภูมิโช. \\ อุปวาโณ ปจฺจโย ภิกฺขุ, เทฺว จ สมณพฺราหฺมณาติ.}}

\csromanheader{1}{4. กฬารขตฺติยวคฺค}
\csromanheader{2}{1. ภูตสุตฺต}
\proseitem{31}{เอกํ สมยํ ภควา สาวตฺถิยํ วิหรติ.\csromanspacer{} ตตฺร โข ภควา อายสฺมนฺตํ สาริปุตฺตํ อามนฺเตสิ “วุตฺตมิทํ สาริปุตฺต ปารายเน\footnote{ปารายเณ (สี)} อชิตปเญฺห\csromandash{}}

\csromangathameasure{เย จ สงฺขาตธมฺมาเส, เย จ เสกฺขา ปุถู อิธ. \\ เตสํ เม นิปโก อิริยํ, ปุฏฺโฐ ปพฺรูหิ มาริสาติ.}
\csromangathagroup{\csromangathastanza{เย จ สงฺขาตธมฺมาเส, เย จ เสกฺขา ปุถู อิธ. \\ เตสํ เม นิปโก อิริยํ, ปุฏฺโฐ ปพฺรูหิ มาริสาติ.}}

\prose{อิมสฺส นุ โข สาริปุตฺต สํขิตฺเตน ภาสิตสฺส กถํ วิตฺถาเรน อตฺโถ ทฏฺฐพฺโพ”ติ.\csromanspacer{} เอวํ วุตฺเต อายสฺมา สาริปุตฺโต ตุณฺหี อโหสิ.\csromanspacer{} ทุติยมฺปิ โข ภควา อายสฺมนฺตํ สาริปุตฺตํ อามนฺเตสิ ฯเปฯ.\csromanspacer{} ทุติยมฺปิ โข อายสฺมา สาริปุตฺโต ตุณฺหี อโหสิ.\csromanspacer{} ตติยมฺปิ โข ภควา อายสฺมนฺตํ สาริปุตฺตํ อามนฺเตสิ “วุตฺตมิทํ สาริปุตฺต ปารายเน อชิตปเญฺห\csromandash{}}

\csromangathameasure{เย จ สงฺขาตธมฺมาเส, เย จ เสกฺขา ปุถู อิธ. \\ เตสํ เม นิปโก อิริยํ, ปุฏฺโฐ ปพฺรูหิ มาริสาติ.}
\csromangathagroup{\csromangathastanza{เย จ สงฺขาตธมฺมาเส, เย จ เสกฺขา ปุถู อิธ. \\ เตสํ เม นิปโก อิริยํ, ปุฏฺโฐ ปพฺรูหิ มาริสาติ.}}

\prose{อิมสฺส นุ โข สาริปุตฺต สํขิตฺเตน ภาสิตสฺส กถํ วิตฺถาเรน อตฺโถ ทฏฺฐพฺโพ”ติ.\csromanspacer{} ตติยมฺปิ โข อายสฺมา สาริปุตฺโต ตุณฺหี อโหสิ.}

\prose{“ภูตมิทนฺ”ติ สาริปุตฺต ปสฺสสีติ. “ภูตมิทนฺ”ติ ภนฺเต ยถาภูตํ สมฺมปฺปญฺญาย ปสฺสติ, “ภูตมิทนฺ”ติ ยถาภูตํ สมฺมปฺปญฺญาย ทิสฺวา ภูตสฺส นิพฺพิทาย วิราคาย นิโรธาย ปฏิปนฺโน โหติ. “ตทาหารสมฺภวนฺ”ติ ยถาภูตํ สมฺมปฺปญฺญาย ปสฺสติ, “ตทาหารสมฺภวนฺ”ติ ยถาภูตํ สมฺมปฺปญฺญาย ทิสฺวา อาหารสมฺภวสฺส นิพฺพิทาย วิราคาย นิโรธาย}

\vspace{\tipitakaparskip}
\csromancenter{นิทานวคฺคสํยุตฺตปาฬิ ปฏิปนฺโน โหติ. “ตทาหารนิโรธา ยํ ภูตํ ตํ นิโรธธมฺมนฺ”ติ ยถาภูตํ สมฺมปฺปญฺญาย ปสฺสติ, “ตทาหารนิโรธา ยํ ภูตํ ตํ นิโรธธมฺมนฺ”ติ ยถาภูตํ สมฺมปฺปญฺญาย ทิสฺวา นิโรธธมฺมสฺส นิพฺพิทาย วิราคาย นิโรธาย ปฏิปนฺโน โหติ.\csromanspacer{} เอวํ โข ภนฺเต เสกฺโข โหติ.}
\prose{\csromanfolio{282}กถญฺจ ภนฺเต สงฺขาตธมฺโม โหติ. “ภูตมิทนฺ”ติ ภนฺเต ยถาภูตํ สมฺมปฺปญฺญาย ปสฺสติ, “ภูตมิทนฺ”ติ ยถาภูตํ สมฺมปฺปญฺญาย ทิสฺวา ภูตสฺส นิพฺพิทา วิราคา นิโรธา อนุปาทา วิมุตฺโต โหติ. “ตทาหารสมฺภวนฺ”ติ ยถาภูตํ สมฺมปฺปญฺญาย ปสฺสติ, “ตทาหารสมฺภวนฺ”ติ ยถาภูตํ สมฺมปฺปญฺญาย ทิสฺวา อาหารสมฺภวสฺส นิพฺพิทา วิราคา นิโรธา อนุปาทา วิมุตฺโต โหติ. “ตทาหารนิโรธา ยํ ภูตํ ตํ นิโรธธมฺมนฺ”ติ ยถาภูตํ สมฺมปฺปญฺญาย ปสฺสติ, “ตทาหารนิโรธา ยํ ภูตํ ตํ นิโรธธมฺมนฺ”ติ ยถาภูตํ สมฺมปฺปญฺญาย ทิสฺวา นิโรธธมฺมสฺส นิพฺพิทา วิราคา นิโรธา อนุปาทา วิมุตฺโต โหติ.\csromanspacer{} เอวํ โข ภนฺเต สงฺขาตธมฺโม โหติ.\csromanspacer{} อิติ โข ภนฺเต ยํ ตํ วุตฺตํ ปารายเน อชิตปเญฺห\csromandash{}}

\csromangathameasure{เย จ สงฺขาตธมฺมาเส, เย จ เสกฺขา ปุถู อิธ. \\ เตสํ เม นิปโก อิริยํ, ปุฏฺโฐ ปพฺรูหิ มาริสาติ.}
\csromangathagroup{\csromangathastanza{เย จ สงฺขาตธมฺมาเส, เย จ เสกฺขา ปุถู อิธ. \\ เตสํ เม นิปโก อิริยํ, ปุฏฺโฐ ปพฺรูหิ มาริสาติ.}}

\prose{อิมสฺส ขฺวาหํ ภนฺเต สํขิตฺเตน ภาสิตสฺส เอวํ วิตฺถาเรน อตฺถํ อาชานามีติ.}

\prose{สาธุ สาธุ สาริปุตฺต, “ภูตมิทนฺ”ติ สาริปุตฺต ยถาภูตํ สมฺมปฺปญฺญาย ปสฺสติ, “ภูตมิทนฺ”ติ ยถาภูตํ สมฺมปฺปญฺญาย ทิสฺวา ภูตสฺส นิพฺพิทาย วิราคาย นิโรธาย ปฏิปนฺโน โหติ. “ตทาหารสมฺภวนฺ”ติ ยถาภูตํ สมฺมปฺปญฺญาย ปสฺสติ, “ตทาหารสมฺภวนฺ”ติ ยถาภูตํ สมฺมปฺปญฺญาย ทิสฺวา อาหารสมฺภวสฺส นิพฺพิทาย วิราคาย นิโรธาย ปฏิปนฺโน โหติ. “ตทาหารนิโรธา ยํ ภูตํ ตํ นิโรธธมฺมนฺ”ติ ยถาภูตํ สมฺมปฺปญฺญาย ปสฺสติ, “ตทาหารนิโรธา ยํ ภูตํ ตํ นิโรธธมฺมนฺ”ติ ยถาภูตํ สมฺมปฺปญฺญาย ทิสฺวา นิโรธธมฺมสฺส นิพฺพิทาย วิราคาย นิโรธาย ปฏิปนฺโน โหติ.\csromanspacer{} เอวํ โข สาริปุตฺต เสกฺโข โหติ.}

\prose{กถญฺจ สาริปุตฺต สงฺขาตธมฺโม โหติ. “ภูตมิทนฺ”ติ สาริปุตฺต ยถาภูตํ สมฺมปฺปญฺญาย ปสฺสติ, “ภูตมิทนฺ”ติ ยถาภูตํ สมฺมปฺปญฺญาย}

\csromanmatikaanchor{mtk.33}
\csromantocmarkref{subsection}{7. ปจฺจยสุตฺต}{mtk.33}
\bookmark[dest={mtk.33},level=2]{7. ปจฺจยสุตฺต}
\vspace{\tipitakaparskip}
\csromancenter{นิทานสํยุตฺต ทิสฺวา ภูตสฺส นิพฺพิทา วิราคา นิโรธา อนุปาทา วิมุตฺโต โหติ. “ตทาหารสมฺภวนฺ”ติ ยถาภูตํ สมฺมปฺปญฺญาย ปสฺสติ, “ตทาหารสมฺภวนฺ”ติ ยถาภูตํ สมฺมปฺปญฺญาย ทิสฺวา อาหารสมฺภวสฺส นิพฺพิทา วิราคา นิโรธา อนุปาทา วิมุตฺโต โหติ. “ตทาหารนิโรธา ยํ ภูตํ ตํ นิโรธธมฺมนฺ”ติ ยถาภูตํ สมฺมปฺปญฺญาย ปสฺสติ, “ตทาหารนิโรธา ยํ ภูตํ ตํ นิโรธธมฺมนฺ”ติ ยถาภูตํ สมฺมปฺปญฺญาย ทิสฺวา นิโรธธมฺมสฺส นิพฺพิทา วิราคา นิโรธา อนุปาทา วิมุตฺโต โหติ.\csromanspacer{} เอวํ โข สาริปุตฺต สงฺขาตธมฺโม โหติ.\csromanspacer{} อิติ โข สาริปุตฺต ยํ ตํ วุตฺตํ ปารายเน อชิตปเญฺห\csromandash{}}
\vspace{0.5\baselineskip}
\csromangathameasure{เย จ สงฺขาตธมฺมาเส, เย จ เสกฺขา ปุถู อิธ. \\ เตสํ เม นิปโก อิริยํ, ปุฏฺโฐ ปพฺรูหิ มาริสาติ.}
\csromangathagroup{\csromangathastanza{\csromanfolio{283}เย จ สงฺขาตธมฺมาเส, เย จ เสกฺขา ปุถู อิธ. \\ เตสํ เม นิปโก อิริยํ, ปุฏฺโฐ ปพฺรูหิ มาริสาติ.}}

\prose{อิมสฺส โข สาริปุตฺต สํขิตฺเตน ภาสิตสฺส เอวํ วิตฺถาเรน อตฺโถ ทฏฺฐพฺโพติ.\csromanspacer{}\csromanspacer{}ปฐมํ.}

\csromanheader{2}{2. กฬารสุตฺต}
\proseitem{32}{สาวตฺถิยํ วิหรติ.\csromanspacer{} อถ โข กฬารขตฺติโย ภิกฺขุ เยนายสฺมา สาริปุตฺโต เตนุปสงฺกมิ, อุปสงฺกมิตฺวา อายสฺมตา สาริปุตฺเตน สทฺธิํ สมฺโมทิ, สมฺโมทนียํ กถํ สารณียํ วีติสาเรตฺวา เอกมนฺตํ นิสีทิ, เอกมนฺตํ นิสินฺโน โข กฬารขตฺติโย ภิกฺขุ อายสฺมนฺตํ สาริปุตฺตํ เอตทโวจ “โมฬิยผคฺคุโน อาวุโส สาริปุตฺต ภิกฺขุ สิกฺขํ ปจฺจกฺขาย หีนายาวตฺโต”ติ.\csromanspacer{} น หิ นูน โส อายสฺมา อิมสฺมิํ ธมฺมวินเย อสฺสาสมลตฺถาติ.\csromanspacer{} เตน หายสฺมา สาริปุตฺโต อิมสฺมิํ ธมฺมวินเย อสฺสาสํ ปตฺโตติ.\csromanspacer{} น ขฺวาหํ อาวุโส กงฺขามีติ.\csromanspacer{} อายติํ ปนาวุโสติ.\csromanspacer{} น ขฺวาหํ อาวุโส วิจิกิจฺฉามีติ.\csromanspacer{} อถ โข กฬารขตฺติโย ภิกฺขุ อุฏฺฐายาสนา เยน ภควา เตนุปสงฺกมิ, อุปสงฺกมิตฺวา ภควนฺตํ อภิวาเทตฺวา เอกมนฺตํ นิสีทิ, เอกมนฺตํ นิสินฺโน โข กฬารขตฺติโย ภิกฺขุ ภควนฺตํ เอตทโวจ “อายสฺมตา ภนฺเต สาริปุตฺเตน อญฺญา พฺยากตา ‘ขีณา ชาติ, วุสิตํ พฺรหฺมจริยํ, กตํ กรณียํ, นาปรํ อิตฺถตฺตายาติ ปชานามี’ติ”.}

\prose{อถ โข ภควา อญฺญตรํ ภิกฺขุํ อามนฺเตสิ “เอหิ ตฺวํ ภิกฺขุ มม วจเนน สาริปุตฺตํ อามนฺเตหิ ‘สตฺถา ตํ อาวุโส สาริปุตฺต}

\vspace{\tipitakaparskip}
\csromancenter{นิทานวคฺคสํยุตฺตปาฬิ อามนฺเตตี’ติ”. “เอวํ ภนฺเต”ติ โข โส ภิกฺขุ ภควโต ปฏิสฺสุตฺวา เยนายสฺมา สาริปุตฺโต เตนุปสงฺกมิ, อุปสงฺกมิตฺวา อายสฺมนฺตํ สาริปุตฺตํ เอตทโวจ “สตฺถา ตํ อาวุโส สาริปุตฺต อามนฺเตตี”ติ. “เอวํ อาวุโส”ติ โข อายสฺมา สาริปุตฺโต ตสฺส ภิกฺขุโน ปฏิสฺสุตฺวา เยน ภควา เตนุปสงฺกมิ, อุปสงฺกมิตฺวา ภควนฺตํ อภิวาเทตฺวา เอกมนฺตํ นิสีทิ, เอกมนฺตํ นิสินฺนํ โข อายสฺมนฺตํ สาริปุตฺตํ ภควา เอตทโวจ “สจฺจํ กิร ตยา สาริปุตฺต อญฺญา พฺยากตา ‘ขีณา ชาติ, วุสิตํ พฺรหฺมจริยํ, กตํ กรณียํ, นาปรํ อิตฺถตฺตายาติ ปชานามี’ติ”.\csromanspacer{} น โข ภนฺเต เอเตหิ ปเทหิ เอเตหิ พฺยญฺชเนหิ อตฺโถ\footnote{อตฺโถ จ (สฺยา, กํ, ก)} วุตฺโตติ.\csromanspacer{} เยน เกนจิปิ สาริปุตฺต ปริยาเยน กุลปุตฺโต อญฺญํ พฺยากโรติ, อถ โข พฺยากตํ พฺยากตโต ทฏฺฐพฺพนฺติ.\csromanspacer{} นนุ อหมฺปิ ภนฺเต เอวํ วทามิ “น โข ภนฺเต เอเตหิ ปเทหิ เอเตหิ พฺยญฺชเนหิ อตฺโถ วุตฺโต”ติ.}
\csromanmatikaanchor{mtk.34}
\csromantocmarkref{subsection}{8. ภิกฺขุสุตฺต}{mtk.34}
\bookmark[dest={mtk.34},level=2]{8. ภิกฺขุสุตฺต}
\prose{\csromanfolio{284}สเจ ตํ สาริปุตฺต เอวํ ปุจฺเฉยฺยุํ “กถํ ชานตา ปน ตยา อาวุโส สาริปุตฺต กถํ ปสฺสตา อญฺญา พฺยากตา ‘ขีณา ชาติ, วุสิตํ พฺรหฺมจริยํ, กตํ กรณียํ, นาปรํ อิตฺถตฺตายาติ ปชานามี’ติ”, เอวํ ปุฏฺโฐ ตฺวํ สาริปุตฺต กินฺติ พฺยากเรยฺยาสีติ.\csromanspacer{} สเจ มํ ภนฺเต เอวํ ปุจฺเฉยฺยุํ “กถํ ชานตา ปน ตยา อาวุโส สาริปุตฺต กถํ ปสฺสตา อญฺญา พฺยากตา ‘ขีณา ชาติ, วุสิตํ พฺรหฺมจริยํ, กตํ กรณียํ, นาปรํ อิตฺถตฺตายาติ ปชานามี’ติ”, เอวํ ปุฏฺโฐหํ\footnote{ปุฏฺโฐ อหํ (สฺยา, กํ), ปุฏฺฐาหํ (อิ, ก)} ภนฺเต เอวํ พฺยากเรยฺยํ “ยํนิทานา อาวุโส ชาติ, ตสฺส นิทานสฺส ขยา ขีณสฺมิํ ‘ขีณามฺหี’ติ วิทิตํ. ‘ขีณามฺหี’ติ วิทิตฺวา ‘ขีณา ชาติ, วุสิตํ พฺรหฺมจริยํ, กตํ กรณียํ, นาปรํ อิตฺถตฺตายา’ติ ปชานามี”ติ.\csromanspacer{} เอวํ ปุฏฺโฐหํ ภนฺเต เอวํ พฺยากเรยฺยนฺติ.}

\prose{สเจ ปน ตํ สาริปุตฺต เอวํ ปุจฺเฉยฺยุํ “ชาติ ปนาวุโส สาริปุตฺต กิํนิทานา กิํสมุทยา กิํชาติกา กิํปภวา”ติ, เอวํ ปุฏฺโฐ ตํ สาริปุตฺต กินฺติ พฺยากเรยฺยาสีติ.\csromanspacer{} สเจ มํ ภนฺเต เอวํ ปุจฺเฉยฺยุํ “ชาติ ปนาวุโส สาริปุตฺต กิํนิทานา กิํสมุทยา กิํชาติกา กิํปภวา”ติ, เอวํ ปุฏฺโฐหํ ภนฺเต เอวํ พฺยากเรยฺยํ “ชาติ โข อาวุโส ภวนิทานา}

\vspace{\tipitakaparskip}
\csromancenter{นิทานสํยุตฺต ภวสมุทยา ภวชาติกา ภวปฺปภวา”ติ.\csromanspacer{} เอวํ ปุฏฺโฐหํ ภนฺเต เอวํ พฺยากเรยฺยนฺติ.}
\prose{\csromanfolio{285}สเจ ปน ตํ สาริปุตฺต เอวํ ปุจฺเฉยฺยุํ “ภโว ปนาวุโส สาริปุตฺต กิํนิทาโน กิํสมุทโย กิํชาติโก กิํปภโว”ติ, เอวํ ปุฏฺโฐ ตฺวํ สาริปุตฺต กินฺติ พฺยากเรยฺยาสีติ.\csromanspacer{} สเจ มํ ภนฺเต เอวํ ปุจฺเฉยฺยุํ “ภโว ปนาวุโส สาริปุตฺต กิํนิทาโน กิํสมุทโย กิํชาติโก กิํปภโว”ติ, เอวํ ปุฏฺโฐหํ ภนฺเต เอวํ พฺยากเรยฺยํ “ภโว โข อาวุโส อุปาทานนิทาโน อุปาทานสมุทโย อุปาทานชาติโก อุปาทานปฺปภโว”ติ.\csromanspacer{} เอวํ ปุฏฺโฐหํ ภนฺเต เอวํ พฺยากเรยฺยนฺติ.}

\prose{สเจ ปน ตํ สาริปุตฺต เอวํ ปุจฺเฉยฺยุํ “อุปาทานํ ปนาวุโส ฯเปฯ.\csromanspacer{} สเจ ปน ตํ สาริปุตฺต เอวํ ปุจฺเฉยฺยุํ “ตณฺหา ปนาวุโส สาริปุตฺต กิํนิทานา กิํสมุทยา กิํชาติกา กิํปภวา”ติ, เอวํ ปุฏฺโฐ ตฺวํ สาริปุตฺต กินฺติ พฺยากเรยฺยาสีติ.\csromanspacer{} สเจ มํ ภนฺเต เอวํ ปุจฺเฉยฺยุํ “ตณฺหา ปนาวุโส สาริปุตฺต กิํนิทานา กิํสมุทยา กิํชาติกา กิํปภวา”ติ, เอวํ ปุฏฺโฐหํ ภนฺเต เอวํ พฺยากเรยฺยํ “ตณฺหา โข อาวุโส เวทนานิทานา เวทนาสมุทยา เวทนาชาติกา เวทนาปภวา”ติ.\csromanspacer{} เอวํ ปุฏฺโฐหํ ภนฺเต เอวํ พฺยากเรยฺยนฺติ.}

\prose{สเจ ปน ตํ สาริปุตฺต เอวํ ปุจฺเฉยฺยุํ “กถํ ชานโต ปน เต อาวุโส สาริปุตฺต กถํ ปสฺสโต ยา เวทนาสุ นนฺที, สา น อุปฏฺฐาสี”ติ, เอวํ ปุฏฺโฐ ตฺวํ สาริปุตฺต กินฺติ พฺยากเรยฺยาสีติ.\csromanspacer{} สเจ มํ ภนฺเต เอวํ ปุจฺเฉยฺยุํ “กถํ ชานโต ปน เต อาวุโส สาริปุตฺต กถํ ปสฺสโต ยา เวทนาสุ นนฺที, สา น อุปฏฺฐาสี”ติ, เอวํ ปุฏฺโฐหํ ภนฺเต เอวํ พฺยากเรยฺยํ “ติสฺโส โข อิมา อาวุโส เวทนา.\csromanspacer{} กตมา ติสฺโส, สุขา เวทนา ทุกฺขา เวทนา อทุกฺขมสุขา เวทนา.\csromanspacer{} อิมา โข อาวุโส ติสฺโส เวทนา อนิจฺจา, ยทนิจฺจํ ตํ ทุกฺขนฺติ วิทิตํ\footnote{วิทิตา (ฏีกา)}, ยา เวทนาสุ นนฺที, สา น อุปฏฺฐาสี”ติ.\csromanspacer{} เอวํ ปุฏฺโฐหํ ภนฺเต เอวํ พฺยากเรยฺยนฺติ.}

\prose{สาธุ สาธุ สาริปุตฺต, อยมฺปิ โข สาริปุตฺต ปริยาโย, เอตสฺเสว อตฺถสฺส สํขิตฺเตน เวยฺยากรณาย “ยํ กิญฺจิ เวทยิตํ, ตํ ทุกฺขสฺมินฺ”ติ.}

\csromanmatikaanchor{mtk.35}
\csromantocmarkref{subsection}{9. สมณพฺราหฺมณสุตฺต}{mtk.35}
\bookmark[dest={mtk.35},level=2]{9. สมณพฺราหฺมณสุตฺต}
\gambhira{นิทานวคฺคสํยุตฺตปาฬิ}
\prose{\csromanfolio{286}สเจ ปน ตํ สาริปุตฺต เอวํ ปุจฺเฉยฺยุํ “กถํ วิโมกฺขา ปน ตยา อาวุโส สาริปุตฺต อญฺญา พฺยากตา ‘ขีณา ชาติ, วุสิตํ พฺรหฺมจริยํ, กตํ กรณียํ, นาปรํ อิตฺถตฺตายาติ ปชานามี’ติ”, เอวํ ปุฏฺโฐ ตฺวํ สาริปุตฺต กินฺติ พฺยากเรยฺยาสีติ.\csromanspacer{} สเจ มํ ภนฺเต เอวํ ปุจฺเฉยฺยุํ “กถํ วิโมกฺขา ปน ตยา อาวุโส สาริปุตฺต อญฺญา พฺยากตา ‘ขีณา ชาติ, วุสิตํ พฺรหฺมจริยํ, กตํ กรณียํ.\csromanspacer{} นาปรํ อิตฺถตฺตายาติ ปชานามี’ติ”, เอวํ ปุฏฺโฐหํ ภนฺเต เอวํ พฺยากเรยฺยํ “อชฺฌตฺตํ วิโมกฺขา ขฺวาหํ อาวุโส สพฺพุปาทานกฺขยา ตถา สโต วิหรามิ, ยถา สตํ วิหรนฺตํ อาสวา นานุสฺสวนฺติ, อตฺตานญฺจ นาวชานามี”ติ.\csromanspacer{} เอวํ ปุฏฺโฐหํ ภนฺเต เอวํ พฺยากเรยฺยนฺติ.}

\prose{สาธุ สาธุ สาริปุตฺต, อยมฺปิ โข สาริปุตฺต ปริยาโย, เอตสฺเสว อตฺถสฺส สํขิตฺเตน เวยฺยากรณาย “เย อาสวา สมเณน วุตฺตา, เตสฺวาหํ น กงฺขามิ, เต เม ปหีนาติ น วิจิกิจฺฉามี”ติ.\csromanspacer{} อิทมโวจ ภควา, อิทํ วตฺวา สุคโต อุฏฺฐายาสนา วิหารํ ปาวิสิ.}

\prose{ตตฺร โข อายสฺมา สาริปุตฺโต อจิรปกฺกนฺตสฺส ภควโต ภิกฺขู อามนฺเตสิ “ปุพฺเพ อปฺปฏิสํวิทิตํ มํ อาวุโส ภควา ปฐมํ ปญฺหํ อปุจฺฉิ, ตสฺส เม อโหสิ ทนฺธายิตตฺตํ.\csromanspacer{} ยโต จ โข เม อาวุโส ภควา ปฐมํ ปญฺหํ อนุโมทิ, ตสฺส มยฺหํ อาวุโส เอตทโหสิ ‘ทิวสํ เจปิ มํ ภควา เอตมตฺถํ ปุจฺเฉยฺย อญฺญมญฺเญหิ ปเทหิ อญฺญมญฺเญหิ ปริยาเยหิ, ทิวสมฺปาหํ ภควโต เอตมตฺถํ พฺยากเรยฺยํ อญฺญมญฺเญหิ ปเทหิ อญฺญมญฺเญหิ ปริยาเยหิ.\csromanspacer{} รตฺติํ เจปิ มํ ภควา เอตมตฺถํ ปุจฺเฉยฺย อญฺญมญฺเญหิ ปเทหิ อญฺญมญฺเญหิ ปริยาเยหิ, รตฺติมฺปาหํ ภควโต เอตมตฺถํ พฺยากเรยฺยํ อญฺญมญฺเญหิ ปเทหิ อญฺญมญฺเญหิ ปริยาเยหิ.\csromanspacer{} รตฺตินฺทิวํ\footnote{รตฺติทิวํ (ก)} เจปิ มํ ภควา เอตมตฺถํ ปุจฺเฉยฺย อญฺญมญฺเญหิ ปเทหิ อญฺญมญฺเญหิ ปริยาเยหิ, รตฺตินฺทิวมฺปาหํ ภควโต เอตมตฺถํ พฺยากเรยฺยํ อญฺญมญฺเญหิ ปเทหิ อญฺญมญฺเญหิ ปริยาเยหิ.\csromanspacer{} เทฺว รตฺตินฺทิวานิ เจปิ มํ ภควา เอตมตฺถํ ปุจฺเฉยฺย ฯเปฯ.\csromanspacer{} เทฺว รตฺตินฺทิวานิปาหํ ภควโต เอตมตฺถํ พฺยากเรยฺยํ ฯเปฯ.\csromanspacer{} ตีณิ รตฺตินฺทิวานิ เจปิ มํ ภควา เอตมตฺถํ ปุจฺเฉยฺย ฯเปฯ.\csromanspacer{} ตีณิ รตฺตินฺทิวานิปาหํ ภควโต เอตมตฺถํ พฺยากเรยฺยํ ฯเปฯ.\csromanspacer{} จตฺตาริ รตฺตินฺทิวานิ เจปิ มํ ภควา เอตมตฺถํ ปุจฺเฉยฺย ฯเปฯ.\csromanspacer{} จตฺตาริ}

\csromanmatikaanchor{mtk.36}
\csromantocmarkref{subsection}{10. ทุติยสมณพฺราหฺมณสุตฺต}{mtk.36}
\bookmark[dest={mtk.36},level=2]{10. ทุติยสมณพฺราหฺมณสุตฺต}
\csromantocmarkref{section}{อุทฺทานคาถา}{mtk.37}
\bookmark[dest={mtk.37},level=1]{อุทฺทานคาถา}
\vspace{\tipitakaparskip}
\csromancenter{นิทานสํยุตฺต รตฺตินฺทิวานิปาหํ ภควโต เอตมตฺถํ พฺยากเรยฺยํ ฯเปฯ.\csromanspacer{} ปญฺจ รตฺตินฺทิวานิ เจปิ มํ ภควา เอตมตฺถํ ปุจฺเฉยฺย ฯเปฯ.\csromanspacer{} ปญฺจ รตฺตินฺทิวานิปาหํ ภควโต เอตมตฺถํ พฺยากเรยฺยํ ฯเปฯ.\csromanspacer{} ฉ รตฺตินฺทิวานิ เจปิ มํ ภควา เอตมตฺถํ ปุจฺเฉยฺย ฯเปฯ.\csromanspacer{} ฉ รตฺตินฺทิวานิปาหํ ภควโต เอตมตฺถํ พฺยากเรยฺยํ ฯเปฯ.\csromanspacer{} สตฺต รตฺตินฺทิวานิ เจปิ มํ ภควา เอตมตฺถํ ปุจฺเฉยฺย อญฺญมญฺเญหิ ปเทหิ อญฺญมญฺเญหิ ปริยาเยหิ, สตฺต รตฺตินฺทิวานิปาหํ ภควโต เอตมตฺถํ พฺยากเรยฺยํ อญฺญมญฺเญหิ ปเทหิ อญฺญมญฺเญหิ ปริยาเยหี’ติ”.}
\prose{\csromanfolio{287}อถ โข กฬารขตฺติโย ภิกฺขุ อุฏฺฐายาสนา เยน ภควา เตนุปสงฺกมิ, อุปสงฺกมิตฺวา ภควนฺตํ อภิวาเทตฺวา เอกมนฺตํ นิสีทิ, เอกมนฺตํ นิสินฺโน โข กฬารขตฺติโย ภิกฺขุ ภควนฺตํ เอตทโวจ “อายสฺมตา ภนฺเต สาริปุตฺเตน สีหนาโท นทิโต ‘ปุพฺเพ อปฺปฏิสํวิทิตํ มํ อาวุโส ภควา ปฐมํ ปญฺหํ อปุจฺฉิ, ตสฺส เม อโหสิ ทนฺธายิตตฺตํ.\csromanspacer{} ยโต จ โข เม อาวุโส ภควา ปฐมํ ปญฺหํ อนุโมทิ, ตสฺส มยฺหํ อาวุโส เอตทโหสิ ‘ทิวสํ เจปิ มํ ภควา เอตมตฺถํ ปุจฺเฉยฺย อญฺญมญฺเญหิ ปเทหิ อญฺญมญฺเญหิ ปริยาเยหิ, ทิวสมฺปาหํ ภควโต เอตมตฺถํ พฺยากเรยฺยํ อญฺญมญฺเญหิ ปเทหิ อญฺญมญฺเญหิ ปริยาเยหิ.\csromanspacer{} รตฺติํ เจปิ ฯเปฯ.\csromanspacer{} รตฺตินฺทิวํ เจปิ มํ ภควา ฯเปฯ.\csromanspacer{} เทฺว รตฺตินฺทิวานิ เจปิ มํ ภควา ฯเปฯ.\csromanspacer{} ตีณิ.\csromanspacer{} จตฺตาริ.\csromanspacer{} ปญฺจ.\csromanspacer{} ฉ.\csromanspacer{} สตฺต รตฺตินฺทิวานิ เจปิ มํ ภควา เอตมตฺถํ ปุจฺเฉยฺย อญฺญมญฺเญหิ ปเทหิ อญฺญมญฺเญหิ ปริยาเยหิ, สตฺต รตฺตินฺทิวานิปาหํ ภควโต เอตมตฺถํ พฺยากเรยฺยํ อญฺญมญฺเญหิ ปเทหิ อญฺญมญฺเญหิ ปริยาเยหี’ติ”.}

\prose{สา หิ ภิกฺขุ สาริปุตฺตสฺส ธมฺมธาตุ สุปฺปฏิวิทฺธา, ยสฺสา ธมฺมธาตุยา สุปฺปฏิวิทฺธตฺตา ทิวสํ เจปาหํ สาริปุตฺตํ เอตมตฺถํ ปุจฺเฉยฺยํ อญฺญมญฺเญหิ ปเทหิ อญฺญมญฺเญหิ ปริยาเยหิ, ทิวสมฺปิ เม สาริปุตฺโต เอตมตฺถํ พฺยากเรยฺย อญฺญมญฺเญหิ ปเทหิ อญฺญมญฺเญหิ ปริยาเยหิ.\csromanspacer{} รตฺติํ เจปาหํ สาริปุตฺตํ เอตมตฺถํ ปุจฺเฉยฺยํ อญฺญมญฺเญหิ ปเทหิ อญฺญมญฺเญหิ ปริยาเยหิ, รตฺติมฺปิ เม สาริปุตฺโต เอตมตฺถํ พฺยากเรยฺย ฯเปฯ.\csromanspacer{} รตฺตินฺทิวํ เจปาหํ สาริปุตฺตํ เอตมตฺถํ ปุจฺเฉยฺยํ, รตฺตินฺทิวมฺปิ เม สาริปุตฺโต เอตมตฺถํ พฺยากเรยฺย.\csromanspacer{} เทฺว รตฺตินฺทิวานิ เจปาหํ สาริปุตฺตํ เอตมตฺถํ ปุจฺเฉยฺยํ, เทฺว รตฺตินฺทิวานิปิ เม สาริปุตฺโต เอตมตฺถํ พฺยากเรยฺย.\csromanspacer{} ตีณิ รตฺตินฺทิวานิ เจปาหํ สาริปุตฺตํ เอตมตฺถํ ปุจฺเฉยฺยํ, ตีณิ รตฺตินฺทิวานิปิ เม สาริปุตฺโต}

\vspace{\tipitakaparskip}
\csromancenter{นิทานวคฺคสํยุตฺตปาฬิ เอตมตฺถํ พฺยากเรยฺย.\csromanspacer{} จตฺตาริ รตฺตินฺทิวานิ เจปาหํ สาริปุตฺตํ เอตมตฺถํ ปุจฺเฉยฺยํ, จตฺตาริ ริตฺตินฺทิวานิปิ เม สาริปุตฺโต เอตมตฺถํ พฺยากเรยฺย.\csromanspacer{} ปญฺจ รตฺตินฺทิวานิ เจปาหํ สาริปุตฺตํ เอตมตฺถํ ปุจฺเฉยฺยํ, ปญฺจ รตฺตินฺทิวานิปิ เม สาริปุตฺโต เอตมตฺถํ พฺยากเรยฺย.\csromanspacer{} ฉ รตฺตินฺทิวานิ เจปาหํ สาริปุตฺตํ เอตมตฺถํ ปุจฺเฉยฺยํ, ฉ ริตฺตินฺทิวานิปิ เม สาริปุตฺโต เอตมตฺถํ พฺยากเรยฺย.\csromanspacer{} สตฺต รตฺตินฺทิวานิ เจปาหํ สาริปุตฺตํ เอตมตฺถํ ปุจฺเฉยฺยํ อญฺญมญฺเญหิ ปเทหิ อญฺญมญฺเญหิ ปริยาเยหิ, สตฺต รตฺตินฺทิวานิปิ เม สาริปุตฺโต เอตมตฺถํ พฺยากเรยฺย อญฺญมญฺเญหิ ปเทหิ อญฺญมญฺเญหิ ปริยาเยหี”ติ.\csromanspacer{}\csromanspacer{}ทุติยํ.}
\csromanheader{2}{3. ญาณวตฺถุสุตฺต}
\proseitem{33}{\csromanfolio{288}สาวตฺถิยํ.\csromanspacer{} จตุจตฺตารีสํ โว ภิกฺขเว ญาณวตฺถูนิ เทเสสฺสามิ, ตํ สุณาถ สาธุกํ มนสิ กโรถ, ภาสิสฺสามีติ. “เอวํ ภนฺเต”ติ โข เต ภิกฺขู ภควโต ปจฺจสฺโสสุํ.\csromanspacer{} ภควา เอตทโวจ\csromandash{}}

\csromanmatikaanchor{mtk.38}
\csromantocmarknopage{section}{4. กฬารขตฺติยวคฺค}{mtk.38}
\bookmark[dest={mtk.38},level=1]{4. กฬารขตฺติยวคฺค}
\prose{กตมานิ\footnote{กตมานิ จ (สฺยา, กํ, อิ, ก)} ภิกฺขเว จตุจตฺตารีสํ ญาณวตฺถูนิ.\csromanspacer{} ชรามรเณ ญาณํ, ชรามรณสมุทเย ญาณํ, ชรามรณนิโรเธ ญาณํ, ชรามรณนิโรธคามินิยา ปฏิปทาย ญาณํ.\csromanspacer{} ชาติยา ญาณํ, ชาติสมุทเย ญาณํ, ชาตินิโรเธ ญาณํ, ชาตินิโรธคามินิยา ปฏิปทาย ญาณํ.\csromanspacer{} ภเว ญาณํ, ภวสมุทเย ญาณํ, ภวนิโรเธ ญาณํ, ภวนิโรธคามินิยา ปฏิปทาย ญาณํ.\csromanspacer{} อุปาทาเน ญาณํ, อุปาทานสมุทเย ญาณํ, อุปาทานนิโรเธ ญาณํ, อุปาทานนิโรธคามินิยา ปฏิปทาย ญาณํ.\csromanspacer{} ตณฺหาย ญาณํ, ตณฺหาสมุทเย ญาณํ, ตณฺหานิโรเธ ญาณํ, ตณฺหานิโรธคามินิยา ปฏิปทาย ญาณํ.\csromanspacer{} เวทนาย ญาณํ, เวทนาสมุทเย ญาณํ, เวทนานิโรเธ ญาณํ, เวทนานิโรธคามินิยา ปฏิปทาย ญาณํ.\csromanspacer{} ผสฺเส ญาณํ ฯเปฯ.\csromanspacer{} สฬายตเน ญาณํ.\csromanspacer{} นามรูเป ญาณํ.\csromanspacer{} วิญฺญาเณ ญาณํ.\csromanspacer{} สงฺขาเรสุ ญาณํ, สงฺขารสมุทเย ญาณํ, สงฺขารนิโรเธ ญาณํ, สงฺขารนิโรธคามินิยา ปฏิปทาย ญาณํ.\csromanspacer{} อิมานิ วุจฺจนฺติ ภิกฺขเว จตุจตฺตารีสํ ญาณวตฺถูนิ.}

\csromanmatikaanchor{mtk.39}
\csromantocmarkref{subsection}{1. ภูตสุตฺต}{mtk.39}
\bookmark[dest={mtk.39},level=2]{1. ภูตสุตฺต}
\prose{กตมญฺจ ภิกฺขเว ชรามรณํ.\csromanspacer{} ยา เตสํ เตสํ สตฺตานํ ตมฺหิ ตมฺหิ สตฺตนิกาเย ชรา ชีรณตา ขณฺฑิจฺจํ ปาลิจฺจํ วลิตฺตจตา อายุโน}

\vspace{\tipitakaparskip}
\csromancenter{นิทานสํยุตฺต สํหานิ อินฺทฺริยานํ ปริปาโก.\csromanspacer{} อยํ วุจฺจติ ชรา.\csromanspacer{} ยา เตสํ เตสํ สตฺตานํ ตมฺหา ตมฺหา สตฺตนิกายา จุติ จวนตา เภโท อนฺตรธานํ มจฺจุ มรณํ กาลกิริยา ขนฺธานํ เภโท กเฬวรสฺส นิกฺเขโป.\csromanspacer{} อิทํ วุจฺจติ มรณํ, อิติ อยญฺจ ชรา อิทญฺจ มรณํ.\csromanspacer{} อิทํ วุจฺจติ ภิกฺขเว ชรามรณํ.}
\prose{\csromanfolio{289}ชาติสมุทยา ชรามรณสมุทโย, ชาตินิโรธา ชรามรณนิโรโธ, อยเมว อริโย อฏฺฐงฺคิโก มคฺโค ชรามรณนิโรธคามินี ปฏิปทา.\csromanspacer{} เสยฺยถิทํ, สมฺมาทิฏฺฐิ ฯเปฯ สมฺมาสมาธิ.}

\prose{ยโต โข ภิกฺขเว อริยสาวโก เอวํ ชรามรณํ ปชานาติ, เอวํ ชรามรณสมุทยํ ปชานาติ, เอวํ ชรามรณนิโรธํ ปชานาติ, เอวํ ชรามรณนิโรธคามินิํ ปฏิปทํ ปชานาติ.\csromanspacer{} อิทมสฺส ธมฺเม ญาณํ, โส อิมินา ธมฺเมน ทิฏฺเฐน วิทิเตน อกาลิเกน ปตฺเตน ปริโยคาเฬฺหน อตีตานาคเตน ยํ เนติ.}

\prose{เย โข เกจิ อตีตมทฺธานํ สมณา วา พฺราหฺมณา วา ชรามรณํ อพฺภญฺญํสุ, ชรามรณสมุทยํ อพฺภญฺญํสุ, ชรามรณนิโรธํ อพฺภญฺญํสุ, ชรามรณนิโรธคามินิํ ปฏิปทํ อพฺภญฺญํสุ, สพฺเพเต เอวเมว อพฺภญฺญํสุ, เสยฺยถาปาหํ เอตรหิ.}

\prose{เยปิ หิ เกจิ อนาคตมทฺธานํ สมณา วา พฺราหฺมณา วา ชรามรณํ อภิชานิสฺสนฺติ, ชรามรณสมุทยํ อภิชานิสฺสนฺติ, ชรามรณนิโรธํ อภิชานิสฺสนฺติ, ชรามรณนิโรธคามินิํ ปฏิปทํ อภิชานิสฺสนฺติ, สพฺเพเต เอวเมว อภิชานิสฺสนฺติ, เสยฺยถาปาหํ เอตรหีติ.\csromanspacer{} อิทมสฺส อนฺวเย ญาณํ.}

\prose{ยโต โข ภิกฺขเว อริยสาวกสฺส อิมานิ เทฺว ญาณานิ ปริสุทฺธานิ โหนฺติ ปริโยทาตานิ ธมฺเม ญาณญฺจ อนฺวเย ญาณญฺจ.\csromanspacer{} อยํ วุจฺจติ ภิกฺขเว อริยสาวโก “ทิฏฺฐิสมฺปนฺโน” อิติปิ, “ทสฺสนสมฺปนฺโน” อิติปิ, “อาคโต อิมํ สทฺธมฺมํ” อิติปิ, “ปสฺสติ อิมํ สทฺธมฺมํ” อิติปิ, “เสกฺเขน ญาเณน สมนฺนาคโต” อิติปิ, “เสกฺขาย วิชฺชาย สมนฺนาคโต” อิติปิ, “ธมฺมโสตํ สมาปนฺโน” อิติปิ, “อริโย นิพฺเพธิกปญฺโญ” อิติปิ, “อมตทฺวารํ อาหจฺจ ติฏฺฐติ “อิติปีติ.}

\prose{กตมา จ ภิกฺขเว ชาติ.\csromanspacer{} กตโม จ ภิกฺขเว ภโว.\csromanspacer{} กตมญฺจ ภิกฺขเว อุปาทานํ.\csromanspacer{} กตมา จ ภิกฺขเว ตณฺหา.\csromanspacer{} กตมา จ ภิกฺขเว เวทนา.\csromanspacer{} กตโม จ ภิกฺขเว ผสฺโส.\csromanspacer{} กตมญฺจ ภิกฺขเว สฬายตนํ.\csromanspacer{} กตมญฺจ}

\vspace{\tipitakaparskip}
\csromancenter{นิทานวคฺคสํยุตฺตปาฬิ ภิกฺขเว นามรูปํ.\csromanspacer{} กตมญฺจ ภิกฺขเว วิญฺญาณํ.\csromanspacer{} กตเม จ ภิกฺขเว สงฺขารา.\csromanspacer{} ตโยเม ภิกฺขเว สงฺขารา, กายสงฺขาโร วจีสงฺขาโร จิตฺตสงฺขาโรติ.\csromanspacer{} อิเม วุจฺจนฺติ ภิกฺขเว สงฺขารา.}
\prose{\csromanfolio{290}อวิชฺชาสมุทยา สงฺขารสมุทโย, อวิชฺชานิโรธา สงฺขารนิโรโธ, อยเมว อริโย อฏฺฐงฺคิโก มคฺโค สงฺขารนิโรธคามินี ปฏิปทา.\csromanspacer{} เสยฺยถิทํ, สมฺมาทิฏฺฐิ ฯเปฯ สมฺมาสมาธิ.}

\prose{ยโต โข ภิกฺขเว อริยสาวโก เอวํ สงฺขาเร ปชานาติ, เอวํ สงฺขารสมุทยํ ปชานาติ, เอวํ สงฺขารนิโรธํ ปชานาติ, เอวํ สงฺขารนิโรธคามินิํ ปฏิปทํ ปชานาติ.\csromanspacer{} อิทมสฺส ธมฺเม ญาณํ, โส อิมินา ธมฺเมน ทิฏฺเฐน วิทิเตน อกาลิเกน ปตฺเตน ปริโยคาเฬฺหน อตีตานาคเตน ยํ เนติ.}

\prose{เย โข เกจิ อตีตมทฺธานํ สมณา วา พฺราหฺมณา วา สงฺขาเร อพฺภญฺญํสุ, สงฺขารสมุทยํ อพฺภญฺญํสุ, สงฺขารนิโรธํ อพฺภญฺญํสุ, สงฺขารนิโรธคามินิํ ปฏิปทํ อพฺภญฺญํสุ, สพฺเพเต เอวเมว อพฺภญฺญํสุ, เสยฺยถาปาหํ เอตรหิ.}

\prose{เย หิปิ เกจิ อนาคตมทฺธานํ สมณา วา พฺราหฺมณา วา สงฺขาเร อภิชานิสฺสนฺติ, สงฺขารสมุทยํ อภิชานิสฺสนฺติ, สงฺขารนิโรธํ อภิชานิสฺสนฺติ, สงฺขารนิโรธคามินิํ ปฏิปทํ อภิชานิสฺสนฺติ, สพฺเพเต เอวเมว อภิชานิสฺสนฺติ, เสยฺยถาปาหํ เอตรหิ.\csromanspacer{} อิทมสฺส อนฺวเย ญาณํ.}

\prose{ยโต โข ภิกฺขเว อริยสาวกสฺส อิมานิ เทฺว ญาณานิ ปริสุทฺธานิ โหนฺติ ปริโยทาตานิ ธมฺเม ญาณญฺจ อนฺวเย ญาณญฺจ.\csromanspacer{} อยํ วุจฺจติ ภิกฺขเว อริยสาวโก “ทิฏฺฐิสมฺปนฺโน” อิติปิ, “ทสฺสนสมฺปนฺโน” อิติปิ, “อาคโต อิมํ สทฺธมฺมํ” อิติปิ, “ปสฺสติ อิมํ สทฺธมฺมํ” อิติปิ, “เสกฺเขน ญาเณน สมนฺนาคโต” อิติปิ, “เสกฺขาย วิชฺชาย สมนฺนาคโต” อิติปิ, “ธมฺมโสตํ สมาปนฺโน” อิติปิ, “อริโย นิพฺเพธิกปญฺโญ” อิติปิ, “อมตทฺวารํ อาหจฺจ ติฏฺฐติ” อิติปีติ.\csromanspacer{}\csromanspacer{}ตติยํ.}

\csromanheader{2}{4. ทุติยญาณวตฺถุสุตฺต}
\proseitem{34}{สาวตฺถิยํ วิหรติ.\csromanspacer{} สตฺตสตฺตริ โว ภิกฺขเว ญาณวตฺถูนิ เทเสสฺสามิ, ตํ สุณาถ สาธุกํ มนสิ กโรถ, ภาสิสฺสามีติ. “เอวํ ภนฺเต”ติ โข เต ภิกฺขู ภควโต ปจฺจสฺโสสุํ.\csromanspacer{} ภควา เอตทโวจ\csromandash{}}

\csromanmatikaanchor{mtk.40}
\csromantocmarkref{subsection}{2. กฬารสุตฺต}{mtk.40}
\bookmark[dest={mtk.40},level=2]{2. กฬารสุตฺต}
\csromanheader{2}{1. นิทานสํยุตฺต}
\prose{\csromanfolio{291}กตมานิ ภิกฺขเว สตฺตสตฺตริ ญาณวตฺถูนิ.\csromanspacer{} ชาติปจฺจยา ชรามรณนฺติ ญาณํ, อสติ ชาติยา นตฺถิ ชรามรณนฺติ ญาณํ.\csromanspacer{} อตีตมฺปิ อทฺธานํ ชาติปจฺจยา ชรามรณนฺติ ญาณํ, อสติ ชาติยา นตฺถิ ชรามรณนฺติ ญาณํ.\csromanspacer{} อนาคตมฺปิ อทฺธานํ ชาติปจฺจยา ชรามรณนฺติ ญาณํ, อสติ ชาติยา นตฺถิ ชรามรณนฺติ ญาณํ.\csromanspacer{} ยมฺปิสฺส ตํ ธมฺมฏฺฐิติญาณํ, ตมฺปิ ขยธมฺมํ วยธมฺมํ วิราคธมฺมํ นิโรธธมฺมนฺติ ญาณํ.}

\prose{ภวปจฺจยา ชาตีติ ญาณํ ฯเปฯ.\csromanspacer{} อุปาทานปจฺจยา ภโวติ ญาณํ.\csromanspacer{} ตณฺหาปจฺจยา อุปาทานนฺติ ญาณํ.\csromanspacer{} เวทนาปจฺจยา ตณฺหาติ ญาณํ.\csromanspacer{} ผสฺสปจฺจยา เวทนาติ ญาณํ.\csromanspacer{} สฬายตนปจฺจยา ผสฺโสติ ญาณํ.\csromanspacer{} นามรูปปจฺจยา สฬายตนนฺติ ญาณํ.\csromanspacer{} วิญฺญาณปจฺจยา นามรูปนฺติ ญาณํ, สงฺขารปจฺจยา วิญฺญาณนฺติ ญาณํ.\csromanspacer{} อวิชฺชาปจฺจยา สงฺขาราติ ญาณํ, อสติ อวิชฺชาย นตฺถิ สงฺขาราติ ญาณํ.\csromanspacer{} อตีตมฺปิ อทฺธานํ อวิชฺชาปจฺจยา สงฺขาราติ ญาณํ, อสติ อวิชฺชาย นตฺถิ สงฺขาราติ ญาณํ.\csromanspacer{} อนาคตมฺปิ อทฺธานํ อวิชฺชาปจฺจยา สงฺขาราติ ญาณํ, อสติ อวิชฺชาย นตฺถิ สงฺขาราติ ญาณํ.\csromanspacer{} ยมฺปิสฺส ตํ ธมฺมฏฺฐิติญาณํ, ตมฺปิ ขยธมฺมํ วยธมฺมํ วิราคธมฺมํ นิโรธธมฺมนฺติ ญาณํ.\csromanspacer{} อิมานิ วุจฺจนฺติ ภิกฺขเว สตฺตสตฺตริ ญาณวตฺถูนีติ.\csromanspacer{}\csromanspacer{}จตุตฺถํ.}

\csromanheader{2}{5. อวิชฺชาปจฺจยสุตฺต}
\proseitem{35}{สาวตฺถิยํ วิหรติ.\csromanspacer{} อวิชฺชาปจฺจยา ภิกฺขเว สงฺขารา, สงฺขารปจฺจยา วิญฺญาณํ ฯเปฯ เอวเมตสฺส เกวลสฺส ทุกฺขกฺขนฺธสฺส สมุทโย โหตีติ.\csromanspacer{} เอวํ วุตฺเต อญฺญตโร ภิกฺขุ ภควนฺตํ เอตทโวจ “กตมํ นุ โข ภนฺเต ชรามรณํ, กสฺส จ ปนิทํ ชรามรณนฺ”ติ. “โน กลฺโล ปโญฺห”ติ ภควา อโวจ, “กตมํ ชรามรณํ, กสฺส จ ปนิทํ ชรามรณนฺ”ติ อิติ วา ภิกฺขุ โย วเทยฺย, “อญฺญํ ชรามรณํ, อญฺญสฺส จ ปนิทํ ชรามรณนฺ”ติ อิติ วา ภิกฺขุ โย วเทยฺย.\csromanspacer{} อุภยเมตํ เอกตฺถํ พฺยญฺชนเมว นานํ. “ตํ ชีวํ ตํ สรีรนฺ”ติ วา ภิกฺขุ ทิฏฺฐิยา สติ พฺรหฺมจริยวาโส น โหติ, “อญฺญํ ชีวํ อญฺญํ สรีรนฺ”ติ วา ภิกฺขุ ทิฏฺฐิยา สติ พฺรหฺมจริยวาโส น โหติ.\csromanspacer{} เอเต เต ภิกฺขุ อุโภ อนฺเต อนุปคมฺม มชฺเฌน ตถาคโต ธมฺมํ เทเสติ “ชาติปจฺจยา ชรามรณนฺ”ติ.}

\prose{“กตมา นุ โข ภนฺเต ชาติ, กสฺส จ ปนายํ ชาตี”ติ. “โน กลฺโล ปโญฺห”ติ ภควา อโวจ, “กตมา ชาติ, กสฺส จ ปนายํ ชาตี”ติ}

\vspace{\tipitakaparskip}
\csromancenter{นิทานวคฺคสํยุตฺตปาฬิ อิติ วา ภิกฺขุ โย วเทยฺย, “อญฺญา ชาติ, อญฺญสฺส จ ปนายํ ชาตี”ติ อิติ วา ภิกฺขุ โย วเทยฺย.\csromanspacer{} อุภยเมตํ เอกตฺถํ พฺยญฺชนเมว นานํ. “ตํ ชีวํ ตํ สรีรนฺ”ติ วา ภิกฺขุ ทิฏฺฐิยา สติ พฺรหฺมจริยวาโส น โหติ, “อญฺญํ ชีวํ อญฺญํ สรีรนฺ”ติ วา ภิกฺขุ ทิฏฺฐิยา สติ พฺรหฺมจริยวาโส น โหติ.\csromanspacer{} เอเต เต ภิกฺขุ อุโภ อนฺเต อนุปคมฺม มชฺเฌน ตถาคโต ธมฺมํ เทเสติ “ภวปจฺจยา ชาตี”ติ.}
\prose{\csromanfolio{292}“กตโม นุ โข ภนฺเต ภโว, กสฺส จ ปนายํ ภโว”ติ. “โน กลฺโล ปโญฺห”ติ ภควา อโวจ, “กตโม ภโว, กสฺส จ ปนายํ ภโว”ติ อิติ วา ภิกฺขุ โย วเทยฺย, “อญฺโญ ภโว, อญฺญสฺส จ ปนายํ ภโว”ติ อิติ วา ภิกฺขุ โย วเทยฺย.\csromanspacer{} อุภยเมตํ เอกตฺถํ พฺยญฺชนเมว นานํ. “ตํ ชีวํ ตํ สรีรนฺ”ติ วา ภิกฺขุ ทิฏฺฐิยา สติ พฺรหฺมจริยวาโส น โหติ, “อญฺญํ ชีวํ อญฺญํ สรีรนฺ”ติ วา ภิกฺขุ ทิฏฺฐิยา สติ พฺรหฺมจริยวาโส น โหติ.\csromanspacer{} เอเต เต ภิกฺขุ อุโภ อนฺเต อนุปคมฺม มชฺเฌน ตถาคโต ธมฺมํ เทเสติ “อุปาทานปจฺจยา ภโว”ติ ฯเปฯ “ตณฺหาปจฺจยา อุปาทานนฺ”ติ. “เวทนาปจฺจยา ตณฺหา”ติ. “ผสฺสปจฺจยา เวทนา”ติ. “สฬายตนปจฺจยา ผสฺโส”ติ. “นามรูปปจฺจยา สฬายตนนฺ”ติ. “วิญฺญาณปจฺจยา นามรูปนฺ”ติ. “สงฺขารปจฺจยา วิญฺญาณนฺ”ติ.}

\prose{“กตเม นุ โข ภนฺเต สงฺขารา, กสฺส จ ปนิเม สงฺขารา”ติ. “โน กลฺโล ปโญฺห”ติ ภควา อโวจ, “กตเม สงฺขารา, กสฺส จ ปนิเม สงฺขารา”ติ อิติ วา ภิกฺขุ โย วเทยฺย, “อญฺเญ สงฺขารา, อญฺญสฺส จ ปนิเม สงฺขารา”ติ อิติ วา ภิกฺขุ โย วเทยฺย.\csromanspacer{} อุภยเมตํ เอกตฺถํ พฺยญฺชนเมว นานํ. “ตํ ชีวํ ตํ สรีรนฺ”ติ วา ภิกฺขุ ทิฏฺฐิยา สติ พฺรหฺมจริยวาโส น โหติ, “อญฺญํ ชีวํ อญฺญํ สรีรนฺ”ติ วา ภิกฺขุ ทิฏฺฐิยา สติ พฺรหฺมจริยวาโส น โหติ.\csromanspacer{} เอเต เต ภิกฺขุ อุโภ อนฺเต อนุปคมฺม มชฺเฌน ตถาคโต ธมฺมํ เทเสติ “อวิชฺชาปจฺจยา สงฺขารา”ติ.}

\prose{อวิชฺชาย เตฺวว ภิกฺขุ อเสสวิราคนิโรธา, ยานิสฺส ตานิ วิสูกายิกานิ วิเสวิตานิ วิปฺผนฺทิตานิ กานิจิ กานิจิ. “กตมํ ชรามรณํ, กสฺส จ ปนิทํ ชรามรณํ” อิติ วา. “อญฺญํ ชรามรณํ, อญฺญสฺส จ ปนิทํ ชรามรณํ” อิติ วา. “ตํ ชีวํ ตํ สรีรํ” อิติ วา. “อญฺญํ ชีวํ อญฺญํ สรีรํ” อิติ วา.\csromanspacer{} สพฺพานิสฺส ตานิ ปหีนานิ ภวนฺติ อุจฺฉินฺนมูลานิ ตาลาวตฺถุกตานิ อนภาวํกตานิ อายติํ อนุปฺปาทธมฺมานิ.}

\csromanheader{2}{1. นิทานสํยุตฺต}
\prose{\csromanfolio{293}อวิชฺชาย เตฺวว ภิกฺขุ อเสสวิราคนิโรธา, ยานิสฺส ตานิ วิสูกายิกานิ วิเสวิตานิ วิปฺผนฺทิตานิ กานิจิ กานิจิ. “กตมา ชาติ, กสฺส จ ปนายํ ชาติ” อิติ วา. “อญฺญา ชาติ, อญฺญสฺส จ ปนายํ ชาติ” อิติ วา. “ตํ ชีวํ ตํ สรีรํ” อิติ วา. “อญฺญํ ชีวํ อญฺญํ สรีรํ” อิติ วา.\csromanspacer{} สพฺพานิสฺส ตานิ ปหีนานิ ภวนฺติ อุจฺฉินฺนมูลานิ ตาลาวตฺถุกตานิ อนภาวํกตานิ อายติํ อนุปฺปาทธมฺมานิ.}

\prose{อวิชฺชาย เตฺวว ภิกฺขุ อเสสวิราคนิโรธา, ยานิสฺส ตานิ วิสูกายิกานิ วิเสวิตานิ วิปฺผนฺทิตานิ กานิจิ กานิจิ.\csromanspacer{} กตโม ภโว ฯเปฯ.\csromanspacer{} กตมํ อุปาทานํ.\csromanspacer{} กตมา ตณฺหา.\csromanspacer{} กตมา เวทนา.\csromanspacer{} กตโม ผสฺโส.\csromanspacer{} กตมํ สฬายตนํ.\csromanspacer{} กตมํ นามรูปํ.\csromanspacer{} กตมํ วิญฺญาณํ ฯเปฯ.}

\prose{อวิชฺชาย เตฺวว ภิกฺขุ อเสสวิราคนิโรธา, ยานิสฺส ตานิ วิสูกายิกานิ วิเสวิตานิ วิปฺผนฺทิตานิ กานิจิ กานิจิ. “กตเม สงฺขารา, กสฺส จ ปนิเม สงฺขารา” อิติ วา. “อญฺเญ สงฺขารา, อญฺญสฺส จ ปนิเม สงฺขารา” อิติ วา. “ตํ ชีวํ ตํ สรีรํ” อิติ วา. “อญฺญํ ชีวํ อญฺญํ สรีรํ” อิติ วา.\csromanspacer{} สพฺพานิสฺส ตานิ ปหีนานิ ภวนฺติ อุจฺฉินฺนมูลานิ ตาลาวตฺถุกตานิ อนภาวํกตานิ อายติํ อนุปฺปาทธมฺมานีติ.\csromanspacer{}\csromanspacer{}ปญฺจมํ.}

\csromanheader{2}{6. ทุติย-อวิชฺชาปจฺจยสุตฺต}
\proseitem{36}{สาวตฺถิยํ วิตรติ.\csromanspacer{} อวิชฺชาปจฺจยา ภิกฺขเว สงฺขารา, สงฺขารปจฺจยา วิญฺญาณํ ฯเปฯ เอวเมตสฺส เกวลสฺส ทุกฺขกฺขนฺธสฺส สมุทโย โหติ.}

\prose{“กตมํ ชรามรณํ, กสฺส จ ปนิทํ ชรามรณนฺ”ติ อิติ วา ภิกฺขเว โย วเทยฺย, “อญฺญํ ชรามรณํ, อญฺญสฺส จ ปนิทํ ชรามรณนฺ”ติ อิติ วา ภิกฺขเว โย วเทยฺย.\csromanspacer{} อุภยเมตํ เอกตฺถํ พฺยญฺชนเมว นานํ. “ตํ ชีวํ ตํ สรีรํ” อิติ วา ภิกฺขเว ทิฏฺฐิยา สติ พฺรหฺมจริยวาโส น โหติ, “อญฺญํ ชีวํ อญฺญํ สรีรํ” อิติ วา ภิกฺขเว ทิฏฺฐิยา สติ พฺรหฺมจริยวาโส น โหติ.\csromanspacer{} เอเต เต ภิกฺขเว อุโภ อนฺเต อนุปคมฺม มชฺเฌน ตถาคโต ธมฺมํ เทเสติ “ชาติปจฺจยา ชรามรณนฺ”ติ.}

\gambhira{นิทานวคฺคสํยุตฺตปาฬิ}
\prose{\csromanfolio{294}กตมํ ชาติ ฯเปฯ.\csromanspacer{} กตโม ภโว.\csromanspacer{} กตมํ อุปาทานํ.\csromanspacer{} กตมา ตณฺหา.\csromanspacer{} กตมา เวทนา.\csromanspacer{} กตโม ผสฺโส.\csromanspacer{} กตมํ สฬายตนํ.\csromanspacer{} กตมํ นามรูปํ.\csromanspacer{} กตมํ วิญฺญาณํ. “กตเม สงฺขารา, กสฺส จ ปนิเม สงฺขารา”ติ อิติ วา ภิกฺขเว โย วเทยฺย, “อญฺเญ สงฺขารา, อญฺญสฺส จ ปนิเม สงฺขารา” อิติ วา ภิกฺขเว โย วเทยฺย.\csromanspacer{} อุภยเมตํ เอกตฺถํ พฺยญฺชนเมว นานํ. “ตํ ชีวํ ตํ สรีรํ” อิติ วา ภิกฺขเว ทิฏฺฐิยา สติ พฺรหฺมจริยวาโส น โหติ, “อญฺญํ ชีวํ อญฺญํ สรีรํ” อิติ วา ภิกฺขเว ทิฏฺฐิยา สติ พฺรหฺมจริยวาโส น โหติ.\csromanspacer{} เอเต เต ภิกฺขเว อุโภ อนฺเต อนุปคมฺม มชฺเฌน ตถาคโต ธมฺมํ เทเสติ “อวิชฺชาปจฺจยา สงฺขารา”ติ.}

\csromanmatikaanchor{mtk.41}
\csromantocmarkref{subsection}{3. ญาณวตฺถุสุตฺต}{mtk.41}
\bookmark[dest={mtk.41},level=2]{3. ญาณวตฺถุสุตฺต}
\prose{อวิชฺชาย เตฺวว ภิกฺขเว อเสสวิราคนิโรธา, ยานิสฺส ตานิ วิสูกายิกานิ วิเสวิตานิ วิปฺผนฺทิตานิ กานิจิ กานิจิ. “กตมํ ชรามรณํ, กสฺส จ ปนิทํ ชรามรณํ” อิติ วา. “อญฺญํ ชรามรณํ, อญฺญสฺส จ ปนิทํ ชรามรณํ” อิติ วา. “ตํ ชีวํ ตํ สรีรํ” อิติ วา. “อญฺญํ ชีวํ อญฺญํ สรีรํ” อิติ วา.\csromanspacer{} สพฺพานิสฺส ตานิ ปหีนานิ ภวนฺติ อุจฺฉินฺนมูลานิ ตาลาวตฺถุกตานิ อนภาวํกตานิ อายติํ อนุปฺปาทธมฺมานิ.}

\prose{อวิชฺชาย เตฺวว ภิกฺขเว อเสสวิราคนิโรธา, ยานิสฺส ตานิ วิสูกายิกานิ วิเสวิตานิ วิปฺผนฺทิตานิ กานิจิ กานิจิ.\csromanspacer{} กตมา ชาติ ฯเปฯ.\csromanspacer{} กตโม ภโว.\csromanspacer{} กตมํ อุปาทานํ.\csromanspacer{} กตมา ตณฺหา.\csromanspacer{} กตมา เวทนา.\csromanspacer{} กตโม ผสฺโส.\csromanspacer{} กตมํ สฬายตนํ.\csromanspacer{} กตมํ นามรูปํ.\csromanspacer{} กตมํ วิญฺญาณํ. “กตเม สงฺขารา, กสฺส จ ปนิเม สงฺขารา” อิติ วา. “อญฺเญ สงฺขารา, อญฺญสฺส จ ปนิเม สงฺขารา” อิติ วา. “ตํ ชีวํ ตํ สรีรํ” อิติ วา. “อญฺญํ ชีวํ อญฺญํ สรีรํ” อิติ วา.\csromanspacer{} สพฺพานิสฺส ตานิ ปหีนานิ ภวนฺติ อุจฺฉินฺนมูลานิ ตาลาวตฺถุกตานิ อนภาวํกตานิ อายติํ อนุปฺปาทธมฺมานีติ.\csromanspacer{}\csromanspacer{}ฉฏฺฐํ.}

\csromanheader{2}{7. นตุมฺหสุตฺต}
\proseitem{37}{สาวตฺถิยํ วิหรติ.\csromanspacer{} นายํ ภิกฺขเว กาโย ตุมฺหากํ, นปิ อญฺเญสํ.\csromanspacer{} ปุราณมิทํ ภิกฺขเว กมฺมํ อภิสงฺขตํ อภิสญฺเจตยิตํ เวทนิยํ ทฏฺฐพฺพํ.}

\csromanheader{2}{1. นิทานสํยุตฺต}
\prose{\csromanfolio{295}ตตฺร โข ภิกฺขเว สุตวา อริยสาวโก ปฏิจฺจสมุปฺปาทญฺเญว สาธุกํ โยนิโส มนสิ กโรติ “อิติ อิมสฺมิํ สติ อิทํ โหติ, อิมสฺสุปฺปาทา อิทํ อุปฺปชฺชติ.\csromanspacer{} อิมสฺมิํ อสติ อิทํ น โหติ, อิมสฺส นิโรธา อิทํ นิรุชฺฌติ.\csromanspacer{} ยทิทํ อวิชฺชาปจฺจยา สงฺขารา, สงฺขารปจฺจยา วิญฺญาณํ ฯเปฯ เอวเมตสฺส เกวลสฺส ทุกฺขกฺขนฺธสฺส สมุทโย โหติ.\csromanspacer{} อวิชฺชาย เตฺวว อเสสวิราคนิโรธา สงฺขารนิโรโธ, สงฺขารนิโรธา วิญฺญาณนิโรโธ ฯเปฯ เอวเมตสฺส เกวลสฺส ทุกฺขกฺขนฺธสฺส นิโรโธ โหตี”ติ.\csromanspacer{}\csromanspacer{}สตฺตมํ.}

\csromanheader{2}{8. เจตนาสุตฺต}
\proseitem{38}{สาวตฺถิยํ วิหรติ.\csromanspacer{} ยญฺจ ภิกฺขเว เจเตติ, ยญฺจ ปกปฺเปติ, ยญฺจ อนุเสติ.\csromanspacer{} อารมฺมณเมตํ\footnote{อารมณเมตํ (?)} โหติ วิญฺญาณสฺส ฐิติยา, อารมฺมเณ สติ ปติฏฺฐา วิญฺญาณสฺส โหติ, ตสฺมิํ ปติฏฺฐิเต วิญฺญาเณ วิรูเฬฺห อายติํ ปุนพฺภวาภินิพฺพตฺติ โหติ, อายติํ ปุนพฺภวาภินิพฺพตฺติยา สติ อายติํ ชาติ ชรามรณํ โสกปริเทวทุกฺขโทมนสฺสุปายาสา สมฺภวนฺติ.\csromanspacer{} เอวเมตสฺส เกวลสฺส ทุกฺขกฺขนฺธสฺส สมุทโย โหติ.}

\prose{โน เจ ภิกฺขเว เจเตติ, โน เจ ปกปฺเปติ, อถ เจ อนุเสติ.\csromanspacer{} อารมฺมณเมตํ โหติ วิญฺญาณสฺส ฐิติยา, อารมฺมเณ สติ ปติฏฺฐา วิญฺญาณสฺส โหติ, ตสฺมิํ ปติฏฺฐิเต วิญฺญาเณ วิรูเฬฺห อายติํ ปุนพฺภวาภินิพฺพตฺติ โหติ, อายติํ ปุนพฺภวาภินิพฺพตฺติยา สติ อายติํ ชาติ ชรามรณํ โสกปริเทวทุกฺขโทมนสฺสุปายาสา สมฺภวนฺติ.\csromanspacer{} เอวเมตสฺส เกวลสฺส ทุกฺขกฺขนฺธสฺส สมุทโย โหติ.}

\prose{ยโต จ โข ภิกฺขเว โน เจว เจเตติ, โน จ ปกปฺเปติ, โน จ อนุเสติ.\csromanspacer{} อารมฺมณเมตํ น โหติ วิญฺญาณสฺส ฐิติยา, อารมฺมเณ อสติ ปติฏฺฐา วิญฺญาณสฺส น โหติ, ตทปฺปติฏฺฐิเต วิญฺญาเณ อวิรูเฬฺห อายติํ ปุนพฺภวาภินิพฺพตฺติ น โหติ, อายติํ ปุนพฺภวาภินิพฺพตฺติยา อสติ อายติํ ชาติ ชรามรณํ โสกปริเทวทุกฺขโทมนสฺสุปายาสา นิรุชฺฌนฺติ.\csromanspacer{} เอวเมตสฺส เกวลสฺส ทุกฺขกฺขนฺธสฺส นิโรโธ โหตีติ.\csromanspacer{}\csromanspacer{}อฏฺฐมํ.}

\csromanheader{2}{นิทานวคฺคสํยุตฺตปาฬิ \textbf{9. ทุติยเจตนาสุตฺต}}
\proseitem{39}{\csromanfolio{296}สาวตฺถิยํ วิหรติ.\csromanspacer{} ยญฺจ ภิกฺขเว เจเตติ, ยญฺจ ปกปฺเปติ, ยญฺจ อนุเสติ.\csromanspacer{} อารมฺมณเมตํ โหติ วิญฺญาณสฺส ฐิติยา, อารมฺมเณ สติ ปติฏฺฐา วิญฺญาณสฺส โหติ, ตสฺมิํ ปติฏฺฐิเต วิญฺญาเณ วิรูเฬฺห นามรูปสฺส อวกฺกนฺติ โหติ.\csromanspacer{} นามรูปปจฺจยา สฬายตนํ, สฬายตนปจฺจยา ผสฺโส, ผสฺสปจฺจยา เวทนา ฯเปฯ ตณฺหา.\csromanspacer{} อุปาทานํ.\csromanspacer{} ภโว.\csromanspacer{} ชาติ.\csromanspacer{} ชรามรณํ.\csromanspacer{} โสกปริเทวทุกฺขโทมนสฺสุปายาสา สมฺภวนฺติ.\csromanspacer{} เอวเมตสฺส เกวลสฺส ทุกฺขกฺขนฺธสฺส สมุทโย โหติ.}

\prose{โน เจ ภิกฺขเว เจเตติ, โน เจ ปกปฺเปติ, อถ เจ อนุเสติ.\csromanspacer{} อารมฺมณเมตํ โหติ วิญฺญาณสฺส ฐิติยา, อารมฺมเณ สติ ปติฏฺฐา วิญฺญาณสฺส โหติ, ตสฺมิํ ปติฏฺฐิเต วิญฺญาเณ วิรูเฬฺห นามรูปสฺส อวกฺกนฺติ โหติ.\csromanspacer{} นามรูปปจฺจยา สฬายตนํ ฯเปฯ เอวเมตสฺส เกวลสฺส ทุกฺขกฺขนฺธสฺส สมุทโย โหติ.}

\prose{ยโต จ โข ภิกฺขเว โน เจว เจเตติ, โน จ ปกปฺเปติ, โน จ อนุเสติ.\csromanspacer{} อารมฺมณเมตํ น โหติ วิญฺญาณสฺส ฐิติยา, อารมฺมเณ อสติ ปติฏฺฐา วิญฺญาณสฺส น โหติ, ตทปฺปติฏฺฐิเต วิญฺญาเณ อวิรูเฬฺห นามรูปสฺส อวกฺกนฺติ น โหติ.\csromanspacer{} นามรูปนิโรธา สฬายตนนิโรโธ ฯเปฯ เอวเมตสฺส เกวลสฺส ทุกฺขกฺขนฺธสฺส นิโรโธ โหตีติ.\csromanspacer{}\csromanspacer{}นวมํ.}

\csromanheader{2}{10. ตติยเจตนาสุตฺต}
\proseitem{40}{สาวตฺถิยํ วิหรติ.\csromanspacer{} ยญฺจ ภิกฺขเว เจเตติ, ยญฺจ ปกปฺเปติ, ยญฺจ อนุเสติ.\csromanspacer{} อารมฺมณเมตํ โหติ วิญฺญาณสฺส ฐิติยา, อารมฺมเณ สติ ปติฏฺฐา วิญฺญาณสฺส โหติ, ตสฺมิํ ปติฏฺฐิเต วิญฺญาเณ วิรูเฬฺห นติ โหติ, นติยา สติ อาคติคติ โหติ, อาคติคติยา สติ จุตูปปาโต โหติ, จุตูปปาเต สติ อายติํ ชาติ ชรามรณํ โสกปริเทวทุกฺขโทมนสฺสุปายาสา สมฺภวนฺติ.\csromanspacer{} เอวเมตสฺส เกวลสฺส ทุกฺขกฺขนฺธสฺส สมุทโย โหติ.}

\prose{โน เจ ภิกฺขเว เจเตติ, โน เจ ปกปฺเปติ, อถ เจ อนุเสติ.\csromanspacer{} อารมฺมณเมตํ โหติ วิญฺญาณสฺส ฐิติยา, อารมฺมเณ สติ ปติฏฺฐา วิญฺญาณสฺส โหติ, ตสฺมิํ ปติฏฺฐิเต วิญฺญาเณ วิรูเฬฺห นติ โหติ,}

\csromanmatikaanchor{mtk.42}
\csromantocmarkref{subsection}{4. ทุติยญาณวตฺถุสุตฺต}{mtk.42}
\bookmark[dest={mtk.42},level=2]{4. ทุติยญาณวตฺถุสุตฺต}
\vspace{\tipitakaparskip}
\csromancenter{นิทานสํยุตฺต นติยา สติ อาคติคติ โหติ, อาคติคติยา สติ จุตูปปาโต โหติ, จุตูปปาเต สติ อายติํ ชาติ ชรามรณํ โสกปริเทวทุกฺขโทมนสฺสุปายาสา สมฺภวนฺติ.\csromanspacer{} เอวเมตสฺส เกวลสฺส ทุกฺขกฺขนฺธสฺส สมุทโย โหติ.}
\csromanmatikaanchor{mtk.43}
\csromanmatikaanchor{mtk.49}
\csromantocmarkref{subsection}{5. อวิชฺชาปจฺจยสุตฺต}{mtk.43}
\bookmark[dest={mtk.43},level=2]{5. อวิชฺชาปจฺจยสุตฺต}
\prose{\csromanfolio{297}ยโต จ โข ภิกฺขเว โน เจว เจเตติ, โน จ ปกปฺเปติ, โน จ อนุเสติ.\csromanspacer{} อารมฺมณเมตํ น โหติ วิญฺญาณสฺส ฐิติยา, อารมฺมเณ อสติ ปติฏฺฐา วิญฺญาณสฺส น โหติ, ตทปฺปติฏฺฐิเต วิญฺญาเณ อวิรูเฬฺห นติ น โหติ, นติยา อสติ อาคติคติ น โหติ, อาคติคติยา อสติ จุตูปปาโต น โหติ, จุตูปปาเต อสติ อายติํ ชาติ ชรามรณํ โสกปริเทวทุกฺขโทมนสฺสุปายาสา นิรุชฺฌนฺติ.\csromanspacer{} เอวเมตสฺส เกวลสฺส ทุกฺขกฺขนฺธสฺส นิโรโธ โหตีติ.\csromanspacer{}\csromanspacer{}ทสมํ.}

\vspace{\tipitakaparskip}
\csromancenter{กฬารขตฺติยวคฺโค จตุตฺโถ.}
\titlehead{\textbf{ตสฺสุทฺทานํ}}
\csromangathameasure{ภูตมิทํ กฬารญฺจ, ทุเว จ ญาณวตฺถูนิ. \\ อวิชฺชาปจฺจยา จ เทฺว, นตุมฺหา เจตนา ตโยติ.}
\csromangathagroup{\csromangathastanza{ภูตมิทํ กฬารญฺจ, ทุเว จ ญาณวตฺถูนิ. \\ อวิชฺชาปจฺจยา จ เทฺว, นตุมฺหา เจตนา ตโยติ.}}

\csromanheader{1}{5. คหปติวคฺค}
\csromanheader{2}{1. ปญฺจเวรภยสุตฺต}
\proseitem{41}{สาวตฺถิยํ วิหรติ.\csromanspacer{} อถ โข อนาถปิณฺฑิโก คหปติ เยน ภควา เตนุปสงฺกมิ, อุปสงฺกมิตฺวา ภควนฺตํ อภิวาเทตฺวา เอกมนฺตํ นิสีทิ, เอกมนฺตํ นิสินฺนํ โข อนาถปิณฺฑิกํ คหปติํ ภควา เอตทโวจ\csromandash{}}

\prose{ยโต โข คหปติ อริยสาวกสฺส ปญฺจ ภยานิ เวรานิ วูปสนฺตานิ โหนฺติ, จตูหิ จ โสตาปตฺติยงฺเคหิ สมนฺนาคโต โหติ, อริโย จสฺส ญาโย ปญฺญาย สุทิฏฺโฐ โหติ สุปฺปฏิวิทฺโธ, โส อากงฺขมาโน อตฺตนาว อตฺตานํ พฺยากเรยฺย “ขีณนิรโยมฺหิ ขีณติรจฺฉานโยนิ ขีณเปตฺติวิสโย ขีณาปายทุคฺคติวินิปาโต, โสตาปนฺโนหมสฺมิ อวินิปาตธมฺโม นิยโต สมฺโพธิปรายโน”ติ.}

\gambhira{นิทานวคฺคสํยุตฺตปาฬิ}
\prose{\csromanfolio{298}กตมานิ ปญฺจ ภยานิ เวรานิ วูปสนฺตานิ โหนฺติ.\csromanspacer{} ยํ คหปติ ปาณาติปาตี ปาณาติปาตปจฺจยา ทิฏฺฐธมฺมิกมฺปิ ภยํ เวรํ ปสวติ, สมฺปรายิกมฺปิ ภยํ เวรํ ปสวติ, เจตสิกมฺปิ ทุกฺขํ โทมนสฺสํ ปฏิสํเวทยติ.\csromanspacer{} ปาณาติปาตา ปฏิวิรตสฺส เอวํ ตํ ภยํ เวรํ วูปสนฺตํ โหติ. (1)}

\prose{ยํ คหปติ อทินฺนาทายี อทินฺนาทานปจฺจยา ทิฏฺฐธมฺมิกมฺปิ ภยํ เวรํ ปสวติ, สมฺปรายิกมฺปิ ภยํ เวรํ ปสวติ, เจตสิกมฺปิ ทุกฺขํ โทมนสฺสํ ปฏิสํเวทยติ.\csromanspacer{} อทินฺนาทานา ปฏิวิรตสฺส เอวํ ตํ ภยํ เวรํ วูปสนฺตํ โหติ. (2)}

\prose{ยํ คหปติ กาเมสุมิจฺฉาจารี กาเมสุมิจฺฉาจารปจฺจยา ทิฏฺฐธมฺมิกมฺปิ ภยํ เวรํ ปสวติ, สมฺปรายิกมฺปิ ภยํ เวรํ ปสวติ, เจตสิกมฺปิ ทุกฺขํ โทมนสฺสํ ปฏิสํเวทยติ.\csromanspacer{} กาเมสุมิจฺฉาจารา ปฏิวิรตสฺส เอวํ ตํ ภยํ เวรํ วูปสนฺตํ โหติ. (3)}

\prose{ยํ คหปติ มุสาวาที มุสาวาทปจฺจยา ทิฏฺฐธมฺมิกมฺปิ ภยํ เวรํ ปสวติ, สมฺปรายิกมฺปิ ภยํ เวรํ ปสวติ, เจตสิกมฺปิ ทุกฺขํ โทมนสฺสํ ปฏิสํเวทยติ.\csromanspacer{} มุสาวาทา ปฏิวิรตสฺส เอวํ ตํ ภยํ เวรํ วูปสนฺตํ โหติ. (4)}

\prose{ยํ คหปติ สุราเมรยมชฺชปมาทฏฺฐายี สุราเมรยมชฺชปมาทฏฺฐานปจฺจยา ทิฏฺฐธมฺมิกมฺปิ ภยํ เวรํ ปสวติ, สมฺปรายิกมฺปิ ภยํ เวรํ ปสวติ, เจตสิกมฺปิ ทุกฺขํ โทมนสฺสํ ปฏิสํเวทยติ.\csromanspacer{} สุราเมรยมชฺชปมาทฏฺฐานา ปฏิวิรตสฺส เอวํ ตํ ภยํ เวรํ วูปสนฺตํ โหติ.\csromanspacer{} อิมานิ ปญฺจ ภยานิ เวรานิ วูปสนฺตานิ โหนฺติ. (5)}

\prose{กตเมหิ จตูหิ โสตาปตฺติยงฺเคหิ สมนฺนาคโต โหติ.\csromanspacer{} อิธ คหปติ อริยสาวโก พุทฺเธ อเวจฺจปฺปสาเทน สมนฺนาคโต โหติ “อิติปิ โส ภควา อรหํ สมฺมาสมฺพุทฺโธ วิชฺชาจรณสมฺปนฺโน สุคโต โลกวิทู อนุตฺตโร ปุริสทมฺมสารถิ สตฺถา เทวมนุสฺสานํ พุทฺโธ ภควา”ติ. (1)}

\csromanmatikaanchor{mtk.44}
\csromantocmarkref{subsection}{6. ทุติย-อวิชฺชาปจฺจยสุตฺต}{mtk.44}
\bookmark[dest={mtk.44},level=2]{6. ทุติย-อวิชฺชาปจฺจยสุตฺต}
\prose{ธมฺเม อเวจฺจปฺปสาเทน สมนฺนาคโต โหติ “สฺวากฺขาโต ภควตา ธมฺโม สนฺทิฏฺฐิโก อกาลิโก เอหิปสฺสิโก โอปเนยฺยิโก ปจฺจตฺตํ เวทิตพฺโพ วิญฺญูหี”ติ. (2)}

\csromanheader{2}{1. นิทานสํยุตฺต}
\prose{\csromanfolio{299}สํเฆ อเวจฺจปฺปสาเทน สมนฺนาคโต โหติ “สุปฺปฏิปนฺโน ภควโต สาวกสํโฆ อุชุปฺปฏิปนฺโน ภควโต สาวกสํโฆ ญายปฺปฏิปนฺโน ภควโต สาวกสํโฆ สามีจิปฺปฏิปนฺโน ภควโต สาวกสํโฆ, ยทิทํ จตฺตาริ ปุริสยุคานิ อฏฺฐ ปุริสปุคฺคลา, เอส ภควโต สาวกสํโฆ อาหุเนยฺโย ปาหุเนยฺโย ทกฺขิเณยฺโย อญฺชลิกรณีโย อนุตฺตรํ ปุญฺญกฺเขตฺตํ โลกสฺสา”ติ. (3)}

\prose{อริยกนฺเตหิ สีเลหิ สมนฺนาคโต โหติ อขณฺเฑหิ อจฺฉิทฺเทหิ อสพเลหิ อกมฺมาเสหิ ภุชิสฺเสหิ วิญฺญุปฺปสตฺเถหิ อปรามฏฺเฐหิ สมาธิสํวตฺตนิเกหิ.\csromanspacer{} อิเมหิ จตูหิ โสตาปตฺติยงฺเคหิ สมนฺนาคโต โหติ. (4)}

\prose{กตโม จสฺส อริโย ญาโย ปญฺญาย สุทิฏฺโฐ โหติ สุปฺปฏิวิทฺโธ.\csromanspacer{} อิธ คหปติ อริยสาวโก ปฏิจฺจสมุปฺปาทญฺเญว สาธุกํ โยนิโส มนสิ กโรติ “อิติ อิมสฺมิํ สติ อิทํ โหติ, อิมสฺมิํ อสติ อิทํ น โหติ, อิมสฺสุปฺปาทา อิทํ อุปฺปชฺชติ, อิมสฺส นิโรธา อิทํ นิรุชฺฌติ, ยทิทํ อวิชฺชาปจฺจยา สงฺขารา, สงฺขารปจฺจยา วิญฺญาณํ ฯเปฯ เอวเมตสฺส เกวลสฺส ทุกฺขกฺขนฺธสฺส สมุทโย โหติ.\csromanspacer{} อวิชฺชาย เตฺวว อเสสวิราคนิโรธา สงฺขารนิโรโธ, สงฺขารนิโรธา วิญฺญาณนิโรโธ ฯเปฯ เอวเมตสฺส เกวลสฺส ทุกฺขกฺขนฺธสฺส นิโรโธ โหตี”ติ.\csromanspacer{} อยมสฺส อริโย ญาโย ปญฺญาย สุทิฏฺโฐ โหติ สุปฺปฏิวิทฺโธ.}

\prose{ยโต โข คหปติ อริยสาวกสฺส อิมานิ ปญฺจ ภยานิ เวรานิ วูปสนฺตานิ โหนฺติ.\csromanspacer{} อิเมหิ จ จตูหิ โสตาปตฺติยงฺเคหิ สมนฺนาคโต โหติ.\csromanspacer{} อยญฺจสฺส อริโย ญาโย ปญฺญาย สุทิฏฺโฐ โหติ สุปฺปฏิวิทฺโธ.\csromanspacer{} โส อากงฺขมาโน อตฺตนาว อตฺตานํ พฺยากเรยฺย “ขีณนิรโยมฺหิ ขีณติรจฺฉานโยนิ ขีณเปตฺติวิสโย ขีณาปายทุคฺคติวินิปาโต, โสตาปนฺโนหมสฺมิ อวินิปาตธมฺโม นิยโต สมฺโพธิปรายโน”ติ.\csromanspacer{}\csromanspacer{}ปฐมํ.}

\csromanheader{2}{2. ทุติยปญฺจเวรภยสุตฺต}
\csromanmatikaanchor{mtk.45}
\csromantocmarkref{subsection}{7. นตุมฺหสุตฺต}{mtk.45}
\bookmark[dest={mtk.45},level=2]{7. นตุมฺหสุตฺต}
\proseitem{42}{สาวตฺถิยํ วิหรติ.\csromanspacer{} ยโต โข ภิกฺขเว อริยสาวกสฺส ปญฺจ ภยานิ เวรานิ วูปสนฺตานิ โหนฺติ.\csromanspacer{} จตูหิ จ โสตาปตฺติยงฺเคหิ สมนฺนาคโต โหติ, อริโย จสฺส ญาโย ปญฺญาย สุทิฏฺโฐ}

\vspace{\tipitakaparskip}
\csromancenter{นิทานวคฺคสํยุตฺตปาฬิ โหติ สุปฺปฏิวิทฺโธ, โส อากงฺขมาโน อตฺตนาว อตฺตานํ พฺยากเรยฺย “ขีณนิรโยมฺหิ ฯเปฯ อวินิปาตธมฺโม นิยโต สมฺโพธิปรายโน”ติ.}
\prose{\csromanfolio{300}กตมานิ ปญฺจ ภยานิ เวรานิ วูปสนฺตานิ โหนฺติ.\csromanspacer{} ยํ ภิกฺขเว ปาณาติปาตี ฯเปฯ.\csromanspacer{} ยํ ภิกฺขเว อทินฺนาทายี.\csromanspacer{} ยํ ภิกฺขเว กาเมสุมิจฺฉาจารี.\csromanspacer{} ยํ ภิกฺขเว มุสาวาที.\csromanspacer{} ยํ ภิกฺขเว สุราเมรยมชฺชปมาทฏฺฐายี ฯเปฯ.\csromanspacer{} อิมานิ ปญฺจ ภยานิ เวรานิ วูปสนฺตานิ โหนฺติ.}

\prose{กตเมหิ จตูหิ โสตาปตฺติยงฺเคหิ สมนฺนาคโต โหติ.\csromanspacer{} อิธ ภิกฺขเว อริยสาวโก พุทฺเธ ฯเปฯ ธมฺเม.\csromanspacer{} สํเฆ.\csromanspacer{} อริยกนฺเตหิ สีเลหิ สมนฺนาคโต โหติ.\csromanspacer{} อิเมหิ จตูหิ โสตาปตฺติยงฺเคหิ สมนฺนาคโต โหติ.}

\csromanmatikaanchor{mtk.46}
\csromantocmarkref{subsection}{8. เจตนาสุตฺต}{mtk.46}
\bookmark[dest={mtk.46},level=2]{8. เจตนาสุตฺต}
\prose{กตโม จสฺส อริโย ญาโย ปญฺญาย สุทิฏฺโฐ โหติ สุปฺปฏิวิทฺโธ.\csromanspacer{} อิธ ภิกฺขเว อริยสาวโก ปฏิจฺจสมุปฺปาทญฺเญว สาธุกํ โยนิโส มนสิ กโรติ ฯเปฯ อยมสฺส อริโย ญาโย ปญฺญาย สุทิฏฺโฐ โหติ สุปฺปฏิวิทฺโธ.}

\prose{ยโต โข ภิกฺขเว อริยสาวกสฺส อิมานิ ปญฺจ ภยานิ เวรานิ วูปสนฺตานิ โหนฺติ.\csromanspacer{} อิเมหิ จ จตูหิ โสตาปตฺติยงฺเคหิ สมนฺนาคโต โหติ, อยญฺจสฺส อริโย ญาโย ปญฺญาย สุทิฏฺโฐ โหติ สุปฺปฏิวิทฺโธ, โส อากงฺขมาโน อตฺตนาว อตฺตานํ พฺยากเรยฺย “ขีณนิรโยมฺหิ ขีณติรจฺฉานโยนิ ขีณเปตฺติวิสโย ขีณาปายทุคฺคติวินิปาโต, โสตาปนฺโนหมสฺมิ อวินิปาตธมฺโม นิยโต สมฺโพธิปรายโน”ติ.\csromanspacer{}\csromanspacer{}ทุติยํ.}

\csromanheader{2}{3. ทุกฺขสุตฺต}
\proseitem{43}{สาวตฺถิยํ วิหรติ.\csromanspacer{} ทุกฺขสฺส ภิกฺขเว สมุทยญฺจ อตฺถงฺคมญฺจ เทเสสฺสามิ, ตํ สุณาถ สาธุกํ มนสิ กโรถ, ภาสิสฺสามีติ. “เอวํ ภนฺเต”ติ โข เต ภิกฺขู ภควโต ปจฺจสฺโสสุํ.\csromanspacer{} ภควา เอตทโวจ\csromandash{}}

\csromanmatikaanchor{mtk.47}
\csromantocmarkref{subsection}{9. ทุติยเจตนาสุตฺต}{mtk.47}
\bookmark[dest={mtk.47},level=2]{9. ทุติยเจตนาสุตฺต}
\prose{กตโม จ ภิกฺขเว ทุกฺขสฺส สมุทโย.\csromanspacer{} จกฺขุญฺจ ปฏิจฺจ รูเป จ อุปฺปชฺชติ จกฺขุวิญฺญาณํ, ติณฺณํ สงฺคติ ผสฺโส, ผสฺสปจฺจยา เวทนา, เวทนาปจฺจยา ตณฺหา.\csromanspacer{} อยํ โข ภิกฺขเว ทุกฺขสฺส สมุทโย.}

\csromanheader{2}{1. นิทานสํยุตฺต}
\prose{\csromanfolio{301}โสตญฺจ ปฏิจฺจ สทฺเท จ อุปฺปชฺชติ โสตวิญฺญาณํ ฯเปฯ.\csromanspacer{} ฆานญฺจ ปฏิจฺจ คนฺเธ จ ฯเปฯ.\csromanspacer{} ชิวฺหญฺจ ปฏิจฺจ รเส จ ฯเปฯ.\csromanspacer{} กายญฺจ ปฏิจฺจ โผฏฺฐพฺเพ จ ฯเปฯ.\csromanspacer{} มนญฺจ ปฏิจฺจ ธมฺเม จ อุปฺปชฺชติ มโนวิญฺญาณํ, ติณฺณํ สงฺคติ ผสฺโส, ผสฺสปจฺจยา เวทนา, เวทนาปจฺจยา ตณฺหา.\csromanspacer{} อยํ โข ภิกฺขเว ทุกฺขสฺส สมุทโย.}

\prose{กตโม จ ภิกฺขเว ทุกฺขสฺส อตฺถงฺคโม.\csromanspacer{} จกฺขุญฺจ ปฏิจฺจ รูเป จ อุปฺปชฺชติ จกฺขุวิญฺญาณํ, ติณฺณํ สงฺคติ ผสฺโส, ผสฺสปจฺจยา เวทนา, เวทนาปจฺจยา ตณฺหา.\csromanspacer{} ตสฺสาเยว ตณฺหาย อเสสวิราคนิโรธา อุปาทานนิโรโธ, อุปาทานนิโรธา ภวนิโรโธ, ภวนิโรธา ชาตินิโรโธ, ชาตินิโรธา ชรามรณํ โสกปริเทวทุกฺขโทมนสฺสุปายาสา นิรุชฺฌนฺติ.\csromanspacer{} เอวเมตสฺส เกวลสฺส ทุกฺขกฺขนฺธสฺส นิโรโธ โหติ.\csromanspacer{} อยํ โข ภิกฺขเว ทุกฺขสฺส อตฺถงฺคโม.}

\csromanmatikaanchor{mtk.48}
\csromantocmarkref{subsection}{10. ตติยเจตนาสุตฺต}{mtk.48}
\bookmark[dest={mtk.48},level=2]{10. ตติยเจตนาสุตฺต}
\csromantocmarkref{section}{อุทฺทานคาถา}{mtk.49}
\bookmark[dest={mtk.49},level=1]{อุทฺทานคาถา}
\prose{โสตญฺจ ปฏิจฺจ สทฺเท จ อุปฺปชฺชติ โสตวิญฺญาณํ ฯเปฯ.\csromanspacer{} ฆานญฺจ ปฏิจฺจ คนฺเธ จ ฯเปฯ.\csromanspacer{} ชิวฺหญฺจ ปฏิจฺจ รเส จ ฯเปฯ.\csromanspacer{} กายญฺจ ปฏิจฺจ โผฏฺฐพฺเพ จ ฯเปฯ.\csromanspacer{} มนญฺจ ปฏิจฺจ ธมฺเม จ อุปฺปชฺชติ มโนวิญฺญาณํ, ติณฺณํ สงฺคติ ผสฺโส, ผสฺสปจฺจยา เวทนา, เวทนาปจฺจยา ตณฺหา.\csromanspacer{} ตสฺสาเยว ตณฺหาย อเสสวิราคนิโรธา อุปาทานนิโรโธ, อุปาทานนิโรธา ภวนิโรโธ, ภวนิโรธา ชาตินิโรโธ, ชาตินิโรธา ชรามรณํ โสกปริเทวทุกฺขโทมนสฺสุปายาสา นิรุชฺฌนฺติ.\csromanspacer{} เอวเมตสฺส เกวลสฺส ทุกฺขกฺขนฺธสฺส นิโรโธ โหติ.\csromanspacer{} อยํ โข ภิกฺขเว ทุกฺขสฺส อตฺถงฺคโมติ.\csromanspacer{}\csromanspacer{}ตติยํ.}

\csromanheader{2}{4. โลกสุตฺต}
\proseitem{44}{สาวตฺถิยํ วิหรติ.\csromanspacer{} โลกสฺส ภิกฺขเว สมุทยญฺจ อตฺถงฺคมญฺจ เทเสสฺสามิ, ตํ สุณาถ สาธุกํ มนสิ กโรถ, ภาสิสฺสามีติ. “เอวํ ภนฺเต”ติ โข เต ภิกฺขู ภควโต ปจฺจสฺโสสุํ.\csromanspacer{} ภควา เอตทโวจ\csromandash{}}

\prose{กตโม จ ภิกฺขเว โลกสฺส สมุทโย.\csromanspacer{} จกฺขุญฺจ ปฏิจฺจ รูเป จ อุปฺปชฺชติ จกฺขุวิญฺญาณํ, ติณฺณํ สงฺคติ ผสฺโส, ผสฺสปจฺจยา เวทนา, เวทนาปจฺจยา ตณฺหา, ตณฺหาปจฺจยา อุปาทานํ, อุปาทานปจฺจยา ภโว, ภวปจฺจยา ชาติ, ชาติปจฺจยา ชรามรณํ โสกปริเทวทุกฺขโทมนสฺสุปายาสา สมฺภวนฺติ.\csromanspacer{} อยํ โข ภิกฺขเว โลกสฺส สมุทโย.}

\gambhira{นิทานวคฺคสํยุตฺตปาฬิ}
\prose{\csromanfolio{302}โสตญฺจ ปฏิจฺจ สทฺเท จ ฯเปฯ.\csromanspacer{} ฆานญฺจ ปฏิจฺจ คนฺเธ จ.\csromanspacer{} ชิวฺหญฺจ ปฏิจฺจ รเส จ.\csromanspacer{} กายญฺจ ปฏิจฺจ โผฏฺฐพฺเพ จ.\csromanspacer{} มนญฺจ ปฏิจฺจ ธมฺเม จ อุปฺปชฺชติ มโนวิญฺญาณํ, ติณฺณํ สงฺคติ ผสฺโส, ผสฺสปจฺจยา เวทนา ฯเปฯ ชาติปจฺจยา ชรามรณํ โสกปริเทวทุกฺขโทมนสฺสุปายาสา สมฺภวนฺติ.\csromanspacer{} อยํ โข ภิกฺขเว โลกสฺส สมุทโย.}

\prose{กตโม จ ภิกฺขเว โลกสฺส อตฺถงฺคโม.\csromanspacer{} จกฺขุญฺจ ปฏิจฺจ รูเป จ อุปฺปชฺชติ จกฺขุวิญฺญาณํ, ติณฺณํ สงฺคติ ผสฺโส, ผสฺสปจฺจยา เวทนา, เวทนาปจฺจยา ตณฺหา.\csromanspacer{} ตสฺสาเยว ตณฺหาย อเสสวิราคนิโรธา อุปาทานนิโรโธ, อุปาทานนิโรธา ภวนิโรโธ ฯเปฯ เอวเมตสฺส เกวลสฺส ทุกฺขกฺขนฺธสฺส นิโรโธ โหติ.\csromanspacer{} อยํ โข ภิกฺขเว โลกสฺส อตฺถงฺคโม.}

\prose{โสตญฺจ ปฏิจฺจ สทฺเท จ ฯเปฯ.\csromanspacer{} ฆานญฺจ ปฏิจฺจ คนฺเธ จ.\csromanspacer{} ชิวฺหญฺจ ปฏิจฺจ รเส จ.\csromanspacer{} กายญฺจ ปฏิจฺจ โผฏฺฐพฺเพ จ.\csromanspacer{} มนญฺจ ปฏิจฺจ ธมฺเม จ อุปฺปชฺชติ มโนวิญฺญาณํ, ติณฺณํ สงฺคติ ผสฺโส, ผสฺสปจฺจยา เวทนา, เวทนาปจฺจยา ตณฺหา.\csromanspacer{} ตสฺสาเยว ตณฺหาย อเสสวิราคนิโรธา อุปาทานนิโรโธ, อุปาทานนิโรธา ภวนิโรโธ ฯเปฯ เอวเมตสฺส เกวลสฺส ทุกฺขกฺขนฺธสฺส นิโรโธ โหติ.\csromanspacer{} อยํ โข ภิกฺขเว โลกสฺส อตฺถงฺคโมติ.\csromanspacer{}\csromanspacer{}จตุตฺถํ.}

\csromanmatikaanchor{mtk.50}
\csromantocmarknopage{section}{5. คหปติวคฺค}{mtk.50}
\bookmark[dest={mtk.50},level=1]{5. คหปติวคฺค}
\csromanheader{2}{5. ญาติกสุตฺต}
\csromanmatikaanchor{mtk.51}
\csromantocmarkref{subsection}{1. ปญฺจเวรภยสุตฺต}{mtk.51}
\bookmark[dest={mtk.51},level=2]{1. ปญฺจเวรภยสุตฺต}
\proseitem{45}{เอวํ เม สุตํ\csromandash{}เอกํ สมยํ ภควา ญาติเก วิหรติ คิญฺชกาวสเถ.\csromanspacer{} อถ โข ภควา รโหคโต ปฏิสลฺลาโน อิมํ ธมฺมปริยายํ อภาสิ\csromandash{}}

\prose{จกฺขุญฺจ ปฏิจฺจ รูเป จ อุปฺปชฺชติ จกฺขุวิญฺญาณํ, ติณฺณํ สงฺคติ ผสฺโส, ผสฺสปจฺจยา เวทนา, เวทนาปจฺจยา ตณฺหา, ตณฺหาปจฺจยา อุปาทานํ ฯเปฯ เอวเมตสฺส เกวลสฺส ทุกฺขกฺขนฺธสฺส สมุทโย โหติ.}

\prose{โสตญฺจ ปฏิจฺจ สทฺเท จ ฯเปฯ.\csromanspacer{} ฆานญฺจ ปฏิจฺจ คนฺเธ จ.\csromanspacer{} ชิวฺหญฺจ ปฏิจฺจ รเส จ.\csromanspacer{} กายญฺจ ปฏิจฺจ โผฏฺฐพฺเพ จ.\csromanspacer{} มนญฺจ ปฏิจฺจ ธมฺเม จ อุปฺปชฺชติ มโนวิญฺญาณํ, ติณฺณํ สงฺคติ ผสฺโส, ผสฺสปจฺจยา เวทนา, เวทนาปจฺจยา ตณฺหา, ตณฺหาปจฺจยา อุปาทานํ ฯเปฯ เอวเมตสฺส เกวลสฺส ทุกฺขกฺขนฺธสฺส สมุทโย โหติ.}

\csromanheader{2}{1. นิทานสํยุตฺต}
\prose{\csromanfolio{303}จกฺขุญฺจ ปฏิจฺจ รูเป จ อุปฺปชฺชติ จกฺขุวิญฺญาณํ, ติณฺณํ สงฺคติ ผสฺโส, ผสฺสปจฺจยา เวทนา, เวทนาปจฺจยา ตณฺหา.\csromanspacer{} ตสฺสาเยว ตณฺหาย อเสสวิราคนิโรธา อุปาทานนิโรโธ, อุปาทานนิโรธา ภวนิโรโธ ฯเปฯ เอวเมตสฺส เกวลสฺส ทุกฺขกฺขนฺธสฺส นิโรโธ โหติ.}

\prose{โสตญฺจ ปฏิจฺจ สทฺเท จ ฯเปฯ.\csromanspacer{} มนญฺจ ปฏิจฺจ ธมฺเม จ อุปฺปชฺชติ มโนวิญฺญาณํ, ติณฺณํ สงฺคติ ผสฺโส, ผสฺสปจฺจยา เวทนา, เวทนาปจฺจยา ตณฺหา.\csromanspacer{} ตสฺสาเยว ตณฺหาย อเสสวิราคนิโรธา อุปาทานนิโรโธ, อุปาทานนิโรธา ภวนิโรโธ ฯเปฯ เอวเมตสฺส เกวลสฺส ทุกฺขกฺขนฺธสฺส นิโรโธ โหตีติ.}

\prose{เตน โข ปน สมเยน อญฺญตโร ภิกฺขุ ภควโต อุปสฺสุติ\footnote{อุปสฺสุติํ (สี, อิ)} ฐิโต โหติ.\csromanspacer{} อทฺทสา โข ภควา ตํ ภิกฺขุํ อุปสฺสุติ ฐิตํ, ทิสฺวาน ตํ ภิกฺขุํ เอตทโวจ “อสฺโสสิ โน ตฺวํ ภิกฺขุ อิมํ ธมฺมปริยายนฺ”ติ.\csromanspacer{} เอวํ ภนฺเตติ.\csromanspacer{} อุคฺคณฺหาหิ ตฺวํ ภิกฺขุ อิมํ ธมฺมปริยายํ, ปริยาปุณาหิ ตฺวํ ภิกฺขุ อิมํ ธมฺมปริยายํ, ธาเรหิ ตฺวํ ภิกฺขุ อิมํ ธมฺมปริยายํ, อตฺถสํหิโต อยํ\footnote{อตฺถสํหิโตยํ (สี, สฺยา, กํ), อตฺถสํหิตายํ (อิ, ก)} ภิกฺขุ ธมฺมปริยาโย อาทิพฺรหฺมจริยโกติ.\csromanspacer{}\csromanspacer{}ปญฺจมํ.}

\csromanheader{2}{6. อญฺญตรพฺราหฺมณสุตฺต}
\proseitem{46}{สาวตฺถิยํ วิหรติ.\csromanspacer{} อถ โข อญฺญตโร พฺราหฺมโณ เยน ภควา เตนุปสงฺกมิ, อุปสงฺกมิตฺวา ภควตา สทฺธิํ สมฺโมทิ, สมฺโมทนียํ กถํ สารณียํ วีติสาเรตฺวา เอกมนฺตํ นิสีทิ, เอกมนฺตํ นิสินฺโน โข โส พฺราหฺมโณ ภควนฺตํ เอตทโวจ\csromandash{}}

\prose{กิํ นุ โข โภ โคตม โส กโรติ, โส ปฏิสํเวทยตีติ. “โส กโรติ, โส ปฏิสํเวทยตี”ติ โข พฺราหฺมณ อยเมโก อนฺโต.}

\prose{กิํ ปน โภ โคตม อญฺโญ กโรติ, อญฺโญ ปฏิสํเวทยตีติ. “อญฺโญ กโรติ, อญฺโญ ปฏิสํเวทยตี”ติ โข พฺราหฺมณ อยํ ทุติโย อนฺโต.\csromanspacer{} เอเต เต พฺราหฺมณ อุโภ อนฺเต อนุปคมฺม มชฺเฌน ตถาคโต ธมฺมํ เทเสติ\csromandash{} อวิชฺชาปจฺจยา สงฺขารา, สงฺขารปจฺจยา วิญฺญาณํ ฯเปฯ เอวเมตสฺส เกวลสฺส ทุกฺขกฺขนฺธสฺส สมุทโย โหติ.}

\vspace{\tipitakaparskip}
\csromancenter{นิทานวคฺคสํยุตฺตปาฬิ อวิชฺชาย เตฺวว อเสสวิราคนิโรธา สงฺขารนิโรโธ, สงฺขารนิโรธา ฯเปฯ เอวเมตสฺส เกวลสฺส ทุกฺขกฺขนฺธสฺส นิโรโธ โหตีติ.}
\prose{\csromanfolio{304}เอวํ วุตฺเต โส พฺราหฺมโณ ภควนฺตํ เอตทโวจ “อภิกฺกนฺตํ โภ โคตม, อภิกฺกนฺตํ โภ โคตม ฯเปฯ อุปาสกํ มํ ภวํ โคตโม ธาเรตุ อชฺชตคฺเค ปาณุเปตํ สรณํ คตนฺ”ติ.\csromanspacer{}\csromanspacer{}ฉฏฺฐํ.}

\csromanheader{2}{7. ชาณุสฺโสณิสุตฺต}
\proseitem{47}{สาวตฺถิยํ วิหรติ.\csromanspacer{} อถ โข ชาณุสฺโสณิ พฺราหฺมโณ เยน ภควา เตนุปสงฺกมิ, อุปสงฺกมิตฺวา ภควตา สทฺธิํ ฯเปฯ เอกมนฺตํ นิสินฺโน โข ชาณุสฺโสณิ พฺราหฺมโณ ภควนฺตํ เอตทโวจ\csromandash{}}

\prose{กิํ นุ โข โภ โคตม สพฺพมตฺถีติ. “สพฺพมตฺถี”ติ โข พฺราหฺมณ อยเมโก อนฺโต.}

\csromanmatikaanchor{mtk.52}
\csromantocmarkref{subsection}{2. ทุติยปญฺจเวรภยสุตฺต}{mtk.52}
\bookmark[dest={mtk.52},level=2]{2. ทุติยปญฺจเวรภยสุตฺต}
\prose{กิํ ปน โภ โคตม สพฺพํ นตฺถีติ. “สพฺพํ นตฺถี”ติ โข พฺราหฺมณ อยํ ทุติโย อนฺโต.\csromanspacer{} เอเต เต พฺราหฺมณ อุโภ อนฺเต อนุปคมฺม มชฺเฌน ตถาคโต ธมฺมํ เทเสติ\csromandash{}อวิชฺชาปจฺจยา สงฺขารา, สงฺขารปจฺจยา วิญฺญาณํ ฯเปฯ เอวเมตสฺส เกวลสฺส ทุกฺขกฺขนฺธสฺส สมุทโย โหติ.\csromanspacer{} อวิชฺชาย เตฺวว อเสสวิราคนิโรธา สงฺขารนิโรโธ, สงฺขารนิโรธา วิญฺญาณนิโรโธ ฯเปฯ เอวเมตสฺส เกวลสฺส ทุกฺขกฺขนฺธสฺส นิโรโธ โหตีติ.}

\prose{เอวํ ปุตฺเต ชาณุสฺโสณิ พฺราหฺมโณ ภควนฺตํ เอตทโวจ “อภิกฺกนฺตํ โภ โคตม ฯเปฯ ปาณุเปตํ สรณํ คตนฺ”ติ.\csromanspacer{}\csromanspacer{}สตฺตมํ.}

\csromanheader{2}{8. โลกายติกสุตฺต}
\proseitem{48}{สาวตฺถิยํ วิหรติ.\csromanspacer{} อถ โข โลกายติโก พฺราหฺมโณ เยน ภควา ฯเปฯ เอกมนฺตํ นิสินฺโน โข โลกายติโก พฺราหฺมโณ ภควนฺตํ เอตทโวจ\csromandash{}}

\prose{กิํ นุ โข โภ โคตม สพฺพมตฺถีติ. “สพฺพมตฺถี”ติ โข พฺราหฺมณ เชฏฺฐเมตํ โลกายตํ.}

\csromanheader{2}{1. นิทานสํยุตฺต}
\prose{\csromanfolio{305}กิํ ปน โภ โคตม สพฺพํ นตฺถีติ. “สพฺพํ นตฺถี”ติ โข พฺราหฺมณ ทุติยเมตํ โลกายตํ.}

\csromanmatikaanchor{mtk.53}
\csromantocmarkref{subsection}{3. ทุกฺขสุตฺต}{mtk.53}
\bookmark[dest={mtk.53},level=2]{3. ทุกฺขสุตฺต}
\prose{กิํ นุ โข โภ โคตม สพฺพเมกตฺตนฺติ. “สพฺพเมกตฺตนฺ”ติ โข พฺราหฺมณ ตติยเมตํ โลกายตํ.}

\prose{กิํ ปน โภ โคตม สพฺพํ ปุถุตฺตนฺติ. “สพฺพํ ปุถุตฺตนฺ”ติ โข พฺราหฺมณ จตุตฺถเมตํ โลกายตํ.\csromanspacer{} เอเต เต พฺราหฺมณ อุโภ อนฺเต อนุปคมฺม มชฺเฌน ตถาคโต ธมฺมํ เทเสติ\csromandash{}อวิชฺชาปจฺจยา สงฺขารา, สงฺขารปจฺจยา วิญฺญาณํ ฯเปฯ เอวเมตสฺส เกวลสฺส ทุกฺขกฺขนฺธสฺส สมุทโย โหติ.\csromanspacer{} อวิชฺชาย เตฺวว อเสสวิราคนิโรธา สงฺขารนิโรโธ, สงฺขารนิโรธา วิญฺญาณนิโรโธ ฯเปฯ เอวเมตสฺส เกวลสฺส ทุกฺขกฺขนฺธสฺส นิโรโธ โหตีติ.}

\prose{เอวํ วุตฺเต โลกายติโก พฺราหฺมโณ ภควนฺตํ เอตทโวจ “อภิกฺกนฺตํ โภ โคตม ฯเปฯ อชฺชตคฺเค ปาณุเปตํ สรณํ คตนฺ”ติ.\csromanspacer{}\csromanspacer{}อฏฺฐมํ.}

\csromanheader{2}{9. อริยสาวกสุตฺต}
\proseitem{49}{สาวตฺถิยํ วิหรติ.\csromanspacer{} น ภิกฺขเว สุตวโต อริยสาวกสฺส เอวํ โหติ “กิํ นุ โข\csromandash{}กิสฺมิํ สติ กิํ โหติ, กิสฺสุปฺปาทา กิํ อุปฺปชฺชติ, (กิสฺมิํ สติ สงฺขารา โหนฺติ, กิสฺมิํ สติ วิญฺญาณํ โหติ)\footnote{( ) เอตฺถนฺตเร ปาฐา เกสุจิ โปตฺถเกสุ น ทิสฺสนฺตีติ สี-อิ-โปตฺถเกสุ ทสฺสิตา. ตถา สติ อนนฺตรสุตฺตฏีกาย สเมติ.} กิสฺมิํ สติ นามรูปํ โหติ, กิสฺมิํ สติ สฬายตนํ โหติ, กิสฺมิํ สติ ผสฺโส โหติ, กิสฺมิํ สติ เวทนา โหติ, กิสฺมิํ สติ ตณฺหา โหติ, กิสฺมิํ สติ อุปาทานํ โหติ, กิสฺมิํ สติ ภโว โหติ, กิสฺมิํ สติ ชาติ โหติ, กิสฺมิํ สติ ชรามรณํ โหตี”ติ.}

\prose{อถ โข ภิกฺขเว สุตวโต อริยสาวกสฺส อปรปฺปจฺจยา ญาณเมเวตฺถ โหติ “อิมสฺมิํ สติ อิทํ โหติ, อิมสฺสุปฺปาทา อิทํ อุปฺปชฺชติ, (อวิชฺชาย สติ สงฺขารา โหนฺติ, สงฺขาเรสุ สติ วิญฺญาณํ โหติ,)1 วิญฺญาเณ สติ นามรูปํ โหติ, นามรูเป สติ สฬายตนํ}

\vspace{\tipitakaparskip}
\csromancenter{นิทานวคฺคสํยุตฺตปาฬิ โหติ, สฬายตเน สติ ผสฺโส โหติ, ผสฺเส สติ เวทนา โหติ, เวทนาย สติ ตณฺหา โหติ, ตณฺหาย สติ อุปาทานํ โหติ, อุปาทาเน สติ ภโว โหติ, ภเว สติ ชาติ โหติ, ชาติยา สติ ชรามรณํ โหตี”ติ.\csromanspacer{} โส เอวํ ปชานาติ “เอวมยํ โลโก สมุทยตี”ติ.}
\csromanmatikaanchor{mtk.54}
\csromantocmarkref{subsection}{4. โลกสุตฺต}{mtk.54}
\bookmark[dest={mtk.54},level=2]{4. โลกสุตฺต}
\prose{\csromanfolio{306}น ภิกฺขเว สุตวโต อริยสาวกสฺส เอวํ โหติ “กิํ นุ โข\csromandash{}กิสฺมิํ อสติ กิํ น โหติ, กิสฺส นิโรธา กิํ นิรุชฺฌติ, (กิสฺมิํ อสติ สงฺขารา น โหนฺติ, กิสฺมิํ อสติ วิญฺญาณํ น โหติ,)\footnote{( ) เอตฺถนฺตเร ปาฐาปิ ตตฺถ ตเถว ทสฺสิตา.} กิสฺมิํ อสติ นามรูปํ น โหติ, กิสฺมิํ อสติ สฬายตนํ น โหติ, กิสฺมิํ อสติ ผสฺโส น โหติ, กิสฺมิํ อสติ เวทนา น โหติ, กิสฺมิํ อสติ ตณฺหา น โหติ, กิสฺมิํ อสติ อุปาทานํ น โหติ, กิสฺมิํ อสติ ภโว น โหติ, กิสฺมิํ อสติ ชาติ น โหติ, กิสฺมิํ อสติ ชรามรณํ น โหตี”ติ.}

\prose{อถ โข ภิกฺขเว สุตวโต อริยสาวกสฺส อปรปฺปจฺจยา ญาณเมเวตฺถ โหติ “อิมสฺมิํ อสติ อิทํ น โหติ, อิมสฺส นิโรธา อิทํ นิรุชฺฌติ, (อวิชฺชาย อสติ สงฺขารา น โหนฺติ, สงฺขาเรสุ อสติ วิญฺญาณํ น โหติ,)1 วิญฺญาเณ อสติ นามรูปํ น โหติ, นามรูเป อสติ สฬายตนํ น โหติ ฯเปฯ ภโว น โหติ.\csromanspacer{} ชาติ น โหติ.\csromanspacer{} ชาติยา อสติ ชรามรณํ น โหตี”ติ.\csromanspacer{} โส เอวํ ปชานาติ “เอวมยํ โลโก นิรุชฺฌตี”ติ.}

\prose{ยโต โข ภิกฺขเว อริยสาวโก เอวํ โลกสฺส สมุทยญฺจ อตฺถงฺคมญฺจ ยถาภูตํ ปชานาติ.\csromanspacer{} อยํ วุจฺจติ ภิกฺขเว อริยสาวโก ทิฏฺฐิสมฺปนฺโน อิติปิ ฯเปฯ อมตทฺวารํ อาหจฺจ ติฏฺฐติ อิติปีติ.\csromanspacer{}\csromanspacer{}นวมํ.}

\csromanheader{2}{10. ทุติย-อริยสาวกสุตฺต}
\proseitem{50}{สาวตฺถิยํ วิหรติ.\csromanspacer{} น ภิกฺขเว สุตวโต อริยสาวกสฺส เอวํ โหติ “กิํ นุ โข\csromandash{}กิสฺมิํ สติ กิํ โหติ, กิสฺสุปฺปาทา กิํ อุปฺปชฺชติ, กิสฺมิํ สติ สงฺขารา โหนฺติ, กิสฺมิํ สติ วิญฺญาณํ โหติ, กิสฺมิํ สติ นามรูปํ โหติ, กิสฺมิํ สติ สฬายตนํ โหติ, กิสฺมิํ}

\vspace{\tipitakaparskip}
\csromancenter{นิทานสํยุตฺต สติ ผสฺโส โหติ, กิสฺมิํ สติ เวทนา โหติ, กิสฺมิํ สติ ตณฺหา โหติ, กิสฺมิํ สติ อุปาทานํ โหติ, กิสฺมิํ สติ ภโว โหติ, กิสฺมิํ สติ ชาติ โหติ, กิสฺมิํ สติ ชรามรณํ โหตี”ติ.}
\prose{\csromanfolio{307}อถ โข ภิกฺขเว สุตวโต อริยสาวกสฺส อปรปฺปจฺจยา ญาณเมเวตฺถ โหติ “อิมสฺมิํ สติ อิทํ โหติ, อิมสฺสุปฺปาทา อิทํ อุปฺปชฺชติ, อวิชฺชาย สติ สงฺขารา โหนฺติ, สงฺขาเรสุ สติ วิญฺญาณํ โหติ, วิญฺญาเณ สติ นามรูปํ โหติ, นามรูเป สติ สฬายตนํ โหติ, สฬายตเน สติ ผสฺโส โหติ, ผสฺเส สติ เวทนา โหติ, เวทนาย สติ ตณฺหา โหติ, ตณฺหาย สติ อุปาทานํ โหติ, อุปาทาเน สติ ภโว โหติ, ภเว สติ ชาติ โหติ, ชาติยา สติ ชรามรณํ โหตี”ติ.\csromanspacer{} โส เอวํ ปชานาติ “เอวมยํ โลโก สมุทยตี”ติ.}

\csromanmatikaanchor{mtk.55}
\csromantocmarkref{subsection}{5. ญาติกสุตฺต}{mtk.55}
\bookmark[dest={mtk.55},level=2]{5. ญาติกสุตฺต}
\prose{น ภิกฺขเว สุตวโต อริยสาวกสฺส เอวํ โหติ “กิํ นุ โข\csromandash{}กิสฺมิํ อสติ กิํ น โหติ, กิสฺส นิโรธา กิํ นิรุชฺฌติ, กิสฺมิํ อสติ สงฺขารา น โหนฺติ, กิสฺมิํ อสติ วิญฺญาณํ น โหติ, กิสฺมิํ อสติ นามรูปํ น โหติ, กิสฺมิํ อสติ สฬายตนํ น โหติ, กิสฺมิํ อสติ ผสฺโส น โหติ, กิสฺมิํ อสติ เวทนา น โหติ, กิสฺมิํ อสติ ตณฺหา น โหติ ฯเปฯ อุปาทานํ.\csromanspacer{} ภโว.\csromanspacer{} ชาติ.\csromanspacer{} กิสฺมิํ อสติ ชรามรณํ น โหตี”ติ.}

\prose{อถ โข ภิกฺขเว สุตวโต อริยสาวกสฺส อปรปฺปจฺจยา ญาณเมเวตฺถ โหติ “อิมสฺมิํ อสติ อิทํ น โหติ, อิมสฺส นิโรธา อิทํ นิรุชฺฌติ, อวิชฺชาย อสติ สงฺขารา น โหนฺติ, สงฺขาเรสุ อสติ วิญฺญาณํ น โหติ, วิญฺญาเณ อสติ นามรูปํ น โหติ, นามรูเป อสติ สฬายตนํ น โหติ ฯเปฯ ชาติยา อสติ ชรามรณํ น โหตี”ติ.\csromanspacer{} โส เอวํ ปชานาติ “เอวมยํ โลโก นิรุชฺฌตี”ติ.}

\prose{ยโต โข ภิกฺขเว อริยสาวโก เอวํ โลกสฺส สมุทยญฺจ อตฺถงฺคมญฺจ ยถาภูตํ ปชานาติ.\csromanspacer{} อยํ วุจฺจติ ภิกฺขเว อริยสาวโก ทิฏฺฐิสมฺปนฺโน อิติปิ, ทสฺสนสมฺปนฺโน อิติปิ, อาคโต อิมํ สทฺธมฺมํ อิติปิ, ปสฺสติ อิมํ สทฺธมฺมํ อิติปิ, เสกฺเขน ญาเณน สมนฺนาคโต อิติปิ,}

\vspace{\tipitakaparskip}
\csromancenter{นิทานวคฺคสํยุตฺตปาฬิ เสกฺขาย วิชฺชาย สมนฺนาคโต อิติปิ, ธมฺมโสตํ สมาปนฺโน อิติปิ, อริโย นิพฺเพธิกปญฺโญ อิติปิ, อมตทฺวารํ อาหจฺจ ติฏฺฐติ อิติปีติ.\csromanspacer{}\csromanspacer{}ทสมํ.}
\csromancenter{คหปติวคฺโค ปญฺจโม.}
\titlehead{\textbf{ตสฺสุทฺทานํ}}
\csromanmatikaanchor{mtk.56}
\csromanmatikaanchor{mtk.61}
\csromantocmarkref{subsection}{6. อญฺญตรพฺราหฺมณสุตฺต}{mtk.56}
\bookmark[dest={mtk.56},level=2]{6. อญฺญตรพฺราหฺมณสุตฺต}
\csromangathameasure{เทฺว ปญฺจเวรภยา วุตฺตา, ทุกฺขํ โลโก จ ญาติกํ. \\ อญฺญตรํ ชาณุสฺโสณิ จ, โลกายติเกน อฏฺฐมํ. \\ เทฺว อริยสาวกา วุตฺตา, วคฺโค เตน ปวุจฺจตีติ.}
\csromangathagroup{\csromangathastanza{\csromanfolio{308}เทฺว ปญฺจเวรภยา วุตฺตา, ทุกฺขํ โลโก จ ญาติกํ. \\ อญฺญตรํ ชาณุสฺโสณิ จ, โลกายติเกน อฏฺฐมํ. \\ เทฺว อริยสาวกา วุตฺตา, วคฺโค เตน ปวุจฺจตีติ.}}

\csromanheader{1}{6. ทุกฺขวคฺค}
\csromanheader{2}{1. ปริวีมํสนสุตฺต}
\proseitem{51}{เอวํ เม สุตํ\csromandash{}เอกํ สมยํ ภควา สาวตฺถิยํ วิหรติ เชตวเน อนาถปิณฺฑิกสฺส อาราเม.\csromanspacer{} ตตฺร โข ภควา ภิกฺขู อามนฺเตสิ “ภิกฺขโว”ติ. “ภทนฺเต”ติ เต ภิกฺขู ภควโต ปจฺจสฺโสสุํ.\csromanspacer{} ภควา เอตทโวจ\csromandash{}}

\prose{กิตฺตาวตา นุ โข ภิกฺขเว ภิกฺขุ ปริวีมํสมาโน ปริวีมํเสยฺย สพฺพโส สมฺมา ทุกฺขกฺขยายาติ.\csromanspacer{} ภควํมูลกา โน ภนฺเต ธมฺมา ภควํเนตฺติกา ภควํปฏิสรณา, สาธุ วต ภนฺเต ภควนฺตํเยว ปฏิภาตุ เอตสฺส ภาสิตสฺส อตฺโถ, ภควโต สุตฺวา ภิกฺขู ธาเรสฺสนฺตีติ.\csromanspacer{} เตน หิ ภิกฺขเว สุณาถ สาธุกํ มนสิ กโรถ, ภาสิสฺสามีติ. “เอวํ ภนฺเต”ติ โข เต ภิกฺขู ภควโต ปจฺจสฺโสสุํ.\csromanspacer{} ภควา เอตทโวจ\csromandash{}}

\prose{อิธ ภิกฺขเว ภิกฺขุ ปริวีมํสมาโน ปริวีมํสติ “ยํ โข อิทํ อเนกวิธํ นานปฺปการกํ ทุกฺขํ โลเก อุปฺปชฺชติ ชรามรณํ, อิทํ นุ โข ทุกฺขํ กิํนิทานํ กิํสมุทยํ กิํชาติกํ กิํปภวํ, กิสฺมิํ สติ ชรามรณํ โหติ, กิสฺมิํ อสติ ชรามรณํ น โหตี”ติ.\csromanspacer{} โส ปริวีมํสมาโน เอวํ ปชานาติ “ยํ โข อิทํ อเนกวิธํ นานปฺปการกํ ทุกฺขํ โลเก}

\vspace{\tipitakaparskip}
\csromancenter{นิทานสํยุตฺต อุปฺปชฺชติ ชรามรณํ, อิทํ โข ทุกฺขํ ชาตินิทานํ ชาติสมุทยํ ชาติชาติกํ ชาติปฺปภวํ, ชาติยา สติ ชรามรณํ โหติ, ชาติยา อสติ ชรามรณํ น โหตี”ติ.}
\prose{\csromanfolio{309}โส ชรามรณญฺจ ปชานาติ, ชรามรณสมุทยญฺจ ปชานาติ, ชรามรณนิโรธญฺจ ปชานาติ, ยา จ ชรามรณนิโรธสารุปฺปคามินี ปฏิปทา, ตญฺจ ปชานาติ, ตถา ปฏิปนฺโน จ โหติ อนุธมฺมจารี.\csromanspacer{} อยํ วุจฺจติ ภิกฺขเว ภิกฺขุ สพฺพโส สมฺมา ทุกฺขกฺขยาย ปฏิปนฺโน ชรามรณนิโรธาย.}

\csromanmatikaanchor{mtk.57}
\csromantocmarkref{subsection}{7. ชานุสฺโสณิสุตฺต}{mtk.57}
\bookmark[dest={mtk.57},level=2]{7. ชานุสฺโสณิสุตฺต}
\prose{อถาปรํ ปริวีมํสมาโน ปริวีมํสติ “ชาติ ปนายํ กิํนิทานา กิํสมุทยา กิํชาติกา กิํปภวา, กิสฺมิํ สติ ชาติ โหติ, กิสฺมิํ อสติ ชาติ น โหตี”ติ.\csromanspacer{} โส ปริวีมํสมาโน เอวํ ปชานาติ “ชาติ ภวนิทานา ภวสมุทยา ภวชาติกา ภวปฺปภวา, ภเว สติ ชาติ โหติ, ภเว อสติ ชาติ น โหตี”ติ.}

\prose{โส ชาติญฺจ ปชานาติ, ชาติสมุทยญฺจ ปชานาติ, ชาตินิโรธญฺจ ปชานาติ, ยา จ ชาตินิโรธสารุปฺปคามินี ปฏิปทา, ตญฺจ ปชานาติ, ตถา ปฏิปนฺโน จ โหติ อนุธมฺมจารี.\csromanspacer{} อยํ วุจฺจติ ภิกฺขเว ภิกฺขุ สพฺพโส สมฺมา ทุกฺขกฺขยาย ปฏิปนฺโน ชาตินิโรธาย.}

\prose{อถาปรํ ปริวีมํสมาโน ปริวีมํสติ “ภโว ปนายํ กิํนิทาโน ฯเปฯ อุปาทานํ ปนิทํ กิํนิทานํ.\csromanspacer{} ตณฺหา ปนายํ กิํนิทานา.\csromanspacer{} เวทนา.\csromanspacer{} ผสฺโส.\csromanspacer{} สฬายตนํ ปนิทํ กิํนิทานํ.\csromanspacer{} นามรูปํ ปนิทํ.\csromanspacer{} วิญฺญาณํ ปนิทํ.\csromanspacer{} สงฺขารา ปนิเม กิํนิทานา กิํสมุทยา กิํชาติกา กิํปภวา, กิสฺมิํ สติ สงฺขารา โหนฺติ, กิสฺมิํ อสติ สงฺขารา น โหนฺตี”ติ.\csromanspacer{} โส ปริวีมํสมาโน เอวํ ปชานาติ “สงฺขารา อวิชฺชานิทานา อวิชฺชาสมุทยา อวิชฺชาชาติกา อวิชฺชาปภวา, อวิชฺชาย สติ สงฺขารา โหนฺติ, อวิชฺชาย อสติ สงฺขารา น โหนฺตี”ติ.}

\prose{โส สงฺขาเร จ ปชานาติ, สงฺขารสมุทยญฺจ ปชานาติ, สงฺขารนิโรธญฺจ ปชานาติ, ยา จ สงฺขารนิโรธสารุปฺปคามินี ปฏิปทา, ตญฺจ ปชานาติ.\csromanspacer{} ตถา ปฏิปนฺโน จ โหติ อนุธมฺมจารี.\csromanspacer{} อยํ วุจฺจติ ภิกฺขเว ภิกฺขุ สพฺพโส สมฺมา ทุกฺขกฺขยาย ปฏิปนฺโน สงฺขารนิโรธาย.}

\gambhira{นิทานวคฺคสํยุตฺตปาฬิ}
\csromanmatikaanchor{mtk.58}
\csromantocmarkref{subsection}{8. โลกายติกสุตฺต}{mtk.58}
\bookmark[dest={mtk.58},level=2]{8. โลกายติกสุตฺต}
\prose{\csromanfolio{310}อวิชฺชาคโต ยํ ภิกฺขเว ปุริสปุคฺคโล ปุญฺญํ เจ สงฺขารํ อภิสงฺขโรติ, ปุญฺญูปคํ โหติ วิญฺญาณํ.\csromanspacer{} อปุญฺญํ เจ สงฺขารํ อภิสงฺขโรติ, อปุญฺญูปคํ โหติ วิญฺญาณํ.\csromanspacer{} อาเนญฺชํ เจ สงฺขารํ อภิสงฺขโรติ, อาเนญฺชูปคํ โหติ วิญฺญาณํ.\csromanspacer{} ยโต โข ภิกฺขเว ภิกฺขุโน อวิชฺชา ปหีนา โหติ, วิชฺชา อุปฺปนฺนา.\csromanspacer{} โส อวิชฺชาวิราคา วิชฺชุปฺปาทา เนว ปุญฺญาภิสงฺขารํ อภิสงฺขโรติ, น อปุญฺญาภิสงฺขารํ อภิสงฺขโรติ, น อาเนญฺชาภิสงฺขารํ อภิสงฺขโรติ.\csromanspacer{} อนภิสงฺขโรนฺโต อนภิสญฺเจตยนฺโต น กิญฺจิ โลเก อุปาทิยติ, อนุปาทิยํ น ปริตสฺสติ, อปริตสฺสํ ปจฺจตฺตญฺเญว ปรินิพฺพายติ, “ขีณา ชาติ, วุสิตํ พฺรหฺมจริยํ, กตํ กรณียํ, นาปรํ อิตฺถตฺตายา”ติ ปชานาติ.}

\prose{โส สุขํ เจ เวทนํ เวทยติ, “สา อนิจฺจา”ติ ปชานาติ, “อนชฺโฌสิตา”ติ ปชานาติ, “อนภินนฺทิตา”ติ ปชานาติ.\csromanspacer{} ทุกฺขํ เจ เวทนํ เวทยติ, “สา อนิจฺจา”ติ ปชานาติ, “อนชฺโฌสิตา”ติ ปชานาติ, “อนภินนฺทิตา”ติ ปชานาติ.\csromanspacer{} อทุกฺขมสุขํ เจ เวทนํ เวทยติ, “สา อนิจฺจา”ติ ปชานาติ, “อนชฺโฌสิตา”ติ ปชานาติ, “อนภินนฺทิตา”ติ ปชานาติ.\csromanspacer{} โส สุขํ เจ เวทนํ เวทยติ, วิสํยุตฺโต นํ\footnote{ตํ เวทนํ (สี, อิ), เวทนํ (ก)} เวทยติ.\csromanspacer{} ทุกฺขํ เจ เวทนํ เวทยติ, วิสํยุตฺโต นํ เวทยติ.\csromanspacer{} อทุกฺขมสุขํ เจ เวทนํ เวทยติ, วิสํยุตฺโต นํ เวทยติ.}

\prose{โส กายปริยนฺติกํ เวทนํ เวทยมาโน “กายปริยนฺติกํ เวทนํ เวทยามี”ติ ปชานาติ, ชีวิตปริยนฺติกํ เวทนํ เวทยมาโน “ชีวิตปริยนฺติกํ เวทนํ เวทยามี”ติ ปชานาติ. “กายสฺส เภทา อุทฺธํ ชีวิตปริยาทานา อิเธว สพฺพเวทยิตานิ อนภินนฺทิตานิ สีตีภวิสฺสนฺติ สรีรานิ อวสิสฺสนฺตี”ติ ปชานาติ.}

\prose{เสยฺยถาปิ ภิกฺขเว ปุริโส กุมฺภการปากา อุณฺหํ กุมฺภํ อุทฺธริตฺวา สเม ภูมิภาเค ปฏิสิสฺเสยฺย\footnote{ปฏิวิเสยฺย (สี), ปติฏฺฐเปยฺย (สฺยา, กํ, อิ), ปฏิเสเวยฺย (ฏีกา)}.\csromanspacer{} ตตฺร ยายํ อุสฺมา, สา ตตฺเถว วูปสเมยฺย, กปลฺลานิ อวสิสฺเสยฺยุํ.\csromanspacer{} เอวเมว โข ภิกฺขเว ภิกฺขุ กายปริยนฺติกํ เวทนํ เวทยมาโน “กายปริยนฺติกํ เวทนํ เวทยามี”ติ ปชานาติ, ชีวิตปริยนฺติกํ เวทนํ เวทยมาโน}

\vspace{\tipitakaparskip}
\csromancenter{นิทานสํยุตฺต “ชีวิตปริยนฺติกํ เวทนํ เวทยามี”ติ ปชานาติ. “กายสฺส เภทา อุทฺธํ ชีวิตปริยาทานา อิเธว สพฺพเวทยิตานิ อนภินนฺทิตานิ สีตีภวิสฺสนฺติ, สรีรานิ อวสิสฺสนฺตี”ติ ปชานาติ.}
\prose{\csromanfolio{311}ตํ กิํ มญฺญถ ภิกฺขเว, อปิ นุ โข ขีณาสโว ภิกฺขุ ปุญฺญาภิสงฺขารํ วา อภิสงฺขเรยฺย, อปุญฺญาภิสงฺขารํ วา อภิสงฺขเรยฺย, อาเนญฺชาภิสงฺขารํ วา อภิสงฺขเรยฺยาติ.\csromanspacer{} โน เหตํ ภนฺเต.\csromanspacer{} สพฺพโส วา ปน สงฺขาเรสุ อสติ สงฺขารนิโรธา อปิ นุ โข วิญฺญาณํ ปญฺญาเยถาติ.\csromanspacer{} โน เหตํ ภนฺเต.\csromanspacer{} สพฺพโส วา ปน วิญฺญาเณ อสติ วิญฺญาณนิโรธา อปิ นุ โข นามรูปํ ปญฺญาเยถาติ.\csromanspacer{} โน เหตํ ภนฺเต.\csromanspacer{} สพฺพโส วา ปน นามรูเป อสติ นามรูปนิโรธา อปิ นุ โข สฬายตนํ ปญฺญาเยถาติ.\csromanspacer{} โน เหตํ ภนฺเต.\csromanspacer{} สพฺพโส วา ปน สฬายตเน อสติ สฬายตนนิโรธา อปิ นุ โข ผสฺโส ปญฺญาเยถาติ.\csromanspacer{} โน เหตํ ภนฺเต.\csromanspacer{} สพฺพโส วา ปน ผสฺเส อสติ ผสฺสนิโรธา อปิ นุ โข เวทนา ปญฺญาเยถาติ.\csromanspacer{} โน เหตํ ภนฺเต.\csromanspacer{} สพฺพโส วา ปน เวทนาย อสติ เวทนานิโรธา อปิ นุ โข ตณฺหา ปญฺญาเยถาติ.\csromanspacer{} โน เหตํ ภนฺเต.\csromanspacer{} สพฺพโส วา ปน ตณฺหาย อสติ ตณฺหานิโรธา อปิ นุ โข อุปาทานํ ปญฺญาเยถาติ.\csromanspacer{} โน เหตํ ภนฺเต.\csromanspacer{} สพฺพโส วา ปน อุปาทาเน อสติ อุปาทานนิโรธา อปิ นุ โข ภโว ปญฺญาเยถาติ.\csromanspacer{} โน เหตํ ภนฺเต.\csromanspacer{} สพฺพโส วา ปน ภเว อสติ ภวนิโรธา อปิ นุ โข ชาติ ปญฺญาเยถาติ.\csromanspacer{} โน เหตํ ภนฺเต.\csromanspacer{} สพฺพโส วา ปน ชาติยา อสติ ชาตินิโรธา อปิ นุ โข ชรามรณํ ปญฺญาเยถาติ.\csromanspacer{} โน เหตํ ภนฺเต.}

\prose{สาธุ สาธุ ภิกฺขเว, เอวเมตํ ภิกฺขเว, เนตํ อญฺญถา, สทฺทหถ เม ตํ ภิกฺขเว อธิมุจฺจถ, นิกฺกงฺขา เอตฺถ โหถ นิพฺพิจิกิจฺฉา, เอเสวนฺโต ทุกฺขสฺสาติ.\csromanspacer{}\csromanspacer{}ปฐมํ.}

\csromanheader{2}{2. อุปาทานสุตฺต}
\csromanmatikaanchor{mtk.59}
\csromantocmarkref{subsection}{9. อริยสาวกสุตฺต}{mtk.59}
\bookmark[dest={mtk.59},level=2]{9. อริยสาวกสุตฺต}
\proseitem{52}{สาวตฺถิยํ วิหรติ.\csromanspacer{} อุปาทานิเยสุ ภิกฺขเว ธมฺเมสุ อสฺสาทานุปสฺสิโน วิหรโต ตณฺหา ปวฑฺฒติ, ตณฺหาปจฺจยา อุปาทานํ, อุปาทานปจฺจยา ภโว, ภวปจฺจยา ชาติ, ชาติปจฺจยา ชรามรณํ โสกปริเทวทุกฺขโทมนสฺสุปายาสา สมฺภวนฺติ, เอวเมตสฺส เกวลสฺส ทุกฺขกฺขนฺธสฺส สมุทโย โหติ.}

\vspace{\tipitakaparskip}
\csromancenter{นิทานวคฺคสํยุตฺตปาฬิ เสยฺยถาปิ ภิกฺขเว ทสนฺนํ วา กฏฺฐวาหานํ วีสาย วา กฏฺฐวาหานํ ติํสาย วา กฏฺฐวาหานํ จตฺตารีสาย วา กฏฺฐวาหานํ มหา- อคฺคิกฺขนฺโธ ชเลยฺย, ตตฺร ปุริโส กาเลน กาลํ สุกฺขานิ เจว ติณานิ ปกฺขิเปยฺย, สุกฺขานิ จ โคมยานิ ปกฺขิเปยฺย, สุกฺขานิ จ กฏฺฐานิ ปกฺขิเปยฺย.\csromanspacer{} เอวํ หิ โส ภิกฺขเว มหา-อคฺคิกฺขนฺโธ ตทาหาโร ตทุปาทาโน จิรํ ทีฆมทฺธานํ ชเลยฺย.\csromanspacer{} เอวเมว โข ภิกฺขเว อุปาทานิเยสุ ธมฺเมสุ อสฺสาทานุปสฺสิโน วิหรโต ตณฺหา ปวฑฺฒติ, ตณฺหาปจฺจยา อุปาทานํ ฯเปฯ เอวเมตสฺส เกวลสฺส ทุกฺขกฺขนฺธสฺส สมุทโย โหติ.}
\prose{\csromanfolio{312}อุปาทานิเยสุ ภิกฺขเว ธมฺเมสุ อาทีนวานุปสฺสิโน วิหรโต ตณฺหา นิรุชฺฌติ, ตณฺหานิโรธา อุปาทานนิโรโธ, อุปาทานนิโรธา ภวนิโรโธ, ภวนิโรธา ชาตินิโรโธ, ชาตินิโรธา ชรามรณํ โสกปริเทวทุกฺขโทมนสฺสุปายาสา นิรุชฺฌนฺติ, เอวเมตสฺส เกวลสฺส ทุกฺขกฺขนฺธสฺส นิโรโธ โหติ.}

\prose{เสยฺยถาปิ ภิกฺขเว ทสนฺนํ วา กฏฺฐวาหานํ วีสาย วา ติํสาย วา จตฺตารีสาย วา กฏฺฐวาหานํ มหา-อคฺคิกฺขนฺโธ ชเลยฺย, ตตฺร ปุริโส น กาเลน กาลํ สุกฺขานิ เจว ติณานิ ปกฺขิเปยฺย, น สุกฺขานิ จ โคมยานิ ปกฺขิเปยฺย, น สุกฺขานิ จ กฏฺฐานิ ปกฺขิเปยฺย.\csromanspacer{} เอวํ หิ โส ภิกฺขเว มหา-อคฺคิกฺขนฺโธ ปุริมสฺส จ อุปาทานสฺส ปริยาทานา อญฺญสฺส จ อนุปหารา\footnote{อนุปาหารา (อิ)} อนาหาโร นิพฺพาเยยฺย.\csromanspacer{} เอวเมว โข ภิกฺขเว อุปาทานิเยสุ ธมฺเมสุ อาทีนวานุปสฺสิโน วิหรโต ตณฺหา นิรุชฺฌติ, ตณฺหานิโรธา อุปาทานนิโรโธ ฯเปฯ เอวเมตสฺส เกวลสฺส ทุกฺขกฺขนฺธสฺส นิโรโธ โหตีติ.\csromanspacer{}\csromanspacer{}ทุติยํ.}

\csromanheader{2}{3. สํโยชนสุตฺต}
\proseitem{53}{สาวตฺถิยํ วิหรติ.\csromanspacer{} สํโยชนิเยสุ ภิกฺขเว ธมฺเมสุ อสฺสาทานุปสฺสิโน วิหรโต ตณฺหา ปวฑฺฒติ, ตณฺหาปจฺจยา อุปาทานํ, อุปาทานปจฺจยา ภโว, ภวปจฺจยา ชาติ, ชาติปจฺจยา ชรามรณํ โสกปริเทวทุกฺขโทมนสฺสุปายาสา สมฺภวนฺติ, เอวเมตสฺส เกวลสฺส ทุกฺขกฺขนฺธสฺส สมุทโย โหติ.}

\csromanheader{2}{1. นิทานสํยุตฺต}
\csromanmatikaanchor{mtk.60}
\csromantocmarkref{subsection}{10. ทุติย-อริยสาวกสุตฺต}{mtk.60}
\bookmark[dest={mtk.60},level=2]{10. ทุติย-อริยสาวกสุตฺต}
\csromantocmarkref{section}{อุทฺทานคาถา}{mtk.61}
\bookmark[dest={mtk.61},level=1]{อุทฺทานคาถา}
\prose{\csromanfolio{313}เสยฺยถาปิ ภิกฺขเว เตลญฺจ ปฏิจฺจ วฏฺฏิญฺจ ปฏิจฺจ เตลปฺปทีโป ฌาเยยฺย, ตตฺร ปุริโส กาเลน กาลํ เตลํ อาสิญฺเจยฺย วฏฺฏิํ อุปสํหเรยฺย.\csromanspacer{} เอวํ หิ โส ภิกฺขเว เตลปฺปทีโป ตทาหาโร ตทุปาทาโน จิรํ ทีฆมทฺธานํ ชเลยฺย.\csromanspacer{} เอวเมว โข ภิกฺขเว สํโยชนิเยสุ ธมฺเมสุ อสฺสาทานุปสฺสิโน วิหรโต ตณฺหา ปวฑฺฒติ, ตณฺหาปจฺจยา อุปาทานํ, อุปาทานปจฺจยา ภโว, ภวปจฺจยา ชาติ, ชาติปจฺจยา ชรามรณํ โสกปริเทวทุกฺขโทมนสฺสุปายาสา สมฺภวนฺติ, เอวเมตสฺส เกวลสฺส ทุกฺขกฺขนฺธสฺส สมุทโย โหติ.}

\prose{สํโยชนิเยสุ ภิกฺขเว ธมฺเมสุ อาทีนวานุปสฺสิโน วิหรโต ตณฺหา นิรุชฺฌติ, ตณฺหานิโรธา อุปาทานนิโรโธ, อุปาทานนิโรธา ภวนิโรโธ, ภวนิโรธา ชาตินิโรโธ, ชาตินิโรธา ชรามรณํ โสกปริเทวทุกฺขโทมนสฺสุปายาสา นิรุชฺฌนฺติ, เอวเมตสฺส เกวลสฺส ทุกฺขกฺขนฺธสฺส นิโรโธ โหติ.}

\prose{เสยฺยถาปิ ภิกฺขเว เตลญฺจ ปฏิจฺจ วฏฺฏิญฺจ ปฏิจฺจ เตลปฺปทีโป ฌาเยยฺย.\csromanspacer{} ตตฺร ปุริโส น กาเลน กาลํ เตลํ อาสิญฺเจยฺย, น วฏฺฏิํ อุปสํหเรยฺย.\csromanspacer{} เอวํ หิ โส ภิกฺขเว เตลปฺปทีโป ปุริมสฺส จ อุปาทานสฺส ปริยาทานา อญฺญสฺส จ อนุปหารา อนาหาโร นิพฺพาเยยฺย.\csromanspacer{} เอวเมว โข ภิกฺขเว สํโยชนิเยสุ ธมฺเมสุ อาทีนวานุปสฺสิโน วิหรโต ตณฺหา นิรุชฺฌติ, ตณฺหานิโรธา อุปาทานนิโรโธ ฯเปฯ เอวเมตสฺส เกวลสฺส ทุกฺขกฺขนฺธสฺส นิโรโธ โหตีติ.\csromanspacer{}\csromanspacer{}ตติยํ.}

\csromanheader{2}{4. ทุติยสํโยชนสุตฺต}
\proseitem{54}{สาวตฺถิยํ วิหรติ.\csromanspacer{} เสยฺยถาปิ ภิกฺขเว เตลญฺจ ปฏิจฺจ วฏฺฏิญฺจ ปฏิจฺจ เตลปฺปทีโป ฌาเยยฺย.\csromanspacer{} ตตฺร ปุริโส กาเลน กาลํ เตลํ อาสิญฺเจยฺย, วฏฺฏิํ อุปสํหเรยฺย.\csromanspacer{} เอวํ หิ โส ภิกฺขเว เตลปฺปทีโป ตทาหาโร ตทุปาทาโน จิรํ ทีฆมทฺธานํ ชเลยฺย.\csromanspacer{} เอวเมว โข ภิกฺขเว สํโยชนิเยสุ ธมฺเมสุ อสฺสาทานุปสฺสิโน วิหรโต ตณฺหา ปวฑฺฒติ, ตณฺหาปจฺจยา อุปาทานํ ฯเปฯ เอวเมตสฺส เกวลสฺส ทุกฺขกฺขนฺธสฺส สมุทโย โหติ.}

\prose{เสยฺยถาปิ ภิกฺขเว เตลญฺจ ปฏิจฺจ วฏฺฏิญฺจ ปฏิจฺจ เตลปฺปทีโป ฌาเยยฺย.\csromanspacer{} ตตฺร ปุริโส น กาเลน กาลํ เตลํ อาสิญฺเจยฺย, น วฏฺฏิํ}

\vspace{\tipitakaparskip}
\csromancenter{นิทานวคฺคสํยุตฺตปาฬิ อุปสํหเรยฺย.\csromanspacer{} เอวํ หิ โส ภิกฺขเว เตลปฺปทีโป ปุริมสฺส จ อุปาทานสฺส ปริยาทานา อญฺญสฺส จ อนุปหารา อนาหาโร นิพฺพาเยยฺย.\csromanspacer{} เอวเมว โข ภิกฺขเว สํโยชนิเยสุ ธมฺเมสุ อาทีนวานุปสฺสิโน วิหรโต ตณฺหา นิรุชฺฌติ, ตณฺหานิโรธา อุปาทานนิโรโธ ฯเปฯ เอวเมตสฺส เกวลสฺส ทุกฺขกฺขนฺธสฺส นิโรโธ โหตีติ.\csromanspacer{}\csromanspacer{}จตุตฺถํ.}
\csromanheader{2}{5. มหารุกฺขสุตฺต}
\proseitem{55}{\csromanfolio{314}สาวตฺถิยํ วิหรติ.\csromanspacer{} อุปาทานิเยสุ ภิกฺขเว ธมฺเมสุ อสฺสาทานุปสฺสิโน วิหรโต ตณฺหา ปวฑฺฒติ, ตณฺหาปจฺจยา อุปาทานํ, อุปาทานปจฺจยา ภโว ฯเปฯ เอวเมตสฺส เกวลสฺส ทุกฺขกฺขนฺธสฺส สมุทโย โหติ.}

\prose{เสยฺยถาปิ ภิกฺขเว มหารุกฺโข.\csromanspacer{} ตสฺส ยานิ เจว มูลานิ อโธคมานิ ยานิ จ ติริยงฺคมานิ, สพฺพานิ ตานิ อุทฺธํ โอชํ อภิหรนฺติ.\csromanspacer{} เอวํ หิ โส ภิกฺขเว มหารุกฺโข ตทาหาโร ตทุปาทาโน จิรํ ทีฆมทฺธานํ ติฏฺเฐยฺย.\csromanspacer{} เอวเมว โข ภิกฺขเว อุปาทานิเยสุ ธมฺเมสุ อสฺสาทานุปสฺสิโน วิหรโต ตณฺหา ปวฑฺฒติ, ตณฺหาปจฺจยา อุปาทานํ ฯเปฯ เอวเมตสฺส เกวลสฺส ทุกฺขกฺขนฺธสฺส สมุทโย โหติ.}

\prose{อุปาทานิเยสุ ภิกฺขเว ธมฺเมสุ อาทีนวานุปสฺสิโน วิหรโต ตณฺหา นิรุชฺฌติ, ตณฺหานิโรธา อุปาทานนิโรโธ, อุปาทานนิโรธา ภวนิโรโธ ฯเปฯ เอวเมตสฺส เกวลสฺส ทุกฺขกฺขนฺธสฺส นิโรโธ โหติ.}

\csromanmatikaanchor{mtk.62}
\csromantocmarknopage{section}{6. ทุกฺขวคฺค}{mtk.62}
\bookmark[dest={mtk.62},level=1]{6. ทุกฺขวคฺค}
\prose{เสยฺยถาปิ ภิกฺขเว มหารุกฺโข.\csromanspacer{} อถ ปุริโส อาคจฺเฉยฺย กุทฺทาลปิฏกํ\footnote{กุทาลปิฏกํ (อญฺญตฺถ)} อาทาย, โส ตํ รุกฺขํ มูเล ฉินฺเทยฺย, มูลํ ฉินฺทิตฺวา ปลิขเณยฺย\footnote{ปลิํ ขเณยฺย (อิ, ก)}, ปลิขณิตฺวา มูลานิ อุทฺธเรยฺย อนฺตมโส อุสีรนาฬิมตฺตานิปิ.\csromanspacer{} โส ตํ รุกฺขํ ขณฺฑาขณฺฑิกํ ฉินฺเทยฺย, ขณฺฑาขณฺฑิกํ ฉินฺทิตฺวา ผาเลยฺย, ผาเลตฺวา สกลิกํ สกลิกํ กเรยฺย, สกลิกํ สกลิกํ กริตฺวา วาตาตเป วิโสเสยฺย, วาตาตเป วิโสเสตฺวา อคฺคินา ฑเหยฺย, อคฺคินา ฑเหตฺวา มสิํ กเรยฺย, มสิํ กริตฺวา มหาวาเต วา โอผุเณยฺย\footnote{โอปุเนยฺย (สี, อิ), โอผุเนยฺย (สฺยา, กํ, ก)}, นทิยา วา สีฆโสตาย ปวาเหยฺย.\csromanspacer{} เอวํ หิ โส ภิกฺขเว มหารุกฺโข}

\csromanmatikaanchor{mtk.63}
\csromantocmarkref{subsection}{1. ปริวีมํสนสุตฺต}{mtk.63}
\bookmark[dest={mtk.63},level=2]{1. ปริวีมํสนสุตฺต}
\vspace{\tipitakaparskip}
\csromancenter{นิทานสํยุตฺต อุจฺฉินฺนมูโล อสฺส, ตาลาวตฺถุกโต อนภาวํกโต\footnote{อนภาวกโต (สี), อนภาวงฺคโต (สฺยา, กํ)} อายติํ อนุปฺปาทธมฺโม.\csromanspacer{} เอวเมว โข ภิกฺขเว อุปาทานิเยสุ ธมฺเมสุ อาทีนวานุปสฺสิโน วิหรโต ตณฺหา นิรุชฺฌติ, ตณฺหานิโรธา อุปาทานนิโรโธ, อุปาทานนิโรธา ภวนิโรโธ ฯเปฯ เอวเมตสฺส เกวลสฺส ทุกฺขกฺขนฺธสฺส นิโรโธ โหตีติ.\csromanspacer{}\csromanspacer{}ปญฺจมํ.}
\csromanheader{2}{6. ทุติยมหารุกฺขสุตฺต}
\proseitem{56}{\csromanfolio{315}สาวตฺถิยํ วิหรติ.\csromanspacer{} เสยฺยถาปิ ภิกฺขเว มหารุกฺโข.\csromanspacer{} ตสฺส ยานิ เจว มูลานิ อโธคมานิ ยานิ จ ติริยงฺคมานิ, สพฺพานิ ตานิ อุทฺธํ โอชํ อภิหรนฺติ, เอวํ หิ โส ภิกฺขเว มหารุกฺโข ตทาหาโร ตทุปาทาโน จิรํ ทีฆมทฺธานํ ติฏฺเฐยฺย.\csromanspacer{} เอวเมว โข ภิกฺขเว อุปาทานิเยสุ ธมฺเมสุ อสฺสาทานุปสฺสิโน วิหรโต ตณฺหา ปวฑฺฒติ, ตณฺหาปจฺจยา อุปาทานํ ฯเปฯ เอวเมตสฺส เกวลสฺส ทุกฺขกฺขนฺธสฺส สมุทโย โหติ.}

\prose{เสยฺยถาปิ ภิกฺขเว มหารุกฺโข.\csromanspacer{} อถ ปุริโส อาคจฺเฉยฺย กุทฺทาลปิฏกํ อาทาย, โส ตํ รุกฺขํ มูเล ฉินฺเทยฺย, มูเล เฉตฺวา ปลิขเณยฺย, ปลิขณิตฺวา มูลานิ อุทฺธเรยฺย ฯเปฯ นทิยา วา สีฆโสตาย ปวาเหยฺย.\csromanspacer{} เอวํ หิ โส ภิกฺขเว มหารุกฺโข อุจฺฉินฺนมูโล อสฺส, ตาลาวตฺถุกโต อนภาวํกโต อายติํ อนุปฺปาทธมฺโม.\csromanspacer{} เอวเมว โข ภิกฺขเว อุปาทานิเยสุ ธมฺเมสุ อาทีนวานุปสฺสิโน วิหรโต ตณฺหา นิรุชฺฌติ, ตณฺหานิโรธา อุปาทานนิโรโธ ฯเปฯ เอวเมตสฺส เกวลสฺส ทุกฺขกฺขนฺธสฺส นิโรโธ โหตีติ.\csromanspacer{}\csromanspacer{}ฉฏฺฐํ.}

\csromanheader{2}{7. ตรุณรุกฺขสุตฺต}
\proseitem{57}{สาวตฺถิยํ วิหรติ.\csromanspacer{} สํโยชนิเยสุ ภิกฺขเว ธมฺเมสุ อสฺสาทานุปสฺสิโน วิหรโต ตณฺหา ปวฑฺฒติ, ตณฺหาปจฺจยา อุปาทานํ ฯเปฯ เอวเมตสฺส เกวลสฺส ทุกฺขกฺขนฺธสฺส สมุทโย โหติ.}

\prose{เสยฺยถาปิ ภิกฺขเว ตรุโณ รุกฺโข.\csromanspacer{} ตสฺส ปุริโส กาเลน กาลํ มูลานิ ปลิมชฺเชยฺย\footnote{ปลิสนฺเนยฺย (สี), ปลิสชฺเชยฺย (สฺยา, กํ, อิ), ปลิปฏฺเฐยฺย (ก), ปลิสนฺเทยฺย, ปลิพนฺเธยฺย (ฏีกานุรูปํ)}, กาเลน กาลํ ปํสุํ ทเทยฺย, กาเลน}

\vspace{\tipitakaparskip}
\csromancenter{นิทานวคฺคสํยุตฺตปาฬิ กาลํ อุทกํ ทเทยฺย.\csromanspacer{} เอวํ หิ โส ภิกฺขเว ตรุโณ รุกฺโข ตทาหาโร ตทุปาทาโน วุทฺธิํ วิรูฬฺหิํ เวปุลฺลํ อาปชฺเชยฺย.\csromanspacer{} เอวเมว โข ภิกฺขเว สํโยชนิเยสุ ธมฺเมสุ อสฺสาทานุปสฺสิโน วิหรโต ตณฺหา ปวฑฺฒติ, ตณฺหาปจฺจยา อุปาทานํ ฯเปฯ เอวเมตสฺส เกวลสฺส ทุกฺขกฺขนฺธสฺส สมุทโย โหติ.}
\prose{\csromanfolio{316}สํโยชนิเยสุ ภิกฺขเว ธมฺเมสุ อาทีนวานุปสฺสิโน วิหรโต ตณฺหา นิรุชฺฌติ, ตณฺหานิโรธา อุปาทานนิโรโธ ฯเปฯ เอวเมตสฺส เกวลสฺส ทุกฺขกฺขนฺธสฺส นิโรโธ โหติ.}

\prose{เสยฺยถาปิ ภิกฺขเว ตรุโณ รุกฺโข.\csromanspacer{} อถ ปุริโส อาคจฺเฉยฺย กุทฺทาลปิฏกํ อาทาย ฯเปฯ นทิยา วา สีฆโสตาย ปวาเหยฺย.\csromanspacer{} เอวํ หิ โส ภิกฺขเว ตรุโณ รุกฺโข อุจฺฉินฺนมูโล อสฺส, ตาลาวตฺถุกโต อนภาวํกโต อายติํ อนุปฺปาทธมฺโม.\csromanspacer{} เอวเมว โข ภิกฺขเว สํโยชนิเยสุ ธมฺเมสุ อาทีนวานุปสฺสิโน วิหรโต ตณฺหา นิรุชฺฌติ, ตณฺหานิโรธา อุปาทานนิโรโธ ฯเปฯ เอวเมตสฺส เกวลสฺส ทุกฺขกฺขนฺธสฺส นิโรโธ โหตีติ.\csromanspacer{}\csromanspacer{}สตฺตมํ.}

\csromanheader{2}{8. นามรูปสุตฺต}
\proseitem{58}{สาวตฺถิยํ วิหรติ.\csromanspacer{} สํโยชนิเยสุ ภิกฺขเว ธมฺเมสุ อสฺสาทานุปสฺสิโน วิหรโต นามรูปสฺส อวกฺกนฺติ โหติ, นามรูปปจฺจยา สฬายตนํ ฯเปฯ เอวเมตสฺส เกวลสฺส ทุกฺขกฺขนฺธสฺส สมุทโย โหติ.}

\prose{เสยฺยถาปิ ภิกฺขเว มหารุกฺโข.\csromanspacer{} ตสฺส ยานิ เจว มูลานิ อโธคมานิ ยานิ จ ติริยงฺคมานิ, สพฺพานิ ตานิ อุทฺธํ โอชํ อภิหรนฺติ.\csromanspacer{} เอวํ หิ โส ภิกฺขเว มหารุกฺโข ตทาหาโร ตทุปาทาโน จิรํ ทีฆมทฺธานํ ติฏฺเฐยฺย.\csromanspacer{} เอวเมว โข ภิกฺขเว สํโยชนิเยสุ ธมฺเมสุ อสฺสาทานุปสฺสิโน วิหรโต นามรูปสฺส อวกฺกนฺติ โหติ ฯเปฯ.}

\prose{สํโยชนิเยสุ ภิกฺขเว ธมฺเมสุ อาทีนวานุปสฺสิโน วิหรโต นามรูปสฺส อวกฺกนฺติ น โหติ, นามรูปนิโรธา สฬายตนนิโรโธ ฯเปฯ เอวเมตสฺส เกวลสฺส ทุกฺขกฺขนฺธสฺส นิโรโธ โหติ.}

\csromanheader{2}{1. นิทานสํยุตฺต}
\prose{\csromanfolio{317}เสยฺยถาปิ ภิกฺขเว มหารุกฺโข.\csromanspacer{} อถ ปุริโส อาคจฺเฉยฺย กุทฺทาลปิฏกํ อาทาย ฯเปฯ อายติํ อนุปฺปาทธมฺโม.\csromanspacer{} เอวเมว โข ภิกฺขเว สํโยชนิเยสุ ธมฺเมสุ อาทีนวานุปสฺสิโน วิหรโต นามรูปสฺส อวกฺกนฺติ น โหติ, นามรูปนิโรธา สฬายตนนิโรโธ ฯเปฯ เอวเมตสฺส เกวลสฺส ทุกฺขกฺขนฺธสฺส นิโรโธ โหตีติ.\csromanspacer{}\csromanspacer{}อฏฺฐมํ.}

\csromanheader{2}{9. วิญฺญาณสุตฺต}
\proseitem{59}{สาวตฺถิยํ วิหรติ.\csromanspacer{} สํโยชนิเยสุ ภิกฺขเว ธมฺเมสุ อสฺสาทานุปสฺสิโน วิหรโต วิญฺญาณสฺส อวกฺกนฺติ โหติ, วิญฺญาณปจฺจยา นามรูปํ ฯเปฯ เอวเมตสฺส เกวลสฺส ทุกฺขกฺขนฺธสฺส สมุทโย โหติ.}

\csromanmatikaanchor{mtk.64}
\csromantocmarkref{subsection}{2. อุปาทานสุตฺต}{mtk.64}
\bookmark[dest={mtk.64},level=2]{2. อุปาทานสุตฺต}
\prose{เสยฺยถาปิ ภิกฺขเว มหารุกฺโข.\csromanspacer{} ตสฺส ยานิ เจว มูลานิ ฯเปฯ.\csromanspacer{} เอวเมว โข ภิกฺขเว สํโยชนิเยสุ ธมฺเมสุ อสฺสาทานุปสฺสิโน วิหรโต วิญฺญาณสฺส อวกฺกนฺติ โหติ ฯเปฯ.}

\prose{สํโยชนิเยสุ ภิกฺขเว ธมฺเมสุ อาทีนวานุปสฺสิโน วิหรโต วิญฺญาณสฺส อวกฺกนฺติ น โหติ.\csromanspacer{} วิญฺญาณนิโรธา นามรูปนิโรโธ ฯเปฯ เอวเมตสฺส เกวลสฺส ทุกฺขกฺขนฺธสฺส นิโรโธ โหติ.}

\prose{เสยฺยถาปิ ภิกฺขเว มหารุกฺโข.\csromanspacer{} อถ ปุริโส อาคจฺเฉยฺย กุทฺทาลปิฏกํ อาทาย ฯเปฯ อายติํ อนุปฺปาทธมฺโม.\csromanspacer{} เอวเมว โข ภิกฺขเว สํโยชนิเยสุ ธมฺเมสุ อาทีนวานุปสฺสิโน วิหรโต วิญฺญาณสฺส อวกฺกนฺติ น โหติ, วิญฺญาณสฺส นิโรธา นามรูปนิโรโธ ฯเปฯ เอวเมตสฺส เกวลสฺส ทุกฺขกฺขนฺธสฺส นิโรโธ โหตีติ.\csromanspacer{}\csromanspacer{}นวมํ.}

\csromanheader{2}{10. นิทานสุตฺต}
\proseitem{60}{เอกํ สมยํ ภควา กุรูสุ วิหรติ กมฺมาสธมฺมํ นาม กุรูนํ นิคโม.\csromanspacer{} อถ โข อายสฺมา อานนฺโท เยน ภควา เตนุปสงฺกมิ, อุปสงฺกมิตฺวา ภควนฺตํ อภิวาเทตฺวา เอกมนฺตํ นิสีทิ, เอกมนฺตํ นิสินฺโน โข อายสฺมา อานนฺโท ภควนฺตํ เอตทโวจ “อจฺฉริยํ ภนฺเต, อพฺภุตํ ภนฺเต, ยาว คมฺภีโร จายํ ภนฺเต ปฏิจฺจสมุปฺปาโท คมฺภีราวภาโส จ, อถ จ ปน เม อุตฺตานกุตฺตานโก วิย ขายตี”ติ.}

\csromanmatikaanchor{mtk.65}
\csromantocmarkref{subsection}{3. สํโยชนสุตฺต}{mtk.65}
\bookmark[dest={mtk.65},level=2]{3. สํโยชนสุตฺต}
\gambhira{นิทานวคฺคสํยุตฺตปาฬิ}
\prose{\csromanfolio{318}มา เหวํ อานนฺท มา เหวํ อานนฺท\footnote{มา เหวํ อานนฺท อวจ มา เหวํ อานนฺท อวจ (ที 2 มหานิทานสุตฺเต) 2. คุฬาคุณฺฐิกชาตา (สี), คุฬีคุณฺฐิกชาตา (สฺยา, กํ)}, คมฺภีโร จายํ อานนฺท ปฏิจฺจสมุปฺปาโท คมฺภีราวภาโส จ, เอตสฺส อานนฺท ธมฺมสฺส อนนุโพธา อปฺปฏิเวธา เอวมยํ ปชา ตนฺตากุลกชาตา กุลคณฺฐิกชาตา2, มุญฺชปพฺพชภูตา\footnote{มุญฺชพพฺพชภูตา (สี)} อปายํ ทุคฺคติํ วินิปาตํ สํสารํ นาติวตฺตติ.}

\prose{อุปาทานิเยสุ อานนฺท ธมฺเมสุ อสฺสาทานุปสฺสิโน วิหรโต ตณฺหา ปวฑฺฒติ, ตณฺหาปจฺจยา อุปาทานํ, อุปาทานปจฺจยา ภโว, ภวปจฺจยา ชาติ, ชาติปจฺจยา ชรามรณํ โสกปริเทวทุกฺขโทมนสฺสุปายาสา สมฺภวนฺติ, เอวเมตสฺส เกวลสฺส ทุกฺขกฺขนฺธสฺส สมุทโย โหติ.}

\prose{เสยฺยถาปิ อานนฺท มหารุกฺโข.\csromanspacer{} ตสฺส ยานิ เจว มูลานิ อโธคมานิ ยานิ จ ติริยงฺคมานิ, สพฺพานิ ตานิ อุทฺธํ โอชํ อภิหรนฺติ.\csromanspacer{} เอวํ หิ โส อานนฺท มหารุกฺโข ตทาหาโร ตทุปาทาโน จิรํ ทีฆมทฺธานํ ติฏฺเฐยฺย.\csromanspacer{} เอวเมว โข อานนฺท อุปาทานิเยสุ ธมฺเมสุ อสฺสาทานุปสฺสิโน วิหรโต ตณฺหา ปวฑฺฒติ, ตณฺหาปจฺจยา อุปาทานํ, อุปาทานปจฺจยา ภโว ฯเปฯ เอวเมตสฺส เกวลสฺส ทุกฺขกฺขนฺธสฺส สมุทโย โหติ.}

\prose{อุปาทานิเยสุ อานนฺท ธมฺเมสุ อาทีนวานุปสฺสิโน วิหรโต ตณฺหา นิรุชฺฌติ, ตณฺหานิโรธา อุปาทานนิโรโธ, อุปาทานนิโรธา ภวนิโรโธ ฯเปฯ เอวเมตสฺส เกวลสฺส ทุกฺขกฺขนฺธสฺส นิโรโธ โหติ.}

\prose{เสยฺยถาปิ อานนฺท มหารุกฺโข.\csromanspacer{} อถ ปุริโส อาคจฺเฉยฺย กุทฺทาลปิฏกํ อาทาย, โส ตํ รุกฺขํ มูเล ฉินฺเทยฺย, มูเล เฉตฺวา ปลิขเณยฺย, ปลิขณิตฺวา มูลานิ อุทฺธเรยฺย อนฺตมโส อุสีรนาฬิมตฺตานิปิ.\csromanspacer{} โส ตํ รุกฺขํ ขณฺฑาขณฺฑิกํ ฉินฺเทยฺย, ขณฺฑาขณฺฑิกํ ฉินฺทิตฺวา ผาเลยฺย, ผาเลตฺวา สกลิกํ สกลิกํ กเรยฺย, สกลิกํ สกลิกํ กริตฺวา วาตาตเป วิโสเสยฺย, วาตาตเป วิโสเสตฺวา อคฺคินา ฑเหยฺย, อคฺคินา ฑเหตฺวา มสิํ กเรยฺย, มสิํ กริตฺวา มหาวาเต วา โอผุเณยฺย, นทิยา วา สีฆโสตาย ปวาเหยฺย.\csromanspacer{} เอวํ หิ โส อานนฺท มหารุกฺโข อุจฺฉินฺนมูโล อสฺส, ตาลาวตฺถุกโต อนภาวํกโต อายติํ}

\csromanmatikaanchor{mtk.66}
\csromantocmarkref{subsection}{4. ทุติยสํโยชนสุตฺต}{mtk.66}
\bookmark[dest={mtk.66},level=2]{4. ทุติยสํโยชนสุตฺต}
\vspace{\tipitakaparskip}
\csromancenter{นิทานสํยุตฺต อนุปฺปาทธมฺโม.\csromanspacer{} เอวเมว โข อานนฺท อุปาทานิเยสุ ธมฺเมสุ อาทีนวานุปสฺสิโน วิหรโต ตณฺหา นิรุชฺฌติ, ตณฺหานิโรธา อุปาทานนิโรโธ, อุปาทานนิโรธา ภวนิโรโธ, ภวนิโรธา ชาตินิโรโธ, ชาตินิโรธา ชรามรณํ โสกปริเทวทุกฺขโทมนสฺสุปายาสา นิรุชฺฌนฺติ, เอวเมตสฺส เกวลสฺส ทุกฺขกฺขนฺธสฺส นิโรโธ โหตีติ.\csromanspacer{}\csromanspacer{}ทสมํ.}
\csromancenter{ทุกฺขวคฺโค ฉฏฺโฐ.}
\titlehead{\textbf{ตสฺสุทฺทานํ}}
\csromanmatikaanchor{mtk.67}
\csromanmatikaanchor{mtk.73}
\csromantocmarkref{subsection}{5. มหารุกฺขสุตฺต}{mtk.67}
\bookmark[dest={mtk.67},level=2]{5. มหารุกฺขสุตฺต}
\csromangathameasure{ปริวีมํสนุปาทานํ, เทฺว จ สํโยชนานิ จ. \\ มหารุกฺเขน เทฺว วุตฺตา, ตรุเณน จ สตฺตมํ. \\ นามรูปญฺจ วิญฺญาณํ, นิทาเนน จ เต ทสาติ.}
\csromangathagroup{\csromangathastanza{\csromanfolio{319}ปริวีมํสนุปาทานํ, เทฺว จ สํโยชนานิ จ. \\ มหารุกฺเขน เทฺว วุตฺตา, ตรุเณน จ สตฺตมํ. \\ นามรูปญฺจ วิญฺญาณํ, นิทาเนน จ เต ทสาติ.}}

\csromanheader{1}{7. มหาวคฺค}
\csromanheader{2}{1. อสฺสุตวาสุตฺต}
\proseitem{61}{เอวํ เม สุตํ\csromandash{}เอกํ สมยํ ภควา สาวตฺถิยํ วิหรติ เชตวเน อนาถปิณฺฑิกสฺส อาราเม.\csromanspacer{} อสฺสุตวา ภิกฺขเว ปุถุชฺชโน อิมสฺมิํ จาตุมหาภูติกสฺมิํ\footnote{จาตุมฺมหาภูติกสฺมิํ (สี, สฺยา, กํ)} กายสฺมิํ นิพฺพินฺเทยฺยปิ วิรชฺเชยฺยปิ วิมุจฺเจยฺยปิ.\csromanspacer{} ตํ กิสฺส เหตุ, ทิสฺสติ ภิกฺขเว\footnote{ทิสฺสติ หิ ภิกฺขเว (สี, สฺยา, กํ)} อิมสฺส จาตุมหาภูติกสฺส กายสฺส อาจโยปิ อปจโยปิ อาทานมฺปิ นิกฺเขปนมฺปิ.\csromanspacer{} ตสฺมา ตตฺราสฺสุตวา ปุถุชฺชโน นิพฺพินฺเทยฺยปิ วิรชฺเชยฺยปิ วิมุจฺเจยฺยปิ.}

\prose{ยญฺจ โข เอตํ ภิกฺขเว วุจฺจติ จิตฺตํ อิติปิ มโน อิติปิ วิญฺญาณํ อิติปิ.\csromanspacer{} ตตฺราสฺสุตวา ปุถุชฺชโน นาลํ นิพฺพินฺทิตุํ, นาลํ วิรชฺชิตุํ, นาลํ วิมุจฺจิตุํ.\csromanspacer{} ตํ กิสฺส เหตุ, ทีฆรตฺตํ เหตํ ภิกฺขเว อสฺสุตวโต ปุถุชฺชนสฺส อชฺโฌสิตํ มมายิตํ ปรามฏฺฐํ “เอตํ มม, เอโสหมสฺมิ, เอโส เม อตฺตา”ติ.\csromanspacer{} ตสฺมา ตตฺราสฺสุตวา ปุถุชฺชโน นาลํ นิพฺพินฺทิตุํ, นาลํ วิรชฺชิตุํ, นาลํ วิมุจฺจิตุํ.}

\gambhira{นิทานวคฺคสํยุตฺตปาฬิ}
\prose{\csromanfolio{320}วรํ ภิกฺขเว อสฺสุตวา ปุถุชฺชโน อิมํ จาตุมหาภูติกํ กายํ อตฺตโต อุปคจฺเฉยฺย, น เตฺวว จิตฺตํ.\csromanspacer{} ตํ กิสฺส เหตุ, ทิสฺสตายํ ภิกฺขเว จาตุมหาภูติโก กาโย เอกมฺปิ วสฺสํ ติฏฺฐมาโน เทฺวปิ วสฺสานิ ติฏฺฐมาโน ตีณิปิ วสฺสานิ ติฏฺฐมาโน จตฺตาริปิ วสฺสานิ ติฏฺฐมาโน ปญฺจปิ วสฺสานิ ติฏฺฐมาโน ทสปิ วสฺสานิ ติฏฺฐมาโน วีสติปิ วสฺสานิ ติฏฺฐมาโน ติํสมฺปิ วสฺสานิ ติฏฺฐมาโน จตฺตารีสมฺปิ วสฺสานิ ติฏฺฐมาโน ปญฺญาสมฺปิ วสฺสานิ ติฏฺฐมาโน วสฺสสตมฺปิ ติฏฺฐมาโน ภิยฺโยปิ ติฏฺฐมาโน.}

\csromanmatikaanchor{mtk.68}
\csromantocmarkref{subsection}{6. ทุติยมหารุกฺขสุตฺต}{mtk.68}
\bookmark[dest={mtk.68},level=2]{6. ทุติยมหารุกฺขสุตฺต}
\prose{ยญฺจ โข เอตํ ภิกฺขเว วุจฺจติ จิตฺตํ อิติปิ มโน อิติปิ วิญฺญาณํ อิติปิ, ตํ รตฺติยา จ ทิวสสฺส จ อญฺญเทว อุปฺปชฺชติ อญฺญํ นิรุชฺฌติ.\csromanspacer{} เสยฺยถาปิ ภิกฺขเว มกฺกโฏ อรญฺเญ ปวเน จรมาโน สาขํ คณฺหติ, ตํ มุญฺจิตฺวา อญฺญํ คณฺหติ, ตํ มุญฺจิตฺวา อญฺญํ คณฺหติ.\csromanspacer{} เอวเมว โข ภิกฺขเว ยมิทํ วุจฺจติ จิตฺตํ อิติปิ มโน อิติปิ วิญฺญาณํ อิติปิ, ตํ รตฺติยา จ ทิวสสฺส จ อญฺญเทว อุปฺปชฺชติ อญฺญํ นิรุชฺฌติ.}

\prose{ตตฺร ภิกฺขเว สุตวา อริยสาวโก ปฏิจฺจสมุปฺปาทํเยว สาธุกํ โยนิโส มนสิ กโรติ “อิติ อิมสฺมิํ สติ อิทํ โหติ, อิมสฺสุปฺปาทา อิทํ อุปฺปชฺชติ.\csromanspacer{} อิมสฺมิํ อสติ อิทํ น โหติ, อิมสฺส นิโรธา อิทํ นิรุชฺฌติ.\csromanspacer{} ยทิทํ อวิชฺชาปจฺจยา สงฺขารา, สงฺขารปจฺจยา วิญฺญาณํ ฯเปฯ เอวเมตสฺส เกวลสฺส ทุกฺขกฺขนฺธสฺส สมุทโย โหติ.\csromanspacer{} อวิชฺชาย เตฺวว อเสสวิราคนิโรธา สงฺขารนิโรโธ, สงฺขารนิโรธา วิญฺญาณนิโรโธ ฯเปฯ เอวเมตสฺส เกวลสฺส ทุกฺขกฺขนฺธสฺส นิโรโธ โหตี”ติ.}

\prose{เอวํ ปสฺสํ ภิกฺขเว สุตวา อริยสาวโก รูปสฺมิมฺปิ นิพฺพินฺทติ, เวทนายปิ นิพฺพินฺทติ, สญฺญายปิ นิพฺพินฺทติ, สงฺขาเรสุปิ นิพฺพินฺทติ, วิญฺญาณสฺมิมฺปิ นิพฺพินฺทติ นิพฺพินฺทํ วิรชฺชติ, วิราคา วิมุจฺจติ, วิมุตฺตสฺมิํ “วิมุตฺตมฺ”อิติ ญาณํ โหติ, “ขีณา ชาติ, วุสิตํ พฺรหฺมจริยํ, กตํ กรณียํ, นาปรํ อิตฺถตฺตายา”ติ ปชานาตีติ.\csromanspacer{}\csromanspacer{}ปฐมํ.}

\csromanmatikaanchor{mtk.69}
\csromantocmarkref{subsection}{7. ตรุณรุกฺขสุตฺต}{mtk.69}
\bookmark[dest={mtk.69},level=2]{7. ตรุณรุกฺขสุตฺต}
\csromanheader{2}{2. ทุติย-อสฺสุตวาสุตฺต}
\proseitem{62}{สาวตฺถิยํ วิหรติ.\csromanspacer{} อสฺสุตวา ภิกฺขเว ปุถุชฺชโน อิมสฺมิํ จาตุมหาภูติกสฺมิํ กายสฺมิํ นิพฺพินฺเทยฺยปิ วิรชฺเชยฺยปิ วิมุจฺเจยฺยปิ.\csromanspacer{} ตํ กิสฺส เหตุ, ทิสฺสติ ภิกฺขเว อิมสฺส จาตุมหาภูติกสฺส กายสฺส}

\vspace{\tipitakaparskip}
\csromancenter{นิทานสํยุตฺต อาจโยปิ อปจโยปิ อาทานมฺปิ นิกฺเขปนมฺปิ.\csromanspacer{} ตสฺมา ตตฺราสฺสุตวา ปุถุชฺชโน นิพฺพินฺเทยฺยปิ วิรชฺเชยฺยปิ วิมุจฺเจยฺยปิ.\csromanspacer{} ยญฺจ โข เอตํ ภิกฺขเว วุจฺจติ จิตฺตํ อิติปิ มโน อิติปิ วิญฺญาณํ อิติปิ.\csromanspacer{} ตตฺราสฺสุตวา ปุถุชฺชโน นาลํ นิพฺพินฺทิตุํ, นาลํ วิรชฺชิตุํ, นาลํ วิมุจฺจิตุํ.\csromanspacer{} ตํ กิสฺส เหตุ, ทีฆรตฺตํ เหตํ ภิกฺขเว อสฺสุตวโต ปุถุชฺชนสฺส อชฺโฌสิตํ มมายิตํ ปรามฏฺฐํ “เอตํ มม, เอโสหมสฺมิ, เอโส เม อตฺตา”ติ.\csromanspacer{} ตสฺมา ตตฺราสฺสุตวา ปุถุชฺชโน นาลํ นิพฺพินฺทิตุํ, นาลํ วิรชฺชิตุํ, นาลํ วิมุจฺจิตุํ.}
\prose{\csromanfolio{321}วรํ ภิกฺขเว อสฺสุตวา ปุถุชฺชโน อิมํ จาตุมหาภูติกํ กายํ อตฺตโต อุปคจฺเฉยฺย, น เตฺวว จิตฺตํ.\csromanspacer{} ตํ กิสฺส เหตุ, ทิสฺสตายํ ภิกฺขเว จาตุมหาภูติโก กาโย เอกมฺปิ วสฺสํ ติฏฺฐมาโน เทฺวปิ วสฺสานิ ติฏฺฐมาโน ตีณิปิ วสฺสานิ ติฏฺฐมาโน จตฺตาริปิ วสฺสานิ ติฏฺฐมาโน ปญฺจปิ วสฺสานิ ติฏฺฐมาโน ทสปิ วสฺสานิ ติฏฺฐมาโน วีสติปิ วสฺสานิ ติฏฺฐมาโน ติํสมฺปิ วสฺสานิ ติฏฺฐมาโน จตฺตารีสมฺปิ วสฺสานิ ติฏฺฐมาโน ปญฺญาสมฺปิ วสฺสานิ ติฏฺฐมาโน วสฺสสตมฺปิ ติฏฺฐมาโน ภิยฺโยปิ ติฏฺฐมาโน.\csromanspacer{} ยญฺจ โข เอตํ ภิกฺขเว วุจฺจติ จิตฺตํ อิติปิ มโน อิติปิ วิญฺญาณํ อิติปิ, ตํ รตฺติยา จ ทิวสสฺส จ อญฺญเทว อุปฺปชฺชติ, อญฺญํ นิรุชฺฌติ.}

\prose{ตตฺร ภิกฺขเว สุตวา อริยสาวโก ปฏิจฺจสมุปฺปาทํเยว สาธุกํ โยนิโส มนสิ กโรติ “อิติ อิมสฺมิํ สติ อิทํ โหติ, อิมสฺสุปฺปาทา อิทํ อุปฺปชฺชติ.\csromanspacer{} อิมสฺมิํ อสติ อิทํ น โหติ, อิมสฺส นิโรธา อิทํ นิรุชฺฌตี”ติ.\csromanspacer{} สุขเวทนิยํ ภิกฺขเว ผสฺสํ ปฏิจฺจ อุปฺปชฺชติ สุขเวทนา, ตสฺเสว สุขเวทนิยสฺส ผสฺสสฺส นิโรธา ยํ ตชฺชํ เวทยิตํ สุขเวทนิยํ ผสฺสํ ปฏิจฺจ อุปฺปนฺนา สุขเวทนา, สา นิรุชฺฌติ สา วูปสมฺมติ.\csromanspacer{} ทุกฺขเวทนิยํ ภิกฺขเว ผสฺสํ ปฏิจฺจ อุปฺปชฺชติ ทุกฺขเวทนา, ตสฺเสว ทุกฺขเวทนิยสฺส ผสฺสสฺส นิโรธา ยํ ตชฺชํ เวทยิตํ ทุกฺขเวทนิยํ ผสฺสํ ปฏิจฺจ อุปฺปนฺนา ทุกฺขเวทนา, สา นิรุชฺฌติ สา วูปสมฺมติ.\csromanspacer{} อทุกฺขมสุขเวทนิยํ ภิกฺขเว ผสฺสํ ปฏิจฺจ อุปฺปชฺชติ อทุกฺขมสุขเวทนา, ตสฺเสว อทุกฺขมสุขเวทนิยสฺส ผสฺสสฺส นิโรธา ยํ ตชฺชํ เวทยิตํ อทุกฺขมสุขเวทนิยํ ผสฺสํ ปฏิจฺจ อุปฺปนฺนา อทุกฺขมสุขเวทนา, สา นิรุชฺฌติ สา วูปสมฺมติ.}

\gambhira{นิทานวคฺคสํยุตฺตปาฬิ}
\csromanmatikaanchor{mtk.70}
\csromantocmarkref{subsection}{8. นามรูปสุตฺต}{mtk.70}
\bookmark[dest={mtk.70},level=2]{8. นามรูปสุตฺต}
\prose{\csromanfolio{322}เสยฺยถาปิ ภิกฺขเว ทฺวินฺนํ กฏฺฐานํ สํฆฏฺฏนสโมธานา อุสฺมา ชายติ, เตโช อภินิพฺพตฺตติ, เตสํเยว ทฺวินฺนํ กฏฺฐานํ นานากตวินิพฺโภคา\footnote{นานาภาวา วินิกฺเขปา (สี, อิ) ม 3 (286) ปิฏฺเฐปิ.}, ยา ตชฺชา อุสฺมา, สา นิรุชฺฌติ สา วูปสมฺมติ.\csromanspacer{} เอวเมว โข ภิกฺขเว สุขเวทนิยํ ผสฺสํ ปฏิจฺจ อุปฺปชฺชติ สุขเวทนา, ตสฺเสว สุขเวทนิยสฺส ผสฺสสฺส นิโรธา ยํ ตชฺชํ เวทยิตํ สุขเวทนิยํ ผสฺสํ ปฏิจฺจ อุปฺปนฺนา สุขเวทนา, สา นิรุชฺฌติ สา วูปสมฺมติ ฯเปฯ.\csromanspacer{} ทุกฺขเวทนิยํ ผสฺสํ ปฏิจฺจ.\csromanspacer{} อทุกฺขมสุขเวทนิยํ ผสฺสํ ปฏิจฺจ อุปฺปชฺชติ อทุกฺขมสุขเวทนา, ตสฺเสว อทุกฺขมสุขเวทนิยสฺส ผสฺสสฺส นิโรธา ยํ ตชฺชํ เวทยิตํ อทุกฺขมสุขเวทนิยํ ผสฺสํ ปฏิจฺจ อุปฺปนฺนา อทุกฺขมสุขเวทนา, สา นิรุชฺฌติ สา วูปสมฺมติ.}

\prose{เอวํ ปสฺสํ ภิกฺขเว สุตวา อริยสาวโก ผสฺเสปิ นิพฺพินฺทติ, เวทนายปิ นิพฺพินฺทติ, สญฺญายปิ นิพฺพินฺทติ, สงฺขาเรสุปิ นิพฺพินฺทติ, วิญฺญาณสฺมิมฺปิ นิพฺพินฺทติ, นิพฺพินฺทํ วิรชฺชติ, วิราคา วิมุจฺจติ, วิมุตฺตสฺมิํ “วิมุตฺตมฺ”อิติ ญาณํ โหติ, “ขีณา ชาติ, วุสิตํ พฺรหฺมจริยํ, กตํ กรณียํ, นาปรํ อิตฺถตฺตายา”ติ ปชานาตีติ.\csromanspacer{}\csromanspacer{}ทุติยํ.}

\csromanheader{2}{3. ปุตฺตมํสูปมสุตฺต}
\proseitem{63}{สาวตฺถิยํ.\csromanspacer{} จตฺตาโรเม ภิกฺขเว อาหารา ภูตานํ วา สตฺตานํ ฐิติยา สมฺภเวสีนํ วา อนุคฺคหาย.\csromanspacer{} กตเม จตฺตาโร, กพฬีกาโร อาหาโร โอฬาริโก วา สุขุโม วา ผสฺโส ทุติโย มโนสญฺเจตนา ตติยา วิญฺญาณํ จตุตฺถํ.\csromanspacer{} อิเม โข ภิกฺขเว จตฺตาโร อาหารา ภูตานํ วา สตฺตานํ ฐิติยา สมฺภเวสีนํ วา อนุคฺคหาย.}

\prose{กถญฺจ ภิกฺขเว กพฬีกาโร อาหาโร ทฏฺฐพฺโพ.\csromanspacer{} เสยฺยถาปิ ภิกฺขเว เทฺว ชายมฺปติกา\footnote{ชยมฺปติกา (สี, อิ) ฏีกา โอโลเกตพฺพา.} ปริตฺตํ สมฺพลํ อาทาย กนฺตารมคฺคํ ปฏิปชฺเชยฺยุํ.\csromanspacer{} เตสมสฺส เอกปุตฺตโก ปิโย มนาโป.\csromanspacer{} อถ โข เตสํ ภิกฺขเว ทฺวินฺนํ ชายมฺปติกานํ กนฺตารคตานํ ยา ปริตฺตา สมฺพลมตฺตา, สา ปริกฺขยํ ปริยาทานํ คจฺเฉยฺย.\csromanspacer{} สิยา จ เนสํ กนฺตาราวเสโส อนติณฺโณ.\csromanspacer{} อถ โข เตสํ ภิกฺขเว ทฺวินฺนํ ชายมฺปติกานํ เอวมสฺส “อมฺหากํ โข ยา ปริตฺตา สมฺพลมตฺตา สา ปริกฺขีณา ปริยาทิณฺณา\footnote{ปริยาทินฺนา (สฺยา, กํ)},}

\vspace{\tipitakaparskip}
\csromancenter{นิทานสํยุตฺต อตฺถิ จายํ กนฺตาราวเสโส อนิตฺติณฺโณ\footnote{อนิตฺถิณฺโณ (สฺยา, กํ), อนติณฺโณ (ก)}.\csromanspacer{} ยํนูน มยํ อิมํ เอกปุตฺตกํ ปิยํ มนาปํ วธิตฺวา วลฺลูรญฺจ โสณฺฑิกญฺจ กริตฺวา ปุตฺตมํสานิ ขาทนฺตา เอวํ ตํ กนฺตาราวเสสํ นิตฺถเรยฺยาม, มา สพฺเพว ตโย วินสฺสิมฺหา”ติ.\csromanspacer{} อถ โข เต ภิกฺขเว เทฺว ชายมฺปติกา ตํ เอกปุตฺตกํ ปิยํ มนาปํ วธิตฺวา วลฺลูรญฺจ โสณฺฑิกญฺจ กริตฺวา ปุตฺตมํสานิ ขาทนฺตา เอวํ ตํ กนฺตาราวเสสํ นิตฺถเรยฺยุํ, เต ปุตฺตมํสานิ เจว ขาเทยฺยุํ, อุเร จ ปฏิปิเสยฺยุํ “กหํ เอกปุตฺตก กหํ เอกปุตฺตกา”ติ.}
\csromanmatikaanchor{mtk.71}
\csromantocmarkref{subsection}{9. วิญฺญาณสุตฺต}{mtk.71}
\bookmark[dest={mtk.71},level=2]{9. วิญฺญาณสุตฺต}
\prose{\csromanfolio{323}ตํ กิํ มญฺญถ ภิกฺขเว อปิ นุ เต ทวาย วา อาหารํ อาหาเรยฺยุํ, มทาย วา อาหารํ อาหาเรยฺยุํ, มณฺฑนาย วา อาหารํ อาหาเรยฺยุํ, วิภูสนาย วา อาหารํ อาหาเรยฺยุนฺติ.\csromanspacer{} โน เหตํ ภนฺเต.\csromanspacer{} นนุ เต ภิกฺขเว ยาวเทว กนฺตารสฺส นิตฺถรณตฺถาย อาหารํ อาหาเรยฺยุนฺติ.\csromanspacer{} เอวํ ภนฺเต.\csromanspacer{} เอวเมว ขฺวาหํ ภิกฺขเว “กพฬีกาโร อาหาโร ทฏฺฐพฺโพ”ติ วทามิ.\csromanspacer{} กพฬีกาเร ภิกฺขเว อาหาเร ปริญฺญาเต ปญฺจ กามคุณิโก ราโค ปริญฺญาโต โหติ.\csromanspacer{} ปญฺจ กามคุณิเก ราเค ปริญฺญาเต นตฺถิ ตํ สํโยชนํ เยน สํโยชเนน สํยุตฺโต อริยสาวโก ปุน อิมํ โลกํ อาคจฺเฉยฺย.}

\prose{กถญฺจ ภิกฺขเว ผสฺสาหาโร ทฏฺฐพฺโพ.\csromanspacer{} เสยฺยถาปิ ภิกฺขเว คาวี นิจฺจมฺมา กุฏฺฏํ เจ\footnote{กุฑฺฑญฺเจ (สี, สฺยา, กํ, อิ)} นิสฺสาย ติฏฺเฐยฺย, เย กุฏฺฏนิสฺสิตา ปาณา, เต นํ ขาเทยฺยุํ.\csromanspacer{} รุกฺขํ เจ นิสฺสาย ติฏฺเฐยฺย, เย รุกฺขนิสฺสิตา ปาณา, เต นํ ขาเทยฺยุํ.\csromanspacer{} อุทกํ เจ นิสฺสาย ติฏฺเฐยฺย, เย อุทกนิสฺสิตา ปาณา, เต นํ ขาเทยฺยุํ.\csromanspacer{} อากาสํ เจ นิสฺสาย ติฏฺเฐยฺย, เย อากาสนิสฺสิตา ปาณา, เต นํ ขาเทยฺยุํ.\csromanspacer{} ยํ ยเทว หิ สา ภิกฺขเว คาวี นิจฺจมฺมา นิสฺสาย ติฏฺเฐยฺย, เย ตนฺนิสฺสิตา\footnote{เย ตนฺนิสฺสิตา ตนฺนิสฺสิตา (สี, สฺยา, กํ, อิ)} ปาณา, เต นํ ขาเทยฺยุํ.\csromanspacer{} เอวเมว ขฺวาหํ ภิกฺขเว “ผสฺสาหาโร ทฏฺฐพฺโพ”ติ วทามิ.\csromanspacer{} ผสฺเส ภิกฺขเว อาหาเร ปริญฺญาเต ติสฺโส เวทนา ปริญฺญาตา โหนฺติ.\csromanspacer{} ตีสุ เวทนาสุ ปริญฺญาตาสุ อริยสาวกสฺส นตฺถิ กิญฺจิ อุตฺตริกรณียนฺติ\footnote{อุตฺตริํกรณียนฺติ (สี, อิ)} วทามิ.}

\gambhira{นิทานวคฺคสํยุตฺตปาฬิ}
\prose{\csromanfolio{324}กถญฺจ ภิกฺขเว มโนสญฺเจตนาหาโร ทฏฺฐพฺโพ.\csromanspacer{} เสยฺยถาปิ ภิกฺขเว องฺคารกาสุ สาธิกโปริสา ปุณฺณา องฺคารานํ วีตจฺจิกานํ วีตธูมานํ.\csromanspacer{} อถ ปุริโส อาคจฺเฉยฺย ชีวิตุกาโม อมริตุกาโม สุขกาโม ทุกฺขปฺปฏิกูโล, ตเมนํ เทฺว พลวนฺโต ปุริสา นานาพาหาสุ คเหตฺวา ตํ องฺคารกาสุํ อุปกฑฺเฒยฺยุํ.\csromanspacer{} อถ โข ภิกฺขเว ตสฺส ปุริสสฺส อารกาวสฺส เจตนา อารกา ปตฺถนา อารกา ปณิธิ.\csromanspacer{} ตํ กิสฺส เหตุ, เอวํ หิ ภิกฺขเว ตสฺส ปุริสสฺส โหติ “อิมํ จาหํ องฺคารกาสุํ ปปติสฺสามิ, ตโตนิทานํ มรณํ วา นิคจฺฉามิ มรณมตฺตํ วา ทุกฺขนฺ”ติ.\csromanspacer{} เอวเมว ขฺวาหํ ภิกฺขเว “มโนสญฺเจตนาหาโร ทฏฺฐพฺโพ”ติ วทามิ.\csromanspacer{} มโนสญฺเจตนาย ภิกฺขเว อาหาเร ปริญฺญาเต ติสฺโส ตณฺหา ปริญฺญาตา โหนฺติ.\csromanspacer{} ตีสุ ตณฺหาสุ ปริญฺญาตาสุ อริยสาวกสฺส นตฺถิ กิญฺจิ อุตฺตริกรณียนฺติ วทามิ.}

\prose{กถญฺจ ภิกฺขเว วิญฺญาณาหาโร ทฏฺฐพฺโพ.\csromanspacer{} เสยฺยถาปิ ภิกฺขเว โจรํ อาคุจาริํ คเหตฺวา รญฺโญ ทสฺเสยฺยุํ “อยํ เต เทว โจโร อาคุจารี, อิมสฺส ยํ อิจฺฉสิ, ตํ ทณฺฑํ ปเณหี”ติ.\csromanspacer{} ตเมนํ ราชา เอวํ วเทยฺย “คจฺฉถ โภ อิมํ ปุริสํ ปุพฺพณฺหสมยํ สตฺติสเตน หนถา”ติ.\csromanspacer{} ตเมนํ ปุพฺพณฺหสมยํ สตฺติสเตน หเนยฺยุํ.\csromanspacer{} อถ ราชา มชฺฌนฺหิกสมยํ เอวํ วเทยฺย “อมฺโภ กถํ โส ปุริโส”ติ.\csromanspacer{} ตเถว เทว ชีวตีติ.\csromanspacer{} ตเมนํ ราชา เอวํ วเทยฺย “คจฺฉถ โภ ตํ ปุริสํ มชฺฌนฺหิกสมยํ สตฺติสเตน หนถา”ติ.\csromanspacer{} ตเมนํ มชฺฌนฺหิกสมยํ สตฺติสเตน หเนยฺยุํ.\csromanspacer{} อถ ราชา สายนฺหสมยํ เอวํ วเทยฺย “อมฺโภ กถํ โส ปุริโส”ติ.\csromanspacer{} ตเถว เทว ชีวตีติ.\csromanspacer{} ตเมนํ ราชา เอวํ วเทยฺย “คจฺฉถ โภ ตํ ปุริสํ สายนฺหสมยํ สตฺติสเตน หนถา”ติ.\csromanspacer{} ตเมนํ สายนฺหสมยํ สตฺติสเตน หเนยฺยุํ.\csromanspacer{} ตํ กิํ มญฺญถ ภิกฺขเว, อปิ นุ โส ปุริโส ทิวสํ ตีหิ สตฺติสเตหิ หญฺญมาโน ตโตนิทานํ ทุกฺขํ โทมนสฺสํ ปฏิสํเวทิเยถาติ.\csromanspacer{} เอกิสฺสาปิ ภนฺเต สตฺติยา หญฺญมาโน ตโตนิทานํ ทุกฺขํ โทมนสฺสํ ปฏิสํเวทิเยถ, โก ปน วาโท ตีหิ สตฺติสเตหิ หญฺญมาโนติ.\csromanspacer{} เอวเมว ขฺวาหํ ภิกฺขเว “วิญฺญาณาหาโร ทฏฺฐพฺโพ”ติ วทามิ.\csromanspacer{} วิญฺญาเณ ภิกฺขเว อาหาเร ปริญฺญาเต นามรูปํ ปริญฺญาตํ โหติ, นามรูเป ปริญฺญาเต อริยสาวกสฺส นตฺถิ กิญฺจิ อุตฺตริกรณียนฺติ วทามีติ.\csromanspacer{}\csromanspacer{}ตติยํ.}

\csromanmatikaanchor{mtk.72}
\csromantocmarkref{subsection}{10. นิทานสุตฺต}{mtk.72}
\bookmark[dest={mtk.72},level=2]{10. นิทานสุตฺต}
\csromantocmarkref{section}{อุทฺทานคาถา}{mtk.73}
\bookmark[dest={mtk.73},level=1]{อุทฺทานคาถา}
\csromanheader{2}{1. นิทานสํยุตฺต \textbf{4. อตฺถิราคสุตฺต}}
\proseitem{64}{\csromanfolio{325}สาวตฺถิยํ วิหรติ.\csromanspacer{} จตฺตาโรเม ภิกฺขเว อาหารา ภูตานํ วา สตฺตานํ ฐิติยา สมฺภเวสีนํ วา อนุคฺคหาย.\csromanspacer{} กตเม จตฺตาโร, กพฬีกาโร อาหาโร โอฬาริโก วา สุขุโม วา ผสฺโส ทุติโย มโนสญฺเจตนา ตติยา วิญฺญาณํ จตุตฺถํ.\csromanspacer{} อิเม โข ภิกฺขเว จตฺตาโร อาหารา ภูตานํ วา สตฺตานํ ฐิติยา สมฺภเวสีนํ วา อนุคฺคหาย.}

\prose{กพฬีกาเร เจ ภิกฺขเว อาหาเร อตฺถิ ราโค อตฺถิ นนฺที อตฺถิ ตณฺหา, ปติฏฺฐิตํ ตตฺถ วิญฺญาณํ วิรูฬฺหํ.\csromanspacer{} ยตฺถ ปติฏฺฐิตํ วิญฺญาณํ วิรูฬฺหํ, อตฺถิ ตตฺถ นามรูปสฺส อวกฺกนฺติ.\csromanspacer{} ยตฺถ อตฺถิ นามรูปสฺส อวกฺกนฺติ, อตฺถิ ตตฺถ สงฺขารานํ วุทฺธิ.\csromanspacer{} ยตฺถ อตฺถิ สงฺขารานํ วุทฺธิ, อตฺถิ ตตฺถ อายติํ ปุนพฺภวาภินิพฺพตฺติ.\csromanspacer{} ยตฺถ อตฺถิ อายติํ ปุนพฺภวาภินิพฺพตฺติ, อตฺถิ ตตฺถ อายติํ ชาติชรามรณํ.\csromanspacer{} ยตฺถ อตฺถิ อายติํ ชาติชรามรณํ, สโสกํ ตํ ภิกฺขเว สทรํ ส-อุปายาสนฺติ วทามิ.}

\prose{ผสฺเส เจ ภิกฺขเว อาหาเร ฯเปฯ.\csromanspacer{} มโนสญฺเจตนาย เจ ภิกฺขเว อาหาเร.\csromanspacer{} วิญฺญาเณ เจ ภิกฺขเว อาหาเร อตฺถิ ราโค อตฺถิ นนฺที อตฺถิ ตณฺหา, ปติฏฺฐิตํ ตตฺถ วิญฺญาณํ วิรูฬฺหํ.\csromanspacer{} ยตฺถ ปติฏฺฐิตํ วิญฺญาณํ วิรูฬฺหํ, อตฺถิ ตตฺถ นามรูปสฺส อวกฺกนฺติ.\csromanspacer{} ยตฺถ อตฺถิ นามรูปสฺส อวกฺกนฺติ, อตฺถิ ตตฺถ สงฺขารานํ วุทฺธิ.\csromanspacer{} ยตฺถ อตฺถิ สงฺขารานํ วุทฺธิ, อตฺถิ ตตฺถ อายติํ ปุนพฺภวาภินิพฺพตฺติ.\csromanspacer{} ยตฺถ อตฺถิ อายติํ ปุนพฺภวาภินิพฺพตฺติ, อตฺถิ ตตฺถ อายติํ ชาติชรามรณํ.\csromanspacer{} ยตฺถ อตฺถิ อายติํ ชาติชรามรณํ, สโสกํ ตํ ภิกฺขเว สทรํ ส-อุปายาสนฺติ วทามิ.}

\prose{เสยฺยถาปิ ภิกฺขเว รชโก วา จิตฺตการโก วา สติ รชนาย วา ลาขาย วา หลิทฺทิยา วา นีลิยา วา มญฺชิฏฺฐาย วา สุปริมฏฺเฐ วา ผลเก ภิตฺติยา วา ทุสฺสปฏฺเฏ วา อิตฺถิรูปํ วา ปุริสรูปํ วา อภินิมฺมิเนยฺย สพฺพงฺคปจฺจงฺคํ.\csromanspacer{} เอวเมว โข ภิกฺขเว กพฬีกาเร เจ อาหาเร อตฺถิ ราโค อตฺถิ นนฺที อตฺถิ ตณฺหา, ปติฏฺฐิตํ ตตฺถ วิญฺญาณํ วิรูฬฺหํ.\csromanspacer{} ยตฺถ ปติฏฺฐิตํ วิญฺญาณํ วิรูฬฺหํ, อตฺถิ ตตฺถ นามรูปสฺส อวกฺกนฺติ.\csromanspacer{} ยตฺถ อตฺถิ นามรูปสฺส อวกฺกนฺติ, อตฺถิ ตตฺถ สงฺขารานํ วุทฺธิ.\csromanspacer{} ยตฺถ อตฺถิ สงฺขารานํ วุทฺธิ, อตฺถิ ตตฺถ อายติํ ปุนพฺภวาภินิพฺพตฺติ.\csromanspacer{} ยตฺถ อตฺถิ อายติํ ปุนพฺภวาภินิพฺพตฺติ, อตฺถิ}

\vspace{\tipitakaparskip}
\csromancenter{นิทานวคฺคสํยุตฺตปาฬิ ตตฺถ อายติํ ชาติชรามรณํ.\csromanspacer{} ยตฺถ อตฺถิ อายติํ ชาติชรามรณํ, สโสกํ ตํ ภิกฺขเว สทรํ ส-อุปายาสนฺติ วทามิ.}
\prose{\csromanfolio{326}ผสฺเส เจ ภิกฺขเว อาหาเร ฯเปฯ.\csromanspacer{} มโนสญฺเจตนาย เจ ภิกฺขเว อาหาเร.\csromanspacer{} วิญฺญาเณ เจ ภิกฺขเว อาหาเร อตฺถิ ราโค อตฺถิ นนฺที อตฺถิ ตณฺหา, ปติฏฺฐิตํ ตตฺถ วิญฺญาณํ วิรูฬฺหํ.\csromanspacer{} ยตฺถ ปติฏฺฐิตํ วิญฺญาณํ วิรูฬฺหํ, อตฺถิ ตตฺถ นามรูปสฺส อวกฺกนฺติ.\csromanspacer{} ยตฺถ อตฺถิ นามรูปสฺส อวกฺกนฺติ, อตฺถิ ตตฺถ สงฺขารานํ วุทฺธิ.\csromanspacer{} ยตฺถ อตฺถิ สงฺขารานํ วุทฺธิ, อตฺถิ ตตฺถ อายติํ ปุนพฺภวาภินิพฺพตฺติ.\csromanspacer{} ยตฺถ อตฺถิ อายติํ ปุนพฺภวาพฺภินิพฺพตฺติ, อตฺถิ ตตฺถ อายติํ ชาติชรามรณํ.\csromanspacer{} ยตฺถ อตฺถิ อายติํ ชาติชรามรณํ, สโสกํ ตํ ภิกฺขเว สทรํ ส-อุปายาสนฺติ วทามิ.}

\prose{กพฬีกาเร เจ ภิกฺขเว อาหาเร นตฺถิ ราโค นตฺถิ นนฺที นตฺถิ ตณฺหา, อปฺปติฏฺฐิตํ ตตฺถ วิญฺญาณํ อวิรูฬฺหํ.\csromanspacer{} ยตฺถ อปฺปติฏฺฐิตํ วิญฺญาณํ อวิรูฬฺหํ, นตฺถิ ตตฺถ นามรูปสฺส อวกฺกนฺติ.\csromanspacer{} ยตฺถ นตฺถิ นามรูปสฺส อวกฺกนฺติ, นตฺถิ ตตฺถ สงฺขารานํ วุทฺธิ.\csromanspacer{} ยตฺถ นตฺถิ สงฺขารานํ วุทฺธิ, นตฺถิ ตตฺถ อายติํ ปุนพฺภวาภินิพฺพตฺติ.\csromanspacer{} ยตฺถ นตฺถิ อายติํ ปุนพฺภวาภินิพฺพตฺติ, นตฺถิ ตตฺถ อายติํ ชาติชรามรณํ.\csromanspacer{} ยตฺถ นตฺถิ อายติํ ชาติชรามรณํ, อโสกํ ตํ ภิกฺขเว อทรํ อนุปายาสนฺติ วทามิ.}

\prose{ผสฺเส เจ ภิกฺขเว อาหาเร ฯเปฯ.\csromanspacer{} มโนสญฺเจตนาย เจ ภิกฺขเว อาหาเร.\csromanspacer{} วิญฺญาเณ เจ ภิกฺขเว อาหาเร นตฺถิ ราโค นตฺถิ นนฺที นตฺถิ ตณฺหา, อปฺปติฏฺฐิตํ ตตฺถ วิญฺญาณํ อวิรูฬฺหํ.\csromanspacer{} ยตฺถ อปฺปติฏฺฐิตํ วิญฺญาณํ อวิรูฬฺหํ, นตฺถิ ตตฺถ นามรูปสฺส อวกฺกนฺติ.\csromanspacer{} ยตฺถ นตฺถิ นามรูปสฺส อวกฺกนฺติ, นตฺถิ ตตฺถ สงฺขารานํ วุทฺธิ.\csromanspacer{} ยตฺถ นตฺถิ สงฺขารานํ วุทฺธิ, นตฺถิ ตตฺถ อายติํ ปุนพฺภวาภินิพฺพตฺติ.\csromanspacer{} ยตฺถ นตฺถิ อายติํ ปุนพฺภวาภินิพฺพตฺติ, นตฺถิ ตตฺถ อายติํ ชาติชรามรณํ.\csromanspacer{} ยตฺถ นตฺถิ อายติํ ชาติชรามรณํ, อโสกํ ตํ ภิกฺขเว อทรํ อนุปายาสนฺติ วทามิ.}

\prose{เสยฺยถาปิ ภิกฺขเว กูฏาคารํ วา กูฏาคารสาลํ วา อุตฺตราย วา ทกฺขิณาย วา ปาจีนาย วา วาตปานา สูริเย อุคฺคจฺฉนฺเต วาตปาเนน รสฺมิ ปวิสิตฺวา กฺวาสฺส ปติฏฺฐิตา\footnote{กตฺถ ปติฏฺฐิตา (ก)} ติ.\csromanspacer{} ปจฺฉิมายํ ภนฺเต ภิตฺติยนฺติ.\csromanspacer{} ปจฺฉิมา เจ ภิกฺขเว ภิตฺติ นาสฺส กฺวาสฺส ปติฏฺฐิตาติ.\csromanspacer{} ปถวิยํ ภนฺเตติ.\csromanspacer{} ปถวี เจ}

\vspace{\tipitakaparskip}
\csromancenter{นิทานสํยุตฺต ภิกฺขเว นาสฺส กฺวาสฺส ปติฏฺฐิตาติ.\csromanspacer{} อาปสฺมิํ ภนฺเตติ.\csromanspacer{} อาโป เจ ภิกฺขเว นาสฺส กฺวาสฺส ปติฏฺฐิตาติ.\csromanspacer{} อปฺปติฏฺฐิตา ภนฺเตติ.\csromanspacer{} เอวเมว โข ภิกฺขเว กพฬีกาเร เจ อาหาเร นตฺถิ ราโค นตฺถิ นนฺที นตฺถิ ตณฺหา ฯเปฯ.}
\prose{\csromanfolio{327}ผสฺเส เจ ภิกฺขเว อาหาเร.\csromanspacer{} มโนสญฺเจตนาย เจ ภิกฺขเว อาหาเร.\csromanspacer{} วิญฺญาเณ เจ ภิกฺขเว อาหาเร นตฺถิ ราโค นตฺถิ นนฺที นตฺถิ ตณฺหา, อปฺปติฏฺฐิตํ ตตฺถ วิญฺญาณํ อวิรูฬฺหํ.\csromanspacer{} ยตฺถ อปฺปติฏฺฐิตํ วิญฺญาณํ อวิรูฬฺหํ, นตฺถิ ตตฺถ นามรูปสฺส อวกฺกนฺติ.\csromanspacer{} ยตฺถ นตฺถิ นามรูปสฺส อวกฺกนฺติ, นตฺถิ ตตฺถ สงฺขารานํ วุทฺธิ.\csromanspacer{} ยตฺถ นตฺถิ สงฺขารานํ วุทฺธิ, นตฺถิ ตตฺถ อายติํ ปุนพฺภวาภินิพฺพตฺติ.\csromanspacer{} ยตฺถ นตฺถิ อายติํ ปุนพฺภวาภินิพฺพตฺติ, นตฺถิ ตตฺถ อายติํ ชาติชรามรณํ.\csromanspacer{} ยตฺถ นตฺถิ อายติํ ชาติชรามรณํ, อโสกํ ตํ ภิกฺขเว อทรํ อนุปายาสนฺติ วทามีติ.\csromanspacer{}\csromanspacer{}จตุตฺถํ.}

\csromanmatikaanchor{mtk.74}
\csromantocmarknopage{section}{7. มหาวคฺค}{mtk.74}
\bookmark[dest={mtk.74},level=1]{7. มหาวคฺค}
\csromanheader{2}{5. นครสุตฺต}
\csromanmatikaanchor{mtk.75}
\csromantocmarkref{subsection}{1. อสฺสุตวาสุตฺต}{mtk.75}
\bookmark[dest={mtk.75},level=2]{1. อสฺสุตวาสุตฺต}
\proseitem{65}{สาวตฺถิยํ วิหรติ.\csromanspacer{} ปุพฺเพ เม ภิกฺขเว สมฺโพธา อนภิสมฺพุทฺธสฺส โพธิสตฺตสฺเสว สโต เอตทโหสิ “กิจฺฉา วตายํ โลโก อาปนฺโน ชายติ จ ชียติ จ มียติ จ จวติ จ อุปปชฺชติ จ, อถ จ ปนิมสฺส ทุกฺขสฺส นิสฺสรณํ นปฺปชานาติ ชรามรณสฺส.\csromanspacer{} กุทาสฺสุ นาม อิมสฺส ทุกฺขสฺส นิสฺสรณํ ปญฺญายิสฺสติ ชรามรณสฺสา”ติ.\csromanspacer{} ตสฺส มยฺหํ ภิกฺขเว เอตทโหสิ “กิมฺหิ นุ โข สติ ชรามรณํ โหติ, กิํปจฺจยา ชรามรณนฺ”ติ.\csromanspacer{} ตสฺส มยฺหํ ภิกฺขเว โยนิโส มนสิการา อหุ ปญฺญาย อภิสมโย “ชาติยา โข สติ ชรามรณํ โหติ, ชาติปจฺจยา ชรามรณนฺ”ติ.}

\prose{ตสฺส มยฺหํ ภิกฺขเว เอตทโหสิ “กิมฺหิ นุ โข สติ ชาติ โหติ ฯเปฯ ภโว โหติ.\csromanspacer{} อุปาทานํ โหติ.\csromanspacer{} ตณฺหา โหติ.\csromanspacer{} เวทนา โหติ.\csromanspacer{} ผสฺโส โหติ.\csromanspacer{} สฬายตนํ โหติ.\csromanspacer{} นามรูปํ โหติ.\csromanspacer{} กิํปจฺจยา นามรูปนฺ”ติ.\csromanspacer{} ตสฺส มยฺหํ ภิกฺขเว โยนิโส มนสิการา อหุ ปญฺญาย อภิสมโย “วิญฺญาเณ โข สติ นามรูปํ โหติ, วิญฺญาณปจฺจยา นามรูปนฺ”ติ.\csromanspacer{} ตสฺส มยฺหํ ภิกฺขเว เอตทโหสิ “กิมฺหิ นุ โข สติ วิญฺญาณํ โหติ กิํปจฺจยา วิญฺญาณนฺ”ติ.\csromanspacer{} ตสฺส มยฺหํ ภิกฺขเว โยนิโส มนสิการา อหุ ปญฺญาย อภิสมโย “นามรูเป โข สติ วิญฺญาณํ โหติ, นามรูปปจฺจยา วิญฺญาณนฺ”ติ.}

\gambhira{นิทานวคฺคสํยุตฺตปาฬิ}
\prose{\csromanfolio{328}ตสฺส มยฺหํ ภิกฺขเว เอตทโหสิ\csromandash{}ปจฺจุทาวตฺตติ โข อิทํ วิญฺญาณํ นามรูปมฺหา น ปรํ คจฺฉติ.\csromanspacer{} เอตฺตาวตา ชาเยถ วา ชีเยถ วา มีเยถ วา จเวถ วา อุปปชฺเชถ วา, ยทิทํ นามรูปปจฺจยา วิญฺญาณํ, วิญฺญาณปจฺจยา นามรูปํ, นามรูปปจฺจยา สฬายตนํ, สฬายตนปจฺจยา ผสฺโส ฯเปฯ เอวเมตสฺส เกวลสฺส ทุกฺขกฺขนฺธสฺส สมุทโย โหติ. “สมุทโย สมุทโย”ติ โข เม ภิกฺขเว ปุพฺเพ อนนุสฺสุเตสุ ธมฺเมสุ จกฺขุํ อุทปาทิ, ญาณํ อุทปาทิ, ปญฺญา อุทปาทิ, วิชฺชา อุทปาทิ, อาโลโก อุทปาทิ.}

\prose{ตสฺส มยฺหํ ภิกฺขเว เอตทโหสิ “กิมฺหิ นุ โข อสติ ชรามรณํ น โหติ, กิสฺส นิโรธา ชรามรณนิโรโธ”ติ.\csromanspacer{} ตสฺส มยฺหํ ภิกฺขเว โยนิโส มนสิการา อหุ ปญฺญาย อภิสมโย “ชาติยา โข อสติ ชรามรณํ น โหติ, ชาตินิโรธา ชรามรณนิโรโธ”ติ.\csromanspacer{} ตสฺส มยฺหํ ภิกฺขเว เอตทโหสิ “กิมฺหิ นุ โข อสติ ชาติ น โหติ ฯเปฯ ภโว น โหติ.\csromanspacer{} อุปาทานํ น โหติ.\csromanspacer{} ตณฺหา น โหติ.\csromanspacer{} เวทนา น โหติ.\csromanspacer{} ผสฺโส น โหติ.\csromanspacer{} สฬายตนํ น โหติ.\csromanspacer{} นามรูปํ น โหติ, กิสฺส นิโรธา นามรูปนิโรโธ”ติ.\csromanspacer{} ตสฺส มยฺหํ ภิกฺขเว โยนิโส มนสิการา อหุ ปญฺญาย อภิสมโย “วิญฺญาเณ โข อสติ นามรูปํ น โหติ, วิญฺญาณนิโรธา นามรูปนิโรโธ”ติ.}

\prose{ตสฺส มยฺหํ ภิกฺขเว เอตทโหสิ “กิมฺหิ นุ โข อสติ วิญฺญาณํ น โหติ, กิสฺส นิโรธา วิญฺญาณนิโรโธ”ติ.\csromanspacer{} ตสฺส มยฺหํ ภิกฺขเว โยนิโส มนสิการา อหุ ปญฺญาย อภิสมโย “นามรูเป โข อสติ วิญฺญาณํ น โหติ, นามรูปนิโรธา วิญฺญาณนิโรโธ”ติ.}

\prose{ตสฺส มยฺหํ ภิกฺขเว เอตทโหสิ\csromandash{}อธิคโต โข มฺยายํ มคฺโค โพธาย, ยทิทํ นามรูปนิโรธา วิญฺญาณนิโรโธ, วิญฺญาณนิโรธา นามรูปนิโรโธ, นามรูปนิโรธา สฬายตนนิโรโธ, สฬายตนนิโรธา ผสฺสนิโรโธ ฯเปฯ เอวเมตสฺส เกวลสฺส ทุกฺขกฺขนฺธสฺส นิโรโธ โหติ. “นิโรโธ นิโรโธ”ติ โข เม ภิกฺขเว ปุพฺเพ อนนุสฺสุเตสุ ธมฺเมสุ จกฺขุํ อุทปาทิ, ญาณํ อุทปาทิ, ปญฺญา อุทปาทิ, วิชฺชา อุทปาทิ, อาโลโก อุทปาทิ.}

\prose{เสยฺยถาปิ ภิกฺขเว ปุริโส อรญฺเญ ปวเน จรมาโน ปสฺเสยฺย ปุราณํ มคฺคํ ปุราณญฺชสํ ปุพฺพเกหิ มนุสฺเสหิ อนุยาตํ, โส}

\csromanmatikaanchor{mtk.76}
\csromantocmarkref{subsection}{2. ทุติย-อสฺสุตวาสุตฺต}{mtk.76}
\bookmark[dest={mtk.76},level=2]{2. ทุติย-อสฺสุตวาสุตฺต}
\vspace{\tipitakaparskip}
\csromancenter{นิทานสํยุตฺต ตมนุคจฺเฉยฺย, ตมนุคจฺฉนฺโต ปสฺเสยฺย ปุราณํ นครํ ปุราณํ ราชธานิํ, ปุพฺพเกหิ มนุสฺเสหิ อชฺฌาวุฏฺฐํ\footnote{อชฺฌาวุตฺถํ (สี, สฺยา, กํ, อิ)} อารามสมฺปนฺนํ วนสมฺปนฺนํ โปกฺขรณีสมฺปนฺนํ อุทฺธาปวนฺตํ\footnote{อุทฺทาปวนฺตํ (สี, สฺยา, กํ, อิ)} รมณียํ.\csromanspacer{} อถ โข โส ภิกฺขเว ปุริโส รญฺโญ วา ราชมหามตฺตสฺส วา อาโรเจยฺย “ยคฺเฆ ภนฺเต ชาเนยฺยาสิ ‘อหํ อทฺทสํ อรญฺเญ ปวเน จรมาโน ปุราณํ มคฺคํ ปุราณญฺชสํ ปุพฺพเกหิ มนุสฺเสหิ อนุยาตํ, ตมนุคจฺฉิํ, ตมนุคจฺฉนฺโต อทฺทสํ ปุราณํ นครํ ปุราณํ ราชธานิํ, ปุพฺพเกหิ มนุสฺเสหิ อชฺฌาวุฏฺฐํ อารามสมฺปนฺนํ วนสมฺปนฺนํ โปกฺขรณีสมฺปนฺนํ อุทฺธาปวนฺตํ รมณียํ, ตํ ภนฺเต นครํ มาเปหี’ติ”.\csromanspacer{} อถ โข โส ภิกฺขเว ราชา วา ราชมหามตฺโต วา ตํ นครํ มาเปยฺย, ตทสฺส นครํ อปเรน สมเยน อิทฺธํ เจว ผีตํ จ พาหุชญฺญํ อากิณฺณมนุสฺสํ วุทฺธิเวปุลฺลปฺปตฺตํ.\csromanspacer{} เอวเมว ขฺวาหํ ภิกฺขเว อทฺทสํ ปุราณํ มคฺคํ ปุราณญฺชสํ ปุพฺพเกหิ สมฺมาสมฺพุทฺเธหิ อนุยาตํ.}
\prose{\csromanfolio{329}กตโม จ โส ภิกฺขเว ปุราณมคฺโค ปุราณญฺชโส ปุพฺพเกหิ สมฺมาสมฺพุทฺเธหิ อนุยาโต, อยเมว อริโย อฏฺฐงฺคิโก มคฺโค.\csromanspacer{} เสยฺยถิทํ, สมฺมาทิฏฺฐิ ฯเปฯ สมฺมาสมาธิ.\csromanspacer{} อยํ โข โส ภิกฺขเว ปุราณมคฺโค ปุราณญฺชโส ปุพฺพเกหิ สมฺมาสมฺพุทฺเธหิ อนุยาโต.\csromanspacer{} ตมนุคจฺฉิํ, ตมนุคจฺฉนฺโต ชรามรณํ อพฺภญฺญาสิํ, ชรามรณสมุทยํ อพฺภญฺญาสิํ, ชรามรณนิโรธํ อพฺภญฺญาสิํ, ชรามรณนิโรธคามินิํ ปฏิปทํ อพฺภญฺญาสิํ.\csromanspacer{} ตมนุคจฺฉิํ, ตมนุคจฺฉนฺโต ชาติํ อพฺภญฺญาสิํ ฯเปฯ ภวํ อพฺภญฺญาสิํ.\csromanspacer{} อุปาทานํ อพฺภญฺญาสิํ.\csromanspacer{} ตณฺหํ อพฺภญฺญาสิํ.\csromanspacer{} เวทนํ อพฺภญฺญาสิํ.\csromanspacer{} ผสฺสํ อพฺภญฺญาสิํ.\csromanspacer{} สฬายตนํ อพฺภญฺญาสิํ.\csromanspacer{} นามรูปํ อพฺภญฺญาสิํ.\csromanspacer{} วิญฺญาณํ อพฺภญฺญาสิํ.\csromanspacer{} ตมนุคจฺฉิํ, ตมนุคจฺฉนฺโต สงฺขาเร อพฺภญฺญาสิํ, สงฺขารสมุทยํ อพฺภญฺญาสิํ, สงฺขารนิโรธํ อพฺภญฺญาสิํ, สงฺขารนิโรธคามินิํ ปฏิปทํ อพฺภญฺญาสิํ.\csromanspacer{} ตทภิญฺญาย อาจิกฺขิํ ภิกฺขูนํ ภิกฺขุนีนํ อุปาสกานํ อุปาสิกา- นํ, ตยิทํ ภิกฺขเว พฺรหฺมจริยํ อิทฺธํ เจว ผีตํ จ วิตฺถาริกํ พาหุชญฺญํ ปุถุภูตํ ยาว เทวมนุสฺเสหิ สุปฺปกาสิตนฺติ.\csromanspacer{}\csromanspacer{}ปญฺจมํ.}

\csromanheader{2}{6. สมฺมสสุตฺต}
\proseitem{66}{เอวํ เม สุตํ\csromandash{}เอกํ สมยํ ภควา กุรูสุ วิหรติ กมฺมาสทมฺมํ นาม กุรูนํ นิคโม.\csromanspacer{} ตตฺร โข ภควา ภิกฺขู อามนฺเตสิ “ภิกฺขโว”ติ.}

\vspace{\tipitakaparskip}
\csromancenter{นิทานวคฺคสํยุตฺตปาฬิ “ภทนฺเต”ติ เต ภิกฺขู ภควโต ปจฺจสฺโสสุํ.\csromanspacer{} ภควา เอตทโวจ “สมฺมสถ โน ตุเมฺห ภิกฺขเว อนฺตรํ สมฺมสนฺ”ติ\footnote{อนฺตรา สมฺมสนนฺติ (สี)}.\csromanspacer{} เอวํ วุตฺเต อญฺญตโร ภิกฺขุ ภควนฺตํ เอตทโวจ “อหํ โข ภนฺเต สมฺมสามิ อนฺตรํ สมฺมสนฺ”ติ.\csromanspacer{} ยถา กถํ ปน ตฺวํ ภิกฺขุ สมฺมสสิ อนฺตรํ สมฺมสนฺติ.\csromanspacer{} อถ โข โส ภิกฺขุ พฺยากาสิ.\csromanspacer{} ยถา โส ภิกฺขุ พฺยากาสิ, น โส ภิกฺขุ ภควโต จิตฺตํ อาราเธสิ.}
\prose{\csromanfolio{330}เอวํ วุตฺเต อายสฺมา อานนฺโท ภควนฺตํ เอตทโวจ “เอตสฺส ภควา กาโล, เอตสฺส สุคต กาโล, ยํ ภควา อนฺตรํ สมฺมสํ ภาเสยฺย, ภควโต สุตฺวา ภิกฺขู ธาเรสฺสนฺตี”ติ.\csromanspacer{} เตนหานนฺท สุณาถ สาธุกํ มนสิ กโรถ, ภาสิสฺสามีติ. “เอวํ ภนฺเต”ติ โข เต ภิกฺขู ภควโต ปจฺจสฺโสสุํ.\csromanspacer{} ภควา เอตทโวจ\csromandash{}}

\prose{อิธ ภิกฺขเว ภิกฺขุ สมฺมสมาโน สมฺมสติ อนฺตรํ สมฺมสํ\footnote{สมฺมสนํ (สี)} “ยํ โข อิทํ อเนกวิธํ นานปฺปการกํ ทุกฺขํ โลเก อุปฺปชฺชติ ชรามรณํ, อิทํ โข ทุกฺขํ กิํนิทานํ กิํสมุทยํ กิํชาติกํ กิํปภวํ, กิสฺมิํ สติ ชรามรณํ โหติ, กิสฺมิํ อสติ ชรามรณํ น โหตี”ติ.\csromanspacer{} โส สมฺมสมาโน เอวํ ชานาติ “ยํ โข อิทํ อเนกวิธํ นานปฺปการกํ ทุกฺขํ โลเก อุปฺปชฺชติ ชรามรณํ, อิทํ โข ทุกฺขํ อุปธินิทานํ อุปธิสมุทยํ อุปธิชาติกํ อุปธิปภวํ, อุปธิสฺมิํ สติ ชรามรณํ โหติ, อุปธิสฺมิํ อสติ ชรามรณํ น โหตี”ติ.\csromanspacer{} โส ชรามรณญฺจ ปชานาติ, ชรามรณสมุทยญฺจ ปชานาติ, ชรามรณนิโรธญฺจ ปชานาติ, ยา จ ชรามรณนิโรธสารุปฺปคามินี ปฏิปทา ตญฺจ ปชานาติ.\csromanspacer{} ตถาปฏิปนฺโน จ โหติ อนุธมฺมจารี.\csromanspacer{} อยํ วุจฺจติ ภิกฺขเว ภิกฺขุ สพฺพโส สมฺมา ทุกฺขกฺขยาย ปฏิปนฺโน ชรามรณนิโรธาย.}

\prose{อถาปรํ สมฺมสมาโน สมฺมสติ อนฺตรํ สมฺมสํ “อุปธิ ปนายํ กิํนิทาโน กิํสมุทโย กิํชาติโก กิํปภโว, กิสฺมิํ สติ อุปธิ โหติ, กิสฺมิํ อสติ อุปธิ น โหตี”ติ.\csromanspacer{} โส สมฺมสมาโน เอวํ ชานาติ “อุปธิ ตณฺหานิทาโน ตณฺหาสมุทโย ตณฺหาชาติโก ตณฺหาปภโว, ตณฺหาย สติ อุปธิ โหติ, ตณฺหาย อสติ อุปธิ น โหตี”ติ.\csromanspacer{} โส อุปธิญฺจ ปชานาติ, อุปธิสมุทยญฺจ ปชานาติ,}

\csromanmatikaanchor{mtk.77}
\csromantocmarkref{subsection}{3. ปุตฺตมํสูปมสุตฺต}{mtk.77}
\bookmark[dest={mtk.77},level=2]{3. ปุตฺตมํสูปมสุตฺต}
\vspace{\tipitakaparskip}
\csromancenter{นิทานสํยุตฺต อุปธินิโรธญฺจ ปชานาติ, ยา จ อุปธินิโรธสารุปฺปคามินี ปฏิปทา ตญฺจ ปชานาติ.\csromanspacer{} ตถาปฏิปนฺโน จ โหติ อนุธมฺมจารี.\csromanspacer{} อยํ วุจฺจติ ภิกฺขเว ภิกฺขุ สพฺพโส สมฺมา ทุกฺขกฺขยาย ปฏิปนฺโน อุปธินิโรธาย.}
\prose{\csromanfolio{331}อถาปรํ สมฺมสมาโน สมฺมสติ อนฺตรํ สมฺมสํ “ตณฺหา ปนายํ กตฺถ อุปฺปชฺชมานา อุปฺปชฺชติ, กตฺถ นิวิสมานา นิวิสตี”ติ.\csromanspacer{} โส สมฺมสมาโน เอวํ ชานาติ\csromandash{}ยํ โข โลเก ปิยรูปํ สาตรูปํ, เอตฺเถสา ตณฺหา อุปฺปชฺชมานา อุปฺปชฺชติ, เอตฺถ นิวิสมานา นิวิสติ.\csromanspacer{} กิญฺจ โลเก ปิยรูปํ สาตรูปํ, จกฺขุํ โลเก ปิยรูปํ สาตรูปํ, เอตฺเถสา ตณฺหา อุปฺปชฺชมานา อุปฺปชฺชติ, เอตฺถ นิวิสมานา นิวิสติ.\csromanspacer{} โสตํ โลเก ปิยรูปํ สาตรูปํ ฯเปฯ.\csromanspacer{} ฆานํ โลเก ปิยรูปํ สาตรูปํ.\csromanspacer{} ชิวฺหา โลเก ปิยรูปํ สาตรูปํ.\csromanspacer{} กาโย โลเก ปิยรูปํ สาตรูปํ.\csromanspacer{} มโน โลเก ปิยรูปํ สาตรูปํ, เอตฺเถสา ตณฺหา อุปฺปชฺชมานา อุปฺปชฺชติ, เอตฺถ นิวิสมานา นิวิสติ.}

\prose{เย หิ เกจิ ภิกฺขเว อตีตมทฺธานํ สมณา วา พฺราหฺมณา วา ยํ โลเก ปิยรูปํ สาตรูปํ, ตํ นิจฺจโต อทฺทกฺขุํ สุขโต อทฺทกฺขุํ อตฺตโต อทฺทกฺขุํ อาโรคฺยโต อทฺทกฺขุํ เขมโต อทฺทกฺขุํ, เต ตณฺหํ วฑฺเฒสุํ.\csromanspacer{} เย ตณฺหํ วฑฺเฒสุํ, เต อุปธิํ วฑฺเฒสุํ.\csromanspacer{} เย อุปธิํ วฑฺเฒสุํ, เต ทุกฺขํ วฑฺเฒสุํ.\csromanspacer{} เย ทุกฺขํ วฑฺเฒสุํ, เต น ปริมุจฺจิํสุ ชาติยา ชราย มรเณน โสเกหิ ปริเทเวหิ ทุกฺเขหิ โทมนสฺเสหิ อุปายาเสหิ, น ปริมุจฺจิํสุ ทุกฺขสฺมาติ วทามิ.}

\prose{เยปิ หิ เกจิ ภิกฺขเว อนาคตมทฺธานํ สมณา วา พฺราหฺมณา วา ยํ โลเก ปิยรูปํ สาตรูปํ, ตํ นิจฺจโต ทกฺขิสฺสนฺติ\footnote{ทกฺขินฺติ (สี)} สุขโต ทกฺขิสฺสนฺติ อตฺตโต ทกฺขิสฺสนฺติ อาโรคฺยโต ทกฺขิสฺสนฺติ เขมโต ทกฺขิสฺสนฺติ, เต ตณฺหํ วฑฺฒิสฺสนฺติ.\csromanspacer{} เย ตณฺหํ วฑฺฒิสฺสนฺติ, เต อุปธิํ วฑฺฒิสฺสนฺติ.\csromanspacer{} เย อุปธิํ วฑฺฒิสฺสนฺติ, เต ทุกฺขํ วฑฺฒิสฺสนฺติ.\csromanspacer{} เย ทุกฺขํ วฑฺฒิสฺสนฺติ, เต น ปริมุจฺจิสฺสนฺติ ชาติยา ชราย มรเณน โสเกหิ ปริเทเวหิ ทุกฺเขหิ โทมนสฺเสหิ อุปายาเสหิ, น ปริมุจฺจิสฺสนฺติ ทุกฺขสฺมาติ วทามิ.}

\prose{เยปิ หิ เกจิ ภิกฺขเว เอตรหิ สมณา วา พฺราหฺมณา วา ยํ โลเก ปิยรูปํ สาตรูปํ, ตํ นิจฺจโต ปสฺสนฺติ สุขโต ปสฺสนฺติ อตฺตโต ปสฺสนฺติ อาโรคฺยโต ปสฺสนฺติ เขมโต ปสฺสนฺติ, เต ตณฺหํ}

\vspace{\tipitakaparskip}
\csromancenter{นิทานวคฺคสํยุตฺตปาฬิ วฑฺเฒนฺติ.\csromanspacer{} เย ตณฺหํ วฑฺเฒนฺติ, เต อุปธิํ วฑฺเฒนฺติ.\csromanspacer{} เย อุปธิํ วฑฺเฒนฺติ, เต ทุกฺขํ วฑฺเฒนฺติ.\csromanspacer{} เย ทุกฺขํ วฑฺเฒนฺติ, เต น ปริมุจฺจนฺติ ชาติยา ชราย มรเณน โสเกหิ ปริเทเวหิ ทุกฺเขหิ โทมนสฺเสหิ อุปายาเสหิ, น ปริมุจฺจนฺติ ทุกฺขสฺมาติ วทามิ.}
\prose{\csromanfolio{332}เสยฺยถาปิ ภิกฺขเว อาปานียกํโส วณฺณสมฺปนฺโน คนฺธสมฺปนฺโน รสสมฺปนฺโน, โส จ โข วิเสน สํสฏฺโฐ.\csromanspacer{} อถ ปุริโส อาคจฺเฉยฺย ฆมฺมาภิตตฺโต ฆมฺมปเรโต กิลนฺโต ตสิโต ปิปาสิโต, ตเมนํ เอวํ วเทยฺยุํ “อยํ เต อมฺโภ ปุริส อาปานียกํโส วณฺณสมฺปนฺโน คนฺธสมฺปนฺโน รสสมฺปนฺโน, โส จ โข วิเสน สํสฏฺโฐ.\csromanspacer{} สเจ อากงฺขสิ ปิว, ปิวโต หิ โข ตํ ฉาเทสฺสติ วณฺเณนปิ คนฺเธนปิ รเสนปิ, ปิวิตฺวา จ ปน ตโตนิทานํ มรณํ วา นิคจฺฉสิ มรณมตฺตํ วา ทุกฺขนฺ”ติ.\csromanspacer{} โส ตํ อาปานียกํสํ สหสา อปฺปฏิสงฺขา ปิเวยฺย นปฺปฏินิสฺสชฺเชยฺย.\csromanspacer{} โส ตโตนิทานํ มรณํ วา นิคจฺเฉยฺย มรณมตฺตํ วา ทุกฺขํ.\csromanspacer{} เอวเมว โข ภิกฺขเว เย หิ เกจิ อตีตมทฺธานํ สมณา วา พฺราหฺมณา วา ยํ โลเก ปิยรูปํ ฯเปฯ.\csromanspacer{} อนาคตมทฺธานํ ฯเปฯ.\csromanspacer{} เอตรหิ สมณา วา พฺราหฺมณา วา ยํ โลเก ปิยรูปํ สาตรูปํ, ตํ นิจฺจโต ปสฺสนฺติ สุขโต ปสฺสนฺติ อตฺตโต ปสฺสนฺติ อาโรคฺยโต ปสฺสนฺติ เขมโต ปสฺสนฺติ, เต ตณฺหํ วฑฺเฒนฺติ.\csromanspacer{} เย ตณฺหํ วฑฺเฒนฺติ, เต อุปธิํ วฑฺเฒนฺติ.\csromanspacer{} เย อุปธิํ วฑฺเฒนฺติ, เต ทุกฺขํ วฑฺเฒนฺติ.\csromanspacer{} เย ทุกฺขํ วฑฺเฒนฺติ, เต น ปริมุจฺจนฺติ ชาติยา ชราย มรเณน โสเกหิ ปริเทเวหิ ทุกฺเขหิ โทมนสฺเสหิ อุปายาเสหิ, น ปริมุจฺจนฺติ ทุกฺขสฺมาติ วทามิ.}

\prose{เย จ โข เกจิ ภิกฺขเว อตีตมทฺธานํ สมณา วา พฺราหฺมณา วา ยํ โลเก ปิยรูปํ สาตรูปํ, ตํ อนิจฺจโต อทฺทกฺขุํ ทุกฺขโต อทฺทกฺขุํ อนตฺตโต อทฺทกฺขุํ โรคโต อทฺทกฺขุํ ภยโต อทฺทกฺขุํ, เต ตณฺหํ ปชหิํสุ.\csromanspacer{} เย ตณฺหํ ปชหิํสุ, เต อุปธิํ ปชหิํสุ.\csromanspacer{} เย อุปธิํ ปชหิํสุ, เต ทุกฺขํ ปชหิํสุ.\csromanspacer{} เย ทุกฺขํ ปชหิํสุ, เต ปริมุจฺจิํสุ ชาติยา ชราย มรเณน โสเกหิ ปริเทเวหิ ทุกฺเขหิ โทมนสฺเสหิ อุปายาเสหิ, ปริมุจฺจิํสุ ทุกฺขสฺมาติ วทามิ.}

\prose{เยปิ หิ เกจิ ภิกฺขเว อนาคตมทฺธานํ สมณา วา พฺราหฺมณา วา ยํ โลเก ปิยรูปํ สาตรูปํ, ตํ อนิจฺจโต ทกฺขิสฺสนฺติ ทุกฺขโต ทกฺขิสฺสนฺติ}

\csromanmatikaanchor{mtk.78}
\csromantocmarkref{subsection}{4. อตฺถิราคสุตฺต}{mtk.78}
\bookmark[dest={mtk.78},level=2]{4. อตฺถิราคสุตฺต}
\vspace{\tipitakaparskip}
\csromancenter{นิทานสํยุตฺต อนตฺตโต ทกฺขิสฺสนฺติ โรคโต ทกฺขิสฺสนฺติ ภยโต ทกฺขิสฺสนฺติ, เต ตณฺหํ ปชหิสฺสนฺติ.\csromanspacer{} เย ตณฺหํ ปชหิสฺสนฺติ ฯเปฯ ปริมุจฺจิสฺสนฺติ ทุกฺขสฺมาติ วทามิ.}
\prose{\csromanfolio{333}เยปิ หิ เกจิ ภิกฺขเว เอตรหิ สมณา วา พฺราหฺมณา วา ยํ โลเก ปิยรูปํ สาตรูปํ, ตํ อนิจฺจโต ปสฺสนฺติ ทุกฺขโต ปสฺสนฺติ อนตฺตโต ปสฺสนฺติ โรคโต ปสฺสนฺติ ภยโต ปสฺสนฺติ, เต ตณฺหํ ปชหนฺติ.\csromanspacer{} เย ตณฺหํ ปชหนฺติ, เต อุปธิํ ปชหนฺติ.\csromanspacer{} เย อุปธิํ ปชหนฺติ, เต ทุกฺขํ ปชหนฺติ.\csromanspacer{} เย ทุกฺขํ ปชหนฺติ, เต ปริมุจฺจนฺติ ชาติยา ชราย มรเณน โสเกหิ ปริเทเวหิ ทุกฺเขหิ โทมนสฺเสหิ อุปายาเสหิ, ปริมุจฺจนฺติ ทุกฺขสฺมาติ วทามิ.}

\prose{เสยฺยถาปิ ภิกฺขเว อาปานียกํโส วณฺณสมฺปนฺโน คนฺธสมฺปนฺโน รสสมฺปนฺโน, โส จ โข วิเสน สํสฏฺโฐ.\csromanspacer{} อถ ปุริโส อาคจฺเฉยฺย ฆมฺมาภิตตฺโต ฆมฺมปเรโต กิลนฺโต ตสิโต ปิปาสิโต, ตเมนํ เอวํ วเทยฺยุํ “อยํ เต อมฺโภ ปุริส อาปานียกํโส วณฺณสมฺปนฺโน คนฺธสมฺปนฺโน รสสมฺปนฺโน, โส จ โข วิเสน สํสฏฺโฐ.\csromanspacer{} สเจ อากงฺขสิ ปิว, ปิวโต หิ โข ตํ ฉาเทสฺสติ วณฺเณนปิ คนฺเธนปิ รเสนปิ, ปิวิตฺวา จ ปน ตโตนิทานํ มรณํ วา นิคจฺฉสิ มรณมตฺตํ วา ทุกฺขนฺ”ติ.\csromanspacer{} อถ โข ภิกฺขเว ตสฺส ปุริสสฺส เอวมสฺส “สกฺกา โข เม อยํ สุราปิปาสิตา\footnote{สุราปิปาสา (?)} ปานีเยน วา วิเนตุํ ทธิมณฺฑเกน วา วิเนตุํ ภฏฺฐโลณิกาย\footnote{มฏฺฐโลณิกาย (สี, สฺยา, กํ, อิ)} วา วิเนตุํ โลณโสวีรเกน วา วิเนตุํ, น เตฺววาหํ ตํ ปิเวยฺยํ, ยํ มม อสฺส ทีฆรตฺตํ หิตาย สุขายา”ติ.\csromanspacer{} โส ตํ อาปานียกํสํ ปฏิสงฺขา น ปิเวยฺย ปฏินิสฺสชฺเชยฺย.\csromanspacer{} โส ตโตนิทานํ น มรณํ วา นิคจฺเฉยฺย มรณมตฺตํ วา ทุกฺขํ.\csromanspacer{} เอวเมว โข ภิกฺขเว เย หิ เกจิ อตีตมทฺธานํ สมณา วา พฺราหฺมณา วา ยํ โลเก ปิยรูปํ สาตรูปํ, ตํ อนิจฺจโต อทฺทกฺขุํ ทุกฺขโต อทฺทกฺขุํ อนตฺตโต อทฺทกฺขุํ โรคโต อทฺทกฺขุํ ภยโต อทฺทกฺขุํ, เต ตณฺหํ ปชหิํสุ.\csromanspacer{} เย ตณฺหํ ปชหิํสุ, เต อุปธิํ ปชหิํสุ.\csromanspacer{} เย อุปธิํ ปชหิํสุ, เต ทุกฺขํ ปชหิํสุ.\csromanspacer{} เย ทุกฺขํ ปชหิํสุ, เต ปริมุจฺจิํสุ ชาติยา ชราย มรเณน โสเกหิ ปริเทเวหิ ทุกฺเขหิ โทมนสฺเสหิ อุปายาเสหิ, ปริมุจฺจิํสุ ทุกฺขสฺมาติ วทามิ.}

\gambhira{นิทานวคฺคสํยุตฺตปาฬิ}
\prose{\csromanfolio{334}เยปิ หิ เกจิ ภิกฺขเว อนาคตมทฺธานํ ฯเปฯ.\csromanspacer{} เอตรหิ สมณา วา พฺราหฺมณา วา ยํ โลเก ปิยรูปํ สาตรูปํ, ตํ อนิจฺจโต ปสฺสนฺติ ทุกฺขโต ปสฺสนฺติ อนตฺตโต ปสฺสนฺติ โรคโต ปสฺสนฺติ ภยโต ปสฺสนฺติ, เต ตณฺหํ ปชหนฺติ.\csromanspacer{} เย ตณฺหํ ปชหนฺติ, เต อุปธิํ ปชหนฺติ.\csromanspacer{} เย อุปธิํ ปชหนฺติ, เต ทุกฺขํ ปชหนฺติ.\csromanspacer{} เย ทุกฺขํ ปชหนฺติ, เต ปริมุจฺจนฺติ ชาติยา ชราย มรเณน โสเกหิ ปริเทเวหิ ทุกฺเขหิ โทมนสฺเสหิ อุปายาเสหิ, ปริมุจฺจนฺติ ทุกฺขสฺมาติ วทามีติ.\csromanspacer{}\csromanspacer{}ฉฏฺฐํ.}

\csromanheader{2}{7. นฬกลาปีสุตฺต}
\proseitem{67}{เอกํ สมยํ อายสฺมา จ สาริปุตฺโต อายสฺมา จ มหาโกฏฺฐิโก\footnote{มหาโกฏฺฐิโต (สี, สฺยา, กํ, อิ)} พาราณสิยํ วิหรนฺติ อิสิปตเน มิคทาเย.\csromanspacer{} อถ โข อายสฺมา มหาโกฏฺฐิโก สายนฺหสมยํ ปฏิสลฺลานา วุฏฺฐิโต เยนายสฺมา สาริปุตฺโต เตนุปสงฺกมิ, อุปสงฺกมิตฺวา อายสฺมตา สาริปุตฺเตน สทฺธิํ สมฺโมทิ, สมฺโมทนียํ กถํ สารณียํ วีติสาเรตฺวา เอกมนฺตํ นิสีทิ, เอกมนฺตํ นิสินฺโน โข อายสฺมา มหาโกฏฺฐิโก อายสฺมนฺตํ สาริปุตฺตํ เอตทโวจ “กิํ นุ โข อาวุโส สาริปุตฺต สยํกตํ ชรามรณํ, ปรํกตํ ชรามรณํ, สยํกตญฺจ ปรํกตญฺจ ชรามรณํ, อุทาหุ อสยํการํ อปรํการํ อธิจฺจสมุปฺปนฺนํ ชรามรณนฺ”ติ.\csromanspacer{} น โข อาวุโส โกฏฺฐิก สยํกตํ ชรามรณํ, น ปรํกตํ ชรามรณํ, น สยํกตญฺจ ปรํกตญฺจ ชรามรณํ, นาปิ อสยํการํ อปรํการํ อธิจฺจสมุปฺปนฺนํ ชรามรณํ, อปิ จ ชาติปจฺจยา ชรามรณนฺติ.}

\prose{กิํ นุ โข อาวุโส สาริปุตฺต สยํกตา ชาติ, ปรํกตา ชาติ.\csromanspacer{} สยํกตา จ ปรํกตา จ ชาติ, อุทาหุ อสยํการา อปรํการา อธิจฺจสมุปฺปนฺนา ชาตีติ.\csromanspacer{} น โข อาวุโส โกฏฺฐิก สยํกตา ชาติ, น ปรํกตา ชาติ, น สยํกตา จ ปรํกตา จ ชาติ, นาปิ อสยํการา อปรํการา อธิจฺจสมุปฺปนฺนา ชาติ, อปิ จ ภวปจฺจยา ชาตีติ.}

\prose{กิํ นุ โข อาวุโส สาริปุตฺต สยํกโต ภโว ฯเปฯ สยํกตํ อุปาทานํ.\csromanspacer{} สยํกตา ตณฺหา.\csromanspacer{} สยํกตา เวทนา.\csromanspacer{} สยํกโต ผสฺโส.\csromanspacer{} สยํกตํ สฬายตนํ.\csromanspacer{} สยํกตํ นามรูปํ, ปรํกตํ นามรูปํ, สยํกตญฺจ ปรํกตญฺจ นามรูปํ, อุทาหุ อสยํการํ อปรํการํ อธิจฺจสมุปฺปนฺนํ}

\vspace{\tipitakaparskip}
\csromancenter{นิทานสํยุตฺต นามรูปนฺติ.\csromanspacer{} น โข อาวุโส โกฏฺฐิก สยํกตํ นามรูปํ, น ปรํกตํ นามรูปํ, น สยํกตญฺจ ปรํกตญฺจ นามรูปํ, นาปิ อสยํการํ อปรํการํ อธิจฺจสมุปฺปนฺนํ นามรูปํ, อปิ จ วิญฺญาณปจฺจยา นามรูปนฺติ.}
\prose{\csromanfolio{335}กิํ นุ โข อาวุโส สาริปุตฺต สยํกตํ วิญฺญาณํ, ปรํกตํ วิญฺญาณํ, สยํกตญฺจ ปรํกตญฺจ วิญฺญาณํ, อุทาหุ อสยํการํ อปรํการํ อธิจฺจสมุปฺปนฺนํ วิญฺญาณนฺติ.\csromanspacer{} น โข อาวุโส โกฏฺฐิก สยํกตํ วิญฺญาณํ, น ปรํกตํ วิญฺญาณํ, น สยํกตญฺจ ปรํกตญฺจ วิญฺญาณํ, นาปิ อสยํการํ อปรํการํ อธิจฺจสมุปฺปนฺนํ วิญฺญาณํ, อปิ จ นามรูปปจฺจยา วิญฺญาณนฺติ.}

\prose{อิทาเนว โข มยํ อายสฺมโต สาริปุตฺตสฺส ภาสิตํ เอวํ อาชานาม.\csromanspacer{} น ขฺวาวุโส โกฏฺฐิก สยํกตํ นามรูปํ, น ปรํกตํ นามรูปํ, น สยํกตญฺจ ปรํกตญฺจ นามรูปํ, นาปิ อสยํการํ อปรํการํ อธิจฺจสมุปฺปนฺนํ นามรูปํ, อปิ จ วิญฺญาณปจฺจยา นามรูปนฺติ.}

\csromanmatikaanchor{mtk.79}
\csromantocmarkref{subsection}{5. นครสุตฺต}{mtk.79}
\bookmark[dest={mtk.79},level=2]{5. นครสุตฺต}
\prose{อิทาเนว จ ปน มยํ อายสฺมโต สาริปุตฺตสฺส ภาสิตํ เอวํ อาชานาม.\csromanspacer{} น ขฺวาวุโส โกฏฺฐิก สยํกตํ วิญฺญาณํ, น ปรํกตํ วิญฺญาณํ, น สยํกตญฺจ ปรํกตญฺจ วิญฺญาณํ, นาปิ อสยํการํ อปรํการํ อธิจฺจสมุปฺปนฺนํ วิญฺญาณํ, อปิ จ นามรูปปจฺจยา วิญฺญาณนฺติ.}

\prose{ยถ กถํ ปนาวุโส สาริปุตฺต อิมสฺส ภาสิตสฺส อตฺโถ ทฏฺฐพฺโพติ.\csromanspacer{} เตนหาวุโส อุปมํ เต กริสฺสามิ, อุปมายปิเธกจฺเจ วิญฺญู ปุริสา ภาสิตสฺส อตฺถํ ชานนฺติ.\csromanspacer{} เสยฺยถาปิ อาวุโส เทฺว นฬกลาปิโย อญฺญมญฺญํ นิสฺสาย ติฏฺเฐยฺยุํ.\csromanspacer{} เอวเมว โข อาวุโส นามรูปปจฺจยา วิญฺญาณํ, วิญฺญาณปจฺจยา นามรูปํ, นามรูปปจฺจยา สฬายตนํ, สฬายตนปจฺจยา ผสฺโส ฯเปฯ เอวเมตสฺส เกวลสฺส ทุกฺขกฺขนฺธสฺส สมุทโย โหติ.\csromanspacer{} ตาสํ เจ อาวุโส นฬกลาปีนํ เอกํ อากฑฺเฒยฺย เอกา ปปเตยฺย, อปรํ เจ อากฑฺเฒยฺย อปรา ปปเตยฺย.\csromanspacer{} เอวเมว โข อาวุโส นามรูปนิโรธา วิญฺญาณนิโรโธ, วิญฺญาณนิโรธา นามรูปนิโรโธ, นามรูปนิโรธา สฬายตนนิโรโธ, สฬายตนนิโรธา ผสฺสนิโรโธ ฯเปฯ เอวเมตสฺส เกวลสฺส ทุกฺขกฺขนฺธสฺส นิโรโธ โหตีติ.}

\gambhira{นิทานวคฺคสํยุตฺตปาฬิ}
\prose{\csromanfolio{336}อจฺฉริยํ อาวุโส สาริปุตฺต, อพฺภุตํ อาวุโส สาริปุตฺต, ยาวสุภาสิตํ จิทํ อายสฺมตา สาริปุตฺเตน.\csromanspacer{} อิทญฺจ ปน มยํ อายสฺมโต สาริปุตฺตสฺส ภาสิตํ อิเมหิ ฉตฺติํสาย วตฺถูหิ อนุโมทาม\csromandash{}}

\prose{ชรามรณสฺส เจ อาวุโส ภิกฺขุ นิพฺพิทาย วิราคาย นิโรธาย ธมฺมํ เทเสติ, “ธมฺมกถิโก ภิกฺขู”ติ อลํ วจนาย.\csromanspacer{} ชรามรณสฺส เจ อาวุโส ภิกฺขุ นิพฺพิทาย วิราคาย นิโรธาย ปฏิปนฺโน โหติ, “ธมฺมานุธมฺมปฺปฏิปนฺโน ภิกฺขู”ติ อลํ วจนาย.\csromanspacer{} ชรามรณสฺส เจ อาวุโส ภิกฺขุ นิพฺพิทา วิราคา นิโรธา อนุปาทา วิมุตฺโต โหติ, “ทิฏฺฐธมฺมนิพฺพานปฺปตฺโต ภิกฺขู”ติ อลํ วจนาย.\csromanspacer{} ชาติยา เจ.\csromanspacer{} ภวสฺส เจ.\csromanspacer{} อุปาทานสฺส เจ.\csromanspacer{} ตณฺหาย เจ.\csromanspacer{} ผสฺสสฺส เจ.\csromanspacer{} สฬายตนสฺส เจ.\csromanspacer{} นามรูปสฺส เจ.\csromanspacer{} วิญฺญาณสฺส เจ.\csromanspacer{} สงฺขารานํ เจ.\csromanspacer{} อวิชฺชาย เจ อาวุโส ภิกฺขุ นิพฺพิทาย วิราคาย นิโรธาย ธมฺมํ เทเสติ, “ธมฺมกถิโก ภิกฺขู”ติ อลํ วจนาย.\csromanspacer{} อวิชฺชาย เจ อาวุโส ภิกฺขุ นิพฺพิทาย วิราคาย นิโรธาย ปฏิปนฺโน โหติ, “ธมฺมานุธมฺมปฺปฏิปนฺโน ภิกฺขู”ติ อลํ วจนาย.\csromanspacer{} อวิชฺชาย เจ อาวุโส ภิกฺขุ นิพฺพิทา วิราคา นิโรธา อนุปาทา วิมุตฺโต โหติ, “ทิฏฺฐธมฺมนิพฺพานปฺปตฺโต ภิกฺขู”ติ อลํ วจนายาติ.\csromanspacer{}\csromanspacer{}สตฺตมํ.}

\csromanheader{2}{8. โกสมฺพิสุตฺต}
\proseitem{68}{เอกํ สมยํ อายสฺมา จ มุสิโล\footnote{มูสิโล (สี), มุสีโล (อิ)} อายสฺมา จ ปวิฏฺโฐ\footnote{สวิฏฺโฐ (สี, อิ)} อายสฺมา จ นารโท อายสฺมา จ อานนฺโท โกสมฺพิยํ วิหรนฺติ โฆสิตาราเม.\csromanspacer{} อถ โข อายสฺมา ปวิฏฺโฐ อายสฺมนฺตํ มุสิลํ เอตทโวจ “อญฺญเตฺรว อาวุโส มุสิล สทฺธาย อญฺญตฺร รุจิยา อญฺญตฺร อนุสฺสวา อญฺญตฺร อาการปริวิตกฺกา อญฺญตฺร ทิฏฺฐินิชฺฌานกฺขนฺติยา อตฺถายสฺมโต มุสิลสฺส ปจฺจตฺตเมว ญาณํ ‘ชาติปจฺจยา ชรามรณนฺ’ติ”.\csromanspacer{} อญฺญเตฺรว อาวุโส ปวิฏฺฐ สทฺธาย อญฺญตฺร รุจิยา อญฺญตฺร อนุสฺสวา อญฺญตฺร อาการปริวิตกฺกา อญฺญตฺร ทิฏฺฐินิชฺฌานกฺขนฺติยา อหเมตํ ชานามิ อหเมตํ ปสฺสามิ “ชาติปจฺจยา ชรามรณนฺ”ติ.}

\prose{อญฺญเตฺรว อาวุโส มุสิล สทฺธาย อญฺญตฺร รุจิยา อญฺญตฺร อนุสฺสวา อญฺญตฺร อาการปริวิตกฺกา อญฺญตฺร ทิฏฺฐินิชฺฌานกฺขนฺติยา อตฺถายสฺมโต}

\vspace{\tipitakaparskip}
\csromancenter{นิทานสํยุตฺต มุสิลสฺส ปจฺจตฺตเมว ญาณํ “ภวปจฺจยา ชาตี”ติ ฯเปฯ “อุปาทานปจฺจยา ภโว”ติ. “ตณฺหาปจฺจยา อุปาทานนฺ”ติ. “เวทนาปจฺจยา ตณฺหา”ติ. “ผสฺสปจฺจยา เวทนา”ติ. “สฬายตนปจฺจยา ผสฺโส”ติ. “นามรูปปจฺจยา สฬายตนนฺ”ติ. “วิญฺญาณปจฺจยา นามรูปนฺ”ติ. “สงฺขารปจฺจยา วิญฺญาณนฺ”ติ. “อวิชฺชาปจฺจยา สงฺขารา”ติ.\csromanspacer{} อญฺญเตฺรว อาวุโส ปวิฏฺฐ สทฺธาย อญฺญตฺร รุจิยา อญฺญตฺร อนุสฺสวา อญฺญตฺร อาการปริวิตกฺกา อญฺญตฺร ทิฏฺฐินิชฺฌานกฺขนฺติยา อหเมตํ ชานามิ อหเมตํ ปสฺสามิ “อวิชฺชาปจฺจยา สงฺขารา”ติ.}
\prose{\csromanfolio{337}อญฺญเตฺรว อาวุโส มุสิล สทฺธาย อญฺญตฺร รุจิยา อญฺญตฺร อนุสฺสวา อญฺญตฺร อาการปริวิตกฺกา อญฺญตฺร ทิฏฺฐินิชฺฌานกฺขนฺติยา อตฺถายสฺมโต มุสิลสฺส ปจฺจตฺตเมว ญาณํ “ชาตินิโรธา ชรามรณนิโรโธ”ติ.\csromanspacer{} อญฺญเตฺรว อาวุโส ปวิฏฺฐ สทฺธาย อญฺญตฺร รุจิยา อญฺญตฺร อนุสฺสวา อญฺญตฺร อาการปริวิตกฺกา อญฺญตฺร ทิฏฺฐินิชฺฌานกฺขนฺติยา อหเมตํ ชานามิ อหเมตํ ปสฺสามิ “ชาตินิโรธา ชรามรณนิโรโธ”ติ.}

\prose{อญฺญเตฺรว อาวุโส มุสิล สทฺธาย อญฺญตฺร รุจิยา อญฺญตฺร อนุสฺสวา อญฺญตฺร อาการปริวิตกฺกา อญฺญตฺร ทิฏฺฐินิชฺฌานกฺขนฺติยา อตฺถายสฺมโต มุสิลสฺส ปจฺจตฺตเมว ญาณํ “ภวนิโรธา ชาตินิโรโธ”ติ ฯเปฯ “อุปาทานนิโรธา ภวนิโรโธ”ติ. “ตณฺหานิโรธา อุปาทานนิโรโธ”ติ. “เวทนานิโรธา ตณฺหานิโรโธ”ติ. “ผสฺสนิโรธา เวทนานิโรโธ”ติ. “สฬายตนนิโรธา ผสฺสนิโรโธ”ติ. “นามรูปนิโรธา สฬายตนนิโรโธ”ติ. “วิญฺญาณนิโรธา นามรูปนิโรโธ”ติ. “สงฺขารนิโรธา วิญฺญาณนิโรโธ”ติ. “อวิชฺชานิโรธา สงฺขารนิโรโธ”ติ.\csromanspacer{} อญฺญเตฺรว อาวุโส ปวิฏฺฐ สทฺธาย อญฺญตฺร รุจิยา อญฺญตฺร อนุสฺสวา อญฺญตฺร อาการปริวิตกฺกา อญฺญตฺร ทิฏฺฐินิชฺฌานกฺขนฺติยา อหเมตํ ชานามิ อหเมตํ ปสฺสามิ “อวิชฺชานิโรธา สงฺขารนิโรโธ”ติ.}

\csromanmatikaanchor{mtk.80}
\csromantocmarkref{subsection}{6. สมฺมสสุตฺต}{mtk.80}
\bookmark[dest={mtk.80},level=2]{6. สมฺมสสุตฺต}
\prose{อญฺญเตฺรว อาวุโส มุสิล สทฺธาย อญฺญตฺร รุจิยา อญฺญตฺร อนุสฺสวา อญฺญตฺร อาการปริวิตกฺกา อญฺญตฺร ทิฏฺฐินิชฺฌานกฺขนฺติยา อตฺถายสฺมโต มุสิลสฺส ปจฺจตฺตเมว ญาณํ “ภวนิโรโธ นิพฺพานนฺ”ติ.\csromanspacer{} อญฺญเตฺรว อาวุโส ปวิฏฺฐ สทฺธาย อญฺญตฺร รุจิยา อญฺญตฺร อนุสฺสวา}

\vspace{\tipitakaparskip}
\csromancenter{นิทานวคฺคสํยุตฺตปาฬิ อญฺญตฺร อาการปริวิตกฺกา อญฺญตฺร ทิฏฺฐินิชฺฌานกฺขนฺติยา อหเมตํ ชานามิ อหเมตํ ปสฺสามิ “ภวนิโรโธ นิพฺพานนฺ”ติ.}
\prose{\csromanfolio{338}เตนหายสฺมา มุสิโล อรหํ ขีณาสโวติ.\csromanspacer{} เอวํ วุตฺเต อายสฺมา มุสิโล ตุณฺหี อโหสิ.\csromanspacer{} อถ โข อายสฺมา นารโท อายสฺมนฺตํ ปวิฏฺฐํ เอตทโวจ “สาธาวุโส ปวิฏฺฐ อหํ เอตํ ปญฺหํ ลเภยฺยํ.\csromanspacer{} มํ เอตํ ปญฺหํ ปุจฺฉ, อหํ เต เอตํ ปญฺหํ พฺยากริสฺสามี”ติ.\csromanspacer{} ลภตายสฺมา นารโท เอตํ ปญฺหํ.\csromanspacer{} ปุจฺฉามหํ อายสฺมนฺตํ นารทํ เอตํ ปญฺหํ, พฺยากโรตุ จ เม อายสฺมา นารโท เอตํ ปญฺหํ.}

\prose{อญฺญเตฺรว อาวุโส นารท สทฺธาย อญฺญตฺร รุจิยา อญฺญตฺร อนุสฺสวา อญฺญตฺร อาการปริวตกฺกา อญฺญตฺร ทิฏฺฐินิชฺฌานกฺขนฺติยา อตฺถายสฺมโต นารทสฺส ปจฺจตฺตเมว ญาณํ “ชาติปจฺจยา ชรามรณนฺ”ติ.\csromanspacer{} อญฺญเตฺรว อาวุโส ปวิฏฺฐ สทฺธาย อญฺญตฺร รุจิยา อญฺญตฺร อนุสฺสวา อญฺญตฺร อาการปริวิตกฺกา อญฺญตฺร ทิฏฺฐินิชฺฌานกฺขนฺติยา อหเมตํ ชานามิ อหเมตํ ปสฺสามิ “ชาติปจฺจยา ชรามรณนฺ”ติ.}

\prose{อญฺญเตฺรว อาวุโส นารท สทฺธาย อญฺญตฺร รุจิยา อญฺญตฺร อนุสฺสวา อญฺญตฺร อาการปริวิตกฺกา อญฺญตฺร ทิฏฺฐินิชฺฌานกฺขนฺติยา อตฺถายสฺมโต นารทสฺส ปจฺจตฺตเมว ญาณํ “ภวปจฺจยา ชาติ” ฯเปฯ “อวิชฺชาปจฺจยา สงฺขารา”ติ.\csromanspacer{} อญฺญเตฺรว อาวุโส ปวิฏฺฐ สทฺธาย อญฺญตฺร รุจิยา อญฺญตฺร อนุสฺสวา อญฺญตฺร อาการปริวิตกฺกา อญฺญตฺร ทิฏฺฐินิชฺฌานกฺขนฺติยา อหเมตํ ชานามิ อหเมตํ ปสฺสามิ “อวิชฺชาปจฺจยา สงฺขารา”ติ.}

\prose{อญฺญเตฺรว อาวุโส นารท สทฺธาย อญฺญตฺร รุจิยา อญฺญตฺร อนุสฺสวา อญฺญตฺร อาการปริวิตกฺกา อญฺญตฺร ทิฏฺฐินิชฺฌานกฺขนฺติยา อตฺถายสฺมโต นารทสฺส ปจฺจตฺตเมว ญาณํ “ชาตินิโรธา ชรามรณนิโรโธ”ติ.\csromanspacer{} อญฺญเตฺรว อาวุโส ปวิฏฺฐ สทฺธาย อญฺญตฺร รุจิยา อญฺญตฺร อนุสฺสวา อญฺญตฺร อาการปริวิตกฺกา อญฺญตฺร ทิฏฺฐินิชฺฌานกฺขนฺติยา อหเมตํ ชานามิ อหเมตํ ปสฺสามิ “ชาตินิโรธา ชรามรณนิโรโธ”ติ.}

\prose{อญฺญเตฺรว อาวุโส นารท สทฺธาย อญฺญตฺร รุจิยา อญฺญตฺร อนุสฺสวา อญฺญตฺร อาการปริวิตกฺกา อญฺญตฺร ทิฏฺฐินิชฺฌานกฺขนฺติยา อตฺถายสฺมโต}

\vspace{\tipitakaparskip}
\csromancenter{นิทานสํยุตฺต นารทสฺส ปจฺจตฺตเมว ญาณํ “ภวนิโรธา ชาตินิโรโธ”ติ ฯเปฯ “อวิชฺชานิโรธา สงฺขารนิโรโธ”ติ.\csromanspacer{} อญฺญเตฺรว อาวุโส ปวิฏฺฐ สทฺธาย อญฺญตฺร รุจิยา อญฺญตฺร อนุสฺสวา อญฺญตฺร อาการปริวิตกฺกา อญฺญตฺร ทิฏฺฐินิชฺฌานกฺขนฺติยา อหเมตํ ชานามิ อหเมตํ ปสฺสามิ “อวิชฺชานิโรธา สงฺขารนิโรโธ”ติ.}
\prose{\csromanfolio{339}อญฺญเตฺรว อาวุโส นารท สทฺธาย อญฺญตฺร รุจิยา อญฺญตฺร อนุสฺสวา อญฺญตฺร อาการปริวิตกฺกา อญฺญตฺร ทิฏฺฐินิชฺฌานกฺขนฺติยา อตฺถายสฺมโต นารทสฺส ปจฺจตฺตเมว ญาณํ “ภวนิโรโธ นิพฺพานนฺ”ติ.\csromanspacer{} อญฺญเตฺรว อาวุโส ปวิฏฺฐ สทฺธาย อญฺญตฺร รุจิยา อญฺญตฺร อนุสฺสวา อญฺญตฺร อาการปริวิตกฺกา อญฺญตฺร ทิฏฺฐินิชฺฌานกฺขนฺติยา อหเมตํ ชานามิ อหเมตํ ปสฺสามิ “ภวนิโรโธ นิพฺพานนฺ”ติ.}

\prose{เตนหายสฺมา นารโท อรหํ ขีณาสโวติ. “ภวนิโรโธ นิพฺพานนฺ”ติ โข เม อาวุโส ยถาภูตํ สมฺมปฺปญฺญาย สุทิฏฺฐํ, น จมฺหิ อรหํ ขีณาสโว.\csromanspacer{} เสยฺยถาปิ อาวุโส กนฺตารมคฺเค อุทปาโน, ตตฺร เนวสฺส รชฺชุ น อุทกวารโก.\csromanspacer{} อถ ปุริโส อาคจฺเฉยฺย ฆมฺมาภิตตฺโต ฆมฺมปเรโต กิลนฺโต ตสิโต ปิปาสิโต, โส ตํ อุทปานํ โอโลเกยฺย.\csromanspacer{} ตสฺส “อุทกนฺ”ติ หิ โข ญาณํ อสฺส, น จ กาเยน ผุสิตฺวา วิหเรยฺย.\csromanspacer{} เอวเมว โข อาวุโส “ภวนิโรโธ นิพฺพานนฺ”ติ ยถาภูตํ สมฺมปฺปญฺญาย สุทิฏฺฐํ, น จมฺหิ อรหํ ขีณาสโวติ.}

\prose{เอวํ วุตฺเต อายสฺมา อานนฺโท อายสฺมนฺตํ ปวิฏฺฐํ เอตทโวจ “เอวํวาที\footnote{เอวํวาทิํ (?)} ตฺวํ อาวุโส ปวิฏฺฐ อายสฺมนฺตํ นารทํ กิํ วเทสี”ติ.\csromanspacer{} เอวํวาทาหํ อาวุโส อานนฺท อายสฺมนฺตํ นารทํ น กิญฺจิ วทามิ อญฺญตฺร กลฺยาณา อญฺญตฺร กุสลาติ.\csromanspacer{}\csromanspacer{}อฏฺฐมํ.}

\csromanheader{2}{9. อุปยนฺติสุตฺต}
\proseitem{69}{เอวํ เม สุตํ\csromandash{}เอกํ สมยํ ภควา สาวตฺถิยํ วิหรติ เชตวเน อนาถปิณฺฑิกสฺส อาราเม.\csromanspacer{} ตตฺร โข ฯเปฯ.\csromanspacer{} มหาสมุทฺโท ภิกฺขเว อุปยนฺโต มหานทิโย อุปยาเปติ, มหานทิโย อุปยนฺติโย}

\vspace{\tipitakaparskip}
\csromancenter{นิทานวคฺคสํยุตฺตปาฬิ กุนฺนทิโย อุปยาเปนฺติ, กุนฺนทิโย อุปยนฺติโย มหาโสพฺเภ อุปยาเปนฺติ, มหาโสพฺภา อุปยนฺตา กุโสพฺเภ อุปยาเปนฺติ.\csromanspacer{} เอวเมว โข ภิกฺขเว อวิชฺชา อุปยนฺตี สงฺขาเร อุปยาเปติ, สงฺขารา อุปยนฺตา วิญฺญาณํ อุปยาเปนฺติ, วิญฺญาณํ อุปยนฺตํ นามรูปํ อุปยาเปติ, นามรูปํ อุปยนฺตํ สฬายตนํ อุปยาเปติ, สฬายตนํ อุปยนฺตํ ผสฺสํ อุปยาเปติ, ผสฺโส อุปยนฺโต เวทนํ อุปยาเปติ, เวทนา อุปยนฺตี ตณฺหํ อุปยาเปติ, ตณฺหา อุปยนฺตี อุปาทานํ อุปยาเปติ, อุปาทานํ อุปยนฺตํ ภวํ อุปยาเปติ, ภโว อุปยนฺโต ชาติํ อุปยาเปติ, ชาติ อุปยนฺตี ชรามรณํ อุปยาเปติ.}
\prose{\csromanfolio{340}มหาสมุทฺโท ภิกฺขเว อปยนฺโต มหานทิโย อปยาเปติ, มหานทิโย อปยนฺติโย กุนฺนทิโย อปยาเปนฺติ, กุนฺนทิโย อปยนฺติโย มหาโสพฺเภ อปยาเปนฺติ, มหาโสพฺภา อปยนฺตา กุโสพฺเภ อปยาเปนฺติ.\csromanspacer{} เอวเมว โข ภิกฺขเว อวิชฺชา อปยนฺตี สงฺขาเร อปยาเปติ, สงฺขารา อปยนฺตา วิญฺญาณํ อปยาเปนฺติ, วิญฺญาณํ อปยนฺตํ นามรูปํ อปยาเปติ, นามรูปํ อปยนฺตํ สฬายตนํ อปยาเปติ, สฬายตนํ อปยนฺตํ ผสฺสํ อปยาเปติ, ผสฺโส อปยนฺโต เวทนํ อปยาเปติ, เวทนา อปยนฺตี ตณฺหํ อปยาเปติ.\csromanspacer{} ตณฺหา อปยนฺตี อุปาทานํ อปยาเปติ, อุปาทานํ อปยนฺตํ ภวํ อปยาเปติ, ภโว อปยนฺโต ชาติํ อปยาเปติ, ชาติ อปยนฺตี ชรามรณํ อปยาเปตีติ.\csromanspacer{}\csromanspacer{}นวมํ.}

\csromanheader{2}{10. สุสิมสุตฺต}
\proseitem{70}{เอวํ เม สุตํ\csromandash{}เอกํ สมยํ ภควา ราชคเห วิหรติ เวฬุวเน กลนฺทกนิวาเป.\csromanspacer{} เตน โข ปน สมเยน ภควา สกฺกโต โหติ ครุกโต มานิโต ปูชิโต อปจิโต, ลาภี จีวรปิณฺฑปาตเสนาสนคิลานปฺปจฺจยเภสชฺชปริกฺขารานํ.\csromanspacer{} ภิกฺขุสํโฆปิ สกฺกโต โหติ ครุกโต มานิโต ปูชิโต อปจิโต, ลาภี จีวรปิณฺฑปาตเสนาสนคิลานปฺปจฺจยเภสชฺชปริกฺขารานํ.\csromanspacer{} อญฺญติตฺถิยา ปน ปริพฺพาชกา อสกฺกตา โหนฺติ อครุกตา อมานิตา อปูชิตา อนปจิตา, น ลาภิโน จีวรปิณฺฑปาตเสนาสนคิลานปฺปจฺจยเภสชฺชปริกฺขารานํ.}

\csromanheader{2}{1. นิทานสํยุตฺต}
\prose{\csromanfolio{341}เตน โข ปน สมเยน สุสิโม\footnote{สุสีโม (สี, ก)} ปริพฺพาชโก ราชคเห ปฏิวสติ มหติยา ปริพฺพาชกปริสาย สทฺธิํ.\csromanspacer{} อถ โข สุสิมสฺส ปริพฺพาชกสฺส ปริสา สุสิมํ ปริพฺพาชกํ เอตทโวจุํ “เอหิ ตฺวํ อาวุโส สุสิม สมเณ โคตเม พฺรหฺมจริยํ จร, ตฺวํ ธมฺมํ ปริยาปุณิตฺวา อเมฺห วาเจยฺยาสิ\footnote{วาเจสฺสสิ (อิ, ก)}, ตํ มยํ ธมฺมํ ปริยาปุณิตฺวา คิหีนํ ภาสิสฺสาม.\csromanspacer{} เอวํ มยมฺปิ สกฺกตา ภวิสฺสาม ครุกตา มานิตา ปูชิตา อปจิตา, ลาภิโน จีวรปิณฺฑปาตเสนาสนคิลานปฺปจฺจยเภสชฺชปริกฺขารานนฺ”ติ. “เอวมาวุโส”ติ โข สุสิโม ปริพฺพาชโก สกาย ปริสาย ปฏิสฺสุณิตฺวา เยนายสฺมา อานนฺโท เตนุปสงฺกมิ, อุปสงฺกมิตฺวา อายสฺมตา อานนฺเทน สทฺธิํ สมฺโมทิ, สมฺโมทนียํ กถํ สารณียํ วีติสาเรตฺวา เอกมนฺตํ นิสีทิ, เอกมนฺตํ นิสินฺโน โข สุสิโม ปริพฺพาชโก อายสฺมนฺตํ อานนฺทํ เอตทโวจ “อิจฺฉามหํ อาวุโส อานนฺท อิมสฺมิํ ธมฺมวินเย พฺรหฺมจริยํ จริตุนฺ”ติ.}

\prose{อถ โข อายสฺมา อานนฺโท สุสิมํ ปริพฺพาชกํ อาทาย เยน ภควา เตนุปสงฺกมิ, อุปสงฺกมิตฺวา ภควนฺตํ อภิวาเทตฺวา เอกมนฺตํ นิสีทิ, เอกมนฺตํ นิสินฺโน โข อายสฺมา อานนฺโท ภควนฺตํ เอตทโวจ “อยํ ภนฺเต สุสิโม ปริพฺพาชโก เอวมาห ‘อิจฺฉามหํ อาวุโส อานนฺท อิมสฺมิํ ธมฺมวินเย พฺรหฺมจริยํ จริตุนฺ’ติ”.\csromanspacer{} เตนหานนฺท สุสิมํ ปพฺพาเชถาติ.\csromanspacer{} อลตฺถ โข สุสิโม ปริพฺพาชโก ภควโต สนฺติเก ปพฺพชฺชํ, อลตฺถ อุปสมฺปทํ.}

\csromanmatikaanchor{mtk.81}
\csromantocmarkref{subsection}{7. นฬกลาปีสุตฺต}{mtk.81}
\bookmark[dest={mtk.81},level=2]{7. นฬกลาปีสุตฺต}
\prose{เตน โข ปน สมเยน สมฺพหุเลหิ ภิกฺขูหิ ภควโต สนฺติเก อญฺญา พฺยากตา โหติ “ขีณา ชาติ, วุสิตํ พฺรหฺมจริยํ, กตํ กรณียํ, นาปรํ อิตฺถตฺตายาติ ปชานามา”ติ.\csromanspacer{} อสฺโสสิ โข อายสฺมา สุสิโม “สมฺพหุเลหิ กิร ภิกฺขูหิ ภควโต สนฺติเก อญฺญา พฺยากตา ‘ขีณา ชาติ, วุสิตํ พฺรหฺมจริยํ, กตํ กรณียํ, นาปรํ อิตฺถตฺตายาติ ปชานามา’ติ”.\csromanspacer{} อถ โข อายสฺมา สุสิโม เยน เต ภิกฺขู เตนุปสงฺกมิ, อุปสงฺกมิตฺวา เตหิ ภิกฺขูหิ สทฺธิํ สมฺโมทิ, สมฺโมทนียํ กถํ สารณียํ วีติสาเรตฺวา เอกมนฺตํ นิสีทิ, เอกมนฺตํ นิสินฺโน โข อายสฺมา สุสิโม เต ภิกฺขู เอตทโวจ “สจฺจํ}

\vspace{\tipitakaparskip}
\csromancenter{นิทานวคฺคสํยุตฺตปาฬิ กิรายสฺมนฺเตหิ ภควโต สนฺติเก อญฺญา พฺยากตา ‘ขีณา ชาติ, วุสิตํ พฺรหฺมจริยํ, กตํ กรณียํ, นาปรํ อิตฺถตฺตายาติ ปชานามา’ติ”.\csromanspacer{} เอวมาวุโสติ.}
\prose{\csromanfolio{342}อปิ ปน\footnote{อปิ นุ (สี, สฺยา, กํ) เอวมุปริปิ.} ตุเมฺห อายสฺมนฺโต เอวํ ชานนฺตา เอวํ ปสฺสนฺตา อเนกวิหิตํ อิทฺธิวิธํ ปจฺจนุโภถ\csromandash{}เอโกปิ หุตฺวา พหุธา โหถ, พหุธาปิ หุตฺวา เอโก โหถ, อาวิภาวํ ติโรภาวํ ติโรกุฏฺฏํ ติโรปาการํ ติโรปพฺพตํ อสชฺชมานา คจฺฉถ เสยฺยถาปิ อากาเส, ปถวิยาปิ อุมฺมุชฺชนิมุชฺชํ กโรถ เสยฺยถาปิ อุทเก, อุทเกปิ อภิชฺชมาเน คจฺฉถ เสยฺยถาปิ ปถวิยํ, อากาเสปิ ปลฺลงฺเกน กมถ เสยฺยถาปิ ปกฺขี สกุโณ, อิเมปิ จนฺทิมสูริเย เอวํมหิทฺธิเก เอวํมหานุภาเว ปาณินา ปริมสถ ปริมชฺชถ, ยาว พฺรหฺมโลกาปิ กาเยน วสํ วตฺเตถาติ.\csromanspacer{} โน เหตํ อาวุโส.}

\prose{อปิ ปน ตุเมฺห อายสฺมนฺโต เอวํ ชานนฺตา เอวํ ปสฺสนฺตา ทิพฺพาย โสตธาตุยา วิสุทฺธาย อติกฺกนฺตมานุสิกาย อุโภ สทฺเท สุณาถ ทิพฺเพ จ มานุเส จ เย ทูเร สนฺติเก จาติ.\csromanspacer{} โน เหตํ อาวุโส.}

\prose{อปิ ปน ตุเมฺห อายสฺมนฺโต เอวํ ชานนฺตา เอวํ ปสฺสนฺตา ปรสตฺตานํ ปรปุคฺคลานํ เจตสา เจโต ปริจฺจ ปชานาถ\csromandash{}สราคํ วา จิตฺตํ “สราคํ จิตฺตนฺ”ติ ปชานาถ, วีตราคํ วา จิตฺตํ “วีตราคํ จิตฺตนฺ”ติ ปชานาถ.\csromanspacer{} สโทสํ วา จิตฺตํ “สโทสํ จิตฺตนฺ”ติ ปชานาถ, วีตโทสํ วา จิตฺตํ “วีตโทสํ จิตฺตนฺ”ติ ปชานาถ.\csromanspacer{} สโมหํ วา จิตฺตํ “สโมหํ จิตฺตนฺ”ติ ปชานาถ, วีตโมหํ วา จิตฺตํ “วีตโมหํ จิตฺตนฺ”ติ ปชานาถ, สํขิตฺตํ วา จิตฺตํ “สํขิตฺตํ จิตฺตนฺ”ติ ปชานาถ, วิกฺขิตฺตํ วา จิตฺตํ “วิกฺขิตฺตํ จิตฺตนฺ”ติ ปชานาถ.\csromanspacer{} มหคฺคตํ วา จิตฺตํ “มหคฺคตํ จิตฺตนฺ”ติ ปชานาถ.\csromanspacer{} อมหคฺคตํ วา จิตฺตํ “อมหคฺคตํ จิตฺตนฺ”ติ ปชานาถ.\csromanspacer{} ส-อุตฺตรํ วา จิตฺตํ “ส-อุตฺตรํ จิตฺตนฺ”ติ ปชานาถ, อนุตฺตรํ วา จิตฺตํ “อนุตฺตรํ จิตฺตนฺ”ติ ปชานาถ.\csromanspacer{} สมาหิตํ วา จิตฺตํ “สมาหิตํ จิตฺตนฺ”ติ ปชานาถ, อสมาหิตํ วา จิตฺตํ “อสมาหิตํ จิตฺตนฺ”ติ ปชานาถ.\csromanspacer{} วิมุตฺตํ วา จิตฺตํ “วิมุตฺตํ จิตฺตนฺ”ติ}

\vspace{\tipitakaparskip}
\csromancenter{นิทานสํยุตฺต ปชานาถ, อวิมุตฺตํ วา จิตฺตํ “อวิมุตฺตํ จิตฺตนฺ”ติ ปชานาถาติ.\csromanspacer{} โน เหตํ อาวุโส.}
\prose{\csromanfolio{343}อปิ ปน ตุเมฺห อายสฺมนฺโต เอวํ ชานนฺตา เอวํ ปสฺสนฺตา อเนกวิหิตํ ปุพฺเพนิวาสํ อนุสฺสรถ\csromandash{}เสยฺยถิทํ, เอกมฺปิ ชาติํ เทฺวปิ ชาติโย ติสฺโสปิ ชาติโย จตสฺโสปิ ชาติโย ปญฺจปิ ชาติโย ทสปิ ชาติโย วีสมฺปิ ชาติโย ติํสมฺปิ ชาติโย จตฺตารีสมฺปิ ชาติโย ปญฺญาสมฺปิ ชาติโย ชาติสตมฺปิ ชาติสหสฺสมฺปิ ชาติสตสหสฺสมฺปิ อเนเกปิ สํวฏฺฏกปฺเป อเนเกปิ วิวฏฺฏกปฺเป อเนเกปิ สํวฏฺฏวิวฏฺฏกปฺเป “อมุตฺราสิํ เอวํนาโม เอวํโคตฺโต เอวํวณฺโณ เอวมาหาโร เอวํสุขทุกฺขปฏิสํเวที เอวมายุปริยนฺโต, โส ตโต จุโต อมุตฺร อุทปาทิํ.\csromanspacer{} ตตฺราปาสิํ, เอวํนาโม เอวํโคตฺโต เอวํวณฺโณ เอวมาหาโร เอวํสุขทุกฺขปฏิสํเวที เอวมายุปริยนฺโต, โส ตโต จุโต อิธูปปนฺโน”ติ.\csromanspacer{} อิติ สาการํ ส-อุทฺเทสํ อเนกวิหิตํ ปุพฺเพนิวาสํ อนุสฺสรถาติ.\csromanspacer{} โน เหตํ อาวุโส.}

\prose{อปิ ปน ตุเมฺห อายสฺมนฺโต เอวํ ชานนฺตา เอวํ ปสฺสนฺตา ทิพฺเพน จกฺขุนา วิสุทฺเธน อติกฺกนฺตมานุสเกน สตฺเต ปสฺสถ จวมาเน อุปปชฺชมาเน หีเน ปณีเต สุวณฺเณ ทุพฺพณฺเณ สุคเต ทุคฺคเต, ยถากมฺมูปเค สตฺเต ปชานาถ\csromandash{}“อิเม วต โภนฺโต สตฺตา กายทุจฺจริเตน สมนฺนาคตา วจีทุจฺจริเตน สมนฺนาคตา มโนทุจฺจริเตน สมนฺนาคตา, อริยานํ อุปวาทกา มิจฺฉาทิฏฺฐิกา มิจฺฉาทิฏฺฐิกมฺมสมาทานา, เต กายสฺส เภทา ปรํ มรณา อปายํ ทุคฺคติํ วินิปาตํ นิรยํ อุปปนฺนา.\csromanspacer{} อิเม วา ปน โภนฺโต สตฺตา กายสุจริเตน สมนฺนาคตา วจีสุจริเตน สมนฺนาคตา มโนสุจริเตน สมนฺนาคตา, อริยานํ อนุปวาทกา สมฺมาทิฏฺฐิกา สมฺมาทิฏฺฐิกมฺมสมาทานา, เต กายสฺส เภทา ปรํ มรณา สุคติํ สคฺคํ โลกํ อุปปนฺนา”ติ.\csromanspacer{} อิติ ทิพฺเพน จกฺขุนา วิสุทฺเธน อติกฺกนฺตมานุสเกน สตฺเต ปสฺสถ จวมาเน อุปปชฺชมาเน หีเน ปณีเต สุวณฺเณ ทุพฺพณฺเณ สุคเต ทุคฺคเต, ยถากมฺมูปเค สตฺเต ปชานาถาติ.\csromanspacer{} โน เหตํ อาวุโส.}

\prose{อปิ ปน ตุเมฺห อายสฺมนฺโต เอวํ ชานนฺตา เอวํ ปสฺสนฺตา เย เต สนฺตา วิโมกฺขา อติกฺกมฺม รูเป อารุปฺปา, เต กาเยน ผุสิตฺวา วิหรถาติ.\csromanspacer{} โน เหตํ อาวุโส.}

\gambhira{นิทานวคฺคสํยุตฺตปาฬิ}
\prose{\csromanfolio{344}เอตฺถ ทานิ อายสฺมนฺโต อิทญฺจ เวยฺยากรณํ อิเมสญฺจ ธมฺมานํ อสมาปตฺติ, อิทํ โน อาวุโส กถนฺติ.\csromanspacer{} ปญฺญาวิมุตฺตา โข มยํ อาวุโส สุสิมาติ.}

\prose{น ขฺวาหํ อิมสฺส อายสฺมนฺตานํ สํขิตฺเตน ภาสิตสฺส วิตฺถาเรน อตฺถํ อาชานามิ.\csromanspacer{} สาธุ เม อายสฺมนฺโต ตถา ภาสนฺตุ, ยถาหํ อิมสฺส อายสฺมนฺตานํ สํขิตฺเตน ภาสิตสฺส วิตฺถาเรน อตฺถํ อาชาเนยฺยนฺติ.\csromanspacer{} อาชาเนยฺยาสิ วา ตฺวํ อาวุโส สุสิม, น วา ตฺวํ อาชาเนยฺยาสิ.\csromanspacer{} อถ โข ปญฺญาวิมุตฺตา มยนฺติ.}

\csromanmatikaanchor{mtk.82}
\csromantocmarkref{subsection}{8. โกสมฺพิสุตฺต}{mtk.82}
\bookmark[dest={mtk.82},level=2]{8. โกสมฺพิสุตฺต}
\prose{อถ โข อายสฺมา สุสิโม อุฏฺฐายาสนา เยน ภควา เตนุปสงฺกมิ, อุปสงฺกมิตฺวา ภควนฺตํ อภิวาเทตฺวา เอกมนฺตํ นิสีทิ, เอกมนฺตํ นิสินฺโน โข อายสฺมา สุสิโม ยาวตโก เตหิ ภิกฺขูหิ สทฺธิํ อโหสิ กถาสลฺลาโป, ตํ สพฺพํ ภควโต อาโรเจสิ.\csromanspacer{} ปุพฺเพ โข สุสิม ธมฺมฏฺฐิติญาณํ, ปจฺฉา นิพฺพาเน ญาณนฺติ.}

\prose{น ขฺวาหํ ภนฺเต อิมสฺส ภควตา\footnote{ภควโต (อิ)} สํขิตฺเตน ภาสิตสฺส วิตฺถาเรน อตฺถํ อาชานามิ.\csromanspacer{} สาธุ เม ภนฺเต ภควา ตถา ภาสตุ, ยถาหํ อิมสฺส ภควตา1 สํขิตฺเตน ภาสิตสฺส วิตฺถาเรน อตฺถํ อาชาเนยฺยนฺติ.\csromanspacer{} อาชาเนยฺยาสิ วา ตฺวํ สุสิม, น วา ตฺวํ อาชาเนยฺยาสิ.\csromanspacer{} อถ โข ธมฺมฏฺฐิติญาณํ ปุพฺเพ, ปจฺฉา นิพฺพาเน ญาณํ.}

\prose{ตํ กิํ มญฺญสิ สุสิม, รูปํ นิจฺจํ วา อนิจฺจํ วาติ.\csromanspacer{} อนิจฺจํ ภนฺเต.\csromanspacer{} ยํ ปนานิจฺจํ, ทุกฺขํ วา ตํ สุขํ วาติ.\csromanspacer{} ทุกฺขํ ภนฺเต.\csromanspacer{} ยํ ปนานิจฺจํ ทุกฺขํ วิปริณามธมฺมํ, กลฺลํ นุ ตํ สมนุปสฺสิตุํ “เอตํ มม, เอโสหมสฺมิ, เอโส เม อตฺตา”ติ.\csromanspacer{} โน เหตํ ภนฺเต.\csromanspacer{} เวทนา นิจฺจา วา อนิจฺจา วาติ.\csromanspacer{} อนิจฺจา ภนฺเต.\csromanspacer{} ยํ ปนานิจฺจํ, ทุกฺขํ วา ตํ สุขํ วาติ.\csromanspacer{} ทุกฺขํ ภนฺเต.\csromanspacer{} ยํ ปนานิจฺจํ ทุกฺขํ วิปริณามธมฺมํ, กลฺลํ นุ ตํ สมนุปสฺสิตุํ “เอตํ มม เอโสหมสฺมิ, เอโส เม อตฺตา”ติ.\csromanspacer{} โน เหตํ ภนฺเต.\csromanspacer{} สญฺญา นิจฺจา วา อนิจฺจา วาติ.\csromanspacer{} อนิจฺจา ภนฺเต ฯเปฯ.\csromanspacer{} สงฺขารา นิจฺจา วา อนิจฺจา วาติ.\csromanspacer{} อนิจฺจา ภนฺเต.\csromanspacer{} ยํ ปนานิจฺจํ, ทุกฺขํ วา ตํ สุขํ วาติ.\csromanspacer{} ทุกฺขํ ภนฺเต.\csromanspacer{} ยํ ปนานิจฺจํ ทุกฺขํ วิปริณามธมฺมํ, กลฺลํ นุ ตํ สมนุปสฺสิตุํ “เอตํ มม, เอโสหมสฺมิ, เอโส เม อตฺตา”ติ.\csromanspacer{} โน เหตํ ภนฺเต.\csromanspacer{} วิญฺญาณํ นิจฺจํ วา อนิจฺจํ วาติ.\csromanspacer{} อนิจฺจํ ภนฺเต.\csromanspacer{} ยํ}

\vspace{\tipitakaparskip}
\csromancenter{นิทานสํยุตฺต ปนานิจฺจํ, ทุกฺขํ วา ตํ สุขํ วาติ.\csromanspacer{} ทุกฺขํ ภนฺเต.\csromanspacer{} ยํ ปนานิจฺจํ ทุกฺขํ วิปริณามธมฺมํ, กลฺลํ นุ ตํ สมนุปสฺสิตุํ “เอตํ มม, เอโสหมสฺมิ, เอโส เม อตฺตา”ติ.\csromanspacer{} โน เหตํ ภนฺเต.}
\prose{\csromanfolio{345}ตสฺมาติห สุสิม ยํ กิญฺจิ รูปํ อตีตานาคตปจฺจุปฺปนฺนํ อชฺฌตฺตํ วา พหิทฺธา วา โอฬาริกํ วา สุขุมํ วา หีนํ วา ปณีตํ วา ยํ ทูเร สนฺติเก วา, สพฺพํ รูปํ “เนตํ มม, เนโสหมสฺมิ, น เมโส อตฺตา”ติ เอวเมตํ ยถาภูตํ สมฺมปฺปญฺญาย ทฏฺฐพฺพํ.\csromanspacer{} ยา กาจิ เวทนา อตีตานาคตปจฺจุปฺปนฺนา อชฺฌตฺตํ วา พหิทฺธา วา โอฬาริกา วา สุขุมา วา หีนา วา ปณีตา วา ยา ทูเร สนฺติเก วา, สพฺพา เวทนา “เนตํ มม, เนโสหมสฺมิ, น เมโส อตฺตา”ติ เอวเมตํ ยถาภูตํ สมฺมปฺปญฺญาย ทฏฺฐพฺพํ.\csromanspacer{} ยา กาจิ สญฺญา ฯเปฯ.\csromanspacer{} เย เกจิ สงฺขารา อตีตานาคตปจฺจุปฺปนฺนา อชฺฌตฺตํ วา พหิทฺธา วา โอฬาริกา วา สุขุมา วา หีนา วา ปณีตา วา เย ทูเร สนฺติเก วา, สพฺเพ สงฺขารา “เนตํ มม, เนโสหมสฺมิ, น เมโส อตฺตา”ติ เอวเมตํ ยถาภูตํ สมฺมปฺปญฺญาย ทฏฺฐพฺพํ.\csromanspacer{} ยํ กิญฺจิ วิญฺญาณํ อตีตานาคตปจฺจุปฺปนฺนํ อชฺฌตฺตํ วา พหิทฺธา วา โอฬาริกํ วา สุขุมํ วา หีนํ วา ปณีตํ วา ยํ ทูเร สนฺติเก วา, สพฺพํ วิญฺญาณํ “เนตํ มม, เนโสหมสฺมิ, น เมโส อตฺตา”ติ เอวเมตํ ยถาภูตํ สมฺมปฺปญฺญาย ทฏฺฐพฺพํ.}

\prose{เอวํ ปสฺสํ สุสิม สุตวา อริยสาวโก รูปสฺมิมฺปิ นิพฺพินฺทติ, เวทนายปิ นิพฺพินฺทติ, สญฺญายปิ นิพฺพินฺทติ, สงฺขาเรสุปิ นิพฺพินฺทติ, วิญฺญาณสฺมิมฺปิ นิพฺพินฺทติ, นิพฺพินฺทํ วิรชฺชติ, วิราคา วิมุจฺจติ, วิมุตฺตสฺมิํ “วิมุตฺตมฺ”อิติ ญาณํ โหติ, “ขีณา ชาติ, วุสิตํ พฺรหฺมจริยํ, กตํ กรณียํ, นาปรํ อิตฺถตฺตายา”ติ ปชานาติ.}

\prose{“ชาติปจฺจยา ชรามรณนฺ”ติ สุสิม ปสฺสสีติ.\csromanspacer{} เอวํ ภนฺเต. “ภวปจฺจยา ชาตี”ติ สุสิม ปสฺสสีติ.\csromanspacer{} เอวํ ภนฺเต. “อุปาทานปจฺจยา ภโว”ติ สุสิม ปสฺสสีติ.\csromanspacer{} เอวํ ภนฺเต. “ตณฺหาปจฺจยา อุปาทานนฺ”ติ สุสิม ปสฺสสีติ.\csromanspacer{} เอวํ ภนฺเต. “เวทนาปจฺจยา ตณฺหา”ติ. “ผสฺสปจฺจยา เวทนา”ติ. “สฬายตนปจฺจยา ผสฺโส”ติ. “นามรูปปจฺจยา สฬายตนนฺ”ติ. “วิญฺญาณปจฺจยา นามรูปนฺ”ติ. “สงฺขารปจฺจยา วิญฺญาณนฺ”ติ. “อวิชฺชาปจฺจยา สงฺขารา”ติ สุสิม ปสฺสสีติ.\csromanspacer{} เอวํ ภนฺเต.}

\gambhira{นิทานวคฺคสํยุตฺตปาฬิ}
\prose{\csromanfolio{346}“ชาตินิโรธา ชรามรณนิโรโธ”ติ สุสิม ปสฺสสีติ.\csromanspacer{} เอวํ ภนฺเต. “ภวนิโรธา ชาตินิโรโธ”ติ สุสิม ปสฺสสีติ.\csromanspacer{} เอวํ ภนฺเต. “อุปาทานนิโรธา ภวนิโรโธ”ติ. “ตณฺหานิโรธา อุปาทานนิโรโธ”ติ. “เวทนานิโรธา ตณฺหานิโรโธ”ติ. “ผสฺสนิโรธา เวทนานิโรโธ”ติ. “สฬายตนนิโรธา ผสฺสนิโรโธ”ติ. “นามรูปนิโรธา สฬายตนนิโรโธ”ติ. “วิญฺญาณนิโรธา นามรูปนิโรโธ”ติ. “สงฺขารนิโรธา วิญฺญาณนิโรโธ”ติ. “อวิชฺชานิโรธา สงฺขารนิโรโธ”ติ สุสิม ปสฺสสีติ.\csromanspacer{} เอวํ ภนฺเต.}

\prose{อปิ ปน ตฺวํ สุสิม เอวํ ชานนฺโต เอวํ ปสฺสนฺโต อเนกวิหิตํ อิทฺธิวิธํ ปจฺจนุโภสิ\csromandash{}เอโกปิ หุตฺวา พหุธา โหสิ, พหุธาปิ หุตฺวา เอโก โหสิ, อาวิภาวํ ติโรภาวํ ติโรกุฏฺฏํ ติโรปาการํ ติโรปพฺพตํ อสชฺชมาโน คจฺฉสิ เสยฺยถาปิ อากาเส, ปถวิยาปิ อุมฺมุชฺชนิมุชฺชํ กโรสิ เสยฺยถาปิ อุทเก, อุทเกปิ อภิชฺชมาโน คจฺฉสิ เสยฺยถาปิ ปถวิยํ, อากาเสปิ ปลฺลงฺเกน กมสิ เสยฺยถาปิ ปกฺขี สกุโณ, อิเมปิ จนฺทิมสูริเย เอวํมหิทฺธิเก เอวํมหานุภาเว ปาณินา ปริมสสิ ปริมชฺชสิ, ยาว พฺรหฺมโลกาปิ กาเยน วสํ วตฺเตสีติ.\csromanspacer{} โน เหตํ ภนฺเต.}

\prose{อปิ ปน ตฺวํ สุสิม เอวํ ชานนฺโต เอวํ ปสฺสนฺโต ทิพฺพาย โสตธาตุยา วิสุทฺธาย อติกฺกนฺตมานุสิกาย อุโภ สทฺเท สุณสิ ทิพฺเพ จ มานุเส จ เย ทูเร สนฺติเก จาติ.\csromanspacer{} โน เหตํ ภนฺเต.}

\prose{อปิ ปน ตฺวํ สุสิม เอวํ ชานนฺโต เอวํ ปสฺสนฺโต ปรสตฺตานํ ปรปุคฺคลานํ เจตสา เจโต ปริจฺจ ปชานาสิ\csromandash{}สราคํ วา จิตฺตํ “สราคํ จิตฺตนฺ”ติ ปชานาสิ ฯเปฯ วิมุตฺตํ วา จิตฺตํ “วิมุตฺตํ จิตฺตนฺ”ติ ปชานาสีติ.\csromanspacer{} โน เหตํ ภนฺเต.}

\prose{อปิ ปน ตฺวํ สุสิม เอวํ ชานนฺโต เอวํ ปสฺสนฺโต อเนกวิหิตํ ปุพฺเพนิวาสํ อนุสฺสรสิ.\csromanspacer{} เสยฺยถิทํ, เอกมฺปิ ชาติํ ฯเปฯ.\csromanspacer{} อิติ สาการํ ส- อุทฺเทสํ อเนกวิหิตํ ปุพฺเพนิวาสํ อนุสฺสรสีติ.\csromanspacer{} โน เหตํ ภนฺเต.}

\prose{อปิ ปน ตฺวํ สุสิม เอวํ ชานนฺโต เอวํ ปสฺสนฺโต ทิพฺเพน จกฺขุนา วิสุทฺเธน อติกฺกนฺตมานุสเกน สตฺเต ปสฺสสิ จวมาเน ฯเปฯ ยถากมฺมูปเค สตฺเต ปชานาสีติ.\csromanspacer{} โน เหตํ ภนฺเต.}

\csromanheader{2}{1. นิทานสํยุตฺต}
\prose{\csromanfolio{347}อปิ ปน ตฺวํ สุสิม เอวํ ชานนฺโต เอวํ ปสฺสนฺโต เย เต สนฺตา วิโมกฺขา อติกฺกมฺม รูเป อารุปฺปา, เต กาเยน ผุสิตฺวา วิหรสีติ.\csromanspacer{} โน เหตํ ภนฺเต.}

\prose{เอตฺถ ทานิ สุสิม อิทญฺจ เวยฺยากรณํ อิเมสญฺจ ธมฺมานํ อสมาปตฺติ, อิทํ โน สุสิม กถนฺติ.}

\csromanmatikaanchor{mtk.83}
\csromantocmarkref{subsection}{9. อุปยนฺติสุตฺต}{mtk.83}
\bookmark[dest={mtk.83},level=2]{9. อุปยนฺติสุตฺต}
\prose{อถ โข อายสฺมา สุสิโม ภควโต ปาเทสุ สิรสา นิปติตฺวา ภควนฺตํ เอตทโวจ “อจฺจโย มํ ภนฺเต อจฺจคมา ยถาพาลํ ยถามูฬฺหํ ยถา-อกุสลํ, ยฺวาหํ เอวํ สฺวากฺขาเต ธมฺมวินเย ธมฺมตฺเถนโก ปพฺพชิโต, ตสฺส เม ภนฺเต ภควา อจฺจยํ อจฺจยโต ปฏิคฺคณฺหาตุ อายติํ สํวรายา”ติ.}

\prose{ตคฺฆ ตฺวํ สุสิม อจฺจโย อจฺจคมา ยถาพาลํ ยถามูฬฺหํ ยถา- อกุสลํ, โย ตฺวํ เอวํ สฺวากฺขาเต ธมฺมวินเย ธมฺมตฺเถนโก ปพฺพชิโต.\csromanspacer{} เสยฺยถาปิ สุสิม โจรํ อาคุจาริํ คเหตฺวา รญฺโญ ทสฺเสยฺยุํ “อยํ เต เทว โจโร อาคุจารี อิมสฺส ยํ อิจฺฉสิ, ตํ ทณฺฑํ ปเณหี”ติ.\csromanspacer{} ตเมนํ ราชา เอวํ วเทยฺย “คจฺฉถ โภ อิมํ ปุริสํ ทฬฺหาย รชฺชุยา ปจฺฉาพาหํ คาฬฺหพนฺธนํ พนฺธิตฺวา ขุรมุณฺฑํ กริตฺวา ขรสฺสเรน ปณเวน รถิยาย รถิยํ สิงฺฆาฏเกน สิงฺฆาฏกํ ปริเนตฺวา ทกฺขิเณน ทฺวาเรน นิกฺขาเมตฺวา ทกฺขิณโต นครสฺส สีสํ ฉินฺทถา”ติ.\csromanspacer{} ตเมนํ รญฺโญ ปุริสา ทฬฺหาย รชฺชุยา ปจฺฉาพาหํ คาฬฺหพนฺธนํ พนฺธิตฺวา ขุรมุณฺฑํ กริตฺวา ขรสฺสเรน ปณเวน รถิยาย รถิยํ สิงฺฆาฏเกน สิงฺฆาฏกํ ปริเนตฺวา ทกฺขิเณน ทฺวาเรน นิกฺขาเมตฺวา ทกฺขิณโต นครสฺส สีสํ ฉินฺเทยฺยุํ.\csromanspacer{} ตํ กิํ มญฺญสิ สุสิม, อปิ นุ โส ปุริโส ตโตนิทานํ ทุกฺขํ โทมนสฺสํ ปฏิสํเวทิเยถาติ.\csromanspacer{} เอวํ ภนฺเต.}

\prose{ยํ โข โส สุสิม ปุริโส ตโตนิทานํ ทุกฺขํ โทมนสฺสํ ปฏิสํเวทิเยถ\footnote{ปฏิสํเวทิเยถ วา, น วา ปฏิสํเวทิเยถ (ก)}.\csromanspacer{} ยา เอวํ สฺวากฺขาเต ธมฺมวินเย ธมฺมตฺเถนกสฺส ปพฺพชฺชา, อยํ ตโต ทุกฺขวิปากตรา จ กฏุกวิปากตรา จ, อปิ จ วินิปาตาย สํวตฺตติ.\csromanspacer{} ยโต จ โข ตฺวํ สุสิม อจฺจยํ อจฺจยโต ทิสฺวา ยถาธมฺมํ ปฏิกโรสิ, ตํ เต มยํ ปฏิคฺคณฺหาม, วุทฺธิ เหสา สุสิม อริยสฺส}

\vspace{\tipitakaparskip}
\csromancenter{นิทานวคฺคสํยุตฺตปาฬิ วินเย, โย อจฺจยํ อจฺจยโต ทิสฺวา ยถาธมฺมํ ปฏิกโรติ, อายติํ จ\footnote{อายติํ (สฺยา, กํ)} สํวรํ อาปชฺชตีติ.\csromanspacer{}\csromanspacer{}ทสมํ.}
\csromanmatikaanchor{mtk.84}
\csromantocmarkref{subsection}{10. สุสิมสุตฺต}{mtk.84}
\bookmark[dest={mtk.84},level=2]{10. สุสิมสุตฺต}
\csromancenter{มหาวคฺโค สตฺตโม.}
\titlehead{\textbf{ตสฺสุทฺทานํ}}
\csromanmatikaanchor{mtk.85}
\csromantocmarkref{section}{อุทฺทานคาถา}{mtk.85}
\bookmark[dest={mtk.85},level=1]{อุทฺทานคาถา}
\csromangathameasure{เทฺว อสฺสุตวตา วุตฺตา, ปุตฺตมํเสน จาปรํ. \\ อตฺถิราโค จ นครํ, สมฺมสํ นฬกลาปิยํ. \\ โกสมฺพี อุปยนฺติ จ, ทสโม สุสิเมน จาติ.}
\csromangathagroup{\csromangathastanza{\csromanfolio{348}เทฺว อสฺสุตวตา วุตฺตา, ปุตฺตมํเสน จาปรํ. \\ อตฺถิราโค จ นครํ, สมฺมสํ นฬกลาปิยํ. \\ โกสมฺพี อุปยนฺติ จ, ทสโม สุสิเมน จาติ\footnote{ทสโม วุตฺโต สุสีเมนาติ (สี)}.}}

\csromanheader{1}{8. สมณพฺราหฺมณวคฺค}
\csromanheader{2}{1. ชรามรณสุตฺต}
\proseitem{71}{เอวํ เม สุตํ\csromandash{}เอกํ สมยํ ภควา สาวตฺถิยํ วิหรติ เชตวเน อนาถปิณฺฑิกสฺส อาราเม.\csromanspacer{} ตตฺร โข ภควา ฯเปฯ.\csromanspacer{} เย หิ เกจิ ภิกฺขเว สมณา วา พฺราหฺมณา วา ชรามรณํ นปฺปชานนฺติ, ชรามรณสมุทยํ นปฺปชานนฺติ, ชรามรณนิโรธํ นปฺปชานนฺติ, ชรามรณนิโรธคามินิํ ปฏิปทํ นปฺปชานนฺติ.\csromanspacer{} น เม เต ภิกฺขเว สมณา วา พฺราหฺมณา วา สมเณสุ วา สมณสมฺมตา พฺราหฺมเณสุ วา พฺราหฺมณสมฺมตา, น จ ปน เต อายสฺมนฺโต สามญฺญตฺถํ วา พฺรหฺมญฺญตฺถํ วา ทิฏฺเฐว ธมฺเม สยํ อภิญฺญา สจฺฉิกตฺวา อุปสมฺปชฺช วิหรนฺติ.}

\prose{เย จ โข เกจิ ภิกฺขเว สมณา วา พฺราหฺมณา วา ชรามรณํ ปชานนฺติ ฯเปฯ ปฏิปทํ ปชานนฺติ.\csromanspacer{} เต โข เม ภิกฺขเว สมณา วา พฺราหฺมณา วา สมเณสุ เจว สมณสมฺมตา พฺราหฺมเณสุ จ พฺราหฺมณสมฺมตา, เต จ ปนายสฺมนฺโต สามญฺญตฺถญฺจ พฺรหฺมญฺญตฺถญฺจ ทิฏฺเฐว ธมฺเม สยํ อภิญฺญา สจฺฉิกตฺวา อุปสมฺปชฺช วิหรนฺตีติ. (สุตฺตนฺโต เอโก.) ปฐมํ.}

\csromanheader{2}{1. นิทานสํยุตฺต \textbf{2-11. ชาติสุตฺตาทิทสก}}
\csromanmatikaanchor{mtk.86}
\csromanmatikaanchor{mtk.89}
\csromantocmarknopage{section}{8. สมณพฺราหฺมณวคฺค}{mtk.86}
\bookmark[dest={mtk.86},level=1]{8. สมณพฺราหฺมณวคฺค}
\proseitem{72}{\csromanfolio{349}สาวตฺถิยํ วิหรติ.\csromanspacer{} ชาติํ นปฺปชานนฺติ ฯเปฯ.}

\prose{ภวํ นปฺปชานนฺติ ฯเปฯ.}

\prose{อุปาทานํ นปฺปชานนฺติ ฯเปฯ.}

\prose{ตณฺหํ นปฺปชานนฺติ ฯเปฯ.}

\prose{เวทนํ นปฺปชานนฺติ ฯเปฯ.}

\prose{ผสฺสํ นปฺปชานนฺติ ฯเปฯ.}

\prose{สฬายตนํ นปฺปชานนฺติ ฯเปฯ.}

\prose{นามรูปํ นปฺปชานนฺติ ฯเปฯ.}

\prose{วิญฺญาณํ นปฺปชานนฺติ ฯเปฯ.}

\prose{สงฺขาเร นปฺปชานนฺติ, สงฺขารสมุทยํ นปฺปชานนฺติ, สงฺขารนิโรธํ นปฺปชานนฺติ, สงฺขารนิโรธคามินิํ ปฏิปทํ นปฺปชานนฺติ ฯเปฯ ปชานนฺติ ฯเปฯ สยํ อภิญฺญา สจฺฉิกตฺวา อุปสมฺปชฺช วิหรนฺตีติ.\csromanspacer{}\csromanspacer{}เอกาทสมํ.}

\vspace{\tipitakaparskip}
\csromancenter{สมณพฺราหฺมณวคฺโค อฏฺฐโม.}
\titlehead{\textbf{ตสฺสุทฺทานํ}}
\csromangathameasure{ปจฺจเยกาทส วุตฺตา, จตุสจฺจวิภชฺชนา. \\ สมณพฺราหฺมณวคฺโค, นิทาเน ภวติ อฏฺฐโม.}
\csromangathagroup{\csromangathastanza{ปจฺจเยกาทส วุตฺตา, จตุสจฺจวิภชฺชนา. \\ สมณพฺราหฺมณวคฺโค, นิทาเน ภวติ อฏฺฐโม.}}

\csromanheader{2}{\textbf{วคฺคุทฺทานํ}}
\csromangathameasure{พุทฺโธ อาหาโร ทสพโล, กฬาโร คหปติปญฺจโม. \\ ทุกฺขวคฺโค มหาวคฺโค, อฏฺฐโม สมณพฺราหฺมโณติ.}
\csromangathagroup{\csromangathastanza{พุทฺโธ อาหาโร ทสพโล, กฬาโร คหปติปญฺจโม. \\ ทุกฺขวคฺโค มหาวคฺโค, อฏฺฐโม สมณพฺราหฺมโณติ.}}

\csromanheader{1}{9. อนฺตรเปยฺยาล}
\csromanheader{2}{1. สตฺถุสุตฺต}
\proseitem{73}{สาวตฺถิยํ วิหรติ.\csromanspacer{} ชรามรณํ ภิกฺขเว อชานตา อปสฺสตา ยถาภูตํ ชรามรเณ ยถาภูตํ ญาณาย สตฺถา ปริเยสิตพฺโพ, ชรามรณสมุทยํ อชานตา อปสฺสตา ยถาภูตํ ชรามรณสมุทเย}

\vspace{\tipitakaparskip}
\csromancenter{นิทานวคฺคสํยุตฺตปาฬิ ยถาภูตํ ญาณาย สตฺถา ปริเยสิตพฺโพ, ชรามรณนิโรธํ อชานตา อปสฺสตา ยถาภูตํ ชรามรณนิโรเธ ยถาภูตํ ญาณาย สตฺถา ปริเยสิตพฺโพ, ชรามรณนิโรธคามินิํ ปฏิปทํ อชานตา อปสฺสตา ยถาภูตํ ชรามรณนิโรธคามินิยา ปฏิปทาย ยถาภูตํ ญาณาย สตฺถา ปริเยสิตพฺโพติ. (สุตฺตนฺโต เอโก.) ปฐมํ. (สพฺเพสํ เปยฺยาโล เอวํ วิตฺถาเรตพฺโพ.)}
\csromanheader{2}{2-11. ทุติยสตฺถุสุตฺตาทิทสก}
\prose{\csromanfolio{350}(2) ชาติํ ภิกฺขเว อชานตา อปสฺสตา ยถาภูตํ ฯเปฯ.}

\prose{(3) ภวํ ภิกฺขเว อชานตา อปสฺสตา ยถาภูตํ ฯเปฯ.}

\prose{(4) อุปาทานํ ภิกฺขเว อชานตา อปสฺสตา ยถาภูตํ ฯเปฯ.}

\prose{(5) ตณฺหํ ภิกฺขเว อชานตา อปสฺสตา ยถาภูตํ ฯเปฯ.}

\prose{(6) เวทนํ ภิกฺขเว อชานตา อปสฺสตา ยถาภูตํ ฯเปฯ.}

\prose{(7) ผสฺสํ ภิกฺขเว อชานตา อปสฺสตา ยถาภูตํ ฯเปฯ.}

\prose{(8) สฬายตนํ ภิกฺขเว อชานตา อปสฺสตา ยถาภูตํ ฯเปฯ.}

\prose{(9) นามรูปํ ภิกฺขเว อชานตา อปสฺสตา ยถาภูตํ ฯเปฯ.}

\prose{(10) วิญฺญาณํ ภิกฺขเว อชานตา อปสฺสตา ยถาภูตํ ฯเปฯ.}

\prose{(11) สงฺขาเร ภิกฺขเว อชานตา อปสฺสตา ยถาภูตํ สงฺขาเรสุ ยถาภูตํ ญาณาย สตฺถา ปริเยสิตพฺโพ, สงฺขารสมุทยํ อชานตา อปสฺสตา ยถาภูตํ สงฺขารสมุทเย ยถาภูตํ ญาณาย สตฺถา ปริเยสิตพฺโพ, สงฺขารนิโรธํ อชานตา อปสฺสตา ยถาภูตํ สงฺขารนิโรเธ ยถาภูตํ ญาณาย สตฺถา ปริเยสิตพฺโพ, สงฺขารนิโรธคามินิํ ปฏิปทํ อชานตา อปสฺสตา ยถาภูตํ สงฺขารนิโรธคามินิยา ปฏิปทาย ยถาภูตํ ญาณาย สตฺถา ปริเยสิตพฺโพติ.\csromanspacer{}\csromanspacer{}เอกาทสมํ. (สพฺเพสํ จตุสจฺจิกํ กาตพฺพํ.)}

\csromanheader{2}{1. นิทานสํยุตฺต \textbf{2-12. สิกฺขาสุตฺตาทิเปยฺยาล-เอกาทสก}}
\csromanmatikaanchor{mtk.87}
\csromanmatikaanchor{mtk.94}
\csromantocmarkref{subsection}{1. ชรามรณสุตฺต}{mtk.87}
\bookmark[dest={mtk.87},level=2]{1. ชรามรณสุตฺต}
\prose{\csromanfolio{351}(2) ชรามรณํ ภิกฺขเว อชานตา อปสฺสตา ยถาภูตํ ชรามรเณ ยถาภูตํ ญาณาย สิกฺขา กรณียา.}

\vspace{\tipitakaparskip}
\csromancenter{(เปยฺยาโล.\csromanspacer{} จตุสจฺจิกํ กาตพฺพํ.)}
\prose{(3) ชรามรณํ ภิกฺขเว อชานตา ฯเปฯ.\csromanspacer{} โยโค กรณีโย ฯเปฯ.}

\prose{(4) ชรามรณํ ภิกฺขเว อชานตา ฯเปฯ.\csromanspacer{} ฉนฺโท กรณีโย ฯเปฯ.}

\prose{(5) ชรามรณํ ภิกฺขเว อชานตา ฯเปฯ.\csromanspacer{} อุสฺโสฬฺหี กรณียา ฯเปฯ.}

\prose{(6) ชรามรณํ ภิกฺขเว อชานาตา ฯเปฯ.\csromanspacer{} อปฺปฏิวานี กรณียา ฯเปฯ.}

\csromanmatikaanchor{mtk.88}
\csromantocmarkref{subsection}{2-11. ชาติสุตฺตาทิทสก}{mtk.88}
\bookmark[dest={mtk.88},level=2]{2-11. ชาติสุตฺตาทิทสก}
\csromantocmarkref{section}{อุทฺทานคาถา}{mtk.89}
\bookmark[dest={mtk.89},level=1]{อุทฺทานคาถา}
\prose{(7) ชรามรณํ ภิกฺขเว อชานตา ฯเปฯ.\csromanspacer{} อาตปฺปํ กรณียํ ฯเปฯ.}

\prose{(8) ชรามรณํ ภิกฺขเว อชานตา ฯเปฯ.\csromanspacer{} วีริยํ กรณียํ ฯเปฯ.}

\prose{(9) ชรามรณํ ภิกฺขเว อชานตา ฯเปฯ.\csromanspacer{} สาตจฺจํ กรณียํ ฯเปฯ.}

\prose{(10) ชรามรณํ ภิกฺขเว อชานตา ฯเปฯ.\csromanspacer{} สติ กรณียา ฯเปฯ.}

\prose{(11) ชรามรณํ ภิกฺขเว อชานตา ฯเปฯ.\csromanspacer{} สมฺปชญฺญํ กรณียํ ฯเปฯ.}

\prose{(12) ชรามรณํ ภิกฺขเว อชานตา ฯเปฯ.\csromanspacer{} อปฺปมาโท กรณีโย ฯเปฯ.}

\vspace{\tipitakaparskip}
\csromancenter{อนฺตรเปยฺยาโล นวโม.}
\titlehead{\textbf{ตสฺสุทฺทานํ}}
\csromangathameasure{สตฺถา สิกฺขา จ โยโค จ, ฉนฺโท อุสฺโสฬฺหิปญฺจมี. \\ อปฺปฏิวานิ ยาตปฺปํ, วีริยํ สาตจฺจมุจฺจติ.}
\csromangathagroup{\csromangathastanza{สตฺถา สิกฺขา จ โยโค จ, ฉนฺโท อุสฺโสฬฺหิปญฺจมี. \\ อปฺปฏิวานิ ยาตปฺปํ, วีริยํ สาตจฺจมุจฺจติ.}}

\nitthitam{สติ จ สมฺปชญฺญญฺจ, อปฺปมาเทน ทฺวาทสาติ. สุตฺตนฺตา อนฺตรเปยฺยาลา นิฏฺฐิตา.}
\csromangathameasure{ปเร เต ทฺวาทส โหนฺติ, สุตฺตา ทฺวตฺติํส สตานิ. \\ จตุสจฺเจน เต วุตฺตา, เปยฺยาล-อนฺตรมฺหิ เยติ.}
\csromangathagroup{\csromangathastanza{ปเร เต ทฺวาทส โหนฺติ, สุตฺตา ทฺวตฺติํส สตานิ. \\ จตุสจฺเจน เต วุตฺตา, เปยฺยาล-อนฺตรมฺหิ เยติ\footnote{เปยฺยาลา อนฺตรมฺหิ เยติ (สี, สฺยา, กํ)}.}}

\vspace{\tipitakaparskip}
\csromancenter{อนฺตรเปยฺยาเลสุ อุทฺทานํ สมตฺตํ.}
\csromancenter{\textbf{นิทานสํยุตฺตํ สมตฺตํ.}}
\chapterheadpageread{2. อภิสมยสํยุตฺต}
\csromanheader{1}{1. นขสิขาสุตฺต}
\proseitem{74}{\csromanfolio{352}เอวํ เม สุตํ\csromandash{}เอกํ สมยํ ภควา สาวตฺถิยํ วิหรติ เชตวเน อนาถปิณฺฑิกสฺส อาราเม.\csromanspacer{} อถ โข ภควา ปริตฺตํ นขสิขายํ ปํสุํ อาโรเปตฺวา ภิกฺขู อามนฺเตสิ “ตํ กิํ มญฺญถ ภิกฺขเว, กตมํ นุ โข พหุตรํ, โย วายํ\footnote{โย จายํ (สพฺพตฺถ) ทุติยสุตฺตาทีสุ ปน วาสทฺโทเยว ทิสฺสติ.} มยา ปริตฺโต นขสิขายํ ปํสุ อาโรปิโต, อยํ วา มหาปถวี”ติ.}

\csromanmatikaanchor{mtk.90}
\csromantocmarknopage{section}{9. อนฺตรเปยฺยาล}{mtk.90}
\bookmark[dest={mtk.90},level=1]{9. อนฺตรเปยฺยาล}
\prose{เอตเทว ภนฺเต พหุตรํ ยทิทํ มหาปถวี, อปฺปมตฺตโก ภควตา ปริตฺโต นขสิขายํ ปํสุ อาโรปิโต, เนว สติมํ กลํ อุเปติ, น สหสฺสิมํ กลํ อุเปติ, น สตสหสฺสิมํ กลํ อุเปติ มหาปถวิํ อุปนิธาย ภควตา ปริตฺโต นขสิขายํ ปํสุ อาโรปิโตติ.\csromanspacer{} เอวเมว โข ภิกฺขเว อริยสาวกสฺส ทิฏฺฐิสมฺปนฺนสฺส ปุคฺคลสฺส อภิสเมตาวิโน เอตเทว พหุตรํ ทุกฺขํ ยทิทํ ปริกฺขีณํ ปริยาทิณฺณํ, อปฺปมตฺตกํ อวสิฏฺฐํ, เนว สติมํ กลํ อุเปติ, น สหสฺสิมํ กลํ อุเปติ, น สตสหสฺสิมํ กลํ อุเปติ ปุริมํ ทุกฺขกฺขนฺธํ ปริกฺขีณํ ปริยาทิณฺณํ อุปนิธาย ยทิทํ สตฺตกฺขตฺตุํปรมตา.\csromanspacer{} เอวํ มหตฺถิโย โข ภิกฺขเว ธมฺมาภิสมโย, เอวํ มหตฺถิโย ธมฺมจกฺขุปฏิลาโภติ.\csromanspacer{}\csromanspacer{}ปฐมํ.}

\csromanmatikaanchor{mtk.91}
\csromantocmarkref{subsection}{1. สตฺถุสุตฺต}{mtk.91}
\bookmark[dest={mtk.91},level=2]{1. สตฺถุสุตฺต}
\csromanheader{1}{2. โปกฺขรณีสุตฺต}
\proseitem{75}{สาวตฺถิยํ วิหรติ.\csromanspacer{} เสยฺยถาปิ ภิกฺขเว โปกฺขรณี ปญฺญาสโยชนานิ อายาเมน ปญฺญาสโยชนานิ วิตฺถาเรน ปญฺญาสโยชนานิ อุพฺเพเธน ปุณฺณา อุทกสฺส สมติตฺติกา กากเปยฺยา, ตโต ปุริโส กุสคฺเคน อุทกํ อุทฺธเรยฺย.\csromanspacer{} ตํ กิํ มญฺญถ ภิกฺขเว, กตมํ นุ โข พหุตรํ, ยํ วา กุสคฺเคน อุทกํ อุพฺภตํ, ยํ วา โปกฺขรณิยา อุทกนฺติ.}

\prose{เอตเทว ภนฺเต พหุตรํ ยทิทํ โปกฺขรณิยา อุทกํ, อปฺปมตฺตกํ กุสคฺเคน อุทกํ อุพฺภตํ, เนว สติมํ กลํ อุเปติ, น สหสฺสิมํ กลํ}

\csromanmatikaanchor{mtk.92}
\csromantocmarkref{subsection}{2-11. ทุติยสตฺถุสุตฺตาทิทสก}{mtk.92}
\bookmark[dest={mtk.92},level=2]{2-11. ทุติยสตฺถุสุตฺตาทิทสก}
\vspace{\tipitakaparskip}
\csromancenter{อภิสมยสํยุตฺต อุเปติ, น สตสหสฺสิมํ กลํ อุเปติ โปกฺขรณิยา อุทกํ อุปนิธาย กุสคฺเคน อุทกํ อุพฺภตนฺติ.\csromanspacer{} เอวเมว โข ภิกฺขเว อริยสาวกสฺส ทิฏฺฐิสมฺปนฺนสฺส ปุคฺคลสฺส อภิสเมตาวิโน เอตเทว พหุตรํ ทุกฺขํ ยทิทํ ปริกฺขีณํ ปริยาทิณฺณํ, อปฺปมตฺตกํ อวสิฏฺฐํ, เนว สติมํ กลํ อุเปติ, น สหสฺสิมํ กลํ อุเปติ, น สตสหสฺสิมํ กลํ อุเปติ ปุริมํ ทุกฺขกฺขนฺธํ ปริกฺขีณํ ปริยาทิณฺณํ อุปนิธาย ยทิทํ สตฺตกฺขตฺตุํปรมตา.\csromanspacer{} เอวํ มหตฺถิโย โข ภิกฺขเว ธมฺมาภิสมโย, เอวํ มหตฺถิโย ธมฺมจกฺขุปฏิลาโภติ.\csromanspacer{}\csromanspacer{}ทุติยํ.}
\csromanheader{1}{3. สํเภชฺช-อุทกสุตฺต}
\proseitem{76}{\csromanfolio{353}สาวตฺถิยํ วิหรติ.\csromanspacer{} เสยฺยถาปิ ภิกฺขเว ยตฺถิมา มหานทิโย สํสนฺทนฺติ สเมนฺติ.\csromanspacer{} เสยฺยถิทํ, คงฺคา ยมุนา อจิรวตี สรภู มหี.\csromanspacer{} ตโต ปุริโส เทฺว วา ตีณิ วา อุทกผุสิตานิ อุทฺธเรยฺย.\csromanspacer{} ตํ กิํ มญฺญถ ภิกฺขเว, กตมํ นุ โข พหุตรํ, ยานิ วา เทฺว วา ตีณิ วา อุทกผุสิตานิ อุพฺภตานิ, ยํ วา สํเภชฺช-อุทกนฺติ.}

\prose{เอตเทว ภนฺเต พหุตรํ ยทิทํ สํเภชฺช-อุทกํ, อปฺปมตฺตกานิ เทฺว วา ตีณิ วา อุทกผุสิตานิ อุพฺภตานิ, เนว สติมํ กลํ อุเปนฺติ, น สหสฺสิมํ กลํ อุเปนฺติ, น สตสหสฺสิมํ กลํ อุเปนฺติ สํเภชฺช- อุทกํ อุปนิธาย เทฺว วา ตีณิ วา อุทกผุสิตานิ อุพฺภตานีติ.\csromanspacer{} เอวเมว โข ภิกฺขเว ฯเปฯ ธมฺมจกฺขุปฏิลาโภติ.\csromanspacer{}\csromanspacer{}ตติยํ.}

\csromanheader{1}{4. ทุติยสํเภชฺช-อุทกสุตฺต}
\proseitem{77}{สาวตฺถิยํ วิหรติ.\csromanspacer{} เสยฺยถาปิ ภิกฺขเว ยตฺถิมา มหานทิโย สํสนฺทนฺติ สเมนฺติ.\csromanspacer{} เสยฺยถิทํ, คงฺคา ยมุนา อจิรวตี สรภู มหี.\csromanspacer{} ตํ อุทกํ ปริกฺขยํ ปริยาทานํ คจฺเฉยฺย ฐเปตฺวา เทฺว วา ตีณิ วา อุทกผุสิตานิ.\csromanspacer{} ตํ กิํ มญฺญถ ภิกฺขเว, กตมํ นุ โข พหุตรํ, ยํ วา สํเภชฺช-อุทกํ ปริกฺขีณํ ปริยาทิณฺณํ, ยานิ วา เทฺว วา ตีณิ วา อุทกผุสิตานิ อวสิฏฺฐานีติ.}

\prose{เอตเทว ภนฺเต พหุตรํ สํเภชฺช-อุทกํ ยทิทํ ปริกฺขีณํ ปริยาทิณฺณํ, อปฺปมตฺตกานิ เทฺว วา ตีณิ วา อุทกผุสิตานิ อวสิฏฺฐานิ, เนว สติมํ}

\vspace{\tipitakaparskip}
\csromancenter{นิทานวคฺคสํยุตฺตปาฬิ กลํ อุเปนฺติ, น สหสฺสิมํ กลํ อุเปนฺติ, น สตสหสฺสิมํ กลํ อุเปนฺติ สํเภชฺช-อุทกํ ปริกฺขีณํ ปริยาทิณฺณํ อุปนิธาย เทฺว วา ตีณิ วา อุทกผุสิตานิ อวสิฏฺฐานีติ.\csromanspacer{} เอวเมว โข ภิกฺขเว ฯเปฯ ธมฺมจกฺขุปฏิลาโภติ.\csromanspacer{}\csromanspacer{}จตุตฺถํ.}
\csromanheader{1}{5. ปถวีสุตฺต}
\proseitem{78}{\csromanfolio{354}สาวตฺถิยํ วิหรติ.\csromanspacer{} เสยฺยถาปิ ภิกฺขเว ปุริโส มหาปถวิยา สตฺต โกลฏฺฐิมตฺติโย คุฬิกา อุปนิกฺขิเปยฺย.\csromanspacer{} ตํ กิํ มญฺญถ ภิกฺขเว, กตมํ นุ โข พหุตรํ, ยา วา สตฺต โกลฏฺฐิมตฺติโย คุฬิกา อุปนิกฺขิตฺตา, อยํ\footnote{ยา (สฺยา, ก)} วา มหาปถวีติ.}

\prose{เอตเทว ภนฺเต พหุตรํ ยทิทํ มหาปถวี, อปฺปมตฺติกา สตฺต โกลฏฺฐิมตฺติโย คุฬิกา อุปนิกฺขิตฺตา, เนว สติมํ กลํ อุเปนฺติ, น สหสฺสิมํ กลํ อุเปนฺติ, น สตสหสฺสิมํ กลํ อุเปนฺติ มหาปถวิํ อุปนิธาย สตฺต โกลฏฺฐิมตฺติโย คุฬิกา อุปนิกฺขิตฺตาติ.\csromanspacer{} เอวเมว โข ภิกฺขเว ฯเปฯ ธมฺมจกฺขุปฏิลาโภติ.\csromanspacer{}\csromanspacer{}ปญฺจมํ.}

\csromanmatikaanchor{mtk.93}
\csromantocmarkref{subsection}{2-12. สิกฺขาสุตฺตาทิเปยฺยาล-เอกาทสก}{mtk.93}
\bookmark[dest={mtk.93},level=2]{2-12. สิกฺขาสุตฺตาทิเปยฺยาล-เอกาทสก}
\csromantocmarkref{section}{อุทฺทานคาถา}{mtk.94}
\bookmark[dest={mtk.94},level=1]{อุทฺทานคาถา}
\csromanheader{1}{6. ทุติยปถวีสุตฺต}
\proseitem{79}{สาวตฺถิยํ วิหรติ.\csromanspacer{} เสยฺยถาปิ ภิกฺขเว มหาปถวี ปริกฺขยํ ปริยาทานํ คจฺเฉยฺย ฐเปตฺวา สตฺต โกลฏฺฐิมตฺติโย คุฬิกา.\csromanspacer{} ตํ กิํ มญฺญถ ภิกฺขเว, กตมํ นุ โข พหุตรํ, ยํ วา มหาปถวิยา ปริกฺขีณํ ปริยาทิณฺณํ, ยา วา สตฺต โกลฏฺฐิมตฺติโย คุฬิกา อวสิฏฺฐาติ.}

\prose{เอตเทว ภนฺเต พหุตรํ มหาปถวิยา ยทิทํ ปริกฺขีณํ ปริยาทิณฺณํ, อปฺปมตฺติกา สตฺต โกลฏฺฐิมตฺติโย คุฬิกา อวสิฏฺฐา, เนว สติมํ กลํ อุเปนฺติ, น สหสฺสิมํ กลํ อุเปนฺติ, น สตสหสฺสิมํ กลํ อุเปนฺติ มหาปถวิยา ปริกฺขีณํ ปริยาทิณฺณํ อุปนิธาย สตฺต โกลฏฺฐิมตฺติโย คุฬิกา อวสิฏฺฐาติ.\csromanspacer{} เอวเมว โข ภิกฺขเว ฯเปฯ ธมฺมจกฺขุปฏิลาโภติ.\csromanspacer{}\csromanspacer{}ฉฏฺฐํ.}

\csromanheader{1}{7. สมุทฺทสุตฺต}
\proseitem{80}{สาวตฺถิยํ วิหรติ.\csromanspacer{} เสยฺยถาปิ ภิกฺขเว ปุริโส มหาสมุทฺทโต เทฺว วา ตีณิ วา อุทกผุสิตานิ อุทฺธเรยฺย.\csromanspacer{} ตํ กิํ มญฺญถ}

\vspace{\tipitakaparskip}
\csromancenter{อภิสมยสํยุตฺต ภิกฺขเว, กตมํ นุ โข พหุตรํ ยานิ วา เทฺว วา ตีณิ วา อุทกผุสิตานิ อุพฺภตานิ, ยํ วา มหาสมุทฺเท อุทกนฺติ.}
\prose{\csromanfolio{355}เอตเทว ภนฺเต พหุตรํ ยทิทํ มหาสมุทฺเท อุทกํ, อปฺปมตฺตกานิ เทฺว วา ตีณิ วา อุทกผุสิตานิ อุพฺภตานิ, เนว สติมํ กลํ อุเปนฺติ, น สหสฺสิมํ กลํ อุเปนฺติ, น สตสหสฺสิมํ กลํ อุเปนฺติ มหาสมุทฺเท อุทกํ อุปนิธาย เทฺว วา ตีณิ วา อุทกผุสิตานิ อุพฺภตานีติ.\csromanspacer{} เอวเมว โข ภิกฺขเว ฯเปฯ ธมฺมจกฺขุปฏิลาโภติ.\csromanspacer{}\csromanspacer{}สตฺตมํ.}

\csromanheader{1}{8. ทุติยสมุทฺทสุตฺต}
\proseitem{81}{สาวตฺถิยํ วิหรติ.\csromanspacer{} เสยฺยถาปิ ภิกฺขเว มหาสมุทฺโท ปริกฺขยํ ปริยาทานํ คจฺเฉยฺย ฐเปตฺวา เทฺว วา ตีณิ วา อุทกผุสิตานิ.\csromanspacer{} ตํ กิํ มญฺญถ ภิกฺขเว, กตมํ นุ โข พหุตรํ, ยํ วา มหาสมุทฺเท อุทกํ ปริกฺขีณํ ปริยาทิณฺณํ, ยานิ วา เทฺว วา ตีณิ วา อุทกผุสิตานิ อวสิฏฺฐานีติ.}

\prose{เอตเทว ภนฺเต พหุตรํ มหาสมุทฺเท อุทกํ ยทิทํ ปริกฺขีณํ ปริยาทิณฺณํ, อปฺปมตฺตกานิ เทฺว วา ตีณิ วา อุทกผุสิตานิ อวสิฏฺฐานิ, เนว สติมํ กลํ อุเปนฺติ, น สหสฺสิมํ กลํ อุเปนฺติ, น สตสหสฺสิมํ กลํ อุเปนฺติ มหาสมุทฺเท อุทกํ ปริกฺขีณํ ปริยาทิณฺณํ อุปนิธาย เทฺว วา ตีณิ วา อุทกผุสิตานิ อวสิฏฺฐานีติ.\csromanspacer{} เอวเมว โข ภิกฺขเว ฯเปฯ ธมฺมจกฺขุปฏิลาโภติ.\csromanspacer{}\csromanspacer{}อฏฺฐมํ.}

\csromanheader{1}{9. ปพฺพตสุตฺต}
\proseitem{82}{สาวตฺถิยํ วิหรติ.\csromanspacer{} เสยฺยถาปิ ภิกฺขเว ปุริโส หิมวโต ปพฺพตราชสฺส สตฺต สาสปมตฺติโย ปาสาณสกฺขรา อุปนิกฺขิเปยฺย.\csromanspacer{} ตํ กิํ มญฺญถ ภิกฺขเว, กตมํ นุ โข พหุตรํ, ยา วา สตฺต สาสปมตฺติโย ปาสาณสกฺขรา อุปนิกฺขิตฺตา, โย วา หิมวา\footnote{อุปนิกฺขิตฺตา, หิมวา วา (สี)} ปพฺพตราชาติ.}

\prose{เอตเทว ภนฺเต พหุตรํ ยทิทํ หิมวา ปพฺพตราชา, อปฺปมตฺติกา สตฺต สาสปมตฺติโย ปาสาณสกฺขรา อุปนิกฺขิตฺตา, เนว สติมํ กลํ}

\vspace{\tipitakaparskip}
\csromancenter{นิทานวคฺคสํยุตฺตปาฬิ อุเปนฺติ, น สหสฺสิมํ กลํ อุเปนฺติ, น สตสหสฺสิมํ กลํ อุเปนฺติ หิมวนฺตํ ปพฺพตราชานํ อุปนิธาย สตฺต สาสปมตฺติโย ปาสาณสกฺขรา อุปนิกฺขิตฺตาติ.\csromanspacer{} เอวเมว โข ฯเปฯ ธมฺมจกฺขุปฏิลาโภติ.\csromanspacer{}\csromanspacer{}นวมํ.}
\csromanheader{1}{10. ทุติยปพฺพตสุตฺต}
\proseitem{83}{\csromanfolio{356}สาวตฺถิยํ วิหรติ.\csromanspacer{} เสยฺยถาปิ ภิกฺขเว หิมวา ปพฺพตราชา ปริกฺขยํ ปริยาทานํ คจฺเฉยฺย ฐเปตฺวา สตฺต สาสปมตฺติโย ปาสาณสกฺขรา.\csromanspacer{} ตํ กิํ มญฺญถ ภิกฺขเว, กตมํ นุ โข พหุตรํ, ยํ วา หิมวโต ปพฺพตราชสฺส ปริกฺขีณํ ปริยาทิณฺณํ, ยา วา สตฺต สาสปมตฺติโย ปาสาณสกฺขรา อวสิฏฺฐาติ.}

\prose{เอตเทว ภนฺเต พหุตรํ หิมวโต ปพฺพตราชสฺส ยทิทํ ปริกฺขีณํ ปริยาทิณฺณํ, อปฺปมตฺติกา สตฺต สาสปมตฺติโย ปาสาณสกฺขรา อวสิฏฺฐา, เนว สติมํ กลํ อุเปนฺติ, น สหสฺสิมํ กลํ อุเปนฺติ, น สตสหสฺสิมํ กลํ อุเปนฺติ หิมวโต ปพฺพตราชสฺส ปริกฺขีณํ ปริยาทิณฺณํ อุปนิธาย สตฺต สาสปมตฺติโย ปาสาณสกฺขรา อวสิฏฺฐาติ.}

\prose{เอวเมว โข ภิกฺขเว อริยสาวกสฺส ทิฏฺฐิสมฺปนฺนสฺส ปุคฺคลสฺส อภิสเมตาวิโน เอตเทว พหุตรํ ทุกฺขํ ยทิทํ ปริกฺขีณํ ปริยาทิณฺณํ, อปฺปมตฺตกํ อวสิฏฺฐํ, เนว สติมํ กลํ อุเปติ, น สหสฺสิมํ กลํ อุเปติ, น สตสหสฺสิมํ กลํ อุเปติ ปุริมํ ทุกฺขกฺขนฺธํ ปริกฺขีณํ ปริยาทิณฺณํ อุปนิธาย ยทิทํ สตฺตกฺขตฺตุํปรมตา.\csromanspacer{} เอวํ มหตฺถิโย โข ภิกฺขเว ธมฺมาภิสมโย, เอวํ มหตฺถิโย ธมฺมจกฺขุปฏิลาโภติ.\csromanspacer{}\csromanspacer{}ทสมํ.}

\csromanheader{1}{11. ตติยปพฺพตสุตฺต}
\proseitem{84}{สาวตฺถิยํ วิหรติ.\csromanspacer{} เสยฺยถาปิ ภิกฺขเว ปุริโส สิเนรุสฺส ปพฺพตราชสฺส สตฺต มุคฺคมตฺติโย ปาสาณสกฺขรา อุปนิกฺขิเปยฺย.\csromanspacer{} ตํ กิํ มญฺญถ ภิกฺขเว, กตมํ นุ โข พหุตรํ, ยา วา สตฺต มุคฺคมตฺติโย ปาสาณสกฺขรา อุปนิกฺขิตฺตา, โย วา สิเนรุ\footnote{อุปนิกฺขิตฺตา, สิเนรุ วา (สี)} ปพฺพตราชาติ.}

\csromanmatikaanchor{mtk.95}
\csromantocmarknopage{chapter}{2. อภิสมยสํยุตฺต}{mtk.95}
\bookmark[dest={mtk.95},level=0]{2. อภิสมยสํยุตฺต}
\prose{เอตเทว ภนฺเต พหุตรํ ยทิทํ สิเนรุ ปพฺพตราชา, อปฺปมตฺติกา สตฺต มุคฺคมตฺติโย ปาสาณสกฺขรา อุปนิกฺขิตฺตา, เนว สติมํ กลํ อุเปนฺติ, น}

\csromanmatikaanchor{mtk.96}
\csromantocmarkref{section}{1. นขสิขาสุตฺต}{mtk.96}
\bookmark[dest={mtk.96},level=1]{1. นขสิขาสุตฺต}
\vspace{\tipitakaparskip}
\csromancenter{อภิสมยสํยุตฺต สหสฺสิมํ กลํ อุเปนฺติ, น สตสหสฺสิมํ กลํ อุเปนฺติ สิเนรุํ ปพฺพตราชานํ อุปนิธาย สตฺต มุคฺคมตฺติโย ปาสาณสกฺขรา อุปนิกฺขิตฺตาติ.\csromanspacer{} เอวเมว โข ภิกฺขเว อริยสาวกสฺส ทิฏฺฐิสมฺปนฺนสฺส ปุคฺคลสฺส อธิคมํ อุปนิธาย อญฺญติตฺถิย สมณพฺราหฺมณปริพฺพาชกานํ อธิคโม เนว สติมํ กลํ อุเปติ, น สหสฺสิมํ กลํ อุเปติ, น สตสหสฺสิมํ กลํ อุเปติ.\csromanspacer{} เอวํ มหาธิคโม ภิกฺขเว ทิฏฺฐิสมฺปนฺโน ปุคฺคโล เอวํ มหาภิญฺโญติ.\csromanspacer{}\csromanspacer{}เอกาทสมํ.}
\csromancenter{\textbf{อภิสมยสํยุตฺตํ สมตฺตํ.}}
\titlehead{\textbf{ตสฺสุทฺทานํ}}
\csromanmatikaanchor{mtk.97}
\csromanmatikaanchor{mtk.107}
\csromantocmarkref{section}{2. โปกฺขรณีสุตฺต}{mtk.97}
\bookmark[dest={mtk.97},level=1]{2. โปกฺขรณีสุตฺต}
\csromangathameasure{นขสิขา โปกฺขรณี, สํเภชฺช-อุทเก จ เทฺว. \\ เทฺว ปถวี เทฺว สมุทฺทา, ตโย จ ปพฺพตูปมาติ.}
\csromangathagroup{\csromangathastanza{\csromanfolio{357}นขสิขา โปกฺขรณี, สํเภชฺช-อุทเก จ เทฺว. \\ เทฺว ปถวี เทฺว สมุทฺทา, ตโย จ ปพฺพตูปมาติ.}}

\chapterheadpageread{3. ธาตุสํยุตฺต}
\csromanheader{1}{1. นานตฺตวคฺค}
\csromanheader{2}{1. ธาตุนานตฺตสุตฺต}
\csromanmatikaanchor{mtk.98}
\csromantocmarkref{section}{3. สํเภชฺช-อุทกสุตฺต}{mtk.98}
\bookmark[dest={mtk.98},level=1]{3. สํเภชฺช-อุทกสุตฺต}
\proseitem{85}{\csromanfolio{358}สาวตฺถิยํ วิหรติ.\csromanspacer{} ธาตุนานตฺตํ โว ภิกฺขเว เทเสสฺสามิ, ตํ สุณาถ สาธุกํ มนสิ กโรถ, ภาสิสฺสามีติ. “เอวํ ภนฺเต”ติ โข เต ภิกฺขู ภควโต ปจฺจสฺโสสุํ.\csromanspacer{} ภควา เอตทโวจ\csromandash{}}

\prose{กตมญฺจ ภิกฺขเว ธาตุนานตฺตํ.\csromanspacer{} จกฺขุธาตุ รูปธาตุ จกฺขุวิญฺญาณธาตุ, โสตธาตุ สทฺทธาตุ โสตวิญฺญาณธาตุ, ฆานธาตุ คนฺธธาตุ ฆานวิญฺญาณธาตุ, ชิวฺหาธาตุ รสธาตุ ชิวฺหาวิญฺญาณธาตุ, กายธาตุ โผฏฺฐพฺพธาตุ กายวิญฺญาณธาตุ, มโนธาตุ ธมฺมธาตุ มโนวิญฺญาณธาตุ.\csromanspacer{} อิทํ วุจฺจติ ภิกฺขเว ธาตุนานตฺตนฺติ.\csromanspacer{}\csromanspacer{}ปฐมํ.}

\csromanheader{2}{2. ผสฺสนานตฺตสุตฺต}
\csromanmatikaanchor{mtk.99}
\csromantocmarkref{section}{4. ทุติยสํเภชฺช-อุทกสุตฺต}{mtk.99}
\bookmark[dest={mtk.99},level=1]{4. ทุติยสํเภชฺช-อุทกสุตฺต}
\proseitem{86}{สาวตฺถิยํ วิหรติ.\csromanspacer{} ธาตุนานตฺตํ ภิกฺขเว ปฏิจฺจ อุปฺปชฺชติ ผสฺสนานตฺตํ.\csromanspacer{} กตมญฺจ ภิกฺขเว ธาตุนานตฺตํ.\csromanspacer{} จกฺขุธาตุ โสตธาตุ ฆานธาตุ ชิวฺหาธาตุ กายธาตุ มโนธาตุ.\csromanspacer{} อิทํ วุจฺจติ ภิกฺขเว ธาตุนานตฺตํ.}

\prose{กถญฺจ ภิกฺขเว ธาตุนานตฺตํ ปฏิจฺจ อุปฺปชฺชติ ผสฺสนานตฺตํ.\csromanspacer{} จกฺขุธาตุํ ภิกฺขเว ปฏิจฺจ อุปฺปชฺชติ จกฺขุสมฺผสฺโส.\csromanspacer{} โสตธาตุํ ปฏิจฺจ.\csromanspacer{} ฆานธาตุํ ปฏิจฺจ.\csromanspacer{} ชิวฺหาธาตุํ ปฏิจฺจ.\csromanspacer{} กายธาตุํ ปฏิจฺจ.\csromanspacer{} มโนธาตุํ ปฏิจฺจ อุปฺปชฺชติ มโนสมฺผสฺโส.\csromanspacer{} เอวํ โข ภิกฺขเว ธาตุนานตฺตํ ปฏิจฺจ อุปฺปชฺชติ ผสฺสนานตฺตนฺติ.\csromanspacer{}\csromanspacer{}ทุติยํ.}

\csromanheader{2}{3. โนผสฺสนานตฺตสุตฺต}
\proseitem{87}{สาวตฺถิยํ วิหรติ.\csromanspacer{} ธาตุนานตฺตํ ภิกฺขเว ปฏิจฺจ อุปฺปชฺชติ ผสฺสนานตฺตํ, โน ผสฺสนานตฺตํ ปฏิจฺจ อุปฺปชฺชติ ธาตุนานตฺตํ.\csromanspacer{} กตมญฺจ ภิกฺขเว ธาตุนานตฺตํ.\csromanspacer{} จกฺกฺกุธาตุ ฯเปฯ มโนธาตุ.\csromanspacer{} อิทํ วุจฺจติ ภิกฺขเว ธาตุนานตฺตํ.}

\csromanmatikaanchor{mtk.100}
\csromantocmarkref{section}{5. ปถวีสุตฺต}{mtk.100}
\bookmark[dest={mtk.100},level=1]{5. ปถวีสุตฺต}
\csromanheader{2}{3. ธาตุสํยุตฺต}
\prose{\csromanfolio{359}กถญฺจ ภิกฺขเว ธาตุนานตฺตํ ปฏิจฺจ อุปฺปชฺชติ ผสฺสนานตฺตํ, โน ผสฺสนานตฺตํ ปฏิจฺจ อุปฺปชฺชติ ธาตุนานตฺตํ.\csromanspacer{} จกฺขุธาตุํ ภิกฺขเว ปฏิจฺจ อุปฺปชฺชติ จกฺขุสมฺผสฺโส, โน จกฺขุสมฺผสฺสํ ปฏิจฺจ อุปฺปชฺชติ จกฺขุธาตุ ฯเปฯ.\csromanspacer{} มโนธาตุํ ปฏิจฺจ อุปฺปชฺชติ มโนสมฺผสฺโส, โน มโนสมฺผสฺสํ ปฏิจฺจ อุปฺปชฺชติ มโนธาตุ.\csromanspacer{} เอวํ โข ภิกฺขเว ธาตุนานตฺตํ ปฏิจฺจ อุปฺปชฺชติ ผสฺสนานตฺตํ, โน ผสฺสนานตฺตํ ปฏิจฺจ อุปฺปชฺชติ ธาตุนานตฺตนฺติ.\csromanspacer{}\csromanspacer{}ตติยํ.}

\csromanheader{2}{4. เวทนานานตฺตสุตฺต}
\csromanmatikaanchor{mtk.101}
\csromantocmarkref{section}{6. ทุติยปถวีสุตฺต}{mtk.101}
\bookmark[dest={mtk.101},level=1]{6. ทุติยปถวีสุตฺต}
\proseitem{88}{สาวตฺถิยํ วิหรติ.\csromanspacer{} ธาตุนานตฺตํ ภิกฺขเว ปฏิจฺจ อุปฺปชฺชติ ผสฺสนานตฺตํ, ผสฺสนานตฺตํ ปฏิจฺจ อุปฺปชฺชติ เวทนานานตฺตํ.\csromanspacer{} กตมญฺจ ภิกฺขเว ธาตุนานตฺตํ.\csromanspacer{} จกฺขุธาตุ ฯเปฯ มโนธาตุ.\csromanspacer{} อิทํ วุจฺจติ ภิกฺขเว ธาตุนานตฺตํ.}

\prose{กถญฺจ ภิกฺขเว ธาตุนานตฺตํ ปฏิจฺจ อุปฺปชฺชติ ผสฺสนานตฺตํ ผสฺสนานตฺตํ ปฏิจฺจ อุปฺปชฺชติ เวทนานานตฺตํ.\csromanspacer{} จกฺขุธาตุํ ภิกฺขเว ปฏิจฺจ อุปฺปชฺชติ จกฺขุสมฺผสฺโส, จกฺขุสมฺผสฺสํ ปฏิจฺจ อุปฺปชฺชติ จกฺขุสมฺผสฺสชา เวทนา ฯเปฯ.\csromanspacer{} มโนธาตุํ ปฏิจฺจ อุปฺปชฺชติ มโนสมฺผสฺโส, มโนสมฺผสฺสํ ปฏิจฺจ อุปฺปชฺชติ มโนสมฺผสฺสชา เวทนา.\csromanspacer{} เอวํ โข ภิกฺขเว ธาตุนานตฺตํ ปฏิจฺจ, อุปฺปชฺชติ ผสฺสนานตฺตํ, ผสฺสนานตฺตํ ปฏิจฺจ อุปฺปชฺชติ เวทนานานตฺตนฺติ.\csromanspacer{}\csromanspacer{}จตุตฺถํ.}

\csromanheader{2}{5. ทุติยเวทนานานตฺตสุตฺต}
\csromanmatikaanchor{mtk.102}
\csromantocmarkref{section}{7. สมุทฺทสุตฺต}{mtk.102}
\bookmark[dest={mtk.102},level=1]{7. สมุทฺทสุตฺต}
\proseitem{89}{สาวตฺถิยํ วิหรติ.\csromanspacer{} ธาตุนานตฺตํ ภิกฺขเว ปฏิจฺจ อุปฺปชฺชติ ผสฺสนานตฺตํ, ผสฺสนานตฺตํ ปฏิจฺจ อุปฺปชฺชติ เวทนานานตฺตํ, โน เวทนานานตฺตํ ปฏิจฺจ อุปฺปชฺชติ ผสฺสนานตฺตํ, โน ผสฺสนานตฺตํ ปฏิจฺจ อุปฺปชฺชติ ธาตุนานตฺตํ.\csromanspacer{} กตมญฺจ ภิกฺขเว ธาตุนานตฺตํ.\csromanspacer{} จกฺขุธาตุ ฯเปฯ มโนธาตุ.\csromanspacer{} อิทํ วุจฺจติ ภิกฺขเว ธาตุนานตฺตํ.}

\prose{กถญฺจ ภิกฺขเว ธาตุนานตฺตํ ปฏิจฺจ อุปฺปชฺชติ ผสฺสนานตฺตํ, ผสฺสนานตฺตํ ปฏิจฺจ อุปฺปชฺชติ เวทนานานตฺตํ, โน เวทนานานตฺตํ ปฏิจฺจ อุปฺปชฺชติ ผสฺสนานตฺตํ, โน ผสฺสนานตฺตํ ปฏิจฺจ อุปฺปชฺชติ ธาตุนานตฺตํ.\csromanspacer{} จกฺขุธาตุํ ภิกฺขเว ปฏิจฺจ อุปฺปชฺชติ จกฺขุสมฺผสฺโส, จกฺขุสมฺผสฺสํ ปฏิจฺจ อุปฺปชฺชติ จกฺขุสมฺผสฺสชา เวทนา, โน จกฺขุสมฺผสฺสชํ เวทนํ ปฏิจฺจ อุปฺปชฺชติ จกฺขุสมฺผสฺโส, โน จกฺขุสมฺผสฺสํ ปฏิจฺจ อุปฺปชฺชติ จกฺขุธาตุ ฯเปฯ.\csromanspacer{} มโนธาตุํ ปฏิจฺจ อุปฺปชฺชติ มโนสมฺผสฺโส,}

\vspace{\tipitakaparskip}
\csromancenter{นิทานวคฺคสํยุตฺตปาฬิ มโนสมฺผสฺสํ ปฏิจฺจ อุปฺปชฺชติ มโนสมฺผสฺสชา เวทนา, โน มโนสมฺผสฺสชํ เวทนํ ปฏิจฺจ อุปฺปชฺชติ มโนสมฺผสฺโส, โน มโนสมฺผสฺสํ ปฏิจฺจ อุปฺปชฺชติ มโนธาตุ.\csromanspacer{} เอวํ โข ภิกฺขเว ธาตุนานตฺตํ ปฏิจฺจ อุปฺปชฺชติ ผสฺสนานตฺตํ, ผสฺสนานตฺตํ ปฏิจฺจ อุปฺปชฺชติ เวทนานานตฺตํ, โน เวทนานานตฺตํ ปฏิจฺจ อุปฺปชฺชติ ผสฺสนานตฺตํ, โน ผสฺสนานตฺตํ ปฏิจฺจ อุปฺปชฺชติ ธาตุนานตฺตนฺติ.\csromanspacer{}\csromanspacer{}ปญฺจมํ.}
\csromanheader{2}{6. พาหิรธาตุนานตฺตสุตฺต}
\csromanmatikaanchor{mtk.103}
\csromantocmarkref{section}{8. ทุติยสมุทฺทสุตฺต}{mtk.103}
\bookmark[dest={mtk.103},level=1]{8. ทุติยสมุทฺทสุตฺต}
\proseitem{90}{\csromanfolio{360}สาวตฺถิยํ วิหรติ.\csromanspacer{} ธาตุนานตฺตํ โว ภิกฺขเว เทเสสฺสามิ, ตํ สุณาถ ฯเปฯ.\csromanspacer{} กตมญฺจ ภิกฺขเว ธาตุนานตฺตํ.\csromanspacer{} รูปธาตุ สทฺทธาตุ คนฺธธาตุ รสธาตุ โผฏฺฐพฺพธาตุ ธมฺมธาตุ.\csromanspacer{} อิทํ วุจฺจติ ภิกฺขเว ธาตุนานตฺตนฺติ.\csromanspacer{}\csromanspacer{}ฉฏฺฐํ.}

\csromanheader{2}{7. สญฺญานานตฺตสุตฺต}
\proseitem{91}{สาวตฺถิยํ วิหรติ.\csromanspacer{} ธาตุนานตฺตํ ภิกฺขเว ปฏิจฺจ อุปฺปชฺชติ สญฺญานานตฺตํ, สญฺญานานตฺตํ ปฏิจฺจ อุปฺปชฺชติ สงฺกปฺปนานตฺตํ, สงฺกปฺปนานตฺตํ ปฏิจฺจ อุปฺปชฺชติ ฉนฺทนานตฺตํ, ฉนฺทนานตฺตํ ปฏิจฺจ อุปฺปชฺชติ ปริฬาหนานตฺตํ, ปริฬาหนานตฺตํ ปฏิจฺจ อุปฺปชฺชติ ปริเยสนานานตฺตํ.\csromanspacer{} กตมญฺจ ภิกฺขเว ธาตุนานตฺตํ, รูปธาตุ ฯเปฯ ธมฺมธาตุ.\csromanspacer{} อิทํ วุจฺจติ ภิกฺขเว ธาตุนานตฺตํ.}

\csromanmatikaanchor{mtk.104}
\csromantocmarkref{section}{9. ปพฺพตสุตฺต}{mtk.104}
\bookmark[dest={mtk.104},level=1]{9. ปพฺพตสุตฺต}
\prose{กถญฺจ ภิกฺขเว ธาตุนานตฺตํ ปฏิจฺจ อุปฺปชฺชติ สญฺญานานตฺตํ, สญฺญานานตฺตํ ปฏิจฺจ อุปฺปชฺชติ สงฺกปฺปนานตฺตํ, สงฺกปฺปนานตฺตํ ปฏิจฺจ อุปฺปชฺชติ ฉนฺทนานตฺตํ, ฉนฺทนานตฺตํ ปฏิจฺจ อุปฺปชฺชติ ปริฬาหนานตฺตํ, ปริฬาหนานตฺตํ ปฏิจฺจ อุปฺปชฺชติ ปริเยสนานานตฺตํ.}

\prose{รูปธาตุํ ภิกฺขเว ปฏิจฺจ อุปฺปชฺชติ รูปสญฺญา, รูปสญฺญํ ปฏิจฺจ อุปฺปชฺชติ รูปสงฺกปฺโป, รูปสงฺกปฺปํ ปฏิจฺจ อุปฺปชฺชติ รูปจฺฉนฺโท, รูปจฺฉนฺทํ ปฏิจฺจ อุปฺปชฺชติ รูปปริฬาโห, รูปปริฬาหํ ปฏิจฺจ อุปฺปชฺชติ รูปปริเยสนา ฯเปฯ.\csromanspacer{} ธมฺมธาตุํ ปฏิจฺจ อุปฺปชฺชติ ธมฺมสญฺญา, ธมฺมสญฺญํ ปฏิจฺจ อุปฺปชฺชติ ธมฺมสงฺกปฺโป, ธมฺมสงฺกปฺปํ ปฏิจฺจ อุปฺปชฺชติ ธมฺมจฺฉนฺโท, ธมฺมจฺฉนฺทํ ปฏิจฺจ อุปฺปชฺชติ ธมฺมปริฬาโห, ธมฺมปริฬาหํ ปฏิจฺจ อุปฺปชฺชติ ธมฺมปริเยสนา.}

\prose{เอวํ โข ภิกฺขเว ธาตุนานตฺตํ ปฏิจฺจ อุปฺปชฺชติ สญฺญานานตฺตํ, สญฺญานานตฺตํ ปฏิจฺจ อุปฺปชฺชติ สงฺกปฺปนานตฺตํ, สงฺกปฺปนานตฺตํ ปฏิจฺจ อุปฺปชฺชติ ฉนฺทนานตฺตํ,}

\vspace{\tipitakaparskip}
\csromancenter{ธาตุสํยุตฺต ฉนฺทนานตฺตํ ปฏิจฺจ อุปฺปชฺชติ ปริฬาหนานตฺตํ, ปริฬาหนานตฺตํ ปฏิจฺจ อุปฺปชฺชติ ปริเยสนานานตฺตนฺติ.\csromanspacer{}\csromanspacer{}สตฺตมํ.}
\csromanmatikaanchor{mtk.105}
\csromantocmarkref{section}{10. ทุติยปพฺพตสุตฺต}{mtk.105}
\bookmark[dest={mtk.105},level=1]{10. ทุติยปพฺพตสุตฺต}
\csromanheader{2}{8. โนปริเยสนานานตฺตสุตฺต}
\proseitem{92}{\csromanfolio{361}สาวตฺถิยํ วิหรติ.\csromanspacer{} ธาตุนานตฺตํ ภิกฺขเว ปฏิจฺจ อุปฺปชฺชติ สญฺญานานตฺตํ, สญฺญานานตฺตํ ปฏิจฺจ อุปฺปชฺชติ สงฺกปฺปนานตฺตํ, สงฺกปฺปนานตฺตํ ปฏิจฺจ อุปฺปชฺชติ ฉนฺทนานตฺตํ, ฉนฺทนานตฺตํ ปฏิจฺจ อุปฺปชฺชติ ปริฬาหนานตฺตํ, ปริฬาหนานตฺตํ ปฏิจฺจ อุปฺปชฺชติ ปริเยสนานานตฺตํ, โน ปริเยสนานานตฺตํ ปฏิจฺจ อุปฺปชฺชติ ปริฬาหนานตฺตํ, โน ปริฬาหนานตฺตํ ปฏิจฺจ อุปฺปชฺชติ ฉนฺทนานตฺตํ, โน ฉนฺทนานตฺตํ ปฏิจฺจ อุปฺปชฺชติ สงฺกปฺปนานตฺตํ, โน สงฺกปฺปนานตฺตํ ปฏิจฺจ อุปฺปชฺชติ สญฺญานานตฺตํ, โน สญฺญานานตฺตํ ปฏิจฺจ อุปฺปชฺชติ ธาตุนานตฺตํ.\csromanspacer{} กตมญฺจ ภิกฺขเว ธาตุนานตฺตํ, รูปธาตุ ฯเปฯ ธมฺมธาตุ.\csromanspacer{} อิทํ วุจฺจติ ภิกฺขเว ธาตุนานตฺตํ.}

\prose{กถญฺจ ภิกฺขเว ธาตุนานตฺตํ ปฏิจฺจ อุปฺปชฺชติ สญฺญานานตฺตํ, สญฺญานานตฺตํ ปฏิจฺจ อุปฺปชฺชติ ฯเปฯ ปริเยสนานานตฺตํ.\csromanspacer{} โน ปริเยสนานานตฺตํ ปฏิจฺจ อุปฺปชฺชติ ปริฬาหนานตฺตํ, โน ปริฬาหนานตฺตํ ปฏิจฺจ อุปฺปชฺชติ ฉนฺทนานตฺตํ, โน ฉนฺทนานตฺตํ ปฏิจฺจ อุปฺปชฺชติ สงฺกปฺปนานตฺตํ, โน สงฺกปฺปนานตฺตํ ปฏิจฺจ อุปฺปชฺชติ สญฺญานานตฺตํ, โน สญฺญานานตฺตํ ปฏิจฺจ อุปฺปชฺชติ ธาตุนานตฺตํ.}

\prose{รูปธาตุํ ภิกฺขเว ปฏิจฺจ อุปฺปชฺชติ รูปสญฺญา ฯเปฯ.\csromanspacer{} ธมฺมธาตุํ ปฏิจฺจ อุปฺปชฺชติ ธมฺมสญฺญา, ธมฺมสญฺญํ ปฏิจฺจ อุปฺปชฺชติ ฯเปฯ ธมฺมปริเยสนา, โน ธมฺมปริเยสนํ ปฏิจฺจ อุปฺปชฺชติ ธมฺมปริฬาโห, โน ธมฺมปริฬาหํ ปฏิจฺจ อุปฺปชฺชติ ธมฺมจฺฉนฺโท, โน ธมฺมจฺฉนฺทํ ปฏิจฺจ อุปฺปชฺชติ ธมฺมสงฺกปฺโป, โน ธมฺมสงฺกปฺปํ ปฏิจฺจ อุปฺปชฺชติ ธมฺมสญฺญา.\csromanspacer{} โน ธมฺมสญฺญํ ปฏิจฺจ อุปฺปชฺชติ ธมฺมธาตุ.}

\csromanmatikaanchor{mtk.106}
\csromantocmarkref{section}{11. ตติยปพฺพตสุตฺต}{mtk.106}
\bookmark[dest={mtk.106},level=1]{11. ตติยปพฺพตสุตฺต}
\csromantocmarkref{section}{อุทฺทานคาถา}{mtk.107}
\bookmark[dest={mtk.107},level=1]{อุทฺทานคาถา}
\prose{เอวํ โข ภิกฺขเว ธาตุนานตฺตํ ปฏิจฺจ อุปฺปชฺชติ สญฺญานานตฺตํ, สญฺญานานตฺตํ ปฏิจฺจ อุปฺปชฺชติ ฯเปฯ ปริเยสนานานตฺตํ, โน ปริเยสนานานตฺตํ ปฏิจฺจ อุปฺปชฺชติ ปริฬาหนานตฺตํ, โน ปริฬาหนานตฺตํ ปฏิจฺจ อุปฺปชฺชติ ฉนฺทนานตฺตํ, โน ฉนฺทนานตฺตํ ปฏิจฺจ อุปฺปชฺชติ สงฺกปฺปนานตฺตํ, โน สงฺกปฺปนานตฺตํ ปฏิจฺจ อุปฺปชฺชติ สญฺญานานตฺตํ, โน สญฺญานานตฺตํ ปฏิจฺจ อุปฺปชฺชติ ธาตุนานตฺตนฺติ.\csromanspacer{}\csromanspacer{}อฏฺฐมํ.}

\csromanheader{2}{นิทานวคฺคสํยุตฺตปาฬิ \textbf{9. พาหิรผสฺสนานตฺตสุตฺต}}
\proseitem{93}{\csromanfolio{362}สาวตฺถิยํ วิหรติ.\csromanspacer{} ธาตุนานตฺตํ ภิกฺขเว ปฏิจฺจ อุปฺปชฺชติ สญฺญานานตฺตํ, สญฺญานานตฺตํ ปฏิจฺจ อุปฺปชฺชติ สงฺกปฺปนานตฺตํ, สงฺกปฺปนานตฺตํ ปฏิจฺจ อุปฺปชฺชติ ผสฺสนานตฺตํ, ผสฺสนานตฺตํ ปฏิจฺจ อุปฺปชฺชติ เวทนานานตฺตํ, เวทนานานตฺตํ ปฏิจฺจ อุปฺปชฺชติ ฉนฺทนานตฺตํ, ฉนฺทนานตฺตํ ปฏิจฺจ อุปฺปชฺชติ ปริฬาหนานตฺตํ, ปริฬาหนานตฺตํ ปฏิจฺจ อุปฺปชฺชติ ปริเยสนานานตฺตํ, ปริเยสนานานตฺตํ ปฏิจฺจ อุปฺปชฺชติ ลาภนานตฺตํ.\csromanspacer{} กตมญฺจ ภิกฺขเว ธาตุนานตฺตํ.\csromanspacer{} รูปธาตุ ฯเปฯ ธมฺมธาตุ.\csromanspacer{} อิทํ วุจฺจติ ภิกฺขเว ธาตุนานตฺตํ.}

\prose{กถญฺจ ภิกฺขเว ธาตุนานตฺตํ ปฏิจฺจ อุปฺปชฺชติ สญฺญานานตฺตํ, สญฺญานานตฺตํ ปฏิจฺจ อุปฺปชฺชติ ฯเปฯ ลาภนานตฺตํ.}

\prose{รูปธาตุํ ภิกฺขเว ปฏิจฺจ อุปฺปชฺชติ รูปสญฺญา, รูปสญฺญํ ปฏิจฺจ อุปฺปชฺชติ รูปสงฺกปฺโป, รูปสงฺกปฺปํ ปฏิจฺจ อุปฺปชฺชติ รูปสมฺผสฺโส, รูปสมฺผสฺสํ ปฏิจฺจ อุปฺปชฺชติ รูปสมฺผสฺสชา เวทนา, รูปสมฺผสฺสชํ เวทนํ ปฏิจฺจ อุปฺปชฺชติ รูปจฺฉนฺโท, รูปจฺฉนฺทํ ปฏิจฺจ อุปฺปชฺชติ รูปปริฬาโห, รูปปริฬาหํ ปฏิจฺจ อุปฺปชฺชติ รูปปริเยสนา, รูปปริเยสนํ ปฏิจฺจ อุปฺปชฺชติ รูปลาโภ ฯเปฯ.\csromanspacer{} ธมฺมธาตุํ ปฏิจฺจ อุปฺปชฺชติ ธมฺมสญฺญา, ธมฺมสญฺญํ ปฏิจฺจ อุปฺปชฺชติ ธมฺมสงฺกปฺโป, ธมฺมสงฺกปฺปํ ปฏิจฺจ อุปฺปชฺชติ ธมฺมสมฺผสฺโส, ธมฺมสมฺผสฺสํ ปฏิจฺจ อุปฺปชฺชติ ธมฺมสมฺผสฺสชา เวทนา, ธมฺมสมฺผสฺสชํ เวทนํ ปฏิจฺจ อุปฺปชฺชติ ธมฺมจฺฉนฺโท, ธมฺมจฺฉนฺทํ ปฏิจฺจ อุปฺปชฺชติ ธมฺมปริฬาโห, ธมฺมปริฬาหํ ปฏิจฺจ อุปฺปชฺชติ ธมฺมปริเยสนา, ธมฺมปริเยสนํ ปฏิจฺจ อุปฺปชฺชติ ธมฺมลาโภ.}

\prose{เอวํ โข ภิกฺขเว ธาตุนานตฺตํ ปฏิจฺจ อุปฺปชฺชติ สญฺญานานตฺตํ, สญฺญานานตฺตํ ปฏิจฺจ อุปฺปชฺชติ ฯเปฯ ปริเยสนานานตฺตํ, ปริเยสนานานตฺตํ ปฏิจฺจ อุปฺปชฺชติ ลาภนานตฺตนฺติ.\csromanspacer{}\csromanspacer{}นวมํ.}

\csromanheader{2}{10. ทุติยพาหิรผสฺสนานตฺตสุตฺต}
\csromanmatikaanchor{mtk.108}
\csromantocmarknopage{chapter}{3. ธาตุสํยุตฺต}{mtk.108}
\bookmark[dest={mtk.108},level=0]{3. ธาตุสํยุตฺต}
\proseitem{94}{สาวตฺถิยํ วิหรติ.\csromanspacer{} ธาตุนานตฺตํ ภิกฺขเว ปฏิจฺจ อุปฺปชฺชติ สญฺญานานตฺตํ, สญฺญานานตฺตํ ปฏิจฺจ อุปฺปชฺชติ สงฺกปฺปนานตฺตํ.\csromanspacer{} ผสฺส.\csromanspacer{} เวทนา.\csromanspacer{} ฉนฺท.\csromanspacer{} ปริฬาห.\csromanspacer{} ปริเยสนานานตฺตํ ปฏิจฺจ อุปฺปชฺชติ ลาภนานตฺตํ, โน ลาภนานตฺตํ ปฏิจฺจ อุปฺปชฺชติ ปริเยสนานานตฺตํ, โน ปริเยสนานานตฺตํ ปฏิจฺจ อุปฺปชฺชติ ปริฬาหนานตฺตํ, โน ปริฬาหนานตฺตํ ปฏิจฺจ อุปฺปชฺชติ ฯเปฯ ฉนฺท.\csromanspacer{} เวทนา.\csromanspacer{} ผสฺส.\csromanspacer{} สงฺกปฺป.\csromanspacer{} สญฺญานานตฺตํ, โน สญฺญานานตฺตํ}

\csromanmatikaanchor{mtk.109}
\csromantocmarknopage{section}{1. นานตฺตวคฺค}{mtk.109}
\bookmark[dest={mtk.109},level=1]{1. นานตฺตวคฺค}
\vspace{\tipitakaparskip}
\csromancenter{ธาตุสํยุตฺต ปฏิจฺจ อุปฺปชฺชติ ธาตุนานตฺตํ.\csromanspacer{} กตมญฺจ ภิกฺขเว ธาตุนานตฺตํ.\csromanspacer{} รูปธาตุ ฯเปฯ ธมฺมธาตุ.\csromanspacer{} อิทํ วุจฺจติ ภิกฺขเว ธาตุนานตฺตํ.}
\csromanmatikaanchor{mtk.110}
\csromanmatikaanchor{mtk.120}
\csromantocmarkref{subsection}{1. ธาตุนานตฺตสุตฺต}{mtk.110}
\bookmark[dest={mtk.110},level=2]{1. ธาตุนานตฺตสุตฺต}
\prose{\csromanfolio{363}กถญฺจ ภิกฺขเว ธาตุนานตฺตํ ปฏิจฺจ อุปฺปชฺชติ สญฺญานานตฺตํ, สญฺญานานตฺตํ ปฏิจฺจ อุปฺปชฺชติ สงฺกปฺปนานตฺตํ.\csromanspacer{} ผสฺส.\csromanspacer{} เวทนา.\csromanspacer{} ฉนฺท.\csromanspacer{} ปริฬาห.\csromanspacer{} ปริเยสนา.\csromanspacer{} ลาภ.\csromanspacer{} โน ลาภนานตฺตํ ปฏิจฺจ อุปฺปชฺชติ ปริเยสนานานตฺตํ, โน ปริเยสนานานตฺตํ ปฏิจฺจ อุปฺปชฺชติ ปริฬาห.\csromanspacer{} ฉนฺท.\csromanspacer{} เวทนา.\csromanspacer{} ผสฺส.\csromanspacer{} โน สงฺกปฺปนานตฺตํ ปฏิจฺจ อุปฺปชฺชติ สญฺญานานตฺตํ, โน สญฺญานานตฺตํ ปฏิจฺจ อุปฺปชฺชติ ธาตุนานตฺตํ.}

\prose{รูปธาตุํ ภิกฺขเว ปฏิจฺจ อุปฺปชฺชติ รูปสญฺญา ฯเปฯ.\csromanspacer{} ธมฺมธาตุํ ปฏิจฺจ อุปฺปชฺชติ ธมฺมสญฺญา, ธมฺมสญฺญํ ปฏิจฺจ อุปฺปชฺชติ ฯเปฯ ธมฺมปริเยสนา, ธมฺมปริเยสนํ ปฏิจฺจ อุปฺปชฺชติ ธมฺมลาโภ, โน ธมฺมลาภํ ปฏิจฺจ อุปฺปชฺชติ ธมฺมปริเยสนา, โน ธมฺมปริเยสนํ ปฏิจฺจ อุปฺปชฺชติ ธมฺมปริฬาโห, โน ธมฺมปริฬาหํ ปฏิจฺจ อุปฺปชฺชติ ธมฺมจฺฉนฺโท, โน ธมฺมจฺฉนฺทํ ปฏิจฺจ อุปฺปชฺชติ ธมฺมสมฺผสฺสชา เวทนา, โน ธมฺมสมฺผสฺสชํ เวทนํ ปฏิจฺจ อุปฺปชฺชติ ธมฺมสมฺผสฺโส, โน ธมฺมสมฺผสฺสํ ปฏิจฺจ อุปฺปชฺชติ ธมฺมสงฺกปฺโป, โน ธมฺมสงฺกปฺปํ ปฏิจฺจ อุปฺปชฺชติ ธมฺมสญฺญา, โน ธมฺมสญฺญํ ปฏิจฺจ อุปฺปชฺชติ ธมฺมธาตุ.}

\prose{เอวํ โข ภิกฺขเว ธาตุนานตฺตํ ปฏิจฺจ อุปฺปชฺชติ สญฺญานานตฺตํ, สญฺญานานตฺตํ ปฏิจฺจ อุปฺปชฺชติ ฯเปฯ สงฺกปฺป.\csromanspacer{} ผสฺส.\csromanspacer{} เวทนา.\csromanspacer{} ฉนฺท.\csromanspacer{} ปริฬาห.\csromanspacer{} ปริเยสนา.\csromanspacer{} ลาภ.\csromanspacer{} โน ลาภนานตฺตํ ปฏิจฺจ อุปฺปชฺชติ ปริเยสนานานตฺตํ, โน ปริเยสนานานตฺตํ ปฏิจฺจ อุปฺปชฺชติ ปริฬาหนานตฺตํ, โน ปริฬาหนานตฺตํ ปฏิจฺจ อุปฺปชฺชติ ฉนฺทนานตฺตํ, โน ฉนฺทนานตฺตํ ปฏิจฺจ อุปฺปชฺชติ เวทนานานตฺตํ, โน เวทนานานตฺตํ ปฏิจฺจ อุปฺปชฺชติ ผสฺสนานตฺตํ, โน ผสฺสนานตฺตํ ปฏิจฺจ อุปฺปชฺชติ สงฺกปฺปนานตฺตํ, โน สงฺกปฺปนานตฺตํ ปฏิจฺจ อุปฺปชฺชติ สญฺญานานตฺตํ, โน สญฺญานานตฺตํ ปฏิจฺจ อุปฺปชฺชติ ธาตุนานตฺตนฺติ.\csromanspacer{}\csromanspacer{}ทสมํ.}

\csromanmatikaanchor{mtk.111}
\csromantocmarkref{subsection}{2. ผสฺสนานตฺตสุตฺต}{mtk.111}
\bookmark[dest={mtk.111},level=2]{2. ผสฺสนานตฺตสุตฺต}
\vspace{\tipitakaparskip}
\csromancenter{นานตฺตวคฺโค ปฐโม.}
\titlehead{\textbf{ตสฺสุทฺทานํ}}
\csromangathameasure{ธาตุผสฺสญฺจ โน เจตํ, เวทนา อปเร ทุเว. \\ เอตํ อชฺฌตฺตปญฺจกํ, ธาตุสญฺญญฺจ โน เจตํ. \\ ผสฺสสฺส อปเร ทุเว, เอตํ พาหิรปญฺจกนฺติ.}
\csromangathagroup{\csromangathastanza{ธาตุผสฺสญฺจ โน เจตํ, เวทนา อปเร ทุเว. \\ เอตํ อชฺฌตฺตปญฺจกํ, ธาตุสญฺญญฺจ โน เจตํ. \\ ผสฺสสฺส อปเร ทุเว, เอตํ พาหิรปญฺจกนฺติ.}}

\csromanmatikaanchor{mtk.112}
\csromantocmarkref{subsection}{3. โนผสฺสนานตฺตสุตฺต}{mtk.112}
\bookmark[dest={mtk.112},level=2]{3. โนผสฺสนานตฺตสุตฺต}
\csromanheader{1}{นิทานวคฺคสํยุตฺตปาฬิ \textbf{2. ทุติยวคฺค}}
\csromanheader{2}{1. สตฺตธาตุสุตฺต}
\proseitem{95}{\csromanfolio{364}สาวตฺถิยํ วิหรติ.\csromanspacer{} สตฺติมา ภิกฺขเว ธาตุโย.\csromanspacer{} กตมา สตฺต, อาภาธาตุ สุภธาตุ อากาสานญฺจายตนธาตุ วิญฺญาณญฺจายตนธาตุ อากิญฺจญฺญายตนธาตุ เนวสญฺญานาสญฺญายตนธาตุ สญฺญาเวทยิตนิโรธธาตุ.\csromanspacer{} อิมา โข ภิกฺขเว สตฺต ธาตุโยติ.}

\prose{เอวํ วุตฺเต อญฺญตโร ภิกฺขุ ภควนฺตํ เอตทโวจ “ยา จายํ ภนฺเต อาภาธาตุ ยา จ สุภธาตุ ยา จ อากาสานญฺจายตนธาตุ ยา จ วิญฺญาณญฺจายตนธาตุ ยา จ อากิญฺจญฺญายตนธาตุ ยา จ เนวสญฺญานาสญฺญายตนธาตุ ยา จ สญฺญาเวทยิตนิโรธธาตุ, อิมา นุ โข ภนฺเต ธาตุโย กิํ ปฏิจฺจ ปญฺญายนฺตี”ติ.}

\csromanmatikaanchor{mtk.113}
\csromantocmarkref{subsection}{4. เวทนานานตฺตสุตฺต}{mtk.113}
\bookmark[dest={mtk.113},level=2]{4. เวทนานานตฺตสุตฺต}
\prose{ยายํ ภิกฺขุ อาภาธาตุ, อยํ ธาตุ อนฺธการํ ปฏิจฺจ ปญฺญายติ.\csromanspacer{} ยายํ ภิกฺขุ สุภธาตุ, อยํ ธาตุ อสุภํ ปฏิจฺจ ปญฺญายติ.\csromanspacer{} ยายํ ภิกฺขุ อากาสานญฺจายตนธาตุ, อยํ ธาตุ รูปํ ปฏิจฺจ ปญฺญายติ.\csromanspacer{} ยายํ ภิกฺขุ วิญฺญาณญฺจายตนธาตุ, อยํ ธาตุ อากาสานญฺจายตนํ ปฏิจฺจ ปญฺญายติ.\csromanspacer{} ยายํ ภิกฺขุ อากิญฺจญฺญายตนธาตุ, อยํ ธาตุ วิญฺญาณญฺจายตนํ ปฏิจฺจ ปญฺญายติ.\csromanspacer{} ยายํ ภิกฺขุ เนวสญฺญานาสญฺญายตนธาตุ, อยํ ธาตุ อากิญฺจญฺญายตนํ ปฏิจฺจ ปญฺญายติ.\csromanspacer{} ยายํ ภิกฺขุ สญฺญาเวทยิตนิโรธธาตุ, อยํ ธาตุ นิโรธํ ปฏิจฺจ ปญฺญายตีติ.}

\prose{ยา จายํ ภนฺเต อาภาธาตุ ยา จ สุภธาตุ ยา จ อากาสานญฺจายตนธาตุ ยา จ วิญฺญาณญฺจายตนธาตุ ยา จ อากิญฺจญฺญายตนธาตุ ยา จ เนวสญฺญานาสญฺญายตนธาตุ ยา จ สญฺญาเวทยิตนิโรธธาตุ, อิมา นุ โข ภนฺเต ธาตุโย กถํ สมาปตฺติ ปตฺตพฺพาติ.}

\prose{ยา จายํ ภิกฺขุ อาภาธาตุ ยา จ สุภธาตุ ยา จ อากาสานญฺจายตนธาตุ ยา จ วิญฺญาณญฺจายตนธาตุ ยา จ อากิญฺจญฺญายตนธาตุ, อิมา ธาตุโย สญฺญาสมาปตฺติ ปตฺตพฺพา.\csromanspacer{} ยายํ ภิกฺขุ เนวสญฺญานาสญฺญายตนธาตุ, อยํ ธาตุ สงฺขาราวเสสสมาปตฺติ}

\csromanmatikaanchor{mtk.114}
\csromantocmarkref{subsection}{5. ทุติยเวทนานานตฺตสุตฺต}{mtk.114}
\bookmark[dest={mtk.114},level=2]{5. ทุติยเวทนานานตฺตสุตฺต}
\vspace{\tipitakaparskip}
\csromancenter{ธาตุสํยุตฺต ปตฺตพฺพา.\csromanspacer{} ยายํ ภิกฺขุ สญฺญาเวทยิตนิโรธธาตุ, อยํ ธาตุ นิโรธสมาปตฺติ ปตฺตพฺพาติ.\csromanspacer{}\csromanspacer{}ปฐมํ.}
\csromanheader{2}{2. สนิทานสุตฺต}
\proseitem{96}{\csromanfolio{365}สาวตฺถิยํ วิหรติ.\csromanspacer{} สนิทานํ ภิกฺขเว อุปฺปชฺชติ กามวิตกฺโก โน อนิทานํ, สนิทานํ อุปฺปชฺชติ พฺยาปาทวิตกฺโก โน อนิทานํ, สนิทานํ อุปฺปชฺชติ วิหิํสาวิตกฺโก โน อนิทานํ.}

\prose{กถญฺจ ภิกฺขเว สนิทานํ อุปฺปชฺชติ กามวิตกฺโก โน อนิทานํ, สนิทานํ อุปฺปชฺชติ พฺยาปาทวิตกฺโก โน อนิทานํ, สนิทานํ อุปฺปชฺชติ วิหิํสาวิตกฺโก โน อนิทานํ.\csromanspacer{} กามธาตุํ ภิกฺขเว ปฏิจฺจ อุปฺปชฺชติ กามสญฺญา, กามสญฺญํ ปฏิจฺจ อุปฺปชฺชติ กามสงฺกปฺโป, กามสงฺกปฺปํ ปฏิจฺจ อุปฺปชฺชติ กามจฺฉนฺโท, กามจฺฉนฺทํ ปฏิจฺจ อุปฺปชฺชติ กามปริฬาโห, กามปริฬาหํ ปฏิจฺจ อุปฺปชฺชติ กามปริเยสนา, กามปริเยสนํ ภิกฺขเว ปริเยสมาโน อสฺสุตวา ปุถุชฺชโน ตีหิ ฐาเนหิ มิจฺฉา ปฏิปชฺชติ กาเยน วาจาย มนสา.}

\csromanmatikaanchor{mtk.115}
\csromantocmarkref{subsection}{6. พาหิรธาตุนานตฺตสุตฺต}{mtk.115}
\bookmark[dest={mtk.115},level=2]{6. พาหิรธาตุนานตฺตสุตฺต}
\prose{พฺยาปาทธาตุํ ภิกฺขเว ปฏิจฺจ อุปฺปชฺชติ พฺยาปาทสญฺญา, พฺยาปาทสญฺญํ ปฏิจฺจ อุปฺปชฺชติ พฺยาปาทสงฺกปฺโป ฯเปฯ พฺยาปาทจฺฉนฺโท.\csromanspacer{} พฺยาปาทปริฬาโห.\csromanspacer{} พฺยาปาทปริเยสนา.\csromanspacer{} พฺยาปาทปริเยสนํ ภิกฺขเว ปริเยสมาโน อสฺสุตวา ปุถุชฺชโน ตีหิ ฐาเนหิ มิจฺฉา ปฏิปชฺชติ กาเยน วาจาย มนสา.}

\prose{วิหิํสาธาตุํ ภิกฺขเว ปฏิจฺจ อุปฺปชฺชติ วิหิํสาสญฺญา.\csromanspacer{} วิหิํสาสญฺญํ ปฏิจฺจ อุปฺปชฺชติ วิหิํสาสงฺกปฺโป ฯเปฯ วิหิํสาฉนฺโท.\csromanspacer{} วิหิํสาปริฬาโห.\csromanspacer{} วิหิํสาปริเยสนา.\csromanspacer{} วิหิํสาปริเยสนํ ภิกฺขเว ปริเยสมาโน อสฺสุตวา ปุถุชฺชโน ตีหิ ฐาเนหิ มิจฺฉา ปฏิปชฺชติ กาเยน วาจาย มนสา.}

\csromanmatikaanchor{mtk.116}
\csromantocmarkref{subsection}{7. สญฺญานานตฺตสุตฺต}{mtk.116}
\bookmark[dest={mtk.116},level=2]{7. สญฺญานานตฺตสุตฺต}
\prose{เสยฺยถาปิ ภิกฺขเว ปุริโส อาทิตฺตํ ติณุกฺกํ สุกฺเข ติณทาเย นิกฺขิเปยฺย, โน เจ หตฺเถหิ จ ปาเทหิ จ ขิปฺปเมว นิพฺพาเปยฺย.\csromanspacer{} เอวํ หิ ภิกฺขเว เย ติณกฏฺฐนิสฺสิตา ปาณา, เต อนยพฺยสนํ อาปชฺเชยฺยุํ.\csromanspacer{} เอวเมว โข ภิกฺขเว โย หิ โกจิ สมโณ วา พฺราหฺมโณ วา อุปฺปนฺนํ วิสมคตํ สญฺญํ น ขิปฺปเมว ปชหติ วิโนเทติ พฺยนฺตีกโรติ}

\vspace{\tipitakaparskip}
\csromancenter{นิทานวคฺคสํยุตฺตปาฬิ อนภาวํ คเมติ, โส ทิฏฺเฐ เจว ธมฺเม ทุกฺขํ วิหรติ สวิฆาตํ ส- อุปายาสํ สปริฬาหํ, กายสฺส จ เภทา ปรํ มรณา ทุคฺคติ ปาฏิกงฺขา.}
\prose{\csromanfolio{366}สนิทานํ ภิกฺขเว อุปฺปชฺชติ เนกฺขมฺมวิตกฺโก โน อนิทานํ, สนิทานํ อุปฺปชฺชติ อพฺยาปาทวิตกฺโก โน อนิทานํ, สนิทานํ อุปฺปชฺชติ อวิหิํสาวิตกฺโก โน อนิทานํ.}

\prose{กถญฺจ ภิกฺขเว สนิทานํ อุปฺปชฺชติ เนกฺขมฺมวิตกฺโก โน อนิทานํ, สนิทานํ อุปฺปชฺชติ อพฺยาปาทวิตกฺโก โน อนิทานํ, สนิทานํ อุปฺปชฺชติ อวิหิํสาวิตกฺโก โน อนิทานํ.\csromanspacer{} เนกฺขมฺมธาตุํ ภิกฺขเว ปฏิจฺจ อุปฺปชฺชติ เนกฺขมฺมสญฺญา, เนกฺขมฺมสญฺญํ ปฏิจฺจ อุปฺปชฺชติ เนกฺขมฺมสงฺกปฺโป, เนกฺขมฺมสงฺกปฺปํ ปฏิจฺจ อุปฺปชฺชติ เนกฺขมฺมจฺฉนฺโท, เนกฺขมฺมจฺฉนฺทํ ปฏิจฺจ อุปฺปชฺชติ เนกฺขมฺมปริฬาโห, เนกฺขมฺมปริฬาหํ ปฏิจฺจ อุปฺปชฺชติ เนกฺขมฺมปริเยสนา, เนกฺขมฺมปริเยสนํ ภิกฺขเว ปริเยสมาโน สุตวา อริยสาวโก ตีหิ ฐาเนหิ สมฺมา ปฏิปชฺชติ กาเยน วาจาย มนสา.}

\prose{อพฺยาปาทธาตุํ ภิกฺขเว ปฏิจฺจ อุปฺปชฺชติ อพฺยาปาทสญฺญา.\csromanspacer{} อพฺยาปาทสญฺญํ ปฏิจฺจ อุปฺปชฺชติ อพฺยาปาทสงฺกปฺโป ฯเปฯ อพฺยาปาทจฺฉนฺโท.\csromanspacer{} อพฺยาปาทปริฬาโห.\csromanspacer{} อพฺยาปาทปริเยสนา.\csromanspacer{} อพฺยาปาทปริเยสนํ ภิกฺขเว ปริเยสมาโน สุตวา อริยสาวโก ตีหิ ฐาเนหิ สมฺมา ปฏิปชฺชติ กาเยน วาจาย มนสา.}

\prose{อวิหิํสาธาตุํ ภิกฺขเว ปฏิจฺจ อุปฺปชฺชติ อวิหิํสาสญฺญา, อวิหิํสาสญฺญํ ปฏิจฺจ อุปฺปชฺชติ อวิหิํสาสงฺกปฺโป, อวิหิํสาสงฺกปฺปํ ปฏิจฺจ อุปฺปชฺชติ อวิหิํสาฉนฺโท, อวิหิํสาฉนฺทํ ปฏิจฺจ อุปฺปชฺชติ อวิหิํสาปริฬาโห, อวิหิํสาปริฬาหํ ปฏิจฺจ อุปฺปชฺชติ อวิหิํสาปริเยสนา, อวิหิํสาปริเยสนํ ภิกฺขเว ปริเยสมาโน สุตวา อริยสาวโก ตีหิ ฐาเนหิ สมฺมา ปฏิปชฺชติ กาเยน วาจาย มนสา.}

\csromanmatikaanchor{mtk.117}
\csromantocmarkref{subsection}{8. โนปริเยสนานานตฺตสุตฺต}{mtk.117}
\bookmark[dest={mtk.117},level=2]{8. โนปริเยสนานานตฺตสุตฺต}
\prose{เสยฺยถาปิ ภิกฺขเว ปุริโส อาทิตฺตํ ติณุกฺกํ สุกฺเข ติณทาเย นิกฺขิเปยฺย, ตเมนํ หตฺเถหิ จ ปาเทหิ จ ขิปฺปเมว นิพฺพาเปยฺย.\csromanspacer{} เอวํ หิ ภิกฺขเว เย ติณกฏฺฐนิสฺสิตา ปาณา, เต น อนยพฺยสนํ อาปชฺเชยฺยุํ.\csromanspacer{} เอวเมว โข ภิกฺขเว โย หิ โกจิ สมโณ วา พฺราหฺมโณ วา อุปฺปนฺนํ วิสมคตํ สญฺญํ ขิปฺปเมว ปชหติ วิโนเทติ}

\vspace{\tipitakaparskip}
\csromancenter{ธาตุสํยุตฺต พฺยนฺตีกโรติ อนภาวํ คเมติ, โส ทิฏฺเฐ เจว ธมฺเม สุขํ วิหรติ อวิฆาตํ อนุปายาสํ อปริฬาหํ, กายสฺส จ เภทา ปรํ มรณา สุคติ ปาฏิกงฺขาติ.\csromanspacer{}\csromanspacer{}ทุติยํ.}
\csromanheader{2}{3. คิญฺชกาวสถสุตฺต}
\proseitem{97}{\csromanfolio{367}เอกํ สมยํ ภควา ญาติเก วิหรติ คิญฺชกาวสเถ.\csromanspacer{} ตตฺร โข ภควา ภิกฺขู อามนฺเตสิ “ภิกฺขโว”ติ. “ภทนฺเต”ติ เต ภิกฺขู ภควโต ปจฺจสฺโสสุํ.\csromanspacer{} ภควา เอตทโวจ\csromandash{}}

\prose{“ธาตุํ ภิกฺขเว ปฏิจฺจ อุปฺปชฺชติ สญฺญา อุปฺปชฺชติ ทิฏฺฐิ อุปฺปชฺชติ วิตกฺโก”ติ.\csromanspacer{} เอวํ วุตฺเต อายสฺมา กจฺจาโน\footnote{สทฺโธ กจฺจาโน (ก)} ภควนฺตํ เอตทโวจ “ยายํ ภนฺเต ทิฏฺฐิ ‘อสมฺมาสมฺพุทฺเธสุ สมฺมาสมฺพุทฺธา’ติ, อยํ นุ โข ภนฺเต ทิฏฺฐิ กิํ ปฏิจฺจ ปญฺญายตี”ติ.}

\csromanmatikaanchor{mtk.118}
\csromantocmarkref{subsection}{9. พาหิรผสฺสนานตฺตสุตฺต}{mtk.118}
\bookmark[dest={mtk.118},level=2]{9. พาหิรผสฺสนานตฺตสุตฺต}
\prose{มหติ โข เอสา กจฺจาน ธาตุ ยทิทํ อวิชฺชาธาตุ.\csromanspacer{} หีนํ กจฺจาน ธาตุํ ปฏิจฺจ อุปฺปชฺชติ หีนา สญฺญา หีนา ทิฏฺฐิ หีโน วิตกฺโก หีนา เจตนา หีนา ปตฺถนา หีโน ปณิธิ หีโน ปุคฺคโล หีนา วาจา หีนํ อาจิกฺขติ เทเสติ ปญฺญเปติ ปฏฺฐเปติ วิวรติ วิภชติ อุตฺตานีกโรติ, หีนา ตสฺส อุปปตฺตีติ วทามิ.}

\prose{มชฺฌิมํ กจฺจาน ธาตุํ ปฏิจฺจ อุปฺปชฺชติ มชฺฌิมา สญฺญา มชฺฌิมา ทิฏฺฐิ มชฺฌิโม วิตกฺโก มชฺฌิมา เจตนา มชฺฌิมา ปตฺถนา มชฺฌิโม ปณิธิ มชฺฌิโม ปุคฺคโล มชฺฌิมา วาจา มชฺฌิมํ อาจิกฺขติ เทเสติ ปญฺญเปติ ปฏฺฐเปติ วิวรติ วิภชติ อุตฺตานีกโรติ, มชฺฌิมา ตสฺส อุปปตฺตีติ วทามิ.}

\prose{ปณีตํ กจฺจาน ธาตุํ ปฏิจฺจ อุปฺปชฺชติ ปณีตา สญฺญา ปณีตา ทิฏฺฐิ ปณีโต วิตกฺโก ปณีตา เจตนา ปณีตา ปตฺถนา ปณีโต ปณิธิ ปณีโต ปุคฺคโล ปณีตา วาจา ปณีตํ อาจิกฺขติ เทเสติ ปญฺญเปติ ปฏฺฐเปติ วิวรติ วิภชติ อุตฺตานีกโรติ, ปณีตา ตสฺส อุปปตฺตีติ วทามีติ.\csromanspacer{}\csromanspacer{}ตติยํ.}

\csromanheader{2}{นิทานวคฺคสํยุตฺตปาฬิ \textbf{4. หีนาธิมุตฺติกสุตฺต}}
\proseitem{98}{\csromanfolio{368}สาวตฺถิยํ วิหรติ.\csromanspacer{} ธาตุโสว\footnote{ธาตุโส (สี, อิ) อยญฺจ ปฐมารมฺภวาเกฺยเยว, น สพฺพตฺถ. ตีสุ ปน อทฺธาสุ จ อุปมาสํสนฺทนนิคมนฏฺฐาเน จ อิทํ ปาฐนานตฺตํ นตฺถิ.} ภิกฺขเว สตฺตา สํสนฺทนฺติ สเมนฺติ, หีนาธิมุตฺติกา หีนาธิมุตฺติเกหิ สทฺธิํ สํสนฺทนฺติ สเมนฺติ, กลฺยาณาธิมุตฺติกา กลฺยาณาธิมุตฺติเกหิ สทฺธิํ สํสนฺทนฺติ สเมนฺติ.}

\csromanmatikaanchor{mtk.119}
\csromantocmarkref{subsection}{10. ทุติยพาหิรผสฺสนานตฺตสุตฺต}{mtk.119}
\bookmark[dest={mtk.119},level=2]{10. ทุติยพาหิรผสฺสนานตฺตสุตฺต}
\csromantocmarkref{section}{อุทฺทานคาถา}{mtk.120}
\bookmark[dest={mtk.120},level=1]{อุทฺทานคาถา}
\prose{อตีตมฺปิ โข\footnote{โขสทฺโท สี-สฺยา-กํ-อิ-โปตฺถเกสุ นตฺถิ.} ภิกฺขเว อทฺธานํ ธาตุโสว\footnote{¢ทิเสสุ ฐาเนสุ ปาฐนานตฺตํ นตฺถิ.} สตฺตา สํสนฺทิํสุ สมิํสุ, หีนาธิมุตฺติกา หีนาธิมุตฺติเกหิ สทฺธิํ สํสนฺทิํสุ สมิํสุ, กลฺยาณาธิมุตฺติกา กลฺยาณาธิมุตฺติเกหิ สทฺธิํ สํสนฺทิํสุ สมิํสุ.}

\prose{อนาคตมฺปิ โข2 ภิกฺขเว อทฺธานํ ธาตุโสว3 สตฺตา สํสนฺทิสฺสนฺติ สเมสฺสนฺติ, หีนาธิมุตฺติกา หีนาธิมุตฺติเกหิ สทฺธิํ สํสนฺทิสฺสนฺติ สเมสฺสนฺติ, กลฺยาณาธิมุตฺติกา กลฺยาณาธิมุตฺติเกหิ สทฺธิํ สํสนฺทิสฺสนฺติ สเมสฺสนฺติ.}

\prose{เอตรหิปิ โข2 ภิกฺขเว ปจฺจุปฺปนฺนํ อทฺธานํ ธาตุโสว3 สตฺตา สํสนฺทนฺติ สเมนฺติ, หีนาธิมุตฺติกา หีนาธิมุตฺติเกหิ สทฺธิํ สํสนฺทนฺติ สเมนฺติ, กลฺยาณาธิมุตฺติกา กลฺยาณาธิมุตฺติเกหิ สทฺธิํ สํสนฺทนฺติ สเมนฺตีติ.\csromanspacer{}\csromanspacer{}จตุตฺถํ.}

\csromanheader{2}{5. จงฺกมสุตฺต}
\proseitem{99}{เอกํ สมยํ ภควา ราชคเห วิหรติ คิชฺฌกูเฏ ปพฺพเต.\csromanspacer{} เตน โข ปน สมเยน อายสฺมา สาริปุตฺโต สมฺพหุเลหิ ภิกฺขูหิ สทฺธิํ ภควโต อวิทูเร จงฺกมติ, อายสฺมาปิ โข มหาโมคฺคลฺลาโน สมฺพหุเลหิ ภิกฺขูหิ สทฺธิํ ภควโต อวิทูเร จงฺกมติ, อายสฺมาปิ โข มหากสฺสโป สมฺพหุเลหิ ภิกฺขูหิ สทฺธิํ ภควโต อวิทูเร จงฺกมติ, อายสฺมาปิ โข อนุรุทฺโธ สมฺพหุเลหิ ภิกฺขูหิ สทฺธิํ ภควโต อวิทูเร จงฺกมติ, อายสฺมาปิ โข ปุณฺโณ มนฺตานิปุตฺโต สมฺพหุเลหิ ภิกฺขูหิ สทฺธิํ ภควโต อวิทูเร จงฺกมติ, อายสฺมาปิ โข อุปาลิ สมฺพหุเลหิ ภิกฺขูหิ สทฺธิํ ภควโต อวิทูเร จงฺกมติ, อายสฺมาปิ โข อานนฺโท สมฺพหุเลหิ ภิกฺขูหิ สทฺธิํ ภควโต อวิทูเร จงฺกมติ, เทวทตฺโตปิ โข สมฺพหุเลหิ ภิกฺขูหิ สทฺธิํ ภควโต อวิทูเร จงฺกมติ.}

\csromanheader{2}{3. ธาตุสํยุตฺต}
\prose{\csromanfolio{369}อถ โข ภควา ภิกฺขู อามนฺเตสิ\csromandash{}ปสฺสถ โน ตุเมฺห ภิกฺขเว สาริปุตฺตํ สมฺพหุเลหิ ภิกฺขูหิ สทฺธิํ จงฺกมนฺตนฺติ.\csromanspacer{} เอวํ ภนฺเต.\csromanspacer{} สพฺเพ โข เอเต ภิกฺขเว ภิกฺขู มหาปญฺญา.\csromanspacer{} ปสฺสถ โน ตุเมฺห ภิกฺขเว โมคฺคลฺลานํ สมฺพหุเลหิ ภิกฺขูหิ สทฺธิํ จงฺกมนฺตนฺติ.\csromanspacer{} เอวํ ภนฺเต.\csromanspacer{} สพฺเพ โข เอเต ภิกฺขเว ภิกฺขู มหิทฺธิกา.\csromanspacer{} ปสฺสถ โน ตุเมฺห ภิกฺขเว กสฺสปํ สมฺพหุเลหิ ภิกฺขูหิ สทฺธิํ จงฺกมนฺตนฺติ.\csromanspacer{} เอวํ ภนฺเต.\csromanspacer{} สพฺเพ โข เอเต ภิกฺขเว ภิกฺขู ธุตวาทา.\csromanspacer{} ปสฺสถ โน ตุเมฺห ภิกฺขเว อนุรุทฺธํ สมฺพหุเลหิ ภิกฺขูหิ สทฺธิํ จงฺกมนฺตนฺติ.\csromanspacer{} เอวํ ภนฺเต.\csromanspacer{} สพฺเพ โข เอเต ภิกฺขเว ภิกฺขู ทิพฺพจกฺขุกา.\csromanspacer{} ปสฺสถ โน ตุเมฺห ภิกฺขเว ปุณฺณํ มนฺตานิปุตฺตํ สมฺพหุเลหิ ภิกฺขูหิ สทฺธิํ จงฺกมนฺตนฺติ.\csromanspacer{} เอวํ ภนฺเต.\csromanspacer{} สพฺเพ โข เอเต ภิกฺขเว ภิกฺขู ธมฺมกถิกา.\csromanspacer{} ปสฺสถ โน ตุเมฺห ภิกฺขเว อุปาลิํ สมฺพหุเลหิ ภิกฺขูหิ สทฺธิํ จงฺกมนฺตนฺติ.\csromanspacer{} เอวํ ภนฺเต.\csromanspacer{} สพฺเพ โข เอเต ภิกฺขเว ภิกฺขู วินยธรา.\csromanspacer{} ปสฺสถ โน ตุเมฺห ภิกฺขเว อานนฺทํ สมฺพหุเลหิ ภิกฺขูหิ สทฺธิํ จงฺกมนฺตนฺติ.\csromanspacer{} เอวํ ภนฺเต.\csromanspacer{} สพฺเพ โข เอเต ภิกฺขเว ภิกฺขู พหุสฺสุตา.\csromanspacer{} ปสฺสถ โน ตุเมฺห ภิกฺขเว เทวทตฺตํ สมฺพหุเลหิ ภิกฺขูหิ สทฺธิํ จงฺกมนฺตนฺติ.\csromanspacer{} เอวํ ภนฺเต.\csromanspacer{} สพฺเพ โข เอเต ภิกฺขเว ภิกฺขู ปาปิจฺฉา.}

\prose{ธาตุโสว ภิกฺขเว สตฺตา สํสนฺทนฺติ สเมนฺติ, หีนาธิมุตฺติกา หีนาธิมุตฺติเกหิ สทฺธิํ สํสนฺทนฺติ สเมนฺติ, กลฺยาณาธิมุตฺติกา กลฺยาณาธิมุตฺติเกหิ สทฺธิํ สํสนฺทนฺติ สเมนฺติ.\csromanspacer{} อตีตมฺปิ โข ภิกฺขเว อทฺธานํ ธาตุโสว สตฺตา สํสนฺทิํสุ สมิํสุ, หีนาธิมุตฺติกา หีนาธิมุตฺติเกหิ สทฺธิํ สํสนฺทิํสุ สมิํสุ, กลฺยาณาธิมุตฺติกา กลฺยาณาธิมุตฺติเกหิ สทฺธิํ สํสนฺทิํสุ สมิํสุ.}

\prose{อนาคตมฺปิ โข ภิกฺขเว อทฺธานํ ธาตุโสว สตฺตา สํสนฺทิสฺสนฺติ สเมสฺสนฺติ, หีนาธิมุตฺติกา หีนาธิมุตฺติเกหิ สทฺธิํ สํสนฺทิสฺสนฺติ สเมสฺสนฺติ, กลฺยาณาธิมุตฺติกา กลฺยาณาธิมุตฺติเกหิ สทฺธิํ สํสนฺทิสฺสนฺติ สเมสฺสนฺติ.}

\csromanmatikaanchor{mtk.121}
\csromantocmarknopage{section}{2. ทุติยวคฺค}{mtk.121}
\bookmark[dest={mtk.121},level=1]{2. ทุติยวคฺค}
\prose{เอตรหิปิ โข ภิกฺขเว ปจฺจุปฺปนฺนํ อทฺธานํ ธาตุโสว สตฺตา สํสนฺทนฺติ สเมนฺติ, หีนาธิมุตฺติกา หีนาธิมุตฺติเกหิ สทฺธิํ สํสนฺทนฺติ สเมนฺติ, กลฺยาณาธิมุตฺติกา กลฺยาณาธิมุตฺติเกหิ สทฺธิํ สํสนฺทนฺติ สเมนฺตีติ.\csromanspacer{}\csromanspacer{}ปญฺจมํ.}

\csromanmatikaanchor{mtk.122}
\csromantocmarkref{subsection}{1. สตฺตธาตุสุตฺต}{mtk.122}
\bookmark[dest={mtk.122},level=2]{1. สตฺตธาตุสุตฺต}
\csromanheader{2}{นิทานวคฺคสํยุตฺตปาฬิ \textbf{6. สคาถาสุตฺต}}
\proseitem{100}{\csromanfolio{370}สาวตฺถิยํ วิหรติ.\csromanspacer{} ธาตุโสว ภิกฺขเว สตฺตา สํสนฺทนฺติ สเมนฺติ, หีนาธิมุตฺติกา หีนาธิมุตฺติเกหิ สทฺธิํ สํสนฺทนฺติ สเมนฺติ.\csromanspacer{} อตีตมฺปิ โข ภิกฺขเว อทฺธานํ ธาตุโสว สตฺตา สํสนฺทิํสุ สมิํสุ, หีนาธิมุตฺติกา หีนาธิมุตฺติเกหิ สทฺธิํ สํสนฺทิํสุ สมิํสุ.}

\prose{อนาคตมฺปิ โข ภิกฺขเว อทฺธานํ ธาตุโสว สตฺตา สํสนฺทิสฺสนฺติ สเมสฺสนฺติ, หีนาธิมุตฺติกา หีนาธิมุตฺติเกหิ สทฺธิํ สํสนฺทิสฺสนฺติ สเมสฺสนฺติ.}

\prose{เอตรหิปิ โข ภิกฺขเว ปจฺจุปฺปนฺนํ อทฺธานํ ธาตุโสว สตฺตา สํสนฺทนฺติ สเมนฺติ, หีนาธิมุตฺติกา หีนาธิมุตฺติเกหิ สทฺธิํ สํสนฺทนฺติ สเมนฺติ.}

\prose{เสยฺยถาปิ ภิกฺขเว คูโถ คูเถน สํสนฺทติ สเมติ, มุตฺตํ มุตฺเตน สํสนฺทติ สเมติ, เขโฬ เขเฬน สํสนฺทติ สเมติ, ปุพฺโพ ปุพฺเพน สํสนฺทติ สเมติ, โลหิตํ โลหิเตน สํสนฺทติ สเมติ.\csromanspacer{} เอวเมว โข ภิกฺขเว ธาตุโสว\footnote{สพฺพตฺถปิ เอวเมว ทิสฺสติ.} สตฺตา สํสนฺทนฺติ สเมนฺติ, หีนาธิมุตฺติกา หีนาธิมุตฺติเกหิ สทฺธิํ สํสนฺทนฺติ สเมนฺติ.\csromanspacer{} อตีตมฺปิ โข อทฺธานํ ฯเปฯ.\csromanspacer{} อนาคตมฺปิ โข อทฺธานํ ฯเปฯ.\csromanspacer{} เอตรหิปิ โข ปจฺจุปฺปนฺนํ อทฺธานํ ธาตุโสว สตฺตา สํสนฺทนฺติ สเมนฺติ, หีนาธิมุตฺติกา หีนาธิมุตฺติเกหิ สทฺธิํ สํสนฺทนฺติ สเมนฺติ.}

\prose{ธาตุโสว ภิกฺขเว สตฺตา สํสนฺทนฺติ สเมนฺติ, กลฺยาณาธิมุตฺติกา กลฺยาณาธิมุตฺติเกหิ สทฺธิํ สํสนฺทนฺติ สเมนฺติ.\csromanspacer{} อตีตมฺปิ โข ภิกฺขเว อทฺธานํ ธาตุโสว สตฺตา สํสนฺทิํสุ สมิํสุ, กลฺยาณาธิมุตฺติกา กลฺยาณาธิมุตฺติเกหิ สทฺธิํ สํสนฺทิํสุ สมิํสุ.}

\prose{อนาคตมฺปิ โข ภิกฺขเว อทฺธานํ ฯเปฯ.\csromanspacer{} เอตรหิปิ โข ภิกฺขเว ปจฺจุปฺปนฺนํ อทฺธานํ ธาตุโสว สตฺตา สํสนฺทนฺติ สเมนฺติ, กลฺยาณาธิมุตฺติกา กลฺยาณาธิมุตฺติเกหิ สทฺธิํ สํสนฺทนฺติ สเมนฺติ.}

\csromanmatikaanchor{mtk.123}
\csromantocmarkref{subsection}{2. สนิทานสุตฺต}{mtk.123}
\bookmark[dest={mtk.123},level=2]{2. สนิทานสุตฺต}
\prose{เสยฺยถาปิ ภิกฺขเว ขีรํ ขีเรน สํสนฺทติ สเมติ, เตลํ เตเลน สํสนฺทติ สเมติ, สปฺปิ สปฺปินา สํสนฺทติ สเมติ, มธุ มธุนา สํสนฺทติ สเมติ, ผาณิตํ ผาณิเตน สํสนฺทติ สเมติ.\csromanspacer{} เอวเมว โข ภิกฺขเว}

\vspace{\tipitakaparskip}
\csromancenter{ธาตุสํยุตฺต ธาตุโสว สตฺตา สํสนฺทนฺติ สเมนฺติ, กลฺยาณาธิมุตฺติกา กลฺยาณาธิมุตฺติเกหิ สทฺธิํ สํสนฺทนฺติ สเมนฺติ.\csromanspacer{} อตีตมฺปิ โข อทฺธานํ.\csromanspacer{} อนาคตมฺปิ โข อทฺธานํ.\csromanspacer{} เอตรหิปิ โข ปจฺจุปฺปนฺนํ อทฺธานํ ธาตุโสว สตฺตา สํสนฺทนฺติ สเมนฺติ, กลฺยาณาธิมุตฺติกา กลฺยาณาธิมุตฺติเกหิ สทฺธิํ สํสนฺทนฺติ สเมนฺตีติ.}
\prose{\csromanfolio{371}อิทมโวจ ภควา, อิทํ วตฺวาน สุคโต อถาปรํ เอตทโวจ สตฺถา\csromandash{}}

\csromangathameasure{“สํสคฺคา วนโถ ชาโต, อสํสคฺเคน ฉิชฺชติ. \\ ปริตฺตํ ทารุมารุยฺห, ยถา สีเท มหณฺณเว. \\ เอวํ กุสีตมาคมฺม, สาธุชีวิปิ สีทติ. \\ ตสฺมา ตํ ปริวชฺเชยฺย, กุสีตํ หีนวีริยํ. \\ ปวิวิตฺเตหิ อริเยหิ, ปหิตตฺเตหิ ฌายีหิ. \\ นิจฺจํ อารทฺธวิริเยหิ, ปณฺฑิเตหิ สหาวเส”ติ.}
\csromangathagroup{\csromangathastanza{“สํสคฺคา วนโถ ชาโต, อสํสคฺเคน ฉิชฺชติ. \\ ปริตฺตํ ทารุมารุยฺห, ยถา สีเท มหณฺณเว.}\csromangathastanza{เอวํ กุสีตมาคมฺม, สาธุชีวิปิ สีทติ. \\ ตสฺมา ตํ ปริวชฺเชยฺย, กุสีตํ หีนวีริยํ.}\csromangathastanza{ปวิวิตฺเตหิ อริเยหิ, ปหิตตฺเตหิ ฌายีหิ\footnote{ฌายิหิ (สี), ฌายิภิ (สฺยา, กํ)}. \\ นิจฺจํ อารทฺธวิริเยหิ, ปณฺฑิเตหิ สหาวเส”ติ.}}

\vspace{\tipitakaparskip}
\csromancenter{ฉฏฺฐํ.}
\csromanheader{2}{7. อสฺสทฺธสํสนฺทนสุตฺต}
\proseitem{101}{สาวตฺถิยํ วิหรติ.\csromanspacer{} ธาตุโสว ภิกฺขเว สตฺตา สํสนฺทนฺติ สเมนฺติ.\csromanspacer{} อสฺสทฺธา อสฺสทฺเธหิ สทฺธิํ สํสนฺทนฺติ สเมนฺติ, อหิริกา อหิริเกหิ สทฺธิํ สํสนฺทนฺติ สเมนฺติ, อโนตฺตปฺปิโน อโนตฺตปฺปีหิ สทฺธิํ สํสนฺทนฺติ สเมนฺติ, อปฺปสฺสุตา อปฺปสฺสุเตหิ สทฺธิํ สํสนฺทนฺติ สเมนฺติ, กุสีตา กุสีเตหิ สทฺธิํ สํสนฺทนฺติ สเมนฺติ, มุฏฺฐสฺสติโน มุฏฺฐสฺสตีหิ สทฺธิํ สํสนฺทนฺติ สเมนฺติ, ทุปฺปญฺญา ทุปฺปญฺเญหิ สทฺธิํ สํสนฺทนฺติ สเมนฺติ.}

\prose{อตีตมฺปิ โข ภิกฺขเว อทฺธานํ ธาตุโสว สตฺตา สํสนฺทิํสุ สมิํสุ.\csromanspacer{} อสฺสทฺธา อสฺสทฺเธหิ สทฺธิํ สํสนฺทิํสุ สมิํสุ, อหิริกา อหิริเกหิ สทฺธิํ สํสนฺทิํสุ สมิํสุ, อโนตฺตปฺปิโน อโนตฺตปฺปีหิ สทฺธิํ สํสนฺทิํสุ สมิํสุ, อปฺปสฺสุตา อปฺปสฺสุเตหิ สทฺธิํ สํสนฺทิํสุ สมิํสุ, กุสีตา กุสีเตหิ สทฺธิํ สํสนฺทิํสุ สมิํสุ, มุฏฺฐสฺสติโน มุฏฺฐสฺสตีหิ สทฺธิํ สํสนฺทิํสุ สมิํสุ, ทุปฺปญฺญา ทุปฺปญฺเญหิ สทฺธิํ สํสนฺทิํสุ สมิํสุ.}

\gambhira{นิทานวคฺคสํยุตฺตปาฬิ}
\prose{\csromanfolio{372}อนาคตมฺปิ โข ภิกฺขเว อทฺธานํ ธาตุโสว สตฺตา สํสนฺทิสฺสนฺติ สเมสฺสนฺติ.\csromanspacer{} อสฺสทฺธา อสฺสทฺเธหิ สทฺธิํ สํสนฺทิสฺสนฺติ สเมสฺสนฺติ, อหิริกา อหิริเกหิ สทฺธิํ สํสนฺทิสฺสนฺติ สเมสฺสนฺติ, อโนตฺตปฺปิโน อโนตฺตปฺปีหิ สทฺธิํ ฯเปฯ อปฺปสฺสุตา อปฺปสฺสุเตหิ สทฺธิํ ฯเปฯ กุสีตา กุสีเตหิ สทฺธิํ ฯเปฯ มุฏฺฐสฺสติโน มุฏฺฐสฺสตีหิ สทฺธิํ ฯเปฯ ทุปฺปญฺญา ทุปฺปญฺเญหิ สทฺธิํ สํสนฺทิสฺสนฺติ สเมสฺสนฺติ.}

\prose{เอตรหิปิ โข ภิกฺขเว ปจฺจุปฺปนฺนํ อทฺธานํ ธาตุโสว สตฺตา สํสนฺทนฺติ สเมนฺติ.\csromanspacer{} อสฺสทฺธา อสฺสทฺเธหิ สทฺธิํ สํสนฺทนฺติ สเมนฺติ, อหิริกา อหิริเกหิ สทฺธิํ ฯเปฯ อโนตฺตปฺปิโน อโนตฺตปฺปีหิ สทฺธิํ อปฺปสฺสุตา อปฺปสฺสุเตหิ สทฺธิํ ฯเปฯ กุสีตา กุสีเตหิ สทฺธิํ ฯเปฯ มุฏฺฐสฺสติโน มุฏฺฐสฺสตีหิ สทฺธิํ สํสนฺทนฺติ สเมนฺติ, ทุปฺปญฺญา ทุปฺปญฺเญหิ สทฺธิํ สํสนฺทนฺติ สเมนฺติ.}

\csromanmatikaanchor{mtk.124}
\csromantocmarkref{subsection}{3. คิญฺชกาวสถสุตฺต}{mtk.124}
\bookmark[dest={mtk.124},level=2]{3. คิญฺชกาวสถสุตฺต}
\prose{ธาตุโสว ภิกฺขเว สตฺตา สํสนฺทนฺติ สเมนฺติ.\csromanspacer{} สทฺธา สทฺเธหิ สทฺธิํ สํสนฺทนฺติ สเมนฺติ, หิริมนา หิริมเนหิ สทฺธิํ สํสนฺทนฺติ สเมนฺติ, โอตฺตปฺปิโน โอตฺตปฺปีหิ สทฺธิํ สํสนฺทนฺติ สเมนฺติ, พหุสฺสุตา พหุสฺสุเตหิ สทฺธิํ สํสนฺทนฺติ สเมนฺติ, อารทฺธวีริยา อารทฺธวีริเยหิ สทฺธิํ สํสนฺทนฺติ สเมนฺติ, อุปฏฺฐิตสฺสติโน อุปฏฺฐิตสฺสตีหิ สทฺธิํ สํสนฺทนฺติ สเมนฺติ, ปญฺญวนฺโต ปญฺญวนฺเตหิ สทฺธิํ สํสนฺทนฺติ สเมนฺติ.\csromanspacer{} อตีตมฺปิ โข ภิกฺขเว อทฺธานํ ฯเปฯ.\csromanspacer{} อนาคตมฺปิ โข ภิกฺขเว ฯเปฯ.\csromanspacer{} เอตรหิปิ โข ภิกฺขเว ปจฺจุปฺปนฺนํ อทฺธานํ ธาตุโสว สตฺตา สํสนฺทนฺติ สเมนฺติ.\csromanspacer{} สทฺธา สทฺเธหิ สทฺธิํ ฯเปฯ ปญฺญวนฺโต ปญฺญวนฺเตหิ สทฺธิํ สํสนฺทนฺติ สเมนฺตีติ.\csromanspacer{}\csromanspacer{}สตฺตมํ.}

\csromanheader{2}{8. อสฺสทฺธมูลกสุตฺต}
\proseitem{102}{สาวตฺถิยํ วิหรติ.\csromanspacer{} ธาตุโสว ภิกฺขเว สตฺตา สํสนฺทนฺติ สเมนฺติ.\csromanspacer{} อสฺสทฺธา อสฺสทฺเธหิ สทฺธิํ สํสนฺทนฺติ สเมนฺติ, อหิริกา อหิริเกหิ สทฺธิํ สํสนฺทนฺติ สเมนฺติ, ทุปฺปญฺญา ทุปฺปญฺเญหิ สทฺธิํ สํสนฺทนฺติ สเมนฺติ.\csromanspacer{} สทฺธา สทฺเธหิ สทฺธิํ สํสนฺทนฺติ สเมนฺติ, หิริมนา หิริมเนหิ สทฺธิํ สํสนฺทนฺติ สเมนฺติ, ปญฺญวนฺโต ปญฺญวนฺเตหิ สทฺธิํ สํสนฺทนฺติ สเมนฺติ.\csromanspacer{} อตีตมฺปิ โข ภิกฺขเว อทฺธานํ ธาตุโสว สตฺตา สํสนฺทิํสุ สมิํสุ ฯเปฯ.\csromanspacer{} อนาคตมฺปิ โข ภิกฺขเว อทฺธานํ ธาตุโสว สตฺตา สํสนฺทิสฺสนฺติ สเมสฺสนฺติ ฯเปฯ.}

\csromanheader{2}{3. ธาตุสํยุตฺต}
\prose{\csromanfolio{373}เอตรหิปิ โข ภิกฺขเว ปจฺจุปฺปนฺนํ อทฺธานํ ธาตุโสว สตฺตา สํสนฺทนฺติ สเมนฺติ.\csromanspacer{} อสฺสทฺธา อสฺสทฺเธหิ สทฺธิํ สํสนฺทนฺติ สเมนฺติ, อหิริกา อหิริเกหิ สทฺธิํ สํสนฺทนฺติ สเมนฺติ, ทุปฺปญฺญา ทุปฺปญฺเญหิ สทฺธิํ สํสนฺทนฺติ สเมนฺติ.\csromanspacer{} สทฺธา สทฺเธหิ สทฺธิํ สํสนฺทนฺติ สเมนฺติ, หิริมนา หิริมเนหิ สทฺธิํ สํสนฺทนฺติ สเมนฺติ, ปญฺญวนฺโต ปญฺญวนฺเตหิ สทฺธิํ สํสนฺทนฺติ สเมนฺตีติ. (1)}

\prose{ธาตุโสว ภิกฺขเว สตฺตา สํสนฺทนฺติ สเมนฺติ.\csromanspacer{} อสฺสทฺธา อสฺสทฺเธหิ สทฺธิํ สํสนฺทนฺติ สเมนฺติ, อโนตฺตปฺปิโน อโนตฺตปฺปีหิ สทฺธิํ สํสนฺทนฺติ สเมนฺติ, ทุปฺปญฺญา ทุปฺปญฺเญหิ สทฺธิํ สํสนฺทนฺติ สเมนฺติ.\csromanspacer{} สทฺธา สทฺเธหิ สทฺธิํ สํสนฺทนฺติ สเมนฺติ, โอตฺตปฺปิโน โอตฺตปฺปีหิ สทฺธิํ สํสนฺทนฺติ สเมนฺติ, ปญฺญวนฺโต ปญฺญวนฺเตหิ สทฺธิํ สํสนฺทนฺติ สเมนฺติ ฯเปฯ ปฐมวาโร วิย วิตฺถาเรตพฺโพ. (2)}

\csromanmatikaanchor{mtk.125}
\csromantocmarkref{subsection}{4. หีนาธิมุตฺติกสุตฺต}{mtk.125}
\bookmark[dest={mtk.125},level=2]{4. หีนาธิมุตฺติกสุตฺต}
\prose{ธาตุโสว ภิกฺขเว ฯเปฯ.\csromanspacer{} อสฺสทฺธา อสฺสทฺเธหิ สทฺธิํ สํสนฺทนฺติ สเมนฺติ, อปฺปสฺสุตา อปฺปสฺสุเตหิ สทฺธิํ สํสนฺทนฺติ สเมนฺติ, ทุปฺปญฺญา ทุปฺปญฺเญหิ สทฺธิํ สํสนฺทนฺติ สเมนฺติ.\csromanspacer{} สทฺธา สทฺเธหิ สทฺธิํ สํสนฺทนฺติ สเมนฺติ, พหุสฺสุตา พหุสฺสุเตหิ สทฺธิํ สํสนฺทนฺติ สเมนฺติ, ปญฺญวนฺโต ปญฺญวนฺเตหิ สทฺธิํ สํสนฺทนฺติ สเมนฺติ ฯเปฯ. (3)}

\prose{ธาตุโสว ภิกฺขเว ฯเปฯ.\csromanspacer{} อสฺสทฺธา อสฺสทฺเธหิ สทฺธิํ สํสนฺทนฺติ สเมนฺติ, กุสีตา กุสีเตหิ สทฺธิํ สํสนฺทนฺติ สเมนฺติ, ทุปฺปญฺญา ทุปฺปญฺเญหิ สทฺธิํ สํสนฺทนฺติ สเมนฺติ.\csromanspacer{} สทฺธา สทฺเธหิ สทฺธิํ สํสนฺทนฺติ สเมนฺติ, อารทฺธวีริยา อารทฺธวีริเยหิ สทฺธิํ สํสนฺทนฺติ สเมนฺติ, ปญฺญวนฺโต ปญฺญวนฺเตหิ สทฺธิํ สํสนฺทนฺติ สเมนฺติ ฯเปฯ. (4)}

\prose{ธาตุโสว ภิกฺขเว ฯเปฯ.\csromanspacer{} อสฺสทฺธา อสฺสทฺเธหิ สทฺธิํ สํสนฺทนฺติ สเมนฺติ, มุฏฺฐสฺสติโน มุฏฺฐสฺสตีหิ สทฺธิํ สํสนฺทนฺติ สเมนฺติ, ทุปฺปญฺญา ทุปฺปญฺเญหิ สทฺธิํ สํสนฺทนฺติ สเมนฺติ.\csromanspacer{} สทฺธา สทฺเธหิ สทฺธิํ สํสนฺทนฺติ สเมนฺติ, อุปฏฺฐิตสฺสติโน อุปฏฺฐิตสฺสตีหิ สทฺธิํ สํสนฺทนฺติ สเมนฺติ, ปญฺญวนฺโต ปญฺญวนฺเตหิ สทฺธิํ สํสนฺทนฺติ สเมนฺตีติ ฯเปฯ.\csromanspacer{}\csromanspacer{}(5) อฏฺฐมํ.}

\csromanheader{2}{9. อหิริกมูลกสุตฺต}
\proseitem{103}{สาวตฺถิยํ วิหรติ.\csromanspacer{} ธาตุโสว ฯเปฯ.\csromanspacer{} อหิริกา อหิริเกหิ สทฺธิํ สํสนฺทนฺติ สเมนฺติ, อโนตฺตปฺปิโน อโนตฺตปฺปีหิ สทฺธิํ สํสนฺทนฺติ}

\csromanmatikaanchor{mtk.126}
\csromantocmarkref{subsection}{5. จงฺกมสุตฺต}{mtk.126}
\bookmark[dest={mtk.126},level=2]{5. จงฺกมสุตฺต}
\vspace{\tipitakaparskip}
\csromancenter{นิทานวคฺคสํยุตฺตปาฬิ สเมนฺติ, ทุปฺปญฺญา ทุปฺปญฺเญหิ สทฺธิํ สํสนฺทนฺติ สเมนฺติ.\csromanspacer{} หิริมนา หิริมเนหิ สทฺธิํ สํสนฺทนฺติ สเมนฺติ, โอตฺตปฺปิโน โอตฺตปฺปีหิ สทฺธิํ สํสนฺทนฺติ สเมนฺติ, ปญฺญวนฺโต ปญฺญวนฺเตหิ สทฺธิํ สํสนฺทนฺติ สเมนฺติ ฯเปฯ. (1)}
\prose{\csromanfolio{374}อหิริกา อหิริเกหิ สทฺธิํ สํสนฺทนฺติ สเมนฺติ, อปฺปสฺสุตา อปฺปสฺสุเตหิ สทฺธิํ สํสนฺทนฺติ สเมนฺติ, ทุปฺปญฺญา ทุปฺปญฺเญหิ สทฺธิํ สํสนฺทนฺติ สเมนฺติ.\csromanspacer{} หิริมนา หิริมเนหิ สทฺธิํ สํสนฺทนฺติ สเมนฺติ, พหุสฺสุตา พหุสฺสุเตหิ สทฺธิํ สํสนฺทนฺติ สเมนฺติ, ปญฺญวนฺโต ปญฺญวนฺเตหิ สทฺธิํ สํสนฺทนฺติ สเมนฺติ ฯเปฯ. (2)}

\prose{อหิริกา อหิริเกหิ สทฺธิํ สํสนฺทนฺติ สเมนฺติ, กุสีตา กุสีเตหิ สทฺธิํ สํสนฺทนฺติ สเมนฺติ, ทุปฺปญฺญา ทุปฺปญฺเญหิ สทฺธิํ สํสนฺทนฺติ สเมนฺติ.\csromanspacer{} หิริมนา หิริมเนหิ สทฺธิํ สํสนฺทนฺติ สเมนฺติ, อารทฺธวีริยา อารทฺธวีริเยหิ สทฺธิํ สํสนฺทนฺติ สเมนฺติ, ปญฺญวนฺโต ปญฺญวนฺเตหิ สทฺธิํ สํสนฺทนฺติ สเมนฺติ ฯเปฯ. (3)}

\prose{อหิริกา อหิริเกหิ สทฺธิํ สํสนฺทนฺติ สเมนฺติ, มุฏฺฐสฺสติโน มุฏฺฐสฺสตีหิ สทฺธิํ สํสนฺทนฺติ สเมนฺติ, ทุปฺปญฺญา ทุปฺปญฺเญหิ สทฺธิํ สํสนฺทนฺติ สเมนฺติ.\csromanspacer{} หิริมนา หิริมเนหิ สทฺธิํ สํสนฺทนฺติ สเมนฺติ, อุปฏฺฐิตสฺสติโน อุปฏฺฐิตสฺสตีหิ สทฺธิํ สํสนฺทนฺติ สเมนฺติ, ปญฺญวนฺโต ปญฺญวนฺเตหิ สทฺธิํ สํสนฺทนฺติ สเมนฺตีติ ฯเปฯ.\csromanspacer{}\csromanspacer{}(4) นวมํ.}

\csromanheader{2}{10. อโนตฺตปฺปมูลกสุตฺต}
\proseitem{104}{สาวตฺถิยํ วิหรติ.\csromanspacer{} ธาตุโสว ภิกฺขเว สตฺตา สํสนฺทนฺติ สเมนฺติ.\csromanspacer{} อโนตฺตปฺปิโน อโนตฺตปฺปีหิ สทฺธิํ สํสนฺทนฺติ สเมนฺติ, อปฺปสฺสุตา อปฺปสฺสุเตหิ สทฺธิํ สํสนฺทนฺติ สเมนฺติ, ทุปฺปญฺญา ทุปฺปญฺเญหิ สทฺธิํ สํสนฺทนฺติ สเมนฺติ.\csromanspacer{} โอตฺตปฺปิโน โอตฺตปฺปีหิ สทฺธิํ สํสนฺทนฺติ สเมนฺติ, พหุสฺสุตา พหุสฺสุเตหิ สทฺธิํ สํสนฺทนฺติ สเมนฺติ, ปญฺญวนฺโต ปญฺญวนฺเตหิ สทฺธิํ สํสนฺทนฺติ สเมนฺติ ฯเปฯ. (1) อโนตฺตปฺปิโน อโนตฺตปฺปีหิ สทฺธิํ สํสนฺทนฺติ สเมนฺติ, กุสีตา กุสีเตหิ สทฺธิํ สํสนฺทนฺติ สเมนฺติ, ทุปฺปญฺญา ทุปฺปญฺเญหิ สทฺธิํ สํสนฺทนฺติ สเมนฺติ.\csromanspacer{} โอตฺตปฺปิโน โอตฺตปฺปีหิ สทฺธิํ สํสนฺทนฺติ สเมนฺติ, อารทฺธวีริยา อารทฺธวีริเยหิ สทฺธิํ สํสนฺทนฺติ สเมนฺติ, ปญฺญวนฺโต ปญฺญวนฺเตหิ สทฺธิํ สํสนฺทนฺติ สเมนฺติ ฯเปฯ. (2)}

\csromanmatikaanchor{mtk.127}
\csromantocmarkref{subsection}{6. สคาถาสุตฺต}{mtk.127}
\bookmark[dest={mtk.127},level=2]{6. สคาถาสุตฺต}
\csromanheader{2}{3. ธาตุสํยุตฺต}
\prose{\csromanfolio{375}อโนตฺตปฺปิโน อโนตฺตปฺปีหิ สทฺธิํ สํสนฺทนฺติ สเมนฺติ, มุฏฺฐสฺสติโน มุฏฺฐสฺสตีหิ สทฺธิํ สํสนฺทนฺติ สเมนฺติ, ทุปฺปญฺญา ทุปฺปญฺเญหิ สทฺธิํ สํสนฺทนฺติ สเมนฺติ.\csromanspacer{} โอตฺตปฺปิโน โอตฺตปฺปีหิ สทฺธิํ สํสนฺทนฺติ สเมนฺติ, อุปฏฺฐิตสฺสติโน อุปฏฺฐิตสฺสตีหิ สทฺธิํ สํสนฺทนฺติ สเมนฺติ, ปญฺญวนฺโต ปญฺญวนฺเตหิ สทฺธิํ สํสนฺทนฺติ สเมนฺตีติ ฯเปฯ.\csromanspacer{}\csromanspacer{}(3) ทสมํ.}

\csromanheader{2}{11. อปฺปสฺสุตมูลกสุตฺต}
\proseitem{105}{สาวตฺถิยํ วิหรติ.\csromanspacer{} ธาตุโสว ภิกฺขเว สตฺตา สํสนฺทนฺติ สเมนฺติ.\csromanspacer{} อปฺปสฺสุตา อปฺปสฺสุเตหิ สทฺธิํ สํสนฺทนฺติ สเมนฺติ, กุสีตา กุสีเตหิ สทฺธิํ สํสนฺทนฺติ สเมนฺติ, ทุปฺปญฺญา ทุปฺปญฺเญหิ สทฺธิํ สํสนฺทนฺติ สเมนฺติ.\csromanspacer{} พหุสฺสุตา พหุสฺสุเตหิ สทฺธิํ สํสนฺทนฺติ สเมนฺติ, อารทฺธวีริยา อารทฺธวีริเยหิ สทฺธิํ สํสนฺทนฺติ สเมนฺติ, ปญฺญวนฺโต ปญฺญวนฺเตหิ สทฺธิํ สํสนฺทนฺติ สเมนฺติ ฯเปฯ. (1)}

\prose{อปฺปสฺสุตา อปฺปสฺสุเตหิ สทฺธิํ สํสนฺทนฺติ สเมนฺติ, มุฏฺฐสฺสติโน มุฏฺฐสฺสตีหิ สทฺธิํ สํสนฺทนฺติ สเมนฺติ, ทุปฺปญฺญา ทุปฺปญฺเญหิ สทฺธิํ สํสนฺทนฺติ สเมนฺติ.\csromanspacer{} พหุสฺสุตา พหุสฺสุเตหิ สทฺธิํ สํสนฺทนฺติ สเมนฺติ, อุปฏฺฐิตสฺสติโน อุปฏฺฐิตสฺสตีหิ สทฺธิํ สํสนฺทนฺติ สเมนฺติ, ปญฺญวนฺโต ปญฺญวนฺเตหิ สทฺธิํ สํสนฺทนฺติ สเมนฺตีติ ฯเปฯ.\csromanspacer{}\csromanspacer{}(2) เอกาทสมํ.}

\csromanheader{2}{12. กุสีตมูลกสุตฺต}
\proseitem{106}{สาวตฺถิยํ วิหรติ.\csromanspacer{} ธาตุโสว ภิกฺขเว สตฺตา สํสนฺทนฺติ สเมนฺติ.\csromanspacer{} กุสีตา กุสีเตหิ สทฺธิํ สํสนฺทนฺติ สเมนฺติ, มุฏฺฐสฺสติโน มุฏฺฐสฺสตีหิ สทฺธิํ สํสนฺทนฺติ สเมนฺติ, ทุปฺปญฺญา ทุปฺปญฺเญหิ สทฺธิํ สํสนฺทนฺติ สเมนฺติ.\csromanspacer{} อารทฺธวีริยา อารทฺธวีริเยหิ สทฺธิํ สํสนฺทนฺติ สเมนฺติ, อุปฏฺฐิตสฺสติโน อุปฏฺฐิตสฺสตีหิ สทฺธิํ สํสนฺทนฺติ สเมนฺติ, ปญฺญวนฺโต ปญฺญวนฺเตหิ สทฺธิํ สํสนฺทนฺติ สเมนฺตีติ ฯเปฯ.\csromanspacer{}\csromanspacer{}ทฺวาทสมํ. (สพฺพตฺถ อตีตานาคตปจฺจุปฺปนฺนํ กาตพฺพํ.)}

\vspace{\tipitakaparskip}
\csromancenter{ทุติโย วคฺโค.}
\gambhira{นิทานวคฺคสํยุตฺตปาฬิ \textbf{ตสฺสุทฺทานํ}}
\csromanmatikaanchor{mtk.128}
\csromanmatikaanchor{mtk.134}
\csromantocmarkref{subsection}{7. อสฺสทฺธสํสนฺทนสุตฺต}{mtk.128}
\bookmark[dest={mtk.128},level=2]{7. อสฺสทฺธสํสนฺทนสุตฺต}
\csromangathameasure{สตฺติมา สนิทานญฺจ, คิญฺชกาวสเถน จ.}
\csromangathagroup{\csromangathastanza{\csromanfolio{376}สตฺติมา สนิทานญฺจ, คิญฺชกาวสเถน จ.}}

\csromanheader{2}{หีนาธิมุตฺติ จงฺกมํ, สคาถา อสฺสทฺธสตฺตมํ.}
\csromangathameasure{อสฺสทฺธมูลกา ปญฺจ, จตฺตาโร อหิริกมูลกา. \\ อโนตฺตปฺปมูลกา ตีณิ, ทุเว อปฺปสฺสุเตน จ. \\ กุสีตํ เอกกํ วุตฺตํ, สุตฺตนฺตา ตีณิ ปญฺจกา. \\ พาวีสติ วุตฺตา สุตฺตา, ทุติโย วคฺโค ปวุจฺจตีติ.}
\csromangathagroup{\csromangathastanza{อสฺสทฺธมูลกา ปญฺจ, จตฺตาโร อหิริกมูลกา. \\ อโนตฺตปฺปมูลกา ตีณิ, ทุเว อปฺปสฺสุเตน จ.}\csromangathastanza{กุสีตํ เอกกํ วุตฺตํ, สุตฺตนฺตา ตีณิ ปญฺจกา. \\ พาวีสติ วุตฺตา สุตฺตา, ทุติโย วคฺโค ปวุจฺจตีติ.}}

\csromanheader{1}{3. กมฺมปถวคฺค}
\csromanheader{2}{1. อสมาหิตสุตฺต}
\proseitem{107}{สาวตฺถิยํ วิหรติ.\csromanspacer{} ธาตุโสว ภิกฺขเว สตฺตา สํสนฺทนฺติ สเมนฺติ.\csromanspacer{} อสฺสทฺธา อสฺสทฺเธหิ สทฺธิํ สํสนฺทนฺติ สเมนฺติ, อหิริกา อหิริเกหิ สทฺธิํ สํสนฺทนฺติ สเมนฺติ, อโนตฺตปฺปิโน อโนตฺตปฺปีหิ สทฺธิํ สํสนฺทนฺติ สเมนฺติ, อสมาหิตา อสมาหิเตหิ สทฺธิํ สํสนฺทนฺติ สเมนฺติ, ทุปฺปญฺญา ทุปฺปญฺเญหิ สทฺธิํ สํสนฺทนฺติ สเมนฺติ.}

\prose{สทฺธา สทฺเธหิ สทฺธิํ สํสนฺทนฺติ สเมนฺติ, หิริมนา หิริมเนหิ สทฺธิํ สํสนฺทนฺติ สเมนฺติ, โอตฺตปฺปิโน โอตฺตปฺปีหิ สทฺธิํ สํสนฺทนฺติ สเมนฺติ, สมาหิตา สมาหิเตหิ สทฺธิํ สํสนฺทนฺติ สเมนฺติ, ปญฺญวนฺโต ปญฺญวนฺเตหิ สทฺธิํ สํสนฺทนฺติ สเมนฺตีติ.\csromanspacer{}\csromanspacer{}ปฐมํ.}

\csromanheader{2}{2. ทุสฺสีลสุตฺต}
\proseitem{108}{สาวตฺถิยํ วิหรติ.\csromanspacer{} ธาตุโสว ภิกฺขเว สตฺตา สํสนฺทนฺติ สเมนฺติ.\csromanspacer{} อสฺสทฺธา อสฺสทฺเธหิ สทฺธิํ สํสนฺทนฺติ สเมนฺติ, อหิริกา อหิริเกหิ สทฺธิํ สํสนฺทนฺติ สเมนฺติ, อโนตฺตปฺปิโน อโนตฺตปฺปีหิ สทฺธิํ สํสนฺทนฺติ สเมนฺติ, ทุสฺสีลา ทุสฺสีเลหิ สทฺธิํ สํสนฺทนฺติ สเมนฺติ, ทุปฺปญฺญา ทุปฺปญฺเญหิ สทฺธิํ สํสนฺทนฺติ สเมนฺติ.}

\prose{สทฺธา สทฺเธหิ สทฺธิํ สํสนฺทนฺติ สเมนฺติ, หิริมนา หิริมเนหิ สทฺธิํ สํสนฺทนฺติ สเมนฺติ, โอตฺตปฺปิโน โอตฺตปฺปีหิ สทฺธิํ สํสนฺทนฺติ สเมนฺติ,}

\vspace{\tipitakaparskip}
\csromancenter{ธาตุสํยุตฺต สีลวนฺโต สีลวนฺเตหิ สทฺธิํ สํสนฺทนฺติ สเมนฺติ, ปญฺญวนฺโต ปญฺญวนฺเตหิ สทฺธิํ สํสนฺทนฺติ สเมนฺตีติ.\csromanspacer{}\csromanspacer{}ทุติยํ.}
\csromanmatikaanchor{mtk.129}
\csromantocmarkref{subsection}{8. อสฺสทฺธมูลกสุตฺต}{mtk.129}
\bookmark[dest={mtk.129},level=2]{8. อสฺสทฺธมูลกสุตฺต}
\csromanheader{2}{3. ปญฺจสิกฺขาปทสุตฺต}
\proseitem{109}{\csromanfolio{377}สาวตฺถิยํ วิหรติ.\csromanspacer{} ธาตุโสว ภิกฺขเว สตฺตา สํสนฺทนฺติ สเมนฺติ.\csromanspacer{} ปาณาติปาติโน ปาณาติปาตีหิ สทฺธิํ สํสนฺทนฺติ สเมนฺติ, อทินฺนาทายิโน อทินฺนาทายีหิ สทฺธิํ สํสนฺทนฺติ สเมนฺติ, กาเมสุมิจฺฉาจาริโน กาเมสุมิจฺฉาจารีหิ สทฺธิํ สํสนฺทนฺติ สเมนฺติ, มุสาวาทิโน มุสาวาทีหิ สทฺธิํ สํสนฺทนฺติ สเมนฺติ, สุราเมรยมชฺชปฺปมาทฏฺฐายิโน สุราเมรยมชฺชปฺปมาทฏฺฐายีหิ สทฺธิํ สํสนฺทนฺติ สเมนฺติ.}

\prose{ปาณาติปาตา ปฏิวิรตา ปาณาติปาตา ปฏิวิรเตหิ สทฺธิํ สํสนฺทนฺติ สเมนฺติ, อทินฺนาทานา ปฏิวิรตา อทินฺนาทานา ปฏิวิรเตหิ สทฺธิํ สํสนฺทนฺติ สเมนฺติ, กาเมสุมิจฺฉาจารา ปฏิวิรตา กาเมสุมิจฺฉาจารา ปฏิวิรเตหิ สทฺธิํ สํสนฺทนฺติ สเมนฺติ, มุสาวาทา ปฏิวิรตา มุสาวาทา ปฏิวิรเตหิ สทฺธิํ สํสนฺทนฺติ สเมนฺติ, สุราเมรยมชฺชปฺปมาทฏฺฐานา ปฏิวิรตา สุราเมรยมชฺชปฺปมาทฏฺฐานา ปฏิวิรเตหิ สทฺธิํ สํสนฺทนฺติ สเมนฺตีติ.\csromanspacer{}\csromanspacer{}ตติยํ.}

\csromanheader{2}{4. สตฺตกมฺมปถสุตฺต}
\proseitem{110}{สาวตฺถิยํ วิหรติ.\csromanspacer{} ธาตุโสว ภิกฺขเว สตฺตา สํสนฺทนฺติ สเมนฺติ.\csromanspacer{} ปาณาติปาติโน ปาณาติปาตีหิ สทฺธิํ สํสนฺทนฺติ สเมนฺติ, อทินฺนาทายิโน อทินฺนาทายีหิ สทฺธิํ สํสนฺทนฺติ สเมนฺติ, กาเมสุมิจฺฉาจาริโน กาเมสุมิจฺฉาจารีหิ สทฺธิํ สํสนฺทนฺติ สเมนฺติ, มุสาวาทิโน มุสาวาทีหิ สทฺธิํ สํสนฺทนฺติ สเมนฺติ, ปิสุณวาจา ปิสุณวาเจหิ สทฺธิํ สํสนฺทนฺติ สเมนฺติ, ผรุสวาจา ผรุสวาเจหิ สทฺธิํ สํสนฺทนฺติ สเมนฺติ, สมฺผปฺปลาปิโน สมฺผปฺปลาปีหิ สทฺธิํ สํสนฺทนฺติ สเมนฺติ.}

\prose{ปาณาติปาตา ปฏิวิรตา ฯเปฯ อทินฺนาทานา ปฏิวิรตา.\csromanspacer{} กาเมสุมิจฺฉาจารา ปฏิวิรตา.\csromanspacer{} มุสาวาทา ปฏิวิรตา.\csromanspacer{} ปิสุณาย วาจาย ปฏิวิรตา ปิสุณาย วาจาย ปฏิวิรเตหิ สทฺธิํ สํสนฺทนฺติ สเมนฺติ, ผรุสาย วาจาย ปฏิวิรตา ผรุสาย วาจาย ปฏิวิรเตหิ สทฺธิํ สํสนฺทนฺติ สเมนฺติ, สมฺผปฺปลาปา ปฏิวิรตา สมฺผปฺปลาปา ปฏิวิรเตหิ สทฺธิํ สํสนฺทนฺติ สเมนฺตีติ.\csromanspacer{}\csromanspacer{}จตุตฺถํ.}

\csromanheader{2}{นิทานวคฺคสํยุตฺตปาฬิ \textbf{5. ทสกมฺมปถสุตฺต}}
\proseitem{111}{\csromanfolio{378}สาวตฺถิยํ วิหรติ.\csromanspacer{} ธาตุโสว ภิกฺขเว สตฺตา สํสนฺทนฺติ สเมนฺติ.\csromanspacer{} ปาณาติปาติโน ปาณาติปาตีหิ สทฺธิํ สํสนฺทนฺติ สเมนฺติ, อทินฺนาทายิโน ฯเปฯ.\csromanspacer{} กาเมสุมิจฺฉาจาริโน.\csromanspacer{} มุสาวาทิโน.\csromanspacer{} ปิสุณวาจา.\csromanspacer{} ผรุสวาจา.\csromanspacer{} สมฺผปฺปลาปิโน สมฺผปฺปลาปีหิ สทฺธิํ สํสนฺทนฺติ สเมนฺติ, อภิชฺฌาลุโน อภิชฺฌาลูหิ สทฺธิํ สํสนฺทนฺติ สเมนฺติ, พฺยาปนฺนจิตฺตา พฺยาปนฺนจิตฺเตหิ สทฺธิํ สํสนฺทนฺติ สเมนฺติ, มิจฺฉาทิฏฺฐิกา มิจฺฉาทิฏฺฐิเกหิ สทฺธิํ สํสนฺทนฺติ สเมนฺติ.}

\csromanmatikaanchor{mtk.130}
\csromantocmarkref{subsection}{9. อหิริกมูลกสุตฺต}{mtk.130}
\bookmark[dest={mtk.130},level=2]{9. อหิริกมูลกสุตฺต}
\prose{ปาณาติปาตา ปฏิวิรตา ปาณาติปาตา ปฏิวิรเตหิ สทฺธิํ สํสนฺทนฺติ สเมนฺติ, อทินฺนาทานา ปฏิวิรตา ฯเปฯ.\csromanspacer{} กาเมสุมิจฺฉาจารา ปฏิวิรตา.\csromanspacer{} มุสาวาทา ปฏิวิรตา.\csromanspacer{} ปิสุณาย วาจาย.\csromanspacer{} ผรุสาย วาจาย.\csromanspacer{} สมฺผปฺปลาปา ปฏิวิรตา สมฺผปฺปลาปา ปฏิวิรเตหิ สทฺธิํ สํสนฺทนฺติ สเมนฺติ, อนภิชฺฌาลุโน อนภิชฺฌาลูหิ สทฺธิํ สํสนฺทนฺติ สเมนฺติ, อพฺยาปนฺนจิตฺตา อพฺยาปนฺนจิตฺเตหิ สทฺธิํ สํสนฺทนฺติ สเมนฺติ, สมฺมาทิฏฺฐิกา สมฺมาทิฏฺฐิเกหิ สทฺธิํ สํสนฺทนฺติ สเมนฺตีติ.\csromanspacer{}\csromanspacer{}ปญฺจมํ.}

\csromanheader{2}{6. อฏฺฐงฺคิกสุตฺต}
\proseitem{112}{สาวตฺถิยํ วิหรติ.\csromanspacer{} ธาตุโสว ภิกฺขเว สตฺตา สํสนฺทนฺติ สเมนฺติ.\csromanspacer{} มิจฺฉาทิฏฺฐิกา มิจฺฉาทิฏฺฐิเกหิ สทฺธิํ สํสนฺทนฺติ สเมนฺติ, มิจฺฉาสงฺกปฺปา ฯเปฯ มิจฺฉาวาจา.\csromanspacer{} มิจฺฉากมฺมนฺตา.\csromanspacer{} มิจฺฉา-อาชีวา.\csromanspacer{} มิจฺฉาวายามา.\csromanspacer{} มิจฺฉาสติโน.\csromanspacer{} มิจฺฉาสมาธิโน มิจฺฉาสมาธีหิ สทฺธิํ สํสนฺทนฺติ สเมนฺติ.\csromanspacer{} สมฺมาทิฏฺฐิกา สมฺมาทิฏฺฐิเกหิ สทฺธิํ สํสนฺทนฺติ สเมนฺติ, สมฺมาสงฺกปฺปา ฯเปฯ สมฺมาวาจา.\csromanspacer{} สมฺมากมฺมนฺตา.\csromanspacer{} สมฺมา-อาชีวา.\csromanspacer{} สมฺมาวายามา.\csromanspacer{} สมฺมาสติโน.\csromanspacer{} สมฺมาสมาธิโน สมฺมาสมาธีหิ สทฺธิํ สํสนฺทนฺติ สเมนฺตีติ.\csromanspacer{}\csromanspacer{}ฉฏฺฐํ.}

\csromanheader{2}{7. ทสงฺคสุตฺต}
\proseitem{113}{สาวตฺถิยํ วิหรติ.\csromanspacer{} ธาตุโสว ภิกฺขเว สตฺตา สํสนฺทนฺติ สเมนฺติ.\csromanspacer{} มิจฺฉาทิฏฺฐิกา มิจฺฉาทิฏฺฐิเกหิ สทฺธิํ สํสนฺทนฺติ สเมนฺติ มิจฺฉาสงฺกปฺปา ฯเปฯ มิจฺฉาวาจา.\csromanspacer{} มิจฺฉากมฺมนฺตา.\csromanspacer{} มิจฺฉา-อาชีวา.\csromanspacer{} มิจฺฉาวายามา.}

\vspace{\tipitakaparskip}
\csromancenter{ธาตุสํยุตฺต มิจฺฉาสติโน.\csromanspacer{} มิจฺฉาสมาธิโน มิจฺฉาสมาธีหิ สทฺธิํ สํสนฺทนฺติ สเมนฺติ, มิจฺฉาญาณิโน มิจฺฉาญาณีหิ สทฺธิํ สํสนฺทนฺติ สเมนฺติ, มิจฺฉาวิมุตฺติโน มิจฺฉาวิมุตฺตีหิ สทฺธิํ สํสนฺทนฺติ สเมนฺติ.}
\csromanmatikaanchor{mtk.131}
\csromanmatikaanchor{mtk.143}
\csromantocmarkref{subsection}{10. อโนตฺตปฺปมูลกสุตฺต}{mtk.131}
\bookmark[dest={mtk.131},level=2]{10. อโนตฺตปฺปมูลกสุตฺต}
\prose{\csromanfolio{379}สมฺมาทิฏฺฐิกา สมฺมาทิฏฺฐิเกหิ สทฺธิํ สํสนฺทนฺติ สเมนฺติ, สมฺมาสงฺกปฺปา ฯเปฯ สมฺมาวาจา.\csromanspacer{} สมฺมากมฺมนฺตา.\csromanspacer{} สมฺมา-อาชีวา.\csromanspacer{} สมฺมาวายามา.\csromanspacer{} สมฺมาสติโน.\csromanspacer{} สมฺมาสมาธิโน.\csromanspacer{} สมฺมาญาณิโน สมฺมาญาณีหิ สทฺธิํ สํสนฺทนฺติ สเมนฺติ, สมฺมาวิมุตฺติโน สมฺมาวิมุตฺตีหิ สทฺธิํ สํสนฺทนฺติ สเมนฺตีติ.\csromanspacer{}\csromanspacer{}สตฺตมํ. (สพฺพตฺถ อตีตานาคตปจฺจุปฺปนฺนํ กาตพฺพํ.)}

\csromanheader{2}{\textbf{สตฺตนฺนํ สุตฺตนฺตานํ อุทฺทานํ}}
\csromangathameasure{อสมาหิตํ ทุสฺสีลํ, ปญฺจสิกฺขาปทานิ จ. \\ สตฺตกมฺมปถา วุตฺตา, ทสกมฺมปเถน จ. \\ ฉฏฺฐํ อฏฺฐงฺคิโก วุตฺโต, ทสงฺเคน จ สตฺตมํ. \\ กมฺมปถวคฺโค ตติโย.}
\csromangathagroup{\csromangathastanza{อสมาหิตํ ทุสฺสีลํ, ปญฺจสิกฺขาปทานิ จ. \\ สตฺตกมฺมปถา วุตฺตา, ทสกมฺมปเถน จ.}\csromangathastanza{ฉฏฺฐํ อฏฺฐงฺคิโก วุตฺโต, ทสงฺเคน จ สตฺตมํ. \\ กมฺมปถวคฺโค ตติโย.}}

\chapterhead{4. จตุตฺถวคฺค}
\csromanmatikaanchor{mtk.132}
\csromantocmarkref{subsection}{11. อปฺปสฺสุตมูลกสุตฺต}{mtk.132}
\bookmark[dest={mtk.132},level=2]{11. อปฺปสฺสุตมูลกสุตฺต}
\csromanheader{2}{1. จตุธาตุสุตฺต}
\proseitem{114}{เอกํ สมยํ ภควา สาวตฺถิยํ วิหรติ เชตวเน อนาถปิณฺฑิกสฺส อาราเม.\csromanspacer{} จตสฺโส อิมา ภิกฺขเว ธาตุโย.\csromanspacer{} กตมา จตสฺโส, ปถวีธาตุ อาโปธาตุ เตโชธาตุ วาโยธาตุ.\csromanspacer{} อิมา โข ภิกฺขเว จตสฺโส ธาตุโยติ.\csromanspacer{}\csromanspacer{}ปฐมํ.}

\csromanheader{2}{2. ปุพฺเพสมฺโพธสุตฺต}
\csromanmatikaanchor{mtk.133}
\csromantocmarkref{subsection}{12. กุสีตมูลกสุตฺต}{mtk.133}
\bookmark[dest={mtk.133},level=2]{12. กุสีตมูลกสุตฺต}
\csromantocmarkref{section}{อุทฺทานคาถา}{mtk.134}
\bookmark[dest={mtk.134},level=1]{อุทฺทานคาถา}
\proseitem{115}{สาวตฺถิยํ วิหรติ.\csromanspacer{} ปุพฺเพว เม ภิกฺขเว สมฺโพธา อนภิสมฺพุทฺธสฺส โพธิสตฺตสฺเสว สโต เอตทโหสิ “โก นุ โข ปถวีธาตุยา อสฺสาโท โก อาทีนโว กิํ นิสฺสรณํ, โก อาโปธาตุยา อสฺสาโท โก อาทีนโว กิํ นิสฺสรณํ, โก}

\vspace{\tipitakaparskip}
\csromancenter{นิทานวคฺคสํยุตฺตปาฬิ เตโชธาตุยา อสฺสาโท โก อาทีนโว กิํ นิสฺสรณํ, โก วาโยธาตุยา อสฺสาโท โก อาทีนโว กิํ นิสฺสรณนฺ”ติ.}
\prose{\csromanfolio{380}ตสฺส มยฺหํ ภิกฺขเว เอตทโหสิ “ยํ โข ปถวีธาตุํ ปฏิจฺจ อุปฺปชฺชติ สุขํ โสมนสฺสํ, อยํ ปถวีธาตุยา อสฺสาโท.\csromanspacer{} ยํ\footnote{ยา (สี)} ปถวีธาตุ อนิจฺจา ทุกฺขา วิปริณามธมฺมา, อยํ ปถวีธาตุยา อาทีนโว.\csromanspacer{} โย ปถวีธาตุยา ฉนฺทราควินโย ฉนฺทราคปฺปหานํ, อิทํ ปถวีธาตุยา นิสฺสรณํ.\csromanspacer{} ยํ อาโปธาตุํ ปฏิจฺจ ฯเปฯ.\csromanspacer{} ยํ เตโชธาตุํ ปฏิจฺจ ฯเปฯ.\csromanspacer{} ยํ วาโยธาตุํ ปฏิจฺจ อุปฺปชฺชติ สุขํ โสมนสฺสํ, อยํ วาโยธาตุยา อสฺสาโท.\csromanspacer{} ยํ วาโยธาตุ อนิจฺจา ทุกฺขา วิปริณามธมฺมา, อยํ วาโยธาตุยา อาทีนโว.\csromanspacer{} โย วาโยธาตุยา ฉนฺทราควินโย ฉนฺทราคปฺปหานํ, อิทํ วาโยธาตุยา นิสฺสรณํ.}

\prose{ยาวกีวญฺจาหํ ภิกฺขเว อิมาสํ จตุนฺนํ ธาตูนํ เอวํ อสฺสาทญฺจ อสฺสาทโต อาทีนวญฺจ อาทีนวโต นิสฺสรณญฺจ นิสฺสรณโต ยถาภูตํ น อพฺภญฺญาสิํ, เนว ตาวาหํ ภิกฺขเว สเทวเก โลเก สมารเก สพฺรหฺมเก สสฺสมณพฺราหฺมณิยา ปชาย สเทวมนุสฺสาย อนุตฺตรํ สมฺมาสมฺโพธิํ อภิสมฺพุทฺโธติ\footnote{อภิสมฺพุทฺโธ (สี, สฺยา, กํ)} ปจฺจญฺญาสิํ.}

\prose{ยโต จ ขฺวาหํ ภิกฺขเว อิมาสํ จตุนฺนํ ธาตูนํ เอวํ อสฺสาทญฺจ อสฺสาทโต อาทีนวญฺจ อาทีนวโต นิสฺสรณญฺจ นิสฺสรณโต ยถาภูตํ อพฺภญฺญาสิํ, อถาหํ ภิกฺขเว สเทวเก โลเก สมารเก สพฺรหฺมเก สสฺสมณพฺราหฺมณิยา ปชาย สเทวมนุสฺสาย อนุตฺตรํ สมฺมาสมฺโพธิํ อภิสมฺพุทฺโธติ ปจฺจญฺญาสิํ.\csromanspacer{} ญาณญฺจ ปน เม ทสฺสนํ อุทปาทิ, อกุปฺปา เม วิมุตฺติ\footnote{เจโตวิมุตฺติ (สี, อิ, ก)}, อยมนฺติมา ชาติ, นตฺถิ ทานิ ปุนพฺภโว”ติ.\csromanspacer{}\csromanspacer{}ทุติยํ.}

\csromanheader{2}{3. อจริํสุตฺต}
\proseitem{116}{สาวตฺถิยํ วิหรติ.\csromanspacer{} ปถวีธาตุยาหํ ภิกฺขเว อสฺสาทปริเยสนํ อจริํ, โย ปถวีธาตุยา อสฺสาโท ตทชฺฌคมํ, ยาวตา ปถวีธาตุยา อสฺสาโท ปญฺญาย เม โส สุทิฏฺโฐ.\csromanspacer{} ปถวีธาตุยาหํ ภิกฺขเว อาทีนวปริเยสนํ อจริํ, โย ปถวีธาตุยา อาทีนโว}

\vspace{\tipitakaparskip}
\csromancenter{ธาตุสํยุตฺต ตทชฺฌคมํ, ยาวตา ปถวีธาตุยา อาทีนโว ปญฺญาย เม โส สุทิฏฺโฐ.\csromanspacer{} ปถวีธาตุยาหํ ภิกฺขเว นิสฺสรณปริเยสนํ อจริํ, ยํ ปถวีธาตุยา นิสฺสรณํ ตทชฺฌคมํ, ยาวตา ปถวีธาตุยา นิสฺสรณํ ปญฺญาย เม ตํ สุทิฏฺฐํ.}
\csromanmatikaanchor{mtk.135}
\csromantocmarknopage{section}{3. กมฺมปถวคฺค}{mtk.135}
\bookmark[dest={mtk.135},level=1]{3. กมฺมปถวคฺค}
\prose{\csromanfolio{381}อาโปธาตุยาหํ ภิกฺขเว ฯเปฯ.\csromanspacer{} เตโชธาตุยาหํ ภิกฺขเว.\csromanspacer{} วาโยธาตุยาหํ ภิกฺขเว อสฺสาทปริเยสนํ อจริํ, โย วาโยธาตุยา อสฺสาโท ตทชฺฌคมํ, ยาวตา วาโยธาตุยา อสฺสาโท ปญฺญาย เม โส สุทิฏฺโฐ.\csromanspacer{} วาโยธาตุยาหํ ภิกฺขเว อาทีนวปริเยสนํ อจริํ, โย วาโยธาตุยา อาทีนโว ตทชฺฌคมํ, ยาวตา วาโยธาตุยา อาทีนโว ปญฺญาย เม โส สุทิฏฺโฐ.\csromanspacer{} วาโยธาตุยาหํ ภิกฺขเว นิสฺสรณปริเยสนํ อจริํ, ยํ วาโยธาตุยา นิสฺสรณํ ตทชฺฌคมํ, ยาวตา วาโยธาตุยา นิสฺสรณํ ปญฺญาย เม ตํ สุทิฏฺฐํ.}

\csromanmatikaanchor{mtk.136}
\csromantocmarkref{subsection}{1. อสมาหิตสุตฺต}{mtk.136}
\bookmark[dest={mtk.136},level=2]{1. อสมาหิตสุตฺต}
\prose{ยาวกีวญฺจาหํ ภิกฺขเว อิมาสํ จตุนฺนํ ธาตูนํ อสฺสาทญฺจ อสฺสาทโต อาทีนวญฺจ อาทีนวโต นิสฺสรณญฺจ นิสฺสรณโต ยถาภูตํ น อพฺภญฺญาสิํ, เนว ตาวาหํ ภิกฺขเว สเทวเก โลเก สมารเก สพฺรหฺมเก สสฺสมณพฺราหฺมณิยา ปชาย สเทวมนุสฺสาย อนุตฺตรํ สมฺมาสมฺโพธิํ อภิสมฺพุทฺโธติ ปจฺจญฺญาสิํ.}

\prose{ยโต จ ขฺวาหํ ภิกฺขเว อิมาสํ จตุนฺนํ ธาตูนํ อสฺสาทญฺจ อสฺสาทโต อาทีนวญฺจ อาทีนวโต นิสฺสรณญฺจ นิสฺสรณโต ยถาภูตํ อพฺภญฺญาสิํ, อถาหํ ภิกฺขเว สเทวเก โลเก สมารเก สพฺรหฺมเก สสฺสมณพฺราหฺมณิยา ปชาย สเทวมนุสฺสาย อนุตฺตรํ สมฺมาสมฺโพธิํ อภิสมฺพุทฺโธติ ปจฺจญฺญาสิํ.\csromanspacer{} ญณญฺจ ปน เม ทสฺสนํ อุทปาทิ, อกุปฺปา เม วิมุตฺติ, อยมนฺติมา ชาติ, นตฺถิ ทานิ ปุนพฺภโวติ.\csromanspacer{}\csromanspacer{}ตติยํ.}

\csromanheader{2}{4. โนเจทํสุตฺต}
\csromanmatikaanchor{mtk.137}
\csromantocmarkref{subsection}{2. ทุสฺสีลสุตฺต}{mtk.137}
\bookmark[dest={mtk.137},level=2]{2. ทุสฺสีลสุตฺต}
\proseitem{117}{สาวตฺถิยํ วิหรติ.\csromanspacer{} โน เจทํ ภิกฺขเว ปถวีธาตุยา อสฺสาโท อภวิสฺส, นยิทํ สตฺตา ปถวีธาตุยา สารชฺเชยฺยุํ.}

\vspace{\tipitakaparskip}
\csromancenter{นิทานวคฺคสํยุตฺตปาฬิ ยสฺมา จ โข ภิกฺขเว อตฺถิ ปถวีธาตุยา อสฺสาโท, ตสฺมา สตฺตา ปถวีธาตุยา สารชฺชนฺติ.\csromanspacer{} โน เจทํ ภิกฺขเว ปถวีธาตุยา อาทีนโว อภวิสฺส, นยิทํ สตฺตา ปถวีธาตุยา นิพฺพินฺเทยฺยุํ.\csromanspacer{} ยสฺมา จ โข ภิกฺขเว อตฺถิ ปถวีธาตุยา อาทีนโว, ตสฺมา สตฺตา ปถวีธาตุยา นิพฺพินฺทนฺติ.\csromanspacer{} โน เจทํ ภิกฺขเว ปถวีธาตุยา นิสฺสรณํ อภวิสฺส, นยิทํ สตฺตา ปถวีธาตุยา นิสฺสเรยฺยุํ.\csromanspacer{} ยสฺมา จ โข ภิกฺขเว อตฺถิ ปถวีธาตุยา นิสฺสรณํ, ตสฺมา สตฺตา ปถวีธาตุยา นิสฺสรนฺติ.}
\prose{\csromanfolio{382}โน เจทํ ภิกฺขเว อาโปธาตุยา อสฺสาโท อภวิสฺส.\csromanspacer{} โน เจทํ ภิกฺขเว เตโชธาตุยา ฯเปฯ.\csromanspacer{} โน เจทํ ภิกฺขเว วาโยธาตุยา อสฺสาโท อภวิสฺส, นยิทํ สตฺตา วาโยธาตุยา สารชฺเชยฺยุํ.\csromanspacer{} ยสฺมา จ โข ภิกฺขเว อตฺถิ วาโยธาตุยา อสฺสาโท, ตสฺมา สตฺตา วาโยธาตุยา สารชฺชนฺติ.\csromanspacer{} โน เจทํ ภิกฺขเว วาโยธาตุยา อาทีนโว อภวิสฺส, นยิทํ สตฺตา วาโยธาตุยา นิพฺพินฺเทยฺยุํ.\csromanspacer{} ยสฺมา จ โข ภิกฺขเว อตฺถิ วาโยธาตุยา อาทีนโว, ตสฺมา สตฺตา วาโยธาตุยา นิพฺพินฺทนฺติ.\csromanspacer{} โน เจทํ ภิกฺขเว วาโยธาตุยา นิสฺสรณํ อภวิสฺส, นยิทํ สตฺตา วาโยธาตุยา นิสฺสเรยฺยุํ.\csromanspacer{} ยสฺมา จ โข ภิกฺขเว อตฺถิ วาโยธาตุยา นิสฺสรณํ, ตสฺมา สตฺตา วาโยธาตุยา นิสฺสรนฺติ.}

\prose{ยาวกีวญฺจิเม ภิกฺขเว สตฺตา อิมาสํ จตุนฺนํ ธาตูนํ อสฺสาทญฺจ อสฺสาทโต อาทีนวญฺจ อาทีนวโต นิสฺสรณญฺจ นิสฺสรณโต ยถาภูตํ น อพฺภญฺญํสุ, เนว ตาวิเม ภิกฺขเว สตฺตา สเทวกา โลกา สมารกา สพฺรหฺมกา สสฺสมณพฺราหฺมณิยา ปชาย สเทวมนุสฺสาย นิสฺสฏา วิสํยุตฺตา วิปฺปมุตฺตา วิมริยาทิกเตน เจตสา วิหริํสุ.}

\csromanmatikaanchor{mtk.138}
\csromantocmarkref{subsection}{3. ปญฺจสิกฺขาปทสุตฺต}{mtk.138}
\bookmark[dest={mtk.138},level=2]{3. ปญฺจสิกฺขาปทสุตฺต}
\prose{ยโต จ โข ภิกฺขเว สตฺตา อิมาสํ จตุนฺนํ ธาตูนํ อสฺสาทญฺจ อสฺสาทโต อาทีนวญฺจ อาทีนวโต นิสฺสรณญฺจ นิสฺสรณโต ยถาภูตํ อพฺภญฺญํสุ, อถ ภิกฺขเว สตฺตา สเทวกา โลกา สมารกา สพฺรหฺมกา สสฺสมณพฺราหฺมณิยา ปชาย สเทวมนุสฺสาย นิสฺสฏา วิสํยุตฺตา วิปฺปมุตฺตา วิมริยาทิกเตน เจตสา วิหรนฺตีติ.\csromanspacer{}\csromanspacer{}จตุตฺถํ.}

\csromanheader{2}{3. ธาตุสํยุตฺต \textbf{5. เอกนฺตทุกฺขสุตฺต}}
\proseitem{118}{\csromanfolio{383}สาวตฺถิยํ วิหรติ.\csromanspacer{} ปถวีธาตุ เจ\footnote{จ (สี, สฺยา, กํ)} หิทํ ภิกฺขเว เอกนฺตทุกฺขา อภวิสฺส ทุกฺขานุปติตา ทุกฺขาวกฺกนฺตา อนวกฺกนฺตา สุเขน, นยิทํ สตฺตา ปถวีธาตุยา สารชฺเชยฺยุํ.\csromanspacer{} ยสฺมา จ โข ภิกฺขเว ปถวีธาตุ สุขา สุขานุปติตา สุขาวกฺกนฺตา อนวกฺกนฺตา ทุกฺเขน, ตสฺมา สตฺตา ปถวีธาตุยา สารชฺชนฺติ.}

\csromanmatikaanchor{mtk.139}
\csromantocmarkref{subsection}{4. สตฺตกมฺมปถสุตฺต}{mtk.139}
\bookmark[dest={mtk.139},level=2]{4. สตฺตกมฺมปถสุตฺต}
\prose{อาโปธาตุ เจ หิทํ ภิกฺขเว ฯเปฯ.\csromanspacer{} เตโชธาตุ เจ หิทํ ภิกฺขเว.\csromanspacer{} วาโยธาตุ เจ หิทํ ภิกฺขเว เอกนฺตทุกฺขา อภวิสฺส ทุกฺขานุปติตา ทุกฺขาวกฺกนฺตา อนวกฺกนฺตา สุเขน, นยิทํ สตฺตา วาโยธาตุยา สารชฺเชยฺยุํ.\csromanspacer{} ยสฺมา จ โข ภิกฺขเว วาโยธาตุ สุขา สุขานุปติตา สุขาวกฺกนฺตา อนวกฺกนฺตา ทุกฺเขน, ตสฺมา สตฺตา วาโยธาตุยา สารชฺชนฺติ.}

\prose{ปถวีธาตุ เจ หิทํ ภิกฺขเว เอกนฺตสุขา อภวิสฺส สุขานุปติตา สุขาวกฺกนฺตา อนวกฺกนฺตา ทุกฺเขน, นยิทํ สตฺตา ปถวีธาตุยา นิพฺพินฺเทยฺยุํ.\csromanspacer{} ยสฺมา จ โข ภิกฺขเว ปถวีธาตุ ทุกฺขา ทุกฺขานุปติตา ทุกฺขาวกฺกนฺตา อนวกฺกนฺตา สุเขน, ตสฺมา สตฺตา ปถวีธาตุยา นิพฺพินฺทนฺติ.}

\prose{อาโปธาตุ เจ หิทํ ภิกฺขเว ฯเปฯ.\csromanspacer{} เตโชธาตุ เจ หิทํ ภิกฺขเว.\csromanspacer{} วาโยธาตุ เจ หิทํ ภิกฺขเว เอกนฺตสุขา อภวิสฺส สุขานุปติตา สุขาวกฺกนฺตา อนวกฺกนฺตา ทุกฺเขน, นยิทํ สตฺตา วาโยธาตุยา นิพฺพินฺเทยฺยุํ.\csromanspacer{} ยสฺมา จ โข ภิกฺขเว วาโยธาตุ ทุกฺขา ทุกฺขานุปติตา ทุกฺขาวกฺกนฺตา อนวกฺกนฺตา สุเขน, ตสฺมา สตฺตา วาโยธาตุยา นิพฺพินฺทนฺตีติ.\csromanspacer{}\csromanspacer{}ปญฺจมํ.}

\csromanmatikaanchor{mtk.140}
\csromantocmarkref{subsection}{5. ทสกมฺมปถสุตฺต}{mtk.140}
\bookmark[dest={mtk.140},level=2]{5. ทสกมฺมปถสุตฺต}
\csromanheader{2}{6. อภินนฺทสุตฺต}
\proseitem{119}{สาวตฺถิยํ วิหรติ.\csromanspacer{} โย ภิกฺขเว ปถวีธาตุํ อภินนฺทติ, ทุกฺขํ โส อภินนฺทติ.\csromanspacer{} โย ทุกฺขํ อภินนฺทติ, อปริมุตฺโต โส ทุกฺขสฺมาติ วทามิ.\csromanspacer{} โย อาโปธาตุํ อภินนฺทติ ฯเปฯ.\csromanspacer{} โย เตโชธาตุํ.\csromanspacer{} โย วาโยธาตุํ อภินนฺทติ, ทุกฺขํ โส อภินนฺทติ.\csromanspacer{} โย ทุกฺขํ อภินนฺทติ, อปริมุตฺโต โส ทุกฺขสฺมาติ วทามิ.}

\gambhira{นิทานวคฺคสํยุตฺตปาฬิ}
\csromanmatikaanchor{mtk.141}
\csromantocmarkref{subsection}{6. อฏฺฐงฺคิกสุตฺต}{mtk.141}
\bookmark[dest={mtk.141},level=2]{6. อฏฺฐงฺคิกสุตฺต}
\prose{\csromanfolio{384}โย จ โข ภิกฺขเว ปถวีธาตุํ นาภินนฺทติ, ทุกฺขํ โส นาภินนฺทติ.\csromanspacer{} โย ทุกฺขํ นาภินนฺทติ, ปริมุตฺโต โส ทุกฺขสฺมาติ วทามิ.\csromanspacer{} โย อาโปธาตุํ ฯเปฯ.\csromanspacer{} โย เตโชธาตุํ.\csromanspacer{} โย วาโยธาตุํ นาภินนฺทติ, ทุกฺขํ โส นาภินนฺทติ.\csromanspacer{} โย ทุกฺขํ นาภินนฺทติ, ปริมุตฺโต โส ทุกฺขสฺมาติํ วทามีติ.\csromanspacer{}\csromanspacer{}ฉฏฺฐํ.}

\csromanheader{2}{7. อุปฺปาทสุตฺต}
\csromanmatikaanchor{mtk.142}
\csromantocmarkref{subsection}{7. ทสงฺคสุตฺต}{mtk.142}
\bookmark[dest={mtk.142},level=2]{7. ทสงฺคสุตฺต}
\csromantocmarkref{section}{อุทฺทานคาถา}{mtk.143}
\bookmark[dest={mtk.143},level=1]{อุทฺทานคาถา}
\proseitem{120}{สาวตฺถิยํ วิหรติ.\csromanspacer{} โย ภิกฺขเว ปถวีธาตุยา อุปฺปาโท ฐิติ อภินิพฺพตฺติ ปาตุภาโว, ทุกฺขสฺเสโส อุปฺปาโท โรคานํ ฐิติ ชรามรณสฺส ปาตุภาโว.\csromanspacer{} โย อาโปธาตุยา ฯเปฯ.\csromanspacer{} โย เตโชธาตุยา.\csromanspacer{} โย วาโยธาตุยา อุปฺปาโท ฐิติ อภินิพฺพตฺติ ปาตุภาโว, ทุกฺขสฺเสโส อุปฺปาโท โรคานํ ฐิติ ชรามรณสฺส ปาตุภาโว.}

\prose{โย จ โข ภิกฺขเว ปถวีธาตุยา นิโรโธ วูปสโม อตฺถงฺคโม, ทุกฺขสฺเสโส นิโรโธ โรคานํ วูปสโม ชรามรณสฺส อตฺถงฺคโม.\csromanspacer{} โย อาโปธาตุยา ฯเปฯ.\csromanspacer{} โย เตโชธาตุยา.\csromanspacer{} โย วาโยธาตุยา นิโรโธ วูปสโม อตฺถงฺคโม, ทุกฺขสฺเสโส นิโรโธ โรคานํ วูปสโม ชรามรณสฺส อตฺถงฺคโมติ.\csromanspacer{}\csromanspacer{}สตฺตมํ.}

\csromanheader{2}{8. สมณพฺราหฺมณสุตฺต}
\proseitem{121}{สาวตฺถิยํ วิหรติ.\csromanspacer{} จตสฺโส อิมา ภิกฺขเว ธาตุโย.\csromanspacer{} กตมา จตสฺโส, ปถวีธาตุ อาโปธาตุ เตโชธาตุ วาโยธาตุ.\csromanspacer{} เย หิ เกจิ ภิกฺขเว สมณา วา พฺราหฺมณา วา อิมาสํ จตุนฺนํ ธาตูนํ อสฺสาทญฺจ อาทีนวญฺจ นิสฺสรณญฺจ ยถาภูตํ นปฺปชานนฺติ.\csromanspacer{} น เม เต ภิกฺขเว สมณา วา พฺราหฺมณา วา สมเณสุ วา สมณสมฺมตา, พฺราหฺมเณสุ วา พฺราหฺมณสมฺมตา.\csromanspacer{} น จ ปน เต อายสฺมนฺโต สามญฺญตฺถํ วา พฺรหฺมญฺญตฺถํ วา ทิฏฺเฐว ธมฺเม สยํ อภิญฺญา สจฺฉิกตฺวา อุปสมฺปชฺช วิหรนฺติ.}

\prose{เย จ โข เกจิ ภิกฺขเว สมณา วา พฺราหฺมณา วา อิมาสํ จตุนฺนํ ธาตูนํ อสฺสาทญฺจ อาทีนวญฺจ นิสฺสรณญฺจ ยถาภูตํ ปชานนฺติ.\csromanspacer{} เต จ}

\vspace{\tipitakaparskip}
\csromancenter{ธาตุสํยุตฺต โข เม ภิกฺขเว สมณา วา พฺราหฺมณา วา สมเณสุ เจว สมณสมฺมตา, พฺราหฺมเณสุ จ พฺราหฺมณสมฺมตา.\csromanspacer{} เต จ ปนายสฺมนฺโต สามญฺญตฺถญฺจ พฺรหฺมญฺญตฺถญฺจ ทิฏฺเฐว ธมฺเม สยํ อภิญฺญา สจฺฉิกตฺวา อุปสมฺปชฺช วิหรนฺตีติ.\csromanspacer{}\csromanspacer{}อฏฺฐมํ.}
\csromanheader{2}{9. ทุติยสมณพฺราหฺมณสุตฺต}
\proseitem{122}{\csromanfolio{385}สาวตฺถิยํ วิหรติ.\csromanspacer{} จตสฺโส อิมา ภิกฺขเว ธาตุโย.\csromanspacer{} กตมา จตสฺโส, ปถวีธาตุ อาโปธาตุ เตโชธาตุ วาโยธาตุ.\csromanspacer{} เย หิ เกจิ ภิกฺขเว สมณา วา พฺราหฺมณา วา อิมาสํ จตุนฺนํ ธาตูนํ สมุทยญฺจ อตฺถงฺคมญฺจ อสฺสาทญฺจ อาทีนวญฺจ นิสฺสรณญฺจ ยถาภูตํ นปฺปชานนฺติ ฯเปฯ. (วิตฺถาเรตพฺพํ.) ปชานนฺติ ฯเปฯ สยํ อภิญฺญา สจฺฉิกตฺวา อุปสมฺปชฺช วิหรนฺตีติ.\csromanspacer{}\csromanspacer{}นวมํ.}

\csromanmatikaanchor{mtk.144}
\csromanmatikaanchor{mtk.145}
\csromantocmarknopage{section}{4. จตุตฺถวคฺค}{mtk.144}
\bookmark[dest={mtk.144},level=1]{4. จตุตฺถวคฺค}
\csromantocmarkref{subsection}{1. จตุธาตุสุตฺต}{mtk.145}
\bookmark[dest={mtk.145},level=2]{1. จตุธาตุสุตฺต}
\csromanheader{2}{10. ตติยสมณพฺราหฺมณสุตฺต}
\proseitem{123}{สาวตฺถิยํ วิหรติ.\csromanspacer{} เย หิ เกจิ ภิกฺขเว สมณา วา พฺราหฺมณา วา ปถวีธาตุํ นปฺปชานนฺติ, ปถวีธาตุสมุทยํ นปฺปชานนฺติ, ปถวีธาตุนิโรธํ นปฺปชานนฺติ, ปถวีธาตุนิโรธคามินิํ ปฏิปทํ นปฺปชานนฺติ ฯเปฯ อาโปธาตุํ นปฺปชานนฺติ.\csromanspacer{} เตโชธาตุํ นปฺปชานนฺติ.\csromanspacer{} วาโยธาตุํ นปฺปชานนฺติ, วาโยธาตุสมุทยํ นปฺปชานนฺติ, วาโยธาตุนิโรธํ นปฺปชานนฺติ, วาโยธาตุนิโรธคามินิํ ปฏิปทํ นปฺปชานนฺติ.\csromanspacer{} น เม เต ภิกฺขเว สมณา วา พฺราหฺมณา วา สมเณสุ วา สมณสมฺมตา, พฺราหฺมเณสุ วา พฺราหฺมณสมฺมตา.\csromanspacer{} น จ ปน เต อายสฺมนฺโต สามญฺญตฺถํ วา พฺรหฺมญฺญตฺถํ วา ทิฏฺเฐว ธมฺเม สยํ อภิญฺญา สจฺฉิกตฺวา อุปสมฺปชฺช วิหรนฺติ.}

\csromanmatikaanchor{mtk.146}
\csromantocmarkref{subsection}{2. ปุพฺเพสมฺโพธสุตฺต}{mtk.146}
\bookmark[dest={mtk.146},level=2]{2. ปุพฺเพสมฺโพธสุตฺต}
\prose{เย จ โข เกจิ ภิกฺขเว สมณา วา พฺราหฺมณา วา ปถวีธาตุํ ปชานนฺติ, ปถวีธาตุสมุทยํ ปชานนฺติ, ปถวีธาตุนิโรธํ ปชานนฺติ, ปถวีธาตุนิโรธคามินิํ ปฏิปทํ ปชานนฺติ.\csromanspacer{} เย จ โข เกจิ ภิกฺขเว สมณา วา พฺราหฺมณา วา ฯเปฯ อาโปธาตุํ ปชานนฺติ.\csromanspacer{} เตโชธาตุํ ปชานนฺติ.\csromanspacer{} วาโยธาตุํ ปชานนฺติ, วาโยธาตุสมุทยํ ปชานนฺติ, วาโยธาตุนิโรธํ ปชานนฺติ, วาโยธาตุนิโรธคามินิํ ปฏิปทํ ปชานนฺติ.}

\vspace{\tipitakaparskip}
\csromancenter{นิทานวคฺคสํยุตฺตปาฬิ เต จ โข เม ภิกฺขเว สมณา วา พฺราหฺมณา วา สมเณสุ เจว สมณสมฺมตา, พฺราหฺมเณสุ จ พฺราหฺมณสมฺมตา.\csromanspacer{} เต จ ปนายสฺมนฺโต สามญฺญตฺถญฺจ พฺรหฺมญฺญตฺถญฺจ ทิฏฺเฐว ธมฺเม สยํ อภิญฺญา สจฺฉิกตฺวา อุปสมฺปชฺช วิหรนฺตีติ.\csromanspacer{}\csromanspacer{}ทสมํ.}
\csromancenter{จตุตฺโถ วคฺโค.}
\titlehead{\textbf{ตสฺสุทฺทานํ}}
\csromanmatikaanchor{mtk.147}
\csromanmatikaanchor{mtk.155}
\csromantocmarkref{subsection}{3. อจริํสุตฺต}{mtk.147}
\bookmark[dest={mtk.147},level=2]{3. อจริํสุตฺต}
\csromangathameasure{จตสฺโส ปุพฺเพ อจริํ, โนเจทํ จ ทุกฺเขน จ. \\ อภินนฺทญฺจ อุปฺปาโท, ตโย สมณพฺราหฺมณาติ. \\ \textbf{ธาตุสํยุตฺตํ สมตฺตํ.}}
\csromangathagroup{\csromangathastanza{\csromanfolio{386}จตสฺโส ปุพฺเพ อจริํ, โนเจทํ จ ทุกฺเขน จ. \\ อภินนฺทญฺจ อุปฺปาโท, ตโย สมณพฺราหฺมณาติ. \\ \textbf{ธาตุสํยุตฺตํ สมตฺตํ.}}}

\chapterheadpageread{4. อนมตคฺคสํยุตฺต}
\chapterhead{1. ปฐมวคฺค}
\csromanheader{2}{1. ติณกฏฺฐสุตฺต}
\proseitem{124}{\csromanfolio{387}เอวํ เม สุตํ\csromandash{}เอกํ สมยํ ภควา สาวตฺถิยํ วิหรติ เชตวเน อนาถปิณฺฑิกสฺส อาราเม.\csromanspacer{} ตตฺร โข ภควา ภิกฺขู อามนฺเตสิ “ภิกฺขโว”ติ. “ภทนฺเต”ติ เต ภิกฺขู ภควโต ปจฺจสฺโสสุํ.\csromanspacer{} ภควา เอตทโวจ\csromandash{}}

\prose{อนมตคฺโคยํ\footnote{อนมตคฺคายํ (อิ, ก)} ภิกฺขเว สํสาโร, ปุพฺพา โกฏิ น ปญฺญายติ อวิชฺชานีวรณานํ สตฺตานํ ตณฺหาสํโยชนานํ สนฺธาวตํ สํสรตํ.\csromanspacer{} เสยฺยถาปิ ภิกฺขเว ปุริโส ยํ อิมสฺมิํ ชมฺพุทีเป ติณกฏฺฐสาขา- ปลาสํ, ตํ เฉตฺวา\footnote{ตจฺเฉตฺวา (พหูสุ)} เอกชฺฌํ สํหริตฺวา จตุรงฺคุลํ จตุรงฺคุลํ ฆฏิกํ กตฺวา นิกฺขิเปยฺย อยํ เม มาตา, ตสฺสา เม มาตุ อยํ มาตา”ติ.\csromanspacer{} อปริยาทินฺนาว\footnote{อปริยาทิณฺณาว (สี)} ภิกฺขเว ตสฺส ปุริสสฺส มาตุมาตโร อสฺสุ, อถ อิมสฺมิํ ชมฺพุทีเป ติณกฏฺฐสาขาปลาสํ ปริกฺขยํ ปริยาทานํ คจฺเฉยฺย.\csromanspacer{} ตํ กิสฺส เหตุ, อนมตคฺโคยํ ภิกฺขเว สํสาโร, ปุพฺพา โกฏิ น ปญฺญายติ อวิชฺชานีวรณานํ สตฺตานํ ตณฺหาสํโยชนานํ สนฺธาวตํ สํสรตํ.\csromanspacer{} เอวํ ทีฆรตฺตํ โว ภิกฺขเว ทุกฺขํ ปจฺจนุภูตํ ติพฺพํ ปจฺจนุภูตํ พฺยสนํ ปจฺจนุภูตํ กฏสี\footnote{กฏสิ (สี, อิ, ก) กฏา ฉวา สยนฺติ เอตฺถาติ กฏสี.} วฑฺฒิตา, ยาวญฺจิทํ ภิกฺขเว อลเมว สพฺพสงฺขาเรสุ นิพฺพินฺทิตุํ, อลํ วิรชฺชิตุํ, อลํ วิมุจฺจิตุนฺติ.\csromanspacer{}\csromanspacer{}ปฐมํ.}

\csromanheader{2}{2. ปถวีสุตฺต}
\csromanmatikaanchor{mtk.148}
\csromantocmarkref{subsection}{4. โนเจทํ สุตฺต}{mtk.148}
\bookmark[dest={mtk.148},level=2]{4. โนเจทํ สุตฺต}
\proseitem{125}{สาวตฺถิยํ วิหรติ.\csromanspacer{} อนมตคฺโคยํ ภิกฺขเว สํสาโร, ปุพฺพา โกฏิ น ปญฺญายติ อวิชฺชานีวรณานํ สตฺตานํ ตณฺหาสํโยชนานํ สนฺธาวตํ สํสรตํ.\csromanspacer{} เสยฺยถาปิ ภิกฺขเว ปุริโส อิมํ มหาปถวิํ โกลฏฺฐิมตฺตํ โกลฏฺฐิมตฺตํ มตฺติกาคุฬิกํ กริตฺวา นิกฺขิเปยฺย “อยํ เม ปิตา, ตสฺส เม ปิตุ อยํ ปิตา”ติ.\csromanspacer{} อปริยาทินฺนาว ภิกฺขเว ตสฺส ปุริสสฺส ปิตุปิตโร อสฺสุ, อถายํ มหาปถวี ปริกฺขยํ ปริยาทานํ}

\vspace{\tipitakaparskip}
\csromancenter{นิทานวคฺคสํยุตฺตปาฬิ คจฺเฉยฺย.\csromanspacer{} ตํ กิสฺส เหตุ, อนมตคฺโคยํ ภิกฺขเว สํสาโร, ปุพฺพา โกฏิ น ปญฺญายติ อวิชฺชานีวรณานํ สตฺตานํ ตณฺหาสํโยชนานํ สนฺธาวตํ สํสรตํ.\csromanspacer{} เอวํ ทีฆรตฺตํ โว ภิกฺขเว ทุกฺขํ ปจฺจนุภูตํ ติพฺพํ ปจฺจนุภูตํ พฺยสนํ ปจฺจนุภูตํ กฏสี วฑฺฒิตา, ยาวญฺจิทํ ภิกฺขเว อลเมว สพฺพสงฺขาเรสุ นิพฺพินฺทิตุํ, อลํ วิรชฺชิตุํ, อลํ วิมุจฺจิตุนฺติ.\csromanspacer{}\csromanspacer{}ทุติยํ.}
\csromanheader{2}{3. อสฺสุสุตฺต}
\proseitem{126}{\csromanfolio{388}สาวตฺถิยํ วิหรติ.\csromanspacer{} อนมตคฺโคยํ ภิกฺขเว สํสาโร, ปุพฺพา โกฏิ น ปญฺญายติ อวิชฺชานีวรณานํ สตฺตานํ ตณฺหาสํโยชนานํ สนฺธาวตํ สํสรตํ.\csromanspacer{} ตํ กิํ มญฺญถ ภิกฺขเว, กตมํ นุ โข พหุตรํ, ยํ วา โว อิมินา ทีเฆน อทฺธุนา สนฺธาวตํ สํสรตํ อมนาปสมฺปโยคา มนาปวิปฺปโยคา กนฺทนฺตานํ โรทนฺตานํ\footnote{รุทนฺตานํ (สี)} อสฺสุ ปสฺสนฺนํ\footnote{ปสฺสนฺทํ (ก-สี), ปสนฺทํ (สฺยา, กํ), ปสนฺนํ (อิ, ก)} ปคฺฆริตํ, ยํ วา จตูสุ มหาสมุทฺเทสุ อุทกนฺติ.\csromanspacer{} ยถา โข มยํ ภนฺเต ภควตา ธมฺมํ เทสิตํ อาชานาม, เอตเทว ภนฺเต พหุตรํ, ยํ โน อิมินา ทีเฆน อทฺธุนา สนฺธาวตํ สํสรตํ อมนาปสมฺปโยคา มนาปวิปฺปโยคา กนฺทนฺตานํ โรทนฺตานํ อสฺสุ ปสฺสนฺนํ ปคฺฆริตํ, น เตฺวว จตูสุ มหาสมุทฺเทสุ อุทกนฺติ.}

\prose{สาธุ สาธุ ภิกฺขเว, สาธุ โข เม ตุเมฺห ภิกฺขเว เอวํ ธมฺมํ เทสิตํ อาชานาถ, เอตเทว ภิกฺขเว พหุตรํ, ยํ โว อิมินา ทีเฆน อทฺธุนา สนฺธาวตํ สํสรตํ อมนาปสมฺปโยคา มนาปวิปฺปโยคา กนฺทนฺตานํ โรทนฺตานํ อสฺสุ ปสฺสนฺนํ ปคฺฆริตํ, น เตฺวว จตูสุ มหาสมุทฺเทสุ อุทกํ.\csromanspacer{} ทีฆรตฺตํ โว ภิกฺขเว มาตุมรณํ ปจฺจนุภูตํ, เตสํ โว มาตุมรณํ ปจฺจนุโภนฺตานํ อมนาปสมฺปโยคา มนาปวิปฺปโยคา กนฺทนฺตานํ โรทนฺตานํ อสฺสุ ปสฺสนฺนํ ปคฺฆริตํ, น เตฺวว จตูสุ มหาสมุทฺเทสุ อุทกํ.\csromanspacer{} ทีฆรตฺตํ โว ภิกฺขเว ปิตุมรณํ ปจฺจนุภูตํ ฯเปฯ ภาตุมรณํ ปจฺจนุภูตํ.\csromanspacer{} ภคินิมรณํ ปจฺจนุภูตํ.\csromanspacer{} ปุตฺตมรณํ ปจฺจนุภูตํ.\csromanspacer{} ธีตุมรณํ ปจฺจนุภูตํ.\csromanspacer{} ญาติพฺยสนํ ปจฺจนุภูตํ.\csromanspacer{} โภคพฺยสนํ ปจฺจนุภูตํ.\csromanspacer{} ทีฆรตฺตํ โว ภิกฺขเว โรคพฺยสนํ ปจฺจนุภูตํ, เตสํ โว โรคพฺยสนํ ปจฺจนุโภนฺตานํ อมนาปสมฺปโยคา มนาปวิปฺปโยคา กนฺทนฺตานํ โรทนฺตานํ อสฺสุ ปสฺสนฺนํ ปคฺฆริตํ, น เตฺวว จตูสุ มหาสมุทฺเทสุ อุทกํ.\csromanspacer{} ตํ กิสฺส เหตุ, อนมตคฺโคยํ ภิกฺขเว}

\vspace{\tipitakaparskip}
\csromancenter{อมตคฺคสํยุตฺต สํสาโร ฯเปฯ.\csromanspacer{} ยาวญฺจิทํ ภิกฺขเว อลเมว สพฺพสงฺขาเรสุ นิพฺพินฺทิตุํ, อลํ วิรชฺชิตุํ, อลํ วิมุจฺจิตุนฺติ.\csromanspacer{}\csromanspacer{}ตติยํ.}
\csromanmatikaanchor{mtk.149}
\csromantocmarkref{subsection}{5. เอกนฺตทุกฺขสุตฺต}{mtk.149}
\bookmark[dest={mtk.149},level=2]{5. เอกนฺตทุกฺขสุตฺต}
\csromanheader{2}{4. ขีรสุตฺต}
\proseitem{127}{\csromanfolio{389}สาวตฺถิยํ วิหรติ.\csromanspacer{} อนมตคฺโคยํ ภิกฺขเว สํสาโร, ปุพฺพา โกฏิ น ปญฺญายติ อวิชฺชานีวรณานํ สตฺตานํ ตณฺหาสํโยชนานํ สนฺธาวตํ สํสรตํ.\csromanspacer{} ตํ กิํ มญฺญถ ภิกฺขเว, กตมํ นุ โข พหุตรํ, ยํ วา โว อิมินา ทีเฆน อทฺธุนา สนฺธาวตํ สํสรตํ มาตุถญฺญํ ปีตํ, ยํ วา จตูสุ มหาสมุทฺเทสุ อุทกนฺติ.\csromanspacer{} ยถา โข มยํ ภนฺเต ภควตา ธมฺมํ เทสิตํ อาชานาม, เอตเทว ภนฺเต พหุตรํ, ยํ โน อิมินา ทีเฆน อทฺธุนา สนฺธาวตํ สํสรตํ มาตุถญฺญํ ปีตํ, น เตฺวว จตูสุ มหาสมุทฺเทสุ อุทกนฺติ.}

\prose{สาธุ สาธุ ภิกฺขเว, สาธุ โข เม ตุเมฺห ภิกฺขเว เอวํ ธมฺมํ เทสิตํ อาชานาถ, เอตเทว ภิกฺขเว พหุตรํ, ยํ โว อิมินา ทีเฆน อทฺธุนา สนฺธาวตํ สํสรตํ มาตุถญฺญํ ปีตํ, น เตฺวว จตูสุ มหาสมุทฺเทสุ อุทกํ.\csromanspacer{} ตํ กิสฺส เหตุ, อนมตคฺโคยํ ภิกฺขเว สํสาโร ฯเปฯ อลํ วิมุจฺจิตุนฺติ.\csromanspacer{}\csromanspacer{}จตุตฺถํ.}

\csromanheader{2}{5. ปพฺพตสุตฺต}
\proseitem{128}{สาวตฺถิยํ วิหรติ ฯเปฯ อาราเม.\csromanspacer{} อถ โข อญฺญตโร ภิกฺขุ เยน ภควา เตนุปสงฺกมิ, อุปสงฺกมิตฺวา ภควนฺตํ อภิวาเทตฺวา เอกมนฺตํ นิสีทิ, เอกมนฺตํ นิสินฺโน โข โส ภิกฺขุ ภควนฺตํ เอตทโวจ “กีวทีโฆ นุ โข ภนฺเต กปฺโป”ติ.\csromanspacer{} ทีโฆ โข ภิกฺขุ กปฺโป, โส น สุกโร สงฺขาตุํ “เอตฺตกานิ วสฺสานิ” อิติ วา “เอตฺตกานิ วสฺสสตานิ” อิติ วา “เอตฺตกานิ วสฺสสหสฺสานิ” อิติ วา “เอตฺตกานิ วสฺสสตสหสฺสานิ” อิติ วาติ.}

\csromanmatikaanchor{mtk.150}
\csromantocmarkref{subsection}{6. อภินนฺทสุตฺต}{mtk.150}
\bookmark[dest={mtk.150},level=2]{6. อภินนฺทสุตฺต}
\prose{สกฺกา ปน ภนฺเต อุปมํ กาตุนฺติ. “สกฺกา ภิกฺขู”ติ ภควา อโวจ, เสยฺยถาปิ ภิกฺขุ มหาเสโล ปพฺพโต โยชนํ อายาเมน โยชนํ วิตฺถาเรน โยชนํ อุพฺเพเธน อจฺฉินฺโน อสุสิโร เอกคฺฆโน, ตเมนํ ปุริโส วสฺสสตสฺส วสฺสสตสฺส อจฺจเยน กาสิเกน}

\vspace{\tipitakaparskip}
\csromancenter{นิทานวคฺคสํยุตฺตปาฬิ วตฺเถน สกิํ สกิํ ปริมชฺเชยฺย.\csromanspacer{} ขิปฺปตรํ โข โส ภิกฺขุ มหาเสโล ปพฺพโต อิมินา อุปกฺกเมน ปริกฺขยํ ปริยาทานํ คจฺเฉยฺย, น เตฺวว กปฺโป.\csromanspacer{} เอวํ ทีโฆ ภิกฺขุ กปฺโป, เอวํ ทีฆานํ โข ภิกฺขุ กปฺปานํ เนโก กปฺโป สํสิโต, เนกํ กปฺปสตํ สํสิตํ, เนกํ กปฺปสหสฺสํ สํสิตํ, เนกํ กปฺปสตสหสฺสํ สํสิตํ.\csromanspacer{} ตํ กิสฺส เหตุ, อนมตคฺโคยํ ภิกฺขุ สํสาโร, ปุพฺพา โกฏิ ฯเปฯ.\csromanspacer{} ยาวญฺจิทํ ภิกฺขุ อลเมว สพฺพสงฺขาเรสุ นิพฺพินฺทิตุํ, อลํ วิรชฺชิตุํ, อลํ วิมุจฺจิตุนฺติ.\csromanspacer{}\csromanspacer{}ปญฺจมํ.}
\csromanheader{2}{6. สาสปสุตฺต}
\proseitem{129}{\csromanfolio{390}สาวตฺถิยํ วิหรติ.\csromanspacer{} อถ โข อญฺญตโร ภิกฺขุ เยน ภควา ฯเปฯ เอกมนฺตํ นิสินฺโน โข โส ภิกฺขุ ภควนฺตํ เอตทโวจ “กีวทีโฆ นุ โข ภนฺเต กปฺโป”ติ.\csromanspacer{} ทีโฆ โข ภิกฺขุ กปฺโป, โส น สุกโร สงฺขาตุํ “เอตฺตกานิ วสฺสานิ” อิติ วา ฯเปฯ “เอตฺตกานิ วสฺสสตสหสฺสานิ” อิติ วาติ.}

\csromanmatikaanchor{mtk.151}
\csromantocmarkref{subsection}{7. อุปฺปาทสุตฺต}{mtk.151}
\bookmark[dest={mtk.151},level=2]{7. อุปฺปาทสุตฺต}
\prose{สกฺกา ปน ภนฺเต อุปมํ กาตุนฺติ. “สกฺกา ภิกฺขู”ติ ภควา อโวจ, เสยฺยถาปิ ภิกฺขุ อายสํ นครํ โยชนํ อายาเมน โยชนํ วิตฺถาเรน โยชนํ อุพฺเพเธน ปุณฺณํ สาสปานํ คุฬิกาพทฺธํ\footnote{จูฬิกาพทฺธํ (สี, อิ)}, ตโต ปุริโส วสฺสสตสฺส วสฺสสตสฺส อจฺจเยน เอกเมกํ สาสปํ อุทฺธเรยฺย.\csromanspacer{} ขิปฺปตรํ โข โส ภิกฺขุ มหาสาสปราสิ อิมินา อุปกฺกเมน ปริกฺขยํ ปริยาทานํ คจฺเฉยฺย, น เตฺวว กปฺโป.\csromanspacer{} เอวํ ทีโฆ โข ภิกฺขุ กปฺโป, เอวํ ทีฆานํ โข ภิกฺขุ กปฺปานํ เนโก กปฺโป สํสิโต, เนกํ กปฺปสตํ สํสิตํ, เนกํ กปฺปสหสฺสํ สํสิตํ, เนกํ กปฺปสตสหสฺสํ สํสิตํ.\csromanspacer{} ตํ กิสฺส เหตุ, อนมตคฺโคยํ ภิกฺขุ สํสาโร ฯเปฯ อลํ วิมุจฺจิตุนฺติ.\csromanspacer{}\csromanspacer{}ฉฏฺฐํ.}

\csromanheader{2}{7. สาวกสุตฺต}
\proseitem{130}{สาวตฺถิยํ วิหรติ.\csromanspacer{} อถ โข สมฺพหุลา ภิกฺขู เยน ภควา ฯเปฯ เอกมนฺตํ นิสินฺนา โข เต ภิกฺขู ภควนฺตํ เอตทโวจุํ “กีวพหุกา นุ โข ภนฺเต กปฺปา อพฺภตีตา อติกฺกนฺตา”ติ.\csromanspacer{} พหุกา โข}

\csromanmatikaanchor{mtk.152}
\csromantocmarkref{subsection}{8. สมณพฺราหฺมณสุตฺต}{mtk.152}
\bookmark[dest={mtk.152},level=2]{8. สมณพฺราหฺมณสุตฺต}
\vspace{\tipitakaparskip}
\csromancenter{อมตคฺคสํยุตฺต ภิกฺขเว กปฺปา อพฺภตีตา อติกฺกนฺตา, เต น สุกรา สงฺขาตุํ “เอตฺตกา กปฺปา” อิติ วา “เอตฺตกานิ กปฺปสตานิ” อิติ วา “เอตฺตกานิ กปฺปสหสฺสานิ” อิติ วา “เอตฺตกานิ กปฺปสตสหสฺสานิ” อิติ วาติ.}
\prose{\csromanfolio{391}สกฺกา ปน ภนฺเต อุปมํ กาตุนฺติ. “สกฺกา ภิกฺขเว”ติ ภควา อโวจ, อิธสฺสุ ภิกฺขเว จตฺตาโร สาวกา วสฺสสตายุกา วสฺสสตชีวิโน, เต ทิวเส ทิวเส กปฺปสตสหสฺสํ กปฺปสตสหสฺสํ อนุสฺสเรยฺยุํ.\csromanspacer{} อนนุสฺสริตาว ภิกฺขเว เตหิ กปฺปา อสฺสุ, อถ โข เต จตฺตาโร สาวกา วสฺสสตายุกา วสฺสสตชีวิโน วสฺสสตสฺส อจฺจเยน กาลํ กเรยฺยุํ.\csromanspacer{} เอวํ พหุกา โข ภิกฺขเว กปฺปา อพฺภตีตา อติกฺกนฺตา, เต น สุกรา สงฺขาตุํ “เอตฺตกา กปฺปา” อิติ วา “เอตฺตกานิ กปฺปสตานิ” อิติ วา “เอตฺตกานิ กปฺปสหสฺสานิ” อิติ วา “เอตฺตกานิ กปฺปสตสหสฺสานิ” อิติ วาติ.\csromanspacer{} ตํ กิสฺส เหตุ, อนมตคฺโคยํ ภิกฺขเว สํสาโร ฯเปฯ อลํ วิมุจฺจิตุนฺติ.\csromanspacer{}\csromanspacer{}สตฺตมํ.}

\csromanheader{2}{8. คงฺคาสุตฺต}
\proseitem{131}{ราชคเห วิหรติ เวฬุวเน.\csromanspacer{} อถ โข อญฺญตโร พฺราหฺมโณ เยน ภควา เตนุปสงฺกมิ, อุปสงฺกมิตฺวา ภควตา สทฺธิํ สมฺโมทิ, สมฺโมทนียํ กถํ สารณียํ วีติสาเรตฺวา เอกมนฺตํ นิสีทิ, เอกมนฺตํ นิสินฺโน โข โส พฺราหฺมโณ ภควนฺตํ เอตทโวจ “กีวพหุกา นุ โข โภ โคตม กปฺปา อพฺภตีตา อติกฺกนฺตา”ติ.\csromanspacer{} พหุกา โข พฺราหฺมณ กปฺปา อพฺภตีตา อติกฺกนฺตา, เต น สุกรา สงฺขาตุํ “เอตฺตกา กปฺปา” อิติ วา “เอตฺตกานิ กปฺปสตานิ” อิติ วา “เอตฺตกานิ กปฺปสหสฺสานิ” อิติ วา “เอตฺตกานิ กปฺปสตสหสฺสานิ” อิติ วาติ.}

\csromanmatikaanchor{mtk.153}
\csromantocmarkref{subsection}{9. ทุติยสมณพฺราหฺมณสุตฺต}{mtk.153}
\bookmark[dest={mtk.153},level=2]{9. ทุติยสมณพฺราหฺมณสุตฺต}
\prose{สกฺกา ปน โภ โคตม อุปมํ กาตุนฺติ. “สกฺกา พฺราหฺมณา”ติ ภควา อโวจ, เสยฺยถาปิ พฺราหฺมณ ยโต จายํ นที ปภวติ, ยตฺถ จ มหาสมุทฺทํ อปฺเปติ.\csromanspacer{} ยา เอตสฺมิํ อนฺตเร วาลิกา, สา น สุกรา สงฺขาตุํ “เอตฺตกา วาลิกา” อิติ วา “เอตฺตกานิ วาลิกสตานิ” อิติ วา “เอตฺตกานิ วาลิกสหสฺสานิ” อิติ วา “เอตฺตกานิ วาลิกสตสหสฺสานิ” อิติ วาติ.\csromanspacer{} ตโต พหุตรา โข พฺราหฺมณ กปฺปา อพฺภตีตา อติกฺกนฺตา, เต น สุกรา สงฺขาตุํ “เอตฺตกา กปฺปา”}

\vspace{\tipitakaparskip}
\csromancenter{นิทานวคฺคสํยุตฺตปาฬิ อิติ วา “เอตฺตกานิ กปฺปสตานิ” อิติ วา “เอตฺตกานิ กปฺปสหสฺสานิ” อิติ วา “เอตฺตกานิ กปฺปสตสหสฺสานิ” อิติ วาติ.\csromanspacer{} ตํ กิสฺส เหตุ, อนมตคฺโคยํ พฺราหฺมณ สํสาโร, ปุพฺพา โกฏิ น ปญฺญายติ อวิชฺชานีวรณานํ สตฺตานํ ตณฺหาสํโยชนานํ สนฺธาวตํ สํสรตํ.\csromanspacer{} เอวํ ทีฆรตฺตํ โข พฺราหฺมณ ทุกฺขํ ปจฺจนุภูตํ ติพฺพํ ปจฺจนุภูตํ พฺยสนํ ปจฺจนุภูตํ กฏสี วฑฺฒิตา, ยาวญฺจิทํ พฺราหฺมณ อลเมว สพฺพสงฺขาเรสุ นิพฺพินฺทิตุํ, อลํ วิรชฺชิตุํ, อลํ วิมุจฺจิตุนฺติ.}
\csromanmatikaanchor{mtk.154}
\csromantocmarkref{subsection}{10. ตติยสมณพฺราหฺมณสุตฺต}{mtk.154}
\bookmark[dest={mtk.154},level=2]{10. ตติยสมณพฺราหฺมณสุตฺต}
\csromantocmarkref{section}{อุทฺทานคาถา}{mtk.155}
\bookmark[dest={mtk.155},level=1]{อุทฺทานคาถา}
\prose{\csromanfolio{392}เอวํ วุตฺเต โส พฺราหฺมโณ ภควนฺตํ เอตทโวจ “อภิกฺกนฺตํ โภ โคตม, อภิกฺกนฺตํ โภ โคตม ฯเปฯ อุปาสกํ มํ ภวํ โคตโม ธาเรตุ อชฺชตคฺเค ปาณุเปตํ สรณํ คตนฺ”ติ.\csromanspacer{}\csromanspacer{}อฏฺฐมํ.}

\csromanheader{2}{9. ทณฺฑสุตฺต}
\proseitem{132}{สาวตฺถิยํ วิหรติ.\csromanspacer{} อนมตคฺโคยํ ภิกฺขเว สํสาโร, ปุพฺพา โกฏิ น ปญฺญายติ อวิชฺชานีวรณานํ สตฺตานํ ตณฺหาสํโยชนานํ สนฺธาวตํ สํสรตํ.\csromanspacer{} เสยฺยถาปิ ภิกฺขเว ทณฺโฑ อุปริเวหาสํ ขิตฺโต สกิมฺปิ มูเลน นิปตติ, สกิมฺปิ มชฺเฌน นิปตติ, สกิมฺปิ อนฺเตน นิปตติ.\csromanspacer{} เอวเมว โข ภิกฺขเว อวิชฺชานีวรณา สตฺตา ตณฺหาสํโยชนา สนฺธาวนฺตา สํสรนฺตา สกิมฺปิ อสฺมา โลกา ปรํ โลกํ คจฺฉนฺติ, สกิมฺปิ ปรสฺมา โลกา อิมํ โลกํ อาคจฺฉนฺติ.\csromanspacer{} ตํ กิสฺส เหตุ, อนมตคฺโคยํ ภิกฺขเว สํสาโร ฯเปฯ อลํ วิมุจฺจิตุนฺติ.\csromanspacer{}\csromanspacer{}นวมํ.}

\csromanheader{2}{10. ปุคฺคลสุตฺต}
\proseitem{133}{เอกํ สมยํ ภควา ราชคเห วิหรติ คิชฺฌกูเฏ ปพฺพเต.\csromanspacer{} ตตฺร โข ภควา ภิกฺขู อามนฺเตสิ “ภิกฺขโว”ติ. “ภทนฺเต”ติ เต ภิกฺขู ภควโต ปจฺจสฺโสสุํ.\csromanspacer{} ภควา เอตทโวจ\csromandash{}}

\prose{อนมตคฺโคยํ ภิกฺขเว สํสาโร ฯเปฯ.\csromanspacer{} เอกปุคฺคลสฺส ภิกฺขเว กปฺปํ สนฺธาวโต สํสรโต สิยา เอวํ มหา อฏฺฐิกงฺกโล อฏฺฐิปุญฺโช อฏฺฐิราสิ, ยถายํ เวปุลฺโล ปพฺพโต.\csromanspacer{} สเจ สํหารโก อสฺส, สมฺภตญฺจ น วินสฺเสยฺย.\csromanspacer{} ตํ กิสฺส เหตุ, อนมตคฺโคยํ ภิกฺขเว สํสาโร ฯเปฯ อลํ วิมุจฺจิตุนฺติ.}

\csromanheader{2}{4. อมตคฺคสํยุตฺต}
\csromanmatikaanchor{mtk.156}
\csromanmatikaanchor{mtk.168}
\csromantocmarknopage{chapter}{4. อนมตคฺคสํยุตฺต}{mtk.156}
\bookmark[dest={mtk.156},level=0]{4. อนมตคฺคสํยุตฺต}
\csromantocmarknopage{section}{1. ปฐมวคฺค}{mtk.157}
\bookmark[dest={mtk.157},level=1]{1. ปฐมวคฺค}
\prose{\csromanfolio{393}อิทมโวจ ภควา, อิทํ วตฺวาน สุคโต อถาปรํ เอตทโวจ สตฺถา\csromandash{}}

\csromanmatikaanchor{mtk.158}
\csromantocmarkref{subsection}{1. ติณกฏฺฐสุตฺต}{mtk.158}
\bookmark[dest={mtk.158},level=2]{1. ติณกฏฺฐสุตฺต}
\csromangathameasure{“เอกสฺเสเกน กปฺเปน, ปุคฺคลสฺสฏฺฐิสญฺจโย. \\ สิยา ปพฺพตสโม ราสิ, อิติ วุตฺตํ มเหสินา. \\ โส โข ปนายํ อกฺขาโต, เวปุลฺโล ปพฺพโต มหา. \\ อุตฺตโร คิชฺฌกูฏสฺส, มคธานํ คิริพฺพเช. \\ ยโต จ อริยสจฺจานิ, สมฺมปฺปญฺญาย ปสฺสติ. \\ ทุกฺขํ ทุกฺขสมุปฺปาทํ, ทุกฺขสฺส จ อติกฺกมํ. \\ อริยํ จฏฺฐงฺคิกํ มคฺคํ, ทุกฺขูปสมคามินํ. \\ ส สตฺตกฺขตฺตุํปรมํ, สนฺธาวิตฺวาน ปุคฺคโล. \\ ทุกฺขสฺสนฺตกโร โหติ, สพฺพสํโยชนกฺขยา”ติ.}
\csromangathagroup{\csromangathastanza{“เอกสฺเสเกน กปฺเปน, ปุคฺคลสฺสฏฺฐิสญฺจโย. \\ สิยา ปพฺพตสโม ราสิ, อิติ วุตฺตํ มเหสินา.}\csromangathastanza{โส โข ปนายํ อกฺขาโต, เวปุลฺโล ปพฺพโต มหา. \\ อุตฺตโร คิชฺฌกูฏสฺส, มคธานํ คิริพฺพเช.}\csromangathastanza{ยโต จ อริยสจฺจานิ, สมฺมปฺปญฺญาย ปสฺสติ. \\ ทุกฺขํ ทุกฺขสมุปฺปาทํ, ทุกฺขสฺส จ อติกฺกมํ.}\csromangathastanza{อริยํ จฏฺฐงฺคิกํ มคฺคํ, ทุกฺขูปสมคามินํ. \\ ส สตฺตกฺขตฺตุํปรมํ, สนฺธาวิตฺวาน ปุคฺคโล. \\ ทุกฺขสฺสนฺตกโร โหติ, สพฺพสํโยชนกฺขยา”ติ.}}

\vspace{\tipitakaparskip}
\csromancenter{ทสมํ.}
\csromanmatikaanchor{mtk.159}
\csromantocmarkref{subsection}{2. ปถวีสุตฺต}{mtk.159}
\bookmark[dest={mtk.159},level=2]{2. ปถวีสุตฺต}
\csromancenter{ปฐโม วคฺโค.}
\titlehead{\textbf{ตสฺสุทฺทานํ}}
\csromangathameasure{ติณกฏฺฐญฺจ ปถวี, อสฺสุ ขีรญฺจ ปพฺพตํ. \\ สาสปา สาวกา คงฺคา, ทณฺโฑ จ ปุคฺคเลน จาติ.}
\csromangathagroup{\csromangathastanza{ติณกฏฺฐญฺจ ปถวี, อสฺสุ ขีรญฺจ ปพฺพตํ. \\ สาสปา สาวกา คงฺคา, ทณฺโฑ จ ปุคฺคเลน จาติ.}}

\csromanmatikaanchor{mtk.160}
\csromantocmarkref{subsection}{3. อสฺสุสุตฺต}{mtk.160}
\bookmark[dest={mtk.160},level=2]{3. อสฺสุสุตฺต}
\chapterhead{2. ทุติยวคฺค}
\csromanheader{2}{1. ทุคฺคตสุตฺต}
\proseitem{134}{เอกํ สมยํ ภควา สาวตฺถิยํ วิหรติ.\csromanspacer{} ตตฺร โข ฯเปฯ.\csromanspacer{} อนมตคฺโคยํ ภิกฺขเว สํสาโร ปุพฺพา โกฏิ น ปญฺญายติ อวิชฺชานีวรณานํ สตฺตานํ ตณฺหาสํโยชนานํ สนฺธาวตํ สํสรตํ.\csromanspacer{} ยํ ภิกฺขเว ปสฺเสยฺยาถ ทุคฺคตํ ทุรูเปตํ, นิฏฺฐเมตฺถ คนฺตพฺพํ “อเมฺหหิปิ เอวรูปํ ปจฺจนุภูตํ อิมินา ทีเฆน อทฺธุนา”ติ.\csromanspacer{} ตํ กิสฺส เหตุ ฯเปฯ.\csromanspacer{} ยาวญฺจิทํ ภิกฺขเว อลเมว สพฺพสงฺขาเรสุ นิพฺพินฺทิตุํ, อลํ วิรชฺชิตุํ, อลํ วิมุจฺจิตุนฺติ.\csromanspacer{}\csromanspacer{}ปฐมํ.}

\csromanheader{2}{นิทานวคฺคสํยุตฺตปาฬิ \textbf{2. สุขิตสุตฺต}}
\csromanmatikaanchor{mtk.161}
\csromantocmarkref{subsection}{4. ขีรสุตฺต}{mtk.161}
\bookmark[dest={mtk.161},level=2]{4. ขีรสุตฺต}
\proseitem{135}{\csromanfolio{394}สาวตฺถิยํ วิหรติ.\csromanspacer{} อนมตคฺโคยํ ภิกฺขเว สํสาโร ฯเปฯ.\csromanspacer{} ยํ ภิกฺขเว ปสฺเสยฺยาถ สุขิตํ สุสชฺชิตํ, นิฏฺฐเมตฺถ คนฺตพฺพํ “อเมฺหหิปิ เอวรูปํ ปจฺจนุภูตํ อิมินา ทีเฆน อทฺธุนา”ติ.\csromanspacer{} ตํ กิสฺส เหตุ, อนมตคฺโคยํ ภิกฺขเว สํสาโร, ปุพฺพา โกฏิ น ปญฺญายติ ฯเปฯ อลํ วิมุจฺจิตุนฺติ.\csromanspacer{}\csromanspacer{}ทุติยํ.}

\csromanheader{2}{3. ติํสมตฺตสุตฺต}
\proseitem{136}{ราชคเห วิหรติ เวฬุวเน.\csromanspacer{} อถ โข ติํสมตฺตา ปาเวยฺยกา\footnote{ปาเฐยฺยกา (กตฺถจิ) วิ 3 กถินกฺขนฺธเกปิ.} ภิกฺขู สพฺเพ อารญฺญิกา สพฺเพ ปิณฺฑปาติกา สพฺเพ ปํสุกูลิกา สพฺเพ เตจีวริกา สพฺเพ สสํโยชนา เยน ภควา เตนุปสงฺกมิํสุ, อุปสงฺกมิตฺวา ภควนฺตํ อภิวาเทตฺวา เอกมนฺตํ นิสีทิํสุ.\csromanspacer{} อถ โข ภควโต เอตทโหสิ “อิเม โข ติํสมตฺตา ปาเวยฺยกา ภิกฺขู สพฺเพ อารญฺญิกา สพฺเพ ปิณฺฑปาติกา สพฺเพ ปํสุกูลิกา สพฺเพ เตจีวริกา สพฺเพ สสํโยชนา.\csromanspacer{} ยํนูนาหํ อิเมสํ ตถา ธมฺมํ เทเสยฺยํ, ยถา เนสํ อิมสฺมิํเยว อาสเน อนุปาทาย อาสเวหิ จิตฺตานิ วิมุจฺเจยฺยุนฺ”ติ.\csromanspacer{} อถ โข ภควา ภิกฺขู อามนฺเตสิ “ภิกฺขโว”ติ. “ภทนฺเต”ติ เต ภิกฺขู ภควโต ปจฺจสฺโสสุํ.\csromanspacer{} ภควา เอตทโวจ\csromandash{}}

\csromanmatikaanchor{mtk.162}
\csromantocmarkref{subsection}{5. ปพฺพตสุตฺต}{mtk.162}
\bookmark[dest={mtk.162},level=2]{5. ปพฺพตสุตฺต}
\prose{อนมตคฺโคยํ ภิกฺขเว สํสาโร, ปุพฺพา โกฏิ น ปญฺญายติ อวิชฺชานีวรณานํ สตฺตานํ ตณฺหาสํโยชนานํ สนฺธาวตํ สํสรตํ.\csromanspacer{} ตํ กิํ มญฺญถ ภิกฺขเว, กตมํ นุ โข พหุตรํ, ยํ วา โว อิมินา ทีเฆน อทฺธุนา สนฺธาวตํ สํสรตํ สีสจฺฉินฺนานํ โลหิตํ ปสฺสนฺนํ ปคฺฆริตํ, ยํ วา จตูสุ มหาสมุทฺเทสุ อุทกนฺติ.\csromanspacer{} ยถา โข มยํ ภนฺเต ภควตา ธมฺมํ เทสิตํ อาชานาม “เอตเทว ภนฺเต พหุตรํ, ยํ โน อิมินา ทีเฆน อทฺธุนา สนฺธาวตํ สํสรตํ สีสจฺฉินฺนานํ โลหิตํ ปสฺสนฺนํ ปคฺฆริตํ, น เตฺวว จตูสุ มหาสมุทฺเทสุ อุทกนฺ”ติ.}

\prose{สาธุ สาธุ ภิกฺขเว, สาธุ โข เม ตุเมฺห ภิกฺขเว เอวํ ธมฺมํ เทสิตํ อาชานาถ, เอตเทว ภิกฺขเว พหุตรํ, ยํ โว อิมินา ทีเฆน อทฺธุนา สนฺธาวตํ สํสรตํ สีสจฺฉินฺนานํ โลหิตํ ปสฺสนฺนํ ปคฺฆริตํ, น เตฺวว จตูสุ}

\vspace{\tipitakaparskip}
\csromancenter{อมตคฺคสํยุตฺต มหาสมุทฺเทสุ อุทกํ.\csromanspacer{} ทีฆรตฺตํ โว ภิกฺขเว คุนฺนํ สตํ โคภูตานํ สีสจฺฉินฺนานํ โลหิตํ ปสฺสนฺนํ ปคฺฆริตํ, น เตฺวว จตูสุ มหาสมุทฺเทสุ อุทกํ.\csromanspacer{} ทีฆรตฺตํ โว ภิกฺขเว มหิํสานํ\footnote{มหิสานํ (สี, อิ)} สตํ มหิํสภูตานํ สีสจฺฉินฺนานํ โลหิตํ ปสฺสนฺนํ ปคฺฆริตํ ฯเปฯ.\csromanspacer{} ทีฆรตฺตํ โว ภิกฺขเว อุรพฺภานํ สตํ อุรพฺภภูตานํ ฯเปฯ.\csromanspacer{} อชานํ สตํ อชภูตานํ.\csromanspacer{} มิคานํ สตํ มิคภูตานํ.\csromanspacer{} กุกฺกุฏานํ สตํ กุกฺกุฏภูตานํ.\csromanspacer{} สูกรานํ สตํ สูกรภูตานํ.\csromanspacer{} ทีฆรตฺตํ โว ภิกฺขเว “โจรา คามฆาตา”ติ คเหตฺวา สีสจฺฉินฺนานํ โลหิตํ ปสฺสนฺนํ ปคฺฆริตํ.\csromanspacer{} ทีฆรตฺตํ โว ภิกฺขเว “โจรา ปาริปนฺถิกา”ติ คเหตฺวา สีสจฺฉินฺนานํ โลหิตํ ปสฺสนฺนํ ปคฺฆริตํ.\csromanspacer{} ทีฆรตฺตํ โว ภิกฺขเว “โจรา ปารทาริกา”ติ คเหตฺวา สีสจฺฉินฺนานํ โลหิตํ ปสฺสนฺนํ ปคฺฆริตํ, น เตฺวว จตูสุ มหาสมุทฺเทสุ อุทกํ.\csromanspacer{} ตํ กิสฺส เหตุ, อนมตคฺโคยํ ภิกฺขเว สํสาโร ฯเปฯ อลํ วิมุจฺจิตุนฺติ.}
\prose{\csromanfolio{395}อิทมโวจ ภควา, อตฺตมนา เต ภิกฺขู ภควโต ภาสิตํ อภินนฺทุนฺติ.\csromanspacer{} อิมสฺมิํ จ ปน เวยฺยากรณสฺมิํ ภญฺญมาเน ติํสมตฺตานํ ปาเวยฺยกานํ ภิกฺขูนํ อนุปาทาย อาสเวหิ จิตฺตานิ วิมุจฺจิํสูติ.\csromanspacer{}\csromanspacer{}ตติยํ.}

\csromanmatikaanchor{mtk.163}
\csromantocmarkref{subsection}{6. สาสปสุตฺต}{mtk.163}
\bookmark[dest={mtk.163},level=2]{6. สาสปสุตฺต}
\csromanheader{2}{4. มาตุสุตฺต}
\proseitem{137}{สาวตฺถิยํ วิหรติ.\csromanspacer{} อนมตคฺโคยํ ภิกฺขเว สํสาโร ฯเปฯ.\csromanspacer{} น โส ภิกฺขเว สตฺโต สุลภรูโป, โย นมาตาภูตปุพฺโพ อิมินา ทีเฆน อทฺธุนา.\csromanspacer{} ตํ กิสฺส เหตุ, อนมตคฺโคยํ ภิกฺขเว สํสาโร ฯเปฯ อลํ วิมุจฺจิตุนฺติ.\csromanspacer{}\csromanspacer{}จตุตฺถํ.}

\csromanheader{2}{5. ปิตุสุตฺต}
\csromanmatikaanchor{mtk.164}
\csromantocmarkref{subsection}{7. สาวกสุตฺต}{mtk.164}
\bookmark[dest={mtk.164},level=2]{7. สาวกสุตฺต}
\proseitem{138}{สาวตฺถิยํ วิหรติ.\csromanspacer{} อนมตคฺโคยํ ภิกฺขเว สํสาโร ฯเปฯ.\csromanspacer{} น โส ภิกฺขเว สตฺโต สุลภรูโป, โย นปิตาภูตปุพฺโพ ฯเปฯ อลํ วิมุจฺจิตุนฺติ.\csromanspacer{}\csromanspacer{}ปญฺจมํ.}

\csromanheader{2}{6. ภาตุสุตฺต}
\proseitem{139}{สาวตฺถิยํ วิหรติ ฯเปฯ.\csromanspacer{} น โส ภิกฺขเว สตฺโต สุลภรูโป, โย นภาตาภูตปุพฺโพ ฯเปฯ อลํ วิมุจฺจิตุนฺติ.\csromanspacer{}\csromanspacer{}ฉฏฺฐํ.}

\csromanheader{2}{นิทานวคฺคสํยุตฺตปาฬิ \textbf{7. ภคินิสุตฺต}}
\csromanmatikaanchor{mtk.165}
\csromantocmarkref{subsection}{8. คงฺคาสุตฺต}{mtk.165}
\bookmark[dest={mtk.165},level=2]{8. คงฺคาสุตฺต}
\proseitem{140}{\csromanfolio{396}สาวตฺถิยํ วิหรติ ฯเปฯ.\csromanspacer{} น โส ภิกฺขเว สตฺโต สุลภรูโป, โย นภคินิภูตปุพฺโพ ฯเปฯ อลํ วิมุจฺจิตุนฺติ.\csromanspacer{}\csromanspacer{}สตฺตมํ.}

\csromanheader{2}{8. ปุตฺตสุตฺต}
\proseitem{141}{สาวตฺถิยํ วิหรติ ฯเปฯ.\csromanspacer{} น โส ภิกฺขเว สตฺโต สุลภรูโป, โย นปุตฺตภูตปุพฺโพ ฯเปฯ อลํ วิมุจฺจิตุนฺติ.\csromanspacer{}\csromanspacer{}อฏฺฐมํ.}

\csromanheader{2}{9. ธีตุสุตฺต}
\proseitem{142}{สาวตฺถิยํ วิหรติ.\csromanspacer{} อนมตคฺโคยํ ภิกฺขเว สํสาโร, ปุพฺพา โกฏิ น ปญฺญายติ อวิชฺชานีวรณานํ สตฺตานํ ตณฺหาสํโยชนานํ สนฺธาวตํ สํสรตํ.\csromanspacer{} น โส ภิกฺขเว สตฺโต สุลภรูโป, โย นธีตาภูตปุพฺโพ อิมินา ทีเฆน อทฺธุนา.\csromanspacer{} ตํ กิสฺส เหตุ, อนมตคฺโคยํ ภิกฺขเว สํสาโร, ปุพฺพา โกฏิ น ปญฺญายติ อวิชฺชานีวรณานํ สตฺตานํ ตณฺหาสํโยชนานํ สนฺธาวตํ สํสรตํ.\csromanspacer{} เอวํ ทีฆรตฺตํ โว ภิกฺขเว ทุกฺขํ ปจฺจนุภูตํ ติพฺพํ ปจฺจนุภูตํ พฺยสนํ ปจฺจนุภูตํ กฏสี วฑฺฒิตา, ยาวญฺจิทํ ภิกฺขเว อลเมว สพฺพสงฺขาเรสุ นิพฺพินฺทิตุํ, อลํ วิรชฺชิตุํ, อลํ วิมุจฺจิตุนฺติ.\csromanspacer{}\csromanspacer{}นวมํ.}

\csromanmatikaanchor{mtk.166}
\csromantocmarkref{subsection}{9. ทณฺฑสุตฺต}{mtk.166}
\bookmark[dest={mtk.166},level=2]{9. ทณฺฑสุตฺต}
\csromanheader{2}{10. เวปุลฺลปพฺพตสุตฺต}
\proseitem{143}{เอกํ สมยํ ภควา ราชคเห วิหรติ คิชฺฌกูเฏ ปพฺพเต.\csromanspacer{} ตตฺร โข ภควา ภิกฺขู อามนฺเตสิ “ภิกฺขโว”ติ. “ภทนฺเต”ติ เต ภิกฺขู ภควโต ปจฺจสฺโสสุํ.\csromanspacer{} ภควา เอตทโวจ\csromandash{}}

\csromanmatikaanchor{mtk.167}
\csromantocmarkref{subsection}{10. ปุคฺคลสุตฺต}{mtk.167}
\bookmark[dest={mtk.167},level=2]{10. ปุคฺคลสุตฺต}
\csromantocmarkref{section}{อุทฺทานคาถา}{mtk.168}
\bookmark[dest={mtk.168},level=1]{อุทฺทานคาถา}
\prose{อนมตคฺโคยํ ภิกฺขเว สํสาโร, ปุพฺพา โกฏิ น ปญฺญายติ อวิชฺชานีวรณานํ สตฺตานํ ตณฺหาสํโยชนานํ สนฺธาวตํ สํสรตํ.\csromanspacer{} ภูตปุพฺพํ ภิกฺขเว อิมสฺส เวปุลฺลสฺส ปพฺพตสฺส “ปาจีนวํโส”เตฺวว สมญฺญา อุทปาทิ.\csromanspacer{} เตน โข ปน ภิกฺขเว สมเยน มนุสฺสานํ “ติวรา” เตฺวว สมญฺญา อุทปาทิ, ติวรานํ ภิกฺขเว มนุสฺสานํ จตฺตารีส วสฺสสหสฺสานิ อายุปฺปมาณํ อโหสิ, ติวรา ภิกฺขเว มนุสฺสา ปาจีนวํสํ ปพฺพตํ จตูเหน อาโรหนฺติ, จตูเหน โอโรหนฺติ.\csromanspacer{} เตน โข ปน}

\vspace{\tipitakaparskip}
\csromancenter{อมตคฺคสํยุตฺต ภิกฺขเว สมเยน กกุสนฺโธ ภควา อรหํ สมฺมาสมฺพุทฺโธ โลเก อุปฺปนฺโน โหติ, กกุสนฺธสฺส ภิกฺขเว ภควโต อรหโต สมฺมาสมฺพุทฺธสฺส วิธุรสญฺชีวํ นาม สาวกยุคํ อโหสิ อคฺคํ ภทฺทยุคํ.\csromanspacer{} ปสฺสถ ภิกฺขเว สา เจวิมสฺส ปพฺพตสฺส สมญฺญา อนฺตรหิตา, เต จ มนุสฺสา กาลงฺกตา, โส จ ภควา ปรินิพฺพุโต.\csromanspacer{} เอวํ อนิจฺจา ภิกฺขเว สงฺขารา, เอวํ อทฺธุวา ภิกฺขเว สงฺขารา, เอวํ อนสฺสาสิกา ภิกฺขเว สงฺขารา, ยาวญฺจิทํ ภิกฺขเว อลเมว สพฺพสงฺขาเรสุ นิพฺพินฺทิตุํ, อลํ วิรชฺชิตุํ, อลํ วิมุจฺจิตุํ.}
\prose{\csromanfolio{397}ภูตปุพฺพํ ภิกฺขเว อิมสฺส เวปุลฺลสฺส ปพฺพตสฺส “วงฺกโก”เตฺวว สมญฺญา อุทปาทิ.\csromanspacer{} เตน โข ปน ภิกฺขเว สมเยน มนุสฺสานํ “โรหิตสฺสา”เตฺวว สมญฺญา อุทปาทิ, โรหิตสฺสานํ ภิกฺขเว มนุสฺสานํ ติํส วสฺสสหสฺสานิ อายุปฺปมาณํ อโหสิ, โรหิตสฺสา ภิกฺขเว มนุสฺสา วงฺกกํ ปพฺพตํ ตีเหน อาโรหนฺติ, ตีเหน โอโรหนฺติ.\csromanspacer{} เตน โข ปน ภิกฺขเว สมเยน โกณาคมโน ภควา อรหํ สมฺมาสมฺพุทฺโธ โลเก อุปฺปนฺโน โหติ, โกณาคมนสฺส ภิกฺขเว ภควโต อรหโต สมฺมาสมฺพุทฺธสฺส ภิยฺโยสุตฺตรํ นาม สาวกยุคํ อโหสิ อคฺคํ ภทฺทยุคํ.\csromanspacer{} ปสฺสถ ภิกฺขเว สา เจวิมสฺส ปพฺพตสฺส สมญฺญา อนฺตรหิตา, เต จ มนุสฺสา กาลงฺกตา, โส จ ภควา ปรินิพฺพุโต.\csromanspacer{} เอวํ อนิจฺจา ภิกฺขเว สงฺขารา ฯเปฯ อลํ วิมุจฺจิตุํ.}

\prose{ภูตปุพฺพํ ภิกฺขเว อิมสฺส เวปุลฺลสฺส ปพฺพตสฺส “สุปสฺโส”เตฺวว\footnote{สุผสฺโสเตฺวว (สี)} สมญฺญา อุทปาทิ.\csromanspacer{} เตน โข ปน ภิกฺขเว สมเยน มนุสฺสานํ “สุปฺปิยา”เตฺวว\footnote{อปฺปิยาเตฺวว (สี)} สมญฺญา อุทปาทิ, สุปฺปิยานํ ภิกฺขเว มนุสฺสานํ วีสติ วสฺสสหสฺสานิ อายุปฺปมาณํ อโหสิ, สุปฺปิยา ภิกฺขเว มนุสฺสา สุปสฺสํ ปพฺพตํ ทฺวีเหน อาโรหนฺติ, ทฺวีเหน โอโรหนฺติ.\csromanspacer{} เตน โข ปน ภิกฺขเว สมเยน กสฺสโป ภควา อรหํ สมฺมาสมฺพุทฺโธ โลเก อุปฺปนฺโน โหติ, กสฺสปสฺส ภิกฺขเว ภควโต อรหโต สมฺมาสมฺพุทฺธสฺส ติสฺสภารทฺวาชํ นาม สาวกยุคํ อโหสิ อคฺคํ ภทฺทยุคํ.\csromanspacer{} ปสฺสถ ภิกฺขเว สา เจวิมสฺส ปพฺพตสฺส สมญฺญา อนฺตรหิตา, เต จ มนุสฺสา กาลงฺกตา, โส จ ภควา ปรินิพฺพุโต.\csromanspacer{} เอวํ อนิจฺจา ภิกฺขเว สงฺขารา, เอวํ อทฺธุวา ภิกฺขเว สงฺขารา ฯเปฯ อลํ วิมุจฺจิตุํ.}

\vspace{\tipitakaparskip}
\csromancenter{นิทานวคฺคสํยุตฺตปาฬิ เอตรหิ โข ปน ภิกฺขเว อิมสฺส เวปุลฺลสฺส ปพฺพตสฺส “เวปุลฺโล”เตฺวว สมญฺญา อุทปาทิ.\csromanspacer{} เอตรหิ โข ปน ภิกฺขเว อิเมสํ มนุสฺสานํ “มาคธกา”เตฺวว สมญฺญา อุทปาทิ, มาคธกานํ ภิกฺขเว มนุสฺสานํ อปฺปกํ อายุปฺปมาณํ ปริตฺตํ ลหุกํ\footnote{ลหุสํ (สี)}.\csromanspacer{} โย จิรํ ชีวติ, โส วสฺสสตํ อปฺปํ วา ภิยฺโย.\csromanspacer{} มาคธกา ภิกฺขเว มนุสฺสา เวปุลฺลํ ปพฺพตํ มุหุตฺเตน อาโรหนฺติ, มุหุตฺเตน โอโรหนฺติ.\csromanspacer{} เอตรหิ โข ปนาหํ ภิกฺขเว อรหํ สมฺมาสมฺพุทฺโธ โลเก อุปฺปนฺโน, มยฺหํ โข ปน ภิกฺขเว สาริปุตฺตโมคฺคลฺลานํ นาม สาวกยุคํ อคฺคํ ภทฺทยุคํ.\csromanspacer{} ภวิสฺสติ ภิกฺขเว โส สมโย, ยา อยญฺเจวิมสฺส ปพฺพตสฺส สมญฺญา อนฺตรธายิสฺสติ, อิเม จ มนุสฺสา กาลํ กริสฺสนฺติ, อหญฺจ ปรินิพฺพายิสฺสามิ.\csromanspacer{} เอวํ อนิจฺจา ภิกฺขเว สงฺขารา, เอวํ อทฺธุวา ภิกฺขเว สงฺขารา, เอวํ อนสฺสาสิกา ภิกฺขเว สงฺขารา, ยาวญฺจิทํ ภิกฺขเว อลเมว สพฺพสงฺขาเรสุ นิพฺพินฺทิตุํ, อลํ วิรชฺชิตุํ, อลํ วิมุจฺจิตุนฺติ.}
\csromanmatikaanchor{mtk.169}
\csromanmatikaanchor{mtk.180}
\csromantocmarknopage{section}{2. ทุติยวคฺค}{mtk.169}
\bookmark[dest={mtk.169},level=1]{2. ทุติยวคฺค}
\prose{\csromanfolio{398}อิทมโวจ ภควา, อิทํ วตฺวาน สุคโต อถาปรํ เอตทโวจ สตฺถา\csromandash{}}

\csromangathameasure{“ปาจีนวํโส ติวรานํ, โรหิตสฺสาน วงฺกโก. \\ สุปฺปิยานํ สุปสฺโสติ, มาคธานญฺจ เวปุลฺโล. \\ อนิจฺจา วต สงฺขารา, อุปฺปาทวยธมฺมิโน. \\ อุปฺปชฺชิตฺวา นิรุชฺฌนฺติ, เตสํ วูปสโม สุโข”ติ.}
\csromangathagroup{\csromangathastanza{“ปาจีนวํโส ติวรานํ, โรหิตสฺสาน วงฺกโก. \\ สุปฺปิยานํ สุปสฺโสติ, มาคธานญฺจ เวปุลฺโล.}\csromangathastanza{อนิจฺจา วต สงฺขารา, อุปฺปาทวยธมฺมิโน. \\ อุปฺปชฺชิตฺวา นิรุชฺฌนฺติ, เตสํ วูปสโม สุโข”ติ.}}

\vspace{\tipitakaparskip}
\csromancenter{ทสมํ.}
\csromancenter{ทุติโย วคฺโค.}
\titlehead{\textbf{ตสฺสุทฺทานํ}}
\csromangathameasure{ทุคฺคตํ สุขิตญฺเจว, ติํส มาตาปิเตน จ. \\ ภาตา ภคินี ปุตฺโต จ, ธีตา เวปุลฺลปพฺพตํ. \\ \textbf{อนมตคฺคสํยุตฺตํ สมตฺตํ.}}
\csromangathagroup{\csromangathastanza{ทุคฺคตํ สุขิตญฺเจว, ติํส มาตาปิเตน จ. \\ ภาตา ภคินี ปุตฺโต จ, ธีตา เวปุลฺลปพฺพตํ. \\ \textbf{อนมตคฺคสํยุตฺตํ สมตฺตํ.}}}

\csromanheader{2}{5. กสฺสปสํยุตฺต}
\csromanmatikaanchor{mtk.170}
\csromantocmarkref{subsection}{1. ทุคฺคตสุตฺต}{mtk.170}
\bookmark[dest={mtk.170},level=2]{1. ทุคฺคตสุตฺต}
\csromanheader{3}{1. สนฺตุฏฺฐสุตฺต}
\proseitem{144}{\csromanfolio{399}สาวตฺถิยํ วิหรติ.\csromanspacer{} สนฺตุฏฺฐายํ\footnote{สนฺตุฏฺโฐยํ (สี)} ภิกฺขเว กสฺสโป อิตรีตเรน จีวเรน, อิตรีตรจีวรสนฺตุฏฺฐิยา จ วณฺณวาที, น จ จีวรเหตุ อเนสนํ อปฺปติรูปํ อาปชฺชติ, อลทฺธา จ จีวรํ น ปริตสฺสติ, ลทฺธา จ จีวรํ อคธิโต\footnote{อคถิโต (สี)} อมุจฺฉิโต อนชฺฌาปนฺโน อาทีนวทสฺสาวี นิสฺสรณปญฺโญ ปริภุญฺชติ.}

\csromanmatikaanchor{mtk.171}
\csromantocmarkref{subsection}{2. สุขิตสุตฺต}{mtk.171}
\bookmark[dest={mtk.171},level=2]{2. สุขิตสุตฺต}
\prose{สนฺตุฏฺฐายํ ภิกฺขเว กสฺสโป อิตรีตเรน ปิณฺฑปาเตน, อิตรีตรปิณฺฑปาตสนฺตุฏฺฐิยา จ วณฺณวาที, น จ ปิณฺฑปาตเหตุ อเนสนํ อปฺปติรูปํ อาปชฺชติ, อลทฺธา จ ปิณฺฑปาตํ น ปริตสฺสติ, ลทฺธา จ ปิณฺฑปาตํ อคธิโต อมุจฺฉิโต อนชฺฌาปนฺโน อาทีนวทสฺสาวี นิสฺสรณปญฺโญ ปริภุญฺชติ.}

\prose{สนฺตุฏฺฐายํ ภิกฺขเว กสฺสโป อิตรีตเรน เสนาสเนน, อิตรีตรเสนาสนสนฺตุฏฺฐิยา จ วณฺณวาที, น จ เสนาสนเหตุ อเนสนํ อปฺปติรูปํ อาปชฺชติ, อลทฺธา จ เสนาสนํ น ปริตสฺสติ, ลทฺธา จ เสนาสนํ อคธิโต อมุจฺฉิโต อนชฺฌาปนฺโน อาทีนวทสฺสาวี นิสฺสรณปญฺโญ ปริภุญฺชติ.}

\csromanmatikaanchor{mtk.172}
\csromantocmarkref{subsection}{3. ติํสมตฺตสุตฺต}{mtk.172}
\bookmark[dest={mtk.172},level=2]{3. ติํสมตฺตสุตฺต}
\prose{สนฺตุฏฺฐายํ ภิกฺขเว กสฺสโป อิตรีตเรน คิลานปฺปจฺจยเภสชฺช- ปริกฺขาเรน, อิตรีตรคิลานปฺปจฺจยเภสชฺชปริกฺขารสนฺตุฏฺฐิยา จ วณฺณวาที, น จ คิลานปฺปจฺจยเภสชฺชปริกฺขารเหตุ อเนสนํ อปฺปติรูปํ อาปชฺชติ, อลทฺธา จ คิลานปฺปจฺจยเภสชฺชปริกฺขารํ น ปริตสฺสติ, ลทฺธา จ คิลานปฺปจฺจยเภสชฺชปริกฺขารํ อคธิโต อมุจฺฉิโต อนชฺฌาปนฺโน อาทีนวทสฺสาวี นิสฺสรณปญฺโญ ปริภุญฺชติ.}

\prose{ตสฺมาติห ภิกฺขเว เอวํ สิกฺขิตพฺพํ\csromandash{}สนฺตุฏฺฐา ภวิสฺสาม อิตรีตเรน จีวเรน, อิตรีตรจีวรสนฺตุฏฺฐิยา จ วณฺณวาทิโน, น จ จีวรเหตุ อเนสนํ อปฺปติรูปํ อาปชฺชิสฺสาม, อลทฺธา จ จีวรํ น จ ปริตสฺสิสฺสาม, ลทฺธา จ จีวรํ อคธิตา อมุจฺฉิตา อนชฺฌาปนฺนา อาทีนวทสฺสาวิโน นิสฺสรณปญฺญา ปริภุญฺชิสฺสาม. (เอวํ สพฺพํ กาตพฺพํ.)}

\gambhira{นิทานวคฺคสํยุตฺตปาฬิ}
\prose{\csromanfolio{400}สนฺตุฏฺฐา ภวิสฺสาม อิตรีตเรน ปิณฺฑปาเตน ฯเปฯ.\csromanspacer{} สนฺตุฏฺฐา ภวิสฺสาม อิตรีตเรน เสนาสเนน ฯเปฯ.\csromanspacer{} สนฺตุฏฺฐา ภวิสฺสาม อิตรีตเรน คิลานปฺปจฺจยเภสชฺชปริกฺขาเรน, อิตรีตรคิลานปฺปจฺจยเภสชฺช- ปริกฺขารสนฺตุฏฺฐิยา จ วณฺณวาทิโน, น จ คิลานปฺปจฺจยเภสชฺช- ปริกฺขารเหตุ อเนสนํ อปฺปติรูปํ อาปชฺชิสฺสาม, อลทฺธา จ คิลานปฺปจฺจยเภสชฺชปริกฺขารํ น ปริตสฺสิสฺสาม, ลทฺธา จ คิลานปฺปจฺจยเภสชฺชปริกฺขารํ อคธิตา อมุจฺฉิตา อนชฺฌาปนฺนา อาทีนวทสฺสาวิโน นิสฺสรณปญฺญา ปริภุญฺชิสฺสามาติ.\csromanspacer{} เอวํ หิ โว ภิกฺขเว สิกฺขิตพฺพํ.\csromanspacer{} กสฺสเปน วา หิ โว ภิกฺขเว โอวทิสฺสามิ โย วา ปนสฺส\footnote{โย วา ปน (สี), โย วา (อิ)} กสฺสปสทิโส, โอวทิเตหิ จ ปน โว ตถตฺตาย ปฏิปชฺชิตพฺพนฺติ.\csromanspacer{}\csromanspacer{}ปฐมํ.}

\csromanheader{3}{2. อโนตฺตปฺปีสุตฺต}
\proseitem{145}{เอวํ เม สุตํ\csromandash{}เอกํ สมยํ อายสฺมา จ มหากสฺสโป อายสฺมา จ สาริปุตฺโต พาราณสิยํ วิหรนฺติ อิสิปตเน มิคทาเย.\csromanspacer{} อถ โข อายสฺมา สาริปุตฺโต สายนฺหสมยํ ปฏิสลฺลานา วุฏฺฐิโต เยนายสฺมา มหากสฺสโป เตนุปสงฺกมิ, อุปสงฺกมิตฺวา อายสฺมตา มหากสฺสเปน สทฺธิํ สมฺโมทิ, สมฺโมทนียํ กถํ สารณียํ วีติสาเรตฺวา เอกมนฺตํ นิสีทิ, เอกมนฺตํ นิสินฺโน โข อายสฺมา สาริปุตฺโต อายสฺมนฺตํ มหากสฺสปํ เอตทโวจ “วุจฺจติ หิทํ อาวุโส กสฺสป อนาตาปี อโนตฺตปฺปี อภพฺโพ สมฺโพธาย อภพฺโพ นิพฺพานาย อภพฺโพ อนุตฺตรสฺส โยคกฺเขมสฺส อธิคมาย, อาตาปี จ โข โอตฺตปฺปี ภพฺโพ สมฺโพธาย ภพฺโพ นิพฺพานาย ภพฺโพ อนุตฺตรสฺส โยคกฺเขมสฺส อธิคมายา”ติ.}

\csromanmatikaanchor{mtk.173}
\csromantocmarkref{subsection}{4. มาตุสุตฺต}{mtk.173}
\bookmark[dest={mtk.173},level=2]{4. มาตุสุตฺต}
\prose{กิตฺตาวตา นุ โข อาวุโส อนาตาปี โหติ อโนตฺตปฺปี อภพฺโพ สมฺโพธาย อภพฺโพ นิพฺพานาย อภพฺโพ อนุตฺตรสฺส โยคกฺเขมสฺส อธิคมาย, กิตฺตาวตา จ ปนาวุโส อาตาปี โหติ โอตฺตปฺปี ภพฺโพ สมฺโพธาย ภพฺโพ นิพฺพานาย ภพฺโพ อนุตฺตรสฺส โยคกฺเขมสฺส อธิคมายาติ.\csromanspacer{} อิธาวุโส ภิกฺขุ “อนุปฺปนฺนา เม ปาปกา อกุสลา ธมฺมา อุปฺปชฺชมานา อนตฺถาย สํวตฺเตยฺยุนฺ”ติ น อาตปฺปํ กโรติ, “อุปฺปนฺนา เม ปาปกา อกุสลา ธมฺมา อปฺปหียมานา อนตฺถาย}

\csromanheader{2}{5. กสฺสปสํยุตฺต}
\csromanmatikaanchor{mtk.174}
\csromantocmarkref{subsection}{5. ปิตุสุตฺต}{mtk.174}
\bookmark[dest={mtk.174},level=2]{5. ปิตุสุตฺต}
\prose{\csromanfolio{401}สํวตฺเตยฺยุนฺ”ติ น อาตปฺปํ กโรติ, “อนุปฺปนฺนา เม กุสลา ธมฺมา นุปฺปชฺชมานา อนตฺถาย สํวตฺเตยฺยุนฺ”ติ น อาตปฺปํ กโรติ, “อุปฺปนฺนา เม กุสลา ธมฺมา นิรุชฺฌมานา อนตฺถาย สํวตฺเตยฺยุนฺ”ติ น อาตปฺปํ กโรติ.\csromanspacer{} เอวํ โข อาวุโส อนาตาปี โหติ.}

\prose{กถญฺจาวุโส อโนตฺตปฺปี โหติ, อิธาวุโส ภิกฺขุ “อนุปฺปนฺนา เม ปาปกา อกุสลา ธมฺมา อุปฺปชฺชมานา อนตฺถาย สํวตฺเตยฺยุนฺ”ติ น โอตฺตปฺปติ, “อุปฺปนฺนา เม ปาปกา อกุสลา ธมฺมา อปฺปหียมานา อนตฺถาย สํวตฺเตยฺยุนฺ”ติ น โอตฺตปฺปติ, “อนุปฺปนฺนา เม กุสลา ธมฺมา นุปฺปชฺชมานา อนตฺถาย สํวตฺเตยฺยุนฺ”ติ น โอตฺตปฺปติ, “อุปฺปนฺนา เม กุสลา ธมฺมา นิรุชฺฌมานา อนตฺถาย สํวตฺเตยฺยุนฺ”ติ น โอตฺตปฺปติ.\csromanspacer{} เอวํ โข อาวุโส อโนตฺตปฺปี โหติ.\csromanspacer{} เอวํ โข อาวุโส อนาตาปี อโนตฺตปฺปี อภพฺโพ สมฺโพธาย อภพฺโพ นิพฺพานาย อภพฺโพ อนุตฺตรสฺส โยคกฺเขมสฺส อธิคมาย.}

\csromanmatikaanchor{mtk.175}
\csromantocmarkref{subsection}{6. ภาตุสุตฺต}{mtk.175}
\bookmark[dest={mtk.175},level=2]{6. ภาตุสุตฺต}
\prose{กถญฺจาวุโส อาตาปี โหติ, อิธาวุโส ภิกฺขุ “อนุปฺปนฺนา เม ปาปกา อกุสลา ธมฺมา อุปฺปชฺชมานา อนตฺถาย สํวตฺเตยฺยุนฺ”ติ อาตปฺปํ กโรติ, “อุปฺปนฺนา เม ปาปกา อกุสลา ธมฺมา อปฺปหียมานา อนตฺถาย สํวตฺเตยฺยุนฺ”ติ อาตปฺปํ กโรติ, “อนุปฺปนฺนา เม กุสลา ธมฺมา ฯเปฯ อาตปฺปํ กโรติ.\csromanspacer{} เอวํ โข อาวุโส อาตาปี โหติ.}

\prose{กถญฺจาวุโส โอตฺตปฺปี โหติ, อิธาวุโส ภิกฺขุ “อนุปฺปนฺนา เม ปาปกา อกุสลา ธมฺมา อุปฺปชฺชมานา อนตฺถาย สํวตฺเตยฺยุนฺ”ติ โอตฺตปฺปติ, “อุปฺปนฺนา เม ปาปกา อกุสลา ธมฺมา อปฺปหียมานา อนตฺถาย สํวตฺเตยฺยุนฺ”ติ โอตฺตปฺปติ. “อนุปฺปนฺนา เม กุสลา ธมฺมา อนุปฺปชฺชมานา อนตฺถาย สํวตฺเตยฺยุนฺ”ติ โอตฺตปฺปติ, “อุปฺปนฺนา เม กุสลา ธมฺมา นิรุชฺฌมานา อนตฺถาย สํวตฺเตยฺยุนฺ”ติ โอตฺตปฺปติ.\csromanspacer{} เอวํ โข อาวุโส โอตฺตปฺปี โหติ.\csromanspacer{} เอวํ โข อาวุโส อาตาปี โอตฺตปฺปี ภพฺโพ สมฺโพธาย ภพฺโพ นิพฺพานาย ภพฺโพ อนุตฺตรสฺส โยคกฺเขมสฺส อธิคมายาติ.\csromanspacer{}\csromanspacer{}ทุติยํ.}

\csromanmatikaanchor{mtk.176}
\csromantocmarkref{subsection}{7. ภคินิสุตฺต}{mtk.176}
\bookmark[dest={mtk.176},level=2]{7. ภคินิสุตฺต}
\csromanheader{3}{3. จนฺทูปมสุตฺต}
\proseitem{146}{สาวตฺถิยํ วิหรติ.\csromanspacer{} จนฺทูปมา ภิกฺขเว กุลานิ อุปสงฺกมถ อปกสฺเสว กายํ อปกสฺส จิตฺตํ นิจฺจนวกา กุเลสุ อปฺปคพฺภา\footnote{อปฺปคพฺพา (ก)}.}

\csromanmatikaanchor{mtk.177}
\csromantocmarkref{subsection}{8. ปุตฺตสุตฺต}{mtk.177}
\bookmark[dest={mtk.177},level=2]{8. ปุตฺตสุตฺต}
\vspace{\tipitakaparskip}
\csromancenter{นิทานวคฺคสํยุตฺตปาฬิ เสยฺยถาปิ ภิกฺขเว ปุริโส ชรุทปานํ วา โอโลเกยฺย ปพฺพตวิสมํ วา นทีวิทุคฺคํ วา อปกสฺเสว กายํ อปกสฺส จิตฺตํ.\csromanspacer{} เอวเมว โข ภิกฺขเว จนฺทูปมา กุลานิ อุปสงฺกมถ อปกสฺเสว กายํ อปกสฺส จิตฺตํ นิจฺจนวกา กุเลสุ อปฺปคพฺภา.}
\prose{\csromanfolio{402}กสฺสโป ภิกฺขเว จนฺทูปโม กุลานิ อุปสงฺกมติ อปกสฺเสว กายํ อปกสฺส จิตฺตํ นิจฺจนวโก กุเลสุ อปฺปคพฺโภ.\csromanspacer{} ตํ กิํ มญฺญถ ภิกฺขเว, กถํรูโป ภิกฺขุ อรหติ กุลานิ อุปสงฺกมิตุนฺติ.\csromanspacer{} ภควํมูลกา โน ภนฺเต ธมฺมา ภควํเนตฺติกา ภควํปฏิสรณา.\csromanspacer{} สาธุ วต ภนฺเต ภควนฺตํเยว ปฏิภาตุ เอตสฺส ภาสิตสฺส อตฺโถ, ภควโต สุตฺวา ภิกฺขู ธาเรสฺสนฺตีติ.}

\csromanmatikaanchor{mtk.178}
\csromantocmarkref{subsection}{9. ธีตุสุตฺต}{mtk.178}
\bookmark[dest={mtk.178},level=2]{9. ธีตุสุตฺต}
\prose{อถ โข ภควา อากาเส ปาณิํ จาเลสิ.\csromanspacer{} เสยฺยถาปิ ภิกฺขเว อยํ อากาเส ปาณิ น สชฺชติ น คยฺหติ น พชฺฌติ.\csromanspacer{} เอวเมว โข ภิกฺขเว ยสฺส กสฺสจิ ภิกฺขุโน กุลานิ อุปสงฺกมโต กุเลสุ จิตฺตํ น สชฺชติ น คยฺหติ น พชฺฌติ “ลภนฺตุ ลาภกามา, ปุญฺญกามา กโรนฺตุ ปุญฺญานี”ติ.\csromanspacer{} ยถาสเกน ลาเภน อตฺตมโน โหติ สุมโน, เอวํ ปเรสํ ลาเภน อตฺตมโน โหติ สุมโน.\csromanspacer{} เอวรูโป โข ภิกฺขเว ภิกฺขุ อรหติ กุลานิ อุปสงฺกมิตุํ.}

\prose{กสฺสปสฺส ภิกฺขเว กุลานิ อุปสงฺกมโต กุเลสุ จิตฺตํ น สชฺชติ น คยฺหติ น พชฺฌติ “ลภนฺตุ ลาภกามา, ปุญฺญกามา กโรนฺตุ ปุญฺญานี”ติ.\csromanspacer{} ยถาสเกน ลาเภน อตฺตมโน โหติ สุมโน, เอวํ ปเรสํ ลาเภน อตฺตมโน โหติ สุมโน.}

\csromanmatikaanchor{mtk.179}
\csromantocmarkref{subsection}{10. เวปุลฺลปพฺพตสุตฺต}{mtk.179}
\bookmark[dest={mtk.179},level=2]{10. เวปุลฺลปพฺพตสุตฺต}
\csromantocmarkref{section}{อุทฺทานคาถา}{mtk.180}
\bookmark[dest={mtk.180},level=1]{อุทฺทานคาถา}
\prose{ตํ กิํ มญฺญถ ภิกฺขเว, กถํรูปสฺส ภิกฺขุโน อปริสุทฺธา ธมฺมเทสนา โหติ, กถํรูปสฺส ภิกฺขุโน ปริสุทฺธา ธมฺมเทสนา โหตีติ.\csromanspacer{} ภควํมูลกา โน ภนฺเต ธมฺมา ภควํเนตฺติกา ภควํปฏิสรณา.\csromanspacer{} สาธุ วต ภนฺเต ภควนฺตํเยว ปฏิภาตุ เอตสฺส ภาสิตสฺส อตฺโถ, ภควโต สุตฺวา ภิกฺขู ธาเรสฺสนฺตีติ.\csromanspacer{} เตน หิ ภิกฺขเว สุณาถ สาธุกํ มนสิ กโรถ, ภาสิสฺสามีติ. “เอวํ ภนฺเต”ติ โข เต ภิกฺขู ภควโต ปจฺจสฺโสสุํ.\csromanspacer{} ภควา เอตทโวจ\csromandash{}}

\prose{โย หิ โกจิ ภิกฺขเว ภิกฺขุ เอวํจิตฺโต ปเรสํ ธมฺมํ เทเสติ “อโห วต เม ธมฺมํ สุเณยฺยุํ, สุตฺวา จ ปน ธมฺมํ ปสีเทยฺยุํ, ปสนฺนา}

\csromanheader{2}{5. กสฺสปสํยุตฺต}
\prose{\csromanfolio{403}จ เม ปสนฺนาการํ กเรยฺยุนฺ”ติ.\csromanspacer{} เอวรูปสฺส โข ภิกฺขเว ภิกฺขุโน อปริสุทฺธา ธมฺมเทสนา โหติ.}

\prose{โย จ โข ภิกฺขเว ภิกฺขุ เอวํจิตฺโต ปเรสํ ธมฺมํ เทเสติ “สฺวากฺขาโต ภควโต ธมฺโม สนฺทิฏฺฐิโก อกาลิโก เอหิปสฺสิโก โอปเนยฺยิโก ปจฺจตฺตํ เวทิตพฺโพ วิญฺญูหีติ\footnote{วิญฺญูหิ (?)}.\csromanspacer{} อโห วต เม ธมฺมํ สุเณยฺยุํ, สุตฺวา จ ปน ธมฺมํ อาชาเนยฺยุํ อาชานิตฺวา จ ปน ตถตฺตาย ปฏิปชฺเชยฺยุนฺ”ติ.\csromanspacer{} อิติ ธมฺมสุธมฺมตํ ปฏิจฺจ ปเรสํ ธมฺมํ เทเสติ, การุญฺญํ ปฏิจฺจ อนุทฺทยํ\footnote{อนุทยํ (พหูสุ) ทฺวิตฺตการณํ ปน คเวสิตพฺพํ.} ปฏิจฺจ อนุกมฺปํ อุปาทาย ปเรสํ ธมฺมํ เทเสติ, เอวรูปสฺส โข ภิกฺขเว ภิกฺขุโน ปริสุทฺธา ธมฺมเทสนา โหติ.}

\prose{กสฺสโป ภิกฺขเว เอวํจิตฺโต ปเรสํ ธมฺมํ เทเสติ “สฺวากฺขาโต ภควตา ธมฺโม สนฺทิฏฺฐิโก อกาลิโก เอหิปสฺสิโก โอปเนยฺยิโก ปจฺจตฺตํ เวทิตพฺโพ วิญฺญูหีติ.\csromanspacer{} อโห วต เม ธมฺมํ สุเณยฺยุํ, สุตฺวา จ ปน ธมฺมํ อาชาเนยฺยุํ, อาชานิตฺวา จ ปน ตถตฺตาย ปฏิปชฺเชยฺยุนฺ”ติ.\csromanspacer{} อิติ ธมฺมสุธมฺมตํ ปฏิจฺจ ปเรสํ ธมฺมํ เทเสติ, การุญฺญํ ปฏิจฺจ อนุทฺทยํ ปฏิจฺจ อนุกมฺปํ อุปาทาย ปเรสํ ธมฺมํ เทเสติ.\csromanspacer{} กสฺสเปน วา หิ โว ภิกฺขเว โอวทิสฺสามิ โย วา ปนสฺส กสฺสปสทิโส, โอวทิเตหิ จ ปน โว ตถตฺตาย ปฏิปชฺชิตพฺพนฺติ.\csromanspacer{}\csromanspacer{}ตติยํ.}

\csromanheader{3}{4. กุลูปกสุตฺต}
\proseitem{147}{สาวตฺถิยํ วิหรติ.\csromanspacer{} ตํ กิํ มญฺญถ ภิกฺขเว, กถํรูโป ภิกฺขุ อรหติ กุลูปโก โหตุํ, กถํรูโป ภิกฺขุ น อรหติ กุลูปโก โหตุนฺติ.\csromanspacer{} ภควํมูลกา โน ภนฺเต ธมฺมา ฯเปฯ.\csromanspacer{} ภควา เอตทโวจ\csromandash{}}

\prose{โย หิ โกจิ ภิกฺขเว ภิกฺขุ เอวํจิตฺโต กุลานิ อุปสงฺกมติ “เทนฺตุเยว เม มา นาทํสุ, พหุกญฺเญว เม เทนฺตุ มา โถกํ, ปณีตญฺเญว เม เทนฺตุ มา ลูขํ, สีฆญฺเญว เม เทนฺตุ มา ทนฺธํ, สกฺกจฺจญฺเญว เม เทนฺตุ มา อสกฺกจฺจนฺ”ติ.\csromanspacer{} ตสฺส เจ ภิกฺขเว ภิกฺขุโน เอวํจิตฺตสฺส กุลานิ อุปสงฺกมโต น เทนฺติ, เตน ภิกฺขุ สนฺทียติ.\csromanspacer{} โส ตโตนิทานํ ทุกฺขํ โทมนสฺสํ ปฏิสํเวทยติ.\csromanspacer{} โถกํ เทนฺติ โน พหุกํ ฯเปฯ ลูขํ เทนฺติ โน ปณีตํ, ทนฺธํ เทนฺติ โน สีฆํ, เตน ภิกฺขุ สนฺทียติ.\csromanspacer{} โส}

\vspace{\tipitakaparskip}
\csromancenter{นิทานวคฺคสํยุตฺตปาฬิ ตโตนิทานํ ทุกฺขํ โทมนสฺสํ ปฏิสํเวทยติ.\csromanspacer{} อสกฺกจฺจํ เทนฺติ โน สกฺกจฺจํ, เตน ภิกฺขุ สนฺทียติ.\csromanspacer{} โส ตโตนิทานํ ทุกฺขํ โทมนสฺสํ ปฏิสํเวทยติ.\csromanspacer{} เอวรูโป โข ภิกฺขเว ภิกฺขุ น อรหติ กุลูปโก โหตุํ.}
\prose{\csromanfolio{404}โย จ โข ภิกฺขเว ภิกฺขุ เอวํจิตฺโต กุลานิ อุปสงฺกมติ “ตํ กุเตตฺถ ลพฺภา ปรกุเลสุ, เทนฺตุเยว เม มา นาทํสุ, พหุกญฺเญว เม เทนฺตุ มา โถกํ, ปณีตญฺเญว เม เทนฺตุ มา ลูขํ, สีฆญฺเญว เม เทนฺตุ มา ทนฺธํ, สกฺกจฺจญฺเญว เม เทนฺตุ มา อสกฺกจฺจนฺ”ติ.\csromanspacer{} ตสฺส เจ ภิกฺขเว ภิกฺขุโน เอวํจิตฺตสฺส กุลานิ อุปสงฺกมโต น เทนฺติ, เตน ภิกฺขุ น สนฺทียติ.\csromanspacer{} โส น ตโตนิทานํ ทุกฺขํ โทมนสฺสํ ปฏิสํเวทยติ.\csromanspacer{} โถกํ เทนฺติ โน พหุกํ, เตน ภิกฺขุ น สนฺทียติ.\csromanspacer{} โส น ตโตนิทานํ ทุกฺขํ โทมนสฺสํ ปฏิสํเวทยติ.\csromanspacer{} ลูขํ เทนฺติ โน ปณีตํ, เตน ภิกฺขุ น สนฺทียติ.\csromanspacer{} โส น ตโตนิทานํ ทุกฺขํ โทมนสฺสํ ปฏิสํเวทยติ.\csromanspacer{} ทนฺธํ เทนฺติ โน สีฆํ, เตน ภิกฺขุ น สนฺทียติ.\csromanspacer{} โส น ตโตนิทานํ ทุกฺขํ โทมนสฺสํ ปฏิสํเวทยติ.\csromanspacer{} อสกฺกจฺจํ เทนฺติ โน สกฺกจฺจํ, เตน ภิกฺขุ น สนฺทียติ.\csromanspacer{} โส น ตโตนิทานํ ทุกฺขํ โทมนสฺสํ ปฏิสํเวทยติ.\csromanspacer{} เอวรูโป โข ภิกฺขเว ภิกฺขุ อรหติ กุลูปโก โหตุํ.}

\prose{กสฺสโป ภิกฺขเว เอวํจิตฺโต กุลานิ อุปสงฺกมติ “ตํ กุเตตฺถ ลพฺภา ปรกุเลสุ, เทนฺตุเยว เม มา นาทํสุ, พหุกญฺเญว เม เทนฺตุ มา โถกํ, ปณีตญฺเญว เม เทนฺตุ มา ลูขํ, สีฆญฺเญว เม เทนฺตุ มา ทนฺธํ, สกฺกจฺจญฺเญว เม เทนฺตุ มา อสกฺกจฺจนฺ”ติ.\csromanspacer{} ตสฺส เจ ภิกฺขเว กสฺสปสฺส เอวํจิตฺตสฺส กุลานิ อุปสงฺกมโต น เทนฺติ, เตน กสฺสโป น สนฺทียติ.\csromanspacer{} โส น ตโตนิทานํ ทุกฺขํ โทมนสฺสํ ปฏิสํเวทยติ.\csromanspacer{} โถกํ เทนฺติ โน พหุกํ, เตน กสฺสโป น สนฺทียติ, โส น ตโตนิทานํ ทุกฺขํ โทมนสฺสํ ปฏิสํเวทยติ.\csromanspacer{} ลูขํ เทนฺติ โน ปณีตํ, เตน กสฺสโป น สนฺทียติ.\csromanspacer{} โส น ตโตนิทานํ ทุกฺขํ โทนมสฺสํ ปฏิสํเวทยติ.\csromanspacer{} ทนฺธํ เทนฺติ โน สีฆํ, เตน กสฺสโป น สนฺทียติ.\csromanspacer{} โส น ตโตนิทานํ ทุกฺขํ โทมนสฺสํ ปฏิสํเวทยติ.\csromanspacer{} อสกจฺจํ เทนฺติ โน สกฺกจฺจํ, เตน กสฺสโป น สนฺทียติ.\csromanspacer{} โส น ตโตนิทานํ ทุกฺขํ โทมนสฺสํ ปฏิสํเวทยติ.\csromanspacer{} กสฺสเปน วา หิ โว ภิกฺขเว โอวทิสฺสามิ โย วา ปนสฺส กสฺสปสทิโส, โอวทิเตหิ จ ปน โว ตถตฺตาย ปฏิปชฺชิตพฺพนฺติ.\csromanspacer{}\csromanspacer{}จตุตฺถํ.}

\csromanheader{2}{5. กสฺสปสํยุตฺต}
\csromanheader{3}{5. ชิณฺณสุตฺต}
\proseitem{148}{\csromanfolio{405}เอวํ เม สุตํ\csromandash{}ราชคเห เวฬุวเน.\csromanspacer{} อถ โข อายสฺมา มหากสฺสโป เยน ภควา เตนุปสงฺกมิ, อุปสงฺกมิตฺวา ภควนฺตํ อภิวาเทตฺวา เอกมนฺตํ นิสีทิ, เอกมนฺตํ นิสินฺนํ โข อายสฺมนฺตํ มหากสฺสปํ ภควา เอตทโวจ “ชิณฺโณสิ ทานิ ตฺวํ กสฺสป ครุกานิ จ เต อิมานิ สาณานิ ปํสุกูลานิ นิพฺพสนานิ, ตสฺมาติห ตฺวํ กสฺสป คหปตานิ\footnote{คหปติกานิ (สี)} เจว จีวรานิ ธาเรหิ, นิมนฺตนานิ จ ภุญฺชาหิ, มม จ สนฺติเก วิหราหี”ติ.}

\csromanmatikaanchor{mtk.181}
\csromantocmarknopage{subsection}{5. กสฺสปสํยุตฺต}{mtk.181}
\bookmark[dest={mtk.181},level=2]{5. กสฺสปสํยุตฺต}
\prose{อหํ โข ภนฺเต ทีฆรตฺตํ อารญฺญิโก เจว อารญฺญิกตฺตสฺส จ วณฺณวาที, ปิณฺฑปาติโก เจว ปิณฺฑปาติกตฺตสฺส จ วณฺณวาที, ปํสุกูลิโก เจว ปํสุกูลิกตฺตสฺส จ วณฺณวาที, เตจีวริโก เจว เตจีวริกตฺตสฺส จ วณฺณวาที, อปฺปิจฺโฉ เจว อปฺปิจฺฉตาย จ วณฺณวาที, สนฺตุฏฺโฐ เจว สนฺตุฏฺฐิยา จ วณฺณวาที, ปวิวิตฺโต เจว ปวิเวกสฺส จ วณฺณวาที, อสํสฏฺโฐ เจว อสํสคฺคสฺส จ วณฺณวาที, อารทฺธวีริโย เจว วีริยารมฺภสฺส\footnote{วีริยารพฺภสฺส (ก)} จ วณฺณวาทีติ.}

\csromanmatikaanchor{mtk.182}
\csromantocmarkref{subsubsection}{1. สนฺตุฏฺฐสุตฺต}{mtk.182}
\bookmark[dest={mtk.182},level=3]{1. สนฺตุฏฺฐสุตฺต}
\prose{กิํ\footnote{กํ (ก)} ปน ตฺวํ กสฺสป อตฺถวสํ สมฺปสฺสมาโน ทีฆรตฺตํ อารญฺญิโก เจว อารญฺญิกตฺตสฺส จ วณฺณวาที, ปิณฺฑปาติโก เจว ฯเปฯ ปํสุกูลิโก เจว.\csromanspacer{} เตจีวริโก เจว.\csromanspacer{} อปฺปิจฺโฉ เจว.\csromanspacer{} สนฺตุฏฺโฐ เจว.\csromanspacer{} ปวิวิตฺโต เจว.\csromanspacer{} อสํสฏฺโฐ เจว.\csromanspacer{} อารทฺธวีริโย เจว วีริยารมฺภสฺส จ วณฺณวาทีติ.}

\prose{เทฺว ขฺวาหํ ภนฺเต อตฺถวเส สมฺปสฺสมาโน ทีฆรตฺตํ อารญฺญิโก เจว อารญฺญิกตฺตสฺส จ วณฺณวาที, ปิณฺฑปาติโก เจว ฯเปฯ ปํสุกูลิโก เจว.\csromanspacer{} เตจีวริโก เจว.\csromanspacer{} อปฺปิจฺโฉ เจว.\csromanspacer{} สนฺตุฏฺโฐ เจว.\csromanspacer{} ปวิวิตฺโต เจว.\csromanspacer{} อสํสฏฺโฐ เจว.\csromanspacer{} อารทฺธวีริโย เจว วีริยารมฺภสฺส จ วณฺณวาที.}

\prose{อตฺตโน จ ทิฏฺฐธมฺมสุขวิหารํ สมฺปสฺสมาโน ปจฺฉิมญฺจ ชนตํ อนุกมฺปมาโน อปฺเปว นาม ปจฺฉิมา ชนตา ทิฏฺฐานุคติํ อาปชฺเชยฺยุํ\footnote{อาปชฺเชยฺย (สี, สฺยา, กํ)}.\csromanspacer{} เย กิร เต อเหสุํ พุทฺธานุพุทฺธสาวกา, เต ทีฆรตฺตํ อารญฺญิกา เจว อเหสุํ อารญฺญิกตฺตสฺส จ วณฺณวาทิโน ฯเปฯ ปิณฺฑปาติกา เจว}

\vspace{\tipitakaparskip}
\csromancenter{นิทานวคฺคสํยุตฺตปาฬิ อเหสุํ ฯเปฯ ปํสุกูลิกา เจว อเหสุํ.\csromanspacer{} เตจีวริกา เจว อเหสุํ.\csromanspacer{} อปฺปิจฺฉา เจว อเหสุํ.\csromanspacer{} สนฺตุฏฺฐา เจว อเหสุํ.\csromanspacer{} ปวิวิตฺตา เจว อเหสุํ.\csromanspacer{} อสํสฏฺฐา เจว อเหสุํ.\csromanspacer{} อารทฺธวีริยา เจว อเหสุํ วีริยารมฺภสฺส จ วณฺณวาทิโน”ติ.\csromanspacer{} เต ตถตฺตาย ปฏิปชฺชิสฺสนฺติ, เตสํ ตํ ภวิสฺสติ ทีฆรตฺตํ หิตาย สุขาย.}
\prose{\csromanfolio{406}อิเม ขฺวาหํ ภนฺเต เทฺว อตฺถวเส สมฺปสฺสมาโน ทีฆรตฺตํ อารญฺญิโก เจว อารญฺญิกตฺตสฺส จ วณฺณวาที, ปิณฺฑปาติโก เจว ฯเปฯ ปํสุกูลิโก เจว.\csromanspacer{} เตจีวริโก เจว.\csromanspacer{} อปฺปิจฺโฉ เจว.\csromanspacer{} สนฺตุฏฺโฐ เจว.\csromanspacer{} ปวิวิตฺโต เจว.\csromanspacer{} อสํสฏฺโฐ เจว.\csromanspacer{} อารทฺธวีริโย เจว วีริยารมฺภสฺส จ วณฺณวาทีติ.}

\prose{สาธุ สาธุ กสฺสป, พหุชนหิตาย กิร ตฺวํ กสฺสป ปฏิปนฺโน พหุชนสุขาย โลกานุกมฺปาย อตฺถาย หิตาย สุขาย เทวมนุสฺสานํ, ตสฺมาติห ตฺวํ กสฺสป สาณานิ เจว ปํสุกูลานิ ธาเรหิ นิพฺพสนานิ, ปิณฺฑาย จ จราหิ, อรญฺเญ จ วิหราหีติ.\csromanspacer{}\csromanspacer{}ปญฺจมํ.}

\csromanheader{3}{6. โอวาทสุตฺต}
\proseitem{149}{ราชคเห เวฬุวเน.\csromanspacer{} อถ โข อายสฺมา มหากสฺสโป เยน ภควา เตนุปสงฺกมิ, อุปสงฺกมิตฺวา ภควนฺตํ อภิวาเทตฺวา เอกมนฺตํ นิสีทิ เอกมนฺตํ นิสินฺนํ โข อายสฺมนฺตํ มหากสฺสปํ ภควา เอตทโวจ “โอวท กสฺสป ภิกฺขู, กโรหิ กสฺสป ภิกฺขูนํ ธมฺมิํ กถํ, อหํ วา กสฺสป ภิกฺขู โอวเทยฺยํ ตฺวํ วา, อหํ วา ภิกฺขูนํ ธมฺมิํ กถํ กเรยฺยํ ตฺวํ วา”ติ.}

\csromanmatikaanchor{mtk.183}
\csromantocmarkref{subsubsection}{2. อโนตฺตปฺปีสุตฺต}{mtk.183}
\bookmark[dest={mtk.183},level=3]{2. อโนตฺตปฺปีสุตฺต}
\prose{ทุพฺพจา โข ภนฺเต เอตรหิ ภิกฺขู, โทวจสฺสกรเณหิ ธมฺเมหิ สมนฺนาคตา, อกฺขมา อปฺปทกฺขิณคฺคาหิโน อนุสาสนิํ, อิธาหํ ภนฺเต อทฺทสํ ภณฺฑญฺจ\footnote{ภณฺฑุญฺจ (สี)} นาม ภิกฺขุํ อานนฺทสฺส สทฺธิวิหาริํ อภิชิกญฺจ\footnote{อาภิญฺชิกญฺจ (สี, ก), อาภิชฺชิกญฺจ (สฺยา, กํ)} นาม ภิกฺขุํ อนุรุทฺธสฺส สทฺธิวิหาริํ อญฺญมญฺญํ สุเตน อจฺจาวทนฺเต “เอหิ ภิกฺขุ, โก พหุตรํ ภาสิสฺสติ, โก สุนฺทรตรํ ภาสิสฺสติ, โก จิรตรํ ภาสิสฺสตี”ติ.}

\csromanheader{2}{5. กสฺสปสํยุตฺต}
\prose{\csromanfolio{407}อถ โข ภควา อญฺญตรํ ภิกฺขุํ อามนฺเตสิ “เอหิ ตฺวํ ภิกฺขุ มม วจเนน ภณฺฑญฺจ ภิกฺขุํ อานนฺทสฺส สทฺธิวิหาริํ อภิชิกญฺจ ภิกฺขุํ อนุรุทฺธสฺส สทฺธิวิหาริํ อามนฺเตหิ ‘สตฺถา อายสฺมนฺเต อามนฺเตตี’ติ”. “เอวํ ภนฺเต”ติ โข โส ภิกฺขุ ภควโต ปฏิสฺสุตฺวา เยน เต ภิกฺขู เตนุปสงฺกมิ, อุปสงฺกมิตฺวา เต ภิกฺขู เอตทโวจ “สตฺถา อายสฺมนฺเต อามนฺเตตี”ติ.}

\prose{“เอวมาวุโส”ติ โข เต ภิกฺขู ตสฺส ภิกฺขุโน ปฏิสฺสุตฺวา เยน ภควา เตนุปสงฺกมิํสุ, อุปสงฺกมิตฺวา ภควนฺตํ อภิวาเทตฺวา เอกมนฺตํ นิสีทิํสุ, เอกมนฺตํ นิสินฺเน โข เต ภิกฺขู ภควา เอตทโวจ “สจฺจํ กิร ตุเมฺห ภิกฺขเว อญฺญมญฺญํ สุเตน อจฺจาวทถ ‘เอหิ ภิกฺขุ, โก พหุตรํ ภาสิสฺสติ, โก สุนฺทรตรํ ภาสิสฺสติ, โก จิรตรํ ภาสิสฺสตี’ติ”.\csromanspacer{} เอวํ ภนฺเต.\csromanspacer{} กิํ นุ โข เม ตุเมฺห ภิกฺขเว เอวํ ธมฺมํ เทสิตํ อาชานาถ “เอถ ตุเมฺห ภิกฺขเว อญฺญมญฺญํ สุเตน อจฺจาวทถ ‘เอหิ ภิกฺขุ, โก พหุตรํ ภาสิสฺสติ, โก สุนฺทรตรํ ภาสิสฺสติ, โก จิรตรํ ภาสิสฺสตี’ติ”.\csromanspacer{} โน เหตํ ภนฺเต.\csromanspacer{} โน เจ กิร เม ตุเมฺห ภิกฺขเว เอวํ ธมฺมํ เทสิตํ อาชานาถ, อถ กิญฺจรหิ ตุเมฺห โมฆปุริสา กิํชานนฺตา กิํปสฺสนฺตา เอวํ สฺวากฺขาเต ธมฺมวินเย ปพฺพชิตา สมานา อญฺญมญฺญํ สุเตน อจฺจาวทถ “เอหิ ภิกฺขุ, โก พหุตรํ ภาสิสฺสติ, โก สุนฺทรตรํ ภาสิสฺสติ, โก จิรตรํ ภาสิสฺสตี”ติ.}

\prose{อถ โข เต ภิกฺขู ภควโต ปาเทสุ สิรสา นิปติตฺวา ภควนฺตํ เอตทโวจุํ “อจฺจโย โน ภนฺเต อจฺจคมา ยถาพาเล ยถามูเฬฺห ยถา- อกุสเล\footnote{ยถา พาเล ยถา มูเฬฺห ยถา อกุสเล (อิ), ยถาพาลํ ยถามูฬฺหํ ยถา-อกุสลํ (?)}, เย มยํ เอวํ สฺวากฺขาเต ธมฺมวินเย ปพฺพชิตา สมานา อญฺญมญฺญํ สุเตน อจฺจาวทิมฺห ‘เอหิ ภิกฺขุ, โก พหุตรํ ภาสิสฺสติ, โก สุนฺทรตรํ ภาสิสฺสติ, โก จิรตรํ ภาสิสฺสตี’ติ.\csromanspacer{} เตสํ โน ภนฺเต ภควา อจฺจยํ อจฺจยโต ปฏิคฺคณฺหาตุ อายติํ สํวรายา”ติ.}

\prose{ตคฺฆ ตุเมฺห ภิกฺขเว อจฺจโย อจฺจคมา ยถาพาเล ยถามูเฬฺห ยถา-อกุสเล, เย ตุเมฺห เอวํ สฺวากฺขาเต ธมฺมวินเย ปพฺพชิตา สมานา อญฺญมญฺญํ สุเตน อจฺจาวทิตฺถ “เอหิ ภิกฺขุ, โก พหุตรํ ภาสิสฺสติ,}

\vspace{\tipitakaparskip}
\csromancenter{นิทานวคฺคสํยุตฺตปาฬิ โก สุนฺทรตรํ ภาสิสฺสติ, โก จิรตรํ ภาสิสฺสตี”ติ.\csromanspacer{} ยโต จ โข ตุเมฺห ภิกฺขเว อจฺจยํ อจฺจยโต ทิสฺวา ยถาธมฺมํ ปฏิกโรถ, ตํ โว มยํ\footnote{มยํ อจฺจยํ (สี)} ปฏิคฺคณฺหาม, วุทฺธิ เหสา ภิกฺขเว อริยสฺส วินเย, โย อจฺจยํ อจฺจยโต ทิสฺวา ยถาธมฺมํ ปฏิกโรติ, อายติํ จ สํวรํ อาปชฺชตีติ.\csromanspacer{}\csromanspacer{}ฉฏฺฐํ.}
\csromanheader{3}{7. ทุติย-โอวาทสุตฺต}
\csromanmatikaanchor{mtk.184}
\csromantocmarkref{subsubsection}{3. จนฺทูปมสุตฺต}{mtk.184}
\bookmark[dest={mtk.184},level=3]{3. จนฺทูปมสุตฺต}
\proseitem{150}{\csromanfolio{408}ราชคเห วิหรติ เวฬุวเน\footnote{สาวตฺถิ, ตตฺร-เอตทโวจ (สี)}.\csromanspacer{} อถ โข อายสฺมา มหากสฺสโป เยน ภควา เตนุปสงฺกมิ ฯเปฯ เอกมนฺตํ นิสินฺนํ โข อายสฺมนฺตํ มหากสฺสปํ ภควา เอตทโวจ “โอวท กสฺสป ภิกฺขู, กโรหิ กสฺสป ภิกฺขูนํ ธมฺมิํ กถํ, อหํ วา กสฺสป ภิกฺขู โอวเทยฺยํ ตฺวํ วา, อหํ วา ภิกฺขูนํ ธมฺมิํ กถํ กเรยฺยํ ตฺวํ วา”ติ.}

\prose{ทุพฺพจา โข ภนฺเต เอตรหิ ภิกฺขู โทวจสฺสกรเณหิ ธมฺเมหิ สมนฺนาคตา อกฺขมา อปฺปทกฺขิณคฺคาหิโน อนุสาสนิํ, ยสฺส กสฺสจิ ภนฺเต สทฺธา นตฺถิ กุสเลสุ ธมฺเมสุ.\csromanspacer{} หิรี\footnote{หิริ (สพฺพตฺถ)} นตฺถิ กุสเลสุ ธมฺเมสุ.\csromanspacer{} โอตฺตปฺปํ นตฺถิ กุสเลสุ ธมฺเมสุ.\csromanspacer{} วีริยํ นตฺถิ กุสเลสุ ธมฺเมสุ.\csromanspacer{} ปญฺญา นตฺถิ กุสเลสุ ธมฺเมสุ, ตสฺส ยา รตฺติ วา ทิวโส วา อาคจฺฉติ\footnote{อาคจฺฉนฺติ (สี)}, หานิเยว ปาฏิกงฺขา กุสเลสุ ธมฺเมสุ โน วุทฺธิ.}

\prose{เสยฺยถาปิ ภนฺเต กาฬปกฺเข จนฺทสฺส ยา รตฺติ วา ทิวโส วา อาคจฺฉติ, หายเตว วณฺเณน, หายติ มณฺฑเลน, หายติ อาภาย, หายติ อาโรหปริณาเหน.\csromanspacer{} เอวเมว โข ภนฺเต ยสฺส กสฺสจิ สทฺธา นตฺถิ กุสเลสุ ธมฺเมสุ ฯเปฯ.\csromanspacer{} หิรี นตฺถิ.\csromanspacer{} โอตฺตปฺปํ นตฺถิ.\csromanspacer{} วีริยํ นตฺถิ.\csromanspacer{} ปญฺญา นตฺถิ กุสเลสุ ธมฺเมสุ, ตสฺส ยา รตฺติ วา ทิวโส วา อาคจฺฉติ, หานิเยว ปาฏิกงฺขา กุสเลสุ ธมฺเมสุ โน วุทฺธิ.}

\prose{“อสฺสทฺโธ ปุริสปุคฺคโล”ติ ภนฺเต ปริหานเมตํ. “อหิริโก ปุริสปุคฺคโล”ติ ภนฺเต ปริหานเมตํ. “อโนตฺตปฺปี ปุริสปุคฺคโล”ติ ภนฺเต ปริหานเมตํ. “กุสีโต ปุริสปุคฺคโล”ติ ภนฺเต ปริหานเมตํ. “ทุปฺปญฺโญ ปุริสปุคฺคโล”ติ ภนฺเต ปริหานเมตํ. “โกธโน}

\csromanheader{2}{5. กสฺสปสํยุตฺต}
\prose{\csromanfolio{409}ปุริสปุคฺคโล”ติ ภนฺเต ปริหานเมตํ. “อุปนาหี ปุริสปุคฺคโล”ติ ภนฺเต ปริหานเมตํ. “น สนฺติ ภิกฺขู โอวาทกา”ติ ภนฺเต ปริหานเมตํ.}

\prose{ยสฺส กสฺสจิ ภนฺเต สทฺธา อตฺถิ กุสเลสุ ธมฺเมสุ.\csromanspacer{} หิรี อตฺถิ กุสเลสุ ธมฺเมสุ.\csromanspacer{} โอตฺตปฺปํ อตฺถิ กุสเลสุ ธมฺเมสุ.\csromanspacer{} วีริยํ อตฺถิ กุสเลสุ ธมฺเมสุ.\csromanspacer{} ปญฺญา อตฺถิ กุสเลสุ ธมฺเมสุ, ตสฺส ยา รตฺติ วา ทิวโส วา อาคจฺฉติ, วุทฺธิเยว ปาฏิกงฺขา กุสเลสุ ธมฺเมสุ โน ปริหานิ.}

\prose{เสยฺยถาปิ ภนฺเต ชุณฺหปกฺเข จนฺทสฺส ยา รตฺติ วา ทิวโส วา อาคจฺฉติ, วฑฺฒเตว วณฺเณน, วฑฺฒติ มณฺฑเลน, วฑฺฒติ อาภาย, วฑฺฒติ อาโรหปริณาเหน.\csromanspacer{} เอวเมว โข ภนฺเต ยสฺส กสฺสจิ สทฺธา อตฺถิ กุสเลสุ ธมฺเมสุ.\csromanspacer{} หิรี อตฺถิ ฯเปฯ.\csromanspacer{} โอตฺตปฺปํ อตฺถิ.\csromanspacer{} วีริยํ อตฺถิ.\csromanspacer{} ปญฺญา อตฺถิ กุสเลสุ ธมฺเมสุ, ตสฺส ยา รตฺติ วา ทิวโส วา อาคจฺฉติ, วุทฺธิเยว ปาฏิกงฺขา กุสเลสุ ธมฺเมสุ โน ปริหานิ.}

\prose{“สทฺโธ ปุริสปุคฺคโล”ติ ภนฺเต อปริหานเมตํ. “หิริมา ปุริสปุคฺคโล”ติ ภนฺเต อปริหานเมตํ. “โอตฺตปฺปี ปุริสปุคฺคโล”ติ ภนฺเต อปริหานเมตํ. “อารทฺธวีริโย ปุริสปุคฺคโล”ติ ภนฺเต อปริหานเมตํ. “ปญฺญวา ปุริสปุคฺคโล”ติ ภนฺเต อปริหานเมตํ. “อกฺโกธโน ปุริสปุคฺคโล”ติ ภนฺเต อปริหานเมตํ. “อนุปนาหี ปุริสปุคฺคโล”ติ ภนฺเต อปริหานเมตํ. “สนฺติ ภิกฺขู โอวาทกา”ติ ภนฺเต อปริหานเมตนฺติ.}

\prose{สาธุ สาธุ กสฺสป, ยสฺส กสฺสจิ กสฺสป สทฺธา นตฺถิ กุสเลสุ ธมฺเมสุ.\csromanspacer{} หิรี นตฺถิ ฯเปฯ.\csromanspacer{} โอตฺตปฺปํ นตฺถิ.\csromanspacer{} วีริยํ นตฺถิ.\csromanspacer{} ปญฺญา นตฺถิ กุสเลสุ ธมฺเมสุ, ตสฺส ยา รตฺติ วา ทิวโส วา อาคจฺฉติ, หานิเยว ปาฏิกงฺขา กุสเลสุ ธมฺเมสุ โน วุทฺธิ.}

\prose{เสยฺยถาปิ กสฺสป กาฬปกฺเข จนฺทสฺส ยา รตฺติ วา ทิวโส วา อาคจฺฉติ, หายเตว วณฺเณน ฯเปฯ หายติ อาโรหปริณาเหน.\csromanspacer{} เอวเมว โข กสฺสป ยสฺส กสฺสจิ สทฺธา นตฺถิ กุสเลสุ ธมฺเมสุ ฯเปฯ.\csromanspacer{} หิรี นตฺถิ.\csromanspacer{} โอตฺตปฺปํ นตฺถิ.\csromanspacer{} วีริยํ นตฺถิ.\csromanspacer{} ปญฺญา นตฺถิ กุสเลสุ ธมฺเมสุ, ตสฺส ยา รตฺติ วา ทิวโส วา อาคจฺฉติ, หานิเยว ปาฏิกงฺขา กุสเลสุ ธมฺเมสุ โน วุทฺธิ. “อสฺสทฺโธ ปุริสปุคฺคโล”ติ กสฺสป}

\vspace{\tipitakaparskip}
\csromancenter{นิทานวคฺคสํยุตฺตปาฬิ ปริหานเมตํ.\csromanspacer{} อหิริโก ฯเปฯ.\csromanspacer{} อโนตฺตปฺปี.\csromanspacer{} กุสีโต.\csromanspacer{} ทุปฺปญฺโญ.\csromanspacer{} โกธโน. “อุปนาหี ปุริสปุคฺคโล”ติ กสฺสป ปริหานเมตํ. “น สนฺติ ภิกฺขู โอวาทกา”ติ กสฺสป ปริหานเมตํ.}
\csromanmatikaanchor{mtk.185}
\csromantocmarkref{subsubsection}{4. กุลูปกสุตฺต}{mtk.185}
\bookmark[dest={mtk.185},level=3]{4. กุลูปกสุตฺต}
\prose{\csromanfolio{410}ยสฺส กสฺสจิ กสฺสป สทฺธา อตฺถิ กุสเลสุ ธมฺเมสุ ฯเปฯ.\csromanspacer{} หิรี อตฺถิ.\csromanspacer{} โอตฺตปฺปํ อตฺถิ.\csromanspacer{} วีริยํ อตฺถิ.\csromanspacer{} ปญฺญา อตฺถิ กุสเลสุ ธมฺเมสุ, ตสฺส ยา รตฺติ วา ทิวโส วา อาคจฺฉติ, วุทฺธิเยว ปาฏิกงฺขา กุสเลสุ ธมฺเมสุ โน ปริหานิ.}

\prose{เสยฺยถาปิ กสฺสป ชุณฺหปกฺเข จนฺทสฺส ยา รตฺติ วา ทิวโส วา อาคจฺฉติ, วฑฺฒเตว วณฺเณน, วฑฺฒติ มณฺฑเลน, วฑฺฒติ อาภาย, วฑฺฒติ อาโรหปริณาเหน.\csromanspacer{} เอวเมว โข กสฺสป ยสฺส กสฺสจิ สทฺธา อตฺถิ กุสเลสุ ธมฺเมสุ.\csromanspacer{} หิรี อตฺถิ.\csromanspacer{} โอตฺตปฺปํ อตฺถิ.\csromanspacer{} วีริยํ อตฺถิ.\csromanspacer{} ปญฺญา อตฺถิ กุสเลสุ ธมฺเมสุ, ตสฺส ยา รตฺติ วา ทิวโส วา อาคจฺฉติ, วุทฺธิเยว ปาฏิกงฺขา กุสเลสุ ธมฺเมสุ โน ปริหานิ.}

\prose{“สทฺโธ ปุริสปุคฺคโล”ติ กสฺสป อปริหานเมตํ.\csromanspacer{} หิริมา ฯเปฯ.\csromanspacer{} โอตฺตปฺปี.\csromanspacer{} อารทฺธวีริโย.\csromanspacer{} ปญฺญวา.\csromanspacer{} อกฺโกธโน. “อนุปนาหี ปุริสปุคฺคโล”ติ กสฺสป อปริหานเมตํ. “สนฺติ ภิกฺขู โอวาทกา”ติ กสฺสป อปริหานเมตนฺติ.\csromanspacer{}\csromanspacer{}สตฺตมํ.}

\csromanheader{3}{8. ตติย-โอวาทสุตฺต}
\proseitem{151}{ราชคเห กลนฺทกนิวาเป\footnote{สาวตฺถิ, อาราเม (สี)}.\csromanspacer{} อถ โข อายสฺมา มหากสฺสโป เยน ภควา เตนุปสงฺกมิ, อุปสงฺกมิตฺวา ภควนฺตํ อภิวาเทตฺวา เอกมนฺตํ นิสีทิ, เอกมนฺตํ นิสินฺนํ โข อายสฺมนฺตํ มหากสฺสปํ ภควา เอตทโวจ “โอวท กสฺสป ภิกฺขู, กโรหิ กสฺสป ภิกฺขูนํ ธมฺมิํ กถํ, อหํ วา กสฺสป ภิกฺขู โอวเทยฺยํ ตฺวํ วา, อหํ วา ภิกฺขูนํ ธมฺมิํ กถํ กเรยฺยํ ตฺวํ วา”ติ.}

\prose{ทุพฺพจา โข ภนฺเต เอตรหิ ภิกฺขู โทวจสฺสกรเณหิ ธมฺเมหิ สมนฺนาคตา อกฺขมา อปฺปทกฺขิณคฺคาหิโน อนุสาสนินฺติ.\csromanspacer{} ตถา หิ ปน กสฺสป ปุพฺเพ เถรา ภิกฺขู อารญฺญิกา เจว อเหสุํ อารญฺญิกตฺตสฺส จ วณฺณวาทิโน, ปิณฺฑปาติกา เจว อเหสุํ ปิณฺฑปาติกตฺตสฺส จ}

\csromanheader{2}{5. กสฺสปสํยุตฺต}
\csromanmatikaanchor{mtk.186}
\csromantocmarkref{subsubsection}{5. ชิณฺณสุตฺต}{mtk.186}
\bookmark[dest={mtk.186},level=3]{5. ชิณฺณสุตฺต}
\prose{\csromanfolio{411}วณฺณวาทิโน, ปํสุกูลิกา เจว อเหสุํ ปํสุกูลิกตฺตสฺส จ วณฺณวาทิโน, เตจีวริกา เจว อเหสุํ เตจีวริกตฺตสฺส จ วณฺณวาทิโน, อปฺปิจฺฉา เจว อเหสุํ อปฺปิจฺฉตาย จ วณฺณวาทิโน, สนฺตุฏฺฐา เจว อเหสุํ สนฺตุฏฺฐิยา จ วณฺณวาทิโน, ปวิวิตฺตา เจว อเหสุํ ปวิเวกสฺส จ วณฺณวาทิโน, อสํสฏฺฐา เจว อเหสุํ อสํสคฺคสฺส จ วณฺณวาทิโน, อารทฺธวีริยา เจว อเหสุํ วีริยารมฺภสฺส จ วณฺณวาทิโน.}

\prose{ตตฺร โย โหติ ภิกฺขุ อารญฺญิโก เจว อารญฺญิกตฺตสฺส จ วณฺณวาที, ปิณฺฑปาติโก เจว ปิณฺฑปาติกตฺตสฺส จ วณฺณวาที, ปํสุกูลิโก เจว ปํสุกูลิกตฺตสฺส จ วณฺณวาที, เตจีวริโก เจว เตจีวริกตฺตสฺส จ วณฺณวาที, อปฺปิจฺโฉ เจว อปฺปิจฺฉตาย จ วณฺณวาที, สนฺตุฏฺโฐ เจว สนฺตุฏฺฐิยา จ วณฺณวาที, ปวิวิตฺโต เจว ปวิเวกสฺส จ วณฺณวาที, อสํสฏฺโฐ เจว อสํสคฺคสฺส จ วณฺณวาที, อารทฺธวีริโย เจว วีริยารมฺภสฺส จ วณฺณวาที, ตํ เถรา ภิกฺขู อาสเนน นิมนฺเตนฺติ “เอหิ ภิกฺขุ, โก นามายํ ภิกฺขุ, ภทฺทโก วตายํ ภิกฺขุ, สิกฺขากาโม วตายํ ภิกฺขุ, เอหิ ภิกฺขุ, อิทํ อาสนํ, นิสีทาหี”ติ.}

\prose{ตตฺร กสฺสป นวานํ ภิกฺขูนํ เอวํ โหติ “โย กิร โส โหติ ภิกฺขุ อารญฺญิโก เจว อารญฺญิกตฺตสฺส จ วณฺณวาที, ปิณฺฑปาติโก เจว ฯเปฯ ปํสุกูลิโก เจว.\csromanspacer{} เตจีวริโก เจว.\csromanspacer{} อปฺปิจฺโฉ เจว.\csromanspacer{} สนฺตุฏฺโฐ เจว.\csromanspacer{} ปวิวิตฺโต เจว.\csromanspacer{} อสํสฏฺโฐ เจว.\csromanspacer{} อารทฺธวีริโย เจว วีริยารมฺภสฺส จ วณฺณวาที, ตํ เถรา ภิกฺขู อาสเนน นิมนฺเตนฺติ ‘เอหิ ภิกฺขุ, โก นามายํ ภิกฺขุ, ภทฺทโก วตายํ ภิกฺขุ, สิกฺขากาโม วตายํ ภิกฺขุ, เอหิ ภิกฺขุ, อิทํ อาสนํ, นิสีทาหี’ติ”.\csromanspacer{} เต ตถตฺตาย ปฏิปชฺชนฺติ, เตสํ ตํ โหติ ทีฆรตฺตํ หิตาย สุขาย.}

\prose{เอตรหิ ปน กสฺสป เถรา ภิกฺขู น เจว อารญฺญิกา น จ อารญฺญิกตฺตสฺส วณฺณวาทิโน, น เจว ปิณฺฑปาติกา น จ ปิณฺฑปาติกตฺตสฺส วณฺณวาทิโน, น เจว ปํสุกูลิกา น จ ปํสุกูลิกตฺตสฺส วณฺณวาทิโน, น เจว เตจีวริกา น จ เตจีวริกตฺตสฺส วณฺณวาทิโน, น เจว อปฺปิจฺฉา น จ อปฺปิจฺฉตาย วณฺณวาทิโน, น เจว สนฺตุฏฺฐา น จ สนฺตุฏฺฐิยา วณฺณวาทิโน, น เจว ปวิวิตฺตา น จ ปวิเวกสฺส วณฺณวาทิโน, น เจว อสํสฏฺฐา น จ อสํสคฺคสฺส วณฺณวาทิโน, น เจว อารทฺธวีริยา น จ วีริยารมฺภสฺส วณฺณวาทิโน.}

\gambhira{นิทานวคฺคสํยุตฺตปาฬิ}
\prose{\csromanfolio{412}ตตฺร โย โหติ ภิกฺขุ ญาโต ยสสฺสี ลาภี จีวรปิณฺฑปาตเสนาสน- คิลานปฺปจฺจยเภสชฺชปริกฺขารานํ, ตํ เถรา ภิกฺขู อาสเนน นิมนฺเตนฺติ “เอหิ ภิกฺขุ, โก นามายํ ภิกฺขุ, ภทฺทโก วตายํ ภิกฺขุ, สพฺรหฺมจาริกาโม วตายํ ภิกฺขุ, เอหิ ภิกฺขุ, อิทํ อาสนํ, นิสีทาหี”ติ.}

\prose{ตตฺร กสฺสป นวานํ ภิกฺขูนํ เอวํ โหติ “โย กิร โส โหติ ภิกฺขุ ญาโต ยสสฺสี ลาภี จีวร ปิณฺฑปาต เสนาสน คิลานปฺปจฺจยเภสชฺช ปริกฺขารานํ, ตํ เถรา ภิกฺขู อาสเนน นิมนฺเตนฺติ ‘เอหิ ภิกฺขุ, โก นามายํ ภิกฺขุ, ภทฺทโก วตายํ ภิกฺขุ, สพฺรหฺมจาริกาโม วตายํ ภิกฺขุ, เอหิ ภิกฺขุ, อิทํ อาสนํ, นิสีทาหี’ติ”.\csromanspacer{} เต ตถตฺตาย ปฏิปชฺชนฺติ, เตสํ ตํ โหติ ทีฆรตฺตํ อหิตาย ทุกฺขาย.\csromanspacer{} ยํ หิ ตํ กสฺสป สมฺมา วทมาโน วเทยฺย อุปทฺทุตา พฺรหฺมจารี พฺรหฺมจารูปทฺทเวน อภิปตฺถนา\footnote{อภิภวนา (สี)} พฺรหฺมจารี พฺรหฺมจาริ- อภิปตฺถเนนาติ\footnote{พฺรหฺมจาริ-อภิภวเนนาติ (สี)}.\csromanspacer{} เอตรหิ ตํ กสฺสป สมฺมา วทมาโน วเทยฺย “อุปทฺทุตา พฺรหฺมจารี พฺรหฺมจารูปทฺทเวน อภิปตฺถนา พฺรหฺมจารี พฺรหฺมจาริ-อภิปตฺถเนนา”ติ.\csromanspacer{}\csromanspacer{}อฏฺฐมํ.}

\csromanheader{3}{9. ฌานาภิญฺญสุตฺต}
\proseitem{152}{สาวตฺถิยํ วิหรติ.\csromanspacer{} อหํ ภิกฺขเว ยาวเท อากงฺขามิ วิวิจฺเจว กาเมหิ วิวิจฺจ อกุสเลหิ ธมฺเมหิ สวิตกฺกํ สวิจารํ วิเวกชํ ปีติสุขํ ปฐมํ ฌานํ อุปสมฺปชฺช วิหรามิ.\csromanspacer{} กสฺสโปปิ ภิกฺขเว ยาวเท อากงฺขติ วิวิจฺเจว กาเมหิ วิวิจฺจ อกุสเลหิ ธมฺเมหิ สวิตกฺกํ สวิจารํ วิเวกชํ ปีติสุขํ ปฐมํ ฌานํ อุปสมฺปชฺช วิหรติ.}

\csromanmatikaanchor{mtk.187}
\csromantocmarkref{subsubsection}{6. โอวาทสุตฺต}{mtk.187}
\bookmark[dest={mtk.187},level=3]{6. โอวาทสุตฺต}
\prose{อหํ ภิกฺขเว ยาวเท อากงฺขามิ วิตกฺกวิจารานํ วูปสมา อชฺฌตฺตํ สมฺปสาทนํ เจตโส เอโกทิภาวํ อวิตกฺกํ อวิจารํ สมาธิชํ ปีติสุขํ ทุติยํ ฌานํ อุปสมฺปชฺช วิหรามิ.\csromanspacer{} กสฺสโปปิ ภิกฺขเว ยาวเท อากงฺขติ วิตกฺกวิจารานํ วูปสมา ฯเปฯ ทุติยํ ฌานํ อุปสมฺปชฺช วิหรติ.}

\prose{อหํ ภิกฺขเว ยาวเท อากงฺขามิ ปีติยา จ วิราคา อุเปกฺขโก จ วิหรามิ สโต จ สมฺปชาโน สุขญฺจ กาเยน ปฏิสํเวเทมิ, ยํ ตํ อริยา อาจิกฺขนฺติ “อุเปกฺขโก สติมา สุขวิหารี”ติ, ตติยํ ฌานํ อุปสมฺปชฺช วิหรามิ.\csromanspacer{} กสฺสโปปิ ภิกฺขเว ยาวเท อากงฺขติ}

\csromanheader{2}{5. กสฺสปสํยุตฺต}
\prose{\csromanfolio{413}ปีติยา จ วิราคา อุเปกฺขโก จ วิหรติ สโต จ สมฺปชาโน สุขญฺจ กาเยน ปฏิสํเวเทติ, ยํ ตํ อริยา อาจิกฺขนฺติ “อุเปกฺขโก สติมา สุขวิหารี”ติ, ตติยํ ฌานํ อุปสมฺปชฺช วิหรติ.}

\prose{อหํ ภิกฺขเว ยาวเท อากงฺขามิ สุขสฺส จ ปหานา ทุกฺขสฺส จ ปหานา ปุพฺเพว โสมนสฺสโทมนสฺสานํ อตฺถงฺคมา อทุกฺขมสุขํ อุเปกฺขาสติปาริสุทฺธิํ จตุตฺถํ ฌานํ อุปสมฺปชฺช วิหรามิ.\csromanspacer{} กสฺสโปปิ ภิกฺขเว ยาวเท อากงฺขติ สุขสฺส จ ปหานา ฯเปฯ จตุตฺถํ ฌานํ อุปสมฺปชฺช วิหรติ.}

\prose{อหํ ภิกฺขเว ยาวเท อากงฺขามิ สพฺพโส รูปสญฺญานํ สมติกฺกมา ปฏิฆสญฺญานํ อตฺถงฺคมา นานตฺตสญฺญานํ อมนสิการา “อนนฺโต อากาโส”ติ อากาสานญฺจายตนํ อุปสมฺปชฺช วิหรามิ.\csromanspacer{} กสฺสโปปิ ภิกฺขเว ยาวเท อากงฺขติ สพฺพโส รูปสญฺญานํ สมติกฺกมา ฯเปฯ อากาสานญฺจายตนํ อุปสมฺปชฺช วิหรติ.}

\prose{อหํ ภิกฺขเว ยาวเท อากงฺขามิ สพฺพโส อากาสานญฺจายตนํ สมติกฺกมฺม “อนนฺตํ วิญฺญาณนฺ”ติ วิญฺญาณญฺจายตนํ อุปสมฺปชฺช วิหรามิ.\csromanspacer{} กสฺสโปปิ ภิกฺขเว ยาวเท อากงฺขติ สพฺพโส อากาสานญฺจายตนํ สมติกฺกมฺม “อนนฺตํ วิญฺญาณนฺ”ติ วิญฺญาณญฺจายตนํ อุปสมฺปชฺช วิหรติ.}

\prose{อหํ ภิกฺขเว ยาวเท อากงฺขามิ สพฺพโส วิญฺญาณญฺจายตนํ สมติกฺกมฺม “นตฺถิ กิญฺจี”ติ อากิญฺจญฺญายตนํ อุปสมฺปชฺช วิหรามิ.\csromanspacer{} กสฺสโปปิ ภิกฺขเว ยาวเท อากงฺขติ ฯเปฯ อากิญฺจญฺญายตนํ อุปสมฺปชฺช วิหรติ.}

\prose{อหํ ภิกฺขเว ยาวเท อากงฺขามิ สพฺพโส อากิญฺจญฺญายตนํ สมติกฺกมฺม เนวสญฺญานาสญฺญายตนํ อุปสมฺปชฺช วิหรามิ.\csromanspacer{} กสฺสโปปิ ภิกฺขเว ยาวเท อากงฺขติ ฯเปฯ เนวสญฺญานาสญฺญายตนํ อุปสมฺปชฺช วิหรติ.}

\csromanmatikaanchor{mtk.188}
\csromantocmarkref{subsubsection}{7. ทุติย-โอวาทสุตฺต}{mtk.188}
\bookmark[dest={mtk.188},level=3]{7. ทุติย-โอวาทสุตฺต}
\prose{อหํ ภิกฺขเว ยาวเท อากงฺขามิ สพฺพโส เนวสญฺญานาสญฺญายตนํ สมติกฺกมฺม สญฺญาเวทยิตนิโรธํ อุปสมฺปชฺช วิหรามิ.\csromanspacer{} กสฺสโปปิ ภิกฺขเว ฯเปฯ สญฺญาเวทยิตนิโรธํ อุปสมฺปชฺช วิหรติ.}

\gambhira{นิทานวคฺคสํยุตฺตปาฬิ}
\prose{\csromanfolio{414}อหํ ภิกฺขเว ยาวเท อากงฺขามิ อเนกวิหิตํ อิทฺธิวิธํ ปจฺจนุโภมิ, เอโกปิ หุตฺวา พหุธา โหมิ, พหุธาปิ หุตฺวา เอโก โหมิ, อาวิภาวํ ติโรภาวํ ติโรกุฏฺฏํ ติโรปาการํ ติโรปพฺพตํ อสชฺชมาโน คจฺฉามิ เสยฺยถาปิ อากาเส, ปถวิยาปิ อุมฺมุชฺชนิมุชฺชํ กโรมิ เสยฺยถาปิ อุทเก, อุทเกปิ อภิชฺชมาเน คจฺฉามิ เสยฺยถาปิ ปถวิยํ, อากาเสปิ ปลฺลงฺเกน กมามิ เสยฺยถาปิ ปกฺขี สกุโณ, อิเมปิ จนฺทิมสูริเย เอวํมหิทฺธิเก เอวํมหานุภาเว ปาณินา ปริมสามิ ปริมชฺชามิ, ยาว พฺรหฺมโลกาปิ กาเยน วสํ วตฺเตมิ.\csromanspacer{} กสฺสโปปิ ภิกฺขเว ยาวเท อากงฺขติ อเนกวิหิตํ อิทฺธิวิธํ ปจฺจนุโภติ ฯเปฯ ยาว พฺรหฺมโลกาปิ กาเยน วสํ วตฺเตติ.}

\prose{อหํ ภิกฺขเว ยาวเท อากงฺขามิ ทิพฺพาย โสตธาตุยา วิสุทฺธาย อติกฺกนฺตมานุสิกาย อุโภ สทฺเท สุณามิ ทิพฺเพ จ มานุเส จ เย ทูเร สนฺติเก จ.\csromanspacer{} กสฺสโปปิ ภิกฺขเว ยาวเท อากงฺขติ ทิพฺพาย โสตธาตุยา ฯเปฯ ทูเร สนฺติเก จ.}

\prose{อหํ ภิกฺขเว ยาวเท อากงฺขามิ ปรสตฺตานํ ปรปุคฺคลานํ เจตสา เจโต ปริจฺจ ปชานามิ, สราคํ วา จิตฺตํ “สราคํ จิตฺตนฺ”ติ ปชานามิ, วีตราคํ วา จิตฺตํ “วีตราคํ จิตฺตนฺ”ติ ปชานามิ, สโทสํ วา จิตฺตํ ฯเปฯ วีตโทสํ วา จิตฺตํ.\csromanspacer{} สโมหํ วา จิตฺตํ.\csromanspacer{} วีตโมหํ วา จิตฺตํ.\csromanspacer{} สํขิตฺตํ วา จิตฺตํ.\csromanspacer{} วิกฺขิตฺตํ วา จิตฺตํ.\csromanspacer{} มหคฺคตํ วา จิตฺตํ.\csromanspacer{} อมหคฺคตํ วา จิตฺตํ.\csromanspacer{} ส-อุตฺตรํ วา จิตฺตํ.\csromanspacer{} อนุตฺตรํ วา จิตฺตํ.\csromanspacer{} สมาหิตํ วา จิตฺตํ.\csromanspacer{} อสมาหิตํ วา จิตฺตํ.\csromanspacer{} วิมุตฺตํ วา จิตฺตํ.\csromanspacer{} อวิมุตฺตํ วา จิตฺตํ “อวิมุตฺตํ จิตฺตนฺ”ติ ปชานามิ.\csromanspacer{} กสฺสโปปิ ภิกฺขเว ยาวเท อากงฺขติ ปรสตฺตานํ ปรปุคฺคลานํ เจตสา เจโต ปริจฺจ ปชานาติ, สราคํ วา จิตฺตํ “สราคํ จิตฺตนฺ”ติ ปชานาติ ฯเปฯ อวิมุตฺตํ วา จิตฺตํ “อวิมุตฺตํ จิตฺตนฺ”ติ ปชานาติ.}

\prose{อหํ ภิกฺขเว ยาวเท อากงฺขามิ อเนกวิหิตํ ปุพฺเพนิวาสํ อนุสฺสรามิ.\csromanspacer{} เสยฺยถิทํ, เอกมฺปิ ชาติํ เทฺวปิ ชาติโย ติสฺโสปิ ชาติโย จตสฺโสปิ ชาติโย ปญฺจปิ ชาติโย ทสปิ ชาติโย วีสมฺปิ ชาติโย ติํสมฺปิ ชาติโย จตฺตาลีสมฺปิ ชาติโย ปญฺญาสมฺปิ ชาติโย ชาติสตมฺปิ ชาติสหสฺสมฺปิ ชาติสตสหสฺสมฺปิ อเนเกปิ สํวฏฺฏกปฺเป อเนเกปิ วิวฏฺฏกปฺเป อเนเกปิ สํวฏฺฏวิวฏฺฏกปฺเป “อมุตฺราสิํ เอวํนาโม เอวํโคตฺโต เอวํวณฺโณ เอวมาหาโร เอวํสุขทุกฺขปฺปฏิสํเวที}

\csromanheader{2}{5. กสฺสปสํยุตฺต}
\prose{\csromanfolio{415}เอวมายุปริยนฺโต, โส ตโต จุโต อมุตฺร อุทปาทิํ, ตตฺราปาสิํ เอวํนาโม เอวํโคตฺโต เอวํวณฺโณ เอวมาหาโร เอวํสุขทุกฺขปฺปฏิสํเวที เอวมายุปริยนฺโต, โส ตโต จุโต อิธูปปนฺโน”ติ, อิติ สาการํ ส-อุทฺเทสํ อเนกวิหิตํ ปุพฺเพนิวาสํ อนุสฺสรามิ.\csromanspacer{} กสฺสโปปิ ภิกฺขเว ยาวเท อากงฺขติ อเนกวิหิตํ ปุพฺเพนิวาสํ อนุสฺสรติ.\csromanspacer{} เสยฺยถิทํ, เอกมฺปิ ชาติํ ฯเปฯ อิติ สาการํ ส-อุทฺเทสํ อเนกวิหิตํ ปุพฺเพนิวาสํ อนุสฺสรติ.}

\prose{อหํ ภิกฺขเว ยาวเท อากงฺขามิ ทิพฺเพน จกฺขุนา วิสุทฺเธน อติกฺกนฺตมานุสเกน สตฺเต ปสฺสามิ จวมาเน อุปปชฺชมาเน หีเน ปณีเต สุวณฺเณ ทุพฺพณฺเณ สุคเต ทุคฺคเต, ยถากมฺมูปเค สตฺเต ปชานามิ “อิเม วต โภนฺโต สตฺตา กายทุจฺจริเตน สมนฺนาคตา วจีทุจฺจริเตน สมนฺนาคตา มโนทุจฺจริเตน สมนฺนาคตา อริยานํ อุปวาทกา มิจฺฉาทิฏฺฐิกา มิจฺฉาทิฏฺฐิกมฺมสมาทานา, เต กายสฺส เภทา ปรํ มรณา อปายํ ทุคฺคติํ วินิปาตํ นิรยํ อุปปนฺนา.\csromanspacer{} อิเม วา ปน โภนฺโต สตฺตา กายสุจริเตน สมนฺนาคตา วจีสุจริเตน สมนฺนาคตา มโนสุจริเตน สมนฺนาคตา อริยานํ อนุปวาทกา สมฺมาทิฏฺฐิกา สมฺมาทิฏฺฐิกมฺมสมาทานา, เต กายสฺส เภทา ปรํ มรณา สุคติํ สคฺคํ โลกํ อุปปนฺนา”ติ.\csromanspacer{} อิติ ทิพฺเพน จกฺขุนา วิสุทฺเธน อติกฺกนฺตมานุสเกน สตฺเต ปสฺสามิ จวมาเน อุปปชฺชมาเน หีเน ปณีเต สุวณฺเณ ทุพฺพณฺเณ สุคเต ทุคฺคเต, ยถากมฺมูปเค สตฺเต ปชานามิ.\csromanspacer{} กสฺสโปปิ ภิกฺขเว ยาวเท อากงฺขติ ทิพฺเพน จกฺขุนา วิสุทฺเธน อติกฺกนฺตมานุสเกน สตฺเต ปสฺสติ จวมาเน ฯเปฯ ยถากมฺมูปเค สตฺเต ปชานาติ.}

\prose{อหํ ภิกฺขเว อาสวานํ ขยา อนาสวํ เจโตวิมุตฺติํ ปญฺญาวิมุตฺติํ ทิฏฺเฐว ธมฺเม สยํ อภิญฺญา สจฺฉิกตฺวา อุปสมฺปชฺช วิหรามิ.\csromanspacer{} กสฺสโปปิ ภิกฺขเว อาสวานํ ขยา อนาสวํ เจโตวิมุตฺติํ ปญฺญาวิมุตฺติํ ทิฏฺเฐว ธมฺเม สยํ อภิญฺญา สจฺฉิกตฺวา อุปสมฺปชฺช วิหรตีติ.\csromanspacer{}\csromanspacer{}นวมํ.}

\csromanheader{3}{10. อุปสฺสยสุตฺต}
\proseitem{153}{เอวํ เม สุตํ\csromandash{}เอกํ สมยํ อายสฺมา มหากสฺสโป สาวตฺถิยํ วิหรติ เชตวเน อนาถปิณฺฑิกสฺส อาราเม.\csromanspacer{} อถ โข อายสฺมา อานนฺโท ปุพฺพณฺหสมยํ นิวาเสตฺวา ปตฺตจีวรมาทาย}

\vspace{\tipitakaparskip}
\csromancenter{นิทานวคฺคสํยุตฺตปาฬิ เยนายสฺมา มหากสฺสโป เตนุปสงฺกมิ, อุปสงฺกมิตฺวา อายสฺมนฺตํ มหากสฺสปํ เอตทโวจ “อายาม ภนฺเต กสฺสป เยน อญฺญตโร ภิกฺขุนุปสฺสโย เตนุปสงฺกมิสฺสามา”ติ.\csromanspacer{} คจฺฉ ตฺวํ อาวุโส อานนฺท พหุกิจฺโจ ตฺวํ พหุกรณีโยติ.\csromanspacer{} ทุติยมฺปิ โข อายสฺมา อานนฺโท อายสฺมนฺตํ มหากสฺสปํ เอตทโวจ “อายาม ภนฺเต กสฺสป เยน อญฺญตโร ภิกฺขุนุปสฺสโย เตนุปสงฺกมิสฺสามา”ติ.\csromanspacer{} คจฺฉ ตฺวํ อาวุโส อานนฺท พหุกิจฺโจ ตฺวํ พหุกรณีโยติ.\csromanspacer{} ตติยมฺปิ โข อายสฺมา อานนฺโท อายสฺมนฺตํ มหากสฺสปํ เอตทโวจ “อายาม ภนฺเต กสฺสป เยน อญฺญตโร ภิกฺขุนุปสฺสโย เตนุปสงฺกมิสฺสามา”ติ.}
\prose{\csromanfolio{416}อถ โข อายสฺมา มหากสฺสโป ปุพฺพณฺหสมยํ นิวาเสตฺวา ปตฺตจีวรมาทาย อายสฺมตา อานนฺเทน ปจฺฉาสมเณน เยน อญฺญตโร ภิกฺขุนุปสฺสโย เตนุปสงฺกมิ, อุปสงฺกมิตฺวา ปญฺญตฺเต อาสเน นิสีทิ.\csromanspacer{} อถ โข สมฺพหุลา ภิกฺขุนิโย เยนายสฺมา มหากสฺสโป เตนุปสงฺกมิํสุ, อุปสงฺกมิตฺวา อายสฺมนฺตํ มหากสฺสปํ อภิวาเทตฺวา เอกมนฺตํ นิสีทิํสุ, เอกมนฺตํ นิสินฺนา โข ตา ภิกฺขุนิโย อายสฺมา มหากสฺสโป ธมฺมิยา กถาย สนฺทสฺเสสิ สมาทเปสิ สมุตฺเตเชสิ สมฺปหํเสสิ.\csromanspacer{} อถ โข อายสฺมา มหากสฺสโป ตา ภิกฺขุนิโย ธมฺมิยา กถาย สนฺทสฺเสตฺวา สมาทเปตฺวา สมุตฺเตเชตฺวา สมฺปหํเสตฺวา อุฏฺฐายาสนา ปกฺกามิ.}

\prose{อถ โข ถุลฺลติสฺสา ภิกฺขุนี อนตฺตมนา อนตฺตมนวาจํ นิจฺฉาเรสิ “กิํ ปน อยฺโย มหากสฺสโป อยฺยสฺส อานนฺทสฺส เวเทหมุนิโน สมฺมุขา ธมฺมํ ภาสิตพฺพํ มญฺญติ.\csromanspacer{} เสยฺยถาปิ นาม สูจิวาณิชโก สูจิการสฺส สนฺติเก สูจิํ วิกฺเกตพฺพํ มญฺเญยฺย.\csromanspacer{} เอวเมว อยฺโย มหากสฺสโป อยฺยสฺส อานนฺทสฺส เวเทหมุนิโน สมฺมุขา ธมฺมํ ภาสิตพฺพํ มญฺญตี”ติ.}

\prose{อสฺโสสิ โข อายสฺมา มหากสฺสโป ถุลฺลติสฺสาย ภิกฺขุนิยา อิมํ วาจํ ภาสมานาย.\csromanspacer{} อถ โข อายสฺมา มหากสฺสโป อายสฺมนฺตํ อานนฺทํ เอตทโวจ “กิํ นุ โข อาวุโส อานนฺท อหํ สูจิวาณิชโก ตฺวํ สูจิกาโร, อุทาหุ อหํ สูจิกาโร ตฺวํ สูจิวาณิชโก”ติ.}

\csromanmatikaanchor{mtk.189}
\csromantocmarkref{subsubsection}{8. ตติย-โอวาทสุตฺต}{mtk.189}
\bookmark[dest={mtk.189},level=3]{8. ตติย-โอวาทสุตฺต}
\csromanheader{2}{5. กสฺสปสํยุตฺต}
\prose{\csromanfolio{417}ขม ภนฺเต กสฺสป, พาโล มาตุคาโมติ.\csromanspacer{} อาคเมติ ตฺวํ อาวุโส อานนฺท, มา เต สํโฆ อุตฺตริ อุปปริกฺขิ.}

\prose{ตํ กิํ มญฺญสิ อาวุโส อานนฺท, อปิ นุ ตฺวํ ภควโต สมฺมุขา ภิกฺขุสํเฆ อุปนีโต “อหํ ภิกฺขเว ยาวเท อากงฺขามิ วิวิจฺเจว กาเมหิ วิวิจฺจ อกุสเลหิ ธมฺเมหิ สวิตกฺกํ สวิจารํ วิเวกชํ ปีติสุขํ ปฐมํ ฌานํ อุปสมฺปชฺช วิหรามิ.\csromanspacer{} อานนฺโทปิ ภิกฺขเว ยาวเท อากงฺขติ วิวิจฺเจว กาเมหิ วิวิจฺจ อกุสเลหิ ธมฺเมหิ สวิตกฺกํ สวิจารํ วิเวกชํ ปีติสุขํ ปฐมํ ฌานํ อุปสมฺปชฺช วิหรตี”ติ.\csromanspacer{} โน เหตํ ภนฺเต.}

\prose{อหํ โข อาวุโส ภควโต สมฺมุขา ภิกฺขุสํเฆ อุปนีโต “อหํ ภิกฺขเว ยาวเท อากงฺขามิ วิวิจฺเจว กาเมหิ วิวิจฺจ อกุสเลหิ ธมฺเมหิ สวิตกฺกํ สวิจารํ วิเวกชํ ปีติสุขํ ปฐมํ ฌานํ อุปสมฺปชฺช วิหรามิ.\csromanspacer{} กสฺสโปปิ ภิกฺขเว ยาวเท อากงฺขติ วิวิจฺเจว กาเมหิ วิวิจฺจ อกุสเลหิ ธมฺเมหิ ฯเปฯ ปฐมํ ฌานํ อุปสมฺปชฺช วิหรตี”ติ ฯเปฯ. (นวนฺนํ อนุปุพฺพวิหารสมาปตฺตีนํ ปญฺจนฺนญฺจ อภิญฺญานํ เอวํ วิตฺถาโร เวทิตพฺโพ.)}

\prose{ตํ กิํ มญฺญสิ อาวุโส อานนฺท, อปิ นุ ตฺวํ ภควโต สมฺมุขา ภิกฺขุสํเฆ อุปนีโต “อหํ ภิกฺขเว อาสวานํ ขยา อนาสวํ เจโตวิมุตฺติํ ปญฺญาวิมุตฺติํ ทิฏฺเฐว ธมฺเม สยํ อภิญฺญา สจฺฉิกตฺวา อุปสมฺปชฺช วิหรามิ.\csromanspacer{} อานนฺโทปิ ภิกฺขเว อาสวานํ ขยา อนาสวํ เจโตวิมุตฺติํ ปญฺญาวิมุตฺติํ ทิฏฺเฐว ธมฺเม สยํ อภิญฺญา สจฺฉิกตฺวา อุปสมฺปชฺช วิหรตี”ติ.\csromanspacer{} โน เหตํ ภนฺเต.}

\prose{อหํ โข อาวุโส ภควโต สมฺมุขา ภิกฺขุสํเฆ อุปนีโต “อหํ ภิกฺขเว อาสวานํ ขยา อนาสวํ เจโตวิมุตฺติํ ปญฺญาวิมุตฺติํ ทิฏฺเฐว ธมฺเม สยํ อภิญฺญา สจฺฉิกตฺวา อุปสมฺปชฺช วิหรามิ.\csromanspacer{} กสฺสโปปิ ภิกฺขเว อาสวานํ ขยา อนาสวํ เจโตวิมุตฺติํ ปญฺญาวิมุตฺติํ ทิฏฺเฐว ธมฺเม สยํ อภิญฺญา สจฺฉิกตฺวา อุปสมฺปชฺช วิหรตี”ติ.}

\prose{สตฺตรตนํ วา อาวุโส นาคํ อฑฺฒฏฺฐมรตนํ วา ตาลปตฺติกาย ฉาเทตพฺพํ มญฺเญยฺย, โย เม ฉ อภิญฺญา ฉาเทตพฺพํ มญฺเญยฺยาติ.}

\prose{จวิตฺถ จ ปน ถุลฺลติสฺสา ภิกฺขุนี พฺรหฺมจริยมฺหาติ.\csromanspacer{}\csromanspacer{}ทสมํ.}

\csromanheader{3}{นิทานวคฺคสํยุตฺตปาฬิ \textbf{11. จีวรสุตฺต}}
\proseitem{154}{\csromanfolio{418}เอกํ สมยํ อายสฺมา มหากสฺสโป ราชคเห วิหรติ เวฬุวเน กลนฺทกนิวาเป.\csromanspacer{} เตน โข ปน สมเยน อายสฺมา อานนฺโท ทกฺขิณคิริสฺมิํ จาริกํ จรติ มหตา ภิกฺขุสํเฆน สทฺธิํ. เตน โข ปน สมเยน อายสฺมโต อานนฺทสฺส ติํสมตฺตา สทฺธิวิหาริโน ภิกฺขู สิกฺขํ ปจฺจกฺขาย หีนายวตฺตา ภวนฺติ เยภุยฺเยน กุมารภูตา.\csromanspacer{} อถ โข อายสฺมา อานนฺโท ทกฺขิณคิริสฺมิํ ยถาภิรนฺตํ จาริกํ จริตฺวา เยน ราชคหํ เวฬุวนํ กลนฺทกนิวาโป เยนายสฺมา มหากสฺสโป เตนุปสงฺกมิ, อุปสงฺกมิตฺวา อายสฺมนฺตํ มหากสฺสปํ อภิวาเทตฺวา เอกมนฺตํ นิสีทิ, เอกมนฺตํ นิสินฺนํ โข อายสฺมนฺตํ อานนฺทํ อายสฺมา มหากสฺสโป เอตทโวจ “กติ นุ โข อาวุโส อานนฺท อตฺถวเส ปฏิจฺจ ภควตา กุเลสุ ติกโภชนํ ปญฺญตฺตนฺ”ติ.}

\csromanmatikaanchor{mtk.190}
\csromantocmarkref{subsubsection}{9. ฌานาภิญฺญสุตฺต}{mtk.190}
\bookmark[dest={mtk.190},level=3]{9. ฌานาภิญฺญสุตฺต}
\prose{ตโย โข ภนฺเต กสฺสป อตฺถวเส ปฏิจฺจ ภควตา กุเลสุ ติกโภชนํ ปญฺญตฺตํ.\csromanspacer{} ทุมฺมงฺกูนํ ปุคฺคลานํ นิคฺคหาย เปสลานํ ภิกฺขูนํ ผาสุวิหาราย, มา ปาปิจฺฉา ปกฺขํ นิสฺสาย สํฆํ ภินฺเทยฺยุํ\footnote{วิ 4 สํฆเภทกกฺขนฺธเก วชิรพุทฺธิยํ อญฺญถา สมฺพนฺโธ ทสฺสิโต.}, กุลานุทฺทยตาย จ.\csromanspacer{} อิเม โข ภนฺเต กสฺสป ตโย อตฺถวเส ปฏิจฺจ ภควตา กุเลสุ ติกโภชนํ ปญฺญตฺตนฺติ.}

\prose{อถ กิญฺจรหิ ตฺวํ อาวุโส อานนฺท อิเมหิ นเวหิ ภิกฺขูหิ อินฺทฺริเยสุ อคุตฺตทฺวาเรหิ โภชเน อมตฺตญฺญูหิ ชาคริยํ อนนุยุตฺเตหิ สทฺธิํ จาริกํ จรสิ, สสฺสฆาตํ มญฺเญ จรสิ, กุลูปฆาตํ มญฺเญ จรสิ, โอลุชฺชติ\footnote{อุลฺลุชฺชติ (สี, อฏฺฐกถาสุ จ)} โข เต อาวุโส อานนฺท ปริสา, ปลุชฺชนฺติ โข เต อาวุโส นวปฺปายา.\csromanspacer{} น วายํ กุมารโก มตฺตมญฺญาสีติ.}

\prose{อปิ เม ภนฺเต กสฺสป สิรสฺมิํ ปลิตานิ ชาตานิ, อถ จ ปน มยํ อชฺชาปิ อายสฺมโต มหากสฺสปสฺส กุมารกวาทา น มุจฺจามาติ.\csromanspacer{} ตถา หิ ปน ตฺวํ อาวุโส อานนฺท อิเมหิ นเวหิ ภิกฺขูหิ อินฺทฺริเยสุ อคุตฺตทฺวาเรหิ โภชเน อมตฺตญฺญูหิ ชาคริยํ อนนุยุตฺเตหิ สทฺธิํ จาริกํ จรสิ, สสฺสฆาตํ มญฺเญ จรสิ, กุลูปฆาตํ มญฺเญ จรสิ,}

\csromanheader{2}{5. กสฺสปสํยุตฺต}
\prose{\csromanfolio{419}โอลุชฺชติ โข เต อาวุโส อานนฺท ปริสา, ปลุชฺชนฺติ โข เต อาวุโส\footnote{ปลุชฺชติ โข เต อาวุโส อานนฺท ปริสา (ก-สี)} นวปฺปายา, น วายํ กุมารโก มตฺตมญฺญาสีติ.}

\prose{อสฺโสสิ โข ถุลฺลนนฺทา ภิกฺขุนี “อยฺเยน กิร มหากสฺสเปน อยฺโย อานนฺโท เวเทหมุนิ กุมารกวาเทน อปสาทิโต”ติ.}

\prose{อถ โข ถุลฺลนนฺทา ภิกฺขุนี อนตฺตมนา อนตฺตมนวาจํ นิจฺฉาเรสิ “กิํ ปน อยฺโย มหากสฺสโป อญฺญติตฺถิยปุพฺโพ สมาโน อยฺยํ อานนฺทํ เวเทหิมุนิํ กุมารกวาเทน อปสาเทตพฺพํ มญฺญตี”ติ.\csromanspacer{} อสฺโสสิ โข อายสฺมา มหากสฺสโป ถุลฺลนนฺทาย ภิกฺขุนิยา อิมํ วาจํ ภาสมานาย.}

\prose{อถ โข อายสฺมา มหากสฺสโป อายสฺมนฺตํ อานนฺทํ เอตทโวจ\csromandash{} ตคฺฆาวุโส อานนฺท ถุลฺลนนฺทาย ภิกฺขุนิยา สหสา อปฺปฏิสงฺขา วาจา ภาสิตา, ยตฺวาหํ อาวุโส เกสมสฺสุํ โอหาเรตฺวา กาสายานิ วตฺถานิ อจฺฉาเทตฺวา อคารสฺมา อนคาริยํ ปพฺพชิโต, นาภิชานามิ อญฺญํ สตฺถารํ อุทฺทิสิตา\footnote{อุทฺทิสิตุํ (สี, อิ, ก)} อญฺญตฺร เตน ภควตา อรหตา สมฺมาสมฺพุทฺเธน.\csromanspacer{} ปุพฺเพ เม อาวุโส อคาริกภูตสฺส สโต เอตทโหสิ “สมฺพาโธ ฆราวาโส รชาปโถ\footnote{รโชปโถ (สี)} อพฺโภกาโส ปพฺพชฺชา.\csromanspacer{} นยิทํ สุกรํ อคารํ อชฺฌาวสตา เอกนฺตปริปุณฺณํ เอกนฺตปริสุทฺธํ สงฺขลิขิตํ พฺรหฺมจริยํ จริตุํ, ยํนูนาหํ เกสมสฺสุํ โอหาเรตฺวา กาสายานิ วตฺถานิ อจฺฉาเทตฺวา อคารสฺมา อนคาริยํ ปพฺพเชยฺยนฺ”ติ.\csromanspacer{} โส ขฺวาหํ อาวุโส อปเรน สมเยน ปฏปิโลติกานํ สงฺฆาฏิํ กาเรตฺวา\footnote{กริตฺวา (สี, สฺยา, กํ, อิ)} เย โลเก อรหนฺโต, เต อุทฺทิสฺส เกสมสฺสุํ โอหาเรตฺวา กาสายานิ วตฺถานิ อจฺฉาเทตฺวา อคารสฺมา อนคาริยํ ปพฺพชิํ.}

\prose{โส เอวํ ปพฺพชิโต สมาโน อทฺธานมคฺคปฺปฏิปนฺโน อทฺทสํ ภควนฺตํ อนฺตรา จ ราชคหํ อนฺตรา จ นาฬนฺทํ พหุปุตฺเต เจติเย นิสินฺนํ, ทิสฺวาน เม เอตทโหสิ “สตฺถารญฺจ วตาหํ ปสฺเสยฺยํ ภควนฺตเมว ปสฺเสยฺยํ, สุคตญฺจ วตาหํ ปสฺเสยฺยํ ภควนฺตเมว ปสฺเสยฺยํ, สมฺมาสมฺพุทฺธญฺจ วตาหํ ปสฺเสยฺยํ ภควนฺตเมว ปสฺเสยฺยนฺ”ติ.\csromanspacer{} โส ขฺวาหํ อาวุโส ตตฺเถว ภควโต ปาเทสุ สิรสา นิปติตฺวา ภควนฺตํ เอตทโวจํ “สตฺถา เม ภนฺเต ภควา, สาวโกหมสฺมิ, สตฺถา เม ภนฺเต ภควา,}

\vspace{\tipitakaparskip}
\csromancenter{นิทานวคฺคสํยุตฺตปาฬิ สาวโกหมสฺมี”ติ.\csromanspacer{} เอวํ วุตฺเต มํ อาวุโส ภควา เอตทโวจ “โย โข กสฺสป เอวํ สพฺพเจตสา สมนฺนาคตํ สาวกํ อชานญฺเญว วเทยฺย ‘ชานามี’ติ, อปสฺสญฺเญว วเทยฺย ‘ปสฺสามี’ติ, มุทฺธาปิ ตสฺส วิปเตยฺย.\csromanspacer{} อหํ โข ปน กสฺสป ชานญฺเญว วทามิ ‘ชานามี’ติ, ปสฺสญฺเญว วทามิ ‘ปสฺสามี’ติ”.}
\prose{\csromanfolio{420}ตสฺมาติห เต กสฺสป เอวํ สิกฺขิตพฺพํ “ติพฺพํ เม หิโรตฺตปฺปํ, ปจฺจุปฏฺฐิตํ ภวิสฺสติ เถเรสุ นเวสุ มชฺฌิเมสู”ติ, เอวํ หิ เต กสฺสป สิกฺขิตพฺพํ.}

\prose{ตสฺมาติห เต กสฺสป เอวํ สิกฺขิตพฺพํ “ยํ กิญฺจิ ธมฺมํ สุณิสฺสามิ กุสลูปสํหิตํ, สพฺพํ ตํ อฏฺฐิํ กตฺวา มนสิ กริตฺวา สพฺพเจตสา สมนฺนาหริตฺวา โอหิตโสโต ธมฺมํ สุณิสฺสามี”ติ, เอวํ หิ เต กสฺสป สิกฺขิตพฺพํ.}

\prose{ตสฺมาติห เต กสฺสป เอวํ สิกฺขิตพฺพํ “สาตสหคตา จ เม กายคตาสติ น วิชหิสฺสตี”ติ, เอวํ หิ เต กสฺสป สิกฺขิตพฺพนฺติ.}

\prose{อถ โข มํ อาวุโส ภควา อิมินา โอวาเทน โอวทิตฺวา อุฏฺฐายาสนา ปกฺกามิ, สตฺตาหเมว ขฺวาหํ อาวุโส สรโณ\footnote{สาโณ (สี)} รฏฺฐปิณฺฑํ ภุญฺชิํ อฏฺฐมิยา อญฺญา อุทปาทิ.}

\prose{อถ โข อาวุโส ภควา มคฺคา โอกฺกมฺม เยน อญฺญตรํ รุกฺขมูลํ เตนุปสงฺกมิ, อถ ขฺวาหํ อาวุโส ปฏปิโลติกานํ สงฺฆาฏิํ จตุคฺคุณํ ปญฺญาเปตฺวา ภควนฺตํ เอตทโวจํ “อิธ ภนฺเต ภควา นิสีทตุ, ยํ มมสฺส ทีฆรตฺตํ หิตาย สุขายา”ติ.\csromanspacer{} นิสีทิ โข อาวุโส ภควา ปญฺญตฺเต อาสเน, นิสชฺช โข มํ อาวุโส ภควา เอตทโวจ “มุทุกา โข ตฺยายํ กสฺสป ปฏปิโลติกานํ สงฺฆาฏี”ติ.\csromanspacer{} ปฏิคฺคณฺหาตุ เม ภนฺเต ภควา ปฏปิโลติกานํ สงฺฆาฏิํ อนุกมฺปํ อุปาทายาติ.\csromanspacer{} ธาเรสฺสสิ ปน เม ตฺวํ กสฺสป สาณานิ ปํสุกูลานิ นิพฺพสนานีติ.\csromanspacer{} ธาเรสฺสามหํ ภนฺเต ภควโต สาณานิ ปํสุกูลานิ นิพฺพสนานีติ.\csromanspacer{} โส ขฺวาหํ อาวุโส ปฏปิโลติกานํ สงฺฆาฏิํ ภควโต ปาทาสิํ, อหํ ปน ภควโต สาณานิ ปํสุกูลานิ นิพฺพสนานิ ปฏิปชฺชิํ.}

\prose{ยํ หิ ตํ อาวุโส สมฺมา วทมาโน วเทยฺย “ภควโต ปุตฺโต โอรโส มุขโต ชาโต ธมฺมโช ธมฺมนิมฺมิโต ธมฺมทายาโท,}

\csromanheader{2}{5. กสฺสปสํยุตฺต}
\prose{\csromanfolio{421}ปฏิคฺคหิตานิ\footnote{ปฏิคฺคเหตา (สี)} สาณานิ ปํสุกูลานิ นิพฺพสนานี”ติ.\csromanspacer{} มมํ ตํ สมฺมา วทมาโน วเทยฺย “ภควโต ปุตฺโต โอรโส มุขโต ชาโต ธมฺมโช ธมฺมนิมฺมิโต ธมฺมทายาโท, ปฏิคฺคหิตานิ สาณานิ ปํสุกูลานิ นิพฺพสนานี”ติ.}

\prose{อหํ โข อาวุโส ยาวเท อากงฺขามิ วิวิจฺเจว กาเมหิ วิวิจฺจ อกุสเลหิ ธมฺเมหิ สวิตกฺกํ สวิจารํ วิเวกชํ ปีติสุขํ ปฐมํ ฌานํ อุปสมฺปชฺช วิหรามิ, อหํ โข อาวุโส ยาวเท อากงฺขามิ ฯเปฯ. (นวนฺนํ อนุปุพฺพวิหารสมาปตฺตีนํ ปญฺจนฺนญฺจ อภิญฺญานํ เอวํ วิตฺถาโร เวทิตพฺโพ.)}

\prose{อหํ โข อาวุโส อาสวานํ ขยา อนาสวํ เจโตวิมุตฺติํ ปญฺญาวิมุตฺติํ ทิฏฺเฐว ธมฺเม สยํ อภิญฺญา สจฺฉิกตฺวา อุปสมฺปชฺช วิหรามิ, สตฺตรตนํ วา อาวุโส นาคํ อฑฺฒฏฺฐมรตนํ วา ตาลปตฺติกาย ฉาเทตพฺพํ มญฺเญยฺย, โย เม ฉ อภิญฺญา ฉาเทตพฺพํ มญฺเญยฺยาติ.}

\prose{จวิตฺถ จ ปน ถุลฺลนนฺทา ภิกฺขุนี พฺรหฺมจริยมฺหาติ.\csromanspacer{}\csromanspacer{}เอกาทสมํ.}

\csromanmatikaanchor{mtk.191}
\csromantocmarkref{subsubsection}{10. อุปสฺสยสุตฺต}{mtk.191}
\bookmark[dest={mtk.191},level=3]{10. อุปสฺสยสุตฺต}
\csromanheader{3}{12. ปรํมรณสุตฺต}
\proseitem{155}{เอกํ สมยํ อายสฺมา จ มหากสฺสโป อายสฺมา จ สาริปุตฺโต พาราณสิยํ วิหรนฺติ อิสิปตเน มิคทาเย.\csromanspacer{} อถ โข อายสฺมา สาริปุตฺโต สายนฺหสมยํ ปฏิสลฺลานา วุฏฺฐิโต เยนายสฺมา มหากสฺสโป เตนุปสงฺกมิ, อุปสงฺกมิตฺวา อายสฺมตา มหากสฺสเปน สทฺธิํ สมฺโมทิ, สมฺโมทนียํ กถํ สารณียํ วีติสาเรตฺวา เอกมนฺตํ นิสีทิ, เอกมนฺตํ นิสินฺโน โข อายสฺมา สาริปุตฺโต อายสฺมนฺตํ มหากสฺสปํ เอตทโวจ “กิํ นุ โข อาวุโส กสฺสป โหติ ตถาคโต ปรํ มรณา”ติ.\csromanspacer{} อพฺยากตํ โข เอตํ อาวุโส ภควตา “โหติ ตถาคโต ปรํ มรณา”ติ.\csromanspacer{} กิํ ปนาวุโส น โหติ ตถาคโต ปรํ มรณาติ.\csromanspacer{} เอตมฺปิ โข อาวุโส อพฺยากตํ ภควตา “น โหติ ตถาคโต ปรํ มรณา”ติ.\csromanspacer{} กิํ นุ โข อาวุโส โหติ จ น จ โหติ ตถาคโต ปรํ มรณาติ.\csromanspacer{} อพฺยากตํ โข เอตํ อาวุโส ภควตา “โหติ จ น จ โหติ ตถาคโต ปรํ มรณา”ติ.\csromanspacer{} กิํ ปนาวุโส เนว โหติ น น โหติ ตถาคโต ปรํ มรณาติ.\csromanspacer{} เอตมฺปิ โข อาวุโส อพฺยากตํ ภควตา “เนว โหติ น น โหติ}

\vspace{\tipitakaparskip}
\csromancenter{นิทานวคฺคสํยุตฺตปาฬิ ตถาคโต ปรํ มรณา”ติ.\csromanspacer{} กสฺมา เจตํ อาวุโส อพฺยากตํ ภควตาติ.\csromanspacer{} น เหตํ อาวุโส อตฺถสํหิตํ นาทิพฺรหฺมจริยกํ น นิพฺพิทาย น วิราคาย น นิโรธาย น อุปสมาย น อภิญฺญาย น สมฺโพธาย น นิพฺพานาย สํวตฺตติ, ตสฺมา ตํ อพฺยากตํ ภควตาติ.}
\prose{\csromanfolio{422}อถ กิญฺจรหาวุโส พฺยากตํ ภควตาติ. “อิทํ ทุกฺขนฺ”ติ โข อาวุโส พฺยากตํ ภควตา, “อยํ ทุกฺขสมุทโย”ติ พฺยากตํ ภควตา, “อยํ ทุกฺขนิโรโธ”ติ พฺยากตํ ภควตา, “อยํ ทุกฺขนิโรธคามินี ปฏิปทา”ติ พฺยากตํ ภควตาติ.\csromanspacer{} กสฺมา เจตํ อาวุโส พฺยากตํ ภควตาติ.\csromanspacer{} เอตํ หิ อาวุโส อตฺถสํหิตํ เอตํ อาทิพฺรหฺมจริยกํ เอตํ นิพฺพิทาย วิราคาย นิโรธาย อุปสมาย อภิญฺญาย สมฺโพธาย นิพฺพานาย สํวตฺตติ, ตสฺมา ตํ พฺยากตํ ภควตาติ.\csromanspacer{}\csromanspacer{}ทฺวาทสมํ.}

\csromanheader{3}{13. สทฺธมฺมปฺปติรูปกสุตฺต}
\proseitem{156}{เอวํ เม สุตํ\csromandash{}เอกํ สมยํ ภควา สาวตฺถิยํ วิหรติ เชตวเน อนาถปิณฺฑิกสฺส อาราเม.\csromanspacer{} อถ โข อายสฺมา มหากสฺสโป เยน ภควา เตนุปสงฺกมิ, อุปสงฺกมิตฺวา ภควนฺตํ อภิวาเทตฺวา เอกมนฺตํ นิสีทิ, เอกมนฺตํ นิสินฺโน โข อายสฺมา มหากสฺสโป ภควนฺตํ เอตทโวจ “โก นุ โข ภนฺเต เหตุ โก ปจฺจโย, เยน ปุพฺเพ อปฺปตรานิ เจว สิกฺขาปทานิ อเหสุํ, พหุตรา จ ภิกฺขู อญฺญาย สณฺฐหิํสุ.\csromanspacer{} โก ปน ภนฺเต เหตุ โก ปจฺจโย, เยเนตรหิ พหุตรานิ เจว สิกฺขาปทานิ อปฺปตรา จ ภิกฺขู อญฺญาย สณฺฐหนฺตี”ติ.\csromanspacer{} เอวญฺเจตํ กสฺสป โหติ สตฺเตสุ หายมาเนสุ สทฺธมฺเม อนฺตรธายมาเน พหุตรานิ เจว สิกฺขาปทานิ โหนฺติ, อปฺปตรา จ ภิกฺขู อญฺญาย สณฺฐหนฺติ.\csromanspacer{} น ตาว กสฺสป สทฺธมฺมสฺส อนฺตรธานํ โหติ, ยาว น สทฺธมฺมปฺปติรูปกํ โลเก อุปฺปชฺชติ.\csromanspacer{} ยโต จ โข กสฺสป สทฺธมฺมปฺปติรูปกํ โลเก อุปฺปชฺชติ, อถ สทฺธมฺมสฺส อนฺตรธานํ โหติ.}

\prose{เสยฺยถาปิ กสฺสป น ตาว ชาตรูปสฺส อนฺตรธานํ โหติ, ยาว น ชาตรูปปฺปติรูปกํ โลเก อุปฺปชฺชติ.\csromanspacer{} ยโต จ โข กสฺสป ชาตรูปปฺปติรูปกํ โลเก อุปฺปชฺชติ, อถ ชาตรูปสฺส อนฺตรธานํ โหติ.\csromanspacer{} เอวเมว โข กสฺสป น ตาว สทฺธมฺมสฺส อนฺตรธานํ โหติ,}

\csromanheader{2}{5. กสฺสปสํยุตฺต}
\csromanmatikaanchor{mtk.192}
\csromanmatikaanchor{mtk.195}
\csromantocmarkref{subsubsection}{11. จีวรสุตฺต}{mtk.192}
\bookmark[dest={mtk.192},level=3]{11. จีวรสุตฺต}
\prose{\csromanfolio{423}ยาว น สทฺธมฺมปฺปติรูปกํ โลเก อุปฺปชฺชติ.\csromanspacer{} ยโต จ โข กสฺสป สทฺธมฺมปฺปติรูปกํ โลเก อุปฺปชฺชติ, อถ สทฺธมฺมสฺส อนฺตรธานํ โหติ.}

\prose{น โข กสฺสป ปถวีธาตุ สทฺธมฺมํ อนฺตรธาเปติ, น อาโปธาตุ สทฺธมฺมํ อนฺตรธาเปติ, น เตโชธาตุ สทฺธมฺมํ อนฺตรธาเปติ, น วาโยธาตุ สทฺธมฺมํ อนฺตรธาเปติ, อถ โข อิเธว เต อุปฺปชฺชนฺติ โมฆปุริสา, เย อิมํ สทฺธมฺมํ อนฺตรธาเปนฺติ.\csromanspacer{} เสยฺยถาปิ กสฺสป นาวา อาทิเกเนว โอปิลวติ, น โข กสฺสป เอวํ สทฺธมฺมสฺส อนฺตรธานํ โหติ.}

\prose{ปญฺจ โขเม กสฺสป โอกฺกมนิยา ธมฺมา สทฺธมฺมสฺส สมฺโมสาย อนฺตรธานาย สํวตฺตนฺติ.\csromanspacer{} กตเม ปญฺจ, อิธ กสฺสป ภิกฺขู ภิกฺขุนิโย อุปาสกา อุปาสิกาโย สตฺถริ อคารวา วิหรนฺติ อปฺปติสฺสา, ธมฺเม อคารวา วิหรนฺติ อปฺปติสฺสา, สํเฆ อคารวา วิหรนฺติ อปฺปติสฺสา, สิกฺขาย อคารวา วิหรนฺติ อปฺปติสฺสา, สมาธิสฺมิํ อคารวา วิหรนฺติ อปฺปติสฺสา.\csromanspacer{} อิเม โข กสฺสป ปญฺจ โอกฺกมนิยา ธมฺมา สทฺธมฺมสฺส สมฺโมสาย อนฺตรธานาย สํวตฺตนฺติ.}

\prose{ปญฺจ โขเม กสฺสป ธมฺมา สทฺธมฺมสฺส ฐิติยา อสมฺโมสาย อนนฺตรธานาย สํวตฺตนฺติ.\csromanspacer{} กตเม ปญฺจ, อิธ กสฺสป ภิกฺขู ภิกฺขุนิโย อุปาสกา อุปาสิกาโย สตฺถริ สคารวา วิหรนฺติ สปฺปติสฺสา, ธมฺเม สคารวา วิหรนฺติ สปฺปติสฺสา, สํเฆ สคารวา วิหรนฺติ สปฺปติสฺสา, สิกฺขาย สคารวา วิหรนฺติ สปฺปติสฺสา, สมาธิสฺมิํ สคารวา วิหรนฺติ สปฺปติสฺสา.\csromanspacer{} อิเม โข กสฺสป ปญฺจ ธมฺมา สทฺธมฺมสฺส ฐิติยา อสมฺโมสาย อนนฺตรธานาย สํวตฺตนฺตีติ.\csromanspacer{}\csromanspacer{}เตรสมํ.}

\vspace{\tipitakaparskip}
\csromancenter{\textbf{กสฺสปสํยุตฺตํ สมตฺตํ.}}
\titlehead{\textbf{ตสฺสุทฺทานํ}}
\csromangathameasure{สนฺตุฏฺฐญฺจ อโนตฺตปฺปี, จนฺทูปมํ กุลูปกํ. \\ ชิณฺณํ ตโย จ โอวาทา, ฌานาภิญฺญา อุปสฺสยํ. \\ จีวรํ ปรํมรณํ, สทฺธมฺมปฺปติรูปกนฺติ.}
\csromangathagroup{\csromangathastanza{สนฺตุฏฺฐญฺจ อโนตฺตปฺปี, จนฺทูปมํ กุลูปกํ. \\ ชิณฺณํ ตโย จ โอวาทา, ฌานาภิญฺญา อุปสฺสยํ. \\ จีวรํ ปรํมรณํ, สทฺธมฺมปฺปติรูปกนฺติ.}}

\csromanheader{2}{6. ลาภสกฺการสํยุตฺต}
\chapterhead{1. ปฐมวคฺค}
\csromanheader{5}{1. ทารุณสุตฺต}
\proseitem{157}{\csromanfolio{424}เอวํ เม สุตํ\csromandash{}เอกํ สมยํ ภควา สาวตฺถิยํ วิหรติ เชตวเน อนาถปิณฺฑิกสฺส อาราเม.\csromanspacer{} ตตฺร โข ภควา ภิกฺขู อามนฺเตสิ “ภิกฺขโว”ติ. “ภทนฺเต”ติ เต ภิกฺขู ภควโต ปจฺจสฺโสสุํ.\csromanspacer{} ภควา เอตทโวจ\csromandash{}}

\prose{ทารุโณ ภิกฺขเว ลาภสกฺการสิโลโก กฏุโก ผรุโส อนฺตรายิโก อนุตฺตรสฺส โยคกฺเขมสฺส อธิคมาย.\csromanspacer{} ตสฺมาติห ภิกฺขเว เอวํ สิกฺขิตพฺพํ “อุปฺปนฺนํ ลาภสกฺการสิโลกํ ปชหิสฺสาม, น จ โน อุปฺปนฺโน ลาภสกฺการสิโลโก จิตฺตํ ปริยาทาย ฐสฺสตี”ติ.\csromanspacer{} เอวํ หิ โว ภิกฺขเว สิกฺขิตพฺพนฺติ.\csromanspacer{}\csromanspacer{}ปฐมํ.}

\csromanheader{5}{2. พฬิสสุตฺต}
\proseitem{158}{สาวตฺถิยํ วิหรติ.\csromanspacer{} ทารุโณ ภิกฺขเว ลาภสกฺการสิโลโก กฏุโก ผรุโส อนฺตรายิโก อนุตฺตรสฺส โยคกฺเขมสฺส อธิคมาย.\csromanspacer{} เสยฺยถาปิ ภิกฺขเว พาฬิสิโก อามิสคตํ พฬิสํ คมฺภีเร อุทกรหเท ปกฺขิเปยฺย, ตเมนํ อญฺญตโร อามิสจกฺขุ มจฺโฉ คิเลยฺย.\csromanspacer{} เอวํ หิ โส ภิกฺขเว มจฺโฉ คิลพฬิโส พาฬิสิกสฺส อนยํ อาปนฺโน พฺยสนํ อาปนฺโน ยถากามกรณีโย พาฬิสิกสฺส.}

\prose{“พาฬิสิโก”ติ โข ภิกฺขเว มารสฺเสตํ ปาปิมโต อธิวจนํ. “พฬิสนฺ”ติ โข ภิกฺขเว ลาภสกฺการสิโลกสฺเสตํ อธิวจนํ.\csromanspacer{} โย หิ โกจิ ภิกฺขเว ภิกฺขุ อุปฺปนฺนํ ลาภสกฺการสิโลกํ อสฺสาเทติ นิกาเมติ.\csromanspacer{} อยํ วุจฺจติ ภิกฺขเว ภิกฺขุ คิลพฬิโส มารสฺส อนยํ อาปนฺโน พฺยสนํ อาปนฺโน ยถากามกรณีโย ปาปิมโต.\csromanspacer{} เอวํ ทารุโณ โข}

\vspace{\tipitakaparskip}
\csromancenter{ลาภสกฺการสํยุตฺต ภิกฺขเว ลาภสกฺการสิโลโก กฏุโก ผรุโส อนฺตรายิโก อนุตฺตรสฺส โยคกฺเขมสฺส อธิคมาย.\csromanspacer{} ตสฺมาติห ภิกฺขเว เอวํ สิกฺขิตพฺพํ “อุปฺปนฺนํ ลาภสกฺการสิโลกํ ปชหิสฺสาม, น จ โน อุปฺปนฺโน ลาภสกฺการสิโลโก จิตฺตํ ปริยาทาย ฐสฺสตี”ติ.\csromanspacer{} เอวํ หิ โว ภิกฺขเว สิกฺขิตพฺพนฺติ.\csromanspacer{}\csromanspacer{}ทุติยํ.}
\csromanheader{5}{3. กุมฺมสุตฺต}
\proseitem{159}{\csromanfolio{425}สาวตฺถิยํ วิหรติ.\csromanspacer{} ทารุโณ ภิกฺขเว ลาภสกฺการสิโลโก ฯเปฯ อธิคมาย.\csromanspacer{} ภูตปุพฺพํ ภิกฺขเว อญฺญตรสฺมิํ อุทกรหเท มหากุมฺมกุลํ จิรนิวาสิ อโหสิ.\csromanspacer{} อถ โข ภิกฺขเว อญฺญตโร กุมฺโม อญฺญตรํ กุมฺมํ เอตทโวจ “มา โข ตฺวํ ตาต กุมฺม เอตํ ปเทสํ อคมาสี”ติ.\csromanspacer{} อคมาสิ โข ภิกฺขเว โส กุมฺโม ตํ ปเทสํ, ตเมนํ ลุทฺโท ปปตาย วิชฺฌิ.\csromanspacer{} อถ โข ภิกฺขเว โส กุมฺโม เยน โส กุมฺโม เตนุปสงฺกมิ.\csromanspacer{} อทฺทสา โข ภิกฺขเว โส กุมฺโม ตํ กุมฺมํ ทูรโตว อาคจฺฉนฺตํ, ทิสฺวาน ตํ กุมฺมํ เอตทโวจ “กจฺจิ ตฺวํ ตาต กุมฺม น ตํ ปเทสํ อคมาสี”ติ.\csromanspacer{} อคมาสิํ ขฺวาหํ ตาต กุมฺม ตํ ปเทสนฺติ.\csromanspacer{} กจฺจิ ปนาสิ ตาต กุมฺม อกฺขโต อนุปหโตติ.\csromanspacer{} อกฺขโต โขมฺหิ ตาต กุมฺม อนุปหโต, อตฺถิ จ เม อิทํ สุตฺตกํ ปิฏฺฐิโต ปิฏฺฐิโต อนุพนฺธนฺติ.\csromanspacer{} ตคฺฆสิ ตาต กุมฺม ขโต ตคฺฆ อุปหโต, เอเตน หิ เต ตาต กุมฺม สุตฺตเกน ปิตโร จ ปิตามหา จ อนยํ อาปนฺนา พฺยาสนํ อาปนฺนา, คจฺฉ ทานิ ตฺวํ ตาต กุมฺม, น ทานิ ตฺวํ อมฺหากนฺติ.}

\prose{“ลุทฺโท”ติ โข ภิกฺขเว มารสฺเสตํ ปาปิมโต อธิวจนํ. “ปปตา”ติ โข ภิกฺขเว ลาภสกฺการสิโลกสฺเสตํ อธิวจนํ. “สุตฺตกนฺ”ติ โข ภิกฺขเว นนฺทีราคสฺเสตํ อธิวจนํ.\csromanspacer{} โย หิ โกจิ ภิกฺขเว ภิกฺขุ อุปฺปนฺนํ ลาภสกฺการสิโลกํ อสฺสาเทติ นิกาเมติ.\csromanspacer{} อยํ วุจฺจติ ภิกฺขเว ภิกฺขุ คิทฺโธ ปปตาย\footnote{ภิกฺขุ ปปตาย (สฺยา, กํ), ภิกฺขุ วิทฺโธ ปปตาย (?)} อนยํ อาปนฺโน พฺยสนํ อาปนฺโน ยถากามกรณีโย ปาปิมโต.\csromanspacer{} เอวํ ทารุโณ โข ภิกฺขเว ลาภสกฺการสิโลโก ฯเปฯ.\csromanspacer{} เอวํ หิ โว ภิกฺขเว สิกฺขิตพฺพนฺติ.\csromanspacer{}\csromanspacer{}ตติยํ.}

\gambhira{นิทานวคฺคสํยุตฺตปาฬิ}
\csromanheader{5}{4. ทีฆโลมิกสุตฺต}
\proseitem{160}{\csromanfolio{426}สาวตฺถิยํ วิหรติ.\csromanspacer{} ทารุโณ ภิกฺขเว ลาภสกฺการสิโลโก ฯเปฯ อธิคมาย.\csromanspacer{} เสยฺยถาปิ ภิกฺขเว ทีฆโลมิกา เอฬกา กณฺฏกคหนํ ปวิเสยฺย.\csromanspacer{} สา ตตฺร ตตฺร สชฺเชยฺย, ตตฺร ตตฺร คเยฺหยฺย\footnote{คจฺเฉยฺย (สี), คเณฺหยฺย (สฺยา, กํ, อิ, ก)}, ตตฺร ตตฺร พชฺเฌยฺย, ตตฺร ตตฺร อนยพฺยสนํ อาปชฺเชยฺย.\csromanspacer{} เอวเมว โข ภิกฺขเว อิเธกจฺโจ ภิกฺขุ ลาภสกฺการสิโลเกน อภิภูโต ปริยาทิณฺณจิตฺโต ปุพฺพณฺหสมยํ นิวาเสตฺวา ปตฺตจีวรมาทาย คามํ วา นิคมํ วา ปิณฺฑาย ปวิสติ.\csromanspacer{} โส ตตฺร ตตฺร สชฺชติ, ตตฺร ตตฺร คยฺหติ, ตตฺร ตตฺร พชฺฌติ, ตตฺร ตตฺร อนยพฺยสนํ อาปชฺชติ.\csromanspacer{} เอวํ ทารุโณ โข ภิกฺขเว ลาภสกฺการสิโลโก ฯเปฯ.\csromanspacer{} เอวํ หิ โว ภิกฺขเว สิกฺขิตพฺพนฺติ.\csromanspacer{}\csromanspacer{}จตุตฺถํ.}

\csromanheader{5}{5. มีฬฺหกสุตฺต}
\proseitem{161}{สาวตฺถิยํ วิหรติ.\csromanspacer{} ทารุโณ ภิกฺขเว ลาภสกฺการสิโลโก ฯเปฯ อธิคมาย.\csromanspacer{} เสยฺยถาปิ ภิกฺขเว มีฬฺหกา คูถาที คูถปูรา ปุณฺณา คูถสฺส, ปุรโต จสฺส มหาคูถปุญฺโช, สา เตน อญฺญา มีฬฺหกา อติมญฺเญยฺย “อหมฺหิ คูถาที คูถปูรา ปุณฺณา คูถสฺส, ปุรโต จ มฺยายํ มหาคูถปุญฺโช”ติ.\csromanspacer{} เอวเมว โข ภิกฺขเว อิเธกจฺโจ ภิกฺขุ ลาภสกฺการสิโลเกน อภิภูโต ปริยาทิณฺณจิตฺโต ปุพฺพณฺหสมยํ นิวาเสตฺวา ปตฺตจีวรมาทาย คามํ วา นิคมํ วา ปิณฺฑาย ปวิสติ.\csromanspacer{} โส ตตฺถ ภุตฺตาวี จ โหติ ยาวทตฺโถ, นิมนฺติโต จ สฺวาตนาย, ปิณฺฑปาโต จสฺส ปูโร, โส อารามํ คนฺตฺวา ภิกฺขุคณสฺส มชฺเฌ วิกตฺถติ “ภุตฺตาวี จมฺหิ ยาวทตฺโถ, นิมนฺติโต จมฺหิ สฺวาตนาย, ปิณฺฑปาโต จ มฺยายํ ปูโร, ลาภี จมฺหิ จีวรปิณฺฑปาตเสนาสนคิลานปฺปจฺจยเภสชฺชปริกฺขารานํ, อิเม ปนญฺเญ ภิกฺขู อปฺปปุญฺญา อปฺเปสกฺขา น ลาภิโน จีวรปิณฺฑปาตเสนาสนคิลานปฺปจฺจยเภสชฺชปริกฺขารานนฺ”ติ.\csromanspacer{} โส เตน ลาภสกฺการสิโลเกน อภิภูโต ปริยาทิณฺณจิตฺโต อญฺเญ เปสเล ภิกฺขู อติมญฺญติ, ตํ หิ ตสฺส ภิกฺขเว โมฆปุริสสฺส โหติ ทีฆรตฺตํ อหิตาย ทุกฺขาย.\csromanspacer{} เอวํ ทารุโณ โข ภิกฺขเว ลาภสกฺการสิโลโก ฯเปฯ.\csromanspacer{} เอวํ หิ โว ภิกฺขเว สิกฺขิตพฺพนฺติ.\csromanspacer{}\csromanspacer{}ปญฺจมํ.}

\csromanheader{5}{6. ลาภสกฺการสํยุตฺต \textbf{6. อสนิสุตฺต}}
\proseitem{162}{\csromanfolio{427}สาวตฺถิยํ วิหรติ.\csromanspacer{} ทารุโณ ภิกฺขเว ลาภสกฺการสิโลโก ฯเปฯ อธิคมาย.1 กํ ภิกฺขเว อสนิวิจกฺกํ อาคจฺฉตุ1, เสขํ\footnote{อสนิวิจกฺกํ, ตํ เสขํ (อิ, ก), อสนิวิจกฺกํ, เสขํ (สฺยา, กํ), อสนิวิจกฺกํ อาคจฺฉตุ, กํ เสขํ (?)} อปฺปตฺตมานสํ ลาภสกฺการสิโลโก อนุปาปุณาตุ\footnote{อนุปาปุณาติ (อิ, ก)}.}

\prose{“อสนิวิจกฺกนฺ”ติ โข ภิกฺขเว ลาภสกฺการสิโลกสฺเสตํ อธิวจนํ.\csromanspacer{} เอวํ ทารุโณ โข ภิกฺขเว ลาภสกฺการสิโลโก ฯเปฯ.\csromanspacer{} เอวํ หิ โว ภิกฺขเว สิกฺขิตพฺพนฺติ.\csromanspacer{}\csromanspacer{}ฉฏฺฐํ.}

\csromanheader{5}{7. ทิทฺธสุตฺต}
\proseitem{163}{สาวตฺถิยํ วิหรติ.\csromanspacer{} ทารุโณ ภิกฺขเว ลาภสกฺการสิโลโก ฯเปฯ อธิคมาย.\csromanspacer{} กํ ภิกฺขเว ทิทฺธคเตน วิสลฺเลน สลฺเลน\footnote{ทิฏฺฐิคเตน วิสลฺเลน (ก-สี), ทิฏฺฐิคเตน สลฺเลน (สฺยา, กํ), ทิฏฺฐิคเตน วิสลฺเลน สลฺเลน (ก), ทิฏฺฐคเตน วิสลฺเลน สลฺเลน (อิ)} วิชฺฌตุ, เสขํ\footnote{วิชฺฌตุ, ตํ เสขํ (สี), วิชฺฌติ, ตํ เสขํ (อิ, ก)} อปฺปตฺตมานสํ ลาภสกฺการสิโลโก อนุปาปุณาตุ3.}

\prose{“สลฺลนฺ”ติ โข ภิกฺขเว ลาภสกฺการสิโลกสฺเสตํ อธิวจนํ.\csromanspacer{} เอวํ ทารุโณ โข ภิกฺขเว ลาภสกฺการสิโลโก ฯเปฯ.\csromanspacer{} เอวํ หิ โว ภิกฺขเว สิกฺขิตพฺพนฺติ.\csromanspacer{}\csromanspacer{}สตฺตมํ.}

\csromanmatikaanchor{mtk.193}
\csromantocmarkref{subsubsection}{12. ปรํมรณสุตฺต}{mtk.193}
\bookmark[dest={mtk.193},level=3]{12. ปรํมรณสุตฺต}
\csromanheader{5}{8. สิงฺคาลสุตฺต}
\proseitem{164}{สาวตฺถิยํ วิหรติ.\csromanspacer{} ทารุโณ ภิกฺขเว ลาภสกฺการสิโลโก ฯเปฯ อธิคมาย.\csromanspacer{} อสฺสุตฺถ โน ตุเมฺห ภิกฺขเว รตฺติยา}

\gambhira{นิทานวคฺคสํยุตฺตปาฬิ}
\prose{\csromanfolio{428}ปจฺจูสสมยํ ชรสิงฺคาลสฺส\footnote{สิงฺคาลสฺส (ก), ชรสิคาลสฺส (สี, สฺยา, กํ)} วสฺสมานสฺสาติ.\csromanspacer{} เอวํ ภนฺเต.\csromanspacer{} เอโส โข ภิกฺขเว ชรสิงฺคาโล อุกฺกณฺฑเกน\footnote{อุกฺกณฺฏเกน (สี), อุกฺกณฺณเกน (สฺยา, กํ, อิ)} นาม โรคชาเตน ผุฏฺโฐ เนว พิลคโต รมติ, น รุกฺขมูลคโต รมติ, น อชฺโฌกาสคโต รมติ.\csromanspacer{} เยน เยน คจฺฉติ, ยตฺถ ยตฺถ ติฏฺฐติ, ยตฺถ ยตฺถ นิสีทติ, ยตฺถ ยตฺถ นิปชฺชติ, ตตฺถ ตตฺถ อนยพฺยสนํ อาปชฺชติ.\csromanspacer{} เอวเมว โข ภิกฺขเว อิเธกจฺโจ ภิกฺขุ ลาภสกฺการสิโลเกน อภิภูโต ปริยาทิณฺณจิตฺโต เนว สุญฺญาคารคโต รมติ, น รุกฺขมูลคโต รมติ, น อชฺโฌกาสคโต รมติ.\csromanspacer{} เยน เยน คจฺฉติ, ยตฺถ ยตฺถ ติฏฺฐติ, ยตฺถ ยตฺถ นิสีทติ, ยตฺถ ยตฺถ นิปชฺชติ, ตตฺถ ตตฺถ อนยพฺยสนํ อาปชฺชติ.\csromanspacer{} เอวํ ทารุโณ โข ภิกฺขเว ลาภสกฺการสิโลโก ฯเปฯ.\csromanspacer{} เอวํ หิ โว ภิกฺขเว สิกฺขิตพฺพนฺติ.\csromanspacer{}\csromanspacer{}อฏฺฐมํ.}

\csromanmatikaanchor{mtk.194}
\csromantocmarkref{subsubsection}{13. สทฺธมฺมปฺปติรูปกสุตฺต}{mtk.194}
\bookmark[dest={mtk.194},level=3]{13. สทฺธมฺมปฺปติรูปกสุตฺต}
\csromantocmarkref{section}{อุทฺทานคาถา}{mtk.195}
\bookmark[dest={mtk.195},level=1]{อุทฺทานคาถา}
\csromanheader{5}{9. เวรมฺภสุตฺต}
\proseitem{165}{สาวตฺถิยํ วิหรติ.\csromanspacer{} ทารุโณ ภิกฺขเว ลาภสกฺการสิโลโก ฯเปฯ อธิคมาย.\csromanspacer{} อุปริ ภิกฺขเว อากาเส เวรมฺภา\footnote{เวรมฺพา (สี, อิ)} นาม วาตา วายนฺติ.\csromanspacer{} ตตฺถ โย ปกฺขี คจฺฉติ, ตเมนํ เวรมฺภา วาตา ขิปนฺติ, ตสฺส เวรมฺภวาตกฺขิตฺตสฺส อญฺเญเนว ปาทา คจฺฉนฺติ, อญฺเญน ปกฺขา คจฺฉนฺติ, อญฺเญน สีสํ คจฺฉติ, อญฺเญน กาโย คจฺฉติ.\csromanspacer{} เอวเมว โข ภิกฺขเว อิเธกจฺโจ ภิกฺขุ ลาภสกฺการสิโลเกน อภิภูโต ปริยาทิณฺณจิตฺโต ปุพฺพณฺหสมยํ นิวาเสตฺวา ปตฺตจีวรมาทาย คามํ วา นิคมํ วา ปิณฺฑาย ปวิสติ, อรกฺขิเตเนว กาเยน อรกฺขิตาย วาจาย อรกฺขิเตน จิตฺเตน อนุปฏฺฐิตาย สติยา อสํวุเตหิ อินฺทฺริเยหิ, โส ตตฺถ ปสฺสติ มาตุคามํ ทุนฺนิวตฺถํ วา ทุปฺปารุตํ วา.\csromanspacer{} ตสฺส มาตุคามํ ทิสฺวา ทุนฺนิวตฺถํ วา ทุปฺปารุตํ วา ราโค จิตฺตํ อนุทฺธํเสติ, โส ราคานุทฺธํสิเตน จิตฺเตน สิกฺขํ ปจฺจกฺขาย หีนายาวตฺตติ.\csromanspacer{} ตสฺส อญฺเญ จีวรํ หรนฺติ, อญฺเญ ปตฺตํ หรนฺติ, อญฺเญ นิสีทนํ หรนฺติ, อญฺเญ สูจิฆรํ หรนฺติ เวรมฺภวาตกฺขิตฺตสฺเสว สกุณสฺส.\csromanspacer{} เอวํ ทารุโณ โข ภิกฺขเว ลาภสกฺการสิโลโก ฯเปฯ.\csromanspacer{} เอวํ หิ โว ภิกฺขเว สิกฺขิตพฺพนฺติ.\csromanspacer{}\csromanspacer{}นวมํ.}

\csromanheader{5}{6. ลาภสกฺการสํยุตฺต \textbf{10. สคาถกสุตฺต}}
\csromanmatikaanchor{mtk.196}
\csromanmatikaanchor{mtk.208}
\csromantocmarknopage{subsection}{6. ลาภสกฺการสํยุตฺต}{mtk.196}
\bookmark[dest={mtk.196},level=2]{6. ลาภสกฺการสํยุตฺต}
\proseitem{166}{\csromanfolio{429}สาวตฺถิยํ วิหรติ.\csromanspacer{} ทารุโณ ภิกฺขเว ลาภสกฺการสิโลโก ฯเปฯ อธิคมาย.\csromanspacer{} อิธาหํ ภิกฺขเว เอกจฺจํ ปุคฺคลํ ปสฺสามิ สกฺกาเรน อภิภูตํ ปริยาทิณฺณจิตฺตํ กายสฺส เภทา ปรํ มรณา อปายํ ทุคฺคติํ วินิปาตํ นิรยํ อุปปนฺนํ.\csromanspacer{} อิธ ปนาหํ ภิกฺขเว เอกจฺจํ ปุคฺคลํ ปสฺสามิ อสกฺกาเรน อภิภูตํ ปริยาทิณฺณจิตฺตํ กายสฺส เภทา ปรํ มรณา อปายํ ทุคฺคติํ วินิปาตํ นิรยํ อุปปนฺนํ.\csromanspacer{} อิธ ปนาหํ ภิกฺขเว เอกจฺจํ ปุคฺคลํ ปสฺสามิ สกฺกาเรน จ อสกฺกาเรน จ ตทุภเยน อภิภูตํ ปริยาทิณฺณจิตฺตํ กายสฺส เภทา ปรํ มรณา อปายํ ทุคฺคติํ วินิปาตํ นิรยํ อุปปนฺนํ.\csromanspacer{} เอวํ ทารุโณ โข ภิกฺขเว ลาภสกฺการสิโลโก ฯเปฯ.\csromanspacer{} เอวํ หิ โว ภิกฺขเว สิกฺขิตพฺพนฺติ.}

\prose{อิทมโวจ ภควา, อิทํ วตฺวาน สุคโต อถาปรํ เอตทโวจ สตฺถา\csromandash{}}

\csromangathameasure{“ยสฺส สกฺกริยมานสฺส, อสกฺกาเรน จูภยํ. \\ สมาธิ น วิกมฺปติ, อปฺปมาณวิหาริโน. \\ ตํ ฌายินํ สาตติกํ, สุขุมํทิฏฺฐิวิปสฺสกํ. \\ อุปาทานกฺขยารามํ, อาหุ สปฺปุริโส อิตี”ติ.}
\csromangathagroup{\csromangathastanza{“ยสฺส สกฺกริยมานสฺส, อสกฺกาเรน จูภยํ. \\ สมาธิ น วิกมฺปติ, อปฺปมาณวิหาริโน\footnote{อปฺปมาทวิหาริโน (อิ, ก) อปฺปมาโณติ เหตฺถ ผลสมาธิ, น สติ.}.}\csromangathastanza{ตํ ฌายินํ สาตติกํ, สุขุมํทิฏฺฐิวิปสฺสกํ. \\ อุปาทานกฺขยารามํ, อาหุ สปฺปุริโส อิตี”ติ.}}

\vspace{\tipitakaparskip}
\csromancenter{ทสมํ.}
\csromancenter{ปฐโม วคฺโค.}
\titlehead{\textbf{ตสฺสุทฺทานํ}}
\csromangathameasure{ทารุโณ พฬิสํ กุมฺมํ, ทีฆโลมิ จ มีฬฺหกํ. \\ อสนิ ทิทฺธํ สิงฺคาลํ, เวรมฺเภน สคาถกนฺติ.}
\csromangathagroup{\csromangathastanza{ทารุโณ พฬิสํ กุมฺมํ, ทีฆโลมิ จ มีฬฺหกํ. \\ อสนิ ทิทฺธํ สิงฺคาลํ, เวรมฺเภน สคาถกนฺติ.}}

\gambhira{นิทานวคฺคสํยุตฺตปาฬิ}
\chapterhead{2. ทุติยวคฺค}
\csromanmatikaanchor{mtk.197}
\csromanmatikaanchor{mtk.198}
\csromantocmarknopage{paragraph}{1. ปฐมวคฺค}{mtk.197}
\bookmark[dest={mtk.197},level=4]{1. ปฐมวคฺค}
\csromantocmarkref{subparagraph}{1. ทารุณสุตฺต}{mtk.198}
\bookmark[dest={mtk.198},level=5]{1. ทารุณสุตฺต}
\csromanheader{5}{1. สุวณฺณปาติสุตฺต}
\proseitem{167}{\csromanfolio{430}สาวตฺถิยํ วิหรติ.\csromanspacer{} ทารุโณ ภิกฺขเว ลาภสกฺการสิโลโก ฯเปฯ อธิคมาย.\csromanspacer{} อิธาหํ ภิกฺขเว เอกจฺจํ ปุคฺคลํ เอวํ เจตสา เจโต ปริจฺจ ปชานามิ “น จายมายสฺมา สุวณฺณปาติยาปิ รูปิยจุณฺณปริปูราย เหตุ สมฺปชานมุสา ภาเสยฺยา”ติ, ตเมนํ ปสฺสามิ อปเรน สมเยน ลาภสกฺการสิโลเกน อภิภูตํ ปริยาทิณฺณจิตฺตํ สมฺปชานมุสา ภาสนฺตํ.\csromanspacer{} เอวํ ทารุโณ โข ภิกฺขเว ลาภสกฺการสิโลโก ฯเปฯ.\csromanspacer{} เอวํ หิ โว ภิกฺขเว สิกฺขิตพฺพนฺติ.\csromanspacer{}\csromanspacer{}ปฐมํ.}

\csromanheader{5}{2. รูปิยปาติสุตฺต}
\csromanmatikaanchor{mtk.199}
\csromantocmarkref{subparagraph}{2. พฬิสสุตฺต}{mtk.199}
\bookmark[dest={mtk.199},level=5]{2. พฬิสสุตฺต}
\proseitem{168}{สาวตฺถิยํ วิหรติ.\csromanspacer{} ทารุโณ ภิกฺขเว ลาภสกฺการสิโลโก ฯเปฯ.\csromanspacer{} อิธาหํ ภิกฺขเว เอกจฺจํ ปุคฺคลํ เอวํ เจตสา เจโต ปริจฺจ ปชานามิ “น จายมายสฺมา รูปิยปาติยาปิ สุวณฺณจุณฺณปริปูราย เหตุ สมฺปชานมุสา ภาเสยฺยา”ติ, ตเมนํ ปสฺสามิ อปเรน สมเยน ลาภสกฺการสิโลเกน อภิภูตํ ปริยาทิณฺณจิตฺตํ, สมฺปชานมุสา ภาสนฺตํ.\csromanspacer{} เอวํ ทารุโณ โข ภิกฺขเว ลาภสกฺการสิโลโก ฯเปฯ.\csromanspacer{} เอวํ หิ โว ภิกฺขเว สิกฺขิตพฺพนฺติ.\csromanspacer{}\csromanspacer{}ทุติยํ.}

\csromanheader{5}{3-10. สุวณฺณนิกฺขสุตฺตาทิ-อฏฺฐก}
\proseitem{169}{สาวตฺถิยํ วิหรติ.\csromanspacer{} อิธาหํ ภิกฺขเว เอกจฺจํ ปุคฺคลํ เอวํ เจตสา เจโต ปริจฺจ ปชานามิ “น จายมายสฺมา สุวณฺณนิกฺขสฺสาปิ เหตุ ฯเปฯ สุวณฺณนิกฺขสตสฺสาปิ เหตุ.\csromanspacer{} สิงฺคีนิกฺขสฺสาปิ เหตุ.\csromanspacer{} สิงฺคีนิกฺขสตสฺสาปิ เหตุ.\csromanspacer{} ปถวิยาปิ ชาตรูปปริปูราย เหตุ.\csromanspacer{} อามิสกิญฺจิกฺขเหตุปิ.\csromanspacer{} ชีวิตเหตุปิ.\csromanspacer{} ชนปทกลฺยาณิยาปิ เหตุ สมฺปชานมุสา ภาเสยฺยา”ติ, ตเมนํ ปสฺสามิ อปเรน สมเยน ลาภสกฺการสิโลเกน อภิภูตํ ปริยาทิณฺณจิตฺตํ สมฺปชานมุสา ภาสนฺตํ.\csromanspacer{} เอวํ ทารุโณ โข ภิกฺขเว ลาภสกฺการสิโลโก ฯเปฯ.\csromanspacer{} เอวํ หิ โว ภิกฺขเว สิกฺขิตพฺพนฺติ.\csromanspacer{}\csromanspacer{}ทสมํ.}

\vspace{\tipitakaparskip}
\csromancenter{ทุติโย วคฺโค.}
\csromanmatikaanchor{mtk.200}
\csromantocmarkref{subparagraph}{3. กุมฺมสุตฺต}{mtk.200}
\bookmark[dest={mtk.200},level=5]{3. กุมฺมสุตฺต}
\titlehead{6. ลาภสกฺการสํยุตฺต \textbf{ตสฺสุทฺทานํ}}
\csromanmatikaanchor{mtk.201}
\csromanmatikaanchor{mtk.213}
\csromantocmarkref{subparagraph}{4. ทีฆโลมิกสุตฺต}{mtk.201}
\bookmark[dest={mtk.201},level=5]{4. ทีฆโลมิกสุตฺต}
\csromangathameasure{เทฺว ปาติ เทฺว สุวณฺณา จ, สิงฺคีหิ อปเร ทุเว. \\ ปถวี กิญฺจิกฺขชีวิตํ, ชนปทกลฺยาณิยา ทสาติ.}
\csromangathagroup{\csromangathastanza{\csromanfolio{431}เทฺว ปาติ เทฺว สุวณฺณา จ, สิงฺคีหิ อปเร ทุเว. \\ ปถวี กิญฺจิกฺขชีวิตํ, ชนปทกลฺยาณิยา ทสาติ.}}

\chapterhead{3. ตติยวคฺค}
\csromanheader{5}{1. มาตุคามสุตฺต}
\proseitem{170}{สาวตฺถิยํ วิหรติ.\csromanspacer{} ทารุโณ ภิกฺขเว ลาภสกฺการสิโลโก ฯเปฯ.\csromanspacer{} น ตสฺส ภิกฺขเว มาตุคาโม เอโก เอกสฺส จิตฺตํ ปริยาทาย ติฏฺฐติ, ยสฺส ลาภสกฺการสิโลโก จิตฺตํ ปริยาทาย ติฏฺฐติ.\csromanspacer{} เอวํ ทารุโณ โข ภิกฺขเว ลาภสกฺการสิโลโก ฯเปฯ.\csromanspacer{} เอวํ หิ โว ภิกฺขเว สิกฺขิตพฺพนฺติ.\csromanspacer{}\csromanspacer{}ปฐมํ.}

\csromanheader{5}{2. กลฺยาณีสุตฺต}
\csromanmatikaanchor{mtk.202}
\csromantocmarkref{subparagraph}{5. มีฬฺหกสุตฺต}{mtk.202}
\bookmark[dest={mtk.202},level=5]{5. มีฬฺหกสุตฺต}
\proseitem{171}{สาวตฺถิยํ วิหรติ.\csromanspacer{} ทารุโณ ภิกฺขเว ลาภสกฺการสิโลโก ฯเปฯ.\csromanspacer{} น ตสฺส ภิกฺขเว ชนปทกลฺยาณี เอกา เอกสฺส จิตฺตํ ปริยาทาย ติฏฺฐติ, ยสฺส ลาภสกฺการสิโลโก จิตฺตํ ปริยาทาย ติฏฺฐติ.\csromanspacer{} เอวํ ทารุโณ โข ภิกฺขเว ลาภสกฺการสิโลโก ฯเปฯ.\csromanspacer{} เอวํ หิ โว ภิกฺขเว สิกฺขิตพฺพนฺติ.\csromanspacer{}\csromanspacer{}ทุติยํ.}

\csromanheader{5}{3. เอกปุตฺตกสุตฺต}
\csromanmatikaanchor{mtk.203}
\csromantocmarkref{subparagraph}{6. อสนิสุตฺต}{mtk.203}
\bookmark[dest={mtk.203},level=5]{6. อสนิสุตฺต}
\proseitem{172}{สาวตฺถิยํ วิหรติ.\csromanspacer{} ทารุโณ ภิกฺขเว ลาภสกฺการสิโลโก ฯเปฯ.\csromanspacer{} สทฺธา ภิกฺขเว อุปาสิกา เอกปุตฺตกํ ปิยํ มนาปํ เอวํ สมฺมา อายาจมานา อายาเจยฺย “ตาทิโส ตาต ภวาหิ, ยาทิโส จิตฺโต จ คหปติ หตฺถโก จ อาฬวโก”ติ.\csromanspacer{} เอสา ภิกฺขเว ตุลา เอตํ ปมาณํ มม สาวกานํ อุปาสกานํ, ยทิทํ จิตฺโต จ คหปติ หตฺถโก จ อาฬวโก.\csromanspacer{} สเจ โข ตฺวํ ตาต อคารสฺมา อนคาริยํ ปพฺพชสิ, ตาทิโส ตาต ภวาหิ, ยาทิสา สาริปุตฺตโมคฺคลฺลานาติ.\csromanspacer{} เอสา ภิกฺขเว ตุลา เอตํ ปมาณํ มม สาวกานํ ภิกฺขูนํ, ยทิทํ}

\gambhira{นิทานวคฺคสํยุตฺตปาฬิ}
\prose{\csromanfolio{432}สาริปุตฺตโมคฺคลฺลานา.\csromanspacer{} มา จ โข ตฺวํ ตาต เสขํ อปฺปตฺตมานสํ ลาภสกฺการสิโลโก อนุปาปุณาตูติ.\csromanspacer{} ตญฺเจ ภิกฺขเว ภิกฺขุํ เสขํ อปฺปตฺตมานสํ ลาภสกฺการสิโลโก อนุปาปุณาติ, โส ตสฺส โหติ อนฺตรายาย.\csromanspacer{} เอวํ ทารุโณ โข ภิกฺขเว ลาภสกฺการสิโลโก ฯเปฯ.\csromanspacer{} เอวํ หิ โว ภิกฺขเว สิกฺขิตพฺพนฺติ.\csromanspacer{}\csromanspacer{}ตติยํ.}

\csromanmatikaanchor{mtk.204}
\csromantocmarkref{subparagraph}{7. ทิทฺธสุตฺต}{mtk.204}
\bookmark[dest={mtk.204},level=5]{7. ทิทฺธสุตฺต}
\csromanheader{5}{4. เอกธีตุสุตฺต}
\proseitem{173}{สาวตฺถิยํ วิหรติ.\csromanspacer{} ทารุโณ ภิกฺขเว ลาภสกฺการสิโลโก ฯเปฯ.\csromanspacer{} สทฺธา ภิกฺขเว อุปาสิกา เอกํ ธีตรํ ปิยํ มนาปํ เอวํ สมฺมา อายาจมานา อายาเจยฺย “ตาทิสา อยฺเย ภวาหิ, ยาทิสา ขุชฺชุตฺตรา จ อุปาสิกา เวฬุกณฺฑกิยา\footnote{เวฬุกณฺฑกี (สี, ฉกฺกงฺคุตฺตเรปิ)} จ นนฺทมาตา”ติ.\csromanspacer{} เอสา ภิกฺขเว ตุลา เอตํ ปมาณํ มม สาวิกานํ อุปาสิกานํ, ยทิทํ ขุชฺชุตฺตรา จ อุปาสิกา เวฬุกณฺฑกิยา จ นนฺทมาตา.\csromanspacer{} สเจ โข ตฺวํ อยฺเย อคารสฺมา อนคาริยํ ปพฺพชสิ, ตาทิสา อยฺเย ภวาหิ, ยาทิสา เขมา จ ภิกฺขุนี อุปฺปลวณฺณา จาติ.\csromanspacer{} เอสา ภิกฺขเว ตุลา เอตํ ปมาณํ มม สาวิกานํ ภิกฺขุนีนํ, ยทิทํ เขมา จ ภิกฺขุนี อุปฺปลวณฺณา จ.\csromanspacer{} มา จ โข ตฺวํ อยฺเย เสขํ อปฺปตฺตมานสํ ลาภสกฺการสิโลโก อนุปาปุณาตูติ.\csromanspacer{} ตญฺเจ ภิกฺขเว ภิกฺขุนิํ เสขํ อปฺปตฺตมานสํ ลาภสกฺการสิโลโก อนุปาปุณาติ, โส ตสฺสา โหติ อนฺตรายาย.\csromanspacer{} เอวํ ทารุโณ โข ภิกฺขเว ลาภสกฺการสิโลโก ฯเปฯ.\csromanspacer{} เอวํ หิ โว ภิกฺขเว สิกฺขิตพฺพนฺติ.\csromanspacer{}\csromanspacer{}จตุตฺถํ.}

\csromanheader{5}{5. สมณพฺราหฺมณสุตฺต}
\csromanmatikaanchor{mtk.205}
\csromantocmarkref{subparagraph}{8. สิงฺคาลสุตฺต}{mtk.205}
\bookmark[dest={mtk.205},level=5]{8. สิงฺคาลสุตฺต}
\proseitem{174}{สาวตฺถิยํ วิหรติ.\csromanspacer{} เย หิ เกจิ ภิกฺขเว สมณา วา พฺราหฺมณา วา ลาภสกฺการสิโลกสฺส อสฺสาทญฺจ อาทีนวญฺจ นิสฺสรณญฺจ ยถาภูตํ นปฺปชานนฺติ.\csromanspacer{} น เม เต ภิกฺขเว สมณา วา พฺราหฺมณา วา สมเณสุ วา สมณสมฺมตา พฺราหฺมเณสุ วา พฺราหฺมณสมฺมตา, น จ ปน เต อายสฺมนฺตา สามญฺญตฺถํ วา พฺรหฺมญฺญตฺถํ วา ทิฏฺเฐว ธมฺเม สยํ อภิญฺญา สจฺฉิกตฺวา อุปสมฺปชฺช วิหรนฺติ.\csromanspacer{} เย จ โข เกจิ ภิกฺขเว สมณา วา พฺราหฺมณา วา ลาภสกฺการสิโลกสฺส อสฺสาทญฺจ}

\vspace{\tipitakaparskip}
\csromancenter{ลาภสกฺการสํยุตฺต อาทีนวญฺจ นิสฺสรณญฺจ ยถาภูตํ ปชานนฺติ, เต จ โข เม ภิกฺขเว สมณา วา พฺราหฺมณา วา สมเณสุ เจว สมณสมฺมตา พฺราหฺมเณสุ จ พฺราหฺมณสมฺมตา, เต จ ปนายสฺมนฺโต สามญฺญตฺถญฺจ พฺรหฺมญฺญตฺถญฺจ ทิฏฺเฐว ธมฺเม สยํ อภิญฺญา สจฺฉิกตฺวา อุปสมฺปชฺช วิหรนฺตีติ.\csromanspacer{}\csromanspacer{}ปญฺจมํ.}
\csromanheader{5}{6. ทุติยสมณพฺราหฺมณสุตฺต}
\proseitem{175}{\csromanfolio{433}สาวตฺถิยํ วิหรติ.\csromanspacer{} เย หิ เกจิ ภิกฺขเว สมณา วา พฺราหฺมณา วา ลาภสกฺการสิโลกสฺส สมุทยญฺจ อตฺถงฺคมญฺจ อสฺสาทญฺจ อาทีนวญฺจ นิสฺสรณญฺจ ยาถาภูตํ นปฺปชานนฺติ ฯเปฯ ปชานนฺติ ฯเปฯ สยํ อภิญฺญา สจฺฉิกตฺวา อุปสมฺปชฺช วิหรนฺตีติ.\csromanspacer{}\csromanspacer{}ฉฏฺฐํ.}

\csromanmatikaanchor{mtk.206}
\csromantocmarkref{subparagraph}{9. เวรมฺภสุตฺต}{mtk.206}
\bookmark[dest={mtk.206},level=5]{9. เวรมฺภสุตฺต}
\csromanheader{5}{7. ตติยสมณพฺราหฺมณสุตฺต}
\proseitem{176}{สาวตฺถิยํ วิหรติ.\csromanspacer{} เย หิ เกจิ ภิกฺขเว สมณา วา พฺราหฺมณา วา ลาภสกฺการสิโลกํ ยถาภูตํ นปฺปชานนฺติ, ลาภสกฺการสิโลกสมุทยํ นปฺปชานนฺติ, ลาภสกฺการสิโลกนิโรธํ นปฺปชานนฺติ, ลาภสกฺการสิโลกนิโรธคามินิํ ปฏิปทํ นปฺปชานนฺติ ฯเปฯ ปชานนฺติ ฯเปฯ สยํ อภิญฺญา สจฺฉิกตฺวา อุปสมฺปชฺช วิหรนฺตีติ.\csromanspacer{}\csromanspacer{}สตฺตมํ.}

\csromanmatikaanchor{mtk.207}
\csromantocmarkref{subparagraph}{10. สคาถกสุตฺต}{mtk.207}
\bookmark[dest={mtk.207},level=5]{10. สคาถกสุตฺต}
\csromantocmarkref{section}{อุทฺทานคาถา}{mtk.208}
\bookmark[dest={mtk.208},level=1]{อุทฺทานคาถา}
\csromanheader{5}{8. ฉวิสุตฺต}
\proseitem{177}{สาวตฺถิยํ วิหรติ.\csromanspacer{} ทารุโณ ภิกฺขเว ลาภสกฺการสิโลโก, ลาภสกฺการสิโลโก ภิกฺขเว ฉวิํ ฉินฺทติ, ฉวิํ เฉตฺวา จมฺมํ ฉินฺทติ, จมฺมํ เฉตฺวา มํสํ ฉินฺทติ, มํสํ เฉตฺวา นฺหารุํ ฉินฺทติ, นฺหารุํ เฉตฺวา อฏฺฐิํ ฉินฺทติ, อฏฺฐิํ เฉตฺวา อฏฺฐิมิญฺชํ อาหจฺจ ติฏฺฐติ.\csromanspacer{} เอวํ ทารุโณ โข ภิกฺขเว ลาภสกฺการสิโลโก ฯเปฯ.\csromanspacer{} เอวํ หิ โว ภิกฺขเว สิกฺขิตพฺพนฺติ.\csromanspacer{}\csromanspacer{}อฏฺฐมํ.}

\csromanheader{5}{9. รชฺชุสุตฺต}
\proseitem{178}{สาวตฺถิยํ วิหรติ.\csromanspacer{} ทารุโณ ภิกฺขเว ลาภสกฺการสิโลโก, ลาภสกฺการสิโลโก ภิกฺขเว ฉวิํ ฉินฺทติ, ฉวิํ เฉตฺวา จมฺมํ ฉินฺทติ, จมฺมํ เฉตฺวา มํสํ ฉินฺทติ, มํสํ เฉตฺวา นฺหารุํ ฉินฺทติ, นฺหารุํ เฉตฺวา อฏฺฐิํ ฉินฺทติ, อฏฺฐิํ เฉตฺวา อฏฺฐิมิญฺชํ อาหจฺจ ติฏฺฐติ.}

\gambhira{นิทานวคฺคสํยุตฺตปาฬิ}
\csromanmatikaanchor{mtk.209}
\csromanmatikaanchor{mtk.225}
\csromantocmarknopage{paragraph}{2. ทุติยวคฺค}{mtk.209}
\bookmark[dest={mtk.209},level=4]{2. ทุติยวคฺค}
\prose{\csromanfolio{434}เสยฺยถาปิ ภิกฺขเว พลวา ปุริโส ทฬฺหาย วาฬรชฺชุยา ชงฺฆํ เวเฐตฺวา ฆํเสยฺย.\csromanspacer{} สา ฉวิํ ฉินฺเทยฺย, ฉวิํ เฉตฺวา จมฺมํ ฉินฺเทยฺย, จมฺมํ เฉตฺวา มํสํ ฉินฺเทยฺย, มํสํ เฉตฺวา นฺหารุํ ฉินฺเทยฺย, นฺหารุํ เฉตฺวา อฏฺฐิํ ฉินฺเทยฺย, อฏฺฐิํ เฉตฺวา อฏฺฐิมิญฺชํ อาหจฺจ ติฏฺเฐยฺย.\csromanspacer{} เอวเมว โข ภิกฺขเว ลาภสกฺการสิโลโก ฉวิํ ฉินฺทติ, ฉวิํ เฉตฺวา จมฺมํ ฉินฺทติ, จมฺมํ เฉตฺวา มํสํ ฉินฺทติ, มํสํ เฉตฺวา นฺหารุํ ฉินฺทติ, นฺหารุํ เฉตฺวา อฏฺฐิํ ฉินฺทติ, อฏฺฐิํ เฉตฺวา อฏฺฐิมิญฺชํ อาหจฺจ ติฏฺฐติ.\csromanspacer{} เอวํ ทารุโณ โข ภิกฺขเว ลาภสกฺการสิโลโก ฯเปฯ.\csromanspacer{} เอวํ หิ โว ภิกฺขเว สิกฺขิตพฺพนฺติ.\csromanspacer{}\csromanspacer{}นวมํ.}

\csromanheader{5}{10. ภิกฺขุสุตฺต}
\proseitem{179}{สาวตฺถิยํ วิหรติ.\csromanspacer{} โยปิ โส ภิกฺขเว ภิกฺขุ อรหํ ขีณาสโว, ตสฺสปาหํ ลาภสกฺการสิโลโก อนฺตรายาย วทามีติ.\csromanspacer{} เอวํ วุตฺเต อายสฺมา อานนฺโท ภควนฺตํ เอตทโวจ “กิสฺส ปน ภนฺเต ขีณาสวสฺส ภิกฺขุโน ลาภสกฺการสิโลโก อนฺตรายายา”ติ.\csromanspacer{} ยา หิสฺส สา อานนฺท อกุปฺปา เจโตวิมุตฺติ, นาหํ ตสฺสา ลาภสกฺการสิโลกํ อนฺตรายาย วทามิ.\csromanspacer{} เย จ ขฺวสฺส อานนฺท อปฺปมตฺตสฺส อาตาปิโน ปหิตตฺตสฺส วิหรโต ทิฏฺฐธมฺมสุขวิหารา อธิคตา, เตสาหมสฺส ลาภสกฺการสิโลกํ อนฺตรายาย วทามิ.\csromanspacer{} เอวํ ทารุโณ โข อานนฺท ลาภสกฺการสิโลโก กฏุโก ผรุโส อนฺตรายิโก อนุตฺตรสฺส โยคกฺเขมสฺส อธิคมาย.\csromanspacer{} ตสฺมาติหานนฺท เอวํ สิกฺขิตพฺพํ “อุปฺปนฺนํ ลาภสกฺการสิโลกํ ปชหิสฺสาม, น จ โน อุปฺปนฺโน ลาภสกฺการสิโลโก จิตฺตํ ปริยาทาย ฐสฺสตี”ติ.\csromanspacer{} เอวํ หิ โว อานนฺท สิกฺขิตพฺพนฺติ.\csromanspacer{}\csromanspacer{}ทสมํ.}

\vspace{\tipitakaparskip}
\csromancenter{ตติโย วคฺโค.}
\titlehead{\textbf{ตสฺสุทฺทานํ}}
\csromangathameasure{มาตุคาโม จ กลฺยาณี, ปุตฺโต จ เอกธีตุ จ. \\ สมณพฺราหฺมณา ตีณิ, ฉวิ รชฺชุ จ ภิกฺขุนาติ.}
\csromangathagroup{\csromangathastanza{มาตุคาโม จ กลฺยาณี, ปุตฺโต จ เอกธีตุ จ. \\ สมณพฺราหฺมณา ตีณิ, ฉวิ รชฺชุ จ ภิกฺขุนาติ.}}

\csromanmatikaanchor{mtk.210}
\csromantocmarkref{subparagraph}{1. สุวณฺณปาติสุตฺต}{mtk.210}
\bookmark[dest={mtk.210},level=5]{1. สุวณฺณปาติสุตฺต}
\csromanheader{1}{6. ลาภสกฺการสํยุตฺต \textbf{4. จตุตฺถวคฺค}}
\csromanheader{5}{1. ภินฺทิสุตฺต}
\csromanmatikaanchor{mtk.211}
\csromantocmarkref{subparagraph}{2. รูปิยปาติสุตฺต}{mtk.211}
\bookmark[dest={mtk.211},level=5]{2. รูปิยปาติสุตฺต}
\proseitem{180}{\csromanfolio{435}สาวตฺถิยํ วิหรติ.\csromanspacer{} ทารุโณ ภิกฺขเว ลาภสกฺการสิโลโก, ลาภสกฺการสิโลเกน อภิภูโต ปริยาทิณฺณจิตฺโต ภิกฺขเว เทวทตฺโต สํฆํ ภินฺทิ.\csromanspacer{} เอวํ ทารุโณ โข ภิกฺขเว ลาภสกฺการสิโลโก ฯเปฯ สิกฺขิตพฺพนฺติ.\csromanspacer{}\csromanspacer{}ปฐมํ.}

\csromanheader{5}{2. กุสลมูลสุตฺต}
\csromanmatikaanchor{mtk.212}
\csromantocmarkref{subparagraph}{3-10. สุวณฺณนิกฺขสุตฺตาทิ-อฏฺฐก}{mtk.212}
\bookmark[dest={mtk.212},level=5]{3-10. สุวณฺณนิกฺขสุตฺตาทิ-อฏฺฐก}
\csromantocmarkref{section}{อุทฺทานคาถา}{mtk.213}
\bookmark[dest={mtk.213},level=1]{อุทฺทานคาถา}
\csromantocmarknopage{paragraph}{3. ตติยวคฺค}{mtk.214}
\bookmark[dest={mtk.214},level=4]{3. ตติยวคฺค}
\proseitem{181}{สาวตฺถิยํ วิหรติ.\csromanspacer{} ทารุโณ ภิกฺขเว ลาภสกฺการสิโลโก, ลาภสกฺการสิโลเกน อภิภูตสฺส ปริยาทิณฺณจิตฺตสฺส ภิกฺขเว เทวทตฺตสฺส กุสลมูลํ สมุจฺเฉทมคมา.\csromanspacer{} เอวํ ทารุโณ โข ภิกฺขเว ลาภสกฺการสิโลโก ฯเปฯ สิกฺขิตพฺพนฺติ.\csromanspacer{} ทุติยํ.}

\csromanheader{5}{3. กุสลธมฺมสุตฺต}
\proseitem{182}{สาวตฺถิยํ วิหรติ.\csromanspacer{} ทารุโณ ภิกฺขเว ลาภสกฺการสิโลโก, ลาภสกฺการสิโลเกน อภิภูตสฺส ปริยาทิณฺณจิตฺตสฺส ภิกฺขเว เทวทตฺตสฺส กุสโล ธมฺโม สมุจฺเฉทมคมา.\csromanspacer{} เอวํ ทารุโณ โข ภิกฺขเว ลาภสกฺการสิโลโก ฯเปฯ สิกฺขิตพฺพนฺติ.\csromanspacer{}\csromanspacer{}ตติยํ.}

\csromanheader{5}{4. สุกฺกธมฺมสุตฺต}
\proseitem{183}{สาวตฺถิยํ วิหรติ.\csromanspacer{} ทารุโณ ภิกฺขเว ลาภสกฺการสิโลโก, ลาภสกฺการสิโลเกน อภิภูตสฺส ปริยาทิณฺณจิตฺตสฺส ภิกฺขเว เทวทตฺตสฺส สุกฺโก ธมฺโม สมุจฺเฉทมคมา.\csromanspacer{} เอวํ ทารุโณ โข ภิกฺขเว ลาภสกฺการสิโลโก ฯเปฯ สิกฺขิตพฺพนฺติ.\csromanspacer{}\csromanspacer{}จตุตฺถํ.}

\csromanheader{5}{5. อจิรปกฺกนฺตสุตฺต}
\csromanmatikaanchor{mtk.215}
\csromantocmarkref{subparagraph}{1. มาตุคามสุตฺต}{mtk.215}
\bookmark[dest={mtk.215},level=5]{1. มาตุคามสุตฺต}
\proseitem{184}{เอกํ สมยํ ภควา ราชคเห วิหรติ คิชฺฌกูเฏ ปพฺพเต อจิรปกฺกนฺเต เทวทตฺเต.\csromanspacer{} ตตฺร โข ภควา เทวทตฺตํ อารพฺภ ภิกฺขู}

\gambhira{นิทานวคฺคสํยุตฺตปาฬิ}
\csromanmatikaanchor{mtk.216}
\csromantocmarkref{subparagraph}{2. กลฺยาณีสุตฺต}{mtk.216}
\bookmark[dest={mtk.216},level=5]{2. กลฺยาณีสุตฺต}
\prose{\csromanfolio{436}อามนฺเตสิ “อตฺตวธาย ภิกฺขเว เทวทตฺตสฺส ลาภสกฺการสิโลโก อุทปาทิ, ปราภวาย เทวทตฺตสฺส ลาภสกฺการสิโลโก อุทปาทิ.}

\prose{เสยฺยถาปิ ภิกฺขเว กทลี อตฺตวธาย ผลํ เทติ, ปราภวาย ผลํ เทติ.\csromanspacer{} เอวเมว โข ภิกฺขเว อตฺตวธาย เทวทตฺตสฺส ลาภสกฺการสิโลโก อุทปาทิ, ปราภวาย เทวทตฺตสฺส ลาภสกฺการสิโลโก อุทปาทิ.}

\csromanmatikaanchor{mtk.217}
\csromantocmarkref{subparagraph}{3. เอกปุตฺตกสุตฺต}{mtk.217}
\bookmark[dest={mtk.217},level=5]{3. เอกปุตฺตกสุตฺต}
\prose{เสยฺยถาปิ ภิกฺขเว เวฬุ อตฺตวธาย ผลํ เทติ, ปราภวาย ผลํ เทติ.\csromanspacer{} เอวเมว โข ภิกฺขเว อตฺตวธาย เทวทตฺตสฺส ลาภสกฺการสิโลโก อุทปาทิ, ปราภวาย เทวทตฺตสฺส ลาภสกฺการสิโลโก อุทปาทิ.}

\prose{เสยฺยถาปิ ภิกฺขเว นโฬ อตฺตวธาย ผลํ เทติ, ปราภวาย ผลํ เทติ.\csromanspacer{} เอวเมว โข ภิกฺขเว อตฺตวธาย เทวทตฺตสฺส ลาภสกฺการสิโลโก อุทปาทิ, ปราภวาย เทวทตฺตสฺส ลาภสกฺการสิโลโก อุทปาทิ.}

\prose{เสยฺยถาปิ ภิกฺขเว อสฺสตรี อตฺตวธาย คพฺภํ คณฺหาติ, ปราภวาย คพฺภํ คณฺหาติ.\csromanspacer{} เอวเมว โข ภิกฺขเว อตฺตวธาย เทวทตฺตสฺส ลาภสกฺการสิโลโก อุทปาทิ, ปราภวาย เทวทตฺตสฺส ลาภสกฺการสิโลโก อุทปาทิ.\csromanspacer{} เอวํ ทารุโณ โข ภิกฺขเว ลาภสกฺการสิโลโก.\csromanspacer{} เอวํ หิ โว ภิกฺขเว สิกฺขิตพฺพนฺ”ติ.}

\prose{อิทมโวจ ภควา, อิทํ วตฺวาน สุคโต อถาปรํ เอตทโวจ สตฺถา\csromandash{}}

\csromanmatikaanchor{mtk.218}
\csromantocmarkref{subparagraph}{4. เอกธีตุสุตฺต}{mtk.218}
\bookmark[dest={mtk.218},level=5]{4. เอกธีตุสุตฺต}
\csromangathameasure{“ผลํ เว กทลิํ หนฺติ, ผลํ เวฬุํ ผลํ นฬํ. \\ สกฺกาโร กาปุริสํ หนฺติ, คพฺโภ อสฺสตริํ ยถา”ติ.}
\csromangathagroup{\csromangathastanza{“ผลํ เว กทลิํ หนฺติ, ผลํ เวฬุํ ผลํ นฬํ. \\ สกฺกาโร กาปุริสํ หนฺติ, คพฺโภ อสฺสตริํ ยถา”ติ.}}

\vspace{\tipitakaparskip}
\csromancenter{ปญฺจมํ.}
\csromanmatikaanchor{mtk.219}
\csromantocmarkref{subparagraph}{5. สมณพฺราหฺมณสุตฺต}{mtk.219}
\bookmark[dest={mtk.219},level=5]{5. สมณพฺราหฺมณสุตฺต}
\csromanheader{5}{6. ปญฺจรถสตสุตฺต}
\proseitem{185}{ราชคเห วิหรติ เวฬุวเน กลนฺทกนิวาเป.\csromanspacer{} เตน โข ปน สมเยน เทวทตฺตสฺส อชาตสตฺตุกุมาโร ปญฺจหิ รถสเตหิ สายํ}

\vspace{\tipitakaparskip}
\csromancenter{ลาภสกฺการสํยุตฺต ปาตํ อุปฏฺฐานํ คจฺฉติ, ปญฺจ จ ถาลิปากสตานิ ภตฺตาภิหาโร อภิหรียติ.\csromanspacer{} อถ โข สมฺพหุลา ภิกฺขู เยน ภควา เตนุปสงฺกมิํสุ, อุปสงฺกมิตฺวา ภควนฺตํ อภิวาเทตฺวา เอกมนฺตํ นิสีทิํสุ, เอกมนฺตํ นิสินฺนา โข เต ภิกฺขู ภควนฺตํ เอตทโวจุํ “เทวทตฺตสฺส ภนฺเต อชาตสตฺตุกุมาโร ปญฺจหิ รถสเตหิ สายํ ปาตํ อุปฏฺฐานํ คจฺฉติ, ปญฺจ จ ถาลิปากสตานิ ภตฺตาภิหาโร อภิหรียตี”ติ.\csromanspacer{} มา ภิกฺขเว เทวทตฺตสฺส ลาภสกฺการสิโลกํ ปิหยิตฺถ, ยาวกีวญฺจ ภิกฺขเว เทวทตฺตสฺส อชาตสตฺตุกุมาโร ปญฺจหิ รถสเตหิ สายํ ปาตํ อุปฏฺฐานํ คมิสฺสติ, ปญฺจ จ ถาลิปากสตานิ ภตฺตาภิหาโร อาหรียิสฺสติ.\csromanspacer{} หานิเยว ภิกฺขเว เทวทตฺตสฺส ปาฏิกงฺขา กุสเลสุ ธมฺเมสุ, โน วุฑฺฒิ.}
\csromanmatikaanchor{mtk.220}
\csromantocmarkref{subparagraph}{6. ทุติยสมณพฺราหฺมณสุตฺต}{mtk.220}
\bookmark[dest={mtk.220},level=5]{6. ทุติยสมณพฺราหฺมณสุตฺต}
\prose{\csromanfolio{437}เสยฺยถาปิ ภิกฺขเว จณฺฑสฺส กุกฺกุรสฺส นาสาย ปิตฺตํ ภินฺเทยฺยุํ, เอวํ หิ โส ภิกฺขเว กุกฺกุโร ภิยฺโยโส มตฺตาย จณฺฑตโร อสฺส.\csromanspacer{} เอวเมว ภิกฺขเว ยาวกีวญฺจ เทวทตฺตสฺส อชาตสตฺตุกุมาโร ปญฺจหิ รถสเตหิ สายํ ปาตํ อุปฏฺฐานํ คมิสฺสติ, ปญฺจ จ ถาลิปากสตานิ ภตฺตาภิหาโร อาหรียิสฺสติ.\csromanspacer{} หานิเยว ภิกฺขเว เทวทตฺตสฺส ปาฏิกงฺขา กุสเลสุ ธมฺเมสุ, โน วุทฺธิ.\csromanspacer{} เอวํ ทารุโณ โข ภิกฺขเว ลาภสกฺการสิโลโก ฯเปฯ.\csromanspacer{} เอวํ หิ โว ภิกฺขเว สิกฺขิตพฺพนฺติ.\csromanspacer{}\csromanspacer{}ฉฏฺฐํ.}

\csromanheader{5}{7. มาตุสุตฺต}
\csromanmatikaanchor{mtk.221}
\csromantocmarkref{subparagraph}{7. ตติยสมณพฺราหฺมณสุตฺต}{mtk.221}
\bookmark[dest={mtk.221},level=5]{7. ตติยสมณพฺราหฺมณสุตฺต}
\proseitem{186}{สาวตฺถิยํ วิหรติ.\csromanspacer{} ทารุโณ ภิกฺขเว ลาภสกฺการสิโลโก กฏุโก ผรุโส อนฺตรายิโก อนุตฺตรสฺส โยคกฺเขมสฺส อธิคมาย.\csromanspacer{} อิธาหํ ภิกฺขเว เอกจฺจํ ปุคฺคลํ เอวํ เจตสา เจโต ปริจฺจ ปชานามิ “น จายมายสฺมา มาตุปิ เหตุ สมฺปชานมุสา ภาเสยฺยา”ติ, ตเมนํ ปสฺสามิ อปเรน สมเยน ลาภสกฺการสิโลเกน อภิภูตํ ปริยาทิณฺณจิตฺตํ สมฺปชานมุสา ภาสนฺตํ.\csromanspacer{} เอวํ ทารุโณ โข ภิกฺขเว ลาภสกฺการสิโลโก กฏุโก ผรุโส อนฺตรายิโก อนุตฺตรสฺส โยคกฺเขมสฺส อธิคมาย, ตสฺมาติห ภิกฺขเว เอวํ สิกฺขิตพฺพํ “อุปฺปนฺนํ ลาภสกฺการสิโลกํ ปชหิสฺสาม, น จ โน อุปฺปนฺโน ลาภสกฺการสิโลโก จิตฺตํ ปริยาทาย ฐสฺสตี”ติ.\csromanspacer{} เอวํ หิ โว ภิกฺขเว สิกฺขิตพฺพนฺติ.\csromanspacer{}\csromanspacer{}สตฺตมํ.}

\gambhira{นิทานวคฺคสํยุตฺตปาฬิ}
\csromanmatikaanchor{mtk.222}
\csromantocmarkref{subparagraph}{8. ฉวิสุตฺต}{mtk.222}
\bookmark[dest={mtk.222},level=5]{8. ฉวิสุตฺต}
\csromanheader{5}{8-13. ปิตุสุตฺตาทิฉกฺก}
\csromanmatikaanchor{mtk.223}
\csromanmatikaanchor{mtk.235}
\csromantocmarkref{subparagraph}{9. รชฺชุสุตฺต}{mtk.223}
\bookmark[dest={mtk.223},level=5]{9. รชฺชุสุตฺต}
\proseitem{187}{\csromanfolio{438}สาวตฺถิยํ วิหรติ.\csromanspacer{} ทารุโณ ภิกฺขเว ลาภสกฺการสิโลโก กฏุโก ผรุโส อนฺตรายิโก อนุตฺตรสฺส โยคกฺเขมสฺส อธิคมาย.\csromanspacer{} อิธาหํ ภิกฺขเว เอกจฺจํ ปุคฺคลํ เอวํ เจตสา เจโต ปริจฺจ ปชานามิ “น จายมายสฺมา ปิตุปิ เหตุ ฯเปฯ. (วิตฺถาเรตพฺพํ.) ภาตุปิ เหตุ.\csromanspacer{} ภคินิยาปิ เหตุ.\csromanspacer{} ปุตฺตสฺสปิ เหตุ.\csromanspacer{} ธีตุยาปิ เหตุ.\csromanspacer{} ปชาปติยาปิ เหตุ สมฺปชานมุสา ภาเสยฺยา”ติ, ตเมนํ ปสฺสามิ อปเรน สมเยน ลาภสกฺการสิโลเกน อภิภูตํ ปริยาทิณฺณจิตฺตํ สมฺปชานมุสา ภาสนฺตํ.\csromanspacer{} เอวํ ทารุโณ โข ภิกฺขเว ลาภสกฺการสิโลโก กฏุโก ผรุโส อนฺตรายิโก อนุตฺตรสฺส โยคกฺเขมสฺส อธิคมาย.\csromanspacer{} ตสฺมาติห ภิกฺขเว เอวํ สิกฺขิตพฺพํ “อุปฺปนฺนํ ลาภสกฺการสิโลกํ ปชหิสฺสาม, น จ โน อุปฺปนฺโน ลาภสกฺการสิโลโก จิตฺตํ ปริยาทาย ฐสฺสตี”ติ.\csromanspacer{} เอวํ หิ โว ภิกฺขเว สิกฺขิตพฺพนฺติ.\csromanspacer{} เตรสมํ.}

\vspace{\tipitakaparskip}
\csromancenter{จตุตฺโถ วคฺโค.}
\titlehead{\textbf{ตสฺสุทฺทานํ}}
\csromangathameasure{ภินฺทิ มูลํ ทุเว ธมฺมา, ปกฺกนฺตํ รถ มาตริ. \\ ปิตา ภาตา จ ภคินี, ปุตฺโต ธีตา ปชาปตีติ. \\ \textbf{ลาภสกฺการสํยุตฺตํ สมตฺตํ.}}
\csromangathagroup{\csromangathastanza{ภินฺทิ มูลํ ทุเว ธมฺมา, ปกฺกนฺตํ รถ มาตริ. \\ ปิตา ภาตา จ ภคินี, ปุตฺโต ธีตา ปชาปตีติ. \\ \textbf{ลาภสกฺการสํยุตฺตํ สมตฺตํ.}}}

\csromanmatikaanchor{mtk.224}
\csromantocmarkref{subparagraph}{10. ภิกฺขุสุตฺต}{mtk.224}
\bookmark[dest={mtk.224},level=5]{10. ภิกฺขุสุตฺต}
\csromantocmarkref{section}{อุทฺทานคาถา}{mtk.225}
\bookmark[dest={mtk.225},level=1]{อุทฺทานคาถา}
\csromanheader{5}{7. ราหุลสํยุตฺต}
\chapterhead{1. ปฐมวคฺค}
\csromanheader{6}{1. จกฺขุสุตฺต}
\proseitem{188}{\csromanfolio{439}เอวํ เม สุตํ\csromandash{}เอกํ สมยํ ภควา สาวตฺถิยํ วิหรติ เชตวเน อนาถปิณฺฑิกสฺส อาราเม.\csromanspacer{} อถ โข อายสฺมา ราหุโล เยน ภควา เตนุปสงฺกมิ, อุปสงฺกมิตฺวา ภควนฺตํ อภิวาเทตฺวา เอกมนฺตํ นิสีทิ, เอกมนฺตํ นิสินฺโน โข อายสฺมา ราหุโล ภควนฺตํ เอตทโวจ “สาธุ เม ภนฺเต ภควา สํขิตฺเตน ธมฺมํ เทเสตุ, ยมหํ ภควโต ธมฺมํ สุตฺวา เอโก วูปกฏฺโฐ อปฺปมตฺโต อาตาปี ปหิตตฺโต วิหเรยฺยนฺ”ติ.}

\prose{ตํ กิํ มญฺญสิ ราหุล, จกฺขุํ นิจฺจํ วา อนิจฺจํ วาติ.\csromanspacer{} อนิจฺจํ ภนฺเต.\csromanspacer{} ยํ ปนานิจฺจํ, ทุกฺขํ วา ตํ สุขํ วาติ.\csromanspacer{} ทุกฺขํ ภนฺเต.\csromanspacer{} ยํ ปนานิจฺจํ ทุกฺขํ วิปริณามธมฺมํ, กลฺลํ นุ ตํ สมนุปสฺสิตุํ “เอตํ มม, เอโสหมสฺมิ, เอโส เม อตฺตา”ติ.\csromanspacer{} โน เหตํ ภนฺเต.}

\prose{โสตํ นิจฺจํ วา อนิจฺจํ วาติ.\csromanspacer{} อนิจฺจํ ภนฺเต ฯเปฯ.\csromanspacer{} ฆานํ นิจฺจํ วา อนิจฺจํ วาติ.\csromanspacer{} อนิจฺจํ ภนฺเต.\csromanspacer{} ชิวฺหา นิจฺจา วา อนิจฺจา วาติ.\csromanspacer{} อนิจฺจา ภนฺเต.\csromanspacer{} กาโย นิจฺโจ วา อนิจฺโจ วาติ.\csromanspacer{} อนิจฺโจ ภนฺเต.\csromanspacer{} มโน นิจฺโจ วา อนิจฺโจ วาติ.\csromanspacer{} อนิจฺโจ ภนฺเต.\csromanspacer{} ยํ ปนานิจฺจํ ทุกฺขํ วา ตํ สุขํ วาติ.\csromanspacer{} ทุกฺขํ ภนฺเต.\csromanspacer{} ยํ ปนานิจฺจํ ทุกฺขํ วิปริณามธมฺมํ, กลฺลํ นุ ตํ สมนุปสฺสิตุํ “เอตํ มม, เอโสหมสฺมิ, เอโส เม อตฺตา”ติ.\csromanspacer{} โน เหตํ ภนฺเต.}

\csromanmatikaanchor{mtk.226}
\csromanmatikaanchor{mtk.227}
\csromantocmarknopage{paragraph}{4. จตุตฺถวคฺค}{mtk.226}
\bookmark[dest={mtk.226},level=4]{4. จตุตฺถวคฺค}
\csromantocmarkref{subparagraph}{1. ภินฺทิสุตฺต}{mtk.227}
\bookmark[dest={mtk.227},level=5]{1. ภินฺทิสุตฺต}
\prose{เอวํ ปสฺสํ ราหุล สุตวา อริยสาวโก จกฺขุสฺมิมฺปิ นิพฺพินฺทติ ฯเปฯ.\csromanspacer{} โสตสฺมิมฺปิ นิพฺพินฺทติ.\csromanspacer{} ฆานสฺมิมฺปิ นิพฺพินฺทติ.\csromanspacer{} ชิวฺหายปิ นิพฺพินฺทติ.\csromanspacer{} กายสฺมิมฺปิ นิพฺพินฺทติ.\csromanspacer{} มนสฺมิมฺปิ นิพฺพินฺทติ, นิพฺพินฺทํ วิรชฺชติ, วิราคา วิมุจฺจติ, วิมุตฺตสฺมิํ “วิมุตฺตมฺ”อิติ ญาณํ โหติ, “ขีณา ชาติ, วุสิตํ พฺรหฺมจริยํ, กตํ กรณียํ, นาปรํ อิตฺถตฺตายา”ติ ปชานาตีติ. (เอเตน เปยฺยาเลน ทส สุตฺตนฺตา กาตพฺพา.) .\csromanspacer{} ปฐมํ.}

\gambhira{นิทานวคฺคสํยุตฺตปาฬิ}
\csromanmatikaanchor{mtk.228}
\csromantocmarkref{subparagraph}{2. กุสลมูลสุตฺต}{mtk.228}
\bookmark[dest={mtk.228},level=5]{2. กุสลมูลสุตฺต}
\csromanheader{6}{2. รูปสุตฺต}
\proseitem{189}{\csromanfolio{440}สาวตฺถิยํ วิหรติ.\csromanspacer{} ตํ กิํ มญฺญสิ ราหุล, รูปา นิจฺจา วา อนิจฺจา วาติ.\csromanspacer{} อนิจฺจา ภนฺเต ฯเปฯ.\csromanspacer{} สทฺทา.\csromanspacer{} คนฺธา.\csromanspacer{} รสา.\csromanspacer{} โผฏฺฐพฺพา.\csromanspacer{} ธมฺมา นิจฺจา วา อนิจฺจา วาติ.\csromanspacer{} อนิจฺจา ภนฺเต.\csromanspacer{} เอวํ ปสฺสํ ราหุล สุตวา อริยสาวโก รูเปสุปิ นิพฺพินฺทติ.\csromanspacer{} สทฺเทสุปิ นิพฺพินฺทติ.\csromanspacer{} คนฺเธสุปิ นิพฺพินฺทติ.\csromanspacer{} รเสสุปิ นิพฺพินฺทติ.\csromanspacer{} โผฏฺฐพฺเพสุปิ นิพฺพินฺทติ.\csromanspacer{} ธมฺเมสุปิ นิพฺพินฺทติ, นิพฺพินฺทํ วิรชฺชติ ฯเปฯ ปชานาตีติ.\csromanspacer{} ทุติยํ.}

\csromanmatikaanchor{mtk.229}
\csromantocmarkref{subparagraph}{3. กุสลธมฺมสุตฺต}{mtk.229}
\bookmark[dest={mtk.229},level=5]{3. กุสลธมฺมสุตฺต}
\csromanheader{6}{3. วิญฺญาณสุตฺต}
\proseitem{190}{สาวตฺถิยํ วิหรติ.\csromanspacer{} ตํ กิํ มญฺญสิ ราหุล, จกฺขุวิญฺญาณํ นิจฺจํ วา อนิจฺจํ วาติ.\csromanspacer{} อนิจฺจํ ภนฺเต.\csromanspacer{} โสตวิญฺญาณํ ฯเปฯ.\csromanspacer{} ฆานวิญฺญาณํ.\csromanspacer{} ชิวฺหาวิญฺญาณํ.\csromanspacer{} กายวิญฺญาณํ.\csromanspacer{} มโนวิญฺญาณํ นิจฺจํ วา อนิจฺจํ วาติ.\csromanspacer{} อนิจฺจํ ภนฺเต.\csromanspacer{} เอวํ ปสฺสํ ราหุล สุตวา อริยสาวโก จกฺขุวิญฺญาณสฺมิมฺปิ นิพฺพินฺทติ ฯเปฯ.\csromanspacer{} โสตวิญฺญาณสฺมิมฺปิ นิพฺพินฺทติ.\csromanspacer{} ฆานวิญฺญาณสฺมิมฺปิ นิพฺพินฺทติ.\csromanspacer{} ชิวฺหาวิญฺญาณสฺมิมฺปิ นิพฺพินฺทติ.\csromanspacer{} กายวิญฺญาณสฺมิมฺปิ นิพฺพินฺทติ.\csromanspacer{} มโนวิญฺญาณสฺมิมฺปิ นิพฺพินฺทติ, นิพฺพินฺทํ วิรชฺชติ ฯเปฯ ปชานาตีติ.\csromanspacer{}\csromanspacer{}ตติยํ.}

\csromanmatikaanchor{mtk.230}
\csromantocmarkref{subparagraph}{4. สุกฺกธมฺมสุตฺต}{mtk.230}
\bookmark[dest={mtk.230},level=5]{4. สุกฺกธมฺมสุตฺต}
\csromanheader{6}{4. สมฺผสฺสสุตฺต}
\proseitem{191}{สาวตฺถิยํ วิหรติ.\csromanspacer{} ตํ กิํ มญฺญสิ ราหุล, จกฺขุสมฺผสฺโส นิจฺโจ วา อนิจฺโจ วาติ.\csromanspacer{} อนิจฺโจ ภนฺเต.\csromanspacer{} โสตสมฺผสฺโส ฯเปฯ.\csromanspacer{} ฆานสมฺผสฺโส.\csromanspacer{} ชิวฺหาสมฺผสฺโส.\csromanspacer{} กายสมฺผสฺโส.\csromanspacer{} มโนสมฺผสฺโส นิจฺโจ วา อนิจฺโจ วาติ.\csromanspacer{} อนิจฺโจ ภนฺเต.\csromanspacer{} เอวํ ปสฺสํ ราหุล สุตวา อริยสาวโก จกฺขุสมฺผสฺสสฺมิมฺปิ นิพฺพินฺทติ ฯเปฯ.\csromanspacer{} โสตสมฺผสฺสสฺมิมฺปิ นิพฺพินฺทติ.\csromanspacer{} ฆานสมฺผสฺสสฺมิมฺปิ นิพฺพินฺทติ.\csromanspacer{} ชิวฺหาสมฺผสฺสสฺมิมฺปิ นิพฺพินฺทติ.\csromanspacer{} กายสมฺผสฺสสฺมิมฺปิ นิพฺพินฺทติ.\csromanspacer{} มโนสมฺผสฺสสฺมิมฺปิ นิพฺพินฺทติ.\csromanspacer{} นิพฺพินฺทํ วิรชฺชติ ฯเปฯ ปชานาตีติ.\csromanspacer{}\csromanspacer{}จตุตฺถํ.}

\csromanmatikaanchor{mtk.231}
\csromantocmarkref{subparagraph}{5. อจิรปกฺกนฺตสุตฺต}{mtk.231}
\bookmark[dest={mtk.231},level=5]{5. อจิรปกฺกนฺตสุตฺต}
\csromanheader{6}{5. เวทนาสุตฺต}
\proseitem{192}{สาวตฺถิยํ วิหรติ.\csromanspacer{} ตํ กิํ มญฺญสิ ราหุล, จกฺขุสมฺผสฺสชา เวทนา นิจฺจา วา อนิจฺจา วาติ.\csromanspacer{} อนิจฺจา ภนฺเต.\csromanspacer{} โสตสมฺผสฺสชา เวทนา ฯเปฯ.\csromanspacer{} ฆานสมฺผสฺสชา เวทนา.\csromanspacer{} ชิวฺหาสมฺผสฺสชา เวทนา.\csromanspacer{} กายสมฺผสฺสชา เวทนา.\csromanspacer{} มโนสมฺผสฺสชา เวทนา นิจฺจา วา อนิจฺจา วาติ.}

\csromanheader{2}{7. ราหุลสํยุตฺต}
\prose{\csromanfolio{441}อนิจฺจา ภนฺเต.\csromanspacer{} เอวํ ปสฺสํ ราหุล สุตวา อริยสาวโก จกฺขุสมฺผสฺสชาย เวทนายปิ นิพฺพินฺทติ ฯเปฯ.\csromanspacer{} โสต.\csromanspacer{} ฆาน.\csromanspacer{} ชิวฺหา.\csromanspacer{} กาย.\csromanspacer{} มโนสมฺผสฺสชาย เวทนายปิ นิพฺพินฺทติ ฯเปฯ ปชานาตีติ.\csromanspacer{}\csromanspacer{}ปญฺจมํ.}

\csromanheader{6}{6. สญฺญาสุตฺต}
\proseitem{193}{สาวตฺถิยํ วิหรติ.\csromanspacer{} ตํ กิํ มญฺญสิ ราหุล, รูปสญฺญา นิจฺจา วา อนิจฺจา วาติ.\csromanspacer{} อนิจฺจา ภนฺเต.\csromanspacer{} สทฺทสญฺญา ฯเปฯ.\csromanspacer{} คนฺธสญฺญา.\csromanspacer{} รสสญฺญา.\csromanspacer{} โผฏฺฐพฺพสญฺญา.\csromanspacer{} ธมฺมสญฺญา นิจฺจา วา อนิจฺจา วาติ.\csromanspacer{} อนิจฺจา ภนฺเต.\csromanspacer{} เอวํ ปสฺสํ ราหุล สุตวา อริยสาวโก รูปสญฺญายปิ นิพฺพินฺทติ ฯเปฯ.\csromanspacer{} สทฺทสญฺญายปิ นิพฺพินฺทติ.\csromanspacer{} คนฺธสญฺญายปิ นิพฺพินฺทติ.\csromanspacer{} รสสญฺญายปิ นิพฺพินฺทติ.\csromanspacer{} โผฏฺฐพฺพสญฺญายปิ นิพฺพินฺทติ.\csromanspacer{} ธมฺมสญฺญายปิ นิพฺพินฺทติ ฯเปฯ ปชานาตีติ.\csromanspacer{}\csromanspacer{}ฉฏฺฐํ.}

\csromanheader{6}{7. สญฺเจตนาสุตฺต}
\proseitem{194}{สาวตฺถิยํ วิหรติ.\csromanspacer{} ตํ กิํ มญฺญสิ ราหุล, รูปสญฺเจตนา นิจฺจา วา อนิจฺจา วาติ.\csromanspacer{} อนิจฺจา ภนฺเต.\csromanspacer{} สทฺทสญฺเจตนา ฯเปฯ.\csromanspacer{} คนฺธสญฺเจตนา.\csromanspacer{} รสสญฺเจตนา.\csromanspacer{} โผฏฺฐพฺพสญฺเจตนา.\csromanspacer{} ธมฺมสญฺเจตนา นิจฺจา วา อนิจฺจา วาติ.\csromanspacer{} อนิจฺจา ภนฺเต.\csromanspacer{} เอวํ ปสฺสํ ราหุล สุตวา อริยสาวโก รูปสญฺเจตนายปิ นิพฺพินฺทติ ฯเปฯ.\csromanspacer{} สทฺทสญฺเจตนายปิ นิพฺพินฺทติ.\csromanspacer{} คนฺธสญฺเจตนายปิ นิพฺพินฺทติ.\csromanspacer{} รสสญฺเจตนายปิ นิพฺพินฺทติ.\csromanspacer{} โผฏฺฐพฺพสญฺเจตนายปิ นิพฺพินฺทติ.\csromanspacer{} ธมฺมสญฺเจตนายปิ นิพฺพินฺทติ ฯเปฯ ปชานาตีติ.\csromanspacer{}\csromanspacer{}สตฺตมํ.}

\csromanheader{6}{8. ตณฺหาสุตฺต}
\proseitem{195}{สาวตฺถิยํ วิหรติ.\csromanspacer{} ตํ กิํ มญฺญสิ ราหุล, รูปตณฺหา นิจฺจา วา อนิจฺจา วาติ.\csromanspacer{} อนิจฺจา ภนฺเต.\csromanspacer{} สทฺทตณฺหา ฯเปฯ.\csromanspacer{} คนฺธตณฺหา.\csromanspacer{} รสตณฺหา.\csromanspacer{} โผฏฺฐพฺพตณฺหา.\csromanspacer{} ธมฺมตณฺหา นิจฺจา วา อนิจฺจา วาติ.\csromanspacer{} อนิจฺจา ภนฺเต.\csromanspacer{} เอวํ ปสฺสํ ราหุล สุตวา อริยสาวโก รูปตณฺหายปิ นิพฺพินฺทติ ฯเปฯ.\csromanspacer{} สทฺทตณฺหายปิ นิพฺพินฺทติ.\csromanspacer{} คนฺธตณฺหายปิ นิพฺพินฺทติ.\csromanspacer{} รสตณฺหายปิ นิพฺพินฺทติ.\csromanspacer{} โผฏฺฐพฺพตณฺหายปิ นิพฺพินฺทติ.\csromanspacer{} ธมฺมตณฺหายปิ นิพฺพินฺทติ ฯเปฯ ปชานาตีติ.\csromanspacer{}\csromanspacer{}อฏฺฐมํ.}

\gambhira{นิทานวคฺคสํยุตฺตปาฬิ}
\csromanmatikaanchor{mtk.232}
\csromantocmarkref{subparagraph}{6. ปญฺจรถสตสุตฺต}{mtk.232}
\bookmark[dest={mtk.232},level=5]{6. ปญฺจรถสตสุตฺต}
\csromanheader{6}{9. ธาตุสุตฺต}
\csromanmatikaanchor{mtk.233}
\csromanmatikaanchor{mtk.248}
\csromantocmarkref{subparagraph}{7. มาตุสุตฺต}{mtk.233}
\bookmark[dest={mtk.233},level=5]{7. มาตุสุตฺต}
\proseitem{196}{\csromanfolio{442}สาวตฺถิยํ วิหรติ.\csromanspacer{} ตํ กิํ มญฺญสิ ราหุล, ปถวีธาตุ นิจฺจา วา อนิจฺจา วาติ.\csromanspacer{} อนิจฺจา ภนฺเต.\csromanspacer{} อาโปธาตุ ฯเปฯ.\csromanspacer{} เตโชธาตุ.\csromanspacer{} วาโยธาตุ.\csromanspacer{} อากาสธาตุ.\csromanspacer{} วิญฺญาณธาตุ นิจฺจา วา อนิจฺจา วาติ.\csromanspacer{} อนิจฺจา ภนฺเต.\csromanspacer{} เอวํ ปสฺสํ ราหุล สุตวา อริยสาวโก ปถวีธาตุยาปิ นิพฺพินฺทติ ฯเปฯ.\csromanspacer{} อาโปธาตุยาปิ นิพฺพินฺทติ.\csromanspacer{} เตโชธาตุยาปิ นิพฺพินฺทติ.\csromanspacer{} วาโยธาตุยาปิ นิพฺพินฺทติ.\csromanspacer{} อากาสธาตุยาปิ นิพฺพินฺทติ.\csromanspacer{} วิญฺญาณธาตุยาปิ นิพฺพินฺทติ ฯเปฯ ปชานาตีติ.\csromanspacer{}\csromanspacer{}นวมํ.}

\csromanheader{6}{10. ขนฺธสุตฺต}
\proseitem{197}{สาวตฺถิยํ วิหรติ.\csromanspacer{} ตํ กิํ มญฺญสิ ราหุล, รูปํ นิจฺจํ วา อนิจฺจํ วาติ.\csromanspacer{} อนิจฺจํ ภนฺเต.\csromanspacer{} เวทนา ฯเปฯ.\csromanspacer{} สญฺญา.\csromanspacer{} สงฺขารา.\csromanspacer{} วิญฺญาณํ นิจฺจํ วา อนิจฺจํ วาติ.\csromanspacer{} อนิจฺจํ ภนฺเต.\csromanspacer{} เอวํ ปสฺสํ ราหุล สุตวา อริยสาวโก รูปสฺมิมฺปิ นิพฺพินฺทติ ฯเปฯ.\csromanspacer{} เวทนายปิ นิพฺพินฺทติ.\csromanspacer{} สญฺญายปิ นิพฺพินฺทติ.\csromanspacer{} สงฺขาเรสุปิ นิพฺพินฺทติ.\csromanspacer{} วิญฺญาณสฺมิมฺปิ นิพฺพินฺทติ, นิพฺพินฺทํ วิรชฺชติ, วิราคา วิมุจฺจติ, วิมุตฺตสฺมิํ “วิมุตฺตมฺ”อิติ ญาณํ โหติ, “ขีณา ชาติ, วุสิตํ พฺรหฺมจริยํ, กตํ กรณียํ, นาปรํ อิตฺถตฺตายา”ติ ปชานาตีติ.\csromanspacer{}\csromanspacer{}ทสมํ.}

\vspace{\tipitakaparskip}
\csromancenter{ปฐโม วคฺโค.}
\titlehead{\textbf{ตสฺสุทฺทานํ}}
\csromangathameasure{จกฺขุ รูปญฺจ วิญฺญาณํ, สมฺผสฺโส เวทนาย จ. \\ สญฺญา สญฺเจตนา ตณฺหา, ธาตุ ขนฺเธน เต ทสาติ.}
\csromangathagroup{\csromangathastanza{จกฺขุ รูปญฺจ วิญฺญาณํ, สมฺผสฺโส เวทนาย จ. \\ สญฺญา สญฺเจตนา ตณฺหา, ธาตุ ขนฺเธน เต ทสาติ.}}

\csromanmatikaanchor{mtk.234}
\csromantocmarkref{subparagraph}{8-13. ปิตุสุตฺตาทิฉกฺก}{mtk.234}
\bookmark[dest={mtk.234},level=5]{8-13. ปิตุสุตฺตาทิฉกฺก}
\csromantocmarkref{section}{อุทฺทานคาถา}{mtk.235}
\bookmark[dest={mtk.235},level=1]{อุทฺทานคาถา}
\chapterhead{2. ทุติยวคฺค}
\csromanheader{6}{1. จกฺขุสุตฺต}
\proseitem{198}{เอวํ เม สุตํ\csromandash{}เอกํ สมยํ ภควา สาวตฺถิยํ วิหรติ.\csromanspacer{} อถ โข อายสฺมา ราหุโล เยน ภควา เตนุปสงฺกมิ, อุปสงฺกมิตฺวา ภควนฺตํ อภิวาเทตฺวา เอกมนฺตํ นิสีทิ, เอกมนฺตํ นิสินฺนํ โข อายสฺมนฺตํ}

\csromanheader{2}{7. ราหุลสํยุตฺต}
\prose{\csromanfolio{443}ราหุลํ ภควา เอตทโวจ “ตํ กิํ มญฺญสิ ราหุล, จกฺขุํ นิจฺจํ วา อนิจฺจํ วาติ.\csromanspacer{} อนิจฺจํ ภนฺเต.\csromanspacer{} ยํ ปนานิจฺจํ, ทุกฺขํ วา ตํ สุขํ วาติ.\csromanspacer{} ทุกฺขํ ภนฺเต.\csromanspacer{} ยํ ปนานิจฺจํ ทุกฺขํ วิปริณามธมฺมํ, กลฺลํ นุ ตํ สมนุปสฺสิตุํ “เอตํ มม, เอโสหมสฺมิ, เอโส เม อตฺตา”ติ.\csromanspacer{} โน เหตํ ภนฺเต.\csromanspacer{} โสตํ ฯเปฯ.\csromanspacer{} ฆานํ.\csromanspacer{} ชิวฺหา.\csromanspacer{} กาโย.\csromanspacer{} มโน นิจฺโจ วา อนิจฺโจ วาติ.\csromanspacer{} อนิจฺโจ ภนฺเต.\csromanspacer{} ยํ ปนานิจฺจํ, ทุกฺขํ วา ตํ สุขํ วาติ.\csromanspacer{} ทุกฺขํ ภนฺเต.\csromanspacer{} ยํ ปนานิจฺจํ ทุกฺขํ วิปริณามธมฺมํ, กลฺลํ นุ ตํ สมนุปสฺสิตุํ “เอตํ มม, เอโสหมสฺมิ, เอโส เม อตฺตา”ติ.\csromanspacer{} โน เหตํ ภนฺเต.\csromanspacer{} เอวํ ปสฺสํ ราหุล สุตวา อริยสาวโก จกฺขุสฺมิมฺปิ นิพฺพินฺทติ ฯเปฯ.\csromanspacer{} โสตสฺมิมฺปิ นิพฺพินฺทติ.\csromanspacer{} ฆานสฺมิมฺปิ นิพฺพินฺทติ.\csromanspacer{} ชิวฺหายปิ นิพฺพินฺทติ.\csromanspacer{} กายสฺมิมฺปิ นิพฺพินฺทติ.\csromanspacer{} มนสฺมิมฺปิ นิพฺพินฺทติ.\csromanspacer{} นิพฺพินฺทํ วิรชฺชติ, วิราคา วิมุจฺจติ, วิมุตฺตสฺมิํ “วิมุตฺตมฺ”อิติ ญาณํ โหติ, “ขีณา ชาติ, วุสิตํ พฺรหฺมจริยํ, กตํ กรณียํ, นาปรํ อิตฺถตฺตายา”ติ ปชานาตีติ. (เอเตน เปยฺยาเลน ทส สุตฺตนฺตา กาตพฺพา.) .\csromanspacer{} ปฐมํ.}

\csromanheader{6}{2-10. รูปาทิสุตฺตนวก}
\csromanmatikaanchor{mtk.236}
\csromantocmarknopage{subparagraph}{7. ราหุลสํยุตฺต}{mtk.236}
\bookmark[dest={mtk.236},level=5]{7. ราหุลสํยุตฺต}
\proseitem{199}{สาวตฺถิยํ วิหรติ.\csromanspacer{} ตํ กิํ มญฺญสิ ราหุล, รูปา นิจฺจา วา อนิจฺจา วาติ.\csromanspacer{} อนิจฺจา ภนฺเต ฯเปฯ.\csromanspacer{} สทฺทา.\csromanspacer{} คนฺธา.\csromanspacer{} รสา.\csromanspacer{} โผฏฺฐพฺพา.\csromanspacer{} ธมฺมา.}

\prose{จกฺขุวิญฺญาณํ ฯเปฯ.\csromanspacer{} โสตวิญฺญาณํ.\csromanspacer{} ฆานวิญฺญาณํ.\csromanspacer{} ชิวฺหาวิญฺญาณํ.\csromanspacer{} กายวิญฺญาณํ.\csromanspacer{} มโนวิญฺญาณํ.}

\csromanmatikaanchor{mtk.237}
\csromanmatikaanchor{mtk.238}
\csromantocmarknopage{subparagraph}{1. ปฐมวคฺค}{mtk.237}
\bookmark[dest={mtk.237},level=5]{1. ปฐมวคฺค}
\csromantocmarkref{subparagraph}{1. จกฺขุสุตฺต}{mtk.238}
\bookmark[dest={mtk.238},level=5]{1. จกฺขุสุตฺต}
\prose{จกฺขุสมฺผสฺโส ฯเปฯ.\csromanspacer{} โสตสมฺผสฺโส.\csromanspacer{} ฆานสมฺผสฺโส.\csromanspacer{} ชิวฺหาสมฺผสฺโส.\csromanspacer{} กายสมฺผสฺโส.\csromanspacer{} มโนสมฺผสฺโส.}

\prose{จกฺขุสมฺผสฺสชา เวทนา ฯเปฯ.\csromanspacer{} โสตสมฺผสฺสชา เวทนา.\csromanspacer{} ฆานสมฺผสฺสชา เวทนา.\csromanspacer{} ชิวฺหาสมฺผสฺสชา เวทนา.\csromanspacer{} กายสมฺผสฺสชา เวทนา.\csromanspacer{} มโนสมฺผสฺสชา เวทนา.}

\prose{รูปสญฺญา ฯเปฯ.\csromanspacer{} สทฺทสญฺญา.\csromanspacer{} คนฺธสญฺญา.\csromanspacer{} รสสญฺญา.\csromanspacer{} โผฏฺฐพฺพสญฺญา.\csromanspacer{} ธมฺมสญฺญา.}

\prose{รูปสญฺเจตนา ฯเปฯ.\csromanspacer{} สทฺทสญฺเจตนา.\csromanspacer{} คนฺธสญฺเจตนา.\csromanspacer{} รสสญฺเจตนา.\csromanspacer{} โผฏฺฐพฺพสญฺเจตนา.\csromanspacer{} ธมฺมสญฺเจตนา.}

\gambhira{นิทานวคฺคสํยุตฺตปาฬิ}
\prose{\csromanfolio{444}รูปตณฺหา ฯเปฯ.\csromanspacer{} สทฺทตณฺหา.\csromanspacer{} คนฺธตณฺหา.\csromanspacer{} รสตณฺหา.\csromanspacer{} โผฏฺฐพฺพตณฺหา.\csromanspacer{} ธมฺมตณฺหา.}

\csromanmatikaanchor{mtk.239}
\csromantocmarkref{subparagraph}{2. รูปสุตฺต}{mtk.239}
\bookmark[dest={mtk.239},level=5]{2. รูปสุตฺต}
\prose{ปถวีธาตุ ฯเปฯ.\csromanspacer{} อาโปธาตุ.\csromanspacer{} เตโชธาตุ.\csromanspacer{} วาโยธาตุ.\csromanspacer{} อากาสธาตุ.\csromanspacer{} วิญฺญาณธาตุ.}

\prose{รูปํ ฯเปฯ.\csromanspacer{} เวทนา.\csromanspacer{} สญฺญา.\csromanspacer{} สงฺขารา.\csromanspacer{} วิญฺญาณํ นิจฺจํ วา อนิจฺจํ วาติ.\csromanspacer{} อนิจฺจํ ภนฺเต ฯเปฯ.\csromanspacer{} เอวํ ปสฺสํ ราหุล ฯเปฯ.\csromanspacer{} นาปรํ อิตฺถตฺตายาติ ปชานาตีติ.\csromanspacer{}\csromanspacer{}ทสมํ.}

\csromanmatikaanchor{mtk.240}
\csromantocmarkref{subparagraph}{3. วิญฺญาณสุตฺต}{mtk.240}
\bookmark[dest={mtk.240},level=5]{3. วิญฺญาณสุตฺต}
\csromanheader{6}{11. อนุสยสุตฺต}
\proseitem{200}{สาวตฺถิยํ วิหรติ.\csromanspacer{} อถ โข อายสฺมา ราหุโล เยน ภควา เตนุปสงฺกมิ, อุปสงฺกมิตฺวา ภควนฺตํ อภิวาเทตฺวา เอกมนฺตํ นิสีทิ, เอกมนฺตํ นิสินฺโน โข อายสฺมา ราหุโล ภควนฺตํ เอตทโวจ “กถํ นุ โข ภนฺเต ชานโต กถํ ปสฺสโต อิมสฺมิํ จ สวิญฺญาณเก กาเย พหิทฺธา จ สพฺพนิมิตฺเตสุ อหงฺการมมงฺการมานานุสยา น โหนฺตี”ติ.\csromanspacer{} ยํ กิญฺจิ ราหุล รูปํ อตีตานาคตปจฺจุปฺปนฺนํ อชฺฌตฺตํ วา พหิทฺธา วา โอฬาริกํ วา สุขุมํ วา หีนํ วา ปณีตํ วา ยํ ทูเร สนฺติเก วา, สพฺพํ รูปํ “เนตํ มม, เนโสหมสฺมิ, น เมโส อตฺตา”ติ เอวเมตํ ยถาภูตํ สมฺมปฺปญฺญาย ปสฺสติ.\csromanspacer{} ยา กาจิ เวทนา ฯเปฯ.\csromanspacer{} ยา กาจิ สญฺญา.\csromanspacer{} เย เกจิ สงฺขารา.\csromanspacer{} ยํ กิญฺจิ วิญฺญาณํ อตีตานาคตปจฺจุปฺปนํ อชฺฌตฺตํ วา พหิทฺธา วา โอฬาริกํ วา สุขุมํ วา หีนํ วา ปณีตํ วา ยํ ทูเร สนฺติเก วา, สพฺพํ วิญฺญาณํ “เนตํ มม, เนโสหมสฺมิ, น เมโส อตฺตา”ติ เอวเมตํ ยถาภูตํ สมฺมปฺปญฺญาย ปสฺสติ.\csromanspacer{} เอวํ โข ราหุล ชานโต เอวํ ปสฺสโต อิมสฺมิํ จ สวิญฺญาณเก กาเย พหิทฺธา จ สพฺพนิมิตฺเตสุ อหงฺการมมงฺการมานานุสยา น โหนฺตีติ.\csromanspacer{}\csromanspacer{}เอกาทสมํ.}

\csromanmatikaanchor{mtk.241}
\csromantocmarkref{subparagraph}{4. สมฺผสฺสสุตฺต}{mtk.241}
\bookmark[dest={mtk.241},level=5]{4. สมฺผสฺสสุตฺต}
\csromanheader{6}{12. อปคตสุตฺต}
\proseitem{201}{สาวตฺถินิทานํ.\csromanspacer{} อถ โข อายสฺมา ราหุโล เยน ภควา เตนุปสงฺกมิ, อุปสงฺกมิตฺวา ภควนฺตํ อภิวาเทตฺวา เอกมนฺตํ นิสีทิ, เอกมนฺตํ นิสินฺโน โข อายสฺมา ราหุโล ภควนฺตํ เอตทโวจ “กถํ นุ โข ภนฺเต ชานโต กถํ ปสฺสโต อิมสฺมิํ จ สวิญฺญาณเก}

\csromanmatikaanchor{mtk.242}
\csromantocmarkref{subparagraph}{5. เวทนาสุตฺต}{mtk.242}
\bookmark[dest={mtk.242},level=5]{5. เวทนาสุตฺต}
\csromanheader{2}{7. ราหุลสํยุตฺต}
\csromanmatikaanchor{mtk.243}
\csromanmatikaanchor{mtk.254}
\csromantocmarkref{subparagraph}{6. สญฺญาสุตฺต}{mtk.243}
\bookmark[dest={mtk.243},level=5]{6. สญฺญาสุตฺต}
\prose{\csromanfolio{445}กาเย พหิทฺธา จ สพฺพนิมิตฺเตสุ อหงฺการมมงฺการมานาปคตํ มานสํ โหติ วิธา สมติกฺกนฺตํ สนฺตํ สุวิมุตฺตนฺ”ติ.\csromanspacer{} ยํ กิญฺจิ ราหุล รูปํ อตีตานาคตปจฺจุปฺปนฺนํ อชฺฌตฺตํ วา พหิทฺธา วา โอฬาริกํ วา สุขุมํ วา หีนํ วา ปณีตํ วา ยํ ทูเร สนฺติเก วา, สพฺพํ รูปํ “เนตํ มม, เนโสหมสฺมิ, น เมโส อตฺตา”ติ เอวเมตํ ยถาภูตํ สมฺมปฺปญฺญาย ทิสฺวา อนุปาทา วิมุตฺโต โหติ.}

\prose{ยา กาจิ เวทนา ฯเปฯ.\csromanspacer{} ยา กาจิ สญฺญา.\csromanspacer{} เย เกจิ สงฺขารา.\csromanspacer{} ยํ กิญฺจิ วิญฺญาณํ อตีตานาคตปจฺจุปฺปนฺนํ อชฺฌตฺตํ วา พหิทฺธา วา โอฬาริกํ วา สุขุมํ วา หีนํ วา ปณีตํ วา ยํ ทูเร สนฺติเก วา, สพฺพํ วิญฺญาณํ “เนตํ มม, เนโสหมสฺมิ, น เมโส อตฺตา”ติ เอวเมตํ ยถาภูตํ สมฺมปฺปญฺญาย ทิสฺวา อนุปาทาวิมุตฺโต โหติ.\csromanspacer{} เอวํ โข ราหุล ชานโต เอวํ ปสฺสโต อิมสฺมิํ จ สวิญฺญาณเก กาเย พหิทฺธา จ สพฺพนิมิตฺเตสุ อหงฺการมมงฺการมานาปคตํ มานสํ โหติ วิธา สมติกฺกนฺตํ สนฺตํ สุวิมุตฺตนฺ”ติ.\csromanspacer{}\csromanspacer{}ทฺวาทสมํ.}

\vspace{\tipitakaparskip}
\csromancenter{ทุติโย วคฺโค.}
\titlehead{\textbf{ตสฺสุทฺทานํ}}
\csromanmatikaanchor{mtk.244}
\csromantocmarkref{subparagraph}{7. สญฺเจตนาสุตฺต}{mtk.244}
\bookmark[dest={mtk.244},level=5]{7. สญฺเจตนาสุตฺต}
\csromangathameasure{จกฺขุ รูปญฺจ วิญฺญาณํ, สมฺผสฺโส เวทนาย จ. \\ สญฺญา สญฺเจตนา ตณฺหา, ธาตุ ขนฺเธน เต ทส. \\ อนุสยํ อปคตญฺเจว, วคฺโค เตน ปวุจฺจตีติ. \\ \textbf{ราหุลสํยุตฺตํ สมตฺตํ.}}
\csromangathagroup{\csromangathastanza{จกฺขุ รูปญฺจ วิญฺญาณํ, สมฺผสฺโส เวทนาย จ. \\ สญฺญา สญฺเจตนา ตณฺหา, ธาตุ ขนฺเธน เต ทส.}\csromangathastanza{อนุสยํ อปคตญฺเจว, วคฺโค เตน ปวุจฺจตีติ. \\ \textbf{ราหุลสํยุตฺตํ สมตฺตํ.}}}

\csromanheader{6}{8. ลกฺขณสํยุตฺต}
\csromanmatikaanchor{mtk.245}
\csromantocmarkref{subparagraph}{8. ตณฺหาสุตฺต}{mtk.245}
\bookmark[dest={mtk.245},level=5]{8. ตณฺหาสุตฺต}
\chapterhead{1. ปฐมวคฺค}
\csromanheader{6}{1. อฏฺฐิสุตฺต}
\proseitem{202}{\csromanfolio{446}เอวํ เม สุตํ\csromandash{}เอกํ สมยํ ภควา ราชคเห วิหรติ เวฬุวเน กลนฺทกนิวาเป.\csromanspacer{} เตน โข ปน สมเยน อายสฺมา จ ลกฺขโณ อายสฺมา จ มหาโมคฺคลฺลาโน\footnote{มหาโมคฺคลาโน (ก)} คิชฺฌกูเฏ ปพฺพเต วิหรนฺติ.\csromanspacer{} อถ โข อายสฺมา มหาโมคฺคลฺลาโน ปุพฺพณฺหสมยํ นิวาเสตฺวา ปตฺตจีวรมาทาย เยนายสฺมา ลกฺขโณ เตนุปสงฺกมิ, อุปสงฺกมิตฺวา อายสฺมนฺตํ ลกฺขณํ เอตทโวจ “อายามาวุโส\footnote{เอหิ อาวุโส (สฺยา, กํ, ก)} ลกฺขณ ราชคหํ ปิณฺฑาย ปวิสิสฺสามา”ติ. “เอวมาวุโส”ติ โข อายสฺมา ลกฺขโณ อายสฺมโต มหาโมคฺคลฺลานสฺส ปจฺจสฺโสสิ.\csromanspacer{} อถ โข อายสฺมา มหาโมคฺคลฺลาโน คิชฺฌกูฏา ปพฺพตา โอโรหนฺโต อญฺญตรสฺมิํ ปเทเส สิตํ ปาตฺวากาสิ.\csromanspacer{} อถ โข อายสฺมา ลกฺขโณ อายสฺมนฺตํ มหาโมคฺคลฺลานํ เอตทโวจ “โก นุ โข อาวุโส โมคฺคลฺลาน เหตุ, โก ปจฺจโย สิตสฺส ปาตุกมฺมายา”ติ.\csromanspacer{} อกาโล โข อาวุโส ลกฺขณ เอตสฺส ปญฺหสฺส, ภควโต มํ สนฺติเก เอตํ ปญฺหํ ปุจฺฉาติ.}

\csromanmatikaanchor{mtk.246}
\csromantocmarkref{subparagraph}{9. ธาตุสุตฺต}{mtk.246}
\bookmark[dest={mtk.246},level=5]{9. ธาตุสุตฺต}
\prose{อถ โข อายสฺมา จ ลกฺขโณ อายสฺมา จ มหาโมคฺคลฺลาโน ราชคเห ปิณฺฑาย จริตฺวา ปจฺฉาภตฺตํ ปิณฺฑปาตปฏิกฺกนฺตา เยน ภควา เตนุปสงฺกมิํสุ, อุปสงฺกมิตฺวา ภควนฺตํ อภิวาเทตฺวา เอกมนฺตํ นิสีทิํสุ, เอกมนฺตํ นิสินฺโน โข อายสฺมา ลกฺขโณ อายสฺมนฺตํ มหาโมคฺคลฺลานํ เอตทโวจ “อิธายสฺมา มหาโมคฺคลฺลาโน คิชฺฌกูฏา ปพฺพตา โอโรหนฺโต อญฺญตรสฺมิํ ปเทเส สิตํ ปาตฺวากาสิ, โก นุ โข อาวุโส โมคฺคลฺลาน เหตุ, โก ปจฺจโย สิตสฺส ปาตุกมฺมายา”ติ.}

\prose{อิธาหํ อาวุโส คิชฺฌกูฏา ปพฺพตา โอโรหนฺโต อทฺทสํ อฏฺฐิกสงฺขลิกํ เวหาสํ คจฺฉนฺติํ, ตเมนํ คิชฺฌาปิ กากาปิ กุลลาปิ}

\csromanmatikaanchor{mtk.247}
\csromantocmarkref{subparagraph}{10. ขนฺธสุตฺต}{mtk.247}
\bookmark[dest={mtk.247},level=5]{10. ขนฺธสุตฺต}
\csromantocmarkref{section}{อุทฺทานคาถา}{mtk.248}
\bookmark[dest={mtk.248},level=1]{อุทฺทานคาถา}
\vspace{\tipitakaparskip}
\csromancenter{ลกฺขณสํยุตฺต อนุปติตฺวา อนุปติตฺวา ผาสุฬนฺตริกาหิ วิตุเทนฺติ วิตจฺเฉนฺติ วิราเชนฺติ\footnote{วิตุเทนฺติ (สี), วิตจฺเฉนฺติ วิภเชนฺติ (อิ, ก)}.\csromanspacer{} สา สุทํ อฏฺฏสฺสรํ กโรติ, ตสฺส มยฺหํ อาวุโส เอตทโหสิ “อจฺฉริยํ วต โภ, อพฺภุตํ วต โภ, เอวรูโปปิ นาม สตฺโต ภวิสฺสติ, เอวรูโปปิ นาม ยกฺโข ภวิสฺสติ, เอวรูโปปิ นาม อตฺตภาวปฏิลาโภ ภวิสฺสตี”ติ.}
\prose{\csromanfolio{447}อถ โข ภควา ภิกฺขู อามนฺเตสิ “จกฺขุภูตา วต ภิกฺขเว สาวกา วิหรนฺติ, ญาณภูตา วต ภิกฺขเว สาวกา วิหรนฺติ, ยตฺร หิ นาม สาวโก เอวรูปํ ญสฺสติ วา ทกฺขติ วา สกฺขิํ วา กริสฺสติ, ปุพฺเพว เม โส ภิกฺขเว สตฺโต ทิฏฺโฐ อโหสิ, อปิ จาหํ น พฺยากาสิํ, อหญฺเจตํ\footnote{อหเมเวตํ (สี)} พฺยากเรยฺยํ, ปเร จ เม\footnote{ปเร เม (สี) 4. วิราเชนฺติ (สี, สฺยา, กํ), วิภเชนฺติ (อิ, ก)} น สทฺทเหยฺยุํ.\csromanspacer{} เย เม น สทฺทเหยฺยุํ, เตสํ ตํ อสฺส ทีฆรตฺตํ อหิตาย ทุกฺขาย.\csromanspacer{} เอโส ภิกฺขเว สตฺโต อิมสฺมิํ เยว ราชคเห โคฆาตโก อโหสิ, โส ตสฺส กมฺมสฺส วิปาเกน พหูนิ วสฺสานิ พหูนิ วสฺสสตานิ พหูนิ วสฺสสหสฺสานิ พหูนิ วสฺสสตสหสฺสานิ นิรเย ปจฺจิตฺวา ตสฺเสว กมฺมสฺส วิปากาวเสเสน เอวรูปํ อตฺตภาวปฏิลาภํ ปฏิสํเวทยตี”ติ. (สพฺเพสํ สุตฺตนฺตานํ เอเสว เปยฺยาโล.) .\csromanspacer{} ปฐมํ.}

\csromanheader{6}{2. เปสิสุตฺต}
\proseitem{203}{อิธาหํ อาวุโส คิชฺฌกูฏา ปพฺพตา โอโรหนฺโต อทฺทสํ มํสเปสิํ เวหาสํ คจฺฉนฺติํ, ตเมนํ คิชฺฌาปิ กากาปิ กุลลาปิ อนุปติตฺวา อนุปติตฺวา วิตจฺเฉนฺติ วิราเชนฺติ4.\csromanspacer{} สา สุทํ อฏฺฏสฺสรํ กโรติ ฯเปฯ.\csromanspacer{} เอโส ภิกฺขเว สตฺโต อิมสฺมิํเยว ราชคเห โคฆาตโก อโหสิ ฯเปฯ.\csromanspacer{}\csromanspacer{}ทุติยํ.}

\csromanheader{6}{3. ปิณฺฑสุตฺต}
\proseitem{204}{อิธาหํ อาวุโส คิชฺฌกูฏา ปพฺพตา โอโรหนฺโต อทฺทสํ มํสปิณฺฑํ เวหาสํ คจฺฉนฺตํ, ตเมนํ คิชฺฌาปิ กากาปิ กุลลาปิ อนุปติตฺวา อนุปติตฺวา วิตจฺเฉนฺติ วิราเชนฺติ.\csromanspacer{} สา สุทํ อฏฺฏสฺสรํ กโรติ ฯเปฯ.\csromanspacer{} เอโส ภิกฺขเว สตฺโต อิมสฺมิํเยว ราชคเห สากุณิโก อโหสิ ฯเปฯ.\csromanspacer{}\csromanspacer{}ตติยํ.}

\csromanmatikaanchor{mtk.249}
\csromanmatikaanchor{mtk.250}
\csromantocmarknopage{subparagraph}{2. ทุติยวคฺค}{mtk.249}
\bookmark[dest={mtk.249},level=5]{2. ทุติยวคฺค}
\csromantocmarkref{subparagraph}{1. จกฺขุสุตฺต}{mtk.250}
\bookmark[dest={mtk.250},level=5]{1. จกฺขุสุตฺต}
\gambhira{นิทานวคฺคสํยุตฺตปาฬิ}
\csromanheader{6}{4. นิจฺฉวิสุตฺต}
\proseitem{205}{\csromanfolio{448}อิธาหํ อาวุโส คิชฺฌกูฏา ปพฺพตา โอโรหนฺโต อทฺทสํ นิจฺฉวิํ ปุริสํ เวหาสํ คจฺฉนฺตํ, ตเมนํ คิชฺฌาปิ กากาปิ กุสลาปิ อนุปติตฺวา อนุปติตฺวา วิตจฺเฉนฺติ วิราเชนฺติ.\csromanspacer{} โส สุทํ อฏฺฏสฺสรํ กโรติ ฯเปฯ.\csromanspacer{} เอโส ภิกฺขเว สตฺโต อิมสฺมิํเยว ราชคเห โอรพฺภิโก อโหสิ ฯเปฯ.\csromanspacer{}\csromanspacer{}จตุตฺถํ.}

\csromanheader{6}{5. อสิโลมสุตฺต}
\csromanmatikaanchor{mtk.251}
\csromantocmarkref{subparagraph}{2-10. รูปาทิสุตฺตนวก}{mtk.251}
\bookmark[dest={mtk.251},level=5]{2-10. รูปาทิสุตฺตนวก}
\proseitem{206}{อิธาหํ อาวุโส คิชฺฌกูฏา ปพฺพตา โอโรหนฺโต อทฺทสํ อสิโลมํ ปุริสํ เวหาสํ คจฺฉนฺตํ, ตสฺส เต อสี อุปฺปติตฺวา อุปฺปติตฺวา ตสฺเสว กาเย นิปตนฺติ.\csromanspacer{} โส สุทํ อฏฺฏสฺสรํ กโรติ ฯเปฯ.\csromanspacer{} เอโส ภิกฺขเว สตฺโต อิมสฺมิํเยว ราชคเห สูกริโก อโหสิ ฯเปฯ.\csromanspacer{} ปญฺจมํ.}

\csromanheader{6}{6. สตฺติสุตฺต}
\proseitem{207}{อิธาหํ อาวุโส คิชฺฌกูฏา ปพฺพตา โอโรหนฺโต อทฺทสํ สตฺติโลมํ ปุริสํ เวหาสํ คจฺฉนฺตํ, ตสฺส ตา สตฺติโย อุปฺปติตฺวา อุปฺปติตฺวา ตสฺเสว กาเย นิปตนฺติ.\csromanspacer{} โส สุทํ อฏฺฏสฺสรํ กโรติ ฯเปฯ.\csromanspacer{} เอโส ภิกฺขเว สตฺโต อิมสฺมิํเยว ราชคเห มาควิโก อโหสิ ฯเปฯ.\csromanspacer{} ฉฏฺฐํ.}

\csromanheader{6}{7. อุสุโลมสุตฺต}
\proseitem{208}{อิธาหํ อาวุโส คิชฺฌกูฏา ปพฺพตา โอโรหนฺโต อทฺทสํ อุสุโลมํ ปุริสํ เวหาสํ คจฺฉนฺตํ, ตสฺส เต อุสู อุปฺปติตฺวา อุปฺปติตฺวา ตสฺเสว กาเย นิปตนฺติ.\csromanspacer{} โส สุทํ อฏฺฏสฺสรํ กโรติ ฯเปฯ.\csromanspacer{} เอโส ภิกฺขเว สตฺโต อิมสฺมิํเยว ราชคเห การณิโก อโหสิ ฯเปฯ.\csromanspacer{} สตฺตมํ.}

\csromanheader{6}{8. สูจิโลมสุตฺต}
\proseitem{209}{อิธาหํ อาวุโส คิชฺฌกูฏา ปพฺพตา โอโรหนฺโต อทฺทสํ สูจิโลมํ ปุริสํ เวหาสํ คจฺฉนฺตํ, ตสฺส ตา สูจิโย อุปฺปติตฺวา อุปฺปติตฺวา ตสฺเสว กาเย นิปตนฺติ.\csromanspacer{} โส สุทํ อฏฺฏสฺสรํ กโรติ ฯเปฯ.}

\vspace{\tipitakaparskip}
\csromancenter{ลกฺขณสํยุตฺต เอโส ภิกฺขเว สตฺโต อิมสฺมิํเยว ราชคเห สูโต\footnote{สารถิโก (ก, วินเยปิ)} อโหสิ ฯเปฯ.\csromanspacer{}\csromanspacer{}อฏฺฐมํ.}
\csromanheader{6}{9. ทุติยสูจิโลมสุตฺต}
\csromanmatikaanchor{mtk.252}
\csromanmatikaanchor{mtk.267}
\csromantocmarkref{subparagraph}{11. อนุสยสุตฺต}{mtk.252}
\bookmark[dest={mtk.252},level=5]{11. อนุสยสุตฺต}
\proseitem{210}{\csromanfolio{449}อิธาหํ อาวุโส คิชฺฌกูฏา ปพฺพตา โอโรหนฺโต อทฺทสํ สูจิโลมํ ปุริสํ เวหาสํ คจฺฉนฺตํ, ตสฺส ตา สูจิโย สีเส ปวิสิตฺวา มุขโต นิกฺขมนฺติ, มุเข ปวิสิตฺวา อุรโต นิกฺขมนฺติ, อุเร ปวิสิตฺวา อุทรโต นิกฺขมนฺติ, อุทเร ปวิสิตฺวา อูรูหิ นิกฺขมนฺติ, อูรูสุ ปวิสิตฺวา ชงฺฆาติ นิกฺขมนฺติ, ชงฺฆาสุ ปวิสิตฺวา ปาเทหิ นิกฺขมนฺติ.\csromanspacer{} โส สุทํ อฏฺฏสฺสรํ กโรติ ฯเปฯ.\csromanspacer{} เอโส ภิกฺขเว สตฺโต อิมสฺมิํเยว ราชคเห สูจโก อโหสิ ฯเปฯ.\csromanspacer{}\csromanspacer{}นวมํ.}

\csromanheader{6}{10. กุมฺภณฺฑสุตฺต}
\proseitem{211}{อิธาหํ อาวุโส คิชฺฌกูฏา ปพฺพตา โอโรหนฺโต อทฺทสํ กุมฺภณฺฑํ ปุริสํ เวหาสํ คจฺฉนฺตํ, โส คจฺฉนฺโตปิ เตว อณฺเฑ ขนฺเธ อาโรเปตฺวา คจฺฉติ, นิสีทนฺโตปิ เตเสฺวว อณฺเฑสุ นิสีทติ, ตเมนํ คิชฺฌาปิ กากาปิ กุลลาปิ อนุปติตฺวา อนุปติตฺวา วิตจฺเฉนฺติ วิราเชนฺติ.\csromanspacer{} โส สุทํ อฏฺฏสฺสรํ กโรติ ฯเปฯ.\csromanspacer{} เอโส ภิกฺขเว สตฺโต อิมสฺมิํเยว ราชคเห คามกูฏโก อโหสิ ฯเปฯ.\csromanspacer{}\csromanspacer{}ทสมํ.}

\vspace{\tipitakaparskip}
\csromancenter{ปฐโม วคฺโค.}
\csromanmatikaanchor{mtk.253}
\csromantocmarkref{subparagraph}{12. อปคตสุตฺต}{mtk.253}
\bookmark[dest={mtk.253},level=5]{12. อปคตสุตฺต}
\csromantocmarkref{section}{อุทฺทานคาถา}{mtk.254}
\bookmark[dest={mtk.254},level=1]{อุทฺทานคาถา}
\titlehead{\textbf{ตสฺสุทฺทานํ}}
\csromangathameasure{อฏฺฐิ เปสิ อุโภ คาวฆาตกา, \\ ปิณฺโฑ สากุณิโย นิจฺฉโวรพฺภิ. \\ อสิ สูกริโก สตฺติมาควิ, \\ อุสุ การณิโก สูจิ สารถิ. \\ โย จ สิพฺพิยติ สูจโก หิ โส, \\ อณฺฑภาริ อหุ คามกูฏโกติ.}
\csromangathagroup{\csromangathastanza{อฏฺฐิ เปสิ อุโภ คาวฆาตกา, \\ ปิณฺโฑ สากุณิโย นิจฺฉโวรพฺภิ. \\ อสิ สูกริโก สตฺติมาควิ, \\ อุสุ การณิโก สูจิ สารถิ. \\ โย จ สิพฺพิยติ สูจโก หิ โส, \\ อณฺฑภาริ อหุ คามกูฏโกติ.}}

\gambhira{นิทานวคฺคสํยุตฺตปาฬิ}
\chapterhead{2. ทุติยวคฺค}
\csromanheader{6}{1. สสีสกสุตฺต}
\proseitem{212}{\csromanfolio{450}เอวํ เม สุตํ\csromandash{}เอกํ สมยํ ราชคเห เวฬุวเน.\csromanspacer{} อิธาหํ อาวุโส คิชฺฌกูฏา ปพฺพตา โอโรหนฺโต อทฺทสํ ปุริสํ คูถกูเป สสีสกํ นิมุคฺคํ ฯเปฯ.\csromanspacer{} เอโส ภิกฺขเว สตฺโต อิมสฺมิํเยว ราชคเห ปารทาริโก อโหสิ ฯเปฯ.\csromanspacer{}\csromanspacer{}ปฐมํ.}

\csromanheader{6}{2. คูถขาทสุตฺต}
\proseitem{213}{อิธาหํ อาวุโส คิชฺฌกูฏา ปพฺพตา โอโรหนฺโต อทฺทสํ ปุริสํ คูถกูเป นิมุคฺคํ อุโภหิ หตฺเถหิ คูถํ ขาทนฺตํ ฯเปฯ.\csromanspacer{} เอโส ภิกฺขเว สตฺโต อิมสฺมิํเยว ราชคเห ทุฏฺฐพฺราหฺมโณ อโหสิ.\csromanspacer{} โส กสฺสปสฺส สมฺมาสมฺพุทฺธสฺส ปาวจเน ภิกฺขุสํฆํ ภตฺเตน นิมนฺเตตฺวา โทณิโย\footnote{โทณิยา (สฺยา, กํ, อิ, ก)} คูถสฺส ปูราเปตฺวา เอตทโวจ “อโห โภนฺโต ยาวทตฺถํ ภุญฺชนฺตุ เจว หรนฺตุ จา”ติ ฯเปฯ.\csromanspacer{}\csromanspacer{}ทุติยํ.}

\csromanheader{6}{3. นิจฺฉวิตฺถิสุตฺต}
\csromanmatikaanchor{mtk.255}
\csromantocmarknopage{subparagraph}{8. ลกฺขณสํยุตฺต}{mtk.255}
\bookmark[dest={mtk.255},level=5]{8. ลกฺขณสํยุตฺต}
\proseitem{214}{อิธาหํ อาวุโส คิชฺฌกูฏา ปพฺพตา โอโรหนฺโต อทฺทสํ นิจฺฉวิํ อิตฺถิํ เวหาสํ คจฺฉนฺติํ, ตเมนํ คิชฺฌาปิ กากาปิ กุลลาปิ อนุปติตฺวา อนุปติตฺวา วิตจฺเฉนฺติ วิราเชนฺติ.\csromanspacer{} สา สุทํ อฏฺฏสฺสรํ กโรติ ฯเปฯ.\csromanspacer{} เอสา ภิกฺขเว อิตฺถี อิมสฺมิํเยว ราชคเห อติจารินี อโหสิ ฯเปฯ.\csromanspacer{}\csromanspacer{}ตติยํ.}

\csromanheader{6}{4. มงฺคุลิตฺถิสุตฺต}
\csromanmatikaanchor{mtk.256}
\csromanmatikaanchor{mtk.257}
\csromantocmarknopage{subparagraph}{1. ปฐมวคฺค}{mtk.256}
\bookmark[dest={mtk.256},level=5]{1. ปฐมวคฺค}
\csromantocmarkref{subparagraph}{1. อฏฺฐิสุตฺต}{mtk.257}
\bookmark[dest={mtk.257},level=5]{1. อฏฺฐิสุตฺต}
\proseitem{215}{อิธาหํ อาวุโส คิชฺฌกูฏา ปพฺพตา โอโรหนฺโต อทฺทสํ อิตฺถิํ ทุคฺคนฺธํ มงฺคุลิํ เวหาสํ คจฺฉนฺติํ, ตเมนํ คิชฺฌาปิ กากาปิ กุลลาปิ อนุปติตฺวา อนุปติตฺวา วิตจฺเฉนฺติ วิราเชนฺติ.\csromanspacer{} สา สุทํ อฏฺฏสฺสรํ กโรติ ฯเปฯ.\csromanspacer{} เอสา ภิกฺขเว อิตฺถี อิมสฺมิํเยว ราชคเห อิกฺขณิกา อโหสิ ฯเปฯ.\csromanspacer{}\csromanspacer{}จตุตฺถํ.}

\csromanheader{6}{8. ลกฺขณสํยุตฺต \textbf{5. โอกิลินีสุตฺต}}
\proseitem{216}{\csromanfolio{451}อิธาหํ อาวุโส คิชฺฌกูฏา ปพฺพตา โอโรหนฺโต อทฺทสํ อิตฺถิํ อุปฺปกฺกํ โอกิลินิํ โอกิรินิํ เวหาสํ คจฺฉนฺติํ.\csromanspacer{} สา สุทํ อฏฺฏสฺสรํ กโรติ ฯเปฯ.\csromanspacer{} เอสา ภิกฺขเว อิตฺถี กลิงฺคสฺส รญฺโญ อคฺคมเหสี อโหสิ, สา อิสฺสาปกตา สปตฺติํ องฺคารกฏาเหน โอกิริ ฯเปฯ.\csromanspacer{}\csromanspacer{}ปญฺจมํ.}

\csromanheader{6}{6. อสีสกสุตฺต}
\proseitem{217}{อิธาหํ อาวุโส คิชฺฌกูฏา ปพฺพตา โอโรหนฺโต อทฺทสํ อสีสกํ กพนฺธํ\footnote{กวนฺธํ (สี, อิ)} เวหาสํ คจฺฉนฺตํ, ตสฺส อุเร อกฺขีนิ เจว โหนฺติ มุขญฺจ.\csromanspacer{} ตเมนํ คิชฺฌาปิ กากาปิ กุลลาปิ อนุปติตฺวา อนุปติตฺวา วิตจฺเฉนฺติ วิราเชนฺติ.\csromanspacer{} โส สุทํ อฏฺฏสฺสรํ กโรติ ฯเปฯ.\csromanspacer{} เอโส ภิกฺขเว สตฺโต อิมสฺมิํเยว ราชคเห หาริโก นาม โจรฆาตโก อโหสิ ฯเปฯ.\csromanspacer{}\csromanspacer{}ฉฏฺฐํ.}

\csromanheader{6}{7. ปาปภิกฺขุสุตฺต}
\csromanmatikaanchor{mtk.258}
\csromantocmarkref{subparagraph}{2. เปสิสุตฺต}{mtk.258}
\bookmark[dest={mtk.258},level=5]{2. เปสิสุตฺต}
\proseitem{218}{อิธาหํ อาวุโส คิชฺฌกูฏา ปพฺพตา โอโรหนฺโต อทฺทสํ ภิกฺขุํ เวหาสํ คจฺฉนฺตํ, ตสฺส สงฺฆาฏิปิ อาทิตฺตา สมฺปชฺชลิตา สโชติภูตา\footnote{สญฺโชติภูตา (สฺยา, กํ)}, ปตฺโตปิ อาทิตฺโต สมฺปชฺชลิโต สโชติภูโต, กายพนฺธนมฺปิ อาทิตฺตํ สมฺปชฺชลิตํ สโชติภูตํ, กาโยปิ อาทิตฺโต สมฺปชฺชลิโต สโชติภูโต.\csromanspacer{} โส สุทํ อฏฺฏสฺสรํ กโรติ ฯเปฯ.\csromanspacer{} เอโส ภิกฺขเว ภิกฺขุ กสฺสปสฺส สมฺมาสมฺพุทฺธสฺส ปาวจเน ปาปภิกฺขุ อโหสิ ฯเปฯ.\csromanspacer{}\csromanspacer{}สตฺตมํ.}

\csromanheader{6}{8. ปาปภิกฺขุนีสุตฺต}
\csromanmatikaanchor{mtk.259}
\csromantocmarkref{subparagraph}{3. ปิณฺฑสุตฺต}{mtk.259}
\bookmark[dest={mtk.259},level=5]{3. ปิณฺฑสุตฺต}
\proseitem{219}{อทฺทสํ ภิกฺขุนิํ เวหาสํ คจฺฉนฺติํ, ตสฺสา สงฺฆาฏิปิ อาทิตฺตา ฯเปฯ ปาปภิกฺขุนี อโหสิ ฯเปฯ.\csromanspacer{}\csromanspacer{}อฏฺฐมํ.}

\gambhira{นิทานวคฺคสํยุตฺตปาฬิ}
\csromanheader{6}{9. ปาปสิกฺขมานสุตฺต}
\csromanmatikaanchor{mtk.260}
\csromantocmarkref{subparagraph}{4. นิจฺฉวิสุตฺต}{mtk.260}
\bookmark[dest={mtk.260},level=5]{4. นิจฺฉวิสุตฺต}
\proseitem{220}{\csromanfolio{452}อทฺทสํ สิกฺขมานํ เวหาสํ คจฺฉนฺติํ, ตสฺสา สงฺฆาฏิปิ อาทิตฺตา ฯเปฯ ปาปสิกฺขมานา อโหสิ ฯเปฯ.\csromanspacer{}\csromanspacer{}นวมํ.}

\csromanheader{6}{10. ปาปสามเณรสุตฺต}
\csromanmatikaanchor{mtk.261}
\csromantocmarkref{subparagraph}{5. อสิโลมสุตฺต}{mtk.261}
\bookmark[dest={mtk.261},level=5]{5. อสิโลมสุตฺต}
\proseitem{221}{อทฺทสํ สามเณรํ เวหาสํ คจฺฉนฺตํ, ตสฺส สงฺฆาฏิปิ อาทิตฺตา ฯเปฯ ปาปสามเณโร อโหสิ ฯเปฯ.\csromanspacer{}\csromanspacer{}ทสมํ.}

\csromanheader{6}{11. ปาปสามเณรีสุตฺต}
\csromanmatikaanchor{mtk.262}
\csromantocmarkref{subparagraph}{6. สตฺติสุตฺต}{mtk.262}
\bookmark[dest={mtk.262},level=5]{6. สตฺติสุตฺต}
\proseitem{222}{อิธาหํ อาวุโส คิชฺฌกูฏา ปพฺพตา โอโรหนฺโต อทฺทสํ สามเณริํ เวหาสํ คจฺฉนฺติํ, ตสฺสา สงฺฆาฏิปิ อาทิตฺตา สมฺปชฺชลิตา สโชติภูตา, ปตฺโตปิ อาทิตฺโต สมฺปชฺชลิโต สโชติภูโต, กายพนฺธนมฺปิ อาทิตฺตํ สมฺปชฺชลิตํ สโชติภูตํ, กาโยปิ อาทิตฺโต สมฺปชฺชลิโต สโชติภูโต.\csromanspacer{} สา สุทํ อฏฺฏสฺสรํ กโรติ.\csromanspacer{} ตสฺส มยฺหํ อาวุโส เอตทโหสิ “อจฺฉริยํ วต โภ, อพฺภุตํ วต โภ, เอวรูโปปิ นาม สตฺโต ภวิสฺสติ, เอวรูโปปิ นาม ยกฺโข ภวิสฺสติ, เอวรูโปปิ นาม อตฺตภาวปฏิลาโภ ภวิสฺสตี”ติ.}

\prose{อถ โข ภควา ภิกฺขู อามนฺเตสิ “จกฺขุภูตา วต ภิกฺขเว สาวกา วิหรนฺติ, ญาณภูตา วต ภิกฺขเว สาวกา วิหรนฺติ, ยตฺร หิ นาม สาวโก เอวรูปํ ญสฺสติ วา ทกฺขติ วา สกฺขิํ วา กริสฺสติ.\csromanspacer{} ปุพฺเพว เม สา ภิกฺขเว สามเณรี ทิฏฺฐา อโหสิ.\csromanspacer{} อปิ จาหํ น พฺยากาสิํ, อหญฺเจตํ พฺยากเรยฺยํ, ปเร จ เม น สทฺทเหยฺยุํ.\csromanspacer{} เย เม น สทฺทเหยฺยุํ, เตสํ ตํ อสฺส ทีฆรตฺตํ อหิตาย ทุกฺขาย.\csromanspacer{} เอสา ภิกฺขเว สามเณรี กสฺสปสฺส สมฺมาสมฺพุทฺธสฺส ปาวจเน ปาปสามเณรี อโหสิ, สา ตสฺส กมฺมสฺส วิปาเกน พหูนิ วสฺสานิ พหูนิ วสฺสสตานิ พหูนิ วสฺสสหสฺสานิ พหูนิ วสฺสสตสหสฺสานิ นิรเย ปจฺจิตฺวา ตสฺเสว กมฺมสฺส วิปากาวเสเสน เอวรูปํ อตฺตภาวปฏิลาภํ ปฏิสํเวทยตี”ติ.\csromanspacer{}\csromanspacer{}เอกาทสมํ.}

\csromanmatikaanchor{mtk.263}
\csromantocmarkref{subparagraph}{7. อุสุโลมสุตฺต}{mtk.263}
\bookmark[dest={mtk.263},level=5]{7. อุสุโลมสุตฺต}
\vspace{\tipitakaparskip}
\csromancenter{ทุติโย วคฺโค.}
\titlehead{8. ลกฺขณสํยุตฺต \textbf{ตสฺสุทฺทานํ}}
\csromanmatikaanchor{mtk.264}
\csromanmatikaanchor{mtk.280}
\csromantocmarkref{subparagraph}{8. สูจิโลมสุตฺต}{mtk.264}
\bookmark[dest={mtk.264},level=5]{8. สูจิโลมสุตฺต}
\csromangathameasure{กูเป นิมุคฺโค หิ โส ปารทาริโก, \\ คูถขาทิ อหุ ทุฏฺฐพฺราหฺมโณ. \\ นิจฺฉวิตฺถิ อติจารินี อหุ, \\ มงฺคุลิตฺถิ อหุ อิกฺขณิตฺถิกา. \\ โอกิลินิ สปตฺตงฺคาโรกิริ, \\ สีสจฺฉินฺโน อหุ โจรฆาตโก. \\ ภิกฺขุ ภิกฺขุนี สิกฺขมานา, \\ สามเณโร อถ สามเณริกา. \\ กสฺสปสฺส วินยสฺมิํ ปพฺพชฺชํ, \\ ปาปกมฺมํ กริํสุ ตาวเทติ. \\ \textbf{ลกฺขณสํยุตฺตํ สมตฺตํ.}}
\csromangathagroup{\csromangathastanza{\csromanfolio{453}กูเป นิมุคฺโค หิ โส ปารทาริโก, \\ คูถขาทิ อหุ ทุฏฺฐพฺราหฺมโณ. \\ นิจฺฉวิตฺถิ อติจารินี อหุ, \\ มงฺคุลิตฺถิ อหุ อิกฺขณิตฺถิกา.}\csromangathastanza{โอกิลินิ สปตฺตงฺคาโรกิริ, \\ สีสจฺฉินฺโน อหุ โจรฆาตโก. \\ ภิกฺขุ ภิกฺขุนี สิกฺขมานา, \\ สามเณโร อถ สามเณริกา.}\csromangathastanza{กสฺสปสฺส วินยสฺมิํ ปพฺพชฺชํ, \\ ปาปกมฺมํ กริํสุ ตาวเทติ. \\ \textbf{ลกฺขณสํยุตฺตํ สมตฺตํ.}}}

\csromanmatikaanchor{mtk.265}
\csromantocmarkref{subparagraph}{9. ทุติยสูจิโลมสุตฺต}{mtk.265}
\bookmark[dest={mtk.265},level=5]{9. ทุติยสูจิโลมสุตฺต}
\csromanheader{6}{9. โอปมฺมสํยุตฺต}
\csromanheader{6}{1. กูฏสุตฺต}
\csromanmatikaanchor{mtk.266}
\csromantocmarkref{subparagraph}{10. กุมฺภณฺฑสุตฺต}{mtk.266}
\bookmark[dest={mtk.266},level=5]{10. กุมฺภณฺฑสุตฺต}
\csromantocmarkref{section}{อุทฺทานคาถา}{mtk.267}
\bookmark[dest={mtk.267},level=1]{อุทฺทานคาถา}
\proseitem{223}{\csromanfolio{454}เอวํ เม สุตํ\csromandash{}เอกํ สมยํ ภควา สาวตฺถิยํ วิหรติ เชตวเน อนาถปิณฺฑิกสฺส อาราเม.\csromanspacer{} ตตฺร โข ภควา ฯเปฯ เอตทโวจ\csromandash{} เสยฺยถาปิ ภิกฺขเว กูฏาคารสฺส ยา กาจิ โคปานสิโย, สพฺพา ตา กูฏงฺคมา กูฏสโมสรณา, กูฏสมุคฺฆาตา สพฺพา ตา สมุคฺฆาตํ คจฺฉนฺติ.\csromanspacer{} เอวเมว โข ภิกฺขเว เย เกจิ อกุสลา ธมฺมา, สพฺเพ เต อวิชฺชามูลกา อวิชฺชาสโมสรณา, อวิชฺชาสมุคฺฆาตา สพฺเพ เต สมุคฺฆาตํ คจฺฉนฺติ.\csromanspacer{} ตสฺมาติห ภิกฺขเว เอวํ สิกฺขิตพฺพํ “อปฺปมตฺตา วิหริสฺสามา”ติ.\csromanspacer{} เอวํ หิ โว ภิกฺขเว สิกฺขิตพฺพนฺติ.\csromanspacer{}\csromanspacer{}ปฐมํ.}

\csromanheader{6}{2. นขสิขสุตฺต}
\proseitem{224}{สาวตฺถิยํ วิหรติ.\csromanspacer{} อถ โข ภควา ปริตฺตํ นขสิขายํ ปํสุํ อาโรเปตฺวา ภิกฺขู อามนฺเตสิ “ตํ กิํ มญฺญถ ภิกฺขเว, กตมํ นุ โข พหุตรํ, โย วายํ\footnote{โย จายํ (พหูสุ)} มยา ปริตฺโต นขสิขายํ ปํสุ อาโรปิโต, อยํ วา\footnote{ยา จายํ (สฺยา, ก)} มหาปถวี”ติ.\csromanspacer{} เอตเทว ภนฺเต พหุตรํ ยทิทํ มหาปถวี, อปฺปมตฺตโก ยํ ภควตา ปริตฺโต นขสิขายํ ปํสุ อาโรปิโต, สงฺขมฺปิ น อุเปติ อุปนิธิมฺปิ น อุเปติ กลภาคมฺปิ น อุเปติ มหาปถวิํ อุปนิธาย ภควตา ปริตฺโต นขสิขายํ ปํสุ อาโรปิโตติ.\csromanspacer{} เอวเมว โข ภิกฺขเว อปฺปกา เต สตฺตา, เย มนุสฺเสสุ ปจฺจาชายนฺติ, อถ โข เอเตเยว พหุตรา สตฺตา, เย อญฺญตฺร มนุสฺเสหิ ปจฺจาชายนฺติ.\csromanspacer{} ตสฺมาติห ภิกฺขเว เอวํ สิกฺขิตพฺพํ “อปฺปมตฺตา วิหริสฺสามา”ติ.\csromanspacer{} เอวํ หิ โว ภิกฺขเว สิกฺขิตพฺพนฺติ.\csromanspacer{}\csromanspacer{}ทุติยํ.}

\csromanheader{6}{3. กุลสุตฺต}
\proseitem{225}{สาวตฺถิยํ วิหรติ.\csromanspacer{} เสยฺยถาปิ ภิกฺขเว ยานิ กานิจิ กุลานิ พหุตฺถิกานิ อปฺปปุริสานิ, ตานิ สุปฺปธํสิยานิ โหนฺติ โจเรหิ}

\vspace{\tipitakaparskip}
\csromancenter{โอปมฺมสํยุตฺต กุมฺภตฺเถนเกหิ.\csromanspacer{} เอวเมว โข ภิกฺขเว ยสฺส กสฺสจิ ภิกฺขุโน เมตฺตาเจโตวิมุตฺติ อภาวิตา อพหุลีกตา, โส สุปฺปธํสิโย โหติ อมนุสฺเสหิ.\csromanspacer{} เสยฺยถาปิ ภิกฺขเว ยานิ กานิจิ กุลานิ อปฺปิตฺถิกานิ พหุปุริสานิ, ตานิ ทุปฺปธํสิยานิ โหนฺติ โจเรหิ กุมฺภตฺเถนเกหิ.\csromanspacer{} เอวเมว โข ภิกฺขเว ยสฺส กสฺสจิ ภิกฺขุโน เมตฺตาเจโตวิมุตฺติ ภาวิตา พหุลีกตา, โส ทุปฺปธํสิโย โหติ อมนุสฺเสหิ.\csromanspacer{} ตสฺมาติห ภิกฺขเว เอวํ สิกฺขิตพฺพํ “เมตฺตา โน เจโตวิมุตฺติ ภาวิตา ภวิสฺสติ พหุลีกตา ยานีกตา วตฺถุกตา อนุฏฺฐิตา ปริจิตา สุสมารทฺธา”ติ.\csromanspacer{} เอวํ หิ โว ภิกฺขเว สิกฺขิตพฺพนฺติ.\csromanspacer{}\csromanspacer{}ตติยํ.}
\csromanheader{6}{4. โอกฺขาสุตฺต}
\csromanmatikaanchor{mtk.268}
\csromanmatikaanchor{mtk.269}
\csromantocmarknopage{subparagraph}{2. ทุติยวคฺค}{mtk.268}
\bookmark[dest={mtk.268},level=5]{2. ทุติยวคฺค}
\csromantocmarkref{subparagraph}{1. สสีสกสุตฺต}{mtk.269}
\bookmark[dest={mtk.269},level=5]{1. สสีสกสุตฺต}
\proseitem{226}{\csromanfolio{455}สาวตฺถิยํ วิหรติ.\csromanspacer{} โย ภิกฺขเว ปุพฺพณฺหสมยํ โอกฺขาสตํ ทานํ ทเทยฺย, โย มชฺฌนฺหิกสมยํ โอกฺขาสตํ ทานํ ทเทยฺย, โย สายนฺหสมยํ โอกฺขาสตํ ทานํ ทเทยฺย.\csromanspacer{} โย วา ปุพฺพณฺหสมยํ อนฺตมโส คทฺทุหนมตฺตมฺปิ เมตฺตจิตฺตํ ภาเวยฺย, โย วา มชฺฌนฺหิกสมยํ อนฺตมโส คทฺทุหนมตฺตมฺปิ เมตฺตจิตฺตํ ภาเวยฺย, โย วา สายนฺหสมยํ อนฺตมโส คทฺทุหนมตฺตมฺปิ เมตฺตจิตฺตํ ภาเวยฺย, อิทํ ตโต มหปฺผลตรํ.\csromanspacer{} ตสฺมาติห ภิกฺขเว เอวํ สิกฺขิตพฺพํ “เมตฺตา โน เจโตวิมุตฺติ ภาวิตา ภวิสฺสติ พหุลีกตา ยานีกตา วตฺถุกตา อนุฏฺฐิตา ปริจิตา สุสมารทฺธา”ติ.\csromanspacer{} เอวํ หิ โว ภิกฺขเว สิกฺขิตพฺพนฺติ.\csromanspacer{}\csromanspacer{}จตุตฺถํ.}

\csromanheader{6}{5. สตฺติสุตฺต}
\csromanmatikaanchor{mtk.270}
\csromantocmarkref{subparagraph}{2. คูถขาทสุตฺต}{mtk.270}
\bookmark[dest={mtk.270},level=5]{2. คูถขาทสุตฺต}
\proseitem{227}{สาวตฺถิยํ วิหรติ.\csromanspacer{} เสยฺยถาปิ ภิกฺขเว สตฺติ ติณฺหผลา, อถ ปุริโส อาคจฺเฉยฺย “อหํ อิมํ สตฺติํ ติณฺหผลํ ปาณินา วา มุฏฺฐินา วา ปฏิเลณิสฺสามิ ปฏิโกฏฺฏิสฺสามิ ปฏิวฏฺเฏสฺสามี”ติ.\csromanspacer{} ตํ กิํ มญฺญถ ภิกฺขเว, ภพฺโพ นุ โข โส ปุริโส อมุํ สตฺติํ ติณฺหผลํ ปาณินา วา มุฏฺฐินา วา ปฏิเลเณตุํ ปฏิโกฏฺเฏตุํ ปฏิวฏฺเฏตุนฺติ.\csromanspacer{} โน เหตํ ภนฺเต.\csromanspacer{} ตํ กิสฺส เหตุ, อสุ หิ ภนฺเต สตฺติ ติณฺหผลา น สุกรา ปาณินา วา มุฏฺฐินา วา ปฏิเลเณตุํ ปฏิโกฏฺเฏตุํ ปฏิวฏฺเฏตุํ.\csromanspacer{} ยาวเทว จ ปน โส ปุริโส กิลมถสฺส วิฆาตสฺส ภาคี อสฺสาติ.}

\gambhira{นิทานวคฺคสํยุตฺตปาฬิ}
\csromanmatikaanchor{mtk.271}
\csromantocmarkref{subparagraph}{3. นิจฺฉวิตฺถิสุตฺต}{mtk.271}
\bookmark[dest={mtk.271},level=5]{3. นิจฺฉวิตฺถิสุตฺต}
\prose{\csromanfolio{456}เอวเมว โข ภิกฺขเว ยสฺส กสฺสจิ ภิกฺขุโน เมตฺตาเจโตวิมุตฺติ ภาวิตา พหุลีกตา ยานีกตา วตฺถุกตา อนุฏฺฐิตา ปริจิตา สุสมารทฺธา.\csromanspacer{} ตสฺส เจ อมนุสฺโส จิตฺตํ ขิปิตพฺพํ มญฺเญยฺย.\csromanspacer{} อถ โข เสฺวว อมนุสฺโส กิลมถสฺส วิฆาตสฺส ภาคี อสฺส.\csromanspacer{} ตสฺมาติห ภิกฺขเว เอวํ สิกฺขิตพฺพํ “เมตฺตา โน เจโตวิมุตฺติ ภาวิตา ภวิสฺสติ พหุลีกตา ยานีกตา วตฺถุกตา อนุฏฺฐิตา ปริจิตา สุสมารทฺธา”ติ.\csromanspacer{} เอวํ หิ โว ภิกฺขเว สิกฺขิตพฺพนฺติ.\csromanspacer{}\csromanspacer{}ปญฺจมํ.}

\csromanheader{6}{6. ธนุคฺคหสุตฺต}
\csromanmatikaanchor{mtk.272}
\csromantocmarkref{subparagraph}{4. มงฺคุลิตฺถิสุตฺต}{mtk.272}
\bookmark[dest={mtk.272},level=5]{4. มงฺคุลิตฺถิสุตฺต}
\proseitem{228}{สาวตฺถิยํ วิหรติ.\csromanspacer{} เสยฺยถาปิ ภิกฺขเว จตฺตาโร ทฬฺหธมฺมา ธนุคฺคหา สุสิกฺขิตา กตหตฺถา กตูปาสนา จตุทฺทิสา ฐิตา อสฺสุ, อถ ปุริโส อาคจฺเฉยฺย “อหํ อิเมสํ จตุนฺนํ ทฬฺหธมฺมานํ ธนุคฺคหานํ สุสิกฺขิตานํ กตหตฺถานํ กตูปาสนานํ จตุทฺทิสา กณฺเฑ ขิตฺเต อปฺปติฏฺฐิเต ปถวิยํ คเหตฺวา อาหริสฺสามี”ติ.\csromanspacer{} ตํ กิํ มญฺญถ ภิกฺขเว “ชวโน ปุริโส ปรเมน ชเวน สมนฺนาคโต”ติ อลํ วจนายาติ.}

\prose{เอกสฺส เจปิ ภนฺเต ทฬฺหธมฺมสฺส ธนุคฺคหสฺส สุสิกฺขิตสฺส กตหตฺถสฺส กตูปาสนสฺส กณฺฑํ ขิตฺตํ อปฺปติฏฺฐิตํ ปถวิยํ คเหตฺวา อาหเรยฺย “ชวโน ปุริโส ปรเมน ชเวน สมนฺนาคโต”ติ อลํ วจนาย, โก ปน วาโท จตุนฺนํ ทฬฺหธมฺมานํ ธนุคฺคหานํ สุสิกฺขิตานํ กตหตฺถานํ กตูปาสนานนฺติ.}

\csromanmatikaanchor{mtk.273}
\csromantocmarkref{subparagraph}{5. โอกิลินีสุตฺต}{mtk.273}
\bookmark[dest={mtk.273},level=5]{5. โอกิลินีสุตฺต}
\prose{ยถา จ ภิกฺขเว ตสฺส ปุริสสฺส ชโว, ยถา จ จนฺทิมสูริยานํ ชโว, ตโต สีฆตโร.\csromanspacer{} ยถา จ ภิกฺขเว ตสฺส ปุริสสฺส ชโว, ยถา จ จนฺทิมสูริยานํ ชโว, ยถา จ ยา เทวตา จนฺทิมสูริยานํ ปุรโต ธาวนฺติ, ตาสํ เทวตานํ ชโว, ( )\footnote{(ตโต สีฆตโร. ยถา จ ภิกฺขเว ตสฺส ปุริสสฺส ชโว, ยถา จ จนฺทิมสุริยานํ ชโว, ยถา จ ยา เทวตา จนฺทิมสุริยานํ ปุรโต ธาวนฺติ, ตาสํ เทวตานํ ชโว,) (สี, สฺยา, กํ)} ตโต สีฆตรํ อายุสงฺขารา ขียนฺติ.\csromanspacer{} ตสฺมาติห ภิกฺขเว เอวํ สิกฺขิตพฺพํ “อปฺปมตฺตา วิหริสฺสามา”ติ.\csromanspacer{} เอวํ หิ โว ภิกฺขเว สิกฺขิตพฺพนฺติ.\csromanspacer{}\csromanspacer{}ฉฏฺฐํ.}

\csromanheader{6}{9. โอปมฺมสํยุตฺต \textbf{7. อาณิสุตฺต}}
\csromanmatikaanchor{mtk.274}
\csromantocmarkref{subparagraph}{6. อสีสกสุตฺต}{mtk.274}
\bookmark[dest={mtk.274},level=5]{6. อสีสกสุตฺต}
\proseitem{229}{\csromanfolio{457}สาวตฺถิยํ วิหรติ.\csromanspacer{} ภูตปุพฺพํ ภิกฺขเว ทสารหานํ อานโก\footnote{อาณโก (สี)} นาม มุทิงฺโค อโหสิ, ตสฺส ทสารหา อานเก ฆฏิเต อญฺญํ อาณิํ โอทหิํสุ, อหุ โข โส ภิกฺขเว สมโย ยํ อานกสฺส มุทิงฺคสฺส โปราณํ โปกฺขรผลกํ อนฺตรธายิ, อาณิสงฺฆาโฏว อวสิสฺสิ.\csromanspacer{} เอวเมว โข ภิกฺขเว ภวิสฺสนฺติ ภิกฺขู อนาคตมทฺธานํ เย เต สุตฺตนฺตา ตถาคตภาสิตา คมฺภีรา คมฺภีรตฺถา โลกุตฺตรา สุญฺญตปฺปฏิสํยุตฺตา, เตสุ ภญฺญมาเนสุ น สุสฺสูสิสฺสนฺติ, น โสตํ โอทหิสฺสนฺติ, น อญฺญา จิตฺตํ อุปฏฺฐาเปสฺสนฺติ, น จ เต ธมฺเม อุคฺคเหตพฺพํ ปริยาปุณิตพฺพํ มญฺญิสฺสนฺติ.}

\prose{เย ปน เต สุตฺตนฺตา กวิกตา กาเวยฺยา จิตฺตกฺขรา จิตฺตพฺยญฺชนา พาหิรกา สาวกภาสิตา, เตสุ ภญฺญมาเนสุ สุสฺสูสิสฺสนฺติ, โสตํ โอทหิสฺสนฺติ, อญฺญา จิตฺตํ อุปฏฺฐาเปสฺสนฺติ, เต จ ธมฺเม อุคฺคเหตพฺพํ ปริยาปุณิตพฺพํ มญฺญิสฺสนฺติ.\csromanspacer{} เอวเมเตสํ ภิกฺขเว สุตฺตนฺตานํ ตถาคตภาสิตานํ คมฺภีรานํ คมฺภีรตฺถานํ โลกุตฺตรานํ สุญฺญตปฺปฏิสํยุตฺตานํ อนฺตรธานํ ภวิสฺสติ, ตสฺมาติห ภิกฺขเว เอวํ สิกฺขิตพฺพํ “เย เต สุตฺตนฺตา ตถาคตภาสิตา คมฺภีรา คมฺภีรตฺถา โลกุตฺตรา สุญฺญตปฺปฏิสํยุตฺตา, เตสุ ภญฺญมาเนสุ สุสฺสูสิสฺสาม, โสตํ โอทหิสฺสาม, อญฺญา จิตฺตํ อุปฏฺฐาเปสฺสาม, เต จ ธมฺเม อุคฺคเหตพฺพํ ปริยาปุณิตพฺพํ มญฺญิสฺสามา”ติ.\csromanspacer{} เอวํ หิ โว ภิกฺขเว สิกฺขิตพฺพนฺติ.\csromanspacer{}\csromanspacer{}สตฺตมํ.}

\csromanmatikaanchor{mtk.275}
\csromantocmarkref{subparagraph}{7. ปาปภิกฺขุสุตฺต}{mtk.275}
\bookmark[dest={mtk.275},level=5]{7. ปาปภิกฺขุสุตฺต}
\csromanheader{6}{8. กลิงฺครสุตฺต}
\proseitem{230}{เอวํ เม สุตํ\csromandash{}เอกํ สมยํ ภควา เวสาลิยํ วิหรติ มหาวเน กูฏาคารสาลายํ.\csromanspacer{} ตตฺร โข ภควา ภิกฺขู อามนฺเตสิ “ภิกฺขโว”ติ. “ภทนฺเต”ติ เต ภิกฺขู ภควโต ปจฺจสฺโสสุํ.\csromanspacer{} ภควา เอตทโวจ\csromandash{}}

\csromanmatikaanchor{mtk.276}
\csromantocmarkref{subparagraph}{8. ปาปภิกฺขุนีสุตฺต}{mtk.276}
\bookmark[dest={mtk.276},level=5]{8. ปาปภิกฺขุนีสุตฺต}
\prose{กลิงฺครูปธานา ภิกฺขเว เอตรหิ ลิจฺฉวี วิหรนฺติ, อปฺปมตฺตา อาตาปิโน อุปาสนสฺมิํ, เตสํ ราชา มาคโธ อชาตสตฺตุ เวเทหิปุตฺโต น ลภติ โอตารํ, น ลภติ อารมฺมณํ.\csromanspacer{} ภวิสฺสนฺติ}

\gambhira{นิทานวคฺคสํยุตฺตปาฬิ}
\prose{\csromanfolio{458}ภิกฺขเว อนาคตมทฺธานํ ลิจฺฉวี สุขุมาลา\footnote{สุกุมาลา (สี), สุขุมา (ก)} มุทุตลุนหตฺถปาทา\footnote{มุทุตลาหตฺถปาทา (สฺยา, กํ)}, เต มุทุกาสุ เสยฺยาสุ ตูลพิพฺโพหนาสุ\footnote{ตูลพิมฺโพหนาสุ (สฺยา, กํ, อิ), ตูลพิมฺโพหนาทีสุ (สี), ตูลพิพฺโพหนาทีสุ (ก)} ยาวสูริยุคฺคมนา เสยฺยํ กปฺปิสฺสนฺติ, เตสํ ราชา มาคโธ อชาตสตฺตุ เวเทหิปุตฺโต ลจฺฉติ โอตารํ, ลจฺฉติ อารมฺมณํ.}

\csromanmatikaanchor{mtk.277}
\csromantocmarkref{subparagraph}{9. ปาปสิกฺขมานสุตฺต}{mtk.277}
\bookmark[dest={mtk.277},level=5]{9. ปาปสิกฺขมานสุตฺต}
\prose{กลิงฺครูปธานา ภิกฺขเว เอตรหิ ภิกฺขู วิหรนฺติ, อปฺปมตฺตา อาตาปิโน ปธานสฺมิํ, เตสํ มาโร ปาปิมา น ลภติ โอตารํ, น ลภติ อารมฺมณํ.\csromanspacer{} ภวิสฺสนฺติ ภิกฺขเว อนาคตมทฺธานํ ภิกฺขู สุขุมา มุทุตลุนหตฺถปาทา, เต มุทุกาสุ เสยฺยาสุ ตูลพิพฺโพหนาสุ ยาวสูริยุคฺคมนา เสยฺยํ กปฺปิสฺสนฺติ, เตสํ มาโร ปาปิมา ลจฺฉติ โอตารํ, ลจฺฉติ อารมฺมณํ.\csromanspacer{} ตสฺมาติห ภิกฺขเว เอวํ สิกฺขิตพฺพํ “กลิงฺครูปธานา วิหริสฺสาม อปฺปมตฺตา อาตาปิโน ปธานสฺมินฺ”ติ.\csromanspacer{} เอวํ หิ โว ภิกฺขเว สิกฺขิตพฺพนฺติ.\csromanspacer{}\csromanspacer{}อฏฺฐมํ.}

\csromanheader{6}{9. นาคสุตฺต}
\csromanmatikaanchor{mtk.278}
\csromantocmarkref{subparagraph}{10. ปาปสามเณรสุตฺต}{mtk.278}
\bookmark[dest={mtk.278},level=5]{10. ปาปสามเณรสุตฺต}
\proseitem{231}{เอวํ เม สุตํ\csromandash{}เอกํ สมยํ ภควา สาวตฺถิยํ วิหรติ เชตวเน อนาถปิณฺฑิกสฺส อาราเม.\csromanspacer{} เตน โข ปน สมเยน อญฺญตโร นโว ภิกฺขุ อติเวลํ กุลานิ อุปสงฺกมติ.\csromanspacer{} ตเมนํ ภิกฺขู เอวมาหํสุ “มายสฺมา อติเวลํ กุลานิ อุปสงฺกมี”ติ.\csromanspacer{} โส ภิกฺขุ ภิกฺขูหิ วุจฺจมาโน เอวมาห “อิเม หิ นาม เถรา ภิกฺขู กุลานิ อุปสงฺกมิตพฺพํ มญฺญิสฺสนฺติ, กิมงฺคํ\footnote{กิมงฺค (สี)} ปนาหนฺ”ติ.}

\prose{อถ โข สมฺพหุลา ภิกฺขู เยน ภควา เตนุปสงฺกมิํสุ, อุปสงฺกมิตฺวา ภควนฺตํ อภิวาเทตฺวา เอกมนฺตํ นิสีทิํสุ, เอกมนฺตํ นิสินฺนา โข เต ภิกฺขู ภควนฺตํ เอตทโวจุํ “อิธ ภนฺเต อญฺญตโร นโว ภิกฺขุ อติเวลํ กุลานิ อุปสงฺกมติ, ตเมนํ ภิกฺขู เอวมาหํสุ ‘มายสฺมา อติเวลํ กุลานิ อุปสงฺกมี’ติ.\csromanspacer{} โส ภิกฺขุ ภิกฺขูหิ วุจฺจมาโน เอวมาห ‘อิเมหิ นาม เถรา ภิกฺขู กุลานิ อุปสงฺกมิตพฺพํ มญฺญิสฺสนฺติ, กิมงฺคํ ปนาหนฺ’ติ”.}

\csromanmatikaanchor{mtk.279}
\csromantocmarkref{subparagraph}{11. ปาปสามเณรีสุตฺต}{mtk.279}
\bookmark[dest={mtk.279},level=5]{11. ปาปสามเณรีสุตฺต}
\csromantocmarkref{section}{อุทฺทานคาถา}{mtk.280}
\bookmark[dest={mtk.280},level=1]{อุทฺทานคาถา}
\csromanheader{2}{9. โอปมฺมสํยุตฺต}
\prose{\csromanfolio{459}ภูตปุพฺพํ ภิกฺขเว อรญฺญายตเน มหาสรสี.\csromanspacer{} ตํ นาคา อุปนิสฺสาย วิหรนฺติ, เต ตํ สรสิํ โอคาเหตฺวา โสณฺฑาย ภิสมุฬาลํ อพฺพุเหตฺวา\footnote{อพฺภูเหตฺวา (ก), อพฺพาหิตฺวา (วิ 3. 308)} สุวิกฺขาลิตํ วิกฺขาเลตฺวา อกทฺทมํ สงฺขาทิตฺวา\footnote{สงฺขริตฺวา (อิ, ก)} อชฺโฌหรนฺติ.\csromanspacer{} เตสํ ตํ วณฺณาย เจว โหติ พลาย จ, น จ ตโตนิทานํ มรณํ วา นิคจฺฉนฺติ มรณมตฺตํ วา ทุกฺขํ.\csromanspacer{} เตสํเยว โข ปน ภิกฺขเว มหานาคานํ อนุสิกฺขมานา ตรุณา ภิงฺกจฺฉาปา ตํ สรสิํ โอคาเหตฺวา โสณฺฑาย ภิสมุฬาลํ อพฺพุเหตฺวา น สุวิกฺขาลิตํ วิกฺขาเลตฺวา สกทฺทมํ อสงฺขาทิตฺวา อชฺโฌหรนฺติ.\csromanspacer{} เตสํ ตํ เนว วณฺณาย โหติ น พลาย, ตโตนิทานํ มรณํ วา นิคจฺฉนฺติ มรณมตฺตํ วา ทุกฺขํ.}

\prose{เอวเมว โข ภิกฺขเว อิธ เถรา ภิกฺขู ปุพฺพณฺหสมยํ นิวาเสตฺวา ปตฺตจีวรมาทาย คามํ วา นิคมํ วา ปิณฺฑาย ปวิสนฺติ.\csromanspacer{} เต ตตฺถ ธมฺมํ ภาสนฺติ, เตสํ คิหี ปสนฺนาการํ กโรนฺติ.\csromanspacer{} เต ตํ ลาภํ อคธิตา อมุจฺฉิตา อนชฺโฌปนฺนา\footnote{อนชฺฌาปนฺนา (สพฺพตฺถ) ม 1 ปาสราสิสุตฺตวณฺณนา โอโลเกตพฺพา.} อาทีนวทสฺสาวิโน นิสฺสรณปญฺญา ปริภุญฺชนฺติ.\csromanspacer{} เตสํ ตํ วณฺณาย เจว โหติ พลาย จ, น จ ตโตนิทานํ มรณํ วา นิคจฺฉนฺติ มรณมตฺตํ วา ทุกฺขํ.\csromanspacer{} เตสํเยว โข ปน ภิกฺขเว เถรานํ ภิกฺขูนํ อนุสิกฺขมานา นวา ภิกฺขู ปุพฺพณฺหสมยํ นิวาเสตฺวา ปตฺตจีวรมาทาย คามํ วา นิคมํ วา ปิณฺฑาย ปวิสนฺติ.\csromanspacer{} เต ตตฺถ ธมฺมํ ภาสนฺติ, เตสํ คิหี ปสนฺนาการํ กโรนฺติ, เต ตํ ลาภํ คธิตา มุจฺฉิตา อชฺโฌปนฺนา อนาทีนวทสฺสาวิโน อนิสฺสรณปญฺญา ปริภุญฺชนฺติ.\csromanspacer{} เตสํ ตํ เนว วณฺณาย โหติ น พลาย, เต ตโตนิทานํ มรณํ วา นิคจฺฉนฺติ มรณมตฺตํ วา ทุกฺขํ.\csromanspacer{} ตสฺมาติห ภิกฺขเว เอวํ สิกฺขิตพฺพํ “อคธิตา อมุจฺฉิตา อนชฺโฌปนฺนา อาทีนวทสฺสาวิโน นิสฺสรณปญฺญา ตํ ลาภํ ปริภุญฺชิสฺสามา”ติ.\csromanspacer{} เอวํ หิ โว ภิกฺขเว สิกฺขิตพฺพนฺติ.\csromanspacer{}\csromanspacer{}นวมํ.}

\csromanheader{6}{10. พิฬารสุตฺต}
\proseitem{232}{สาวตฺถิยํ วิหรติ.\csromanspacer{} เตน โข ปน สมเยน อญฺญตโร ภิกฺขุ อติเวลํ กุเลสุ จาริตฺตํ อาปชฺชติ.\csromanspacer{} ตเมนํ ภิกฺขู เอวมาหํสุ “มายสฺมา อติเวลํ กุเลสุ จาริตฺตํ อาปชฺชี”ติ.\csromanspacer{} โส ภิกฺขุ ภิกฺขูหิ วุจฺจมาโน น วิรมติ.\csromanspacer{} อถ โข สมฺพหุลา ภิกฺขู เยน ภควา}

\gambhira{นิทานวคฺคสํยุตฺตปาฬิ}
\prose{\csromanfolio{460}เตนุปสงฺกมิํสุ, อุปสงฺกมิตฺวา ภควนฺตํ อภิวาเทตฺวา เอกมนฺตํ นิสีทิํสุ, เอกมนฺตํ นิสินฺนา โข เต ภิกฺขู ภควนฺตํ เอตทโวจุํ “อิธ ภนฺเต อญฺญตโร ภิกฺขุ อติเวลํ กุเลสุ จาริตฺตํ อาปชฺชติ, ตเมนํ ภิกฺขู เอวมาหํสุ ‘มายสฺมา อติเวลํ กุเลสุ จาริตฺตํ อาปชฺชี’ติ, โส ภิกฺขุ ภิกฺขูหิ วุจฺจมาโน น วิรมตี”ติ.}

\prose{ภูตปุพฺพํ ภิกฺขเว พิฬาโร สนฺธิสมลสํกฏีเร ฐิโต อโหสิ มุทุมูสิํ มคฺคยมาโน, ยทายํ มุทุมูสิ โคจราย ปกฺกมิสฺสติ, ตตฺเถว นํ คเหตฺวา ขาทิสฺสามีติ.\csromanspacer{} อถ โข โส ภิกฺขเว มุทุมูสิ โคจราย ปกฺกามิ.\csromanspacer{} ตเมนํ พิฬาโร คเหตฺวา สหสา สงฺขาทิตฺวา อชฺโฌหริ, ตสฺส โส มุทุมูสิ อนฺตมฺปิ ขาทิ, อนฺตคุณมฺปิ ขาทิ, โส ตโตนิทานํ มรณมฺปิ นิคจฺฉิ มรณมตฺตมฺปิ ทุกฺขํ.}

\csromanmatikaanchor{mtk.281}
\csromantocmarknopage{subparagraph}{9. โอปมฺมสํยุตฺต}{mtk.281}
\bookmark[dest={mtk.281},level=5]{9. โอปมฺมสํยุตฺต}
\prose{เอวเมว โข ภิกฺขเว อิเธกจฺโจ ภิกฺขุ ปุพฺพณฺหสมยํ นิวาเสตฺวา ปตฺตจีวรมาทาย คามํ วา นิคมํ วา ปิณฺฑาย ปวิสติ อรกฺขิเตเนว กาเยน อรกฺขิตาย วาจาย อรกฺขิเตน จิตฺเตน อนุปฏฺฐิตาย สติยา อสํวุเตหิ อินฺทฺริเยหิ.\csromanspacer{} โส ตตฺถ ปสฺสติ มาตุคามํ ทุนฺนิวตฺถํ วา ทุปฺปารุตํ วา, ตสฺส มาตุคามํ ทิสฺวา ทุนฺนิวตฺถํ วา ทุปฺปารุตํ วา ราโค จิตฺตํ อนุทฺธํเสติ.\csromanspacer{} โส ราคานุทฺธํเสน จิตฺเตน มรณํ วา นิคจฺฉติ มรณมตฺตํ วา ทุกฺขํ.\csromanspacer{} มรณํ เหตํ ภิกฺขเว อริยสฺส วินเย โย สิกฺขํ ปจฺจกฺขาย หีนายาวตฺตติ.\csromanspacer{} มรณมตฺตํ เหตํ ภิกฺขเว ทุกฺขํ ยทิทํ อญฺญตรํ สํกิลิฏฺฐํ อาปตฺติํ อาปชฺชติ.\csromanspacer{} ยถารูปาย อาปตฺติยา วุฏฺฐานํ ปญฺญายติ.\csromanspacer{} ตสฺมาติห ภิกฺขเว เอวํ สิขิตพฺพํ “รกฺขิเตเนว กาเยน รกฺขิตาย วาจาย รกฺขิเตน จิตฺเตน อุปฏฺฐิตาย สติยา สํวุเตหิ อินฺทฺริเยหิ คามํ วา นิคมํ วา ปิณฺฑาย ปวิสิสฺสามา”ติ.\csromanspacer{} เอวํ หิ โว ภิกฺขเว สิกฺขิตพฺพนฺติ.\csromanspacer{}\csromanspacer{}ทสมํ.}

\csromanmatikaanchor{mtk.282}
\csromantocmarkref{subparagraph}{1. กูฏสุตฺต}{mtk.282}
\bookmark[dest={mtk.282},level=5]{1. กูฏสุตฺต}
\csromanheader{6}{11. สิงฺคาลสุตฺต}
\proseitem{233}{สาวตฺถิยํ วิหรติ.\csromanspacer{} อสฺสุตฺถ โน ตุเมฺห ภิกฺขเว รตฺติยา ปจฺจูสสสฺมยํ ชรสิงฺคาลสฺส วสฺสมานสฺสาติ.\csromanspacer{} เอวํ ภนฺเต.\csromanspacer{} เอโส โข ภิกฺขเว ชรสิงฺคาโล อุกฺกณฺฑเกน นาม โรคชาเตน ผุฏฺโฐ.\csromanspacer{} โส}

\csromanmatikaanchor{mtk.283}
\csromantocmarkref{subparagraph}{2. นขสิขสุตฺต}{mtk.283}
\bookmark[dest={mtk.283},level=5]{2. นขสิขสุตฺต}
\vspace{\tipitakaparskip}
\csromancenter{โอปมฺมสํยุตฺต เยน เยน อิจฺฉติ, เตน เตน คจฺฉติ.\csromanspacer{} ยตฺถ ยตฺถ อิจฺฉติ, ตตฺถ ตตฺถ ติฏฺฐติ.\csromanspacer{} ยตฺถ ยตฺถ อิจฺฉติ, ตตฺถ ตตฺถ นิสีทติ.\csromanspacer{} ยตฺถ ยตฺถ อิจฺฉติ, ตตฺถ ตตฺถ นิปชฺชติ.\csromanspacer{} สีตโกปิ นํ วาโต อุปวายติ.\csromanspacer{} สาธุ ขฺวสฺส ภิกฺขเว ยํ อิเธกจฺโจ สกฺยปุตฺติยปฏิญฺโญ เอวรูปมฺปิ อตฺตภาวปฏิลาภํ ปฏิสํเวทิเยถ.\csromanspacer{} ตสฺมาติห ภิกฺขเว เอวํ สิกฺขิตพฺพํ “อปฺปมตฺตา วิหริสฺสามา”ติ.\csromanspacer{} เอวํ หิ โว ภิกฺขเว สิกฺขิตพฺพนฺติ.\csromanspacer{}\csromanspacer{}เอกาทสมํ.}
\csromanheader{6}{12. ทุติยสิงฺคาลสุตฺต}
\csromanmatikaanchor{mtk.284}
\csromanmatikaanchor{mtk.294}
\csromantocmarkref{subparagraph}{3. กุลสุตฺต}{mtk.284}
\bookmark[dest={mtk.284},level=5]{3. กุลสุตฺต}
\proseitem{234}{\csromanfolio{461}สาวตฺถิยํ วิหรติ.\csromanspacer{} อสฺสุตฺถ โน ตุเมฺห ภิกฺขเว รตฺติยา ปจฺจูสสมยํ ชรสิงฺคาลสฺส วสฺสมานสฺสาติ.\csromanspacer{} เอวํ ภนฺเต.\csromanspacer{} สิยา โข ภิกฺขเว ตสฺมิํ ชรสิงฺคาเล ยา กาจิ กตญฺญุตา กตเวทิตา น เตฺวว อิเธกจฺเจ สกฺยปุตฺติยปฏิญฺเญ สิยา ยา กาจิ กตญฺญุตา กตเวทิตา.\csromanspacer{} ตสฺมาติห ภิกฺขเว เอวํ สิกฺขิตพฺพํ “กตญฺญุโน ภวิสฺสาม กตเวทิโน.\csromanspacer{} น จ โน\footnote{น จ โนติ อิทํ สี-อิ-โปตฺถเกสุ นตฺถิ.} อเมฺหสุ อปฺปกมฺปิ กตํ นสฺสิสฺสตี”ติ\footnote{มา นสฺสิสฺสตีติ (สี, อิ), วินสฺสิสฺสตีติ (สฺยา, กํ)}.\csromanspacer{} เอวํ หิ โว ภิกฺขเว สิกฺขิตพฺพนฺติ.\csromanspacer{}\csromanspacer{}ทฺวาทสมํ.}

\vspace{\tipitakaparskip}
\csromancenter{\textbf{โอปมฺมสํยุตฺตํ สมตฺตํ.}}
\titlehead{\textbf{ตสฺสุทฺทานํ}}
\csromanmatikaanchor{mtk.285}
\csromantocmarkref{subparagraph}{4. โอกฺขาสุตฺต}{mtk.285}
\bookmark[dest={mtk.285},level=5]{4. โอกฺขาสุตฺต}
\csromangathameasure{กูฏํ นขสิขํ กุลํ, โอกฺขา สตฺติ ธนุคฺคโห. \\ อาณิ กลิงฺคโร นาโค, พิฬาโร เทฺว สิงฺคาลกาติ.}
\csromangathagroup{\csromangathastanza{กูฏํ นขสิขํ กุลํ, โอกฺขา สตฺติ ธนุคฺคโห. \\ อาณิ กลิงฺคโร นาโค, พิฬาโร เทฺว สิงฺคาลกาติ.}}

\csromanheader{6}{10. ภิกฺขุสํยุตฺต}
\csromanmatikaanchor{mtk.286}
\csromantocmarkref{subparagraph}{5. สตฺติสุตฺต}{mtk.286}
\bookmark[dest={mtk.286},level=5]{5. สตฺติสุตฺต}
\csromanheader{6}{1. โกลิตสุตฺต}
\proseitem{235}{\csromanfolio{462}เอวํ เม สุตํ\csromandash{}เอกํ สมยํ ภควา สาวตฺถิยํ วิหรติ เชตวเน อนาถปิณฺฑิกสฺส อาราเม.\csromanspacer{} ตตฺร โข อายสฺมา มหาโมคฺคลฺลาโน ภิกฺขู อามนฺเตสิ “อาวุโส ภิกฺขเว”ติ. “อาวุโส”ติ โข เต ภิกฺขู อายสฺมโต มหาโมคฺคลฺลานสฺส ปจฺจสฺโสสุํ.}

\prose{อายสฺมา มหาโมคฺคลฺลาโน เอตทโวจ\csromandash{}อิธ มยฺหํ อาวุโส รโหคตสฺส ปฏิสลฺลีนสฺส เอวํ เจตโส ปริวิตกฺโก อุทปาทิ “อริโย ตุณฺหีภาโว อริโย ตุณฺหีภาโวติ วุจฺจติ.\csromanspacer{} กตโม นุ โข อริโย ตุณฺหีภาโว”ติ.\csromanspacer{} ตสฺส มยฺหํ อาวุโส เอตทโหสิ “อิธ ภิกฺขุ วิตกฺกวิจารานํ วูปสมา อชฺฌตฺตํ สมฺปสาทนํ เจตโส เอโกทิภาวํ อวิตกฺกํ อวิจารํ สมาธิชํ ปีติสุขํ ทุติยํ ฌานํ อุปสมฺปชฺช วิหรติ, อยํ วุจฺจติ อริโย ตุณฺหีภาโว”ติ.\csromanspacer{} โส ขฺวาหํ อาวุโส วิตกฺกวิจารานํ วูปสมา อชฺฌตฺตํ สมฺปสาทนํ เจตโส เอโกทิภาวํ อวิตกฺกํ อวิจารํ สมาธิชํ ปีติสุขํ ทุติยํ ฌานํ อุปสมฺปชฺช วิหริํ, ตสฺส มยฺหํ อาวุโส อิมินา วิหาเรน วิหรโต วิตกฺกสหคตา สญฺญา มนสิการา สมุทาจรนฺติ.}

\prose{อถ โข มํ อาวุโส ภควา อิทฺธิยา อุปสงฺกมิตฺวา เอตทโวจ “โมคฺคลฺลาน โมคฺคลฺลาน มา พฺราหฺมณ อริยํ ตุณฺหีภาวํ ปมาโท.\csromanspacer{} อริเย ตุณฺหีภาเว จิตฺตํ สณฺฐเปหิ, อริเย ตุณฺหีภาเว จิตฺตํ เอโกทิภาวํ กโรหิ, อริเย ตุณฺหีภาเว จิตฺตํ สมาทหา”ติ.\csromanspacer{} โส ขฺวาหํ อาวุโส อปเรน สมเยน วิตกฺกวิจารานํ วูปสมา อชฺฌตฺตํ สมฺปสาทนํ เจตโส เอโกทิภาวํ อวิตกฺกํ อวิจารํ สมาธิชํ ปีติสุขํ ทุติยํ ฌานํ อุปสมฺปชฺช วิหรามิ.\csromanspacer{} ยํ หิ ตํ อาวุโส สมฺมา วทมาโน วเทยฺย “สตฺถารา อนุคฺคหิโต สาวโก มหาภิญฺญตํ ปตฺโต”ติ.\csromanspacer{} มมํ ตํ สมฺมา วทมาโน วเทยฺย “สตฺถารา อนุคฺคหิโต สาวโก มหาภิญฺญตํ ปตฺโต”ติ.\csromanspacer{}\csromanspacer{}ปฐมํ.}

\csromanmatikaanchor{mtk.287}
\csromantocmarkref{subparagraph}{6. ธนุคฺคหสุตฺต}{mtk.287}
\bookmark[dest={mtk.287},level=5]{6. ธนุคฺคหสุตฺต}
\csromanheader{6}{10. ภิกฺขุสํยุตฺต \textbf{2. อุปติสฺสสุตฺต}}
\proseitem{236}{\csromanfolio{463}สาวตฺถินิทานํ.\csromanspacer{} ตตฺร โข อายสฺมา สาริปุตฺโต ภิกฺขู อามนฺเตสิ “อาวุโส ภิกฺขเว”ติ. “อาวุโส”ติ โข เต ภิกฺขู อายสฺมโต สาริปุตฺตสฺส ปจฺจสฺโสสุํ.\csromanspacer{} อายสฺมา สาริปุตฺโต เอตทโวจ\csromandash{}}

\prose{อิธ มยฺหํ อาวุโส รโหคตสฺส ปฏิสลฺลีนสฺส เอวํ เจตโส ปริวิตกฺโก อุทปาทิ “อตฺถิ นุ โข ตํ กิญฺจิ โลกสฺมิํ, ยสฺส เม วิปริณามญฺญถาภาวา อุปฺปชฺเชยฺยุํ โสกปริเทวทุกฺขโทมนสฺสุปายาสา”ติ.\csromanspacer{} ตสฺส มยฺหํ อาวุโส เอตทโหสิ “นตฺถิ โข ตํ กิญฺจิ โลกสฺมิํ, ยสฺส เม วิปริณามญฺญถาภาวา อุปฺปชฺเชยฺยุํ โสกปริเทวทุกฺขโทมนสฺสุปายาสา”ติ.}

\prose{เอวํ วุตฺเต อายสฺมา อานนฺโท อายสฺมนฺตํ สาริปุตฺตํ เอตทโวจ “สตฺถุปิ โข เต อาวุโส สาริปุตฺต วิปริณามญฺญถาภาวา นุปฺปชฺเชยฺยุํ โสกปริเทวทุกฺขโทมนสฺสุปายาสา”ติ.\csromanspacer{} สตฺถุปิ โข เม อาวุโส วิปริณามญฺญถาภาวา นุปฺปชฺเชยฺยุํ โสกปริเทวทุกฺขโทมนสฺสุปายาสา, อปิ จ เม เอวมสฺส “มเหสกฺโข วต โภ สตฺถา อนฺตรหิโต มหิทฺธิโก มหานุภาโว.\csromanspacer{} สเจ หิ ภควา จิรํ ทีฆมทฺธานํ ติฏฺเฐยฺย ตทสฺส พหุชนหิตาย พหุชนสุขาย โลกานุกมฺปาย อตฺถาย หิตาย สุขาย เทวมนุสฺสานนฺ”ติ.\csromanspacer{} ตถา หิ ปนายสฺมโต สาริปุตฺตสฺส ทีฆรตฺตํ อหงฺการมมงฺการมานานุสยา สุสมูหตา.\csromanspacer{} ตสฺมา อายสฺมโต สาริปุตฺตสฺส สตฺถุปิ วิปริณามญฺญถาภาวา นุปฺปชฺเชยฺยุํ โสกปริเทวทุกฺขโทมนสฺสุปายาสาติ.\csromanspacer{}\csromanspacer{}ทุติยํ.}

\csromanmatikaanchor{mtk.288}
\csromantocmarkref{subparagraph}{7. อาณิสุตฺต}{mtk.288}
\bookmark[dest={mtk.288},level=5]{7. อาณิสุตฺต}
\csromanheader{6}{3. ฆฏสุตฺต}
\proseitem{237}{เอวํ เม สุตํ\csromandash{}เอกํ สมยํ ภควา สาวตฺถิยํ วิหรติ เชตวเน อนาถปิณฺฑิกสฺส อาราเม.\csromanspacer{} เตน โข ปน สมเยน อายสฺมา จ สาริปุตฺโต อายสฺมา จ มหาโมคฺคลฺลาโน ราชคเห วิหรนฺติ เวฬุวเน กลนฺทกนิวาเป เอกวิหาเร.\csromanspacer{} อถ โข อายสฺมา สาริปุตฺโต สายนฺหสมยํ ปฏิสลฺลานา วุฏฺฐิโต เยนายสฺมา มหาโมคฺคลฺลาโน เตนุปสงฺกมิ, อุปสงฺกมิตฺวา อายสฺมตา มหาโมคฺคลฺลาเนน สทฺธิํ สมฺโมทิ, สมฺโมทนียํ กถํ สารณียํ วีติสาเรตฺวา}

\gambhira{นิทานวคฺคสํยุตฺตปาฬิ}
\csromanmatikaanchor{mtk.289}
\csromantocmarkref{subparagraph}{8. กลิงฺครสุตฺต}{mtk.289}
\bookmark[dest={mtk.289},level=5]{8. กลิงฺครสุตฺต}
\prose{\csromanfolio{464}เอกมนฺตํ นิสีทิ, เอกมนฺตํ นิสินฺโน โข อายสฺมา สาริปุตฺโต อายสฺมนฺตํ มหาโมคฺคลฺลานํ เอตทโวจ\csromandash{}}

\prose{วิปฺปสนฺนานิ โข เต อาวุโส โมคฺคลฺลาน อินฺทฺริยานิ ปริสุทฺโธ มุขวณฺโณ ปริโยทาโต สนฺเตน นูนายสฺมา มหาโมคฺคลฺลาโน อชฺช วิหาเรน วิหาสีติ.\csromanspacer{} โอฬาริเกน ขฺวาหํ อาวุโส อชฺช วิหาเรน วิหาสิํ, อปิ จ เม อโหสิ “ธมฺมี กถา”ติ.\csromanspacer{} เกน สทฺธิํ ปนายสฺมโต มหาโมคฺคลฺลานสฺส อโหสิ “ธมฺมี กถา”ติ.\csromanspacer{} ภควตา โข เม อาวุโส สทฺธิํ อโหสิ “ธมฺมี กถา”ติ.\csromanspacer{} ทูเร โข อาวุโส ภควา, เอตรหิ สาวตฺถิยํ วิหรติ เชตวเน อนาถปิณฺฑิกสฺส อาราเม.\csromanspacer{} กิํ นุ โข อายสฺมา มหาโมคฺคลฺลาโน ภควนฺตํ อิทฺธิยา อุปสงฺกมิ, อุทาหุ ภควา อายสฺมนฺตํ มหาโมคฺคลฺลานํ อิทฺธิยา อุปสงฺกมีติ.\csromanspacer{} น ขฺวาหํ อาวุโส ภควนฺตํ อิทฺธิยา อุปสงฺกมิํ, นปิ มํ ภควา อิทฺธิยา อุปสงฺกมิ.\csromanspacer{} อปิ จ เม ยาวตา ภควา เอตฺตาวตา ทิพฺพจกฺขุ วิสุชฺฌิ ทิพฺพา จ โสตธาตุ.\csromanspacer{} ภควโตปิ ยาวตาหํ เอตฺตาวตา ทิพฺพจกฺขุ วิสุชฺฌิ ทิพฺพา จ โสตธาตูติ.\csromanspacer{} ยถากถํ ปนายสฺมโต มหาโมคฺคลฺลานสฺส ภควตา สทฺธิํ อโหสิ ธมฺมี กถาติ.}

\prose{อิธาหํ อาวุโส ภควนฺตํ เอตทโวจํ “อารทฺธวีริโย อารทฺธวีริโยติ ภนฺเต วุจฺจติ.\csromanspacer{} กิตฺตาวตา นุ โข ภนฺเต อารทฺธวีริโย โหตี”ติ.\csromanspacer{} เอวํ วุตฺเต มํ อาวุโส ภควา เอตทโวจ “อิธ โมคฺคลฺลาน ภิกฺขุ อารทฺธวีริโย วิหรติ ‘กามํ ตโจ จ นฺหารุ จ อฏฺฐิ จ อวสิสฺสตุ สรีเร อุปสุสฺสตุ มํสโลหิตํ, ยํ ตํ ปุริสถาเมน ปุริสวีริเยน ปุริสปรกฺกเมน ปตฺตพฺพํ, น ตํ อปาปุณิตฺวา วีริยสฺส สณฺฐานํ ภวิสฺสตี’ติ.\csromanspacer{} เอวํ โข โมคฺคลฺลาน อารทฺธวีริโย โหตี”ติ.\csromanspacer{} เอวํ โข เม อาวุโส ภควตา สทฺธิํ อโหสิ ธมฺมี กถาติ.}

\prose{เสยฺยถาปิ อาวุโส หิมวโต ปพฺพตราชสฺส ปริตฺตา ปาสาณสกฺขรา ยาวเทว อุปนิกฺเขปนมตฺตาย.\csromanspacer{} เอวเมว โข มยํ อายสฺมโต มหาโมคฺคลฺลานสฺส ยาวเทว อุปนิกฺเขปนมตฺตาย.\csromanspacer{} อายสฺมา หิ มหาโมคฺคลฺลาโน มหิทฺธิโก มหานุภาโว อากงฺขมาโน กปฺปํ ติฏฺเฐยฺยาติ.}

\csromanheader{2}{10. ภิกฺขุสํยุตฺต}
\prose{\csromanfolio{465}เสยฺยถาปิ อาวุโส มหติยา โลณฆฏาย ปริตฺตา โลณสกฺขราย ยาวเทว อุปนิกฺเขปนมตฺตาย.\csromanspacer{} เอวเมว โข มยํ อายสฺมโต สาริปุตฺตสฺส ยาวเทว อุปนิกฺเขปนมตฺตาย.\csromanspacer{} อายสฺมา หิ สาริปุตฺโต ภควตา อเนกปริยาเยน โถมิโต วณฺณิโต ปสตฺโถ.}

\csromanmatikaanchor{mtk.290}
\csromantocmarkref{subparagraph}{9. นาคสุตฺต}{mtk.290}
\bookmark[dest={mtk.290},level=5]{9. นาคสุตฺต}
\csromangathameasure{“สาริปุตฺโตว ปญฺญาย, สีเลน อุปสเมน จ. \\ โยปิ ปารงฺคโต ภิกฺขุ, เอตาวปรโม สิยา”ติ.}
\csromangathagroup{\csromangathastanza{“สาริปุตฺโตว ปญฺญาย, สีเลน อุปสเมน จ. \\ โยปิ ปารงฺคโต ภิกฺขุ, เอตาวปรโม สิยา”ติ.}}

\prose{อิติห เต อุโภ มหานาคา อญฺญมญฺญสฺส สุภาสิตํ สุลปิตํ สมนุโมทิํสูติ.\csromanspacer{}\csromanspacer{}ตติยํ.}

\csromanheader{6}{4. นวสุตฺต}
\proseitem{238}{สาวตฺถิยํ วิหรติ.\csromanspacer{} เตน โข ปน สมเยน อญฺญตโร นโว ภิกฺขุ ปจฺฉาภตฺตํ ปิณฺฑปาตปฏิกฺกนฺโต วิหารํ ปวิสิตฺวา อปฺโปสฺสุกฺโก ตุณฺหีภูโต สงฺกสายติ, น ภิกฺขูนํ เวยฺยาวจฺจํ กโรติ จีวรการสมเย.\csromanspacer{} อถ โข สมฺพหุลา ภิกฺขู เยน ภควา เตนุปสงฺกมิํสุ, อุปสงฺกมิตฺวา ภควนฺตํ อภิวาเทตฺวา เอกมนฺตํ นิสีทิํสุ, เอกมนฺตํ นิสินฺนา โข เต ภิกฺขู ภควนฺตํ เอตทโวจุํ “อิธ ภนฺเต อญฺญตโร นโว ภิกฺขุ ปจฺฉาภตฺตํ ปิณฺฑปาตปฏิกฺกนฺโต วิหารํ ปวิสิตฺวา อปฺโปสฺสุกฺโก ตุณฺหีภูโต สงฺกสายติ, น ภิกฺขูนํ เวยฺยาวจฺจํ กโรติ จีวรการสมเย”ติ.}

\prose{อถ โข ภควา อญฺญตรํ ภิกฺขุํ อามนฺเตสิ “เอหิ ตฺวํ ภิกฺขุ มม วจเนน ตํ ภิกฺขุํ อามนฺเตหิ ‘สตฺถา ตํ อาวุโส อามนฺเตตี’ติ”. “เอวํ ภนฺเต”ติ โข โส ภิกฺขุ ภควโต ปฏิสฺสุตฺวา เยน โส ภิกฺขุ เตนุปสงฺกมิ, อุปสงฺกมิตฺวา ตํ ภิกฺขุํ เอตทโวจ “สตฺถา ตํ อาวุโส อามนฺเตตี”ติ. “เอวมาวุโส”ติ โข โส ภิกฺขุ ตสฺส ภิกฺขุโน ปฏิสฺสุตฺวา เยน ภควา เตนุปสงฺกมิ, อุปสงฺกมิตฺวา ภควนฺตํ อภิวาเทตฺวา เอกมนฺตํ นิสีทิ, เอกมนฺตํ นิสินฺนํ โข ตํ ภิกฺขุํ ภควา เอตทโวจ “สจฺจํ กิร ตฺวํ ภิกฺขุ ปจฺจาภตฺตํ ปิณฺฑปาตปฏิกฺกนฺโต วิหารํ ปวิสิตฺวา อปฺโปสฺสุกฺโก ตุณฺหีภูโต สงฺกสายสิ, น ภิกฺขูนํ เวยฺยาวจฺจํ กโรสิ จีวรการสมเย”ติ.\csromanspacer{} อหมฺปิ โข ภนฺเต สกํ กิจฺจํ กโรมีติ.}

\gambhira{นิทานวคฺคสํยุตฺตปาฬิ}
\csromanmatikaanchor{mtk.291}
\csromantocmarkref{subparagraph}{10. พิฬารสุตฺต}{mtk.291}
\bookmark[dest={mtk.291},level=5]{10. พิฬารสุตฺต}
\prose{\csromanfolio{466}อถ โข ภควา ตสฺส ภิกฺขุโน เจตสา เจโตปริวิตกฺกมญฺญาย ภิกฺขู อามนฺเตสิ “มา โข ตุเมฺห ภิกฺขเว เอตสฺส ภิกฺขุโน อุชฺฌายิตฺถ, เอโส โข ภิกฺขเว ภิกฺขุ จตุนฺนํ ฌานานํ อาภิเจตสิกานํ ทิฏฺฐธมฺมสุขวิหารานํ นิกามลาภี อกิจฺฉลาภี อกสิรลาภี, ยสฺส จตฺถาย กุลปุตฺตา สมฺมเทว อคารสฺมา อนคาริยํ ปพฺพชนฺติ, ตทนุตฺตรํ พฺรหฺมจริยปริโยสานํ ทิฏฺเฐว ธมฺเม สยํ อภิญฺญา สจฺฉิกตฺวา อุปสมฺปชฺช วิหรตี”ติ.}

\prose{อิทมโวจ ภควา, อิทํ วตฺวาน สุคโต อถาปรํ เอตทโวจ สตฺถา\csromandash{}}

\csromangathameasure{“นยิทํ สิถิลมารพฺภ, นยิทํ อปฺเปน ถามสา. \\ นิพฺพานํ อธิคนฺตพฺพํ, สพฺพทุกฺขปฺปโมจนํ. \\ อยญฺจ ทหโร ภิกฺขุ, อยมุตฺตมปุริโส. \\ ธาเรติ อนฺติมํ เทหํ, เชตฺวา มารํ สวาหินินฺ”ติ.}
\csromangathagroup{\csromangathastanza{“นยิทํ สิถิลมารพฺภ, นยิทํ อปฺเปน ถามสา. \\ นิพฺพานํ อธิคนฺตพฺพํ, สพฺพทุกฺขปฺปโมจนํ.}\csromangathastanza{อยญฺจ ทหโร ภิกฺขุ, อยมุตฺตมปุริโส. \\ ธาเรติ อนฺติมํ เทหํ, เชตฺวา มารํ สวาหินินฺ”ติ.}}

\vspace{\tipitakaparskip}
\csromancenter{จตุตฺถํ.}
\csromanheader{6}{5. สุชาตสุตฺต}
\csromanmatikaanchor{mtk.292}
\csromantocmarkref{subparagraph}{11. สิงฺคาลสุตฺต}{mtk.292}
\bookmark[dest={mtk.292},level=5]{11. สิงฺคาลสุตฺต}
\proseitem{239}{สาวตฺถิยํ วิหรติ.\csromanspacer{} อถ โข อายสฺมา สุชาโต เยน ภควา เตนุปสงฺกมิ.\csromanspacer{} อทฺทสา โข ภควา อายสฺมนฺตํ สุชาตํ ทูรโตว อาคจฺฉนฺตํ, ทิสฺวาน ภิกฺขู อามนฺเตสิ “อุภเยเนวายํ ภิกฺขเว กุลปุตฺโต โสภติ ยญฺจ อภิรูโป ทสฺสนีโย ปาสาทิโก ปรมาย วณฺณโปกฺขรตาย สมนฺนาคโต.\csromanspacer{} ยสฺส จตฺถาย กุลปุตฺตา สมฺมเทว อคารสฺมา อนคาริยํ ปพฺพชนฺติ, ตทนุตฺตรํ พฺรหฺมจริยปริโยสานํ ทิฏฺเฐว ธมฺเม สยํ อภิญฺญา สจฺฉิกตฺวา อุปสมฺปชฺช วิหรตี”ติ.\csromanspacer{} อิทมโวจ ภควา ฯเปฯ สตฺถา\csromandash{}}

\csromangathameasure{“โสภติ วตายํ ภิกฺขุ, อุชุภูเตน เจตสา. \\ วิปฺปยุตฺโต วิสํยุตฺโต, อนุปาทาย นิพฺพุโต. \\ ธาเรติ อนฺติมํ เทหํ, เชตฺวา มารํ สวาหินินฺ”ติ.}
\csromangathagroup{\csromangathastanza{“โสภติ วตายํ ภิกฺขุ, อุชุภูเตน เจตสา. \\ วิปฺปยุตฺโต วิสํยุตฺโต, อนุปาทาย นิพฺพุโต. \\ ธาเรติ อนฺติมํ เทหํ, เชตฺวา มารํ สวาหินินฺ”ติ.}}

\vspace{\tipitakaparskip}
\csromancenter{ปญฺจมํ.}
\csromanmatikaanchor{mtk.293}
\csromantocmarkref{subparagraph}{12. ทุติยสิงฺคาลสุตฺต}{mtk.293}
\bookmark[dest={mtk.293},level=5]{12. ทุติยสิงฺคาลสุตฺต}
\csromantocmarkref{section}{อุทฺทานคาถา}{mtk.294}
\bookmark[dest={mtk.294},level=1]{อุทฺทานคาถา}
\csromanheader{6}{10. ภิกฺขุสํยุตฺต \textbf{6. ลกุณฺฑกภทฺทิยสุตฺต}}
\proseitem{240}{\csromanfolio{467}สาวตฺถิยํ วิหรติ.\csromanspacer{} อถ โข อายสฺมา ลกุณฺฑกภทฺทิโย เยน ภควา เตนุปสงฺกมิ.\csromanspacer{} อทฺทสา โข ภควา อายสฺมนฺตํ ลกุณฺฑกภทฺทิยํ ทูรโตว อาคจฺฉนฺตํ, ทิสฺวาน ภิกฺขู อามนฺเตสิ “ปสฺสถ โน ตุเมฺห ภิกฺขเว เอตํ ภิกฺขุํ อาคจฺฉนฺตํ ทุพฺพณฺณํ ทุทฺทสิกํ โอโกฏิมกํ ภิกฺขูนํ ปริภูตรูปนฺ”ติ.\csromanspacer{} เอวํ ภนฺเต.\csromanspacer{} เอโส โข ภิกฺขเว ภิกฺขุ มหิทฺธิโก มหานุภาโว, น จ สา สมาปตฺติ สุลภรูปา, ยา เตน ภิกฺขุนา อสมาปนฺนปุพฺพา.\csromanspacer{} ยสฺส จตฺถาย กุลปุตฺตา สมฺมเทว อคารสฺมา อนคาริยํ ปพฺพชนฺติ, ตทนุตฺตรํ พฺรหฺมจริยปริโยสานํ ทิฏฺเฐว ธมฺเม สยํ อภิญฺญา สจฺฉิกตฺวา อุปสมฺปชฺช วิหรตีติ.\csromanspacer{} อิทมโวจ ภควา ฯเปฯ สตฺถา\csromandash{}}

\csromangathameasure{“หํสา โกญฺจา มยูรา จ, หตฺถโย ปสทา มิคา. \\ สพฺเพ สีหสฺส ภายนฺติ, นตฺถิ กายสฺมิํ ตุลฺยตา. \\ เอวเมว มนุสฺเสสุ, ทหโร เจปิ ปญฺญวา. \\ โส หิ ตตฺถ มหา โหติ, เนว พาโล สรีรวา”ติ.}
\csromangathagroup{\csromangathastanza{“หํสา โกญฺจา มยูรา จ, หตฺถโย ปสทา มิคา. \\ สพฺเพ สีหสฺส ภายนฺติ, นตฺถิ กายสฺมิํ ตุลฺยตา.}\csromangathastanza{เอวเมว มนุสฺเสสุ, ทหโร เจปิ ปญฺญวา. \\ โส หิ ตตฺถ มหา โหติ, เนว พาโล สรีรวา”ติ.}}

\vspace{\tipitakaparskip}
\csromancenter{ฉฏฺฐํ.}
\csromanmatikaanchor{mtk.295}
\csromantocmarknopage{subparagraph}{10. ภิกฺขุสํยุตฺต}{mtk.295}
\bookmark[dest={mtk.295},level=5]{10. ภิกฺขุสํยุตฺต}
\csromanheader{6}{7. วิสาขสุตฺต}
\csromanmatikaanchor{mtk.296}
\csromantocmarkref{subparagraph}{1. โกลิตสุตฺต}{mtk.296}
\bookmark[dest={mtk.296},level=5]{1. โกลิตสุตฺต}
\proseitem{241}{เอวํ เม สุตํ\csromandash{}เอกํ สมยํ ภควา เวสาลิยํ วิหรติ มหาวเน กูฏาคารสาลายํ.\csromanspacer{} เตน โข ปน สมเยน อายสฺมา วิสาโข ปญฺจาลปุตฺโต อุปฏฺฐานสาลายํ ภิกฺขู ธมฺมิยา กถาย สนฺทสฺเสติ สมาทเปติ สมุตฺเตเชติ สมฺปหํเสติ โปริยา วาจาย วิสฺสฏฺฐาย อเนลคลาย อตฺถสฺส วิญฺญาปนิยา ปริยาปนฺนาย อนิสฺสิตาย.}

\prose{อถ โข ภควา สายนฺหสมยํ ปฏิสลฺลานา วุฏฺฐิโต เยน อุปฏฺฐานสาลา เตนุปสงฺกมิ, อุปสงฺกมิตฺวา ปญฺญตฺเต อาสเน นิสีทิ, นิสชฺช โข ภควา ภิกฺขู อามนฺเตสิ “โก นุ โข ภิกฺขเว อุปฏฺฐานสาลายํ ภิกฺขู ธมฺมิยา กถาย สนฺทสฺเสติ สมาทเปติ สมุตฺเตเชติ สมฺปหํเสติ โปริยา วาจาย วิสฺสฏฺฐาย อเนลคลาย อตฺถสฺส วิญฺญาปนิยา ปริยาปนฺนาย อนิสฺสิตายา”ติ.\csromanspacer{} อายสฺมา ภนฺเต}

\gambhira{นิทานวคฺคสํยุตฺตปาฬิ}
\prose{\csromanfolio{468}วิสาโข ปญฺจาลปุตฺโต อุปฏฺฐานสาลายํ ภิกฺขู ธมฺมิยา กถาย สนฺทสฺเสติ สมาทเปติ สมุตฺเตเชติ สมฺปหํเสติ โปริยา วาจาย วิสฺสฏฺฐาย อเนลคลาย อตฺถสฺส วิญฺญาปนิยา ปริยาปนฺนาย อนิสฺสิตายาติ.\csromanspacer{} อถ โข ภควา อายสฺมนฺตํ วิสาขํ ปญฺจาลปุตฺตํ อามนฺเตสิ “สาธุ สาธุ วิสาข, สาธุ โข ตฺวํ วิสาข ภิกฺขู ธมฺมิยา กถาย สนฺทสฺเสสิ ฯเปฯ อตฺถสฺส วิญฺญาปนิยา ปริยาปนฺนาย อนิสฺสิตายา”ติ.}

\csromanmatikaanchor{mtk.297}
\csromantocmarkref{subparagraph}{2. อุปติสฺสสุตฺต}{mtk.297}
\bookmark[dest={mtk.297},level=5]{2. อุปติสฺสสุตฺต}
\prose{อิทมโวจ ภควา, อิทํ วตฺวาน สุคโต อถาปรํ เอตทโวจ สตฺถา\csromandash{}}

\csromangathameasure{“นาภาสมานํ ชานนฺติ, มิสฺสํ พาเลหิ ปณฺฑิตํ. \\ ภาสมานญฺจ ชานนฺติ, เทเสนฺตํ อมตํ ปทํ. \\ ภาสเย โชตเย ธมฺมํ, ปคฺคเณฺห อิสินํ ธชํ. \\ สุภาสิตธชา อิสโย, ธมฺโม หิ อิสินํ ธโช”ติ.}
\csromangathagroup{\csromangathastanza{“นาภาสมานํ ชานนฺติ, มิสฺสํ พาเลหิ ปณฺฑิตํ. \\ ภาสมานญฺจ ชานนฺติ, เทเสนฺตํ อมตํ ปทํ.}\csromangathastanza{ภาสเย โชตเย ธมฺมํ, ปคฺคเณฺห อิสินํ ธชํ. \\ สุภาสิตธชา อิสโย, ธมฺโม หิ อิสินํ ธโช”ติ.}}

\vspace{\tipitakaparskip}
\csromancenter{สตฺตมํ.}
\csromanmatikaanchor{mtk.298}
\csromantocmarkref{subparagraph}{3. ฆฏสุตฺต}{mtk.298}
\bookmark[dest={mtk.298},level=5]{3. ฆฏสุตฺต}
\csromanheader{6}{8. นนฺทสุตฺต}
\proseitem{242}{สาวตฺถิยํ วิหรติ.\csromanspacer{} อถ โข อายสฺมา นนฺโท ภควโต มาตุจฺฉาปุตฺโต อาโกฏิตปจฺจาโกฏิตานิ จีวรานิ ปารุปิตฺวา อกฺขีนิ อญฺเชตฺวา อจฺฉํ ปตฺตํ คเหตฺวา เยน ภควา เตนุปสงฺกมิ, อุปสงฺกมิตฺวา ภควนฺตํ อภิวาเทตฺวา เอกมนฺตํ นิสีทิ, เอกมนฺตํ นิสินฺนํ โข อายสฺมนฺตํ นนฺทํ ภควา เอตทโวจ “น โข เต ตํ นนฺท ปติรูปํ กุลปุตฺตสฺส สทฺธา อคารสฺมา อนคาริยํ ปพฺพชิตสฺส, ยํ ตฺวํ อาโกฏิตปจฺจาโกฏิตานิ จีวรานิ ปารุเปยฺยาสิ, อกฺขีนิ จ อญฺเชยฺยาสิ, อจฺฉญฺจ ปตฺตํ ธาเรยฺยาสิ.\csromanspacer{} เอตํ โข เต นนฺท ปติรูปํ กุลปุตฺตสฺส สทฺธา อคารสฺมา อนคาริยํ ปพฺพชิตสฺส, ยํ ตฺวํ อารญฺญิโก จ อสฺสสิ, ปิณฺฑปาติโก จ ปํสุกูลิโก จ กาเมสุ จ อนเปกฺโข วิหเรยฺยาสี”ติ.\csromanspacer{} อิทมโวจ ภควา ฯเปฯ สตฺถา\csromandash{}}

\csromangathameasure{“กทาหํ นนฺทํ ปสฺเสยฺยํ, อารญฺญํ ปํสุกูลิกํ. \\ อญฺญาตุญฺเฉน ยาเปนฺตํ, กาเมสุ อนเปกฺขินนฺ”ติ.}
\csromangathagroup{\csromangathastanza{“กทาหํ นนฺทํ ปสฺเสยฺยํ, อารญฺญํ ปํสุกูลิกํ. \\ อญฺญาตุญฺเฉน ยาเปนฺตํ, กาเมสุ อนเปกฺขินนฺ”ติ.}}

\csromanheader{2}{10. ภิกฺขุสํยุตฺต}
\prose{\csromanfolio{469}อถ โข อายสฺมา นนฺโท อปเรน สมเยน อารญฺญิโก จ ปิณฺฑปาติโก จ ปํสุกูลิโก จ กาเมสุ จ อนเปกฺโข วิหาสีติ.\csromanspacer{}\csromanspacer{}อฏฺฐมํ.}

\csromanheader{6}{9. ติสฺสสุตฺต}
\proseitem{243}{สาวตฺถิยํ วิหรติ.\csromanspacer{} อถ โข อายสฺมา ติสฺโส ภควโต ปิตุจฺฉาปุตฺโต เยน ภควา เตนุปสงฺกมิ, อุปสงฺกมิตฺวา ภควนฺตํ อภิวาเทตฺวา เอกมนฺตํ นิสีทิ ทุกฺขี ทุมฺมโน อสฺสูนิ ปวตฺตยมาโน.\csromanspacer{} อถ โข ภควา อายสฺมนฺตํ ติสฺสํ เอตทโวจ “กิํ นุ โข ตฺวํ ติสฺส เอกมนฺตํ นิสินฺโน ทุกฺขี ทุมฺมโน อสฺสูนิ ปวตฺตยมาโน”ติ.\csromanspacer{} ตถา หิ ปน มํ ภนฺเต ภิกฺขู สมนฺตา วาจาสนฺนิโตทเกน\footnote{วาจาย สนฺนิโตทเกน (ก)} สญฺชมฺภริมกํสูติ\footnote{สญฺชพฺภริมกํสูติ (?)}.\csromanspacer{} ตถา หิ ปน ตฺวํ ติสฺส วตฺตา โน จ วจนกฺขโม, น โข เต ตํ ติสฺส ปติรูปํ กุลปุตฺตสฺส สทฺธา อคารสฺมา อนคาริยํ ปพฺพชิตสฺส, ยํ ตฺวํ วตฺตา โน จ วจนกฺขโม, เอตํ โข เต ติสฺส ปติรูปํ กุลปุตฺตสฺส สทฺธา อคารสฺมา อนคาริยํ ปพฺพชิตสฺส, ยํ ตฺวํ วตฺตา จ อสฺส วจนกฺขโม จาติ.}

\prose{อิทมโวจ ภควา, อิทํ วตฺวาน สุคโต อถาปรํ เอตทโวจ สตฺถา\csromandash{}}

\csromangathameasure{“กิํ นุ กุชฺฌสิ มา กุชฺฌิ, อกฺโกโธ ติสฺส เต วรํ. \\ โกธมานมกฺขวินยตฺถํ หิ, ติสฺส พฺรหฺมจริยํ วุสฺสตี”ติ.}
\csromangathagroup{\csromangathastanza{“กิํ นุ กุชฺฌสิ มา กุชฺฌิ, อกฺโกโธ ติสฺส เต วรํ. \\ โกธมานมกฺขวินยตฺถํ หิ, ติสฺส พฺรหฺมจริยํ วุสฺสตี”ติ.}}

\vspace{\tipitakaparskip}
\csromancenter{นวมํ.}
\csromanheader{6}{10. เถรนามกสุตฺต}
\csromanmatikaanchor{mtk.299}
\csromantocmarkref{subparagraph}{4. นวสุตฺต}{mtk.299}
\bookmark[dest={mtk.299},level=5]{4. นวสุตฺต}
\proseitem{244}{เอกํ สมยํ ภควา ราชคเห วิหรติ เวฬุวเน กลนฺทกนิวาเป.\csromanspacer{} เตน โข ปน สมเยน อญฺญตโร ภิกฺขุ เถรนามโก เอกวิหารี เจว โหติ เอกวิหารสฺส จ วณฺณวาที.\csromanspacer{} โส เอโก คามํ ปิณฺฑาย ปวิสติ, เอโก ปฏิกฺกมติ, เอโก รโห นิสีทติ, เอโก จงฺกมํ อธิฏฺฐาติ.\csromanspacer{} อถ โข สมฺพหุลา ภิกฺขู เยน ภควา เตนุปสงฺกมิํสุ, อุปสงฺกมิตฺวา ภควนฺตํ อภิวาเทตฺวา เอกมนฺตํ นิสีทิํสุ, เอกมนฺตํ นิสินฺนา}

\gambhira{นิทานวคฺคสํยุตฺตปาฬิ}
\prose{\csromanfolio{470}โข เต ภิกฺขู ภควนฺตํ เอตทโวจุํ “อิธ ภนฺเต อญฺญตโร ภิกฺขุ เถรนามโก เอกวิหารี เอกวิหารสฺส จ วณฺณวาที”ติ.}

\prose{อถ โข ภควา อญฺญตรํ ภิกฺขุํ อามนฺเตสิ “เอหิ ตฺวํ ภิกฺขุ มม วจเนน เถรํ ภิกฺขุํ อามนฺเตหิ ‘สตฺถา ตํ อาวุโส เถร อามนฺเตตี’ติ”. “เอวํ ภนฺเต”ติ โข โส ภิกฺขุ ภควโต ปฏิสฺสุตฺวา เยนายสฺมา เถโร เตนุปสงฺกมิ, อุปสงฺกมิตฺวา อายสฺมนฺตํ เถรํ เอตทโวจ “สตฺถา ตํ อาวุโส เถร อามนฺเตตี”ติ. “เอวมาวุโส”ติ โข อายสฺมา เถโร ตสฺส ภิกฺขุโน ปฏิสฺสุตฺวา เยน ภควา เตนุปสงฺกมิ, อุปสงฺกมิตฺวา ภควนฺตํ อภิวาเทตฺวา เอกมนฺตํ นิสีทิ, เอกมนฺตํ นิสินฺนํ โข อายสฺมนฺตํ เถรํ ภควา เอตทโวจ “สจฺจํ กิร ตฺวํ เถร เอกวิหารี เอกวิหารสฺส จ วณฺณวาทีติ.\csromanspacer{} เอวํ ภนฺเต.\csromanspacer{} ยถา กถํ ปน ตฺวํ เถร เอกวิหารี เอกวิหารสฺส จ วณฺณวาทีติ.\csromanspacer{} อิธาหํ ภนฺเต เอโก คามํ ปิณฺฑาย ปวิสามิ, เอโก ปฏิกฺกมามิ, เอโก รโห นิสีทามิ, เอโก จงฺกมํ อธิฏฺฐามิ, เอวํ ขฺวาหํ ภนฺเต เอกวิหารี เอกวิหารสฺส จ วณฺณวาทีติ.}

\prose{อตฺเถโส เถร เอกวิหาโร, เนโส นตฺถีติ วทามิ, อปิ จ เถร ยถา เอกวิหาโร วิตฺถาเรน ปริปุณฺโณ โหติ, ตํ สุณาหิ สาธุกํ มนสิ กโรหิ, ภาสิสฺสามีติ. “เอวํ ภนฺเต”ติ โข ฯเปฯ.\csromanspacer{} กถญฺจ เถร เอกวิหาโร วิตฺถาเรน ปริปุณฺโณ โหติ.\csromanspacer{} อิธ เถร ยํ อตีตํ ตํ ปหีนํ, ยํ อนาคตํ ตํ ปฏินิสฺสฏฺฐํ, ปจฺจุปฺปนฺเนสุ จ อตฺตภาวปฏิลาเภสุ ฉนฺทราโค สุปฺปฏิวินีโต.\csromanspacer{} เอวํ โข เถร เอกวิหาโร วิตฺถาเรน ปริปุณฺโณ โหตีติ.}

\prose{อิทมโวจ ภควา, อิทํ วตฺวาน สุคโต อถาปรํ เอตทโวจ สตฺถา\csromandash{}}

\csromangathameasure{“สพฺพาภิภุํ สพฺพวิทุํ สุเมธํ, \\ สพฺเพสุ ธมฺเมสุ อนูปลิตฺตํ. \\ สพฺพญฺชหํ ตณฺหากฺขเย วิมุตฺตํ, \\ ตมหํ นรํ เอกวิหารีติ พฺรูมี”ติ.}
\csromangathagroup{\csromangathastanza{“สพฺพาภิภุํ สพฺพวิทุํ สุเมธํ, \\ สพฺเพสุ ธมฺเมสุ อนูปลิตฺตํ. \\ สพฺพญฺชหํ ตณฺหากฺขเย วิมุตฺตํ, \\ ตมหํ นรํ เอกวิหารีติ พฺรูมี”ติ.}}

\vspace{\tipitakaparskip}
\csromancenter{ทสมํ.}
\csromanheader{6}{10. ภิกฺขุสํยุตฺต \textbf{11. มหากปฺปินสุตฺต}}
\csromanmatikaanchor{mtk.300}
\csromantocmarkref{subparagraph}{5. สุชาตสุตฺต}{mtk.300}
\bookmark[dest={mtk.300},level=5]{5. สุชาตสุตฺต}
\proseitem{245}{\csromanfolio{471}สาวตฺถิยํ วิหรติ.\csromanspacer{} อถ โข อายสฺมา มหากปฺปิโน เยน ภควา เตนุปสงฺกมิ.\csromanspacer{} อทฺทสา โข ภควา อายสฺมนฺตํ มหากปฺปินํ ทูรโตว อาคจฺฉนฺตํ, ทิสฺวาน ภิกฺขู อามนฺเตสิ “ปสฺสถ โน ตุเมฺห ภิกฺขเว เอตํ ภิกฺขุํ อาคจฺฉนฺตํ โอทาตกํ ตนุกํ ตุงฺคนาสิกนฺ”ติ.\csromanspacer{} เอวํ ภนฺเต.\csromanspacer{} เอโส โข ภิกฺขเว ภิกฺขุ มหิทฺธิโก มหานุภาโว, น จ สา สมาปตฺติ สุลภรูปา, ยา เตน ภิกฺขุนา อสมาปนฺนปุพฺพา.\csromanspacer{} ยสฺส จตฺถาย กุลปุตฺตา สมฺมเทว อคารสฺมา อนคาริยํ ปพฺพชนฺติ, ตทนุตฺตรํ พฺรหฺมจริยปริโยสานํ ทิฏฺเฐว ธมฺเม สยํ อภิญฺญา สจฺฉิกตฺวา อุปสมฺปชฺช วิหรตีติ.}

\prose{อิทมโวจ ภควา, อิทํ วตฺวาน สุคโต อถาปรํ เอตทโวจ สตฺถา\csromandash{}}

\csromangathameasure{“ขตฺติโย เสฏฺโฐ ชเนตสฺมิํ, เย โคตฺตปฏิสาริโน. \\ วิชฺชาจรณสมฺปนฺโน, โส เสฏฺโฐ เทวมานุเส. \\ ทิวา ตปติ อาทิจฺโจ, รตฺติมาภาติ จนฺทิมา. \\ สนฺนทฺโธ ขตฺติโย ตปติ, ฌายี ตปติ พฺราหฺมโณ. \\ อถ สพฺพมโหรตฺติํ, พุทฺโธ ตปติ เตชสา”ติ.}
\csromangathagroup{\csromangathastanza{“ขตฺติโย เสฏฺโฐ ชเนตสฺมิํ, เย โคตฺตปฏิสาริโน. \\ วิชฺชาจรณสมฺปนฺโน, โส เสฏฺโฐ เทวมานุเส.}\csromangathastanza{ทิวา ตปติ อาทิจฺโจ, รตฺติมาภาติ จนฺทิมา. \\ สนฺนทฺโธ ขตฺติโย ตปติ, ฌายี ตปติ พฺราหฺมโณ. \\ อถ สพฺพมโหรตฺติํ\footnote{อถ สพฺพมโหรตฺตํ (สี, สฺยา, กํ)}, พุทฺโธ ตปติ เตชสา”ติ.}}

\csromanmatikaanchor{mtk.301}
\csromantocmarkref{subparagraph}{6. ลกุณฺฑกภทฺทิยสุตฺต}{mtk.301}
\bookmark[dest={mtk.301},level=5]{6. ลกุณฺฑกภทฺทิยสุตฺต}
\vspace{\tipitakaparskip}
\csromancenter{เอกาทสมํ.}
\csromanheader{6}{12. สหายกสุตฺต}
\proseitem{246}{สาวตฺถิยํ วิหรติ.\csromanspacer{} อถ โข เทฺว ภิกฺขู สหายกา อายสฺมโต มหากปฺปินสฺส สทฺธิวิหาริโน เยน ภควา เตนุปสงฺกมิํสุ.\csromanspacer{} อทฺทสา โข ภควา เต ภิกฺขู ทูรโตว อาคจฺฉนฺเต, ทิสฺวาน ภิกฺขู อามนฺเตสิ “ปสฺสถ โน ตุเมฺห ภิกฺขเว เอเต ภิกฺขู สหายเก อาคจฺฉนฺเต กปฺปินสฺส สทฺธิวิหาริโน”ติ.\csromanspacer{} เอวํ ภนฺเต.\csromanspacer{} เอเต โข เต ภิกฺขู มหิทฺธิกา มหานุภาวา, น จ สา สมาปตฺติ สุลภรูปา, ยา เตหิ ภิกฺขูหิ อสมาปนฺนปุพฺพา.\csromanspacer{} ยสฺส จตฺถาย กุลปุตฺตา สมฺมเทว อคารสฺมา อนคาริยํ ปพฺพชนฺติ, ตทนุตฺตรํ พฺรหฺมจริยปริเยสานํ ทิฏฺเฐว ธมฺเม สยํ อภิญฺญา สจฺฉิกตฺวา อุปสมฺปชฺช วิหรนฺตีติ.}

\gambhira{นิทานวคฺคสํยุตฺตปาฬิ}
\csromanmatikaanchor{mtk.302}
\csromanmatikaanchor{mtk.308}
\csromantocmarkref{subparagraph}{7. วิสาขสุตฺต}{mtk.302}
\bookmark[dest={mtk.302},level=5]{7. วิสาขสุตฺต}
\prose{\csromanfolio{472}อิทมโวจ ภควา, อิทํ วตฺวาน สุคโต อถาปรํ เอตทโวจ สตฺถา\csromandash{}}

\csromangathameasure{“สหายาวติเม ภิกฺขู, จิรรตฺตํ สเมติกา. \\ สเมติ เนสํ สทฺธมฺโม, ธมฺเม พุทฺธปฺปเวทิเต. \\ สุวินีตา กปฺปิเนน, ธมฺเม อริยปฺปเวทิเต. \\ ธาเรนฺติ อนฺติมํ เทหํ, เชตฺวา มารํ สวาหินินฺ”ติ.}
\csromangathagroup{\csromangathastanza{“สหายาวติเม ภิกฺขู, จิรรตฺตํ สเมติกา. \\ สเมติ เนสํ สทฺธมฺโม, ธมฺเม พุทฺธปฺปเวทิเต.}\csromangathastanza{สุวินีตา กปฺปิเนน, ธมฺเม อริยปฺปเวทิเต. \\ ธาเรนฺติ อนฺติมํ เทหํ, เชตฺวา มารํ สวาหินินฺ”ติ.}}

\vspace{\tipitakaparskip}
\csromancenter{ทฺวาทสมํ.}
\csromancenter{\textbf{ภิกฺขุสํยุตฺตํ สมตฺตํ.}}
\titlehead{\textbf{ตสฺสุทฺทานํ}}
\csromangathameasure{โกลิโต อุปติสฺโส จ, ฆโฏ จาปิ ปวุจฺจติ. \\ นโว สุชาโต ภทฺทิ จ, วิสาโข นนฺโท ติสฺโส จ. \\ เถรนาโม จ กปฺปิโน, สหาเยน จ ทฺวาทสาติ. \\ นิทานวคฺโค ทุติโย.}
\csromangathagroup{\csromangathastanza{โกลิโต อุปติสฺโส จ, ฆโฏ จาปิ ปวุจฺจติ. \\ นโว สุชาโต ภทฺทิ จ, วิสาโข นนฺโท ติสฺโส จ.}\csromangathastanza{เถรนาโม จ กปฺปิโน, สหาเยน จ ทฺวาทสาติ. \\ นิทานวคฺโค ทุติโย.}}

\titlehead{\textbf{ตสฺสุทฺทานํ}}
\csromanmatikaanchor{mtk.303}
\csromantocmarkref{subparagraph}{8. นนฺทสุตฺต}{mtk.303}
\bookmark[dest={mtk.303},level=5]{8. นนฺทสุตฺต}
\csromangathameasure{นิทานาภิสมยธาตุ, อนมตคฺเคน กสฺสปํ. \\ สกฺการราหุลลกฺขโณ, โอปมฺม ภิกฺขุนา วคฺโค. \\ ทุติโย เตน ปวุจฺจตีติ.}
\csromangathagroup{\csromangathastanza{นิทานาภิสมยธาตุ, อนมตคฺเคน กสฺสปํ. \\ สกฺการราหุลลกฺขโณ, โอปมฺม ภิกฺขุนา วคฺโค. \\ ทุติโย เตน ปวุจฺจตีติ.}}

\nitthitam{\textbf{นิทานวคฺคสํยุตฺตปาฬิ นิฏฺฐิตา.}}
\orphannote{เอตฺถ “อถ โข ราชา ปเสนทิ โกสโล จตุรงฺคินิํ เสนํ สนฺนยฺหิตฺวา ราชานํ มาคธํ อชาตสตฺตุํ เวเทหิปุตฺตํ อพฺภุยฺยาสี”ติ-อาทินา ปาเฐน ภวิตพฺพํ. อฏฺฐกถายํ หิ “อพฺภุยฺยาสีติ ปราชเย ครหปฺปตฺโต ฯเปฯ วุตฺตชยการณํ สุตฺวา อภิ-อุยฺยาสี”ติ วุตฺตํ.}

\orphannote{อุปริ ตติยวคฺเค ตติยจตุตฺถสุตฺเตสุ “มา จ โข ตฺวํ ตาต เสขํ ... อนุปาปุณาตู”ติ อาคตํ. เตน นเยน อิธาปิ อตฺโถ คเหตพฺโพ. เอตฺถ หิ กิํสทฺเทน ปฏิกฺเขปตฺโถปิ สกฺกา ญาตุํ, ยถา “สยํ อภิญฺญาย กมุทฺทิเสยฺยนฺ”ติ. ตสฺมา กํ ... อาคจฺฉตูติ เอตฺถ กมปิ ... มา อาคจฺฉตูติ จ, กํ เสขํ ... อนุปาปุณาตูติ เอตฺถ กมปิ เสขํ ... มา ปาปุณาตูติ จ อตฺโถ เวทิตพฺโพ. อฏฺฐกถาฏีกาสุ จ อยเมวตฺโถ ญาปิโต.}

\orphannote{สงฺขริตฺวา (อิ, ก), มํสํ ขาทิตฺวา (สฺยา, กํ), อสํขาทิตฺวา (กตฺถจิ)}

